% Auto-generated from data/segments.json + layout.json (+ transforms) — do not edit by hand.
% volume: 04Vin04
% mode: printing
% segments: 2757
% transform_rules: 8
% toc_marks: 442

% Volume layout defaults (layout.json ``layout``).
\csromanlayoutapply{1}{1.5}{21.6pt}{8pt}{8pt}{65pt}{2.5em}%

\pitaka{\textbf{วินยปิฏก}}
\csromanmatikaanchor{mtk.0}
\csromantocmarknopage{chapter}{จูฬวคฺคปาฬิ}{mtk.0}
\bookmark[dest={mtk.0},level=0]{จูฬวคฺคปาฬิ}
\gambhira{\textbf{จูฬวคฺคปาฬิ}}
\namakkaram{นโม ตสฺส ภควโต อรหโต สมฺมา\-สมฺพุทฺธสฺส.}
\csromanmatikaanchor{mtk.1}
\csromantocmarknopage{chapter}{1. กมฺมกฺขนฺธก}{mtk.1}
\bookmark[dest={mtk.1},level=0]{1. กมฺมกฺขนฺธก}
\chapterhead{1. กมฺม\-กฺขนฺธก}
\csromanmatikaanchor{mtk.2}
\csromantocmarkref{section}{1. ตชฺชนียกมฺม}{mtk.2}
\bookmark[dest={mtk.2},level=1]{1. ตชฺชนียกมฺม}
\csromanheader{1}{1. \textbf{ตชฺชนียกมฺม}}
\proseitem{1}{\csromanfolio{1}เตน สมเยน พุทฺโธ ภควา สาวตฺถิยํ วิหรติ เชตวเน อนาถ\-ปิณฺฑิกสฺส อาราเม. เตน โข ปน สมเยน ปณฺฑุก\-โลหิตกา ภิกฺขู อตฺตนา ภณฺฑนการกา กลหการกา วิวาทการกา ภสฺสการกา สํเฆ อธิกรณ\-การกา, เยปิ จญฺเญ ภิกฺขู ภณฺฑนการกา กลหการกา วิวาทการกา ภสฺสการกา สํเฆ อธิกรณ\-การกา, เต อุปสงฺกมิตฺวา เอวํ วทนฺติ “มา โข ตุเมฺห อายสฺมนฺโต เอโส อเชสิ, พลวาพลวํ ปฏิมนฺเตถ, ตุเมฺห เตน ปณฺฑิตตรา จ พฺยตฺตตรา จ พหุสฺสุตตรา จ อลมตฺตตรา จ\csromansharedfootnote{2d9078b417a17e3d}{อลมตฺถตรา จ (สฺยา, ก)}, มา จสฺส ภายิตฺถ, มยมฺปิ ตุมฺหากํ ปกฺขา ภวิสฺสามา”ติ. เตน อนุปฺปนฺนานิ เจว ภณฺฑนานิ อุปฺปชฺชนฺติ, อุปฺปนฺนานิ จ ภณฺฑนานิ ภิยฺโยภาวาย เวปุลฺลาย สํวตฺตนฺติ. เย เต ภิกฺขู อปฺปิจฺฉา ฯเปฯ เต อุชฺฌายนฺติ ขิยฺยนฺติ วิปาเจนฺติ “กถํ หิ นาม ปณฺฑุก\-โลหิตกา ภิกฺขู อตฺตนา ภณฺฑนการกา กลหการกา วิวาทการกา ภสฺสการกา สํเฆ อธิกรณ\-การกา, เยปิ จญฺเญ ภิกฺขู ภณฺฑนการกา กลหการกา วิวาทการกา ภสฺสการกา สํเฆ อธิกรณ\-การกา, เต อุปสงฺกมิตฺวา เอวํ วกฺขนฺติ ‘มา โข ตุเมฺห อายสฺมนฺโต เอโส อเชสิ, \csromanfolio{2}พลวาพลวํ ปฏิมนฺเตถ, ตุเมฺห เตน ปณฺฑิตตรา จ พฺยตฺตตรา จ พหุสฺสุตตรา จ อลมตฺตตรา จ, มา จสฺส ภายิตฺถ, มยมฺปิ ตุมฺหากํ ปกฺขา ภวิสฺสามา’ติ, เตน อนุปฺปนฺนานิ เจว ภณฺฑนานิ อุปฺปชฺชนฺติ, อุปฺปนฺนานิ จ ภณฺฑนานิ ภิยฺโยภาวาย เวปุลฺลาย สํวตฺตนฺตี”ติ.}

\proseitem{2}{อถ โข เต ภิกฺขู ภควโต เอตมตฺถํ อาโรเจสุํ. อถ โข ภควา เอตสฺมิํ นิทาเน เอตสฺมิํ ปกรเณ ภิกฺขุสํฆํ สนฺนิปาตาเปตฺวา ภิกฺขู ปฏิปุจฺฉิ “สจฺจํ กิร ภิกฺขเว ปณฺฑุก\-โลหิตกา ภิกฺขู อตฺตนา ภณฺฑนการกา กลหการกา วิวาทการกา ภสฺสการกา สํเฆ อธิกรณ\-การกา, เยปิ จญฺเญ ภิกฺขู ภณฺฑนการกา กลหการกา วิวาทการกา ภสฺสการกา สํเฆ อธิกรณ\-การกา, เต อุปสงฺกมิตฺวา เอวํ วทนฺติ ‘มา โข ตุเมฺห อายสฺมนฺโต เอโส อเชสิ, พลวาพลวํ ปฏิมนฺเตถ, ตุเมฺห เตน ปณฺฑิตตรา จ พฺยตฺตตรา จ พหุสฺสุตตรา จ อลมตฺตตรา จ, มา จสฺส ภายิตฺถ, มยมฺปิ ตุมฺหากํ ปกฺขา ภวิสฺสามา’ติ, เตน อนุปฺปนฺนานิ เจว ภณฺฑนานิ อุปฺปชฺชนฺติ, อุปฺปนฺนานิ จ ภณฺฑนานิ ภิยฺโยภาวาย เวปุลฺลาย สํวตฺตนฺตี”ติ. สจฺจํ ภควาติ. วิครหิ พุทฺโธ ภควา “อนนุจฺฉวิกํ ภิกฺขเว เตสํ โมฆปุริสานํ อนนุโลมิกํ อปฺปติรูปํ อสฺสามณกํ อกปฺปิยํ อกรณียํ, กถํ หิ นาม เต ภิกฺขเว โมฆปุริสา อตฺตนา ภณฺฑนการกา ฯเปฯ สํเฆ อธิกรณ\-การกา, เยปิ จญฺเญ ภิกฺขู ภณฺฑนการกา ฯเปฯ สํเฆ อธิกรณ\-การกา, เต อุปสงฺกมิตฺวา เอวํ วกฺขนฺติ ‘มา โข ตุเมฺห อายสฺมนฺโต เอโส อเชสิ, พลวาพลวํ ปฏิมนฺเตถ, ตุเมฺห เตน ปณฺฑิตตรา จ พฺยตฺตตรา จ พหุสฺสุตตรา จ อลมตฺตตรา จ, มา จสฺส ภายิตฺถ, มยมฺปิ ตุมฺหากํ ปกฺขา ภวิสฺสามา’ติ, เตน อนุปฺปนฺนานิ เจว ภณฺฑนานิ อุปฺปชฺชนฺติ, อุปฺปนฺนานิ จ ภณฺฑนานิ ภิยฺโยภาวาย เวปุลฺลาย สํวตฺตนฺติ. เนตํ ภิกฺขเว อปฺปสนฺนานํ วา ปสาทาย ปสนฺนานํ วา ภิยฺโยภาวาย, อถ เขฺวตํ ภิกฺขเว อปฺปสนฺนาน\-ญฺเจว อปฺปสาทาย ปสนฺนานญฺจ เอกจฺจานํ อญฺญถตฺตายา”ติ.}

\prose{อถ โข ภควา เต\csromansharedfootnote{c86862d4b4082119}{ปณฺฑุกโลหิตเก (สฺยา)} ภิกฺขู อเนก\-ปริยาเยน วิครหิตฺวา ทุพฺภรตาย ทุปฺโปสตาย มหิจฺฉตาย อสนฺตุฏฺฐิตาย\csromansharedfootnote{038e43864c470883}{อสนฺตุฏฺฐตาย (สฺยา), อสนฺตุฏฺฐิยา (สี)} สงฺคณิกาย \csromanfolio{3}โกสชฺชสฺส อวณฺณํ ภาสิตฺวา อเนก\-ปริยาเยน สุภรตาย สุโปสตาย อปฺปิจฺฉสฺส สนฺตุฏฺฐสฺส สลฺเลขสฺส ธุตสฺส ปาสาทิกสฺส อปจยสฺส วีริยา\-รมฺภสฺส\csromansharedfootnote{1e491b2e8c0c218f}{วิริยารมฺภสฺส (สี), วีริยารพฺภสฺส (ก)} วณฺณํ ภาสิตฺวา ภิกฺขูนํ ตทนุจฺฉวิกํ ตทนุโลมิกํ ธมฺมิํ กถํ กตฺวา ภิกฺขู อามนฺเตสิ\csromandash{} เตน หิ ภิกฺขเว สํโฆ ปณฺฑุก\-โลหิตกานํ ภิกฺขูนํ ตชฺชนีย\-กมฺมํ กโรตุ. เอวญฺจ ปน ภิกฺขเว กาตพฺพํ, ปฐมํ ปณฺฑุก\-โลหิตกา ภิกฺขู โจเทตพฺพา, โจเทตฺวา สาเรตพฺพา, สาเรตฺวา อาปตฺติํ\csromansharedfootnote{228eef6ef5c661aa}{อาปตฺติ (สี, สฺยา)} อาโรเปตพฺพา, อาปตฺติํ อาโรเปตฺวา พฺยตฺเตน ภิกฺขุนา ปฏิพเลน สํโฆ ญาเปตพฺโพ\csromandash{}}

\proseitem{3}{“สุณาตุ เม ภนฺเต สํโฆ, อิเม ปณฺฑุก\-โลหิตกา ภิกฺขู อตฺตนา ภณฺฑนการกา กลหการกา วิวาทการกา ภสฺสการกา สํเฆ อธิกรณ\-การกา, เยปิ จญฺเญ ภิกฺขู ภณฺฑนการกา กลหการกา วิวาทการกา ภสฺสการกา สํเฆ อธิกรณ\-การกา, เต อุปสงฺกมิตฺวา เอวํ วทนฺติ ‘มา โข ตุเมฺห อายสฺมนฺโต เอโส อเชสิ, พลวาพลวํ ปฏิมนฺเตถ, ตุเมฺห เตน ปณฺฑิตตรา จ พฺยตฺตตฺตตรา จ พหุสฺสุตตรา จ อลมตฺตตรา จ, มา จสฺส ภายิตฺถ, มยมฺปิ ตุมฺหากํ ปกฺขา ภวิสฺสามา’ติ, เตน อนุปฺปนฺนานิ เจว ภณฺฑนานิ อุปฺปชฺชนฺติ, อุปฺปนฺนานิ จ ภณฺฑนานิ ภิยฺโยภาวาย เวปุลฺลาย สํวตฺตนฺติ, ยทิ สํฆสฺส ปตฺตกลฺลํ, สํโฆ ปณฺฑุก\-โลหิตกานํ ภิกฺขูนํ ตชฺชนีย\-กมฺมํ กเรยฺย, เอสา ญตฺติ.}

\prose{สุณาตุ เม ภนฺเต สํโฆ, อิเม ปณฺฑุก\-โลหิตกา ภิกฺขู อตฺตนา ภณฺฑนการกา ฯเปฯ สํเฆ อธิกรณ\-การกา, เยปิ จญฺเญ ภิกฺขู ภณฺฑนการกา ฯเปฯ สํเฆ อธิกรณ\-การกา, เต อุปสงฺกมิตฺวา เอวํ วทนฺติ ‘มา โข ตุเมฺห อายสฺมนฺโต เอโส อเชสิ, พลวาพลวํ ปฏิมนฺเตถ, ตุเมฺห เตน ปณฺฑิตตรา จ พฺยตฺตตรา จ พหุสฺสุตตรา จ อลมตฺตตรา จ, มา จสฺส ภายิตฺถ, มยมฺปิ ตุมฺหากํ ปกฺขา ภวิสฺสามา’ติ, เตน อนุปฺปนฺนานิ เจว ภณฺฑนานิ อุปฺปชฺชนฺติ, อุปฺปนฺนานิ จ ภณฺฑนานิ ภิยฺโยภาวาย เวปุลฺลาย สํวตฺตนฺติ, สํโฆ ปณฺฑุก\-โลหิตกานํ. \csromanfolio{4}ภิกฺขูนํ ตชฺชนีย\-กมฺมํ กโรติ, ยสฺสายสฺมโต ขมติ ปณฺฑุก\-โลหิตกานํ ภิกฺขูนํ ตชฺชนียสฺส กมฺมสฺส กรณํ, โส ตุณฺหสฺส, ยสฺส นกฺขมติ, โส ภาเสยฺย.}

\prose{ทุติยมฺปิ เอตมตฺถํ วทามิ. สุณาตุ เม ภนฺเต สํโฆ, อิเม ปณฺฑุก\-โลหิตกา ภิกฺขู อตฺตนา ภณฺฑนการกา ฯเปฯ สํเฆ อธิกรณ\-การกา, เยปิ จญฺเญ ภิกฺขู ภณฺฑนการกา ฯเปฯ สํเฆ อธิกรณ\-การกา, เต อุปสงฺกมิตฺวา เอวํ วทนฺติ ‘มา โข ตุเมฺห อายสฺมนฺโต เอโส อเชสิ, พลวาพลวํ ปฏิมนฺเตถ, ตุเมฺห เตน ปณฺฑิตตรา จ พฺยตฺตตรา จ พหุสฺสุตตรา จ อลมตฺตตรา จ, มา จสฺส ภายิตฺถ, มยมฺปิ ตุมฺหากํ ปกฺขา ภวิสฺสามา’ติ, เตน อนุปฺปนฺนานิ เจว ภณฺฑนานิ อุปฺปชฺชนฺติ, อุปฺปนฺนานิ จ ภณฺฑนานิ ภิยฺโยภาวาย เวปุลฺลาย สํวตฺตนฺติ, สํโฆ ปณฺฑุก\-โลหิตกานํ ภิกฺขูนํ ตชฺชนีย\-กมฺมํ กโรติ, ยสฺสายสฺมโต ขมติ ปณฺฑุก\-โลหิตกานํ ภิกฺขูนํ ตชฺชนียสฺส กมฺมสฺส กรณํ, โส ตุณฺหสฺส, ยสฺส นกฺขมติ, โส ภาเสยฺย.}

\prose{ตติยมฺปิ เอตมตฺถํ วทามิ. สุณาตุ เม ภนฺเต สํโฆ, อิเม ปณฺฑุก\-โลหิตกา ภิกฺขู อตฺตนา ภณฺฑนการกา ฯเปฯ สํเฆ อธิกรณ\-การกา, เยปิ จญฺเญ ภิกฺขู ภณฺฑนการกา ฯเปฯ สํเฆ อธิกรณ\-การกา, เต อุปสงฺกมิตฺวา เอวํ วทนฺติ ‘มา โข ตุเมฺห อายสฺมนฺโต เอโส อเชสิ, พลวาพลวํ ปฏิมนฺเตถ, ตุเมฺห เตน ปณฺฑิตตรา จ พฺยตฺตตรา จ พหุสฺสุตตรา จ อลมตฺตตรา จ, มา จสฺส ภายิตฺถ, มยมฺปิ ตุมฺหากํ ปกฺขา ภวิสฺสามา’ติ, เตน อนุปฺปนฺนานิ เจว ภณฺฑนานิ อุปฺปชฺชนฺติ, อุปฺปนฺนานิ จ ภณฺฑนานิ ภิยฺโยภาวาย เวปุลฺลาย สํวตฺตนฺติ, สํโฆ ปณฺฑุก\-โลหิตกานํ ภิกฺขูนํ ตชฺชนีย\-กมฺมํ กโรติ, ยสฺสายสฺมโต ขมติ ปณฺฑุก\-โลหิตกานํ ภิกฺขูนํ ตชฺชนียสฺส กมฺมสฺส กรณํ, โส ตุณฺหสฺส, ยสฺส นกฺขมติ, โส ภาเสยฺย.}

\prose{กตํ สํเฆน ปณฺฑุก\-โลหิตกานํ ภิกฺขูนํ ตชฺชนีย\-กมฺมํ, ขมติ สํฆสฺส, ตสฺมา ตุณฺหี, เอวเมตํ ธารยามี”ติ.}

\csromanmatikaanchor{mtk.3}
\csromantocmarkref{subsection}{อธมฺมกมฺมทฺวาทสก}{mtk.3}
\bookmark[dest={mtk.3},level=2]{อธมฺมกมฺมทฺวาทสก}
\csromanheader{2}{อธมฺมกมฺม\-ทฺวาทสก}
\proseitem{4}{\csromanfolio{5}ตีหิ ภิกฺขเว องฺเคหิ สมนฺนาคตํ ตชฺชนีย\-กมฺมํ อธมฺมกมฺมญฺจ\csromansharedfootnote{2497bb441598c21e}{อธมฺมกมฺมญฺเจว (สฺยา)} โหติ อวินยกมฺมญฺจ ทุวูปสนฺตญฺจ. อสมฺมุขา กตํ โหติ, อปฺปฏิปุจฺฉา กตํ โหติ, อปฺปฏิญฺญาย กตํ โหติ. อิเมหิ โข ภิกฺขเว ตีหงฺเคหิ สมนฺนาคตํ ตชฺชนีย\-กมฺมํ อธมฺมกมฺมญฺจ โหติ อวินยกมฺมญฺจ ทุวูปสนฺตญฺจ. (1)}

\prose{อปเรหิปิ ภิกฺขเว ตีหงฺเคหิ สมนฺนาคตํ ตชฺชนีย\-กมฺมํ อธมฺมกมฺมญฺจ โหติ อวินยกมฺมญฺจ ทุวูปสนฺตญฺจ. อนาปตฺติยา กตํ โหติ, อเทสนาคามินิยา อาปตฺติยา กตํ โหติ, เทสิตาย อาปตฺติยา กตํ โหติ. อิเมหิ โข ภิกฺขเว ตีหงฺเคหิ สมนฺนาคตํ ตชฺชนีย\-กมฺมํ อธมฺมกมฺมญฺจ โหติ อวินยกมฺมญฺจ ทุวูปสนฺตญฺจ. (2)}

\prose{อปเรหิปิ ภิกฺขเว ตีหงฺเคหิ สมนฺนาคตํ ตชฺชนีย\-กมฺมํ อธมฺมกมฺมญฺจ โหติ อวินยกมฺมญฺจ ทุวูปสนฺตญฺจ. อโจเทตฺวา กตํ โหติ, อสาเรตฺวา กตํ โหติ, อาปตฺติํ อนาโรเปตฺวา กตํ โหติ. อิเมหิ โข ภิกฺขเว ตีหงฺเคหิ สมนฺนาคตํ ตชฺชนีย\-กมฺมํ อธมฺมกมฺมญฺจ โหติ อวินยกมฺมญฺจ ทุวูปสนฺตญฺจ. (3)}

\prose{อปเรหิปิ ภิกฺขเว ตีหงฺเคหิ สมนฺนาคตํ ตชฺชนีย\-กมฺมํ อธมฺมกมฺมญฺจ โหติ อวินยกมฺมญฺจ ทุวูปสนฺตญฺจ. อสมฺมุขากตํ โหติ, อธมฺเมน กตํ โหติ, วคฺเคน กตํ โหติ. อิเมหิ โข ภิกฺขเว ตีหงฺเคหิ สมนฺนาคตํ ตชฺชนีย\-กมฺมํ อธมฺมกมฺมญฺจ โหติ อวินยกมฺมญฺจ ทุวูปสนฺตญฺจ. (4)}

\prose{อปเรหิปิ ภิกฺขเว ตีหงฺเคหิ สมนฺนาคตํ ตชฺชนีย\-กมฺมํ อธมฺมกมฺมญฺจ โหติ อวินยกมฺมญฺจ ทุวูปสนฺตญฺจ. อปฺปฏิปุจฺฉา กตํ โหติ, อธมฺเมน กตํ โหติ, วคฺเคน กตํ โหติ. อิเมหิ โข ภิกฺขเว ตีหงฺเคหิ สมนฺนาคตํ ตชฺชนีย\-กมฺมํ อธมฺมกมฺมญฺจ โหติ อวินยกมฺมญฺจ ทุวูปสนฺตญฺจ. (5)}

\prose{อปเรหิปิ ภิกฺขเว ตีหงฺเคหิ สมนฺนาคตํ ตชฺชนีย\-กมฺมํ อธมฺมกมฺมญฺจ โหติ อวินยกมฺมญฺจ ทุวูปสนฺตญฺจ. อปฺปฏิญฺญาย กตํ โหติ, อธมฺเมน กตํ โหติ, วคฺเคน กตํ โหติ. อิเมหิ โข ภิกฺขเว ตีหงฺเคหิ สมนฺนาคตํ ตชฺชนีย\-กมฺมํ อธมฺมกมฺมญฺจ โหติ อวินยกมฺมญฺจ ทุวูปสนฺตญฺจ. (6)}

\prose{\csromanfolio{6}อปเรหิปิ ภิกฺขเว ตีหงฺเคหิ สมนฺนาคตํ ตชฺชนีย\-กมฺมํ อธมฺมกมฺมญฺจ โหติ อวินยกมฺมญฺจ ทุวูปสนฺตญฺจ. อนาปตฺติยา กตํ โหติ, อธมฺเมน กตํ โหติ, วคฺเคน กตํ โหติ. อิเมหิ โข ภิกฺขเว ตีหงฺเคหิ สมนฺนาคตํ ตชฺชนีย\-กมฺมํ อธมฺมกมฺมญฺจ โหติ อวินยกมฺมญฺจ ทุวูปสนฺตญฺจ. (7)}

\prose{อปเรหิปิ ภิกฺขเว ตีหงฺเคหิ สมนฺนาคตํ ตชฺชนีย\-กมฺมํ อธมฺมกมฺมญฺจ โหติ อวินยกมฺมญฺจ ทุวูปสนฺตญฺจ. อเทสนาคามินิยา อาปตฺติยา กตํ โหติ, อธมฺเมน กตํ โหติ, วคฺเคน กตํ โหติ. อิเมหิ โข ภิกฺขเว ตีหงฺเคหิ สมนฺนาคตํ ตชฺชนีย\-กมฺมํ อธมฺมกมฺมญฺจ โหติ อวินยกมฺมญฺจ ทุวูปสนฺตญฺจ. (8)}

\prose{อปเรหิปิ ภิกฺขเว ตีหงฺเคหิ สมนฺนาคตํ ตชฺชนีย\-กมฺมํ อธมฺมกมฺมญฺจ โหติ อวินยกมฺมญฺจ ทุวูปสนฺตญฺจ. เทสิตาย อาปตฺติยา กตํ โหติ, อธมฺเมน กตํ โหติ, วคฺเคน กตํ โหติ. อิเมหิ โข ภิกฺขเว ตีหงฺเคหิ สมนฺนาคตํ ตชฺชนีย\-กมฺมํ อธมฺมกมฺมญฺจ โหติ อวินยกมฺมญฺจ ทุวูปสนฺตญฺจ. (9)}

\prose{อปเรหิปิ ภิกฺขเว ตีหงฺเคหิ สมนฺนาคตํ ตชฺชนีย\-กมฺมํ อธมฺมกมฺมญฺจ โหติ อวินยกมฺมญฺจ ทุวูปสนฺตญฺจ. อโจเทตฺวา กตํ โหติ, อธมฺเมน กตํ โหติ, วคฺเคน กตํ โหติ. อิเมหิ โข ภิกฺขเว ตีหงฺเคหิ สมนฺนาคตํ ตชฺชนีย\-กมฺมํ อธมฺมกมฺมญฺจ โหติ อวินยกมฺมญฺจ ทุวูปสนฺตญฺจ. (10)}

\prose{อปเรหิปิ ภิกฺขเว ตีหงฺเคหิ สมนฺนาคตํ ตชฺชนีย\-กมฺมํ อธมฺมกมฺมญฺจ โหติ อวินยกมฺมญฺจ ทุวูปสนฺตญฺจ. อสาเรตฺวา กตํ โหติ, อธมฺเมน กตํ โหติ, วคฺเคน กตํ โหติ. อิเมหิ โข ภิกฺขเว ตีหงฺเคหิ สมนฺนาคตํ ตชฺชนีย\-กมฺมํ อธมฺมกมฺมญฺจ โหติ อวินยกมฺมญฺจ ทุวูปสนฺตญฺจ. (11)}

\prosewithcloser{อปเรหิปิ ภิกฺขเว ตีหงฺเคหิ สมนฺนาคตํ ตชฺชนีย\-กมฺมํ อธมฺมกมฺมญฺจ โหติ อวินยกมฺมญฺจ ทุวูปสนฺตญฺจ. อาปตฺติํ อนาโรเปตฺวา กตํ โหติ, อธมฺเมน กตํ โหติ, วคฺเคน กตํ โหติ. อิเมหิ โข \csromanfolio{7}ภิกฺขเว ตีหงฺเคหิ สมนฺนาคตํ ตชฺชนีย\-กมฺมํ อธมฺมกมฺมญฺจ โหติ อวินยกมฺมญฺจ ทุวูปสนฺตญฺจ. (12)}{อธมฺมกมฺม\-ทฺวาทสกํ นิฏฺฐิตํ.}
\csromanmatikaanchor{mtk.4}
\csromantocmarkref{subsection}{ธมฺมกมฺมทฺวาทสก}{mtk.4}
\bookmark[dest={mtk.4},level=2]{ธมฺมกมฺมทฺวาทสก}
\csromanheader{2}{ธมฺมกมฺม\-ทฺวาทสก}
\proseitem{5}{ตีหิ ภิกฺขเว องฺเคหิ สมนฺนาคตํ ตชฺชนีย\-กมฺมํ ธมฺมกมฺมญฺจ โหติ วินยกมฺมญฺจ สุวูปสนฺตญฺจ. สมฺมุขา กตํ โหติ, ปฏิปุจฺฉา กตํ โหติ, ปฏิญฺญาย กตํ โหติ. อิเมหิ โข ภิกฺขเว ตีหงฺเคหิ สมนฺนาคตํ ตชฺชนีย\-กมฺมํ ธมฺมกมฺมญฺจ โหติ วินยกมฺมญฺจ สุวูปสนฺตญฺจ. (1)}

\prose{อปเรหิปิ ภิกฺขเว ตีหงฺเคหิ สมนฺนาคตํ ตชฺชนีย\-กมฺมํ ธมฺมกมฺมญฺจ โหติ วินยกมฺมญฺจ สุวูปสนฺตญฺจ. อาปตฺติยา กตํ โหติ, เทสนาคามินิยา อาปตฺติยา กตํ โหติ, อเทสิตาย อาปตฺติยา กตํ โหติ. อิเมหิ โข ภิกฺขเว ตีหงฺเคหิ สมนฺนาคตํ ตชฺชนีย\-กมฺมํ ธมฺมกมฺมญฺจ โหติ วินยกมฺมญฺจ สุวูปสนฺตญฺจ. (2)}

\prose{อปเรหิปิ ภิกฺขเว ตีหงฺเคหิ สมนฺนาคตํ ตชฺชนีย\-กมฺมํ ธมฺมกมฺมญฺจ โหติ วินยกมฺมญฺจ สุวูปสนฺตญฺจ. โจเทตฺวา กตํ โหติ, สาเรตฺวา กตํ โหติ, อาปตฺติํ อาโรเปตฺวา กตํ โหติ. อิเมหิ โข ภิกฺขเว ตีหงฺเคหิ สมนฺนาคตํ ตชฺชนีย\-กมฺมํ ธมฺมกมฺมญฺจ โหติ วินยกมฺมญฺจ สุวูปสนฺตญฺจ. (3)}

\prose{อปเรหิปิ ภิกฺขเว ตีหงฺเคหิ สมนฺนาคตํ ตชฺชนีย\-กมฺมํ ธมฺมกมฺมญฺจ โหติ วินยกมฺมญฺจ สุวูปสนฺตญฺจ. สมฺมุขา กตํ โหติ, ธมฺเมน กตํ โหติ, สมคฺเคน กตํ โหติ. อิเมหิ โข ภิกฺขเว ตีหงฺเคหิ สมนฺนาคตํ ตชฺชนีย\-กมฺมํ ธมฺมกมฺมญฺจ โหติ วินยกมฺมญฺจ สุวูปสนฺตญฺจ. (4)}

\prose{อปเรหิปิ ภิกฺขเว ตีหงฺเคหิ สมนฺนาคตํ ตชฺชนีย\-กมฺมํ ธมฺมกมฺมญฺจ โหติ วินยกมฺมญฺจ สุวูปสนฺตญฺจ. ปฏิปุจฺฉา กตํ โหติ, ธมฺเมน กตํ โหติ, สมคฺเคน กตํ โหติ. อิเมหิ โข ภิกฺขเว ตีหงฺเคหิ สมนฺนาคตํ ตชฺชนีย\-กมฺมํ ธมฺมกมฺมญฺจ โหติ วินยกมฺมญฺจ สุวูปสนฺตญฺจ. (5)}

\prose{\csromanfolio{8}อปเรหิปิ ภิกฺขเว ตีหงฺเคหิ สมนฺนาคตํ ตชฺชนีย\-กมฺมํ ธมฺมกมฺมญฺจ โหติ วินยกมฺมญฺจ สุวูปสนฺตญฺจ. ปฏิญฺญาย กตํ โหติ, ธมฺเมน กตํ โหติ, สมคฺเคน กตํ โหติ. อิเมหิ โข ภิกฺขเว ตีหงฺเคหิ สมนฺนาคตํ ตชฺชนีย\-กมฺมํ ธมฺมกมฺมญฺจ โหติ วินยกมฺมญฺจ สุวูปสนฺตญฺจ. (6)}

\prose{อปเรหิปิ ภิกฺขเว ตีหงฺเคหิ สมนฺนาคตํ ตชฺชนีย\-กมฺมํ ธมฺมกมฺมญฺจ โหติ วินยกมฺมญฺจ สุวูปสนฺตญฺจ. อาปตฺติยา กตํ โหติ, ธมฺเมน กตํ โหติ, สมคฺเคน กตํ โหติ. อิเมหิ โข ภิกฺขเว ตีหงฺเคหิ สมนฺนาคตํ ตชฺชนีย\-กมฺมํ ธมฺมกมฺมญฺจ โหติ วินยกมฺมญฺจ สุวูปสนฺตญฺจ. (7)}

\prose{อปเรหิปิ ภิกฺขเว ตีหงฺเคหิ สมนฺนาคตํ ตชฺชนีย\-กมฺมํ ธมฺมกมฺมญฺจ โหติ วินยกมฺมญฺจ สุวูปสนฺตญฺจ. เทสนาคามินิยา อาปตฺติยา กตํ โหติ, ธมฺเมน กตํ โหติ, สมคฺเคน กตํ โหติ. อิเมหิ โข ภิกฺขเว ตีหงฺเคหิ สมนฺนาคตํ ตชฺชนีย\-กมฺมํ ธมฺมกมฺมญฺจ โหติ วินยกมฺมญฺจ สุวูปสนฺตญฺจ. (8)}

\prose{อปเรหิปิ ภิกฺขเว ตีหงฺเคหิ สมนฺนาคตํ ตชฺชนีย\-กมฺมํ ธมฺมกมฺมญฺจ โหติ วินยกมฺมญฺจ สุวูปสนฺตญฺจ. อเทสิตาย อาปตฺติยา กตํ โหติ, ธมฺเมน กตํ โหติ, สมคฺเคน กตํ โหติ. อิเมหิ โข ภิกฺขเว ตีหงฺเคหิ สมนฺนาคตํ ตชฺชนีย\-กมฺมํ ธมฺมกมฺมญฺจ โหติ วินยกมฺมญฺจ สุวูปสนฺตญฺจ. (9)}

\prose{อปเรหิปิ ภิกฺขเว ตีหงฺเคหิ สมนฺนาคตํ ตชฺชนีย\-กมฺมํ ธมฺมกมฺมญฺจ โหติ วินยกมฺมญฺจ สุวูปสนฺตญฺจ. โจเทตฺวา กตํ โหติ, ธมฺเมน กตํ โหติ, สมคฺเคน กตํ โหติ. อิเมหิ โข ภิกฺขเว ตีหงฺเคหิ สมนฺนาคตํ ตชฺชนีย\-กมฺมํ ธมฺมกมฺมญฺจ โหติ วินยกมฺมญฺจ สุวูปสนฺตญฺจ. (10)}

\prose{อปเรหิปิ ภิกฺขเว ตีหงฺเคหิ สมนฺนาคตํ ตชฺชนีย\-กมฺมํ ธมฺมกมฺมญฺจ โหติ วินยกมฺมญฺจ สุวูปสนฺตญฺจ. สาเรตฺวา กตํ โหติ, ธมฺเมน กตํ โหติ, สมคฺเคน กตํ โหติ. อิเมหิ โข ภิกฺขเว ตีหงฺเคหิ สมนฺนาคตํ ตชฺชนีย\-กมฺมํ ธมฺมกมฺมญฺจ โหติ วินยกมฺมญฺจ สุวูปสนฺตญฺจ. (11)}

\csromanmatikaanchor{mtk.5}
\csromantocmarkref{subsection}{อากงฺขมานฉกฺก}{mtk.5}
\bookmark[dest={mtk.5},level=2]{อากงฺขมานฉกฺก}
\prosewithcloser{อปเรหิปิ ภิกฺขเว ตีหงฺเคหิ สมนฺนาคตํ ตชฺชนีย\-กมฺมํ ธมฺมกมฺมญฺจ โหติ วินยกมฺมญฺจ สุวูปสนฺตญฺจ. อาปตฺติํ อาโรเปตฺวา กตํ โหติ, ธมฺเมน กตํ โหติ, สมคฺเคน กตํ โหติ. อิเมหิ โข ภิกฺขเว \csromanfolio{9}ตีหงฺเคหิ สมนฺนาคตํ ตชฺชนีย\-กมฺมํ ธมฺมกมฺมญฺจ โหติ วินยกมฺมญฺจ สุวูปสนฺตญฺจ. (12)}{ธมฺมกมฺม\-ทฺวาทสกํ นิฏฺฐิตํ.}
\proseitem{6}{\csromansymbolfootnote{*}{วิ 5. 219 ปิฏฺเฐปิ.}ตีหิ ภิกฺขเว องฺเคหิ สมนฺนาคตสฺส ภิกฺขุโน อากงฺขมาโน สํโฆ ตชฺชนีย\-กมฺมํ กเรยฺย, ภณฺฑนการโก โหติ กลหการโก วิวาทการโก ภสฺสการโก สํเฆ อธิกรณ\-การโก, พาโล โหติ อพฺยตฺโต อาปตฺติพหุโล อนปทาโน, คิหิสํสฏฺโฐ วิหรติ อนนุโลมิเกหิ คิหิสํสคฺเคหิ. อิเมหิ โข ภิกฺขเว ตีหงฺเคหิ สมนฺนาคตสฺส ภิกฺขุโน อากงฺขมาโน สํโฆ ตชฺชนีย\-กมฺมํ กเรยฺย. (1)}

\prose{อปเรหิปิ ภิกฺขเว ตีหงฺเคหิ สมนฺนาคตสฺส ภิกฺขุโน อากงฺขมาโน สํโฆ ตชฺชนีย\-กมฺมํ กเรยฺย. อธิสีเล สีลวิปนฺโน โหติ, อชฺฌาจาเร อาจารวิปนฺโน โหติ, อติทิฏฺฐิยา ทิฏฺฐิวิปนฺโน โหติ. อิเมหิ โข ภิกฺขเว ตีหงฺเคหิ สมนฺนาคตสฺส ภิกฺขุโน อากงฺขมาโน สํโฆ ตชฺชนีย\-กมฺมํ กเรยฺย. (2)}

\prose{อปเรหิปิ ภิกฺขเว ตีหงฺเคหิ สมนฺนาคตสฺส ภิกฺขุโน อากงฺขมาโน สํโฆ ตชฺชนีย\-กมฺมํ กเรยฺย. พุทฺธสฺส อวณฺณํ ภาสติ, ธมฺมสฺส อวณฺณํ ภาสติ, สํฆสฺส อวณฺณํ ภาสติ. อิเมหิ โข ภิกฺขเว ตีหงฺเคหิ สมนฺนาคตสฺส ภิกฺขุโน อากงฺขมาโน สํโฆ ตชฺชนีย\-กมฺมํ กเรยฺย. (3)}

\prose{ติณฺณํ ภิกฺขเว ภิกฺขูนํ อากงฺขมาโน สํโฆ ตชฺชนีย\-กมฺมํ กเรยฺย. เอโก ภณฺฑนการโก โหติ กลหการโก วิวาทการโก ภสฺสการโก สํเฆ อธิกรณ\-การโก, เอโก พาโล โหติ อพฺยตฺโต อาปตฺติพหุโล อนปทาโน, เอโก คิหิสํสฏฺโฐ วิหรติ อนนุโลมิเกหิ คิหิสํสคฺเคหิ. อิเมสํ โข ภิกฺขเว ติณฺณํ ภิกฺขูนํ อากงฺขมาโน สํโฆ ตชฺชนีย\-กมฺมํ กเรยฺย. (4)}

\prose{\csromanfolio{10}อปเรสมฺปิ ภิกฺขเว ติณฺณํ ภิกฺขูนํ อากงฺขมาโน สํโฆ ตชฺชนีย\-กมฺมํ กเรยฺย. เอโก อธิสีเล สีลวิปนฺโน โหติ, เอโก อชฺฌาจาเร อาจารวิปนฺโน โหติ, เอโก อติทิฏฺฐิยา ทิฏฺฐิวิปนฺโน โหติ. อิเมสํ โข ภิกฺขเว ติณฺณํ ภิกฺขูนํ อากงฺขมาโน สํโฆ ตชฺชนีย\-กมฺมํ กเรยฺย. (5)}

\prosewithcloser{อปเรสมฺปิ ภิกฺขเว ติณฺณํ ภิกฺขูนํ อากงฺขมาโน สํโฆ ตชฺชนีย\-กมฺมํ กเรยฺย. เอโก พุทฺธสฺส อวณฺณํ ภาสติ, เอโก ธมฺมสฺส อวณฺณํ ภาสติ, เอโก สํฆสฺส อวณฺณํ ภาสติ. อิเมสํ โข ภิกฺขเว ติณฺณํ ภิกฺขูนํ อากงฺขมาโน สํโฆ ตชฺชนีย\-กมฺมํ กเรยฺย. (6)}{อากงฺขมาน\-ฉกฺกํ นิฏฺฐิตํ.}
\csromanmatikaanchor{mtk.6}
\csromantocmarkref{subsection}{อฏฺฐารสวตฺต}{mtk.6}
\bookmark[dest={mtk.6},level=2]{อฏฺฐารสวตฺต}
\csromanheader{2}{\textbf{อฏฺฐารสวตฺต}}
\proseitemwithcloser{7}{ตชฺชนียกมฺมกเตน ภิกฺขเว ภิกฺขุนา สมฺมา วตฺติตพฺพํ. ตตฺรายํ สมฺมาวตฺตนา\csromandash{}น อุปสมฺ\-ปาเท\-ตพฺพํ, น นิสฺสโย ทาตพฺโพ, น สามเณโร อุปฏฺฐาเปตพฺโพ, น ภิกฺขุโนวาทก\-สมฺมุติ\csromansharedfootnote{5c6f5c0884ad928c}{สมฺมติ (สฺยา)} สาทิตพฺพา, สมฺมเตนปิ ภิกฺขุนิโย น โอวทิตพฺพา, ยาย อาปตฺติยา สํเฆน ตชฺชนีย\-กมฺมํ กตํ โหติ, สา อาปตฺติ น อาปชฺชิตพฺพา, อญฺญา วา ตาทิสิกา, ตโต วา ปาปิฏฺฐตรา, กมฺมํ น ครหิตพฺพํ, กมฺมิกา น ครหิตพฺพา, น ปกตตฺตสฺส ภิกฺขุโน อุโปสโถ ฐเปตพฺโพ, น ปวารณา ฐเปตพฺพา, น สวจนียํ กาตพฺพํ, น อนุวาโท ปฏฺฐเปตพฺโพ, น โอกาโส กาเรตพฺโพ, น โจเทตพฺโพ, น สาเรตพฺโพ, น ภิกฺขูหิ\csromansharedfootnote{c38a58b2bff6e995}{น ภิกฺขู ภิกฺขูหิ (สฺยา)} สมฺป\-โยเช\-ตพฺพนฺติ.}{ตชฺชนียกมฺเม อฏฺฐารส\-วตฺตํ นิฏฺฐิตํ.}
\csromanmatikaanchor{mtk.7}
\csromantocmarkref{subsection}{นปฺปฏิปฺปสฺสมฺเภตพฺพอฏฺฐารสก}{mtk.7}
\bookmark[dest={mtk.7},level=2]{นปฺปฏิปฺปสฺสมฺเภตพฺพอฏฺฐารสก}
\csromanheader{2}{นปฺปฏิปฺปสฺสมฺเภตพฺพ\-อฏฺฐารสก}
\proseitem{8}{อถ โข สํโฆ ปณฺฑุก\-โลหิตกานํ ภิกฺขูนํ ตชฺชนีย\-กมฺมํ อกาสิ. เต สํเฆน ตชฺชนียกมฺมกตา สมฺมา วตฺตนฺติ, โลมํ \csromanfolio{11}ปาเตนฺติ, เนตฺถารํ วตฺตนฺติ, ภิกฺขู อุปสงฺกมิตฺวา เอวํ วทนฺติ “มยํ อาวุโส สํเฆน ตชฺชนียกมฺมกตา สมฺมา วตฺตาม, โลมํ ปาเตม, เนตฺถารํ วตฺตาม, กถํ นุ โข อเมฺหหิ ปฏิปชฺชิตพฺพนฺ”ติ. ภิกฺขู ภควโต เอตมตฺถํ อาโรเจสุํ ฯเปฯ. เตน หิ ภิกฺขเว สํโฆ ปณฺฑุก\-โลหิตกานํ ภิกฺขูนํ ตชฺชนีย\-กมฺมํ ปฏิปฺปสฺสมฺเภตุ.}

\prose{ปญฺจหิ ภิกฺขเว องฺเคหิ สมนฺนาคตสฺส ภิกฺขุโน ตชฺชนีย\-กมฺมํ นปฺปฏิปฺปสฺสมฺเภตพฺพํ. อุปสมฺปาเทติ, นิสฺสยํ เทติ, สามเณรํ อุปฏฺฐาเปติ, ภิกฺขูโนวาทกสมฺมุติํ สาทิยติ, สมฺมโตปิ ภิกฺขุนิโย โอวทติ. อิเมหิ โข ภิกฺขเว ปญฺจหงฺเคหิ สมนฺนาคตสฺส ภิกฺขุโน ตชฺชนีย\-กมฺมํ นปฺปฏิปฺปสฺสมฺเภตพฺพํ.}

\prose{\csromansymbolfootnote{*}{วิ 5. 317 ปิฏฺเฐปิ.}อปเรหิปิ ภิกฺขเว ปญฺจหงฺเคหิ สมนฺนาคตสฺส ภิกฺขุโน ตชฺชนีย\-กมฺมํ นปฺปฏิปฺปสฺสมฺเภตพฺพํ. ยาย อาปตฺติยา สํเฆน ตชฺชนีย\-กมฺมํ กตํ โหติ, ตํ อาปตฺติํ อาปชฺชติ, อญฺญํ วา ตาทิสิกํ, ตโต วา ปาปิฏฺฐตรํ, กมฺมํ ครหติ, กมฺมิเก ครหติ. อิเมหิ โข ภิกฺขเว ปญฺจหงฺเคหิ สมนฺนาคตสฺส ภิกฺขุโน ตชฺชนีย\-กมฺมํ นปฺปฏิปฺปสฺสมฺเภตพฺพํ.}

\prosewithcloser{อฏฺฐหิ ภิกฺขเว องฺเคหิ สมนฺนาคตสฺส ภิกฺขุโน ตชฺชนีย\-กมฺมํ นปฺปฏิปฺปสฺสมฺเภตพฺพํ. ปกตตฺตสฺส ภิกฺขุโน อุโปสถํ ฐเปติ, ปวารณํ ฐเปติ, สวจนียํ กโรติ, อนุวาทํ ปฏฺฐเปติ, โอกาสํ กาเรติ, โจเทติ, สาเรติ, ภิกฺขูหิ สมฺปโยเชติ\csromansharedfootnote{6f0ef5a358894931}{ภิกฺขู ภิกฺขูหิ สมฺปโยเชติ (สฺยา)}. อิเมหิ โข ภิกฺขเว อฏฺฐหงฺเคหิ สมนฺนาคตสฺส ภิกฺขุโน ตชฺชนีย\-กมฺมํ นปฺปฏิปฺปสฺสมฺเภตพฺพํ.}{นปฺปฏิปฺปสฺสมฺเภตพฺพ\-อฏฺฐารสกํ นิฏฺฐิตํ.}
\csromanmatikaanchor{mtk.8}
\csromantocmarkref{subsection}{ปฏิปฺปสฺสมฺเภตพฺพอฏฺฐารสก}{mtk.8}
\bookmark[dest={mtk.8},level=2]{ปฏิปฺปสฺสมฺเภตพฺพอฏฺฐารสก}
\csromanheader{2}{ปฏิปฺปสฺสมฺเภตพฺพ\-อฏฺฐารสก}
\proseitem{9}{ปญฺจหิ ภิกฺขเว องฺเคหิ สมนฺนาคตสฺส ภิกฺขุโน ตชฺชนีย\-กมฺมํ ปฏิปฺปสฺสมฺเภตพฺพํ. น อุปสมฺปาเทติ, น นิสฺสยํ เทติ, น สามเณรํ อุปฏฺฐาเปติ, น ภิกฺขุโนวาทก\-สมฺมุติํ สาทิยติ, สมฺมโตปิ ภิกฺขุนิโย น โอวทติ. อิเมหิ โข ภิกฺขเว ปญฺจหงฺเคหิ สมนฺนาคตสฺส ภิกฺขุโน ตชฺชนีย\-กมฺมํ ปฏิปฺปสฺสมฺเภตพฺพํ.}

\prose{\csromanfolio{12}อปเรหิปิ ภิกฺขเว ปญฺจหงฺเคหิ สมนฺนาคตสฺส ภิกฺขุโน ตชฺชนีย\-กมฺมํ ปฏิปฺปสฺสมฺเภตพฺพํ. ยาย อาปตฺติยา สํเฆน ตชฺชนีย\-กมฺมํ กตํ โหติ, ตํ อาปตฺติํ น อาปชฺชติ, อญฺญํ วา ตาทิสิกํ, ตโต วา ปาปิฏฺฐตรํ, กมฺมํ น ครหติ, กมฺมิเก น ครหติ. อิเมหิ โข ภิกฺขเว ปญฺจหงฺเคหิ สมนฺนาคตสฺส ภิกฺขุโน ตชฺชนีย\-กมฺมํ ปฏิปฺปสฺสมฺเภตพฺพํ.}

\prosewithcloser{อฏฺฐหิ ภิกฺขเว องฺเคหิ สมนฺนาคตสฺส ภิกฺขุโน ตชฺชนีย\-กมฺมํ ปฏิปฺปสฺสมฺเภตพฺพํ. น ปกตตฺตสฺส ภิกฺขุโน อุโปสถํ ฐเปติ, น ปวารณํ ฐเปติ, น สวจนียํ กโรติ, น อนุวาทํ ปฏฺฐเปติ, น โอกาสํ กาเรติ, น โจเทติ, น สาเรติ, น ภิกฺขูหิ สมฺปโยเชติ. อิเมหิ โข ภิกฺขเว อฏฺฐงฺเคหิ สมนฺนาคตสฺส ภิกฺขุโน ตชฺชนีย\-กมฺมํ ปฏิปฺปสฺสมฺเภตพฺพํ.}{ปฏิปฺปสฺสมฺเภตพฺพ\-อฏฺฐารสกํ นิฏฺฐิตํ.}
\proseitem{10}{เอวญฺจ ปน ภิกฺขเว ปฏิปฺปสฺสมฺเภตพฺพํ, เตหิ ภิกฺขเว ปณฺฑุก\-โลหิตเกหิ ภิกฺขูหิ สํฆํ อุปสงฺกมิตฺวา เอกํสํ อุตฺตราสงฺคํ กริตฺวา วุฑฺฒานํ ภิกฺขูนํ ปาเท วนฺทิตฺวา อุกฺกุฏิกํ นิสีทิตฺวา อญฺชลิํ ปคฺคเหตฺวา เอวมสฺส วจนีโย “มยํ ภนฺเต สํเฆน ตชฺชนียกมฺมกตา สมฺมา วตฺตาม, โลมํ ปาเตม, เนตฺถารํ วตฺตาม, ตชฺชนียสฺส กมฺมสฺส ปฏิปฺปสฺสทฺธิํ ยาจามา”ติ. ทุติยมฺปิ ยาจิตพฺพา. ตติยมฺปิ ยาจิตพฺพา. พฺยตฺเตน ภิกฺขุนา ปฏิพเลน สํโฆ ญาเปตพฺโพ\csromandash{}}

\prose{“สุณาตุ เม ภนฺเต สํโฆ, อิเม ปณฺฑุก\-โลหิตกา ภิกฺขู สํเฆน ตชฺชนียกมฺมกตา สมฺมา วตฺตนฺติ, โลมํ ปาเตนฺติ, เนตฺถารํ วตฺตนฺติ, ตชฺชนียสฺส กมฺมสฺส ปฏิปฺปสฺสทฺธิํ ยาจนฺติ, ยทิ สํฆสฺส ปตฺตกลฺลํ, สํโฆ ปณฺฑุก\-โลหิตกานํ ภิกฺขูนํ ตชฺชนีย\-กมฺมํ ปฏิปฺปสฺสมฺเภยฺย, เอสา ญตฺติ.}

\prose{สุณาตุ เม ภนฺเต สํโฆ, อิเม ปณฺฑุก\-โลหิตกา ภิกฺขู สํเฆน ตชฺชนียกมฺมกตา สมฺมา วตฺตนฺติ, โลมํ ปาเตนฺติ, เนตฺถารํ วตฺตนฺติ, ตชฺชนียสฺส กมฺมสฺส ปฏิปฺปสฺสทฺธิํ ยาจนฺติ, สํโฆ ปณฺฑุก\-โลหิตกานํ ภิกฺขูนํ ตชฺชนีย\-กมฺมํ ปฏิปฺปสฺสมฺเภติ, ยสฺสายสฺมโต ขมติ ปณฺฑุก\-โลหิตกานํ ภิกฺขูนํ ตชฺชนียสฺส กมฺมสฺส ปฏิปฺปสฺสทฺธิ, โส ตุณฺหสฺส, ยสฺส นกฺขมติ, โส ภาเสยฺย.}

\prose{\csromanfolio{13}ทุติยมฺปิ เอตมตฺถํ วทามิ. สุณาตุ เม ภนฺเต สํโฆ, อิเม ปณฺฑุก\-โลหิตกา ภิกฺขู สํเฆน ตชฺชนียกมฺมกตา สมฺมา วตฺตนฺติ, โลมํ ปาเตนฺติ, เนตฺถารํ วตฺตนฺติ, ตชฺชนียสฺส กมฺมสฺส ปฏิปฺปสฺสทฺธิํ ยาจนฺติ, สํโฆ ปณฺฑุก\-โลหิตกานํ ภิกฺขูนํ ตชฺชนีย\-กมฺมํ ปฏิปฺปสฺสมฺเภติ, ยสฺสายสฺมโต ขมติ ปณฺฑุก\-โลหิตกานํ ภิกฺขูนํ ตชฺชนียสฺส กมฺมสฺส ปฏิปฺปสฺสทฺธิ, โส ตุณฺหสฺส, ยสฺส นกฺขมติ, โส ภาเสยฺย.}

\prose{ตติยมฺปิ เอตมตฺถํ วทามิ. สุณาตุ เม ภนฺเต สํโฆ, อิเม ปณฺฑุก\-โลหิตกา ภิกฺขู สํเฆน ตชฺชนียกมฺมกตา สมฺมา วตฺตนฺติ, โลมํ ปาเตนฺติ, เนตฺถารํ วตฺตนฺติ, ตชฺชนียสฺส กมฺมสฺส ปฏิปฺปสฺสทฺธิํ ยาจนฺติ, สํโฆ ปณฺฑุก\-โลหิตกานํ ภิกฺขูนํ ตชฺชนีย\-กมฺมํ ปฏิปฺปสฺสมฺเภติ, ยสฺสายสฺมโต ขมติ ปณฺฑุก\-โลหิตกานํ ภิกฺขูนํ ตชฺชนียสฺส กมฺมสฺส ปฏิปฺปสฺสทฺธิ, โส ตุณฺหสฺส, ยสฺส นกฺขมติ, โส ภาเสยฺย.}

\prose{ปฏิปฺปสฺสทฺธํ สํเฆน ปณฺฑุก\-โลหิตกานํ ภิกฺขูนํ ตชฺชนีย\-กมฺมํ, ขมติ สํฆสฺส, ตสฺมา ตุณฺหี, เอวเมตํ ธารยามี”ติ.}

\vspace{\tipitakaparskip}
\csromancenter{ตชฺชนีย\-กมฺมํ นิฏฺฐิตํ ปฐมํ.}
\csromanmatikaanchor{mtk.9}
\csromantocmarkref{section}{2. นิยสฺสกมฺม}{mtk.9}
\bookmark[dest={mtk.9},level=1]{2. นิยสฺสกมฺม}
\csromanheader{1}{2. \textbf{นิยสฺสกมฺม}}
\proseitem{11}{เตน โข ปน สมเยน อายสฺมา เสยฺยสโก พาโล โหติ อพฺยตฺโต อาปตฺติพหุโล อนปทาโน, คิหิสํสฏฺโฐ วิหรติ อนนุโลมิเกหิ คิหิสํสคฺเคหิ, อปิสฺสุ ภิกฺขู ปกตา\csromansharedfootnote{b587f91c09675564}{ปกตตฺตา (สี, สฺยา)} ปริวาสํ เทนฺตา มูลาย ปฏิกสฺสนฺตา มานตฺตํ เทนฺตา อพฺเภนฺตา. เย เต ภิกฺขู อปฺปิจฺฉา ฯเปฯ เต อุชฺฌายนฺติ ขิยฺยนฺติ วิปาเจนฺติ “กถํ หิ นาม อายสฺมา เสยฺยสโก พาโล ภวิสฺสติ อพฺยตฺโต อาปตฺติพหุโล อนปทาโน, คิหิสํสฏฺโฐ วิหริสฺสติ อนนุโลมิเกหิ คิหิสํสคฺเคหิ, อปิสฺสุ ภิกฺขู ปกตา ปริวาสํ เทนฺตา มูลาย ปฏิกสฺสนฺตา มานตฺตํ เทนฺตา อพฺเภนฺตา”ติ. อถ โข เต ภิกฺขู ภควโต เอตมตฺถํ อาโรเจสุํ.}

\prose{\csromanfolio{14}อถ โข ภควา เอตสฺมิํ นิทาเน เอตสฺมิํ ปกรเณ ภิกฺขุสํฆํ สนฺนิปาตาเปตฺวา ภิกฺขู ปฏิปุจฺฉิ “สจฺจํ กิร ภิกฺขเว เสยฺยสโก ภิกฺขุ พาโล อพฺยตฺโต อาปตฺติพหุโล อนปทาโน, คิหิสํสฏฺโฐ วิหรติ อนนุโลมิเกหิ คิหิสํสคฺเคหิ, อปิสฺสุ ภิกฺขู ปกตา ปริวาสํ เทนฺตา มูลาย ปฏิกสฺสนฺตา มานตฺตํ เทนฺตา อพฺเภนฺตา”ติ. สจฺจํ ภควาติ. วิครหิ พุทฺโธ ภควา อนนุจฺฉวิกํ ภิกฺขเว ตสฺส โมฆปุริสสฺส อนนุโลมิกํ อปฺปติรูปํ อสฺสามณกํ อกปฺปิยํ อกรณียํ. กถํ หิ นาม โส ภิกฺขเว โมฆปุริโส พาโล ภวิสฺสติ อพฺยตฺโต อาปตฺติพหุโล อนปทาโน, คิหิสํสฏฺโฐ วิหริสฺสติ อนนุโลมิเกหิ คิหิสํสคฺเคหิ, อปิสฺสุ ภิกฺขู ปกตา ปริวาสํ เทนฺตา มูลาย ปฏิกสฺสนฺตา มานตฺตํ เทนฺตา อพฺเภนฺตา. เนตํ ภิกฺขเว อปฺปสนฺนานํ วา ปสาทาย ฯเปฯ วิครหิตฺวา ฯเปฯ ธมฺมิํ กถํ กตฺวา ภิกฺขู อามนฺเตสิ\csromandash{} เตน หิ ภิกฺขเว สํโฆ เสยฺยสกสฺส ภิกฺขุโน นิยสฺสกมฺมํ\csromansharedfootnote{d9e228f72afe7dd1}{นิยสกมฺมํ (ก)} กโรตุ “นิสฺสาย เต วตฺถพฺพนฺ”ติ. เอวญฺจ ปน ภิกฺขเว กาตพฺพํ, ปฐมํ เสยฺยสโก ภิกฺขุ โจเทตพฺโพ, โจเทตฺวา สาเรตพฺโพ, สาเรตฺวา อาปตฺติํ อาโรเปตพฺโพ\csromansharedfootnote{a5ea3b2af66e8b9a}{อาปตฺติ อาโรเปตพฺพา (สี, สฺยา)}, อาปตฺติํ อาโรเปตฺวา พฺยตฺเตน ภิกฺขุนา ปฏิพเลน สํโฆ ญาเปตพฺโพ\csromandash{}}

\proseitem{12}{“สุณาตุ เม ภนฺเต สํโฆ, อยํ เสยฺยสโก ภิกฺขุ, พาโล อพฺยตฺโต อาปตฺติพหุโล อนปทาโน, คิหิสํสฏฺโฐ วิหรติ อนนุโลมิเกหิ คิหิสํสคฺเคหิ, อปิสฺสุ ภิกฺขู ปกตา ปริวาสํ เทนฺตา มูลาย ปฏิกสฺสนฺตา มานตฺตํ เทนฺตา อพฺเภนฺตา, ยทิ สํฆสฺส ปตฺตกลฺลํ, สํโฆ เสยฺยสกสฺส ภิกฺขุโน นิยสฺสกมฺมํ กเรยฺย ‘นิสฺสาย เต วตฺถพฺพนฺ’ติ, เอสา ญตฺติ.}

\prose{สุณาตุ เม ภนฺเต สํโฆ, อยํ เสยฺยสโก ภิกฺขุ พาโล อพฺยตฺโต อาปตฺติพหุโล อนปทาโน, คิหิสํสฏฺโฐ วิหรติ อนนุโลมิเกหิ คิหิสํสคฺเคหิ, อปิสฺสุ ภิกฺขู ปกตา ปริวาสํ เทนฺตา มูลาย ปฏิกสฺสนฺตา มานตฺตํ เทนฺตา อพฺเภนฺตา, สํโฆ เสยฺยสกสฺส ภิกฺขุโน นิยสฺสกมฺมํ กโรติ ‘นิสฺสาย เต วตฺถพฺพนฺ’ติ, ยสฺสายสฺมโต ขมติ เสยฺยสกสฺส ภิกฺขุโน นิยสฺสสฺส กมฺมสฺส กรณํ ‘นิสฺสาย เต วตฺถพฺพนฺ’ติ, โส ตุณฺหสฺส, ยสฺส นกฺขมติ, โส ภาเสยฺย.}

\prose{\csromanfolio{15}ทุติยมฺปิ เอตมตฺถํ วทามิ ฯเปฯ. ตติยมฺปิ เอตมตฺถํ วทามิ. สุณาตุ เม ภนฺเต สํโฆ, อยํ เสยฺยสโก ภิกฺขุ พาโล อพฺยตฺโต อาปตฺติพหุโล อนปทาโน, คิหิสํสฏฺโฐ วิหรติ อนนุโลมิเกหิ คิหิสํสคฺเคหิ, อปิสฺสุ ภิกฺขู ปกตา ปริวาสํ เทนฺตา มูลาย ปฏิกสฺสนฺตา มานตฺตํ เทนฺตา อพฺเภนฺตา, สํโฆ เสยฺยสกสฺส ภิกฺขุโน นิยสฺสกมฺมํ กโรติ ‘นิสฺสาย เต วตฺถพฺพนฺ’ติ, ยสฺสายสฺมโต ขมติ เสยฺยสกสฺส ภิกฺขุโน นิยสฺสสฺส กมฺมสฺส กรณํ ‘นิสฺสาย เต วตฺถพฺพนฺ’ติ, โส ตุณฺหสฺส, ยสฺส นกฺขมติ, โส ภาเสยฺย.}

\prose{กตํ สํเฆน เสยฺยสกสฺส ภิกฺขุโน นิยสฺสกมฺมํ ‘นิสฺสาย เต วตฺถพฺพนฺ’ติ, ขมติ สํฆสฺส, ตสฺมา ตุณฺหี, เอวเมตํ ธารยามี”ติ.}

\csromanmatikaanchor{mtk.10}
\csromantocmarkref{subsection}{อธมฺมกมฺมทฺวาทสก}{mtk.10}
\bookmark[dest={mtk.10},level=2]{อธมฺมกมฺมทฺวาทสก}
\csromanheader{2}{อธมฺมกมฺม\-ทฺวาทสก}
\proseitem{13}{ตีหิ ภิกฺขเว องฺเคหิ สมนฺนาคตํ นิยสฺสกมฺมํ อธมฺมกมฺมญฺจ โหติ อวินยกมฺมญฺจ ทุวูปสนฺตญฺจ. อสมฺมุขา กตํ โหติ, อปฺปฏิปุจฺฉา กตํ โหติ, อปฺปฏิญฺญาย กตํ โหติ. อิเมหิ โข ภิกฺขเว ตีหงฺเคหิ สมนฺนาคตํ นิยสฺสกมฺมํ อธมฺมกมฺมญฺจ โหติ อวินยกมฺมญฺจ ทุวูปสนฺตญฺจ. (1)}

\prose{อปเรหิปิ ภิกฺขเว ตีหงฺเคหิ สมนฺนาคตํ นิยสฺสกมฺมํ อธมฺมกมฺมญฺจ โหติ อวินยกมฺมญฺจ ทุวูปสนฺตญฺจ. อนาปตฺติยา กตํ โหติ, อเทสนาคามินิยา อาปตฺติยา กตํ โหติ, เทสิตาย อาปตฺติยา กตํ โหติ. อิเมหิ โข ภิกฺขเว ตีหงฺเคหิ สมนฺนาคตํ นิยสฺสกมฺมํ อธมฺมกมฺมญฺจ โหติ อวินยกมฺมญฺจ ทุวูปสนฺตญฺจ. (2)}

\prose{อปเรหิปิ ภิกฺขเว ตีหงฺเคหิ สมนฺนาคตํ นิยสฺสกมฺมํ อธมฺมกมฺมญฺจ โหติ อวินยกมฺมญฺจ ทุวูปสนฺตญฺจ. อโจเทตฺวา กตํ โหติ, อสาเรตฺวา กตํ โหติ, อาปตฺติํ อนาโรเปตฺวา กตํ โหติ. อิเมหิ โข ภิกฺขเว ตีหงฺเคหิ สมนฺนาคตํ นิยสฺสกมฺมํ อธมฺมกมฺมญฺจ โหติ อวินยกมฺมญฺจ ทุวูปสนฺตญฺจ. (3)}

\prose{อปเรหิปิ ภิกฺขเว ฯเปฯ. อสมฺมุขา กตํ โหติ, อธมฺเมน กตํ โหติ, วคฺเคน กตํ โหติ. อิเมหิ โข ภิกฺขเว ฯเปฯ. (4)}

\prose{\csromanfolio{16}อปเรหิปิ ภิกฺขเว ฯเปฯ. อปฺปฏิปุจฺฉา กตํ โหติ, อธมฺเมน กตํ โหติ, วคฺเคน กตํ โหติ. อิเมหิ โข ภิกฺขเว ฯเปฯ. (5)}

\prose{อปเรหิปิ ภิกฺขเว ฯเปฯ. อปฺปฏิญฺญาย กตํ โหติ, อธมฺเมน กตํ โหติ, วคฺเคน กตํ โหติ. อิเมหิ โข ภิกฺขเว ฯเปฯ. (6)}

\prose{อปเรหิปิ ภิกฺขเว ฯเปฯ. อนาปตฺติยา กตํ โหติ, อธมฺเมน กตํ โหติ, วคฺเคน กตํ โหติ. อิเมหิ โข ภิกฺขเว ฯเปฯ. (7)}

\prose{อปเรหิปิ ภิกฺขเว ฯเปฯ. อเทสนาคามินิยา อาปตฺติยา กตํ โหติ, อธมฺเมน กตํ โหติ, วคฺเคน กตํ โหติ. อิเมหิ โข ภิกฺขเว ฯเปฯ. (8)}

\prose{อปเรหิปิ ภิกฺขเว ฯเปฯ. เทสิตาย อาปตฺติยา กตํ โหติ, อธมฺเมน กตํ โหติ, วคฺเคน กตํ โหติ. อิเมหิ โข ภิกฺขเว ฯเปฯ. (9)}

\prose{อปเรหิปิ ภิกฺขเว ฯเปฯ. อโจเทตฺวา กตํ โหติ, อธมฺเมน กตํ โหติ, วคฺเคน กตํ โหติ. อิเมหิ โข ภิกฺขเว ฯเปฯ. (10)}

\prose{อปเรหิปิ ภิกฺขเว ฯเปฯ. อสาเรตฺวา กตํ โหติ, อธมฺเมน กตํ โหติ, วคฺเคน กตํ โหติ. อิเมหิ โข ภิกฺขเว ฯเปฯ. (11)}

\prose{อปเรหิปิ ภิกฺขเว ตีหงฺเคหิ สมนฺนาคตํ นิยสฺสกมฺมํ อธมฺมกมฺมญฺจ โหติ อวินยกมฺมญฺจ ทุวูปสนฺตญฺจ. อาปตฺติํ อนาโรเปตฺวา กตํ โหติ, อธมฺเมน กตํ โหติ, วคฺเคน กตํ โหติ. อิเมหิ โข ภิกฺขเว ตีหงฺเคหิ สมนฺนาคตํ นิยสฺสกมฺมํ อธมฺมกมฺมญฺจ โหติ อวินยกมฺมญฺจ ทุวูปสนฺตญฺจ. (12)}

\vspace{\tipitakaparskip}
\csromancenter{อธมฺมกมฺม\-ทฺวาทสกํ ทิฏฺฐิตํ.}
\csromanmatikaanchor{mtk.11}
\csromantocmarkref{subsection}{ธมฺมกมฺมทฺวาทสก}{mtk.11}
\bookmark[dest={mtk.11},level=2]{ธมฺมกมฺมทฺวาทสก}
\csromanheader{2}{ธมฺมกมฺม\-ทฺวาทสก}
\proseitem{14}{ตีหิ ภิกฺขเว องฺเคหิ สมนฺนาคตํ นิยสฺสกมฺมํ ธมฺมกมฺมญฺจ โหติ วินยกมฺมญฺจ สุวูปสนฺตญฺจ. สมฺมุขา กตํ โหติ, ปฏิปุจฺฉา กตํ โหติ, ปฏิญฺญาย กตํ โหติ. อิเมหิ โข ภิกฺขเว ตีหงฺเคหิ สมนฺนาคตํ นิยสฺสกมฺมํ ธมฺมกมฺมญฺจ โหติ วินยกมฺมญฺจ สุวูปสนฺตญฺจ. (1)}

\prose{\csromanfolio{17}อปเรหิปิ ภิกฺขเว ตีหงฺเคหิ สมนฺนาคตํ นิยสฺสกมฺมํ ธมฺมกมฺมญฺจ โหติ วินยกมฺมญฺจ สุวูปสนฺตญฺจ. อาปตฺติยา กตํ โหติ, เทสนาคามินิยา อาปตฺติยา กตํ โหติ, อเทสิตาย อาปตฺติยา กตํ โหติ. อิเมหิ โข ภิกฺขเว ตีหงฺเคหิ สมนฺนาคตํ นิยสฺสกมฺมํ ธมฺมกมฺมญฺจ โหติ วินยกมฺมญฺจ สุวูปสนฺตญฺจ. (2)}

\prose{อปเรหิปิ ภิกฺขเว ตีหงฺเคหิ สมนฺนาคตํ นิยสฺสกมฺมํ ธมฺมกมฺมญฺจ โหติ วินยกมฺมญฺจ สุวูปสนฺตญฺจ. โจเทตฺวา กตํ โหติ, สาเรตฺวา กตํ โหติ, อาปตฺติํ อาโรเปตฺวา กตํ โหติ. อิเมหิ โข ภิกฺขเว ตีหงฺเคหิ สมนฺนาคตํ นิยสฺสกมฺมํ ธมฺมกมฺมญฺจ โหติ วินยกมฺมญฺจ สุวูปสนฺตญฺจ. (3)}

\prose{อปเรหิปิ ภิกฺขเว ฯเปฯ. สมฺมุขา กตํ โหติ, ธมฺเมน กตํ โหติ, สมคฺเคน กตํ โหติ. อิเมหิ โข ภิกฺขเว ฯเปฯ. (4)}

\prose{อปเรหิปิ ภิกฺขเว ฯเปฯ. ปฏิปุจฺฉา กตํ โหติ, ธมฺเมน กตํ โหติ, สมคฺเคน กตํ โหติ. อิเมหิ โข ภิกฺขเว ฯเปฯ. (5)}

\prose{อปเรหิปิ ภิกฺขเว ฯเปฯ. ปฏิญฺญายกตํ โหติ, ธมฺเมน กตํ โหติ, สมคฺเคน กตํ โหติ. อิเมหิ โข ภิกฺขเว ฯเปฯ. (6)}

\prose{อปเรหิปิ ภิกฺขเว ฯเปฯ. อาปตฺติยา กตํ โหติ, ธมฺเมน กตํ โหติ, สมคฺเคน กตํ โหติ. อิเมหิ โข ภิกฺขเว ฯเปฯ. (7)}

\prose{อปเรหิปิ ภิกฺขเว ฯเปฯ. เทสนาคามินิยา อาปตฺติยา กตํ โหติ, ธมฺเมน กตํ โหติ, สมคฺเคน กตํ โหติ. อิเมหิ โข ภิกฺขเว ฯเปฯ. (8)}

\prose{อปเรหิปิ ภิกฺขเว ฯเปฯ. อเทสิตาย อาปตฺติยา กตํ โหติ, ธมฺเมน กตํ โหติ, สมคฺเคน กตํ โหติ. อิเมหิ โข ภิกฺขเว ฯเปฯ. (9)}

\prose{อปเรหิปิ ภิกฺขเว ฯเปฯ. โจเทตฺวา กตํ โหติ, ธมฺเมน กตํ โหติ, สมคฺเคน กตํ โหติ. อิเมหิ โข ภิกฺขเว ฯเปฯ. (10)}

\prose{อปเรหิปิ ภิกฺขเว ฯเปฯ. สาเรตฺวา กตํ โหติ, ธมฺเมน กตํ โหติ, สมคฺเคน กตํ โหติ. อิเมหิ โข ภิกฺขเว ฯเปฯ. (11)}

\csromanmatikaanchor{mtk.12}
\csromantocmarkref{subsection}{อากงฺขมานฉกฺก}{mtk.12}
\bookmark[dest={mtk.12},level=2]{อากงฺขมานฉกฺก}
\prosewithcloser{\csromanfolio{18}อปเรหิปิ ภิกฺขเว ตีหงฺเคหิ สมนฺนาคตํ นิยสฺสกมฺมํ ธมฺมกมฺมญฺจ โหติ วินยกมฺมญฺจ สุวูปสนฺตญฺจ. อาปตฺติํ อาโรเปตฺวา กตํ โหติ, ธมฺเมน กตํ โหติ, สมคฺเคน กตํ โหติ. อิเมหิ โข ภิกฺขเว ตีหงฺเคหิ สมนฺนาคตํ นิยสฺสกมฺมํ ธมฺมกมฺมญฺจ โหติ วินยกมฺมญฺจ สุวูปสนฺตญฺจ. (12)}{ธมฺมกมฺม\-ทฺวาทสกํ นิฏฺฐิตํ.}
\proseitem{15}{\csromansymbolfootnote{*}{วิ 5. 219 ปิฏฺเฐปิ.}ตีหิ ภิกฺขเว องฺเคหิ สมนฺนาคตสฺส ภิกฺขุโน อากงฺขมาโน สํโฆ นิยสฺสกมฺมํ กเรยฺย. ภณฺฑนการโก โหติ กลหการโก วิวาทการโก ภสฺสการโก สํเฆ อธิกรณ\-การโก, พาโล โหติ อพฺยตฺโต อาปตฺติพหุโล อนปทาโน, คิหิสํสฏฺโฐ วิหรติ อนนุโลมิเกหิ คิหิสํสคฺเคหิ. อิเมหิ โข ภิกฺขเว ตีหงฺเคหิ สมนฺนาคตสฺส ภิกฺขุโน อากงฺขมาโน สํโฆ นิยสฺสกมฺมํ กเรยฺย. (1)}

\prose{อปเรหิปิ ภิกฺขเว ตีหงฺเคหิ สมนฺนาคตสฺส ภิกฺขุโน อากงฺขมาโน สํโฆ นิยสฺสกมฺมํ กเรยฺย. อธิสีเล สีลวิปนฺโน โหติ, อชฺฌาจาเร อาจารวิปนฺโน โหติ, อติทิฏฺฐิยา ทิฏฺฐิวิปนฺโน โหติ. อิเมหิ โข ภิกฺขเว ตีหงฺเคหิ สมนฺนาคตสฺส ภิกฺขุโน อากงฺขมาโน สํโฆ นิยสฺสกมฺมํ กเรยฺย. (2)}

\prose{อปเรหิปิ ภิกฺขเว ตีหงฺเคหิ สมนฺนาคตสฺส ภิกฺขุโน อากงฺขมาโน สํโฆ นิยสฺสกมฺมํ กเรยฺย. พุทฺธสฺส อวณฺณํ ภาสติ, ธมฺมสฺส อวณฺณํ ภาสติ, สํฆสฺส อวณฺณํ ภาสติ. อิเมหิ โข ภิกฺขเว ตีหงฺเคหิ สมนฺนาคตสฺส ภิกฺขุโน อากงฺขมาโน สํโฆ นิยสฺสกมฺมํ กเรยฺย. (3)}

\prose{ติณฺณํ ภิกฺขเว ภิกฺขูนํ อากงฺขมาโน สํโฆ นิยสฺสกมฺมํ กเรยฺย. เอโก ภณฺฑนการโก โหติ กลหการโก วิวาทการโก ภสฺสการโก สํเฆ อธิกรณ\-การโก, เอโก พาโล โหติ อพฺยตฺโต อาปตฺติพหุโล อนปทาโน, เอโก คิหิสํสฏฺโฐ วิหรติ \csromanfolio{19}อนนุโลมิเกหิ คิหิสํสคฺเคหิ. อิเมสํ โข ภิกฺขเว ติณฺณํ ภิกฺขูนํ อากงฺขมาโน สํโฆ นิยสฺสกมฺมํ กเรยฺย. (4)}

\prose{อปเรสมฺปิ ภิกฺขเว ติณฺณํ ภิกฺขูนํ อากงฺขมาโน สํโฆ นิยสฺสกมฺมํ กเรยฺย. เอโก อธิสีเล สีลวิปนฺโน โหติ, เอโก อชฺฌาจาเร อาจารวิปนฺโน โหติ, เอโก อติทิฏฺฐิยา ทิฏฺฐิวิปนฺโน โหติ. อิเมสํ โข ภิกฺขเว ติณฺณํ ภิกฺขูนํ อากงฺขมาโน สํโฆ นิยสฺสกมฺมํ กเรยฺย. (5)}

\prosewithcloser{อปเรสมฺปิ ภิกฺขเว ติณฺณํ ภิกฺขูนํ อากงฺขมาโน สํโฆ นิยสฺสกมฺมํ กเรยฺย. เอโก พุทฺธสฺส อวณฺณํ ภาสติ, เอโก ธมฺมสฺส อวณฺณํ ภาสติ, เอโก สํฆสฺส อวณฺณํ ภาสติ. อิเมสํ โข ภิกฺขเว ติณฺณํ ภิกฺขูนํ อากงฺขมาโน สํโฆ นิยสฺสกมฺมํ กเรยฺย. (6)}{อากงฺขมาน\-ฉกฺกํ นิฏฺฐิตํ.}
\csromanmatikaanchor{mtk.13}
\csromantocmarkref{subsection}{อฏฺฐารสวตฺต}{mtk.13}
\bookmark[dest={mtk.13},level=2]{อฏฺฐารสวตฺต}
\csromanheader{2}{\textbf{อฏฺฐารสวตฺต}}
\proseitemwithcloser{16}{นิยสฺสกมฺมกเตน ภิกฺขเว ภิกฺขุนา สมฺมา วตฺติตพฺพํ. ตตฺรายํ สมฺมาวตฺตนา\csromandash{}น อุปสมฺ\-ปาเท\-ตพฺพํ, น นิสฺสโย ทาตพฺโพ, น สามเณโร อุปฏฺฐาเปตพฺโพ, น ภิกฺขุโนวาทก\-สมฺมุติ สาทิตพฺพา, สมฺมเตนปิ ภิกฺขุนิโย น โอวทิตพฺพา, ยาย อาปตฺติยา สํเฆน นิยสฺสกมฺมํ กตํ โหติ, สา อาปตฺติ น อาปชฺชิตพฺพา, อญฺญา วา ตาทิสิกา, ตโต วา ปาปิฏฺฐตรา, กมฺมํ น ครหิตพฺพํ, กมฺมิกา น ครหิตพฺพา, น ปกตตฺตสฺส ภิกฺขุโน อุโปสโถ ฐเปตพฺโพ, น ปวารณา ฐเปตพฺพา, น สวจนียํ กาตพฺพํ, น อนุวาโท ปฏฺฐเปตพฺโพ, น โอกาโส การาเปตพฺโพ, น โจเทตพฺโพ, น สาเรตพฺโพ, น ภิกฺขูหิ สมฺป\-โยเช\-ตพฺพนฺติ.}{นิยสฺสกมฺเม อฏฺฐารส\-วตฺตํ นิฏฺฐิตํ.}
\proseitem{17}{อถ โข สํโฆ เสยฺยสกสฺส ภิกฺขุโน นิยสฺสกมฺมํ อกาสิ “นิสฺสาย เต วตฺถพฺพนฺ”ติ. โส สํเฆน นิยสฺสกมฺมกโต \csromanfolio{20}\textbf{กลฺยาณมิตฺเต เสวมาโน ภชมาโน ปยิรุปาสมาโน อุทฺทิสาเปนฺโต} ปริปุจฺฉนฺโต พหุสฺสุโต โหติ อาคตาคโม ธมฺมธโร วินยธโร มาติกาธโร ปณฺฑิโต วิยตฺโต เมธาวี ลชฺชี กุกฺกุจฺจโก สิกฺขากาโม, สมฺมา วตฺตติ, โลมํ ปาเตติ, เนตฺถารํ วตฺตติ, ภิกฺขู อุปสงฺกมิตฺวา เอวํ วเทติ “อหํ อาวุโส สํเฆน นิยสฺสกมฺมกโต สมฺมา วตฺตามิ, โลมํ ปาเตมิ, เนตฺถารํ วตฺตามิ, กถํ นุ โข มยา ปฏิปชฺชิตพฺพนฺ”ติ. ภิกฺขู ภควโต เอตมตฺถํ อาโรเจสุํ ฯเปฯ. เตน หิ ภิกฺขเว สํโฆ เสยฺยสกสฺส ภิกฺขุโน นิยสฺสกมฺมํ ปฏิปฺปสฺสมฺเภตุ.}

\csromanmatikaanchor{mtk.14}
\csromantocmarkref{subsection}{นปฺปฏิปฺปสฺสมฺเภตพฺพอฏฺฐารสก}{mtk.14}
\bookmark[dest={mtk.14},level=2]{นปฺปฏิปฺปสฺสมฺเภตพฺพอฏฺฐารสก}
\csromanheader{2}{นปฺปฏิปฺปสฺสมฺเภตพฺพ\-อฏฺฐารสก}
\proseitem{18}{ปญฺจหิ ภิกฺขเว องฺเคหิ สมนฺนาคตสฺส ภิกฺขุโน นิยสฺสกมฺมํ นปฺปฏิปฺปสฺสมฺเภตพฺพํ. อุปสมฺปาเทติ, นิสฺสยํ เทติ, สามเณรํ อุปฏฺฐาเปติ, ภิกฺขุโนวาทก\-สมฺมุติํ สาทิยติ, สมฺมโตปิ ภิกฺขุนิโย โอวทติ. อิเมหิ โข ภิกฺขเว ปญฺจหงฺเคหิ สมนฺนาคตสฺส ภิกฺขุโน นิยสฺสกมฺมํ นปฺปฏิปฺปสฺสมฺเภตพฺพํ.}

\prose{\csromansymbolfootnote{*}{วิ 5. 317 ปิฏฺเฐปิ.}อปเรหิปิ ภิกฺขเว ปญฺจหงฺเคหิ สมนฺนาคตสฺส ภิกฺขุโน นิยสฺสกมฺมํ นปฺปฏิปฺปสฺสมฺเภตพฺพํ. ยาย อาปตฺติยา สํเฆน นิยสฺสกมฺมํ กตํ โหติ, ตํ อาปตฺติํ อาปชฺชติ, อญฺญํ วา ตาทิสิกํ, ตโต วา ปาปิฏฺฐตรํ, กมฺมํ ครหติ, กมฺมิเก ครหติ. อิเมหิ โข ภิกฺขเว ปญฺจหงฺเคหิ สมนฺนาคตสฺส ภิกฺขุโน นิยสฺสกมฺมํ นปฺปฏิปฺปสฺสมฺเภตพฺพํ.}

\prosewithcloser{อฏฺฐหิ ภิกฺขเว องฺเคหิ สมนฺนาคตสฺส ภิกฺขุโน นิยสฺสกมฺมํ นปฺปฏิปฺปสฺสมฺเภตพฺพํ. ปกตตฺตสฺส ภิกฺขุโน อุโปสถํ ฐเปติ, ปวารณํ ฐเปติ, สวจนียํ กโรติ, อนุวาทํ ปฏฺฐเปติ, โอกาสํ กาเรติ, โจเทติ, สาเรติ, ภิกฺขูหิ สมฺปโยเชติ. อิเมหิ โข ภิกฺขเว อฏฺฐหงฺเคหิ สมนฺนาคตสฺส ภิกฺขุโน นิยสฺสกมฺมํ นปฺปฏิปฺปสฺสมฺเภตพฺพํ.}{นปฺปฏิปฺปสฺสมฺเภตพฺพ\-อฏฺฐารสกํ นิฏฺฐิตํ.}
\csromanmatikaanchor{mtk.15}
\csromantocmarkref{subsection}{ปฏิปฺปสฺสมฺเภตพฺพอฏฺฐารสก}{mtk.15}
\bookmark[dest={mtk.15},level=2]{ปฏิปฺปสฺสมฺเภตพฺพอฏฺฐารสก}
\csromanheader{2}{ปฏิปฺปสฺสมฺเภตพฺพ\-อฏฺฐารสก}
\proseitem{19}{\csromanfolio{21}ปญฺจหิ ภิกฺขเว องฺเคหิ สมนฺนาคตสฺส ภิกฺขุโน นิยสฺสกมฺมํ ปฏิปฺปสฺสมฺเภตพฺพํ. น อุปสมฺปาเทติ, น นิสฺสยํ เทติ, น สามเณรํ อุปฏฺฐาเปติ, น ภิกฺขุโนวาทก\-สมฺมุติํ สาทิยติ, สมฺมโตปิ ภิกฺขุนิโย น โอวทติ. อิเมหิ โข ภิกฺขเว ปญฺจหงฺเคหิ สมนฺนาคตสฺส ภิกฺขุโน นิยสฺสกมฺมํ ปฏิปฺปสฺสมฺเภตพฺพํ.}

\prose{อปเรหิปิ ภิกฺขเว ปญฺจหงฺเคหิ สมนฺนาคตสฺส ภิกฺขุโน นิยสฺสกมฺมํ ปฏิปฺปสฺสมฺเภตพฺพํ. ยาย อาปตฺติยา สํเฆน นิยสฺสกมฺมํ กตํ โหติ, ตํ อาปตฺติํ น อาปชฺชติ, อญฺญํ วา ตาทิสิกํ, ตโต วา ปาปิฏฺฐตรํ, กมฺมํ น ครหติ, กมฺมิเก น ครหติ. อิเมหิ โข ภิกฺขเว ปญฺจหงฺเคหิ สมนฺนาคตสฺส ภิกฺขุโน นิยสฺสกมฺมํ ปฏิปฺปสฺสมฺเภตพฺพํ.}

\prosewithcloser{อฏฺฐหิ ภิกฺขเว องฺเคหิ สมนฺนาคตสฺส ภิกฺขุโน นิยสฺสกมฺมํ ปฏิปฺปสฺสมฺเภตพฺพํ. น ปกตตฺตสฺส ภิกฺขุโน อุโปสถํ ฐเปติ, น ปวารณํ ฐเปติ, น สวจนียํ กโรติ, น อนุวาทํ ปฏฺฐเปติ, น โอกาสํ กาเรติ, น โจเทติ, น สาเรติ, น ภิกฺขูหิ สมฺปโยเชติ. อิเมหิ โข ภิกฺขเว อฏฺฐหงฺเคหิ สมนฺนาคตสฺส ภิกฺขุโน นิยสฺสกมฺมํ ปฏิปฺปสฺสมฺเภตพฺพํ.}{ปฏิปฺปสฺสมฺเภตพฺพ\-อฏฺฐารสกํ นิฏฺฐิตํ.}
\proseitem{20}{เอวญฺจ ปน ภิกฺขเว ปฏิปฺปสฺสมฺเภตพฺพํ, เตน ภิกฺขเว เสยฺยสเกน ภิกฺขุนา สํฆํ อุปสงฺกมิตฺวา เอกํสํ อุตฺตราสงฺคํ กริตฺวา วุฑฺฒานํ ภิกฺขูนํ ปาเท วนฺทิตฺวา อุกฺกุฏิกํ นิสีทิตฺวา อญฺชลิํ ปคฺคเหตฺวา เอวมสฺส วจนีโย “อหํ ภนฺเต สํเฆน นิยสฺสกมฺมกโต สมฺมา วตฺตามิ, โลมํ ปาเตมิ, เนตฺถารํ วตฺตามิ, นิยสฺสสฺส กมฺมสฺส ปฏิปฺปสฺสทฺธิํ ยาจามี”ติ. ทุติยมฺปิ ยาจิตพฺพา. ตติยมฺปิ ยาจิตพฺพา. พฺยตฺเตน ภิกฺขุนา ปฏิพเลน สํโฆ ญาเปตพฺโพ\csromandash{}}

\prose{“สุณาตุ เม ภนฺเต สํโฆ, อยํ เสยฺยสโก ภิกฺขุ สํเฆน นิยสฺสกมฺมกโต สมฺมา วตฺตติ, โลมํ ปาเตติ, เนตฺถารํ วตฺตติ, นิยสฺสสฺส กมฺมสฺส ปฏิปฺปสฺสทฺธิํ ยาจติ, ยทิ สํฆสฺส ปตฺตกลฺลํ, สํโฆ เสยฺยสกสฺส ภิกฺขุโน นิยสฺสกมฺมํ ปฏิปฺปสฺสมฺเภยฺย, เอสา ญตฺติ.}

\prose{\csromanfolio{22}สุณาตุ เม ภนฺเต สํโฆ, อยํ เสยฺยสโก ภิกฺขุ สํเฆน นิยสฺสกมฺมกโต สมฺมา วตฺตติ, โลมํ ปาเตติ, เนตฺถารํ วตฺตติ, นิยสฺสสฺส กมฺมสฺส ปฏิปฺปสฺสทฺธิํ ยาจติ, สํโฆ เสยฺยสกสฺส ภิกฺขุโน นิยสฺสกมฺมํ ปฏิปฺปสฺสมฺเภติ, ปสฺสายสฺมโต ขมติ เสยฺยสกสฺส ภิกฺขุโน นิยสฺสสฺส กมฺมสฺส ปฏิปฺปสฺสทฺธิ, โส ตุณฺหสฺส, ยสฺส นกฺขมติ, โส ภาเสยฺย.}

\prose{ทุติยมฺปิ เอตมตฺถํ วทามิ. สุณาตุ เม ภนฺเต สํโฆ, อยํ เสยฺยสโก ภิกฺขุ สํเฆน นิยสฺสกมฺมกโต สมฺมา วตฺตติ, โลมํ ปาเตติ, เนตฺถารํ วตฺตติ, นิยสฺสสฺส กมฺมสฺส ปฏิปฺปสฺสทฺธิํ ยาจติ, สํโฆ เสยฺยสกสฺส ภิกฺขุโน นิยสฺสกมฺมํ ปฏิปฺปสฺสมฺเภติ, ยสฺสายสฺมโต ขมติ เสยฺยสกสฺส ภิกฺขุโน นิยสฺสสฺส กมฺมสฺส ปฏิปฺปสฺสทฺธิ, โส ตุณฺหสฺส, ยสฺส นกฺขมติ, โส ภาเสยฺย.}

\prose{ตติยมฺปิ เอตมตฺถํ วทามิ. สุณาตุ เม ภนฺเต สํโฆ, อยํ เสยฺยสโก ภิกฺขุ สํเฆน นิยสฺสกมฺมกโต สมฺมา วตฺตติ, โลมํ ปาเตติ, เนตฺถารํ วตฺตติ, นิยสฺสสฺส กมฺมสฺส ปฏิปฺปสฺสทฺธิํ ยาจติ, สํโฆ เสยฺยสกสฺส ภิกฺขุโน นิยสฺสกมฺมํ ปฏิปฺปสฺสมฺเภติ, ยสฺสายสฺมโต ขมติ เสยฺยสกสฺส ภิกฺขุโน นิยสฺสสฺส กมฺมสฺส ปฏิปฺปสฺสทฺธิ, โส ตุณฺหสฺส, ยสฺส นกฺขมติ, โส ภาเสยฺย.}

\prose{ปฏิปฺปสฺสทฺธํ สํเฆน เสยฺยสกสฺส ภิกฺขุโน นิยสฺสกมฺมํ, ขมติ สํฆสฺส, ตสฺมา ตุณฺหี, เอวเมตํ ธารยามี”ติ.}

\vspace{\tipitakaparskip}
\csromancenter{นิยสฺสกมฺมํ นิฏฺฐิตํ ทุติยํ.}
\csromanmatikaanchor{mtk.16}
\csromantocmarkref{section}{3. ปพฺพาชนียกมฺม}{mtk.16}
\bookmark[dest={mtk.16},level=1]{3. ปพฺพาชนียกมฺม}
\csromanheader{1}{3. ปพฺพาชนีย\-กมฺม}
\proseitem{21}{\csromansymbolfootnote{*}{อิทํ วตฺถุ วิ 1. 274 ปิฏฺฐาทีสุปิ อาคตํ.}เตน โข ปน สมเยน อสฺสชิ\-ปุนพฺพสุกา นาม\csromansharedfootnote{e177e442d5af50ad}{นาม ภิกฺขู (ก)} กีฏาคิริสฺมิํ อาวาสิกา โหนฺติ อลชฺชิโน ปาปภิกฺขู, เต\csromansymbolfootnote{+}{วิ 4. 284 ปิฏฺฐาทีสุปิ.}เอวรูปํ อนาจารํ อาจรนฺติ, มาลาวจฺฉํ โรเปนฺติปิ โรปาเปนฺติปิ, สิญฺจนฺติปิ สิญฺจาเปนฺติปิ, โอจินนฺติปิ โอจินาเปนฺติปิ, คนฺเถนฺติปิ คนฺถาเปนฺติปิ, เอกโตวณฺฏิกมาลํ \csromanfolio{23}กโรนฺติปิ การาเปนฺติปิ, อุภโตวณฺฏิกมาลํ กโรนฺติปิ การาเปนฺติปิ, มญฺชริกํ กโรนฺติปิ การาเปนฺติปิ, วิธูติกํ\csromansharedfootnote{994634950fce6b19}{วิธุติกํ (สฺยา)} กโรนฺติปิ การาเปนฺติปิ, วฏํสกํ กโรนฺติปิ การาเปนฺติปิ, อาเวฬํ กโรนฺติปิ การาเปนฺติปิ, อุรจฺฉทํ กโรนฺติปิ การาเปนฺติปิ. เต กุลิตฺถีนํ กุลธีตานํ กุลกุมารีนํ กุลสุณฺหานํ กุลทาสีนํ เอกโตวณฺฏิกมาลํ หรนฺติปิ หราเปนฺติปิ, อุภโตวณฺฏิกมาลํ หรนฺติปิ หราเปนฺติปิ, มญฺชริกํ หรนฺติปิ หราเปนฺติปิ, วิธูติกํ หรนฺติปิ หราเปนฺติปิ, วฏํสกํ หรนฺติปิ หราเปนฺติปิ, อาเวฬํ หรนฺติปิ หราเปนฺติปิ, อุรจฺฉทํ หรนฺติปิ หราเปนฺติปิ. เต กุลิตฺถีหิ กุลธีตาหิ กุลกุมารีหิ กุลสุณฺหาหิ กุลทาสีหิ สทฺธิํ เอกภาชเนปิ ภุญฺชนฺติ, เอกถาลเกปิ ปิวนฺติ, เอกาสเนปิ นิสีทนฺติ, เอกมญฺเจปิ ตุวฏฺเฏนฺติ, เอกตฺถรณาปิ ตุวฏฺเฏนฺติ, เอกปาวุรณาปิ ตุวฏฺเฏนฺติ, เอก\-ตฺถรณ\-ปาวุรณาปิ ตุวฏฺเฏนฺติ, วิกาเลปิ ภุญฺชนฺติ, มชฺชมฺปิ ปิวนฺติ, มาลา\-คนฺธ\-วิเลปนมฺปิ ธาเรนฺติ, นจฺจนฺติปิ คายนฺติปิ วาเทนฺติปิ ลาเสนฺติปิ, นจฺจนฺติยาปิ นจฺจนฺติ นจฺจนฺติยาปิ คายนฺติ นจฺจนฺติยาปิ วาเทนฺติ นจฺจนฺติยาปิ ลาเสนฺติ, คายนฺติยาปิ นจฺจนฺติ คายนฺติยาปิ คายนฺติ คายนฺติยาปิ วาเทนฺติ คายนฺติยาปิ ลาเสนฺติ, วาเทนฺติยาปิ นจฺจนฺติ วาเทนฺติยาปิ คายนฺติ วาเทนฺติยาปิ วาเทนฺติ วาเทนฺติยาปิ ลาเสนฺติ, ลาเสนฺติยาปิ นจฺจนฺติ ลาเสนฺติยาปิ คายนฺติ ลาเสนฺติยาปิ วาเทนฺติ ลาเสนฺติยาปิ ลาเสนฺติ. อฏฺฐปเทปิ กีฬนฺติ, ทสปเทปิ กีฬนฺติ, อากาเสปิ กีฬนฺติ, ปริหาร\-ปเถปิ กีฬนฺติ, สนฺติกายปิ กีฬนฺติ, ขลิกายปิ กีฬนฺติ, ฆฏิกายปิ กีฬนฺติ, สลากหตฺเถนปิ กีฬนฺติ, อกฺเขนปิ กีฬนฺติ, ปงฺคจีเรนปิ กีฬนฺติ, วงฺกเกนปิ กีฬนฺติ, โมกฺขจิกายปิ กีฬนฺติ, จิงฺคุลเกนปิ กีฬนฺติ, ปตฺตาฬฺหเกนปิ กีฬนฺติ, รถเกนปิ กีฬนฺติ, ธนุเกนปิ กีฬนฺติ, อกฺขริกายปิ กีฬนฺติ, มเนสิกายปิ กีฬนฺติ, ยถาวชฺเชนปิ กีฬนฺติ, หตฺถิสฺมิมฺปิ สิกฺขนฺติ, อสฺสสฺมิมฺปิ สิกฺขนฺติ, รถสฺมิมฺปิ สิกฺขนฺติ, ธนุสฺมิมฺปิ สิกฺขนฺติ, ถรุสฺมิมฺปิ สิกฺขนฺติ, หตฺถิสฺสปิ ปุรโต ธาวนฺติ, อสฺสสฺสปิ ปุรโต ธาวนฺติ, รถสฺสปิ ปุรโต\csromansharedfootnote{3c8b67f5c4ccb131}{ปุรโต ธาวนฺติ (สฺยา)} ธาวนฺติปิ อาธาวนฺติปิ, อุสฺเสเฬนฺติปิ, อปฺโผเฏนฺติปิ, นิพฺพุชฺฌนฺติปิ, มุฏฺฐีหิปิ ยุชฺฌนฺติ, รงฺคมชฺเฌปิ สงฺฆาฏิํ ปตฺถริตฺวา นจฺจกิํ\csromansharedfootnote{f605615cc9978f89}{นจฺจนฺติํ (สี, สฺยา)} เอวํ วทนฺติ “อิธ ภคินิ นจฺจสฺสู”ติ, นลาฏิกมฺปิ เทนฺติ, วิวิธมฺปิ อนาจารํ อาจรนฺติ.}

\proseitem{22}{\csromanfolio{24}เตน โข ปน สมเยน อญฺญตโร ภิกฺขุ กาสีสุ วสฺสํวุฏฺโฐ\csromansharedfootnote{1331ed676cedffb6}{วสฺสํวุตฺโถ (สี, สฺยา)} สาวตฺถิํ คจฺฉนฺโต ภควนฺตํ ทสฺสนาย เยน กีฏาคิริ ตทวสริ. อถ โข โส ภิกฺขุ ปุพฺพณฺหสมยํ นิวาเสตฺวา ปตฺต\-จีวรมา\-ทาย กีฏาคิริํ ปิณฺฑาย ปาวิสิ ปาสาทิเกน อภิกฺกนฺเตน ปฏิกฺกนฺเตน อาโลกิเตน วิโลกิเตน สมิญฺชิเตน\csromansharedfootnote{dfd385415aa74a30}{สมฺมิญฺชิเตน (สี, สฺยา, กํ)} ปสาริเตน โอกฺขิตฺตจกฺขุ อิริยาปถ\-สมฺปนฺโน. มนุสฺสา ตํ ภิกฺขุํ ปสฺสิตฺวา เอวมาหํสุ “กฺวายํ อพลพโล วิย มนฺทมนฺโท วิย ภากุฏิก\-ภากุฏิโก วิย, โก อิมสฺส อุปคตสฺส ปิณฺฑกมฺปิ ทสฺสติ, อมฺหากํ ปน อยฺยา อสฺสชิ\-ปุนพฺพสุกา สณฺหา สขิลา สุขสมฺภาสา มิหิต\-ปุพฺพงฺคมา เอหิ\-สฺวาคต\-วาทิโน อพฺภากุฏิกา อุตฺตานมุขา ปุพฺพภาสิโน, เตสํ โข นาม ปิณฺโฑ ทาตพฺโพ”ติ.}

\prose{อทฺทสา โข อญฺญตโร อุปาสโก ตํ ภิกฺขุํ กีฏาคิริสฺมิํ ปิณฺฑาย จรนฺตํ, ทิสฺวาน เยน โส ภิกฺขุ เตนุปสงฺกมิ, อุปสงฺกมิตฺวา ตํ ภิกฺขุํ อภิวาเทตฺวา เอตทโวจ “อปิ ภนฺเต ปิณฺโฑ ลพฺภตี”ติ. น โข อาวุโส ปิณฺโฑ ลพฺภตีติ. เอหิ ภนฺเต, ฆรํ คมิสฺสามาติ. อถ โข โส อุปาสโก ตํ ภิกฺขุํ ฆรํ เนตฺวา โภเชตฺวา เอตทโวจ “กหํ ภนฺเต อยฺโย คมิสฺสตี”ติ. สาวตฺถิํ โข อหํ อาวุโส คมิสฺสามิ ภควนฺตํ ทสฺสนายาติ. เตน หิ ภนฺเต มม วจเนน ภควโต ปาเท สิรสา วนฺท, เอวญฺจ วเทหิ “ทุฏฺโฐ ภนฺเต กีฏาคิริสฺมิํ อาวาโส, อสฺสชิ\-ปุนพฺพสุกา นาม กีฏาคิริสฺมิํ อาวาสิกา อลชฺชิโน ปาปภิกฺขู, เต เอวรูปํ อนาจารํ อาจรนฺติ, มาลาวจฺฉํ โรเปนฺติปิ โรปาเปนฺติปิ, สิญฺจนฺติปิ สิญฺจาเปนฺติปิ, โอจินนฺติปิ โอจินาเปนฺติปิ, คนฺเถนฺติปิ คนฺถาเปนฺติปิ, เอกโตวณฺฏิกมาลํ กโรนฺติปิ การาเปนฺติปิ, อุภโตวณฺฏิกมาลํ กโรนฺติปิ การาเปนฺติปิ, มญฺชริกํ กโรนฺติปิ การาเปนฺติปิ, วิธูติกํ กโรนฺติปิ การาเปนฺติปิ, วฏํสกํ กโรนฺติปิ การาเปนฺติปิ, อาเวฬํ กโรนฺติปิ การาเปนฺติปิ, อุรจฺฉทํ กโรนฺติปิ การาเปนฺติปิ. เต กุลิตฺถีนํ กุลธีตานํ กุลกุมารีนํ กุลสุณฺหานํ กุลทาสีนํ เอกโตวณฺฏิกมาลํ หรนฺติปิ หราเปนฺติปิ, อุภโตวณฺฏิกมาลํ หรนฺติปิ หราเปนฺติปิ, มญฺชริกํ หรนฺติปิ หราเปนฺติปิ, \csromanfolio{25}วิธูติกํ หรนฺติปิ หราเปนฺติปิ, วฏํสกํ หรนฺติปิ หราเปนฺติปิ, อาเวฬํ หรนฺติปิ หราเปนฺติปิ, อุรจฺฉทํ หรนฺติปิ หราเปนฺติปิ. เต กุลิตฺถีหิ กุลธีตาหิ กุลกุมารีหิ กุลสุณฺหาหิ กุลทาสีหิ สทฺธิํ เอกภาชเนปิ ภุญฺชนฺติ, เอกถาลเกปิ ปิวนฺติ, เอกาสเนปิ นิสีทนฺติ, เอกมญฺเจปิ ตุวฏฺเฏนฺติ, เอกตฺถรณาปิ ตุวฏฺเฏนฺติ, เอกปาวุรณาปิ ตุวฏฺเฏนฺติ, เอก\-ตฺถรณ\-ปาวุรณาปิ ตุวฏฺเฏนฺติ, วิกาเลปิ ภุญฺชนฺติ, มชฺชมฺปิ ปิวนฺติ, มาลา\-คนฺธ\-วิเลปนมฺปิ ธาเรนฺติ, นจฺจนฺติปิ คายนฺติปิ วาเทนฺติปิ ลาเสนฺติปิ, นจฺจนฺติยาปิ นจฺจนฺติ นจฺจนฺติยาปิ คายนฺติ นจฺจนฺติยาปิ วาเทนฺติ นจฺจนฺติยาปิ ลาเสนฺติ ฯเปฯ (จกฺกํ กาตพฺพํ). ลาเสนฺติยาปิ นจฺจนฺติ ลาเสนฺติยาปิ คายนฺติ ลาเสนฺติยาปิ วาเทนฺติ ลาเสนฺติยาปิ ลาเสนฺติ, อฏฺฐปเทปิ กีฬนฺติ, ทสปเทปิ กีฬนฺติ, อากาเสปิ กีฬนฺติ, ปริหาร\-ปเถปิ กีฬนฺติ, สนฺติกายปิ กีฬนฺติ, ขลิกายปิ กีฬนฺติ, ฆฏิกายปิ กีฬนฺติ, สลากหตฺเถนปิ กีฬนฺติ, อกฺเขนปิ กีฬนฺติ, ปงฺคจีเรนปิ กีฬนฺติ, วงฺกเกนปิ กีฬนฺติ, โมกฺขจิกายปิ กีฬนฺติ, จิงฺคุลเกนปิ กีฬนฺติ, ปตฺตาฬฺหเกนปิ กีฬนฺติ, รถเกนปิ กีฬนฺติ, ธนุเกนปิ กีฬนฺติ, อกฺขริกายปิ กีฬนฺติ, มเนสิกายปิ กีฬนฺติ, ยถาวชฺเชนปิ กีฬนฺติ, หตฺถิสฺมิมฺปิ สิกฺขนฺติ, อสฺสสฺมิมฺปิ สิกฺขนฺติ, รถสฺมิมฺปิ สิกฺขนฺติ, ธนุสฺมิมฺปิ สิกฺขนฺติ, ถรุสฺมิมฺปิ สิกฺขนฺติ, หตฺถิสฺสปิ ปุรโต ธาวนฺติ, อสฺสสฺสปิ ปุรโต ธาวนฺติ, รถสฺสปิ ปุรโต ธาวนฺติปิ อาธาวนฺติปิ, อุสฺเสเฬนฺติปิ, อปฺโผเฏนฺติปิ, นิพฺพุชฺฌนฺติปิ, มุฏฺฐีหิปิ ยุชฺฌนฺติ, รงฺคมชฺเฌปิ สงฺฆาฏิํ ปตฺถริตฺวา นจฺจกิํ เอวํ วทนฺติ “อิธ ภคินิ นจฺจสฺสู”ติ, นลาฏิกมฺปิ เทนฺติ, วิวิธมฺปิ อนาจารํ อาจรนฺติ. เยปิ เต ภนฺเต มนุสฺสา ปุพฺเพ สทฺธา อเหสุํ ปสนฺนา, เตปิ เอตรหิ อสฺสทฺธา อปฺปสนฺนา, ยานิปิ ตานิ สํฆสฺส ปุพฺเพ ทานปถานิ, ตานิปิ เอตรหิ อุปจฺฉินฺนานิ, ริญฺจนฺติ เปสลา ภิกฺขู, นิวสนฺติ ปาปภิกฺขู, สาธุ ภนฺเต ภควา กีฏาคิริํ ภิกฺขู ปหิเณยฺย, ยถายํ กีฏาคิริสฺมิํ อาวาโส สณฺฐเหยฺยา”ติ.}

\prose{“เอวมาวุโส”ติ โข โส ภิกฺขุ ตสฺส อุปาสกสฺส ปฏิสฺสุณิตฺวา อุฏฺฐายาสนา เยน สาวตฺถิ เตน ปกฺกามิ, อนุปุพฺเพน เยน สาวตฺถิ เชตวนํ อนาถ\-ปิณฺฑิกสฺส อาราโม, เยน ภควา เตนุปสงฺกมิ, อุปสงฺกมิตฺวา ภควนฺตํ อภิวาเทตฺวา เอกมนฺตํ นิสีทิ, อาจิณฺณํ โข ปเนตํ พุทฺธานํ ภควนฺตานํ อาคนฺตุเกหิ ภิกฺขูหิ สทฺธิํ ปฏิสมฺโมทิตุํ. \csromanfolio{26}อถ โข ภควา ตํ ภิกฺขุํ เอตทโวจ “กจฺจิ ภิกฺขุ ขมนียํ, กจฺจิ ยาปนียํ, กจฺจิสิ อปฺป\-กิลมเถน อทฺธานํ อาคโต, กุโต จ ตฺวํ ภิกฺขุ อาคจฺฉสี”ติ. ขมนียํ ภควา, ยาปนียํ ภควา, อปฺป\-กิลมเถน จ อหํ ภนฺเต อทฺธานํ อาคโต, อิธาหํ ภนฺเต กาสีสุ วสฺสํวุฏฺโฐ สาวตฺถิํ อาคจฺฉนฺโต ภควนฺตํ ทสฺสนาย เยน กีฏาคิริ ตทวสริํ. อถ ขฺวาหํ ภนฺเต ปุพฺพณฺหสมยํ นิวาเสตฺวา ปตฺต\-จีวรมา\-ทาย กีฏาคิริํ ปิณฺฑาย ปาวิสิํ, อทฺทสา โข มํ ภนฺเต อญฺญตโร อุปาสโก กีฏาคิริสฺมิํ ปิณฺฑาย จรนฺตํ, ทิสฺวาน เยนาหํ เตนุปสงฺกมิ, อุปสงฺกมิตฺวา มํ อภิวาเทตฺวา เอตทโวจ “อปิ ภนฺเต ปิณฺโฑ ลพฺภตี”ติ, น โข อาวุโส ปิณฺโฑ ลพฺภตีติ, เอหิ ภนฺเต, ฆรํ คมิสฺสามาติ, อถ โข ภนฺเต โส อุปาสโก มํ ฆรํ เนตฺวา โภเชตฺวา เอตทโวจ “กหํ ภนฺเต อยฺโย คมิสฺสาตี”ติ, สาวตฺถิํ โข อหํ อาวุโส คมิสฺสามิ ภควนฺตํ ทสฺสนายาติ, เตน หิ ภนฺเต มม วจเนน ภควโต ปาเท สิรสา วนฺท, เอวญฺจ วเทหิ “ทุฏฺโฐ ภนฺเต กีฏาคิริสฺมิํ อาวาโส, อสฺสชิ\-ปุนพฺพสุกา นาม กีฏาคิริสฺมิํ อาวาสิกา อลชฺชิโน ปาปภิกฺขู, เต เอวรูปํ อนาจารํ อาจรนฺติ, มาลาวจฺฉํ โรเปนฺติปิ โรปาเปนฺติปิ ฯเปฯ วิวิธมฺปิ อนาจารํ อาจรนฺติ. เยปิ เต ภนฺเต มนุสฺสา ปุพฺเพ สทฺธา อเหสุํ ปสนฺนา, เตปิ เอตรหิ อสฺสทฺธา อปฺปสนฺนา, ยานิปิ ตานิ สํฆสฺส ปุพฺเพ ทานปถานิ, ตานิปิ เอตรหิ อุปจฺฉินฺนานิ, ริญฺจนฺติ เปสลา ภิกฺขู, นิวสนฺติ ปาปภิกฺขู, สาธุ ภนฺเต ภควา กีฏาคิริํ ภิกฺขู ปหิเณยฺย, ยถายํ กีฏาคิริสฺมิํ อาวาโส สณฺฐเหยฺยา”ติ, ตโต อหํ ภควา อาคจฺฉามีติ.}

\proseitem{23}{อถ โข ภควา เอตสฺมิํ นิทาเน เอตสฺมิํ ปกรเณ ภิกฺขุสํฆํ สนฺนิปาตาเปตฺวา ภิกฺขู ปฏิปุจฺฉิ “สจฺจํ กิร ภิกฺขเว อสฺสชิ\-ปุนพฺพสุกา นาม กีฏาคิริสฺมิํ อาวาสิกา อลชฺชิโน ปาปภิกฺขู, เต เอวรูปํ อนาจารํ อาจรนฺติ, มาลาวจฺฉํ โรเปนฺติปิ ฯเปฯ วิวิธมฺปิ อนาจารํ อาจรนฺติ. เยปิ เต มนุสฺสา ปุพฺเพ สทฺธา อเหสุํ ปสนฺนา, เตปิ เอตรหิ อสฺสทฺธา อปฺปสนฺนา, ยานิปิ ตานิ สํฆสฺส ปุพฺเพ ทานปถานิ, ตานิปิ เอตรหิ อุปจฺฉินฺนานิ, ริญฺจนฺติ เปสลา ภิกฺขู, นิวสนฺติ ปาปภิกฺขู”ติ. สจฺจํ ภควาติ. วิครหิ พุทฺโธ ภควา อนนุจฺฉวิกํ ฯเปฯ. กถํ หิ นาม เต \csromanfolio{27}ภิกฺขเว โมฆปุริสา เอวรุปํ อนาจารํ อาจริสฺสนฺติ, มาลาวจฺฉํ โรเปสฺสนฺติปิ โรปาเปสฺสนฺติปิ, สิญฺจิสฺสนฺติปิ สิญฺจาเปสฺสนฺติปิ, โอจินิสฺสนฺติปิ โอจินาเปสฺสนฺติปิ, คนฺเถสฺสนฺติปิ คนฺถาเปสฺสนฺติปิ, เอกโตวณฺฏิกมาลํ กริสฺสนฺติปิ การาเปสฺสนฺติปิ, อุภโตวณฺฏิกมาลํ กริสฺสนฺติปิ การาเปสฺสนฺติปิ, มญฺชริกํ กริสฺสนฺติปิ การาเปสฺสนฺติปิ, วิธูติกํ กริสฺสนฺติปิ การาเปสฺสนฺติปิ, วฏํสกํ กริสฺสนฺติปิ การาเปสฺสนฺติปิ, อาเวฬํ กริสฺสนฺติปิ การาเปสฺสนฺติปิ, อุรจฺฉทํ กริสฺสนฺติปิ การาเปสฺสนฺติปิ. เต กุลิตฺถีนํ กุลธีตานํ กุลกุมารีนํ กุลสุณฺหานํ กุลทาสีนํ เอกโตวณฺฏิกมาลํ หริสฺสนฺติปิ หราเปสฺสนฺติปิ, อุภโตวณฺฏิกมาลํ หริสฺสนฺติปิ หราเปสฺสนฺติปิ, มญฺชริกํ หริสฺสนฺติปิ หราเปสฺสนฺติปิ, วิธูติกํ หริสฺสนฺติปิ หราเปสฺสนฺติปิ, วฏํสกํ หริสฺสนฺติปิ หราเปสฺสนฺติปิ, อาเวฬํ หริสฺสนฺติปิ หราเปสฺสนฺติปิ, อุรจฺฉทํ หริสฺสนฺติปิ หราเปสฺสนฺติปิ. เต กุลิตฺถีหิ กุลธีตาหิ กุลกุมารีหิ กุลสุณฺหาหิ กุลทาสีหิ สทฺธิํ เอกภาชเนปิ ภุญฺชิสฺสนฺติ, เอกถาลเกปิ ปิวิสฺสนฺติ, เอกาสเนปิ นิสีทิสฺสนฺติ, เอกมญฺเจปิ ตุวฏฺฏิสฺสนฺติ, เอกตฺถรณาปิ ตุวฏฺฏิสฺสนฺติ, เอกปาวุรณาปิ ตุวฏฺฏิสฺสนฺติ, เอก\-ตฺถรณ\-ปาวุรณาปิ ตุวฏฺฏิสฺสนฺติ, วิกาเลปิ ภุญฺชิสฺสนฺติ, มชฺชมฺปิ ปิวิสฺสนฺติ, มาลา\-คนฺธ\-วิเลปนมฺปิ ธาเรสฺสนฺติ, นจฺจิสฺสนฺติปิ คายิสฺสนฺติปิ วาเทสฺสนฺติปิ ลาเสสฺสนฺติปิ, นจฺจนฺติยาปิ นจฺจิสฺสนฺติ นจฺจนฺติยาปิ คายิสฺสนฺติ นจฺจนฺติยาปิ วาเทสฺสนฺติ นจฺจนฺติยาปิ ลาเสสฺสนฺติ, คายนฺติยาปิ นจฺจิสฺสนฺติ คายนฺติยาปิ คายิสฺสนฺติ คายนฺติยาปิ วาเทสฺสนฺติ คายนฺติยาปิ ลาเสสฺสนฺติ, วาเทนฺติยาปิ นจฺจิสฺสนฺติ วาเทนฺติยาปิ คายิสฺสนฺติ วาเทนฺติยาปิ วาเทสฺสนฺติ วาเทนฺติยาปิ ลาเสสฺสนฺติ, ลาเสนฺติยาปิ นจฺจิสฺสนฺติ ลาเสนฺติยาปิ คายิสฺสนฺติ, ลาเสนฺติยาปิ วาเทสฺสนฺติ, ลาเสนฺติยาปิ ลาเสสฺสนฺติ, อฏฺฐปเทปิ กีฬิสฺสนฺติ, ทสปเทปิ กีฬิสฺสนฺติ, อากาเสปิ กีฬิสฺสนฺติ, ปริหาร\-ปเถปิ กีฬิสฺสนฺติ, สนฺติกายปิ กีฬิสฺสนฺติ, ขลิกายปิ กีฬิสฺสนฺติ, ฆฏิกายปิ กีฬิสฺสนฺติ, สลากหตฺเถนปิ กีฬิสฺสนฺติ, อกฺเขนปิ กีฬิสฺสนฺติ, ปงฺคจีเรนปิ กีฬิสฺสนฺติ, วงฺกเกนปิ กีฬิสฺสนฺติ, โมกฺขจิกายปิ กีฬิสฺสนฺติ, จิงฺคุลเกนปิ กีฬิสฺสนฺติ, ปตฺตาฬฺหเกนปิ กีฬิสฺสนฺติ, รถเกนปิ กีฬิสฺสนฺติ, ธนุเกนปิ กีฬิสฺสนฺติ, อกฺขริกายปิ กีฬิสฺสนฺติ, มเนสิกายปิ กีฬิสฺสนฺติ, ยถาวชฺเชนปิ กีฬิสฺสนฺติ, หตฺถิสฺมิมฺปิ สิกฺขิสฺสนฺติ, อสฺสสฺมิมฺปิ สิกฺขิสฺสนฺติ, รถสฺมิมฺปิ สิกฺขิสฺสนฺติ, ธนุสฺมิมฺปิ สิกฺขิสฺสนฺติ, ถรุสฺมิมฺปิ สิกฺขิสฺสนฺติ, หตฺถิสฺสปิ ปุรโต \csromanfolio{28}ธาวิสฺสนฺติ, อสฺสสฺสปิ ปุรโต มาวิสฺสนฺติ, รถสฺสปิ ปุรโต\csromansharedfootnote{5947dce77d2f62c8}{ปุรโต ธาวิสฺสนฺติ (สฺยา)} ธาวิสฺสนฺติปิ อาธาวิสฺสนฺติปิ, อุสฺเสเฬสฺสนฺติปิ, อปฺโผเฏสฺสนฺติปิ, นิพฺพุชฺฌิสฺสนฺติปิ, มุฏฺฐีหิปิ ยุชฺฌิสฺสนฺติ, รงฺคมชฺเฌปิ สงฺฆาฏิํ ปตฺถริตฺวา นจฺจกิํ เอวํ วกฺขนฺติ\csromansharedfootnote{014df73c22346d6b}{วทิสฺสนฺติ (ก)} “อิธ ภคินิ นจฺจสฺสู”ติ, นลาฏิกมฺปิ ทสฺสนฺติ, วิวิธมฺปิ อนาจารํ อาจาริสฺสนฺติ, เนตํ ภิกฺขเว อปฺปสนฺนานํ วา ปสาทาย ฯเปฯ วิครหิตฺวา ฯเปฯ ธมฺมิํ กถํ กตฺวา สาริปุตฺต\-โมคฺคลฺลาเน อามนฺเตสิ “คจฺฉถ ตุเมฺห สาริปุตฺตา, กีฏาคิริํ คนฺตฺวา อสฺสชิ\-ปุนพฺพสุกานํ ภิกฺขูนํ กีฏาคิริสฺมา ปพฺพาชนีย\-กมฺมํ กโรถ, ตุมฺหากํ เอเต สทฺธิวิหาริโน”ติ.}

\prose{กถํ มยํ ภนฺเต อสฺสชิ\-ปุนพฺพสุกานํ ภิกฺขูนํ กีฏาคิริสฺมา ปพฺพาชนีย\-กมฺมํ กโรม, จณฺฑา เต ภิกฺขู ผรุสาติ. เตน หิ ตุเมฺห สาริปุตฺตา พหุเกหิ ภิกฺขูหิ สทฺธิํ คจฺฉถาติ. “เอวํ ภนฺเต”ติ โข สาริปุตฺต\-โมคฺคลฺลานา ภควโต ปจฺจสฺโสสุํ. เอวญฺจ ปน ภิกฺขเว กาตพฺพํ, ปฐมํ อสฺสชิ\-ปุนพฺพสุกา ภิกฺขู โจเทตพฺพา, โจเทตฺวา สาเรตพฺพา, สาเรตฺวา อาปตฺติํ อาโรเปตพฺพา, อาปตฺติํ อาโรเปตฺวา พฺยตฺเตน ภิกฺขุนา ปฏิพเลน สํโฆ ญาเปตพฺโพ\csromandash{}}

\proseitem{24}{“สุณาตุ เม ภนฺเต สํโฆ, อิเม อสฺสชิ\-ปุนพฺพสุกา ภิกฺขู กุลทูสกา ปาปสมาจารา, อิเมสํ ปาปกา สมาจารา ทิสฺสนฺติ เจว สุยฺยนฺติ จ, กุลานิ จ อิเมหิ ทุฏฺฐานิ ทิสฺสนฺติ เจว สุยฺยนฺติ จ, ยทิ สํฆสฺส ปตฺตกลฺลํ, สํโฆ อสฺสชิ\-ปุนพฺพสุกานํ ภิกฺขูนํ กีฏาคิริสฺมา ปพฺพาชนีย\-กมฺมํ กเรยฺย ‘น อสฺสชิปุนพฺพสุเกหิ ภิกฺขูหิ กีฏาคิริสฺมิํ วตฺถพฺพนฺ’ติ, เอสา ญตฺติ.}

\prose{สุณาตุ เม ภนฺเต สํโฆ, อิเม อสฺสชิ\-ปุนพฺพสุกา ภิกฺขู กุลทูสกา ปาปสมาจารา, อิเมสํ ปาปกา สมาจารา ทิสฺสนฺติ เจว สุยฺยนฺติ จ, กุลานิ จ อิเมหิ ทุฏฺฐานิ ทิสฺสนฺติ เจว สุยฺยนฺติ จ, สํโฆ อสฺสชิ\-ปุนพฺพสุกานํ ภิกฺขูนํ กีฏาคิริสฺมา ปพฺพาชนีย\-กมฺมํ กโรติ ‘น อสฺสชิปุนพฺพสุเกหิ ภิกฺขูหิ กีฏาคิริสฺมิํ วตฺถพฺพนฺ’ติ, ยสฺสายสฺมโต ขมติ อสฺสชิ\-ปุนพฺพสุกานํ ภิกฺขูนํ กีฏาคิริสฺมา ปพฺพาชนียสฺส กมฺมสฺส กรณํ ‘น อสฺสชิปุนพฺพสุเกหิ ภิกฺขูหิ กีฏาคิริสฺมิํ วตฺถพฺพนฺ’ติ, โส ตุณฺหสฺส, ยสฺส นกฺขมติ, โส ภาเสยฺย.}

\prose{\csromanfolio{29}ทุติยมฺปิ เอตมตฺถํ วทามิ ฯเปฯ. ตติยมฺปิ เอตมตฺถํ วทามิ. สุณาตุ เม ภนฺเต สํโฆ, อิเม อสฺสชิ\-ปุนพฺพสุกา ภิกฺขู กุลทูสกา ปาปสมาจารา, อิเมสํ ปาปกา สมาจารา ทิสฺสนฺติ เจว สุยฺยนฺติ จ, กุลานิ จ อิเมหิ ทุฏฺฐานิ ทิสฺสนฺติ เจว สุยฺยนฺติ จ, สํโฆ อสฺสชิ\-ปุนพฺพสุกานํ ภิกฺขูนํ กีฏาคิริสฺมา ปพฺพาชนีย\-กมฺมํ กโรติ ‘น อสฺสชิปุนพฺพสุเกหิ ภิกฺขูหิ กีฏาคิริสฺมิํ วตฺตพฺพนฺ’ติ, ยสฺสายสฺมโต ขมติ อสฺสชิ\-ปุนพฺพสุกานํ ภิกฺขูนํ กีฏาคิริสฺมา ปพฺพาชนียสฺส กมฺมสฺส กรณํ ‘น อสฺสชิปุนพฺพสุเกหิ ภิกฺขูหิ กีฏาคิริสฺมิํ วตฺตพฺพนฺ’ติ, โส ตุณฺหสฺส, ยสฺส นกฺขมติ, โส ภาเสยฺย.}

\prose{กตํ สํเฆน อสฺสชิ\-ปุนพฺพสุกานํ ภิกฺขูนํ กีฏาคิริสฺมา ปพฺพาชนีย\-กมฺมํ ‘น อสฺสชิปุนพฺพสุเกหิ ภิกฺขูหิ กีฏาคิริสฺมิํ วตฺตพฺพนฺ’ติ, ขมติ สํฆสฺส, ตสฺมา ตุณฺหี, เอวเมตํ ธารยามี”ติ.}

\csromanmatikaanchor{mtk.17}
\csromantocmarkref{subsection}{อธมฺมกมฺมทฺวาทสก}{mtk.17}
\bookmark[dest={mtk.17},level=2]{อธมฺมกมฺมทฺวาทสก}
\csromanheader{2}{อธมฺมกมฺม\-ทฺวาทสก}
\proseitem{25}{ตีหิ ภิกฺขเว องฺเคหิ สมนฺนาคตํ ปพฺพาชนีย\-กมฺมํ อธมฺมกมฺมญฺจ โหติ อวินยกมฺมญฺจ ทุวูปสนฺตญฺจ. อสมฺมุขา กตํ โหติ, อปฺปฏิปุจฺฉา กตํ โหติ, อปฺปฏิญฺญาย กตํ โหติ. อิเมหิ โข ภิกฺขเว ตีหงฺเคหิ สมนฺนาคตํ ปพฺพาชนีย\-กมฺมํ อธมฺมกมฺมญฺจ โหติ อวินยกมฺมญฺจ ทุวูปสนฺตญฺจ.}

\prosewithcloser{อปเรหิปิ ภิกฺขเว ตีหงฺเคหิ สมนฺนาคตํ ปพฺพาชนีย\-กมฺมํ อธมฺมกมฺมญฺจ โหติ อวินยกมฺมญฺจ ทุวูปสนฺตญฺจ. อนาปตฺติยา กตํ โหติ, อเทสนาคามินิยา อาปตฺติยา กตํ โหติ, เทสิตาย อาปตฺติยา กตํ โหติ ฯเปฯ. อโจเทตฺวา กตํ โหติ, อสาเรตฺวา กตํ โหติ, อาปตฺติํ อนาโรเปตฺวา กตํ โหติ ฯเปฯ. อสมฺมุขา กตํ โหติ, อธมฺเมน กตํ โหติ, วคฺเคน กตํ โหติ ฯเปฯ. อปฺปฏิปุจฺฉา กตํ โหติ, อธมฺเมน กตํ โหติ, วคฺเคน กตํ โหติ ฯเปฯ. อปฺปฏิญฺญาย กตํ โหติ, อธมฺเมน กตํ โหติ, วคฺเคน กตํ โหติ ฯเปฯ. อนาปตฺติยา กตํ โหติ, อธมฺเมน กตํ โหติ, วคฺเคน กตํ โหติ ฯเปฯ. อเทสนาคามินิยา อาปตฺติยา กตํ โหติ, อธมฺเมน กตํ โหติ, วคฺเคน กตํ โหติ ฯเปฯ. เทสิตาย อาปตฺติยา กตํ โหติ, อธมฺเมน กตํ โหติ, วคฺเคน กตํ โหติ ฯเปฯ. อโจเทตฺวา กตํ โหติ, อธมฺเมน กตํ โหติ, วคฺเคน กตํ \csromanfolio{30}โหติ ฯเปฯ. อสาเรตฺวา กตํ โหติ, อธมฺเมน กตํ โหติ, วคฺเคน กตํ โหติ ฯเปฯ. อาปตฺติํ อนาโรเปตฺวา กตํ โหติ, อธมฺเมน กตํ โหติ, วคฺเคน กตํ โหติ. อิเมหิ โข ภิกฺขเว ตีหงฺเคหิ สมนฺนาคตํ ปพฺพาชนีย\-กมฺมํ อธมฺมกมฺมญฺจ โหติ อวินยกมฺมญฺจ ทุวูปสนฺตญฺจ.}{อธมฺมกมฺม\-ทฺวาทสกํ นิฏฺฐิตํ.}
\csromanmatikaanchor{mtk.18}
\csromantocmarkref{subsection}{ธมฺมกมฺมทฺวาทสก}{mtk.18}
\bookmark[dest={mtk.18},level=2]{ธมฺมกมฺมทฺวาทสก}
\csromanheader{2}{ธมฺมกมฺม\-ทฺวาทสก}
\proseitem{26}{ตีหิ ภิกฺขเว องฺเคหิ สมนฺนาคตํ ปพฺพาชนีย\-กมฺมํ ธมฺมกมฺมญฺจ โหติ วินยกมฺมญฺจ สุวูปสนฺตญฺจ. สมฺมุขา กตํ โหติ, ปฏิปุจฺฉา กตํ โหติ, ปฏิญฺญาย กตํ โหติ. อิเมหิ โข ภิกฺขเว ตีหงฺเคหิ สมนฺนาคตํ ปพฺพาชนีย\-กมฺมํ ธมฺมกมฺมญฺจ โหติ วินยกมฺมญฺจ สุวูปสนฺตญฺจ.}

\prosewithcloser{อปเรหิปิ ภิกฺขเว ตีหงฺเคหิ สมนฺนาคตํ ปพฺพาชนีย\-กมฺมํ ธมฺมกมฺมญฺจ โหติ วินยกมฺมญฺจ สุวูปสนฺตญฺจ. อาปตฺติยา กตํ โหติ, เทสนาคามินิยา อาปตฺติยา กตํ โหติ, อเทสิตาย อาปตฺติยา กตํ โหติ ฯเปฯ. โจเทตฺวา กตํ โหติ, สาเรตฺวา กตํ โหติ, อาปตฺติํ อาโรเปตฺวา กตํ โหติ ฯเปฯ. สมฺมุขา กตํ โหติ, ธมฺเมน กตํ โหติ, สมคฺเคน กตํ โหติ ฯเปฯ. ปฏิปุจฺฉา กตํ โหติ, ธมฺเมน กตํ โหติ, สมคฺเคน กตํ โหติ ฯเปฯ. ปฏิญฺญาย กตํ โหติ, ธมฺเมน กตํ โหติ, สมคฺเคน กตํ โหติ ฯเปฯ. อาปตฺติยา กตํ โหติ, ธมฺเมน กตํ โหติ, สมคฺเคน กตํ โหติ ฯเปฯ. เทสนาคามินิยา อาปตฺติยา กตํ โหติ, ธมฺเมน กตํ โหติ, สมคฺเคน กตํ โหติ ฯเปฯ. อเทสิตาย อาปตฺติยา กตํ โหติ, ธมฺเมน กตํ โหติ, สมคฺเคน กตํ โหติ ฯเปฯ. โจเทตฺวา กตํ โหติ, ธมฺเมน กตํ โหติ, สมคฺเคน กตํ โหติ ฯเปฯ. สาเรตฺวา กตํ โหติ, ธมฺเมน กตํ โหติ, สมคฺเคน กตํ โหติ ฯเปฯ. อาปตฺติํ อาโรเปตฺวา กตํ โหติ, ธมฺเมน \csromanfolio{31}กตํ โหติ, สมคฺเคน กตํ โหติ. อิเมหิ โข ภิกฺขเว ตีหงฺเคหิ สมนฺนาคตํ ปพฺพาชนีย\-กมฺมํ ธมฺมกมฺมญฺจ โหติ วินยกมฺมญฺจ สุวูปสนฺตญฺจ.}{ธมฺมกมฺม\-ทฺวาทสกํ นิฏฺฐิตํ.}
\csromanmatikaanchor{mtk.19}
\csromantocmarkref{subsection}{อากงฺขมานจุทฺทสก}{mtk.19}
\bookmark[dest={mtk.19},level=2]{อากงฺขมานจุทฺทสก}
\csromanheader{2}{อากงฺขมาน\-จุทฺทสก}
\proseitem{27}{\csromansymbolfootnote{*}{วิ 5. 219 ปิฏฺเฐปิ.}ตีหิ ภิกฺขเว องฺเคหิ สมนฺนาคตสฺส ภิกฺขุโน อากงฺขมาโน สํโฆ ปพฺพาชนีย\-กมฺมํ กเรยฺย. ภณฺฑนการโก โหติ กลหการโก วิวาทการโก ภสฺสการโก สํเฆ อธิกรณ\-การโก, พาโล โหติ อพฺยตฺโต อาปตฺติพหุโล อนปทาโน, คิหิสํสฏฺโฐ วิหรติ อนนุโลมิเกหิ คิหิสํสคฺเคหิ. อิเมหิ โข ภิกฺขเว ตีหงฺเคหิ สมนฺนาคตสฺส ภิกฺขุโน อากงฺขมาโน สํโฆ ปพฺพาชนีย\-กมฺมํ กเรยฺย. (1)}

\prose{อปเรหิปิ ภิกฺขเว ตีหงฺเคหิ สมนฺนาคตสฺส ภิกฺขุโน อากงฺขมาโน สํโฆ ปพฺพาชนีย\-กมฺมํ กเรยฺย. อธิสีเล สีลวิปนฺโน โหติ, อชฺฌาจาเร อาจารวิปนฺโน โหติ, อติทิฏฺฐิยา ทิฏฺฐิวิปนฺโน โหติ. อิเมหิ โข ภิกฺขเว ตีหงฺเคหิ สมนฺนาคตสฺส ภิกฺขุโน อากงฺขมาโน สํโฆ ปพฺพาชนีย\-กมฺมํ กเรยฺย. (2)}

\prose{อปเรหิปิ ภิกฺขเว ตีหงฺเคหิ สมนฺนาคตสฺส ภิกฺขุโน อากงฺขมาโน สํโฆ ปพฺพาชนีย\-กมฺมํ กเรยฺย. พุทฺธสฺส อวณฺณํ ภาสติ, ธมฺมสฺส อวณฺณํ ภาสติ, สํฆสฺส อวณฺณํ ภาสติ. อิเมหิ โข ภิกฺขเว ตีหงฺเคหิ สมนฺนาคตสฺส ภิกฺขุโน อากงฺขมาโน สํโฆ ปพฺพาชนีย\-กมฺมํ กเรยฺย. (3)}

\prose{อปเรหิปิ ภิกฺขเว ตีหงฺเคหิ สมนฺนาคตสฺส ภิกฺขุโน อากงฺขมาโน สํโฆ ปพฺพาชนีย\-กมฺมํ กเรยฺย. กายิเกน ทเวน สมนฺนาคโต โหติ, วาจสิเกน ทเวน สมนฺนาคโต โหติ, กายิก\-วาจสิเกน ทเวน สมนฺนาคโต โหติ. อิเมหิ โข ภิกฺขเว ตีหงฺเคหิ สมนฺนาคตสฺส ภิกฺขุโน อากงฺขมาโน สํโฆ ปพฺพาชนีย\-กมฺมํ กเรยฺย. (4)}

\prose{\csromanfolio{32}อปเรหิปิ ภิกฺขเว ตีหงฺเคหิ สมนฺนาคตสฺส ภิกฺขุโน อากงฺขมาโน สํโฆ ปพฺพาชนีย\-กมฺมํ กเรยฺย. กายิเกน อนาจาเรน สมนฺนาคโต โหติ, วาจสิเกน อนาจาเรน สมนฺนาคโต โหติ, กายิก\-วาจสิเกน อนาจาเรน สมนฺนาคโต โหติ. อิเมหิ โข ภิกฺขเว ตีหงฺเคหิ สมนฺนาคตสฺส ภิกฺขุโน อากงฺขมาโน สํโฆ ปพฺพาชนีย\-กมฺมํ กเรยฺย. (5)}

\prose{อปเรหิปิ ภิกฺขเว ตีหงฺเคหิ สมนฺนาคตสฺส ภิกฺขุโน อากงฺขมาโน สํโฆ ปพฺพาชนีย\-กมฺมํ กเรยฺย. กายิเกน อุปฆาติเกน สมนฺนาคโต โหติ, วาจสิเกน อุปฆาติเกน สมนฺนาคโต โหติ, กายิก\-วาจสิเกน อุปฆาติเกน สมนฺนาคโต โหติ. อิเมหิ โข ภิกฺขเว ตีหงฺเคหิ สมนฺนาคตสฺส ภิกฺขุโน อากงฺขมาโน สํโฆ ปพฺพาชนีย\-กมฺมํ กเรยฺย. (6)}

\prose{อปเรหิปิ ภิกฺขเว ตีหงฺเคหิ สมนฺนาคตสฺส ภิกฺขุโน อากงฺขมาโน สํโฆ ปพฺพาชนีย\-กมฺมํ กเรยฺย. กายิเกน มิจฺฉาชีเวน สมนฺนาคโต โหติ, วาจสิเกน มิจฺฉาชีเวน สมนฺนาคโต โหติ, กายิก\-วาจสิเกน มิจฺฉาชีเวน สมนฺนาคโต โหติ. อิเมหิ โข ภิกฺขเว ตีหงฺเคหิ สมนฺนาคตสฺส ภิกฺขุโน อากงฺขมาโน สํโฆ ปพฺพาชนีย\-กมฺมํ กเรยฺย. (7)}

\prose{ติณฺณํ ภิกฺขเว ภิกฺขูนํ อากงฺขมาโน สํโฆ ปพฺพาชนีย\-กมฺมํ กเรยฺย. เอโก ภณฺฑนการโก โหติ กลหการโก วิวาทการโก ภสฺสการโก สํเฆ อธิกรณ\-การโก, เอโก พาโล โหติ อพฺยตฺโต อาปตฺติพหุโล อนปทาโน, เอโก คิหิสํสฏฺโฐ วิหรติ อนนุโลมิเกหิ คิหิสํสคฺเคหิ. อิเมสํ โข ภิกฺขเว ติณฺณํ ภิกฺขูนํ อากงฺขมาโน สํโฆ ปพฺพาชนีย\-กมฺมํ กเรยฺย. (8)}

\prose{อปเรสมฺปิ ภิกฺขเว ติณฺณํ ภิกฺขูนํ อากงฺขมาโน สํโฆ ปพฺพาชนีย\-กมฺมํ กเรยฺย. เอโก อธิสีเล สีลวิปนฺโน โหติ, เอโก อชฺฌาจาเร อาจารวิปนฺโน โหติ, เอโก อติทิฏฺฐิยา ทิฏฺฐิวิปนฺโน โหติ. อิเมสํ โข ภิกฺขเว ติณฺณํ ภิกฺขูนํ อากงฺขมาโน สํโฆ ปพฺพาชนีย\-กมฺมํ กเรยฺย. (9)}

\prose{\csromanfolio{33}อปเรสมฺปิ ภิกฺขเว ติณฺณํ ภิกฺขูนํ อากงฺขมาโน สํโฆ ปพฺพาชนีย\-กมฺมํ กเรยฺย. เอโก พุทฺธสฺส อวณฺณํ ภาสติ, เอโก ธมฺมสฺส อวณฺณํ ภาสติ, เอโก สํฆสฺส อวณฺณํ ภาสติ. อิเมสํ โข ภิกฺขเว ติณฺณํ ภิกฺขูนํ อากงฺขมาโน สํโฆ ปพฺพาชนีย\-กมฺมํ กเรยฺย. (10)}

\prose{อปเรสมฺปิ ภิกฺขเว ติณฺณํ ภิกฺขูนํ อากงฺขมาโน สํโฆ ปพฺพาชนีย\-กมฺมํ กเรยฺย. เอโก กายิเกน ทเวน สมนฺนาคโต โหติ, เอโก วาจสิเกน ทเวน สมนฺนาคโต โหติ, เอโก กายิก\-วาจสิเกน ทเวน สมนฺนาคโต โหติ. อิเมสํ โข ภิกฺขเว ติณฺณํ ภิกฺขูนํ อากงฺขมาโน สํโฆ ปพฺพาชนีย\-กมฺมํ กเรยฺย. (11)}

\prose{อปเรสมฺปิ ภิกฺขเว ติณฺณํ ภิกฺขูนํ อากงฺขมาโน สํโฆ ปพฺพาชนีย\-กมฺมํ กเรยฺย. เอโก กายิเกน อนาจาเรน สมนฺนาคโต โหติ, เอโก วาจสิเกน อนาจาเรน สมนฺนาคโต โหติ, เอโก กายิก\-วาจสิเกน อนาจาเรน สมนฺนาคโต โหติ. อิเมสํ โข ภิกฺขเว ติณฺณํ ภิกฺขูนํ อากงฺขมาโน สํโฆ ปพฺพาชนีย\-กมฺมํ กเรยฺย. (12)}

\prose{อปเรสมฺปิ ภิกฺขเว ติณฺณํ ภิกฺขูนํ อากงฺขมาโน สํโฆ ปพฺพาชนีย\-กมฺมํ กเรยฺย. เอโก กายิเกน อุปฆาติเกน สมนฺนาคโต โหติ, เอโก วาจสิเกน อุปฆาติเกน สมนฺนาคโต โหติ. อิเมสํ โข ภิกฺขเว ติณฺณํ ภิกฺขูนํ อากงฺขมาโน สํโฆ ปพฺพาชนีย\-กมฺมํ กเรยฺย. (13)}

\prosewithcloser{อปเรสมฺปิ ภิกฺขเว ติณฺณํ ภิกฺขูนํ อากงฺขมาโน สํโฆ ปพฺพาชนิย\-กมฺมํ กเรยฺย. เอโก กายิเกน มิจฺฉาชีเวน สมนฺนาคโต โหติ, เอโก วาจสิเกน มิจฺฉาชีเวน สมนฺนาคโต โหติ, เอโก กายิก\-วาจสิเกน มิจฺฉาชีเวน สมนฺนาคโต โหติ. อิเมสํ โข ภิกฺขเว ติณฺณํ ภิกฺขูนํ อากงฺขมาโน สํโฆ ปพฺพาชนีย\-กมฺมํ กเรยฺย. (14)}{อากงฺขมาน\-จุทฺทสกํ นิฏฺฐิตํ.}
\csromanmatikaanchor{mtk.20}
\csromantocmarkref{subsection}{อฏฺฐารสวตฺต}{mtk.20}
\bookmark[dest={mtk.20},level=2]{อฏฺฐารสวตฺต}
\csromanheader{2}{\textbf{อฏฺฐารสวตฺต}}
\proseitemwithcloser{28}{\csromanfolio{34}ปพฺพาชนีย\-กมฺมกเตน ภิกฺขเว ภิกฺขุนา สมฺมา วตฺติตพฺพํ. ตตฺรายํ สมฺมาวตฺตนา\csromandash{}น อุปสมฺ\-ปาเท\-ตพฺพํ, น นิสฺสโย ทาตพฺโพ, น สามเณโร อุปฏฺฐาเปตพฺโพ, น ภิกฺขุโนวาทก\-สมฺมุติ สาทิตพฺพา, สมฺมเตนปิ ภิกฺขุนิโย น โอวทิตพฺพา, ยาย อาปตฺติยา สํเฆน ปพฺพาชนีย\-กมฺมํ กตํ โหติ, สา อาปตฺติ น อาปชฺชิตพฺพา, อญฺญา วา ตาทิสิกา, ตโต วา ปาปิฏฺฐตรา, กมฺมํ น ครหิตพฺพํ, กมฺมิกา น ครหิตพฺพา, น ปกตตฺตสฺส ภิกฺขุโน อุโปสโถ ฐเปตพฺโพ, น ปวารณา ฐเปตพฺพา, น สวจนียํ กาตพฺพํ, น อนุวาโท อฏฺฐเปตพฺโพ, น โอกาโส กาเรตพฺโพ, น โจเทตพฺโพ, น สาเรตพฺโพ, น ภิกฺขูหิ สมฺป\-โยเช\-ตพฺพนฺติ.}{ปพฺพาชนีย\-กมฺเม อฏฺฐารส\-วตฺตํ นิฏฺฐิตํ.}
\proseitem{29}{อถ โข สาริปุตฺต\-โมคฺคลฺลาน\-ปฺปมุโข ภิกฺขุสํโฆ กีฏาคิริํ คนฺตฺวา อสฺสชิ\-ปุนพฺพสุกานํ ภิกฺขูนํ กีฏาคิริสฺมา ปพฺพาชนีย\-กมฺมํ อกาสิ “น อสฺสชิปุนพฺพสุเกหิ ภิกฺขูหิ กีฏาคิริสฺมิํ วตฺถพฺพนฺ”ติ. เต สํเฆน ปพฺพาชนีย\-กมฺมกตา น สมฺมา วตฺตนฺติ, น โลมํ ปาเตนฺติ, น เนตฺถารํ วตฺตนฺติ, น ภิกฺขู ขมาเปนฺติ, อกฺโกสนฺติ, ปริภาสนฺติ, ฉนฺทคามิตา โทสคามิตา โมหคามิตา ภยคามิตา ปาเปนฺติ, ปกฺกมนฺติปิ วิพฺภมนฺติปิ. เย เต ภิกฺขู อปฺปิจฺฉา ฯเปฯ เต อุชฺฌายนฺติ ขิยฺยนฺติ วิปาเจนฺติ “กถํ หิ นาม อสฺสชิ\-ปุนพฺพสุกา ภิกฺขู สํเฆน ปพฺพาชนีย\-กมฺมกตา น สมฺมา วตฺติสฺสนฺติ, น โลมํ ปาเตสฺสนฺติ, น เนตฺถารํ วตฺติสฺสนฺติ, น ภิกฺขู ขมาเปสฺสนฺติ, อกฺโกสิสฺสนฺติ, ปริภาสิสฺสนฺติ, ฉนฺทคามิตา โทสคามิตา โมหคามิตา ภยคามิตา ปาเปสฺสนฺติ, ปกฺกมิสฺสนฺติปิ วิพฺภมิสฺสนฺติปี”ติ. อถ โข เต ภิกฺขู ภควโต เอตมตฺถํ อาโรเจสุํ.}

\prose{อถ โข ภควา เอตสฺมิํ นิทาเน เอตสฺมิํ ปกรเณ ภิกฺขุสํฆํ สนฺนิปาตาเปตฺวา ภิกฺขู ปฏิปุจฺฉิ “สจฺจํ กิร ภิกฺขเว อสฺสชิ\-ปุนพฺพสุกา ภิกฺขู สํเฆน ปพฺพาชนีย\-กมฺมกตา น สมฺมา วตฺตนฺติ, น โลมํ ปาเตนฺติ, น เนตฺถารํ วตฺตนฺติ, น ภิกฺขู ขมาเปนฺติ, อกฺโกสนฺติ, ปริภาสนฺติ, ฉนฺทคามิตา โทสคามิตา โมหคามิตา ภยคามิตา ปาเปนฺติ, ปกฺกมนฺติปิ วิพฺภมนฺติปี”ติ. สจฺจํ ภควาติ. วิครหิ พุทฺโธ ภควา \csromanfolio{35}อนนุจฺฉวิกํ ฯเปฯ. กถํ หิ นาม เต ภิกฺขเว โมฆปุริสา สํเฆน ปพฺพาชนีย\-กมฺมกตา น สมฺมา วตฺติสฺสนฺติ, น โลมํ ปาเตสฺสนฺติ, น เนตฺถารํ วตฺติสฺสนฺติ, น ภิกฺขู ขมาเปสฺสนฺติ, อกฺโกสิสฺสนฺติ, ปริภาสิสฺสนฺติ, ฉนฺทคามิตา โทสคามิตา โมหคามิตา ภยคามิตา ปาเปสฺสนฺติ, ปกฺกมิสฺสนฺติปิ วิพฺภมิสฺสนฺติปิ. เนตํ ภิกฺขเว อปฺปสนฺนานํ วา ปสาทาย ฯเปฯ วิครหิตฺวา ฯเปฯ ธมฺมิํ กถํ กตฺวา ภิกฺขู อามนฺเตสิ\csromandash{} เตน หิ ภิกฺขเว สํโฆ ปพฺพาชนีย\-กมฺมํ ปฏิปฺปสฺสมฺเภตุ.}

\csromanmatikaanchor{mtk.21}
\csromantocmarkref{subsection}{นปฺปฏิปฺปสฺสมฺเภตพฺพอฏฺฐารสก}{mtk.21}
\bookmark[dest={mtk.21},level=2]{นปฺปฏิปฺปสฺสมฺเภตพฺพอฏฺฐารสก}
\csromanheader{2}{นปฺปฏิปฺปสฺสมฺเภตพฺพ\-อฏฺฐารสก}
\proseitem{30}{ปญฺจหิ ภิกฺขเว องฺเคหิ สมนฺนาคตสฺส ภิกฺขุโน ปพฺพาชนีย\-กมฺมํ นปฺปฏิปฺปสฺสมฺเภตพฺพํ. อุปสมฺปาเทติ, นิสฺสยํ เทติ, สามเณรํ อุปฏฺฐาเปติ, ภิกฺขุโนวาทก\-สมฺมุติํ สาทิยติ, สมฺมโตปิ ภิกฺขุนิโย โอวทติ. อิเมหิ โข ภิกฺขเว ปญฺจหงฺเคหิ สมนฺนาคตสฺส ภิกฺขุโน ปพฺพาชนีย\-กมฺมํ นปฺปฏิปฺปสฺสมฺเภตพฺพํ.}

\prose{\csromansymbolfootnote{*}{วิ 5. 317 ปิฏฺเฐปิ.}อปเรหิปิ ภิกฺขเว ปญฺจหงฺเคหิ สมนฺนาคตสฺส ภิกฺขุโน ปพฺพาชนีย กมฺมํ นปฺปฏิปฺปสฺสมฺเภตพฺพํ. ยาย อาปตฺติยา สํเฆน ปพฺพาชนีย\-กมฺมํ กตํ โหติ, ตํ อาปตฺติํ อาปชฺชติ, อญฺญํ วา ตาทิสิกํ, ตโต วา ปาปิฏฺฐตรํ, กมฺมํ ครหติ, กมฺมิเก ครหติ. อิเมหิ โข ภิกฺขเว ปญฺจหงฺเคหิ สมนฺนาคตสฺส ภิกฺขุโน ปพฺพาชนีย\-กมฺมํ นปฺปฏิปฺปสฺสมฺเภตพฺพํ.}

\prosewithcloser{อฏฺฐหิ ภิกฺขเว องฺเคหิ สมนฺนาคตสฺส ภิกฺขุโน ปพฺพาชนีย\-กมฺมํ นปฺปฏิปฺปสฺสมฺเภตพฺพํ. ปกตตฺตสฺส ภิกฺขุโน อุโปสถํ ฐเปติ, ปวารณํ ฐเปติ, สวจนียํ กโรติ, อนุวาทํ ปฏฺฐเปติ, โอกาสํ กาเรติ, โจเทติ, สาเรติ, ภิกฺขูหิ สมฺปโยเชติ. อิเมหิ โข ภิกฺขเว อฏฺฐหงฺเคหิ สมนฺนาคตสฺส ภิกฺขุโน ปพฺพาชนีย\-กมฺมํ นปฺปฏิปฺปสฺสมฺเภตพฺพํ.}{ปพฺพาชนีย\-กมฺเม นปฺปฏิปฺปสฺสมฺเภตพฺพ\-อฏฺฐารสกํ นิฏฺฐิตํ.}
\csromanmatikaanchor{mtk.22}
\csromantocmarkref{subsection}{ปฏิปฺปสฺสมฺเภตพฺพอฏฺฐารสก}{mtk.22}
\bookmark[dest={mtk.22},level=2]{ปฏิปฺปสฺสมฺเภตพฺพอฏฺฐารสก}
\csromanheader{2}{ปฏิปฺปสฺสมฺเภตพฺพ\-อฏฺฐารสก}
\proseitem{31}{ปญฺจหิ ภิกฺขเว องฺเคหิ สมนฺนาคตสฺส ภิกฺขุโน ปพฺพาชนีย\-กมฺมํ ปฏิปฺปสฺสมฺเภตพฺพํ. น อุปสมฺปาเทติ, น นิสฺสยํ เทติ, น สามเณรํ \csromanfolio{36}อุปฏฺฐาเปติ, น ภิกฺขุโนวาทก\-สมฺมุติํ สาทิยติ, สมฺมโตปิ ภิกฺขุนิโย น โอวทติ. อิเมหิ โข ภิกฺขเว ปญฺจหงฺเคหิ สมนฺนาคตสฺส ภิกฺขุโน ปพฺพาชนีย\-กมฺมํ ปฏิปฺปสฺสมฺเภตพฺพํ.}

\prose{อปเรหิปิ ภิกฺขเว ปญฺจหงฺเคหิ สมนฺนาคตสฺส ภิกฺขุโน ปพฺพาชนีย\-กมฺมํ ปฏิปฺปสฺสมฺเภตพฺพํ. ยาย อาปตฺติยา สํเฆน ปพฺพาชนีย\-กมฺมํ กตํ โหติ, ตํ อาปตฺติํ น อาปชฺชติ, อญฺญํ วา ตาทิสิกํ, ตโต วา ปาปิฏฺฐตรํ, กมฺมํ น ครหติ, กมฺมิเก น ครหติ. อิเมหิ โข ภิกฺขเว ปญฺจหงฺเคหิ สมนฺนาคตสฺส ภิกฺขุโน ปพฺพาชนีย\-กมฺมํ ปฏิปฺปสฺสมฺเภตพฺพํ.}

\prosewithcloser{อฏฺฐหิ ภิกฺขเว องฺเคหิ สมนฺนาคตสฺส ภิกฺขุโน ปพฺพาชนีย\-กมฺมํ ปฏิปฺปสฺสมฺเภตพฺพํ. น ปกตตฺตสฺส ภิกฺขุโน อุโปสถํ ฐเปติ, น ปวารณํ ฐเปติ, น สวจนียํ กโรติ, น อนุวาทํ ปฏฺฐเปติ, น โอกาสํ กาเรติ, น โจเทติ, น สาเรติ, น ภิกฺขูหิ สมฺปโยเชติ. อิเมหิ โข ภิกฺขเว อฏฺฐหงฺเคหิ สมนฺนาคตสฺส ภิกฺขุโน ปพฺพาชนีย\-กมฺมํ ปฏิปฺปสฺสมฺเภตพฺพํ.}{ปพฺพาชนีย\-กมฺเม ปฏิปฺปสฺสมฺเภตพฺพ\-อฏฺฐารสกํ นิฏฺฐิตํ.}
\proseitem{32}{เอวญฺจ ปน ภิกฺขเว ปฏิปฺปสฺสมฺเภตพฺพํ, เตน ภิกฺขเว ปพฺพาชนีย\-กมฺมกเตน ภิกฺขุนา สํฆํ อุปสงฺกมิตฺวา เอกํสํ อุตฺตราสงฺคํ กริตฺวา วุฑฺฒานํ ภิกฺขูนํ ปาเท วนฺทิตฺวา อุกฺกุฏิกํ นิสีทิตฺวา อญฺชลิํ ปคฺคเหตฺวา เอวมสฺส วจนีโย “อหํ ภนฺเต สํเฆน ปพฺพาชนีย\-กมฺมกโต สมฺมา วตฺตามิ, โลมํ ปาเตมิ, เนตฺถารํ วตฺตามิ, ปพฺพาชนียสฺส กมฺมสฺส ปฏิปฺปสฺสทฺธิํ ยาจามี”ติ. ทุติยมฺปิ ยาจิตพฺพา. ตติยมฺปิ ยาจิตพฺพา. พฺยตฺเตน ภิกฺขุนา ปฏิพเลน สํโฆ ญาเปตพฺโพ\csromandash{}}

\prose{“สุณาตุ เม ภนฺเต สํโฆ, อยํ อิตฺตนฺนาโม ภิกฺขุ สํเฆน ปพฺพาชนีย\-กมฺมกโต สมฺมา วตฺตติ, โลมํ ปาเตติ, เนตฺถารํ วตฺตติ, ปพฺพาชนียสฺส กมฺมสฺส ปฏิปฺปสฺสทฺธิํ ยาจติ, ยทิ สํฆสฺส ปตฺตกลฺลํ, สํโฆ อิตฺถนฺนามสฺส ภิกฺขุโน ปพฺพาชนีย\-กมฺมํ ปฏิปฺปสฺสมฺเภยฺย, เอสา ญตฺติ.}

\prose{สุณาตุ เม ภนฺเต สํโฆ, อยํ อิตฺถนฺนาโม ภิกฺขุ สํเฆน ปพฺพาชนีย\-กมฺมกโต สมฺมา วตฺตติ, โลมํ ปาเตติ, เนตฺถารํ วตฺตติ, \csromanfolio{37}ปพฺพาชนียสฺส กมฺมสฺส ปฏิปฺปสฺสทฺธิํ ยาจติ, สํโฆ อิตฺถนฺนามสฺส ภิกฺขุโน ปพฺพาชนีย\-กมฺมํ ปฏิปฺปสฺสมฺเภติ, ยสฺสายสฺมโต ขมติ อิตฺถนฺนามสฺส ภิกฺขุโน ปพฺพาชนียสฺส กมฺมสฺส ปฏิปฺปสฺสทฺธิ, โส ตุณฺหสฺส, ยสฺส นกฺขมติ, โส ภาเสยฺย.}

\prose{ทุติยมฺปิ เอตมตฺถํ วทามิ ฯเปฯ. ตติยมฺปิ เอตมตฺถํ วทามิ. สุณาตุ เม ภนฺเต สํโฆ, อยํ อิตฺถนฺนาโม ภิกฺขุ สํเฆน ปพฺพาชนีย\-กมฺมกโต สมฺมา วตฺตติ, โลมํ ปาเตติ, เนตฺถารํ วตฺตติ, ปพฺพาชนียสฺส กมฺมสฺส ปฏิปฺปสฺสทฺธิํ ยาจติ, สํโฆ อิตฺถนฺนามสฺส ภิกฺขุโน ปพฺพาชนีย\-กมฺมํ ปฏิปฺปสฺสมฺเภติ, ยสฺสายสฺมโต ขมติ อิตฺถนฺนามสฺส ภิกฺขุโน ปพฺพาชนียสฺส กมฺมสฺส ปฏิปฺปสฺสทฺธิ, โส ตุณฺหสฺส, ยสฺส นกฺขมติ, โส ภาเสยฺย.}

\prose{ปฏิปฺปสฺสทฺธํ สํเฆน อิตฺถนฺนามสฺส ภิกฺขุโน ปพฺพาชนีย\-กมฺมํ, ขมติ สํฆสฺส, ตสฺมา ตุณฺหี, เอวเมตํ ธารยามี”ติ.}

\vspace{\tipitakaparskip}
\csromancenter{ปพฺพาชนีย\-กมฺมํ นิฏฺฐิตํ ตติยํ.}
\csromanmatikaanchor{mtk.23}
\csromantocmarkref{section}{4. ปฏิสารณียกมฺม}{mtk.23}
\bookmark[dest={mtk.23},level=1]{4. ปฏิสารณียกมฺม}
\csromanheader{1}{4. ปฏิสารณีย\-กมฺม}
\proseitem{33}{เตน โข ปน สมเยน อายสฺมา สุธมฺโม มจฺฉิกาสณฺเฑ จิตฺตสฺส คหปติโน อาวาสิโก โหติ นวกมฺมิโก ธุวภตฺติโก. ยทา จิตฺโต คหปติ สํฆํ วา คณํ วา ปุคฺคลํ วา นิมนฺเตตุกาโม โหติ, ตทา น อายสฺมนฺตํ สุธมฺมํ อนปโลเกตฺวา สํฆํ วา คณํ วา ปุคฺคลํ วา นิมนฺเตติ.}

\prose{เตน โข ปน สมเยน สมฺพหุลา เถรา ภิกฺขู อายสฺมา จ สาริปุตฺโต อายสฺมา จ มหาโมคฺคลฺลาโน อายสฺมา จ มหากจฺจาโน อายสฺมา จ มหาโกฏฺฐิโก อายสฺมา จ มหากปฺปิโน อายสฺมา จ มหาจุนฺโท อายสฺมา จ อนุรุทฺโธ อายสฺมา จ เรวโต อายสฺมา จ อุปาลิ อายสฺมา จ อานนฺโท อายสฺมา จ ราหุโล กาสีสุ จาริกํ จรมานา เยน มจฺฉิกาสณฺโฑ ตทวสรุํ.}

\prose{อสฺโสสิ โข จิตฺโต คหปติ “เถรา กิร ภิกฺขู มจฺฉิกาสณฺฑํ อนุปฺปตฺตา”ติ. อถ โข จิตฺโต คหปติ เยน เถรา ภิกฺขู \csromanfolio{38}เตนุปสงฺกมิ, อุปสงฺกมิตฺวา เถเร ภิกฺขู อภิวาเทตฺวา เอกมนฺตํ นิสีทิ, เอกมนฺตํ นิสินฺนํ โข จิตฺตํ คหปติํ อายสฺมา สาริปุตฺโต ธมฺมิยา กถาย สนฺทสฺเสสิ สมาทเปสิ สมุตฺเตเชสิ สมฺปหํเสสิ.}

\prose{อถ โข จิตฺโต คหปติ อายสฺมตา สาริปุตฺเตน ธมฺมิยา กถาย สนฺทสฺสิโต สมาทปิโต สมุตฺเตชิโต สมฺปหํสิโต เถเร ภิกฺขู เอตทโวจ “อธิวาเสนฺตุ เม ภนฺเต เถรา สฺวาตนาย อาคนฺตุกภตฺตนฺ”ติ. อธิวาเสสุํ โข เถรา ภิกฺขู\csromansharedfootnote{5954c91a583fb6f1}{อธิวาเสสุํ โข เต เถรา ภิกฺขู (สฺยา)} ตุณฺหีภาเวน.}

\prose{อถ โข จิตฺโต คหปติ เถรานํ ภิกฺขูนํ อธิวาสนํ วิทิตฺวา อุฏฺฐายาสนา เถเร ภิกฺขู อภิวาเทตฺวา ปทกฺขิณํ กตฺวา เยนายสฺมา สุธมฺโม เตนุปสงฺกมิ, อุปสงฺกมิตฺวา อายสฺมนฺตํ สุธมฺมํ อภิวาเทตฺวา เอกมนฺตํ อฏฺฐาสิ, เอกมนฺตํ ฐิโต โข จิตฺโต คหปติ อายสฺมนฺตํ สุธมฺมํ เอตทโวจ “อธิวาเสตุ เม ภนฺเต อยฺโย สุธมฺโม สฺวาตนาย ภตฺตํ สทฺธิํ เถเรหี”ติ. อถ โข อายสฺมา สุธมฺโม “ปุพฺเพ ขฺวายํ จิตฺโต คหปติ ยทา สํฆํ วา คณํ วา ปุคฺคลํ วา นิมนฺเตตุกาโม, น มํ อนปโลเกตฺวา สํฆํ วา คณํ วา ปุคฺคลํ วา นิมนฺเตติ, โสทานิ มํ อนปโลเกตฺวา เถเร ภิกฺขู นิมนฺเตสิ, ทุฏฺโฐทานายํ จิตฺโต คหปติ อนเปกฺโข วิรตฺตรูโป มยี”ติ จิตฺตํ คหปติํ เอตทโวจ “อลํ คหปติ, นาธิวาเสมี”ติ. ทุติยมฺปิ โข ฯเปฯ. ตติยมฺปิ โข จิตฺโต คหปติ อายสฺมนฺตํ สุธมฺมํ เอตทโวจ “อธิวาเสตุ เม ภนฺเต อยฺโย สุธมฺโม สฺวาตนาย ภตฺตํ สทฺธิํ เถเรหี”ติ. อลํ คหปติ, นาธิวาเสมีติ. อถ โข จิตฺโต คหปติ “กิํ เม กริสฺสติ อยฺโย สุธมฺโม อธิวาเสนฺโต วา อนธิวาเสนฺโต วา”ติ อายสฺมนฺตํ สุธมฺมํ อภิวาเทตฺวา ปทกฺขิณํ กตฺวา ปกฺกามิ.}

\proseitem{34}{อถ โข จิตฺโต คหปติ ตสฺสา รตฺติยา อจฺจเยน เถรานํ ภิกฺขูนํ ปณีตํ ขาทนียํ โภชนียํ ปฏิยาทาเปสิ. อถ โข อายสฺมา สุธมฺโม “ยํนูนาหํ จิตฺตสฺส คหปติโน เถรานํ ภิกฺขูนํ ปฏิยตฺตํ ปสฺเสยฺยนฺ”ติ ปุพฺพณฺหสมยํ นิวาเสตฺวา ปตฺต\-จีวรมา\-ทาย เยน จิตฺตสฺส คหปติโน นิเวสนํ เตนุปสงฺกมิ, อุปสงฺกมิตฺวา ปญฺญตฺเต อาสเน นิสีทิ. อถ โข จิตฺโต คหปติ เยนายสฺมา สุธมฺโม เตนุปสงฺกมิ, \csromanfolio{39}อุปสงฺกมิตฺวา อายสฺมนฺตํ สุธมฺมํ อภิวาเทตฺวา เอกมนฺตํ นิสีทิ, เอกมนฺตํ นิสินฺนํ โข จิตฺตํ คหปติํ อายสฺมา สุธมฺโม เอตทโวจ “ปหูตํ โข เต อิทํ คหปติ ขาทนียํ โภชนียํ ปฏิยตฺตํ, เอกา จ โข อิธ นตฺถิ ยทิทํ ติลสํคุฬิกา”ติ. พหุมฺหิ วต ภนฺเต รตเน พุทฺธวจเน\csromansharedfootnote{2af135e9024e8d7d}{ภนฺเต พุทฺธวจเน (สฺยา)} วิชฺชมาเน อยฺเยน สุธมฺเมน ยเทว กิญฺจิ ภาสิตํ ยทิทํ ติลสํคุฬิกาติ. ภูตปุพฺพํ ภนฺเต ทกฺขิณาปถกา วาณิชา ปุรตฺถิมํ ชนปทํ อคมํสุ วาณิชฺชาย, เต ตโต กุกฺกุฏิํ อาเนสุํ. อถ โข สา ภนฺเต กุกฺกุฏี กาเกน สทฺธิํ สํวาสํ กปฺเปสิ, สา โปตกํ ชเนสิ, ยทา โข โส ภนฺเต กุกฺกุฏโปตโก กากวสฺสํ วสฺสิตุกาโม โหติ, “กากกุกฺกุฏี”ติ วสฺสติ. ยทา กุกฺกุฏิวสฺสํ วสฺสิตุกาโม โหติ, “กุกฺกุฏิกากา”ติ วสฺสติ. เอวเมว โข ภนฺเต พหุมฺหิ รตเน พุทฺธวจเน วิชฺชมาเน อยฺเยน สุธมฺเมน ยเทว กิญฺจิ ภาสิตํ ยทิทํ ติลสํคุฬิกาติ. อกฺโกสสิ มํ ตฺวํ คหปติ, ปริภาสสิ มํ ตฺวํ คหปติ, เอโส เต คหปติ อาวาโส, ปกฺกมิสฺสามีติ. นาหํ ภนฺเต อยฺยํ สุธมฺมํ อกฺโกสามิ ปริภาสามิ\csromansharedfootnote{1b0f35fe9cf1c737}{น ปริภาสามิ (สี, สฺยา)}, วสตุ ภนฺเต อยฺโย สุธมฺโม มจฺฉิกาสณฺเฑ รมณียํ อมฺพาฏกวนํ, อหํ อยฺยสฺส สุธมฺมสฺส อุสฺสุกฺกํ กริสฺสามิ จีวร\-ปิณฺฑปาต\-เสนาสน\-คิลานปฺปจฺจย\-เภสชฺช\-ปริกฺขารานนฺติ. ทุติยมฺปิ โข ฯเปฯ. ตติยมฺปิ โข อายสฺมา สุธมฺโม จิตฺตํ คหปติํ เอตทโวจ “อกฺโกสสิ มํ ตฺวํ คหปติ, ปริภาสสิ มํ ตฺวํ คหปติ, เอโส เต คหปติ อาวาโส, ปกฺกมิสฺสามี”ติ. กหํ ภนฺเต อยฺโย สุธมฺโม คมิสฺสตีติ. สาวตฺถิํ โข อหํ คหปติ คมิสฺสามิ ภควนฺตํ ทสฺสนายาติ. เตน หิ ภนฺเต ยญฺจ อตฺตนา ภณิตํ, ยญฺจ มยา ภณิตํ, ตํ สพฺพํ ภควโต อาโรเจหิ, อนจฺฉริยํ โข ปเนตํ ภนฺเต ยํ อยฺโย สุธมฺโม ปุนเทว มจฺฉิกาสณฺฑํ ปจฺจาคจฺเฉยฺยาติ.}

\proseitem{35}{อถ โข อายสฺมา สุธมฺโม เสนาสนํ สํสาเมตฺวา ปตฺต\-จีวรมา\-ทาย เยน สาวตฺถิ เตน ปกฺกามิ, อนุปุพฺเพน เยน สาวตฺถิ เชตวนํ อนาถ\-ปิณฺฑิกสฺส อาราโม, เยน ภควา เตนุปสงฺกมิ, อุปสงฺกมิตฺวา ภควนฺตํ อภิวาเทตฺวา เอกมนฺตํ นิสีทิ, เอกมนฺตํ นิสินฺโน โข อายสฺมา สุธมฺโม ยญฺจ อตฺตนา ภณิตํ, ยญฺจ จิตฺเตน คหปตินา ภณิตํ, ตํ สพฺพํ ภควโต อาโรเจสิ.}

\prose{\csromanfolio{40}วิครหิ พุทฺโธ ภควา อนนุจฺฉวิกํ โมฆปุริส อนนุโลมิกํ อปฺปติรูปํ อสฺสามณกํ อกปฺปิยํ อกรณียํ. กถํ หิ นาม ตฺวํ โมฆปุริส จิตฺตํ คหปติํ สทฺธํ ปสนฺนํ ทายกํ การกํ สํฆุปฏฺฐากํ หีเนน ขุํเสสาสิ, หีเนน วมฺเภสฺสสิ. เนตํ โมฆปุริส อปฺปสนฺนานํ วา ปสาทาย ฯเปฯ วิครหิตฺวา ฯเปฯ ธมฺมิํ กถํ กตฺวา ภิกฺขู อามนฺเตสิ\csromandash{} เตน หิ ภิกฺขเว สํโฆ สุธมฺมสฺส ภิกฺขุโน ปฏิสารณีย\-กมฺมํ กโรตุ “จิตฺโต เต คหปติ ขมาเปตพฺโพ”ติ. เอวญฺจ ปน ภิกฺขเว กาตพฺพํ, ปฐมํ สุธมฺโม ภิกฺขุ โจเทตพฺโพ, โจเทตฺวา สาเรตพฺโพ, สาเรตฺวา อาปตฺติํ อาโรเปตพฺโพ, อาปตฺติํ อาโรเปตฺวา พฺยตฺเตน ภิกฺขุนา ปฏิพเลน สํโฆ ญาเปตพฺโพ\csromandash{}}

\proseitem{36}{“สุณาตุ เม ภนฺเต สํโฆ, อยํ สุธมฺโม ภิกฺขุ จิตฺตํ คหปติํ สทฺธํ ปสนฺนํ ทายกํ การกํ สํฆุปฏฺฐากํ หีเนน ขุํเสติ, หีเนน วมฺเภติ, ยทิ สํฆสฺส ปตฺตกลฺลํ, สํโฆ สุธมฺมสฺส ภิกฺขุโน ปฏิสารณีย\-กมฺมํ กเรยฺย ‘จิตฺโต เต คตปติ ขมาเปตพฺโพ’ติ, เอสา ญตฺติ.}

\prose{สุณาตุ เม ภนฺเต สํโฆ, อยํ สุธมฺโม ภิกฺขุ จิตฺตํ คหปติํ สทฺธํ ปสนฺนํ ทายกํ การกํ สํฆุปฏฺฐากํ หีเนน ขุํเสติ, หีเนน วมฺเภติ, สํโฆ สุธมฺมสฺส ภิกฺขุโน ปฏิสารณีย\-กมฺมํ กโรติ ‘จิตฺโต เต คหปติ ขมาเปตพฺโพ’ติ, ยสฺสายสฺมโต ขมติ สุธมฺมสฺส ภิกฺขุโน ปฏิสารณียสฺส กมฺมสฺส กรณํ ‘จิตฺโต เต คหปติ ขมาเปตพฺโพ’ติ, โส ตุณฺหสฺส, ยสฺส นกฺขมติ, โส ภาเสยฺย.}

\prose{ทุติยมฺปิ เอตมตฺถํ วทามิ ฯเปฯ. ตติยมฺปิ เอตมตฺถํ วทามิ. สุณาตุ เม ภนฺเต สํโฆ, อยํ สุธมฺโม ภิกฺขุ จิตฺตํ คหปติํ สทฺธํ ปสนฺนํ ทายกํ การกํ สํฆุปฏฺฐากํ หีเนน ขุํเสติ, หีเนน วมฺเภติ, สํโฆ สุธมฺมสฺส ภิกฺขุโน ปฏิสารณีย\-กมฺมํ กโรติ ‘จิตฺโต เต คหปติ ขมาเปตพฺโพ’ติ, ยสฺสายสฺมโต ขมติ สุธมฺมสฺส ภิกฺขุโน ปฏิสารณียสฺส กมฺมสฺส กรณํ ‘จิตฺโต เต คหปติ ขมาเปตพฺโพ’ติ, โส ตุณฺหสฺส, ยสฺส นกฺขมติ, โส ภาเสยฺย.}

\prose{กตํ สํเฆน สุธมฺมสฺส ภิกฺขุโน ปฏิสารณีย\-กมฺมํ ‘จิตฺโต เต คหปติ ขมาเปตพฺโพ’ติ, ขมติ สํฆสฺส, ตสฺมา ตุณฺหี, เอวเมตํ ธารยามี”ติ.}

\csromanmatikaanchor{mtk.24}
\csromantocmarkref{subsection}{อธมฺมกมฺมทฺวาทสก}{mtk.24}
\bookmark[dest={mtk.24},level=2]{อธมฺมกมฺมทฺวาทสก}
\csromanheader{2}{อธมฺมกมฺม\-ทฺวาทสก}
\proseitem{37}{\csromanfolio{41}ตีหิ ภิกฺขเว องฺเคหิ สมนฺนาคตํ ปฏิสารณีย\-กมฺมํ อธมฺมกมฺมญฺจ โหติ อวินยกมฺมญฺจ ทุวูปสนฺตญฺจ. อสมฺมุขา กตํ โหติ, อปฺปฏิปุจฺฉา กตํ โหติ, อปฺปฏิญฺญาย กตํ โหติ. อิเมหิ โข ภิกฺขเว ตีหงฺเคหิ สมนฺนาคตํ ปฏิสารณีย\-กมฺมํ อธมฺมกมฺมญฺจ โหติ อวินยกมฺมญฺจ ทุวูปสนฺตญฺจ.}

\prosewithcloser{อปเรหิปิ ภิกฺขเว ตีหงฺเคหิ สมนฺนาคตํ ปฏิสารณีย\-กมฺมํ อธมฺมกมฺมญฺจ โหติ อวินยกมฺมญฺจ ทุวูปสนฺตญฺจ. อนาปตฺติยา กตํ โหติ, อเทสนาคามินิยา อาปตฺติยา กตํ โหติ, เทสิตาย อาปตฺติยา กตํ โหติ ฯเปฯ. อโจเทตฺวา กตํ โหติ, อสาเรตฺวา กตํ โหติ, อาปตฺติํ อนาโรเปตฺวา กตํ โหติ ฯเปฯ. อสมฺมุขา กตํ โหติ, อธมฺเมน กตํ โหติ, วคฺเคน กตํ โหติ ฯเปฯ. อปฺปฏิปุจฺฉา กตํ โหติ, อธมฺเมน กตํ โหติ, วคฺเคน กตํ โหติ ฯเปฯ. อปฺปฏิญฺญาย กตํ โหติ, อธมฺเมน กตํ โหติ, วคฺเคน กตํ โหติ ฯเปฯ. อนาปตฺติยา กตํ โหติ, อธมฺเมน กตํ โหติ, วคฺเคน กตํ โหติ ฯเปฯ. อเทสนาคามินิยา อาปตฺติยา กตํ โหติ, อธมฺเมน กตํ โหติ, วคฺเคน กตํ โหติ ฯเปฯ. เทสิตาย อาปตฺติยา กตํ โหติ, อธมฺเมน กตํ โหติ, วคฺเคน กตํ โหติ ฯเปฯ. อโจเทตฺวา กตํ โหติ, อธมฺเมน กตํ โหติ, วคฺเคน กตํ โหติ ฯเปฯ. อสาเรตฺวา กตํ โหติ, อธมฺเมน กตํ โหติ, วคฺเคน กตํ โหติ ฯเปฯ. อาปตฺติํ อนาโรเปตฺวา กตํ โหติ, อธมฺเมน กตํ โหติ, วคฺเคน กตํ โหติ. อิเมหิ โข ภิกฺขเว ตีหงฺเคหิ สมนฺนาคตํ ปฏิสารณีย\-กมฺมํ อธมฺมกมฺมญฺจ โหติ อวินยกมฺมญฺจ ทุวูปสนฺตญฺจ.}{ปฏิสารณีย\-กมฺเม อธมฺมกมฺม\-ทฺวาทสกํ นิฏฺฐิตํ.}
\csromanmatikaanchor{mtk.25}
\csromantocmarkref{subsection}{ธมฺมกมฺมทฺวาทสก}{mtk.25}
\bookmark[dest={mtk.25},level=2]{ธมฺมกมฺมทฺวาทสก}
\csromanheader{2}{ธมฺมกมฺม\-ทฺวาทสก}
\proseitem{38}{ตีหิ ภิกฺขเว องฺเคหิ สมนฺนาคตํ ปฏิสารณีย\-กมฺมํ ธมฺมกมฺมญฺจ โหติ วินยกมฺมญฺจ สุวูปสนฺตญฺจ. สมฺมุขา กตํ โหติ, ปฏิปุจฺฉา กตํ โหติ, ปฏิญฺญาย กตํ โหติ. อิเมหิ โข ภิกฺขเว ตีหงฺเคหิ สมนฺนาคตํ ปฏิสารณีย\-กมฺมํ ธมฺมกมฺมญฺจ โหติ วินยกมฺมญฺจ สุวูปสนฺตญฺจ.}

\prosewithcloser{\csromanfolio{42}อปเรหิปิ ภิกฺขเว ตีหงฺเคหิ สมนฺนาคตํ ปฏิสารณีย\-กมฺมํ ธมฺมกมฺมญฺจ โหติ วินยกมฺมญฺจ สุวูปสนฺตญฺจ. อาปตฺติยา กตํ โหติ, เทสนาคามินิยา อาปตฺติยา กตํ โหติ, อเทสิตาย อาปตฺติยา กตํ โหติ ฯเปฯ. โจเทตฺวา กตํ โหติ, สาเรตฺวา กตํ โหติ, อาปตฺติํ อาโรเปตฺวา กตํ โหติ ฯเปฯ. สมฺมุขา กตํ โหติ, ธมฺเมน กตํ โหติ, สมคฺเคน กตํ โหติ ฯเปฯ. ปฏิปุจฺฉา กตํ โหติ, ธมฺเมน กตํ โหติ, สมคฺเคน กตํ โหติ ฯเปฯ. ปฏิญฺญาย กตํ โหติ, ธมฺเมน กตํ โหติ, สมคฺเคน กตํ โหติ ฯเปฯ. อาปตฺติยา กตํ โหติ, ธมฺเมน กตํ โหติ, สมคฺเคน กตํ โหติ ฯเปฯ. เทสนาคามินิยา อาปตฺติยา กตํ โหติ, ธมฺเมน กตํ โหติ, สมคฺเคน กตํ โหติ ฯเปฯ. อเทสิตาย อาปตฺติยา กตํ โหติ, ธมฺเมน กตํ โหติ, สมคฺเคน กตํ โหติ ฯเปฯ. โจเทตฺวา กตํ โหติ, ธมฺเมน กตํ โหติ, สมคฺเคน กตํ โหติ ฯเปฯ. สาเรตฺวา กตํ โหติ, ธมฺเมน กตํ โหติ, สมคฺเคน กตํ โหติ ฯเปฯ. อาปตฺติํ อาโรเปตฺวา กตํ โหติ, ธมฺเมน กตํ โหติ, สมคฺเคน กตํ โหติ. อิเมหิ โข ภิกฺขเว ตีหงฺเคหิ สมนฺนาคตํ ปฏิสารณีย\-กมฺมํ ธมฺมกมฺมญฺจ โหติ วินยกมฺมญฺจ สุวูปสนฺตญฺจ.}{ปฏิสารณีย\-กมฺเม ธมฺมกมฺม\-ทฺวาทสกํ นิฏฺฐิตํ.}
\csromanmatikaanchor{mtk.26}
\csromantocmarkref{subsection}{อากงฺขมานจตุกฺก}{mtk.26}
\bookmark[dest={mtk.26},level=2]{อากงฺขมานจตุกฺก}
\csromanheader{2}{อากงฺขมาน\-จตุกฺก}
\proseitem{39}{ปญฺจหิ ภิกฺขเว องฺเคหิ สมนฺนาคตสฺส ภิกฺขุโน อากงฺขมาโน สํโฆ ปฏิสารณีย\-กมฺมํ กเรยฺย. คิหีนํ อลาภาย ปริสกฺกติ คิหีนํ อนตฺถาย ปริสกฺกติ, คิหีนํ อนาวาสาย\csromansharedfootnote{f821a118c0b374cc}{อวาสาย (สี)} ปริสกฺกติ, คิหี อกฺโกสติ ปริภาสติ, คิหี คิหีหิ เภเทติ. อิเมหิ โข ภิกฺขเว ปญฺจหงฺเคหิ สมนฺนาคตสฺส ภิกฺขุโน อากงฺขมาโน สํโฆ ปฏิสารณีย\-กมฺมํ กเรยฺย. (1)}

\prose{อปเรหิปิ ภิกฺขเว ปญฺจหงฺเคหิ สมนฺนาคตสฺส ภิกฺขุโน อากงฺขมาโน สํโฆ ปฏิสารณีย\-กมฺมํ กเรยฺย. คิหีนํ พุทฺธสฺส อวณฺณํ \csromanfolio{43}ภาสติ, คิหีนํ ธมฺมสฺส อวณฺณํ ภาสติ, คิหีนํ สํฆสฺส อวณฺณํ ภาสติ, คิหี หีเนน ขุํเสติ หีเนน วมฺเภติ, คิหีนํ ธมฺมิกํ ปฏิสฺสวํ น สจฺจาเปติ. อิเมหิ โข ภิกฺขเว ปญฺจหงฺเคหิ สมนฺนาคตสฺส ภิกฺขุโน อากงฺขมาโน สํโฆ ปฏิสารณีย\-กมฺมํ กเรยฺย. (2)}

\prose{ปญฺจนฺนํ ภิกฺขเว ภิกฺขูนํ อากงฺขมาโน สํโฆ ปฏิสารณีย\-กมฺมํ กเรยฺย. เอโก คิหีนํ อลาภาย ปริสกฺกติ, เอโก คิหีนํ อนตฺถาย ปริสกฺกติ, เอโก คิหีนํ อนาวาสาย ปริสกฺกติ, เอโก คิหี อกฺโกสติ ปริภาสติ, เอโก คิหี คิหีหิ เภเทติ. อิเมสํ โข ภิกฺขเว ปญฺจนฺนํ ภิกฺขูนํ อากงฺขมาโน สํโฆ ปฏิสารณีย\-กมฺมํ กเรยฺย. (3)}

\prosewithcloser{อปเรสมฺปิ ภิกฺขเว ปญฺจนฺนํ ภิกฺขูนํ อากงฺขมาโน สํโฆ ปฏิสารณีย\-กมฺมํ กเรยฺย. เอโก คิหีนํ พุทฺธสฺส อวณฺณํ ภาสติ, เอโก คิหีนํ ธมฺมสฺส อวณฺณํ ภาสติ, เอโก คิหีนํ สํฆสฺส อวณฺณํ ภาสติ, เอโก คิหี หีเนน ขุํเสติ หีเนน วมฺเภติ, เอโก คิหีนํ ธมฺมิกํ ปฏิสฺสวํ น สจฺจาเปติ. อิเมสํ โข ภิกฺขเว ปญฺจนฺนํ ภิกฺขูนํ อากงฺขมาโน สํโฆ ปฏิสารณีย\-กมฺมํ กเรยฺย. (4)}{อากงฺขมาน\-จตุกฺกํ นิฏฺฐิตํ.}
\csromanmatikaanchor{mtk.27}
\csromantocmarkref{subsection}{อฏฺฐารสวตฺต}{mtk.27}
\bookmark[dest={mtk.27},level=2]{อฏฺฐารสวตฺต}
\csromanheader{2}{\textbf{อฏฺฐารสวตฺต}}
\proseitemwithcloser{40}{ปฏิสารณีย\-กมฺมกเตน ภิกฺขเว ภิกฺขุนา สมฺมา วตฺติตพฺพํ. ตตฺรายํ สมฺมาวตฺตนา\csromandash{}น อุปสมฺ\-ปาเท\-ตพฺพํ, น นิสฺสโย ทาตพฺโพ, น สามเณโร อุปฏฺฐาเปตพฺโพ, น ภิกฺขุโนวาทก\-สมฺมุติ สาทิตพฺพา, สมฺมเตนปิ ภิกฺขุนิโย น โอวทิตพฺพา, ยาย อาปตฺติยา สํเฆน ปฏิสารณีย\-กมฺมํ กตํ โหติ, สา อาปตฺติ น อาปชฺชิตพฺพา, อญฺญา วา ตาทิสิกา, ตโต วา ปาปิฏฺฐตรา, กมฺมํ น ครหิตพฺพํ, กมฺมิกา น ครหิตพฺพา, น ปกตตฺตสฺส ภิกฺขุโน อุโปสโถ ฐเปตพฺโพ, น ปวารณา ฐเปตพฺพา, น สวจนียํ กาตพฺพํ, น อนุวาโท ปฏฺฐเปตพฺโพ, น โอกาโส กาเรตพฺโพ, น โจเทตพฺโพ, น สาเรตพฺโพ, น ภิกฺขูหิ สมฺป\-โยเช\-ตพฺพนฺติ.}{ปฏิสารณีย\-กมฺเม อฏฺฐารส\-วตฺตํ นิฏฺฐิตํ.}
\proseitem{41}{\csromanfolio{44}อถ โข สํโฆ สุธมฺมสฺส ภิกฺขุโน ปฏิสารณีย\-กมฺมํ อกาสิ ‘จิตฺโต เต คหปติ ขมาเปตพฺโพ’ติ. โส สํเฆน ปฏิสารณีย\-กมฺมกโต มจฺฉิกาสณฺฑํ คนฺตฺวา มงฺกุภูโต นาสกฺขิ จิตฺตํ คหปติํ ขมาเปตุํ, ปุนเทว สาวตฺถิํ ปจฺจาคญฺฉิ. ภิกฺขู เอวมาหํสุ “ขมาปิโต ตยา อาวุโส สุธมฺม จิตฺโต คหปตี”ติ. อิธาหํ อาวุโส มจฺฉิกาสณฺฑํ คนฺตฺวา มงฺกุภูโต นาสกฺขิํ จิตฺตํ คหปติํ ขมาเปตุนฺติ. ภิกฺขู ภควโต เอตมตฺถํ อาโรเจสุํ ฯเปฯ. เตน หิ ภิกฺขเว สํโฆ สุธมฺมสฺส ภิกฺขุโน อนุทูตํ เทตุ จิตฺตํ คหปติํ ขมาเปตุํ. เอวญฺจ ปน ภิกฺขเว ทาตพฺโพ, ปฐมํ ภิกฺขุ ยาจิตพฺโพ, ยาจิตฺวา พฺยตฺเตน ภิกฺขุนา ปฏิพเลน สํโฆ ญาเปตพฺโพ\csromandash{}}

\prose{“สุณาตุ เม ภนฺเต สํโฆ, ยทิ สํฆสฺส ปตฺตกลฺลํ, สํโฆ อิตฺถนฺนามํ ภิกฺขุํ สุธมฺมสฺส ภิกฺขุโน อนุทูตํ ทเทยฺย จิตฺตํ คหปติํ ขมาเปตุํ, เอสา ญตฺติ.}

\prose{สุณาตุ เม ภนฺเต สํโฆ, สํโฆ อิตฺถนฺนามํ ภิกฺขุํ สุธมฺมสฺส ภิกฺขุโน อนุทูตํ เทติ จิตฺตํ คหปติํ ขมาเปตุํ, ยสฺสายสฺมโต ขมติ อิตฺถนฺนามสฺส ภิกฺขุโน สุธมฺมสฺส ภิกฺขุโน อนุทูตสฺส ทานํ จิตฺตํ คหปติํ ขมาเปตุํ, โส ตุณฺหสฺส, ยสฺส นกฺขมติ, โส ภาเสยฺย.}

\prose{ทินฺโน สํเฆน อิตฺถนฺนาโม ภิกฺขุ สุธมฺมสฺส ภิกฺขุโน อนุทูโต จิตฺตํ คหปติํ ขมาเปตุํ, ขมติ สํฆสฺส, ตสฺมา ตุณฺหี, เอวเมตํ ธารยามี”ติ.}

\proseitem{42}{เตน ภิกฺขเว สุธมฺเมน ภิกฺขุนา อนุทูเตน ภิกฺขุนา สทฺธิํ มจฺฉิกาสณฺฑํ คนฺตฺวา จิตฺโต คหปติ ขมาเปตพฺโพ “ขม คหปติ, ปสาเทมิ ตนฺ”ติ. เอวญฺเจ วุจฺจมาโน ขมติ, อิจฺเจตํ กุสลํ, โน เจ ขมติ, อนุทูเตน ภิกฺขุนา วตฺตพฺโพ “ขม คหปติ อิมสฺส ภิกฺขุโน, ปสาเทติ ตนฺ”ติ. เอวญฺเจ วุจฺจมาโน ขมติ, อิจฺเจตํ กุสลํ, โน เจ ขมติ, อนุทูเตน ภิกฺขุนา วตฺตพฺโพ “ขม คหปติ อิมสฺส ภิกฺขุโน, อหํ ตํ ปสาเทมี”ติ. เอวญฺเจ วุจฺจมาโน ขมติ, อิจฺเจตํ กุสลํ, โน เจ ขมติ, อนุทูเตน ภิกฺขุนา วตฺตพฺโพ “ขม คหปติ อิมสฺส ภิกฺขุโน สํฆสฺส วจเนนา”ติ. เอวญฺเจ วุจฺจมาโน ขมติ, อิจฺเจตํ กุสลํ, โน \csromanfolio{45}เจ ขมติ, อนุทูเตน ภิกฺขุนา สุธมฺโม ภิกฺขุ\csromansharedfootnote{7fbc75a74c60cbd9}{สุธมฺมํ ภิกฺขุ ... สา อาปตฺติ เทสาเปตพฺพาติ (สี, สฺยา)} จิตฺตสฺส คหปติโน ทสฺสนูปเจรํ อวิชหาเปตฺวา สวนูปจารํ อวิชหาเปตฺวา เอกํสํ อุตฺตราสงฺคํ การาเปตฺวา อุกฺกุฏิกํ นิสีทาเปตฺวา อญฺชลิํ ปคฺคณฺหาเปตฺวา ตํ อาปตฺติํ เทสาเปตพฺโพติ\csromansharedfootnote{7fbc75a74c60cbd9}{สุธมฺมํ ภิกฺขุ ... สา อาปตฺติ เทสาเปตพฺพาติ (สี, สฺยา)}.}

\prose{อถ โข อายสฺมา สุธมฺโม อนุทูเตน ภิกฺขุนา สทฺธิํ มจฺฉิกาสณฺฑํ คนฺตฺวา จิตฺตํ คหปติํ ขมาเปสิ, โส สมฺมา วตฺตติ, โลมํ ปาเตติ, เนตฺถารํ วตฺตติ, ภิกฺขู อุปสงฺกํตฺวา เอวํ วเทติ “อหํ อาวุโส สํเฆน ปฏิสารณีย\-กมฺมกโต สมฺมา วตฺตามิ, โลมํ ปาเตมิ, เนตฺถารํ วตฺตามิ, กถํ นุ โข มยา ปฏิปชฺชิตพฺพนฺ”ติ. ภิกฺขู ภควโต เอตมตฺถํ อาโรเจสุํ ฯเปฯ. เตน หิ ภิกฺขเว สํโฆ สุธมฺมสฺส ภิกฺขุโน ปฏิสารณีย\-กมฺมํ ปฏิปฺปสฺสมฺเภตุ.}

\csromanmatikaanchor{mtk.28}
\csromantocmarkref{subsection}{นปฺปฏิปฺปสฺสมฺเภตพฺพอฏฺฐารสก}{mtk.28}
\bookmark[dest={mtk.28},level=2]{นปฺปฏิปฺปสฺสมฺเภตพฺพอฏฺฐารสก}
\csromanheader{2}{นปฺปฏิปฺปสฺสมฺเภตพฺพ\-อฏฺฐารสก}
\proseitem{43}{ปญฺจหิ ภิกฺขเว องฺเคหิ สมนฺนาคตสฺส ภิกฺขุโน ปฏิสารณีย\-กมฺมํ นปฺปฏิปฺปสฺสมฺเภตพฺพํ. อุปสมฺปาเทติ, นิสฺสยํ เทติ, สามเณรํ อุปฏฺฐาเปติ, ภิกฺขุโนวาทก\-สมฺมุติํ สาทิยติ, สมฺมโตปิ ภิกฺขุนิโย โอวทติ. อิเมหิ โข ภิกฺขเว ปญฺจหงฺเคหิ สมนฺนาคตสฺส ภิกฺขุโน ปฏิสารณีย\-กมฺมํ นปฺปฏิปฺปสฺสมฺเภตพฺพํ.}

\prose{\csromansymbolfootnote{*}{วิ 5. 317 ปิฏฺเฐปิ.}อปเรหิปิ ภิกฺขเว ปญฺจหงฺเคหิ สมนฺนาคตสฺส ภิกฺขุโน ปฏิสารณีย\-กมฺมํ นปฺปฏิปฺปสฺสมฺเภตพฺพํ. ยาย อาปตฺติยา สํเฆน ปฏิสารณีย\-กมฺมํ กตํ โหติ, ตํ อาปตฺติํ อาปชฺชติ, อญฺญํ วา ตาทิสิกํ, ตโต วา ปาปิฏฺฐตรํ, กมฺมํ ครหติ, กมฺมิเก ครหติ. อิเมหิ โข ภิกฺขเว ปญฺจหงฺเคหิ สมนฺนาคตสฺส ภิกฺขุโน ปฏิสารณีย\-กมฺมํ นปฺปฏิปฺปสฺสมฺเภตพฺพํ.}

\prosewithcloser{อฏฺฐหิ ภิกฺขเว องฺเคหิ สมนฺนาคตสฺส ภิกฺขุโน ปฏิสารณีย\-กมฺมํ นปฺปฏิปฺปสฺสมฺเภตพฺพํ. ปกตตฺตสฺส ภิกฺขุโน อุโปสถํ ฐเปติ, ปวารณํ ฐเปติ, สวจนียํ กโรติ, อนุวาทํ ปฏฺฐเปติ, โอกาสํ กาเรติ, \csromanfolio{46}โจเทติ, สาเรติ, ภิกฺขูหิ สมฺปโยเชติ. อิเมหิ โข ภิกฺขเว อฏฺฐหงฺเคหิ สมนฺนาคตสฺส ภิกฺขุโน ปฏิสารณีย\-กมฺมํ นปฺปฏิปฺปสฺสมฺเภตพฺพํ.}{ปฏิสารณีย\-กมฺเม นปฺปฏิปฺปสฺสมฺเภตพฺพ\-อฏฺฐารสกํ นิฏฺฐิตํ.}
\csromanmatikaanchor{mtk.29}
\csromantocmarkref{subsection}{ปฏิปฺปสฺสมฺเภตพฺพอฏฺฐารสก}{mtk.29}
\bookmark[dest={mtk.29},level=2]{ปฏิปฺปสฺสมฺเภตพฺพอฏฺฐารสก}
\csromanheader{2}{ปฏิปฺปสฺสมฺเภตพฺพ\-อฏฺฐารสก}
\proseitem{44}{ปญฺจหิ ภิกฺขเว องฺเคหิ สมนฺนาคตสฺส ภิกฺขุโน ปฏิสารณีย\-กมฺมํ ปฏิปฺปสฺสมฺเภตพฺพํ. น อุปสมฺปาเทติ, น นิสฺสยํ เทติ, น สามเณรํ อุปฏฺฐาเปติ, น ภิกฺขุโนวาทก\-สมฺมุติํ สาทิยติ, สมฺมโตปิ ภิกฺขุนิโย น โอวทติ. อิเมหิ โข ภิกฺขเว ปญฺจหงฺเคหิ สมนฺนาคตสฺส ภิกฺขุโน ปฏิสารณีย\-กมฺมํ ปฏิปฺปสฺสมฺเภตพฺพํ.}

\prose{อปเรหิปิ ภิกฺขเว ปญฺจหงฺเคหิ สมนฺนาคตสฺส ภิกฺขุโน ปฏิสารณีย\-กมฺมํ ปฏิปฺปสฺสมฺเภตพฺพํ. ยาย อาปตฺติยา สํเฆน ปฏิสารณีย\-กมฺมํ กตํ โหติ, ตํ อาปตฺติํ น อาปชฺชติ, อญฺญํ วา ตาทิสิกํ, ตโต วา ปาปิฏฺฐตรํ, กมฺมํ น ครหติ, กมฺมิเก น ครหติ. อิเมหิ โข ภิกฺขเว ปญฺจหงฺเคหิ สมนฺนาคตสฺส ภิกฺขุโน ปฏิสารณีย\-กมฺมํ ปฏิปฺปสฺสมฺเภตพฺพํ.}

\prosewithcloser{อฏฺฐหิ ภิกฺขเว องฺเคหิ สมนฺนาคตสฺส ภิกฺขุโน ปฏิสารณีย\-กมฺมํ ปฏิปฺปสฺสมฺเภตพฺพํ. น ปกตตฺตสฺส ภิกฺขุโน อุโปสถํ ฐเปติ, น ปวารณํ ฐเปติ, น สวจนียํ กโรติ, น อนุวาทํ ปฏฺฐเปติ, น โอกาสํ กาเรติ, น โจเทติ, น สาเรติ, น ภิกฺขูหิ สมฺปโยเชติ. อิเมหิ โข ภิกฺขเว อฏฺฐหงฺเคหิ สมนฺนาคตสฺส ภิกฺขุโน ปฏิสารณีย\-กมฺมํ ปฏิปฺปสฺสมฺเภตพฺพํ.}{ปฏิสารณีย\-กมฺเม ปฏิปฺปสฺสมฺเภตพฺพ\-อฏฺฐารสกํ นิฏฺฐิตํ.}
\proseitem{45}{เอวญฺจ ปน ภิกฺขเว ปฏิปฺปสฺสมฺเภตพฺพํ, เตน\csromansharedfootnote{4ba2f8233ccda3c0}{เตน หิ (สฺยา, ก)} ภิกฺขเว สุธมฺเมน ภิกฺขุนา สํฆํ อุปสงฺกมิตฺวา เอกํสํ อุตฺตราสงฺคํ กริตฺวา วุฑฺฒานํ ภิกฺขูนํ ปาเท วนฺทิตฺวา อุกฺกุฏิกํ นิสีทิตฺวา อญฺชลิํ ปคฺคเหตฺวา เอวมสฺส วจนีโย “อหํ ภนฺเต สํเฆน ปฏิสารณีย\-กมฺมกโต สมฺมา วตฺตามิ, โลมํ ปาเตมิ, เนตฺถารํ วตฺตามิ, ปฏิสารณียสฺส กมฺมสฺส ปฏิปฺปสฺสทฺธิํ \csromanfolio{47}ยาจามี”ติ. ทุติยมฺปิ ยาจิตพฺพา. ตติยมฺปิ ยาจิตพฺพา. พฺยตฺเตน ภิกฺขุนา ปฏิพเลน สํโฆ ญาเปตพฺโพ\csromandash{}}

\prose{“สุณาตุ เม ภนฺเต สํโฆ, อยํ สุธมฺโม ภิกฺขุ สํเฆน ปฏิสารณีย\-กมฺมกโต สมฺมา วตฺตติ, โลมํ ปาเตติ, เนตฺถารํ วตฺตติ, ปฏิสารณียสฺส กมฺมสฺส ปฏิปฺปสฺสทฺธิํ ยาจติ, ยทิ สํฆสฺส ปตฺตกลฺลํ, สํโฆ สุธมฺมสฺส ภิกฺขุโน ปฏิสารณีย\-กมฺมํ ปฏิปฺปสฺสมฺเภยฺย, เอสา ญตฺติ.}

\prose{สุณาตุ เม ภนฺเต สํโฆ, อยํ สุธมฺโม ภิกฺขุ สํเฆน ปฏิสารณีย\-กมฺมกโต สมฺมา วตฺตติ, โลมํ ปาเตติ, เนตฺถารํ วตฺตติ, ปฏิสารณียสฺส กมฺมสฺส ปฏิปฺปสฺสทฺธิํ ยาจติ, สํโฆ สุธมฺมสฺส ภิกฺขุโน ปฏิสารณีย\-กมฺมํ ปฏิปฺปสฺสมฺเภติ, ยสฺสายสฺมโต ขมติ สุธมฺมสฺส ภิกฺขุโน ปฏิสารณียสฺส กมฺมสฺส ปฏิปฺปสฺสทฺธิ, โส ตุณฺหสฺส, ยสฺส นกฺขมติ, โส ภาเสยฺย.}

\prose{ทุติยมฺปิ เอตมตฺถํ วทามิ ฯเปฯ. ตติยมฺปิ เอตมตฺถํ วทามิ. สุณาตุ เม ภนฺเต สํโฆ, อยํ สุธมฺโม ภิกฺขุ สํเฆน ปฏิสารณีย\-กมฺมกโต สมฺมา วตฺตติ, โลมํ ปาเตติ, เนตฺถารํ วตฺตติ, ปฏิสารณียสฺส กมฺมสฺส ปฏิปฺปสฺสทฺธิํ ยาจติ, สํโฆ สุธมฺมสฺส ภิกฺขุโน ปฏิสารณีย\-กมฺมํ ปฏิปฺปสฺสมฺเภติ, ยสฺสายสฺมโต ขมติ สุธมฺมสฺส ภิกฺขุโน ปฏิสารณียสฺส กมฺมสฺส ปฏิปฺปสฺสทฺธิ, โส ตุณฺหสฺส, ยสฺส นกฺขมติ, โส ภาเสยฺย.}

\prose{ปฏิปฺปสฺสทฺธํ สํเฆน สุธมฺมสฺส ภิกฺขุโน ปฏิสารณีย\-กมฺมํ, ขมติ สํฆสฺส, ตสฺมา ตุณฺหี, เอวเมตํ ธารยามี”ติ.}

\vspace{\tipitakaparskip}
\csromancenter{ปฏิสารณีย\-กมฺมํ นิฏฺฐิตํ จตุตฺถํ.}
\csromanmatikaanchor{mtk.30}
\csromantocmarkref{section}{5. อาปตฺติยา อทสฺสเน อุกฺเขปนียกมฺม}{mtk.30}
\bookmark[dest={mtk.30},level=1]{5. อาปตฺติยา อทสฺสเน อุกฺเขปนียกมฺม}
\csromanheader{1}{5. อาปตฺติยา อทสฺสเน อุกฺเขปนีย\-กมฺม}
\proseitem{46}{เตน สมเยน พุทฺโธ ภควา โกสมฺพิยํ วิหรติ โฆสิตาราเม. เตน โข ปน สมเยน อายสฺมา ฉนฺโน อาปตฺติํ อาปชฺชิตฺวา น อิจฺฉติ อาปตฺติํ ปสฺสิตุํ. เย เต ภิกฺขู อปฺปิจฺฉา ฯเปฯ เต \csromanfolio{48}อุชฺฌายนฺติ ขิยฺยนฺติ วิปาเจนฺติ “กถํ หิ นาม อายสฺมา ฉนฺโน อาปตฺติํ อาปชฺชิตฺวา น อิจฺฉิสฺสติ อาปตฺติํ ปสฺสิตุนฺ”ติ. อถ โข เต ภิกฺขู ภควโต เอตมตฺถํ อาโรเจสุํ.}

\prose{อถ โข ภควา เอตสฺมิํ นิทาเน เอตสฺมิํ ปกรเณ ภิกฺขุสํฆํ สนฺนิปาตาเปตฺวา ภิกฺขู ปฏิปุจฺฉิ “สจฺจํ กิร ภิกฺขเว ฉนฺโน ภิกฺขุ อาปตฺติํ อาปชฺชิตฺวา น อิจฺฉติ อาปตฺติํ ปสฺสิตุนฺ”ติ. สจฺจํ ภควาติ. วิครหิ พุทฺโธ ภควา อนนุจฺฉวิกํ ฯเปฯ. กถํ หิ นาม โส ภิกฺขเว โมฆปุริโส อาปตฺติํ อาปชฺชิตฺวา น อิจฺฉิสฺสติ อาปตฺติํ ปสฺสิตุํ, เนตํ ภิกฺขเว อปฺปสนฺนานํ วา ปสฺสทาย ฯเปฯ วิครหิตฺวา ฯเปฯ ธมฺมิํ กถํ กตฺวา ภิกฺขู อามนฺเตสิ\csromandash{} เตน หิ ภิกฺขเว สํโฆ ฉนฺนสฺส ภิกฺขุโน อาปตฺติยา อทสฺสเน อุกฺเขปนีย\-กมฺมํ กโรตุ อสมฺโภคํ สํเฆน. เอวญฺจ ปน ภิกฺขเว กาตพฺพํ, ปฐมํ ฉนฺโน ภิกฺขุ โจเทตพฺโพ, โจเทตฺวา สาเรตพฺโพ, สาเรตฺวา อาปตฺติํ อาโรเปตพฺโพ, อาปตฺติํ อาโรเปตฺวา พฺยตฺเตน ภิกฺขุนา ปฏิพเลน สํโฆ ญาเปตพฺโพ\csromandash{}}

\proseitem{47}{“สุณาตุ เม ภนฺเต สํโฆ, อยํ ฉนฺโน ภิกฺขุ อาปตฺติํ อาปชฺชิตฺวา น อิจฺฉติ อาปตฺติํ ปสฺสิตุํ, ยทิ สํฆสฺส ปตฺตกลฺลํ, สํโฆ ฉนฺนสฺส ภิกฺขุโน อาปตฺติยา อทสฺสเน อุกฺเขปนีย\-กมฺมํ กเรยฺย อสมฺโภคํ สํเฆน, เอสา ญตฺติ.}

\prose{สุณาตุ เม ภนฺเต สํโฆ, อยํ ฉนฺโน ภิกฺขุ อาปตฺติํ อาปชฺชิตฺวา น อิจฺฉติ อาปตฺติํ ปสฺสิตุํ, สํโฆ ฉนฺนสฺส ภิกฺขุโน อาปตฺติยา อทสฺสเน อุกฺเขปนีย\-กมฺมํ กโรติ อสมฺโภคํ สํเฆน, ยสฺสายสฺมโต ขมติ ฉนฺนสฺส ภิกฺขุโน อาปตฺติยา อทสฺสเน อุกฺเขปนียสฺส กมฺมสฺส กรณํ อสมฺโภคํ สํเฆน, โส ตุณฺหสฺส, ยสฺส นกฺขมติ, โส ภาเสยฺย.}

\prose{ทุติยมฺปิ เอตมตฺถํ วทามิ ฯเปฯ. ตติยมฺปิ เอตมตฺถํ วทามิ. สุณาตุ เม ภนฺเต สํโฆ, อยํ ฉนฺโน ภิกฺขุ อาปตฺติํ อาปชฺชิตฺวา น อิจฺฉติ อาปตฺติํ ปสฺสิตุํ. สํโฆ ฉนฺนสฺส ภิกฺขุโน อาปตฺติยา อทสฺสเน อุกฺเขปนีย\-กมฺมํ กโรติ อสมฺโภคํ สํเฆน, ยสฺสายสฺมโต ขมติ ฉนฺนสฺส ภิกฺขุโน อาปตฺติยา อทสฺสเน อุกฺเขปนียสฺส กมฺมสฺส \csromanfolio{49}กรณํ อสมฺโภคํ สํเฆน, โส ตุณฺหสฺส, ยสฺส นกฺขมติ, โส ภาเสยฺย.}

\prose{กตํ สํเฆน ฉนฺนสฺส ภิกฺขุโน อาปตฺติยา อทสฺสเน อุกฺเขปนีย\-กมฺมํ อสมฺโภคํ สํเฆน, ขมติ สํฆสฺส, ตสฺมา ตุณฺหี, เอวเมตํ ธารยามี”ติ.}

\prose{อาวาสปรมฺปรญฺจ ภิกฺขเว สํสถ “ฉนฺโน ภิกฺขุ สํเฆน อาปตฺติยา\csromansharedfootnote{d0af2f6c05f86da0}{ฉนฺโน ภิกฺขุ อาปตฺติยา (สี, ก)} อทสฺสเน อุกฺเขปนีย\-กมฺมกโต อสมฺโภคํ สํเฆนา”ติ.}

\csromanmatikaanchor{mtk.31}
\csromantocmarkref{subsection}{อธมฺมกมฺมทฺวาทสก}{mtk.31}
\bookmark[dest={mtk.31},level=2]{อธมฺมกมฺมทฺวาทสก}
\csromanheader{2}{อธมฺมกมฺม\-ทฺวาทสก}
\proseitem{48}{ตีหิ ภิกฺขเว องฺเคหิ สมนฺนาคตํ อาปตฺติยา อทสฺสเน อุกฺเขปนีย\-กมฺมํ อธมฺมกมฺมญฺจ โหติ อวินยกมฺมญฺจ ทุวูปสนฺตญฺจ. อสมฺมุขา กตํ โหติ, อปฺปฏิปุจฺฉา กตํ โหติ, อปฺปฏิญฺญาย กตํ โหติ. อิเมหิ โข ภิกฺขเว ตีหงฺเคหิ สมนฺนาคตํ อาปตฺติยา อทสาเน อุกฺเขปนีย\-กมฺมํ อธมฺมกมฺมญฺจ โหติ อวินยกมฺมญฺจ ทุวูปสนฺตญฺจ.}

\prosewithcloser{อปเรหิปิ ภิกฺขเว ตีหงฺเคหิ สมนฺนาคตํ อาปตฺติยา อทสฺสเน อุกฺเขปนีย\-กมฺมํ อธมฺมกมฺมญฺจ โหติ อวินยกมฺมญฺจ ทุวูปสนฺตญฺจ. อนาปตฺติยา กตํ โหติ, อเทสนาคามินิยา อาปตฺติยา กตํ โหติ, เทสิตาย อาปตฺติยา กตํ โหติ ฯเปฯ. อโจเทตฺวา กตํ โหติ, อสาเรตฺวา กตํ โหติ, อาปตฺติํ อนาโรเปตฺวา กตํ โหติ ฯเปฯ. อสมฺมุขา กตํ โหติ, อธมฺเมน กตํ โหติ, วคฺเคน กตํ โหติ ฯเปฯ. อปฺปฏิปุจฺฉา กตํ โหติ, อธมฺเมน กตํ โหติ, วคฺเคน กตํ โหติ ฯเปฯ. อปฺปฏิญฺญาย กตํ โหติ, อธมฺเมน กตํ โหติ, วคฺเคน กตํ โหติ ฯเปฯ. อนาปตฺติยา กตํ โหติ, อธมฺเมน กตํ โหติ, วคฺเคน กตํ โหติ ฯเปฯ. อเทสนาคามินิยา อาปตฺติยา กตํ โหติ, อธมฺเมน กตํ โหติ, วคฺเคน กตํ โหติ ฯเปฯ. เทสิตาย อาปตฺติยา กตํ โหติ, อธมฺเมน กตํ โหติ, วคฺเคน กตํ โหติ ฯเปฯ. อโจเทตฺวา กตํ โหติ, อธมฺเมน กตํ โหติ, วคฺเคน กตํ โหติ ฯเปฯ. อสาเรตฺวา กตํ โหติ, อธมฺเมน กตํ โหติ, วคฺเคน กตํ โหติ ฯเปฯ. อาปตฺติํ \csromanfolio{50}อนาโรเปตฺวา กตํ โหติ, อธมฺเมน กตํ โหติ, วคฺเคน กตํ โหติ. อิเมหิ โข ภิกฺขเว ตีหงฺเคหิ สมนฺนาคตํ อาปตฺติยา อทสฺสเน อุกฺเขปนีย\-กมฺมํ อธมฺมกมฺมญฺจ โหติ อวินยกมฺมญฺจ ทุวูปสนฺตญฺจ.}{อาปตฺติยา อทสฺสเน อุกฺเขปนีย\-กมฺเม อธมฺมกมฺม\-ทฺวาทสกํ นิฏฺฐิตํ.}
\csromanmatikaanchor{mtk.32}
\csromantocmarkref{subsection}{ธมฺมกมฺมทฺวาทสก}{mtk.32}
\bookmark[dest={mtk.32},level=2]{ธมฺมกมฺมทฺวาทสก}
\csromanheader{2}{ธมฺมกมฺม\-ทฺวาทสก}
\proseitem{49}{ตีหิ ภิกฺขเว องฺเคหิ สมนฺนาคตํ อาปตฺติยา อทสฺสเน อุกฺเขปนีย\-กมฺมํ ธมฺมกมฺมญฺจ โหติ วินยกมฺมญฺจ สุวูปสนฺตญฺจ. สมฺมุขา กตํ โหติ, ปฏิปุจฺฉา กตํ โหติ, ปฏิญฺญาย กตํ โหติ. อิเมหิ โข ภิกฺขเว ตีหงฺเคหิ สมนฺนาคตํ อาปตฺติยา อทสฺสเน อุกฺเขปนีย\-กมฺมํ ธมฺมกมฺมญฺจ โหติ วินยกมฺมญฺจ สุวูปสนฺตญฺจ.}

\prosewithcloser{อปเรหิปิ ภิกฺขเว ตีหงฺเคหิ สมนฺนาคตํ ฯเปฯ. อาปตฺติยา กตํ โหติ, เทสนาคามินิยา อาปตฺติยา กตํ โหติ, อเทสิตาย อาปตฺติยา กตํ โหติ ฯเปฯ. โจเทตฺวา กตํ โหติ, สาเรตฺวา กตํ โหติ, อาปตฺติํ อาโรเปตฺวา กตํ โหติ ฯเปฯ. สมฺมุขา กตํ โหติ, ธมฺเมน กตํ โหติ, สมคฺเคน กตํ โหติ ฯเปฯ. ปฏิปุจฺฉา กตํ โหติ, ธมฺเมน กตํ โหติ, สมคฺเคน กตํ โหติ ฯเปฯ. ปฏิญฺญาย กตํ โหติ, ธมฺเมน กตํ โหติ, สมคฺเคน กตํ โหติ ฯเปฯ. อาปตฺติยา กตํ โหติ, ธมฺเมน กตํ โหติ, สมคฺเคน กตํ โหติ ฯเปฯ. เทสนาคามินิยา อาปตฺติยา กตํ โหติ, ธมฺเมน กตํ โหติ, สมคฺเคน กตํ โหติ ฯเปฯ. อเทสิตาย อาปตฺติยา กตํ โหติ, ธมฺเมน กตํ โหติ, สมคฺเคน กตํ โหติ ฯเปฯ. โจเทตฺวา กตํ โหติ, ธมฺเมน กตํ โหติ, สมคฺเคน กตํ โหติ ฯเปฯ. สาเรตฺวา กตํ โหติ, ธมฺเมน กตํ โหติ, สมคฺเคน กตํ โหติ ฯเปฯ. อาปตฺติํ อาโรเปตฺวา กตํ โหติ, ธมฺเมน กตํ โหติ, สมคฺเคน กตํ โหติ. อิเมหิ โข ภิกฺขเว ตีหงฺเคหิ สมนฺนาคตํ อาปตฺติยา อทสฺสเน อุกฺเขปนีย\-กมฺมํ ธมฺมกมฺมญฺจ โหติ วินยกมฺมญฺจ สุวูปสนฺตญฺจ.}{อาปตฺติยา อทสฺสเน อุกฺเขปนีย\-กมฺเม ธมฺมกมฺม\-ทฺวาทสกํ นิฏฺฐิตํ.}
\csromanmatikaanchor{mtk.33}
\csromantocmarkref{subsection}{อากงฺขมานฉกฺก}{mtk.33}
\bookmark[dest={mtk.33},level=2]{อากงฺขมานฉกฺก}
\proseitem{50}{\csromanfolio{51}\csromansymbolfootnote{*}{วิ 5. 220 ปิฏฺเฐปิ.}ตีหิ ภิกฺขเว องฺเคหิ สมนฺนาคตสฺส ภิกฺขุโน อากงฺขมาโน สํโฆ อาปตฺติยา อทสฺสเน อุกฺเขปนีย\-กมฺมํ กเรยฺย. ภณฺฑนการโก โหติ กลหการโก วิวาทการโก ภสฺสการโก สํเฆ อธิกรณ\-การโก, พาโล โหติ อพฺยตฺโต อาปตฺติพหุโล อนปทาโน, คิหิสํสฏฺโฐ วิหรติ อนนุโลมิเกหิ คิหิสํสคฺเคหิ. อิเมหิ โข ภิกฺขเว ตีหงฺเคหิ สมนฺนาคตสฺส ภิกฺขุโน อากงฺขมาโน สํโฆ อาปตฺติยา อทสฺสเน อุกฺเขปนีย\-กมฺมํ กเรยฺย.}

\prose{อปเรหิปิ ภิกฺขเว ตีหงฺเคหิ สมนฺนาคตสฺส ภิกฺขุโน อากงฺขมาโน สํโฆ อาปตฺติยา อทสฺสเน อุกฺเขปนีย\-กมฺมํ กเรยฺย. อธิสีเล สีลวิปนฺโน โหติ, อชฺฌาจาเร อาจารวิปนฺโน โหติ, อติทิฏฺฐิยา ทิฏฺฐิวิปนฺโน โหติ. อิเมหิ โข ภิกฺขเว ตีหงฺเคหิ สมนฺนาคตสฺส ภิกฺขุโน อากงฺขมาโน สํโฆ อาปตฺติยา อทสฺสเน อุกฺเขปนีย\-กมฺมํ กเรยฺย.}

\prose{อปเรหิปิ ภิกฺขเว ตีหงฺเคหิ สมนฺนาคตสฺส ภิกฺขุโน อากงฺขมาโน สํโฆ อาปตฺติยา อทสฺสเน อุกฺเขปนีย\-กมฺมํ กเรยฺย. พุทฺธสฺส อวณฺณํ ภาสติ, ธมฺมสฺส อวณฺณํ ภาสติ, สํฆสฺส อวณฺณํ ภาสติ. อิเมหิ โข ภิกฺขเว ตีหงฺเคหิ สมนฺนาคตสฺส ภิกฺขุโน อากงฺขมาโน สํโฆ อาปตฺติยา อทสฺสเน อุกฺเขปนีย\-กมฺมํ กเรยฺย.}

\prose{ติณฺณํ ภิกฺขเว ภิกฺขูนํ อากงฺขมาโน สํโฆ อาปตฺติยา อทสฺสเน อุกฺเขปนีย\-กมฺมํ กเรยฺย. เอโก ภณฺฑนการโก โหติ กลหการโก วิวาทการโก ภสฺสการโก สํเฆ อธิกรณ\-การโก, เอโก พาโล โหติ อพฺยตฺโต อาปตฺติพหุโล อนปทาโน, เอโก คิหิสํสฏฺโฐ วิหรติ อนนุโลมิเกหิ คิหิสํสคฺเคหิ. อิเมสํ โข ภิกฺขเว ติณฺณํ ภิกฺขูนํ อากงฺขมาโน สํโฆ อาปตฺติยา อทสฺสเน อุกฺเขปนีย\-กมฺมํ กเรยฺย.}

\prose{อปเรสมฺปิ ภิกฺขเว ติณฺณํ ภิกฺขูนํ อากงฺขมาโน สํโฆ อาปตฺติยา อทสฺสเน อุกฺเขปนีย\-กมฺมํ กเรยฺย. เอโก อธิสีเล สีลวิปนฺโน \csromanfolio{52}โหติ, เอโก อชฺฌาจาเร อาจารวิปนฺโน โหติ, เอโก อติทิฏฺฐิยา ทิฏฺฐิวิปนฺโน โหติ. อิเมสํ โข ภิกฺขเว ติณฺณํ ภิกฺขูนํ อากงฺขมาโน สํโฆ อาปตฺติยา อทสฺสเน อุกฺเขปนีย\-กมฺมํ กเรยฺย.}

\prose{อปเรสมฺปิ ภิกฺขเว ติณฺณํ ภิกฺขูนํ อากงฺขมาโน สํโฆ อาปตฺติยา อทสฺสเน อุกฺเขปนีย\-กมฺมํ กเรยฺย. เอโก พุทฺธสฺส อวณฺณํ ภาสติ, เอโก ธมฺมสฺส อวณฺณํ ภาสติ, เอโก สํฆสฺส อวณฺณํ ภาสติ. อิเมสํ โข ภิกฺขเว ติณฺณํ ภิกฺขูนํ อากงฺขมาโน สํโฆ อาปตฺติยา อทสฺสเน อุกฺเขปนีย\-กมฺมํ กเรยฺย.}

\csromanheader{2}{อาปตฺติยา อทสฺสเน อุกฺเขปนีย\-กมฺเม}
\nitthitam{อากงฺขมาน\-ฉกฺกํ นิฏฺฐิตํ.}
\csromanmatikaanchor{mtk.34}
\csromantocmarkref{subsection}{เตจตฺตาลีสวตฺต}{mtk.34}
\bookmark[dest={mtk.34},level=2]{เตจตฺตาลีสวตฺต}
\csromanheader{2}{เตจตฺตาลีส\-วตฺต}
\proseitem{51}{อาปตฺติยา อทสฺสเน อุกฺเขปนีย\-กมฺมกเตน ภิกฺขเว ภิกฺขุนา สมฺมา วตฺติตพฺพํ. ตตฺรายํ สมฺมาวตฺตนา\csromandash{}น อุปสมฺ\-ปาเท\-ตพฺพํ, น นิสฺสโย ทาตพฺโพ, น สามเณโร อุปฏฺฐาเปตพฺโพ, น ภิกฺขุโนวาทก\-สมฺมุติ สาทิตพฺพา, สมฺมเตนปิ ภิกฺขุนิโย น โอวทิตพฺพา, ยาย อาปตฺติยา สํเฆน อาปตฺติยา อทสฺสเน อุกฺเขปนีย\-กมฺมํ กตํ โหติ, สา อาปตฺติ น อาปชฺชิตพฺพา, อญฺญา วา ตาทิสิกา, ตโต วา ปาปิฏฺฐตรา, กมฺมํ น ครหิตพฺพํ, กมฺมิกา น ครหิตพฺพา, น ปกตตฺตสฺส ภิกฺขุโน อภิวาทนํ ปจฺจุฏฺฐานํ อญฺชลิกมฺมํ สามีจิกมฺมํ อาสนาภิหาโร เสยฺยาภิหาโร ปาโททกํ ปาทปีฐํ ปาทกถลิกํ ปตฺตจีวร\-ปฺปฏิคฺคหณํ นหาเน ปิฏฺฐิ\-ปริกมฺมํ สาทิตพฺพํ, น ปกตตฺโต ภิกฺขุ สีลวิปตฺติยา อนุทฺธํเส\-ตพฺโพ, น ปตฺติยา อนุทฺธํเส\-ตพฺโพ, น ทิฏฺฐิ\-วิปตฺติยา อนุทฺธํเส\-ตพฺโพ, น อาชีววิปตฺติยา ทฺธํเสตพฺโพ, น ภิกฺขุ ภิกฺขูหิ เภเทตพฺโพ\csromansharedfootnote{27562eb5c6815e97}{น ภิกฺขู ภิกฺขูหิ เภเทตพฺพา (สฺยา)}, น คิหิทฺธโช ธาเรตพฺโพ, น ติตฺถิยทฺธโช ธาเรตพฺโพ, น ติตฺถิยา เสวิตพฺพา, ภิกฺขู เสวิตพฺพา,ภิกฺขุ\-สิกฺขาย สิกฺขิตพฺพํ\csromansharedfootnote{1575ef6bbe232546}{ภิกฺขุสิกฺขา สิกฺขิตพฺพา (สฺยา)}, น ปกตตฺเตน ภิกฺขุนา สทฺธิํ เอกจฺฉนฺเน อาวาเส วตฺถพฺพํ, น เอกจฺฉนฺเน \csromanfolio{53}อนาวาเส วตฺถพฺพํ, น เอกจฺฉนฺเน อาวาเส วา อนาวาเส วา วตฺถพฺพํ, ปกตตฺตํ ภิกฺขุํ ทิสฺวา อาสนา วุฏฺฐาตพฺพํ, น ปกตตฺโต ภิกฺขุ อาสาเทตพฺโพ อนฺโต วา พหิ วา, น ปกตตฺตสฺส ภิกฺขุโน อุโปสโถ ฐเปตพฺโพ, น ปวารณา ฐเปตพฺพา, น สวจนียํ กาตพฺพํ, น อนุวาโท ปฏฺฐเปตพฺโพ, น โอกาโส กาเรตพฺโพ, น โจเทตพฺโพ, น สาเรตพฺโพ, น ภิกฺขูหิ สมฺป\-โยเช\-ตพฺพนฺติ.}

\csromanheader{2}{อาปตฺติยา อทสฺสเน อุกฺเขปนีย\-กมฺเม}
\nitthitam{เตจตฺตาลีส\-วตฺตํ นิฏฺฐิตํ.}
\proseitem{52}{อถ โข สํโฆ ฉนฺนสฺส ภิกฺขุโน อาปตฺติยา อทสฺสเน อุกฺเขปนีย\-กมฺมํ อกาสิ อสมฺโภคํ สํเฆน. โส สํเฆน อาปตฺติยา อทสฺสเน อุกฺเขปนีย\-กมฺมกโต ตมฺหา อาวาสา อญฺญํ อาวาสํ อคมาสิ. ตตฺถ ภิกฺขู เนว อภิวาเทสุํ, น ปจฺจุฏฺเฐสุํ, น อญฺชลิกมฺมํ น สามีจิกมฺมํ อกํสุ, น สกฺกริํสุ, น ครุํ กริํสุ\csromansharedfootnote{296ff41d10e20064}{น ครุกริํสุ (ก)}, น มาเนสุํ, น ปูเชสุํ. โส ภิกฺขูหิ อสกฺกริยมาโน อครุกริยมาโน อมานิยมาโน อปูชิยมาโน อสกฺการปกโต ตมฺหาปิ อาวาสา อญฺญํ อาวาสํ อคมาสิ. ตตฺถปิ ภิกฺขู เนว อภิวาเทสุํ, น ปจฺจุฏฺเฐสุํ, น อญฺชลิกมฺมํ น สามีจิกมฺมํ อกํสุ, น สกฺกริํสุ, น ครุํ กริํสุ, น มาเนสุํ, น ปูเชสุํ. โส ภิกฺขูหิ อสกฺกริยมาโน อครุกริยมาโน อมานิยมาโน อปูชิยมาโน อสกฺการปกโต ตมฺหาปิ อาวาสา อญฺญํ อาวาสํ อคมาสิ. ตตฺถปิ ภิกฺขู เนว อภิวาเทสุํ, น ปจฺจุฏฺเฐสุํ, น อญฺชลิกมฺมํ น สามีจิกมฺมํ อกํสุ, น สกฺกริํสุ, น ครุํ กริํสุ, น มาเนสุํ, น ปูเชสุํ. โส ภิกฺขูหิ อสกฺกริยมาโน อครุกริยมาโน อมานิยมาโน อปูชิยมาโน อสกฺการปกโต ปุนเทว โกสมฺพิํ ปจฺจาคญฺฉิ. โส สมฺมา วตฺตติ, โลมํ ปาเตติ, เนตฺถารํ วตฺตติ, ภิกฺขู อุปสงฺกมิตฺวา เอวํ วเทติ “อหํ อาวุโส สํเฆน อาปตฺติยา อทสฺสเน อุกฺเขปนีย\-กมฺมกโต สมฺมา วตฺตามิ, โลมํ ปาเตมิ, เนตฺถารํ วตฺตามิ, กถํ นุ โข มยา ปฏิปชฺชิตพฺพนฺ”ติ. ภิกฺขู ภควโต เอตมตฺถํ อาโรเจสุํ ฯเปฯ. เตน หิ ภิกฺขเว สํโฆ ฉนฺนสฺส ภิกฺขุโน อาปตฺติยา อทสฺสเน อุกฺเขปนีย\-กมฺมํ ปฏิปฺปสฺสมฺเภตุ.}

\csromanmatikaanchor{mtk.35}
\csromantocmarkref{subsection}{นปฺปฏิปฺปสฺสมฺเภตพฺพเตจตฺตาลีสก}{mtk.35}
\bookmark[dest={mtk.35},level=2]{นปฺปฏิปฺปสฺสมฺเภตพฺพเตจตฺตาลีสก}
\csromanheader{2}{1. นปฺปฏิปฺปสฺสมฺเภตพฺพ\-เตจตฺตาลีสก}
\proseitem{53}{\csromanfolio{54}ปญฺจหิ ภิกฺขเว องฺเคหิ สมนฺนาคตสฺส ภิกฺขุโน อาปตฺติยา อทสฺสเน อุกฺเขปนีย\-กมฺมํ นปฺปฏิปฺปสฺสมฺเภตพฺพํ. อุปสมฺปาเทติ, นิสฺสยํ เทติ, สามเณรํ อุปฏฺฐาเปติ, ภิกฺขุโนวาทก\-สมฺมุติํ สาทิยติ, สมฺมโตปิ ภิกฺขุนิโย โอวทติ. อิเมหิ โข ภิกฺขเว ปญฺจหงฺเคหิ สมนฺนาคตสฺส ภิกฺขุโน อาปตฺติยา อทสฺสเน อุกฺเขปนีย\-กมฺมํ นปฺปฏิปฺปสฺสมฺเภตพฺพํ. (5)}

\prose{\csromansymbolfootnote{*}{วิ 5. 317 ปิฏฺเฐปิ.}อปเรหิปิ ภิกฺขเว ปญฺจหงฺเคหิ สมนฺนาคตสฺส ภิกฺขุโน อาปตฺติยา อทสฺสเน อุกฺเขปนีย\-กมฺมํ นปฺปฏิปฺปสฺสมฺเภตพฺพํ. ยาย อาปตฺติยา สํเฆน อาปตฺติยา อทสฺสเน อุกฺเขปนีย\-กมฺมํ กตํ โหติ, ตํ อาปตฺติํ อาปชฺชติ, อญฺญํ วา ตาทิสิกํ, ตโต วา ปาปิฏฺฐตรํ, กมฺมํ ครหติ, กมฺมิเก ครหติ. อิเมหิ โข ภิกฺขเว ปญฺจหงฺเคหิ สมนฺนาคตสฺส ภิกฺขุโน อาปตฺติยา อทสฺสเน อุกฺเขปนีย\-กมฺมํ นปฺปฏิปฺปสฺสมฺเภตพฺพํ. (10)}

\prose{อปเรหิปิ ภิกฺขเว ปญฺจหงฺเคหิ สมนฺนาคตสฺส ภิกฺขุโน อาปตฺติยา อทสฺสเน อุกฺเขปนีย\-กมฺมํ นปฺปฏิปฺปสฺสมฺเภตพฺพํ. ปกตตฺตสฺส ภิกฺขุโน อภิวาทนํ ปจฺจุฏฺฐานํ อญฺชลิกมฺมํ สามีจิกมฺมํ อาสนาภิหารํ สาทิยติ. อิเมหิ โข ภิกฺขเว ปญฺจหงฺเคหิ สมนฺนาคตสฺส ภิกฺขุโน อาปตฺติยา อทสฺสเน อุกฺเขปนีย\-กมฺมํ นปฺปฏิปฺปสฺสมฺเภตพฺพํ. (15)}

\prose{อปเรหิปิ ภิกฺขเว ปญฺจหงฺเคหิ สมนฺนาคตสฺส ภิกฺขุโน อาปตฺติยา อทสฺสเน อุกฺเขปนีย\-กมฺมํ นปฺปฏิปฺปสฺสมฺเภตพฺพํ. ปกตตฺตสฺส ภิกฺขุโน เสยฺยาภิหารํ ปาโททกํ ปาทปิฐํ ปาทกถลิกํ ปตฺตจีวร\-ปฺปฏิคฺคหณํ นหาเน ปิฏฺฐิ\-ปริกมฺมํ สาทิยติ. อิเมหิ โข ภิกฺขเว ปญฺจหงฺเคหิ สมนฺนาคตสฺส ภิกฺขุโน อาปตฺติยา อทสฺสเน อุกฺเขปนีย\-กมฺมํ นปฺปฏิปฺปสฺสมฺเภตพฺพํ. (20)}

\prose{อปเรหิปิ ภิกฺขเว ปญฺจหงฺเคหิ สมนฺนาคตสฺส ภิกฺขุโน อาปตฺติยา อทสฺสเน อุกฺเขปนีย\-กมฺมํ นปฺปฏิปฺปสฺสมฺเภตพฺพํ. ปกตตฺตํ ภิกฺขุํ สีลวิปตฺติยา อนุทฺธํเสติ, อาจารวิปตฺติยา อนุทฺธํเสติ, ทิฏฺฐิ\-วิปตฺติยา อนุทฺธํเสติ, อาชีววิปตฺติยา อนุทฺธํเสติ, ภิกฺขุํ\csromansharedfootnote{80ea68b0b5d0714a}{ภิกฺขู (สี, สฺยา)} ภิกฺขูหิ เภเทติ.}

\prose{\csromanfolio{55}อิเมหิ โข ภิกฺขเว ปญฺจหงฺเคหิ สมนฺนาคตสฺส ภิกฺขุโน อาปตฺติยา อทสฺสเน อุกฺเขปนีย\-กมฺมํ นปฺปฏิปฺปสฺสมฺเภตพฺพํ. (25)}

\prose{อปเรหิปิ ภิกฺขเว ปญฺจหงฺเคหิ สมนฺนาคตสฺส ภิกฺขุโน อาปตฺติยา อทสฺสเน อุกฺเขปนีย\-กมฺมํ นปฺปฏิปฺปสฺสมฺเภตพฺพํ. คิหิทฺธชํ ธาเรติ, ติตฺถิยทฺธชํ ธาเรติ, ติตฺถิเย เสวติ, ภิกฺขู น เสวติ, ภิกฺขุ\-สิกฺขาย น สิกฺขติ. อิเมหิ โข ภิกฺขเว ปญฺจหงฺเคหิ สมนฺนาคตสฺส ภิกฺขุโน อาปตฺติยา อทสฺสเน อุกฺเขปนีย\-กมฺมํ นปฺปฏิปฺปสฺสมฺเภตพฺพํ. (30)}

\prose{อปเรหิปิ ภิกฺขเว ปญฺจหงฺเคหิ สมนฺนาคตสฺส ภิกฺขุโน อาปตฺติยา อทสฺสเน อุกฺเขปนีย\-กมฺมํ นปฺปฏิปฺปสฺสมฺเภตพฺพํ. ปกตตฺเตน ภิกฺขุนา สทฺธิํ เอกจฺฉนฺเน อาวาเส วสติ, เอกจฺฉนฺเน อนาวาเส วสติ, เอกจฺฉนฺเน อาวาเส วา อนาวาเส วา วสติ, ปกตตฺตํ ภิกฺขุํ ทิสฺวา อาสนา น วุฏฺฐาติ, ปกตตฺตํ ภิกฺขุํ อาสาเทติ อนฺโต วา พหิ วา. อิเมหิ โข ภิกฺขเว ปญฺจหงฺเคหิ สมนฺนาคตสฺส ภิกฺขุโน อาปตฺติยา อทสฺสเน อุกฺเขปนีย\-กมฺมํ นปฺปฏิปฺปสฺสมฺเภตพฺพํ. (35)}

\prose{อฏฺฐหิ ภิกฺขเว องฺเคหิ สมนฺนาคตสฺส ภิกฺขุโน อาปตฺติยา อทสฺสเน อุกฺเขปนีย\-กมฺมํ นปฺปฏิปฺปสฺสมฺเภตพฺพํ. ปกตตฺตสฺส ภิกฺขุโน อุโปสถํ ฐเปติ, ปวารณํ ฐเปติ, สวจนียํ กโรติ, อนุวาทํ ปฏฺฐเปติ, โอกาสํ กาเรติ, โจเทติ, สาเรติ, ภิกฺขูหิ สมฺปโยเชติ. อิเมหิ โข ภิกฺขเว อฏฺฐหงฺเคหิ สมนฺนาคตสฺส ภิกฺขุโน อาปตฺติยา อทสฺสเน อุกฺเขปนีย\-กมฺมํ นปฺปฏิปฺปสฺสมฺเภตพฺพํ. (43)}

\csromanheader{2}{อาปตฺติยา อทสฺสเน อุกฺเขปนีย\-กมฺเม}
\nitthitam{นปฺปฏิปฺปสฺสมฺเภตพฺพ\-เตจตฺตาลีสกํ นิฏฺฐิตํ.}
\csromanmatikaanchor{mtk.36}
\csromantocmarkref{subsection}{ปฏิปฺปสฺสมฺเภตพฺพเตจตฺตาลีสก}{mtk.36}
\bookmark[dest={mtk.36},level=2]{ปฏิปฺปสฺสมฺเภตพฺพเตจตฺตาลีสก}
\csromanheader{2}{ปฏิปฺปสฺสมฺเภตพฺพ\-เตจตฺตาลีสก}
\proseitem{54}{ปญฺจหิ ภิกฺขเว องฺเคหิ สมนฺนาคตสฺส ภิกฺขุโน อาปตฺติยา อทสฺสเน อุกฺเขปนีย\-กมฺมํ ปฏิปฺปสฺสมฺเภตพฺพํ. น อุปสมฺปาเทติ, น นิสฺสยํ เทติ, น สามเณรํ อุปฏฺฐาเปติ, น ภิกฺขุโนวาทก\-สมฺมุติํ สาทิยติ, สมฺมโตปิ ภิกฺขุนิโย น โอวทติ. อิเมหิ โข ภิกฺขเว ปญฺจหงฺเคหิ \csromanfolio{56}สมนฺนาคตสฺส ภิกฺขุโน อาปตฺติยา อทสฺสเน อุกฺเขปนีย\-กมฺมํ ปฏิปฺปสฺสมฺเภตพฺพํ. (5)}

\prose{อปเรหิปิ ภิกฺขเว ปญฺจหงฺเคหิ สมนฺนาคตสฺส ภิกฺขุโน อาปตฺติยา อทสฺสเน อุกฺเขปนีย\-กมฺมํ ปฏิปฺปสฺสมฺเภตพฺพํ. ยาย อาปตฺติยา สํเฆน อาปตฺติยา อทสฺสเน อุกฺเขปนีย\-กมฺมํ กตํ โหติ, ตํ อาปตฺติํ น อาปชฺชติ, อญฺญํ วา ตาทิสิกํ, ตโต วา ปาปิฏฺฐตรํ, กมฺมํ น ครหติ, กมฺมิเก น ครหติ. อิเมหิ โข ภิกฺขเว ปญฺจหงฺเคหิ สมนฺนาคตสฺส ภิกฺขุโน อาปตฺติยา อทสฺสเน อุกฺเขปนีย\-กมฺมํ ปฏิปฺปสฺสมฺเภตพฺพํ. (10)}

\prose{อปเรหิปิ ภิกฺขเว ปญฺจหงฺเคหิ สมนฺนาคตสฺส ภิกฺขุโน อาปตฺติยา อทสฺสเน อุกฺเขปนีย\-กมฺมํ ปฏิปฺปสฺสมฺเภตพฺพํ. น ปกตตฺตสฺส ภิกฺขุโน อภิวาทนํ ปจฺจุฏฺฐานํ อญฺชลิกมฺมํ สามีจิกมฺมํ อาสนาภิหารํ สาทิยติ. อิเมหิ โข ภิกฺขเว ปญฺจหงฺเคหิ สมนฺนาคตสฺส ภิกฺขุโน อาปตฺติยา อทสฺสเน อุกฺเขปนีย\-กมฺมํ ปฏิปฺปสฺสมฺเภตพฺพํ. (15)}

\prose{อปเรหิปิ ภิกฺขเว ปญฺจหงฺเคหิ สมนฺนาคตสฺส ภิกฺขุโน อาปตฺติยา อทสฺสเน อุกฺเขปนีย\-กมฺมํ ปฏิปฺปสฺสมฺเภตพฺพํ. น ปกตตฺตสฺส ภิกฺขุโน เสยฺยาภิหารํ ปาโททกํ ปาทปีฐํ ปาทกถลิกํ ปตฺตจีวร\-ปฺปฏิคฺคหณํ นหาเน ปิฏฺฐิ\-ปริกมฺมํ สาทิยติ. อิเมหิ โข ภิกฺขเว ปญฺจหงฺเคหิ สมนฺนาคตสฺส ภิกฺขุโน อาปตฺติยา อทสฺสเน อุกฺเขปนีย\-กมฺมํ ปฏิปฺปสฺสมฺเภตพฺพํ. (20)}

\prose{อปเรหิปิ ภิกฺขเว ปญฺจหงฺเคหิ สมนฺนาคตสฺส ภิกฺขุโน อาปตฺติยา อทสฺสเน อุกฺเขปนีย\-กมฺมํ ปฏิปฺปสฺสมฺเภตพฺพํ. น ปกตตฺตํ ภิกฺขุํ สีลวิปตฺติยา อนุทฺธํเสติ, น อาจารวิปตฺติยา อนุทฺธํเสติ, น ทิฏฺฐิ\-วิปตฺติยา อนุทฺธํเสติ, น อาชีววิปตฺติยา อนุทฺธํเสติ, น ภิกฺขุํ ภิกฺขูหิ เภเทติ. อิเมหิ โข ภิกฺขเว ปญฺจหงฺเคหิ สมนฺนาคตสฺส ภิกฺขุโน อาปตฺติยา อทสฺสเน อุกฺเขปนีย\-กมฺมํ ปฏิปฺปสฺสมฺเภตพฺพํ. (25)}

\prose{อปเรหิปิ ภิกฺขเว ปญฺจหงฺเคหิ สมนฺนาคตสฺส ภิกฺขุโน อาปตฺติยา อทสฺสเน อุกฺเขปนีย\-กมฺมํ ปฏิปฺปสฺสมฺเภตพฺพํ. น คิหิทฺธชํ ธาเรติ, น ติตฺถิยทฺธชํ ธาเรติ, น ติตฺถิเย เสวติ, ภิกฺขู เสวติ, \csromanfolio{57}ภิกฺขุ\-สิกฺขาย สิกฺขติ. อิเมหิ โข ภิกฺขเว ปญฺจหงฺเคหิ สมนฺนาคตสฺส ภิกฺขุโน อาปตฺติยา อทสฺสเน อุกฺเขปนีย\-กมฺมํ ปฏิปฺปสฺสมฺเภตพฺพํ. (30)}

\prose{อปเรหิปิ ภิกฺขเว ปญฺจหงฺเคหิ สมนฺนาคตสฺสภิกฺขุโน อาปตฺติยา อทสฺสเน อุกฺเขปนีย\-กมฺมํ ปฏิปฺปสฺสมฺเภตพฺพํ. น ปกตตฺเตน ภิกฺขุนา สทฺธิํ เอกจฺฉนฺเน อาวาเส วสติ, น เอกจฺฉนฺเน อนาวาเส วสติ, น เอกจฺฉนฺเน อาวาเส วา อนาวาเส วา วสติ, ปกตตฺตํ ภิกฺขุํ ทิสฺวา อาสนา วุฏฺฐาติ, น ปกตตฺตํ ภิกฺขุํ อาสาเทติ อนฺโต วา พหิ วา. อิเมหิ โข ภิกฺขเว ปญฺจหงฺเคหิ สมนฺนาคตสฺส ภิกฺขุโน อาปตฺติยา อทสฺสเน อุกฺเขปนีย\-กมฺมํ ปฏิปฺปสฺสมฺเภตพฺพํ. (35)}

\prose{อฏฺฐหิ ภิกฺขเว องฺเคหิ สมนฺนาคตสฺส ภิกฺขุโน อาปตฺติยา อทสฺสเน อุกฺเขปนีย\-กมฺมํ ปฏิปฺปสฺสมฺเภตพฺพํ. น ปกตตฺตสฺส ภิกฺขุโน อุโปสถํ ฐเปติ, น ปวารณํ ฐเปติ, น สวจนียํ กโรติ, น อนุวาทํ ปฏฺฐเปติ, น โอกาสํ กาเรติ, น โจเทติ, น สาเรติ, น ภิกฺขูหิ สมฺปโยเชติ. อิเมหิ โข ภิกฺขเว อฏฺฐหงฺเคหิ สมนฺนาคตสฺส ภิกฺขุโน อาปตฺติยา อทสฺสเน อุกฺเขปนีย\-กมฺมํ ปฏิปฺปสฺสมฺเภตพฺพํ. (43)}

\csromanheader{2}{อาปตฺติยา อทสฺสเน อุกฺเขปนีย\-กมฺเม}
\nitthitam{ปฏิปฺปสฺสมฺเภตพฺพ\-เตจตฺตาลีสกํ นิฏฺฐิตํ.}
\proseitem{55}{เอวญฺจ ปน ภิกฺขเว ปฏิปฺปสฺสมฺเภตพฺพํ, เตน ภิกฺขเว ฉนฺเนน ภิกฺขุนา สํฆํ อุปสงฺกมิตฺวา เอกํสํ อุตฺตราสงฺคํ กริตฺวา วุฑฺฒานํ ภิกฺขูนํ ปาเท วนฺทิตฺวา อุกฺกุฏิกํ นิสีทิตฺวา อญฺชลิํ ปคฺคเหตฺวา เอวมสฺส วจนีโย “อหํ ภนฺเต สํเฆน อาปตฺติยา อทสฺสเน อุกฺเขปนีย\-กมฺมกโต สมฺมา วตฺตามิ, โลมํ ปาเตมิ, เนตฺถารํ วตฺตามิ, อาปตฺติยา อทสฺสเน อุกฺเขปนียสฺส กมฺมสฺส ปฏิปฺปสฺสทฺธิํ ยาจามี”ติ. ทุติยมฺปิ ยาจิตพฺพา. ตติยมฺปิ ยาจิตพฺพา. พฺยตฺเตน ภิกฺขุนา ปฏิพเลน สํโฆ ญาเปตพฺโพ\csromandash{}}

\prose{“สุณาตุ เม ภนฺเต สํโฆ, อยํ ฉนฺโน ภิกฺขุ สํเฆน อาปตฺติยา อทสฺสเน อุกฺเขปนีย\-กมฺมกโต สมฺมา วตฺตติ, โลมํ ปาเตติ, เนตฺถารํ วตฺตติ, อาปตฺติยา อทสฺสเน อุกฺเขปนียสฺส กมฺมสฺส ปฏิปฺปสฺสทฺธิํ \csromanfolio{58}ยาจติ, ยทิ สํฆสฺส ปตฺตกลฺลํ, สํโฆ ฉนฺนสฺส ภิกฺขุโน อาปตฺติยา อทสฺสเน อุกฺเขปนีย\-กมฺมํ ปฏิปฺปสฺสมฺเภยฺย, เอสา ญตฺติ.}

\prose{สุณาตุ เม ภนฺเต สํโฆ, อยํ ฉนฺโน ภิกฺขุ สํเฆน อาปตฺติยา อทสฺสเน อุกฺเขปนีย\-กมฺมกโต สมฺมา วตฺตติ, โลมํ ปาเตติ, เนตฺถารํ, วตฺตติ, อาปตฺติยา อทสฺสเน อุกฺเขปนียสฺส กมฺมสฺส ปฏิปฺปสฺสทฺธิํ ยาจติ, สํโฆ ฉนฺนสฺส ภิกฺขุโน อาปตฺติยา อทสฺสเน อุกฺเขปนีย\-กมฺมํ ปฏิปฺปสฺสมฺเภติ, ยสฺสายสฺมโต ขมติ ฉนฺนสฺส ภิกฺขุโน อาปตฺติยา อทสฺสเน อุกฺเขปนียสฺส กมฺมสฺส ปฏิปฺปสฺสทฺธิ, โส ตุณฺหสฺส, ยสฺส นกฺขมติ, โส ภาเสยฺย.}

\prose{ทุติยมฺปิ เอตมตฺถํ วทามิ ฯเปฯ. ตติยมฺปิ เอตมตฺถํ วทามิ. สุณาตุ เม ภนฺเต สํโฆ, อยํ ฉนฺโน ภิกฺขุ สํเฆน อาปตฺติยา อทสฺสเน อุกฺเขปนีย\-กมฺมกโต สมฺมา วตฺตติ, โลมํ ปาเตติ, เนตฺถารํ วตฺตติ, อาปตฺติยา อทสฺสเน อุกฺเขปนียสฺส กมฺมสฺส ปฏิปฺปสฺสทฺธิํ ยาจติ, สํโฆ ฉนฺนสฺส ภิกฺขุโน อาปตฺติยา อทสฺสเน อุกฺเขปนีย\-กมฺมํ ปฏิปฺปสฺสมฺเภติ, ยสฺสายสฺมโต ขมติ ฉนฺนสฺส ภิกฺขุโน อาปตฺติยา อทสฺสเน อุกฺเขปนียสฺส กมฺมสฺส ปฏิปฺปสฺสทฺธิ, โส ตุณฺหสฺส, ยสฺส นกฺขมติ, โส ภาเสยฺย.}

\prose{ปฏิปฺปสฺสทฺธํ สํเฆน ฉนฺนสฺส ภิกฺขุโน อาปตฺติยา อทสฺสเน อุกฺเขปนีย\-กมฺมํ, ขมติ สํฆสฺส, ตสฺมา ตุณฺหี, เอวเมตํ ธารยามี”ติ.}

\prose{อาปตฺติยา อทสฺสเน อุกฺเขปนีย\-กมฺมํ นิฏฺฐิตํ ปญฺจมํ.}

\csromanmatikaanchor{mtk.37}
\csromantocmarkref{section}{6. อาปตฺติยา อปฺปฏิกมฺเม อุกฺเขปนียกมฺม ...}{mtk.37}
\bookmark[dest={mtk.37},level=1]{6. อาปตฺติยา อปฺปฏิกมฺเม อุกฺเขปนียกมฺม ...}
\csromanheader{1}{6. อาปตฺติยา อปฺปฏิกมฺเม อุกฺเขปนีย\-กมฺม}
\proseitem{56}{เตน สมเยน พุทฺโธ ภควา โกสมฺพิยํ วิหรติ โฆสิตาราเม. เตน โข ปน สมเยน อายสฺมา ฉนฺโน อาปตฺติํ อาปชฺชิตฺวา น อิจฺฉติ อาปตฺติํ ปฏิกาตุํ. เย เต ภิกฺขู อปฺปิจฺฉา ฯเปฯ เต อุชฺฌายนฺติ ขิยฺยนฺติ วิปาเจนฺติ “กถํ หิ นาม อายสฺมา ฉนฺโน อาปตฺติํ อาปชฺชิตฺวา น อิจฺฉิสฺสติ อาปตฺติํ ปฏิกาตุนฺ”ติ. อถ โข เต ภิกฺขู ภควโต เอตมตฺถํ อาโรเจสุํ.}

\prose{\csromanfolio{59}อถ โข ภควา เอตสฺมิํ นิทาเน เอตสฺมิํ ปกรเณ ภิกฺขุสํฆํ สนฺนิปาตาเปตฺวา ภิกฺขู ปฏิปุจฺฉิ “สจฺจํ กิร ภิกฺขเว ฉนฺโน ภิกฺขุ อาปตฺติํ อาปชฺชิตฺวา น อิจฺฉติ อาปตฺติํ ปฏิกาตุนฺ”ติ. สจฺจํ ภควาติ. วิครหิ พุทฺโธ ภควา อนนุจฺฉวิกํ ฯเปฯ. กถํ หิ นาม โส ภิกฺขเว โมฆปุริโส อาปตฺติํ อาปชฺชิตฺวา น อิจฺฉิสฺสติ อาปตฺติํ ปฏิกาตุํ, เนตํ ภิกฺขเว อปฺปสนฺนานํ วา ปสาทาย ฯเปฯ วิครหิตฺวา ฯเปฯ ธมฺมิํ กถํ กตฺวา ภิกฺขู อามนฺเตสิ\csromandash{} เตน หิ ภิกฺขเว สํโฆ ฉนฺนสฺส ภิกฺขุโน อาปตฺติยา อปฺปฏิกมฺเม อุกฺเขปนีย\-กมฺมํ กโรตุ อสมฺโภคํ สํเฆน, เอวญฺจ ปน ภิกฺขเว กาตพฺพํ, ปฐมํ ฉนฺโน ภิกฺขุ โจเทตพฺโพ, โจเทตฺวา สาเรตพฺโพ, สาเรตฺวา อาปตฺติํ อาโรเปตพฺโพ, อาปตฺติํ อาโรเปตฺวา พฺยตฺเตน ภิกฺขุนา ปฏิพเลน สํโฆ ญาเปตพฺโพ\csromandash{}}

\prose{“สุณาตุ เม ภนฺเต สํโฆ, อยํ ฉนฺโน ภิกฺขุ อาปตฺติํ อาปชฺชิตฺวา น อิจฺฉติ อาปตฺติํ ปฏิกาตุํ, ยทิ สํฆสฺส ปตฺตกลฺลํ, สํโฆ ฉนฺนสฺส ภิกฺขุโน อาปตฺติยา อปฺปฏิกมฺเม อุกฺเขปนีย\-กมฺมํ กเรยฺย อสมฺโภคํ สํเฆน, เอสา ญตฺติ.}

\prose{สุณาตุ เม ภนฺเต สํโฆ, อยํ ฉนฺโน ภิกฺขุ อาปตฺติํ อาปชฺชิตฺวา น อิจฺฉติ อาปตฺติํ ปฏิกาตุํ, สํโฆ ฉนฺนสฺส ภิกฺขุโน อาปตฺติยา อปฺปฏิกมฺเม อุกฺเขปนีย\-กมฺมํ กโรติ อสมฺโภคํ สํเฆน, ยสฺสายสฺมโต ขมติ ฉนฺนสฺส ภิกฺขุโน อาปตฺติยา อปฺปฏิกมฺเม อุกฺเขปนียสฺส กมฺมสฺส กรณํ อสมฺโภคํ สํเฆน, โส ตุณฺหสฺส, ยสฺส นกฺขมติ, โส ภาเสยฺย.}

\prose{ทุติยมฺปิ เอตมตฺถํ วทามิ ฯเปฯ. ตติยมฺปิ เอตมตฺถํ วทามิ. สุณาตุ เม ภนฺเต สํโฆ, อยํ ฉนฺโน ภิกฺขุ อาปตฺติํ อาปชฺชิตฺวา น อิจฺฉติ อาปตฺติํ ปฏิกาตุํ, สํโฆ ฉนฺนสฺส ภิกฺขุโน อาปตฺติยา อปฺปฏิกมฺเม อุกฺเขปนีย\-กมฺมํ กโรติ อสมฺโภคํ สํเฆน, ยสฺสายสฺมโต ขมติ ฉนฺนสฺส ภิกฺขุโน อาปตฺติยา อปฺปฏิกมฺเม อุกฺเขปนียสฺส กมฺมสฺส กรณํ อสมฺโภคํ สํเฆน, โส ตุณฺหสฺส, ยสฺส นกฺขมติ, โส ภาเสยฺย.}

\prose{กตํ สํเฆน ฉนฺนสฺส ภิกฺขุโน อาปตฺติยา อปฺปฏิกมฺเม อุกฺเขปนีย\-กมฺมํ อสมฺโภคํ สํเฆน, ขมติ สํฆสฺส, ตสฺมา ตุณฺหี, เอวเมตํ ธารยามี”ติ.}

\prose{\csromanfolio{60}อาวาสปรมฺปรญฺจ ภิกฺขเว สํสถ “ฉนฺโน ภิกฺขุ สํเฆน อาปตฺติยา\csromansharedfootnote{d0af2f6c05f86da0}{ฉนฺโน ภิกฺขุ อาปตฺติยา (สี, ก)} อปฺปฏิกมฺเม อุกฺเขปนีย\-กมฺมกโต อสมฺโภคํ สํเฆนา”ติ.}

\csromanmatikaanchor{mtk.38}
\csromantocmarkref{subsection}{อธมฺมกมฺมทฺวาทสก}{mtk.38}
\bookmark[dest={mtk.38},level=2]{อธมฺมกมฺมทฺวาทสก}
\csromanheader{2}{อธมฺมกมฺม\-ทฺวาทสก}
\proseitem{57}{ตีหิ ภิกฺขเว องฺเคหิ สมนฺนาคตํ อาปตฺติยา อปฺปฏิกมฺเม อุกฺเขปนีย\-กมฺมํ อธมฺมกมฺมญฺจ โหติ อวินยกมฺมญฺจ ทุวูปสนฺตญฺจ. อสมฺมุขา กตํ โหติ, อปฺปฏิปุจฺฉา กตํ โหติ, อปฺปฏิญฺญาย กตํ โหติ. อิเมหิ โข ภิกฺขเว ตีหงฺเคหิ สมนฺนาคตํ อาปตฺติยา อปฺปฏิกมฺเม อุกฺเขปนีย\-กมฺมํ อธมฺมกมฺมญฺจ โหติ อวินยกมฺมญฺจ ทุวูปสนฺตญฺจ.}

\prose{อปเรหิปิ ภิกฺขเว ตีหงฺเคหิ สมนฺนาคตํ อาปตฺติยา อปฺปฏิกมฺเม อุกฺเขปนีย\-กมฺมํ อธมฺมกมฺมญฺจ โหติ อวินยกมฺมญฺจ ทุวูปสนฺตญฺจ. อนาปตฺติยา กตํ โหติ, อเทสนาคามินิยา อาปตฺติยา กตํ โหติ, เทสิตาย อาปตฺติยา กตํ โหติ ฯเปฯ. อโจเทตฺวา กตํ โหติ, อสาเรตฺวา กตํ โหติ, อาปตฺติํ อนาโรเปตฺวา กตํ โหติ ฯเปฯ. อสมฺมุขา กตํ โหติ, อธมฺเมน กตํ โหติ, วคฺเคน กตํ โหติ ฯเปฯ. อปฺปฏิปุจฺฉา กตํ โหติ, อธมฺเมน กตํ โหติ, วคฺเคน กตํ โหติ ฯเปฯ. อปฺปฏิญฺญาย กตํ โหติ, อธมฺเมน กตํ โหติ, วคฺเคน กตํ โหติ ฯเปฯ. อนาปตฺติยา กตํ โหติ, อธมฺเมน กตํ โหติ, วคฺเคน กตํ โหติ ฯเปฯ. อเทสนาคามินิยา อาปตฺติยา กตํ โหติ, อธมฺเมน กตํ โหติ, วคฺเคน กตํ โหติ ฯเปฯ. เทสิตาย อาปตฺติยา กตํ โหติ, อธมฺเมน กตํ โหติ, วคฺเคน กตํ โหติ ฯเปฯ. อโจเทตฺวา กตํ โหติ, อธมฺเมน กตํ โหติ, วคฺเคน กตํ โหติ ฯเปฯ. อสาเรตฺวา กตํ โหติ, อธมฺเมน กตํ โหติ, วคฺเคน กตํ โหติ ฯเปฯ. อาปตฺติํ อนาโรเปตฺวา กตํ โหติ, อธมฺเมน กตํ โหติ, วคฺเคน กตํ โหติ. อิเมหิ โข ภิกฺขเว หีหงฺเคหิ สมนฺนาคตํ อาปตฺติยา อปฺปฏิกมฺเม อุกฺเขปนีย\-กมฺมํ อธมฺมกมฺมญฺจ โหติ อวินยกมฺมญฺจ ทุวูปสนฺตญฺจ.}

\csromanheader{2}{อาปตฺติยา อปฺปฏิกมฺเม อุกฺเขปนีย\-กมฺเม}
\nitthitam{อธมฺมกมฺม\-ทฺวาทสกํ นิฏฺฐิตํ.}
\csromanmatikaanchor{mtk.39}
\csromantocmarkref{subsection}{ธมฺมกมฺมทฺวาทสก}{mtk.39}
\bookmark[dest={mtk.39},level=2]{ธมฺมกมฺมทฺวาทสก}
\csromanheader{2}{ธมฺมกมฺม\-ทฺวาทสก}
\proseitem{58}{\csromanfolio{61}ตีหิ ภิกฺขเว องฺเคหิ สมนฺนาคตํ อาปตฺติยา อปฺปฏิกมฺเม อุกฺเขปนีย\-กมฺมํ ธมฺมกมฺมญฺจ โหติ วินยกมฺมญฺจ สุวูปสนฺตญฺจ. สมฺมุขา กตํ โหติ, ปฏิปุจฺฉา กตํ โหติ, ปฏิญฺญาย กตํ โหติ. อิเมหิ โข ภิกฺขเว ตีหงฺเคหิ สมนฺนาคตํ อาปตฺติยา อปฺปฏิกมฺเม อุกฺเขปนีย\-กมฺมํ ธมฺมกมฺมญฺจ โหติ วินยกมฺมญฺจ สุวูปสนฺตญฺจ.}

\prose{อปเรหิปิ ภิกฺขเว ตีหงฺเคหิ สมนฺนาคตํ อาปตฺติยา อปฺปฏิกมฺเม อุกฺเขปนีย\-กมฺมํ ธมฺมกมฺมญฺจ โหติ วินยกมฺมญฺจ สุวูปสนฺตญฺจ. อาปตฺติยา กตํ โหติ, เทสนาคามินิยา อาปตฺติยา กตํ โหติ, อเทสิตาย อาปตฺติยา กตํ โหติ ฯเปฯ. โจเทตฺวา กตํ โหติ, สาเรตฺวา กตํ โหติ, อาปตฺติํ อาโรเปตฺวา กตํ โหติ ฯเปฯ. สมฺมุขา กตํ โหติ, ธมฺเมน กตํ โหติ. สมคฺเคน กตํ โหติ ฯเปฯ. ปฏิปุจฺฉา กตํ โหติ, ธมฺเมน กตํ โหติ, สมคฺเคน กตํ โหติ ฯเปฯ. ปฏิญฺญาย กตํ โหติ, ธมฺเมน กตํ โหติ, สมคฺเคน กตํ โหติ ฯเปฯ. อาปตฺติยา กตํ โหติ, ธมฺเมน กตํ โหติ, สมคฺเคน กตํ โหติ ฯเปฯ. เทสนาคามินิยา อาปตฺติยา กตํ โหติ, ธมฺเมน กตํ โหติ, สมคฺเคน กตํ โหติ ฯเปฯ. อเทสิตาย อาปตฺติยา กตํ โหติ, ธมฺเมน กตํ โหติ, สมคฺเคน กตํ โหติ ฯเปฯ. โจเทตฺวา กตํ โหติ, ธมฺเมน กตํ โหติ, สมคฺเคน กตํ โหติ ฯเปฯ. สาเรตฺวา กตํ โหติ, ธมฺเมน กตํ โหติ, สมคฺเคน กตํ โหติ ฯเปฯ. อาปตฺติํ อาโรเปตฺวา กตํ โหติ, ธมฺเมน กตํ โหติ, สมคฺเคน กตํ โหติ. อิเมหิ โข ภิกฺขเว ตีหงฺเคหิ สมนฺนาคตํ อาปตฺติยา อปฺปฏิกมฺเม อุกฺเขปนีย\-กมฺมํ ธมฺมกมฺมญฺจ โหติ วินยกมฺมญฺจ สุวูปสนฺตญฺจ.}

\csromanheader{2}{อาปตฺติยา อปฺปฏิกมฺเม อุกฺเขปนีย\-กมฺเม}
\nitthitam{ธมฺมกมฺม\-ทฺวาทสกํ นิฏฺฐิตํ.}
\csromanmatikaanchor{mtk.40}
\csromantocmarkref{subsection}{อากงฺขมานฉกฺก}{mtk.40}
\bookmark[dest={mtk.40},level=2]{อากงฺขมานฉกฺก}
\proseitem{59}{\csromanfolio{62}\csromansymbolfootnote{*}{วิ 5. 220 ปิฏฺเฐปิ.}ตีหิ ภิกฺขเว องฺเคหิ สมนฺนาคตสฺส ภิกฺขุโน อากงฺขมาโน สํโฆ อาปตฺติยา อปฺปฏิกมฺเม อุกฺเขปนีย\-กมฺมํ กเรยฺย. ภณฺฑนการโก โหติ กลหการโก วิวาทการโก ภสฺสการโก สํเฆ อธิกรณ\-การโก, พาโล โหติ อพฺยตฺโต อาปตฺติพหุโล อนปทาโน, คิหิสํสฏฺโฐ วิหรติ อนนุโลมิเกหิ คิหิสํสคฺเคหิ. อิเมหิ โข ภิกฺขเว ตีหงฺเคหิ สมนฺนาคตสฺส ภิกฺขุโน อากงฺขมาโน สํโฆ อาปตฺติยา อปฺปฏิกมฺเม อุกฺเขปนีย\-กมฺมํ กเรยฺย.}

\prose{อปเรหิปิ ภิกฺขเว ตีหงฺเคหิ สมนฺนาคตสฺส ภิกฺขุโน อากงฺขมาโน สํโฆ อาปตฺติยา อปฺปฏิกมฺเม อุกฺเขปนีย\-กมฺมํ กเรยฺย. อธิสีเล สีลวิปนฺโน โหติ, อชฺฌาจาเร อาจารวิปนฺโน โหติ, อติทิฏฺฐิยา ทิฏฺฐิวิปนฺโน โหติ. อิเมหิ โข ภิกฺขเว ตีหงฺเคหิ สมนฺนาคตสฺส ภิกฺขุโน อากงฺขมาโน สํโฆ อาปตฺติยา อปฺปฏิกมฺเม อุกฺเขปนีย\-กมฺมํ กเรยฺย.}

\prose{อปเรหิปิ ภิกฺขเว ตีหงฺเคหิ สมนฺนาคตสฺส ภิกฺขุโน อากงฺขมาโน สํโฆ อาปตฺติยา อปฺปฏิกมฺเม อุกฺเขปนีย\-กมฺมํ กเรยฺย. พุทฺธสฺส อวณฺณํ ภาสติ, ธมฺมสฺส อวณฺณํ ภาสติ, สํฆสฺส อวณฺณํ ภาสติ. อิเมหิ โข ภิกฺขเว ตีหงฺเคหิ สมนฺนาคตสฺส ภิกฺขุโน อากงฺขมาโน สํโฆ อาปตฺติยา อปฺปฏิกมฺเม อุกฺเขปนีย\-กมฺมํ กเรยฺย.}

\prose{ติณฺณํ ภิกฺขเว ภิกฺขูนํ อากงฺขมาโน สํโฆ อาปตฺติยา อปฺปฏิกมฺเม อุกฺเขปนีย\-กมฺมํ กเรยฺย. เอโก ภณฺฑนการโก โหติ กลหการโก วิวาทการโก ภสฺสการโก สํเฆ อธิกรณ\-การโก, เอโก พาโล โหติ อพฺยตฺโต อาปตฺติพหุโล อนปทาโน, เอโก คิหิสํสฏฺโฐ วิหรติ อนนุโลมิเกหิ คิหิ สํสคฺเคหิ. อิเมสํ โข ภิกฺขเว ติณฺณํ ภิกฺขูนํ อากงฺขมาโน สํโฆ อาปตฺติยา อปฺปฏิกมฺเม อุกฺเขปนีย\-กมฺมํ กเรยฺย.}

\prose{อปเรสมฺปิ ภิกฺขเว ติณฺณํ ภิกฺขูนํ อากงฺขมาโน สํโฆ อาปตฺติยา อปฺปฏิกมฺเม อุกฺเขปนีย\-กมฺมํ กเรยฺย. เอโก อธิสีเล \csromanfolio{63}สีลวิปนฺโน โหติ. เอโก อชฺฌาจาเร อาจารวิปนฺโน โหติ, เอโก อติทิฏฺฐิยา ทิฏฺฐิวิปนฺโน โหติ. อิเมสํ โข ภิกฺขเว ติณฺณํ ภิกฺขูนํ อากงฺขมาโน สํโฆ อาปตฺติยา อปฺปฏิกมฺเม อุกฺเขปนีย\-กมฺมํ กเรยฺย.}

\prose{อปเรสมฺปิ ภิกฺขเว ติณฺณํ ภิกฺขูนํ อากงฺขมาโน สํโฆ อาปตฺติยา อปฺปฏิกมฺเม อุกฺเขปนีย\-กมฺมํ กเรยฺย. เอโก พุทฺธสฺส อวณฺณํ ภาสติ, เอโก ธมฺมสฺส อวณฺณํ ภาสติ, เอโก สํฆสฺส อวณฺณํ ภาสติ. อิเมสํ โข ภิกฺขเว ติณฺณํ ภิกฺขูนํ อากงฺขมาโน สํโฆ อาปตฺติยา อปฺปฏิกมฺเม อุกฺเขปนีย\-กมฺมํ กเรยฺย.}

\csromanheader{2}{อาปตฺติยา อปฺปฏิกมฺเม อุกฺเขปนีย\-กมฺเม}
\nitthitam{อากงฺขมาน\-ฉกฺกํ นิฏฺฐิตํ.}
\csromanmatikaanchor{mtk.41}
\csromantocmarkref{subsection}{เตจตฺตาลีสวตฺต}{mtk.41}
\bookmark[dest={mtk.41},level=2]{เตจตฺตาลีสวตฺต}
\csromanheader{2}{เตจตฺตาลีส\-วตฺต}
\proseitem{60}{อาปตฺติยา อปฺปฏิกมฺเม อุกฺเขปนีย\-กมฺมกเตน ภิกฺขเว ภิกฺขุนา สมฺมา วตฺติตพฺพํ. ตตฺรยํ สมฺมา วตฺตนา\csromandash{}น อุปสมฺ\-ปาเท\-ตพฺพํ, น นิสฺสโย ทาตพฺโพ, น สามเณโร อุปฏฺฐาเปตพฺโพ, น ภิกฺขุโนวาทก\-สมฺมุติ สาทิตพฺพา, สมฺมเตนปิ ภิกฺขุนิโย น โอวทิตพฺพา, ยาย อาปตฺติยา สํเฆน อาปตฺติยา อปฺปฏิกมฺเม อุกฺเขปนีย\-กมฺมํ กตํ โหติ, สา อาปตฺติ น อาปชฺชิตพฺพา, อญฺญา วา ตาทิสิกา, ตโต วา ปาปิฏฺฐตรา, กมฺมํ น ครหิตพฺพํ, กมฺมิกา น ครหิตพฺพา, น ปกตตฺตสฺส ภิกฺขุโน อภิวาทนํ ปจฺจุฏฺฐานํ อญฺชลิกมฺมํ สามิจิกมฺมํ อาสนาภิหาโร เสยฺยาภิหาโร ปาโททกํ ปาทปีฐํ ปาทกถลิกํ ปตฺตจีวร\-ปฺปฏิคฺคหณํ นหาเน ปิฏฺฐิ\-ปริกมฺมํ สาทิตพฺพํ, น ปกตตฺโต ภิกฺขุ สีลวิปตฺติยา อนุทฺธํเส\-ตพฺโพ, น อาจารวิปตฺติยา อนุทฺธํเส\-ตพฺโพ, น ทิฏฺฐิ\-วิปตฺติยา อนุทฺธํเส\-ตพฺโพ, น อาชีววิปตฺติยา อนุทฺธํเส\-ตพฺโพ, น ภิกฺขุ ภิกฺขูหิ เภเทตพฺโพ, น คิหิทฺธโช ธาเรตพฺโพ, น ติตฺถิยทฺธโช ธาเรตพฺโพ, น ติตฺถิยา เสวิตพฺพา, ภิกฺขู เสวิตพฺพา, ภิกฺขุ\-สิกฺขาย สิกฺขิตพฺพํ, น ปกตตฺเตน ภิกฺขุนา สทฺธิํ เอกจฺฉนฺเน อาวาเส วตฺถพฺพํ, น เอกจฺฉนฺเน อนาวาเส วตฺถพฺพํ, น เอกจฺฉนฺเน อาวาเส \csromanfolio{64}วา อนาวาเส วา วตฺถพฺพํ, ปกตตฺตํ ภิกฺขุํ ทิสฺวา อาสนา วุฏฺฐาตพฺพํ, น ปกตตฺโต ภิกฺขุ อาสาเทตพฺโพ อนฺโต วา พหิ วา, น ปกตตฺตสฺส ภิกฺขุโน อุโปสโถ ฐเปตพฺโพ, น ปวารณา ฐเปตพฺพา, น สวจนียํ กาตพฺพํ, น อนุวาโท ปฏฺฐเปตพฺโพ, น โอกาโส กาเรตพฺโพ, น โจเทตพฺโพ, น สาเรตพฺโพ, น ภิกฺขูหิ สมฺป\-โยเช\-ตพฺพนฺติ.}

\csromanheader{2}{อาปตฺติยา อปฺปฏิกมฺเม อุกฺเขปนีย\-กมฺเม}
\nitthitam{เตจตฺตาลีส\-วตฺตํ นิฏฺฐิตํ.}
\proseitem{61}{อถ โข สํโฆ ฉนฺนสฺส ภิกฺขุโน อาปตฺติยา อปฺปฏิกมฺเม อุกฺเขปนีย\-กมฺมํ อกาสิ อสมฺโภคํ สํเฆน. โส สํเฆน อาปตฺติยา อปฺปฏิกมฺเม อุกฺเขปนีย\-กมฺมกโต ตมฺหา อาวาสา อญฺญํ อาวาสํ อคมาสิ. ตตฺถ ภิกฺขู เนว อภิวาเทสุํ, น ปจฺจุฏฺเฐสุํ, น อญฺชลิกมฺมํ, น สามีจิกมฺมํ อกํสุ, น สกฺกริสุ, น ครุํ กริํสุ, น มาเนสุํ, น ปูเชสุํ. โส ภิกฺขูหิ อสกฺกริยมาโน อครุกริยมาโน อมานิยมาโน อปูชิยมาโน อสกฺการปกโต ตมฺหาปิ อาวาสา อญฺญํ อาวาสํ อคมาสิ. ตตฺถปิ ภิกฺขู เนว อภิวาเทสุํ, น ปจฺจุฏฺเฐสุํ, น อญฺชลิกมฺมํ, น สามีจิกมฺมํ อกํสุ, น สกฺกริํสุ, น ครุํ กริํสุ, น มาเนสุํ, น ปูเชสุํ. โส ภิกฺขูหิ อสกฺกริยมาโน อครุกริยมาโน อมานิยมาโน อปูชิยมาโน อสกฺการปกโต ตมฺหาปิ อาวาสา อญฺญํ อาวาสํ อคมาสิ. ตตฺถปิ ภิกฺขู เนว อภิวาเทสุํ, น ปจฺจุฏฺเฐสุํ, น อญฺชลิกมฺมํ, น สามีจิกมฺมํ อกํสุ, น สกฺกริํสุ, น ครุํ กริํสุ, น มาเนสุํ, น ปูเชสุํ. โส ภิกฺขูหิ อสกฺกริยมาโน อครุกริยมาโน อมานิยมาโน อปูชิยมาโน อสกฺการปกโต ปุนเทว โกสมฺพิํ ปจฺจาคญฺฉิ. โส สมฺมา วตฺตติ, โลมํ ปาเตติ, เนตฺถารํ วตฺตติ, ภิกฺขู อุปสงฺกมิตฺวา เอวํ วเทติ “อหํ อาวุโส สํเฆน อาปตฺติยา อปฺปฏิกมฺเม อุกฺเขปนีย\-กมฺมกโต สมฺมา วตฺตามิ, โลมํ ปาเตมิ, เนตฺถารํ วตฺตามิ, กถํ นุ ขา มยา ปฏิปชฺชิตพฺพนฺ”ติ. ภิกฺขู ภควโต เอตมตฺถํ อาโรเจสุํ ฯเปฯ. เตน หิ ภิกฺขเว สํโฆ ฉนฺนสฺส ภิกฺขุโน อาปตฺติยา อปฺปฏิกมฺเม อุกฺเขปนีย\-กมฺมํ ปฏิปฺปสฺสมฺเภตุ.}

\csromanmatikaanchor{mtk.42}
\csromantocmarkref{subsection}{นปฺปฏิปฺปสฺสมฺเภตพฺพเตจตฺตาลีสก}{mtk.42}
\bookmark[dest={mtk.42},level=2]{นปฺปฏิปฺปสฺสมฺเภตพฺพเตจตฺตาลีสก}
\csromanheader{2}{นปฺปฏิปฺปสฺสมฺเภตพฺพ\-เตจตฺตาลีสก}
\proseitem{62}{\csromanfolio{65}ปญฺจหิ ภิกฺขเว องฺเคหิ สมนฺนาคตสฺส ภิกฺขุโน อาปตฺติยา อปฺปฏิกมฺเม อุกฺเขปนีย\-กมฺมํ นปฺปฏิปฺปสฺสมฺเภตพฺพํ. อุปสมฺปาเทติ, นิสฺสยํ เทติ, สามเณรํ อุปฏฺฐาเปติ, ภิกฺขุโนวาทก\-สมฺมุติํ สาทิยติ, สมฺมโตปิ ภิกฺขุนิโย โอวทติ. อิเมหิ โข ภิกฺขเว ปญฺจหงฺเคหิ สมนฺนาคตสฺส ภิกฺขุโน อาปตฺติยา อปฺปฏิกมฺเม อุกฺเขปนีย\-กมฺมํ นปฺปฏิปฺปสฺสมฺเภตพฺพํ.}

\prose{\csromansymbolfootnote{*}{วิ 5. 317 ปิฏฺเฐปิ.}อปเรหิปิ ภิกฺขเว ปญฺจหงฺเคหิ สมนฺนาคตสฺส ภิกฺขุโน อาปตฺติยา อปฺปฏิกมฺเม อุกฺเขปนีย\-กมฺมํ นปฺปฏิปฺปสฺสมฺเภตพฺพํ. ยาย อาปตฺติยา สํเฆน อาปตฺติยา อปฺปฏิกมฺเม อุกฺเขปนีย\-กมฺมํ กตํ โหติ, ตํ อาปตฺติํ อาปชฺชติ, อญฺญํ วา ตาทิสิกํ, ตโต วา ปาปิฏฺฐตรํ, กมฺมํ ครหติ, กมฺมิเก ครหติ ฯเปฯ. ปกตตฺตสฺส ภิกฺขุโน อภิวาทนํ ปจฺจุฏฺฐานํ อญฺชลิกมฺมํ สามีจิกมฺมํ อาสนาภิหารํ สาทิยติ ฯเปฯ. ปกตตฺตสฺส ภิกฺขุโน เสยฺยาภิหารํ ปาโททกํ ปาทปีฐํ ปาทกถลิกํ ปตฺตจีวร\-ปฺปฏิคฺคหณํ นหาเน ปิฏฺฐิ\-ปริกมฺมํ สาทิยติ ฯเปฯ. ปกตตฺตํ ภิกฺขุํ สีลวิปตฺติยา อนุทฺธํเสติ, อาจารวิปตฺติยา อนุทฺธํเสติ, ทิฏฺฐิ\-วิปตฺติยา อนุทฺธํเสติ, อาชีววิปตฺติยา อนุทฺธํเสติ, ภิกฺขุํ ภิกฺขูหิ เภเทติ ฯเปฯ. คิหิทฺธชํ ธาเรติ, ติตฺถิยทฺธชํ ธาเรติ, ติตฺถิเย เสวติ, ภิกฺขู น เสวติ, ภิกฺขุ\-สิกฺขาย น สิกฺขติ ฯเปฯ. ปกตตฺเตน ภิกฺขุนา สทฺธิํ เอกจฺฉนฺเน อาวาเส วสติ, เอกจฺฉนฺเน อนาวาเส วสติ, เอกจฺฉนฺเน อาวาเส วา อนาวาเส วา วสติ, ปกตตฺตํ ภิกฺขุํ ทิสฺวา อาสนา น วุฏฺฐาติ, ปกตตฺตํ ภิกฺขุํ อาสาเทติ อนฺโต วา พหิ วา. อิเมหิ โข ภิกฺขเว ปญฺจหงฺเคหิ สมนฺนาคตสฺส ภิกฺขุโน อาปตฺติยา อปฺปฏิกมฺเม อุกฺเขปนีย\-กมฺมํ นปฺปฏิปฺปสฺสมฺเภตพฺพํ.}

\prose{อฏฺฐหิ ภิกฺขเว องฺเคหิ สมนฺนาคตสฺส ภิกฺขุโน อาปตฺติยา อปฺปฏิกมฺเม อุกฺเขปนีย\-กมฺมํ นปฺปฏิปฺปสฺสมฺเภตพฺพํ. ปกตตฺตสฺส ภิกฺขุโน อุโปสถํ ฐเปติ, ปวารณํ ฐเปติ, สวจนียํ กโรติ, อนุวาทํ ปฏฺฐเปติ, โอกาสํ กาเรติ, โจเทติ, สาเรติ, ภิกฺขูหิ \csromanfolio{66}สมฺปโยเชติ. อิเมหิ โข ภิกฺขเว อฏฺฐหงฺเคหิ สมนฺนาคตสฺส ภิกฺขุโน อาปตฺติยา อปฺปฏิกมฺเม อุกฺเขปนีย\-กมฺมํ นปฺปฏิปฺปสฺสมฺเภตพฺพํ.}

\csromanheader{2}{อาปตฺติยา อปฺปฏิกมฺเม อุกฺเขปนีย\-กมฺเม}
\nitthitam{นปฺปฏิปฺปสฺสมฺเภตพฺพ\-เตจตฺตาลีสกํ นิฏฺฐิตํ.}
\csromanmatikaanchor{mtk.43}
\csromantocmarkref{subsection}{ปฏิปฺปสฺสมฺเภตพฺพเตจตฺตาลีสก}{mtk.43}
\bookmark[dest={mtk.43},level=2]{ปฏิปฺปสฺสมฺเภตพฺพเตจตฺตาลีสก}
\csromanheader{2}{ปฏิปฺปสฺสมฺเภตพฺพ\-เตจตฺตาลีสก}
\proseitem{63}{ปญฺจหิ ภิกฺขเว องฺเคหิ สมนฺนาคตสฺส ภิกฺขุโน อาปตฺติยา อปฺปฏิกมฺเม อุกฺเขปนีย\-กมฺมํ ปฏิปฺปสฺสมฺเภตพฺพํ. น อุปสมฺปาเทติ, น นิสฺสยํ เทติ, น สามเณรํ อุปฏฺฐาเปติ, น ภิกฺขุโนวาทก\-สมฺมุติํ สาทิยติ, สมฺมโตปิ ภิกฺขุนิโย น โอวทติ. อิเมหิ โข ภิกฺขเว ปญฺจหงฺเคหิ สมนฺนาคตสฺส ภิกฺขุโน อาปตฺติยา อปฺปฏิกมฺเม อุกฺเขปนีย\-กมฺมํ ปฏิปฺปสฺสมฺเภตพฺพํ.}

\prose{อปเรหิปิ ภิกฺขเว ปญฺจหงฺเคหิ สมนฺนาคตสฺส ภิกฺขุโน อาปตฺติยา อปฺปฏิกมฺเม อุกฺเขปนีย\-กมฺมํ ปฏิปฺปสฺสมฺเภตพฺพํ. ยาย อาปตฺติยา สํเฆน อาปตฺติยา อปฺปฏิกมฺเม อุกฺเขปนีย\-กมฺมํ กตํ โหติ, ตํ อาปตฺติํ น อาปชฺชติ, อญฺญํ วา ตาทิสิกํ, ตโต วา ปาปิฏฺฐตรํ, กมฺมํ น ครหติ, กมฺมิเก น ครหติ ฯเปฯ. น ปกตตฺตสฺส ภิกฺขุโน อภิวาทนํ ปจฺจุฏฺฐานํ อญฺชลิกมฺมํ สามีจิกมฺมํ อาสนาภิหารํ สาทิยติ ฯเปฯ. น ปกตตฺตสฺส ภิกฺขุโน เสยฺยาภิหารํ ปาโททกํ ปาทปีฐํ ปาทกถลิกํ ปตฺตจีวร\-ปฺปฏิคฺคหณํ นหาเน ปิฏฺฐิ\-ปริกมฺมํ สาทิยติ ฯเปฯ. น ปกตตฺตํ ภิกฺขุํ สีลวิปตฺติยา อนุทฺธํเสติ, น อาจารวิปตฺติยา อนุทฺธํเสติ, น ทิฏฺฐิ\-วิปตฺติยา อนุทฺธํเสติ, น อาชีววิปตฺติยา อนุทฺธํเสติ, น ภิกฺขุํ ภิกฺขูหิ เภเทติ ฯเปฯ. น คิหิทฺธชํ ธาเรติ, น ติตฺถิยทฺธชํ ธาเรติ, น ติตฺถิเย เสวติ, ภิกฺขู เสวติ, ภิกฺขุ\-สิกฺขาย สิกฺขติ ฯเปฯ. น ปกตตฺเตน ภิกฺขุนา สทฺธิํ เอกจฺฉนฺเน อาวาเส วสติ, น เอกจฺฉนฺเน อนาวาเส วสติ, น เอกจฺฉนฺเน อาวาเส วา อนาวาเส วา วสติ, ปกตตฺตํ ภิกฺขุํ ทิสฺวา อาสนา วุฏฺฐาติ, น ปกตตฺตํ ภิกฺขุํ อาสาเทติ อนฺโต วา พหิ วา. อิเมหิ โข ภิกฺขเว ปญฺจหงฺเคหิ สมนฺนาคตสฺส ภิกฺขุโน อาปตฺติยา อปฺปฏิกมฺเม อุกฺเขปนีย\-กมฺมํ ปฏิปฺปสฺสมฺเภตพฺพํ.}

\prose{\csromanfolio{67}อฏฺฐหิ ภิกฺขเว องฺเคหิ สมนฺนาคตสฺส ภิกฺขุโน อาปตฺติยา อปฺปฏิกมฺเม อุกฺเขปนีย\-กมฺมํ ปฏิปฺปสฺสมฺเภตพฺพํ. น ปกตตฺตสฺส ภิกฺขุโน อุโปสถํ ฐเปติ, น ปวารณํ ฐเปติ, น สวจนียํ กโรติ, น อนุวาทํ ปฏฺฐเปติ, น โอกาสํ กาเรติ, น โจเทติ, น สาเรติ, น ภิกฺขูหิ สมฺปโยเชติ. อิเมหิ โข ภิกฺขเว อฏฺฐหงฺเคหิ สมนฺนาคตสฺส ภิกฺขุโน อาปตฺติยา อปฺปฏิกมฺเม อุกฺเขปนีย\-กมฺมํ ปฏิปฺปสฺสมฺเภตพฺพํ.}

\csromanheader{2}{อาปตฺติยา อปฺปฏิกมฺเม อุกฺเขปนีย\-กมฺเม}
\nitthitam{ปฏิปฺปสฺสมฺเภตพฺพ\-เตจตฺตาลีสกํ นิฏฺฐิตํ.}
\proseitem{64}{เอวญฺจ ปน ภิกฺขเว ปฏิปฺปสฺสมฺเภตพฺพํ, เตน ภิกฺขเว ฉนฺเนน ภิกฺขุนา สํฆํ อุปสงฺกมิตฺวา เอกํสํ อุตฺตราสงฺคํ กริตฺวา วุฑฺฒานํ ภิกฺขูนํ ปาเท วนฺทิตฺวา อุกฺกุฏิกํ นิสีทิตฺวา อญฺชลิํ ปคฺคเหตฺวา เอวมสฺส วจนีโย “อหํ ภนฺเต สํเฆน อาปตฺติยา อปฺปฏิกมฺเม อุกฺเขปนีย\-กมฺมกโต สมฺมา วตฺตามิ, โลมํ ปาเตมิ, เนตฺถารํ วตฺตามิ, อาปตฺติยา อปฺปฏิกมฺเม อุกฺเขปนียสฺส กมฺมสฺส ปฏิปฺปสฺสทฺธิํ ยาจามี”ติ. ทุติยมฺปิ ยาจิตพฺพา. ตติยมฺปิ ยาจิตพฺพา. พฺยตฺเตน ภิกฺขุนา ปฏิพเลน สํโฆ ญาเปตพฺโพ\csromandash{}}

\prose{“สุณาตุ เม ภนฺเต สํโฆ, อยํ ฉนฺโน ภิกฺขุ สํเฆน อาปตฺติยา อปฺปฏิกมฺเม อุกฺเขปนีย\-กมฺมกโต สมฺมา วตฺตติ, โลมํ ปาเตติ, เนตฺถารํ วตฺตติ, อาปตฺติยา อปฺปฏิกมฺเม อุกฺเขปนียสฺส กมฺมสฺส ปฏิปฺปสฺสทฺธิํ ยาจติ, ยทิ สํฆสฺส ปตฺตกลฺลํ, สํโฆ ฉนฺนสฺส ภิกฺขุโน อาปตฺติยา อปฺปฏิกมฺเม อุกฺเขปนีย\-กมฺมํ ปฏิปฺปสฺสมฺเภยฺย, เอสา ญตฺติ.}

\prose{สุณาตุ เม ภนฺเต สํโฆ, อยํ ฉนฺโน ภิกฺขุ สํเฆน อาปตฺติยา อปฺปฏิกมฺเม อุกฺเขปนีย\-กมฺมกโต สมฺมา วตฺตติ, โลมํ ปาเตติ, เนตฺถารํ วตฺตติ, อาปตฺติยา อปฺปฏิกมฺเม อุกฺเขปนียสฺส กมฺมสฺส ปฏิปฺปสฺสทฺธิํ ยาจติ, สํโฆ ฉนฺนสฺส ภิกฺขุโน อาปตฺติยา อปฺปฏิกมฺเม อุกฺเขปนีย\-กมฺมํ ปฏิปฺปสฺสมฺเภติ, ยสฺสายสฺมโต ขมติ ฉนฺนสฺส ภิกฺขุโน อาปตฺติยา อปฺปฏิกมฺเม อุกฺเขปนียสฺส กมฺมสฺส ปฏิปฺปสฺสทฺธิ, โส ตุณฺหสฺส, ยสฺส นกฺขมติ, โส ภาเสยฺย.}

\csromanmatikaanchor{mtk.44}
\csromantocmarkref{section}{7. ปาปิกาย ทิฏฺฐิยา อปฺปฏินิสฺสคฺเค อุกฺเขปนียกมฺม}{mtk.44}
\bookmark[dest={mtk.44},level=1]{7. ปาปิกาย ทิฏฺฐิยา อปฺปฏินิสฺสคฺเค อุกฺเขปนียกมฺม}
\prose{\csromanfolio{68}ทุติยมฺปิ เอตมตฺถํ วทามิ ฯเปฯ. ตติยมฺปิ เอตมตฺถํ วทามิ. สุณาตุ เม ภนฺเต สํโฆ, อยํ ฉนฺโน ภิกฺขุ สํเฆน อาปตฺติยา อปฺปฏิกมฺเม อุกฺเขปนีย\-กมฺมกโต สมฺมา วตฺตติ, โลมํ ปาเตติ, เนตฺถารํ วตฺตติ, อาปตฺติยา อปฺปฏิกมฺเม อุกฺเขปนียสฺส กมฺมสฺส ปฏิปฺปสฺสทฺธิํ ยาจติ, สํโฆ ฉนฺนสฺส ภิกฺขุโน อาปตฺติยา อปฺปฏิกมฺเม อุกฺเขปนีย\-กมฺมํ ปฏิปฺปสฺสมฺเภติ, ยสฺสายสฺมโต ขมติ ฉนฺนสฺส ภิกฺขุโน อาปตฺติยา อปฺปฏิกมฺเม อุกฺเขปนียสฺส กมฺมสฺส ปฏิปฺปสฺสทฺธิ, โส ตุณฺหสฺส, ยสฺส นกฺขมติ, โส ภาเสยฺย.}

\prose{ปฏิปฺปสฺสทฺธํ สํเฆน ฉนฺนสฺส ภิกฺขุโน อาปตฺติยา อปฺปฏิกมฺเม อุกฺเขปนีย\-กมฺมํ, ขมติ สํฆสฺส, ตสฺมา ตุณฺหี, เอวเมตํ ธารยามี”ติ.}

\prose{อาปตฺติยา อปฺปฏิกมฺเม อุกฺเขปนีย\-กมฺมํ นิฏฺฐิตํ ฉฏฺฐํ.}

\prose{7. ปาปิกาย ทิฏฺฐิยา อปฺปฏินิสฺสคฺเค อุกฺเขปนีย\-กมฺม}

\proseitem{65}{\csromansymbolfootnote{*}{อิทํ วตฺถุ วิ 2. 175; ม 1. 182 ปิฏฺฐาทีสุปิ อาคตํ.}เตน สมเยน พุทฺโธ ภควา สาวตฺถิยํ วิหรติ เชตวเน อนาถ\-ปิณฺฑิกสฺส อาราเม. เตน โข ปน สมเยน อริฏฺฐสฺส นาม ภิกฺขุโน คทฺธพาธิ\-ปุพฺพสฺส\csromansharedfootnote{ce519d550967bef5}{คนฺธพาธิปุพฺพสฺส (ก)} เอวรูปํ ปาปกํ ทิฏฺฐิคตํ อุปฺปนฺนํ โหติ “ตถาหํ ภควตา ธมฺมํ เทสิตํ อาชานามิ, ยถา เยเม อนฺตรายิกา ธมฺมา วุตฺตา ภควตา, เต ปฏิเสวโต นาลํ อนฺตรายายา”ติ. อสฺโสสุํ โข สมฺพหุลา ภิกฺขู “อริฏฺฐสฺส นาม กิร ภิกฺขุโน คทฺธพาธิ\-ปุพฺพสฺส เอวรูปํ ปาปกํ ทิฏฺฐิคตํ อุปฺปนฺนํ ‘ตถาหํ ภควตา ธมฺมํ เทสิตํ อาชานามิ, ยถา เยเม อนฺตรายิกา ธมฺมา วุตฺตา ภควตา, เต ปฏิเสวโต นาลํ อนฺตรายายา’ติ”. อถ โข เต ภิกฺขู เยน อริฏฺโฐ ภิกฺขุ คทฺธพาธิ\-ปุพฺโพ เตนุ\-ปสงฺกมิํสุ, อุปสงฺกมิตฺวา อริฏฺฐํ ภิกฺขุํ คทฺธพาธิ\-ปุพฺพํ เอตทโวจุํ “สจฺจํ กิร เต อาวุโส อริฏฺฐ เอวรูปํ ปาปกํ ทิฏฺฐิคตํ อุปฺปนฺนํ ‘ตถาหํ ภควตา ธมฺมํ เทสิตํ อาชานามิ, ยถา เยเม อนฺตรายิกา ธมฺมา วุตฺตา ภควตา, เต ปฏิเสวโต นาลํ อนฺตรายายา’ติ”. เอวํพฺยา โข อหํ อาวุโส ภควตา ธมฺมํ เทสิตํ อาชานามิ, ยถา \csromanfolio{69}เยเม อนฺตรายิกา ธมฺมา วุตฺตา ภควตา, เต ปฏิเสวโต นาลํ อนฺตรายายาติ.}

\prose{มาวุโส อริฏฺฐ เอวํ อวจ, มา ภควนฺตํ อพฺภาจิกฺขิ, น หิ สาธุ ภควโต อพฺภกฺขานํ\csromansharedfootnote{792e546ef2a1bad2}{อพฺภาจิกฺขนํ (ก)}, น หิ ภควา เอวํ วเทยฺย, อเนกปริยาเยนา\-วุโส อริฏฺฐ อนฺตรายิกา ธมฺมา อนฺตรายิกา วุตฺตา ภควตา, อลญฺจ ปน เต ปฏิเสวโต อนฺตรายาย. อปฺปสฺสาทา กามา วุตฺตา ภควตา พหุทุกฺขา พหุปายาสา\csromansharedfootnote{d5d359bc39d938e3}{พหูปายายา (สี, สฺยา)}, อาทีนโว เอตฺถ ภิยฺโย. อฏฺฐิ\-กงฺกลู\-ปมา กามา วุตฺตา ภควตา พหุทุกฺขา พหุปายาสา, อาทีนโว เอตฺถ ภิยฺโย. มํสเปสูปมา กามา วุตฺตา ภควตา ฯเปฯ. ติณุกฺกูปมา กามา วุตฺตา ภควตา. องฺคารกาสูปมา กามา วุตฺตา ภควตา. สุปินกูปมา กามา วุตฺตา ภควตา. ยาจิตกูปมา กามา วุตฺตา ภควตา. รุกฺข\-ผลู\-ปมา กามา วุตฺตา ภควตา. อสิสูนูปมา กามา วุตฺตา ภควตา. สตฺติสูลูปมา กามา วุตฺตา ภควตา. สปฺปสิรูปมา กามา วุตฺตา ภควตา พหุทุกฺขา พหุปายาสา, อาทีนโว เอตฺถ ภิยฺโยติ.}

\prose{เอวมฺปิ โข อริฏฺโฐ ภิกฺขุ คทฺธพาธิ\-ปุพฺโพ เตหิ ภิกฺขูหิ วุจฺจมาโน ตเถว ตํ ปาปกํ ทิฏฺฐิคตํ ถามสา ปรามาสา อภินิวิสฺส โวหรติ “เอวํพฺยา โข อหํ อาวุโส ภควตา ธมฺมํ เทสิตํ อาชานามิ, ยถา เยเม อนฺตรายิกา ธมฺมา วุตฺตา ภควตา, เต ปฏิเสวโต นาลํ อนฺตรายายา”ติ. ยโต จ โข เต ภิกฺขู นาสกฺขิํสุ อริฏฺฐํ ภิกฺขุํ คทฺธพาธิ\-ปุพฺพํ เอตสฺมา ปาปกา ทิฏฺฐิคตา วิเวเจตุํ, อถ โข เต ภิกฺขู เยน ภควา เตนุ\-ปสงฺกมิํสุ, อุปสงฺกมิตฺวา ภควโต เอตมตฺถํ อาโรเจสุํ. อถ โข ภควา เอตสฺมิํ นิทาเน เอตสฺมิํ ปกรเณ ภิกฺขุสํฆํ สนฺนิปาตาเปตฺวา อริฏฺฐํ ภิกฺขุํ คทฺธพาธิ\-ปุพฺพํ ปฏิปุจฺฉิ “สจฺจํ กิร เต อริฏฺฐ เอวรูปํ ปาปกํ ทิฏฺฐิคตํ อุปฺปนฺนํ ‘ตถาหํ ภควตา ธมฺมํ เทสิตํ อาชานามิ, ยถา เยเม อนฺตรายิกา ธมฺมา วุตฺตา ภควตา, เต ปฏิเสวโต นาลํ อนฺตรายายา’ติ”. เอวํพฺยา โข อหํ ภนฺเต ภควตา ธมฺมํ เทสิตํ อาชานามิ, ยถา เยเม อนฺตรายิกา ธมฺมา วุตฺตา ภควตา, เต ปฏิเสวโต นาลํ อนฺตรายายาติ.}

\prose{\csromanfolio{70}กสฺส นุ โข นาม ตฺวํ โมฆปุริส มยา เอวํ ธมฺมํ เทสิตํ อาชานาสิ. นนุ มยา โมฆปุริส อเนก\-ปริยาเยน อนฺตรายิกา ธมฺมา อนฺตรายิกา วุตฺตา, อลญฺจ ปน เต ปฏิเสวโต อนฺตรายาย. อปฺปสฺสาทา กามา วุตฺตา มยา พหุทุกฺขา พหุปายาสา, อาทีนโว เอตฺถ ภิยฺโย. อฏฺฐิ\-กงฺกลู\-ปมา กามา วุตฺตา มยา พหุทุกฺขา พหุปายาสา, อาทีนโว เอตฺถ ภิยฺโย. มํสเปสูปมา กามา วุตฺตา มยา ฯเปฯ. ติณุกฺกูปมา กามา วุตฺตา มยา. องฺคารกาสูปมา กามา วุตฺตา มยา. สุปินกูปมา กามา วุตฺตา มยา. ยาจิตกูปมา กามา วุตฺตา มยา. รูกฺขผลูปมา กามา วุตฺตา มยา. อสิสูนูปมา กามา วุตฺตา มยา. สตฺติสูลูปมา กามา วุตฺตา มยา. สปฺปสิรูปมา กามา วุตฺถา มยา พหุทุกฺขา พหุปายาสา, อาทีนโว เอตฺถ ภิยฺโย. อถ จ ปน ตฺวํ โมฆปุริส อตฺตนา ทุคฺคหิเตน\csromansharedfootnote{eeda9ad4c66dbcce}{ทุคฺคหิเตน ทิฏฺฐิคเตน (สฺยา)} อเมฺห เจว อพฺภาจิกฺขสิ, อตฺตานญฺจ ขณสิ, พหุญฺจ อปุญฺญํ ปสวสิ, ตญฺหิ เต โมฆปุริส ภวิสฺสติ ทีฆรตฺตํ อหิตาย ทุกฺขาย. เนตํ โมฆปุริส อปฺปสนฺนานํ วา ปสาทาย ฯเปฯ วิครหิตฺวา ฯเปฯ ธมฺมิํ กถํ กตฺวา ภิกฺขู อามนฺเตสิ\csromandash{} เตน หิ ภิกฺขเว สํโฆ อริฏฺฐสฺส ภิกฺขุโน คทฺธพาธิ\-ปุพฺพสฺส ปาปิกาย ทิฏฺฐิยา อปฺปฏินิสฺสคฺเค อุกฺเขปนีย\-กมฺมํ กโรตุ อสมฺโภคํ สํเฆน. เอวญฺจ ปน ภิกฺขเว กาตพฺพํ, ปฐมํ อริฏฺโฐ ภิกฺขุ โจเทตพฺโพ, โจเทตฺวา สาเรตพฺโพ, สาเรตฺวา อาปตฺติํ อาโรเปตพฺโพ, อาปตฺติํ อาโรเปตฺวา พฺยตฺเตน ภิกฺขุนา ปฏิพเลน สํโฆ ญาเปตพฺโพ\csromandash{}}

\proseitem{66}{“สุณาตุ เม ภนฺเต สํโฆ, อริฏฺฐสฺส ภิกฺขุโน คทฺธพาธิ\-ปุพฺพสฺส เอวรูปํ ปาปกํ ทิฏฺฐิคตํ อุปฺปนฺนํ ‘ตถาหํ ภควตา ธมฺมํ เทสิตํ อาชานามิ, ยถา เยเม อนฺตรายิกา ธมฺมา วุตฺตา ภควตา, เต ปฏิเสวโต นาลํ อนฺตรายายา’ติ. โส ตํ ทิฏฺฐิํ น ปฏินิสฺสชฺชติ, ยทิ สํฆสฺส ปตฺตกลฺลํ, สํโฆ อริฏฺฐสฺส ภิกฺขุโน คทฺธพาธิ\-ปุพฺพสฺส ปาปิกาย ทิฏฺฐิยา อปฺปฏินิสฺสคฺเค อุกฺเขปนีย\-กมฺมํ กเรยฺย อสมฺโภคํ สํเฆน, เอสา ญตฺติ.}

\prose{สุณาตุ เม ภนฺเต สํโฆ, อริฏฺฐสฺส ภิกฺขุโน คทฺธพาธิ\-ปุพฺพสฺส เอวรูปํ ปาปกํ ทิฏฺฐิคตํ อุปฺปนฺนํ ‘ตถาหํ ภควตา ธมฺมํ เทสิตํ อาชานามิ, \csromanfolio{71}ยถา เยเม อนฺตรายิกา ธมฺมา วุตฺตา ภควตา, เต ปฏิเสวโต นาลํ อนฺตรายายา’ติ. โส ตํ ทิฏฺฐิํ น ปฏินิสฺสชฺชติ, สํโฆ อริฏฺฐสฺส ภิกฺขุโน คทฺธพาธิ\-ปุพฺพสฺส ปาปิกาย ทิฏฺฐิยา อปฺปฏินิสฺสคฺเค อุกฺเขปนีย\-กมฺมํ กโรติ อสมฺโภคํ สํเฆน, ยสฺสายสฺมโต ขมติ อริฏฺฐสฺส ภิกฺขุโน คทฺธพาธิ\-ปุพฺพสฺส ปาปิกาย ทิฏฺฐิยา อปฺปฏินิสฺสคฺเค อุกฺเขปนียสฺส กมฺมสฺส กรณํ อสมฺโภคํ สํเฆน, โส ตุณฺหสฺส, ยสฺส นกฺขมติ, โส ภาเสยฺย.}

\prose{ทุติยมฺปิ เอตมตฺวํ วทามิ ฯเปฯ. ตติยมฺปิ เอตมตฺถํ วทามิ. สุณาตุ เม ภนฺเต สํโฆ, อริฏฺฐสฺส ภิกฺขุโน คทฺธพาธิ\-ปุพฺพสฺส เอวรูปํ ปาปกํ ทิฏฺฐิคตํ อุปฺปนฺนํ ‘ตถาหํ ภควตา ธมฺมํ เทสิตํ อาชานามิ, ยถา เยเม อนฺตรายิกา ธมฺมา วุตฺตา ภควตา, เต ปฏิเสวโต นาลํ อนฺตรายายา’ติ, โส ตํ ทิฏฺฐิํ น ปฏินิสฺสชฺชติ, สํโฆ อริฏฺฐสฺส ภิกฺขุโน คทฺธพาธิ\-ปุพฺพสฺส ปาปิกาย ทิฏฺฐิยา อปฺปฏินิสฺสคฺเค อุกฺเขปนีย\-กมฺมํ กโรติ อสมฺโภคํ สํเฆน, ยสฺสายสฺมโต ขมติ อริฏฺฐสฺส ภิกฺขุโน คทฺธพาธิ\-ปุพฺพสฺส ปาปิกาย ทิฏฺฐิยา อปฺปฏินิสฺสคฺเค อุกฺเขปนียสฺส กมฺมสฺส กรณํ อสมฺโภคํ สํเฆน, โส ตุณฺหสฺส, ยสฺส นกฺขมติ, โส ภาเสยฺย.}

\prose{กตํ สํเฆน อริฏฺฐสฺส ภิกฺขุโน คทฺธพาธิ\-ปุพฺพสฺส ปาปิกาย ทิฏฺฐิยา อปฺปฏินิสฺสคฺเค อุกฺเขปนีย\-กมฺมํ อสมฺโภคํ สํเฆน, ขมติ สํฆสฺส, ตสฺมา ตุณฺหี, เอวเมตํ ธารยามี”ติ.}

\prose{อาวาสปรมฺปรญฺจ ภิกฺขเว สํสถ “อริฏฺโฐ ภิกฺขุ คทฺธพาธิ\-ปุพฺโพ สํเฆน ปาปิกาย\csromansharedfootnote{b4a36321f5652c90}{คทฺธพาธิปุพฺโพ ปาปิกาย (สี, ก)} ทิฏฺฐิยา อปฺปฏินิสฺสคฺเค อุกฺเขปนีย\-กมฺมกโต อสมฺโภคํ สํเฆนา”ติ.}

\csromanmatikaanchor{mtk.45}
\csromantocmarkref{subsection}{อธมฺมกมฺมทฺวาทสก}{mtk.45}
\bookmark[dest={mtk.45},level=2]{อธมฺมกมฺมทฺวาทสก}
\csromanheader{2}{อธมฺมกมฺม\-ทฺวาทสก}
\proseitem{67}{ตีหิ ภิกฺขเว องฺเคหิ สมนฺนาคตํ ปาปิกาย ทิฏฺฐิยา อปฺปฏินิสฺสคฺเค อุกฺเขปนีย\-กมฺมํ อธมฺมกมฺมญฺจ โหติ อวินยกมฺมญฺจ ทุวูปสนฺตญฺจ. อสมฺมุขา กตํ โหติ, อปฺปฏิปุจฺฉา กตํ โหติ, อปฺปฏิญฺญาย กตํ โหติ. อิเมหิ โข ภิกฺขเว ตีหงฺเคหิ สมนฺนาคตํ \csromanfolio{72}ปาปิกาย ทิฏฺฐิยา อปฺปฏินิสฺสคฺเค อุกฺเขปนีย\-กมฺมํ อธมฺมกมฺมญฺจ โหติ อวินยกมฺมญฺจ ทุวูปสนฺตญฺจ.}

\prose{อปเรหิปิ ภิกฺขเว ตีหงฺเคหิ สมนฺนาคตํ ปาปิกาย ทิฏฺฐิยา อปฺปฏินิสฺสคฺเค อุกฺเขปนีย\-กมฺมํ อธมฺมกมฺมญฺจ โหติ อวินยกมฺมญฺจ ทุวูปสนฺตญฺจ. อนาปตฺติยา กตํ โหติ, อเทสนาคามินิยา อาปตฺติยา กตํ โหติ, เทสิตาย อาปตฺติยา กตํ โหติ ฯเปฯ. อโจเทตฺวา กตํ โหติ, อสาเรตฺวา กตํ โหติ, อาปตฺติํ อนาโรเปตฺวา กตํ โหติ ฯเปฯ. อสมฺมุขา กตํ โหติ, อธมฺเมน กตํ โหติ, วคฺเคน กตํ โหติ ฯเปฯ. อปฺปฏิปุจฺฉา กตํ โหติ, อธมฺเมน กตํ โหติ, วคฺเคน กตํ โหติ ฯเปฯ. อปฺปฏิญฺญาย กตํ โหติ, อธมฺเมน กตํ โหติ, วคฺเคน กตํ โหติ ฯเปฯ. อนาปตฺติยา กตํ โหติ, อธมฺเมน กตํ โหติ, วคฺเคน กตํ โหติ ฯเปฯ. อเทสนาคามินิยา อาปตฺติยา กตํ โหติ, อธมฺเมน กตํ โหติ, วคฺเคน กตํ โหติ ฯเปฯ. เทสิตาย อาปตฺติยา กตํ โหติ, อธมฺเมน กตํ โหติ, วคฺเคน กตํ โหติ ฯเปฯ. อโจเทตฺวา กตํ โหติ, อธมฺเมน กตํ โหติ, วคฺเคน กตํ โหติ ฯเปฯ. อาสาเรตฺวา กตํ โหติ, อธมฺเมน กตํ โหติ, วคฺเคน กตํ โหติ ฯเปฯ. อาปตฺติํ อนาโรเปตฺวา กตํ โหติ, อธมฺเมน กตํ โหติ, วคฺเคน กตํ โหติ. อิเมหิ โข ภิกฺขเว ตีหงฺเคหิ สมนฺนาคตํ ปาปิกาย ทิฏฺฐิยา อปฺปฏินิสฺสคฺเค อุกฺเขปนีย\-กมฺมํ อธมฺมกมฺมญฺจ โหติ อวินยกมฺมญฺจ ทุวูปสนฺตญฺจ.}

\csromanheader{2}{ปาปิกาย ทิฏฺฐิยา อปฺปฏินิสฺสคฺเค อุกฺเขปนีย\-กมฺเม}
\nitthitam{อธมฺมกมฺม\-ทฺวาทสกํ นิฏฺฐิตํ.}
\csromanmatikaanchor{mtk.46}
\csromantocmarkref{subsection}{ธมฺมกมฺมทฺวาทสก}{mtk.46}
\bookmark[dest={mtk.46},level=2]{ธมฺมกมฺมทฺวาทสก}
\csromanheader{2}{ธมฺมกมฺม\-ทฺวาทสก}
\proseitem{68}{ตีหิ ภิกฺขเว องฺเคหิ สมนฺนาคตํ ปาปิกาย ทิฏฺฐิยา อปฺปฏินิสฺสคฺเค อุกฺเขปนีย\-กมฺมํ ธมฺมกมฺมญฺจ โหติ วินยกมฺมญฺจ สุวูปสนฺตญฺจ. สมฺมุขา กตํ โหติ, ปฏิปุจฺฉา กตํ โหติ, ปฏิญฺญาย กตํ โหติ. อิเมหิ โข ภิกฺขเว ตีหงฺเคหิ สมนฺนาคตํ ปาปิกาย ทิฏฺฐิยา อปฺปฏินิสฺสคฺเค อุกฺเขปนีย\-กมฺมํ ธมฺมกมฺมญฺจ โหติ วินยกมฺมญฺจ สุวูปสนฺตญฺจ.}

\csromanmatikaanchor{mtk.47}
\csromantocmarkref{subsection}{อากงฺขมานฉกฺก}{mtk.47}
\bookmark[dest={mtk.47},level=2]{อากงฺขมานฉกฺก}
\prosewithcloser{\csromanfolio{73}อปเรหิปิ ภิกฺขเว ตีหงฺเคหิ สมนฺนาคตํ ปาปิกาย ทิฏฺฐิยา อปฺปฏินิสฺสคฺเค อุกฺเขปนีย\-กมฺมํ ธมฺมกมฺมญฺจ โหติ วินยกมฺมญฺจ สุวูปสนฺตญฺจ. อาปตฺติยา กตํ โหติ, เทสนาคามินิยา อาปตฺติยา กตํ โหติ, อเทสิตาย อาปตฺติยา กตํ โหติ ฯเปฯ. โจเทตฺวา กตํ โหติ, สาเรตฺวา กตํ โหติ, อาปตฺติํ อาโรเปตฺวา กตํ โหติ ฯเปฯ. สมฺมุขา กตํ โหติ, ธมฺเมน กตํ โหติ, สมคฺเคน กตํ โหติ ฯเปฯ. ปฏิปุจฺฉา กตํ โหติ, ธมฺเมน กตํ โหติ, สมคฺเคน กตํ โหติ ฯเปฯ. ปฏิญฺญาย กตํ โหติ, ธมฺเมน กตํ โหติ, สมคฺเคน กตํ โหติ ฯเปฯ. อาปตฺติยา กตํ โหติ, ธมฺเมน กตํ โหติ, สมคฺเคน กตํ โหติ ฯเปฯ. เทสนาคามินิยา อาปตฺติยา กตํ โหติ, ธมฺเมน กตํ โหติ, สมคฺเคน กตํ โหติ ฯเปฯ. อเทสิตาย อาปตฺติยา กตํ โหติ, ธมฺเมน กตํ โหติ, สมคฺเคน กตํ โหติ ฯเปฯ. โจเทตฺวา กตํ โหติ, ธมฺเมน กตํ โหติ, สมคฺเคน กตํ โหติ ฯเปฯ. สาเรตฺวา กตํ โหติ, ธมฺเมน กตํ โหติ, สมคฺเคน กตํ โหติ ฯเปฯ. อาปตฺติํ อาโรเปตฺวา กตํ โหติ, ธมฺเมน กตํ โหติ, สมคฺเคน กตํ โหติ. อิเมหิ โข ภิกฺขเว ตีหงฺเคหิ สมนฺนาคตํ ปาปิกาย ทิฏฺฐิยา อปฺปฏินิสฺสคฺเค อุกฺเขปนีย\-กมฺมํ ธมฺมกมฺมญฺจ โหติ วินยกมฺมญฺจ สุวูปสนฺตญฺจ.}{ปาปิกาย ทิฏฺฐายา อปฺปฏินิสฺสคฺเค อุกฺเขปนีย\-กมฺเม ธมฺมกมฺม\-ทฺวาทสกํ นิฏฺฐิตํ.}
\proseitem{69}{\csromansymbolfootnote{*}{วิ 5. 220 ปิฏฺเฐปิ.}ตีหิ ภิกฺขเว องฺเคหิ สมนฺนาคตสฺส ภิกฺขุโน อากงฺขมาโน สํโฆ ปาปิกาย ทิฏฺฐิยา อปฺปฏินิสฺสคฺเค อุกฺเขปนีย\-กมฺมํ กเรยฺย. ภณฺฑนการโก โหติ กลหการโก วิวาทการโก ภสฺสการโก สํเฆ อธิกรณ\-การโก, พาโล โหติ อพฺยตฺโต อาปตฺติพหุโล อนปทาโน, คิหิสํสฏฺโฐ วิหรติ อนนุโลมิเกหิ คิหิสํสคฺเคหิ. อิเมหิ โข ภิกฺขเว ตีหงฺเคหิ สมนฺนาคตสฺส ภิกฺขุโน อากงฺขมาโน สํโฆ ปาปิกาย ทิฏฺฐิยา อปฺปฏินิสฺสคฺเค อุกฺเขปนีย\-กมฺมํ กเรยฺย. (1)}

\prose{\csromanfolio{74}อปเรหิปิ ภิกฺขเว ตีหงฺเคหิ สมนฺนาคตสฺส ภิกฺขุโน อากงฺขมาโน สํโฆ ปาปิกาย ทิฏฺฐิยา อปฺปฏินิสฺสคฺเค อุกฺเขปนีย\-กมฺมํ กเรยฺย. อธิสีเล สีลวิปนฺโน โหติ, อชฺฌาจาเร อาจารวิปนฺโน โหติ, อติทิฏฺฐิยา ทิฏฺฐิวิปนฺโน โหติ. อิเมหิ โข ภิกฺขเว ตีหงฺเคหิ สมนฺนาคตสฺส ภิกฺขุโน อากงฺขมาโน สํโฆ ปาปิกาย ทิฏฺฐิยา อปฺปฏินิสฺสคฺเค อุกฺเขปนีย\-กมฺมํ กเรยฺย. (2)}

\prose{อปเรหิปิ ภิกฺขเว ตีหงฺเคหิ สมนฺนาคตสฺส ภิกฺขุโน อากงฺขมาโน สํโฆ ปาปิกาย ทิฏฺฐิยา อปฺปฏินิสฺสคฺเค อุกฺเขปนีย\-กมฺมํ กเรยฺย. พุทฺธสฺส อวณฺณํ ภาสติ, ธมฺมสฺส อวณฺณํ ภาสติ, สํฆสฺส อวณฺณํ ภาสติ. อิเมหิ โข ภิกฺขเว ตีหงฺเคหิ สมนฺนาคตสฺส ภิกฺขุโน อากงฺขมาโน สํโฆ ปาปิกาย ทิฏฺฐิยา อปฺปฏินิสฺสคฺเค อุกฺเขปนีย\-กมฺมํ กเรยฺย. (3)}

\prose{ติณฺณํ ภิกฺขเว ภิกฺขูนํ อากงฺขมาโน สํโฆ ปาปิกาย ทิฏฺฐิยา อปฺปฏินิสฺสคฺเค อุกฺเขปนีย\-กมฺมํ กเรยฺย. เอโก ภณฺฑนการโก โหติ กลหการโก วิวาทการโก ภสฺสการโก สํเฆ อธิกรณ\-การโก, เอโก พาโล โหติ อพฺยตฺโต อาปตฺติพหุโล อนปทาโน, เอโก คิหิสํสฏฺโฐ วิหรติ อนนุโลมิเกหิ คิหิสํสคฺเคหิ. อิเมสํ โข ภิกฺขเว ติณฺณํ ภิกฺขูนํ อากงฺขมาโน สํโฆ ปาปิกาย ทิฏฺฐิยา อปฺปฏินิสฺสคฺเค อุกฺเขปนีย\-กมฺมํ กเรยฺย. (4)}

\prose{อปเรสมฺปิ ภิกฺขเว ติณฺณํ ภิกฺขูนํ อากงฺขมาโน สํโฆ ปาปิกาย ทิฏฺฐิยา อปฺปฏินิสฺสคฺเค อุกฺเขปนีย\-กมฺมํ กเรยฺย. เอโก อธิสีเล สีลวิปนฺโน โหติ, เอโก อชฺฌาจาเร อาจารวิปนฺโน โหติ, เอโก อติทิฏฺฐิยา ทิฏฺฐิวิปนฺโน โหติ. อิเมสํ โข ภิกฺขเว ติณฺณํ ภิกฺขูนํ อากงฺขมาโน สํโฆ ปาปิกาย ทิฏฺฐิยา อปฺปฏินิสฺสคฺเค อุกฺเขปนีย\-กมฺมํ กเรยฺย. (5)}

\prosewithcloser{อปเรสมฺปิ ภิกฺขเว ติณฺณํ ภิกฺขูนํ อากงฺขมาโน สํโฆ ปาปิกาย ทิฏฺฐิยา อปฺปฏินิสฺสคฺเค อุกฺเขปนีย\-กมฺมํ กเรยฺย. เอโก พุทฺธสฺส อวณฺณํ ภสติ, เอโก ธมฺมสฺส อวณฺณํ ภาสติ, เอโก สํฆสฺส อวณฺณํ \csromanfolio{75}ภาสติ. อิเมสํ โข ภิกฺขเว ติณฺณํ ภิกฺขูนํ อากงฺขมาโน สํโฆ ปาปิกาย ทิฏฺฐิยา อปฺปฏินิสฺสคฺเค อุกฺเขปนีย\-กมฺมํ กเรยฺย.}{(6) ปาปิกาย ทิฏฺฐิยา อปฺปฏินิสฺสคฺเค อุกฺเขปนีย\-กมฺเม อกงฺขมานฉกฺกํ นิฏฺฐิตํ.}
\csromanmatikaanchor{mtk.48}
\csromantocmarkref{subsection}{เตจตฺตาลีสวตฺต}{mtk.48}
\bookmark[dest={mtk.48},level=2]{เตจตฺตาลีสวตฺต}
\csromanheader{2}{เตจตฺตาลีส\-วตฺต}
\proseitemwithcloser{70}{ปาปิกาย ทิฏฺฐิยา อปฺปฏินิสฺสคฺเค อุกฺเขปนีย\-กมฺมกเตน ภิกฺขเว ภิกฺขุนา สมฺมา วตฺติตพฺพํ. ตตฺรายํ สมฺมาวตฺตนา\csromandash{}น อุปสมฺ\-ปาเท\-ตพฺพํ, น นิสฺสโย ทาตพฺโพ, น สามเณโร อุปฏฺฐาเปตพฺโพ, น ภิกฺขุโนวาทก\-สมฺมุติ สาทิตพฺพา, สมฺมเตนปิ ภิกฺขุนิโย น โอวทิตพฺพา, ยาย อาปตฺติยา สํเฆน ปาปิกาย ทิฏฺฐิยา อปฺปฏินิสฺสคฺเค อุกฺเขปนีย\-กมฺมํ กตํ โหติ, สา อาปตฺติ น อาปชฺชิตพฺพา, อญฺญา วา ตาทิสิกา, ตโต วา ปาปิฏฺฐตรา, กมฺมํ น ครหิตพฺพํ, กมฺมิกา น ครหิตพฺพา - ป-. น ปกตตฺตสฺส ภิกฺขุโน อุโปสโถ ฐเปตพฺโพ, น ปวารณา ฐเปตพฺพา, น สวจนียํ กาตพฺพํ, น อนุวาโท ปฏฺฐเปตพฺโพ, น โอกาโส กาเรตพฺโพ, น โจเทตพฺโพ, น สาเรตพฺโพ, น ภิกฺขูหิ สมฺป\-โยเช\-ตพฺพนฺติ.}{ปาปิกาย ทิฏฺฐิยา อปฺปฏินิสฺสคฺเค อุกฺเขปนีย\-กมฺเม เตจตฺตาลีส\-วตฺตํ นิฏฺฐิตํ.}
\proseitem{71}{อถ โข สํโฆ อริฏฺฐสฺส ภิกฺขุโน คทฺธพาธิ\-ปุพฺพสฺส ปาปิกาย ทิฏฺฐิยา อปฺปฏินิสฺสคฺเค อุกฺเขปนีย\-กมฺมํ อกาสิ อสมฺโภคํ สํเฆน. โส สํเฆน ปาปิกาย ทิฏฺฐิยา อปฺปฏินิสฺสคฺเค อุกฺเขปนีย\-กมฺมกโต วิพฺภมิ. เย เต ภิกฺขู อปฺปิจฺฉา ฯเปฯ เต อุชฺฌายนฺติ ขิยฺยนฺติ วิปาเจนฺติ “กถํ หิ นาม อริฏฺโฐ ภิกฺขุ คทฺธพาธิ\-ปุพฺโพ สํเฆน ปาปิกาย ทิฏฺฐิยา อปฺปฏินิสฺสคฺเค อุกฺเขปนีย\-กมฺมกโต วิพฺภมิสฺสตี”ติ. อถ โข เต ภิกฺขู ภควโต เอตมตฺถํ อาโรเจสุํ.}

\prose{อถ โข ภควา เอตสฺมิํ นิทาเน เอตสฺมิํ ปกรเณ ภิกฺขุสํฆํ สนฺนิปาตาเปตฺวา ภิกฺขู ปฏิปุจฺฉิ “สจฺจํ กิร ภิกฺขเว อริฏฺโฐ ภิกฺขุ \csromanfolio{76}คทฺธพาธิ\-ปุพฺโพ สํเฆน ปาปิกาย ทิฏฺฐิยา อปฺปฏินิสฺสคฺเค อุกฺเขปนีย\-กมฺมกโต วิพฺภมตี”ติ\csromansharedfootnote{9e59d08f67e97a2e}{วิพฺภมีติ (สี, ก)}. สจฺจํ ภควาติ. วิครหิ พุทฺโธ ภควา อนนุจฺฉวิกํ ฯเปฯ. กถํ หิ นาม โส ภิกฺขเว โมฆปุริโส สํเฆน ปาปิกาย ทิฏฺฐิยา อปฺปฏินิสฺสคฺเค อุกฺเขปนีย\-กมฺมกโต วิพฺภมิสฺสติ. เนตํ ภิกฺขเว อปฺปสนฺนานํ วา ปสาทาย ฯเปฯ วิครหิตฺวา ฯเปฯ ธมฺมิํ กถํ กตฺวา ภิกฺขู อามนฺเตสิ\csromandash{} เตน หิ ภิกฺขเว สํโฆ ปาปิกาย ทิฏฺฐิยา อปฺปฏินิสฺสคฺเค อุกฺเขปนีย\-กมฺมํ ปฏิปฺปสฺสมฺเภตุ. นปฺปฏิปฺปสฺสมฺเภตพฺพ\-เตจตฺตาลีสก}

\proseitem{72}{ปญฺจหิ ภิกฺขเว องฺเคหิ สมนฺนาคตสฺส ภิกฺขุโน ปาปิกาย ทิฏฺฐิยา อปฺปฏินิสฺสคฺเค อุกฺเขปนีย\-กมฺมํ นปฺปฏิปฺปสฺสมฺเภตพฺพํ. อุปสมฺปาเทติ, นิสฺสยํ เทติ, สามเณรํ อุปฏฺฐาเปติ, ภิกฺขุโนวาทก\-สมฺมุติํ สาทิยติ, สมฺมโตปิ ภิกฺขุนิโย โอวทติ. อิเมหิ โข ภิกฺขเว ปญฺจหงฺเคหิ สมนฺนาคตสฺส ภิกฺขุโน ปาปิกาย ทิฏฺฐิยา อปฺปฏินิสฺสคฺเค อุกฺเขปนีย\-กมฺมํ นปฺปฏิปฺปสฺสมฺเภตพฺพํ.}

\prose{\csromansymbolfootnote{*}{วิ 5. 317 ปิฏฺเฐปิ.}อปเรหิปิ ภิกฺขเว ปญฺจหงฺเคหิ สมนฺนาคตสฺส ภิกฺขุโน ปาปิกาย ทิฏฺฐิยา อปฺปฏินิสฺสคฺเค อุกฺเขปนีย\-กมฺมํ นปฺปฏิปฺปสฺสมฺเภ- ตพฺพํ. ยาย อาปตฺติยา สํเฆน ปาปิกาย ทิฏฺฐิยา อปฺปฏินิสฺสคฺเค อุกฺเขปนีย\-กมฺมํ กตํ โหติ, ตํ อาปตฺติํ อาปชฺชติ, อญฺญํ วา ตาทิสิกํ, ตโต วา ปาปิฏฺฐตรํ, กมฺมํ ครหติ, กมฺมิเก ครหติ. อิเมหิ โข ภิกฺขเว ปญฺจหงฺเคหิ สมนฺนาคตสฺส ภิกฺขุโน ปาปิกาย ทิฏฺฐิยา อปฺปฏินิสฺสคฺเค อุกฺเขปนีย\-กมฺมํ นปฺปฏิปฺปสฺสมฺเภตพฺพํ ฯเปฯ.}

\prose{อฏฺฐหิ ภิกฺขเว องฺเคหิ สมนฺนาคตสฺส ภิกฺขุโน ปาปิกาย ทิฏฺฐิยา อปฺปฏินิสฺสคฺเค อุกฺเขปนีย\-กมฺมํ นปฺปฏิปฺปสฺสมฺเภตพฺพํ. ปกตตฺตสฺส ภิกฺขุโน อุโปสถํ ฐเปติ, ปวารณํ ฐเปติ, สวจนียํ กโรติ, อนุวาทํ ปฏฺฐเปติ, โอกาสํ กาเรติ, โจเทติ, สาเรติ, ภิกฺขูหิ สมฺปโยเชติ. อิเมหิ โข ภิกฺขเว อฏฺฐหงฺเคหิ สมนฺนาคตสฺส ภิกฺขุโน ปาปิกาย ทิฏฺฐิยา อปฺปฏินิสฺสคฺเค อุกฺเขปนีย\-กมฺมํ นปฺปฏิปฺปสฺสมฺเภตพฺพํ. ปาปิกาย ทิฏฺฐิยา อปฺปฏินิสฺสคฺเค อุกฺเขปนีย\-กมฺเม นปฺปฏิปฺปสฺสมฺเภตพฺพ\-เตจตฺตาลีสกํ นิฏฺฐิตํ.}

\csromanmatikaanchor{mtk.49}
\csromanmatikaanchor{mtk.50}
\csromantocmarkref{subsection}{นปฺปฏิปฺปสฺสมฺเภตพฺพเตจตฺตาลีสก}{mtk.49}
\bookmark[dest={mtk.49},level=2]{นปฺปฏิปฺปสฺสมฺเภตพฺพเตจตฺตาลีสก}
\csromantocmarkref{subsection}{ปฏิปฺปสฺสมฺเภตพฺพเตจตฺตาลีสก}{mtk.50}
\bookmark[dest={mtk.50},level=2]{ปฏิปฺปสฺสมฺเภตพฺพเตจตฺตาลีสก}
\csromanheader{2}{ปฏิปฺปสฺสมฺเภตพฺพ\-เตจตฺตาลีสก}
\proseitem{73}{\csromanfolio{77}ปญฺจหิ ภิกฺขเว องฺเคหิ สมนฺนาคตสฺส ภิกฺขุโน ปาปิกาย ทิฏฺฐิยา อปฺปฏินิสฺสคฺเค อุกฺเขปนีย\-กมฺมํ ปฏิปฺปสฺสมฺเภตพฺพํ. น อุปสมฺปาเทติ, น นิสฺสยํ เทติ, น สามเณรํ อุปฏฺฐาเปติ, น ภิกฺขุโนวาทก\-สมฺมุติํ สาทิยติ, สมฺมโตปิ ภิกฺขุนิโย น โอวทติ. อิเมหิ โข ภิกฺขเว ปญฺจหงฺเคหิ สมนฺนาคตสฺส ภิกฺขุโน ปาปิกาย ทิฏฺฐิยา อปฺปฏินิสฺสคฺเค อุกฺเขปนีย\-กมฺมํ ปฏิปฺปสฺสมฺเภตพฺพํ.}

\prose{อปเรหิปิ ภิกฺขเว ปญฺจหงฺเคหิ สมนฺนาคตสฺส ภิกฺขุโน ปาปิกาย ทิฏฺฐิยา อปฺปฏินิสฺสคฺเค อุกฺเขปนีย\-กมฺมํ ปฏิปฺปสฺสมฺเภตพฺพํ. ยาย อาปตฺติยา สํเฆน ปาปิกาย ทิฏฺฐิยา อปฺปฏินิสฺสคฺเค อุกฺเขปนีย\-กมฺมํ กตํ โหติ, ตํ อาปตฺติํ น อาปชฺชติ, อญฺญํ วา ตาทิสิกํ, ตโต วา ปาปิฏฺฐตรํ, กมฺมํ น ครหติ, กมฺมิเก น ครหติ. อิเมหิ โข ภิกฺขเว ปญฺจหงฺเคหิ สมนฺนาคตสฺส ภิกฺขุโน ปาปิกาย ทิฏฺฐิยา อปฺปฏินิสฺสคฺเค อุกฺเขปนีย\-กมฺมํ ปฏิปฺปสฺสมฺเภตพฺพํ ฯเปฯ.}

\prose{อฏฺฐหิ ภิกฺขเว องฺเคหิ สมนฺนาคตสฺส ภิกฺขุโน ปาปิกาย ทิฏฺฐิยา อปฺปฏินิสฺสคฺเค อุกฺเขปนีย\-กมฺมํ ปฏิปฺปสฺสมฺเภตพฺพํ. น ปกตตฺตสฺส ภิกฺขุโน อุโปสถํ ฐเปติ, น ปวารณํ ฐเปติ, น สวจนียํ กโรติ, น อนุวาทํ ปฏฺฐเปติ, น โอกาสํ กาเรติ, น โจเทติ, น สาเรติ, น ภิกฺขูหิ สมฺปโยเชติ. อิเมหิ โข ภิกฺขเว อฏฺฐหงฺเคหิ สมนฺนาคตสฺส ภิกฺขุโน ปาปิกาย ทิฏฺฐิยา อปฺปฏินิสฺสคฺเค อุกฺเขปนีย\-กมฺมํ ปฏิปฺปสฺสมฺเภตพฺพํ. ปาปิกาย ทิฏฺฐิยา อปฺปฏินิสฺสคฺเค อุกฺเขปนีย\-กมฺเม ปฏิปฺปสฺสมฺเภตพฺพ\-เตจตฺตาลีสกํ นิฏฺฐิตํ.}

\proseitem{74}{เอวญฺจ ปน ภิกฺขเว ปฏิปฺปสฺสมฺเภตพฺพํ, เตน ภิกฺขเว ปาปิกาย ทิฏฺฐิยา อปฺปฏินิสฺสคฺเค อุกฺเขปนีย\-กมฺมกเตน ภิกฺขุนา สํฆํ อุปสงฺกมิตฺวา เอกํสํ อุตฺตราสงฺคํ กริตฺวา วุฑฺฒานํ ภิกฺขูนํ ปาเท วนฺทิตฺวา อุกฺกุฏิกํ นิสีทิตฺวา อญฺชลิํ ปคฺคเหตฺวา เอวมสฺส วจนีโย “อหํ ภนฺเต สํเฆน ปาปิกาย ทิฏฺฐิยา อปฺปฏินิสฺสคฺเค อุกฺเขปนีย\-กมฺมกโต สมฺมา วตฺตามิ, โลมํ ปาเตมิ, เนตฺถารํ วตฺตามิ, ปาปิกาย ทิฏฺฐิยา อปฺปฏินิสฺสคฺเค \csromanfolio{78}อุกฺเขปนียสฺส กมฺมสฺส ปฏิปฺปสฺสทฺธิํ ยาจามี”ติ. ทุติยมฺปิ ยาจิตพฺพา. ตติยมฺปิ ยาจิตพฺพา. พฺยตฺเตน ภิกฺขุนา ปฏิพเลน สํโฆ ญาเปตพฺโพ\csromandash{} “สุณาตุ เม ภนฺเต สํโฆ, อยํ อิตฺถนฺนาโม ภิกฺขุ สํเฆน ปาปิกาย ทิฏฺฐิยา อปฺปฏินิสฺสคฺเค อุกฺเขปนีย\-กมฺมกโต สมฺมา วตฺตติ, โลมํ ปาเตติ, เนตฺถารํ วตฺตติ, ปาปิกาย ทิฏฺฐิยา อปฺปฏินิสฺสคฺเค อุกฺเขปนียสฺส กมฺมสฺส ปฏิปฺปสฺสทฺธิํ ยาจติ, ยทิ สํฆสฺส ปตฺตกลฺลํ, สํโฆ อิตฺถนฺนามสฺส ภิกฺขุโน ปาปิกาย ทิฏฺฐิยา อปฺปฏินิสฺสคฺเค อุกฺเขปนีย\-กมฺมํ ปฏิปฺปสฺสมฺเภยฺย, เอสา ญตฺติ. สุณาตุ เม ภนฺเต สํโฆ, อยํ อิตฺถนฺนาโม ภิกฺขุ สํเฆน ปาปิกาย ทิฏฺฐิยา อปฺปฏินิสฺสคฺเค อุกฺเขปนีย\-กมฺมกโต สมฺมา วตฺตติ, โลมํ ปาเตติ, เนตฺถารํ วตฺตติ, ปาปิกาย ทิฏฺฐิยา อปฺปฏินิสฺสคฺเค อุกฺเขปนียสฺส กมฺมสฺส ปฏิปฺปสฺสทฺธิํ ยาจติ, สํโฆ อิตฺถนฺนามสฺส ภิกฺขุโน ปาปิกาย ทิฏฺฐิยา อปฺปฏินิสฺสคฺเค อุกฺเขปนีย\-กมฺมํ ปฏิปฺปสฺสมฺเภติ, ยสฺสายสฺมโต ขมติ อิตฺถนฺนามสฺส ภิกฺขุโน ปาปิกาย ทิฏฺฐิยา อปฺปฏินิสฺสคฺเค อุกฺเขปนียสฺส กมฺมสฺส ปฏิปฺปสฺสทฺธิ, โส ตุณฺหสฺส, ยสฺส นกฺขมติ, โส ภาเสยฺย.}

\prose{ทุติยมฺปิ เอตมตฺถํ วทามิ ฯเปฯ. ตติยมฺปิ เอตมตฺถํ วทามิ. สุณาตุ เม ภนฺเต สํโฆ, อยํ อิตฺถนฺนาโม ภิกฺขุ สํเฆน ปาปิกาย ทิฏฺฐิยา อปฺปฏินิสฺสคฺเค อุกฺเขปนีย\-กมฺมกโต สมฺมา วตฺตติ, โลมํ ปาเตติ, เนตฺถารํ วตฺตติ, ปาปิกาย ทิฏฺฐิยา อปฺปฏินิสฺสคฺเค อุกฺเขปนียสฺส กมฺมสฺส ปฏิปฺปสฺสทฺธิํ ยาจติ, สํโฆ อิตฺถนฺนามสฺส ภิกฺขุโน ปาปิกาย ทิฏฺฐิยา อปฺปฏินิสฺสคฺเค อุกฺเขปนีย\-กมฺมํ ปฏิปฺปสฺสมฺเภติ, ยสฺสายสฺมโต ขมติ อิตฺถนฺนามสฺส ภิกฺขุโน ปาปิกาย ทิฏฺฐิยา อปฺปฏินิสฺสคฺเค อุกฺเขปนียสฺส กมฺมสฺส ปฏิปฺปสฺสทฺธิ, โส ตุณฺหสฺส, ยสฺส นกฺขมติ, โส ภาเสยฺย.}

\prose{ปฏิปฺปสฺสทฺธํ สํเฆน อิตฺถนฺนามสฺส ภิกฺขุโน ปาปิกาย ทิฏฺฐิยา อปฺปฏินิสฺสคฺเค อุกฺเขปนีย\-กมฺมํ, ขมติ สํฆสฺส, ตสฺมา ตุณฺหี, เอวเมตํ ธารยามี”ติ. ปาปิกาย ทิฏฺฐิยา อปฺปฏินิสฺสคฺเค อุกฺเขปนีย\-กมฺมํ นิฏฺฐิตํ สตฺตมํ. กมฺม\-กฺขนฺธโก ปฐโม. อิมมฺหิ ขนฺธเก วตฺถู สตฺต.}

\vspace{\tipitakaparskip}
\csromancenter{\textbf{ตสฺสุทฺทานํ}}
\csromanmatikaanchor{mtk.51}
\csromantocmarkref{section}{อุทฺทานคาถา}{mtk.51}
\bookmark[dest={mtk.51},level=1]{อุทฺทานคาถา}
\setlength{\csromangathapagewidth}{0pt}
\setlength{\csromangathawakwidth}{0pt}
\csromangathameasuregroup{ปณฺฑุโลหิตกา ภิกฺขู, \\ ตาทิเส อุปสงฺกมฺม, \\ อนุปฺปนฺนาปิ ชายนฺติ\textsuperscript{1}, \\ อปฺปิจฺฉา เปสลา ภิกฺขู, \\ สทฺธมฺมฏฺฐิติโก พุทฺโธ, \\ อาณาเปสิ ตชฺชนีย, \\ อสมฺมุขาปฺปฏิปุจฺฉา, \\ อนาปตฺติ อเทสเน, \\ อโจเทตฺวา อสาเรตฺวา, \\ อสมฺมุขา อธมฺเมน, \\ อปฺปฏิปุจฺฉา อธมฺเมน, \\ อปฺปฏิญฺญาย อธมฺเมน, \\ อนาปตฺติ\textsuperscript{1} อธมฺเมน, \\ อเทสนาคามินิยา, \\ เทสิตาย อธมฺเมน, \\ อโจเทตฺวา อธมฺเมน, \\ อสาเรตฺวา อธมฺเมน, \\ อนาโรเปตฺวา อธมฺเมน, \\ กณฺหวารนเยเนว, \\ สํโฆ อากงฺขมาโน จ, \\ ภณฺฑนํ พาโล สํสฏฺโฐ, \\ อติทิฏฺฐิวิปนฺนสฺส, \\ พุทฺธธมฺมสฺส สํฆสฺส, \\ ติณฺณนฺนมฺปิ จ ภิกฺขูนํ, \\ ภณฺฑนํ การโก เอโก, \\ อธิสีเล อชฺฌาจาเร, \\ พุทฺธธมฺมสฺส สํฆสฺส, \\ ตชฺชนียกมฺมกโต, \\ อุปสมฺปทนิสฺสยา, \\ โอวาทสมฺมเตนาปิ, \\ นาปชฺเช ตญฺจ อาปตฺติํ, \\ กมฺมญฺจ กมฺมิเก จาปิ, \\ อุโปสถํ ปวารณํ, \\ สวจนิํ\textsuperscript{1} อนุวาโท, \\ สารณํ สมฺปโยคญฺจ, \\ อุปสมฺปทนิสฺสยา, \\ โอวาทสมฺมเตนาปิ, \\ ตญฺจาปชฺชติ อาปตฺติํ, \\ กมฺมญฺจ กมฺมิเก จาปิ, \\ อุโปสถํ ปวารณํ, \\ โอกาโส โจทนญฺเจว, \\ อิเมหฏฺฐงฺเคหิ โย ยุตฺโต, \\ กณฺหวารนเยเนว, \\ พาโล อาปตฺติพหุโล, \\ นิยสฺสกมฺมํ สมฺพุทฺโธ, \\ กีฏาคิริสฺมิํ เทฺว ภิกฺขู, \\ อนาจารญฺจ วิวิธํ, \\ ปพฺพาชนียํ สมฺพุทฺโธ, \\ มจฺฉิกาสณฺเฑ สุธมฺโม, \\ ชาติวาเทน ขุํเสติ, \\ ปฏิสารณียกมฺมํ, \\ โกสมฺพิยํ ฉนฺนํ ภิกฺขุํ, \\ อทสฺสเน อุกฺขิปิตุํ, \\ ฉนฺโน ตํเยว อาปตฺติํ, \\ อุกฺเขปนาปฺปฏิกมฺเม, \\ ปาปทิฏฺฐิ อริฏฺฐสฺส, \\ ทิฏฺฐิยาปฺปฏินิสฺสคฺเค\textsuperscript{1}, \\ นิยสฺสกมฺมํ ปพฺพชฺชํ\textsuperscript{1}, \\ อทสฺสนาปฺปฏิกมฺเม, \\ ทวานาจารูปฆาติ, \\ ปพฺพาชนียกมฺมมฺหิ, \\ อลาภาวณฺณา เทฺว ปญฺจ, \\ ปฏิสารณียกมฺมมฺหิ, \\ ตชฺชนียํ นิยสฺสญฺจ, \\ ปพฺพชฺชา\textsuperscript{1} ปฏิสารี จ, \\ ตโย อุกฺเขปนา กมฺมา, \\ ตชฺชนียนเยนาปิ,}{\csromangathabat{ปณฺฑุโลหิตกา ภิกฺขู,}{สยํ ภณฺฑนการกา.} \\ \csromangathabat{ตาทิเส อุปสงฺกมฺม,}{อุสฺสหิํสุ จ ภณฺฑเน.} \\ \csromangathabat{อนุปฺปนฺนาปิ ชายนฺติ\textsuperscript{1},}{อุปฺปนฺนานิปิ วฑฺฒเร\textsuperscript{1}.} \\ \csromangathabat{อปฺปิจฺฉา เปสลา ภิกฺขู,}{อุชฺฌายนฺติ ปทสฺสโต\textsuperscript{1}.} \\ \csromangathabat{สทฺธมฺมฏฺฐิติโก พุทฺโธ,}{สยมฺภู อคฺคปุคฺคโล.} \\ \csromangathabat{อาณาเปสิ ตชฺชนีย,}{กมฺมํ สาวตฺถิยํ ชิโน.} \\ \csromangathabat{อสมฺมุขาปฺปฏิปุจฺฉา,}{ปฺปฏิญฺญาย กตญฺจ ยํ.} \\ \csromangathabat{อนาปตฺติ อเทสเน,}{เทสิตาย กตญฺจ ยํ.} \\ \csromangathabat{อโจเทตฺวา อสาเรตฺวา,}{อนาโรเปตฺวา จ ยํ กตํ.} \\ \csromangathabat{อสมฺมุขา อธมฺเมน,}{วคฺเคน จาปิ\textsuperscript{1} ยํ กตํ.} \\ \csromangathabat{อปฺปฏิปุจฺฉา อธมฺเมน,}{ปุน วคฺเคน\textsuperscript{1} ยํ กตํ.} \\ \csromangathabat{อปฺปฏิญฺญาย อธมฺเมน,}{วคฺเคน จาปิ\textsuperscript{1} ยํ กตํ.} \\ \csromangathabat{อนาปตฺติ\textsuperscript{1} อธมฺเมน,}{วคฺเคน จาปิ\textsuperscript{1} ยํ กตํ.} \\ \csromangathabat{อเทสนาคามินิยา,}{อธมฺเมวคฺคเมว จ.} \\ \csromangathabat{เทสิตาย อธมฺเมน,}{วคฺเคนาปิ ตเถว จ.} \\ \csromangathabat{อโจเทตฺวา อธมฺเมน,}{วคฺเคนาปิ ตเถว จ.} \\ \csromangathabat{อสาเรตฺวา อธมฺเมน,}{วคฺเคนาปิ ตเถว จ.} \\ \csromangathabat{อนาโรเปตฺวา อธมฺเมน,}{วคฺเคนาปิ ตเถว จ.} \\ \csromangathabat{กณฺหวารนเยเนว,}{สุกฺกวารํ วิชานิยา.} \\ \csromangathabat{สํโฆ อากงฺขมาโน จ,}{ยสฺส ตชฺชนียํ กเร.} \\ \csromangathabat{ภณฺฑนํ พาโล สํสฏฺโฐ,}{อธิสีเล อชฺฌาจาเร.} \\ \csromangathabat{อติทิฏฺฐิวิปนฺนสฺส,}{สํโฆ ตชฺชนิยํ กเร.} \\ \csromangathabat{พุทฺธธมฺมสฺส สํฆสฺส,}{อวณฺณํ โย จ ภาสติ.} \\ \csromangathabat{ติณฺณนฺนมฺปิ จ ภิกฺขูนํ,}{สํโฆ ตชฺชนิยํ กเร.} \\ \csromangathabat{ภณฺฑนํ การโก เอโก,}{พาโล สํสคฺคนิสฺสิโต.} \\ \csromangathabat{อธิสีเล อชฺฌาจาเร,}{ตเถว อติทิฏฺฐิยา.} \\ \csromangathabat{พุทฺธธมฺมสฺส สํฆสฺส,}{อวณฺณํ โย จ ภาสติ.} \\ \csromangathabat{ตชฺชนียกมฺมกโต,}{เอวํ สมฺมานุวตฺตนา.} \\ \csromangathabat{อุปสมฺปทนิสฺสยา,}{สามเณรํ อุปฏฺฐนา.} \\ \csromangathabat{โอวาทสมฺมเตนาปิ,}{น กเร ตชฺชนีกโต.} \\ \csromangathabat{นาปชฺเช ตญฺจ อาปตฺติํ,}{ตาทิสญฺจ ตโต ปรํ.} \\ \csromangathabat{กมฺมญฺจ กมฺมิเก จาปิ,}{น ครเห ตถาวิโธ.} \\ \csromangathabat{อุโปสถํ ปวารณํ,}{ปกตตฺตสฺส นฏฺฐเป.} \\ \csromangathabat{สวจนิํ\textsuperscript{1} อนุวาโท,}{โอกาโส โจทเนน จ.} \\ \csromangathabat{สารณํ สมฺปโยคญฺจ,}{น กเรยฺย ตถาวิโธ.} \\ \csromangathabat{อุปสมฺปทนิสฺสยา,}{สามเณรํ อุปฏฺฐนา.} \\ \csromangathabat{โอวาทสมฺมเตนาปิ,}{ปญฺจหงฺเคหิ\textsuperscript{1} น สมฺมติ.} \\ \csromangathabat{ตญฺจาปชฺชติ อาปตฺติํ,}{ตาทิสญฺจ ตโต ปรํ.} \\ \csromangathabat{กมฺมญฺจ กมฺมิเก จาปิ,}{ครหนฺโต น สมฺมติ.} \\ \csromangathabat{อุโปสถํ ปวารณํ,}{สวจนียา จ โนวาโท.} \\ \csromangathabat{โอกาโส โจทนญฺเจว,}{สารณา สมฺปโยชนา.} \\ \csromangathabat{อิเมหฏฺฐงฺเคหิ โย ยุตฺโต,}{ตชฺชนานุปสมฺมติ.} \\ \csromangathabat{กณฺหวารนเยเนว,}{สุกฺกวารํ วิชานิยา.} \\ \csromangathabat{พาโล อาปตฺติพหุโล,}{สํสฏฺโฐปิ จ เสยฺยโส.} \\ \csromangathabat{นิยสฺสกมฺมํ สมฺพุทฺโธ,}{อาณาเปสิ มหามุนิ.} \\ \csromangathabat{กีฏาคิริสฺมิํ เทฺว ภิกฺขู,}{อสฺสชิปุนพฺพสุกา.} \\ \csromangathabat{อนาจารญฺจ วิวิธํ,}{อาจริํสุ อสญฺญตา.} \\ \csromangathabat{ปพฺพาชนียํ สมฺพุทฺโธ,}{กมฺมํ สาวตฺถิยํ ชิโน.} \\ \csromangathabat{มจฺฉิกาสณฺเฑ สุธมฺโม,}{จิตฺตสฺสาวาสิโก อหุ.} \\ \csromangathabat{ชาติวาเทน ขุํเสติ,}{สุธมฺโม จิตฺตุปาสกํ.} \\ \csromangathabat{ปฏิสารณียกมฺมํ,}{อาณาเปสิ ตถาคโต.} \\ \csromangathabat{โกสมฺพิยํ ฉนฺนํ ภิกฺขุํ,}{นิจฺฉนฺตาปตฺติํ ปสฺสิตุํ.} \\ \csromangathabat{อทสฺสเน อุกฺขิปิตุํ,}{อาณาเปสิ ชินุตฺตโม.} \\ \csromangathabat{ฉนฺโน ตํเยว อาปตฺติํ,}{ปฏิกาตุํ น อิจฺฉติ.} \\ \csromangathabat{อุกฺเขปนาปฺปฏิกมฺเม,}{อาณาเปสิ วินายโก.} \\ \csromangathabat{ปาปทิฏฺฐิ อริฏฺฐสฺส,}{อาสิ อญฺญาณนิสฺสิตา.} \\ \csromangathabat{ทิฏฺฐิยาปฺปฏินิสฺสคฺเค\textsuperscript{1},}{อุกฺเขปํ ชินภาสิตํ.} \\ \csromangathabat{นิยสฺสกมฺมํ ปพฺพชฺชํ\textsuperscript{1},}{ตเถว ปฏิสารณี.} \\ \csromangathabat{อทสฺสนาปฺปฏิกมฺเม,}{อนิสฺสคฺเค จ ทิฏฺฐิยา.} \\ \csromangathabat{ทวานาจารูปฆาติ,}{มิจฺฉาอาชีวเมว จ.} \\ \csromangathabat{ปพฺพาชนียกมฺมมฺหิ,}{อติเรกปทา อิเม.} \\ \csromangathabat{อลาภาวณฺณา เทฺว ปญฺจ,}{เทฺว ปญฺจกาติ นามกา\textsuperscript{1}.} \\ \csromangathabat{ปฏิสารณียกมฺมมฺหิ,}{อติเรกปทา อิเม.} \\ \csromangathabat{ตชฺชนียํ นิยสฺสญฺจ,}{ทุเว กมฺมาปิ สาทิสา\textsuperscript{1}.} \\ \csromangathabat{ปพฺพชฺชา\textsuperscript{1} ปฏิสารี จ,}{อตฺถิ ปทาติริตฺตตา.} \\ \csromangathabat{ตโย อุกฺเขปนา กมฺมา,}{สทิสา เต วิภตฺติโต.} \\ \csromangathabat{ตชฺชนียนเยนาปิ,}{เสสกมฺมํ วิชานิยาติ.}}
\csromangathasetleft{ปณฺฑุโลหิตกา ภิกฺขู, \\ ตาทิเส อุปสงฺกมฺม, \\ อนุปฺปนฺนาปิ ชายนฺติ\textsuperscript{1}, \\ อปฺปิจฺฉา เปสลา ภิกฺขู, \\ สทฺธมฺมฏฺฐิติโก พุทฺโธ, \\ อาณาเปสิ ตชฺชนีย, \\ อสมฺมุขาปฺปฏิปุจฺฉา, \\ อนาปตฺติ อเทสเน, \\ อโจเทตฺวา อสาเรตฺวา, \\ อสมฺมุขา อธมฺเมน, \\ อปฺปฏิปุจฺฉา อธมฺเมน, \\ อปฺปฏิญฺญาย อธมฺเมน, \\ อนาปตฺติ\textsuperscript{1} อธมฺเมน, \\ อเทสนาคามินิยา, \\ เทสิตาย อธมฺเมน, \\ อโจเทตฺวา อธมฺเมน, \\ อสาเรตฺวา อธมฺเมน, \\ อนาโรเปตฺวา อธมฺเมน, \\ กณฺหวารนเยเนว, \\ สํโฆ อากงฺขมาโน จ, \\ ภณฺฑนํ พาโล สํสฏฺโฐ, \\ อติทิฏฺฐิวิปนฺนสฺส, \\ พุทฺธธมฺมสฺส สํฆสฺส, \\ ติณฺณนฺนมฺปิ จ ภิกฺขูนํ, \\ ภณฺฑนํ การโก เอโก, \\ อธิสีเล อชฺฌาจาเร, \\ พุทฺธธมฺมสฺส สํฆสฺส, \\ ตชฺชนียกมฺมกโต, \\ อุปสมฺปทนิสฺสยา, \\ โอวาทสมฺมเตนาปิ, \\ นาปชฺเช ตญฺจ อาปตฺติํ, \\ กมฺมญฺจ กมฺมิเก จาปิ, \\ อุโปสถํ ปวารณํ, \\ สวจนิํ\textsuperscript{1} อนุวาโท, \\ สารณํ สมฺปโยคญฺจ, \\ อุปสมฺปทนิสฺสยา, \\ โอวาทสมฺมเตนาปิ, \\ ตญฺจาปชฺชติ อาปตฺติํ, \\ กมฺมญฺจ กมฺมิเก จาปิ, \\ อุโปสถํ ปวารณํ, \\ โอกาโส โจทนญฺเจว, \\ อิเมหฏฺฐงฺเคหิ โย ยุตฺโต, \\ กณฺหวารนเยเนว, \\ พาโล อาปตฺติพหุโล, \\ นิยสฺสกมฺมํ สมฺพุทฺโธ, \\ กีฏาคิริสฺมิํ เทฺว ภิกฺขู, \\ อนาจารญฺจ วิวิธํ, \\ ปพฺพาชนียํ สมฺพุทฺโธ, \\ มจฺฉิกาสณฺเฑ สุธมฺโม, \\ ชาติวาเทน ขุํเสติ, \\ ปฏิสารณียกมฺมํ, \\ โกสมฺพิยํ ฉนฺนํ ภิกฺขุํ, \\ อทสฺสเน อุกฺขิปิตุํ, \\ ฉนฺโน ตํเยว อาปตฺติํ, \\ อุกฺเขปนาปฺปฏิกมฺเม, \\ ปาปทิฏฺฐิ อริฏฺฐสฺส, \\ ทิฏฺฐิยาปฺปฏินิสฺสคฺเค\textsuperscript{1}, \\ นิยสฺสกมฺมํ ปพฺพชฺชํ\textsuperscript{1}, \\ อทสฺสนาปฺปฏิกมฺเม, \\ ทวานาจารูปฆาติ, \\ ปพฺพาชนียกมฺมมฺหิ, \\ อลาภาวณฺณา เทฺว ปญฺจ, \\ ปฏิสารณียกมฺมมฺหิ, \\ ตชฺชนียํ นิยสฺสญฺจ, \\ ปพฺพชฺชา\textsuperscript{1} ปฏิสารี จ, \\ ตโย อุกฺเขปนา กมฺมา, \\ ตชฺชนียนเยนาปิ,}
\csromangathagroupwithclosermajor{\csromangathastanza{\csromanfolio{79}\csromangathabat{ปณฺฑุโลหิตกา ภิกฺขู,}{สยํ ภณฺฑนการกา.} \\ \csromangathabat{ตาทิเส อุปสงฺกมฺม,}{อุสฺสหิํสุ จ ภณฺฑเน.}}\csromangathastanza{\csromangathabat{อนุปฺปนฺนาปิ ชายนฺติ\csromansharedfootnote{6af4fb1483616b64}{อนุปฺปนฺนานิ ชายนฺติ (สี, สฺยา) 2. อุปฺปนฺนานิ ปวฑฺฒเร (สี), อุปฺปนฺนาปิ ปวฑฺฒนฺติ (ก) 3. ปรีสโต (สฺยา), ปรสฺสโต (สี)},}{อุปฺปนฺนานิปิ วฑฺฒเร\csromansharedfootnote{854b7c611d9a86ec}{อุปฺปนฺนานิ ปวฑฺฒเร (สี), อุปฺปนฺนาปิ ปวฑฺฒนฺติ (ก)}.} \\ \csromangathabat{อปฺปิจฺฉา เปสลา ภิกฺขู,}{อุชฺฌายนฺติ ปทสฺสโต\csromansharedfootnote{13c2fcfb310b0016}{ปรีสโต (สฺยา), ปรสฺสโต (สี)}.}}\csromangathastanza{\csromangathabat{สทฺธมฺมฏฺฐิติโก พุทฺโธ,}{สยมฺภู อคฺคปุคฺคโล.} \\ \csromangathabat{อาณาเปสิ ตชฺชนีย,}{กมฺมํ สาวตฺถิยํ ชิโน.}}\csromangathastanza{\csromangathabat{อสมฺมุขาปฺปฏิปุจฺฉา,}{ปฺปฏิญฺญาย กตญฺจ ยํ.} \\ \csromangathabat{อนาปตฺติ อเทสเน,}{เทสิตาย กตญฺจ ยํ.}}\csromangathastanza{\csromangathabat{อโจเทตฺวา อสาเรตฺวา,}{อนาโรเปตฺวา จ ยํ กตํ.} \\ \csromangathabat{อสมฺมุขา อธมฺเมน,}{วคฺเคน จาปิ\csromansharedfootnote{23bbcc3a0d37d488}{วคฺเคนาปิ จ (สี, สฺยา)} ยํ กตํ.}}\csromangathastanza{\csromangathabat{อปฺปฏิปุจฺฉา อธมฺเมน,}{ปุน วคฺเคน\csromansharedfootnote{23bbcc3a0d37d488}{วคฺเคนาปิ จ (สี, สฺยา)} ยํ กตํ.} \\ \csromangathabat{อปฺปฏิญฺญาย อธมฺเมน,}{วคฺเคน จาปิ\csromansharedfootnote{23bbcc3a0d37d488}{วคฺเคนาปิ จ (สี, สฺยา)} ยํ กตํ.}}\csromangathastanza{\csromangathabat{อนาปตฺติ\csromansharedfootnote{1da5a778fafaf3a7}{อนาปตฺติยา (สี, สฺยา)} อธมฺเมน,}{วคฺเคน จาปิ\csromansharedfootnote{23bbcc3a0d37d488}{วคฺเคนาปิ จ (สี, สฺยา)} ยํ กตํ.} \\ \csromangathabat{อเทสนาคามินิยา,}{อธมฺเมวคฺคเมว จ.}}\csromangathastanza{\csromangathabat{เทสิตาย อธมฺเมน,}{วคฺเคนาปิ ตเถว จ.} \\ \csromangathabat{อโจเทตฺวา อธมฺเมน,}{วคฺเคนาปิ ตเถว จ.}}\csromangathastanza{\csromangathabat{อสาเรตฺวา อธมฺเมน,}{วคฺเคนาปิ ตเถว จ.} \\ \csromangathabat{อนาโรเปตฺวา อธมฺเมน,}{วคฺเคนาปิ ตเถว จ.}}\csromangathastanza{\csromangathabat{กณฺหวารนเยเนว,}{สุกฺกวารํ วิชานิยา.} \\ \csromangathabat{สํโฆ อากงฺขมาโน จ,}{ยสฺส ตชฺชนียํ กเร.}}\csromangathastanza{\csromangathabat{ภณฺฑนํ พาโล สํสฏฺโฐ,}{อธิสีเล อชฺฌาจาเร.} \\ \csromangathabat{อติทิฏฺฐิวิปนฺนสฺส,}{สํโฆ ตชฺชนิยํ กเร.}}\csromangathastanza{\csromanfolio{80}\csromangathabat{พุทฺธธมฺมสฺส สํฆสฺส,}{อวณฺณํ โย จ ภาสติ.} \\ \csromangathabat{ติณฺณนฺนมฺปิ จ ภิกฺขูนํ,}{สํโฆ ตชฺชนิยํ กเร.}}\csromangathastanza{\csromangathabat{ภณฺฑนํ การโก เอโก,}{พาโล สํสคฺคนิสฺสิโต.} \\ \csromangathabat{อธิสีเล อชฺฌาจาเร,}{ตเถว อติทิฏฺฐิยา.}}\csromangathastanza{\csromangathabat{พุทฺธธมฺมสฺส สํฆสฺส,}{อวณฺณํ โย จ ภาสติ.} \\ \csromangathabat{ตชฺชนียกมฺมกโต,}{เอวํ สมฺมานุวตฺตนา.}}\csromangathastanza{\csromangathabat{อุปสมฺปทนิสฺสยา,}{สามเณรํ อุปฏฺฐนา.} \\ \csromangathabat{โอวาทสมฺมเตนาปิ,}{น กเร ตชฺชนีกโต.}}\csromangathastanza{\csromangathabat{นาปชฺเช ตญฺจ อาปตฺติํ,}{ตาทิสญฺจ ตโต ปรํ.} \\ \csromangathabat{กมฺมญฺจ กมฺมิเก จาปิ,}{น ครเห ตถาวิโธ.}}\csromangathastanza{\csromangathabat{อุโปสถํ ปวารณํ,}{ปกตตฺตสฺส นฏฺฐเป.} \\ \csromangathabat{สวจนิํ\csromansharedfootnote{fdb0e6d198df21a6}{น สวจนิยํ (สี, สฺยา)} อนุวาโท,}{โอกาโส โจทเนน จ.}}\csromangathastanza{\csromangathabat{สารณํ สมฺปโยคญฺจ,}{น กเรยฺย ตถาวิโธ.} \\ \csromangathabat{อุปสมฺปทนิสฺสยา,}{สามเณรํ อุปฏฺฐนา.}}\csromangathastanza{\csromangathabat{โอวาทสมฺมเตนาปิ,}{ปญฺจหงฺเคหิ\csromansharedfootnote{420d591e1f072186}{ปญฺจองฺโค (ก)} น สมฺมติ.} \\ \csromangathabat{ตญฺจาปชฺชติ อาปตฺติํ,}{ตาทิสญฺจ ตโต ปรํ.}}\csromangathastanza{\csromangathabat{กมฺมญฺจ กมฺมิเก จาปิ,}{ครหนฺโต น สมฺมติ.} \\ \csromangathabat{อุโปสถํ ปวารณํ,}{สวจนียา จ โนวาโท.}}\csromangathastanza{\csromangathabat{โอกาโส โจทนญฺเจว,}{สารณา สมฺปโยชนา.} \\ \csromangathabat{อิเมหฏฺฐงฺเคหิ โย ยุตฺโต,}{ตชฺชนานุปสมฺมติ.}}\csromangathastanza{\csromangathabat{กณฺหวารนเยเนว,}{สุกฺกวารํ วิชานิยา.} \\ \csromangathabat{พาโล อาปตฺติพหุโล,}{สํสฏฺโฐปิ จ เสยฺยโส.}}\csromangathastanza{\csromangathabat{นิยสฺสกมฺมํ สมฺพุทฺโธ,}{อาณาเปสิ มหามุนิ.} \\ \csromangathabat{กีฏาคิริสฺมิํ เทฺว ภิกฺขู,}{อสฺสชิ\-ปุนพฺพสุกา.}}\csromangathastanza{\csromanfolio{81}\csromangathabat{อนาจารญฺจ วิวิธํ,}{อาจริํสุ อสญฺญตา.} \\ \csromangathabat{ปพฺพาชนียํ สมฺพุทฺโธ,}{กมฺมํ สาวตฺถิยํ ชิโน.}}\csromangathastanza{\csromangathabat{มจฺฉิกาสณฺเฑ สุธมฺโม,}{จิตฺตสฺสาวาสิโก อหุ.} \\ \csromangathabat{ชาติวาเทน ขุํเสติ,}{สุธมฺโม จิตฺตุปาสกํ.}}\csromangathastanza{\csromangathabat{ปฏิสารณีย\-กมฺมํ,}{อาณาเปสิ ตถาคโต.} \\ \csromangathabat{โกสมฺพิยํ ฉนฺนํ ภิกฺขุํ,}{นิจฺฉนฺตาปตฺติํ ปสฺสิตุํ.}}\csromangathastanza{\csromangathabat{อทสฺสเน อุกฺขิปิตุํ,}{อาณาเปสิ ชินุตฺตโม.} \\ \csromangathabat{ฉนฺโน ตํเยว อาปตฺติํ,}{ปฏิกาตุํ น อิจฺฉติ.}}\csromangathastanza{\csromangathabat{อุกฺเขปนาปฺปฏิกมฺเม,}{อาณาเปสิ วินายโก.} \\ \csromangathabat{ปาปทิฏฺฐิ อริฏฺฐสฺส,}{อาสิ อญฺญาณนิสฺสิตา.}}\csromangathastanza{\csromangathabat{ทิฏฺฐิยาปฺปฏินิสฺสคฺเค\csromansharedfootnote{92e18d5db39e7031}{ทิฏฺฐิอปฺปฏินิสฺสคฺเค (ก)},}{อุกฺเขปํ ชินภาสิตํ.} \\ \csromangathabat{นิยสฺสกมฺมํ ปพฺพชฺชํ\csromansharedfootnote{92e18d5db39e7031}{ทิฏฺฐิอปฺปฏินิสฺสคฺเค (ก)},}{ตเถว ปฏิสารณี.}}\csromangathastanza{\csromangathabat{อทสฺสนาปฺปฏิกมฺเม,}{อนิสฺสคฺเค จ ทิฏฺฐิยา.} \\ \csromangathabat{ทวานาจารูปฆาติ,}{มิจฺฉาอาชีวเมว จ.}}\csromangathastanza{\csromangathabat{ปพฺพาชนียกมฺมมฺหิ,}{อติเรกปทา อิเม.} \\ \csromangathabat{อลาภาวณฺณา เทฺว ปญฺจ,}{เทฺว ปญฺจกาติ นามกา\csromansharedfootnote{8ffcccbffaaa86ab}{เทฺว ปญฺจโกติ นามโก (ก) 4. กมฺเมสุ สทิสํ (ก)}.}}\csromangathastanza{\csromangathabat{ปฏิสารณียกมฺมมฺหิ,}{อติเรกปทา อิเม.} \\ \csromangathabat{ตชฺชนียํ นิยสฺสญฺจ,}{ทุเว กมฺมาปิ สาทิสา\csromansharedfootnote{08baf8dda02a9dc7}{กมฺเมสุ สทิสํ (ก)}.}}\csromangathastanza{\csromangathabat{ปพฺพชฺชา\csromansharedfootnote{11451cc009b8cb24}{ปพฺพาชา (ก)} ปฏิสารี จ,}{อตฺถิ ปทาติริตฺตตา.} \\ \csromangathabat{ตโย อุกฺเขปนา กมฺมา,}{สทิสา เต วิภตฺติโต.} \\ \csromangathabat{ตชฺชนียนเยนาปิ,}{เสสกมฺมํ วิชานิยาติ.}}}{กมฺม\-กฺขนฺธโก นิฏฺฐิโต.}
\csromanmatikaanchor{mtk.52}
\csromantocmarknopage{chapter}{2. ปาริวาสิกกฺขนฺธก}{mtk.52}
\bookmark[dest={mtk.52},level=0]{2. ปาริวาสิกกฺขนฺธก}
\chapterheadpageread{2. ปาริวาสิก\-กฺขนฺธก}
\csromanmatikaanchor{mtk.53}
\csromantocmarkref{section}{1. ปาริวาสิกวตฺต}{mtk.53}
\bookmark[dest={mtk.53},level=1]{1. ปาริวาสิกวตฺต}
\csromanheader{1}{1. ปาริวาสิก\-วตฺต}
\proseitem{75}{\csromanfolio{82}เตน สมเยน พุทฺโธ ภควา สาวตฺถิยํ วิหรติ เชตวเน อนาถ\-ปิณฺฑิกสฺส อาราเม. เตน โข ปน สมเยน ปาริวาสิกา ภิกฺขู สาทิยนฺติ ปกตตฺตานํ ภิกฺขูนํ ภิวาทนํ ปจฺจุฏฺฐานํ อญฺชลิกมฺมํ สามีจิกมฺมํ อาสนาภิหารํ เสยฺยาภิหารํ ปาโททกํ ปาทปีฐํ ปาทกถลิกํ ปตฺตจีวร\-ปฏิคฺคหณํ นหาเน ปฏฺฐิปริกมฺมํ. เย เต ภิกฺขู อปฺปิจฺฉา ฯเปฯ เต อุชฺฌายนฺติ ขิยฺยนฺติ วิปาเจนฺติ “กถํ หิ นาม ปาริวาสิกา ภิกฺขู สาทิยิสฺสนฺติ ปกตตฺตานํ ภิกฺขูนํ อภิวาทนํ ปจฺจุฏฺฐานํ อญฺชลิกมฺมํ สามีจิกมฺมํ อาสนาภิหารํ เสยฺยาภิหารํ ปาโททกํ ปาทปีฐํ ปาทกถลิกํ ปตฺตจีวร\-ปฏิคฺคหณํ นหาเน ปิฏฺฐิปริกมฺมนฺ”ติ. อถ โข เต ภิกฺขู ภควโต เอตมตฺถํ อาโรเจสุํ.}

\prose{อถ โข ภควา เอตสฺมิํ นิทาเน เอตสฺมิํ ปกรเณ ภิกฺขุสํฆํ สนฺนิปาตาเปตฺวา ภิกฺขู ปฏิปุจฺฉิ “สจฺจํ กิร ภิกฺขเว ปาริวาสิกา ภิกฺขู สาทิยนฺติ ปกตตฺตานํ ภิกฺขูนํ อภิวาทนํ ปจฺจุฏฺฐานํ อญฺชลิกมฺมํ สามีจิกมฺมํ อาสนาภิหารํ เสยฺยาภิหารํ ปาโททกํ ปาทปีฐํ ปาทกถลิกํ ปตฺตจีวร\-ปฏิคฺคหณํ นหาเน ปิฏฺฐิปริกมฺมนฺ”ติ. สจฺจํ ภควาติ. วิครหิ พุทฺโธ ภควา อนนุจฺฉวิกํ ฯเปฯ. กถํ หิ นาม ภิกฺขเว ปาริวาสิกา ภิกฺขู สาทิยิสฺสนฺติ ปกตตฺตานํ ภิกฺขูนํ อภิวาทนํ ปจฺจุฏฺฐานํ อญฺชลิกมฺมํ สามีจิกมฺมํ อาสนาภิหารํ เสยฺยาภิหารํ ปาโททกํ ปาทปีฐํ ปาทกถลิกํ ปตฺตจีวร\-ปฏิคฺคหณํ นหาเน ปิฏฺฐิ\-ปริกมฺมํ. เนตํ ภิกฺขเว อปฺปสนฺนานํ วา ปสาทาย ฯเปฯ วิครหิตฺวา ฯเปฯ ธมฺมิํ กถํ กตฺวา ภิกฺขู อามนฺเตสิ\csromandash{}น ภิกฺขเว ปาริวาสิเกน ภิกฺขุนา สาทิตพฺพํ ปกตตฺตานํ ภิกฺขูนํ อภิวาทนํ ปจฺจุฏฺฐานํ อญฺชลิกมฺมํ สามีจิกมฺมํ อาสนาภิหาโร เสยฺยาภิหาโร ปาโททกํ ปาทปีฐํ ปาทกถลิกํ ปตฺตจีวร\-ปฏิคฺคหณํ นหาเน ปิฏฺฐิ\-ปริกมฺมํ, โย สาทิเยยฺย, อาปตฺติ ทุกฺกฏสฺส.}

\prose{อนชานามิ ภิกฺขเว ปาริวาสิกานํ ภิกฺขูนํ มิถุ ยถาวุฑฺฒํ อภิวาทนํ ปจฺจุฏฺฐานํ อญฺชลิกมฺมํ สามีจิกมฺมํ อาสนาภิหารํ เสยฺยาภิหารํ ปาโททกํ ปาทปีฐํ ปาทกถลิกํ ปตฺตจีวร\-ปฏิคฺคหณํ นหาเน ปิฏฺฐิ\-ปริกมฺมํ.}

\prose{\csromanfolio{83}อนุชานามิ ภิกฺขเว ปาริวาสิกานํ ภิกฺขูนํ ปญฺจ ยถาวุฑฺฒํ อุโปสถํ ปวารณํ วสฺสิกสาฏิกํ โอโณชนํ ภตฺตํ. เตน หิ ภิกฺขเว ปาริวาสิกานํ ภิกฺขูนํ วตฺตํ ปญฺญเปสฺสามิ, ยถา ปาริวาสิเกหิ ภิกฺขูหิ วตฺติตพฺพํ.}

\proseitem{76}{ปาริวาสิเกน ภิกฺขเว ภิกฺขุนา สมฺมา วตฺติตพฺพํ. ตตฺรายํ สมฺมาวตฺตนา\csromandash{}}

\prose{น อุปสมฺ\-ปาเท\-ตพฺพํ, น นิสฺสโย ทาตพฺโพ, น สามเณโร อุปฏฺฐาเปตพฺโพ, น ภิกฺขุโนวาทก\-สมฺมุติ สาทิตพฺพา, สมฺมเตนปิ ภิกฺขุนิโย น โอวทิตพฺพา, ยาย อาปตฺติยา สํเฆน ปริวาโส ทินฺโน โหติ, สา อาปตฺติ น อาปชฺชิตพฺพา, อญฺญา วา ตาทิสิกา, ตโต วา ปาปิฏฺฐตรา, กมฺมํ น ครหิตพฺพํ, กมฺมิกา น ครหิตพฺพา, น ปกตตฺตสฺส ภิกฺขุโน อุโปสโถ ฐเปตพฺโพ, น ปวารณา ฐเปตพฺพา, น สวจนียํ กาตพฺพํ, น อนุวาโท ปฏฺฐเปตพฺโพ, น โอกาโส กาเรตพฺโพ, น โจเทตพฺโพ, น สาเรตพฺโพ, น ภิกฺขูหิ สมฺป\-โยเช\-ตพฺพํ.}

\prose{น ภิกฺขเว ปาริวาสิเกน ภิกฺขุนา ปกตตฺตสฺส ภิกฺขุโน ปุรโต คนฺตพฺพํ, น ปุรโต นิสีทิตพฺพํ, โย โหติ สํฆสฺส อาสนปริยนฺโต เสยฺยาปริยนฺโต วิหาร\-ปริยนฺโต, โส ตสฺส ปทาตพฺโพ, เตน จ โส สาทิตพฺโพ.}

\prose{น ภิกฺขเว ปาริวาสิเกน ภิกฺขุนา ปกตตฺเตน ภิกฺขุนา ปุเรสมเณน วา ปจฺฉาสมเณน วา กุลานิ อุปสงฺกมิตพฺพานิ, น อารญฺญิกงฺคํ สมาทาตพฺพํ\csromansharedfootnote{286dc62f0d9e0319}{สมาทิตพฺพํ (ก)}, น ปิณฺฑปาติ\-กงฺคํ สมาทาตพฺพํ, น จ ตปฺปจฺจยา ปิณฺฑปาโต นีหราเปตพฺโพ “มา มํ ชานิํสู”ติ.}

\prose{ปาริวาสิเกน ภิกฺขเว ภิกฺขุนา อาคนฺตุเกน อาโรเจตพฺพํ, อาคนฺตุกสฺส อาโรเจตพฺพํ, อุโปสเถ อาโรเจตพฺพํ, ปวารณาย อาโรเจตพฺพํ, สเจ คิลาโน โหติ, ทูเตนปิ อาโรเจตพฺพํ.}

\prose{น ภิกฺขเว ปาริวาสิเกน ภิกฺขุนา สภิกฺขุกา อาวาสา อภิกฺขุโก อาวาโส คนฺตพฺโพ อญฺญตฺร ปกตตฺเตน อญฺญตฺร อนฺตรายา.}

\prose{\csromanfolio{84}น ภิกฺขเว ปาริวาสิเกน ภิกฺขุนา สภิกฺขุกา อาวาสา อภิกฺขุโก อนาวาโส คนฺตพฺโพ อญฺญตฺร ปกตตฺเตน อญฺญตฺร อนฺตรายา. น ภิกฺขเว ปาริวาสิเกน ภิกฺขุนา สภิกฺขุกา อาวาสา อภิกฺขุโก อาวาโส วา อนาวาโส วา คนฺตพฺโพ อญฺญตฺร ปกตตฺเตน อญฺญตฺร อนฺตรายา.}

\prose{น ภิกฺขเว ปาริวาสิเกน ภิกฺขุนา สภิกฺขุกา อนาวาสา อภิกฺขุโก อาวาโส คนฺตพฺโพ อญฺญตฺร ปกตตฺเตน อญฺญตฺร อนฺตรายา. น ภิกฺขเว ปาริวาสิเกน ภิกฺขุนา สภิกฺขุกา อนาวาสา อภิกฺขุโก อนาวาโส คนฺตพฺโพ อญฺญตฺร ปกตตฺเตน อญฺญตฺร อนฺตรายา. น ภิกฺขเว ปาริวาสิเกน ภิกฺขุนา สภิกฺขุกา อนาวาสา อภิกฺขุโก อาวาโส วา อนาวาโส วา คนฺตพฺโพ อญฺญตฺร ปกตตฺเตน อญฺญตฺร อนฺตรายา.}

\prose{น ภิกฺขเว ปาริวาสิเกน ภิกฺขุนา สภิกฺขุกา อาวาสา วา อนาวาสา วา อภิกฺขุโก อาวาโส คนฺตพฺโพ อญฺญตฺร ปกตตฺเตน อญฺญตฺร อนฺตรายา. น ภิกฺขเว ปาริวาสิเกน ภิกฺขุนา สภิกฺขุกา อาวาสา วา อนาวาสา วา อภิกฺขุโก อนาวาโส คนฺตพฺโพ อญฺญตฺร ปกตตฺเตน อญฺญตฺร อนฺตรายา. น ภิกฺขเว ปาริวาสิเกน ภิกฺขุนา สภิกฺขุกา อาวาสา วา อนาวาสา วา อภิกฺขุโก อาวาโส วา อนาวาโส วา คนฺตพฺโพ อญฺญตฺร ปกตตฺเตน อญฺญตฺร อนฺตรายา.}

\prose{น ภิกฺขเว ปาริวาสิเกน ภิกฺขุนา สภิกฺขุกา อาวาสา สภิกฺขุโก อาวาโส คนฺตพฺโพ ยตฺถสฺสุ ภิกฺขู นานาสํวาสกา อญฺญตฺร ปกตตฺเตน อญฺญตฺร อนฺตรายา. น ภิกฺขเว ปาริวาสิเกน ภิกฺขุนา สภิกฺขุกา อาวาสา สภิกฺขุโก อนาวาโส คนฺตพฺโพ ยตฺถสฺสุ ภิกฺขู นานาสํวาสกา อญฺญตฺร ปกตตฺเตน อญฺญตฺร อนฺตรายา. น ภิกฺขเว ปาริวาสิเกน ภิกฺขุนา สภิกฺขุกา อาวาสา สภิกฺขุโก อาวาโส วา อนาวาโส วา คนฺตพฺโพ ยตฺถสฺสุ ภิกฺขู นานาสํวาสกา อญฺญตฺร ปกตตฺเตน อญฺญตฺร อนฺตรายา.}

\prose{น ภิกฺขเว ปาริวาสิเกน ภิกฺขุนา สภิกฺขุกา อนาวาสา สภิกฺขุโก อาวาโส คนฺตพฺโพ ยตฺถสฺสุ ภิกฺขู นานาสํวาสกา \csromanfolio{85}อญฺญตฺร ปกตตฺเตน อญฺญตฺร อนฺตรายา. น ภิกฺขเว ปาริวาสิเกน ภิกฺขุนา สภิกฺขุกา อนาวาสา สภิกฺขุโก อนาวาโส คนฺตพฺโพ ยตฺถสฺสุ ภิกฺขู นานาสํวาสกา อญฺญตฺร ปกตตฺเตน อญฺญตฺร อนฺตรายา. น ภิกฺขเว ปาริวาสิเกน ภิกฺขุนา สภิกฺขุกา อนาวาสา สภิกฺขุโก อาวาโส วา อนาวาโส วา คนฺตพฺโพ ยตฺถสฺสุ ภิกฺขู นานาสํวาสกา อญฺญตฺร ปกตตฺเตน อญฺญตฺร อนฺตรายา.}

\prose{น ภิกฺขเว ปาริวาสิเกน ภิกฺขุนา สภิกฺขุกา อาวาสา วา อนาวาสา วา สภิกฺขุโก อาวาโส คนฺตพฺโพ ยตฺถสฺสุ ภิกฺขู นานาสํวาสกา อญฺญตฺร ปกตตฺเตน อญฺญตฺร อนฺตรายา. น ภิกฺขเว ปาริวาสิเกน ภิกฺขุนา สภิกฺขุกา อาวาสา วา อนาวาสา วา สภิกฺขุโก อนาวาโส คนฺตพฺโพ ยตฺถสฺสุ ภิกฺขู นานาสํวาสกา อญฺญตฺร ปกตตฺเตน อญฺญตฺร อนฺตรายา. น ภิกฺขเว ปาริวาสิเกน ภิกฺขุนา สภิกฺขุกา อาวาสา วา อนาวาสา วา สภิกฺขุโก อาวาโส วา อนาวาโส วา คนฺตพฺโพ ยตฺถสฺสุ ภิกฺขู นานาสํวาสกา อญฺญตฺร ปกตตฺเตน อญฺญตฺร อนฺตรายา.}

\proseitem{80}{คนฺตพฺโพ ภิกฺขเว ปาริวาสิเกน ภิกฺขุนา สภิกฺขุกา อาวาสา สภิกฺขุโก อาวาโส ยตฺถสฺสุ ภิกฺขู สมาน\-สํวาสกา ยํ ชญฺญา “สกฺโกมิ อชฺเชว คนฺตุนฺ”ติ.}

\prose{คนฺตพฺโพ ภิกฺขเว ปาริวาสิเกน ภิกฺขุนา สภิกฺขุกา อาวาสา สภิกฺขุโก อนาวาโส ยตฺถสฺสุ ภิกฺขู สมาน\-สํวาสกา ยํ ชญฺญา “สกฺโกมิ อชฺเชว คนฺตุนฺ”ติ.}

\prose{คนฺตพฺโพ ภิกฺขเว ปาริวาสิเกน ภิกฺขุนา สภิกฺขุกา อาวาสา สภิกฺขุโก อาวาโส วา อนาวาโส วา ยตฺถสฺสุ ภิกฺขู สมาน\-สํวาสกา ยํ ชญฺญา “สกฺโกมิ อชฺเชว คนฺตุนฺ”ติ.}

\prose{คนฺตพฺโพ ภิกฺขเว ปาริวาสิเกน ภิกฺขุนา สภิกฺขุกา อนาวาสา สภิกฺขุโก อาวาโส ยตฺถสฺสุ ภิกฺขู สมาน\-สํวาสกา ยํ ชญฺญา “สกฺโกมิ อชฺเชว คนฺตุนฺ”ติ.}

\prose{\csromanfolio{86}คนฺตพฺโพ ภิกฺขเว ปาริวาสิเกน ภิกฺขุนา สภิกฺขุกา อนาวาสา สภิกฺขุโก อนาวาโส ยตฺถสฺสุ ภิกฺขู สมาน\-สํวาสกา ยํ ชญฺญา “สกฺโกมิ อชฺเชว คนฺตุนฺ”ติ.}

\prose{คนฺตพฺโพ ภิกฺขเว ปาริวาสิเกน ภิกฺขุนา สภิกฺขุกา อนาวาสา สภิกฺขุโก อาวาโส วา อนาวาโส วา ยตฺถสฺสุ ภิกฺขู สมาน\-สํวาสกา ยํ ชญฺญา “สกฺโกมิ อชฺเชว คนฺตุนฺ”ติ.}

\prose{คนฺตพฺโพ ภิกฺขเว ปาริวาสิเกน ภิกฺขุนา สภิกฺขุกา อาวาสา วา อนาวาสา วา สภิกฺขุโก อาวาโส ยตฺถสฺสุ ภิกฺขู สมาน\-สํวาสกา ยํ ชญฺญา “สกฺโกมิ อชฺเชว คนฺตุนฺ”ติ.}

\prose{คนฺตพฺโพ ภิกฺขเว ปาริวาสิเกน ภิกฺขุนา สภิกฺขุกา อาวาสา วา อนาวาสา วา สภิกฺขุโก อนาวาโส ยตฺถสฺสุ ภิกฺขู สมาน\-สํวาสกา ยํ ชญฺญา “สกฺโกมิ อชฺเชว คนฺตุนฺ”ติ.}

\prose{คนฺตพฺโพ ภิกฺขเว ปาริวาสิเกน ภิกฺขุนา สภิกฺขุกา อาวาสา วา อนาวาสา วา สภิกฺขุโก อาวาโส วา อนาวาโส วา ยตฺถสฺสุ ภิกฺขู สมาน\-สํวาสกา ยํ ชญฺญา “สกฺโกมิ อชฺเชว คนฺตุนฺ”ติ.}

\proseitem{81}{น ภิกฺขเว ปาริวาสิเกน ภิกฺขุนา ปกตตฺเตน ภิกฺขุนา สทฺธิํ เอกจฺฉนฺเน อาวาเส วตฺตพฺพํ, น เอกจฺฉนฺเน อนาวาเส วตฺถพฺพํ, น เอกจฺฉนฺเน อาวาเส วา อนาวาเส วา วตฺถพฺพํ, ปกตตฺตํ ภิกฺขุํ ทิสฺวา อาสนา วุฏฺฐาตพฺพํ, ปกตตฺโต ภิกฺขุ อาสเนน นิมนฺเต ตพฺโพ, น ปกตตฺเตน ภิกฺขุนา สทฺธิํ เอกาสเน นิสีทิตพฺพํ, น นีเจ อาสเน นิสินฺเน อุจฺเจ อาสเน นิสีทิตพฺพํ, น ฉมายํ นิสินฺเน อาสมน นิสีทิตพฺพํ, น เอกจงฺกเม จงฺกมิตพฺพํ, น นีเจ จงฺกเม จงฺกมนฺเต อุจฺเจ จงฺกเม จงฺกมิตพฺพํ, น ฉมายํ จงฺกมนฺเต\csromansharedfootnote{316440f1e7c0e162}{จงฺกมนฺตํ (อฏฺฐกถายํ สํวณฺเณตพฺพปาโฐ)} จงฺเกเม จงฺกมิตพฺพํ.}

\proseitem{82}{น ภิกฺขเว ปาริวาสิเกน ภิกฺขุนา ปาริวาสิเกน วุฑฺฒตเรน ภิกฺขุนา สทฺธิํ ฯเปฯ มูลาย\-ปฏิกสฺสนา\-รเหน ภิกฺขุนา สทฺธิํ ฯเปฯ มานตฺตารเหน ภิกฺขุนา สทฺธิํ ฯเปฯ มานตฺต\-จาริเกน ภิกฺขุนา สทฺธิํ ฯเปฯ อพฺภานารเหน ภิกฺขุนา สทฺธิํ เอกจฺฉนฺเน อาวาเส วตฺถพฺพํ, น เอกจฺฉนฺเน อนาวาเส วตฺถพฺพํ, น เอกจฺฉนฺเน อาวาเส วา อนาวาเส วา วตฺถพฺพํ, น เอกาสเน \csromanfolio{87}นิสีทิตพฺพํ, น นีเจ อาสเน นิสินฺเน อุจฺเจ อาสเน นิสีทิตพฺพํ, น ฉมายํ นิสินฺเน อาสเน นิสีทิตพฺพํ, น เอกจงฺกเม จงฺกมิตพฺพํ, น นีเจ จงฺกเม จงฺกมนฺเต อุจฺเจ จงฺกเม จงฺกมิตพฺพํ, น ฉมายํ จงฺกมนฺเต จงฺกเม จงฺกมิตพฺพํ.}

\prosewithcloser{\csromansymbolfootnote{*}{วิ 3. 442 ปิฏฺเฐ.}ปาริวาสิก\-จตุตฺโถ เจ ภิกฺขเว ปริวาสํ ทเทยฺย, มูลาย ปฏิกสฺเสยฺย, มานตฺตํ ทเทยฺย, ตํวีโส อพฺเภยฺย, อกมฺมํ\csromansharedfootnote{1ca5446e6819855d}{อกมฺมํ ตํ (สฺยา)} น จ กรณียนฺติ.}{จตุนฺนวุติ\-ปาริวาสิก\-วตฺตํ นิฏฺฐิตํ.}
\proseitem{83}{อถ โข อายสฺมา อุปาลิ เยน ภควา เตนุปสงฺกมิ, อุปสงฺกมิตฺวา ภควนฺตํ อภิวาเทตฺวา เอกมนฺตํ นิสีทิ, เอกมนฺตํ นิสินฺโน โข อายสฺมา อุปาลิ ภควนฺตํ เอตทโวจ “กติ นุ โข ภนฺเต ปาริวาสิกสฺส ภิกฺขุโน รตฺติจฺเฉทา”ติ. ตโย โข อุปาลิ ปาริวาสิกสฺส ภิกฺขุโน รตฺติจฺเฉทา, สหวาโส วิปฺปวาโส อนาโรจนา, อิเม โข อุปาลิ ตโย ปาริวาสิกสฺส ภิกฺขุโน รตฺติจฺเฉทาติ.}

\proseitem{84}{เตน โข ปน สมเยน สาวตฺถิยํ มหา\-ภิกฺขุสํโฆ สนฺนิปติโต โหติ, น สกฺโกนฺติ ปาริวาสิกา ภิกฺขู ปริวาสํ โสเธตุํ, ภควโต\csromansharedfootnote{6e7c9ce7663fdcbd}{เต ภิกฺขู ภควโต (สฺยา, เอวมุปริปิ)} เอตมตฺถํ อาโรเจสุํ. อนุชานามิ ภิกฺขเว ปริวาสํ นิกฺขิปิตุํ. เอวญฺจ ปน ภิกฺขเว นิกฺขิปิตพฺโพ, เตน ปาริวาสิเกน ภิกฺขุนา เอกํ ภิกฺขุํ อุปสงฺกมิตฺวา เอกํสํ อุตฺตราสงฺคํ กริตฺวา อุกฺกุฏิกํ นิสีทิตฺวา อญฺชลิํ ปคฺคเหตฺวา เอวมสฺส วจนีโย “ปริวาสํ นิกฺขิปามี”ติ, นิกฺขิตฺโต โหติ ปริวาโส. “วตฺตํ นิกฺขิปามี”ติ, นิกฺขิตฺโต โหติ ปริวาโส.}

\proseitemwithcloser{85}{เตน โข ปน สมเยน สาวตฺถิยา ภิกฺขู ตหํ ตหํ ปกฺกมิํสุ, สกฺโกนฺติ ปาริวาสิกา ภิกฺขู ปริวาสํ โสเธสุํ, ภควโต เอตมตฺถํ อาโรเจสุํ. อนุชานามิ ภิกฺขเว ปริวาสํ สมาทิยิตุํ\csromansharedfootnote{c10f7a693a3a7182}{สมาทาตุํ (สฺยา, กํ)}. เอวญฺจ ปน ภิกฺขเว สมาทิยิตพฺโพ\csromansharedfootnote{6c42faade4de08d7}{สมาทิตพฺโพ (สี, สฺยา, กํ)}, เตน ปาริวาสิเกน ภิกฺขุนา เอกํ ภิกฺขุํ \csromanfolio{88}อุปสงฺกมิตฺวา เอกํสํ อุตฺตราสงฺคํ กริตฺวา อุกฺกุฏิกํ นิสีทิตฺวา อญฺชลิํ ปคฺคเหตฺวา เอวมสฺส วจนีโย “ปริวาสํ สมาทิยามี”ติ, สมาทินฺโน โหติ ปริวาโส. “วตฺตํ สมาทิยามี”ติ, สมาทินฺโน โหติ ปริวาโส.}{ปาริวาสิก\-วตฺตํ นิฏฺฐิตํ.}
\csromanmatikaanchor{mtk.54}
\csromantocmarkref{section}{2. มูลายปฏิกสฺสนารหวตฺต}{mtk.54}
\bookmark[dest={mtk.54},level=1]{2. มูลายปฏิกสฺสนารหวตฺต}
\csromanheader{1}{2. มูลายปฏิกสฺสนารห\-วตฺต}
\proseitem{86}{เตน โข ปน สมเยน มูลาย\-ปฏิกสฺสนา\-รหา ภิกฺขู สาทิยนฺติ ปกตตฺตานํ ภิกฺขูนํ อภิวาทนํ ปจฺจุฏฺฐานํ อญฺชลิกมฺมํ สามีจิกมฺมํ อาสนาภิหารํ เสยฺยาภิหารํ ปาโททกํ ปาทปีฐํ ปาทกถลิกํ ปตฺตจีวร\-ปฏิคฺคหณํ นหาเน ปิฏฺฐิ\-ปริกมฺมํ. เย เต ภิกฺขู อปฺปิจฺฉา ฯเปฯ เต อุชฺฌายนฺติ ขิยฺยนฺติ วิปาเจนฺติ “กถํ หิ นาม มูลาย\-ปฏิกสฺสนา\-รหา ภิกฺขู สาทิยิสฺสนฺติ ปกตตฺตานํ ภิกฺขูนํ อภิวาทนํ ปจฺจุฏฺฐานํ ฯเปฯ นหาเน ปิฏฺฐิปริกมฺมนฺ”ติ. อถ โข เต ภิกฺขู ภควโต เอตมตฺถํ อาโรเจสุํ.}

\prose{อถ โข ภควา เอตสฺมิํ นิทาเน เอตสฺมิํ ปกรเณ ภิกฺขุสํฆํ สนฺนิปาตาเปตฺวา ภิกฺขู ปฏิปุจฺฉิ “สจฺจํ กิร ภิกฺขเว มูลาย\-ปฏิกสฺสนา\-รหา ภิกฺขู สาทิยนฺติ ปกตตฺตานํ ภิกฺขูนํ อภิวาทนํ ปจฺจุฏฺฐานํ ฯเปฯ นหาเน ปิฏฺฐิปริกมฺมนฺ”ติ. สจฺจํ ภควาติ. วิครหิ พุทฺโธ ภควา อนนุจฺฉวิกํ ฯเปฯ. กถํ หิ นาม ภิกฺขเว มูลาย\-ปฏิกสฺสนา\-รหา ภิกฺขู สาทิยิสฺสนฺติ ปกตตฺตานํ ภิกฺขูนํ อภิวาทนํ ปจฺจุฏฺฐานํ ฯเปฯ นหาเน ปิฏฺฐิ\-ปริกมฺมํ. เนตํ ภิกฺขเว อปฺปสนฺนานํ วา ปสฺสาทาย ฯเปฯ วิครหิตฺวา ฯเปฯ ธมฺมิํ กถํ กตฺวา ภิกฺขู อามนฺเตสิ\csromandash{}น ภิกฺขเว มูลาย\-ปฏิกสฺสนา\-รเหน ภิกฺขุนา สาทิตพฺพํ ปกตตฺตานํ ภิกฺขูนํ อภิวาทนํ ปจฺจุฏฺฐานํ อญฺชลิกมฺมํ สามีจิกมฺมํ อาสนาภิหาโร เสยฺยาภิหาโร ปาโททกํ ปาทปีฐํ ปาทกถลิกํ ปตฺตจีวร\-ปฏิคฺคหณํ นหาเน ปิฏฺฐิ\-ปริกมฺมํ, โย สาทิเยยฺย, อาปตฺติ ทุกฺกฏสฺส. อนุชานามิ ภิกฺขเว มูลาย\-ปฏิกสฺสนา\-รหานํ ภิกฺขูนํ มิถุ ยถาวุฑฺฒํ อภิวาทนํ ปจฺจุฏฺฐานํ ฯเปฯ นหาเน ปิฏฺฐิ\-ปริกมฺมํ. อนุชานามิ ภิกฺขเว มูลาย\-ปฏิกสฺสนา\-รหานํ ภิกฺขูนํ ปญฺจ ยถาวุฑฺฒํ \csromanfolio{89}อุโปสถํ ปวารณํ วสฺสิกสาฏิกํ โอโณชนํ ภตฺตํ. เตน หิ ภิกฺขเว มูลาย\-ปฏิกสฺสนา\-รหานํ ภิกฺขูนํ วตฺตํ ปญฺญเปสฺสามิ, ยถา มูลาย\-ปฏิกสฺสนา\-รเหหิ ภิกฺขูหิ วตฺติตพฺพํ.}

\proseitem{87}{มูลาย\-ปฏิกสฺสนา\-รเหน ภิกฺขเว ภิกฺขุนา สมฺมา วตฺติตพฺพํ. ตตฺรายํ สมฺมาวตฺตนา\csromandash{}}

\prose{น อุปสมฺ\-ปาเท\-ตพฺพํ, น นิสฺสโย ทาตพฺโพ, น สามเณโร อุปฏฺฐาเปตพฺโพ, น ภิกฺขุโนวาทก\-สมฺมุติ สาทิตพฺพา, สมฺมเตนปิ ภิกฺขุนิโย น โอวทิตพฺพา, ยาย อาปตฺติยา สํเฆน มูลาย ปฏิกสฺสนา\-รโห กโต โหติ, สา อาปตฺติ น อาปชฺชิตพฺพา, อญฺญา วา ตาทิสิกา, ตโต วา ปาปิฏฺฐตรา, กมฺมํ น ครหิตพฺพํ, กมฺมิกา น ครหิตพฺพา, น ปกตตฺตสฺส ภิกฺขุโน อุโปสโถ ฐเปตพฺโพ, น ปวารณา ฐเปตพฺพา, น สวจนียํ กาตพฺพํ, น อนุวาโท ปฏฺฐเปตพฺโพ, น โอกาโส กาเรตพฺโพ, น โจเทตพฺโพ, น สาเรตพฺโพ, น ภิกฺขูหิ สมฺป\-โยเช\-ตพฺพํ.}

\prose{น ภิกฺขเว มูลาย\-ปฏิกสฺสนา\-รเหน ภิกฺขุนา ปกตตฺตสฺส ภิกฺขุโน ปุรโต คนฺตพฺพํ, น ปุรโต นิสีทิตพฺพํ, โย โหติ สํฆสฺส อาสนปริยนฺโต เสยฺยาปริยนฺโต วิหาร\-ปริยนฺโต, โส ตสฺส ปทาตพฺโพ, เตน จ โส สาทิตพฺโพ.}

\prose{น ภิกฺขเว มูลาย\-ปฏิกสฺสนา\-รเหน ภิกฺขุนา ปกตตฺเตน ภิกฺขุนา ปุเรสมเณน วา ปจฺฉาสมเณน วา กุลานิ อุปสงฺกมิตพฺพานิ, น อารญฺญิกงฺคํ สมาทาตพฺพํ, น ปิณฺฑปาติ\-กงฺคํ สมาทาตพฺพํ, น จ ตปฺปจฺจยา ปิณฺฑปาโต นีหราเปตพฺโพ “มา มํ ชานิํสู”ติ.}

\prose{น ภิกฺขเว มูลาย\-ปฏิกสฺสนา\-รเหน ภิกฺขุนา สภิกฺขุกา อาวาสา อภิกฺขุโก อาวาโส คนฺตพฺโพ อญฺญตฺร ปกตตฺเตน อญฺญตฺร อนฺตรายา.}

\prose{น ภิกฺขเว มูลาย\-ปฏิกสฺสนา\-รเหน ภิกฺขุนา สภิกฺขุกา อาวาสา อภิกฺขุโก อนาวาโส คนฺตพฺโพ อญฺญตฺร ปกตตฺเตน อญฺญตฺร อนฺตรายา.}

\prose{\csromanfolio{90}น ภิกฺขเว มูลาย\-ปฏิกสฺสนา\-รเหน ภิกฺขุนา สภิกฺขุกา อาวาสา อภิกฺขุโก อาวาโส วา อนาวาโส วา คนฺตพฺโพ อญฺญตฺร ปกตตฺเตน อญฺญตฺร อนฺตรายา.}

\prose{น ภิกฺขเว มูลาย\-ปฏิกสฺสนา\-รเหน ภิกฺขุนา สภิกฺขุกา อนาวาสา อภิกฺขุโก อาวาโส คนฺตพฺโพ ฯเปฯ อภิกฺขุโก อนาวาโส คนฺตพฺโพ ฯเปฯ อภิกฺขุโก อาวาโส วา อนาวาโส วา คนฺตพฺโพ อญฺญตฺร ปกตฺเตน อญฺญตฺร อนฺตรายา.}

\prose{น ภิกฺขเว มูลาย\-ปฏิกสฺสนา\-รเหน ภิกฺขุนา สภิกฺขุกา อาวาสา วา อนาวาสา วา อภิกฺขุโก อาวาโส คนฺตพฺโพ ฯเปฯ อภิกฺขุโก อนาวาโส คนฺตพฺโพ ฯเปฯ อภิกฺขุโก อาวาโส วา อนาวาโส วา คนฺตพฺโพ อญฺญตฺร ปกตตฺเตน อญฺญตฺร อนฺตรายา.}

\prose{น ภิกฺขเว มูลาย\-ปฏิกสฺสนา\-รเหน ภิกฺขุนา สภิกฺขุกา อาวาสา สภิกฺขุโก อาวาโส คนฺตพฺโพ ฯเปฯ สภิกฺขุโก อนาวาโส คนฺตพฺโพ ฯเปฯ สภิกฺขุโก อาวาโส วา อนาวาโส วา คนฺตพฺโพ ยตฺถสฺสุ ภิกฺขู นานาสํวาสกา อญฺญตฺร ปกตตฺเตน อญฺญตฺร อนฺตรายา.}

\prose{น ภิกฺขเว มูลาย\-ปฏิกสฺสนา\-รเหน ภิกฺขุนา สภิกฺขุกา อานาวาสา สภิกฺขุโก อาวาโส คนฺตพฺโพ ฯเปฯ สภิกฺขุโก อนาวาโส คนฺตพฺโพ ฯเปฯ สภิกฺขุโก อาวาโส วา อนาวาโส วา คนฺตพฺโพ ยตฺถสฺสุ ภิกฺขู นานาสํวาสกา อญฺญตฺร ปกตตฺเตน อญฺญตฺร อนฺตรายา.}

\prose{น ภิกฺขเว มูลาย\-ปฏิกสฺสนา\-รเหน ภิกฺขุนา สภิกฺขุกา อาวาสา วา อนาวาสา วา สภิกฺขุโก อาวาโส คนฺตพฺโพ ฯเปฯ สภิกฺขุโก อนาวาโส คนฺตพฺโพ ฯเปฯ สภิกฺขุโก อาวาโส วา อนาวาโส วา คนฺตพฺโพ ยตฺถสฺสุ ภิกฺขู นานาสํวาสกา อญฺญตฺร ปกตตฺเตน อญฺญตฺร อนฺตรายา.}

\prose{คนฺตพฺโพ ภิกฺขเว มูลาย\-ปฏิกสฺสนา\-รเหน ภิกฺขุนา สภิกฺขุกา อาวาสา สภิกฺขุโก อาวาโส ฯเปฯ สภิกฺขุโก อนาวาโส ฯเปฯ สภิกฺขุโก อาวาโส วา อนาวาโส วา ยตฺถสฺสุ ภิกฺขู สมาน\-สํวาสกา ยํ ชญฺญา “สกฺโกมิ อชฺเชว คนฺตุนฺ”ติ.}

\prose{\csromanfolio{91}คนฺตพฺโพ ภิกฺขเว มูลาย\-ปฏิกสฺสนา\-รเหน ภิกฺขุนา สภิกฺขุกา อนาวาสา สภิกฺขุโก อาวาโส ฯเปฯ สภิกฺขุโก อนาวาโส ฯเปฯ สภิกฺขุโก อาวาโส วา อนาวาโส วา ยตฺถสฺสุ ภิกฺขู สมาน\-สํวาสกา ยํ ชญฺญา “สกฺโกมิ อชฺเชว คนฺตุนฺ”ติ.}

\prose{คนฺตพฺโพ ภิกฺขเว มูลาย\-ปฏิกสฺสนา\-รเหน ภิกฺขุนา สภิกฺขุกา อาวาสา วา อนาวาสา วา สภิกฺขุโก อาวาโส ฯเปฯ สภิกฺขุโก อนาวาโส ฯเปฯ สภิกฺขุโก อาวาโส วา อนาวาโส วา ยตฺถสฺสุ ภิกฺขู สมาน\-สํวาสกา ยํ ชญฺญา “สกฺโกมิ อชฺเชว คนฺตุนฺ”ติ.}

\prose{น ภิกฺขเว มูลาย\-ปฏิกสฺสนา\-รเหน ภิกฺขุนา ปกตตฺเตน ภิกฺขุนา สทฺธิํ เอกจฺฉนฺเน อาวาเส วตฺถพฺพํ, น เอกจฺฉนฺเน อนาวาเส วตฺถพฺพํ, น เอกจฺฉนฺเน อาวาเส วา อนาวาเส วา วตฺถพฺพํ, ปกตตฺตํ ภิกฺขุํ ทิสฺวา อาสนา วุฏฺฐาตพฺพํ, ปกตตฺโต ภิกฺขุ อาสเนน นิมนฺเตตพฺโพ, น ปกตตฺเตน ภิกฺขุนา สทฺธิํ เอกาสเน นิสีทิตพฺพํ, น นีเจ อาสเน นิสินฺเน อุจฺเจ อาสเน นิสีทิตพฺพํ, น ฉมายํ นิสินฺเน อาสเน นิสีทิตพฺพํ, น เอกจงฺกเม จงฺกมิตพฺพํ, น นีเจ จงฺกเม จงฺกมนฺเต อุจฺเจ จงฺกเม จงฺกมิตพฺพํ, น ฉมายํ จงฺกมนฺเต จงฺกเม จงฺกมิตพฺพํ.}

\prose{น ภิกฺขเว มูลาย\-ปฏิกสฺสนา\-รเหน ภิกฺขุนา ปาริวาสิเกน ภิกฺขุนา สทฺธิํ ฯเปฯ มูลาย\-ปฏิกสฺสนา\-รเหน วุฑฺฒตเรน ภิกฺขุนา สทฺธิํ ฯเปฯ มานตฺตารเหน ภิกฺขุนา สทฺธิํ ฯเปฯ มานตฺต\-จาริเกน ภิกฺขุนา สทฺธิํ ฯเปฯ อพฺภานารเหน ภิกฺขุนา สทฺธิํ เอกจฺฉนฺเน อาวาเส วตฺถพฺพํ, น เอกจฺเฉนฺเน อนาวาเส วตฺตพฺพํ, น เอกจฺฉนฺเน อาวาเส วา อนาวาเส วา วตฺถพฺพํ, น เอกาสเน นิสีทิตพฺพํ, น นีเจ อาสเน นิสินฺเน อุจฺเจ อาสเน นิสีทิตพฺพํ, น ฉมายํ นิสินฺเน อาสเน นิสีทิตพฺพํ, น เอกจงฺกเม จงฺกมิตพฺพํ, น นีเจ จงฺกเม จงฺกมนฺเต อุจฺเจ จงฺกเม จงฺกมิตพฺพํ, น ฉมายํ จงฺกมนฺเต จงฺกเม จงฺกมิตพฺพํ.}

\prosewithcloser{\csromansymbolfootnote{*}{วิ 3. 442 ปิฏฺเฐ.}มูลายปฏิกสฺสนารห\-จตุตฺโถ เจ ภิกฺขเว ปริวาสํ ทเทยฺย, มูลาย ปฏิกสฺเสยฺย, มานตฺตํ ทเทยฺย, ตํวีโส อพฺเภยฺย, อกมฺมํ น จ กรณียนฺติ.}{มูลายปฏิกสฺสนารห\-วตฺตํ นิฏฺฐิตํ.}
\csromanmatikaanchor{mtk.55}
\csromantocmarkref{section}{3. มานตฺตารหวตฺต}{mtk.55}
\bookmark[dest={mtk.55},level=1]{3. มานตฺตารหวตฺต}
\csromanheader{1}{3. มานตฺตารห\-วตฺต}
\proseitem{88}{\csromanfolio{92}เตน โข ปน สมเยน มานตฺตารหา ภิกฺขู สาทิยนฺติ ปกตตฺตานํ ภิกฺขูนํ อภิวาทนํ ปจฺจุฏฺฐานํ อญฺชลิกมฺมํ สามีจิกมฺมํ อาสนาภิหารํ เสยฺยาภิหารํ ปาโททกํ ปาทปีฐํ ปาทกถลิกํ ปตฺตจีวร\-ปฏิคฺคหณํ นหาเน ปิฏฺฐิ\-ปริกมฺมํ. เย เต ภิกฺขู อปฺปิจฺฉา ฯเปฯ เต อุชฺฌายนฺติ ขิยฺยนฺติ วิปาเจนฺติ “กถํ หิ นาม มานตฺตารหา ภิกฺขู สาทิยิสฺสนฺติ ปกตตฺตานํ ภิกฺขูนํ อภิวาทนํ ปจฺจุฏฺฐานํ อญฺชลิกมฺมํ สามีจิกมฺมํ อาสนาภิหารํ เสยฺยาภิหารํ ปาโททกํ ปาทปีฐํ ปาทกถลิกํ ปตฺตจีวร\-ปฏิคฺคหณํ นหาเน ปิฏฺฐิปริกมฺมนฺ”ติ. อถ โข เต ภิกฺขู ภควโต เอตมตฺถํ อาโรเจสุํ ฯเปฯ. สจฺจํ กิร ภิกฺขเว มานตฺตารหา ภิกฺขู สาทิยนฺติ ปกตตฺตานํ ภิกฺขูนํ อภิวาทนํ ปจฺจุฏฺฐานํ ฯเปฯ นหาเน ปิฏฺฐิ\-ปริกมฺมนฺติ. สจฺจํ ภควาติ. วิครหิ พุทฺโธ ภควา อนนุจฺฉวิกํ ฯเปฯ. กถํ หิ นาม ภิกฺขเว มานตฺตารหา ภิกฺขู สาทิยิสฺสนฺติ ปกตตฺตานํ ภิกฺขูนํ อภิวาทนํ ฯเปฯ นหาเน ปิฏฺฐิ\-ปริกมฺมํ. เนตํ ภิกฺขเว อปฺปสนฺนานํ วา ปสาทาย ฯเปฯ วิครหิตฺวา ฯเปฯ ธมฺมิํ กถํ กตฺวา ภิกฺขู อามนฺเตสิ\csromandash{}น ภิกฺขเว มานตฺตารเหน ภิกฺขุนา สาทิตพฺพํ ปกตตฺตานํ ภิกฺขูนํ อภิวาทนํ ปจฺจุฏฺฐานํ ฯเปฯ นหาเน ปิฏฺฐิ\-ปริกมฺมํ, โย สาทิเยยฺย, อาปตฺติ ทุกฺกฏสฺส. อนุชานามิ ภิกฺขเว มานตฺตารหานํ ภิกฺขูนํ มิถุ ยถาวุฑฺฒํ อภิวาทนํ ปจฺจุฏฺฐานํ ฯเปฯ นหาเน ปิฏฺฐิ\-ปริกมฺมํ. อนุชานามิ ภิกฺขเว มานตฺตารหานํ ภิกฺขูนํ ปญฺจ ยถาวุฑฺฒํ อุโปสถํ ปวารณํ วสฺสิกสาฏิกํ โอโณชนํ ภตฺตํ. เตน หิ ภิกฺขเว มานตฺตารหานํ ภิกฺขูนํ วตฺตํ ปญฺญเปสฺสามิ, ยถา มานตฺตารเหหิ ภิกฺขูหิ วตฺติตพฺพํ.}

\proseitem{89}{มานตฺตารเหน ภิกฺขเว ภิกฺขุนา สมฺมา วตฺติตพฺพํ. ตตฺรายํ สมฺมาวตฺตนา\csromandash{}}

\prose{น อุปสมฺ\-ปาเท\-ตพฺพํ ฯเปฯ (ยถา มูลาย ปฏิกสฺสนา, ตถา วิตฺถาเรตพฺพํ.) น ภิกฺขูหิ สมฺป\-โยเช\-ตพฺพํ.}

\prose{น ภิกฺขเว มานตฺตารเหน ภิกฺขุนา ปกตตฺตสฺส ภิกฺขุโน ปุรโต คนฺตพฺพํ, น ปุรโต นิสีทิตพฺพํ, โย โหติ สํฆสฺส อาสนปริยนฺโต เสยฺยาปริยนฺโต วิหาร\-ปริยนฺโต, โส ตสฺส ปทาตพฺโพ, เตน จ โส สาทิตพฺโพ.}

\prose{\csromanfolio{93}น ภิกฺขเว มานตฺตารเหน ภิกฺขุนา ปกตตฺเตน ภิกฺขุนา ปุเรสมเณน วา ปจฺฉาสมเณน วา กุลานิ อุปสงฺกมิตพฺพานิ, น อารญฺญิกงฺคํ สมาทาตพฺพํ, น ปิณฺฑปาติ\-กงฺคํ สมาทาตพฺพํ, น จ ตปฺปจฺจยา ปิณฺฑปาโต นีหราเปตพฺโพ “มา มํ ชานิํสู”ติ.}

\prose{น ภิกฺขเว มานตฺตารเหน ภิกฺขุนา สภิกฺขุกา อาวาสา อภิกฺขุโก อาวาโส คนฺตพฺโพ อญฺญตฺร ปกตตฺเตน อญฺญตฺร อนฺตรายา.}

\prose{น ภิกฺขเว มานตฺตารเหน ภิกฺขุนา สภิกฺขุกา อาวาสา อภิกฺขุโก อนาวาโส คนฺตพฺโพ อญฺญตฺร ปกตตฺเตน อญฺญตฺร อนฺตรายา.}

\prose{น ภิกฺขเว มานตฺตารเหน ภิกฺขุนา สภิกฺขุกา อาวาสา อภิกฺขุโก อาวาโส วา อนาวาโส วา คนฺตพฺโพ อญฺญตฺร ปกตตฺเตน อญฺญตฺร อนฺตรายา.}

\prose{น ภิกฺขเว มานตฺตารเหน ภิกฺขุนา สภิกฺขุกา อนาวาสา อภิกฺขุโก อาวาโส คนฺตพฺโพ ฯเปฯ อภิกฺขุโก อนาวาโส คนฺตพฺโพ ฯเปฯ อภิกฺขุโก อาวาโส วา อนาวาโส วา คนฺตพฺโพ อญฺญตฺร ปกตตฺเตน อญฺญตฺร อนฺตรายา.}

\prose{น ภิกฺขเว มานตฺตารเหน ภิกฺขุนา สภิกฺขุกา อาวาสา วา อนาวาสา วา อภิกฺขุโก อาวาโส คนฺตพฺโพ ฯเปฯ อภิกฺขุโก อนาวาโส คนฺตพฺโพ ฯเปฯ อภิกฺขุโก อาวาโส วา อนาวาโส วา คนฺตพฺโพ อญฺญตฺร ปกตตฺเตน อญฺญตฺร อนฺตรายา.}

\prose{น ภิกฺขเว มานตฺตารเหน ภิกฺขุนา สภิกฺขุกา อาวาสา สภิกฺขุโก อาวาโส คนฺตพฺโพ ฯเปฯ สภิกฺขุโก อนาวาโส คนฺตพฺโพ ฯเปฯ สภิกฺขุโก อาวาโส วา อนาวาโส วา คนฺตพฺโพ ยตฺถสฺสุ ภิกฺขู นานาสํวาสกา อญฺญตฺร ปกตตฺเตน อญฺญตฺร อนฺตรายา.}

\prose{น ภิกฺขเว มานตฺตารเหน ภิกฺขุนา สภิกฺขุกา อนาวาสา สภิกฺขุโก อาวาโส คนฺตพฺโพ ฯเปฯ สภิกฺขุโก อนาวาโส \csromanfolio{94}คนฺตพฺโพ ฯเปฯ สภิกฺขุโก อาวาโส วา อนาวาโส วา คนฺตพฺโพ ยตฺถสฺสุ ภิกฺขู นานาสํวาสกา อญฺญตฺร ปกตตฺเตน อญฺญตฺร อนฺตรายา.}

\prose{น ภิกฺขเว มานตฺตารเหน ภิกฺขุนา สภิกฺขุกา อาวาสา วา อนาวาสา วา สภิกฺขุโก อาวาโส คนฺตพฺโพ ฯเปฯ สภิกฺขุโก อนาวาโส คนฺตพฺโพ ฯเปฯ สภิกฺขุโก อาวาโส วา อนาวาโส วา คนฺตพฺโพ ยตฺถสฺสุ ภิกฺขู นานาสํวาสกา อญฺญตฺร ปกตตฺเตน อญฺญตฺร อนฺตรายา.}

\prose{คนฺตพฺโพ ภิกฺขเว มานตฺตารเหน ภิกฺขุนา สภิกฺขุกา อาวาสา สภิกฺขุโก อาวาโส ฯเปฯ สภิกฺขุโก อนาวาโส ฯเปฯ สภิกฺขุโก อาวาโส วา อนาวาโส วา ยตฺถสฺสุ ภิกฺขู สมาน\-สํวาสกา ยํ ชญฺญา “สกฺโกมิ อชฺเชว คนฺตุนฺ”ติ.}

\prose{คนฺตพฺโพ ภิกฺขเว มานตฺตารเหน ภิกฺขุนา สภิกฺขุกา อนาวาสา สภิกฺขุโก อาวาโส ฯเปฯ สภิกฺขุโก อนาวาโส ฯเปฯ สภิกฺขุโก อาวาโส วา อนาวาโส วา ยตฺถสฺสุ ภิกฺขู สมาน\-สํวาสกา ยํ ชญฺญา “สกฺโกมิ อชฺเชว คนฺตุนฺ”ติ.}

\prose{คนฺตพฺโพ ภิกฺขเว มานตฺตารเหน ภิกฺขุนา สภิกฺขุกา อาวาสา วา อนาวาสา วา สภิกฺขุโก อาวาโส ฯเปฯ สภิกฺขุโก อนาวาโส ฯเปฯ สภิกฺขุโก อาวาโส วา อนาวาโส วา ยตฺถสฺสุ ภิกฺขู สมาน\-สํวาสกา ยํ ชญฺญา “สกฺโกมิ อชฺเชว คนฺตุนฺ”ติ.}

\prose{น ภิกฺขเว มานตฺตารเหน ภิกฺขุนา ปกตตฺเตน ภิกฺขุนา สทฺธิํ เอกจฺฉนฺเน อาวาเส วตฺถพฺพํ, น เอกจฺฉนฺเน อนาวาเส วตฺถพฺพํ, น เอกจฺฉนฺเน อาวาเส วา อนาวาเส วา วตฺถพฺพํ, ปกตตฺตํ ภิกฺขุํ ทิสฺวา อาสนา วุฏฺฐาตพฺพํ, ปกตตฺโต ภิกฺขุ อาสเนน นิมนฺเตตพฺโพ, น ปกตตฺเตน ภิกฺขุนา สทฺธิํ เอกาสเน นิสีทิตพฺพํ, น นีเจ อาสเน นิสินฺเน อุจฺเจ อาสเน นิสีทิตพฺพํ, น ฉมายํ นิสินฺเน อาสเน นิสีทิตพฺพํ, น เอกจงฺกเม จงฺกมิตพฺพํ, น นีเจ จงฺกเม จงฺกมนฺเต อุจฺเจ จงฺกเม จงฺกมิตพฺพํ, น ฉมายํ จงฺกมนฺเต จงฺกเม จงฺกมิตพฺพํ.}

\prose{\csromanfolio{95}น ภิกฺขเว มานตฺตารเหน ภิกฺขุนา ปาริวาสิเกน ภิกฺขุนา สทฺธิํ ฯเปฯ มูลาย\-ปฏิกสฺสนา\-รเหน ภิกฺขุนา สทฺธิํ ฯเปฯ มานตฺตารเหน วุฑฺฒตเรน ภิกฺขุนา สทฺธิํ ฯเปฯ มานตฺต\-จาริเกน ภิกฺขุนา สทฺธิํ ฯเปฯ อพฺภานาเรเหน ภิกฺขุนา สทฺธิํ เอกจฺฉนฺเน อาวาเส วตฺถพฺพํ, น เอกจฺฉนฺเน อนาวาเส วตฺถพฺพํ, น เอกจฺฉนฺเน อาวาเส วา อนาวาเส วา วตฺถพฺพํ, น เอกาสเน นิสีทิตพฺพํ, น นีเจ อาสเน นิสินฺเน อุจฺเจ อาสเน นิสีทิตพฺพํ, น ฉมายํ นิสินฺเน อาสเน นิสีทิตพฺพํ, น เอกจงฺกเม จงฺกมิตพฺพํ, น นีเจ จงฺกเม จงฺกมนฺเต อุจฺเจ จงฺกเม จงฺกมิตพฺพํ, น ฉมายํ จงฺกมนฺเต จงฺกเม จงฺกมิตพฺพํ.}

\prosewithcloser{\csromansymbolfootnote{*}{วิ 3. 442 ปิฏฺเฐ.}มานตฺตารห\-จตุตฺโถ เจ ภิกฺขเว ปริวาสํ ทเทยฺย, มูลาย ปฏิกสฺเสยฺย, มานตฺตํ ทเทยฺย, ตํวีโส อพฺเภยฺย, อกมฺมํ น จ กรณียนฺติ.}{มานตฺตารห\-วตฺตํ นิฏฺฐิตํ.}
\csromanmatikaanchor{mtk.56}
\csromantocmarkref{section}{4. มานตฺตจาริกวตฺต}{mtk.56}
\bookmark[dest={mtk.56},level=1]{4. มานตฺตจาริกวตฺต}
\csromanheader{1}{4. มานตฺตจาริก\-วตฺต}
\proseitem{90}{เตน โข ปน สมเยน มานตฺตจาริกา ภิกฺขู สาทิยนฺติ ปกตตฺตานํ ภิกฺขูนํ อภิวาทนํ ปจฺจุฏฺฐานํ อญฺชลิกมฺมํ สามีจิกมฺมํ อาสนาภิหารํ เสยฺยาภิหารํ ปาโททกํ ปาทปีฐํ ปาทกถลิกํ ปตฺตจีวร\-ปฏิคฺคหณํ นหาเน ปิฏฺฐิ\-ปริกมฺมํ. เย เต ภิกฺขู อปฺปิจฺฉา ฯเปฯ เต อุชฺฌายนฺติ ขิยฺยนฺติ วิปาเจนฺติ “กถํ หิ นาม มานตฺตจาริกา ภิกฺขู สาทิยิสฺสนฺติ ปกตตฺตานํ ภิกฺขูนํ อภิวาทนํ ปจฺจุฏฺฐานํ ฯเปฯ นหาเน ปิฏฺฐิปริกมฺมนฺ”ติ. อถ โข เต ภิกฺขู ภควโต เอตมตฺถํ อาโรเจสุํ.}

\prose{อถ โข ภควา เอตสฺมิํ นิทาเน เอตสฺมิํ ปกรเณ ภิกฺขุสํฆํ สนฺนิปาตาเปตฺวา ภิกฺขู ปฏิปุจฺฉิ “สจฺจํ กิร ภิกฺขเว มานตฺตจาริกา ภิกฺขู สาทิยนฺติ ปกตตฺตานํ ภิกฺขูนํ อภิวาทนํ ปจฺจุฏฺฐานํ ฯเปฯ นหาเน ปิฏฺฐิปริกมฺมนฺ”ติ. สจฺจํ ภควาติ. วิครหิ พุทฺโธ ภควา อนนุจฺฉวิกํ ฯเปฯ. กถํ หิ นาม ภิกฺขเว มานตฺตจาริกา ภิกฺขู สาทิยิสฺสนฺติ ปกตตฺตานํ ภิกฺขูนํ อภิวาทนํ ฯเปฯ นหาเน ปิฏฺฐิ\-ปริกมฺมํ. เนตํ ภิกฺขเว อปฺปสนฺนานํ วา ปสาทาย ฯเปฯ วิครหิตฺวา ฯเปฯ ธมฺมิํ กถํ กตฺวา ภิกฺขู อามนฺเตสิ\csromandash{}}

\prose{\csromanfolio{96}น ภิกฺขเว มานตฺต\-จาริเกน ภิกฺขุนา สาทิตพฺพํ ปกตตฺตานํ ภิกฺขูนํ อภิวาทนํ ปจฺจุฏฺฐานํ อญฺชลิกมฺมํ สามีจิกมฺมํ อาสนาภิหาโร เสยฺยาภิหาโร ปาโททกํ ปาทปีฐํ ปาทกถลิกํ ปตฺตจีวร\-ปฏิคฺคหณํ นหาเน ปิฏฺฐิ\-ปริกมฺมํ, โย สาทิเยยฺย, อาปตฺติ ทุกฺกฏสฺส. อนุชานามิ ภิกฺขเว มานตฺต\-จาริกานํ ภิกฺขูนํ มิถุ ยถาวุฑฺฒํ อภิวาทนํ ปจฺจุฏฺฐานํ ฯเปฯ นหาเน ปิฏฺฐิ\-ปริกมฺมํ. อนุชานามิ ภิกฺขเว มานตฺตจริกานํ ภิกฺขูนํ ปญฺจ ยถาวุฑฺฒํ อุโปสถํ ปวารณํ วสฺสิกสาฏิกํ โอโณชนํ ภตฺตํ. เตน หิ ภิกฺขเว มานตฺต\-จาริกานํ ภิกฺขูนํ วตฺตํ ปญฺญเปสฺสามิ, ยถา มานตฺต\-จาริเกหิ ภิกฺขูหิ วตฺติตพฺพํ.}

\proseitem{91}{มานตฺต\-จาริเกน ภิกฺขเว ภิกฺขุนา สมฺมา วตฺติตพฺพํ. ตตฺรายํ สมฺมาวตฺตนา\csromandash{}}

\prose{น อุปสมฺ\-ปาเท\-ตพฺพํ, น นิสฺสโย ทาตพฺโพ, น สามเณโร อุปฏฺฐาเปตพฺโพ, น ภิกฺขุโนวาทก\-สมฺมุติ สาทิตพฺพา, สมฺมเตนปิ ภิกฺขุนิโย น โอวทิตพฺพา, ยาย อาปตฺติยา สํเฆน มานตฺตํ ทินฺนํ โหติ, สา อาปตฺติ น อาปชฺชิตพฺพา, อญฺญา วา ตาทิสิกา, ตโต วา ปาปิฏฺฐตรา, กมฺมํ น ครหิตพฺพํ, กมฺมิกา น ครหิตพฺพา, น ปกตตฺตสฺส ภิกฺขุโน อุโปสโถ ฐเปตพฺโพ, น ปวารณา ฐเปตพฺพา, น สวจนียํ กาตพฺพํ, น อนุวาโท ปฏฺฐเปตพฺโพ, น โอกาโส กาเรตพฺโพ น โจเทตพฺโพ, น สาเรตพฺโพ, น ภิกฺขูหิ สมฺป\-โยเช\-ตพฺพํ.}

\prose{น ภิกฺขเว มานตฺต\-จาริเกน ภิกฺขุนา ปกตตฺตสฺส ภิกฺขุโน ปุรโต คนฺตพฺพํ, น ปุรโต นิสีทิตพฺพํ, โย โหติ สํฆสฺส อาสนปริยนฺโต เสยฺยาปริยนฺโต วิหาร\-ปริยนฺโต, โส ตสฺส ปทาตพฺโพ, เตน จ โส สาทิตพฺโพ.}

\prose{น ภิกฺขเว มานตฺต\-จาริเกน ภิกฺขุนา ปกตตฺเตน ภิกฺขุนา ปุเรสมเณน วา ปจฺฉาสมเณน วา กุลานิ อุปสงฺกมิตพฺพานิ, น อารญฺญิกงฺคํ สมาทาตพฺพํ, น ปิณฺฑปาติ\-กงฺคํ สมาทาตพฺพํ, น จ ตปฺปจฺจยา ปิณฺฑปาโต นีหราเปตพฺโพ “มา มํ ชานิํสู”ติ.}

\prose{มานตฺต\-จาริเกน ภิกฺขเว ภิกฺขุนา อาคนฺตุเกน อาโรเจตพฺพํ, อาคนฺตุกสฺส อาโรเจตพฺพํ, อุโปสเถ อาโรเจตพฺพํ, ปวารณาย \csromanfolio{97}อาโรเจตพฺพํ, เทวสิกํ อาโรเจตพฺพํ, สเจ คิลาโน โหติ, ทูเตนปิ อาโรเจตพฺพํ.}

\prose{น ภิกฺขเว มานตฺต\-จาริเกน ภิกฺขุนา สภิกฺขุกา อาวาสา อภิกฺขุโก อาวาโส คนฺตพฺโพ อญฺญตฺร สํเฆน อญฺญตฺร อนฺตรายา.}

\prose{น ภิกฺขเว มานตฺต\-จาริเกน ภิกฺขุนา สภิกฺขุกา อาวาสา อภิกฺขุโก อนาวาโส คนฺตพฺโพ อญฺญตฺร สํเฆน อญฺญตฺร อนฺตรายา.}

\prose{น ภิกฺขเว มานตฺต\-จาริเกน ภิกฺขุนา สภิกฺขุกา อาวาสา อภิกฺขุโก อาวาโส วา อนาวาโส วา คนฺตพฺโพ อญฺญตฺร สํเฆน อญฺญตฺร อนฺตรายา.}

\prose{น ภิกฺขเว มานตฺต\-จาริเกน ภิกฺขุนา สภิกฺขุกา อนาวาสา อภิกฺขุโก อาวาโส คนฺตพฺโพ ฯเปฯ อภิกฺขุโก อนาวาโส คนฺตพฺโพ ฯเปฯ อภิกฺขุโก อาวาโส วา อนาวาโส วา คนฺตพฺโพ อญฺญตฺร สํเฆน อญฺญตฺร อนฺตรายา.}

\prose{น ภิกฺขเว มานตฺต\-จาริเกน ภิกฺขุนา สภิกฺขุกา อาวาสา วา อนาวาสา วา อภิกฺขุโก อาวาโส คนฺตพฺโพ ฯเปฯ อภิกฺขุโก อนาวาโส คนฺตพฺโพ ฯเปฯ อภิกฺขุโก อาวาโส วา อนาวาโส วา คนฺตพฺโพ อญฺญตฺร สํเฆน อญฺญตฺร อนฺตรายา.}

\prose{น ภิกฺขเว มานตฺต\-จาริเกน ภิกฺขุนา สภิกฺขุกา อาวาสา สภิกฺขุโก อาวาโส คนฺตพฺโพ ฯเปฯ สภิกฺขุโก อานาวาโส คนฺตพฺโพ ฯเปฯ สภิกฺขุโก อาวาโส วา อนาวาโส วา คนฺตพฺโพ ยตฺถสฺสุ ภิกฺขู นานาสํวาสกา อญฺญตฺร สํเฆน อญฺญตฺร อนฺตรายา.}

\prose{น ภิกฺขเว มานตฺต\-จาริเกน ภิกฺขุนา สภิกฺขุกา อนาวาสา สภิกฺขุโก อาปาโส คนฺตพฺโพ ฯเปฯ สภิกฺขุโก อนาวาโส คนฺตพฺโพ ฯเปฯ สภิกฺขุโก อาวาโส วา อนาวาโส วา คนฺตพฺโพ ยตฺถสฺสุ ภิกฺขู นานาสํวาสกา อญฺญตฺร สํเฆน อญฺญตฺร อนฺตรายา.}

\prose{น ภิกฺขเว มานตฺต\-จาริเกน ภิกฺขุนา สภิกฺขุกา อาวาสา วา อนาวาสา วา สภิกฺขุโก อาวาโส คนฺตพฺโพ ฯเปฯ สภิกฺขุโก อนาวาโส คนฺตพฺโพ ฯเปฯ สภิกฺขุโก อาวาโส วา อนาวาโส วา \csromanfolio{98}คนฺตพฺโพ ยตฺถสฺสุ ภิกฺขู นานาสํวาสกา อญฺญตฺร สํเฆน อญฺญตฺร อนฺตรายา.}

\prose{คนฺตพฺโพ ภิกฺขเว มานตฺต\-จาริเกน ภิกฺขุนา สภิกฺขุกา อาวาสา สภิกฺขุโก อาวาโส ฯเปฯ สภิกฺขุโก อนาวาโส ฯเปฯ สภิกฺขุโก อาวาโส วา อนาวาโส วา ยตฺถสฺสุ ภิกฺขู สมาน\-สํวาสกา ยํ ชญฺญา “สกฺโกมิ อชฺเชว คนฺตุนฺ”ติ.}

\prose{คนฺตพฺโพ ภิกฺขเว มานตฺต\-จาริเกน ภิกฺขุนา สภิกฺขุกา อนาวาสา สภิกฺขุโก อาวาโส ฯเปฯ สภิกฺขุโก อนาวาโส ฯเปฯ สภิกฺขุโก อาวาโส วา อนาวาโส วา ยตฺถสฺสุ ภิกฺขู สมาน\-สํวาสกา ยํ ชญฺญา “สกฺโกมิ อชฺเชว คนฺตุนฺ”ติ.}

\prose{คนฺตพฺโพ ภิกฺขเว มานตฺต\-จาริเกน ภิกฺขุนา สภิกฺขุกา อาวาสา วา อนาวาสา วา สภิกฺขุโก อาวาโส ฯเปฯ สภิกฺขุโก อนาวาโส ฯเปฯ สภิกฺขุโก อาวาโส วา อนาวาโส วา ยตฺถสฺสุ ภิกฺขู สมาน\-สํวาสกา ยํ ชญฺญา “สกฺโกมิ อชฺเชว คนฺตุนฺ”ติ.}

\prose{น ภิกฺขเว มานตฺต\-จาริเกน ภิกฺขุนา ปกตตฺเตน ภิกฺขุนา สทฺธิํ เอกจฺฉนฺเน อาวาเส วตฺถพฺพํ, น เอกจฺฉนฺเน อนาวาเส วตฺถพฺพํ, น เอกจฺฉนฺเน อาวาเส วา อนาวาเส วา วตฺถพฺพํ, ปกตตฺตํ ภิกฺขุํ ทิสฺวา อาสนา วุฏฺฐาตพฺพํ, ปกตตฺโต ภิกฺขุ อาสเนน นิมนฺเตตพฺโพ, น ปกตตฺเตน ภิกฺขุนา สทฺธิํ เอกาสเน นิสีทิตพฺพํ, น นีเจ อาสเน นิสินฺเน อุจฺเจ อาสเน นิสีทิตพฺพํ, น ฉมายํ นิสินฺเน อาสเน นิสีทิตพฺพํ, น เอกจงฺกเม จงฺกมิตพฺพํ, น นีเจ จงฺกเม จงฺกมนฺเต อุจฺเจ จงฺกเม จงฺกมิตพฺพํ, น ฉมายํ จงฺกมนฺเต จงฺกมิตพฺพํ.}

\prose{น ภิกฺขเว มานตฺต\-จาริเกน ภิกฺขุนา ปาริวาสิเกน ภิกฺขุนา สทฺธิํ ฯเปฯ มูลาย\-ปฏิกสฺสนา\-รเหน ภิกฺขุนา สทฺธิํ ฯเปฯ มานตฺตารเหน ภิกฺขุนา สทฺธิํ ฯเปฯ มานตฺต\-จาริเกน วุฑฺฒตเรน ภิกฺขุนา สทฺธิํ ฯเปฯ อพฺภานารเหน ภิกฺขุนา สทฺธิํ เอกจฺฉนฺเน อาวาเส วตฺถพฺพํ, น เอกจฺฉนฺเน อนาวาเส วตฺถพฺพํ, น เอกจฺฉนฺเน อาวาเส วา อนาวเส วา วตฺถพฺพํ, น เอกาสเน นิสีทิตพฺพํ, น นีเจ อาสเน นิสินฺเน อุจฺเจ อาสเน นิสีทิตพฺพํ, น ฉมายํ นิสินฺเน อาสเน นิสีทิตพฺพํ, น เอกจงฺกเม \csromanfolio{99}จงฺกมิตพฺพํ, น นีเจ จงฺกเม จงฺกมนฺเต อุจฺเจ จงฺกเม จงฺกมิตพฺพํ, น ฉมายํ จงฺกมนฺเต จงฺกเม จงฺกมิตพฺพํ.}

\prose{\csromansymbolfootnote{*}{วิ 3. 442 ปิฏฺเฐ.}มานตฺตจาริก\-จตุตฺโถ เจ ภิกฺขเว ปริวาสํ ทเทยฺย, มูลาย ปฏิกสฺเสยฺย, มานตฺตํ ทเทยฺย, ตํวีโส อพฺเภยฺย, อกมฺมํ น จ กรณียนฺติ.}

\proseitem{92}{อถ โข อายสฺมา อุปาลิ เยน ภควา เตนุปสงฺกมิ, อุปสงฺกมิตฺวา ภควนฺตํ อภิวาเทตฺวา เอกมนฺตํ นิสีทิ, เอกมนฺตํ นิสินฺโน โข อายสฺมา อุปาลิ ภควนฺตํ เอตทโวจ “กติ นุ โข ภนฺเต มานตฺต\-จาริกสฺส ภิกฺขุโน รตฺติจฺเฉทา”ติ. จตฺตาโร โข อุปาลิ มานตฺต\-จาริกสฺส ภิกฺขุโน รตฺติจฺเฉทา, สหวาโส วิปฺปวาโส อนาโรจนา อูเน คเณ จรณํ\csromansharedfootnote{a856ea35bd5d134d}{จรณนฺติ (ก)}, อิเม โข อุปาลิ จตฺตาโร มานตฺต\-จาริกสฺส ภิกฺขุโน รตฺติจฺเฉทาติ.}

\proseitem{93}{เตน โข ปน สมเยน สาวตฺถิยํ มหา ภิกฺขุสํโฆ สนฺนิปติโต โหติ, น สกฺโกนฺติ มานตฺตจาริกา ภิกฺขู มานตฺตํ โสเธตุํ, ภควโต เอตมตฺถํ อาโรเจสุํ. อนุชานามิ ภิกฺขเว มานตฺตํ นิกฺขิปิตุํ. เอวญฺจ ปน ภิกฺขเว นิกฺขิปิตพฺพํ, เตน มานตฺต\-จาริเกน ภิกฺขุนา เอกํ ภิกฺขุํ อุปสงฺกมิตฺวา เอกํสํ อุตฺตราสงฺคํ กริตฺวา อุกฺกุฏิกํ นิสีทิตฺวา อญฺชลิํ ปคฺคเหตฺวา เอวมสฺส วจนีโย “มานตฺตํ นิกฺขิปามี”ติ, นิกฺขิตฺตํ โหติ มานตฺตํ. “วตฺตํ นิกฺขิปามี”ติ, นิกฺขิตฺตํ โหติ มานตฺตนฺติ.}

\proseitemwithcloser{94}{เตน โข ปน สมเยน สาวตฺถิยา ภิกฺขู ตหํ ตหํ ปกฺกมิํสุ, สกฺโกนฺติ มานตฺตจาริกา ภิกฺขู มานตฺตํ โสเธตุํ, ภควโต เอตมตฺถํ อาโรเจสุํ. อนุชานามิ ภิกฺขเว มานตฺตํ สมาทิยิตุํ. เอวญฺจ ปน ภิกฺขเว สมาทิยิตพฺพํ, เตน มานตฺต\-จาริเกน ภิกฺขุนา เอกํ ภิกฺขุํ อุปสงฺกมิตฺวา เอกํสํ อุตฺตราสงฺคํ กริตฺวา อุกฺกุฏิกํ นิสีทิตฺวา อญฺชลิํ ปคฺคเหตฺวา เอวมสฺส วจนีโย “มานตฺตํ สมาทิยามี”ติ, สมาทินฺนํ โหติ มานตฺตํ. “วตฺตํ สมาทิยามี”ติ, สมาทินฺนํ โหติ มานตฺตนฺติ.}{มานตฺตจาริก\-วตฺตํ นิฏฺฐิตํ.}
\csromanheader{2}{5. \textbf{อพฺภานารตวตฺต}}
\csromanmatikaanchor{mtk.57}
\csromantocmarkref{section}{5. อพฺภานารหวตฺต}{mtk.57}
\bookmark[dest={mtk.57},level=1]{5. อพฺภานารหวตฺต}
\proseitem{95}{\csromanfolio{100}เตน โข ปน สมเยน อพฺภานารหาภิกฺขู สาทิยนฺติ ปกตตฺตานํ ภิกฺขูนํ อภิวาทนํ ปจฺจุฏฺฐานํ ฯเปฯ นหาเน ปิฏฺฐิ\-ปริกมฺมํ. เย เต ภิกฺขู อปฺปิจฺฉา ฯเปฯ เต อุชฺฌายนฺติ ขิยฺยนฺติ วิปาเจนฺติ “กถํ หิ นาม อพฺภานารหา ภิกฺขู สาทิยิสฺสนฺติ ปกตตฺตานํ ภิกฺขูนํ อภิวาทนํ ปจฺจุฏฺฐานํ ฯเปฯ นหาเน ปิฏฺฐิปริกมฺมนฺ”ติ. อถ โข เต ภิกฺขู ภควโต เอตมตฺถํ อาโรเจสุํ.}

\prose{อถ โข ภควา เอตสฺมิํ นิทาเน เอตสฺมิํ ปกรเณ ภิกฺขุสํฆํ สนฺนิปาตาเปตฺวา ภิกฺขู ปฏิปุจฺฉิ “สจฺจํ กิร ภิกฺขเว อพฺภานารหา ภิกฺขู สาทิยนฺติ ปกตตฺตานํ ภิกฺขูนํ อภิวาทนํ ปจฺจุฏฺฐานํ ฯเปฯ นหาเน ปิฏฺฐิปริกมฺมนฺ”ติ, สจฺจํ ภควาติ. วิครหิ พุทฺโธ ภควา อนนุจฺฉวิกํ ฯเปฯ. กถํ หิ นาม ภิกฺขเว อพฺภานารหา ภิกฺขู สาทิยิสฺสนฺติ ปกตตฺถานํ ภิกฺขูนํ อภิวาทนํ ปจฺจุฏฺฐานํ ฯเปฯ นหาเน ปิฏฺฐิ\-ปริกมฺมํ. เนตํ ภิกฺขเว อปฺปสนฺนานํ วา ปสาทาย ฯเปฯ วิครหิตฺวา ฯเปฯ ธมฺมิํ กถํ กตฺวา ภิกฺขู อามนฺเตสิ\csromandash{} น ภิกฺขเว อพฺภานารเหน ภิกฺขุนา สาทิตพฺพํ ปกตตฺตานํ ภิกฺขูนํ อภิวาทนํ ปจฺจุฏฺฐานํ ฯเปฯ นหาเน ปิฏฺฐิ\-ปริกมฺมํ. โย สาทิเยยฺย, อาปตฺติ ทุกฺกฏสฺส. อนุชานามิ ภิกฺขเว อพฺภานารหานํ ภิกฺขูนํ มิถุ ยถาวุฑฺฒํ อภิวาทนํ ปจฺจุฏฺฐานํ ฯเปฯ นหาเน ปิฏฺฐิ\-ปริกมฺมํ, อนุชานามิ ภิกฺขเว อพฺภานารหานํ ภิกฺขูนํ ปญฺจ ยถาวุฑฺฒํ อุโปสถํ ปวารณํ วสฺสิกสาฏิกํ โอโณชนํ ภตฺตํ. เตน หิ ภิกฺขเว อพฺภานารหานํ ภิกฺขูนํ ปตฺตํ ปญฺญ ปสฺสามิ, ยถา อพฺภานารเหหิ ภิกฺขูหิ วตฺติตพฺพํ.}

\proseitem{96}{อพฺภานารเหน ภิกฺขเว ภิกฺขุนา สมฺมา วตฺติตพฺพํ. ตตฺรายํ สมฺมาวตฺตนา\csromandash{}}

\prose{น อุปสมฺ\-ปาเท\-ตพฺพํ ฯเปฯ (ยถา เหฏฺฐา, ตถา วิตฺถาเรตพฺพํ,) น ภิกฺขูหิ สมฺป\-โยเช\-ตพฺพํ.}

\prose{น ภิกฺขเว อพฺภานารเหน ภิกฺขุนา ปกตตฺตสฺส ภิกฺขุโน ปุรโต คนฺตพฺพํ, น ปุรโต นิสีทิตพฺพํ, โย โหติ สํฆสฺส อาสนปริยนฺโต เสยฺยาปริยนฺโต วิหาร\-ปริยนฺโต, โส ตสฺส ปทาตพฺโพ, เตน จ โส สาทิตพฺโพ.}

\prose{\csromanfolio{101}น ภิกฺขเว อพฺภานารเหน ภิกฺขุนา ปกตตฺเตน ภิกฺขุนา ปุเรสมเณน วา ปจฺฉาสมเณน วา กุลานิ อุปสงฺกมิตพฺพานิ, น อารญฺญิกงฺคํ สมาทาตพฺพํ, น ปิณฺฑปาติ\-กงฺคํ สมาทาตพฺพํ, น จ ตปฺปจฺจยา ปิณฺฑปาโต นีหราเปตพฺโพ “มา มํ ชานิํสู”ติ.}

\prose{น ภิกฺขเว อพฺภานารเหน ภิกฺขุนา สภิกฺขุกา อาวาสา อภิกฺขุโก อาวาโส คนฺตพฺโพ อญฺญตฺร ปกตตฺเตน อญฺญตฺร อนฺตรายา.}

\prose{น ภิกฺขเว อพฺภานารเหน ภิกฺขุนา สภิกฺขุกา อาวาสา อภิกฺขุโก อนาวาโส คนฺตพฺโพ อญฺญตฺร ปกตตฺเตน อญฺญตฺร อนฺตรายา ฯเปฯ (ยถา เหฏฺฐา, ตถา วิตฺถาเรตพฺพา.)}

\prose{คนฺตพฺโพ ภิกฺขเว อพฺภานารเหน ภิกฺขุนา สภิกฺขุกา อาวาสา ฯเปฯ อนาวาสา ฯเปฯ อาวาสา วา อนาวาสา วา สภิกฺขุโก อาวาโส ฯเปฯ สภิกฺขุโก อนาวาโส ฯเปฯ สภิกฺขุโก อาวาโส วา อนาวาโส วา ยตฺถสฺสุ ภิกฺขู สมาน\-สํวาสกา ยํ ชญฺญา “สกฺโกมิ อชฺเชว คนฺตุนฺ”ติ.}

\prose{น ภิกฺขเว อพฺภานารเหน ภิกฺขุนา ปกตตฺเตน ภิกฺขุนา สทฺธิํ เอกจฺฉนฺเน อาวาเส วตฺถพฺพํ, น เอกจฺฉนฺเน อนาวาเส วตฺถพฺพํ, น เอกจฺฉนฺเน อาวาเส วา อนาวาเส วา วตฺถพฺพํ, ปกตตฺตํ ภิกฺขุํ ทิสฺวา อาสนา วุฏฺฐาตพฺพํ, ปกตตฺโต ภิกฺขุ อาสเนน นิมนฺเตตพฺโพ, น ปกตตฺเตน ภิกฺขุนา สทฺธิํ เอกาสเน นิสีทิตพฺพํ, น นีเจ อาสเน นิสินฺเน อุจฺเจ อาสเน นิสีทิตพฺพํ, น ฉมายํ นิสินฺเน อาสเน นิสีทิตพฺพํ, น เอกจงฺกเม จงฺกมิตพฺพํ, น นีเจ จงฺกเม จงฺกมนฺเต อุจฺเจ จงฺกเม จงฺกมิตพฺพํ, น ฉมายํ จงฺกมนฺเต จงฺกเม จงฺกมิตพฺพํ.}

\csromanmatikaanchor{mtk.58}
\csromantocmarkref{section}{อุทฺทานคาถา}{mtk.58}
\bookmark[dest={mtk.58},level=1]{อุทฺทานคาถา}
\prose{น ภิกฺขเว อพฺภานารเหน ภิกฺขุนา ปาริวาสิเกน ภิกฺขุนา สทฺธิํ ฯเปฯ มูลาย\-ปฏิกสฺสนา\-รเหน ภิกฺขุนา สทฺธิํ ฯเปฯ มานตฺตารเหน ภิกฺขุนา สทฺธิํ ฯเปฯ มานตฺต\-จาริเกน ภิกฺขุนา สทฺธิํ ฯเปฯ อพฺภานารเหน วุฑฺฒตเรน ภิกฺขุนา สทฺธิํ เอกจฺฉนฺเน อาวาเส วตฺถพฺพํ, น เอกจฺฉนฺเน อนาวาเส วตฺถพฺพํ, น เอกจฺฉนฺเน อาวาเส วา อนาวาเส วา วตฺถพฺพํ, น เอกาสเน นิสีทิตพฺพํ, น นีเจ อาสเน นิสินฺเน อุจฺเจ อาสเน นิสีทิตพฺพํ, น ฉมายํ นิสินฺเน อาสเน นิสีทิตพฺพํ, น เอกจงฺกเม \csromanfolio{102}จงฺกมิตพฺพํ, น นีเจ จงฺกเม จงฺกมนฺเต อุจฺเจ จงฺกเม จงฺกมิตพฺพํ, น ฉมายํ จงฺกมนฺเต จงฺกเม จงฺกมิตพฺพํ.}

\prose{\csromansymbolfootnote{*}{วิ 3. 442 ปิฏฺเฐ.}อพฺภานารห\-จตุตฺโถ เจ ภิกฺขเว ปริวาสํ ทเทยฺย, มูลาย ปฏิกสฺเสยฺย, มานตฺตํ ทเทยฺย, ตํวีโส อพฺเภยฺย, อกมฺมํ น จ กรณียนฺติ.}

\vspace{\tipitakaparskip}
\csromancenter{อพฺภานารห\-วตฺตํ นิฏฺฐิตํ. ปาริวาสิก\-กฺขนฺธโก ทุติโย. อิมมฺหิ ขนฺธเก วตฺถู ปญฺจ.}
\csromancenter{\textbf{ตสฺสุทฺทานํ}}
\setlength{\csromangathapagewidth}{0pt}
\setlength{\csromangathawakwidth}{0pt}
\csromangathameasuregroup{ปาริวาสิกา สาเทนฺติ, \\ อภิวาทนํ ปจฺจุฏฺฐานํ, \\ อาสนํ เสยฺยาภิหารํ, \\ ปตฺตํ นหาเน ปริกมฺมํ, \\ ทุกฺกฏํ สาทิยนฺตสฺส, \\ อุโปสถํ ปวารณํ, \\ สมฺมา จ วตฺตนา ตตฺถ, \\ โย จ โหติ ปริยนฺโต, \\ อารญฺญปิณฺฑนีหาโร, \\ ปวารณาย ทูเตน, \\ เอกจฺฉนฺเน จ วุฏฺฐานํ, \\ อาสเน นีเจ จงฺกเม, \\ วุฑฺฒตเรน อกมฺมํ, \\ นิกฺขิปนํ สมาทานํ, \\ มูลาย มานตฺตารหา, \\ อพฺภานารเห นโย จาปิ, \\ ปาริวาสิเกสุ ตโย,}{\csromangathabat{ปาริวาสิกา สาเทนฺติ,}{ปกตตฺตาน ภิกฺขุนํ.} \\ \csromangathabat{อภิวาทนํ ปจฺจุฏฺฐานํ,}{อญฺชลิํ จ สามีจิยํ.} \\ \csromangathabat{อาสนํ เสยฺยาภิหารํ,}{ปาโท ปีฐํ กถลิกํ.} \\ \csromangathabat{ปตฺตํ นหาเน ปริกมฺมํ,}{อุชฺฌายนฺติ จ เปสลา.} \\ \csromangathabat{ทุกฺกฏํ สาทิยนฺตสฺส,}{มิถุ ปญฺจ ยถาวุฑฺฒํ\textsuperscript{1}.} \\ \csromangathabat{อุโปสถํ ปวารณํ,}{วสฺสิโกโณชโภชนํ.} \\ \csromangathabat{สมฺมา จ วตฺตนา ตตฺถ,}{ปกตตฺตสฺส คจฺฉนฺตํ.} \\ \csromangathabat{โย จ โหติ ปริยนฺโต,}{ปุเร ปจฺฉา ตเถว จ\textsuperscript{1}.} \\ \csromangathabat{อารญฺญปิณฺฑนีหาโร,}{อาคนฺตุเก อุโปสเถ.} \\ \csromangathabat{ปวารณาย ทูเตน,}{คนฺตพฺโพ จ สภิกฺขุโก.} \\ \csromangathabat{เอกจฺฉนฺเน จ วุฏฺฐานํ,}{ตเถว จ นิมนฺตเย.} \\ \csromangathabat{อาสเน นีเจ จงฺกเม,}{ฉมายํ จงฺกเมน จ.} \\ \csromangathabat{วุฑฺฒตเรน อกมฺมํ,}{รตฺติจฺเฉทา จ โสธนา.} \\ \csromangathabat{นิกฺขิปนํ สมาทานํ,}{วตฺตํว ปาริวาสิเก\textsuperscript{1}.} \\ \csromangathabat{มูลาย มานตฺตารหา,}{ตถา มานตฺตจาริตา.} \\ \csromangathabat{อพฺภานารเห นโย จาปิ,}{สมฺเภทํ นยโต\textsuperscript{1} ปุน.} \\ \csromangathabat{ปาริวาสิเกสุ ตโย,}{จตุ มานตฺตจาริเก.} \\ น สเมนฺติ รตฺติจฺเฉเทสุ\textsuperscript{1} มานตฺเตสุ จ เทวสิ. \\ เทฺว กมฺมา สทิสา เสสา ตโย กมฺมา สมาสมาติ\textsuperscript{1}.}
\csromangathasetleft{ปาริวาสิกา สาเทนฺติ, \\ อภิวาทนํ ปจฺจุฏฺฐานํ, \\ อาสนํ เสยฺยาภิหารํ, \\ ปตฺตํ นหาเน ปริกมฺมํ, \\ ทุกฺกฏํ สาทิยนฺตสฺส, \\ อุโปสถํ ปวารณํ, \\ สมฺมา จ วตฺตนา ตตฺถ, \\ โย จ โหติ ปริยนฺโต, \\ อารญฺญปิณฺฑนีหาโร, \\ ปวารณาย ทูเตน, \\ เอกจฺฉนฺเน จ วุฏฺฐานํ, \\ อาสเน นีเจ จงฺกเม, \\ วุฑฺฒตเรน อกมฺมํ, \\ นิกฺขิปนํ สมาทานํ, \\ มูลาย มานตฺตารหา, \\ อพฺภานารเห นโย จาปิ, \\ ปาริวาสิเกสุ ตโย,}
\csromangathagroupwithclosermajor{\csromangathastanza{\csromangathabat{ปาริวาสิกา สาเทนฺติ,}{ปกตตฺตาน ภิกฺขุนํ.} \\ \csromangathabat{อภิวาทนํ ปจฺจุฏฺฐานํ,}{อญฺชลิํ จ สามีจิยํ.}}\csromangathastanza{\csromangathabat{อาสนํ เสยฺยาภิหารํ,}{ปาโท ปีฐํ กถลิกํ.} \\ \csromangathabat{ปตฺตํ นหาเน ปริกมฺมํ,}{อุชฺฌายนฺติ จ เปสลา.}}\csromangathastanza{\csromangathabat{ทุกฺกฏํ สาทิยนฺตสฺส,}{มิถุ ปญฺจ ยถาวุฑฺฒํ\csromansharedfootnote{c74ea57b33a5f961}{ปุนาปเร (ก)}.} \\ \csromangathabat{อุโปสถํ ปวารณํ,}{วสฺสิโกโณชโภชนํ.}}\csromangathastanza{\csromangathabat{สมฺมา จ วตฺตนา ตตฺถ,}{ปกตตฺตสฺส คจฺฉนฺตํ.} \\ \csromangathabat{โย จ โหติ ปริยนฺโต,}{ปุเร ปจฺฉา ตเถว จ\csromansharedfootnote{ac2d44a68bed7d03}{น ปุเร ปจฺฉาสมเณน (สี, สฺยา)}.}}\csromangathastanza{\csromangathabat{อารญฺญปิณฺฑนีหาโร,}{อาคนฺตุเก อุโปสเถ.} \\ \csromangathabat{ปวารณาย ทูเตน,}{คนฺตพฺโพ จ สภิกฺขุโก.}}\csromangathastanza{\csromangathabat{เอกจฺฉนฺเน จ วุฏฺฐานํ,}{ตเถว จ นิมนฺตเย.} \\ \csromangathabat{อาสเน นีเจ จงฺกเม,}{ฉมายํ จงฺกเมน จ.}}\csromangathastanza{\csromangathabat{วุฑฺฒตเรน อกมฺมํ,}{รตฺติจฺเฉทา จ โสธนา.} \\ \csromangathabat{นิกฺขิปนํ สมาทานํ,}{วตฺตํว ปาริวาสิเก\csromansharedfootnote{b3a76b8e44480180}{รตฺติ วา ปาริวาสิเก (ก), ญาตพฺพํ ปาริวาสิกา (สี, สฺยา)}.}}\csromangathastanza{\csromanfolio{103}\csromangathabat{มูลาย มานตฺตารหา,}{ตถา มานตฺตจาริตา.} \\ \csromangathabat{อพฺภานารเห นโย จาปิ,}{สมฺเภทํ นยโต\csromansharedfootnote{104b7c28a4312556}{สมฺเภทนยโต (สฺยา)} ปุน.}}\csromangathastanza{\csromangathabat{ปาริวาสิเกสุ ตโย,}{จตุ มานตฺตจาริเก.} \\ น สเมนฺติ รตฺติจฺเฉเทสุ\csromansharedfootnote{f1ab73e3996427f0}{รตฺติจฺเฉเท (อิติปิ), รตฺติจฺเฉทา (สฺยา)} มานตฺเตสุ จ เทวสิ. \\ เทฺว กมฺมา สทิสา เสสา ตโย กมฺมา สมาสมาติ\csromansharedfootnote{9e5b1ff0e9095d69}{สมา มตาติ (สี)}.}}{ปาริวาสิก\-กฺขนฺธโก นิฏฺฐิโต.}
\csromanmatikaanchor{mtk.59}
\csromantocmarknopage{chapter}{3. สมุจฺจยกฺขนฺธก}{mtk.59}
\bookmark[dest={mtk.59},level=0]{3. สมุจฺจยกฺขนฺธก}
\chapterheadpageread{3. สมุจฺจย\-กฺขนฺธก}
\csromanmatikaanchor{mtk.60}
\csromantocmarkref{section}{1. สุกฺกวิสฺสฏฺฐิ}{mtk.60}
\bookmark[dest={mtk.60},level=1]{1. สุกฺกวิสฺสฏฺฐิ}
\csromanheader{1}{1. \textbf{สุกฺกวิสฺสฏฺฐิ}}
\proseitem{97}{\csromanfolio{104}เตน สมเยน พุทฺโธ ภควา สาวตฺถิยํ วิหรติ เชตวเน อนาถ\-ปิณฺฑิกสฺส อาราเม. เตน โข ปน สมเยน อายสฺมา อุทายี เอกํ อาปตฺติํ อาปนฺโน โหติ สญฺเจตนิกํ สุกฺก\-วิสฺสฏฺฐิํ อปฺปฏิจฺฉนฺนํ, โส ภิกฺขูนํ อาโรเจสิ “อหํ อาวุโส เอกํ อาปตฺติํ อาปชฺชิํ สญฺเจตนิกํ สุกฺก\-วิสฺสฏฺฐิํ อปฺปฏิจฺฉนฺนํ, กถํ นุ โข มยา ปฏิปชฺชิตพฺพนฺ”ติ. ภควโต\csromansharedfootnote{44917929ccfdd3e3}{เต ภิกฺขู ภควโต (สฺยา)} เอตมตฺถํ อาโรเจสุํ. เตน หิ ภิกฺขเว สํโฆ อุทายิสฺส ภิกฺขุโน เอกิสฺสา อาปตฺติยา สญฺเจตนิกาย สุกฺก\-วิสฺสฏฺฐิยา อปฺปฏิจฺฉนฺนาย ฉารตฺตํ มานตฺตํ เทตุ. เอวญฺจ ปน ภิกฺขเว ทาตพฺพํ\csromandash{}}

\csromanmatikaanchor{mtk.61}
\csromantocmarkref{subsection}{อปฺปฏิจฺฉนฺนมานตฺต}{mtk.61}
\bookmark[dest={mtk.61},level=2]{อปฺปฏิจฺฉนฺนมานตฺต}
\csromanheader{2}{อปฺปฏิจฺฉนฺน\-มานตฺต}
\prose{98 เตน ภิกฺขเว อุทายินา ภิกฺขุนา สํฆํ อุปสงฺกมิตฺวา เอกํสํ อุตฺตราสงฺคํ กริตฺวา วุฑฺฒานํ ภิกฺขูนํ ปาเท วนฺทิตฺวา อุกฺกุฏิกํ นิสีทิตฺวา อญฺชลิํ ปคฺคเหตฺวา เอวมสฺส วจนีโย “อหํ ภนฺเต เอกํ อาปตฺติํ อาปชฺชิํ สญฺเจตนิกํ สุกฺก\-วิสฺสฏฺฐิํ อปฺปฏิจฺฉนฺนํ, โสหํ ภนฺเต สํฆํ เอกิสฺสา อาปตฺติยา สญฺเจตนิกาย สุกฺก\-วิสฺสฏฺฐิยา อปฺปฏิจฺฉนฺนาย ฉารตฺตํ มานตฺตํ ยาจามิ. อหํ ภนฺเต เอกํ อาปตฺติํ อาปชฺชิํ สญฺเจตนิกํ สุกฺก\-วิสฺสฏฺฐิํ อปฺปฏิจฺฉนฺนํ, ทุติยมฺปิ โสหํ\csromansharedfootnote{2e9d61fe46a6a538}{ทุติยมฺปิ (สี, ก)} ภนฺเต สํฆํ เอกิสฺสา อาปตฺติยา สญฺเจตนิกาย สุกฺก\-วิสฺสฏฺฐิยา อปฺปฏิจฺฉนฺนาย ฉารตฺตํ มานตฺตํ ยาจามิ. อหํ ภนฺเต เอกํ อาปตฺติํ อาปชฺชิํ สญฺเจตนิกํ สุกฺก\-วิสฺสฏฺฐิํ อปฺปฏิจฺฉนฺนํ, ตติยมฺปิ โสหํ\csromansharedfootnote{c11bcf7bc8b54786}{ตติยมฺปิ (สี, ก)} ภนฺเต สํฆํ เอกิสฺสา อาปตฺติยา สญฺเจตนิกาย สุกฺก\-วิสฺสฏฺฐิยา อปฺปฏิจฺฉนฺนาย ฉารตฺตํ มานตฺตํ ยาจามี”ติ. พฺยตฺเตน ภิกฺขุนา ปฏิพเลน สํโฆ ญาเปตพฺโพ\csromandash{}}

\proseitem{99}{“สุณาตุ เม ภนฺเต สํโฆ, อยํ อุทายี ภิกฺขุ เอกํ อาปตฺติํ อาปชฺชิ สญฺเจตนิกํ สุกฺก\-วิสฺสฏฺฐิํ อปฺปฏิจฺฉนฺนํ, โส สํฆํ เอกิสฺสา อาปตฺติยา สญฺเจตนิกาย สุกฺก\-วิสฺสฏฺฐิยา อปฺปฏิจฺฉนฺนาย ฉารตฺตํ มานตฺตํ ยาจติ, ยทิ สํฆสฺส ปตฺตกลฺลํ, สํโฆ อุทายิสฺส ภิกฺขุโน \csromanfolio{105}เอกิสฺสา อาปตฺติยา สญฺเจตนิกาย สุกฺก\-วิสฺสฏฺฐิยา อปฺปฏิจฺฉนฺนาย ฉารตฺตํ มานตฺตํ ทเทยฺย, เอสา ญตฺติ.}

\prose{สุณาตุ เม ภนฺเต สํโฆ, อยํ อุทายี ภิกฺขุ เอกํ อาปตฺติํ อาปชฺชิ สญฺเจตนิกํ สุกฺก\-วิสฺสฏฺฐิํ อปฺปฏิจฺฉนฺนํ, โส สํฆํ เอกิสฺสา อาปตฺติยา สญฺเจตนิกาย สุกฺก\-วิสฺสฏฺฐิยา อปฺปฏิจฺฉนฺนาย ฉารตฺตํ มานตฺตํ ยาจติ, สํโฆ อุทายิสฺส ภิกฺขุโน เอกิสฺสา อาปตฺติยา สญฺเจตนิกาย สุกฺก\-วิสฺสฏฺฐิยา อปฺปฏิจฺฉนฺนาย ฉารตฺตํ มานตฺตํ เทติ, ยสฺสายสฺมโต ขมติ อุทายิสฺส ภิกฺขุโน เอกิสฺสา อาปตฺติยา สญฺเจตนิกาย สุกฺก\-วิสฺสฏฺฐิยา อปฺปฏิจฺฉนฺนาย ฉารตฺตํ มานตฺตสฺส ทานํ, โส ตุณฺหสฺส, ยสฺส นกฺขมติ, โส ภาเสยฺย.}

\prose{ทุติยมฺปิ เอตมตฺถํ วทามิ. สุณาตุ เม ภนฺเต สํโฆ, อยํ อุทายี ภิกฺขุ เอกํ อาปตฺติํ อาปชฺชิ สญฺเจตนิกํ สุกฺก\-วิสฺสฏฺฐิํ อปฺปฏิจฺฉนฺนํ, โส สํฆํ เอกิสฺสา อาปตฺติยา สญฺเจตนิกาย สุกฺก\-วิสฺสฏฺฐิยา อปฺปฏิจฺฉนฺนาย ฉารตฺตํ มานตฺตํ ยาจติ, สํโฆ อุทายิสฺส ภิกฺขุโน เอกิสฺสา อาปตฺติยา สญฺเจตนิกาย สุกฺก\-วิสฺสฏฺฐิยา อปฺปฏิจฺฉนฺนาย ฉารตฺตํ มานตฺตํ เทติ, ยสฺสายสฺมโต ขมติ อุทายิสฺส ภิกฺขุโน เอกิสฺสา อาปตฺติยา สญฺเจตนิกาย สุกฺก\-วิสฺสฏฺฐิยา อปฺปฏิจฺฉนฺนาย ฉารตฺตํ มานตฺตสฺส ทานํ, โส ตุณฺหสฺส, ยสฺส นกฺขมติ, โส ภาเสยฺย.}

\prose{ตติยมฺปิ เอตมตฺถํ วทามิ. สุณาตุ เม ภนฺเต สํโฆ, อยํ อุทายี ภิกฺขุ เอกํ อาปตฺติํ อาปชฺชิ สญฺเจตนิกํ สุกฺก\-วิสฺสฏฺฐิํ อปฺปฏิจฺฉนฺนํ, โส สํฆํ เอกิสฺสา อาปตฺติยา สญฺเจตนิกาย สุกฺก\-วิสฺสฏฺฐิยา อปฺปฏิจฺฉนฺนาย ฉารตฺตํ มานตฺตํ ยาจติ, สํโฆ อุทายิสฺส ภิกฺขุโน เอกิสฺสา อาปตฺติยา สญฺเจตนิกาย สุกฺก\-วิสฺสฏฺฐิยา อปฺปฏิจฺฉนฺนาย ฉารตฺตํ มานตฺตํ เทติ, ยสฺสายสฺมโต ขมติ อุทายิสฺส ภิกฺขุโน เอกิสฺสา อาปตฺติยา สญฺเจตนิกาย สุกฺก\-วิสฺสฏฺฐิยา อปฺปฏิจฺฉนฺนาย ฉารตฺตํ มานตฺตสฺส ทานํ, โส ตุณฺหสฺส, ยสฺส นกฺขมติ, โส ภาเสยฺย.}

\prose{ทินฺนํ สํเฆน อุทายิสฺส ภิกฺขุโน เอกิสฺสา อาปตฺติยา สญฺเจตนิกาย สุกฺก\-วิสฺสฏฺฐิยา อปฺปฏิจฺฉนฺนาย ฉารตฺตํ มานตฺตํ, ขมติ สํฆสฺส, ตสฺมา ตุณฺหี, เอวเมตํ ธารยามี”ติ.}

\csromanmatikaanchor{mtk.62}
\csromantocmarkref{subsection}{อปฺปฏิจฺฉนฺนอพฺภาน}{mtk.62}
\bookmark[dest={mtk.62},level=2]{อปฺปฏิจฺฉนฺนอพฺภาน}
\csromanheader{2}{อปฺปฏิจฺฉนฺน\-อพฺภาน}
\proseitem{100}{\csromanfolio{106}โส จิณฺณมานตฺโต ภิกฺขูนํ อาโรเจสิ “อหํ อาวุโส เอกํ อาปตฺติํ อาปชฺชิ สญฺเจตนิกํ สุกฺก\-วิสฺสฏฺฐิํ อปฺปฏิจฺฉนฺนํ, โสหํ สํฆํ เอกิสฺสา อาปตฺติยา สญฺเจตนิกาย สุกฺก\-วิสฺสฏฺฐิยา อปฺปฏิจฺฉนฺนาย ฉารตฺตํ มานตฺตํ ยาจิํ, ตสฺส เม สํโฆ เอกิสฺสา อาปตฺติยา สญฺเจตนิกาย สุกฺก\-วิสฺสฏฺฐิยา อปฺปฏิจฺฉนฺนาย ฉารตฺตํ มานตฺตํ อทาสิ, โสหํ จิณฺณมานตฺโต, กถํ นุ โข มยา ปฏิปชฺชิตพฺพนฺ”ติ. ภควโต เอตมตฺถํ อาโรเจสุํ. เตน หิ ภิกฺขเว สํโฆ อุทายิํ ภิกฺขุํ อพฺเภตุ. เอวญฺจ ปน ภิกฺขเว อพฺเภตพฺโพ\csromandash{}}

\prose{เตน ภิกฺขเว อุทายินา ภิกฺขุนา สํฆํ อุปสงฺกมิตฺวา เอกํสํ อุตฺตราสงฺคํ กริตฺวา วุฑฺฒานํ ภิกฺขูนํ ปาเท วนฺทิตฺวา อุกฺกุฏิกํ นิสีทิตฺวา อญฺชลิํ ปคฺคเหตฺวา เอวมสฺส วจนีโย “อหํ ภนฺเต เอกํ อาปตฺติํ อาปชฺชิ สญฺเจตนิกํ สุกฺก\-วิสฺสฏฺฐิํ อปฺปฏิจฺฉนฺนํ, โสหํ สํฆํ เอกิสฺสา อาปตฺติยา สญฺเจตนิกาย สุกฺก\-วิสฺสฏฺฐิยา อปฺปฏิจฺฉนฺนาย ฉารตฺตํ มานตฺตํ ยาจิํ, ตสฺส เม สํโฆ เอกิสฺสา อาปตฺติยา สญฺเจตนิกาย สุกฺก\-วิสฺสฏฺฐิยา อปฺปฏิจฺฉนฺนาย ฉารตฺตํ มานตฺตํ อทาสิ, โสหํ ภนฺเต จิณฺณมานตฺโต สํฆํ อพฺภานํ ยาจามิ.}

\prose{อหํ ภนฺเต เอกํ อาปตฺติํ อาปชฺชิํ สญฺเจตนิกํ สุกฺก\-วิสฺสฏฺฐิํ อปฺปฏิจฺฉนฺนํ โสหํ สํฆํ เอกิสฺสา อาปตฺติยา สญฺเจตนิกาย สุกฺก\-วิสฺสฏฺฐิยา อปฺปฏิจฺฉนฺนาย ฉารตฺตํ มานตฺตํ ยาจิํ, ตสฺส เม สํโฆ เอกิสฺสา อาปตฺติยา สญฺเจตนิกาย สุกฺก\-วิสฺสฏฺฐิยา อปฺปฏิจฺฉนฺนาย ฉารตฺตํ มานตฺตํ อทาสิ, โสหํ จิณฺณมานตฺโต ทุติยมฺปิ ภนฺเต สํฆํ อพฺภานํ ยาจามิ.}

\prose{อหํ ภนฺเต เอกํ อาปตฺติํ อาปชฺชิํ สญฺเจตนิกํ สุกฺก\-วิสฺสฏฺฐิํ อปฺปฏิจฺฉนฺนํ, โสหํ สํฆํ เอกิสฺสา อาปตฺติยา สญฺเจตนิกาย สุกฺก\-วิสฺสฏฺฐิยา อปฺปฏิจฺฉนฺนาย ฉารตฺตํ มานตฺตํ ยาจิํ, ตสฺส เม สํโฆ เอกิสฺสา อาปตฺติยา สญฺเจตนิกาย สุกฺก\-วิสฺสฏฺฐิยา อปฺปฏิจฺฉนฺนาย ฉารตฺตํ มานตฺตํ อทาสิ, โสหํ จิณฺณมานตฺโต ตติยมฺปิ ภนฺเต สํฆํ อพฺภานํ ยาจามี”ติ. พฺยตฺเตน ภิกฺขุนา ปฏิพเลน สํโฆ ญาเปตพฺโพ\csromandash{}}

\proseitem{101}{“สุณาตุ เม ภนฺเต สํโฆ, อยํ อุทายี ภิกฺขุ เอกํ อาปตฺติํ อาปชฺชิ สญฺเจตนิกํ สุกฺก\-วิสฺสฏฺฐิํ อปฺปฏิจฺฉนฺนํ, โส สํฆํ เอกิสฺสา \csromanfolio{107}อาปตฺติยา สญฺเจตนิกาย สุกฺก\-วิสฺสฏฺฐิยา อปฺปฏิจฺฉนฺนาย ฉารตฺตํ มานตฺตํ ยาจิ, สํโฆ อุทายิสฺส ภิกฺขุโน เอกิสฺสา อาปตฺติยา สญฺเจตนิกายา สุกฺก\-วิสฺสฏฺฐิยา อปฺปฏิจฺฉนฺนาย ฉารตฺตํ มานตฺตํ อทาสิ, โส จิณฺณมานตฺโต สํฆํ อพฺภานํ ยาจติ, ยทิ สํฆสฺส ปตฺตกลฺลํ, สํโฆ อุทายิํ ภิกฺขุํ อพฺเภยฺย, เอสา ญตฺติ.}

\prose{สุณาตุ เม ภนฺเต สํโฆ, อยํ อุทายี ภิกฺขุ เอกํ อาปตฺติํ อาปชฺชิ สญฺเจตนิกํ สุกฺก\-วิสฺสฏฺฐิํ อปฺปฏิจฺฉนฺนํ, โส สํฆํ เอกิสฺสา อาปตฺติยา สญฺเจตนิกาย สุกฺก\-วิสฺสฏฺฐิยา อปฺปฏิจฺฉนฺนาย ฉารตฺตํ มานตฺตํ ยาจิ, สํโฆ อุทายิสฺส ภิกฺขุโน เอกิสฺสา อาปตฺติยา สญฺเจตนิกาย สุกฺก\-วิสฺสฏฺฐิยา อปฺปฏิจฺฉนฺนาย ฉารตฺตํ มานตฺตํ อทาสิ, โส จิณฺณมานตฺโต สํฆํ อพฺภานํ ยาจติ, สํโฆ อุทายิํ ภิกฺขุํ อพฺเภติ, ยสฺสายสฺมโต ขมติ อุทายิสฺส ภิกฺขุโน อพฺภานํ, โส ตุณฺหสฺส, ยสฺส นกฺขมติ, โส ภาเสยฺย.}

\prose{ทุติยมฺปิ เอตมตฺถํ วทามิ. สุณาตุ เม ภนฺเต สํโฆ, อยํ อุทายี ภิกฺขุ เอกํ อาปตฺติํ อาปชฺชิ สญฺเจตนิกํ สุกฺก\-วิสฺสฏฺฐิํ อปฺปฏิจฺฉนฺนํ, โส สํฆํ เอกิสฺสา อาปตฺติยา สญฺเจตนิกาย สุกฺก\-วิสฺสฏฺฐิยา อปฺปฏิจฺฉนฺนาย ฉารตฺตํ มานตฺตํ ยาจิ, สํโฆ อุทายิสฺส ภิกฺขุโน เอกิสฺสา อาปตฺติยา สญฺเจตนิกาย สุกฺก\-วิสฺสฏฺฐิยา อปฺปฏิจฺฉนฺนาย ฉารตฺตํ มานตฺตํ อทาสิ, โส จิณฺณมานตฺโต สํฆํ อพฺภานํ ยาจติ, สํโฆ อุทายิํ ภิกฺขุํ อพฺเภติ, ยสฺสายสฺมโต ขมติ อุทายิสฺส ภิกฺขุโน อพฺภานํ, โส ตุณฺหสฺส, ยสฺส นกฺขมติ, โส ภาเสยฺย.}

\prose{ตติยมฺปิ เอตมตฺถํ วทามิ. สุณาตุ เม ภนฺเต สํโฆ, อยํ อุทายี ภิกฺขุ เอกํ อาปตฺติํ อาปชฺชิ สญฺเจตนิกํ สุกฺก\-วิสฺสฏฺฐิํ อปฺปฏิจฺฉนฺนํ, โส สํฆํ เอกิสฺสา อาปตฺติยา สญฺเจตนิกาย สุกฺก\-วิสฺสฏฺฐิยา อปฺปฏิจฺฉนฺนาย ฉารตฺตํ มานตฺตํ ยาจิ, สํโฆ อุทายิสฺส ภิกฺขุโน เอกิสฺสา อาปตฺติยา สญฺเจตนิกาย สุกฺก\-วิสฺสฏฺฐิยา อปฺปฏิจฺฉนฺนาย ฉารตฺตํ มานตฺตํ อทาสิ, โส จิณฺณมานตฺโต สํฆํ อพฺภานํ ยาจติ, สํโฆ อุทายิํ ภิกฺขุํ อพฺเภติ, ยสฺสายสฺมโต ขมติ อุทายิสฺส ภิกฺขุโน อพฺภานํ, โส ตุณฺหสฺส, ยสฺส นกฺขมติ, โส ภาเสยฺย.}

\prose{\csromanfolio{108}อพฺภิโต สํเฆน อุทายี ภิกฺขุ, ขมติ สํฆสฺส, ตสฺมา ตุณฺหี เอวเมตํ ธารยามี”ติ.}

\csromanmatikaanchor{mtk.63}
\csromantocmarkref{subsection}{เอกาหปฺปฏิจฺฉนฺนปริวาส}{mtk.63}
\bookmark[dest={mtk.63},level=2]{เอกาหปฺปฏิจฺฉนฺนปริวาส}
\csromanheader{2}{\textbf{เอกาหปฺปฏิจฺฉนฺนปริวาส}}
\proseitem{102}{เตน โข ปน สมเยน อายสฺมา อุทายี เอกํ อาปตฺติํ อาปนฺโน โหติ สญฺเจตนิกํ สุกฺก\-วิสฺสฏฺฐิํ เอกาห\-ปฺปฏิจฺฉนฺนํ, โส ภิกฺขูนํ อาโรเจสิ “อหํ อาวุโส เอกํ อาปตฺติํ อาปชฺชิํ สญฺเจตนิกํ สุกฺก\-วิสฺสฏฺฐิํ เอกาห\-ปฺปฏิจฺฉนฺนํ, กถํ นุ โข มยา ปฏิปชฺชิตพฺพนฺ”ติ. ภควโต เอตมตฺถํ อาโรเจสุํ. เตน หิ ภิกฺขเว สํโฆ อุทายิสฺส ภิกฺขุโน เอกิสฺสา อาปตฺติยา สญฺเจตนิกาย สุกฺก\-วิสฺสฏฺฐิยา เอกาห\-ปฺปฏิจฺฉนฺนาย เอกาหปริวาสํ เทตุ. เอวญฺจ ปน ภิกฺขเว ทาตพฺโพ\csromandash{}}

\prose{เตน ภิกฺขเว อุทายินา ภิกฺขุนา สํฆํ อุปสงฺกมิตฺวา เอกํสํ อุตฺตราสงฺคํ กริตฺวา วุฑฺฒานํ ภิกฺขูนํ ปาเท วนฺทิตฺวา อุกฺกุฏิกํ นิสีทิตฺวา อญฺชลิํ ปคฺคเหตฺวา เอวมสฺส วจนีโย “อหํ ภนฺเต เอกํ อาปตฺติํ อาปชฺชิํ สญฺเจตนิกํ สุกฺก\-วิสฺสฏฺฐิํ เอกาห\-ปฺปฏิจฺฉนฺนํ, โสหํ ภนฺเต สํฆํ เอกิสฺสา อาปตฺติยา สญฺเจตนิกาย สุกฺก\-วิสฺสฏฺฐิยา เอกาห\-ปฺปฏิจฺฉนฺนาย เอกาหปริวาสํ ยาจามี”ติ. ทุติยมฺปิ ยาจิตพฺโพ. ตติยมฺปิ ยาจิตพฺโพ. พฺยตฺเตน ภิกฺขุนา ปฏิพเลน สํโฆ ญาเปตพฺโพ\csromandash{}}

\proseitem{103}{“สุณาตุ เม ภนฺเต สํโฆ, อยํ อุทายี ภิกฺขุ เอกํ อาปตฺติํ อาปชฺชิ สญฺเจตนิกํ สุกฺก\-วิสฺสฏฺฐิํ เอกาห\-ปฺปฏิจฺฉนฺนํ, โส สํฆํ เอกิสฺสา อาปตฺติยา สญฺเจตนิกาย สุกฺก\-วิสฺสฏฺฐิยา เอกาห\-ปฺปฏิจฺฉนฺนาย เอกาหปริวาสํ ยาจติ, ยทิ สํฆสฺส ปตฺตกลฺลํ, สํโฆ อุทายิสฺส ภิกฺขุโน เอกิสฺสา อาปตฺติยา สญฺเจตนิกาย สุกฺก\-วิสฺสฏฺฐิยา เอกาห\-ปฺปฏิจฺฉนฺนาย เอกาหปริวาสํ ทเทยฺย, เอสา ญตฺติ.}

\prose{สุณาตุ เม ภนฺเต สํโฆ, อยํ อุทายี ภิกฺขุ เอกํ อาปตฺติํ อาปชฺชิ สญฺเจตนิกํ สุกฺก\-วิสฺสฏฺฐิํ เอกาห\-ปฺปฏิจฺฉนฺนํ, โส สํฆํ เอกิสฺสา อาปตฺติยา สญฺเจตนิกาย สุกฺก\-วิสฺสฏฺฐิยา เอกาห\-ปฺปฏิจฺฉนฺนาย เอกาหปริวาสํ ยาจติ, สํโฆ อุทายิสฺส ภิกฺขุโน เอกิสฺสา อาปตฺติยา สญฺเจตนิกาย สุกฺก\-วิสฺสฏฺฐิยา เอกาห\-ปฺปฏิจฺฉนฺนาย เอกาหปริวาสํ \csromanfolio{109}เทหิ, ยสฺสายสฺมโต ขมติ อุทายิสฺส ภิกฺขุโน เอกิสฺสา อาปตฺติยา สญฺเจตนิกาย สุกฺก\-วิสฺสฏฺฐิยา เอกาห\-ปฺปฏิจฺฉนฺนาย เอกาห\-ปริวาสสฺส ทานํ, โส ตุณฺหสฺส, ยสฺส นกฺขมติ, โส ภาเสยฺย.}

\prose{ทุติยมฺปิ เอตมตฺถํ วทามิ ฯเปฯ. ตติยมฺปิ เอตมตฺถํ วทามิ ฯเปฯ.}

\prose{ทินฺโน สํเฆน อุทายิสฺส ภิกฺขุโน เอกิสฺสา อาปตฺติยา สญฺเจตนิกาย สุกฺก\-วิสฺสฏฺฐิยา เอกาห\-ปฺปฏิจฺฉนฺนาย เอกาหปริวาโส, ขมติ สํฆสฺส, ตสฺมา ตุณฺหี, เอวเมตํ ธารยามี”ติ.}

\csromanmatikaanchor{mtk.64}
\csromantocmarkref{subsection}{เอกาหปฺปฏิจฺฉนฺนมานตฺต}{mtk.64}
\bookmark[dest={mtk.64},level=2]{เอกาหปฺปฏิจฺฉนฺนมานตฺต}
\csromanheader{2}{เอกาหปฺปฏิจฺฉนฺน\-มานตฺต}
\proseitem{104}{โส ปริวุตฺถ\-ปริวาโส ภิกฺขูนํ อาโรเจสิ “อหํ อาวุโส เอกํ อาปตฺติํ อาปชฺชิํ สญฺเจตนิกํ สุกฺก\-วิสฺสฏฺฐิํ เอกาห\-ปฺปฏิจฺฉนฺนํ, โสหํ สํฆํ เอกิสฺสา อาปตฺติยา สญฺเจตนิกาย สุกฺก\-วิสฺสฏฺฐิยา เอกาห\-ปฺปฏิจฺฉนฺนาย เอกาหปริวาสํ ยาจิํ, ตสฺส เม สํโฆ เอกิสฺสา อาปตฺติยา สญฺเจตนิกาย สุกฺก\-วิสฺสฏฺฐิยา เอกาห\-ปฺปฏิจฺฉนฺนาย เอกาหปริวาสํ อทาสิ, โสหํ ปริวุตฺถ\-ปริวาโส, กถํ นุ โข มยา ปฏิปชฺชิตพฺพนฺ”ติ. ภควโต เอตมตฺถํ อาโรเจสุํ. เตน หิ ภิกฺขเว สํโฆ อุทายิสฺส ภิกฺขุโน เอกิสฺสา อาปตฺติยา สญฺเจตนิกาย สุกฺก\-วิสฺสฏฺฐิยา เอกาห\-ปฺปฏิจฺฉนฺนาย ฉารตฺตํ มานตฺตํ เทตุ. เอวญฺจ ปน ภิกฺขเว ทาตพฺพํ\csromandash{}}

\prose{เตน ภิกฺขเว อุทายินา ภิกฺขุนา สํฆํ อุปสงฺกมิตฺวา ฯเปฯ เอวมสฺส วจนีโย “อหํ ภนฺเต เอกํ อาปตฺติํ อาปชฺชิํ สญฺเจตนิกํ สุกฺก\-วิสฺสฏฺฐิํ เอกาห\-ปฺปฏิจฺฉนฺนํ, โสหํ สํฆํ เอกิสฺสา อาปตฺติยา สญฺเจตนิกาย สุกฺก\-วิสฺสฏฺฐิยา เอกาห\-ปฺปฏิจฺฉนฺนาย เอกาหปริวาสํ ยาจิํ, ตสฺส เม สํโฆ เอกิสฺสา อาปตฺติยา สญฺเจตนิกาย สุกฺก\-วิสฺสฏฺฐิยา เอกาห\-ปฺปฏิจฺฉนฺนาย เอกาหปริวาสํ อทาสิ, โสหํ ภนฺเต ปริวุตฺถ\-ปริวาโส สํฆํ เอกิสฺสา อาปตฺติยา สญฺเจตนิกาย สุกฺก\-วิสฺสฏฺฐิยา เอกาห\-ปฺปฏิจฺฉนฺนาย ฉารตฺตํ มานตฺตํ ยาจามี”ติ. ทุติยมฺปิ ยาจิตพฺพํ\csromansharedfootnote{aa4994d44f0a54f7}{ยาจิตพฺโพ (สี, เอวมุปริปิ)}. ตติยมฺปิ ยาจิตพฺพํ\csromansharedfootnote{aa4994d44f0a54f7}{ยาจิตพฺโพ (สี, เอวมุปริปิ)}. พฺยตฺเตน ภิกฺขุนา ปฏิพเลน สํโฆ ญาเปตพฺโพ\csromandash{}}

\proseitem{105}{\csromanfolio{110}“สุณาตุ เม ภนฺเต สํโฆ, อยํ อุทายี ภิกฺขุ เอกํ อาปตฺติํ อาปชฺชิ สญฺเจตนิกํ สุกฺก\-วิสฺสฏฺฐิํ เอกาห\-ปฺปฏิจฺฉนฺนํ, โส สํฆํ เอกิสฺสา อาปตฺติยา สญฺเจตนิกายา สุกฺก\-วิสฺสฏฺฐิยา เอกาห\-ปฺปฏิจฺฉนฺนาย เอกาหปริวาสํ ยาจิ, สํโฆ อุทายิสฺส ภิกฺขุโน เอกิสฺสา อาปตฺติยา สญฺเจตนิกาย สุกฺก\-วิสฺสฏฺฐิยา เอกาห\-ปฺปฏิจฺฉนฺนาย เอกาหปริวาสํ อทาสิ, โส ปริวุตฺถ\-ปริวาโส สํฆํ เอกิสฺสา อาปตฺติยา สญฺเจตนิกาย สุกฺก\-วิสฺสฏฺฐิยา เอกาห\-ปฺปฏิจฺฉนฺนาย ฉารตฺตํ มานตฺตํ ยาจติ, ยทิ สํฆสฺส ปตฺตกลฺลํ, สํโฆ อุทายิสฺส ภิกฺขุโน เอกิสฺสา อาปตฺติยา สญฺเจตนิกาย สุกฺก\-วิสฺสฏฺฐิยา เอกาห\-ปฺปฏิจฺฉนฺนาย ฉารตฺตํ มานตฺตํ ทเทยฺย, เอสา ญตฺติ.}

\prose{สุณาตุ เม ภนฺเต สํโฆ, อยํ อุทายี ภิกฺขุ เอกํ อาปตฺติํ อาปชฺชิ สญฺเจตนิกํ สุกฺก\-วิสฺสฏฺฐิํ เอกาห\-ปฺปฏิจฺฉนฺนํ, โส สํฆํ เอกิสฺสา อาปตฺติยา สญฺเจตนิกาย สุกฺก\-วิสฺสฏฺฐิยา เอกาห\-ปฺปฏิจฺฉนฺนาย เอกาหปริวาสํ ยาจิ, สํโฆ อุทายิสฺส ภิกฺขุโน เอกิสฺสา อาปตฺติยา สญฺเจตนิกาย สุกฺก\-วิสฺสฏฺฐิยา เอกาห\-ปฺปฏิจฺฉนฺนาย เอกาหปริวาสํ อทาสิ, โส ปริวุตฺถ\-ปริวาโส สํฆํ เอกิสฺสา อาปตฺติยา สญฺเจตนิกาย สุกฺก\-วิสฺสฏฺฐิยา เอกาห\-ปฺปฏิจฺฉนฺนาย ฉารตฺตํ มานตฺตํ ยาจติ, สํโฆ อุทายิสฺส ภิกฺขุโน เอกิสฺสา อาปตฺติยา สญฺเจตนิกาย สุกฺก\-วิสฺสฏฺฐิยา เอกาห\-ปฺปฏิจฺฉนฺนาย ฉารตฺตํ มานตฺตํ เทติ, ยสฺสายสฺมโต ขมติ อุทายิสฺส ภิกฺขุโน เอกิสฺสา อาปตฺติยา สญฺเจตนิกาย สุกฺก\-วิสฺสฏฺฐิยา เอกาห\-ปฺปฏิจฺฉนฺนาย ฉารตฺต มานตฺตสฺส ทานํ, โส ตุณฺหสฺส, ยสฺส นกฺขมติ, โส ภาเสยฺย.}

\prose{ทุติยมฺปิ เอตมตฺตํ วทามิ ฯเปฯ. ตติยมฺปิ เอตมตฺตํ วทามิ ฯเปฯ.}

\prose{ทินฺนํ สํเฆน อุทายิสฺส ภิกฺขุโน เอกิสฺสา อาปตฺติยา สญฺเจตนิกาย สุกฺก\-วิสฺสฏฺฐิยา เอกาห\-ปฺปฏิจฺฉนฺนาย ฉารตฺตํ มานตฺตํ, ขมติ สํฆสฺส, ตสฺมา ตุณฺหี, เอวเมตํ ธารยามี”ติ.}

\csromanmatikaanchor{mtk.65}
\csromantocmarkref{subsection}{เอกาหปฺปฏิจฺฉนฺนอพฺภาน}{mtk.65}
\bookmark[dest={mtk.65},level=2]{เอกาหปฺปฏิจฺฉนฺนอพฺภาน}
\csromanheader{2}{เอกาหปฺปฏิจฺฉนฺน\-อพฺภาน}
\proseitem{106}{โส จิณฺณมานตฺโต ภิกฺขูนํ อาโรเจสิ “อหํ อาวุโส เอกํ อาปตฺติํ อาปชฺชิํ สญฺเจตนิกํ สุกฺก\-วิสฺสฏฺฐิํ เอกาหปฺปฏิจฺฉนํ, \csromanfolio{111}โสหํ สํฆํ เอกิสฺสา อาปตฺติยา สญฺเจตนิกาย สุกฺก\-วิสฺสฏฺฐิยา เอกาห\-ปฺปฏิจฺฉนฺนาย เอกาหปริวาสํ ยาจิํ, ตสฺส เม สํโฆ เอกิสฺสา อาปตฺติยา สญฺเจตนิกาย สุกฺก\-วิสฺสฏฺฐิยา เอกาห\-ปฺปฏิจฺฉนฺนาย เอกาหปริวาสํ อทาสิ, โสหํ ปริวุตฺถ\-ปริวาโส สํฆํ เอกิสฺสา อาปตฺติยา สญฺเจตนิกาย สุกฺก\-วิสฺสฏฺฐิยา เอกาห\-ปฺปฏิจฺฉนฺนาย ฉารตฺตํ มานตฺตํ ยาจิํ, ตสฺส เม สํโฆ เอกิสฺสา อาปตฺติยา สญฺเจตนิกาย สุกฺก\-วิสฺสฏฺฐิยา เอกาห\-ปฺปฏิจฺฉนฺนาย ฉารตฺตํ มานตฺตํ อทาสิ, โสหํ จิณฺณมานตฺโต, กถํ นุ โข มยา ปฏิปชฺชิตพฺพนฺ”ติ. ภควโต เอตมตฺถํ อาโรเจสุํ. เตน หิ ภิกฺขเว สํโฆ อุทายิํ ภิกฺขุํ อพฺเภตุ. เอวญฺจ ปน ภิกฺขเว อพฺเภตพฺโพ\csromandash{}}

\prose{เตน ภิกฺขเว อุทายินา ภิกฺขุนา สํฆํ อุปสงฺกมิตฺวา ฯเปฯ เอวมสฺส วจนีโย “อหํ ภนฺเต เอกํ อาปตฺติํ อาปชฺชิํ สญฺเจตนิกํ สุกฺก\-วิสฺสฏฺฐิํ เอกาห\-ปฺปฏิจฺฉนฺนํ, โสหํ สํฆํ เอกิสฺสา อาปตฺติยา สญฺเจตนิกาย สุกฺก\-วิสฺสฏฺฐิยา เอกาห\-ปฺปฏิจฺฉนฺนาย เอกาหปริวาสํ ยาจิํ, ตสฺส เม สํโฆ เอกิสฺสา อาปตฺติยา สญฺเจตนิกาย สุกฺก\-วิสฺสฏฺฐิยา เอกาห\-ปฺปฏิจฺฉนฺนาย เอกาหปริวาสํ อทาสิ, โสหํ ปริวุตฺถ\-ปริวาโส สํฆํ เอกิสฺสา อาปตฺติยา สญฺเจตนิกาย สุกฺก\-วิสฺสฏฺฐิยา เอกาห\-ปฺปฏิจฺฉนฺนาย ฉารตฺตํ มานตฺตํ ยาจิํ, ตสฺส เม สํโฆ เอกิสฺสา อาปตฺติยา สญฺเจตนิกาย สุกฺก\-วิสฺสฏฺฐิยา เอกาห\-ปฺปฏิจฺฉนฺนาย ฉารตฺตํ มานตฺตํ อทาสิ, โสหํ ภนฺเต จิณฺณมานตฺโต สํฆํ อพฺภานํ ยาจามี”ติ. ทุติยมฺปิ ยาจิตพฺพํ. ตติยมฺปิ ยาจิตพฺพํ. พฺยตฺเตน ภิกฺขุนา ปฏิพเลน สํโฆ ญาเปตพฺโพ\csromandash{}}

\proseitem{107}{“สุณาตุ เม ภนฺเต สํโฆ, อยํ อุทายี ภิกฺขุ เอกํ อาปตฺติํ อาปชฺชิ สญฺเจตนิกํ สุกฺกวิสฺสฏฺฐิ เอกาห\-ปฺปฏิจฺฉนฺนํ, โส สํฆํ เอกิสฺสา อาปตฺติยา สญฺเจตนิกาย สุกฺก\-วิสฺสฏฺฐิยา เอกาห\-ปฺปฏิจฺฉนฺนาย เอกาหปริวาสํ ยาจิ, สํโฆ อุทายิสฺส ภิกฺขุโน เอกิสฺสา อาปตฺติยา สญฺเจตนิกาย สุกฺก\-วิสฺสฏฺฐิยา เอกาห\-ปฺปฏิจฺฉนฺนาย เอกาหปริวาสํ อทาสิ, โส ปริวุตฺถ\-ปริวาโส สํฆํ เอกิสฺสา อาปตฺติยา สญฺเจตนิกาย สุกฺก\-วิสฺสฏฺฐิยา เอกาห\-ปฺปฏิจฺฉนฺนาย ฉารตฺตํ มานตฺตํ ยาจิ, สํโฆ อุทายิสฺส ภิกฺขุโน เอกิสฺสา อาปตฺติยา สญฺเจตนิกาย สุกฺก\-วิสฺสฏฺฐิยา เอกาห\-ปฺปฏิจฺฉนฺนาย ฉารตฺตํ มานตฺตํ อทาสิ, โส \csromanfolio{112}จิณฺณมานตฺโต สํฆํ อพฺภานํ ยาจติ, ยทิ สํฆสฺส ปตฺตกลฺลํ, สํโฆ อุทายิํ ภิกฺขุํ อพฺเภยฺย, เอสา ญตฺติ.}

\prose{สุณาตุ เม ภนฺเต สํโฆ, อยํ อุทายี ภิกฺขุ เอกํ อาปตฺติํ อาปชฺชิ สญฺเจตนิกํ สุกฺก\-วิสฺสฏฺฐิํ เอกาห\-ปฺปฏิจฺฉนฺนํ, โส สํฆํ เอกิสฺสา อาปตฺติยา สญฺเจตนิกาย สุกฺก\-วิสฺสฏฺฐิยา เอกาห\-ปฺปฏิจฺฉนฺนาย เอกาหปริวาสํ ยาจิ, สํโฆ อุทายิสฺส ภิกฺขุโน เอกิสฺสา อาปตฺติยา สญฺเจตนิกาย สุกฺก\-วิสฺสฏฺฐิยา เอกาห\-ปฺปฏิจฺฉนฺนาย เอกาหปริวาสํ อทาสิ, โส ปริวุตฺถ\-ปริวาโส สํฆํ เอกิสฺสา อาปตฺติยา สญฺเจตนิกาย สุกฺก\-วิสฺสฏฺฐิยา เอกาห\-ปฺปฏิจฺฉนฺนาย ฉารตฺตํ มานตฺตํ ยาจิ, สํโฆ อุทายิสฺส ภิกฺขุโน เอติสฺสา อาปตฺติยา สญฺเจตนิกาย สุกฺก\-วิสฺสฏฺฐิยา เอกาห\-ปฺปฏิจฺฉนฺนาย ฉารตฺตํ มานตฺตํ อทาสิ, โส จิณฺณมานตฺโต สํฆํ อพฺภานํ ยาจติ สํโฆ อุทายิํ ภิกฺขุํ อพฺเภติ, ยสฺสายสฺมโต ขมติ อุทายิสฺส ภิกฺขุโน อพฺภานํ, โส ตุณฺหสฺส, ยสฺส นกฺขมติ, โส ภาเสยฺย.}

\prose{ทุติยมฺปิ เอตมตฺถํ วทามิ ฯเปฯ. ตติยมฺปิ เอตมตฺถํ วทามิ ฯเปฯ.}

\prose{อพฺภิโต สํเฆน อุทายี ภิกฺขุ, ขมติ สํฆสฺส, ตสฺมา ตุณฺหี เอวเมตํ ธารยามี”ติ.}

\csromanmatikaanchor{mtk.66}
\csromantocmarkref{subsection}{ปญฺจาหปฺปฏิจฺฉนฺนปริวาส}{mtk.66}
\bookmark[dest={mtk.66},level=2]{ปญฺจาหปฺปฏิจฺฉนฺนปริวาส}
\csromanheader{2}{ปญฺจาหปฺปฏิจฺฉนฺน\-ปริวาส}
\proseitem{108}{เตน โข ปน สมเยน อายสฺมา อุทายี เอกํ อาปตฺติํ อาปนฺโน โหติ สญฺเจตนิกํ สุกฺก\-วิสฺสฏฺฐิํ ทฺวีห\-ปฺปฏิจฺฉนฺนํ ฯเปฯ ตีหปฺปฏิจฺฉนฺนํ ฯเปฯ จตูห\-ปฺปฏิจฺฉนฺนํ ฯเปฯ ปญฺจาห\-ปฺปฏิจฺฉนฺนํ. โส ภิกฺขูนํ อาโรเจสิ “อหํ อาวุโส เอกํ อาปตฺติํ อาปชฺชิํ สญฺเจตนิกํ สุกฺก\-วิสฺสฏฺฐิํ ปญฺจาห\-ปฺปฏิจฺฉนฺนํ, กถํ นุ โข มยา ปฏิปชฺชิตพฺพนฺ”ติ. ภควโต เอตมตฺถํ อาโรเจสุํ. เตน หิ ภิกฺขเว สํโฆ อุทายิสฺส ภิกฺขุโน เอกิสฺสา อาปตฺติยา สญฺเจตนิกาย สุกฺก\-วิสฺสฏฺฐิยา ปญฺจาห\-ปฺปฏิจฺฉนฺนาย ปญฺจาห\-ปริวาสํ เทตุ. เอวญฺจ ปน ภิกฺขเว ทาตพฺโพ\csromandash{}}

\csromanmatikaanchor{mtk.67}
\csromantocmarkref{subsection}{ปาริวาสิกมูลายปฏิกสฺสนา}{mtk.67}
\bookmark[dest={mtk.67},level=2]{ปาริวาสิกมูลายปฏิกสฺสนา}
\prose{เตน ภิกฺขเว อุทายินา ภิกฺขุนา สํฆํ อุปสงฺกมิตฺวา เอกํสํ อุตฺตราสงฺคํ กริตฺวา วุฑฺฒานํ ภิกฺขูนํ ปาเท วนฺทิตฺวา อุกฺกุฏิกํ นิสีทิตฺวา อญฺชลิํ ปคฺคเหตฺวา เอวมสฺส วจนีโย “อหํ ภนฺเต เอกํ อาปตฺติํ \csromanfolio{113}อาปชฺชิํ สญฺเจตนิกํ สุกฺก\-วิสฺสฏฺฐิํ ปญฺจาห\-ปฺปฏิจฺฉนฺนํ, โสหํ ภนฺเต สํฆํ เอกิสฺสา อาปตฺติยา สญฺเจตนิกาย สุกฺก\-วิสฺสฏฺฐิยา ปญฺจาห\-ปฺปฏิจฺฉนฺนาย ปญฺจาห\-ปริวาสํ ยาจามี”ติ. ทุติยมฺปิ ยาจิตพฺโพ. ตติยมฺปิ ยาจิตพฺโพ. พฺยตฺเตน ภิกฺขุนา ปฏิพเลน สํโฆ ญาเปตพฺโพ\csromandash{}}

\proseitem{109}{“สุณาตุ เม ภนฺเต สํโฆ, อยํ อุทายี ภิกฺขุ เอกํ อาปตฺติํ อาปชฺชิ สญฺเจตนิกํ สุกฺก\-วิสฺสฏฺฐิํ ปญฺจาห\-ปฺปฏิจฺฉนฺนํ, โส สํฆํ เอกิสฺสา อาปตฺติยา สญฺเจตนิกาย สุกฺก\-วิสฺสฏฺฐิยา ปญฺจาห\-ปฺปฏิจฺฉนฺนาย ปญฺจาห\-ปริวาสํ ยาจติ, ยทิสํฆสฺส ปตฺตกลฺลํ, สํโฆ อุทายิสฺส ภิกฺขุโน เอกิสฺสา อาปตฺติยา สญฺเจตนิกาย สุกฺก\-วิสฺสฏฺฐิยา ปญฺจาห\-ปฺปฏิจฺฉนฺนาย ปญฺจาห\-ปริวาสํ ทเทยฺย, เอสา ญตฺติ.}

\prose{สุณาตุ เม ภนฺเต สํโฆ, อยํ อุทายี ภิกฺขุ เอกํ อาปตฺติํ อาปชฺชิ สญฺเจตนิกํ สุกฺก\-วิสฺสฏฺฐิํ ปญฺจาห\-ปฺปฏิจฺฉนฺนํ, โส สํฆํ เอกิสฺสา อาปตฺติยา สญฺเจตนิกาย สุกฺก\-วิสฺสฏฺฐิยา ปญฺจาห\-ปฺปฏิจฺฉนฺนาย ปญฺจาห\-ปริวาสํ ยาจติ, สํโฆ อุทายิสฺส ภิกฺขุโน เอกิสฺสา อาปตฺติยา สญฺเจตนิกาย สุกฺก\-วิสฺสฏฺฐิยา ปญฺจาห\-ปฺปฏิจฺฉนฺนาย ปญฺจาห\-ปริวาสํ เทติ, ยสฺสายสฺมโต ขมติ อุทายิสฺส ภิกฺขุโน เอกิสฺสา อาปตฺติยา สญฺเจตนิกาย สุกฺก\-วิสฺสฏฺฐิยา ปญฺจาห\-ปฺปฏิจฺฉนฺนาย ปญฺจาห\-ปริวาสสฺส ทานํ, โส ตุณฺหสฺส, ยสฺส นกฺขมติ, โส ภาเสยฺย.}

\prose{ทุติยมฺปิ เอตมตฺถํ วทามิ ฯเปฯ. ตติยมฺปิ เอตมตฺถํ วทามิ ฯเปฯ.}

\prose{ทินฺโน สํเฆน อุทายิสฺส ภิกฺขุโน เอกิสฺสา อาปตฺติยา สญฺเจตนิกาย สุกฺก\-วิสฺสฏฺฐิยา ปญฺจาห\-ปฺปฏิจฺฉนฺนาย ปญฺจาห\-ปริวาโส, ขมติ สํฆสฺส, ตสฺมา ตุณฺหี, เอวเมตํ ธารยามี”ติ.}

\csromanheader{2}{ปริวา\-สิก\-มูลายปฏิกสฺสนา}
\proseitem{110}{โส ปริวสนฺโต อนฺตรา เอกํ อาปตฺติํ อาปชฺชิ สญฺเจตนิกํ สุกฺก\-วิสฺสฏฺฐิํ อปฺปฏิจฺฉนฺนํ, โส ภิกฺขูนํ อาโรเจสิ “อหํ อาวุโส เอกํ อาปตฺติํ อาปชฺชิํ สญฺเจตนิกํ สุกฺก\-วิสฺสฏฺฐิํ ปญฺจาห\-ปฺปฏิจฺฉนฺนํ, โสหํ สํฆํ เอกิสฺสา อาปตฺติยา สญฺเจตนิกาย สุกฺก\-วิสฺสฏฺฐิยา ปญฺจาห\-ปฺปฏิจฺฉนฺนาย ปญฺจาห\-ปริวาสํ ยาจิํ, ตสฺส เม สํโฆ เอกิสฺสา \csromanfolio{114}อาปตฺติยา สญฺเจตนิกาย สุกฺก\-วิสฺสฏฺฐิยา ปญฺจาห\-ปฺปฏิจฺฉนฺนาย ปญฺจาห\-ปริวาสํ อทาสิ, โสหํ ปริวสนฺโต อนฺตรา เอกํ อาปตฺติํ อาปชฺชิํ สญฺเจตนิกํ สุกฺก\-วิสฺสฏฺฐิํ อปฺปฏิจฺฉนฺนํ, กถํ นุ โข มยา ปฏิปชฺชิตพฺพนฺ”ติ. ภควโต เอตมตฺถํ อาโรเจสุํ. เตน หิ ภิกฺขเว สํโฆ อุทายิํ ภิกฺขุํ อนฺตรา เอกิสฺสา อาปตฺติยา สญฺเจตนิกาย สุกฺก\-วิสฺสฏฺฐิยา อปฺปฏิจฺฉนฺนาย มูลาย ปฏิกสฺสตุ. เอวญฺจ ปน ภิกฺขเว มูลาย ปฏิกสฺสิตพฺโพ\csromandash{}}

\prose{เตน ภิกฺขเว อุทายินา ภิกฺขุนา สํฆํ อุปสงฺกมิตฺวา ฯเปฯ เอวมสฺส วจนีโย “อหํ ภนฺเต เอกํ อาปตฺติํ อาปชฺชิํ สญฺเจตนิกํ สุกฺก\-วิสฺสฏฺฐิํ ปญฺจาห\-ปฺปฏิจฺฉนฺนํ, โสหํ สํฆํ เอกิสฺสา อาปตฺติยา สญฺเจตนิกาย สุกฺก\-วิสฺสฏฺฐิยา ปญฺจาห\-ปฺปฏิจฺฉนฺนาย ปญฺจาห\-ปริวาสํ ยาจิํ, ตสฺส เม สํโฆ เอกิสฺสา อาปตฺติยา สญฺเจตนิกาย สุกฺก\-วิสฺสฏฺฐิยา ปญฺจาห\-ปฺปฏิจฺฉนฺนาย ปญฺจาห\-ปริวาสํ อทาสิ, โสหํ ปริวสนฺโต อนฺตรา เอกํ อาปตฺติํ อาปชฺชิํ สญฺเจตนิกํ สุกฺก\-วิสฺสฏฺฐิํ อปฺปฏิจฺฉนฺนํ, โสหํ ภนฺเต สํฆํ อนฺตรา เอกิสฺสา อาปตฺติยา สญฺเจตนิกาย สุกฺก\-วิสฺสฏฺฐิยา อปฺปฏิจฺฉนฺนาย มูลาย ปฏิกสฺสนํ ยาจามี”ติ. ทุติยมฺปิ ยาจิตพฺพา\csromansharedfootnote{8cb3ba29a5083f18}{ยาจิตพฺโพ (สี, เอวมุปริวิ)}. ตติยมฺปิ ยาจิตพฺพา2. พฺยตฺเตน ภิกฺขุนา ปฏิพเลน สํโฆ ญาเปตพฺโพ\csromandash{}}

\proseitem{111}{“สุณาตุ เม ภนฺเต สํโฆ, อยํ อุทายี ภิกฺขุ เอกํ อาปตฺติํ อาปชฺชิ สญฺเจตนิกํ สุกฺก\-วิสฺสฏฺฐิํ ปญฺจาห\-ปฺปฏิจฺฉนฺนํ, โส สํฆํ เอกิสฺสา อาปตฺติยา สญฺเจตนิกาย สุกฺก\-วิสฺสฏฺฐิยา ปญฺจาห\-ปฺปฏิจฺฉนฺนาย ปญฺจาห\-ปริวาสํ ยาจิ, สํโฆ อุทายิสฺส ภิกฺขุโน เอกิสฺสา อาปตฺติยา สญฺเจตนิกาย สุกฺก\-วิสฺสฏฺฐิยา ปญฺจาห\-ปฺปฏิจฺฉนฺนาย ปญฺจาห\-ปริวาสํ อทาสิ, โส ปริวสนฺโต อนฺตรา เอกํ อาปตฺติํ อาปชฺชิ สญฺเจตนิกํ สุกฺก\-วิสฺสฏฺฐิํ อปฺปฏิจฺฉนฺนํ, โส สํฆํ อนฺตรา เอกิสฺสา อาปตฺติยา สญฺเจตนิกาย สุกฺก\-วิสฺสฏฺฐิยา อปฺปฏิจฺฉนฺนาย มูลาย\-ปฏิกสฺสนํ ยาจติ, ยทิ สํฆสฺส ปตฺตกลฺลํ, สํโฆ อุทายิํ ภิกฺขุํ อนฺตรา เอกิสฺสา อาปตฺติยา สญฺเจตนิกาย สุกฺก\-วิสฺสฏฺฐิยา อปฺปฏิจฺฉนฺนาย มูลาย ปฏิกสฺเสยฺย, เอสา ญตฺติ.}

\prose{\csromanfolio{115}สุณาตุ เม ภนฺเต สํโฆ, อยํ อุทายี ภิกฺขุ เอกํ อาปตฺติํ อาปชฺชิ สญฺเจตนิกํ สุกฺก\-วิสฺสฏฺฐิํ ปญฺจาห\-ปฺปฏิจฺฉนฺนํ, โส สํฆํ เอกิสฺสา อาปตฺติยา สญฺเจตนิกาย สุกฺก\-วิสฺสฏฺฐิยา ปญฺจาห\-ปฺปฏิจฺฉนฺนาย ปญฺจาห\-ปริวาสํ ยาจิ, สํโฆ อุทายิสฺส ภิกฺขุโน เอกิสฺสา อาปตฺติยา สญฺเจตนิกาย สุกฺก\-วิสฺสฏฺฐิยา ปญฺจาห\-ปฺปฏิจฺฉนฺนาย ปญฺจาห\-ปริวาสํ อทาสิ, โส ปริวสนฺโต อนฺตรา เอกํ อาปตฺติํ อาปชฺชิ สญฺเจตนิกํ สุกฺก\-วิสฺสฏฺฐิํ อปฺปฏิจฺฉนฺนํ, โส สํฆํ อนฺตรา เอกิสฺสา อาปตฺติยา สญฺเจตนิกาย สุกฺก\-วิสฺสฏฺฐิยา อปฺปฏิจฺฉนฺนาย มูลาย\-ปฏิกสฺสนํ ยาจติ, สํโฆ อุทายิํ ภิกฺขุํ อนฺตรา เอกิสฺสา อาปตฺติยา สญฺเจตนิกาย สุกฺก\-วิสฺสฏฺฐิยา อปฺปฏิจฺฉนฺนาย มูลาย ปฏิกสฺสติ, ยสฺสายสฺมโต ขมติ อุทายิสฺส ภิกฺขุโน อนฺตรา เอกิสฺสา อาปตฺติยา สญฺเจตนิกาย สุกฺก\-วิสฺสฏฺฐิยา อปฺปฏิจฺฉนฺนาย มูลาย\-ปฏิกสฺสนา, โส ตุณฺหสฺส, ยสฺส นกฺขมติ, โส ภาเสยฺย.}

\prose{ทุติยมฺปิ เอตมตฺถํ วทามิ ฯเปฯ. ตติยมฺปิ เอตมตฺถํ วทามิ ฯเปฯ.}

\prose{ปฏิกสฺสิโต สํเฆน อุทายี ภิกฺขุ อนฺตรา เอกิสฺสา อาปตฺติยา สญฺเจตนิกาย สุกฺก\-วิสฺสฏฺฐิยา อปฺปฏิจฺฉนฺนาย มูลาย\-ปฏิกสฺสนา\csromansharedfootnote{39ac7b401fed4ea4}{มูลาย (สฺยา)}, ขมติ สํฆสฺส, ตสฺมา ตุณฺหี, เอวเมตํ ธารยามี”ติ.}

\csromanmatikaanchor{mtk.68}
\csromantocmarkref{subsection}{มานตฺตารหมูลายปฏิกสฺสนา}{mtk.68}
\bookmark[dest={mtk.68},level=2]{มานตฺตารหมูลายปฏิกสฺสนา}
\csromanheader{2}{มานตฺตารห\-มูลาย\-ปฏิกสฺสนา}
\proseitem{112}{โส ปริวุตฺถ\-ปริวาโส มานตฺตารโห อนฺตรา เอกํ อาปตฺติํ อาปชฺชิ สญฺเจตนิกํ สุกฺก\-วิสฺสฏฺฐิํ อปฺปฏิจฺฉนฺนํ, โส ภิกฺขูนํ อาโรเจสิ “อหํ อาวุโส เอกํ อาปตฺติํ อาปชฺชิํ สญฺเจตนิกํ สุกฺก\-วิสฺสฏฺฐิํ ปญฺจาห\-ปฺปฏิจฺฉนฺนํ โสหํ สํฆํ เอกิสฺสา อาปตฺติยา สญฺเจตนิกาย สุกฺก\-วิสฺสฏฺฐิยา ปญฺจาห\-ปฺปฏิจฺฉนฺนาย ปญฺจาห\-ปริวาสํ ยาจิํ, ตสฺส เม สํโฆ เอกิสฺสา อาปตฺติยา สญฺเจตนิกาย สุกฺก\-วิสฺสฏฺฐิยา ปญฺจาห\-ปฺปฏิจฺฉนฺนาย ปญฺจาห\-ปริวาสํ อทาสิ, โสหํ ปริวสนฺโต อนฺตรา เอกํ อาปตฺติํ อาปชฺชิํ สญฺเจตนิกํ สุกฺก\-วิสฺสฏฺฐิํ อปฺปฏิจฺฉนฺนํ, โสหํ สํฆํ อนฺตรา เอกิสฺสา อาปตฺติยา สญฺเจตนิกาย สุกฺก\-วิสฺสฏฺฐิยา อปฺปฏิจฺฉนฺนาย มูลาย\-ปฏิกสฺสนํ ยาจิํ, ตํ มํ สํโฆ อนฺตรา เอกิสฺสา อาปตฺติยา \csromanfolio{116}สญฺเจตนิกาย สุกฺก\-วิสฺสฏฺฐิยา อปฺปฏิจฺฉนฺนาย มูลาย ปฏิกสฺสิ, โสหํ ปริวุตฺถ\-ปริวาโส มานตฺตารโห อนฺตรา เอกํ อาปตฺติํ อาปชฺชิํ สญฺเจตนิกํ สุกฺก\-วิสฺสฏฺฐิํ อปฺปฏิจฺฉนฺนํ, กถํ นุ โข มายา ปฏิปชฺชิตพฺพนฺ”ติ. ภควโต เอตมตฺถํ อาโรเจสุํ. เตน หิ ภิกฺขเว สํโฆ อุทายิํ ภิกฺขุํ อนฺตรา เอกิสฺสา อาปตฺติยา สญฺเจตนิกาย สุกฺก\-วิสฺสฏฺฐิยา อปฺปฏิจฺฉนฺนาย มูลาย ปฏิกสฺสตุ. เอวญฺจ ปน ภิกฺขเว มูลาย ปฏิกสฺสิตพฺโพ\csromandash{}}

\prose{เตน ภิกฺขเว อุทายินา ภิกฺขุนา สํฆํ อุปสงฺกมิตฺวา ฯเปฯ เอวมสฺส วจนีโย “อหํ ภนฺเต เอกํ อาปตฺติํ อาปชฺชิํ สญฺเจตนิกํ สุกฺก\-วิสฺสฏฺฐิํ ปญฺจาห\-ปฺปฏิจฺฉนฺนํ ฯเปฯ โสหํ ปริวุตฺถ\-ปริวาโส มานตฺตารโห อนฺตรา เอกํ อาปตฺติํ อาปชฺชิํ สญฺเจตนิกํ สุกฺก\-วิสฺสฏฺฐิํ อปฺปฏิจฺฉนฺนํ โสหํ สํฆํ อนฺตรา เอกิสฺสา อาปตฺติยา สญฺเจตนิกาย สุกฺก\-วิสฺสฏฺฐิยา อปฺปฏิจฺฉนฺนาย มูลาย\-ปฏิกสฺสนํ ยาจามี”ติ. ทุติยมฺปิ ยาจิตพฺพา. ตติยมฺปิ ยาจิตพฺพา. พฺยตฺเตน ภิกฺขุนา ปฏิพเลน สํโฆ ญาเปตพฺโพ\csromandash{}}

\proseitem{113}{“สุณาตุ เม ภนฺเต สํโฆ, อยํ อุทายี ภิกฺขุ เอกํ อาปตฺติํ อาปชฺชิ สญฺเจตนิกํ สุกฺก\-วิสฺสฏฺฐิํ ปญฺจาห\-ปฺปฏิจฺฉนฺนํ ฯเปฯ โส ปริวุตฺถ\-ปริวาโส มานตฺตารโห อนฺตรา เอกํ อาปตฺติํ อาปชฺชิ สญฺเจตนิกํ สุกฺก\-วิสฺสฏฺฐิํ อปฺปฏิจฺฉนฺนํ, โส สํฆํ อนฺตรา เอกิสฺสา อาปตฺติยา สญฺเจตนิกาย สุกฺก\-วิสฺสฏฺฐิยา อปฺปฏิจฺฉนฺนาย มูลาย\-ปฏิกสฺสนํ ยาจติ, ยทิ สํฆสฺส ปตฺตกลฺลํ, สํโฆ อุทายิํ ภิกฺขุํ อนฺตรา เอกิสฺสา อาปตฺติยา สญฺเจตนิกาย สุกฺก\-วิสฺสฏฺฐิยา อปฺปฏิจฺฉนฺนาย มูลาย ปฏิกสฺเสยฺย, เอสา ญตฺติ.}

\prose{สุณาตุ เม ภนฺเต สํโฆ, อยํ อุทายี ภิกฺขุ เอกํ อาปตฺติํ อาปชฺชิ สญฺเจตนิกํ สุกฺก\-วิสฺสฏฺฐิํ ปญฺจาห\-ปฺปฏิจฺฉนฺนํ ฯเปฯ โส ปริวุตฺถ\-ปริวาโส มานตฺตารโห อนฺตรา เอกํ อาปตฺติํ อาปชฺชิ สญฺเจตนิกํ สุกฺก\-วิสฺสฏฺฐิํ อปฺปฏิจฺฉนฺนํ, โส สํฆํ อนฺตรา เอกิสฺสา อาปตฺติยา สญฺเจตนิกาย สุกฺก\-วิสฺสฏฺฐิยา อปฺปฏิจฺฉนฺนาย มูลาย\-ปฏิกสฺสนํ ยาจติ, สํโฆ อุทายิํ ภิกฺขุํ อนฺตรา เอกิสฺสา อาปตฺติยา สญฺเจตนิกาย สุกฺก\-วิสฺสฏฺฐิยา อปฺปฏิจฺฉนฺนาย มูลาย ปฏิกสฺสติ, ยสฺสายสฺมโต ขมติ อุทายิสฺส ภิกฺขุโน อนฺตรา เอกิสฺสา อาปตฺติยา สญฺเจตนิกาย \csromanfolio{117}สุกฺก\-วิสฺสฏฺฐิยา อปฺปฏิจฺฉนฺนาย มูลาย\-ปฏิกสฺสนา, โส ตุณฺหสฺส, ยสฺส นกฺขมติ, โส ภาเสยฺย.}

\prose{ทุติยมฺปิ เอตมตฺถํ วทามิ ฯเปฯ. ตติยมฺปิ เอตมตฺถํ วทามิ ฯเปฯ.}

\prose{ปฏิกสฺสิโต สํเฆน อุทายี ภิกฺขุ อนฺตรา เอกิสฺสา อาปตฺติยา สญฺเจตนิกาย สุกฺก\-วิสฺสฏฺฐิยา อปฺปฏิจฺฉนฺนาย มูลาย\-ปฏิกสฺสนา, ขมติ สํฆสฺส, ตสฺมา ตุณฺหี, เอวเมตํ ธารยามี”ติ.}

\csromanmatikaanchor{mtk.69}
\csromantocmarkref{subsection}{ติกาปตฺติมานตฺต}{mtk.69}
\bookmark[dest={mtk.69},level=2]{ติกาปตฺติมานตฺต}
\csromanheader{2}{ติกา\-ปตฺติ\-มานตฺต}
\proseitem{114}{โส ปริวุตฺถ\-ปริวาโส ภิกฺขูนํ อาโรเจสิ “อหํ อาวุโส เอกํ อาปตฺติํ อาปชฺชิํ สญฺเจตนิกํ สุกฺก\-วิสฺสฏฺฐิํ ปญฺจาห\-ปฺปฏิจฺฉนฺนํ ฯเปฯ โสหํ ปริวุตฺถ\-ปริวาโส, กถํ นุ โข มยา ปฏิปชฺชิตพฺพนฺ”ติ. ภควโต เอตมตฺถํ อาโรเจสุํ. เตน หิ ภิกฺขเว สํโฆ อุทายิสฺส ภิกฺขุโน ติสฺสนฺนํ อาปตฺตีนํ ฉารตฺตํ มานตฺตํ เทตุ. เอวญฺจ ปน ภิกฺขเว ทาตพฺพํ\csromandash{}}

\prose{เตน ภิกฺขเว อุทายินา ภิกฺขุนา สํฆํ อุปสงฺกมิตฺวา ฯเปฯ เอวมสฺส วจนีโย “อหํ ภนฺเต เอกํ อาปตฺติํ อาปชฺชิํ สญฺเจตนิกํ สุกฺก\-วิสฺสฏฺฐิํ ปญฺจาห\-ปฺปฏิจฺฉนฺนํ, โสหํ สํฆํ เอกิสฺสา อาปตฺติยา สญฺเจตนิกาย สุกฺก\-วิสฺสฏฺฐิยา ปญฺจาห\-ปฺปฏิจฺฉนฺนาย ปญฺจาห\-ปริวาสํ ยาจิํ, ตสฺส เม สํโฆ เอกิสฺสา อาปตฺติยา สญฺเจตนิกาย สุกฺก\-วิสฺสฏฺฐิยา ปญฺจาห\-ปฺปฏิจฺฉนฺนาย ปญฺจาห\-ปริวาสํ อทาสิ ฯเปฯ โสหํ ภนฺเต ปริวุตฺถ\-ปริวาโส สํฆํ ติสฺสนฺนํ อาปตฺตีนํ ฉารตฺถํ มานตฺตํ ยาจามี”ติ. ทุติยมฺปิ ยาจิตพฺพํ. ตติยมฺปิ ยาจิตพฺพํ. พฺยตฺเตน ภิกฺขุนา ปฏิพเลน สํโฆ ญาเปตพฺโพ\csromandash{}}

\proseitem{115}{“สุณาตุ เม ภนฺเต สํโฆ, อยํ อุทายี ภิกฺขุ เอกํ อาปตฺติํ อาปชฺชิ สญฺเจตนิกํ สุกฺก\-วิสฺสฏฺฐิํ ปญฺจาห\-ปฺปฏิจฺฉนฺนํ ฯเปฯ. โส ปริวุตฺถ\-ปริวาโส สํฆํ ติสฺสนฺนํ อาปตฺตีนํ ฉารตฺตํ มานตฺตํ ยาจติ, ยทิ สํฆสฺส ปตฺตกลฺลํ, สํโฆ อุทายิสฺส ภิกฺขุโน ติสฺสนฺนํ อาปตฺตีนํ ฉารตฺตํ มานตฺตํ ทเทยฺย, เอสา ญตฺติ.}

\prose{สุณาตุ เม ภนฺเต สํโฆ, อยํ อุทายี ภิกฺขุ เอกํ อาปตฺติํ อาปชฺชิ สญฺเจตนิกํ สุกฺก\-วิสฺสฏฺฐิํ ปญฺจาห\-ปฺปฏิจฺฉนฺนํ ฯเปฯ. โส ปริวุตฺถ\-ปริวาโส \csromanfolio{118}สํฆํ ติสฺสนฺนํ อาปตฺตีนํ ฉารตฺตํ มานตฺตํ ยาจติ, สํโฆ อุทายิสฺส ภิกฺขุโน ติสฺสนฺนํ อาปตฺตีนํ ฉารตฺตํ มานตฺตํ เทติ, ยสฺสายสฺมโต ขมติ อุทายิสฺส ภิกฺขุโน ติสฺสนฺนํ อาปตฺตีนํ ฉารตฺตํ มานตฺตสฺส ทานํ, โส ตุณฺหสฺส, ยสฺส นกฺขมติ, โส ภาเสยฺย.}

\prose{ทุติยมฺปิ เอตมตฺถํ วทามิ ฯเปฯ. ตติยมฺปิ เอตมตฺถํ วทามิ ฯเปฯ.}

\prose{ทินฺนํ สํเฆน อุทายิสฺส ภิกฺขุโน ติสฺสนฺนํ อาปตฺตีนํ ฉารตฺตํ มานตฺตํ, ขมติ สํฆสฺส, ตสฺมา ตุณฺหี, เอวเมตํ ธารยามี”ติ.}

\csromanmatikaanchor{mtk.70}
\csromantocmarkref{subsection}{มานตฺตจาริกมูลายปฏิกสฺสนา}{mtk.70}
\bookmark[dest={mtk.70},level=2]{มานตฺตจาริกมูลายปฏิกสฺสนา}
\csromanheader{2}{มานตฺตจาริก\-มูลาย\-ปฏิกสฺสนา}
\proseitem{116}{โส มานตฺตํ จรนฺโต อนฺตรา เอกํ อาปตฺติํ อาปชฺชิ สญฺเจตนิกํ สุกฺก\-วิสฺสฏฺฐิํ อปฺปฏิจฺฉนฺนํ, โส ภิกฺขูนํ อาโรเจสิ “อหํ อาวุโส เอกํ อาปตฺติํ อาปชฺชิํ สญฺเจตนิกํ สุกฺก\-วิสฺสฏฺฐิํ ปญฺจาห\-ปฺปฏิจฺฉนฺนํ ฯเปฯ โสหํ มานตฺตํ จรนฺโต อนฺตรา เอกํ อาปตฺติํ อาปชฺชิํ สญฺเจตนิกํ สุกฺก\-วิสฺสฏฺฐิํ อปฺปฏิจฺฉนฺนํ, กถํ นุ โข มยา ปฏิปชฺชิตพฺพนฺ”ติ. ภควโต เอตมตฺถํ อาโรเจสุํ. เตน หิ ภิกฺขเว สํโฆ อุทายิํ ภิกฺขุํ อนฺตรา เอกิสฺสา อาปตฺติยา สญฺเจตนิกาย สุกฺก\-วิสฺสฏฺฐิยา อปฺปฏิจฺฉนฺนาย มูลาย ปฏิกสฺสิตฺวา ฉารตฺตํ มานตฺตํ เทตุ. เอวญฺจ ปน ภิกฺขเว มูลาย ปฏิกสฺสิตพฺโพ, เตน ภิกฺขเว อุทายินา ภิกฺขุนา สํฆํ อุปสงฺกมิตฺวา ฯเปฯ เอวมสฺส วจนีโย “อหํ ภนฺเต เอกํ อาปตฺติํ อาปชฺชิํ สญฺเจตนิกํ สุกฺก\-วิสฺสฏฺฐิํ ปญฺจาห\-ปฺปฏิจฺฉนฺนํ ฯเปฯ โสหํ มานตฺตํ จรนฺโต อนฺตรา เอกํ อาปตฺติํ อาปชฺชิํ สญฺเจตนิกํ สุกฺก\-วิสฺสฏฺฐิํ อปฺปฏิจฺฉนฺนํ โสหํ ภนฺเต สํฆํ อนฺตรา เอกิสฺสา อาปตฺติยา สญฺเจตนิกาย สุกฺก\-วิสฺสฏฺฐิยา อปฺปฏิจฺฉนฺนาย มูลาย\-ปฏิกสฺสนํ ยาจามี”ติ. ทุติยมฺปิ ยาจิตพฺพา. ตติยมฺปิ ยาจิตพฺพา. พฺยตฺเตน ภิกฺขุนา ปฏิพเลน สํโฆ ญาเปตพฺโพ\csromandash{}}

\proseitem{117}{“สุณาตุ เม ภนฺเต สํโฆ, อยํ อุทายี ภิกฺขุ ฯเปฯ มูลาย\-ปฏิกสฺสนํ ยาจติ, ยทิ สํฆสฺส ปตฺตกลฺลํ, สํโฆ อุทายิํ ภิกฺขุํ อนฺตรา เอกิสฺสา อาปตฺติยา สญฺเจตนิกาย สุกฺก\-วิสฺสฏฺฐิยา อปฺปฏิจฺฉนฺนาย มูลาย ปฏิกสฺเสยฺย, เอสา ญตฺติ ฯเปฯ. ปฏิกสฺสิโต สํเฆน อุทายี ภิกฺขุ อนฺตรา เอกิสฺสา อาปตฺติยา สญฺเจตนิกาย \csromanfolio{119}สุกฺก\-วิสฺสฏฺฐิยา อปฺปฏิจฺฉนฺนาย มูลาย\-ปฏิกสฺสนา, ขมติ สํฆสฺส, ตสฺมา ตุณฺหี, เอวเมตํ ธารยามี”ติ.}

\prose{เอวญฺจ ปน ภิกฺขเว ฉารตฺตํ มานตฺตํ ทาตพฺพํ, เตน ภิกฺขเว อุทายินา ภิกฺขุนา ฯเปฯ เอวมสฺส วจนีโย “อหํ ภนฺเต เอกํ อาปตฺติํ อาปชฺชิํ สญฺเจตนิกํ สุกฺก\-วิสฺสฏฺฐิํ ปญฺจาห\-ปฺปฏิจฺฉนฺนํ ฯเปฯ โสหํ มานตฺตํ จรนฺโต อนฺตรา เอกํ อาปตฺติํ อาปชฺชิํ สญฺเจตนิกํ สุกฺก\-วิสฺสฏฺฐิํ อปฺปฏิจฺฉนฺนํ, โสหํ สํฆํ อนฺตรา เอกิสฺสา อาปตฺติยา สญฺเจตนิกาย สุกฺก\-วิสฺสฏฺฐิยา อปฺปฏิจฺฉนฺนาย มูลาย\-ปฏิกสฺสนํ ยาจิํ, ตํ มํ สํโฆ อนฺตรา เอกิสฺสา อาปตฺติยา สญฺเจตนิกายา สุกฺก\-วิสฺสฏฺฐิยา อปฺปฏิจฺฉนฺนาย มูลาย ปฏิกสฺสิ. โสหํ ภนฺเต สํฆํ อนฺตรา เอกิสฺสา อาปตฺติยา สญฺเจตนิกาย สุกฺก\-วิสฺสฏฺฐิยา อปฺปฏิจฺฉนฺนาย ฉารตฺตํ มานตฺตํ ยาจามี”ติ. ทุติยมฺปิ ยาจิตพฺพํ. ตติยมฺปิ ยาจิตพฺพํ. พฺยตฺเตน ภิกฺขุนา ปฏิพเลน สํโฆ ญาเปตพฺโพ\csromandash{}}

\prose{“สุณาตุ เม ภนฺเต สํโฆ อยํ อุทายี ภิกฺขุ ฯเปฯ ฉารตฺตํ มานตฺตํ ยาจติ, ยทิ สํฆสฺส ปตฺตกลฺลํ, สํโฆ อุทายิสฺส ภิกฺขุโน อนฺตรา เอกิสฺสา อาปตฺติยา สญฺเจตนิกาย สุกฺก\-วิสฺสฏฺฐิยา อปฺปฏิจฺฉนฺนาย ฉารตฺตํ มานตฺตํ ทเทยฺย, เอสา ญตฺติ ฯเปฯ. ทินฺนํ สํเฆน อุทายิสฺส ภิกฺขุโน อนฺตรา เอกิสฺสา อาปตฺติยา สญฺเจตนิกาย สุกฺก\-วิสฺสฏฺฐิยา อปฺปฏิจฺฉนฺนาย ฉารตฺตํ มานตฺตํ, ขมติ สํฆสฺส, ตสฺมา ตุณฺหี, เอวเมตํ ธารยามี”ติ.}

\csromanmatikaanchor{mtk.71}
\csromantocmarkref{subsection}{อพฺภานารหมูลายปฏิกสฺสนา}{mtk.71}
\bookmark[dest={mtk.71},level=2]{อพฺภานารหมูลายปฏิกสฺสนา}
\csromanheader{2}{อพฺภานารห\-มูลาย\-ปฏิกสฺสนา}
\proseitem{118}{โส จิณฺณมานตฺโต อพฺภานารโห อนฺตรา เอกํ อาปตฺติํ อาปชฺชิ สญฺเจตนิกํ สุกฺก\-วิสฺสฏฺฐิํ อปฺปฏิจฺฉนฺนํ, โส ภิกฺขูนํ อาโรเจสิ “อหํ อาวุโส เอกํ อาปตฺติํ อาปชฺชิํ สญฺเจตนิกํ สุกฺก\-วิสฺสฏฺฐิํ ปญฺจาห\-ปฺปฏิจฺฉนฺนํ ฯเปฯ โสหํ จิณฺณมานตฺโต อพฺภานารโห อนฺตรา เอกํ อาปตฺติํ อาปชฺชิํ สญฺเจตนิกํ สุกฺก\-วิสฺสฏฺฐิํ อปฺปฏิจฺฉนฺนํ, กถํ นุ โข มยา ปฏิปชฺชิตพฺพนฺ”ติ. ภควโต เอตมตฺถํ อาโรเจสุํ. เตน หิ ภิกฺขเว สํโฆ อุทายิํ ภิกฺขุํ อนฺตรา เอกิสฺสา อาปตฺติยา สญฺเจตนิกาย สุกฺก\-วิสฺสฏฺฐิยา อปฺปฏิจฺฉนฺนาย มูลาย ปฏิกสฺสิตฺวา ฉารตฺตํ มานตฺตํ เทตุ. เอวญฺจ ปน ภิกฺขเว มูลาย ปฏิกสฺสิตพฺโพ ฯเปฯ.}

\prose{\csromanfolio{120}เอวญฺจ ปน ภิกฺขเว ฉารตฺตํ มานตฺตํ ทาตพฺพํ ฯเปฯ.}

\prose{ทินฺนํ สํเฆน อุทายิสฺส ภิกฺขุโน อนฺตรา เอกิสฺสา อาปตฺติยา สญฺเจตนิกาย สุกฺก\-วิสฺสฏฺฐิยา อปฺปฏิจฺฉนฺนาย ฉารตฺตํ มานตฺตํ, ขมติ สํฆสฺส, ตสฺมา ตุณฺหี, เอวเมตํ ธารยามี”ติ.}

\csromanmatikaanchor{mtk.72}
\csromantocmarkref{subsection}{มูลายปฏิกสฺสิตอพฺภาน}{mtk.72}
\bookmark[dest={mtk.72},level=2]{มูลายปฏิกสฺสิตอพฺภาน}
\csromanheader{2}{มูลาย\-ปฏิกสฺสิต\-อพฺภาน}
\proseitem{119}{โส จิณฺณมานตฺโต ภิกฺขูนํ อาโรเจสิ “อหํ อาวุโส เอกํ อาปตฺติํ อาปชฺชิํ สญฺเจตนิกํ สุกฺก\-วิสฺสฏฺฐิํ ปญฺจาห\-ปฺปฏิจฺฉนฺนํ ฯเปฯ โสหํ จิณฺณมานตฺโต, กถํ นุ โข มยา ปฏิปชฺชิตพฺพนฺ”ติ. ภควโต เอตมตฺถํ อาโรเจสุํ. เตน หิ ภิกฺขเว สํโฆ อุทายิํ ภิกฺขุํ อพฺเภตุ. เอวญฺจ ปน ภิกฺขเว อพฺเภตพฺโพ\csromandash{}}

\prose{เตน ภิกฺขเว อุทายินา ภิกฺขุนา สํฆํ อุปสงฺกมิตฺวา ฯเปฯ เอวมสฺส วจนีโย “อหํ ภนฺเต เอกํ อาปตฺติํ อาปชฺชิํ สญฺเจตนิกํ สุกฺก\-วิสฺสฏฺฐิํ ปญฺจาห\-ปฺปฏิจฺฉนฺนํ, โสหํ สํฆํ เอกิสฺสา อาปตฺติยา สญฺเจตนิกาย สุกฺก\-วิสฺสฏฺฐิยา ปญฺจาห\-ปฺปฏิจฺฉนฺนาย ปญฺจาห\-ปริวาสํ ยาจิํ, ตสฺส เม สํโฆ เอกิสฺสา อาปตฺติยา สญฺเจตนิกาย สุกฺก\-วิสฺสฏฺฐิยา ปญฺจาห\-ปฺปฏิจฺฉนฺนาย ปญฺจาห\-ปริวาสํ อทาสิ. โสหํ ปริวสนฺโต อนฺตรา เอกํ อาปตฺติํ อาปชฺชิํ สญฺเจตนิกํ สุกฺก\-วิสฺสฏฺฐิํ อปฺปฏิจฺฉนฺนํ, โสหํ สํฆํ อนฺตรา เอกิสฺสา อาปตฺติยา สญฺเจตนิกาย สุกฺก\-วิสฺสฏฺฐิยา อปฺปฏิจฺฉนฺนาย มูลาย\-ปฏิกสฺสนํ ยาจิํ, ตํ มํ สํโฆ อนฺตรา เอกิสฺสา อาปตฺติยา สญฺเจตนิกาย สุกฺก\-วิสฺสฏฺฐิยา อปฺปฏิจฺฉนฺนาย มูลาย ปฏิกสฺสิ. โสหํ ปริวุตฺถ\-ปริวาโส มานตฺตารโห อนฺตรา เอกํ อาปตฺติํ อาปชฺชิํ สญฺเจตนิกํ สุกฺก\-วิสฺสฏฺฐิํ อปฺปฏิจฺฉนฺนํ, โสหํ สํฆํ อนฺตรา เอกิสฺสา อาปตฺติยา สญฺเจตนิกาย สุกฺก\-วิสฺสฏฺฐิยา อปฺปฏิจฺฉนฺนาย มูลาย\-ปฏิกสฺสนํ ยาจิํ, ตํ มํ สํโฆ อนฺตรา เอกิสฺสา อาปตฺติยา สญฺเจตนิกาย สุกฺก\-วิสฺสฏฺฐิยา อปฺปฏิจฺฉนฺนาย มูลาย ปฏิกสฺสิ. โสหํ ปริวุตฺถวาโส สํฆํ ติสฺสนฺนํ อาปตฺตีนํ ฉารตฺตํ มานตฺตํ ยาจิํ, ตสฺส เม สํโฆ ติสฺสนฺนํ อาปตฺตีนํ ฉารตฺตํ มานตฺตํ อทาสิ. โสหํ มานตฺตํ จรนฺโต อนฺตรา เอกํ อาปตฺติํ อาปชฺชิํ สญฺเจตนิกํ สุกฺก\-วิสฺสฏฺฐิํ อปฺปฏิจฺฉนฺนํ, โสหํ สํฆํ อนฺตรา เอกิสฺสา อาปตฺติยา สญฺเจตนิกาย สุกฺก\-วิสฺสฏฺฐิยา อปฺปฏิจฺฉนฺนาย มูลาย\-ปฏิกสฺสนํ ยาจิํ, ตํ มํ \csromanfolio{121}สํโฆ อนฺตรา เอกิสฺสา อาปตฺติยา สญฺเจตนิกาย สุกฺก\-วิสฺสฏฺฐิยา อปฺปฏิจฺฉนฺนาย มูลาย ปฏิกสฺสิ, โสหํ สํฆํ อนฺตรา เอกิสฺสา อาปตฺติยา สญฺเจตนิกาย สุกฺก\-วิสฺสฏฺฐิยา อปฺปฏิจฺฉนฺนาย ฉารตฺตํ มานตฺตํ ยาจิํ, ตสฺส เม สํโฆ อนฺตรา เอกิสฺสา อาปตฺติยา สญฺเจตนิกาย สุกฺก\-วิสฺสฏฺฐิยา อปฺปฏิจฺฉนฺนาย ฉารตฺตํ มานตฺตํ อทาสิ. โสหํ จิณฺณมานตฺโต อพฺภานารโห อนฺตรา เอกํ อาปตฺติํ อาปชฺชิํ สญฺเจตนิกํ สุกฺก\-วิสฺสฏฺฐิํ อปฺปฏิจฺฉนฺนํ, โสหํ สํฆํ อนฺตรา เอกิสฺสา อาปตฺติยา สญฺเจตนิกาย สุกฺก\-วิสฺสฏฺฐิยา อปฺปฏิจฺฉนฺนาย มูลาย\-ปฏิกสฺสนํ ยาจิํ, ตํ มํ สํโฆ อนฺตรา เอกิสฺสา อาปตฺติยา สญฺเจตนิกาย สุกฺก\-วิสฺสฏฺฐิยา อปฺปฏิจฺฉนฺนาย มูลาย ปฏิกสฺสิ, โสหํ สํฆํ อนฺตรา เอกิสฺสา อาปตฺติยา สญฺเจตนิกาย สุกฺก\-วิสฺสฏฺฐิยา อปฺปฏิจฺฉนฺนาย ฉารตฺตํ มานตฺตํ ยาจิํ, ตสฺส เม สํโฆ อนฺตรา เอกิสฺสา อาปตฺติยา สญฺเจตนิกาย สุกฺก\-วิสฺสฏฺฐิยา อปฺปฏิจฺฉนฺนาย ฉารตฺตํ มานตฺตํ อทาสิ, โสหํ ภนฺเต จิณฺณมานตฺโต สํฆํ อพฺภานํ ยาจามี”ติ.}

\prose{ทุติยมฺปิ ยาจิตพฺพํ. ตติยมฺปิ ยาจิตพฺพํ. พฺยตฺเตน ภิกฺขุนา ปฏิพเลน สํโฆ ญาเปตพฺโพ\csromandash{}}

\proseitem{120}{“สุณาตุ เม ภนฺเต สํโฆ, อยํ อุทายี ภิกฺขุ เอกํ อาปตฺติํ อาปชฺชิ สญฺเจตนิกํ สุกฺก\-วิสฺสฏฺฐิํ ปญฺจาห\-ปฺปฏิจฺฉนฺนํ, โส สํฆํ เอกิสฺสา อาปตฺติยา สญฺเจตนิกาย สุกฺก\-วิสฺสฏฺฐิยา ปญฺจาห\-ปฺปฏิจฺฉนฺนาย ปญฺจาห\-ปริวาสํ ยาจิ, สํโฆ อุทายิสฺส ภิกฺขุโน เอกิสฺสา อาปตฺติยา สญฺเจตนิกาย สุกฺก\-วิสฺสฏฺฐิยา ปญฺจาห\-ปฺปฏิจฺฉนฺนาย ปญฺจาห\-ปริวาสํ อทาสิ. โส ปริวสนฺโต อนฺตรา เอกํ อาปตฺติํ อาปชฺชิ สญฺเจตนิกํ สุกฺก\-วิสฺสฏฺฐิํ อปฺปฏิจฺฉนฺนํ, โส สํฆํ อนฺตรา เอกิสฺสา อาปตฺติยา สญฺเจตนิกาย สุกฺก\-วิสฺสฏฺฐิยา อปฺปฏิจฺฉนฺนาย มูลาย\-ปฏิกสฺสนํ ยาจิ, สํโฆ อุทายิํ ภิกฺขุํ อนฺตรา เอกิสฺสา อาปตฺติยา สญฺเจตนิกาย สุกฺก\-วิสฺสฏฺฐิยา อปฺปฏิจฺฉนฺนาย มูลาย ปฏิกสฺสิ. โส ปริวุตฺถ\-ปริวาโส มานตฺตารโห อนฺตรา เอกํ อาปตฺติํ อาปชฺชิ สญฺเจตนิกํ สุกฺก\-วิสฺสฏฺฐิํ อปฺปฏิจฺฉนฺนํ, โส สํฆํ อนฺตรา เอกิสฺสา อาปตฺติยา สญฺเจตนิกาย สุกฺก\-วิสฺสฏฺฐิยา อปฺปฏิจฺฉนฺนาย มูลาย\-ปฏิกสฺสนํ ยาจิ, สํโฆ อุทายิํ ภิกฺขุํ อนฺตรา เอกิสฺสา อาปตฺติยา \csromanfolio{122}สญฺเจตนิกาย สุกฺก\-วิสฺสฏฺฐิยา อปฺปฏิจฺฉนฺนาย มูลาย ปฏิกสฺสิ. โส ปริวุตฺถ\-ปริวาโส สํฆํ ติสฺสนฺนํ อาปตฺตีนํ ฉารตฺตํ มานตฺตํ ยาจิ, สํโฆ อุทายิสฺส ภิกฺขุโน ติสฺสนฺนํ อาปตฺตีนํ ฉารตฺตํ มานตฺตํ อทาสิ. โส มานตฺตํ จรนฺโต อนฺตรา เอกํ อาปตฺติํ อาปชฺชิ สญฺเจตนิกํ สุกฺก\-วิสฺสฏฺฐิํ อปฺปฏิจฺฉนฺนํ, โส สํฆํ อนฺตรา เอกิสฺสา อาปตฺติยา สญฺเจตนิกาย สุกฺก\-วิสฺสฏฺฐิยา อปฺปฏิจฺฉนฺนาย มูลาย\-ปฏิกสฺสนํ ยาจิ, สํโฆ อุทายิํ ภิกฺขุํ อนฺตรา เอกิสฺสา อาปตฺติยา สญฺเจตนิกาย สุกฺก\-วิสฺสฏฺฐิยา อปฺปฏิจฺฉนฺนาย มูลาย ปฏิกสฺสิ. โส สํฆํ อนฺตรา เอกิสฺสา อาปตฺติยา สญฺเจตนิกาย สุกฺก\-วิสฺสฏฺฐิยา อปฺปฏิจฺฉนฺนาย ฉารตฺตํ มานตฺตํ ยาจิ, สํโฆ อุทายิสฺส ภิกฺขุโน อนฺตรา เอกิสฺสา อาปตฺติยา สญฺเจตนิกาย สุกฺก\-วิสฺสฏฺฐิยา อปฺปฏิจฺฉนฺนาย ฉารตฺตํ มานตฺตํ อทาสิ. โส จิณฺณมานตฺโต อพฺภานารโห อนฺตรา เอกํ อาปตฺติํ อาปชฺชิ สญฺเจตนิกํ สุกฺก\-วิสฺสฏฺฐิํ อปฺปฏิจฺฉนฺนํ, โส สํฆํ อนฺตรา เอกิสฺสา อาปตฺติยา สญฺเจตนิกาย สุกฺก\-วิสฺสฏฺฐิยา อปฺปฏิจฺฉนฺนาย มูลาย ปฏิกสฺสนํ ยาจิ, สํโฆ อุทายิํ ภิกฺขุํ อนฺตรา เอกิสฺสา อาปตฺติยา สญฺเจตนิกาย สุกฺก\-วิสฺสฏฺฐิยา อปฺปฏิจฺฉนฺนาย มูลาย ปฏิกสฺสิ. โส สํฆํ อนฺตรา เอกิสฺสา อาปตฺติยา สญฺเจตนิกาย สุกฺก\-วิสฺสฏฺฐิยา อปฺปฏิจฺฉนฺนาย ฉารตฺตํ มานตฺตํ ยาจิ, สํโฆ อุทายิสฺส ภิกฺขุโน อนฺตรา เอกิสฺสา อาปตฺติยา สญฺเจตนิกาย สุกฺก\-วิสฺสฏฺฐิยา อปฺปฏิจฺฉนฺนาย ฉารตฺตํ มานตฺตํ อทาสิ, โส จิณฺณมานตฺโต สํฆํ อพฺภานํ ยาจติ. ยทิ สํฆสฺส ปตฺตกลฺลํ, สํโฆ อุทายิํ ภิกฺขุํ อพฺเภยฺย, เอสา ญตฺติ.}

\prose{สุณาตุ เม ภนฺเต สํโฆ, อยํ อุทายี ภิกฺขุ เอกํ อาปตฺติํ อาปชฺชิ สญฺเจตนิกํ สุกฺก\-วิสฺสฏฺฐิํ ปญฺจาห\-ปฺปฏิจฺฉนฺนํ ฯเปฯ โส จิณฺณมานตฺโต สํฆํ อพฺภานํ ยาจติ. สํโฆ อุทายิํ ภิกฺขุํ อพฺเภติ, ยสฺสายสฺมโต ขมติ อุทายิสฺส ภิกฺขุโน อพฺภานํ, โส ตุณฺหสฺส, ยสฺส นกฺขมติ โส ภาเสยฺย.}

\prose{ทุติยมฺปิ เอตมตฺถํ วทามิ ฯเปฯ. ตติยมฺปิ เอตมตฺถํ วทามิ ฯเปฯ.}

\prose{อพฺภิโต สํเฆน อุทายี ภิกฺขุ, ขมติ สํฆสฺส, ตสฺมา ตุณฺหี, เอวเมตํ ธารยามี”ติ.}

\csromanheader{2}{\textbf{ปกฺขฏิจฺฉนฺนปริวาส}}
\csromanmatikaanchor{mtk.73}
\csromantocmarkref{subsection}{ปกฺขปฺปฏิจฺฉนฺนปริวาส}{mtk.73}
\bookmark[dest={mtk.73},level=2]{ปกฺขปฺปฏิจฺฉนฺนปริวาส}
\proseitem{121}{\csromanfolio{123}เตน โข ปน สมเยน อายสฺมา อุทายี เอกํ อาปตฺติํ อาปนฺโน โหติ สญฺเจตนิกํ สุกฺก\-วิสฺสฏฺฐิํ ปกฺข\-ปฺปฏิจฺฉนฺนํ, โส ภิกฺขูนํ อาโรเจสิ “อหํ อาวุโส เอกํ อาปตฺติํ อาปชฺชิํ สญฺเจตนิกํ สุกฺก\-วิสฺสฏฺฐิํ ปกฺข\-ปฺปฏิจฺฉนฺนํ, กถํ นุ โข มยา ปฏิปชฺชิตพฺพนฺ”ติ. ภควโต เอตมตฺถํ อาโรเจสุํ. เตน หิ ภิกฺขเว สํโฆ อุทายิสฺส ภิกฺขุโน เอกิสฺสา อาปตฺติยา สญฺเจตนิกาย สุกฺก\-วิสฺสฏฺฐิยา ปกฺข\-ปฺปฏิจฺฉนฺนาย ปกฺข\-ปริวาสํ เทตุ, เอวญฺจ ปน ภิกฺขเว ทาตพฺโพ\csromandash{}}

\prose{เตน ภิกฺขเว อุทายินา ภิกฺขุนา สํฆํ อุปสงฺกมิตฺวา ฯเปฯ เอวมสฺส วจนีโย “อหํ ภนฺเต เอกํ อาปตฺติํ อาปชฺชิํ สญฺเจตนิกํ สุกฺก\-วิสฺสฏฺฐิํ ปกฺข\-ปฺปฏิจฺฉนฺนํ, โสหํ ภนฺเต สํฆํ เอกิสฺสา อาปตฺติยา สญฺเจตนิกาย สุกฺก\-วิสฺสฏฺฐิยา ปกฺข\-ปฺปฏิจฺฉนฺนาย ปกฺข\-ปริวาสํ ยาจามี”ติ. ทุติยมฺปิ ยาจิตพฺโพ. ตติยมฺปิ ยาจิตพฺโพ. พฺยตฺเตน ภิกฺขุนา ปฏิพเลน สํโฆ ญาเปตพฺโพ\csromandash{}}

\proseitem{122}{“สุณาตุ เม ภนฺเต สํโฆ, อยํ อุทายี ภิกฺขุ เอกํ อาปตฺติํ อาปชฺชิ สญฺเจตนิกํ สุกฺก\-วิสฺสฏฺฐิํ ปกฺข\-ปฺปฏิจฺฉนฺนํ, โส สํฆํ เอกิสฺสา อาปตฺติยา สญฺเจตนิกาย สุกฺก\-วิสฺสฏฺฐิยา ปกฺข\-ปฺปฏิจฺฉนฺนาย ปกฺข\-ปริวาสํ ยาจติ, ยทิ สํฆสฺส ปตฺตกลฺลํ, สํโฆ อุทายิสฺส ภิกฺขุโน เอกิสฺสา อาปตฺติยา สญฺเจตนิกาย สุกฺก\-วิสฺสฏฺฐิยา ปกฺข\-ปฺปฏิจฺฉนฺนาย ปกฺข\-ปริวาสํ ทเทยฺย, เอสา ญตฺติ.}

\prose{สุณาตุ เม ภนฺเต สํโฆ, อยํ อุทายี ภิกฺขุ เอกํ อาปตฺติํ อาปชฺชิ สญฺเจตนิกํ สุกฺก\-วิสฺสฏฺฐิํ ปกฺข\-ปฺปฏิจฺฉนฺนํ, โส สํฆํ เอกิสฺสา อาปตฺติยา สญฺเจตนิกาย สุกฺก\-วิสฺสฏฺฐิยา ปกฺข\-ปฺปฏิจฺฉนฺนาย ปกฺข\-ปริวาสํ ยาจติ, สํโฆ อุทายิสฺส ภิกฺขุโน เอกิสฺสา อาปตฺติยา สญฺเจตนิกาย สุกฺก\-วิสฺสฏฺฐิยา ปกฺข\-ปฺปฏิจฺฉนฺนาย ปกฺข\-ปริวาสํ เทติ, ยสฺสายสฺมโต ขมติ อุทายิสฺส ภิกฺขุโน เอกิสฺสา อาปตฺติยา สญฺเจตนิกาย สุกฺก\-วิสฺสฏฺฐิยา ปกฺข\-ปฺปฏิจฺฉนฺนาย ปกฺข\-ปริวาสสฺส ทานํ, โส ตุณฺหสฺส, ยสฺส นกฺขมติ, โส ภาเสยฺย.}

\prose{ทุติยมฺปิ เอตมตฺถํ วทามิ ฯเปฯ. ตติยมฺปิ เอตมตฺถํ วทามิ ฯเปฯ.}

\prose{\csromanfolio{124}ทินฺโน สํเฆน อุทายิสฺส ภิกฺขุโน เอกิสฺสา อาปตฺติยา สญฺเจตนิกาย สุกฺก\-วิสฺสฏฺฐิยา ปกฺข\-ปฺปฏิจฺฉนฺนาย ปกฺขปริวาโส, ขมติ สํฆสฺส, ตสฺมา ตุณฺหี, เอวเมตํ ธารยามี”ติ.}

\csromanmatikaanchor{mtk.74}
\csromantocmarkref{subsection}{ปกฺขปาริวาสิกมูลายปฏิกสฺสนา}{mtk.74}
\bookmark[dest={mtk.74},level=2]{ปกฺขปาริวาสิกมูลายปฏิกสฺสนา}
\csromanheader{2}{ปกฺขปาริวาสิก\-มูลาย\-ปฏิกสฺสนา}
\proseitem{123}{โส ปริวสนฺโต อนฺตรา เอกํ อาปตฺติํ อาปชฺชิ สญฺเจตนิกํ สุกฺก\-วิสฺสฏฺฐิํ ปญฺจาห\-ปฺปฏิจฺฉนฺนํ, โส ภิกฺขูนํ อาโรเจสิ “อหํ อาวุโส เอกํ อาปตฺติํ อาปชฺชิํ สญฺเจตนิกํ สุกฺก\-วิสฺสฏฺฐิํ ปกฺข\-ปฺปฏิจฺฉนฺนํ, โสหํ สํฆํ เอกิสฺสา อาปตฺติยา สญฺเจตนิกาย สุกฺก\-วิสฺสฏฺฐิยา ปกฺข\-ปฺปฏิจฺฉนฺนาย ปกฺข\-ปริวาสํ ยาจิํ, ตสฺส เม สํโฆ เอกิสฺสา อาปตฺติยา สญฺเจตนิกาย สุกฺก\-วิสฺสฏฺฐิยา ปกฺข\-ปฺปฏิจฺฉนฺนาย ปกฺข\-ปริวาสํ อทาสิ, โสหํ ปริวสนฺโต อนฺตรา เอกํ อาปตฺติํ อาปชฺชิํ สญฺเจตนิกํ สุกฺก\-วิสฺสฏฺฐิํ ปญฺจาห\-ปฺปฏิจฺฉนฺนํ, กถํ นุ โข มยา ปฏิปชฺชิตพฺพนฺ”ติ. ภควโต เอตมตฺถํ อาโรเจสุํ. เตน หิ ภิกฺขเว สํโฆ อุทายิํ ภิกฺขุํ อนฺตรา เอกิสฺสา อาปตฺติยา สญฺเจตนิกาย สุกฺก\-วิสฺสฏฺฐิยา ปญฺจาห\-ปฺปฏิจฺฉนฺนาย มูลาย ปฏิกสฺสิตฺวา ปุริมาย อาปตฺติยา สโมธาน\-ปริวาสํ เทตุ, เอวญฺจ ปน ภิกฺขเว มูลาย ปฏิกสฺสิตพฺโพ\csromandash{}}

\prose{เตน ภิกฺขเว อุทายินา ภิกฺขุนา สํฆํ อุปสงฺกมิตฺวา ฯเปฯ เอวมสฺส วจนีโย “อหํ ภนฺเต เอกํ อาปตฺติํ อาปชฺชิํ สญฺเจตนิกํ สุกฺก\-วิสฺสฏฺฐิํ ปกฺข\-ปฺปฏิจฺฉนฺนํ, โสหํ สํฆํ เอกิสฺสา อาปตฺติยา สญฺเจตนิกาย สุกฺก\-วิสฺสฏฺฐิยา ปกฺข\-ปฺปฏิจฺฉนฺนาย ปกฺข\-ปริวาสํ ยาจิํ, ตสฺส เม สํโฆ เอกิสฺสา อาปตฺติยา สญฺเจตนิกาย สุกฺก\-วิสฺสฏฺฐิยา ปกฺข\-ปฺปฏิจฺฉนฺนาย ปกฺข\-ปริวาสํ อทาสิ. โสหํ ปริวสนฺโต อนฺตรา เอกํ อาปตฺติํ อาปชฺชิํ สญฺเจตนิกํ สุกฺก\-วิสฺสฏฺฐิํ ปญฺจาห\-ปฺปฏิจฺฉนฺนํ, โสหํ ภนฺเต สํฆํ อนฺตรา เอกิสฺสา อาปตฺติยา สญฺเจตนิกาย สุกฺก\-วิสฺสฏฺฐิยา ปญฺจาห\-ปฺปฏิจฺฉนฺนาย มูลาย\-ปฏิกสฺสนํ ยาจามี”ติ. ทุติยมฺปิ ยาจิตพฺพา. ตติยมฺปิ ยาจิตพฺพา. พฺยตฺเตน ภิกฺขุนา ปฏิพเลน สํโฆ ญาเปตพฺโพ\csromandash{}}

\proseitem{124}{“สุณาตุ เม ภนฺเต สํโฆ, อยํ อุทายี ภิกฺขุ เอกํ อาปตฺติํ อาปชฺชิ สญฺเจตนิกํ สุกฺก\-วิสฺสฏฺฐิํ ปกฺข\-ปฺปฏิจฺฉนฺนํ, โส สํฆํ เอกิสฺสา อาปตฺติยา สญฺเจตนิกาย สุกฺก\-วิสฺสฏฺฐิยา ปกฺข\-ปฺปฏิจฺฉนฺนาย ปกฺข\-ปริวาสํ \csromanfolio{125}ยาจิ, สํโฆ อุทายิสฺส ภิกฺขุโน เอกิสฺสา อาปตฺติยา สญฺเจตนิกาย สุกฺก\-วิสฺสฏฺฐิยา ปกฺข\-ปฺปฏิจฺฉนฺนาย ปกฺข\-ปริวาสํ อทาสิ. โส ปริวสนฺโต อนฺตรา เอกํ อาปตฺติํ อาปชฺชิ สญฺเจตนิกํ สุกฺก\-วิสฺสฏฺฐิํ ปญฺจาห\-ปฺปฏิจฺฉนฺนํ, โส สํฆํ อนฺตรา เอกิสฺสา อาปตฺติยา สญฺเจตนิกาย สุกฺก\-วิสฺสฏฺฐิยา ปญฺจาห\-ปฺปฏิจฺฉนฺนาย มูลาย\-ปฏิกสฺสนํ ยาจติ. ยทิ สํฆสฺส ปตฺตกลฺลํ, สํโฆ อุทายิํ ภิกฺขุํ อนฺตรา เอกิสฺสา อาปตฺติยา สญฺเจตนิกาย สุกฺก\-วิสฺสฏฺฐิยา ปญฺจาห\-ปฺปฏิจฺฉนฺนาย มูลาย ปฏิกสฺเสยฺย, เอสา ญตฺติ.}

\prose{สุณาตุ เม ภนฺเต สํโฆ, อยํ อุทายี ภิกฺขุ เอกํ อาปตฺติํ อาปชฺชิ สญฺเจตนิกํ สุกฺก\-วิสฺสฏฺฐิํ ปกฺข\-ปฺปฏิจฺฉนฺนํ, โส สํฆํ เอกิสฺสา อาปตฺติยา สญฺเจตนิกาย สุกฺก\-วิสฺสฏฺฐิยา ปกฺข\-ปฺปฏิจฺฉนฺนาย ปกฺข\-ปริวาสํ ยาจิ, สํโฆ อุทายิสฺส ภิกฺขุโน เอกิสฺสา อาปตฺติยา สญฺเจตนิกาย สุกฺก\-วิสฺสฏฺฐิยา ปกฺข\-ปฺปฏิจฺฉนฺนาย ปกฺข\-ปริวาสํ อทาสิ. โส ปริวสนฺโต อนฺตรา เอกํ อาปตฺติํ อาปชฺชิ สญฺเจตนิกํ สุกฺก\-วิสฺสฏฺฐิํ ปญฺจาห\-ปฺปฏิจฺฉนฺนํ, โส สํฆํ อนฺตรา เอกิสฺสา อาปตฺติยา สญฺเจตนิกาย สุกฺก\-วิสฺสฏฺฐิยา ปญฺจาห\-ปฺปฏิจฺฉนฺนาย มูลาย\-ปฏิกสฺสนํ ยาจติ. สํโฆ อุทายิํ ภิกฺขุํ อนฺตรา เอกิสฺสา อาปตฺติยา สญฺเจตนิกาย สุกฺก\-วิสฺสฏฺฐิยา ปญฺจาห\-ปฺปฏิจฺฉนฺนาย มูลาย ปฏิกสฺสติ, ยสฺสายสฺมโต ขมติ อุทายิสฺส ภิกฺขุโน อนฺตรา เอกิสฺสา อาปตฺติยา สญฺเจตนิกาย สุกฺก\-วิสฺสฏฺฐิยา ปญฺจาห\-ปฺปฏิจฺฉนฺนาย มูลาย\-ปฏิกสฺสนา, โส ตุณฺหสฺส, ยสฺส นกฺขมติ, โส ภาเสยฺย.}

\prose{ทุติยมฺปิ เอตมตฺถํ วาทมิ ฯเปฯ. ตติยมฺปิ เอตมตฺถํ วทามิ ฯเปฯ.}

\prose{ปฏิกสฺสิโต สํเฆน อุทายี ภิกฺขุ อนฺตรา เอกิสฺสา อาปตฺติยา สญฺเจตนิกาย สุกฺก\-วิสฺสฏฺฐิยา ปญฺจาห\-ปฺปฏิจฺฉนฺนาย มูลาย\-ปฏิกสฺสนา, ขมติ สํฆสฺส, ตสฺมา ตุณฺหี, เอวเมตํ ธารยามี”ติ.}

\csromanmatikaanchor{mtk.75}
\csromantocmarkref{subsection}{สโมธานปริวาส}{mtk.75}
\bookmark[dest={mtk.75},level=2]{สโมธานปริวาส}
\csromanheader{2}{สโมธาน\-ปริวาส}
\proseitem{125}{เอวญฺจ ปน ภิกฺขเว ปุริมาย อาปตฺติยา สโมธาน\-ปริวาโส ทาตพฺโพ\csromandash{}เตน ภิกฺขเว อุทายินา ภิกฺขุนา สํฆํ อุปสงฺกมิตฺวา ฯเปฯ เอวมสฺส วจนีโย “อหํ ภนฺเต เอกํ อาปตฺติํ อาปชฺชิํ สญฺเจตนิกํ \csromanfolio{126}สุกฺก\-วิสฺสฏฺฐิํ ปกฺข\-ปฺปฏิจฺฉนฺนํ, โสหํ สํฆํ เอกิสฺสา อาปตฺติยา สญฺเจตนิกาย สุกฺก\-วิสฺสฏฺฐิยา ปกฺข\-ปฺปฏิจฺฉนฺนาย ปกฺข\-ปริวาสํ ยาจิํ, ตสฺส เม สํโฆ เอกิสฺสา อาปตฺติยา สญฺเจตนิกาย สุกฺก\-วิสฺสฏฺฐิยา ปกฺข\-ปฺปฏิจฺฉนฺนาย ปกฺข\-ปริวาสํ อทาสิ. โสหํ ปริวสนฺโต อนฺตรา เอกํ อาปตฺติํ อาปชฺชิํ สญฺเจตนิกํ สุกฺก\-วิสฺสฏฺฐิํ ปญฺจาห\-ปฺปฏิจฺฉนฺนํ, โสหํ สํฆํ อนฺตรา เอกิสฺสา อาปตฺติยา สญฺเจตนิกาย สุกฺก\-วิสฺสฏฺฐิยา ปญฺจาห\-ปฺปฏิจฺฉนฺนาย มูลาย\-ปฏิกสฺสนํ ยาจิํ, ตํ มํ สํโฆ อนฺตรา เอกิสฺสา อาปตฺติยา สญฺเจตนิกาย สุกฺก\-วิสฺสฏฺฐิยา ปญฺจาห\-ปฺปฏิจฺฉนฺนาย มูลาย ปฏิกสฺสิ, โสหํ ภนฺเต สํฆํ อนฺตรา เอกิสฺสา อาปตฺติยา สญฺเจตนิกาย สุกฺก\-วิสฺสฏฺฐิยา ปญฺจาห\-ปฺปฏิจฺฉนฺนาย ปุริมาย อาปตฺติยา สโมธาน\-ปริวาสํ ยาจามี”ติ.}

\prose{ทุติยมฺปิ ยาจิตพฺโพ. ตติยมฺปิ ยาจิตพฺโพ. พฺยตฺเตน ภิกฺขุนา ปฏิพเลน สํโฆ ญาเปตพฺโพ\csromandash{}}

\proseitem{126}{“สุณาตุ เม ภนฺเต สํโฆ, อยํ อุทายี ภิกฺขุ เอกํ อาปตฺติํ อาปชฺชิ สญฺเจตนิกํ สุกฺก\-วิสฺสฏฺฐิํ ปกฺข\-ปฺปฏิจฺฉนฺนํ, โส สํฆํ เอกิสฺสา อาปตฺติยา สญฺเจตนิกาย สุกฺก\-วิสฺสฏฺฐิยา ปกฺข\-ปฺปฏิจฺฉนฺนาย ปกฺข\-ปริวาสํ ยาจิ, สํโฆ อุทายิสฺส ภิกฺขุโน เอกิสฺสา อาปตฺติยา สญฺเจตนิกาย สุกฺก\-วิสฺสฏฺฐิยา ปกฺข\-ปฺปฏิจฺฉนฺนาย ปกฺข\-ปริวาสํ อทาสิ. โส ปริวสนฺโต อนฺตรา เอกํ อาปตฺติํ อาปชฺชิ สญฺเจตนิกํ สุกฺก\-วิสฺสฏฺฐิํ ปญฺจาห\-ปฺปฏิจฺฉนฺนํ, โส สํฆํ อนฺตรา เอกิสฺสา อาปตฺติยา สญฺเจตนิกาย สุกฺก\-วิสฺสฏฺฐิยา ปญฺจาห\-ปฺปฏิจฺฉนฺนาย มูลาย\-ปฏิกสฺสนํ ยาจิ, สํโฆ อุทายิํ ภิกฺขุํ อนฺตรา เอกิสฺสา อาปตฺติยา สญฺเจตนิกาย สุกฺก\-วิสฺสฏฺฐิยา ปญฺจาห\-ปฺปฏิจฺฉนฺนาย มูลาย ปฏิกสฺสิ, โส สํฆํ อนฺตรา เอกิสฺสา อาปตฺติยา สญฺเจตนิกาย สุกฺก\-วิสฺสฏฺฐิยา ปญฺจาห\-ปฺปฏิจฺฉนฺนาย ปุริมาย อาปตฺติยา สโมธาน\-ปริวาสํ ยาจติ. ยทิ สํฆสฺส ปตฺตกลฺลํ, สํโฆ อุทายิสฺส ภิกฺขุโน อนฺตรา เอกิสฺสา อาปตฺติยา สญฺเจตนิกาย สุกฺก\-วิสฺสฏฺฐิยา ปญฺจาห\-ปฺปฏิจฺฉนฺนาย ปุริมาย อาปตฺติํยา สโมธาน\-ปริวาสํ ทเทยฺย, เอสา ญตฺติ.}

\prose{สุณาตุ เม ภนฺเต สํโฆ, อยํ อุทายี ภิกฺขุ เอกํ อาปตฺติํ อาปชฺชิ สญฺเจตนิกํ สุกฺก\-วิสฺสฏฺฐิํ ปกฺข\-ปฺปฏิจฺฉนฺนํ, โส สํฆํ เอกิสฺสา \csromanfolio{127}อาปตฺติยา สญฺเจตนิกาย สุกฺก\-วิสฺสฏฺฐิยา ปกฺกปฺปฏิจฺฉนฺนาย ปกฺข\-ปริวาสํ ยาจิ, สํโฆ อุทายิสฺส ภิกฺขุโน เอกิสฺสา อาปตฺติยา สญฺเจตนิกาย สุกฺก\-วิสฺสฏฺฐิยา ปกฺข\-ปฺปฏิจฺฉนฺนาย ปกฺข\-ปริวาสํ อทาสิ. โส ปริวสนฺโต อนฺตรา เอกํ อาปตฺติํ อาปชฺชิ สญฺเจตนิกํ สุกฺก\-วิสฺสฏฺฐิํ ปญฺจาห\-ปฺปฏิจฺฉนฺนํ, โส สํฆํ อนฺตรา เอกิสฺสา อาปตฺติยา สญฺเจตนิกาย สุกฺก\-วิสฺสฏฺฐิยา ปญฺจาห\-ปฺปฏิจฺฉนฺนาย มูลาย\-ปฏิกสฺสนํ ยาจิ, สํโฆ อุทายิํ ภิกฺขุํ อนฺตรา เอกิสฺสา อาปตฺติยา สญฺเจตนิกาย สุกฺก\-วิสฺสฏฺฐิยา ปญฺจาห\-ปฺปฏิจฺฉนฺนาย มูลาย ปฏิกสฺสิ, โส สํฆํ อนฺตรา เอกิสฺสา อาปตฺติยา สญฺเจตนิกาย สุกฺก\-วิสฺสฏฺฐิยา ปญฺจาห\-ปฺปฏิจฺฉนฺนาย ปุริมาย อาปตฺติยา สโมธาน\-ปริวาสํ ยาจติ. สํโฆ อุทายิสฺส ภิกฺขุโน อนฺตรา เอกิสฺสา อาปตฺติยา สญฺเจตนิกาย สุกฺก\-วิสฺสฏฺฐิยา ปญฺจาห\-ปฺปฏิจฺฉนฺนาย ปุริมาย อาปตฺติยา สโมธาน\-ปริวาสํ เทติ, ยสฺสายสฺมโต ขมติ อุทายิสฺส ภิกฺขุโน อนฺตรา เอกิสฺสา อาปตฺติยา สญฺเจตนิกาย สุกฺก\-วิสฺสฏฺฐิยา ปญฺจาห\-ปฺปฏิจฺฉนฺนาย ปุริมาย อาปตฺติยา สโมธาน\-ปริวาสสฺส ทานํ, โส ตุณฺหสฺส, ยสฺส นกฺขมติ, โส ภ ภาเสยฺย.}

\prose{ทุติยมฺปิ เอตมตฺถํ วทามิ ฯเปฯ. ตติยมฺปิ เอตมตฺถํ วทามิ ฯเปฯ.}

\prose{ทินฺโน สํเฆน อุทายิสฺส ภิกฺขุโน อนฺตรา เอกิสฺสา อาปตฺติยา สญฺเจตนิกาย สุกฺก\-วิสฺสฏฺฐิยา ปญฺจาห\-ปฺปฏิจฺฉนฺนาย ปุริมาย อาปตฺติยา สโมธาน\-ปริวาโส, ขมติ สํฆสฺส, ตสฺมา ตุณฺหี, เอวเมตํ ธารยามี”ติ.}

\csromanmatikaanchor{mtk.76}
\csromantocmarkref{subsection}{มานตฺตารหมูลายปฏิกสฺสนาทิ}{mtk.76}
\bookmark[dest={mtk.76},level=2]{มานตฺตารหมูลายปฏิกสฺสนาทิ}
\csromanheader{2}{มานตฺตารห\-มูลายปฏิกสฺสนาทิ}
\proseitem{127}{โส ปริวุตฺถ\-ปริวาโส มานตฺตารโห อนฺตรา เอกํ อาปตฺติํ อาปชฺชิ สญฺเจตนิกํ สุกฺก\-วิสฺสฏฺฐิํ ปญฺจาห\-ปฺปฏิจฺฉนฺนํ. โส ภิกฺขูนํ อาโรเจสิ “อหํ อาวุโส เอกํ อาปตฺติํ อาปชฺชิํ สญฺเจตนิกํ สุกฺก\-วิสฺสฏฺฐิํ ปกฺข\-ปฺปฏิจฺฉนฺนํ ฯเปฯ โสหํ ปริวุตฺถ\-ปริวาโส มานตฺตารโห อนฺตรา เอกํ อาปตฺติํ อาปชฺชิํ สญฺเจตนิกํ สุกฺก\-วิสฺสฏฺฐิํ ปญฺจาห\-ปฺปฏิจฺฉนฺนํ, กถํ นุ โข มยา ปฏิปชฺชิตพฺพนฺ”ติ. ภควโต เอตมตฺถํ อาโรเจสุํ. เตน หิ ภิกฺขเว สํโฆ อุทายิํ ภิกฺขุํ อนฺตรา เอกิสฺสา อาปตฺติยา สญฺเจตนิกาย สุกฺก\-วิสฺสฏฺฐิยา ปญฺจาห\-ปฺปฏิจฺฉนฺนาย มูลาย ปฏิกสฺสิตฺวา \csromanfolio{128}ปุริมาย อาปตฺติยา สโมธาน\-ปริวาสํ เทตุ. เอวญฺจ ปน ภิกฺขเว มูลาย ปฏิกสฺสิตพฺโพ ฯเปฯ.}

\prose{เอวญฺจ ปน ภิกฺขเว ปุริมาย อาปตฺติยา สโมธาน\-ปริวาโส ทาตพฺโพ ฯเปฯ เทติ ฯเปฯ.}

\prose{ทินฺโน สํเฆน อุทายิสฺส ภิกฺขุโน อนฺตรา เอกิสฺสา อาปตฺติยา สญฺเจตนิกาย สุกฺก\-วิสฺสฏฺฐิยา ปญฺจาห\-ปฺปฏิจฺฉนฺนาย ปุริมาย อาปตฺติยา สโมธาน\-ปริวาโส, ขมติ สํฆสฺส, ตสฺมา ตุณฺหี, เอวเมตํ ธารยามี”ติ.}

\csromanmatikaanchor{mtk.77}
\csromantocmarkref{subsection}{ติกาปตฺติมานตฺต}{mtk.77}
\bookmark[dest={mtk.77},level=2]{ติกาปตฺติมานตฺต}
\csromanheader{2}{ติกา\-ปตฺติ\-มานตฺต}
\proseitem{128}{โส ปริวุตฺถ\-ปริวาโส ภิกฺขูนํ อาโรโจสิ “อหํ อาวุโส เอกํ อาปตฺติํ อาปชฺชิํ สญฺเจตนิกํ สุกฺก\-วิสฺสฏฺฐิํ ปกฺข\-ปฺปฏิจฺฉนฺนํ ฯเปฯ โสหํ ปริวุตฺถ\-ปริวาโส, กถํ นุ โข มยา ปฏิปชฺชิตพฺพนฺ”ติ. ภควโต เอตมตฺถํ อาโรเจสุํ. เตน หิ ภิกฺขเว สํโฆ อุทายิสฺส ภิกฺขุโน ติสฺสนฺนํ อาปตฺตีนํ ฉารตฺตํ มานตฺตํ เทตุ. เอวญฺจ ปน ภิกฺขเว ทาตพฺพํ\csromandash{}}

\prose{เตน ภิกฺขเว อุทายิน ภิกฺขุนา สํฆํ อุปสงฺกมิตฺวา ฯเปฯ เอวมสฺส วจนีโย “อหํ ภนฺเต เอกํ อาปตฺติํ อาปชฺชิํ สญฺเจตนิกํ สุกฺก\-วิสฺสฏฺฐิํ ปกฺข\-ปฺปฏิจฺฉนฺนํ ฯเปฯ โสหํ ภนฺเต ปริวุตฺถ\-ปริวาโส สํฆํ ติสฺสนฺนํ อาปตฺตีนํ ฉารตฺตํ มานตฺตํ ยาจามี”ติ. ทุติยมฺปิ ยาจิตพฺพํ. ตติยมฺปิ ยาจิตพฺพํ. พฺยตฺเตน ภิกฺขุนา ปฏิพเลน สํโฆ ญาเปตพฺโพ\csromandash{}}

\proseitem{129}{“สุณาตุ เม ภนฺเต สํโฆ, อยํ อุทายี ภิกฺขุ เอกํ อาปตฺติํ อาปชฺชิ สญฺเจตนิกํ สุกฺก\-วิสฺสฏฺฐิํ ปกฺข\-ปฺปฏิจฺฉนฺนํ ฯเปฯ โส ปริวุตฺถ\-ปริวาโส สํฆํ ติสฺสนฺนํ อาปตฺตีนํ ฉารตฺตํ มานตฺตํ ยาจติ. ยทิ สํฆสฺส ปตฺตกลฺลํ, สํโฆ อุทายิสฺส ภิกฺขุโน ติสฺสนฺนํ อาปตฺตีนํ ฉารตฺตํ มานตฺตํ ทเทยฺย, เอสา ญตฺติ.}

\prose{สุณาตุ เม ภนฺเต สํโฆ, อยํ อุทายี ภิกฺขุ เอกํ อาปตฺติํ อาปชฺชิ สญฺเจตนิกํ สุกฺก\-วิสฺสฏฺฐิํ ปกฺข\-ปฺปฏิจฺฉนฺนํ ฯเปฯ โส ปริวุตฺถ\-ปริวาโส สํฆํ ติสฺสนฺนํ อาปตฺตีนํ ฉารตฺตํ มานตฺตํ ยาจติ. สํโฆ อุทายิสฺส ภิกฺขุโน ติสฺสนฺนํ อาปตฺตีนํ ฉารตฺตํ มานตฺตํ เทติ, ยสฺสายสฺมโต \csromanfolio{129}ขมติ อุทายิสฺส ภิกฺขุโน ติสฺสนฺนํ อาปตฺตีนํ ฉารตฺตํ มานตฺตํ ทานํ, โส ตุณฺหสฺส, ยสฺส นกฺขมติ, โส ภาเสยฺย.}

\prose{ทุติยมฺปิ เอตมตฺถํ วทามิ ฯเปฯ. ตติยมฺปิ เอตมตฺถํ วทามิ ฯเปฯ.}

\prose{ทินฺนํ สํเฆน อุทายิสฺส ภิกฺขุโน ติสฺสนฺนํ อาปตฺตีนํ ฉารตฺตํ มานตฺตํ, ขมติ สํฆสฺส, ตสฺมา ตุณฺหี, เอวเมตํ ธารยามี”ติ.}

\csromanmatikaanchor{mtk.78}
\csromantocmarkref{subsection}{มานตฺตจาริกมูลายปฏิกสฺสนาทิ}{mtk.78}
\bookmark[dest={mtk.78},level=2]{มานตฺตจาริกมูลายปฏิกสฺสนาทิ}
\csromanheader{2}{มานตฺตจาริก\-มูลายปฏิกสฺสนาทิ}
\proseitem{130}{โส มานตฺตํ จรนฺโต อนฺตรา เอกํ อาปตฺติํ อาปชฺชิ สญฺเจตนิกํ สุกฺก\-วิสฺสฏฺฐิํ ปญฺจาห\-ปฺปฏิจฺฉนฺนํ. โส ภิกฺขูนํ อาโรเจสิ “อหํ อาวุโส เอกํ อาปตฺติํ อาปชฺชิํ สญฺเจตนิกํ สุกฺก\-วิสฺสฏฺฐิํ ปกฺข\-ปฺปฏิจฺฉนฺนํ ฯเปฯ โสหํ มานตฺตํ จรนฺโต อนฺตรา เอกํ อาปตฺติํ อาปชฺชิํ สญฺเจตนิกํ สุกฺก\-วิสฺสฏฺฐิํ ปญฺจาห\-ปฺปฏิจฺฉนฺนํ, กถํ นุ โข มยา ปฏิปชฺชิตพฺพนฺ”ติ. ภควโต เอตมตฺถํ อาโรเจสุํ. เตน หิ ภิกฺขเว สํโฆ อุทายิํ ภิกฺขุํ อนฺตรา เอกิสฺสา อาปตฺติยา สญฺเจตนิกาย สุกฺก\-วิสฺสฏฺฐิยา ปญฺจหปฺปฏิจฺฉนฺนาย มูลาย ปฏิกสฺสิตฺวา ปุริมาย อาปตฺติยา สโมธาน\-ปริวาสํ ทตฺวา ฉารตฺตํ มานตฺตํ เทตุ. เอวญฺจ ปน ภิกฺขเว มูลาย ปฏิกสฺสิตพฺโพ ฯเปฯ.}

\prose{เอวญฺจ ปน ภิกฺขเว ปุริมาย อาปตฺติยา สโมธาน\-ปริวาโส ทาตพฺโพ ฯเปฯ.}

\prose{เอวญฺจ ปน ภิกฺขเว ฉารตฺตํ มานตฺตํ ทาตพฺพํ ฯเปฯ เทติ ฯเปฯ.}

\prose{ทินฺนํ สํเฆน อุทายิสฺส ภิกฺขุโน อนฺตรา เอกิสฺสา อาปตฺติยา สญฺเจตนิกาย สุกฺก\-วิสฺสฏฺฐิยา ปญฺจาห\-ปฺปฏิจฺฉนฺนาย ฉารตฺตํ มานตฺตํ, ขมติ สํฆสฺส, ตสฺมา ตุณฺหี, เอวเมตํ ธารยามี”ติ.}

\csromanmatikaanchor{mtk.79}
\csromantocmarkref{subsection}{อพฺภานารหมูลายปฏิกสฺสนาทิ}{mtk.79}
\bookmark[dest={mtk.79},level=2]{อพฺภานารหมูลายปฏิกสฺสนาทิ}
\csromanheader{2}{อพฺภานารห\-มูลายปฏิกสฺสนาทิ}
\proseitem{131}{โส จิณฺณมานตฺโต อพฺภานารโห อนฺตรา เอกํ อาปตฺติํ อาปชฺชิ สญฺเจตนิกํ สุกฺก\-วิสฺสฏฺฐิํ ปญฺจาห\-ปฺปฏิจฺฉนฺนํ. โส ภิกฺขูนํ อาโรเจสิ “อหํ อาวุโส เอกํ อาปตฺติํ อาปชฺชิํ สญฺเจตนิกํ สุกฺก\-วิสฺสฏฺฐิํ ปกฺข\-ปฺปฏิจฺฉนฺนํ ฯเปฯ โสหํ จิณฺณมานตฺโต อพฺภานารโห อนฺตรา เอกํ อาปตฺติํ อาปชฺชิํ สญฺเจตนิกํ สุกฺก\-วิสฺสฏฺฐิํ ปญฺจาห\-ปฺปฏิจฺฉนฺนํ, กถํ นุ โข \csromanfolio{130}มยา ปฏิปชฺชิตพฺพนฺ”ติ. ภควโต เอตมตฺถํ อาโรเจสุํ. เตน หิ ภิกฺขเว สํโฆ อุทายิํ ภิกฺขุํ อนฺตรา เอกิสฺสา อาปตฺติยา สญฺเจตนิกาย สุกฺก\-วิสฺสฏฺฐิยา ปญฺจาห\-ปฺปฏิจฺฉนฺนาย มูลาย ปฏิกสฺสิตฺวา ปุริมาย อาปตฺติยา สโมธาน\-ปริวาสํ ทตฺวา ฉารตฺตํ มานตฺตํ เทตุ. เอวญฺจ ปน ภิกฺขเว มูลาย ปฏิกสฺสิตพฺโพ ฯเปฯ.}

\prose{เอวญฺจ ปน ภิกฺขเว ปุริมาย อาปตฺติยา สโมธาน\-ปริวาโส ทาตพฺโพ ฯเปฯ.}

\prose{เอวญฺจ ปน ภิกฺขเว ฉารตฺตํ มานตฺตํ ทาตพฺพํ ฯเปฯ เทติ ฯเปฯ.}

\prose{ทินฺนํ สํเฆน อุทายิสฺส ภิกฺขุโน อนฺตรา เอกิสฺสา อาปตฺติยา สญฺเจตนิกาย สุกฺก\-วิสฺสฏฺฐิยา ปญฺจาห\-ปฺปฏิจฺฉนฺนาย ฉารตฺตํ มานตฺตํ, ขมติ สํฆสฺส, ตสฺมา ตุณฺหี, เอวเมตํ ธารยามี”ติ.}

\csromanmatikaanchor{mtk.80}
\csromantocmarkref{subsection}{ปกฺขปฺปฏิจฺฉนฺนอพฺภาน}{mtk.80}
\bookmark[dest={mtk.80},level=2]{ปกฺขปฺปฏิจฺฉนฺนอพฺภาน}
\csromanheader{2}{ปกฺข\-ปฺปฏิจฺฉนฺน\-อพฺภาน}
\proseitem{132}{โส จิณฺณมานตฺโต ภิกฺขูนํ อาโรเจสิ “อหํ อาวุโส เอกํ อาปตฺติํ อาปชฺชิํ สญฺเจตนิกํ สุกฺก\-วิสฺสฏฺฐิํ ปกฺข\-ปฺปฏิจฺฉนฺนํ ฯเปฯ โสหํ จิณฺณมานตฺโต, กถํ นุ โข มยา ปฏิปชฺชิตพฺพนฺ”ติ. ภควโต เอตมตฺถํ อาโรเจสุํ. เตน หิ ภิกฺขเว สํโฆ อุทายิํ ภิกฺขุํ อพฺเภตุ. เอวญฺจ ปน ภิกฺขเว อพฺเภตพฺโพ\csromandash{}}

\prose{เตน ภิกฺขเว อุทายินา ภิกฺขุนา สํฆํ อุปสงฺกมิตฺวา เอกํสํ อุตฺตราสงฺคํ กริตฺวา วุฑฺฒานํ ภิกฺขูนํ ปาเท วนฺทิตฺวา อุกฺกุฏิกํ นิสีทิตฺวา อญฺชลิํ ปคฺคเหตฺวา เอวมสฺส วจนีโย “อหํ ภนฺเต เอกํ อาปตฺติํ อาปชฺชิํ สญฺเจตนิกํ สุกฺก\-วิสฺสฏฺฐิํ ปกฺข\-ปฺปฏิจฺฉนฺนํ, โสหํ สํฆํ เอกิสฺสา อาปตฺติยา สญฺเจตนิกาย สุกฺก\-วิสฺสฏฺฐิยา ปกฺข\-ปฺปฏิจฺฉนฺนาย ปกฺข\-ปริวาสํ ยาจิํ, ตสฺส เม สํโฆ เอกิสฺสา อาปตฺติยา สญฺเจตนิกาย สุกฺก\-วิสฺสฏฺฐิยา ปกฺข\-ปฺปฏิจฺฉนฺนาย ปกฺข\-ปริวาสํ อทาสิ. โสหํ ปริวสนฺโต อนฺตรา เอกํ อาปตฺติํ อาปชฺชิํ สญฺเจตนิกํ สุกฺก\-วิสฺสฏฺฐิํ ปญฺจาห\-ปฺปฏิจฺฉนฺนํ, โสหํ สํฆํ อนฺตรา เอกิสฺสา อาปตฺติยา สญฺเจตนิกาย สุกฺก\-วิสฺสฏฺฐิยา ปญฺจาห\-ปฺปฏิจฺฉนฺนาย มูลาย\-ปฏิกสฺสนํ ยาจิํ, ตํ มํ สํโฆ อนฺตรา เอกิสฺสา อาปตฺติยา สญฺเจตนิกาย สุกฺก\-วิสฺสฏฺฐิยา ปญฺจาหปฺปติจฺฉนฺนาย มูลาย ปฏิกสฺสิ, โสหํ สํฆํ อนฺตรา เอกิสฺสา \csromanfolio{131}อาปตฺติยา สญฺเจตนิกาย สุกฺก\-วิสฺสฏฺฐิยา ปญฺจาห\-ปฺปฏิจฺฉนฺนาย ปุริมาย อาปตฺติยา สโมธาน\-ปริวาสํ ยาจิํ, ตสฺส เม สํโฆ อนฺตรา เอกิสฺสา อาปตฺติยา สญฺเจตนิกาย สุกฺก\-วิสฺสฏฺฐิยา ปญฺจาห\-ปฺปฏิจฺฉนฺนาย ปุริมาย อาปตฺติยา สโมธาน\-ปริวาสํ อทาสิ. โสหํ ปริวุตฺถ\-ปริวาโส มานตฺตารโห อนฺตรา เอกํ อาปตฺติํ อาปชฺชิํ สญฺเจตนิกํ สุกฺก\-วิสฺสฏฺฐิํ ปญฺจาห\-ปฺปฏิจฺฉนฺนํ, โสหํ สํฆํ อนฺตรา เอกิสฺสา อาปตฺติยา สญฺเจตนิกาย สุกฺก\-วิสฺสฏฺฐิยา ปญฺจาห\-ปฺปฏิจฺฉนฺนาย มูลาย\-ปฏิกสฺสนํ ยาจิํ, ตํ มํ สํโฆ อนฺตรา เอกิสฺสา อาปตฺติยา สญฺเจตนิกาย สุกฺก\-วิสฺสฏฺฐิยา ปญฺจาห\-ปฺปฏิจฺฉนฺนาย มูลาย ปฏิกสฺสิ, โสหํ สํฆํ อนฺตรา เอกิสฺสา อาปตฺติยา สญฺเจตนิกาย สุกฺก\-วิสฺสฏฺฐิยา ปญฺจาห\-ปฺปฏิจฺฉนฺนาย ปุริมาย อาปตฺติยา สโมธาน\-ปริวาสํ ยาจิํ, ตสฺส เม สํโฆ อนฺตรา เอกิสฺสา อาปตฺติยา สญฺเจตนิกาย สุกฺก\-วิสฺสฏฺฐิยา ปญฺจาห\-ปฺปฏิจฺฉนฺนาย ปุริมาย อาปตฺติยา สโมธาน\-ปริวาสํ อทาสิ. โสหํ ปริวุตฺถ\-ปริวาโส สํฆํ ติสฺสนฺนํ อาปตฺตีนํ ฉารตฺตํ มานตฺตํ ยาจิํ, ตสฺส เม สํโฆ ติสฺสนฺนํ อาปตฺตีนํ ฉารตฺตํ มานตฺตํ อทาสิ. โสหํ มานตฺตํ จรนฺโต อนฺตรา เอกํ อาปตฺติํ อาปชฺชิํ สญฺเจตนิกํ สุกฺก\-วิสฺสฏฺฐิํ ปญฺจาห\-ปฺปฏิจฺฉนฺนํ, โสหํ สํฆํ อนฺตรา เอกิสฺสา อาปตฺติยา สญฺเจตนิกาย สุกฺก\-วิสฺสฏฺฐิยา ปญฺจาห\-ปฺปฏิจฺฉนฺนาย มูลาย\-ปฏิกสฺสนํ ยาจิํ, ตํ มํ สํโฆ อนฺตรา เอกิสฺสา อาปตฺติยา สญฺเจตนิกาย สุกฺก\-วิสฺสฏฺฐิยา ปญฺจาห\-ปฺปฏิจฺฉนฺนาย มูลาย ปฏิกสฺสิ, โสหํ สํฆํ อนฺตรา เอกิสฺสา อาปตฺติยา สญฺเจกนิกาย สุกฺก\-วิสฺสฏฺฐิยา ปญฺจาห\-ปฺปฏิจฺฉนฺนาย ปุริมาย อาปตฺติยา สโมธาน\-ปริวาสํ ยาจิํ, ตสฺส เม สํโฆ อนฺตรา เอกิสฺสา อาปตฺติยา สญฺเจตนิกาย สุกฺก\-วิสฺสฏฺฐิยา ปญฺจาห\-ปฺปฏิจฺฉนฺนาย ปุริมาย อาปตฺติยา สโมธาน\-ปริวาสํ อทาสิ. โสหํ ปริวุตฺถ\-ปริวาโส สํฆํ อนฺตรา เอกิสฺสา อาปตฺติยา สญฺเจตนิกาย สุกฺก\-วิสฺสฏฺฐิยา ปญฺจาห\-ปฺปฏิจฺฉนฺนาย ฉารตฺตํ มานตฺตํ ยาจิํ, ตสฺส เม สํโฆ อนฺตรา เอกิสฺสา อาปตฺติยา สญฺเจตนิกาย สุกฺก\-วิสฺสฏฺฐิยา ปญฺจาห\-ปฺปฏิจฺฉนฺนาย ฉารตฺตํ มานตฺตํ อทาสิ. โสหํ จิณฺณมานตฺโต อพฺภานารโห อนฺตรา เอกํ อาปตฺติํ อาปชฺชิํ สญฺเจตนิกํ สุกฺก\-วิสฺสฏฺฐิํ ปญฺจาห\-ปฺปฏิจฺฉนฺนํ, โสหํ สํฆํ อนฺตรา เอกิสฺสา อาปตฺติยา สญฺเจตนิกาย สุกฺก\-วิสฺสฏฺฐิยา ปญฺจาห\-ปฺปฏิจฺฉนฺนาย มูลาย\-ปฏิกสฺสนํ ยาจิํ, ตํ มํ สํโฆ อนฺตรา เอกิสฺสา อาปตฺติยา สญฺเจตนิกาย สุกฺก\-วิสฺสฏฺฐิยา \csromanfolio{132}ปญฺจาห\-ปฺปฏิจฺฉนฺนาย มูลาย ปฏิกสฺสิ, โสหํ สํฆํ อนฺตรา เอกิสฺสา อาปตฺติยา สญฺเจตนิกาย สุกฺก\-วิสฺสฏฺฐิยา ปญฺจาห\-ปฺปฏิจฺฉนฺนาย ปุริมาย อาปตฺติยา สโมธาน\-ปริวาสํ ยาจิํ, ตสฺส เม สํโฆ อนฺตรา เอกิสฺสา อาปตฺติยา สญฺเจตนิกาย สุกฺก\-วิสฺสฏฺฐิยา ปญฺจาห\-ปฺปฏิจฺฉนฺนาย ปุริมาย อาปตฺติยา สโมธาน\-ปริวาสํ อทาสิ, โสหํ ปริวุตฺถ\-ปริวาโส สํฆํ อนฺตรา เอกิสฺสา อาปตฺติยา สญฺเจตนิกาย สุกฺก\-วิสฺสฏฺฐิยา ปญฺจาห\-ปฺปฏิจฺฉนฺนาย ฉารตฺตํ มานตฺตํ ยาจิํ, ตสฺส เม สํโฆ อนฺตรา เอกิสฺสา อาปตฺติยา สญฺเจตนิกาย สุกฺก\-วิสฺสฏฺฐิยา ปญฺจาห\-ปฺปฏิจฺฉนฺนาย ฉารตฺตํ มานตฺตํ อทาสิ, โสหํ ภนฺเต จิณฺณมานตฺโต สํฆํ อพฺภานํ ยาจามี”ติ.}

\prose{ทุติยมฺปิ ยาจิตพฺพํ. ตติยมฺปิ ยาจิตพฺพํ. พฺยตฺเตน ภิกฺขุนา ปฏิพเลน สํโฆ ญาเปตพฺโพ\csromandash{}}

\proseitem{133}{“สุณาตุ เม ภนฺเต สํโฆ, อยํ อุทายี ภิกฺขุ เอกํ อาปตฺติํ อาปชฺชิ สญฺเจตนิกํ สุกฺก\-วิสฺสฏฺฐิํ ปกฺข\-ปฺปฏิจฺฉนฺนํ, โส สํฆํ เอกิสฺสา อาปตฺติยา สญฺเจตนิกาย สุกฺก\-วิสฺสฏฺฐิยา ปกฺข\-ปฺปฏิจฺฉนฺนาย ปกฺข\-ปริวาสํ ยาจิ, สํโฆ อุทายิสฺส ภิกฺขุโน เอกิสฺสา อาปตฺติยา สญฺเจตนิกาย สุกฺก\-วิสฺสฏฺฐิยา ปกฺข\-ปฺปฏิจฺฉนฺนาย ปกฺข\-ปริวาสํ อทาสิ. โส ปริวสนฺโต อนฺตรา เอกํ อาปตฺติํ อาปชฺชิ สญฺเจตนิกํ สุกฺก\-วิสฺสฏฺฐิํ ปญฺจาห\-ปฺปฏิจฺฉนฺนํ, โส สํฆํ อนฺตรา เอกิสฺสา อาปตฺติยา สญฺเจตนิกาย สุกฺก\-วิสฺสฏฺฐิยา ปญฺจาห\-ปฺปฏิจฺฉนฺนาย มูลาย\-ปฏิกสฺสนํ ยาจิ, สํโฆ อุทายิํ ภิกฺขุํ อนฺตรา เอกิสฺสา อาปตฺติยา สญฺเจตนิกาย สุกฺก\-วิสฺสฏฺฐิยา ปญฺจาห\-ปฺปฏิจฺฉนฺนาย มูลาย ปฏิกสฺสิ, โส สํฆํ อนฺตรา เอกิสฺสา อาปตฺติยา สญฺเจตนิกาย สุกฺก\-วิสฺสฏฺฐิยา ปญฺจาห\-ปฺปฏิจฺฉนฺนาย ปุริมาย อาปตฺติยา สโมธาน\-ปริวาสํ ยาจิ, สํโฆ อุทายิสฺส ภิกฺขุโน อนฺตรา เอกิสฺสา อาปตฺติยา สญฺเจตนิกาย สุกฺก\-วิสฺสฏฺฐิยา ปญฺจาห\-ปฺปฏิจฺฉนฺนาย ปุริมาย อาปตฺติยา สโมธาน\-ปริวาสํ อทาสิ. โส ปริวุตฺถ\-ปริวาโส มานตฺตารโห อนฺตรา เอกํ อาปตฺติํ อาปชฺชิ สญฺเจตนิกํ สุกฺก\-วิสฺสฏฺฐิํ ปญฺจหปฺปฏิจฺฉนฺนํ, โส สํฆํ อนฺตรา เอกิสฺสา อาปตฺติยา สญฺเจตนิกาย สุกฺก\-วิสฺสฏฺฐิยา ปญฺจาห\-ปฺปฏิจฺฉนฺนาย มูลาย\-ปฏิกสฺสนํ ยาจิ, สํโฆ อุทายิํ ภิกฺขุํ อนฺตรา เอกิสฺสา อาปตฺติยา สญฺเจตนิกาย สุกฺก\-วิสฺสฏฺฐิยา ปญฺจหปฺปฏิจฺฉนฺนาย มูลาย ปฏิกสฺสิ, โส สํฆํ อนฺตรา \csromanfolio{133}เอกิสฺสา อาปตฺติยา สญฺเจตนิกาย สุกฺก\-วิสฺสฏฺฐิยา ปญฺจาห\-ปฺปฏิจฺฉนฺนาย ปุริมาย อาปตฺติยา สโมธาน\-ปริวาสํ ยาจิ, สํโฆ อุทายิสฺส ภิกฺขุโน อนฺตรา เอกิสฺสา อาปตฺติยา สญฺเจตนิกาย สุกฺก\-วิสฺสฏฺฐิยา ปญฺจาห\-ปฺปฏิจฺฉนฺนาย ปุริมาย อาปตฺติยา สโมธาน\-ปริวาสํ อทาสิ. โส ปริวุตฺถ\-ปริวาโส สํฆํ ติสฺสนฺนํ อาปตฺตีนํ ฉารตฺตํ มานตฺตํ ยจิ, สํโฆ อุทายิสฺส ภิกฺขุโน ติสฺสนฺนํ อาปตฺตีนํ ฉารตฺตํ มานตฺตํ อทาสิ. โส มานตฺตํ จรนฺโต อนฺตรา เอกํ อาปตฺติํ อาปชฺชิ สญฺเจตนิกํ สุกฺก\-วิสฺสฏฺฐิํ ปญฺจาห\-ปฺปฏิจฺฉนฺนํ, โส สํฆํ อนฺตรา เอกิสฺสา อาปตฺติยา สญฺเจตนิกาย สุกฺก\-วิสฺสฏฺฐิยา ปญฺจาห\-ปฺปฏิจฺฉนฺนาย มูลาย\-ปฏิกสฺสนํ ยาจิ, สํโฆ อุทายิํ ภิกฺขุํ อนฺตรา เอกิสฺสา อาปตฺติยา สญฺเจตนิกาย สุกฺก\-วิสฺสฏฺฐิยา ปญฺจาห\-ปฺปฏิจฺฉนฺนาย มูลาย ปฏิกสฺสิ, โส สํฆํ อนฺตรา เอกิสฺสา อาปตฺติยา สญฺเจตนิกาย สุกฺก\-วิสฺสฏฺฐิยา ปญฺจาห\-ปฺปฏิจฺฉนฺนาย ปุริมาย อาปตฺติยา สโมธาน\-ปริวาสํ ยาจิ, สํโฆ อุทายิสฺส ภิกฺขุโน อนฺตรา เอกิสฺสา อาปตฺติยา สญฺเจตนิกาย สุกฺก\-วิสฺสฏฺฐิยา ปญฺจาห\-ปฺปฏิจฺฉนฺนาย ปุริมาย อาปตฺติยา สโมธาน\-ปริวาสํ อทาสิ, โส ปริวุตฺถ\-ปริวาโส สํฆํ\csromansharedfootnote{baacf61485fb0378}{โส สํฆํ (ก)} อนฺตรา เอกิสฺสา อาปตฺติยา สญฺเจตนิกาย สุกฺก\-วิสฺสฏฺฐิยา ปญฺจาห\-ปฺปฏิจฺฉนฺนาย ฉารตฺตํ มานตฺตํ ยาจิ, สํโฆ อุทายิสฺส ภิกฺขุโน อนฺตรา เอกิสฺสา อาปตฺติยา สญฺเจตนิกาย สุกฺก\-วิสฺสฏฺฐิยา ปญฺจาห\-ปฺปฏิจฺฉนฺนาย ฉารตฺตํ มานตฺตํ อทาสิ. โส จิณฺณมานตฺโต อพฺภานารโห อนฺตรา เอกํ อาปตฺติํ อาปชฺชิ สญฺเจตนิกํ สุกฺก\-วิสฺสฏฺฐิํ ปญฺจาห\-ปฺปฏิจฺฉนฺนํ, โส สํฆํ อนฺตรา เอกิสฺสา อาปตฺติยา สญฺเจตนิกาย สุกฺก\-วิสฺสฏฺฐิยา ปญฺจาห\-ปฺปฏิจฺฉนฺนาย มูลาย\-ปฏิกสฺสนํ ยาจิ, สํโฆ อุทายิํ ภิกฺขุํ อนฺตรา เอกิสฺสา อาปตฺติยา สญฺเจตนิกาย สุกฺก\-วิสฺสฏฺฐิยา ปญฺจาห\-ปฺปฏิจฺฉนฺนาย มูลาย ปฏิกสฺสิ, โส สํฆํ อนฺตรา เอกิสฺสา อาปตฺติยา สญฺเจตนิกาย สุกฺก\-วิสฺสฏฺฐิยา ปญฺจาห\-ปฺปฏิจฺฉนฺนาย ปุริมาย อาปตฺติยา สโมธาน\-ปริวาสํ ยาจิ, สํโฆ อุทายิสฺส ภิกฺขุโน อนฺตรา เอกิสฺสา อาปตฺติยา สญฺเจตนิกาย สุกฺก\-วิสฺสฏฺฐิยา ปญฺจาห\-ปฺปฏิจฺฉนฺนาย ปุริมาย อาปตฺติยา สโมธาน\-ปริวาสํ อทาสิ, โส ปริวุตฺถ\-ปริวาโส สํฆํ อนฺตรา สโมธาน\-ปริวาสํ อทาสิ, โส ปริวุตฺถ\-ปริวาโส สํฆํ อนฺตรา เอกิสฺสา อาปตฺติยา สญฺเจตนิกาย สุกฺก\-วิสฺสฏฺฐิยา ปญฺจาห\-ปฺปฏิจฺฉนฺนาย \csromanfolio{134}ฉารตฺตํ มานตฺตํ ยาจิ, สํโฆ อุทายิสฺส ภิกฺขุโน อนฺตรา เอกิสฺสา อาปตฺติยา สญฺเจตนิกาย สุกฺก\-วิสฺสฏฺฐิยา ปญฺจาห\-ปฺปฏิจฺฉนฺนาย ฉารตฺตํ มานตฺตํ อทาสิ, โส จิณฺณมานตฺโต สํฆํ อพฺภานํ ยาจติ. ยทิ สํฆสฺส ปตฺตกลฺลํ, สํโฆ อุทายิํ อพฺเภยฺย, เอสา ญตฺติ.}

\prose{สุณาตุ เม ภนฺเต สํโฆ, อยํ อุทายี ภิกฺขุ เอกํ อาปตฺติํ อาปชฺชิ สญฺเจตนิกํ สุกฺก\-วิสฺสฏฺฐิํ ปกฺข\-ปฺปฏิจฺฉนฺนํ ฯเปฯ โส จิณฺณมานตฺโต สํฆํ อพฺภานํ ยาจติ. สํโฆ อุทายิํ ภิกฺขุํ อพฺเภติ, ยสฺสายสฺมโต ขมติ อุทายิสฺส ภิกฺขุโน อพฺภานํ, โส ตุณฺหสฺส, ยสฺส นกฺขมติ, โส ภาเสยฺย.}

\prose{ทุติยมฺปิ เอตมตฺตํ วทามิ ฯเปฯ. ตติยมฺปิ เอตมตฺถํ วทามิ ฯเปฯ.}

\prosewithcloser{อพฺภิโต สํเฆน อุทายี ภิกฺขุ, ขมติ สํฆสฺส, ตสฺมา ตุณฺหี, เอวเมตํ ธารยามี”ติ.}{สุกฺกวิสฺสฏฺฐิ สมตฺตา.}
\csromanmatikaanchor{mtk.81}
\csromantocmarkref{section}{2. ปริวาส}{mtk.81}
\bookmark[dest={mtk.81},level=1]{2. ปริวาส}
\csromanheader{1}{2. \textbf{ปริวาส}}
\csromanmatikaanchor{mtk.82}
\csromantocmarkref{subsection}{อคฺฆสโมธานปริวาส}{mtk.82}
\bookmark[dest={mtk.82},level=2]{อคฺฆสโมธานปริวาส}
\csromanheader{2}{อคฺฆ\-สโมธาน\-ปริวาส}
\proseitem{134}{เตน โข ปน สมเยน อญฺญตโร ภิกฺขุ สมฺพหุลา สํฆาทิเสสา อาปตฺติโย อาปนฺโน โหติ, เอกา อาปตฺติ เอกาห\-ปฺปฏิจฺฉนฺนา, เอกา อาปตฺติ ทฺวีห\-ปฺปฏิจฺฉนฺนา, เอกา อาปตฺติ ตีหปฺปฏิจฺฉนฺนา, เอกา อาปตฺติ จตูห\-ปฺปฏิจฺฉนฺนา, เอกา อาปตฺติ ปญฺจาห\-ปฺปฏิจฺฉนฺนา, เอกา อาปตฺติ ฉาห\-ปฺปฏิจฺฉนฺนา, เอกา อาปตฺติ สตฺตาห\-ปฺปฏิจฺฉนฺนา, เอกา อาปตฺติ อฏฺฐาห\-ปฺปฏิจฺฉนฺนา, เอกา อาปตฺติ นวาห\-ปฺปฏิจฺฉนฺนา, เอกา อาปตฺติ ทสาห\-ปฺปฏิจฺฉนฺนา. โส ภิกฺขูนํ อาโรเจสิ “อหํ อาวุโส สมฺพหุลา สํฆาทิเสสา อาปตฺติโย อาปชฺชิํ, เอกา อาปตฺติ เอกาห\-ปฺปฏิจฺฉนฺนา ฯเปฯ เอกา อาปตฺติ ทสาห\-ปฺปฏิจฺฉนฺนา, กถํ นุ โข มยา ปฏิปชฺชิตพฺพนฺ”ติ. ภควโต เอตมตฺถํ อาโรเจสุํ. เตน หิ ภิกฺขเว สํโฆ ตสฺส ภิกฺขุโน ตาสํ อาปตฺตีนํ ยา อาปตฺติ ทสาห\-ปฺปฏิจฺฉนฺนา, ตสฺสา อคฺเฆน สโมธาน\-ปริวาสํ เทตุ. เอวญฺจ ปน ภิกฺขเว ทาตพฺโพ\csromandash{}}

\prose{\csromanfolio{135}เตน ภิกฺขเว ภิกฺขุนา สํฆํ อุปสงฺกมิตฺวา ฯเปฯ เอวมสฺส วจนีโย “อหํ ภนฺเต สมฺพหุลา สํฆาทิเสสา อาปตฺติโย อาปชฺชิํ, เอกา อาปตฺติ เอกาห\-ปฺปฏิจฺฉนฺนา ฯเปฯ เอกา อาปตฺติ ทสาห\-ปฺปฏิจฺฉนฺนา, โสหํ ภนฺเต สํฆํ ตาสํ อาปตฺตีนํ ยา อาปตฺติ ทสาห\-ปฺปฏิจฺฉนฺนา, ตสฺสา อคฺเฆน สโมธาน\-ปริวาสํ ยาจามี”ติ. ทุติยมฺปิ ยาจิตพฺโพ. ตติยมฺปิ ยาจิตพฺโพ. พฺยตฺเตน ภิกฺขุนา ปฏิพเลน สํโฆ ญาเปตพฺโพ\csromandash{}}

\proseitem{135}{“สุณาตุ เม ภนฺเต สํโฆ, อยํ อิตฺถนฺนาโม ภิกฺขุ สมฺพหุลา สํฆาทิเสสา อาปตฺติโย อาปชฺชิ, เอกา อาปตฺติ เอกาห\-ปฺปฏิจฺฉนฺนา ฯเปฯ เอกา อาปตฺติ ทสาห\-ปฺปฏิจฺฉนฺนา, โส สํฆํ ตาสํ อาปตฺตีนํ ยา อาปตฺติ ทสาห\-ปฺปฏิจฺฉนฺนา, ตสฺสา อคฺเฆน สโมธาน\-ปริวาสํ ยาจติ. ยทิ สํฆสฺส ปตฺตกลฺลํ, สํโฆ อิตฺถนฺนามสฺส ภิกฺขุโน ตาสํ อาปตฺตีนํ ยา อาปตฺติ ทสาห\-ปฺปฏิจฺฉนฺนา, ตสฺสา อคฺเฆน สโมธาน\-ปริวาสํ ทเทยฺย, เอสา ญตฺติ.}

\prose{สุณาตุ เม ภนฺเต สํโฆ, อยํ อิตฺถนฺนาโม ภิกฺขุ สมฺพหุลา สํฆาทิเสสา อาปตฺติโย อาปชฺชิ, เอกา อาปตฺติ เอกาห\-ปฺปฏิจฺฉนฺนา ฯเปฯ เอกา อาปตฺติ ทสาห\-ปฺปฏิจฺฉนฺนา, โส สํฆํ ตาสํ อาปตฺตีนํ ยา อาปตฺติ ทสาห\-ปฺปฏิจฺฉนฺนา, ตสฺสา อคฺเฆน สโมธาน\-ปริวาสํ ยาจติ. สํโฆ อิตฺถนฺนามสฺส ภิกฺขุโน ตาสํ อาปตฺตีนํ ยา อาปตฺติ ทสาห\-ปฺปฏิจฺฉนฺนา, ตสฺสา อคฺเฆน สโมธาน\-ปริวาสํ เทติ, ยสฺสายสฺมโต ขมติ อิตฺถนฺนามสฺส ภิกฺขุโน ตาสํ อาปตฺตีนํ ยา อาปตฺติ ทสาห\-ปฺปฏิจฺฉนฺนา, ตสฺสา อคฺเฆน สโมธาน\-ปริวาสสฺส ทานํ, โส ตุณฺหสฺส, ยสฺส นกฺขมติ, โส ภาเสยฺย.}

\prose{ทุติยมฺปิ เอตมตฺถํ วทามิ ฯเปฯ. ตติยมฺปิ เอตมตฺถํ วทามิ ฯเปฯ.}

\prose{ทินฺโน สํเฆน อิตฺถนฺนามสฺส ภิกฺขุโน ตาสํ อาปตฺตีนํ ยา อาปตฺติ ทสาห\-ปฺปฏิจฺฉนฺนา, ตสฺสา อคฺเฆน สโมธาน\-ปริวาโส, ขมติ สํฆสฺส, ตสฺมา ตุณฺหี, เอวเมตํ ธารยามี”ติ.}

\csromanmatikaanchor{mtk.83}
\csromantocmarkref{subsection}{สพฺพจิรปฺปฏิจฺฉนฺนอคฺฆสโมธาน}{mtk.83}
\bookmark[dest={mtk.83},level=2]{สพฺพจิรปฺปฏิจฺฉนฺนอคฺฆสโมธาน}
\csromanheader{2}{สพฺพจิรปฺปฏิจฺฉนฺน\-อคฺฆ\-สโมธาน}
\proseitem{136}{เตน โข ปน สมเยน อญฺญตโร ภิกฺขุ สมฺพหุลา สํฆาทิเสสา อาปตฺติโย อาปนฺโน โหติ, เอกา อาปตฺติ \csromanfolio{136}เอกาห\-ปฺปฏิจฺฉนฺนา, เทฺว อาปตฺติโย ทฺวีห\-ปฺปฏิจฺฉนฺนาโย\csromansharedfootnote{48f89f045589ccde}{ทฺวีหปฺปฏิจฺฉนฺนา (ก, เอวํ ยาวทสาหปฺปฏิจฺฉนฺนา)}, ติสฺโส อาปตฺติโย ตีหปฺปฏิจฺฉนฺนาโย, จตสฺโส อาปตฺติโย จตูห\-ปฺปฏิจฺฉนฺนาโย, ปญฺจ อาปตฺติโย ปญฺจาห\-ปฺปฏิจฺฉนฺนาโย, ฉ อาปตฺติโย ฉาห\-ปฺปฏิจฺฉนฺนาโย, สตฺต อาปตฺติโย สตฺตาห\-ปฺปฏิจฺฉนฺนาโย, อฏฺฐ อาปตฺติโย อฏฺฐาห\-ปฺปฏิจฺฉนฺนาโย, นว อาปตฺติโย นวาห\-ปฺปฏิจฺฉนฺนาโย, ทส อาปตฺติโย ทสาห\-ปฺปฏิจฺฉนฺนาโย, โส ภิกฺขูนํ อาโรเจสิ “อหํ อาวุโส สมฺพหุลา สํฆาทิเสสา อาปตฺติโย อาปชฺชิํ, เอกา อาปตฺติ เอกาห\-ปฺปฏิจฺฉนฺนา ฯเปฯ ทส อาปตฺติโย ทสาห\-ปฺปฏิจฺฉนฺนาโย, กถํ นุ โข มยา ปฏิปชฺชิตพฺพนฺ”ติ. ภควโต เอตมตฺถํ อาโรเจสุํ. เตน หิ ภิกฺขเว สํโฆ ตสฺส ภิกฺขุโน ตาสํ อาปตฺตีนํ ยา อาปตฺติโย สพฺพจิร\-ปฺปฏิจฺฉนฺนาโย, ตาสํ อคฺเฆน สโมธาน\-ปริวาสํ เทตุ. เอวญฺจ ปน ภิกฺขเว ทาตพฺโพ\csromandash{}}

\prose{เตน ภิกฺขเว ภิกฺขุนา สํฆํ อุปสงฺกมิตฺวา ฯเปฯ เอวมสฺส วจนีโย “อหํ ภนฺเต สมฺพหุลา สํฆาทิเสสา อาปตฺติโย อาปชฺชิํ, เอกา อาปตฺติ เอกาห\-ปฺปฏิจฺฉนฺนา ฯเปฯ ทส อาปตฺติโย ทสาห\-ปฺปฏิจฺฉนฺนาโย, โสหํ ภนฺเต สํฆํ ตาสํ อาปตฺตีนํ ยา อาปตฺติโย สพฺพจิรปฺปฏิจฺฉานฺนาโย, ตาสํ อคฺเฆน สโมธาน\-ปริวาสํ ยาจามี”ติ. ทุติยมฺปิ ยาจิตพฺโพ. ตติยมฺปิ ยาจิตพฺโพ. พฺยตฺเตน ภิกฺขุนา ปฏิพเลน ปฏิพเลน สํโฆ ญาเปตพฺโพ\csromandash{}}

\proseitem{137}{“สุณาตุ เม ภนฺเต สํโฆ, อยํ อิตฺถนฺนาโม ภิกฺขุ สมฺพหุลา สํฆาทิเสสา อาปตฺติโย อาปชฺชิ, เอกา อาปตฺติ เอกาห\-ปฺปฏิจฺฉนฺนา ฯเปฯ ทส อาปตฺติโย ทสาห\-ปฺปฏิจฺฉนฺนาโย, โส สํฆํ ตาสํ อาปตฺตีนํ ยา อาปตฺติโย สพฺพจิร\-ปฺปฏิจฺฉนฺนาโย, ตาสํ อคฺเฆน สโมธาน\-ปริวาสํ ยาจติ. ยทิ สํฆสฺส ปตฺตกลฺลํ, สํโฆ อิตฺถนฺนามสฺส ภิกฺขุโน ตาสํ อาปตฺตีนํ ยา อาปตฺติโย สพฺพจิร\-ปฺปฏิจฺฉนฺนาโย, ตาสํ อคฺเฆน สโมธาน\-ปริวาสํ ทเทยฺย, เอสา ญตฺติ.}

\prose{สุณาตุ เม ภนฺเต สํโฆ, อยํ อิตฺถนฺนาโม ภิกฺขุ สมฺพหุลา สํฆาทิเสสา อาปตฺติโย อาปชฺชิ, เอกา อาปตฺติ เอกาห\-ปฺปฏิจฺฉนฺนา. ทส อาปตฺติโย ทสาห\-ปฺปฏิจฺฉนฺนาโย, โส สํฆํ ตาสํ อาปตฺตีนํ \csromanfolio{137}ยา อาปตฺติโย สพฺพจิร\-ปฺปฏิจฺฉนฺนาโย, ตาสํ อคฺเฆน สโมธาน\-ปริวาสํ ยาจติ. สํโฆ อิตฺถนฺนามสฺส ภิกฺขุโน ตาสํ อาปตฺตีนํ ยา อาปตฺติโย สพฺพจิร\-ปฺปฏิจฺฉนฺนาโย, ตาสํ อคฺเฆน สโมธาน\-ปริวาสํ เทติ, ยสฺสายสฺมโต ขมติ อิตฺถนฺนามสฺส ภิกฺขุโน ตาสํ อาปตฺตีนํ ยา อาปตฺติโย สพฺพจิร\-ปฺปฏิจฺฉนฺนาโย, ตาสํ อคฺเฆน สโมธาน\-ปริวาสสฺส ทานํ, โส ตุณฺหสฺส, ยสฺส นกฺขมติ, โส ภาเสยฺย.}

\prose{ทุติยมฺปิ เอตมตฺถํ วทามิ ฯเปฯ. ตติยมฺปิ เอตมตฺถํ วทามิ ฯเปฯ.}

\prose{ทินฺโน สํเฆน อิตฺถนฺนามสฺส ภิกฺขุโน ตาสํ อาปตฺตีนํ ยา อาปตฺติโย สพฺพจิร\-ปฺปฏิจฺฉนฺนาโย, ตาสํ อคฺเฆน สโมธาน\-ปริวาโส, ขมติ สํฆสฺส, ตสฺมา ตุณฺหี, เอวเมตํ ธารยามี”ติ.}

\csromanmatikaanchor{mtk.84}
\csromantocmarkref{subsection}{เทฺวมาสปริวาส}{mtk.84}
\bookmark[dest={mtk.84},level=2]{เทฺวมาสปริวาส}
\csromanheader{2}{เทฺวมาส\-ปริวาส}
\proseitem{138}{เตน โข ปน สมเยน อญฺญตโร ภิกฺขุ เทฺว สํฆาทิเสสา อาปตฺติโย อาปนฺโน โหติ เทฺวมาส\-ปฺปฏิจฺฉนฺนาโย. ตสฺส เอตทโหสิ “อหํ โข เทฺว สํฆาทิเสสา อาปตฺติโย อาปชฺชิํ เทฺวมาส\-ปฺปฏิจฺฉนฺนาโย, ยํนูนาหํ สํฆํ เอกิสฺสา อาปตฺติยา เทฺวมาส\-ปฺปฏิจฺฉนฺนาย เทฺวมาส\-ปริวาสํ ยาเจยฺยนฺ”ติ. โส สํฆํ เอกิสฺสา อาปตฺติยา เทฺวมาส\-ปฺปฏิจฺฉนฺนาย เทฺวมาส\-ปริวาสํ ยาจิ, ตสฺส สํโฆ เอกิสฺสา อาปตฺติยา เทฺวมาส\-ปฺปฏิจฺฉนฺนาย เทฺวมาส\-ปริวาสํ อทาสิ. ตสฺส ปริวสนฺตสฺส ลชฺชีธมฺโม โอกฺกมิ\csromandash{}อหํ โข เทฺว สํฆาทิเสสา อาปตฺติโย อาปชฺชิํ เทฺวมาส\-ปฺปฏิจฺฉนฺนาโย, ตสฺส เม เอตทโหสิ “อหํ โข เทฺว สํฆาทิเสสา อาปตฺติโย อาปชฺชิํ เทฺวมาส\-ปฺปฏิจฺฉนฺนาโย, ยํนูนาหํ สํฆํ เอกิสฺสา อาปตฺติยา เทฺวมาส\-ปฺปฏิจฺฉนฺนาย เทฺวมาส\-ปริวาสํ ยาเจยฺยนฺ”ติ, โสหํ สํฆํ เอกิสฺสา อาปตฺติยา เทฺวมาส\-ปฺปฏิจฺฉนฺนาย เทฺวมาส\-ปริวาสํ ยาจิํ, ตสฺส เม สํโฆ เอกิสฺสา อาปตฺติยา เทฺวมาส\-ปฺปฏิจฺฉนฺนาย เทฺวมาส\-ปริวาสํ อทาสิ, ตสฺส เม ปริวสนฺตสฺส ลชฺชีธมฺโม โอกฺกมิ “ยํนูนาหํ สํฆํ อิตริสฺสาปิ อาปตฺติยา เทฺวมาส\-ปฺปฏิจฺฉนฺนาย เทฺวมาส\-ปริวาสํ ยาเจยฺยนฺ”ติ.}

\prose{โส ภิกฺขูนํ อาโรเจสิ “อหํ อาวุโส เทฺว สํฆาทิเสสา อาปตฺติโย อาปชฺชิํ เทฺวมาส\-ปฺปฏิจฺฉนฺนาโย, ตสฺส เม เอตทโหสิ \csromanfolio{138}“อหํ โข เทฺว สํฆาทิเสสา อาปตฺติโย อาปชฺชิํ เทฺวมาส\-ปฺปฏิจฺฉนฺนาโย, ยํนูนาหํ สํฆํ เอกิสฺสา อาปตฺติยา เทฺวมาส\-ปฺปฏิจฺฉนฺนาย เทฺวมาส\-ปริวาสํ ยาเจยฺยนฺ”ติ. โสหํ สํฆํ เอกิสฺสา อาปตฺติยา เทฺวมาส\-ปฺปฏิจฺฉนฺนาย เทฺวมาส\-ปริวาสํ ยาจิํ, ตสฺส เม สํโฆ เอกิสฺสา อาปตฺติยา เทฺวมาส\-ปฺปฏิจฺฉนฺนาย เทฺวมาส\-ปริวาสํ อทาสิ, ตสฺส เม ปริวสนฺตสฺส ลชฺชีธมฺโม โอกฺกมิ\csromandash{}อหํ โข เทฺว สํฆาทิเสสา อาปตฺติโย อาปชฺชิํ เทฺวมาส\-ปฺปฏิจฺฉนฺนาโย, ตสฺส เม เอตทโหสิ “อหํ โข เทฺว สํฆาทิเสสา อาปตฺติโย อาปชฺชิํ เทฺวมาส\-ปฺปฏิจฺฉนฺนาโย, ยํนูนาหํ สํฆํ เอกิสฺสา อาปตฺติยา เทฺวมาส\-ปฺปฏิจฺฉนฺนาย เทฺวมาส\-ปริวาสํ ยาเจยฺยนฺ”ติ, โสหํ สํฆํ เอกิสฺสา อาปตฺติยา เทฺวมาส\-ปฺปฏิจฺฉนฺนาย เทฺวมาส\-ปริวาสํ ยาจิํ, ตสฺส เม สํโฆ เอกิสฺสา อาปตฺติยา เทฺวมาส\-ปฺปฏิจฺฉนฺนาย เทฺวมาส\-ปริวาสํ อทาสิ, ตสฺส เม ปริวสนฺตสฺส ลชฺชีธมฺโม โอกฺกมิ “ยํนูนาหํ สํฆํ อิตริสฺสาปิ อาปตฺติยา เทฺวมาส\-ปฺปฏิจฺฉนฺนาย เทฺวมาส\-ปริวาสํ ยาเจยฺยนฺ”ติ. กถํ นุ โข มยา ปฏิปชฺชิตพฺพนฺติ. ภควโต เอตมตฺถํ อาโรเจสุํ. เตน หิ ภิกฺขเว สํโฆ ตสฺส ภิกฺขุโน อิตริสฺสาปิ อาปตฺติยา เทฺวมาส\-ปฺปฏิจฺฉนฺนาย เทฺวมาส\-ปริวาสํ เทตุ. เอวญฺจ ปน ภิกฺขเว ทาตพฺโพ\csromandash{}}

\prose{เตน ภิกฺขเว ภิกฺขุนา สํฆํ อุปสงฺกมิตฺวา เอกํสํ อุตฺตราสงฺคํ กริตฺวา วุฑฺฒานํ ภิกฺขูนํ ปาเท วนฺทิตฺวา อุกฺกุฏิกํ นิสีทิตฺวา อญฺชลิํ ปคฺคเหตฺวา เอวมสฺส วจนีโย “อหํ ภนฺเต เทฺว สํฆาทิเสสา อาปตฺติโย อาปชฺชิํ เทฺวมาส\-ปฺปฏิจฺฉนฺนาโย, ตสฺส เม เอตทโหสิ ‘อหํ โข เทฺว สํฆาทิเสสา อาปตฺติโย อาปชฺชิํ เทฺวมาส\-ปฺปฏิจฺฉนฺนาโย, ยํนูนาหํ สํฆํ เอกิสฺสา อาปตฺติยา เทฺวมาส\-ปฺปฏิจฺฉนฺนาย เทฺวมาส\-ปริวาสํ ยาเจยฺยนฺ’ติ, โสหํ สํฆํ เอกิสฺสา อาปตฺติยา เทฺวมาส\-ปฺปฏิจฺฉนฺนาย เทฺวมาส\-ปริวาสํ ยาจิํ, ตสฺส เม สํโฆ เอกิสฺสา อาปตฺติยา เทฺวมาส\-ปฺปฏิจฺฉนฺนาย เทฺวมาส\-ปริวาสํ อทาสิ, ตสฺส เม ปริวสนฺตสฺส ลชฺชีธมฺโม โอกฺกมิ\csromandash{}อหํ โข เทฺว สํฆาทิเสสา อาปตฺติโย อาปชฺชิํ เทฺวมาส\-ปฺปฏิจฺฉนฺนาโย, ตสฺส เม เอตทโหสิ ‘อหํ โข เทฺว สํฆาทิเสสา อาปตฺติโย อาปชฺชิํ เทฺวมาส\-ปฺปฏิจฺฉนฺนาโย, ยํนูนาหํ สํฆํ เอกิสฺสา อาปตฺติยา เทฺวมาส\-ปฺปฏิจฺฉนฺนาย เทฺวมาส\-ปริวาสํ ยาเจยฺยนฺ’ติ, โสหํ สํฆํ เอกิสฺสา อาปตฺติยา เทฺวมาส\-ปฺปฏิจฺฉนฺนาย \csromanfolio{139}เทฺวมาส\-ปริวาสํ ยาจิํ, ตสฺส เม สํโฆ เอกิสฺสา อาปตฺติยา เทฺวมาส\-ปฺปฏิจฺฉนฺนาย เทฺวมาส\-ปริวาสํ อทาสิ, ตสฺส เม ปริวสนฺตสฺส ลชฺชีธมฺโม โอกฺกมิ ‘ยํนูนาหํ สํฆํ อิตริสฺสาปิ อาปตฺติยา เทฺวมาส\-ปฺปฏิจฺฉนฺนาย เทฺวมาส\-ปริวาสํ ยาเจยฺยนฺ’ติ, โสหํ ภนฺเต สํฆํ อิตริสฺสาปิ อาปตฺติยา เทฺวมาส\-ปฺปฏิจฺฉนฺนาย เทฺวมาส\-ปริวาสํ ยาจามี”ติ.}

\prose{ทุติยมฺปิ ยาจิตพฺโพ. ตติยมฺปิ ยาจิตพฺโพ. พฺยตฺเตน ภิกฺขุนา ปฏิพเลน สํโฆ ญาเปตพฺโพ\csromandash{}}

\proseitem{139}{“สุณาตุ เม ภนฺเต สํโฆ, อยํ อิตฺถนฺนาโม ภิกฺขุ เทฺว สํฆาทิเสสา อาปตฺติโย อาปชฺชิ เทฺวมาส\-ปฺปฏิจฺฉนฺนาโย, ตสฺส เอตทโหสิ ‘อหํ โข เทฺว สํฆาทิเสสา อาปตฺติโย อาปชฺชิํ เทฺวมาส\-ปฺปฏิจฺฉนฺนาโย, ยํนูนาหํ สํฆํ เอกิสฺสา อาปตฺติยา เทฺวมาส\-ปฺปฏิจฺฉนฺนาย เทฺวมาส\-ปริวาสํ ยาเจยฺยนฺ’ติ, โส สํฆํ เอกิสฺสา อาปตฺติยา เทฺวมาส\-ปฺปฏิจฺฉนฺนาย เทฺวมาส\-ปริวาสํ ยาจิ, ตสฺส สํโฆ เอกิสฺสา อาปตฺติยา เทฺวมาส\-ปฺปฏิจฺฉนฺนาย เทฺวมาส\-ปริวาสํ อทาสิ, ตสฺส ปริวสนฺตสฺส ลชฺชีธมฺโม โอกฺกมิ\csromandash{}อหํ โข เทฺว สํฆาทิเสสา อาปตฺติโย อาปชฺชิํ เทฺวมาส\-ปฺปฏิจฺฉนฺนาโย, ตสฺส เม เอตทโหสิ ‘อหํ โข เทฺว สํฆาทิเสสา อาปตฺติโย อาปชฺชิํ เทฺวมาส\-ปฺปฏิจฺฉนฺนาโย, ยํนูนาหํ สํฆํ เอกิสฺสา อาปตฺติยา เทฺวมาส\-ปฺปฏิจฺฉนฺนาย เทฺวมาส\-ปริวาสํ ยาเจยฺยนฺ’ติ, โสหํ สํฆํ เอกิสฺสา อาปตฺติยา เทฺวมาส\-ปฺปฏิจฺฉนฺนาย เทฺวมาส\-ปริวาสํ ยาจิํ, ตสฺส เม สํโฆ เอกิสฺสา อาปตฺติยา เทฺวมาส\-ปฺปฏิจฺฉนฺนาย เทฺวมาส\-ปริวาสํ อทาสิ, ตสฺส เม ปริวสนฺตสฺส ลชฺชีธมฺโม โอกฺกมิ ‘ยํนูนาหํ สํฆํ อิตริสฺสาปิ อาปตฺติยา เทฺวมาส\-ปฺปฏิจฺฉนฺนาย เทฺวมาส\-ปริวาสํ ยาเจยฺยนฺ’ติ, โส สํฆํ อิตริสฺสาปิ อาปตฺติยา เทฺวมาส\-ปฺปฏิจฺฉนฺนาย เทฺวมาส\-ปริวาสํ ยาจติ, ยทิ สํฆสฺส ปตฺตกลฺลํ, สํโฆ อิตฺถนฺนามสฺส ภิกฺขุโน อิตริสฺสาปิ อาปตฺติยา เทฺวมาส\-ปฺปฏิจฺฉนฺนาย เทฺวมาส\-ปริวาสํ เทยฺย, เอสา ญตฺติ.}

\prose{สุณาตุ เม ภนฺเต สํโฆ, อยํ อิตฺถนฺนาโม ภิกฺขุ เทฺว สํฆาทิเสสา อาปตฺติโย อาปชฺชิ เทฺวมาส\-ปฺปฏิจฺฉนฺนาโย, ตสฺส เอตทโหสิ ‘อหํ โข เทฺว สํฆาทิเสสา อาปตฺติโย อาปชฺชิํ เทฺวมาส\-ปฺปฏิจฺฉนฺนาโย, ยํนูนาหํ สํฆํ เอกิสฺสา อาปตฺติยา \csromanfolio{140}เทฺวมาส\-ปฺปฏิจฺฉนฺนาย เทฺวมาส\-ปริวาสํ ยาเจยฺยนฺ’ติ, โส สํฆํ เอกิสฺสา อาปตฺติยา เทฺวมาส\-ปฺปฏิจฺฉนฺนาย เทฺวมาส\-ปริวาสํ ยาจิ, ตสฺส สํโฆ เอกิสฺสา อาปตฺติยา เทฺวมาส\-ปฺปฏิจฺฉนฺนาย เทฺวมาส\-ปริวาสํ อทาสิ, ตสฺส ปริวสนฺตสฺส ลชฺชีธมฺโม โอกฺกมิ\csromandash{}อหํ โข เทฺว สํฆาทิเสสา อาปตฺติโย อาปชฺชิํ เทฺวมาส\-ปฺปฏิจฺฉนฺนาโย, ตสฺส เม เอตทโหสิ ‘อหํ โข เทฺว สํฆาทิเสสา อาปตฺติโย อาปชฺชิํ เทฺวมาส\-ปฺปฏิจฺฉนฺนาโย, ยํนูนาหํ สํฆํ เอกิสฺสา อาปตฺติยา เทฺวมาส\-ปฺปฏิจฺฉนฺนาย เทฺวสปริวาสํ ยาเจยฺยนฺ’ติ, โสหํ สํฆํ เอกิสฺสา อาปตฺติยา เทฺวมาส\-ปฺปฏิจฺฉนฺนาย เทฺวมาส\-ปริวาสํ ยาจิํ, ตสฺส เม สํโฆ เอกิสฺสา อาปตฺติยา เทฺวมาส\-ปฺปฏิจฺฉนฺนาย เทฺวมาส\-ปริวาสํ อทาสิ, ตสฺส เม ปริวสนฺตสฺส ลชฺชีธมฺโม โอกฺกมิ ‘ยํนูนาหํ สํฆํ อิตริสฺสาปิ อาปตฺติยา เทฺวมาส\-ปฺปฏิจฺฉนฺนาย เทฺวมาส\-ปริวาสํ ยาเจยฺยนฺ’ติ, โส สํฆํ อิตริสฺสาปิ อาปตฺติยา เทฺวมาส\-ปฺปฏิจฺฉนฺนาย เทฺวมาส\-ปริวาสํ ยาจติ, สํโฆ อิตฺถนฺนามสฺส ภิกฺขุโน อิตริสฺสาปิ อาปตฺติยา เทฺวมาส\-ปฺปฏิจฺฉนฺนาย เทฺวมาส\-ปริวาสํ เทติ, ยสฺสายสฺมโต ขมติ อิตฺถนฺนามสฺส ภิกฺขุโน อิตริสฺสาปิ อาปตฺติยา เทฺวมาส\-ปฺปฏิจฺฉนฺนาย เทฺวมาส\-ปริวาสสฺส ทานํ, โส ตุณฺหสฺส, ยสฺส นกฺขมติ, โส ภาเสยฺย.}

\prose{ทุติยมฺปิ เอตมตฺถํ วทามิ ฯเปฯ. ตติยมฺปิ เอตมตฺถํ วทามิ ฯเปฯ.}

\prose{ทินฺโน สํเฆน อิตฺถนฺนามสฺส ภิกฺขุโน อิตริสฺสาปิ อาปตฺติยา เทฺวมาส\-ปฺปฏิจฺฉนฺนาย เทฺวมสปริวาโส, ขมติ สํฆสฺส, ตสฺมา ตุณฺหี, เอวเมตํ ธารยามี”ติ. เตน ภิกฺขเว ภิกฺขุนา ตทุปาทาย เทฺว มาสา ปริวสิตพฺพา.}

\csromanmatikaanchor{mtk.85}
\csromantocmarkref{subsection}{เทฺว มาสา ปริวสิตพฺพวิธิ}{mtk.85}
\bookmark[dest={mtk.85},level=2]{เทฺว มาสา ปริวสิตพฺพวิธิ}
\csromanheader{2}{\textbf{เทฺว มาสา ปริวสิตพฺพวิธิ}}
\proseitem{140}{อิธ ปน ภิกฺขเว ภิกฺขุ เทฺว สํฆาทิเสสา อาปตฺติโย อาปชฺชติ เทฺวมาส\-ปฺปฏิจฺฉนฺนาโย, ตสฺส เอวํ โหติ “อหํ โข เทฺว สํฆาทิเสสา อาปตฺติโย อาปชฺชิํ เทฺวมาส\-ปฺปฏิจฺฉนฺนาโย, ยํนุนาหํ สํฆํ เอกิสฺสา อาปตฺติยา เทฺวมาส\-ปฺปฏิจฺฉนฺนาย เทฺวมาส\-ปริวาสํ ยาเจยฺยนฺ”ติ, โส สํฆํ เอกิสฺสา อาปตฺติยา เทฺวมาส\-ปฺปฏิจฺฉนฺนาย เทฺวมาส\-ปริวาสํ ยาจติ, ตสฺส สํโฆ เอกิสฺสา อาปตฺติยา เทฺวมาส\-ปฺปฏิจฺฉนฺนาย เทฺวมาส\-ปริวาสํ เทติ, ตสฺส ปริวสนฺตสฺส ลชฺชีธมฺโม \csromanfolio{141}โอกฺกมิ\csromandash{}อหํ โข เทฺว สํฆาทิเสสา อาปตฺติโย อาปชฺชิํ เทฺวมาส\-ปฺปฏิจฺฉนฺนาโย, ตสฺส เม เอตทโหสิ “อหํ โข เทฺว สํฆาทิเสสา อาปตฺติโย อาปชฺชิํ เทฺวมาส\-ปฺปฏิจฺฉนฺนาโย, ยํนูนาหํ สํฆํ เอกิสฺสา อาปตฺติยา เทฺวมาส\-ปฏิจฺฉนฺนาย เทฺวมาส\-ปริวาสํ ยาเจยฺยนฺ”ติ, โสหํ สํฆํ เอกิสฺสา อาปตฺติยา เทฺวมาส\-ปฺปฏิจฺฉนฺนาย เทฺวมาส\-ปริวาสํ ยาจิํ, ตสฺส เม สํโฆ เอกิสฺสา อาปตฺติยา เทฺวมาส\-ปฺปฏิจฺฉนฺนาย เทฺวมาส\-ปริวาสํ อทาสิ, ตสฺส เม ปริวสนฺตสฺส ลชฺชีธมฺโม โอกฺกมิ ‘ยํนุนาหํ สํฆํ อิตริสฺสาปิ อาปตฺติยา เทฺวมาส\-ปฺปฏิจฺฉนฺนาย เทฺวมาส\-ปริวาสํ ยาเจยฺยนฺ”ติ, โส สํฆํ อิตริสฺสาปิ อาปตฺติยา เทฺวมาส\-ปฺปฏิจฺฉนฺนาย เทฺวมาส\-ปริวาสํ ยาจติ, ตสฺส สํโฆ อิตริสฺสาปิ อาปตฺติยา เทฺวมาส\-ปฺปฏิจฺฉนฺนาย เทฺวมาส\-ปริวาสํ เทติ. เตน ภิกฺขเว ภิกฺขุนา ตทุปาทาย เทฺว มาสา ปริวสิตพฺพา.}

\proseitem{141}{อิธ ปน ภิกฺขเว ภิกฺขุ เทฺว สํฆาทิเสสา อาปตฺติโย อาปชฺชติ เทฺวมาส\-ปฺปฏิจฺฉนฺนาโย, เอกํ อาปตฺติํ ชานาติ, เอกํ อาปตฺติํ น ชานาติ. โส สํฆํ ยํ อาปตฺติํ ชานาติ, ตสฺสา อาปตฺติยา เทฺวมาส\-ปฺปฏิจฺฉนฺนาย เทฺวมาส\-ปริวาสํ ยาจติ, ตสฺส สํโฆ ตสฺสา อาปตฺติยา เทฺวมาส\-ปฺปฏิจฺฉนฺนาย เทฺวมาส\-ปริวาสํ เทติ, โส ปริวสนฺโต อิตรมฺปิ อาปตฺติํ ชานาติ, ตสฺส เอวํ โหติ “อหํ โข เทฺว สํฆาทิเสสา อาปตฺติโย อาปชฺชิํ เทฺวมาส\-ปฺปฏิจฺฉนฺนาโย, เอกํ อาปตฺติ ชานิํ, เอกํ อาปตฺติํ น ชานิํ, โสหํ สํฆํ ยํ อาปตฺติํ ชานิํ, ตสฺสา อาปตฺติยา เทฺวมาส\-ปฺปฏิจฺฉนฺนาย เทฺวมาส\-ปริวาสํ ยาจิํ, ตสฺส เม สํโฆ ตสฺสา อาปตฺติยา เทฺวมาส\-ปฺปฏิจฺฉนฺนาย เทฺวมาส\-ปริวาสํ อทาสิ, โสหํ ปริวสนฺโต อิตรมฺปิ อาปตฺติํ ชานามิ, ยํนุนาหํ สํฆํ อิตริสฺสาปิ อาปตฺติยา เทฺวมาส\-ปฺปฏิจฺฉนฺนาย เทฺวมาส\-ปริวาสํ ยาเจยฺยนฺ”ติ, โส สํฆํ อิตริสฺสาปิ อาปตฺติยา เทฺวมาส\-ปฺปฏิจฺฉนฺนาย เทฺวมาส\-ปริวาสํ ยาจติ, ตสฺส สํโฆ อิตริสฺสาปิ อาปตฺติยา เทฺวมาส\-ปฺปฏิจฺฉนฺนาย เทฺวมาส\-ปริวาสํ เทติ. เตน ภิกฺขเว ภิกฺขุนา ตทุปาทาย เทฺวมาสา ปริวสิตพฺพา.}

\proseitem{142}{อิธ ปน ภิกฺขเว ภิกฺขุ เทฺว สํฆาทิเสสา อาปตฺติโย อาปชฺชติ เทฺวมาส\-ปฺปฏิจฺฉนฺนาโย, เอกํ อาปตฺติํ สรติ, เอกํ อาปตฺติํ นสฺสารติ. โส สํฆํ ยํ อาปตฺติํ สรติ, ตสฺสา อาปตฺติยา \csromanfolio{142}เทฺวมาส\-ปฺปฏิจฺฉนฺนาย เทฺวมาส\-ปริวาสํ ยาจติ, ตสฺส สํโฆ ตสฺสา อาปตฺติยา เทฺวมาส\-ปฺปฏิจฺฉนฺนาย เทฺวมาส\-ปริวาสํ เทติ, โส ปริวสนฺโต อิตรมฺปิ อาปตฺติํ สรติ, ตสฺส เอวํ โหติ “อหํ โข เทฺว สํฆาทิเสสา อาปตฺติโย อาปชฺชิํ เทฺวมาส\-ปฺปฏิจฺฉนฺนาโย, เอกํ อาปตฺติํ สริํ, เอกํ อาปตฺติํ นสฺสริํ, โสหํ สํฆํ ยํ อาปตฺติํ สริํ, ตสฺสา อาปตฺติยา เทฺวมาส\-ปฺปฏิจฺฉนฺนาย เทฺวมาส\-ปริวาสํ ยาจิํ, ตสฺส เม สํโฆ ตสฺสา อาปตฺติยา เทฺวมาส\-ปฺปฏิจฺฉนฺนาย เทฺวมาส\-ปริวาสํ อทาสิ, โสหํ ปริวสนฺโต อิตรมฺปิ อาปตฺติํ สรามิ, ยํนูนาหํ สํฆํ อิตริสฺสาปิ อาปตฺติยา เทฺวมาส\-ปฺปฏิจฺฉนฺนาย เทฺวมาส\-ปริวาสํ ยาเจยฺยนฺ”ติ, โส สํฆํ อิตริสฺสาปิ อาปตฺติยา เทฺวมาส\-ปฺปฏิจฺฉนฺนาย เทฺวมาส\-ปริวาสํ ยาจติ, ตสฺส สํโฆ อิตริสฺสาปิ อาปตฺติยา เทฺวมาส\-ปฺปฏิจฺฉนฺนาย เทฺวมาส\-ปริวาสํ เทติ. เตน ภิกฺขเว ภิกฺขุนา ตทุปาทาย เทฺว มาสา ปริวสิตพฺพา.}

\proseitem{143}{อิธ ปน ภิกฺขเว ภิกฺขุ เทฺว สํฆาทิเสสา อาปตฺติโย อาปชฺชติ เทฺวมาส\-ปฺปฏิจฺฉนฺนาโย, เอกาย อาปตฺติยา นิพฺเพมติโก, เอกาย อาปตฺติยา เวมติโก. โส สํฆํ ยาย อาปตฺติยา นิพฺเพมติโก, ตสฺสา อาปตฺติยา เทฺวมาส\-ปฺปฏิจฺฉนฺนาย เทฺวมาส\-ปริวาสํ ยาจติ, ตสฺส สํโฆ ตสฺสา อาปตฺติยา เทฺวมาส\-ปฺปฏิจฺฉนฺนาย เทฺวมาส\-ปริวาสํ เทติ, โส ปริวสนฺโต อิตริสฺสาปิ อาปตฺติยา นิพฺเพมติโก โหติ, ตสฺส เอวํ โหติ “อหํ โข เทฺว สํฆาทิเสสา อาปตฺติโย อาปชฺชิํ เทฺวมาส\-ปฺปฏิจฺฉนฺนาโย, เอกาย อาปตฺติยา นิพฺเพมติโก, เอกาย อาปตฺติยา เวมติโก, โสหํ สํฆํ ยาย อาปตฺติยา นิพฺเพมติโก, ตสฺสา อาปตฺติยา เทฺวมาส\-ปฺปฏิจฺฉนฺนาย เทฺวมาส\-ปริวาสํ ยาจิํ, ตสฺส เม สํโฆ ตสฺสา อาปตฺติยา เทฺวมาส\-ปฺปฏิจฺฉนฺนาย เทฺวมาส\-ปริวาสํ อทาสิ, โสหํ ปริวสนฺโต อิตริสฺสาปิ อาปตฺติยา นิพฺเพมติโก, ยํนูนาหํ สํฆํ อิตริสฺสาปิ อาปตฺติยา เทฺวมาส\-ปฺปฏิจฺฉนฺนาย เทฺวมาส\-ปริวาสํ ยาเจยฺยนฺ”ติ, โส สํฆํ อิตริสฺสาปิ อาปตฺติยา เทฺวมาส\-ปฺปฏิจฺฉนฺนาย เทฺวมาส\-ปริวาสํ ยาจติ, ตสฺส สํโฆ อิตริสฺสาปิ อาปตฺติยา เทฺวมาส\-ปฺปฏิจฺฉนฺนาย เทฺวมาส\-ปริวาสํ เทติ. เตน ภิกฺขเว ภิกฺขุนา ตทุปาทาย เทฺว มาสา ปริวสิตพฺพา.}

\proseitem{144}{\csromanfolio{143}อิธ ปน ภิกฺขเว ภิกฺขุ เทฺว สํฆาทิเสสา อาปตฺติโย อาปชฺชติ เทฺวมาส\-ปฺปฏิจฺฉนฺนาโย, เอกา อาปตฺติ ชาน\-ปฺปฏิจฺฉนฺนา, เอกา อาปตฺติ อชาน\-ปฺปฏิจฺฉนฺนา. โส สํฆํ ตาสํ อาปตฺตีนํ เทฺวมาส\-ปฺปฏิจฺฉนฺนานํ เทฺวมาส\-ปริวาสํ ยาจติ, ตสฺส สํโฆ ตาสํ อาปตฺตีนํ เทฺวมาส\-ปฺปฏิจฺฉนฺนานํ เทฺวมาส\-ปริวาสํ เทติ, ตสฺส ปริวสนฺตสฺส อญฺโญ ภิกฺขุ อาคจฺฉติ พหุสฺสุโต อาคตาคโม ธมฺมธโร วินยธโร มาติกาธโร ปณฺฑิโต วิยตฺโต เมธาวี ลชฺชี กุกฺกุจฺจโก สิกฺขากาโม, โส เอวํ วเทติ “กิํ อยํ อาวุโส ภิกฺขุ อาปนฺโน กิสฺสายํ ภิกฺขุ ปริวสตี”ติ. เต เอวํ วเทนฺติ “อยํ อาวุโส ภิกฺขุ เทฺว สํฆาทิเสสา อาปตฺติโย อาปชฺชิ เทฺวมาส\-ปฺปฏิจฺฉนฺนาโย, เอกา อาปตฺติ ชาน\-ปฺปฏิจฺฉนฺนา, เอกา อาปตฺติ อชาน\-ปฺปฏิจฺฉนฺนา, โส สํฆํ ตาสํ อาปตฺตีนํ เทฺวมาส\-ปฺปฏิจฺฉนฺนานํ เทฺวมาส\-ปริวาสํ ยาจติ, ตสฺส สํโฆ ตาสํ อาปตฺตีนํ เทฺวมาส\-ปฺปฏิจฺฉนฺนานํ เทฺวมาส\-ปริวาสํ อทาสิ, ตาโย อยํ อาวุโส ภิกฺขุ อาปนฺโน, ตาสายํ ภิกฺขุ ปริวสตี”ติ. โส เอวํ วเทติ “ยายํ อาวุโส อาปตฺติ ชาน\-ปฺปฏิจฺฉนฺนา, ธมฺมิกํ ตสฺสา อาปตฺติยา ปริวาสทานํ, ธมฺมตา รุหติ. ยา จ ขฺวายํ อาวุโส อาปตฺติ อชาน\-ปฺปฏิจฺฉนฺนา, อธมฺมิกํ ตสฺสา อาปตฺติยา ปริวาสทานํ, อธมฺมตฺตา น รุหติ. เอกิสฺสา อาวุโส อาปตฺติยา ภิกฺขุ มานตฺตารโห”ติ.}

\proseitem{145}{อิธ ปน ภิกฺขเว ภิกฺขุ เทฺว สํฆาทิเสสา อาปตฺติโย อาปชฺชติ เทฺวมาส\-ปฺปฏิจฺฉนฺนาโย, เอกา อาปตฺติ สรมาน\-ปฺปฏิจฺฉนฺนา, เอกา อาปตฺติ อสฺสรมาน\-ปฺปฏิจฺฉนฺนา. โส สํฆํ ตาสํ อาปตฺตีนํ เทฺวมาส\-ปฺปฏิจฺฉนฺนานํ เทฺวมาส\-ปริวาสํ ยาจติ, ตสฺส สํโฆ ตาสํ อาปตฺตีนํ เทฺวมาส\-ปฺปฏิจฺฉนฺนานํ เทฺวมาส\-ปริวาสํ เทติ, ตสฺส ปริวสนฺตสฺส อญฺโญ ภิกฺขุ อาคจฺฉติ พหุสฺสุโต อาคตาคโม ธมฺมธโร วินยธโร มาติกาธโร ปณฺฑิโต วิยตฺโต เมธาวี ลชฺชี กุกฺกุจฺจโก สิกฺขากาโม, โส เอวํ วเทติ “กิํ อยํ อาวุโส ภิกฺขุ อาปนฺโน, กิสฺสายํ ภิกฺขุ ปริวสตี”ติ. เต เอวํ วเทนฺติ “อยํ อาวุโส ภิกฺขุ เทฺว สํฆาทิเสสา อาปตฺติโย อาปชฺชิ เทฺวมาส\-ปฺปฏิจฺฉนฺนาโย, เอกา อาปตฺติ สรมาน\-ปฺปฏิจฺฉนฺนา, เอกา อาปตฺติ อสฺสรมาน\-ปฺปฏิจฺฉนฺนา, โส สํฆํ ตาสํ อาปตฺตีนํ เทฺวมาส\-ปฺปฏิจฺฉนฺนานํ เทฺวมาส\-ปริวาสํ ยาจิ, ตสฺส สํโฆ ตาสํ อาปตฺตีนํ เทฺวมาส\-ปฺปฏิจฺฉนฺนานํ เทฺวมาส\-ปริวาสํ อทาสิ, \csromanfolio{144}ตาโย อยํ อาวุโส ภิกฺขุ อาปนฺโน, ตาสายํ ภิกฺขุ ปริวสตี”ติ. โส เอวํ วเทติ “ยายํ อาวุโส อาปตฺติ สรมาน\-ปฺปฏิจฺฉนฺนา, ธมฺมิกํ ตสฺสา อาปตฺติยา ปริวาสทานํ, ธมฺมตฺตา รุหติ. ยา จ ขฺวายํ อาวุโส อาปตฺติ อสฺสรมาน\-ปฺปฏิจฺฉนฺนา, อธมฺมิกํ ตสฺส อาปตฺติยา ปริวาสทานํ, อธมฺมตฺตา น รุหติ. เอกิสฺสา อาวุโส อาปตฺติยา ภิกฺขุ มานตฺตารโห”ติ.}

\proseitem{146}{อิธ ปน ภิกฺขเว ภิกฺขุ เทฺว สํฆาทิเสสา อาปตฺติโย อาปชฺชติ เทฺวมาส\-ปฺปฏิจฺฉนฺนาโย, เอกา อาปตฺติ นิพฺเพมติก\-ปฺปฏิจฺฉนฺนา, เอกา อาปตฺติ เวมติก\-ปฺปฏิจฺฉนฺนา. โส สํฆํ ตาสํ อาปตฺตีนํ เทฺวมาส\-ปฺปฏิจฺฉนฺนานํ เทฺวมาส\-ปริวาสํ ยาจติ, ตสฺส สํโฆ ตาสํ อาปตฺตีนํ เทฺวมาส\-ปฺปฏิจฺฉนฺนานํ เทฺวมาส\-ปริวาสํ เทติ, ตสฺส ปริวสนฺตสฺส อญฺโญ ภิกฺขุ อาคจฺฉติ พหุสฺสุโต อาคตาคโม ธมฺมธโร วินยธโร มาติกาธโร ปณฺฑิโต ปิยตฺโต เมธาวี ลชฺชี กุกฺกุจฺจโก สิกฺขากาโม, โส เอวํ วเทติ “กิํ อยํ อาวุโส ภิกฺขุ อาปนฺโน, กิสฺสายํ ติกฺขุ ปริวสตี”ติ. เต เอวํ วเทนฺติ “อยํ อาวุโส ภิกฺขุ เทฺว สํฆาทิเสสา อาปตฺติโย อาปชฺชิ เทฺวมาส\-ปฺปฏิจฺฉนฺนาโย, เอกา อาปตฺติ นิพฺเพมติก\-ปฺปฏิจฺฉนฺนา, เอกา อาปตฺติ เวมติก\-ปฺปฏิจฺฉนฺนา, โส สํฆํ ตาสํ อาปตฺตีนํ เทฺวมาส\-ปฺปฏิจฺฉนฺนานํ เทฺวมาส\-ปริวาสํ ยาจิ, ตสฺส สํโฆ ตาสํ อาปตฺตีนํ เทฺวมาส\-ปฺปฏิจฺฉนฺนานํ เทฺวมาส\-ปริวาสํ อทาสิ, ตาโย อยํ อาวุโส ภิกฺขุ อาปนฺโน, ตาสายํ ภิกฺขุ ปริวสตี”ติ. โส เอวํ วเทติ “ยายํ อาวุโส อาปตฺติ นิพฺเพมติก\-ปฺปฏิจฺฉนฺนา, ธมฺมิกํ ตสฺสา อาปตฺติยา ปริวาสทานํ, ธมฺมตฺตา รุหติ. ยา จ ขฺวายํ อาวุโส อาปตฺติ เวมติก\-ปฺปฏิจฺฉนฺนา, อธมฺมิกํ ตสฺสา อาปตฺติยา ปริวาสทานํ, อธมฺมตฺตา น รุหติ. เอกิสฺสา อาวุโส อาปตฺติยา ภิกฺขุ มานตฺตารโห”ติ.}

\proseitem{147}{เตน โข ปน สมเยน อญฺญตโร ภิกฺขุ เทฺว สํฆาทิเสสา อาปตฺติโย อาปนฺโน โหติ เทฺวมาส\-ปฺปฏิจฺฉนฺนาโย. ตสฺส เอตทโหสิ “อหํ โข เทฺว สํฆาทิเสสา อาปตฺติโย อาปชฺชิํ เทฺวมาส\-ปฺปฏิจฺฉนฺนาโย, ยํนูนาหํ สํฆํ ทฺวินฺนํ อาปตฺตีนํ เทฺวมาส\-ปฺปฏิจฺฉนฺนานํ เอกมาส\-ปริวาสํ ยาเจยฺยนฺ”ติ. โส สํฆํ ทฺวินฺนํ อาปตฺตีนํ เทฺวมาส\-ปฺปฏิจฺฉนฺนานํ \csromanfolio{145}เอกมาส\-ปริวาสํ ยาจิ, ตสฺส สํโฆ ทฺวินฺนํ อาปตฺตีนํ เทฺวมาส\-ปฺปฏิจฺฉนฺนานํ เอกมาส\-ปริวาสํ อทาสิ, ตสฺส ปริวสนฺตสฺส ลชฺชีธมฺโม โอกฺกมิ\csromandash{}อหํ โข เทฺว สํฆาทิเสสา อาปตฺติโย อาปชฺชิํ เทฺวมาส\-ปฺปฏิจฺฉนฺนาโย, ตสฺส เม เอตทโหสิ “อหํ โข เทฺว สํฆาทิเสสา อาปตฺติโย อาปชฺชิํ เทฺวมาส\-ปฺปฏิจฺฉนฺนาโย, ยํนูนาหํ สํฆํ ทฺวินฺนํ อาปตฺตีนํ เทฺวมาส\-ปฺปฏิจฺฉนฺนานํ เอกมาส\-ปริวาสํ ยาเจยฺยนฺ”ติ, โสหํ สํฆํ ทฺวินฺนํ อาปตฺตีนํ เทฺวมาส\-ปฺปฏิจฺฉนฺนานํ เอกมาส\-ปริวาสํ ยาจิํ, ตสฺส เม สํโฆ ทฺวินฺนํ อาปตฺตีนํ เทฺวมาส\-ปฺปฏิจฺฉนฺนานํ เอกมาส\-ปริวาสํ อทาสิ, ตสฺส เม ปริวสนฺตสฺส ลชฺชีธมฺโม โอกฺกมิ “ยํนูนาหํ สํฆํ ทฺวินฺนํ อาปตฺตีนํ เทฺวมาส\-ปฺปฏิจฺฉนฺนานํ อิตรมฺปิ มาสํ ปริวาสํ\csromansharedfootnote{9ef3d38a5c5d4030}{อิตรมฺปิ มาสปริวาสํ (สฺยา, ก, เอวมุปริปิ)} ยาเจยฺยนฺ”ติ.}

\prose{โส ภิกฺขูนํ อาโรเจสิ “อหํ โข อาวุโส เทฺว สํฆาทิเสสา อาปตฺติโย อาปชฺชิํ เทฺวมาส\-ปฺปฏิจฺฉนฺนาโย, ตสฺส เม เอตทโหสิ ‘อหํ โข เทฺว สํฆาทิเสสา อาปตฺติโย อาปชฺชิํ เทฺวมาส\-ปฺปฏิจฺฉนฺนาโย, ยํนูนาหํ สํฆํ ทฺวินฺนํ อาปตฺตีนํ เทฺวมาส\-ปฺปฏิจฺฉนฺนานํ เอกมาส\-ปริวาสํ ยาเจยฺยนฺ’ติ, โสหํ สํฆํ ทฺวินฺนํ อาปตฺตีนํ เทฺวมาส\-ปฺปฏิจฺฉนฺนานํ เอกมาส\-ปริวาสํ ยาจิํ, ตสฺส เม สํโฆ ทฺวินฺนํ อาปตฺตีนํ เทฺวมาส\-ปฺปฏิจฺฉนฺนานํ เอกมาส\-ปริวาสํ อทาสิ, ตสฺส เม ปริวสนฺตสฺส ลชฺชีธมฺโม โอกฺกมิ\csromandash{}อหํ โข เทฺว สํฆาทิเสสา อาปตฺติโย อาปชฺชิํ เทฺวมาส\-ปฺปฏิจฺฉนฺนาโย, ตสฺส เม เอตทโหสิ ‘อหํ โข เทฺว สํฆาทิเสสา อาปตฺติโย อาปชฺชิํ เทฺวมาส\-ปฺปฏิจฺฉนฺนาโย, ยํนูนาหํ สํฆํ ทฺวินฺนํ อาปตฺตีนํ เทฺวมาส\-ปฺปฏิจฺฉนฺนานํ เอกมาส\-ปริวาสํ ยาเจยฺยนฺ’ติ, โสหํ สํฆํ ทฺวินฺนํ อาปตฺตีนํ เทฺวมาส\-ปฺปฏิจฺฉนฺนานํ เอกมาส\-ปริวาสํ ยาจิํ, ตสฺส เม สํโฆ ทฺวินฺนํ อาปตฺตีนํ เทฺวมาส\-ปฺปฏิจฺฉนฺนานํ เอกมาส\-ปริวาสํ อทาสิ, ตสฺส เม ปริวสนฺตสฺส ลชฺชีธมฺโม โอกฺกมิ ‘ยํนูนาหํ สํฆํ ทฺวินฺนํ อาปตฺตีนํ เทฺวมาส\-ปฺปฏิจฺฉนฺนานํ อิตรมฺปิ มาสํ ปริวาสํ ยาเจยฺยนฺ’ติ. กถํ นุ โข มยา ปฏิปชฺชิตพฺพนฺ”ติ. ภควโต เอตมตฺถํ อาโรเจสุํ. เตน หิ ภิกฺขเว สํโฆ ตสฺส ภิกฺขุโน ทฺวินฺนํ อาปตฺตีนํ เทฺวมาส\-ปฺปฏิจฺฉนฺนานํ อิตรมฺปิ มาสํ ปริวาสํ เทตุ. เอวญฺจ ปน ภิกฺขเว ทาตพฺโพ\csromandash{}}

\prose{เตน ภิกฺขเว ภิกฺขุนา สํฆํ อุปสงฺกมิตฺวา ฯเปฯ เอวมสฺส วจนีโย “อหํ ภนฺเต เทฺว สํฆาทิเสสา อาปตฺติโย อาปชฺชิํ \csromanfolio{146}เทฺวมาส\-ปฺปฏิจฺฉนฺนาโย, ตสฺส เม เอตทโหสิ ‘อหํ โข เทฺว สํฆาทิเสสา อาปตฺติโย อาปชฺชิํ เทฺวมาส\-ปฺปฏิจฺฉนฺนาโย, ยํนูนาหํ สํฆํ ทฺวินฺนํ อาปตฺตีนํ เทฺวมาส\-ปฺปฏิจฺฉนฺนานํ เอกมาส\-ปริวาสํ ยาเจยฺยนฺ’ติ, โสหํ สํฆํ ทฺวินฺนํ อาปตฺตีนํ เทฺวมาส\-ปฺปฏิจฺฉนฺนานํ เอกมาส\-ปริวาสํ ยาจิํ, ตสฺส เม สํโฆ ทฺวินฺนํ อาปตฺตีนํ เทฺวมาส\-ปฺปฏิจฺฉนฺนานํ เอกมาส\-ปริวาสํ อทาสิ, ตสฺส เม ปริวสนฺตสฺส ลชฺชีธมฺโม โอกฺกมิ\csromandash{} อหํ โข เทฺว สํฆาทิเสสา อาปตฺติโย อาปชฺชิํ เทฺวมาส\-ปฺปฏิจฺฉนฺนาโย, ตสฺส เม เอตทโหสิ ‘อหํ โข เทฺว สํฆาทิเสสา อาปตฺติโย อาปชฺชิํ เทฺวมาส\-ปฺปฏิจฺฉนฺนาโย ยํนูนาหํ สํฆํ ทฺวินฺนํ อาปตฺตีนํ เทฺวมาส\-ปฺปฏิจฺฉนฺนานํ เอกมาส\-ปริวาสํ ยาเจยฺยนฺ’ติ, โสหํ สํฆํ ทฺวินฺนํ อาปตฺตีนํ เทฺวมาส\-ปฺปฏิจฺฉนฺนานํ เอกมาส\-ปริวาสํ ยาจิํ, ตสฺส เม สํโฆ ทฺวินฺนํ อาปตฺตีนํ เทฺวมาส\-ปฺปฏิจฺฉนฺนานํ เอกมาส\-ปริวาสํ อทาสิ, ตสฺส เม ปริวสนฺตสฺส ลชฺชีธมฺโม โอกฺกมิ ‘ยํนูนาหํ สํฆํ ทฺวินฺนํ อาปตฺตีนํ เทฺวมาส\-ปฺปฏิจฺฉนฺนานํ อิตรมฺปิ มาสํ ปริวาสํ ยาเจยฺยนฺ’ติ, โสหํ ภนฺเต สํฆํ ทฺวินฺนํ อาปตฺตีนํ เทฺวมาส\-ปฺปฏิจฺฉนฺนานํ อิตรมฺปิ มาสํ ปริวาสํ ยาจามี”ติ. ทุติยมฺปิ ยาจิตพฺโพ. ตติยมฺปิ ยาจิตพฺโพ. พฺยตฺเตน ภิกฺขุนา ปฏิพเลน สํโฆ ญาเปตพฺโพ\csromandash{}}

\proseitem{148}{“สุณาตุ เม ภนฺเต สํโฆ, อยํ อิตฺถนฺนาโม ภิกฺขุ เทฺว สํฆาทิเสสา อาปตฺติโย อาปชฺชิ เทฺวมาส\-ปฺปฏิจฺฉนฺนาโย, ตสฺส เอตทโหสิ ‘อหํ โข เทฺว สํฆาทิเสสา อาปตฺติโย อาปชฺชิํ เทฺวมาส\-ปฺปฏิจฺฉนฺนาโย, ยํนูนาหํ สํฆํ ทฺวินฺนํ อาปตฺตีนํ เทฺวมาส\-ปฺปฏิจฺฉนฺนานํ เอกมาส\-ปริวาสํ ยาเจยฺยนฺ’ติ, โส สํฆํ ทฺวินฺนํ อาปตฺตีนํ เทฺว มาสปฺปฏิจฺฉนฺนานํ เอกมาส\-ปริวาสํ ยาจิ, ตสฺส สํโฆ ทฺวินฺนํ อาปตฺตีนํ เทฺวมาส\-ปฺปฏิจฺฉนฺนานํ เอกมาส\-ปริวาสํ อทาสิ, ตสฺส ปริวสนฺตสฺส ลชฺชีธมฺโม โอกฺกมิ\csromandash{}อหํ โข เทฺว สํฆาทิเสสา อาปตฺติโย อาปชฺชิํ เทฺวมาส\-ปฺปฏิจฺฉนฺนาโย, ตสฺส เม เอตทโหสิ ‘อหํ โข เทฺว สํฆาทิเสสา อาปตฺติโย อาปชฺชิํ เทฺวมาส\-ปฺปฏิจฺฉนฺนาโย, ยํนูนาหํ สํฆํ ทฺวินฺนํ อาปตฺตีนํ เทฺวมาส\-ปฺปฏิจฺฉนฺนานํ เอกมาส\-ปริวาสํ ยาเจยฺยนฺ’ติ, โสหํ สํฆํ ทฺวินฺนํ อาปตฺตีนํ เทฺวมาส\-ปฺปฏิจฺฉนฺนานํ เอกมาส\-ปริวาสํ ยาจิํ, ตสฺส เม สํโฆ ทฺวินฺนํ อาปตฺตีนํ เทฺวมาส\-ปฺปฏิจฺฉนฺนานํ เอกมาส\-ปริวาสํ อทาสิ, ตสฺส เม ปริวสนฺตสฺส ลชฺชีธมฺโม โอกฺกมิ ‘ยํนูนาหํ สํฆํ ทฺวินฺนํ อาปตฺตีนํ เทฺวมาส\-ปฺปฏิจฺฉนฺนานํ อิตรมฺปิ มาสํ ปริวาสํ ยาเจยฺยนฺ’ติ, โส สํฆํ ทฺวินฺนํ \csromanfolio{147}อาปตฺตีนํ เทฺวมาส\-ปฺปฏิจฺฉนฺนานํ อิตรมฺปิ มาสํ ปริวาสํ ยาจติ, ยทิ สํฆสฺส ปตฺตกลฺลํ, สํโฆ อิตฺถนฺนามสฺส ภิกฺขุโน ทฺวินฺนํ อาปตฺตีนํ เทฺวมาส\-ปฺปฏิจฺฉนฺนานํ อิตรมฺปิ มาสํ ปริวาสํ ทเทยฺย, เอสา ญตฺติ.}

\prose{สุณาตุ เม ภนฺเต สํโฆ, อยํ อิตฺถนฺนาโม ภิกฺขุ เทฺว สํฆาทิเสสา อาปตฺติโย อาปชฺชิ เทฺวมาส\-ปฺปฏิจฺฉนฺนาโย, ตสฺส เอตทโหสิ ‘อหํ โข เทฺว สํฆาทิเสสา อาปตฺติโย อาปชฺชิํ เทฺวมาส\-ปฺปฏิจฺฉนฺนาโย, ยํนูนาหํ สํฆํ ทฺวินฺนํ อาปตฺตีนํ เทฺวมาส\-ปฺปฏิจฺฉนฺนานํ เอกมาส\-ปริวาสํ ยาเจยฺยนฺ’ติ, โส สํฆํ ทฺวินฺนํ อาปตฺตีนํ เทฺวมาส\-ปฺปฏิจฺฉนฺนานํ เอกมาส\-ปริวาสํ ยาจิ, ตสฺส สํโฆ ทฺวินฺนํ อาปตฺตีนํ เทฺวมาส\-ปฺปฏิจฺฉนฺนานํ เอกมาส\-ปริวาสํ อทาสิ, ตสฺส ปริวสนฺตสฺส ลชฺชีธมฺโม โอกฺกมิ\csromandash{}อหํ โข เทฺว สํฆาทิเสสา อาปตฺติโย อาปชฺชิํ เทฺวมาส\-ปฺปฏิจฺฉนฺนาโย, กสฺส เม เอตทโหสิ ‘อหํ โข เทฺว สํฆาทิเสสา อาปตฺติโย อาปชฺชิํ เทฺวมาส\-ปฺปฏิจฺฉนฺนาโย, ยํนูนาหํ สํฆํ ทฺวินฺนํ อาปตฺตีนํ เทฺวมาส\-ปฺปฏิจฺฉนฺนานํ เอกมาส\-ปริวาสํ ยาเจยฺยนฺ’ติ, โสหํ สํฆํ ทฺวินฺนํ อาปตฺตีนํ เทฺวมาส\-ปฺปฏิจฺฉนฺนานํ เอกมาส\-ปริวาสํ ยาจิํ, ตสฺส เม สํโฆ ทฺวินฺนํ อาปตฺตีนํ เทฺวมาส\-ปฺปฏิจฺฉนฺนานํ เอกมาส\-ปริวาสํ อทาสิ, ตสฺส เม ปริวสนฺตสฺส ลชฺชีธมฺโม โอกฺกมิ ‘ยํนูนาหํ สํฆํ ทฺวินฺนํ อาปตฺตีนํ เทฺวมาส\-ปฺปฏิจฺฉนฺนานํ อิตรมฺปิ มาสํ ปริวาสํ ยาเจยฺยนฺ’ติ, โส สํฆํ ทฺวินฺนํ อาปตฺตีนํ เทฺวมาส\-ปฺปฏิจฺฉนฺนานํ อิตรมฺปิ มาสํ ปริวาสํ ยาจติ, สํโฆ อิตฺถนฺนามสฺส ภิกฺขุโน ทฺวินฺนํ อาปตฺตีนํ เทฺวมาส\-ปฺปฏิจฺฉนฺนานํ อิตรมฺปิ มาสํ ปริวาสํ เทติ, ยสฺสายสฺมโต ขมติ อิตฺถนฺนามสฺส ภิกฺขุโน ทฺวินฺนํ อาปตฺตีนํ เทฺวมาส\-ปฺปฏิจฺฉนฺนานํ อิตรมฺปิ มาสํ ปริวาสสฺส\csromansharedfootnote{84983fe2345f7399}{อิตรมฺปิ มาสปริวาสสฺส (ก), อิตรสฺสปิ มาสปริวาสสฺส (สฺยา)} ทานํ, โส ตุณฺหสฺส, ยสฺส นกฺขมติ, โส ภาเสยฺย.}

\prose{ทุติยมฺปิ เอตมตฺถํ วทามิ ฯเปฯ. ตติยมฺปิ เอตมตฺถํ วทามิ ฯเปฯ.}

\prose{ทินฺโน สํเฆน อิตฺถนฺนามสฺส ภิกฺขุโน ทฺวินฺนํ อาปตฺตีนํ เทฺวมาส\-ปฺปฏิจฺฉนฺนานํ อิตรมฺปิ มาสํ ปริวาโส\csromansharedfootnote{704d9091bbb34c3b}{อิตรมฺปิ มาสปริวาโส (ก), อิตโรปิ มาสปริวาโส (สฺยา)}, ขมติ สํฆสฺส, ตสฺมา ตุณฺหี, เอวเมตํ ธารยามี”ติ.}

\prose{เตน ภิกฺขเว ภิกฺขุนา ปุริมํ อุปาทาย เทฺว มาสา ปริวสิตพฺพา.}

\proseitem{149}{\csromanfolio{148}อิธ ปน ภิกฺขเว ภิกฺขุ เทฺว สํฆาทิเสสา อาปตฺติโย อาปชฺชติ เทฺวมาส\-ปฺปฏิจฺฉนฺนาโย. ตสฺส เอวํ โหติ “อหํ โข เทฺว สํฆาทิเสสา อาปตฺติโย อาปชฺชิํ เทฺวมาส\-ปฺปฏิจฺฉนฺนาโย, ยํนูนาหํ สํฆํ ทฺวินฺนํ อาปตฺตีนํ เทฺวมาส\-ปฺปฏิจฺฉนฺนานํ เอกมาส\-ปริวาสํ ยาเจยฺยนฺ”ติ, โส สํฆํ ทฺวินฺนํ อาปตฺตีนํ เทฺวมาส\-ปฺปฏิจฺฉนฺนานํ เอกมาส\-ปริวาสํ ยาจติ, ตสฺส สํโฆ ทฺวินฺนํ อาปตฺตีนํ เทฺวมาส\-ปฺปฏิจฺฉนฺนานํ เอกมาส\-ปริวาสํ เทติ, ตสฺส ปริวสนฺตสฺส ลชฺชีธมฺโม โอกฺกมิ\csromandash{}อหํ โข เทฺว สํฆาทิเสสา อาปตฺติโย อาปชฺชิํ เทฺวมาส\-ปฺปฏิจฺฉนฺนาโย, ตสฺส เม เอตทโหสิ “อหํ โข เทฺว สํฆาทิเสสา อาปตฺติโย อาปชฺชิํ เทฺวมาส\-ปฺปฏิจฺฉนฺนาโย, ยํนูนาหํ สํฆํ ทฺวินฺนํ อาปตฺตีนํ เทฺวมาส\-ปฺปฏิจฺฉนฺนานํ เอกมาส\-ปริวาสํ ยาเจยฺยนฺ”ติ, โสหํ สํฆํ ทฺวินฺนํ อาปตฺตีนํ เทฺวมาส\-ปฺปฏิจฺฉนฺนานํ เอกมาส\-ปริวาสํ ยาจิํ, ตสฺส เม สํโฆ ทฺวินฺนํ อาปตฺตีนํ เทฺวมาส\-ปฺปฏิจฺฉนฺนานํ เอกมาส\-ปริวาสํ อทาสิ, ตสฺส เม ปริวสนฺตสฺส ลชฺชีธมฺโม โอกฺกมิ “ยํนูนาหํ สํฆํ ทฺวินฺนํ อาปตฺตีนํ เทฺวมาส\-ปฺปฏิจฺฉนฺนานํ อิตรมฺปิ มาสํ ปริวาสํ ยาเจยฺยนฺ”ติ, โส สํฆํ ทฺวินฺนํ อาปตฺตีนํ เทฺวมาส\-ปฺปฏิจฺฉนฺนานํ อิตรมฺปิ มาสํ ปริวาสํ ยาจติ, ตสฺส สํโฆ ทฺวินฺนํ อาปตฺตีนํ เทฺวมาส\-ปฺปฏิจฺฉนฺนานํ อิตรมฺปิ มาสํ ปริวาสํ เทติ. เตน ภิกฺขเว ภิกฺขุนา ปุริมํ อุปาทาย เทฺว มาสา ปริวสิตพฺพา.}

\proseitem{150}{อิธ ปน ภิกฺขเว ภิกฺขุ เทฺว สํฆาทิเสสา อาปตฺติโย อาปชฺชติ เทฺวมาส\-ปฺปฏิจฺฉนฺนาโย, เอกํ มาสํ ชานาติ, เอกํ มาสํ น ชานติ. โส สํฆํ ทฺวินฺนํ อาปตฺตีนํ เทฺวมาส\-ปฺปฏิจฺฉนฺนานํ ยํ มาสํ ชานาติ, ตํ มาสํ ปริวาสํ ยาจติ, ตสฺส สํโฆ ทฺวินฺนํ อาปตฺตีนํ เทฺวมาส\-ปฺปฏิจฺฉนฺนานํ ยํ มาสํ ชานาติ, ตํ มาสํ ปริวาสํ เทติ, โส ปริวสนฺโต อิตรมฺปิ มาสํ ชานาติ, ตสฺส เอวํ โหติ ‘อหํ โข เทฺว สํฆาทิเสสา อาปตฺติโย อาปชฺชิํ เทฺวมาส\-ปฺปฏิจฺฉนฺนาโย, เอกํ มาสํ ชานิํ, เอกํ มาสํ น ชานิํ, โสหํ สํฆํ ทฺวินฺนํ อาปตฺตีนํ เทฺวมาส\-ปฺปฏิจฺฉนฺนานํ ยํ มาสํ ชานิํ, ตํ มาสํ ปริวาสํ ยาจิํ, ตสฺส เม สํโฆ ทฺวินฺนํ อาปตฺตีนํ เทฺวมาส\-ปฺปฏิจฺฉนฺนานํ ยํ มาสํ ชานิํ, ตํ มาสํ ปริวาสํ อทาสิ, โสหํ ปริวสนฺโต อิตรมฺปิ มาสํ ชานามิ, ยํนูนาหํ สํฆํ ทฺวินฺนํ อาปตฺตีนํ เทฺวมาส\-ปฺปฏิจฺฉนฺนานํ อิตรมฺปิ มาสํ ปริวาสํ ยาเจยฺยนฺ’ติ, โส สํฆํ ทฺวินฺนํ อาปตฺตีนํ เทฺวมาส\-ปฺปฏิจฺฉนฺนานํ อิตรมฺปิ มาสํ ปริวาสํ \csromanfolio{149}ยาจติ, ตสฺส สํโฆ ทฺวินฺนํ อาปตฺตีนํ เทฺวมาส\-ปฺปฏิจฺฉนฺนานํ อิตรมฺปิ มาสํ ปริวาสํ เทติ. เตน ภิกฺขเว ภิกฺขุนา ปุริมํ อุปาทาย เทฺว มาสา ปริวสิตพฺพา.}

\proseitem{151}{อิธ ปน ภิกฺขเว ภิกฺขุ เทฺว สํฆาทิเสสา อาปตฺติโย อาปชฺชติ เทฺวมาส\-ปฺปฏิจฺฉนฺนาโย, เอกํ มาสํ สรติ, เอกํ มาสํ นสฺสรติ. โส สํฆํ ทฺวินฺนํ อาปตฺตีนํ เทฺวมาส\-ปฺปฏิจฺฉนฺนานํ ยํ มาสํ สรติ, ตํ มาสํ ปริวาสํ ยาจติ, ตสฺส สํโฆ ทฺวินฺนํ อาปตฺตีนํ เทฺวมาส\-ปฺปฏิจฺฉนฺนานํ ยํ มาสํ สรติ, ตํ มาสํ ปริวาสํ เทติ, โส ปริวสนฺโต อิตรมฺปิ มาสํ สรติ, ตสฺส เอวํ โหติ, “อหํ โข เทฺว สํฆาทิเสสา อาปตฺติโย อาปชฺชิํ เทฺวมาส\-ปฺปฏิจฺฉนฺนาโย, เอกํ มาสํ สริํ, เอกํ มาสํ นสฺสริํ, โสหํ สํฆํ ทฺวินฺนํ อาปตฺตีนํ เทฺวมาส\-ปฺปฏิจฺฉนฺนานํ ยํ มาสํ สริํ, ตํ มาสํ ปริวาสํ ยาจิํ, ตสฺส เม สํโฆ ทฺวินฺนํ อาปตฺตีนํ เทฺวมาส\-ปฺปฏิจฺฉนฺนานํ ยํ มาสํ สริํ, ตํ มาสํ ปริวาสํ อทาสิ, โสหํ ปริวสนฺโต อิตรมฺปิ มาสํ สรามิ, ยํนูนาหํ สํฆํ ทฺวินฺนํ อาปตฺตีนํ เทฺวมาส\-ปฺปฏิจฺฉนฺนานํ อิตรมฺปิ มาสํ ปริวาสํ ยาเจยฺยนฺ”ติ, โส สํฆํ ทฺวินฺนํ อาปตฺตีนํ เทฺวมาส\-ปฺปฏิจฺฉนฺนานํ อิตรมฺปิ มาสํ ปริวาสํ ยาจติ, ตสฺส สํโฆ ทฺวินฺนํ อาปตฺตีนํ เทฺวมาส\-ปฺปฏิจฺฉนฺนานํ อิตรมฺปิ มาสํ ปริวาสํ เทติ. เตน ภิกฺขเว ภิกฺขุนา ปุริมํ อุปาทาย เทฺว มาสา ปริวสิตพฺพา.}

\proseitem{152}{อิธ ปน ภิกฺขเว ภิกฺขุ เทฺว สํฆาทิเสสา อาปตฺติโย อาปชฺชติ เทฺวมาส\-ปฺปฏิจฺฉนฺนาโย, เอกํ มาสํ นิพฺเพมติโก, เอกํ มาสํ เวมติโก. โส สํฆํ ทฺวินฺนํ อาปตฺตีนํ เทฺวมาส\-ปฺปฏิจฺฉนฺนานํ ยํ มาสํ นิพฺเพมติโก, ตํ มาสํ ปริวาสํ ยาจติ, ตสฺส สํโฆ ทฺวินฺนํ อาปตฺตีนํ เทฺวมาส\-ปฺปฏิจฺฉนฺนานํ ยํ มาสํ นิพฺเพมติโก, ตํ มาสํ ปริวาสํ เทติ, โส ปริวสนฺโต อิตรมฺปิ มาสํ นิพฺเพมติโก โหติ, ตสฺส เอวํ โหติ “อหํ โข เทฺว สํฆาทิเสสา อาปตฺติโย อาปชฺชิํ เทฺวมาส\-ปฺปฏิจฺฉนฺนาโย, เอกํ มาสํ นิพฺเพมติโก, เอกํ มาสํ เวมติโก, โสหํ สํฆํ ทฺวินฺนํ อาปตฺตีนํ เทฺวมาส\-ปฺปฏิจฺฉนฺนานํ ยํ มาสํ นิพฺเพมติโก, ตํ มาสํ ปริวาสํ ยาจิํ, ตสฺส เม สํโฆ ทฺวินฺนํ อาปตฺตีนํ เทฺวมาส\-ปฺปฏิจฺฉนฺนานํ ยํ มาสํ นิพฺเพมติโก, ตํ มาสํ ปริวาสํ อทาสิ, โสหํ ปริวสนฺโต อิตรมฺปิ มาสํ นิพฺเพมติโก, ยํนูนาหํ สํฆํ ทฺวินฺนํ อาปตฺตีนํ \csromanfolio{150}เทฺวมาส\-ปฺปฏิจฺฉนฺนานํ อิตรมฺปิ มาสํ ปริวาสํ ยาเจยฺยนฺ”ติ, โส สํฆํ ทฺวินฺนํ อาปตฺตีนํ เทฺวมาส\-ปฺปฏิจฺฉนฺนานํ อิตรมฺปิ มาสํ ปริวาสํ ยาจติ, ตสฺส สํโฆ ทฺวินฺนํ อาปตฺตีนํ เทฺวมาส\-ปฺปฏิจฺฉนฺนานํ อิตรมฺปิ มาสํ ปริวาสํ เทติ. เตน ภิกฺขเว ภิกฺขุนา ปุริมํ อุปาทาย เทฺว มาสา ปริวสิตพฺพา.}

\proseitem{153}{อิธ ปน ภิกฺขเว ภิกฺขุ เทฺว สํฆาทิเสสา อาปตฺติโย อาปชฺชติ เทฺวมาส\-ปฺปฏิจฺฉนฺนาโย, เอโก มาโส ชาน\-ปฺปฏิจฺฉนฺโน, เอโก มาโส อชาน\-ปฺปฏิจฺฉนฺโน. โส สํฆํ ทฺวินฺนํ อาปตฺตีนํ เทฺวมาส\-ปฺปฏิจฺฉนฺนานํ เทฺวมาส\-ปริวาสํ ยาจติ, ตสฺส สํโฆ ทฺวินฺนํ อาปตฺตีนํ เทฺวมาส\-ปฺปฏิจฺฉนฺนานํ เทฺวมาส\-ปริวาสํ เทติ, ตสฺส ปริวสนฺตสฺส อญฺโญ ภิกฺขุ อาคจฺฉติ พหุสฺสุโต อาคตาคโม ธมฺมธโร วินยธโร มาติกาธโร ปณฺฑิโต วิยตฺโต เมธาวี ลชฺชี กุกฺกุจฺจโก สิกฺขากาโม, โส เอวํ วเทติ “กิํ อยํ อาวุโส ภิกฺขุ อาปนฺโน, กิสฺสายํ ภิกฺขุ ปริวสตี”ติ. เต เอวํ วเทนฺติ “อยํ อาวุโส ภิกฺขุ เทฺว สํฆาทิเสสา อาปตฺติโย อาปชฺชิ เทฺวมาส\-ปฺปฏิจฺฉนฺนาโย, เอโก มาโส ชาน\-ปฺปฏิจฺฉนฺโน, เอโก มาโส อชาน\-ปฺปฏิจฺฉนฺโน. โส สํฆํ ทฺวินฺนํ อาปตฺตีนํ เทฺวมาส\-ปฺปฏิจฺฉนฺนานํ เทฺวมาส\-ปริวาสํ ยาจิ, ตสฺส สํโฆ ทฺวินฺนํ อาปตฺตีนํ เทฺวมาส\-ปฺปฏิจฺฉนฺนานํ เทฺวมาส\-ปริวาสํ อทาสิ, ตาโย อยํ อาวุโส ภิกฺขุ อาปนฺโน, ตาสายํ ภิกฺขุ ปริวสตี”ติ. โส เอวํ วเทติ “ยฺวายํ อาวุโส มาโส ชาน\-ปฺปฏิจฺฉนฺโน, ธมฺมิกํ ตสฺส มาสสฺส ปริวาสทานํ, ธมฺมตฺตา รุหติ. โย จ ขฺวายํ อาวุโส มาโส อชาน\-ปฺปฏิจฺฉนฺโน, อธมฺมิกํ ตสฺส มาสสฺส ปริวาสทานํ, อธมฺมตฺตา น รุหติ, เอกสฺส อาวุโส มาสสฺส ภิกฺขุ มานตฺตารโห”ติ.}

\proseitem{154}{อิธ ปน ภิกฺขเว ภิกฺขุ เทฺว สํฆาทิเสสา อาปตฺติโย อาปชฺชติ เทฺวมาส\-ปฺปฏิจฺฉนฺนาโย, เอโก มาโส สรมาน\-ปฺปฏิจฺฉนฺโน, เอโก มาโส อสฺสรมาน\-ปฺปฏิจฺฉนฺโน. โส สํฆํ ทฺวินฺนํ อาปตฺตีนํ เทฺวมาส\-ปฺปฏิจฺฉนฺนานํ เทฺวมาส\-ปริวาสํ ยาจติ, ตสฺส สํโฆ ทฺวินฺนํ อาปตฺตีนํ เทฺวมาส\-ปฺปฏิจฺฉนฺนานํ เทฺวมาส\-ปริวาสํ เทติ, ตสฺส ปริวสนฺตสฺส อญฺโญ ภิกฺขุ อาคจฺฉติ พหุสฺสุโต อาคตาคโม ธมฺมธโร วินยธโร มาติกาธโร ปณฺฑิโต วิยตฺโต เมธาวี ลชฺชี กุกฺกุจฺจโก \csromanfolio{151}สิกฺขากาโม, โส เอวํ วเทติ “กิํ อยํ อาวุโส ภิกฺขุ อาปนฺโน, กิสฺสายํ ภิกฺขุ ปริวสตี”ติ. เต เอวํ วเทนฺติ “อยํ อาวุโส ภิกฺขุ เทฺว สํฆาทิเสสา อาปตฺติโย อาปชฺชิ เทฺวมาส\-ปฺปฏิจฺฉนฺนาโย, เอโก มาโส สรมาน\-ปฺปฏิจฺฉนฺโน, เอโก มาโส อสฺสรมาน\-ปฺปฏิจฺฉนฺโน, โส สํฆํ ทฺวินฺนํ อาปตฺตีนํ เทฺวมาส\-ปฺปฏิจฺฉนฺนานํ เทฺวมาสปริวาทํ ยาจิ, ตสฺส สํโฆ ทฺวินฺนํ อาปตฺตีนํ เทฺวมาส\-ปฺปฏิจฺฉนฺนานํ เทฺวมาส\-ปริวาสํ อทาสิ, ตาโย อยํ อาวุโส ภิกฺขุ อาปนฺโน, ตาสายํ ภิกฺขุ ปริวสตี”ติ. โส เอวํ วเทติ “ยฺวายํ อาวุโส มาโส สรมาน\-ปฺปฏิจฺฉนฺโน, ธมฺมิกํ ตสฺส มาสสฺส ปริวาสทานํ, ธมฺมตฺตา รุหติ. โย จ ขฺวายํ อาวุโส มาโส อสฺสรมาน\-ปฺปฏิจฺฉนฺโน, อธมฺมิกํ ตสฺส มาสสฺส ปริวาสทานํ, อธมฺมตฺตา น รุหติ, เอกสฺส อาวุโส มาสสฺส ภิกฺขุ มานตฺตารโห”ติ.}

\proseitem{155}{อิธ ปน ภิกฺขเว ภิกฺขุ เทฺว สํฆาทิเสสา อาปตฺติโย อาปชฺชติ เทฺวมาส\-ปฺปฏิจฺฉนฺนาโย, เอโก มาโส นิพฺเพมติก\-ปฺปฏิจฺฉนฺโน, เอโก มาโส เวมติก\-ปฺปฏิจฺฉนฺโน. โส สํฆํ ทฺวินฺนํ อาปตฺตีนํ เทฺวมาส\-ปฺปฏิจฺฉนฺนานํ เทฺวมาส\-ปริวาสํ ยาจติ, ตสฺส สํโฆ ทฺวินฺนํ อาปตฺตีนํ เทฺวมาส\-ปฺปฏิจฺฉนฺนานํ เทฺวมาส\-ปริวาสํ เทติ, ตสฺส ปริวสนฺตสฺส อญฺโญ ภิกฺขุ อาคจฺฉติ พหุสฺสุโต อาคตาคโม ธมฺมธโร วินยธโร มาติกาธโร ปณฺฑิโต วิยตฺโต เมธาวี ลชฺชี กุกฺกุจฺจโก สิกฺขากาโม, โส เอวํ วเทติ “กิํ อยํ อาวุโส ภิกฺขุ อาปนฺโน, กิสฺสายํ ภิกฺขุ ปริวสตี”ติ. เต เอวํ วเทนฺติ “อยํ อาวุโส ภิกฺขุ เทฺว สํฆาทิเสสา อาปตฺติโย อาปชฺชิ เทฺวมาส\-ปฺปฏิจฺฉนฺนาโย, เอโก มาโส นิพฺเพมติก\-ปฺปฏิจฺฉนฺโน, เอโก มาโส เวมติก\-ปฺปฏิจฺฉนฺโน, โส สํฆํ ทฺวินฺนํ อาปตฺตีนํ เทฺวมาส\-ปฺปฏิจฺฉนฺนานํ เทฺวมาส\-ปริวาสํ ยาจิ, ตสฺส สํโฆ ทฺวินฺนํ อาปตฺตีนํ เทฺวมาส\-ปฺปฏิจฺฉนฺนานํ เทฺวมาส\-ปริวาสํ อทาสิ, ตาโย อยํ อาวุโส ภิกฺขุ อาปนฺโน, ตาสายํ ภิกฺขุ ปริวสตี”ติ. โส เอวํ วเทติ “ยฺวายํ อาวุโส มาโส นิพฺเพมติก\-ปฺปฏิจฺฉนฺโน, ธมฺมิกํ ตสฺส มาสสฺส ปริวาสทานํ, ธมฺมตฺตา รุหติ. โย จ ขฺวายํ อาวุโส มาโส เวมติก\-ปฺปฏิจฺฉนฺโน, อธมฺมิกํ ตสฺส มาสสฺส ปริวาสทานํ, อธมฺมตฺตา น รุหติ, เอกสฺส อาวุโส มาสสฺส ภิกฺขุ มานตฺตารโห”ติ.}

\csromanmatikaanchor{mtk.86}
\csromantocmarkref{subsection}{สุทฺธนฺตปริวาส}{mtk.86}
\bookmark[dest={mtk.86},level=2]{สุทฺธนฺตปริวาส}
\csromanheader{2}{สุทฺธนฺต\-ปริวาส}
\proseitem{156}{\csromanfolio{152}เตน โข ปน สมเยน อญฺญตโร ภิกฺขุ สมฺพหุลา สํฆาทิเสสา อาปตฺติโย อาปนฺโน โหติ, โส อาปตฺติ\-ปริยนฺตํ น ชานาติ, รตฺติ\-ปริยนฺตํ น ชานาติ, อาปตฺติ\-ปริยนฺตํ นสฺสรติ, รตฺติ\-ปริยนฺตํ นสฺสรติ, อาปตฺติ\-ปริยนฺเต เวมติโก, รตฺติปริยนฺเต เวมติโก, โส ภิกฺขูนํ อาโรเจสิ “อหํ อาวุโส สมฺพหุลา สํฆาทิเสสา อาปตฺติโย อาปชฺชิํ, อาปตฺติ\-ปริยนฺตํ น ชานามิ, รตฺติ\-ปริยนฺตํ น ชานามิ, อาปตฺติ\-ปริยนฺตํ นสฺสรามิ, รตฺติ\-ปริยนฺตํ นสฺสรามิ, อาปตฺติ\-ปริยนฺเต เวมติโก, รตฺติปริยนฺเต เวมติโก, กถํ นุ โข มยา ปฏิปชฺชิตพฺพนฺ”ติ. ภควโต เอตมตฺถํ อาโรเจสุํ. เตน หิ ภิกฺขเว สํโฆ ตสฺส ภิกฺขุโน ตาสํ อาปตฺตีนํ สุทฺธนฺต\-ปริวาสํ เทตุ. เอวญฺจ ปน ภิกฺขเว ทาตพฺโพ\csromandash{}}

\prose{เตน ภิกฺขเว ภิกฺขุนา สํฆํ อุปสงฺกมิตฺวา ฯเปฯ เอวมสฺส วจนีโย “อหํ ภนฺเต สมฺพหุลา สํฆาทิเสสา อาปตฺติโย อาปชฺชิํ, อาปตฺติ\-ปริยนฺตํ น ชานามิ, รตฺติ\-ปริยนฺตํ น ชานามิ, อาปตฺติ\-ปริยนฺตํ นสฺสรามิ, รตฺติ\-ปริยนฺตํ นสฺสรามิ, อาปตฺติ\-ปริยนฺเต เวมติโก, รตฺติปริยนฺเต เวมติโก, โสหํ ภนฺเต สํฆํ ตาสํ อาปตฺตีนํ สุทฺธนฺต\-ปริวาสํ ยาจามี”ติ. ทุติยมฺปิ ยาจิตพฺโพ. ตติยมฺปิ ยาจิตพฺโพ. พฺยตฺเตน ภิกฺขุนา ปฏิพเลน สํโฆ ญาเปตพฺโพ\csromandash{}}

\proseitem{157}{“สุณาตุ เม ภนฺเต สํโฆ, อยํ อิตฺถนฺนาโม ภิกฺขุ สมฺพหุลา สํฆาทิเสสา อาปตฺติโย อาปชฺชิ, อาปตฺติ\-ปริยนฺตํ น ชานาติ, รตฺติ\-ปริยนฺตํ น ชานาติ, อาปตฺติ\-ปริยนฺตํ นสฺสรติ, รตฺติ\-ปริยนฺตํ นสฺสรติ, อาปตฺติ\-ปริยนฺเต เวมติโก, รตฺติปริยนฺเต เวมติโก โส สํฆํ ตาสํ อาปตฺตีนํ สุทฺธนฺต\-ปริวาสํ ยาจติ. ยทิ สํฆสฺส ปตฺตกลฺลํ, สํโฆ อิตฺถนฺนามสฺส ภิกฺขุโน ตาสํ อาปตฺตีนํ สุทฺธนฺต\-ปริวาสํ ทเทยฺย, เอสา ญตฺติ.}

\prose{สุณาตุ เม ภนฺเต สํโฆ, อยํ อิตฺถนฺนาโม ภิกฺขุ สมฺพหุลา สํฆาทิเสสา อาปตฺติโย อาปชฺชิ, อาปตฺติ\-ปริยนฺตํ น ชานาติ, รตฺติ\-ปริยนฺตํ น ชานาติ, อาปตฺติ\-ปริยนฺตํ นสฺสรติ, รตฺติ\-ปริยนฺตํ นสฺสรติ, อาปตฺติ\-ปริยนฺเต เวมติโก, รตฺติปริยนฺเต เวมติโก, โส สํฆํ \csromanfolio{153}ตาสํ อาปตฺตีนํ สุทฺธนฺต\-ปริวาสํ ยาจติ. สํโฆ อิตฺถนฺนามสฺส ภิกฺขุโน ตาสํ อาปตฺตีนํ สุทฺธนฺต\-ปริวาสํ เทติ, ยสฺสยสฺมโต ขมติ อิตฺถนฺนามสฺส ภิกฺขุโน ตาสํ อาปตฺตีนํ สุทฺธนฺต\-ปริวาสสฺส ทานํ, โส ตุณฺหสฺส, ยสฺส นกฺขมติ, โส ภาเสยฺย.}

\prose{ทุติยมฺปิ เอตมตฺถํ วทามิ ฯเปฯ. ตติยมฺปิ เอตมตฺถํ วทามิ ฯเปฯ.}

\prose{ทินฺโน สํเฆน อิตฺถนฺนามสฺส ภิกฺขุโน ตาสํ อาปตฺตีนํ สุทฺธนฺต\-ปริวาโส, ขมติ สํฆสฺส, ตสฺมา ตุณฺหี, เอวเมตํ ธารยามี”ติ.}

\proseitem{158}{เอวํ โข ภิกฺขเว สุทฺธนฺต\-ปริวาโส ทาตพฺโพ, เอวํ ปริวาโส ทาตพฺโพ. กถญฺจ ภิกฺขเว สุทฺธนฺต\-ปริวาโส ทาตพฺโพ. อาปตฺติ\-ปริยนฺตํ น ชานาติ, รตฺติ\-ปริยนฺตํ น ชานาติ, อาปตฺติ\-ปริยนฺตํ นสฺสรติ, รตฺติ\-ปริยนฺตํ นสฺสรติ, อาปตฺติ\-ปริยนฺเต เวมติโก, รตฺติปริยนฺเต เวมติโก. สุทฺธนฺต\-ปริวาโส ทาตพฺโพ.}

\prose{อาปตฺติ\-ปริยนฺตํ ชานาติ, รตฺติ\-ปริยนฺตํ น ชานาติ, อาปตฺติ\-ปริยนฺตํ สรติ, รตฺติ\-ปริยนฺตํ นสฺสรติ, อาปตฺติ\-ปริยนฺเต นิพฺเพมติโก, รตฺติปริยนฺเต เวมาติโก. สุทฺธนฺต\-ปริวาโส ทาตพฺโพ.}

\prose{อาปตฺติ\-ปริยนฺตํ เอกจฺจํ ชานาติ, เอกจฺจํ น ชานาติ, รตฺติ\-ปริยนฺตํ น ชานาติ, อาปตฺติ\-ปริยนฺตํ เอกจฺจํ สรติ, เอกจฺจํ นสฺสรติ, รตฺติ\-ปริยนฺตํ นสฺสรติ, อาปตฺติ\-ปริยนฺเต เอกจฺเจ เวมติโก, เอกจฺเจ นิพฺเพมติโก, รตฺติปริยนฺเต เวมติโก. สุทฺธนฺต\-ปริวาโส ทาตพฺโพ.}

\prose{อาปตฺติ\-ปริยนฺตํ น ชานาติ, รตฺติ\-ปริยนฺตํ เอกจฺจํ ชานาติ, เอกจฺจํ น ชานาติ, อาปตฺติ\-ปริยนฺตํ นสฺสรติ, รตฺติ\-ปริยนฺตํ เอกจฺจํ สรติ, เอกจฺจํ นสฺสรติ, อาปตฺติ\-ปริยนฺเต เวมติโก, รตฺติปริยนฺเต เอกจฺเจ เวมติโก, เอกจฺเจ นิพฺเพมติโก. สุทฺธนฺต\-ปริวาโส ทาตพฺโพ.}

\prose{อาปตฺติ\-ปริยนฺตํ ชานาติ, รตฺติ\-ปริยนฺตํ เอกจฺจํ ชานาติ, เอกจฺจํ น ชานาติ, อาปตฺติ\-ปริยนฺตํ สรติ, รตฺติ\-ปริยนฺตํ เอกจฺจํ สรติ, เอกจฺจํ นสฺสรติ, อาปตฺติ\-ปริยนฺเต นิพฺเพมติโก, รตฺติปริยนฺเต เอกจฺเจ เวมติโก, เอกจฺเจ นิพฺเพมติโก. สุทฺธนฺต\-ปริวาโส ทาตพฺโพ.}

\prose{อาปตฺติ\-ปริยนฺตํ เอกจฺจํ ชานาติ, เอกจฺจํ น ชานาติ, รตฺติ\-ปริยนฺตํ เอกจฺจํ ชานาติ, เอกจฺจํ น ชานาติ, อาปตฺติ\-ปริยนฺตํ เอกจฺจํ สรติ, เอกจฺจํ \csromanfolio{154}นสฺสรติ, รตฺติ\-ปริยนฺตํ เอกจฺจํ สรติ, เอกจฺจํ นสฺสรติ, อาปตฺติ\-ปริยนฺเต เอกจฺเจ เวมติโก, เอกจฺเจ นิพฺเพมติโก, รตฺติปริยนฺเต เอกจฺเจ เวมติโก, เอกจฺเจ นิพฺเพมติโก. สุทฺธนฺต\-ปริวาโส ทาตพฺโพ. เอวํ โข ภิกฺขเว สุทฺธนฺต\-ปริวาโส ทาตพฺโพ.}

\proseitem{159}{กถญฺจ ภิกฺขเว ปริวาโส ทาตพฺโพ, อาปตฺติ\-ปริยนฺตํ ชานาติ, รตฺติ\-ปริยนฺตํ ชานาติ, อาปตฺติ\-ปริยนฺตํ สรติ, รตฺติ\-ปริยนฺตํ สรติ, อาปตฺติ\-ปริยนฺเต นิพฺเพมติโก, รตฺติปริยนฺเต นิพฺเพมติโก. ปริวาโส ทาตพฺโพ.}

\prose{อาปตฺติ\-ปริยนฺตํ น ชานาติ, รตฺติ\-ปริยนฺตํ ชานาติ, อาปตฺติ\-ปริยนฺตํ นสฺสรติ, รตฺติ\-ปริยนฺตํ สรติ, อาปตฺติ\-ปริยนฺเต เวมติโก, รตฺติปริยนฺเต นิพฺเพมติโก. ปริวาโส ทาตพฺโพ.}

\prosewithcloser{อาปตฺติ\-ปริยนฺตํ เอกจฺจํ ชานาติ, เอกจฺจํ น ชานาติ, รตฺติ\-ปริยนฺตํ ชานาติ, อาปตฺติ\-ปริยนฺตํ เอกจฺจํ สรติ, เอกจฺจํ นสฺสรติ, รตฺติ\-ปริยนฺตํ สรติ, อาปตฺติ\-ปริยนฺเต เอกจฺเจ เวมติโก, เอกจฺเจ นิพฺเพมติโก, รตฺติปริยนฺเต นิพฺเพมติโก. ปริวาโส ทาตพฺโพ. เอวํ โข ภิกฺขเว ปริวาโส ทาตพฺโพ.}{ปริวาโส นิฏฺฐิโต.}
\csromanmatikaanchor{mtk.87}
\csromantocmarkref{section}{3. จตฺตาลีสก}{mtk.87}
\bookmark[dest={mtk.87},level=1]{3. จตฺตาลีสก}
\csromanheader{1}{3. \textbf{จตฺตาลีสก}}
\proseitem{160}{เตน โข ปน สมเยน อญฺญตโร ภิกฺขุ ปริวสนฺโต วิพฺภมิ. โส ปุน ปจฺจาคนฺตฺวา ภิกฺขู อุปสมฺปทํ ยาจิ. ภควโต เอตมตฺถํ อาโรเจสุํ.}

\prose{อิธ ปน ภิกฺขเว ภิกฺขุ ปริวสนฺโต วิพฺภมติ, วิพฺภนฺตกสฺส ภิกฺขเว ปริวาโส น รุหติ. โส เจ ปุน อุปสมฺปชฺชติ, ตสฺส ตเทว ปุริมํ ปริวาสทานํ, โย ปริวาโส ทินฺโน สุทินฺโน, โย ปริวุตฺโถ สุปริวุตฺโถ, อวเสโส ปริวสิตพฺโพ. (1)}

\prose{อิธ ปน ภิกฺขเว ภิกฺขุ ปริวสนฺโต สามเณโร โหติ, สามเณรสฺส ภิกฺขเว ปริวาโส น รุหติ. โส เจ ปุน อุปสมฺปชฺชติ, \csromanfolio{155}ตสฺส ตเทว ปุริมํ ปริวาสทานํ, โย ปริวาโส ทินฺโน สุทินฺโน, โย ปริวุตฺโถ สุปริวุตฺโถ, อวเสโส ปริวสิตพฺโพ. (2)}

\prose{อิธ ปน ภิกฺขเว ภิกฺขุ ปริวสนฺโต อุมฺมตฺตโก โหติ, อุมฺมตฺตกสฺส ภิกฺขเว ปริวาโส น รุหติ. โส เจ ปุน อนุมฺมตฺตโก โหติ, ตสฺส ตเทว ปุริมํ ปริวาสทานํ, โย ปริวาโส ทินฺโน สุทินฺโน, โย ปริวุตฺโถ สุปริวุตฺโถ, อวเสโส ปริวสิตพฺโพ. (3)}

\prose{อิธ ปน ภิกฺขเว ภิกฺขุ ปริวสนฺโต ขิตฺตจิตฺโต โหติ, ขิตฺตจิตฺตสฺส ภิกฺขเว ปริวาโส น รุหติ. โส เจ ปุน อขิตฺตจิตฺโต โหติ, ตสฺส ตเทว ปุริมํ ปริวาสทานํ, โย ปริวาโส ทินฺโน สุทินฺโน, โย ปริวุตฺโถ สุปริวุตฺโถ, อวเสโส ปริวสิตพฺโพ. (4)}

\prose{อิธ ปน ภิกฺขเว ภิกฺขุ ปริวสนฺโต เวทนาฏฺโฏ โหติ, เวทนาฏฺฏสฺส ภิกฺขเว ปริวาโส น รุหติ. โส เจ ปุน อเวทนาฏฺโฏ โหติ, ตสฺส ตเทว ปุริมํ ปริวาสทานํ, โย ปริวาโส ทินฺโน สุทินฺโน, โย ปริวุตฺโถ สุปริวุตฺโถ, อวเสโส ปริวสิตพฺโพ. (5)}

\prose{อิธ ปน ภิกฺขเว ภิกฺขุ ปริวสนฺโต อาปตฺติยา อทสฺสเน อุกฺขิปิยฺยติ\csromansharedfootnote{7b11a06397054e38}{อุกฺขิปิยติ (สฺยา), อุกฺขิปียติ (ก)}, อุกฺขิตฺตกสฺส ภิกฺขเว ปริวาโส น รุหติ. โส เจ ปุน โอสาริยฺยติ\csromansharedfootnote{29ac1a602f55b6cb}{โอสาริยติ (สฺยา), โอสารียติ (ก)}, ตสฺส ตเทว ปุริมํ ปริวาสทานํ, โย ปริวาโส ทินฺโน สุทินฺโน, โย ปริวุตฺโถ สุปริวุตฺโถ, อวเสโส ปริวสิตพฺโพ. (6)}

\prose{อิธ ปน ภิกฺขเว ภิกฺขุ ปริวสนฺโต อาปตฺติยา อปฺปฏิกมฺเม อุกฺขิปิยฺยติ, อุกฺขิตฺตกสฺส ภิกฺขเว ปริวาโส น รุหติ. โส เจ ปุน โอสาริยฺยติ, ตสฺส ตเทว ปุริมํ ปริวาสทานํ, โย ปริวาโส ทินฺโน สุทินฺโน, โย ปริวุตฺโถ สุปริวุตฺโถ, อวเสโส ปริวสิตพฺโพ. (7)}

\prose{อิธ ปน ภิกฺขเว ภิกฺขุ ปริวสนฺโต ปาปิกาย ทิฏฺฐิยา อปฺปฏินิสฺสคฺเค อุกฺขิปิยฺยติ, อุกฺขิตฺตกสฺส ภิกฺขเว ปริวาโส น รุหติ. โส เจ ปุน โอสาริยฺยติ, ตสฺส ตเทว ปุริมํ ปริวาสทานํ, โย ปริวาโส ทินฺโน สุทินฺโน, โย ปริวุตฺโถ สุปริวุตฺโถ, อวเสโส ปริวสิตพฺโพ. (8)}

\proseitem{161}{\csromanfolio{156}อิธ ปน ภิกฺขเว ภิกฺขุ มูลาย\-ปฏิกสฺสนา\-รโห วิพฺภมติ, วิพฺภนฺตกสฺส ภิกฺขเว มูลาย\-ปฏิกสฺสนา น รุหติ. โส เจ ปุน อุปสมฺปชฺชติ, ตสฺส ตเทว ปุริมํ ปริวาสทานํ, โย ปริวาโส ทินฺโน สุทินฺโน, โย ปริวุตฺโถ สุปริวุตฺโถ, โส ภิกฺขุ มูลาย ปฏิกสฺสิตพฺโพ. (9)}

\prose{อิธ ปน ภิกฺขเว ภิกฺขุ มูลาย\-ปฏิกสฺสนา\-รโห สามเณโร โหติ ฯเปฯ อุมฺมตฺตโก โหติ ฯเปฯ ขิตฺตจิตฺโต โหติ ฯเปฯ เวทนาฏฺโฏ โหติ ฯเปฯ อาปตฺติยา อทสฺสเน อุกฺขิปิยฺยติ ฯเปฯ อาปตฺติยา อปฺปฏิกมฺเม อุกฺขิปิยฺยติ ฯเปฯ ปาปิกาย ทิฏฺฐิยา อปฺปฏินิสฺสคฺเค อุกฺขิปิยฺยติ, อุกฺขิตฺตกสฺส ภิกฺขเว มูลาย ปฏิกสฺสนา น รุหติ. โส เจ ปุน โอสาริยฺยติ, ตสฺส ตเทว ปุริมํ ปริวาสทานํ, โย ปริวาโส ทินฺโน สุทินฺโน, โย ปริวุตฺโถ สุปริวุตฺโถ, โส ภิกฺขุ มูลาย ปฏิกสฺสิตพฺโพ. (16)}

\proseitem{162}{อิธ ปน ภิกฺขเว ภิกฺขุ มานตฺตารโห วิพฺภมติ, วิพฺภนฺตกสฺส ภิกฺขเว มานตฺตทานํ น รุหติ. โส เจ ปุน อุปสมฺปชฺชติ, ตสฺส ตเทว ปุริมํ ปริวาสทานํ, โย ปริวาโส ทินฺโน สุทินฺโน, โย ปริวุตฺโถ สุปริวุตฺโถ, ตสฺส ภิกฺขุโน มานตฺตํ ทาตพฺพํ. (17)}

\prose{อิธ ปน ภิกฺขเว ภิกฺขุ มานตฺตารโห สามเณโร โหติ ฯเปฯ อุมฺมตฺตโก โหติ ฯเปฯ ขิตฺตจิตฺโต โหติ ฯเปฯ เวทนาฏฺโฏ โหติ ฯเปฯ อาปตฺติยา อทสฺสเน อุกฺขิปิยฺยติ ฯเปฯ อาปตฺติยา อปฺปฏิกมฺเม อุกฺขิปิยฺยติ ฯเปฯ ปาปิกาย ทิฏฺฐิยา อปฺปฏินิสฺสคฺเค อุกฺขิปิยฺยติ, อุกฺขิตฺตกสฺส ภิกฺขเว มานตฺตทานํ น รุหติ. โส เจ ปุน โอสาริยฺยติ, ตสฺส ตเทว ปุริมํ ปริวาสทานํ, โย ปริวาโส ทินฺโน สุทินฺโน, โย ปริวุตฺโถ สุปริวุตฺโถ, ตสฺส ภิกฺขุโน มานตฺตํ ทาตพฺพํ. (24)}

\proseitem{163}{อิธ ปน ภิกฺขเว ภิกฺขุ มานตฺตํ จรนฺโต วิพฺภมติ, วิพฺภนฺตกสฺส ภิกฺขเว มานตฺตจริยา น รุหติ. โส เจ ปุน อุปสมฺปชฺชติ, ตสฺส ตเทว ปุริมํ ปริวาสทานํ, โย ปริวาโส ทินฺโน สุทินฺโน, โย ปริวุตฺโถ สุปริวุตฺโถ, ยํ มานตฺตํ ทินฺนํ สุทินฺนํ, ยํ มานตฺตํ จิณฺณํ สุจิณฺณํ, อวเสสํ จริตพฺพํ. (25)}

\csromanmatikaanchor{mtk.88}
\csromantocmarkref{section}{4. ฉติํสก}{mtk.88}
\bookmark[dest={mtk.88},level=1]{4. ฉติํสก}
\prose{\csromanfolio{157}อิธ ปน ภิกฺขเว ภิกฺขุ มานตฺตํ จรนฺโต สามเณโร โหติ ฯเปฯ อุมฺมตฺตโก โหติ ฯเปฯ ขิตฺตจิตฺโต โหติ ฯเปฯ เวทนาฏฺโฏ โหติ ฯเปฯ อาปตฺติยา อทสฺสเน อุกฺขิปิยฺยติ ฯเปฯ อาปตฺติยา อปฺปฏิกมฺเม อุกฺขิปิยฺยติ ฯเปฯ ปาปิกาย ทิฏฺฐิยา อปฺปฏินิสฺสคฺเค อุกฺขิปิยฺยติ, อุกฺขิตฺตกสฺส ภิกฺขเว มานตฺตจริยา น รุหติ. โส เจ ปุน โอสาริยฺยติ, ตสฺส ตเทว ปุริมํ ปริวาสทานํ, โย ปริวาโส ทินฺโน สุทินฺโน, โย ปริวุตฺโถ สุปริวุตฺโถ, ยํ มานตฺตํ ทินฺนํ สุทินฺนํ, ยํ มานตฺตํ จิณฺณํ สุจิณฺณํ, อวเสสํ จริตพฺพํ. (32)}

\proseitem{164}{อิธ ปน ภิกฺขเว ภิกฺขุ อพฺภานารโห วิพฺภมติ, วิพฺภนฺตกสฺส ภิกฺขเว อพฺภานํ น รุหติ. โส เจ ปุน อุปสมฺปชฺชติ, ตสฺส ตเทว ปุริมํ ปริวาสทานํ, โย ปริวาโส ทินฺโน สุทินฺโน, โย ปริวุตฺโถ สุปริวุตฺโถ, ยํ มานตฺตํ ทินฺนํ สุทินฺนํ, ยํ มานตฺตํ จิณฺณํ สุจิณฺณํ, โส ภิกฺขุ อพฺเภตพฺโพ. (33)}

\prosewithcloser{อิธ ปน ภิกฺขเว ภิกฺขุ อพฺภานารโห สามเณโร โหติ ฯเปฯ อุมฺมตฺตโก โหติ ฯเปฯ ขิตฺตจิตฺโต โหติ ฯเปฯ เวทนาฏฺโฏ โหติ ฯเปฯ อาปตฺติยา อทสฺสเน อุกฺขิปิยฺยติ ฯเปฯ อาปตฺติยา อปฺปฏิกมฺเม อุกฺขิปิยฺยติ ฯเปฯ ปาปิกาย ทิฏฺฐิยา อปฺปฏินิสฺสคฺเค อุกฺขิปิยฺยติ, อุกฺขิตฺตกสฺส ภิกฺขเว อพฺภานํ น รุหติ. โส เจ ปุน โอสาริยฺยติ, ตสฺส ตเทว ปุริมํ ปริวาสทานํ, โย ปริวาโส ทินฺโน สุทินฺโน, โย ปริวุตฺโถ สุปริวุตฺโถ, ยํ มานตฺตํ ทินฺนํ สุทินฺนํ, ยํ มานตฺตํ จิณฺณํ สุจิณฺณํ, โส ภิกฺขุ อพฺเภตพฺโพ. (40)}{จตฺตาลีสกํ สมตฺตํ.}
\csromanheader{2}{4. \textbf{ฉตฺติํสก}}
\proseitem{165}{อิธ ปน ภิกฺขเว ภิกฺขุ ปริวสนฺโต อนฺตรา สมฺพหุลา สํฆาทิเสสา อาปตฺติโย อาปชฺชติ ปริมาณา\csromansharedfootnote{75e56c055ae3fa6c}{ปริมาณาโย (สี, สฺยา)} อปฺปฏิจฺฉนฺนาโย, โส ภิกฺขุ มูลาย ปฏิกสฺสิตพฺโพ. (1)}

\prose{\csromanfolio{158}อิธ ปน ภิกฺขเว ภิกฺขุ ปริวสนฺโต อนฺตรา สมฺพหุลา สํฆาทิเสสา อาปตฺติโย อาปชฺชติ ปริมาณา ปฏิจฺฉนฺนาโย, โส ภิกฺขุ มูลาย ปฏิกสฺสิตพฺโพ, ยถาปฏิจฺฉนฺนาน\-ญฺจสฺส อาปตฺตีนํ ปุริมาย อาปตฺติยา สโมธาน\-ปริวาโส ทาตพฺโพ. (2)}

\prose{อิธ ปน ภิกฺขเว ภิกฺขุ ปริวสนฺโต อนฺตรา สมฺพหุลา สํฆาทิเสสา อาปตฺติโย อาปชฺชติ ปริมาณา ปฏิจฺฉนฺนาโยปิ อปฺปฏิจฺฉนฺนาโยปิ, โส ภิกฺขุ มูลาย ปฏิกสฺสิตพฺโพ, ยถาปฏิจฺฉนฺนาน\-ญฺจสฺส อาปตฺตีนํ ปุริมาย อาปตฺติยา สโมธาน\-ปริวาโส ทาตพฺโพ. (3)}

\prose{อิธ ปน ภิกฺขเว ภิกฺขุ ปริวสนฺโต อนฺตรา สมฺพหุลา สํฆาทิเสสา อาปตฺติโย อาปชฺชติ อปริมาณา\csromansharedfootnote{177e4d7108950826}{อปริมาณาโย (สี, สฺยา)} อปฺปฏิจฺฉนฺนาโย ฯเปฯ อปริมาณา ปฏิจฺฉนฺนาโย ฯเปฯ อปริมาณา ปฏิจฺฉนฺนาโยปิ อปฺปฏิจฺฉนฺนาโยปิ ฯเปฯ ปริมาณาโยปิ อปริมาณาโยปิ อปฺปฏิจฺฉนฺนาโย ฯเปฯ ปริมาณาโยปิ อปริมาณาโยปิ ปฏิจฺฉนฺนาโย ฯเปฯ ปริมาณาโยปิ อปริมาณาโยปิ ปฏิจฺฉนฺนาโยปิ อปฺปฏิจฺฉนฺนาโยปิ, โส ภิกฺขุ มูลาย ปฏิกสฺสิตพฺโพ. ยถาปฏิจฺฉนฺนาน\-ญฺจสฺส อาปตฺตีนํ ปุริมาย อาปตฺติยา สโมธาน\-ปริวาโส ทาตพฺโพ. (9)}

\prosewithcloser{อิธ ปน ภิกฺขเว ภิกฺขุ มานตฺตารโห ฯเปฯ มานตฺตํ จรนฺโต ฯเปฯ (ยถา ปริวาสํ, ตถา วิตฺถาเรตพฺพํ.) อพฺภานารโห อนฺตรา สมฺพหุลา สํฆาทิเสสา อาปตฺติโย อาปชฺชติ ปริมาณา อปฺปฏิจฺฉนฺนาโย ฯเปฯ ปริมาณา ปฏิจฺฉนฺนาโย ฯเปฯ ปริมาณา ปฏิจฺฉนฺนาโยปิ อปฺปฏิจฺฉนฺนาโยปิ ฯเปฯ อปริมาณา อปฺปฏิจฺฉนฺนาโย ฯเปฯ อปริมาณา ปฏิจฺฉนฺนาโย ฯเปฯ อปริมาณา ปฏิจฺฉนฺนาโยปิ อปฺปฏิจฺฉนฺนาโยปิ ฯเปฯ ปริมาณาโยปิ อปริมาณาโยปิ อปฺปฏิจฺฉนฺนาโย ฯเปฯ ปริมาณาโยปิ อปริมาณาโยปิ ปฏิจฺฉนฺนาโย ฯเปฯ ปริมาณาโยปิ อปริมาณาโยปิ ปฏิจฺฉนฺนาโยปิ อปฺปฏิจฺฉนฺนาโยปิ, โส ภิกฺขุ มูลาย ปฏิกสฺสิตพฺโพ. ยถาปฏิจฺฉนฺนาน\-ญฺจสฺส อาปตฺตีนํ ปุริมาย อาปตฺติยา สโมธาน\-ปริวาโส ทาตพฺโพ. (36)}{ฉตฺติํสกํ สมตฺตํ.}
\csromanmatikaanchor{mtk.89}
\csromantocmarkref{section}{5. มานตฺตสต}{mtk.89}
\bookmark[dest={mtk.89},level=1]{5. มานตฺตสต}
\csromanheader{1}{5. \textbf{มานตฺตสต}}
\proseitem{166}{\csromanfolio{159}อิธ ปน ภิกฺขเว ภิกฺขุ สมฺพหุลา สํฆาทิเสสา อาปตฺติโย อาปชฺชิตฺวา อปฺปฏิจฺฉาเทตฺวา วิพฺภมติ. โส ปุน\csromansharedfootnote{608be3c6f737057a}{โส เจ ปุน (ก)} อุปสมฺปนฺโน ตา อาปตฺติโย นจฺฉาเทติ, ตสฺส ภิกฺขเว ภิกฺขุโน มานตฺตํ ทาตพฺพํ. (1)}

\prose{อิธ ปน ภิกฺขเว ภิกฺขุ สมฺพหุลา สํฆาทิเสสา อาปตฺติโย อาปชฺชิตฺวา อปฺปฏิจฺฉาเทตฺวา วิพฺภมติ. โส ปุน อุปสมฺปนฺโน ตา อาปตฺติโย ฉาเทติ, ตสฺส ภิกฺขเว ภิกฺขุโน ปจฺฉิมสฺมิํ อาปตฺติกฺขนฺเธ ยถา\-ปฏิจฺฉนฺเน ปริวาสํ ทตฺวา มานตฺตํ ทาตพฺพํ. (2)}

\prose{อิธ ปน ภิกฺขเว ภิกฺขุ สมฺพหุลา สํฆาทิเสสา อาปตฺติโย อาปชฺชิตฺวา ปฏิจฺฉาเทตฺวา วิพฺภมติ. โส ปุน อุปสมฺปนฺโน ตา อาปตฺติโย นจฺฉาเทติ, ตสฺส ภิกฺขเว ภิกฺขุโน ปุริมสฺมิํ อาปตฺติกฺขนฺเธ ยถา\-ปฏิจฺฉนฺเน ปริวาสํ ทตฺวา มานตฺตํ ทาตพฺพํ. (3)}

\prose{อิธ ปน ภิกฺขเว ภิกฺขุ สมฺพหุลา สํฆาทิเสสา อาปตฺติโย อาปชฺชิตฺวา ปฏิจฺฉาเทตฺวา วิพฺภมติ. โส ปุน อุปสมฺปนฺโน ตา อาปตฺติโย ฉาเทติ, ตสฺส ภิกฺขเว ภิกฺขุโน ปุริมสฺมิํ จ ปจฺฉิมสฺมิํ จ อาปตฺติกฺขนฺเธ ยถา\-ปฏิจฺฉนฺเน ปริวาสํ ทตฺวา มานตฺตํ ทาตพฺพํ. (4)}

\proseitem{167}{อิธ ปน ภิกฺขเว ภิกฺขุ สมฺพหุลา สํฆาทิเสสา อาปตฺติโย อาปชฺชติ, ตสฺส โหนฺติ อาปตฺติโย ปฏิจฺฉนฺนาโยปิ อปฺปฏิจฺฉนฺนาโยปิ, โส วิพฺภมิตฺวา ปุน อุปสมฺปนฺโน ยา อาปตฺติโย ปุพฺเพ ฉาเทสิ, ตา อาปตฺติโย ปจฺฉา นจฺฉาเทติ, ยา อาปตฺติโย ปุพฺเพ นจฺฉาเทสิ, ตา อาปตฺติโย ปจฺฉา นจฺฉาเทติ. ตสฺส ภิกฺขเว ภิกฺขุโน ปุริมสฺมิํ อาปตฺติกฺขนฺเธ ยถา\-ปฏิจฺฉนฺเน ปริวาสํ ทตฺวา มานตฺตํ ทาตพฺพํ. (5)}

\prose{อิธ ปน ภิกฺขเว ภิกฺขุ สมฺพหุลา สํฆาทิเสสา อาปตฺติโย อาปชฺชติ, ตสฺส โหนฺติ อาปตฺติโย ปฏิจฺฉนฺนาโยปิ อปฺปฏิจฺฉนฺนาโยปิ. โส วิพฺภมิตฺวา ปุน อุปสมฺปนฺโน ยา อาปตฺติโย ปุพฺเพ ฉาเทสิ, ตา อาปตฺติโย ปจฺฉา นจฺฉาเทติ. ยา อาปตฺติโย ปุพฺเพ นจฺฉาเทสิ, ตา อาปตฺติโย ปจฺฉา ฉาเทติ. ตสฺส ภิกฺขเว ภิกฺขุโน ปุริมสฺมิํ จ ปจฺฉิมสฺมิํ จ อาปตฺติกฺขนฺเธ ยถา\-ปฏิจฺฉนฺเน ปริวาสํ ทตฺวา มานตฺตํ ทาตพฺพํ. (6)}

\prose{\csromanfolio{160}อิธ ปน ภิกฺขเว ภิกฺขุ สมฺพหุลา สํฆาทิเสสา อาปตฺติโย อาปชฺชติ, ตสฺส โหนฺติ อาปตฺติโย ปฏิจฺฉนฺนาโยปิ อปฺปฏิจฺฉนฺนาโยปิ. โส วิพฺภมิตฺวา ปุน อุปสมฺปนฺโน ยา อาปตฺติโย ปุพฺเพ ฉาเทสิ, ตา อาปตฺติโย ปจฺฉา ฉาเทติ. ยา อาปตฺติโย ปุพฺเพ นจฺฉาเทสิ, ตา อาปตฺติโย ปจฺฉา นจฺฉาเทติ. ตสฺส ภิกฺขเว ภิกฺขุโน ปุริมสฺมิํ จ ปจฺฉิมสฺมิํ จ อาปตฺติกฺขนฺเธ ยถา\-ปฏิจฺฉนฺเน ปริวาสํ ทตฺวา มานตฺตํ ทาตพฺพํ. (7)}

\prose{อิธ ปน ภิกฺขเว ภิกฺขุ สมฺพหุลา สํฆาทิเสสา อาปตฺติโย อาปชฺชติ, ตสฺส โหนฺติ อาปตฺติโย ปฏิจฺฉนฺนาโยปิ อปฺปฏิจฺฉนฺนาโยปิ. โส วิพฺภมิตฺวา ปุน อุปสมฺปนฺโน ยา อาปตฺติโย ปุพฺเพ ฉาเทสิ, ตา อาปตฺติโย ปจฺฉา ฉาเทติ. ยา อาปตฺติโย ปุพฺเพ นจฺฉาเทสิ, ตา อาปตฺติโย ปจฺฉา ฉาเทติ. ตสฺส ภิกฺขเว ภิกฺขุโน ปุริมสฺมิํ จ ปจฺฉิมสฺมิํ จ อาปตฺติกฺขนฺเธ ยถา\-ปฏิจฺฉนฺเน ปริวาสํ ทตฺวา มานตฺตํ ทาตพฺพํ. (8)}

\proseitem{168}{อิธ ปน ภิกฺขเว ภิกฺขุ สมฺพหุลา สํฆาทิเสสา อาปตฺติโย อาปชฺชติ, เอกจฺจา อาปตฺติโย ชานาติ, เอกจฺจา อาปตฺติโย น ชานาติ. ยา อาปตฺติโย ชานาติ, ตา อาปตฺติโย ฉาเทติ. ยา อาปตฺติโย น ชานาติ, ตา อาปตฺติโย นจฺฉาเทติ. โส วิพฺภมิตฺวา ปุน อุปสมฺปนฺโน ยา อาปตฺติโย ปุพฺเพ ชานิตฺวา ฉาเทสิ, ตา อาปตฺติโย ปจฺฉา ชานิตฺวา นจฺฉาเทติ. ยา อาปตฺติโย ปุพฺเพ อชานิตฺวา นจฺฉาเทสิ, ตา อาปตฺติโย ปจฺฉา ชานิตฺวา นจฺฉาเทติ. ตสฺส ภิกฺขเว ภิกฺขุโน ปุริมสฺมิํ อาปตฺติกฺขนฺเธ ยถา\-ปฏิจฺฉนฺเน ปริวาสํ ทตฺวา มานตฺตํ ทาตพฺพํ. (9)}

\prose{อิธ ปน ภิกฺขเว ภิกฺขุ สมฺพหุลา สํฆาทิเสสา อาปตฺติโย อาปชฺชติ, เอกจฺจา อาปตฺติโย ชานาติ, เอกจฺจา อาปตฺติโย น ชานาติ. ยา อาปตฺติโย ชานาติ, ตา อาปตฺติโย ฉาเทติ. ยา อาปตฺติโย น ชานาติ, ตา อาปตฺติโย นจฺฉาเทติ. โส วิพฺภมิตฺวา ปุน อุปสมฺปนฺโน ยา อาปตฺติโย ปุพฺเพ ชานิตฺวา ฉาเทสิ, ตา อาปตฺติโย ปจฺฉา ชานิตฺวา นจฺฉาเทติ. ยา อาปตฺติโย ปุพฺเพ อชานิตฺวา นจฺฉาเทสิ, ตา อาปตฺติโย ปจฺฉา ชานิตฺวา ฉาเทติ. \csromanfolio{161}ตสฺส ภิกฺขเว ภิกฺขุโน ปุริมสฺมิํ จ ปจฺฉิมสฺมิํ จ อาปตฺติกฺขนฺเธ ยถา\-ปฏิจฺฉนฺเน ปริวาสํ ทตฺวา มานตฺตํ ทาตพฺพํ. (10)}

\prose{อิธ ปน ภิกฺขเว ภิกฺขุ สมฺพหุลา สํฆาทิเสสา อาปตฺติโย อาปชฺชติ, เอกจฺจา อาปตฺติโย ชานาติ, เอกจฺจา อาปตฺติโย น ชานาติ. ยา อาปตฺติโย ชานาติ, ตา อาปตฺติโย ฉาเทติ. ยา อาปตฺติโย น ชานาติ, ตา อาปตฺติโย นจฺฉาเทติ. โส วิพฺภมิตฺวา ปุน อุปสมฺปนฺโน ยา อาปตฺติโย ปุพฺเพ ชานิตฺวา ฉาเทสิ, ตา อาปตฺติโย ปจฺฉา ชานิตฺวา ฉาเทติ. ยา อาปตฺติโย ปุพฺเพ อชานิตฺวา นจฺฉาเทสิ, ตา อาปตฺติโย ปจฺฉา ชานิตฺวา นจฺฉาเทติ. ตสฺส ภิกฺขเว ภิกฺขุโน ปุริมสฺมิํ จ ปจฺฉิมสฺมิํ จ อาปตฺติกฺขนฺเธ ยถา\-ปฏิจฺฉนฺเน ปริวาสํ ทตฺวา มานตฺตํ ทาตพฺพํ. (11)}

\prose{อิธ ปน ภิกฺขเว ภิกฺขุ สมฺพหุลา สํฆาทิเสสา อาปตฺติโย อาปชฺชติ, เอกจฺจา อาปตฺติโย ชานาติ, เอกจฺจา อาปตฺติโย น ชานาติ. ยา อาปตฺติโย ชานาติ, ตา อาปตฺติโย ฉาเทติ. ยา อาปตฺติโย น ชานาติ, ตา อาปตฺติโย นจฺฉาเทติ. โส วิพฺภมิตฺวา ปุน อุปสมฺปนฺโน ยา อาปตฺติโย ปุพฺเพ ชานิตฺวา ฉาเทสิ, ตา อาปตฺติโย ปจฺฉา ชานิตฺวา ฉาเทติ. ยา อาปตฺติโย ปุพฺเพ อชานิตฺวา นจฺฉาเทสิ, ตา อาปตฺติโย ปจฺฉา ชานิตฺวา ฉาเทติ. ตสฺส ภิกฺขเว ภิกฺขุโน ปุริมสฺมิํ จ ปจฺฉิมสฺมิํ จ อาปตฺติกฺขนฺเธ ยถา\-ปฏิจฺฉนฺเน ปริวาสํ ทตฺวา มานตฺตํ ทาตพฺพํ. (12)}

\proseitem{196}{อิธ ปน ภิกฺขเว ภิกฺขุ สมฺพหุลา สํฆาทิเสสา อาปตฺติโย อาปชฺชติ, เอกจฺจา อาปตฺติโย สรติ, เอกจฺจา อาปตฺติโย นสฺสรติ. ยา อาปตฺติโย สรติ, ตา อาปตฺติโย ฉาเทติ. ยา อาปตฺติโย นสฺสรติ, ตา อาปตฺติโย นจฺฉาเทติ. โส วิพฺภมิตฺวา ปุน อุปสมฺปนฺโน ยา อาปตฺติโย ปุพฺเพ สริตฺวา ฉาเทสิ, ตา อาปตฺติโย ปจฺฉา สริตฺวา นจฺฉาเทติ. ยา อาปตฺติโย ปุพฺเพ อสฺสริตฺวา นจฺฉาเทสิ, ตา อาปตฺติโย ปจฺฉา สริตฺวา นจฺฉาเทติ. ตสฺส ภิกฺขเว ภิกฺขุโน ปุริมสฺมิํ อาปตฺติกฺขนฺเธ ยถา\-ปฏิจฺฉนฺเน ปริวาสํ ทตฺวา มานตฺตํ ทาตพฺพํ. (13)}

\prose{อิธ ปน ภิกฺขเว ภิกฺขุ สมฺพหุลา สํฆาทิเสสา อาปตฺติโย อาปชฺชติ, เอกจฺจา อาปตฺติโย สรติ, เอกจฺจา อาปตฺติโย นสฺสรติ. \csromanfolio{162}ยา อาปตฺติโย สรติ, ตา อาปตฺติโย ฉาเทติ. ยา อาปตฺติโย นสฺสรติ, ตา อาปตฺติโย นจฺฉาเทติ. โส วิพฺภมิตฺวา ปุน อุปสมฺปนฺโน ยา อาปตฺติโย ปุพฺเพ สริตฺวา ฉาเทสิ, ตา อาปตฺติโย ปจฺฉา สริตฺวา นจฺฉาเทติ. ยา อาปตฺติโย ปุพฺเพ อสฺสริตฺวา นจฺฉาเทสิ, ตา อาปตฺติโย ปจฺฉา สริตฺวา ฉาเทติ. ตสฺส ภิกฺขเว ภิกฺขุโน ปุริมสฺมิํ จ ปจฺฉิมสฺมิํ จ อาปตฺติกฺขนฺเธ ยถาปฏิจฺเฉนฺเน ปริวาสํ ทตฺวา มานตฺตํ ทาตพฺพํ. (14)}

\prose{อิธ ปน ภิกฺขเว ภิกฺขุ สมฺพหุลา สํฆาทิเสสา อาปตฺติโย อาปชฺชติ, เอกจฺจา อาปตฺติโย สรติ, เอกจฺจา อาปตฺติโย นสฺสรติ. ยา อาปตฺติโย สรติ, ตา อาปตฺติโย ฉาเทติ. ยา อาปตฺติโย นสฺสรติ, ตา อาปตฺติโย นจฺฉาเทติ. โส วิพฺภมิตฺวา ปุน อุปสมฺปนฺโน ยา อาปตฺติโย ปุพฺเพ สริตฺวา ฉาเทสิ, ตา อาปตฺติโย ปจฺฉา สริตฺวา ฉาเทติ. ยา อาปตฺติโย ปุพฺเพ อสฺสริตฺวา นจฺฉาเทสิ, ตา อาปตฺติโย ปจฺฉา สริตฺวา นจฺฉาเทติ. ตสฺส ภิกฺขเว ภิกฺขุโน ปุริมสฺมิํ จ ปจฺฉิมสฺมิํ จ อาปตฺติกฺขนฺเธ ยถา\-ปฏิจฺฉนฺเน ปริวาสํ ทตฺวา มานตฺตํ ทาตพฺพํ. (15)}

\prose{อิธ ปน ภิกฺขเว ภิกฺขุ สมฺพหุลา สํฆาทิเสสา อาปตฺติโย อาปชฺชติ, เอกจฺจา อาปตฺติโย สรติ, เอกจฺจา อาปตฺติโย นสฺสรติ. ยา อาปตฺติโย สรติ, ตา อาปตฺติโย ฉาเทติ. ยา อาปตฺติโย นสฺสรติ, ตา อาปตฺติโย นจฺฉาเทติ. โส วิพฺภมิตฺวา ปุน อุปสมฺปนฺโน ยา อาปตฺติโย ปุพฺเพ สริตฺวา ฉาเทสิ, ตา อาปตฺติโย ปจฺฉา สริตฺวา ฉาเทติ. ยา อาปติโย ปุพฺเพ อสฺสริตฺวา นจฺฉาเทสิ, ตา อาปตฺติโย ปจฺฉา สริตฺวา ฉาเทติ. ตสฺส ภิกฺขเว ภิกฺขุโน ปุริมสฺมิํ จ ปจฺฉิมสฺมิํ จ อาปตฺติกฺขนฺเธ ยถา\-ปฏิจฺฉนฺเน ปริวาสํ ทตฺวา มานตฺตํ ทาตพฺพํ. (16)}

\proseitem{170}{อิธ ปน ภิกฺขเว ภิกฺขุ สมฺพหุลา สํฆาทิเสสา อาปตฺติโย อาปชฺชติ, เอกจฺจาสุ อาปตฺตีสุ นิพฺเพมติโก, เอกจฺจาสุ อาปตฺตีสุ เวมติโก. ยาสุ อาปตฺตีสุ นิพฺเพมติโก, ตา อาปตฺติโย ฉาเทติ. ยาสุ อาปตฺตีสุ เวมติโก, ตา อาปตฺติโย นจฺฉาเทติ. โส วิพฺภมิตฺวา ปุน อุปสมฺปนฺโน ยา อาปตฺติโย ปุพฺเพ นิพฺเพมติโก ฉาเทสิ, ตา อาปตฺติโย ปจฺฉา นิพฺเพมติโก นจฺฉาเทติ. ยา \csromanfolio{163}อาปตฺติโย ปุพฺเพ เวมติโก นจฺฉาเทสิ, ตา อาปตฺติโย ปจฺฉา นิพฺเพมติโก นจฺฉาเทติ. ตสฺส ภิกฺขเว ภิกฺขุโน ปุริมสฺมิํ อาปตฺติกฺขนฺเธ ยถา\-ปฏิจฺฉนฺเน ปริวาสํ ทตฺวา มานตฺตํ ทาตพฺพํ. (17)}

\prose{อิธ ปน ภิกฺขเว ภิกฺขุ สมฺพหุลา สํฆาทิเสสา อาปตฺติโย อาปชฺชติ, เอกจฺจาสุ อาปตฺตีสุ นิพฺเพมติโก, เอกจฺจาสุ อาปตฺตีสุ เวมติโก. ยาสุ อาปตฺตีสุ นิพฺเพมติโก, ตา อาปตฺติโย ฉาเทติ. ยาสุ อาปตฺตีสุ เวมติโก, ตา อาปตฺติโย นจฺฉาเทติ. โส วิพฺภมิตฺวา ปุน อุปสมฺปนฺโน ยา อาปตฺติโย ปุพฺเพ นิพฺเพมติโก ฉาเทสิ, ตา อาปตฺติโย ปจฺฉา นิพฺเพมติโก นจฺฉาเทติ. ยา อาปตฺติโย ปุพฺเพ เวมติโก นจฺฉาเทสิ, ตา อาปตฺติโย ปจฺฉา นิพฺเพมติโก ฉาเทติ. ตสฺส ภิกฺขเว ภิกฺขุโน ปุริมสฺมิํ จ ปจฺฉิมสฺมิํ จ อาปตฺติกฺขนฺเธ ยถา\-ปฏิจฺฉนฺเน ปริวาสํ ทตฺวา มานตฺตํ ทาตพฺพํ. (18)}

\prose{อิธ ปน ภิกฺขเว ภิกฺขุ สมฺพหุลา สํฆาทิเสสา อาปตฺติโย อาปชฺชติ, เอกจฺจาสุ อาปตฺตีสุ นิพฺเพมติโก, เอกจฺจาสุ อาปตฺตีสุ เวมติโก. ยาสุ อาปตฺตีสุ นิพฺเพมติโก, ตา อาปตฺติโย ฉาเทติ. ยาสุ อาปตฺตีสุ เวมติโก, ตา อาปตฺติโย นจฺฉาเทติ. โส วิพฺภมิตฺวา ปุน อุปสมฺปนฺโน ยา อาปตฺติโย ปุพฺเพ นิพฺเพมติโก ฉาเทสิ, ตา อาปตฺติโย ปจฺฉา นิพฺเพมติโก ฉาเทติ. ยา อาปตฺติโย ปุพฺเพ เวมติโก นจฺฉาเทสิ, ตา อาปตฺติโย ปจฺฉา นิพฺเพมติโก นจฺฉาเทติ. ตสฺส ภิกฺขเว ภิกฺขุโน ปุริมสฺมิํ จ ปจฺฉิมสฺมิํ จ อาปตฺติกฺขนฺเธ ยถา\-ปฏิจฺฉนฺเน ปริวาสํ ทตฺวา มานตฺตํ ทาตพฺพํ. (19)}

\prose{อิธ ปน ภิกฺขเว ภิกฺขุ สมฺพหุลา สํฆาทิเสสา อาปตฺติโย อาปชฺชติ, เอกจฺจาสุ อาปตฺตีสุ นิพฺเพมติโก, เอกจฺจาสุ อาปตฺตีสุ เวมติโก. ยาสุ อาปตฺตีสุ นิพฺเพมติโก, ตา อาปตฺติโย ฉาเทติ. ยาสุ อาปตฺตีสุ เวมติโก, ตา อาปตฺติโย นจฺฉาเทติ. โส วิพฺภมิตฺวา ปุน อุปสมฺปนฺโน ยา อาปตฺติโย ปุพฺเพ นิพฺเพมติโก ฉาเทสิ, ตา อาปตฺติโย ปจฺฉา นิพฺเพมติโก ฉาเทติ. ยา อาปตฺติโย ปุพฺเพ เวมติโก นจฺฉาเทสิ, ตา อาปตฺติโย ปจฺฉา นิพฺเพมติโก ฉาเทติ. ตสฺส ภิกฺขเว ภิกฺขุโน ปุริมสฺมิํ จ ปจฺฉิมสฺมิํ จ อาปตฺติกฺขนฺเธ ยถา\-ปฏิจฺฉนฺเน ปริวาสํ ทตฺวา มานตฺตํ ทาตพฺพํ. (20)}

\proseitemwithcloser{171}{\csromanfolio{164}อิธ ปน ภิกฺขเว ภิกฺขุ สมฺพหุลา สํฆาทิเสสา อาปตฺติโย อาปชฺชิตฺวา อปฺปฏิจฺฉาเทตฺวา สามเณโร โหติ ฯเปฯ อุมฺมตฺตโก โหติ ฯเปฯ ขิตฺตจิตฺโต โหติ ฯเปฯ (ยถา เหฏฺฐา, ตถา วิตฺถาเรตพฺพํ.) เวทนาฏฺโฏ โหติ ฯเปฯ ตสฺส โหนฺติ อาปตฺติโย ปฏิจฺฉนฺนาโยปิ อปฺปฏิจฺฉนฺนาโยปิ ฯเปฯ เอกจฺจา อาปตฺติโย ชานาติ, เอกจฺจา อาปตฺติโย น ชานาติ ฯเปฯ เอกจฺจา อาปตฺติโย สรติ, เอกจฺจา อาปตฺติโย นสฺสรติ ฯเปฯ เอกจฺจาสุ อาปตฺตีสุ นิพฺเพมติโก, เอกจฺจาสุ อาปตฺตีสุ เวมติโก. ยาสุ อาปตฺตีสุ นิพฺเพมติโก, ตา อาปตฺติโย ฉาเทติ. ยาสุ อาปตฺตีสุ เวมติโก, ตา อาปตฺติโย นจฺฉาเทติ. โส เวทนาฏฺโฏ โหติ. โส ปุน อเวทนาฏฺโฏ หุตฺวา ยา อาปตฺติโย ปุพฺเพ นิพฺเพมติโก ฉาเทสิ, ตา อาปตฺติโย ปจฺฉา นิพฺเพมติโก นจฺฉาเทติ. ยา อาปตฺติโย ปุพฺเพ เวมติโก นจฺฉาเทสิ, ตา อาปตฺติโย ปจฺฉา นิพฺเพมติโก นจฺฉาเทติ ฯเปฯ. ยา อาปตฺติโย ปุพฺเพ นิพฺเพมติโก ฉาเทสิ, ตา อาปตฺติโย ปจฺฉา นิพฺเพมติโก นจฺฉาเทติ. ยา อาปตฺติโย ปุพฺเพ เวมติโก นจฺฉาเทสิ, ตา อาปตฺติโย ปจฺฉา นิพฺเพมติโก ฉาเทติ ฯเปฯ. ยา อาปตฺติโย ปุพฺเพ นิพฺเพมติโก ฉาเทสิ, ตา อาปตฺติโย ปจฺฉา นิพฺเพมติโก ฉาเทติ. ยา อาปตฺติโย ปุพฺเพ เวมติโก นจฺฉาเทสิ, ตา อาปตฺติโย ปจฺฉา นิพฺเพมติโก นจฺฉาเทติ ฯเปฯ. ยา อาปตฺติโย ปุพฺเพ นิพฺเพมติโก ฉาเทสิ, ตา อาปตฺติโย ปจฺฉา นิพฺเพมติโก ฉาเทติ. ยา อาปตฺติโย ปุพฺเพ เวมติโก นจฺฉาเทสิ, ตา อาปตฺติโย ปจฺฉา นิพฺเพมติโก ฉาเทติ. ตสฺส ภิกฺขเว ภิกฺขุโน ปุริมสฺมิํ จ ปจฺฉิมสฺมิํ จ อาปตฺติกฺขนฺเธ ยถา\-ปฏิจฺฉนฺเน ปริวาสํ ทตฺวา มานตฺตํ ทาตพฺพํ. (100)}{มานตฺตสตํ นิฏฺฐิตํ.}
\csromanmatikaanchor{mtk.90}
\csromantocmarkref{section}{6. สมูลายสโมธานปริวาสจตุสฺสต}{mtk.90}
\bookmark[dest={mtk.90},level=1]{6. สมูลายสโมธานปริวาสจตุสฺสต}
\csromanheader{1}{6. สมูลาย\-สโมธานปริวาส\-จตุสฺสต}
\proseitem{172}{อิธ ปน ภิกฺขเว ภิกฺขุ ปริวสนฺโต อนฺตรา สมฺพหุลา สํฆาทิเสสา อาปตฺติโย อาปชฺชิตฺวา อปฺปฏิจฺฉาเทตฺวา วิพฺภมติ, \csromanfolio{165}โส ปุน อุปสมฺปนฺโน ตา อาปตฺติโย นจฺฉาเทติ, โส ภิกฺขุ มูลาย ปฏิกสฺสิตพฺโพ. (1)}

\prose{อิธ ปน ภิกฺขเว ภิกฺขุ ปริวสนฺโต อนฺตรา สมฺพหุลา สํฆาทิเสสา อาปตฺติโย อาปชฺชิตฺวา อปฺปฏิจฺฉาเทตฺวา วิพฺภมติ, โส ปุน อุปสมฺปนฺโน ตา อาปตฺติโย ฉาเทติ, โส ภิกฺขุ มูลาย ปฏิกสฺสิตพฺโพ, ยถาปฏิจฺฉนฺนาน\-ญฺจสฺส อาปตฺตีนํ ปุริมาย อาปตฺติยา สโมธาน\-ปริวาโส ทาตพฺโพ. (2)}

\prose{อิธ ปน ภิกฺขเว ภิกฺขุ ปริวสนฺโต อนฺตรา สมฺพหุลา สํฆาทิเสสา อาปตฺติโย อาปชฺชิตฺวา ปฏิจฺฉาเทตฺวา วิพฺภมติ, โส ปุน อุปสมฺปนฺโน ตา อาปตฺติโย นจฺฉาเทติ, โส ภิกฺขุ มูลาย ปฏิกสฺสิตพฺโพ, ยถาปฏิจฺฉนฺนาน\-ญฺจสฺส อาปตฺตีนํ ปุริมาย อาปตฺติยา สโมธาน\-ปริวาโส ทาตพฺโพ. (3)}

\prose{อิธ ปน ภิกฺขเว ภิกฺขุ ปริวสนฺโต อนฺตรา สมฺพหุลา สํฆาทิเสสา อาปตฺติโย อาปชฺชิตฺวา ปฏิจฺฉาเทตฺวา วิพฺภมติ, โส ปุน อุปสมฺปนฺโน ตา อาปตฺติโย ฉาเทติ, โส ภิกฺขุ มูลาย ปฏิกสฺสิตพฺโพ, ยถาปฏิจฺฉนฺนาน\-ญฺจสฺส อาปตฺตีนํ ปุริมาย อาปตฺติยา สโมธาน\-ปริวาโส ทาตพฺโพ. (4)}

\proseitem{173}{อิธ ปน ภิกฺขเว ภิกฺขุ ปริวสนฺโต อนฺตรา สมฺพหุลา สํฆาทิเสสา อาปตฺติโย อาปชฺชติ, ตสฺส โหนฺติ อาปตฺติโย ปฏิจฺฉนฺนาโยปิ อปฺปฏิจฺฉนฺนาโยปิ. โส วิพฺภมิตฺวา ปุน อุปสมฺปนฺโน ยา อาปตฺติโย ปุพฺเพ ฉาเทสิ, ตา อาปตฺติโย ปจฺฉา นจฺฉาเทติ. ยา อาปตฺติโย ปุพฺเพ นจฺฉาเทสิ, ตา อาปตฺติโย ปจฺฉา นจฺฉาเทติ. โส ภิกฺขุ มูลาย ปฏิกสฺสิตพฺโพ, ยถาปฏิจฺฉนฺนาน\-ญฺจสฺส อาปตฺตีนํ ปุริมาย อาปตฺติยา สโมธาน\-ปริวาโส ทาตพฺโพ. (5)}

\prose{อิธ ปน ภิกฺขเว ภิกฺขุ ปริวสนฺโต อนฺตรา สมฺพหุลา สํฆาทิเสสา อาปตฺติโย อาปชฺชติ, ตสฺส โหนฺติ อาปตฺติโย ปฏิจฺฉนฺนาโยปิ อปฺปฏิจฺฉนฺนาโยปิ. โส วิพฺภมิตฺวา ปุน อุปสมฺปนฺโน ยา อาปตฺติโย ปุพฺเพ ฉาเทสิ, ตา อาปตฺติโย ปจฺฉา นจฺฉาเทติ. \csromanfolio{166}ยา อาปตฺติโย ปุพฺเพ นจฺฉาเทสิ, ตา อาปตฺติโย ปจฺฉา ฉาเทติ. โส ภิกฺขุ มูลาย ปฏิกสฺสิตพฺโพ, ยถาปฏิจฺฉนฺนาน\-ญฺจสฺส อาปตฺตีนํ ปุริมาย อาปตฺติยา สโมธาน\-ปริวาโส ทาตพฺโพ. (6)}

\prose{อิธ ปน ภิกฺขเว ภิกฺขุ ปริวสนฺโต อนฺตรา สมฺพหุลา สํฆาทิเสสา อาปตฺติโย อาปชฺชติ, ตสฺส โหนฺติ อาปตฺติโย ปฏิจฺฉนฺนาโยปิ อปฺปฏิจฺฉนฺนาโยปิ. โส วิพฺภมิตฺวา ปุน อุปสมฺปนฺโน ยา อาปตฺติโย ปุพฺเพ ฉาเทสิ, ตา อาปตฺติโย ปจฺฉา ฉาเทติ. ยา อาปตฺติโย ปุพฺเพ นจฺฉาเทสิ, ตา อาปตฺติโย ปจฺฉา นจฺฉาเทติ. โส ภิกฺขุ มูลาย ปฏิกสฺสิตพฺโพ, ยถาปฏิจฺฉนฺนาน\-ญฺจสฺส อาปตฺตีนํ ปุริมาย อาปตฺติยา สโมธาน\-ปริวาโส ทาตพฺโพ. (7)}

\prose{อิธ ปน ภิกฺขเว ภิกฺขุ ปริวสนฺโต อนฺตรา สมฺพหุลา สํฆาทิเสสา อาปตฺติโย อาปชฺชติ, ตสฺส โหนฺติ อาปตฺติโย ปฏิจฺฉนฺนาโยปิ อปฺปฏิจฺฉนฺนาโยปิ. โส วิพฺภมิตฺวา ปุน อุปสมฺปนฺโน ยา อาปตฺติโย ปุพฺเพ ฉาเทสิ, ตา อาปตฺติโย ปจฺฉา ฉาเทติ. ยา อาปตฺติโย ปุพฺเพ นจฺฉาเทสิ, ตา อาปตฺติโย ปจฺฉา ฉาเทติ. โส ภิกฺขุ มูลาย ปฏิกสฺสิตพฺโพ, ยถาปฏิจฺฉนฺนาน\-ญฺจสฺส อาปตฺตีนํ ปุริมาย อาปตฺติยา สโมธาน\-ปริวาโส ทาตพฺโพ. (8)}

\proseitem{174}{อิธ ปน ภิกฺขเว ภิกฺขุ ปริวสนฺโต อนฺตรา สมฺพหุลา สํฆาทิเสสา อาปตฺติโย อาปชฺชติ, เอกจฺจา อาปตฺติโย ชานาติ, เอกจฺจา อาปตฺติโย น ชานาติ. ยา อาปตฺติโย ชานาติ, ตา อาปตฺติโย ฉาเทติ. ยา อาปตฺติโย น ชานาติ, ตา อาปตฺติโย นจฺฉาเทติ. โส วิพฺภมิตฺวา ปุน อุปสมฺปนฺโน ยา อาปตฺติโย ปุพฺเพ ชานิตฺวา ฉาเทสิ, ตา อาปตฺติโย ปจฺฉา ชานิตฺวา นจฺฉาเทติ. ยา อาปตฺติโย ปุพฺเพ อชานิตฺวา นจฺฉาเทสิ, ตา อาปตฺติโย ปจฺฉา ชานิตฺวา นจฺฉาเทติ. โส ภิกฺขุ มูลาย ปฏิกสฺสิตพฺโพ, ยถาปฏิจฺฉนฺนาน\-ญฺจสฺส อาปตฺตีนํ ปุริมาย อาปตฺติยา สโมธาน\-ปริวาโส ทาตพฺโพ. (9)}

\prose{อิธ ปน ภิกฺขเว ภิกฺขุ ปริวสนฺโต อนฺตรา สมฺพหุลา สํฆาทิเสสา อาปตฺติโย อาปชฺชติ, เอกจฺจา อาปตฺติโย ชานาติ, เอกจฺจา \csromanfolio{167}อาปตฺติโย น ชานาติ. ยา อาปตฺติโย ชานาติ, ตา อาปตฺติโย ฉาเทติ. ยา อาปตฺติโย น ชานาติ, ตา อาปตฺติโย นจฺฉาเทติ. โส วิพฺภมิตฺวา ปุน อุปสมฺปนฺโน ยา อาปตฺติโย ปุพฺเพ ชานิตฺวา ฉาเทสิ, ตา อาปตฺติโย ปจฺฉา ชานิตฺวา นจฺฉาเทติ. ยา อาปตฺติโย ปุพฺเพ อชานิตฺวา นจฺฉาเทสิ, ตา อาปตฺติโย ปจฺฉา ชานิตฺวา ฉาเทติ. โส ภิกฺขุ มูลาย ปฏิกสฺสิตพฺโพ, ยถาปฏิจฺฉนฺนาน\-ญฺจสฺส อาปตฺตีนํ ปุริมาย อาปตฺติยา สโมธาน\-ปริวาโส ทาตพฺโพ. (10)}

\prose{อิธ ปน ภิกฺขุ ปริวสนฺโต อนฺตรา สมฺพหุลา สํฆาทิเสสา อาปตฺติโย อาปชฺชติ, เอกจฺจา อาปตฺติโย ชานาติ, เอกจฺจา อาปตฺติโย น ชานาติ. ยา อาปตฺติโย ชานาติ, ตา อาปตฺติโย ฉาเทติ. ยา อาปตฺติโย น ชานาติ, ตา อาปตฺติโย นจฺฉาเทติ. โส วิพฺภมิตฺวา ปุน อุปสมฺปนฺโน ยา อาปตฺติโย ปุพฺเพ ชานิตฺวา ฉาเทสิ, ตา อาปตฺติโย ปจฺฉา ชานิตฺวา ฉาเทติ. ยา อาปตฺติโย ปุพฺเพ อชานิตฺวา นจฺฉาเทสิ, ตา อาปตฺติโย ปจฺฉา ชานิตฺวา นจฺฉาเทติ. โส ภิกฺขุ มูลาย ปฏิกสฺสิตพฺโพ, ยถาปฏิจฺฉนฺนาน\-ญฺจสฺส อาปตฺตีนํ ปุริมาย อาปตฺติยา สโมธาน\-ปริวาโส ทาตพฺโพ. (11)}

\prose{อิธ ปน ภิกฺขเว ภิกฺขุ ปริวสนฺโต อนฺตรา สมฺพหุลา สํฆาทิเสสา อาปตฺติโย อาปชฺชติ, เอกจฺจา อาปตฺติโย ชานาติ, เอกจฺจา อาปตฺติโย น ชานาติ. ยา อาปตฺติโย ชานาติ, ตา อาปตฺติโย ฉาเทติ. ยา อาปตฺติโย น ชานาติ, ตา อาปตฺติโย นจฺฉาเทติ. โส วิพฺภมิตฺวา ปุน อุปสมฺปนฺโน ยา อาปตฺติโย ปุพฺเพ ชานิตฺวา ฉาเทสิ, ตา อาปตฺติโย ปจฺฉา ชานิตฺวา ฉาเทติ. ยา อาปตฺติโย ปุพฺเพ อชานิตฺวา นจฺฉาเทสิ, ตา อาปตฺติโย ปจฺฉา ชานิตฺวา ฉาเทติ. โส ภิกฺขุ มูลาย ปฏิกสฺสิตพฺโพ, ยถาปฏิจฺฉนฺนาน\-ญฺจสฺส อาปตฺตีนํ ปุริมาย อาปตฺติยา สโมธาน\-ปริวาโส ทาตพฺโพ. (12)}

\proseitem{175}{อิธ ปน ภิกฺขเว ภิกฺขุ ปริวสนฺโต อนฺตรา สมฺพหุลา สํฆาทิเสสา อาปตฺติโย อาปชฺชติ, เอกจฺจา อาปตฺติโย สรติ เอกจฺจา อาปตฺติโย นสฺสรติ. ยา อาปตฺติโย สรติ, ตา \csromanfolio{168}อาปตฺติโย ฉาเทติ. ยา อาปตฺติโย นสฺสรติ, ตา อาปตฺติโย นจฺฉาเทติ. โส วิพฺภมิตฺวา ปุน อุปสมฺปนฺโน ยา อาปตฺติโย ปุพฺเพ สริตฺวา ฉาเทสิ, ตา อาปตฺติโย ปจฺฉา สริตฺวา นจฺฉาเทติ. ยา อาปตฺติโย ปุพฺเพ อสฺสริตฺวา นจฺฉาเทสิ, ตา อาปตฺติโย ปจฺฉา สริตฺวา นจฺฉาเทติ. โส ภิกฺขุ มูลาย ปฏิกสฺสิตพฺโพ, ยถา ปฏิจฺฉนฺนา\-นญฺจสฺส อาปตฺตีนํ ปุริมาย อาปตฺติยา สโมธาน\-ปริวาโส ทาตพฺโพ. (13)}

\prose{อิธ ปน ภิกฺขเว ภิกฺขุ ปริวสนฺโต อนฺตรา สมฺพหุลา สํฆาทิเสสา อาปตฺติโย อาปชฺชติ, เอกจฺจา อาปตฺติโย สรติ, เอกจฺจา อาปตฺติโย นสฺสรติ. ยา อาปตฺติโย สรติ, ตา อาปตฺติโย ฉาเทติ. ยา อาปตฺติโย นสฺสรติ, ตา อาปตฺติโย นจฺฉาเทติ. โส วิพฺภมิตฺวา ปุน อุปสมฺปนฺโน ยา อาปตฺติโย ปุพฺเพ สริตฺวา ฉาเทสิ, ตา อาปตฺติโย ปจฺฉา สริตฺวา นจฺฉาเทติ. ยา อาปตฺติโย ปุพฺเพ อสฺสริตฺวา นจฺฉาเทสิ, ตา อาปตฺติโย ปจฺฉา สริตฺวา ฉาเทติ. โส ภิกฺขุ มูลาย ปฏิกสฺสิตพฺโพ, ยถาปฏิจฺฉนฺนาน\-ญฺจสฺส อาปตฺตีนํ ปุริมาย อาปตฺติยา สโมธาน\-ปริวาโส ทาตพฺโพ. (14)}

\prose{อิธ ปน ภิกฺขเว ภิกฺขุ ปริวสนฺโต อนฺตรา สมฺพหุลา สํฆาทิเสสา อาปตฺติโย อาปชฺชติ, เอกจฺจา อาปตฺติโย สรติ, เอกจฺจา อาปตฺติโย นสฺสรติ. ยา อาปตฺติโย สรติ, ตา อาปตฺติโย ฉาเทติ. ยา อาปตฺติโย นสฺสรติ, ตา อาปตฺติโย นจฺฉาเทติ. โส วิพฺภมิตฺวา ปุน อุปสมฺปนฺโน ยา อาปตฺติโย ปุพฺเพ สริตฺวา ฉาเทสิ, ตา อาปตฺติโย ปจฺฉา สริตฺวา ฉาเทติ. ยา อาปตฺติโย ปุพฺเพ อสฺสริตฺวา นจฺฉาเทสิ, ตา อาปตฺติโย ปจฺฉา สริตฺวา นจฺฉาเทติ. โส ภิกฺขุ มูลาย ปฏิกสฺสิตพฺโพ, ยถาปฏิจฺฉนฺนาน\-ญฺจสฺส อาปตฺตีนํ ปุริมาย อาปตฺติยา สโมธาน\-ปริวาโส ทาตพฺโพ. (15)}

\prose{อิธ ปน ภิกฺขเว ภิกฺขุ ปริวสนฺโต อนฺตรา สมฺพหุลา สํฆาทิเสสา อาปตฺติโย อาปชฺชติ, เอกจฺจา อาปตฺติโย สรติ, เอกจฺจา อาปตฺติโย นสฺสรติ. ยา อาปตฺติโย สรติ, ตา อาปตฺติโย ฉาเทติ. ยา อาปตฺติโย นสฺสรติ, ตา อาปตฺติโย นจฺฉาเทติ. โส วิพฺภมิตฺวา ปุน อุปสมฺปนฺโน ยา อาปตฺติโย ปุพฺเพ สริตฺวา \csromanfolio{169}ฉาเทสิ, ตา อาปตฺติโย ปจฺฉา สริตฺวา ฉาเทติ. ยา อาปตฺติโย ปุพฺเพ อสฺสริตฺวา นจฺฉาเทสิ, ตา อาปตฺติโย ปจฺฉา สริตฺวา ฉาเทติ. โส ภิกฺขุ มูลาย ปฏิกสฺสิตพฺโพ, ยถาปฏิจฺฉนฺนาน\-ญฺจสฺส อาปตฺตีนํ ปุริมาย อาปตฺติยา สโมธาน\-ปริวาโส ทาตพฺโพ. (16)}

\proseitem{176}{อิธ ปน ภิกฺขเว ภิกฺขุ ปริวสนฺโต อนฺตรา สมฺพหุลา สํฆาทิเสสา อาปตฺติโย อาปชฺชติ, เอกจฺจาสุ อาปตฺตีสุ นิพฺเพมติโก, เอกจฺจาสุ อาปตฺตีสุ เวมติโก. ยาสุ อาปตฺตีสุ นิพฺเพมติโก, ตา อาปตฺติโย ฉาเทติ. ยาสุ อาปตฺตีสุ เวมติโก, ตา อาปตฺติโย นจฺฉาเทติ. โส วิพฺภมิตฺวา ปุน อุปสมฺปนฺโน ยา อาปตฺติโย ปุพฺเพ นิพฺเพมติโก ฉาเทสิ, ตา อาปตฺติโย ปจฺฉา นิพฺเพมติโก นจฺฉาเทติ. ยา อาปตฺติโย ปุพฺเพ เวมติโก นจฺฉาเทสิ, ตา อาปตฺติโย ปจฺฉา นิพฺเพมติโก นจฺฉาเทติ. โส ภิกฺขุ มูลาย ปฏิกสฺสิตพฺโพ, ยถาปฏิจฺฉนฺนาน\-ญฺจสฺส อาปตฺตีนํ ปุริมาย อาปตฺติยา สโมธาน\-ปริวาโส ทาตพฺโพ. (17)}

\prose{อิธ ปน ภิกฺขเว ภิกฺขุ ปริวสนฺโต อนฺตรา สมฺพหุลา สํฆาทิเสสา อาปตฺติโย อาปชฺชติ, เอกจฺจาสุ อาปตฺตีสุ นิพฺเพมติโก, เอกจฺจาสุ อาปตฺตีสุ เวมติโก. ยาสุ อาปตฺตีสุ นิพฺเพมติโก, ตา อาปตฺติโย ฉาเทติ. ยาสุ อาปตฺตีสุ เวมติโก, ตา อาปตฺติโย นจฺฉาเทติ. โส วิพฺภมิตฺวา ปุน อุปสมฺปนฺโน ยา อาปตฺติโย ปุพฺเพ นิพฺเพมติโก ฉาเทสิ, ตา อาปตฺติโย ปจฺฉา นิพฺเพมติโก นจฺฉาเทติ. ยา อาปตฺติโย ปุพฺเพ เวมติโก นจฺฉาเทสิ, ตา อาปตฺติโย ปจฺฉา นิพฺเพมติโก ฉาเทติ. โส ภิกฺขุ มูลาย ปฏิกสฺสิตพฺโพ, ยถาปฏิจฺฉนฺนาน\-ญฺจสฺส อาปตฺตีนํ ปุริมาย อาปตฺติยา สโมธาน\-ปริวาโส ทาตพฺโพ. (18)}

\prose{อิธ ปน ภิกฺขเว ภิกฺขุ ปริวสนฺโต อนฺตรา สมฺพหุลา สํฆาทิเสสา อาปตฺติโย อาปชฺชติ, เอกจฺจาสุ อาปตฺตีสุ นิพฺเพมติโก, เอกจฺจาสุ อาปตฺตีสุ เวมติโก. ยาสุ อาปตฺตีสุ นิพฺเพมติโก, ตา อาปตฺติโย ฉาเทติ. ยาสุ อาปตฺตีสุ เวมติโก, ตา อาปตฺติโย นจฺฉาเทติ. โส วิพฺภมิตฺวา ปุน อุปสมฺปนฺโน ยา \csromanfolio{170}อาปตฺติโย ปุพฺเพ นิพฺเพมติโก ฉาเทสิ, ตา อาปตฺติโย ปจฺฉา นิพฺเพมติโก ฉาเทติ. ยา อาปตฺติโย ปุพฺเพ เวมติโก นจฺฉาเทสิ, ตา อาปตฺติโย ปจฺฉา นิพฺเพมติโก นจฺฉาเทติ. โส ภิกฺขุ มูลาย ปฏิกสฺสิตพฺโพ, ยถาปฏิจฺฉนฺนาน\-ญฺจสฺส อาปตฺตีนํ ปุริมาย อาปตฺติยา สโมธาน\-ปริวาโส ทาตพฺโพ. (19)}

\prose{อิธ ปน ภิกฺขเว ภิกฺขุ ปริวสนฺโต อนฺตรา สมฺพหุลา สํฆาทิเสสา อาปตฺติโย อาปชฺชติ, เอกจฺจาสุ อาปตฺตีสุ นิพฺเพมติโก, เอกจฺจาสุ อาปตฺตีสุ เวมติโก. ยาสุ อาปตฺตีสุ นิพฺเพมติโก, ตา อาปตฺติโย ฉาเทติ. ยาสุ อาปตฺตีสุ เวมติโก, ตา อาปตฺติโย นจฺฉาเทติ. โส วิพฺภมิตฺวา\csromansharedfootnote{9fdded679207dafb}{โส ภิกฺขุ วิพฺภมิตฺวา (ก)} ปุน อุปสมฺปนฺโน ยา อาปตฺติโย ปุพฺเพ นิพฺเพมติโก ฉาเทสิ, ตา อาปตฺติโย ปจฺฉา นิพฺเพมติโก ฉาเทติ. ยา อาปตฺติโย ปุพฺเพ เวมติโก นจฺฉาเทสิ, ตา อาปตฺติโย ปจฺฉา นิพฺเพมติโก ฉาเทติ. โส ภิกฺขุ มูลาย ปฏิกสฺสิตพฺโพ, ยถาปฏิจฺฉนฺนาน\-ญฺจสฺส อาปตฺตีนํ ปุริมาย อาปตฺติยา สโมธาน\-ปริวาโส ทาตพฺโพ. (20)}

\proseitem{177}{อิธ ปน ภิกฺขเว ภิกฺขุ ปริวสนฺโต อนฺตรา สมฺพหุลา สํฆาทิเสสา อาปตฺติโย อาปชฺชิตฺวา อปฺปฏิจฺฉาเทตฺวา สามเณโร โหติ ฯเปฯ อุมฺมตฺตโก โหติ ฯเปฯ ขิตฺตจิตฺโต โหติ ฯเปฯ เวทนาฏฺโฏ โหติ ฯเปฯ ตสฺส โหนฺติ อาปตฺติโย ปฏิจฺฉนฺนาโยปิ อปฺปฏิจฺฉนฺนาโยปิ. (ยถา เหฏฺฐา วิตฺถาริตํ, ตถา วิตฺถาเรตพฺพํ) ฯเปฯ เอกจฺจา อาปตฺติโย ชานาติ, เอกจฺจา อาปตฺติโย น ชานาติ ฯเปฯ เอกจฺจา อาปตฺติโย สรติ, เอกจฺจา อาปตฺติโย นสฺสรติ ฯเปฯ เอกจฺจาสุ อาปตฺตีสุ นิพฺเพมติโก, เอกจฺจาสุ อาปตฺตีสุ เวมติโก. ยาสุ อาปตฺตีสุ นิพฺเพมติโก, ตา อาปตฺติโย ฉาเทติ. ยาสุ อาปตฺตีสุ เวมติโก, ตา อาปตฺติโย นจฺฉาเทติ. โส เวทนาฏฺโฏ โหติ, โส ปุน อเวทนาฏฺโฏ หุตฺวา ยา อาปตฺติโย ปุพฺเพ นิพฺเพมติโก ฉาเทสิ, ตา อาปตฺติโย ปจฺฉา นิพฺเพมติโก นจฺฉาเทติ. ยา อาปตฺติโย ปุพฺเพ เวมติโก นจฺฉาเทสิ, ตา อาปตฺติโย ปจฺฉา นิพฺเพมติโก นจฺฉาเทติ ฯเปฯ. ยา อาปตฺติโย ปุพฺเพ นิพฺเพมติโก ฉาเทสิ, ตา อาปตฺติโย ปจฺฉา นิพฺเพมติโก \csromanfolio{171}นจฺฉาเทติ. ยา อาปตฺติโย ปุพฺเพ เวมติโก นจฺฉาเทสิ, ตา อาปตฺติโย ปจฺฉา นิพฺเพมติโก ฉาเทติ ฯเปฯ. ยา อาปตฺติโย ปุพฺเพ นิพฺเพมติโก ฉาเทสิ, ตา อาปตฺติโย ปจฺฉา นิพฺเพมติโก ฉาเทติ. ยา อาปตฺติโย ปุพฺเพ เวมติโก นจฺฉาเทสิ, ตา อาปตฺติโย ปจฺฉา นิพฺเพมติโก นจฺฉาเทติ ฯเปฯ. ยา อาปตฺติโย ปุพฺเพ นิพฺเพมติโก ฉาเทสิ, ตา อาปตฺติโย ปจฺฉา นิพฺเพมติโก ฉาเทติ. ยา อาปตฺติโย ปุพฺเพ เวมติโก นจฺฉาเทสิ, ตา อาปตฺติโย ปจฺฉา นิพฺเพมติโก ฉาเทติ. โส ภิกฺขุ มูลาย ปฏิกสฺสิตพฺโพ, ยถาปฏิจฺฉนฺนาน\-ญฺจสฺส อาปตฺตีนํ ปุริมาย อาปตฺติยา สโมธาน\-ปริวาโส ทาตพฺโพ. (100)}

\proseitem{178}{อิธ ปน ภิกฺขเว ภิกฺขุ มานตฺตารโห ฯเปฯ มานตฺตํ จรนฺโต ฯเปฯ อพฺภานารโห อนฺตรา สมฺพหุลา สํฆาทิเสสา อาปตฺติโย อาปชฺชิตฺวา อปฺปฏิจฺฉาเทตฺวา วิพฺภมติ ฯเปฯ (มานตฺตารโห จ มานตฺตจารี จ อพฺภานารโห จ ยถา ปริวาโส วิตฺถาริโต, ตถา วิตฺถาเรตพฺโพ.) (320)}

\proseitemwithcloser{179}{อิธ ปน ภิกฺขเว ภิกฺขุ อพฺภานารโห อนฺตรา สมฺพหุลา สํฆาทิเสสา อาปตฺติโย อาปชฺชิตฺวา อปฺปฏิจฺฉาเทตฺวา สามเณโร โหติ ฯเปฯ อุมฺมตฺตโก โหติ ฯเปฯ ขิตฺตจิตฺโต โหติ ฯเปฯ เวทนาฏฺโฏ โหติ ฯเปฯ ตสฺส โหนฺติ อาปตฺติโย ปฏิจฺฉนฺนาโยปิ อปฺปฏิจฺฉนฺนาโยปิ ฯเปฯ เอกจฺจา อาปตฺติโย ชานาติ, เอกจฺจา อาปตฺติโย น ชานาติ ฯเปฯ เอกจฺจา อาปตฺติโย สรติ, เอกจฺจา อาปตฺติโย นสฺสรติ ฯเปฯ เอกจฺจาสุ อาปตฺตีสุ นิพฺเพมติโก, เอกจฺจาสุ อาปตฺตีสุ เวมติโก. ยาสุ อาปตฺตีสุ นิพฺเพมติโก, ตา อาปตฺติโย ฉาเทติ. ยาสุ อาปตฺตีสุ เวมติโก, ตา อาปตฺติโย นจฺฉาเทติ. โส เวทนาฏฺโฏ โหติ, โส ปุน อเวทนาฏฺโฏ หุตฺวา ยา อาปตฺติโย ปุพฺเพ นิพฺเพมติโก ฉาเทสิ, ตา อาปตฺติโย ปจฺฉา นิพฺเพมติโก นจฺฉาเทติ. ยา อาปตฺติโย ปุพฺเพ เวมติโก นจฺฉาเทสิ, ตา อาปตฺติโย ปจฺฉา นิพฺเพมติโก นจฺฉาเทติ ฯเปฯ. ยา อาปตฺติโย ปุพฺเพ นิพฺเพมติโก ฉาเทสิ, ตา อาปตฺติโย ปจฺฉา นิพฺเพมติโก นจฺฉาเทติ. ยา อาปตฺติโย ปุพฺเพ เวมติโก \csromanfolio{172}นจฺฉาเทสิ, ตา อาปตฺติโย ปจฺฉา นิพฺเพมติโก ฉาเทติ ฯเปฯ. ยา อาปตฺติโย ปุพฺเพ นิพฺเพมติโก ฉาเทสิ, ตา อาปตฺติโย ปจฺฉา นิพฺเพมติโก ฉาเทติ. ยา อาปตฺติโย ปุพฺเพ เวมติโก นจฺฉาเทสิ, ตา อาปตฺติโย ปจฺฉา นิพฺเพมติโก นจฺฉาเทติ ฯเปฯ. ยา อาปตฺติโย ปุพฺเพ นิพฺเพมติโก ฉาเทสิ, ตา อาปตฺติโย ปจฺฉา นิพฺเพมติโก ฉาเทติ. ยา อาปตฺติโย ปุพฺเพ เวมติโก นจฺฉาเทสิ, ตา อาปตฺติโย ปจฺฉา นิพฺเพมติโก ฉาเทติ. โส ภิกฺขุ มูลาย ปฏิกสฺสิตพฺโพ, ยถาปฏิจฺฉนฺนาน\-ญฺจสฺส อาปตฺตีนํ ปุริมาย อาปตฺติยา สโมธาน\-ปริวาโส ทาตพฺโพ. (400)}{สมูลาย\-สโมธานปริวาส\-จตุสฺสตํ นิฏฺฐิตํ.}
\csromanmatikaanchor{mtk.91}
\csromantocmarkref{section}{7. ปริมาณาทิวารอฏฺฐก}{mtk.91}
\bookmark[dest={mtk.91},level=1]{7. ปริมาณาทิวารอฏฺฐก}
\csromanheader{1}{7. ปริมาณา\-ทิ\-วาร\-อฏฺฐก}
\proseitemwithcloser{180}{อิธ ปน ภิกฺขเว ภิกฺขุ สมฺพหุลา สํฆาทิเสสา อาปตฺติโย อาปชฺชิตฺวา ปริมาณา อปฺปฏิจฺฉาเทตฺวา ฯเปฯ อปริมาณา อปฺปฏิจฺฉาเทตฺวา ฯเปฯ เอกนามา อปฺปฏิจฺฉาเทตฺวา ฯเปฯ นานานามา อปฺปฏิจฺฉาเทตฺวา ฯเปฯ สภาคา อปฺปฏิจฺฉาเทตฺวา ฯเปฯ วิสภาคา อปฺปฏิจฺฉาเทตฺวา ฯเปฯ ววตฺถิตา อปฺปฏิจฺฉาเทตฺวา ฯเปฯ สมฺภินฺนา อปฺปฏิจฺฉาเทตฺวา วิพฺภมติ ฯเปฯ. (ยถา เหฏฺฐา, ตถา วิตฺถาเรตพฺพํ.)}{ปริมาณา\-ทิ\-วาร\-อฏฺฐกํ นิฏฺฐิตํ.}
\csromanmatikaanchor{mtk.92}
\csromantocmarkref{section}{8. เทฺวภิกฺขุวารเอกาทสก}{mtk.92}
\bookmark[dest={mtk.92},level=1]{8. เทฺวภิกฺขุวารเอกาทสก}
\csromanheader{1}{8. เทฺวภิกฺขุวารเอกาทสก}
\proseitem{181}{เทฺว ภิกฺขู สํฆาทิเสสํ อาปนฺนา โหนฺติ, เต สํฆาทิเสเส สํฆาทิเสสทิฏฺฐิโน โหนฺติ, เอโก ฉาเทติ, เอโก นจฺฉาเทติ. โย ฉาเทติ, โส ทุกฺกฏํ เทสาเปตพฺโพ, ยถา\-ปฏิจฺฉนฺเน\csromansharedfootnote{9c815f97492308c6}{ยถาปฏิจฺฉนฺนานํ (สี)} จสฺส ปริวาสํ ทตฺวา อุภินฺนมฺปิ มานตฺตํ ทาตพฺพํ. (1)}

\prose{\csromanfolio{173}เทฺว ภิกฺขู สํฆาทิเสสํ อาปนฺนา โหนฺติ, เต สํฆาทิเสเส เวมติกา โหนฺติ, เอโก ฉาเทติ, เอโก นจฺฉาเทติ. โย ฉาเทติ, โส ทุกฺกฏํ เทสาเปตพฺโพ, ยถา\-ปฏิจฺฉนฺเน จสฺส ปริวาสํ ทตฺวา อุภินฺนมฺปิ มานตฺตํ ทาตพฺพํ. (2)}

\prose{เทฺว ภิกฺขู สํฆาทิเสสํ อาปนฺนา โหนฺติ, เต สํฆาทิเสเส มิสฺสก\-ทิฏฺฐิโน โหนฺติ, เอโก ฉาเทติ, เอโก นจฺฉาเทติ. โย ฉาเทติ, โส ทุกฺกฏํ เทสาเปตพฺโพ, ยถา\-ปฏิจฺฉนฺเน จสฺส ปริวาสํ ทตฺวา อุภินฺนมฺปิ มานตฺตํ ทาตพฺพํ. (3)}

\prose{เทฺว ภิกฺขู มิสฺสกํ อาปนฺนา โหนฺติ, เต มิสฺสเก สํฆาทิเสสทิฏฺฐิโน โหนฺติ, เอโก ฉาเทติ, เอโก นจฺฉาเทติ. โย ฉาเทติ, โส ทุกฺกฏํ เทสาเปตพฺโพ, ยถา\-ปฏิจฺฉนฺเน จสฺส ปริวาสํ ทตฺวา อุภินฺนมฺปิ มานตฺตํ ทาตพฺพํ. (4)}

\prose{เทฺว ภิกฺขู มิสฺสกํ อาปนฺนา โหนฺติ, เต มิสฺสเก มิสฺสก\-ทิฏฺฐิโน โหนฺติ, เอโก ฉาเทติ, เอโก นจฺฉาเทติ. โย ฉาเทติ, โส ทุกฺกฏํ เทสาเปตพฺโพ, ยถา\-ปฏิจฺฉนฺเน จสฺส ปริวาสํ ทตฺวา อุภินฺนมฺปิ มานตฺตํ ทาตพฺพํ. (5)}

\prose{เทฺว ภิกฺขู สุทฺธกํ อาปนฺนา โหนฺติ, เต สุทฺธเก สํฆาทิเสสทิฏฺฐิโน โหนฺติ, เอโก ฉาเทติ, เอโก นจฺฉาเทติ. โย ฉาเทติ, โส ทุกฺกฏํ เทสาเปตพฺโพ, อุโภปิ ยถาธมฺมํ การาเปตพฺพา. (6)}

\prose{เทฺว ภิกฺขู สุทฺธกํ อาปนฺนา โหนฺติ, เต สุทฺธเก สุทฺธก\-ทิฏฺฐิโน โหนฺติ, เอโก ฉาเทติ, เอโก นจฺฉาเทติ. โย ฉาเทติ, โส ทุกฺกฏํ เทสาเปตพฺโพ, อุโภปิ ยถาธมฺมํ การาเปตพฺพา. (7)}

\prose{เทฺว ภิกฺขู สํฆาทิเสสํ อาปนฺนา โหนฺติ, เต สํฆาทิเสเส สํฆาทิเสสทิฏฺฐิโน โหนฺติ, เอกสฺส โหติ “อาโรเจสฺสามี”ติ, เอกสฺส โหติ “น อาโรเจสฺสามี”ติ, โส ปฐมมฺปิ ยามํ ฉาเทติ, ทุติยมฺปิ ยามํ ฉาเทติ, ตติยมฺปิ ยามํ ฉาเทติ, อุฏฺฐิเต อรุเณ ฉนฺนา โหติ อาปตฺติ. โย ฉาเทติ, โส ทุกฺกฏํ เทสาเปตพฺโพ, ยถา\-ปฏิจฺฉนฺเน จสฺส ปริวาสํ ทตฺวา อุภินฺนมฺปิ มานตฺตํ ทาตพฺพํ. (8)}

\prose{\csromanfolio{174}เทฺว ภิกฺขู สํฆาทิเสสํ อาปนฺนา โหนฺติ, เต สํฆาทิเสเส สํฆาทิเสสทิฏฺฐิโน โหนฺติ, เต คจฺฉนฺติ “อาโรเจสฺสามา”ติ, เอกสฺส อนฺตรามคฺเค มกฺขธมฺโม อุปฺปชฺชติ “น อาโรเจสฺสามี”ติ, โส ปฐมมฺปิ ยามํ ฉาเทติ, ทุติยมฺปิ ยามํ ฉาเทติ, ตติยมฺปิ ยามํ ฉาเทติ, อุฏฺฐิเต อรุเณ ฉนฺนา โหติ อาปตฺติ. โย ฉาเทติ, โส ทุกฺกฏํ เทสาเปตพฺโพ, ยถา\-ปฏิจฺฉนฺเน จสฺส ปริวาสํ ทตฺวา อุภินฺนมฺปิ มานตฺตํ ทาตพฺพํ. (9)}

\prose{เทฺว ภิกฺขู สํฆาทิเสสํ อาปนฺนา โหนฺติ, เต สํฆาทิเสเส สํฆาทิเสสทิฏฺฐิโน โหนฺติ, เต อุมฺมตฺตกา โหนฺติ, เต ปจฺฉา อนุมฺมตฺตกา หุตฺวา เอโก ฉาเทติ, เอโก นจฺฉาเทติ. โย ฉาเทติ, โส ทุกฺกฏํ เทสาเปตพฺโพ, ยถา\-ปฏิจฺฉนฺเน จสฺส ปริวาสํ ทตฺวา อุภินฺนมฺปิ มานตฺตํ ทาตพฺพํ. (10)}

\prosewithcloser{เทฺว ภิกฺขู สํฆาทิเสสํ อาปนฺนา โหนฺติ, เต ปาติโมกฺเข อุทฺทิสฺสมาเน เอวํ วทนฺติ “อิทาเนว โข มยํ ชานาม อยมฺปิ กิร ธมฺโม สุตฺตาคโต สุตฺต\-ปริยาปนฺโน อนฺวทฺธมาสํ อุทฺเทสํ อาคจฺฉตี”ติ, เต สํฆาทิเสเส สํฆาทิเสสทิฏฺฐิโน โหนฺติ, เอโก ฉาเทติ, เอโก นจฺฉาเทติ. โย ฉาเทติ, โส ทุกฺกฏํ เทสาเปตพฺโพ, ยถา\-ปฏิจฺฉนฺเน จสฺส ปริวาสํ ทตฺวา อุภินฺนมฺปิ มานตฺตํ ทาตพฺพํ. (11)}{เทฺวภิกฺขุวารเอกาทสกํ นิฏฺฐิตํ.}
\csromanmatikaanchor{mtk.93}
\csromantocmarkref{section}{9. มูลายอวิสุทฺธินวก}{mtk.93}
\bookmark[dest={mtk.93},level=1]{9. มูลายอวิสุทฺธินวก}
\csromanheader{1}{9. มูลาย\-อวิสุทฺธิ\-นวก}
\proseitem{182}{อิธ ปน ภิกฺขเว ภิกฺขุ สมฺพหุลา สํฆาทิเสสา อาปตฺติโย อาปชฺชติ ปริมาณมฺปิ อปริมาณมฺปิ เอกนามมฺปิ นานานามปฺปิ สภาคมฺปิ วิสภาคมฺปิ ววตฺถิตมฺปิ สมฺภินฺนมฺปิ. โส สํฆํ ตาสํ อาปตฺตีนํ สโมธาน\-ปริวาสํ ยาจติ, ตสฺส สํโฆ ตาสํ อาปตฺตีนํ สโมธาน\-ปริวาสํ เทติ, โส ปริวสนฺโต อนฺตรา สมฺพหุลา สํฆาทิเสสา อาปตฺติโย อาปชฺชติ ปริมาณาโย อปฺปฏิจฺฉนฺนาโย. โส สํฆํ อนฺตรา\-อาปตฺตีนํ มูลาย\-ปฏิกสฺสนํ ยาจติ, ตํ สํโฆ อนฺตรา\-อาปตฺตีนํ มูลาย \csromanfolio{175}ปฏิกสฺสติ ธมฺมิเกน กมฺเมน อกุปฺเปน ฐานารเหน, ธมฺเมน สโมธาน\-ปริวาสํ เทติ, อธมฺเมน มานตฺตํ เทติ, อธมฺเมน อพฺเภติ. โส ภิกฺขเว ภิกฺขุ อวิสุทฺโธ ตาหิ อาปตฺตีหิ. (1)}

\prose{อิธ ปน ภิกฺขเว ภิกฺขุ สมฺพหุลา สํฆาทิเสสา อาปตฺติโย อาปชฺชติ ปริมาณมฺปิ อปริมาณมฺปิ เอกนามมฺปิ นานานามมฺปิ สภาคมฺปิ วิสภาคมฺปิ ววตฺถิตมฺปิ สมฺภินฺนมฺปิ. โส สํฆํ ตาสํ อาปตฺตีนํ สโมธาน\-ปริวาสํ ยาจติ, ตสฺส สํโฆ ตาสํ อาปตฺตีนํ สโมธาน\-ปริวาสํ เทติ, โส ปริวสนฺโต อนฺตรา สมฺพหุลา สํฆาทิเสสา อาปตฺติโย อาปชฺชติ ปริมาณาโย ปฏิจฺฉนฺนาโย. โส สํฆํ อนฺตรา\-อาปตฺตีนํ มูลาย\-ปฏิกสฺสนํ ยาจติ, ตํ สํโฆ อนฺตรา\-อาปตฺตีนํ มูลาย\-ปฏิกสฺสติ ธมฺมิเกน กมฺเมน อกุปฺเปน ฐานารเหน, ธมฺเมน สโมธาน\-ปริวาสํ เทติ, อธมฺเมน มานตฺตํ เทติ, อธมฺเมน อพฺเภติ. โส ภิกฺขเว ภิกฺขุ อวิสุทฺโธ ตาหิ อาปตฺตีหิ. (2)}

\prose{อิธ ปน ภิกฺขเว ภิกฺขุ สมฺพหุลา สํฆาทิเสสา อาปตฺติโย อาปชฺชติ ปริมาณมฺปิ อปริมาณมฺปิ เอกนามมฺปิ นานานามมฺปิ สภาคมฺปิ วิสภาคมฺปิ ววตฺถิตมฺปิ สมฺภินฺนมฺปิ. โส สํฆํ ตาสํ อาปตฺตีนํ สโมธาน\-ปริวาสํ ยาจติ, ตสฺส สํโฆ ตาสํ อาปตฺตีนํ สโมธาน\-ปริวาสํ เทติ, โส ปริวสนฺโต อนฺตรา สมฺพหุลา สํฆาทิเสสา อาปตฺติโย อาปชฺชติ ปริมาณาโย ปฏิจฺฉนฺนาโยปิ อปฺปฏิจฺฉนฺนาโยปิ. โส สํฆํ อนฺตรา\-อาปตฺตีนํ มูลาย\-ปฏิกสฺสนํ ยาจติ, ตํ สํโฆ อนฺตรา\-อาปตฺตีนํ มูลาย ปฏิกสฺสติ ธมฺมิเกน กมฺเมน อกุปฺเปน ฐานารเหน, ธมฺเมน สโมธาน\-ปริวาสํ เทติ, อธมฺเมน มานตฺตํ เทติ, อธมฺเมน อพฺเภติ. โส ภิกฺขเว ภิกฺขุ อวิสุทฺโธ ตาหิ อาปตฺตีหิ. (3)}

\prose{อิธ ปน ภิกฺขเว ภิกฺขุ สมฺพหุลา สํฆาทิเสสา อาปตฺติโย อาปชฺชติ ปริมาณมฺปิ อปริมาณมฺปิ เอกนามมฺปิ นานานามมฺปิ สภาคมฺปิ วิสภาคมฺปิ ววตฺถิตมฺปิ สมฺภินฺนมฺปิ. โส สํฆํ ตาสํ อาปตฺตีนํ สโมธาน\-ปริวาสํ ยาจติ, ตสฺส สํโฆ ตาสํ อาปตฺตีนํ สโมธาน\-ปริวาสํ เทติ, โส ปริวสนฺโต อนฺตรา สมฺพหุลา สํฆาทิเสสา อาปตฺติโย อาปชฺชติ อปริมาณาโย อปฺปฏิจฺฉนฺนาโย ฯเปฯ อปริมาณาโย ปฏิจฺฉนฺนาโย ฯเปฯ อปริมาณาโย ปฏิจฺฉนฺนาโยปิ \csromanfolio{176}อปฺปฏิจฺฉนฺนาโยปิ ฯเปฯ ปริมาณาโยปิ อปริมาณาโยปิ อปฺปฏิจฺฉนฺนาโย. โส สํฆํ อนฺตรา อาปตฺตีนํ มูลาย\-ปฏิกสฺสนํ ยาจติ, ตํ สํโฆ อนฺตรา\-อาปตฺตีนํ มูลาย ปฏิกสฺสติ ธมฺมิเกน กมฺเมน อกุปฺเปน ฐานารเหน, ธมฺเมน สโมธาน\-ปริวาสํ เทติ, อธมฺเมน มานตฺตํ เทติ, อธมฺเมน อพฺเภติ. โส ภิกฺขเว ภิกฺขุ อวิสุทฺโธ ตาหิ อาปตฺตีหิ. \mbox{(4-7)}}

\prose{อิธ ปน ภิกฺขเว ภิกฺขุ สมฺพหุลา สํฆาทิเสสา อาปตฺติโย อาปชฺชติ ปริมาณมฺปิ อปริมาณมฺปิ เอกนามมฺปิ นานานามมฺปิ สภาคมฺปิ วิสภาคมฺปิ ววตฺถิตมฺปิ สมฺภินฺนมฺปิ. โส สํฆํ ตาสํ อาปตฺตีนํ สโมธาน\-ปริวาสํ ยาจติ, ตสฺส สํโฆ ตาสํ อาปตฺตีนํ สโมธาน\-ปริวาสํ เทติ, โส ปริวสนฺโต อนฺตรา สมฺพหุลา สํฆาทิเสสา อาปตฺติโย อาปชฺชติ ปริมาณาโยปิ อปริมาณาโยปิ ปฏิจฺฉนฺนาโย. โส สํฆํ อนฺตรา\-อาปตฺตีนํ มูลาย\-ปฏิกสฺสนํ ยาจติ, ตํ สํโฆ อนฺตรา\-อาปตฺตีนํ มูลาย ปฏิกสฺสติ ธมฺมิเกน กมฺเมน อกุปฺเปน ฐานารเหน, ธมฺเมน สโมธาน\-ปริวาสํ เทติ, อธมฺเมน มานตฺตํ เทติ, อธมฺเมน อพฺเภติ. โส ภิกฺขเว ภิกฺขุ อวิสุทฺโธ ตาหิ อาปตฺตีหิ. (8)}

\prosewithcloser{อิธ ปน ภิกฺขเว ภิกฺขุ สมฺพหุลา สํฆาทิเสสา อาปตฺติโย อาปชฺชติ ปริมาณมฺปิ อปริมาณมฺปิ เอกนามมฺปิ นานานามมฺปิ สภาคมฺปิ วิสภาคมฺปิ ววตฺถิตมฺปิ สมฺภินฺนมฺปิ. โส สํฆํ ตาสํ อาปตฺตีนํ สโมธาน\-ปริวาสํ ยาจติ, ตสฺส สํโฆ ตาสํ อาปตฺตีนํ สโมธาน\-ปริวาสํ เทติ, โส ปริวสนฺโต อนฺตรา สมฺพหุลา สํฆาทิเสสา อาปตฺติโย อาปชฺชติ ปริมาณาโยปิ อปริมาณาโยปิ ปฏิจฺฉนฺนาโยปิ อปฺปฏิจฺฉนฺนาโยปิ. โส สํฆํ อนฺตรา\-อาปตฺตีนํ มูลาย\-ปฏิกสฺสนํ ยาจติ, ตํ สํโฆ อนฺตรา\-อาปตฺตีนํ มูลาย ปฏิกสฺสติ ธมฺมิเกน กมฺเมน อกุปฺเปน ฐานารเหน, ธมฺเมน สโมธาน\-ปริวาสํ เทติ, อธมฺเมน มานตฺตํ เทติ, อธมฺเมน อพฺเภติ. โส ภิกฺขเว ภิกฺขุ อวิสุทฺโธ ตาหิ อาปตฺตีหิ. (9)}{มูลาย\-อวิสุทฺธิ\-นวกํ\csromansharedfootnote{5d746c523b6cdbfc}{สมูลาวิสุทฺธินวกํ (ก)} นิฏฺฐิตํ.}
\csromanmatikaanchor{mtk.94}
\csromantocmarkref{section}{10. ทุติยนวก}{mtk.94}
\bookmark[dest={mtk.94},level=1]{10. ทุติยนวก}
\csromanheader{1}{10. \textbf{ทุติยนวก}\csromansharedfootnote{82f43f4d1e96638a}{อิทํ นวกํ โปราณโปตฺถเกสุ อวิสุกฺขิวเสเนว อาคตํ, วุจฺจมานตติยนวเกน จ สํสฏฺฐํ. ตํ ปฏิวิโสธเกหิ อสํสฏฺฐํ กตฺวา วิสุํ ปติฏฺฐาปิตํ. สีหฬสฺยามโปตฺถเกสุ ปน ตํ วิสุทฺธิวเสเนว อาคตํ. ตํ ปเนวํ เวทิตพฺพํ\csromandash{}}}
\csromanmatikaanchor{mtk.95}
\csromantocmarkref{section}{10. มูลายวิสุทฺธินวก (สี, สฺยา)}{mtk.95}
\bookmark[dest={mtk.95},level=1]{10. มูลายวิสุทฺธินวก (สี, สฺยา)}
\proseitem{183}{\csromanfolio{177}อิธ ปน ภิกฺขเว ภิกฺขุ สมฺพหุลา สํฆาทิเสสา อาปตฺติโย อาปชฺชติ ปริมาณมฺปิ อปริมาณมฺปิ ฯเปฯ ววตฺถิตมฺปิ สมฺภินฺนมฺปิ. โส สํฆํ ตาสํ อาปตฺตีนํ สโมธาน\-ปริวาสํ ยาจติ, ตสฺส สํโฆ ตาสํ อาปตฺตีนํ สโมธาน\-ปริวาสํ เทติ, โส ปริวสนฺโต อนฺตรา สมฺพหุลา สํฆาทิเสสา อาปตฺติโย อาปชฺชติ ปริมาณา อปฺปฏิจฺฉนฺนาโย. โส สํฆํ อนฺตรา\-อาปตฺตีนํ มูลาย\-ปฏิกสฺสนํ ยาจติ, ตํ สํโฆ อนฺตรา\-อาปตฺตีนํ มูลาย ปฏิกสฺสติ อธมฺมิเกน กมฺเมน กุปฺเปน อฏฺฐานารเหน, อธมฺเมน สโมธาน\-ปริวาสํ เทติ, ธมฺเมน มานตฺตํ เทติ, ธมฺเมน อพฺเภติ. โส ภิกฺขเว ภิกฺขุ อวิสุทฺโธ ตาหิ อาปตฺตีหิ. (1)}

\csromanmatikaanchor{mtk.96}
\csromantocmarkref{section}{10. มูลายวิสุทฺธินวก (สี, สฺยา)}{mtk.96}
\bookmark[dest={mtk.96},level=1]{10. มูลายวิสุทฺธินวก (สี, สฺยา)}
\prose{อิธ ปน ภิกฺขเว ภิกฺขุ สมฺพหุลา สํฆาทิเสสา อาปตฺติโย อาปชฺชติ ปริมาณมฺปิ อปริมาณมฺปิ ฯเปฯ ววตฺถิตมฺปิ สมฺภินฺนมฺปิ. โส สํฆํ ตาสํ อาปตฺตีนํ สโมธาน\-ปริวาสํ ยาจติ, ตสฺส สํโฆ ตาสํ อาปตฺตีนํ สโมธาน\-ปริวาสํ เทติ, โส ปริวสนฺโต อนฺตรา สมฺพหุลา สํฆาทิเสสา อาปตฺติโย อาปชฺชติ ปริมาณา ปฏิจฺฉนฺนาโย. โส สํฆํ อนฺตรา\-อาปตฺตีนํ มูลาย\-ปฏิกสฺสนํ ยาจติ, ตํ สํโฆ อนฺตรา\-อาปตฺตีนํ มูลาย ปฏิกสฺสติ อธมฺมิเกน กมฺเมน กุปฺเปน \csromanfolio{178}อฏฺฐานารเหน, อธมฺเมน สโมธาน\-ปริวาสํ เทติ, ธมฺเมน มานตฺตํ เทติ, ธมฺเมน อพฺเภติ. โส ภิกฺขเว ภิกฺขุ อวิสุทฺโธ ตาหิ อาปตฺตีหิ. (2)}

\prose{อิธ ปน ภิกฺขเว ภิกฺขุ สมฺพหุลา สํฆาทิเสสา อาปตฺติโย อาปชฺชติ ปริมาณมฺปิ อาปริมาณมฺปิ ฯเปฯ ววตฺถิตมฺปิ สมฺภินฺนมฺปิ. โส สํฆํ ตาสํ อาปตฺตีนํ สโมธาน\-ปริวาสํ ยาจติ, ตสฺส สํโฆ ตาสํ อาปตฺตีนํ สโมธาน\-ปริวาสํ เทติ, โส ปริวสนฺโต อนฺตรา สมฺพหุลา สํฆาทิเสสา อาปตฺติโย อาปชฺชติ ปริมาณา ปฏิจฺฉนฺนาโยปิ อปฺปฏิจฺฉนฺนาโยปิ. โส สํฆํ อนฺตรา\-อาปตฺตีนํ มูลาย\-ปฏิกสฺสนํ ยาจติ, ตํ สํโฆ อนฺตรา\-อาปตฺตีนํ มูลาย ปฏิกสฺสติ อธมฺมิเกน กมฺเมน กุปฺเปน อฏฺฐานารเหน, อธมฺเมน สโมธาน\-ปริวาสํ เทติ, ธมฺเมน มานตฺตํ เทติ, ธมฺเมน อพฺเภติ. โส ภิกฺขเว ภิกฺขุ อวิสุทฺโธ ตาหิ อาปตฺตีหิ. (3)}

\csromanmatikaanchor{mtk.97}
\csromantocmarkref{section}{10. มูลายวิสุทฺธินวก (สี, สฺยา)}{mtk.97}
\bookmark[dest={mtk.97},level=1]{10. มูลายวิสุทฺธินวก (สี, สฺยา)}
\prose{อิธ ปน ภิกฺขเว ภิกฺขุ สมฺพหุลา สํฆาทิเสสา อาปตฺติโย อาปชฺชติ ปริมาณมฺปิ อปริมาณมฺปิ เอกนามมฺปิ นานานามมฺปิ สภาคมฺปิ วิสภาคมฺปิ ววตฺถิตมฺปิ สมฺภินฺนมฺปิ. โส สํฆํ ตาสํ อาปตฺตีนํ สโมธาน\-ปริวาสํ ยาจติ, ตสฺส สํโฆ ตาสํ อาปตฺตีนํ สโมธาน\-ปริวาสํ \csromanfolio{179}เทติ, โส ปริวสนฺโต อนฺตรา สมฺพหุลา สํฆาทิเสสา อาปตฺติโย อาปชฺชติ อปริมาณา อปฺปฏิจฺฉนฺนาโย ฯเปฯ อปริมาณา ปฏิจฺฉนฺนาโย ฯเปฯ อปริมาณา ปฏิจฺฉนฺนาโยปิ อปฺปฏิจฺฉนฺนาโยปิ ฯเปฯ ปริมาณาโยปิ อปริมาณาโยปิ อปฺปฏิจฺฉนฺนาโย. โส สํฆํ อนฺตรา อาปตฺตีนํ มูลาย\-ปฏิกสฺสนํ ยาจติ, ตํ สํโฆ อนฺตรา\-อาปตฺตีนํ มูลาย ปฏิกสฺสติ อธมฺมิเกน กมฺเมน กุปฺเปน อฏฺฐานารเหน, อธมฺเมน สโมธาน\-ปริวาสํ เทติ, ธมฺเมน มานตฺตํ เทติ, ธมฺเมน อพฺเภติ. โส ภิกฺขเว ภิกฺขุ อวิสุทฺโธ ตาหิ อาปตฺตีหิ. \mbox{(4-7)}}

\prose{อิธ ปน ภิกฺขเว ภิกฺขุ สมฺพหุลา สํฆาทิเสสา อาปตฺติโย อาปชฺชติ ปริมาณมฺปิ อาปริมาณมฺปิ ฯเปฯ ววตฺถิตมฺปิ สมฺภินฺนมฺปิ. โส สํฆํ ตาสํ อาปตฺตีนํ สโมธาน\-ปริวาสํ ยาจติ, ตสฺส สํโฆ ตาสํ อาปตฺตีนํ สโมธาน\-ปริวาสํ เทติ, โส ปริวสนฺโต อนฺตรา สมฺพหุลา สํฆาทิเสสา อาปตฺติโย อาปชฺชติ ปริมาณาโยปิ อปริมาณาโยปิ ปฏิจฺฉนฺนาโย. โส สํฆํ อนฺตรา\-อาปตฺตีนํ มูลาย\-ปฏิกสฺสนํ ยาจติ, ตํ สํโฆ อนฺตรา\-อาปตฺตีนํ มูลาย ปฏิกสฺสติ อธมฺมิเกน กมฺเมน กุปฺเปน อฏฺฐานารเหน, อธมฺเมน สโมธาน\-ปริวาสํ เทติ, ธมฺเมน มานตฺตํ เทติ, ธมฺเมน อพฺเภติ. โส ภิกฺขเว ภิกฺขุ อวิสุทฺโธ ตาหิ อาปตฺตีหิ. (8)}

\csromanmatikaanchor{mtk.98}
\csromanmatikaanchor{mtk.99}
\csromantocmarkref{section}{11. ตติยนวก}{mtk.98}
\bookmark[dest={mtk.98},level=1]{11. ตติยนวก}
\csromantocmarkref{section}{10. มูลายวิสุทฺธินวก (สี, สฺยา)}{mtk.99}
\bookmark[dest={mtk.99},level=1]{10. มูลายวิสุทฺธินวก (สี, สฺยา)}
\prosewithcloser{\csromanfolio{180}อิธ ปน ภิกฺขเว ภิกฺขุ สมฺพหุลา สํฆาทิเสสา อาปตฺติโย อาปชฺชติ ปริมาณมฺปิ อปริมาณมฺปิ ฯเปฯ ววตฺถิตมฺปิ สมฺภินฺนมฺปิ. โส สํฆํ ตาสํ อาปตฺตีนํ สโมธาน\-ปริวาสํ ยาจติ, ตสฺส สํโฆ ตาสํ อาปตฺตีนํ สโมธาน\-ปริวาสํ เทติ, โส ปริวสนฺโต อนฺตรา สมฺพหุลา สํฆาทิเสสา อาปตฺติโย อาปชฺชติ ปริมาณาโยปิ อปริมาณาโยปิ ปฏิจฺฉนฺนาโยปิ อปฺปฏิจฺฉนฺนาโยปิ. โส สํฆํ อนฺตรา\-อาปตฺตีนํ มูฏิกสฺสนํ ยาจติ, ตํ สํโฆ อนฺตรา\-อาปตฺตีนํ มูลาย ปฏิกสฺสติ อธมฺมิเกน กมฺเมน กุปฺเปน อฏฺฐานารเหน, อธมฺเมน สโมธาน\-ปริวาสํ เทติ, ธมฺเมน มานตฺตํ เทติ, ธมฺเมน อพฺเภติ. โส ภิกฺขเว ภิกฺขุ อวิสุทฺโธ ตาหิ อาปตฺตีหิ. (9)}{ทุติยนวกํ นิฏฺฐิตํ.}
\csromanheader{1}{11. \textbf{ตติยนวก}}
\proseitem{184}{อิธ ปน ภิกฺขเว ภิกฺขุ สมฺพหุลา สํฆาทิเสสา อาปตฺติโย อาปชฺชติ ปริมาณมฺปิ อปริมาณมฺปิ ฯเปฯ ววตฺถิตมฺปิ สมฺภินฺนมฺปิ. โส สํฆํ ตาสํ อาปตฺตีนํ สโมธาน\-ปริวาสํ ยาจติ, ตสฺส สํโฆ ตาสํ \csromanfolio{181}อาปตฺตีนํ สโมธาน\-ปริวาสํ เทติ, โส ปริวสนฺโต อนฺตรา สมฺพหุลา สํฆาทิเสสา อาปตฺติโย อาปชฺชติ ปริมาณาโย อปฺปฏิจฺฉนฺนาโย. โส สํฆํ อนฺตรา\-อาปตฺตีนํ มูลาย\-ปฏิกสฺสนํ ยาจติ, ตํ สํโฆ อนฺตรา\-อาปตฺตีนํ มูลาย ปฏิกสฺสติ อธมฺมิเกน กมฺเมน กุปฺเปน อฏฺฐานารเหน, อธมฺเมน สโมธาน\-ปริวาสํ เทติ, โส “ปริวสามี”ติ มญฺญมาโน อนฺตรา สมฺพหุลา สํฆาทิเสสา อาปตฺติโย อาปชฺชติ ปริมาณาโย อปฺปฏิจฺฉนฺนาโย. โส ตสฺมิํ ภูมิยํ ฐิโต ปุริมา\-อาปตฺตีนํ อนฺตราอาปตฺติโย สรติ, อปราอาปตฺตีนํ อนฺตราอาปตฺติโย สรติ, ตสฺส เอวํ โหติ “อหํ โข สมฺพหุลา สํฆาทิเสสา อาปตฺติโย อาปชฺชิํ ปริมาณมฺปิ อปริมาณมฺปิ ฯเปฯ ววตฺถิตมฺปิ สมฺภินฺนมฺปิ. โสหํ สํฆํ ตาสํ อาปตฺตีนํ สโมธาน\-ปริวาสํ ยาจิํ, ตสฺส เม สํโฆ ตาสํ อาปตฺตีนํ สโมธาน\-ปริวาสํ อทาสิ, โสหํ ปริวสนฺโต อนฺตรา สมฺพหุลา สํฆาทิเสสา อาปตฺติโย อาปชฺชิํ ปริมาณาโย อปฺปฏิจฺฉนฺนาโย. โสหํ สํฆํ อนฺตรา\-อาปตฺตีนํ มูลาย\-ปฏิกสฺสนํ ยาจิํ, ตํ มํ สํโฆ อนฺตรา\-อาปตฺตีนํ มูลาย ปฏิกสฺสิ อธมฺมิเกน กมฺเมน กุปฺเปน อฏฺฐานารเหน, อธมฺเมน สโมธาน\-ปริวาสํ อทาสิ, โสหํ ‘ปริวสามี’ติ มญฺญมาโน อนฺตรา สมฺพหุลา สํฆาทิเสสา อาปตฺติโย อาปชฺชิํ ปริมาณาโย อปฺปฏิจฺฉนฺนาโย. โสหํ ตสฺมิํ ภูมิยํ ฐิโต ปุริมา\-อาปตฺตีนํ อนฺตราอาปตฺติโย สรามิ, อปราอาปตฺตีนํ อนฺตราอาปตฺติโย สรามิ, ยํนูนาหํ สํฆํ ปุริมา\-อาปตฺตีนํ อนฺตราอาปตฺตีนญฺจ อปราอาปตฺตีนํ อนฺตราอาปตฺตีนญฺจ มูลาย\-ปฏิกสฺสนํ ยาเจยฺยํ ธมฺมิเกน กมฺเมน อกุปฺเปน ฐานารเหน, ธมฺเมน สโมธาน\-ปริวาสํ, ธมฺเมน มานตฺตํ, ธมฺเมน อพฺภานนฺ”ติ. โส สํฆํ ปุริมา\-อาปตฺตีนํ อนฺตราอาปตฺตีนญฺจ อปราอาปตฺตีนํ อนฺตราอาปตฺตีนญฺจ มูลาย\-ปฏิกสฺสนํ ยาจติ ธมฺมิเกน กมฺเมน อกุปฺเปน ฐานารเหน, ธมฺเมน สโมธาน\-ปริวาสํ, ธมฺเมน มานตฺตํ, ธมฺเมน อพฺภานํ. ตํ สํโฆ ปุริมา\-อาปตฺตีนํ อนฺตราอาปตฺตีนญฺจ อปราอาปตฺตีนํ อนฺตราอาปตฺตีนญฺจ มูลาย ปฏิกสฺสติ ธมฺมิเกน กมฺเมน อกุปฺเปน ฐานารเหน, ธมฺเมน สโมธาน\-ปริวาสํ เทติ, ธมฺเมน มานตฺตํ เทติ, ธมฺเมน อพฺเภติ. โส ภิกฺขเว ภิกฺขุ วิสุทฺโธ ตาหิ อาปตฺตีหิ\csromansharedfootnote{0008e706890e60ca}{อยํ ปฐมวาโร สีสฺยาโปตฺถเกสุ ปริปุณฺโณ ทิสฺสติ.}. (1)}

\prose{\csromanfolio{182}อิธ ปน ภิกฺขเว ภิกฺขุ สมฺพหุลา สํฆาทิเสสา อาปตฺติโย อาปชฺชติ ปริมาณมฺปิ อปริมาณมฺปิ ฯเปฯ ววตฺถิตมฺปิ สมฺภินฺนมฺปิ. โส สํฆํ ตาสํ อาปตฺตีนํ สโมธาน\-ปริวาสํ ยาจติ, ตสฺส สํโฆ ตาสํ อาปตฺตีนํ สโมธาน\-ปริวาสํ เทติ, โส ปริวสนฺโต อนฺตรา สมฺพหุลา สํฆาทิเสสา อาปตฺติโย อาปชฺชติ ปริมาณา ปฏิจฺฉนฺนาโย. โส สํฆํ อนฺตรา\-อาปตฺตีนํ มูลาย\-ปฏิกสฺสนํ ยาจติ, ตํ สํโฆ อนฺตรา\-อาปตฺตีนํ มูลาย ปฏิกสฺสติ อธมฺมิเกน กมฺเมน กุปฺเปน อฏฺฐานารเหน, อธมฺเมน สโมธาน\-ปริวาสํ เทติ, โส “ปริวสฺสามี”ติ มญฺญมาโน อนฺตรา สมฺพหุลา สํฆาทิเสสา อาปตฺติโย อาปชฺชติ ปริมาณา ปฏิจฺฉนฺนาโย. โส ตสฺมิํ ภูมิยํ ฐิโต ปุริมา\-อาปตฺตีนํ\csromansharedfootnote{18b7b45249747ab1}{ปุริมานํ อาปตฺตีนํ (ก)} อนฺตราอาปตฺติโย สรติ, อปราอาปตฺตีนํ อนฺตราอาปตฺติโย สรติ, ตสฺส เอวํ โหติ “อหํ โข สมฺพหุลา สํฆาทิเสสา อาปตฺติโย อาปชฺชิํ ปริมาณมฺปิ อปริมาณมฺปิ ฯเปฯ ววตฺถิตมฺปิ สมฺภินฺนมฺปิ. โสหํ สํฆํ ตาสํ อาปตฺตีนํ สโมธาน\-ปริวาสํ ยาจิํ, ตสฺส เม สํโฆ ตาสํ อาปตฺตีนํ สโมธาน\-ปริวาสํ อทาสิ, โสหํ ปริวสนฺโต อนฺตรา สมฺพหุลา สํฆาทิเสสา อาปตฺติโย อาปชฺชิํ ปริมาณา ปฏิจฺฉนฺนาโย. โสหํ สํฆํ อนฺตรา\-อาปตฺตีนํ มูลยปฏิกสฺสนํ ยาจิํ, ตํ มํ สํโฆ อนฺตรา\-อาปตฺตีนํ มูลาย ปฏิกสฺสิ อธมฺมิเกน กมฺเมน กุปฺเปน อฏฺฐานารเหน, อธมฺเมน สโมธาน\-ปริวาสํ อทาสิ, โสหํ ‘ปริวสามี’ติ มญฺญมาโน อนฺตรา สมฺพหุลา สํฆาทิเสสา อาปตฺติโย อาปชฺชิํ ปริมาณา ปฏิจฺฉนฺนาโย โสหํ ตสฺมิํ ภูมิยํ ฐิโต ปุริมา\-อาปตฺตีนํ อนฺตราอาปตฺติโย สรามิ, อปรา อาปตฺตีนํ อนฺตราอาปตฺติโย สรามิ, ยํนูนาหํ สํฆํ ปุริมา\-อาปตฺตีนํ อนฺตรา อาปตฺตีนญฺจ อปรา อาปตฺตีนํ อนฺตรา อาปตฺตีนญฺจ มูลาย\-ปฏิกสฺสนํ ยาเจยฺยํ ธมฺมิเกน กมฺเมน อกุปฺเปน ฐานารเหน, ธมฺเมน สโมธาน\-ปริวาสํ, ธมฺเมน มานตฺตํ, ธมฺเมน อพฺภานนฺ”ติ. โส สํฆํ ปุริมา\-อาปตฺตีนํ อนฺตราอาปตฺตีนญฺจ อปราอาปตฺตีนํ อนฺตราอาปตฺตีนญฺจ มูลาย\-ปฏิกสฺสนํ ยาจติ ธมฺมิเกน กมฺเมน อกุปฺเปน ฐานารเหน, ธมฺเมน สโมธาน\-ปริวาสํ, ธมฺเมน มานตฺตํ, ธมฺเมน อพฺภานํ. ตํ สํโฆ ปุริมา\-อาปตฺตีนํ อนฺตราอาปตฺตีนญฺจ อปราอาปตฺตีนํ อนฺตราอาปตฺตีนญฺจ มูลาย ปฏิกสฺสติ \csromanfolio{183}ธมฺมิเกน กมฺเมน อกุปฺเปน ฐานารเหน, ธมฺเมน สโมธาน\-ปริวาสํ เทติ, ธมฺเมน มานตฺตํ เทติ, ธมฺเมน อพฺเภติ. โส ภิกฺขเว ภิกฺขุ วิสุทฺโธ ตาหิ อาปตฺตีหิ. (2)}

\prose{อิธ ปน ภิกฺขเว ภิกฺขุ สมฺพหุลา สํฆาทิเสสา อาปตฺติโย อาปชฺชติ ปริมาณมฺปิ อปริมาณมฺปิ ฯเปฯ ววตฺถิตมฺปิ สมฺภินฺนมฺปิ. โส สํฆํ ตาสํ อาปตฺตีนํ สโมธาน\-ปริวาสํ ยาจติ, ตสฺส สํโฆ ตาสํ อาปตฺตีนํ สโมธาน\-ปริวาสํ เทติ, โส ปริวสนฺโต อนฺตรา สมฺพหุลา สํฆาทิเสสา อาปตฺติโย อาปชฺชติ ปริมาณาปฏิจฺฉนฺนาโยปิ อปฺปฏิจฺฉนฺนาโยปิ. โส สํฆํ อนฺตรา\-อาปตฺตีนํ มูลาย\-ปฏิกสฺสนํ ยาจติ, ตํ สํโฆ อนฺตรา\-อาปตฺตีนํ มูลาย ปฏิกสฺสติ อธมฺมิเกน กมฺเมน กุปฺเปน อฏฺฐานารเหน, อธมฺเมน สโมธาน\-ปริวาสํ เทติ, โส “ปริวสามี”ติ มญฺญมาโน อนฺตรา สมฺพหุลา สํฆาทิเสสา อาปตฺติโย อาปชฺชติ ปริมาณา ปฏิจฺฉนฺนาโยปิ อปฺปฏิจฺฉนาโยปิ. โส ตสฺมิํ ภูมิยํ ฐิโต ปุริมา\-อาปตฺตีนํ อนฺตราอาปตฺติโย สรติ, อปราอาปตฺตีนํ อนฺตราอาปตฺติโย สรติ, ตสฺส เอวํ โหติ “อหํ โข สมฺพหุลา สํฆาทิเสสา อาปตฺติโย อาปชฺชิํ ปริมาณมฺปิ อปริมาณมฺปิ ฯเปฯ ววตฺถิตมฺปิ สมฺภินฺนมฺปิ. โสหํ สํฆํ ตาสํ อาปตฺตีนํ สโมธาน\-ปริวาสํ ยาจิํ, ตสฺส เม สํโฆ ตาสํ อาปตฺตีนํ สโมธาน\-ปริวาสํ อทาสิ, โสหํ ปริวสนฺโต อนฺตรา สมฺพหุลา สํฆาทิเสสา อาปตฺติโย อาปชฺชิํ ปริมาณา ปฏิจฺฉนฺนาโยปิ อปฺปฏิจฺฉนฺนาโยปิ. โสหํ สํฆํ อนฺตรา\-อาปตฺตีนํ มูลาย\-ปฏิกสฺสนํ ยาจิํ, ตํ มํ สํโฆ อนฺตรา\-อาปตฺตีนํ มูลาย ปฏิกสฺสิ อธมฺมิเกน กมฺเมน กุปฺเปน อฏฺฐานารเหน, อธมฺเมน สโมธาน\-ปริวาสํ อทาสิ, โสหํ ‘ปริวสามี’ติ มญฺญมาโน อนฺตรา สมฺพหุลา สํฆาทิเสสา อาปตฺติโย อาปชฺชิํ ปริมาณา ปฏิจฺฉนฺนาโยปิ อปฺปฏิจฺฉนฺนาโยปิ. โสหํ ตสฺมิํ ภูมิยํ ฐิโต ปุริมา\-อาปตฺตีนํ อนฺตราอาปตฺติโย สรามิ, อปราอาปตฺตีนํ อนฺตราอาปตฺติโย สรามิ, ยํนูนาหํ สํฆํ ปุริมา อาปตฺตีนํ อนฺตรา อาปตฺตีนญฺจ อปรา อาปตฺตีนํ อนฺตรา อาปตฺตีนญฺจ มูลาย\-ปฏิกสฺสนํ ยาเจยฺยํ ธมฺมิเกน กมฺเมน อกุปฺเปน ฐานารเหน, ธมฺเมน สโมธาน\-ปริวาสํ, ธมฺเมน มานตฺตํ, ธมฺเมน อพฺภานนฺ”ติ. โส สํฆํ ปุริมา\-อาปตฺตีนํ อนฺตรา อาปตฺตีนญฺจ อปรา อาปตฺตีนํ อนฺตรา อาปตฺตินญฺจ \csromanfolio{184}มูลาย\-ปฏิกสฺสนํ ยาจติ, ธมฺมิเกน กมฺเมน อกุปฺเปน ฐานารเหน, ธมฺเมน สโมธาน\-ปริวาสํ, ธมฺเมน มานตฺตํ, ธมฺเมน อพฺภานํ. ตํ สํโฆ ปุริมา\-อาปตฺตีนํ อนฺตรา อาปตฺตีนญฺจ อปรา อาปตฺตีนํ อนฺตรา อาปตฺตีนญฺจ มูลาย ปฏิกสฺสติ ธมฺมิเกน กมฺเมน อกุปฺเปน ฐานารเหน, ธมฺเมน สโมธาน\-ปริวาสํ เทติ, ธมฺเมน มานตฺตํ เทติ, ธมฺเมน อพฺเภติ. โส ภิกฺขเว ภิกฺขุ วิสุทฺโธ ตาหิ อาปตฺตีหิ. (3)}

\prose{อิธ ปน ภิกฺขเว ภิกฺขุ สมฺพหุลา สํฆาทิเสสา อาปตฺติโย อาปชฺชติ ปริมาณมฺปิ อปริมาณมฺปิ ฯเปฯ ววตฺถิตมฺปิ สมฺภินฺนมฺปิ. โส สํฆํ ตาสํ อาปตฺตีนํ สโมธาน\-ปริวาสํ ยาจติ, ตสฺส สํโฆ ตาสํ อาปตฺตีนํ สโมธาน\-ปริวาสํ เทติ, โส ปริวสนฺโต อนฺตรา สมฺพหุลา สํฆาทิเสสา อาปตฺติโย อาปชฺชติ อปริมาณาโย อปฺปฏิจฺฉนฺนาโย ฯเปฯ อปริมาณาโย ปฏิจฺฉนฺนาโย ฯเปฯ อปริมาณาโย ปฏิจฺฉนาโยปิ อปฺปฏิจฺฉนฺนาโยปิ ฯเปฯ ปริมาณาโยปิ อปริมาณาโยปิ อปฺปฏิจฺฉนฺนาโย. โส สํฆํ อนฺตรา- อาปตฺตีนํ มูลาย\-ปฏิกสฺสนํ ยาจติ. ตํ สํโฆ อนฺตรา\-อาปตฺตีนํ มูลาย ปฏิกสฺสิ อธมฺมิเกน กมฺเมน กุปฺเปน อฏฺฐานารเหน, อธมฺเมน สโมธาน\-ปริวาสํ เทติ, โส “ปริวสามี”ติ มญฺญมาโน ฯเปฯ. ตํ สํโฆ ปุริมา\-อาปตฺตีนํ อนฺตราอาปตฺตีนญฺจ อปราอาปตฺตีนํ อนฺตราอาปตฺตีนญฺจ มูลาย ปฏิกสฺสติ ธมฺมิเกน กมฺเมน อกุปฺเปน ฐานารเหน, ธมฺเมน สโมธาน\-ปริวาสํ เทติ, ธมฺเมน มานตฺตํ เทติ, ธมฺเมน อพฺเภติ. โส ภิกฺขเว ภิกฺขุ วิสุทฺโธ ตาหิ อาปตฺตีหิ\csromansharedfootnote{6735ce11cb8185d2}{อิเมปิ จตฺตาโร วารา สีสฺยาโปตฺถเกสุ ปริปุณฺณา ทิสฺสนฺติ.}. \mbox{(4-7)}}

\prose{อิธ ปน ภิกฺขเว ภิกฺขุ สมฺพหุลา สํฆาทิเสสา อาปตฺติโย อาปชฺชติ ปริมาณมฺปิ อปริมาณมฺปิ ฯเปฯ ววตฺถิตมฺปิ สมฺภินฺนมฺปิ. โส สํฆํ ตาสํ อาปตฺตีนํ สโมธาน\-ปริวาสํ ยาจติ, ตสฺส สํโฆ ตาสํ อาปตฺตีนํ สโมธาน\-ปริวาสํ เทติ, โส ปริวสนฺโต อนฺตรา สมฺพหุลา สํฆาทิเสสา อาปตฺติโย อาปชฺชติ ปริมาณาโยปิ อปริมาณาโยปิ ปฏิจฺฉนฺนาโย. โส สํฆํ อนฺตรา\-อาปตฺตีนํ มูลาย\-ปฏิกสฺสนํ ยาจติ, ตํ สํโฆ อนฺตรา\-อาปตฺตีนํ มูลาย ปฏิกสฺสติ อธมฺมิเกน กมฺเมน กุปฺเปน อฏฺฐานารเหน, อธมฺเมน สโมธาน\-ปริวาสํ เทติ, โส “ปริวสามี”ติ มญฺญมาโน อนฺตรา \csromanfolio{185}สํฆาทิเสสา อาปตฺติโย อาปชฺชติ ปริมาณาโยปิ อปริมาณาโยปิ ปฏิจฺฉนฺนาโย. โส ตสฺมิํ ภูมิยํ ฐิโต ปุริมา\-อาปตฺตีนํ อนฺตราอาปตฺติโย สรติ, อปรา อาปตฺตีนํ อนฺตรา อาปตฺติโย สรติ, ตสฺส เอวํ โหติ “อหํ โข สมฺพหุลา สํฆาทิเสสา อาปตฺติโย อาปชฺชิํ ปริมาณมฺปิ อปริมาณมฺปิ ฯเปฯ ววตฺถิตมฺปิ สมฺภินฺนมฺปิ. โสหํ สํฆํ ตาสํ อาปตฺตีนํ สโมธาน\-ปริวาสํ ยาจิํ, ตสฺส เม สํโฆ ตาสํ อาปตฺตีนํ สโมธาน\-ปริวาสํ อทาสิ, โสหํ ปริวสนฺโต อนฺตรา สมฺพหุลา สํฆาทิเสสา อาปตฺติโย อาปชฺชิํ ปริมาณาโยปิ อปริมาณาโยปิ ปฏิจฺฉนฺนาโย. โสหํ สํฆํ อนฺตรา\-อาปตฺตีนํ มูลาย ปฏิกสฺสนํ ยาจิํ, ตํ มํ สํโฆ อนฺตรา\-อาปตฺตีนํ มูลาย ปฏิกสฺสิ อธมฺมิเกน กมฺเมน กุปฺเปน อฏฺฐานารเหน, อธมฺเมน สโมธาน\-ปริวาสํ อทาสิ, โสหํ ‘ปริวสามี’ติ มญฺญมาโน อนฺตรา สมฺพหุลา สํฆาทิเสสา อาปตฺติโย อาปชฺชิํ ปริมาณาโยปิ อปริมาณาโยปิ ปฏิจฺฉนฺนาโย. โสหํ ตสฺมิํ ภูมิยํ ฐิโต ปุริมา\-อาปตฺตีนํ อนฺตรา อาปตฺติโย สรามิ, อปราปตฺตีนํ อนฺตราอาปตฺติโย สรามิ, ยํนูนาหํ สํฆํ ปุริมา\-อาปตฺตีนํ อนฺตราปตฺตีนญฺจ อปรา อาปตฺตีนํ อนฺตรา อาปตฺตีนญฺจ มูลาย ปฏิกสฺสนํ ยาเจยฺยํ ธมฺมิเกน กมฺเมน อกุปฺเปน ฐานารเหน, ธมฺเมน สโมธาน\-ปริวาสํ, ธมฺเมน มานตฺตํ, ธมฺเมน อพฺภานนฺ”ติ. โส สํฆํ ปุริมา\-อาปตฺตีนํ อนฺตราปตฺตีนญฺจ อปรา อาปตฺตีนํ อนฺตรา อาปตฺตีนญฺจ มูลาย\-ปฏิกสฺสนํ ยาจติ, ธมฺมิเกน กมฺเมน อกุปฺเปน ฐานารเหน, ธมฺเมน สโมธาน\-ปริวาสํ, ธมฺเมน มานตฺตํ, ธมฺเมน อพฺภานํ. ตํ สํโฆ ปุริมา\-อาปตฺตีนํ อนฺตรา อาปตฺตีนญฺจ อปรา อาปตฺตีนํ อนฺตรา อาปตฺตีนญฺจ มูลาย ปฏิกสฺสติ ธมฺมิเกน กมฺเมน อกุปฺเปน ฐานารเหน, ธมฺเมน สโมธาน\-ปริวาสํ เทติ, ธมฺเมน มานตฺตํ เทติ, ธมฺเมน อพฺเภติ. โส ภิกฺขเว ภิกฺขุ วิสุทฺโธ ตาหิ อาปตฺตีหิ. (8)}

\prosewithcloser{อิธ ปน ภิกฺขเว ภิกฺขุ สมฺพหุลา สํฆาทิเสสา อาปตฺติโย อาปชฺชติ ปริมาณมฺปิ อปริมาณมฺปิ ฯเปฯ ววตฺถิตมฺปิ สมฺภินฺนมฺปิ. โส สํฆํ ตาสํ อาปตฺตีนํ สโมธาน\-ปริวาสํ ยาจติ, ตสฺส สํโฆ ตาสํ อาปตฺตีนํ สโมธาน\-ปริวาสํ เทติ, โส ปริวสนฺโต อนฺตรา สมฺพหุลา สํฆาทิเสสา อาปตฺติโย อาปชฺชติ ปริมาณาโยปิ อปริมาณาโยปิ ปฏิจฺฉนฺนาโยปิ อปฺปฏิจฺฉนฺนาโยปิ. โส สํฆํ \csromanfolio{186}อนฺตรา\-อาปตฺตีนํ มูลาย\-ปฏิกสฺสนํ ยาจติ, ตํ สํโฆ อนฺตรา\-อาปตฺตีนํ มูลาย ปฏิกสฺสติ อธมฺมิเกน กมฺเมน กุปฺเปน อฏฺฐานารเหน, อธมฺเมน สโมธาน\-ปริวาสํ เทติ, โส “ปริวสามี”ติ มญฺญมาโน อนฺตรา สมฺพหุลา สํฆาทิเสสา อาปตฺติโย อาปชฺชติ ปริมาณาโยปิ อปริมาณาโยปิ ปฏิจฺฉนฺนาโยปิ อปฺปฏิจฺฉนฺนาโยปิ. โส ตสฺมิํ ภูมิยํ ฐิโต ปุริมา\-อาปตฺตีนํ อนฺตราอาปตฺติโย สรติ. อปรา อาปตฺตีนํ อนฺตราอาปตฺติโย สรติ, ตสฺส เอวํ โหติ “อหํ โข สมฺพหุลา สํฆาทิเสสา อาปตฺติโย อาปชฺชิํ ปริมาณมฺปิ อปริมาณมฺปิ ฯเปฯ ววตฺถิตมฺปิ สมฺภินฺนมฺปิ. โสหํ สํฆํ ตาสํ อาปตฺตีนํ สโมธาน\-ปริวาสํ ยาจิํ, ตสฺส เม สํโฆ ตาสํ อาปตฺตีนํ สโมธาน\-ปริวาสํ อทาสิ, โสหํ ปริวสนฺโต อนฺตรา สมฺพหุลา สํฆาทิเสสา อาปตฺติโย อาปชฺชิํ ปริมาณาโยปิ อปริมาณาโยปิ ปฏิจฺฉนฺนาโยปิ อปฺปฏิจฺฉนฺนาโยปิ. โสหํ สํฆํ อนฺตรา\-อาปตฺตีนํ มูลาย\-ปฏิกสฺสนํ ยาจิํ, ตํ มํ สํโฆ อนฺตรา\-อาปตฺตีนํ มูลาย ปฏิกสฺสิ อธมฺมิเกน กมฺเมน กุปฺเปน อฏฺฐานารเหน, อธมฺเมน สโมธาน\-ปริวาสํ อทาสิ, โสหํ ‘ปริวสามี’ติ มญฺญมาโน อนฺตรา สมฺพหุลา สํฆาทิเสสา อาปตฺติโย อาปชฺชิํ ปริมาณาโยปิ อปริมาณาโยปิ ปฏิจฺฉนฺนาโยปิ อปฺปฏิจฺฉนฺนาโยปิ. โสหํ ตสฺมิํ ภูมิยํ ฐิโต ปุริมา\-อาปตฺตีนํ อนฺตราอาปตฺติโย สรามิ, อปรา อาปตฺตีนํ อนฺตราอาปตฺติโย สรามิ, ยํนูนาหํ สํฆํ ปุริมา\-อาปตฺตีนํ อนฺตรา อาปตฺตีนญฺจ อปรา อาปตฺตีนํ อนฺตรา อาปตฺตีนญฺจ มูลาย\-ปฏิกสฺสนํ ยาเจยฺยํ ธมฺมิเกน กมฺเมน อกุปฺเปน ฐานารเหน, ธมฺเมน สโมธาน\-ปริวาสํ, ธมฺเมน มานตฺตํ, ธมฺเมน อพฺภานนฺ”ติ. โส สํฆํ ปุริมา\-อาปตฺตีนํ อนฺตรา อาปตฺตีนญฺจ อปรา อาปตฺตีนํ อนฺตรา อาปตฺตีนญฺจ มูลาย\-ปฏิกสฺสนํ ยาจติ ธมฺมิเกน กมฺเมน อกุปฺเปน ฐานารเหน, ธมฺเมน สโมธาน\-ปริวาสํ, ธมฺเมน มานตฺตํ, ธมฺเมน อพฺภานํ. ตํ สํโฆ ปุริมา\-อาปตฺตีนํ อนฺตรา อาปตฺตีนญฺจ อปรา อาปตฺตีนํ อนฺตรา อาปตฺตีนญฺจ มูลาย ปฏิกสฺสติ ธมฺมิเกน กมฺเมน อกุปฺเปน ฐานารเหน, ธมฺเมน สโมธาน\-ปริวาสํ เทติ, ธมฺเมน มานตฺตํ เทติ, ธมฺเมน อพฺเภติ. โส ภิกฺขเว ภิกฺขุ วิสุทฺโธ ตาหิ อาปตฺตีหิ. (9)}{ตติยนวกํ นิฏฺฐิตํ.}
\csromancenter{สมุจฺจย\-กฺขนฺธโก ตติโย.}
\csromancenter{\textbf{ตสฺสุทฺทานํ}}
\csromanmatikaanchor{mtk.100}
\csromantocmarkref{section}{อุทฺทานคาถา}{mtk.100}
\bookmark[dest={mtk.100},level=1]{อุทฺทานคาถา}
\setlength{\csromangathapagewidth}{0pt}
\setlength{\csromangathawakwidth}{0pt}
\csromangathameasuregroup{อปฺปฏิจฺฉนฺนา เอกาห, \\ ปญฺจาหปกฺขทสนฺนํ, \\ สุทฺธนฺโต จ วิพฺภมนฺโต, \\ ตตฺถ สญฺญิโน เทฺว ยถา, \\ มิสฺสกทิฏฺฐิโน เทฺว จ,}{\csromangathabat{อปฺปฏิจฺฉนฺนา เอกาห,}{ทฺวีห ตีห จตูห จ.} \\ \csromangathabat{ปญฺจาหปกฺขทสนฺนํ,}{อาปตฺติมาห มหามุนิ.} \\ \csromangathabat{สุทฺธนฺโต จ วิพฺภมนฺโต,}{ปริมาณมุขํ เทฺว ภิกฺขู.} \\ \csromangathabat{ตตฺถ สญฺญิโน เทฺว ยถา,}{เวมติกา ตเถว จ.} \\ \csromangathabat{มิสฺสกทิฏฺฐิโน เทฺว จ,}{อสุทฺธเกกทิฏฺฐิโน.} \\ เทฺว เจว สุทฺธทิฏฺฐิโน. \\ ตเถว จ เอโก ฉาเทติ, \\ อถ มกฺขมเตน จ. \\ อุมฺมตฺตกเทสนญฺจ, \\ มูลา อฏฺฐารส\textsuperscript{1} วิสุทฺธโต. \\ อาจริยานํ วิภชฺชปทานํ\textsuperscript{1}, \\ ตมฺพปณฺณิทีปปสาทกานํ. \\ มหาวิหารวาสีนํ, \\ วาจนา สทฺธมฺมฏฺฐิติยาติ.}
\csromangathameasurewak{อถ มกฺขมเตน จ. \\ อุมฺมตฺตกเทสนญฺจ, \\ มูลา อฏฺฐารส\textsuperscript{1} วิสุทฺธโต. \\ อาจริยานํ วิภชฺชปทานํ\textsuperscript{1}, \\ ตมฺพปณฺณิทีปปสาทกานํ. \\ มหาวิหารวาสีนํ, \\ วาจนา สทฺธมฺมฏฺฐิติยาติ.}
\csromangathasetleft{อปฺปฏิจฺฉนฺนา เอกาห, \\ ปญฺจาหปกฺขทสนฺนํ, \\ สุทฺธนฺโต จ วิพฺภมนฺโต, \\ ตตฺถ สญฺญิโน เทฺว ยถา, \\ มิสฺสกทิฏฺฐิโน เทฺว จ,}
\csromangathagroupwithclosermajor{\csromangathastanza{\csromanfolio{187}\csromangathabat{อปฺปฏิจฺฉนฺนา เอกาห,}{ทฺวีห ตีห จตูห จ.} \\ \csromangathabat{ปญฺจาหปกฺขทสนฺนํ,}{อาปตฺติมาห มหามุนิ.}}\csromangathastanza{\csromangathabat{สุทฺธนฺโต จ วิพฺภมนฺโต,}{ปริมาณมุขํ เทฺว ภิกฺขู.} \\ \csromangathabat{ตตฺถ สญฺญิโน เทฺว ยถา,}{เวมติกา ตเถว จ.}}\csromangathastanza{\csromangathabat{มิสฺสก\-ทิฏฺฐิโน เทฺว จ,}{อสุทฺธเกกทิฏฺฐิโน.} \\ เทฺว เจว สุทฺธทิฏฺฐิโน. \\ ตเถว จ เอโก ฉาเทติ,}\csromangathastanza{\csromangathawak{อถ มกฺขมเตน จ. \\ อุมฺมตฺตกเทสนญฺจ, \\ มูลา อฏฺฐารส\csromansharedfootnote{a94a9b7d15191ef5}{ปนฺนรส (ก)} วิสุทฺธโต. \\ อาจริยานํ วิภชฺชปทานํ\csromansharedfootnote{a94a9b7d15191ef5}{ปนฺนรส (ก)},}}\csromangathastanza{\csromangathawak{ตมฺพปณฺณิทีปปสาทกานํ. \\ มหาวิหารวาสีนํ, \\ วาจนา สทฺธมฺม\-ฏฺฐิติยาติ.}}}{สมุจฺจย\-กฺขนฺธโก นิฏฺฐิโต.}
\csromanmatikaanchor{mtk.101}
\csromantocmarknopage{chapter}{4. สมถกฺขนฺธก}{mtk.101}
\bookmark[dest={mtk.101},level=0]{4. สมถกฺขนฺธก}
\chapterheadpageread{4. สมถ\-กฺขนฺธก}
\csromanmatikaanchor{mtk.102}
\csromantocmarkref{section}{1. สมฺมุขาวินย}{mtk.102}
\bookmark[dest={mtk.102},level=1]{1. สมฺมุขาวินย}
\csromanheader{1}{1. \textbf{สมฺมุขาวินย}}
\proseitem{185}{\csromanfolio{188}เตน สมเยน พุทฺโธ ภควา สาวตฺถิยํ วิหรติ เชตวเน อนาถ\-ปิณฺฑิกสฺส อาราเม. เตน โข ปน สมเยน ฉพฺพคฺคิยา ภิกฺขู อสมฺมุขีภูตานํ ภิกฺขูนํ กมฺมานิ กโรนฺติ ตชฺชนียมฺปิ นิยสฺสมฺปิ ปพฺพาชนียมฺปิ ปฏิสารณียมฺปิ อุกฺเขปนียมฺปิ. เย เต ภิกฺขู อปฺปิจฺฉา ฯเปฯ เต อุชฺฌายนฺติ ขิยฺยนฺติ วิปาเจนฺติ “กถํ หิ นาม ฉพฺพคฺคิยา ภิกฺขู อสมฺมุขีภูตานํ ภิกฺขูนํ กมฺมานิ กริสฺสนฺติ ตชฺชนียมฺปิ นิยสฺสมฺปิ ปพฺพาชนียมฺปิ ปฏิสารณียมฺปิ อุกฺเขปณียมฺปี”ติ. อถ โข เต ภิกฺขู ภควโต เอตมตฺถํ อาโรเจสุํ ฯเปฯ. สจฺจํ กิร ภิกฺขเว ฉพฺพคฺคิยา ภิกฺขู อสมฺมุขีภูตานํ ภิกฺขูนํ กมฺมานิ กโรนฺติ ตชฺชนียมฺปิ นิยสฺสมฺปิ ปพฺพาชนียมฺปิ ปฏิสารณียมฺปิ อุกฺเขปนียมฺปีติ. สจฺจํ ภควาติ. วิครหิ พุทฺโธ ภควา อนนุจฺฉวิกํ ภิกฺขเว เตสํ โมฆปุริสานํ อนนุโลมิกํ อปฺปติรูปํ อสฺสามณกํ อกปฺปิยํ อกรณียํ. กถํ หิ นาม เต ภิกฺขเว โมฆปุริสา อสมฺมุขีภูตานํ ภิกฺขูนํ กมฺมานิ กริสฺสนฺติ ตชฺชนียมฺปิ นิยสฺสมฺปิ ปพฺพาชนียมฺปิ ปฏิสารณียมฺปิ อุกฺเขปนียมฺปิ. เนตํ ภิกฺขเว อปฺปสนฺนานํ วา ปสาทาย ฯเปฯ วิครหิตฺวา ฯเปฯ ธมฺมิํ กถํ กตฺวา ภิกฺขู อามนฺเตสิ\csromandash{}น ภิกฺขเว อสมฺมุขีภูตานํ ภิกฺขูนํ กมฺมํ กาตพฺพํ ตชฺชนียํ วา นิยสฺสํ วา ปพฺพาชนียํ วา ปฏิสารณียํ วา อุกฺเขปนียํ วา, โย กเรยฺย, อาปตฺติ ทุกฺกฏสฺส.}

\proseitem{186}{อธมฺมวาที ปุคฺคโล, อธมฺมวาที สมฺพหุลา, อธมฺมวาที สํโฆ, ธมฺมวาที ปุคฺคโล, ธมฺมวาที สมฺพหุลา, ธมฺมวาที สํโฆ.}

\csromanmatikaanchor{mtk.103}
\csromantocmarkref{subsection}{กณฺหปกฺขนวก}{mtk.103}
\bookmark[dest={mtk.103},level=2]{กณฺหปกฺขนวก}
\csromanheader{2}{กณฺหปกฺข\-นวก}
\proseitem{187}{อธมฺมวาที ปุคฺคโล ธมฺมวาทิํ ปุคฺคลํ สญฺญาเปติ นิชฺฌาเปติ เปกฺเขติ อนุเปกฺเขติ ทสฺเสติ อนุทสฺเสติ “อยํ ธมฺโม, อยํ วินโย, อิทํ สตฺถุสาสนํ, อิมํ คณฺหาหิ, อิมํ โรเจหี”ติ. เอวญฺเจตํ อธิกรณํ วูปสมฺมติ, อธมฺเมน วูปสมฺมติ สมฺมุขาวินย\-ปติรูปเกน. (1)}

\prose{อธมฺมวาที ปุคฺคโล ธมฺมวาที สมฺพหุเล สญฺญาเปติ นิชฺฌาเปติ เปกฺเขติ อนุเปกฺเขติ ทสฺเสติ อนุทสฺเสติ “อยํ ธมฺโม, อยํ \csromanfolio{189}วินโย, อิทํ สตฺถุสาสนํ, อิมํ คณฺหถ, อิมํ โรเจถา”ติ. เอวญฺเจตํ อธิกรณํ วูปสมฺมติ, อธมฺเมน วูปสมฺมติ สมฺมุขาวินย\-ปติรูปเกน. (2)}

\prose{อธมฺมวาที ปุคฺคโล ธมฺมวาทิํ สํฆํ สญฺญาเปติ นิชฺฌาเปติ เปกฺเขติ อนุเปกฺเขติ ทสฺเสติ อนุทสฺเสติ “อยํ ธมฺโม, อยํ วินโย, อิทํ สตฺถุสาสนํ, อิมํ คณฺหาหิ, อิมํ โรเจหี”ติ\csromansharedfootnote{5f028dd9ae84b516}{อิมํ คณฺหถ, อิมํ โรเจถาติ (ก)}. เอวญฺเจตํ อธิกรณํ วูปสมฺมติ, อธมฺเมน วูปสมฺมติ สมฺมุขาวินย\-ปติรูปเกน. (3)}

\prose{อธมฺมวาที สมฺพหุลา ธมฺมวาทิํ ปุคฺคลํ สญฺญาเปนฺติ นิชฺฌาเปนฺติ เปกฺเขนฺติ อนุเปกฺเขนฺติ ทสฺเสนฺติ อนุทสฺเสนฺติ “อยํ ธมฺโม, อยํ วินโย, อิทํ สตฺถุสาสนํ, อิมํ คณฺหาหิ, อิมํ โรเจหี”ติ. เอวญฺเจตํ อธิกรณํ วูปสมฺมติ, อธมฺเมน วูปสมฺมติ สมฺมุขาวินย\-ปติรูปเกน. (4)}

\prose{อธมฺมวาที สมฺพหุลา ธมฺมวาที สมฺพหุเล สญฺญาเปนฺติ นิชฺฌาเปนฺติ เปกฺเขนฺติ อนุเปกฺเขนฺติ ทสฺเสนฺติ อนุทสฺเสนฺติ “อยํ ธมฺโม, อยํ วินโย, อิทํ สตฺถุสาสนํ, อิมํ คณฺหถ, อิมํ โรเจถา”ติ. เอวญฺเจตํ อธิกรณํ วูปสมฺมติ, อธมฺเมน วูปสมฺมติ สมฺมุขาวินย\-ปติรูปเกน. (5)}

\prose{อธมฺมวาที สมฺพหุลา ธมฺมวาทิํ สํฆํ สญฺญาเปนฺติ นิชฺฌาเปนฺติ เปกฺเขนฺติ อนุเปกฺเขนฺติ ทสฺเสนฺติ อนุทสฺเสนฺติ “อยํ ธมฺโม, อยํ วินโย, อิทํ สตฺถุสาสนํ, อิมํ คณฺหาหิ, อิมํ โรเจหี”ติ. เอวญฺเจตํ อธิกรณํ วูปสมฺมติ, อธมฺเมน วูปสมฺมติ สมฺมุขาวินย\-ปติรูปเกน. (6)}

\prose{อธมฺมวาที สํโฆ ธมฺมวาทิํ ปุคฺคลํ สญฺญาเปติ นิชฺฌาเปติ เปกฺเขติ อนุเปกฺเขติ ทสฺเสติ อนุทสฺเสติ “อยํ ธมฺโม, อยํ วินโย, อิทํ สตฺถุสาสนํ, อิมํ คณฺหาหิ, อิมํ โรเจหี”ติ. เอวญฺเจตํ อธิกรณํ วูปสมฺมติ, อธมฺเมน วูปสมฺมติ สมฺมุขาวินย\-ปติรูปเกน. (7)}

\prose{อธมฺมวาที สํโฆ ธมฺมวาที สมฺพหุเล สญฺญาเปติ นิชฺฌาเปติ เปกฺเขติ อนุเปกฺเขติ ทสฺเสติ อนุทสฺเสติ “อยํ ธมฺโม, อยํ วินโย, อิทํ สตฺถุสาสนํ, อิมํ คณฺหถ, อิมํ โรเจถา”ติ. เอวญฺเจตํ อธิกรณํ วูปสมฺมติ, อธมฺเมน วูปสมฺมติ สมฺมุขาวินย\-ปติรูปเกน. (8)}

\prosewithcloser{อธมฺมวาที สํโฆ ธมฺมวาทิํ สํฆํ สญฺญาเปติ นิชฺฌาเปติ เปกฺเขติ, อนุเปกฺเขติ ทสฺเสติ อนุทสฺเสติ “อยํ ธมฺโม, อยํ วินโย, \csromanfolio{190}อิทํ สตฺถุสาสนํ, อิมํ คณฺหาหิ, อิมํ โรเจหี”ติ. เอวญฺเจตํ อธิกรณํ วูปสมฺมติ, อธมฺเมน วูปสมฺมติ สมฺมุขาวินยปติรุปเกน.}{(9) กณฺหปกฺข\-นวกํ นิฏฺฐิตํ.}
\csromanmatikaanchor{mtk.104}
\csromantocmarkref{subsection}{สุกฺกปกฺขนวก}{mtk.104}
\bookmark[dest={mtk.104},level=2]{สุกฺกปกฺขนวก}
\csromanheader{2}{สุกฺกปกฺข\-นวก}
\proseitem{188}{ธมฺมวาที ปุคฺคโล อธมฺมวาทิํ ปุคฺคลํ สญฺญาเปติ นิชฺฌาเปติ เปกฺเขติ อนุเปกฺเขติ ทสฺเสติ อนุทสฺเสติ “อยํ ธมฺโม, อยํ วินโย, อิทํ สตฺถุสาสนํ, อิมํ คณฺหาหิ, อิมํ โรเจหี”ติ. เอวญฺเจตํ อธิกรณํ วูปสมฺมติ, ธมฺเมน วูปสมฺมติ สมฺมุขา\-วินเยน. (1)}

\prose{ธมฺมวาที ปุคฺคโล อธมฺมวาที สมฺพหุเล สญฺญาเปติ นิชฺฌาเปติ เปกฺเขติ อนุเปกฺเขติ ทสฺเสติ อนุทสฺเสติ “อยํ ธมฺโม, อยํ วินโย, อิทํ สตฺถุสาสนํ, อิมํ คณฺหถ, อิมํ โรเจถา”ติ. เอวญฺเจตํ อธิกรณํ วูปสมฺมติ, ธมฺเมน วูปสมฺมติ สมฺมุขา\-วินเยน. (2)}

\prose{ธมฺมวาที ปุคฺคโล อธมฺมวาทิํ สํฆํ สญฺญาเปติ นิชฺฌาเปติ เปกฺเขติ อนุเปกฺเขติ ทสฺเสติ อนุทสฺเสติ “อยํ ธมฺโม, อยํ วินโย, อิทํ สตฺถุสาสนํ, อิมํ คณฺหาหิ, อิมํ โรเจหี”ติ. เอวญฺเจตํ อธิกรณํ วูปสมฺมติ, ธมฺเมน วูปสมฺมติ สมฺมุขา\-วินเยน. (3)}

\prose{ธมฺมวาที สมฺพหุลา อธมฺมวาทิํ ปุคฺคลํ สญฺญาเปนฺติ นิชฺฌาเปนฺติ เปกฺเขนฺติ อนุเปกฺเขนฺติ ทสฺเสนฺติ อนุทสฺเสนฺติ “อยํ ธมฺโม, อยํ วินโย, อิทํ สตฺถุสาสนํ, อิมํ คณฺหาหิ, อิมํ โรเจหี”ติ. เอวญฺเจตํ อธิกรณํ วูปสมฺมติ ธมฺเมน วูปสมฺมติ สมฺมุขา\-วินเยน. (4)}

\prose{ธมฺมวาที สมฺพหุลา อธมฺมวาที สมฺพหุเล สญฺญาเปนฺติ นิชฺฌาเปนฺติ เปกฺเขนฺติ อนุเปกฺเขนฺติ ทสฺเสนฺติ อนุทสฺเสนฺติ “อยํ ธมฺโม, อยํ วินโย, อิทํ สตฺถุสาสนํ, อิมํ คณฺหถ, อิมํ โรเจถา”ติ. เอวญฺเจตํ อธิกรณํ วูปสมฺมติ, ธมฺเมน วูปสมติ สมฺมุขา\-วินเยน. (5)}

\prose{ธมฺมวาที สมฺพหุลา อธมฺมวาทิํ สํฆํ สญฺญาเปนฺติ นิชฺฌาเปนฺติ เปกฺเขนฺติ อนุเปกฺเขนฺติ ทสฺเสนฺติ อนุทสฺเสนฺติ “อยํ ธมฺโม, อยํ วินโย, อิทํ \csromanfolio{191}สตฺถุสาสนํ, อิมํ คณฺหาหิ, อิมํ โรเจหี”ติ. เอวญฺเจตํ อธิกรณํ วูปสมฺมติ, ธมฺเมน วูปสมฺมติ สมฺมุขา\-วินเยน. (6)}

\prose{ธมฺมวาที สํโฆ อธมฺมวาทิํ ปุคฺคลํ สญฺญาเปติ นิชฺฌาเปติ เปกฺเขติ อนุเปกฺเขติ ทสฺเสติ อนุทสฺเสติ “อยํ ธมฺโม, อยํ วินโย, อิทํ สตฺถุสาสนํ, อิมํ คณฺหาหิ, อิมํ โรเจหี”ติ. เอวญฺเจตํ อธิกรณํ วูปสมฺมติ, ธมฺเมน วูปสมฺมติ สมฺมุขา\-วินเยน. (7)}

\prose{ธมฺมวาที สํโฆ อธมฺมวาที สมฺพหุเล สญฺญาเปติ นิชฺฌาเปติ เปกฺเขติ อนุเปกฺเขติ ทสฺเสติ อนุทสฺเสติ “อยํ ธมฺโม, อยํ วินโย, อิทํ สตฺถุสาสนํ, อิมํ คณฺหถ, อิมํ โรเจถา”ติ. เอวญฺเจตํ อธิกรณํ วูปสมฺมติ, ธมฺเมน วูปสมฺมติ สมฺมุขา\-วินเยน. (8)}

\prosewithcloser{ธมฺมวาที สํโฆ อธมฺมวาทิํ สํฆํ สญฺญาเปติ นิชฺฌาเปติ เปกฺเขติ อนุเปกฺเขติ ทสฺเสติ อนุทสฺเสติ “อยํ ธมฺโม, อยํ วินโย, อิทํ สตฺถุสาสนํ, อิมํ คณฺหาหิ, อิมํ โรเจหี”ติ. เอวญฺเจตํ อธิกรณํ วูปสมฺมติ, ธมฺเมน วูปสมฺมติ สมฺมุขาวินเยนาติ. (9)}{สุกฺกปกฺข\-นวกํ นิฏฺฐิตํ.}
\csromanmatikaanchor{mtk.105}
\csromantocmarkref{section}{2. สติวินย}{mtk.105}
\bookmark[dest={mtk.105},level=1]{2. สติวินย}
\csromanheader{1}{2. \textbf{สติวินย}}
\proseitem{189}{\csromansymbolfootnote{*}{อิทํ วตฺถุ วิ 1. 243 ปิฏฺฐาทีสุปิ อาคตํ.}เตน สมเยน พุทฺโธ ภควา ราชคเห วิหรติ เวฬุวเน กลนฺทก\-นิวาเป. เตน โข ปน สมเยน อายสฺมตา ทพฺเพน มลฺลปุตฺเตน ชาติยา สตฺตวสฺเสน อรหตฺตํ สจฺฉิกตํ โหติ, ยํ กิญฺจิ สาวเกน ปตฺตพฺพํ, สพฺพํ เตน อนุปฺปตฺตํ โหติ, นตฺถิ จสฺส กิญฺจิ อุตฺตริ\csromansharedfootnote{d8552ab5a18ea4eb}{อุตฺตริํ (สี)} กรณียํ, กตสฺส วา ปติจโย. อถ โข อายสฺมโต ทพฺพสฺส มลฺลปุตฺตสฺส รโหคตสฺส ปฏิสลฺลีนสฺส เอวํ เจตสฺโส ปริวิตกฺโก อุทปาทิ “มยา โข ชาติยา สตฺตวสฺเสน อรหตฺตํ สจฺฉิกตํ, ยํ กิญฺจิ สาวเกน ปตฺตพฺพํ, สพฺพํ มยา อนุปฺปตฺตํ, นตฺถิ จ เม กิญฺจิ อุตฺตริกรณียํ, กตสฺส วา ปติจโย, กิํ นุ โข อหํ สํฆสฺส เวยฺยาวจฺจํ กเรยฺยนฺ”ติ.}

\prose{\csromanfolio{192}อถ โข อายสฺมโต ทพฺพสฺส มลฺลปุตฺตสฺส เอตทโหสิ “ยํนูนาหํ สํฆสฺส เสนาสนญฺจ ปญฺญเปยฺยํ ภตฺตานิ จ อุทฺทิเสยฺยนฺ”ติ. อถ โข อายสฺมา ทพฺโพ มลฺลปุตฺโต สายนฺหสมยํ ปฏิสลฺลานา วุฏฺฐิโต เยน ภควา เตนุปสงฺกมิ, อุปสงฺกมิตฺวา ภควนฺตํ อภิวาเทตฺวา เอกมนฺตํ นิสีทิ, เอกมนฺตํ นิสินฺโน โข อายสฺมา ทพฺโพ มลฺลปุตฺโต ภควนฺตํ เอตทโวจ “อิธ มยฺหํ ภนฺเต รโหคตสฺส ปฏิสลฺลีนสฺส เอวํ เจตโส ปริวิตกฺโก อุทปาทิ ‘มยา โข ชาติยา สตฺตวสฺเสน อรหตฺตํ สจฺฉิกตํ, ยํ กิญฺจิ สาวเกน ปตฺตพฺพํ, สพฺพํ มยา อนุปฺปตฺตํ, นตฺถิ จ เม กิญฺจิ อุตฺตริกรณียํ, กตสฺส วา ปติจโย, กิํ นุ โข อหํ สํฆสฺส เวยฺยาวจฺจํ กเรยฺยนฺ’ติ. ตสฺส มยฺหํ ภนฺเต เอตทโหสิ ‘ยํนูนาหํ สํฆสฺส เสนาสนญฺจ ปญฺญเปยฺยํ ภตฺตานิ จ อุทฺทิเสยฺยนฺ’ติ, อิจฺฉามหํ ภนฺเต สํฆสฺส เสนาสนญฺจ ปญฺญาเปตุํ ภตฺตานิ จ อุทฺทิสิตุนฺ”ติ. สาธุ สาธุ ทพฺพ, เตน หิ ตฺวํ ทพฺพ สํฆสฺส เสนาสนญฺจ ปญฺญเปติ ภตฺตานิ จ อุทฺทิสาหีติ\csromansharedfootnote{777f6432c765609b}{อุทฺทิสาติ (วิ 1. 244 ปิฏฺเฐ.)}. “เอวํ ภนฺเต”ติ โข อายสฺมา ทพฺโพ มลฺลปุตฺโต ภควโต ปจฺจสฺโสสิ. อถ โข ภควา เอตสฺมิํ นิทาเน เอตสฺมิํ ปกรเณ ธมฺมิํ กถํ กตฺวา ภิกฺขู อามนฺเตสิ\csromandash{} เตน หิ ภิกฺขเว สํโฆ ทพฺพํ มลฺลปุตฺตํ เสนาสนปญฺญาปกญฺจ ภตฺตุทฺเทสกญฺจ สมฺมนฺนตุ. เอวญฺจ ปน ภิกฺขเว สมฺมนฺนิตพฺโพ, ปฐมํ ทพฺโพ มลฺลปุตฺโต ยาจิตพฺโพ\csromansharedfootnote{1f6593cbc3ad1ae4}{ทพฺโพ ยจิตพฺโพ (สฺยา, ก)}, ยาจิตฺวา พฺยตฺเตน ภิกฺขุนา ปฏิพเลน สํโฆ ญาเปตพฺโพ\csromandash{}}

\proseitem{190}{“สุณาตุ เม ภนฺเต สํโฆ, ยทิ สํฆสฺส ปตฺตกลฺลํ, สํโฆ อายสฺมนฺตํ ทพฺพํ มลฺลปุตฺตํ เสนาสนปญฺญาปกญฺจ ภตฺตุทฺเทสกญฺจ สมฺมนฺเนยฺย, เอสา ญตฺติ.}

\prose{สุณาตุ เม ภนฺเต สํโฆ, สํโฆ อายสฺมนฺตํ ทพฺพํ มลฺลปุตฺตํ เสนาสนปญฺญาปกญฺจ ภตฺตุทฺเทสกญฺจ สมฺมนฺนติ, ยสฺสายสฺมโต ขมติ อายสฺมโต ทพฺพสฺส มลฺลปุตฺตสฺส เสนาสน\-ปญฺญาปกสฺส จ ภุตฺตุทฺเทสกสฺส จ สมฺมุติ, โส ตุณฺหสฺส, ยสฺส นกฺขมติ, โส ภาเสยฺย.}

\prose{สมฺมโต สํเฆน อายสฺมา ทพฺโพ มลฺลปุตฺโต เสนาสน\-ปญฺญาปโก จ ภตฺตุทฺเทสโก จ ขมติ สํฆสฺส, ตสฺมา ตุณฺหี, เอวเมตํ ธารยามี”ติ.}

\proseitem{191}{\csromanfolio{193}สมฺมโต จ ปนายสฺมา ทพฺโพ มลฺลปุตฺโต สภาคานํ ภิกฺขูนํ เอกชฺฌํ เสนาสนํ ปญฺญเปติ, เย เต ภิกฺขู สุตฺตนฺติกา, เตสํ เอกชฺฌํ เสนาสนํ ปญฺญเปติ “เต อญฺญมญฺญํ สุตฺตนฺตํ สงฺคายิสฺสนฺตี”ติ. เย เต ภิกฺขู วินยธรา, เตสํ เอกชฺฌํ เสนาสนํ ปญฺญเปติ “เต อญฺญมญฺญํ วินยํ วินิจฺฉินิ\-สฺสนฺตี’ติ. เย เต ภิกฺขู ธมฺมกถิกา, เตสํ เอกชฺฌํ เสนาสนํ ปญฺญเปติ “เต อญฺญมญฺญํ ธมฺมํ สากจฺฉิสฺสนฺตี”ติ. เย เต ภิกฺขู ฌายิโน, เตสํ เอกชฺฌํ เสนาสนํ ปญฺญเปติ “เต อญฺญมญฺญํ น พฺยาพาธิสฺสนฺตี”ติ. เย เต ภิกฺขู ติรจฺฉาน\-กถิกา กายทฬฺหิ\-พหุลา\csromansharedfootnote{3fb9caade9574027}{กายทฑฺฒิพหุลา (สี)} วิหรนฺติ, เตสมฺปิ เอกชฺฌํ เสนาสนํ ปญฺญเปติ “อิมายปิเม อายสฺมนฺโต รติยา อจฺฉิสฺสนฺตี”ติ. เยปิ เต ภิกฺขู วิกาเล อาคจฺฉนฺติ, เตสมฺปิ เตโชธาตุํ สมาปชฺชิตฺวา เตเนว อาโลเกน เสนาสนํ ปญฺญเปติ, อปิสุ ภิกฺขู สญฺจิจฺจ วิกาเล อาคจฺฉนฺติ “มยํ อายสฺมโต ทพฺพสฺส มลฺลปุตฺตสฺส อิทฺธิ\-ปาฏิหาริยํ ปสฺสิสฺสามา”ติ.}

\prose{เต อายสฺมนฺตํ ทพฺพํ มลฺลปุตฺตํ อุปสงฺกมิตฺวา เอวํ วทนฺติ “อมฺหากํ อาวุโส ทพฺพ เสนาสนํ ปญฺญเปหี”ติ. เต อายสฺมา ทพฺโพ มลฺลปุตฺโต เอวํ วเทติ “กตฺถ อายสฺมนฺตา อิจฺฉนฺติ, กตฺถ ปญฺญเปมี”ติ. เต สญฺจิจฺจ ทูเร อปทิสนฺติ “อมฺหากํ อาวุโส ทพฺพ คิชฺฌกูเฏ ปพฺพเต เสนาสนํ ปญฺญเปหิ, อมฺหากํ อาวุโส โจรปปาเต เสนาสนํ ปญฺญเปหิ, อมฺหากํ อาวุโส อิสิคิลิปสฺเส กาฬสิลายํ เสนาสนํ ปญฺญเปหิ, อมฺหากํ อาวุโส เวภารปสฺเส สตฺตปณฺณิ\-คุหายํ เสนาสนํ ปญฺญเปหิ, อมฺหากํ อาวุโส สีตวเน สปฺป\-โสณฺฑิก\-ปพฺภาเร เสนาสนํ ปญฺญเปหิ, อมฺหากํ อาวุโส โคตมก\-กนฺทรายํ\csromansharedfootnote{a54c7ff3859a28e2}{โคมฏกนฺทรายํ (สฺยา, กํ)} เสนาสนํ ปญฺญเปหิ, อมฺหากํ อาวุโส ตินฺทุก\-กนฺทรายํ เสนาสนํ ปญฺญเปหิ, อมฺหากํ อาวุโส ตโปท\-กนฺทรายํ\csromansharedfootnote{7fe579e5fc849c61}{กโปตกนฺทรายํ (ก)} เสนาสนํ ปญฺญเปหิ, อมฺหากํ อาวุโส ตโปทาราเม เสนาสนํ ปญฺญเปหิ, อมฺหากํ อาวุโส ชีวกมฺพวเน เสนาสนํ ปญฺญเปหิ, อมฺหากํ อาวุโส มทฺทกุจฺฉิมฺหิ มิคทาเย เสนาสนํ ปญฺญเปหี”ติ.}

\prose{\csromanfolio{194}เตสํ อายสฺมา ทพฺโพ มลฺลปุตฺโต เตโชธาตุํ สมาปชฺชิตฺวา องฺคุลิยา ชลมานาย ปุรโต ปุรโต คจฺฉติ. เตปิ เตเนว อาโลเกน อายสฺมโต ทพฺพสฺส มลฺลปุตฺตสฺส ปิฏฺฐิโต ปิฏฺฐิโต คจฺฉนฺติ. เตสํ อายสฺมา ทพฺโพ มลฺลปุตฺโต เอวํ เสนาสนํ ปญฺญเปติ “อยํ มญฺโจ, อิทํ ปีฐํ, อยํ ภิสิ, อิทํ พิพฺโพหนํ\csromansharedfootnote{4f6ab0f5086911e0}{พิมฺโพหนํ (สี, สฺยา, กํ)}, อิทํ วจฺจฏฺฐานํ, อิทํ ปสฺสาวฏฺฐานํ, อิทํ ปานียํ, อิทํ ปริโภชนียํ, อยํ กตฺตรทณฺโฑ, อิทํ สํฆสฺส กติก\-สณฺฐานํ, อิมํ กาลํ ปวิสิตพฺพํ, อิมํ กาลํ นิกฺขมิตพฺพนฺ”ติ. เตสํ อายสฺมา ทพฺโพ มลฺลปุตฺโต เอวํ เสนาสนํ ปญฺญเปตฺวา ปุนเทว เวฬุวนํ ปจฺจาคจฺฉติ.}

\proseitem{192}{เตน โข ปน สมเยน เมตฺติย\-ภูมชกา\csromansharedfootnote{d0eca3f3f24fbbda}{เมตฺติยภุมฺมชกา (สี, สฺยา, กํ)} ภิกฺขู นวกา เจว โหนฺติ อปฺปปุญฺญา จ, ยานิ สํฆสฺส ลามกานิ เสนาสนานิ, ตานิ เตสํ ปาปุณนฺติ ลามกานิ จ ภตฺตานิ. เตน โข ปน สมเยน ราชคเห มนุสฺสา อิจฺฉนฺติ เถรานํ ภิกฺขูนํ อภิสงฺขาริกํ ปิณฺฑปาตํ ทาตุํ สปฺปิมฺปิ เตลมฺปิ อุตฺตริ\-ภงฺคมฺปิ. เมตฺติย\-ภูมชกานํ ปน ภิกฺขูนํ ปากติกํ เทนฺติ ยถารนฺธํ\csromansharedfootnote{5eecfb84f79ae5de}{ยถารทฺธํ (สฺยา)} กณาชกํ พิลงฺคทุติยํ, เต ปจฺฉาภตฺตํ ปิณฺฑปาต\-ปฺปฏิกฺกนฺตา เถเร ภิกฺขู ปุจฺฉนฺติ “ตุมฺหากํ อาวุโส ภตฺตคฺเค กิํ อโหสิ, ตุมฺหากํ กิํ อโหสี”ติ\csromansharedfootnote{d37735a93a7d8785}{กิํ นาโหสิ (สฺยา, กํ)}. เอกจฺเจ เถรา เอวํ วทนฺติ “อมฺหากํ อาวุโส สปฺปิ อโหสิ, เตลํ อโหสิ, อุตฺตริภงฺคํ อโหสี”ติ. เมตฺติย\-ภูมชกา ปน ภิกฺขู เอวํ วทนฺติ “อมฺหากํ อาวุโส น กิญฺจิ อโหสิ, ปากติกํ ยถารนฺธํ กณาชกํ พิลงฺคทุติยนฺ’ติ.}

\prose{เตน โข ปน สมเยน กลฺยาณภตฺติโก คหปติ สํฆสฺส จตุกฺกภตฺตํ เทติ นิจฺจภตฺตํ, โส ภตฺตคฺเค สปุตฺตทาโร อุปติฏฺฐิตฺวา ปริวิสติ, อญฺเญ โอทเนน ปุจฺฉนฺติ, อญฺเญ สูเปน ปุจฺฉนฺติ, อญฺเญ เตเลน ปุจฺฉนฺติ, อญฺเญ อุตฺตริภงฺเคน ปุจฺฉนฺติ. เตน โข ปน สมเยน กลฺยาณภตฺติกสฺส คหปติโน ภตฺตํ สฺวาตนาย เมตฺติย\-ภูมชกานํ ภิกฺขูนํ อุทฺทิฏฺฐํ โหติ. อถ โข กลฺยาณภตฺติโก คหปติ อารามํ อคมาสิ เกนจิเทว กรณีเยน. โส เยนยสฺมา ทพฺโพ มลฺลปุตฺโต เตนุปสงฺกมิ, อุปสงฺกมิตฺวา อายสฺมนฺตํ ทพฺพํ มลฺลปุตฺตํ อภิวาเทตฺวา \csromanfolio{195}เอกมนฺตํ นิสีทิ, เอกมนฺตํ นิสินฺนํ โข กลฺยาณภตฺติกํ คหปติํ อายสฺมา ทพฺโพ มลฺลปุตฺโต ธมฺมิยา กถาย สนฺทสฺเสสิ สมาทเปสิ สมุตฺเตเชสิ สมฺปหํเสสิ. อถ โข กลฺยาณภตฺติโก คหปติ อายสฺมตา ทพฺเพน มลฺลปุตฺเตน ธมฺมิยา กถาย สนฺทสฺสิโต สมาทปิโต สมุตฺเตชิโต สมฺปหํสิโต อายสฺมนฺตํ ทพฺพํ มลฺลปุตฺตํ เอตทโวจ “กสฺส ภนฺเต อมฺหากํ ฆเร สฺวาตนาย ภตฺตํ อุทฺทิฏฺฐนฺ”ติ. เมตฺติย\-ภูมชกานํ โข คหปติ ภิกฺขูนํ ตุมฺหากํ ฆเร สฺวาตนาย ภตฺตํ อุทฺทิฏฺฐนฺติ. อถ โข กลฺยาณภตฺติโก คหปติ อนตฺตมโน อโหสิ “กถํ หิ นาม ปาปภิกฺขู อมฺหากํ ฆเร ภุญฺชิสฺสนฺตี”ติ ฆรํ คนฺตฺวา ทาสิํ อาณาเปสิ “เย เช เสฺว ภตฺติกา อาคจฺฉนฺติ, เต โกฏฺฐเก อาสนํ ปญฺญเปตฺวา กณาชเกน พิลงฺค\-ทุติเยน ปริวิสา”ติ. “เอวํ อยฺยา”ติ โข สา ทาสี กลฺยาณภตฺติกสฺส คหปติโน ปจฺจสฺโสสิ.}

\prose{อถ โข เมตฺติย\-ภูมชกา ภิกฺขู “หิยฺโย โข อาวุโส อมฺหากํ กลฺยาณภตฺติกสฺส คหปติโน ภตฺตํ อุทฺทิฏฺฐํ, เสฺว อเมฺห กลฺยาณภตฺติโก คหปติ สปุตฺตทาโร อุปติฏฺฐิตฺวา ปริวิสิสฺสติ, อญฺเญ โอทเนน ปุจฺฉิสฺสนฺติ, อญฺเญ สูเปน ปุจฺฉิสฺสนฺติ, อญฺเญ เตเลน ปุจฺฉิสฺสนฺติ, อญฺเญ อุตฺตริภงฺเคน ปุจฺฉิสฺสนฺตี”ติ เต เตเนว โสมนสฺเสน น จิตฺตรูปํ รตฺติยา สุปิํสุ. อถ โข เมตฺติย\-ภูมชกา ภิกฺขู ปุพฺพณฺหสมยํ นิวาเสตฺวา ปตฺต\-จีวรมา\-ทาย เยน กลฺยาณภตฺติกสฺส คหปติโน นิเวสนํ เตนุ\-ปสงฺกมิํสุ. อทฺทสา โข สา ทาสี เมตฺติย\-ภูมชเก ภิกฺขู ทูรโตว อาคจฺฉนฺเต, ทิสฺวาน โกฏฺฐเก อาสนํ ปญฺญเปตฺวา เมตฺติย\-ภูมชเก ภิกฺขู เอตทโวจ “นิสีทถ ภนฺเต”ติ. อถ โข เมตฺติย\-ภูมชกานํ ภิกฺขูนํ เอตทโหสิ “นิสฺสํสยํ โข น ตาว ภตฺตํ สิทฺธํ ภวิสฺสติ, ยถา มยํ โกฏฺฐเก นิสีทาปิยามา”ติ\csromansharedfootnote{4ecb2502b92b06da}{นิสีทาเปยฺยามาติ (ก)}. อถ โข สา ทาสี กณาชเกน\csromansharedfootnote{eab9d71f4044edd8}{กาณาชเกน (สฺยา, กํ)} พิลงฺค\-ทุติเยน อุปคญฺฉิ “ภุญฺชถ ภนฺเต”ติ. มยํ โข ภคินิ นิจฺจภตฺติกาติ. ชานามิ อยฺยา “นิจฺจภตฺติกา”ติ, อปิ จาหํ หิยฺโยว คหปตินา อาณตฺตา “เย เช เสฺว ภตฺติกา อาคจฺฉนฺติ, เต โกฏฺฐเก อาสนํ ปญฺญเปตฺวา กณาชเกน พิลงฺค\-ทุติเยน ปริวิสา”ติ, ภุญฺชถ ภนฺเตติ. อถ โข เมตฺติย\-ภูมชกา ภิกฺขู “หิยฺโย โข อาวุโส กลฺยาณภตฺติโก คหปติ อารามํ อคมาสิ ทพฺพสฺส \csromanfolio{196}มลฺลปุตฺตสฺส สนฺติเก, นิสฺสํสยํ โข มยํ ทพฺเพน มลฺลปุตฺเตน คหปติโน อนฺตเร ปริภินฺนา”ติ\csromansharedfootnote{6a0541fdcbc5f3cb}{สนฺติเก ปริภินฺนาติ (สฺยา, กํ)} เต เตเนว โทมนสฺเสน น จิตฺตรูปํ ภุญฺชิํสุ. อถ โข เมตฺติย\-ภูมชกา ภิกฺขู ปจฺฉาภตฺตํ ปิณฺฑปาต\-ปฺปฏิกฺกนฺตา อารามํ คนฺตฺวา ปตฺตจีวรํ ปฏิสาเมตฺวา พหา\-รามโกฏฺฐเก สํฆาฏิปลฺลตฺถิกาย นิสีทิํสุ ตุณฺหีภูตา มงฺกุภูตา ปตฺตกฺขนฺธา อโธมุขา ปชฺฌายนฺตา อปฺปฏิภานา.}

\prose{อถ โข เมตฺติยา ภิกฺขุนี เยน เมตฺติย\-ภูมชกา ภิกฺขู เตนุปสงฺกมิ, อุปสงฺกมิตฺวา เมตฺติย\-ภูมชเก ภิกฺขู เอตทโวจ “วนฺทามิ อยฺยา”ติ. เอวํ วุตฺเต เมตฺติย\-ภูมชกา ภิกฺขู นาลปิํสุ. ทุติยมฺปิ โข ฯเปฯ. ตติยมฺปิ โข เมตฺติยา ภิกฺขุนี เมตฺติย\-ภูมชเก ภิกฺขู เอตทโวจ “วนฺทามิ อยฺยา”ติ. ตติยมฺปิ โข เมตฺติย\-ภูมชกา ภิกฺขู นาลปิํสุ. กฺยาหํ อยฺยานํ อปรชฺฌามิ, กิสฺส มํ อยฺยา นาลปนฺตีติ. ตถา หิ ปน ตฺวํ ภคินิ อเมฺห ทพฺเพน มลฺลปุตฺเตน วิเหฐิยมาเน อชฺฌุเปกฺขสีติ. กฺยาหํ อยฺยา กโรมีติ. สเจ โข ตฺวํ ภคินิ อิจฺเฉยฺยาสิ, อชฺเชว ภควา อายสฺมนฺตํ ทพฺพํ มลฺลปุตฺตํ นาสาเปยฺยาติ. กฺยาหํ อยฺยา กโรมิ, กิํ มยา สกฺกา กาตุนฺติ. เอหิ ตฺวํ ภคินิ, เยน ภควา เตนุปสงฺกม, อุปสงฺกมิตฺวา ภควนฺตํ เอวํ วเทหิ “อิทํ ภนฺเต นจฺฉนฺนํ นปฺปติรูปํ, ยายํ ภนฺเต ทิสา อภยา อนีติกา อนุปทฺทวา, สายํ ทิสา สภยา สอีติกา สอุปทฺทวา, ยโต นิวาตํ ตโต สวาตํ\csromansharedfootnote{4cd960ef6c9105c4}{ตโต ปวาตํ (สี, สฺยา, กํ)}, อุทกํ มญฺเญ อาทิตฺตํ, อยฺเยนมฺหิ ทพฺเพน มลฺลปุตฺเตน ทูสิตา”ติ. “เอวํ อยฺยา”ติ โข เมตฺติยา ภิกฺขุนี เมตฺติย\-ภูมชกานํ ภิกฺขูนํ ปฏิสฺสุตฺวา\csromansharedfootnote{ba2aa24167250915}{ปฏิสฺสุณิตฺวา (สฺยา, กํ)} เยน ภควา เตนุปสงฺกมิ, อุปสงฺกมิตฺวา ภควนฺตํ อภิวาเทตฺวา เอกมนฺตํ อฏฺฐาสิ, เอกมนฺตํ ฐิตา โข เมตฺติยา\csromansharedfootnote{86621a7c9cccefb2}{สา เมตฺติยา (สฺยา, ก)} ภิกฺขุนี ภควนฺตํ เอตทโวจ “อิทํ ภนฺเต นจฺฉนฺนํ นปฺปติรูปํ, ยายํ ภนฺเต ทิสา อภยา อนีติกา อนุปทฺทวา, สายํ ทิสา สภยา สอีติกา สอุปทฺทวา, ยโต นิวาตํ ตโต สวาตํ, อุทกํ มญฺเญ อาทิตฺตํ, อยฺเยนมฺหิ ทพฺเพน มลฺลปุตฺเตน ทูสิตา”ติ.}

\proseitem{193}{อถ โข ภควา เอตสฺมิํ นิทาเน เอตสฺมิํ ปกรเณ ภิกฺขุสํฆํ สนฺนิปาตาเปตฺวา อายสฺมนฺตํ ทพฺพํ มลฺลปุตฺตํ ปฏิปุจฺฉิ “สรสิ ตฺวํ ทพฺพ เอวรูปํ \csromanfolio{197}กตฺตา, ยถายํ ภิกฺขุนี อาหา”ติ. ยถา มํ ภนฺเต ภควา ชานาตีติ. ทุติยมฺปิ โข ภควา อายสฺมนฺตํ ทพฺพํ มลฺลปุตฺตํ เอตทโวจ “สรสิ ตฺวํ ทพฺพ เอวรูปํ กตฺตา, ยถายํ ภิกฺขุนี อาหา”ติ. ยถา มํ ภนฺเต ภควา ชานาตีติ. ตติยมฺปิ โข ภควา อายสฺมนฺตํ ทพฺพํ มลฺลปุตฺตํ เอตทโวจ “สรสิ ตฺวํ ทพฺพ เอวรูปํ กตฺตา, ยถายํ ภิกฺขุนี อาหา”ติ. ยถา มํ ภนฺเต ภควา ชานาตีติ. น โข ทพฺพ ทพฺพา เอวํ นิพฺเพเฐนฺติ, สเจ ตยา กตํ ‘กตนฺ’ติ วเทหิ, สเจ อกตํ ‘อกตนฺ’ติ วเทหีติ. ยโตหํ ภนฺเต ชาโต, นาภิชานามิ สุปินนฺเตนปิ เมถุนํ ธมฺมํ ปฏิเสวิตา, ปเคว ชาคโรติ. อถ โข ภควา ภิกฺขู อามนฺเตสิ “เตน หิ ภิกฺขเว เมตฺติยํ ภิกฺขุนิํ นาเสถ, อิเม จ ภิกฺขู อนุยุญฺชถา”ติ. อิทํ วตฺวา ภควา อุฏฺฐายาสนา วิหารํ ปาวิสิ.}

\prose{อถ โข เต ภิกฺขู เมตฺติยํ ภิกฺขุนิํ นาเสสุํ. อถ โข เมตฺติย\-ภูมชกา ภิกฺขู เต ภิกฺขู เอตทโวจุํ “มาวุโส เมตฺติยํ ภิกฺขุนิํ นาเสถ, น สา กิญฺจิ อปรชฺฌติ, อเมฺหหิ สา อุสฺสาหิตา กุปิเตหิ อนตฺตมเนหิ จาวนาธิปฺปาเยหี”ติ. กิํ ปน ตุเมฺห อาวุโส อายสฺมนฺตํ ทพฺพํ มลฺลปุตฺตํ อมูลิกาย สีลวิปตฺติยา อนุทฺธํเส\-ถาติ. เอวมาวุโสติ. เย เต ภิกฺขู อปฺปิจฺฉา ฯเปฯ เต อุชฺฌายนฺติ ขิยฺยนฺติ วิปาเจนฺติ “กถํ หิ นาม เมตฺติย\-ภูมชกา ภิกฺขู อายสฺมนฺตํ ทพฺพํ มลฺลปุตฺตํ อมูลิกาย สีลวิปตฺติยา อนุทฺธํเสสฺสนฺตี”ติ. อถ โข เต ภิกฺขู ภควโต เอตมตฺถํ อาโรเจสุํ ฯเปฯ. สจฺจํ กิร ภิกฺขเว เมตฺติย\-ภูมชกา ภิกฺขู ทพฺพํ มลฺลปุตฺตํ อมูลิกาย สีลวิปตฺติยา อนุทฺธํเส\-นฺตีติ. สจฺจํ ภควาติ ฯเปฯ วิครหิตฺวา ฯเปฯ ธมฺมิํ กถํ กตฺวา ภิกฺขู อามนฺเตสิ\csromandash{}}

\prose{เตน หิ ภิกฺขเว สํโฆ ทพฺพสฺส มลฺลปุตฺตสฺส สติเวปุลฺล\-ปฺปตฺตสฺส สติวินยํ เทตุ. เอวญฺจ ปน ภิกฺขเว ทาตพฺโพ, เตน ภิกฺขเว ทพฺเพน มลฺลปุตฺเตน สํฆํ อุปสงฺกมิตฺวา เอกํสํ อุตฺตราสงฺคํ กริตฺวา วุฑฺฒานํ ภิกฺขูนํ ปาเท วนฺทิตฺวา อุกฺกุฏิกํ นิสีทิตฺวา อญฺชลิํ ปคฺคเหตฺวา เอวมสฺส วจนีโย “อิเม มํ ภนฺเต เมตฺติย\-ภูมชกา ภิกฺขู อมูลิกาย สีลวิปตฺติยา อนุทฺธํเสนฺติ, โสหํ ภนฺเต สติเวปุลฺล\-ปฺปตฺโต สํฆํ สติวินยํ ยาจามี”ติ. ทุติยมฺปิ ยาจิตพฺโพ. ตติยมฺปิ ยาจิตพฺโพ “อิเม มํ ภนฺเต เมตฺติย\-ภูมชกา ภิกฺขู อมูลิกาย สีลวิปตฺติยา อนุทฺธํเสนฺติ, \csromanfolio{198}โสหํ\csromansharedfootnote{1982bba1e9b58fe6}{โสหํ ภนฺเต (ก)} สติเวปุลฺล\-ปฺปตฺโต ตติยมฺปิ ภนฺเต สํฆํ สติวินยํ ยาจามี”ติ. พฺยตฺเตน ภิกฺขุนา ปฏิพเลน สํโฆ ญาเปตพฺโพ\csromandash{}}

\proseitem{194}{“สุณาตุ เม ภนฺเต สํโฆ, อิเม เมตฺติย\-ภูมชกา ภิกฺขู อายสฺมนฺตํ ทพฺพํ มลฺลปุตฺตํ อมูลิกาย สีลวิปตฺติยา อนุทฺธํเสนฺติ, อายสฺมา ทพฺโพ มลฺลปุตฺโต สติเวปุลฺล\-ปฺปตฺโต สํฆํ สติวินยํ ยาจติ, ยทิ สํฆสฺส ปตฺตกลฺลํ, สํโฆ อายสฺมโต ทพฺพสฺส มลฺลปุตฺตสฺส สติเวปุลฺล\-ปฺปตฺตสฺส สติวินยํ ทเทยฺย, เอสา ญตฺติ.}

\prose{สุณาตุ เม ภนฺเต สํโฆ, อิเม เมตฺติย\-ภูมชกา ภิกฺขู อายสฺมนฺตํ ทพฺพํ มลฺลปุตฺตํ อมูลิกาย สีลวิปตฺติยา อนุทฺธํเสนฺติ, อายสฺมา ทพฺโพ มลฺลปุตฺโต สติเวปุลฺล\-ปฺปตฺโต สํฆํ สติวินยํ ยาจติ, สํโฆ อายสฺมโต ทพฺพสฺส มลฺลปุตฺตสฺส สติเวปุลฺล\-ปฺปตฺตสฺส สติวินยํ เทติ, ยสฺสายสฺมโต ขมติ อายสฺมโต ทพฺพสฺส มลฺลปุตฺตสฺส สติเวปุลฺล\-ปฺปตฺตสฺส สติวินยสฺส ทานํ, โส ตุณฺหสฺส, ยสฺส นกฺขมติ, โส ภาเสยฺย.}

\prose{ทุติยมฺปิ เอตมตฺถํ วทามิ ฯเปฯ. ตติยมฺปิ เอตมตฺถํ วทามิ. สุณาตุ เม ภนฺเต สํโฆ อิเม เมตฺติย\-ภูมชกา ภิกฺขู อายสฺมนฺตํ ทพฺพํ มลฺลปุตฺตํ อมูลิกาย สีลวิปตฺติยา อนุทฺธํเสนฺติ, อายสฺมา ทพฺโพ มลฺลปุตฺโต สติเวปุลฺล\-ปฺปตฺโต สํฆํ สติวินยํ ยาจติ, สํโฆ อายสฺมโต ทพฺพสฺส มลฺลปุตฺตสฺส สติเวปุลฺล\-ปฺปตฺตสฺส สติวินยํ เทติ, ยสฺสายสฺมโต ขมติ อายสฺมโต ทพฺพสฺส มลฺลปุตฺตสฺส สติเวปุลฺล\-ปฺปตฺตสฺส สติวินยสฺส ทานํ, โส ตุณฺหสฺส, ยสฺส นกฺขมติ, โส ภาเสยฺย.}

\prose{ทินฺโน สํเฆน อายสฺมโต ทพฺพสฺส มลฺลปุตฺตสฺส สติเวปุลฺล\-ปฺปตฺตสฺส สติวินโย, ขมติ สํฆสฺส, ตสฺมา ตุณฺหี, เอวเมตํ ธารยามี”ติ.}

\proseitem{195}{ปญฺจิมานิ ภิกฺขเว ธมฺมิกานิ สติวินยสฺส ทานานิ. สุทฺโธ โหติ ภิกฺขุ อนาปตฺติโก, อนุวทนฺติ จ นํ, ยาจติ จ, ตสฺส สํโฆ สติวินยํ เทติ ธมฺเมน สมคฺเคน. อิมานิ โข ภิกฺขเว ปญฺจ ธมฺมิกานิ สติวินยสฺส ทานานีติ.}

\csromanmatikaanchor{mtk.106}
\csromantocmarkref{section}{3. อมูฬฺหวินย}{mtk.106}
\bookmark[dest={mtk.106},level=1]{3. อมูฬฺหวินย}
\csromanheader{1}{3. \textbf{อมูฬฺหวินย}}
\proseitem{196}{\csromanfolio{199}เตน โข ปน สมเยน คคฺโค ภิกฺขุ อุมฺมตฺตโก โหติ จิตฺตวิปริยาสกโต, เตน อุมฺมตฺตเกน จิตฺตวิปริยาสกเตน พหุํ อสฺสามณกํ อชฺฌาจิณฺณํ โหติ ภาสิต\-ปริกฺกนฺตํ. ภิกฺขู คคฺคํ ภิกฺขุํ อุมฺมตฺตเกน จิตฺตวิปริยาสกเตน อชฺฌาจิณฺเณน อาปตฺติยา โจเทนฺติ “สรตายสฺมา เอวรูปิํ อาปตฺติํ อาปชฺชิตา”ติ. โส เอวํ วเทติ “อหํ โข อาวุโส อุมฺมตฺตโก อโหสิํ จิตฺตวิปริยาสกโต, เตน เม อุมฺมตฺตเกน จิตฺตวิปริยาสกเตน พหุํ อสฺสามณกํ อชฺฌาจิณฺณํ ภาสิต\-ปริกฺกนฺตํ, นาหํ ตํ สรามิ, มูเฬฺหน เม เอตํ กตนฺ”ติ. เอวมฺปิ นํ วุจฺจมานา โจเทนฺเตว “สรตายสฺมา เอวรูปิํ อาปตฺติํ อาปชฺชิตา”ติ. เย เต ภิกฺขู อปฺปิจฺฉา ฯเปฯ เต อุชฺฌายนฺติ ขิยฺยนฺติ วิปาเจนฺติ “กถํ หิ นาม ภิกฺขู คคฺคํ ภิกฺขุํ อุมฺมตฺตเกน จิตฺตวิปริยาสกเตน อชฺฌาจิณฺเณน อาปตฺติยา โจเทสฺสนฺติ ‘สรตายสฺมา เอวรูปิํ อาปตฺติํ อาปชฺชิตา’ติ. โส เอวํ วเทติ ‘อหํ โข อาวุโส อุมฺมตฺตโก อโหสิํ จิตฺตวิปริยาสกโต, เตน เม อุมฺมตฺตเกน จิตฺตวิปริยาสกเตน พหุํ อสฺสามณกํ อชฺฌาจิณฺณํ ภาสิต\-ปริกฺกนฺตํ, นาหํ ตํ สรามิ, มูเฬฺหน เม เอตํ กตนฺ’ติ. เอวมฺปิ นํ วุจฺจมานา โจเทนฺเตว สรตายสฺมา เอวรูปิํ อาปตฺติํ อาปชฺชิตา”ติ. อถ โข เต ภิกฺขู ภควโต เอตมตฺถํ อาโรเจสุํ ฯเปฯ. สจฺจํ กิร ภิกฺขเว ฯเปฯ. สจฺจํ ภควาติ ฯเปฯ วิครหิตฺวา ฯเปฯ ธมฺมิํ กถํ กตฺวา ภิกฺขู อามนฺเตสิ\csromandash{}}

\prose{เตน หิ ภิกฺขเว สํโฆ คคฺคสฺส ภิกฺขุโน อมูฬฺหสฺส อมูฬฺหวินยํ เทตุ. เอวญฺจ ปน ภิกฺขเว ทาตพฺโพ, เตน ภิกฺขเว วคฺเคน ภิกฺขุนา สํฆํ อุปสงฺกมิตฺวา เอกํสํ อุตฺตราสงฺคํ กริตฺวา วุฑฺฒานํ ภิกฺขูนํ ปาเท วนฺทิตฺวา อุกฺกุฏิกํ นิสีทิตฺวา อญฺชลิํ ปคฺคเหตฺวา เอวมสฺส วจนีโย “อหํ ภนฺเต อุมฺมตฺตโก อโหสิํ จิตฺตวิปริยาสกโต, เตน เม อุมฺมตฺตเกน จิตฺตวิปริยาสกเตน พหุํ อสฺสามณกํ อชฺฌาจิณฺณํ ภาสิต\-ปริกฺกนฺตํ, มํ ภิกฺขู อุมฺมตฺตเกน จิตฺตวิปริยาสกเตน อชฺฌาจิณฺเณน อาปตฺติยา โจเทนฺติ ‘สรตายสฺมา เอวรูปิํ อาปตฺติํ อาปชฺชิตา’ติ, ตฺยาหํ เอวํ วทามิ ‘อหํ โข อาวุโส อุมฺมตฺตโก อโหสิํ จิตฺตวิปริยาสกโต, เตน เม อุมฺมตฺตเกน จิตฺตวิปริยาสกเตน พหุํ อสฺสามณกํ อชฺฌาจิณฺณํ ภาสิต\-ปริกฺกนฺตํ, นาหํ ตํ สรามิ, มูเฬฺหน เม เอตํ กตนฺ’ติ, เอวมฺปิ มํ วุจฺจมานา \csromanfolio{200}โจเทนฺเตว ‘สรตายสฺมา เอวรูปิํ อาปตฺติํ อาปชฺชิตา’ติ, โสหํ ภนฺเต อมูโฬฺห สํฆํ อมูฬฺหวินยํ ยาจามี”ติ. ทุติยมฺปิ ยาจิตพฺโพ. ตติยมฺปิ ยาจิตพฺโพ “อหํ ภนฺเต อุมฺมตฺตโก อโหสิํ จิตฺตวิปริยาสกโต, เตน เม อุมฺมตฺตเกน จิตฺตวิปริยาสกเตน พหุํ อสฺสามณกํ อชฺฌาจิณฺณํ ภาสิต\-ปริกฺกนฺตํ, มํ ภิกฺขู อุมฺมตฺตเกน จิตฺตวิปริยาสกเตน อชฺฌาจิณฺเณน อาปตฺติยา โจเทนฺติ ‘สรตายสฺมา เอวรูปิํ อาปตฺติํ อาปชฺชิตา’ติ, ตฺยาหํ เอวํ วทามิ ‘อหํ โข อาวุโส อุมฺมตฺตโก อโหสิํ จิตฺตวิปริยาสกโต, เตน เม อุมฺมตฺตเกน จิตฺตวิปริยาสกเตน พหุํ อสฺสามณกํ อชฺฌาจิณฺณํ ภาสิต\-ปริกฺกนฺตํ, นาหํ ตํ สรามิ, มูเฬฺหน เม เอตํ กตนฺ’ติ, เอวมฺปิ มํ วุจฺจมานา โจเทนฺเตว ‘สรตายสฺมา เอวรูปิํ อาปตฺติํ อาปชฺชิตา’ติ, โสหํ อมูโฬฺห\csromansharedfootnote{544fa17fbff11ef0}{โสหํ ภนฺเต อมูโฬฺห (ก)} ตติยมฺปิ ภนฺเต สํฆํ อมูฬฺหวินยํ ยาจามี”ติ. พฺยตฺเตน ภิกฺขุนา ปฏิพเลน สํโฆ ญาเปตพฺโพ\csromandash{}}

\proseitem{197}{“สุณาตุ เม ภนฺเต สํโฆ, อยํ คคฺโค ภิกฺขุ อุมฺมตฺตโก อโหสิ จิตฺตวิปริยาสกโต, เตน อุมฺมตฺตเกน จิตฺตวิปริยาสกเตน พหุํ อสฺสามณกํ อชฺฌาจิณฺณํ ภาสิต\-ปริกฺกนฺตํ, ภิกฺขู คคฺคํ ภิกฺขุํ อุมฺมตฺตเกน จิตฺตวิปริยาสกเตน อชฺฌาจิณฺเณน อาปตฺติยา โจเทนฺติ ‘สรตายสฺมา เอวรูปิํ อาปตฺติํ อาปชฺชิตา’ติ, โส เอวํ วเทติ ‘อหํ โข อาวุโส อุมฺมตฺตโก อโหสิํ จิตฺตวิปริยาสกโต, เตน เม อุมฺมตฺตเกน จิตฺตวิปริยาสกเตน พหุํ อสฺสามณกํ อชฺฌาจิณฺณํ ภาสิต\-ปริกฺกนฺตํ, นาหํ ตํ สรามิ, มูเฬฺหน เม เอตํ กตนฺ’ติ, เอวมฺปิ นํ วุจฺจมานา โจเทนฺเตว ‘สรตายสฺมา เอวรูปิํ อาปตฺติํ อาปชฺชิตา’ติ, โส อมูโฬฺห สํฆํ อมูฬฺหวินยํ ยาจติ, ยทิ สํฆสฺส ปตฺตกลฺลํ, สํโฆ คคฺคสฺส ภิกฺขุโน อมูฬฺหสฺส อมูฬฺหวินยํ ทเทยฺย, เอสา ญตฺติ.}

\prose{สุณาตุ เม ภนฺเต สํโฆ, อยํ คคฺโค ภิกฺขุ อุมฺมตฺตโก อโหสิ จิตฺตวิปริยาสกโต, เตน อุมฺมตฺตเกน จิตฺตวิปริยาสกเตน พหุํ อสฺสามณกํ อชฺฌาจิณฺณํ ภาสิต\-ปริกฺกนฺตํ, ภิกฺขู คคฺคํ ภิกฺขุํ อุมฺมตฺตเกน จิตฺตวิปริยาสกเตน อชฺฌาจิณฺเณน อาปตฺติยา โจเทนฺติ ‘สรตายสฺมา เอวรูปิํ อาปตฺติํ อาปชฺชิตา’ติ, โส เอวํ วเทติ, ‘อหํ โข อาวุโส อุมฺมตฺตโก อโหสิํ จิตฺตวิปริยาสกโต, เตน เม อุมฺมตฺตเกน จิตฺตวิปริยาสกเตน พหุํ อสฺสามณกํ \csromanfolio{201}อชฺฌาจิณฺณํ ภาสิต\-ปริกฺกนฺตํ, นาหํ ตํ สรามิ, มูเฬฺหน เม เอตํ กตนฺ’ติ, เอวมฺปิ นํ วุจฺจมานา โจเทนฺเตว ‘สรตายสฺมา เอวรูปิํ อาปตฺติํ อาปชฺชิตา’ติ, โส อมูโฬฺห สํฆํ อมูฬฺหวินยํ ยาจติ, สํโฆ คคฺคสฺส ภิกฺขุโน อมูฬฺหสฺส อมูฬฺหวินยํ เทติ, ยสฺสายสฺมโต ขมติ คคฺคสฺส ภิกฺขุโน อมูฬฺหสฺส อมูฬฺห\-วินยสฺส ทานํ, โส ตุณฺหสฺส, ยสฺส นกฺขมติ, โส ภาเสยฺย.}

\prose{ทุติยมฺปิ เอตมตฺถํ วทามิ ฯเปฯ. ตติยมฺปิ เอตมตฺถํ วทามิ ฯเปฯ.}

\prose{ทินฺโน สํเฆน คคฺคสฺส ภิกฺขุโน อมูฬฺหสฺส อมูฬฺหวินโย, ขมติ สํฆสฺส, ตสฺมา ตุณฺหี, เอวเมตํ ธารยามี”ติ.}

\proseitem{198}{ตีณิมานิ ภิกฺขเว อธมฺมิกานิ อมูฬฺห\-วินยสฺส ทานานิ, ตีณิ ธมฺมิกานิ. กตมานิ ตีณิ อธมฺมิกานิ อมูฬฺห\-วินยสฺส ทานานิ.}

\prose{อิธ ปน ภิกฺขเว ภิกฺขุ อาปตฺติํ อาปนฺโน โหติ, ตเมนํ โจเทติ สํโฆ วา สมฺพหุลา วา เอกปุคฺคโล วา “สรตายสฺมา เอวรูปิํ อาปตฺติํ อาปชฺชิตา”ติ, โส สรมาโนว เอวํ วเทติ “น โข อหํ อาวุโส สรามิ เอวรูปิํ อาปตฺติํ อาปชฺชิตา”ติ, ตสฺส สํโฆ อมูฬฺหวินยํ เทติ, อธมฺมิกํ อมูฬฺห\-วินยสฺส ทานํ.}

\prose{อิธ ปน ภิกฺขเว ภิกฺขุ อาปตฺติํ อาปนฺโน โหติ, ตเมนํ โจเทติ สํโฆ วา สมฺพหุลา วา เอกปุคฺคโล วา “สรตายสฺมา เอวรูปิํ อาปตฺติํ อาปชฺชิตา”ติ, โส สรมาโนว เอวํ วเทติ “สรามิ โข อหํ อาวุโส ยถาสุปินนฺเตนา”ติ, ตสฺส สํโฆ อมูฬฺหวินยํ เทติ, อธมฺมิกํ อมูฬฺห\-วินยสฺส ทานํ.}

\prose{อิธ ปน ภิกฺขเว ภิกฺขุ อาปตฺติํ อาปนฺโน โหติ, ตเมนํ โจเทติ สํโฆ วา สมฺพหุลา วา เอกปุคฺคโล วา “สรตายสฺมา เอวรูปิํ อาปตฺติํ อาปชฺชิตา”ติ, โส อนุมฺมตฺตโก อุมฺมตฺตกาลยํ กโรติ “อหมฺปิ โข เอวํ กโรมิ, ตุเมฺหปิ เอวํ กโรถ, มยฺหมฺปิ เอตํ กปฺปติ, ตุมฺหากมฺเปตํ กปฺปตี”ติ, ตสฺส สํโฆ อมูฬฺหวินยํ เทติ, อธมฺมิกํ อมูฬฺห\-วินยสฺส ทานํ. อิมานิ ตีณิ อธมฺมิกานิ อมูฬฺห\-วินยสฺส ทานานิ.}

\csromanmatikaanchor{mtk.107}
\csromantocmarkref{section}{4. ปฏิญฺญาตกรณ}{mtk.107}
\bookmark[dest={mtk.107},level=1]{4. ปฏิญฺญาตกรณ}
\proseitem{199}{\csromanfolio{202}กตมานิ ตีณิ ธมฺมิกานิ อมูฬฺห\-วินยสฺส ทานานิ.}

\prose{อิธ ปน ภิกฺขเว ภิกฺขุ อุมฺมตฺตโก โหติ จิตฺตวิปริยาสกโต, เตน อุมฺมตฺตเกน จิตฺตวิปริยาสกเตน พหุํ อสฺสามณกํ อชฺฌาจิณฺณํ โหติ ภาสิต\-ปริกฺกนฺตํ, ตเมนํ โจเทติ สํโฆ วา สมฺพหุลา วา เอกปุคฺคโล วา “สรตายสฺมา เอวรูปิํ อาปตฺติํ อาปชฺชิตา”ติ, โส อสฺสรมาโนว เอวํ วเทติ “น โข อหํ อาวุโส สรามิ เอวรูปิํ อาปตฺติํ อาปชฺชิตา”ติ, ตสฺส สํโฆ อมูฬฺหวินยํ เทติ, ธมฺมิกํ อมูฬฺห\-วินยสฺส ทานํ.}

\prose{อิธ ปน ภิกฺขเว ภิกฺขุ อุมฺมตฺตโก โหติ จิตฺตวิปริยาสกโต, เตน อุมฺมตฺตเกน จิตฺตวิปริยาสกเตน พหุํ อสฺสามณกํ อชฺฌาจิณฺณํ โหติ ภาสิต\-ปริกฺกนฺตํ, ตเมนํ โจเทติ สํโฆ วา สมฺพหุลา วา เอกปุคฺคโล วา “สรตายสฺมา เอวรูปิํ อาปตฺติํ อาปชฺชิตา”ติ, โส อสฺสรมาโนว เอวํ วเทติ “สรามิ โข อหํ อาวุโส ยถา สุปินนฺเตนา”ติ, ตสฺส สํโฆ อมูฬฺหวินยํ เทติ, ธมฺมิกํ อมูฬฺห\-วินยสฺส ทานํ.}

\prose{อิธ ปน ภิกฺขเว ภิกฺขุ อุมฺมตฺตโก โหติ จิตฺตวิปริยาสกโต, เตน อุมฺมตฺตเกน จิตฺตวิปริยาสกเตน พหุํ อสฺสามณกํ อชฺฌาจิณฺณํ โหติ ภาสิต\-ปริกฺกนฺตํ, ตเมนํ โจเทติ สํโฆ วา สมฺพหุลา วา เอกปุคฺคโล วา “สรตายสฺมา เอวรูปิํ อาปตฺติํ อาปชฺชิตา”ติ, โส อุมฺมตฺตโก อุมฺมตฺตกาลยํ กโรติ “อหมฺปิ เอวํ กโรมิ, ตุเมฺหปิ เอวํ กโรถ, มยฺหมฺปิ เอตํ กปฺปติ, ตุมฺหากมฺเปตํ กปฺปตี”ติ, ตสฺส สํโฆ อมูฬฺหวินยํ เทติ, ธมฺมิกํ อมูฬฺห\-วินยสฺส ทานํ. อิมานิ ตีณิ ธมฺมิกานิ อมูฬฺห\-วินยสฺส ทานานีติ.}

\csromanheader{2}{4. ปฏิญฺญาต\-กรณา}
\proseitem{200}{เตน โข ปน สมเยน ฉพฺพคฺคิยา ภิกฺขู อปฺปฏิญฺญาย ภิกฺขูนํ กมฺมานิ กโรนฺติ ตชฺชนียมฺปิ นิยสฺสมฺปิ ปพฺพาชนียมฺปิ ปฏิสารณียมฺปิ อุกฺเขปนียมฺปิ. เย เต ภิกฺขู อปฺปิจฺฉา ฯเปฯ เต อุชฺฌายนฺติ ขิยฺยนฺติ วิปาเจนฺติ “กถํ หิ นาม ฉพฺพคฺคิยา ภิกฺขู อปฺปฏิญฺญาย ภิกฺขูนํ กมฺมานิ กริสฺสนฺติ ตชฺชนียมฺปิ นิยสฺสมฺปิ ปพฺพาชนียมฺปิ ปฏิสารณียมฺปิ อุกฺเขปนียมฺปี”ติ. อถ \csromanfolio{203}โข เต ภิกฺขู ภควโต เอตมตฺถํ อาโรเจสุํ ฯเปฯ. สจฺจํ กิร ภิกฺขเว ฯเปฯ. สจฺจํ ภควาติ ฯเปฯ วิครหิตฺวา ฯเปฯ ธมฺมิํ กถํ กตฺวา ภิกฺขู อามนฺเตสิ\csromandash{}น ภิกฺขเว อปฺปฏิญฺญาย ภิกฺขูนํ กมฺมํ กาตพฺพํ ตชฺชนียํ วา นิยสฺสํ วา ปพฺพาชนียํ วา ปฏิสารณียํ วา อุกฺเขปนียํ วา, โย กเรยฺย, อาปตฺติ ทุกฺกฏสฺส.}

\proseitem{201}{เอวํ โข ภิกฺขเว อธมฺมิกํ โหติ ปฏิญฺญาต\-กรณํ, เอวํ ธมฺมิกํ. กถญฺจ ภิกฺขเว อธมฺมิกํ โหติ ปฏิญฺญาต\-กรณํ.}

\prose{ภิกฺขุ ปาราชิกํ อชฺฌาปนฺโน โหติ, ตเมนํ โจเทติ สํโฆ วา สมฺพหุลา วา เอกปุคฺคโล วา “ปาราชิกํ อายสฺมา อชฺฌาปนฺโน”ติ. โส เอวํ วเทติ “น โข อหํ อาวุโส ปาราชิกํ อชฺฌาปนฺโน สํฆาทิเสสํ อชฺฌาปนฺโน”ติ. ตํ สํโฆ สํฆาทิเสเสน กาเรติ, อธมฺมิกํ ปฏิญฺญาต\-กรณํ.}

\prose{ภิกฺขุ ปาราชิกํ อชฺฌาปนฺโน โหติ, ตเมนํ โจเทติ สํโฆ วา สมฺพหุลา วา เอกปุคฺคโล วา “ปาราชิกํ อายสฺมา อชฺฌาปนฺโน”ติ. โส เอวํ วเทติ “น โข อหํ อาวุโส ปาราชิกํ อชฺฌาปนฺโน, ถุลฺลจฺจยํ ฯเปฯ ปาจิตฺติยํ ฯเปฯ ปาฏิเทสนียํ ฯเปฯ ทุกฺกฏํ ฯเปฯ ทุพฺภาสิตํ อชฺฌาปนฺโน”ติ. ตํ สํโฆ ทุพฺภาสิเตน กาเรติ, อธมฺมิกํ ปฏิญฺญาต\-กรณํ.}

\prose{ภิกฺขุ สํฆาทิเสสํ ฯเปฯ ถุลฺลจฺจยํ ฯเปฯ ปาจิตฺติยํ ฯเปฯ ปาฏิเทสนียํ ฯเปฯ ทุกฺกฏํ ฯเปฯ ทุพฺภาสิตํ อชฺฌาปนฺโน โหติ, ตเมนํ โจเทติ สํโฆ วา สมฺพหุลา วา เอกปุคฺคโล วา “ทุพฺภาสิตํ อายสฺมา อชฺฌาปนฺโน”ติ. โส เอวํ วเทติ “น โข อหํ อาวุโส ทุพฺภาสิตํ อชฺฌาปนฺโน ปาราชิกํ อชฺฌาปนฺโน”ติ. ตํ สํโฆ ปาราชิเกน กาเรติ, อธมฺมิกํ ปฏิญฺญาต\-กรณํ.}

\prose{ภิกฺขุ ทุพฺภาสิตํ อชฺฌาปนฺโน โหติ, ตเมนํ โจเทติ สํโฆ วา สมฺพหุลา วา เอกปุคฺคโล วา “ทุพฺภาสิตํ อายสฺมา อชฺฌาปนฺโน”ติ. โส เอวํ วเทติ “น โข อหํ อาวุโส ทุพฺภาสิตํ อชฺฌาปนฺโน สํฆาทิเสสํ ฯเปฯ ถุลฺลจฺจยํ ฯเปฯ ปาจิตฺติยํ ฯเปฯ ปาฏิเทสนียํ ฯเปฯ ทุกฺกฏํ อชฺฌาปนฺโน”ติ. ตํ สํโฆ ทุกฺกเฏน กาเรติ, อธมฺมิกํ ปฏิญฺญาต\-กรณํ. เอวํ โข ภิกฺขเว อธมฺมิกํ โหติ ปฏิญฺญาต\-กรณํ.}

\prose{\csromanfolio{204}กถญฺจ ภิกฺขเว ธมฺมิกํ โหติ ปฏิญฺญาต\-กรณํ. ภิกฺขุ ปาราชิกํ อชฺฌาปนฺโน โหติ, ตเมนํ โจเทติ สํโฆ วา สมฺพหุลา วา เอกปุคฺคโล วา “ปาราชิกํ อายสฺมา อชฺฌาปนฺโน”ติ. โส เอวํ วเทติ “อาม อาวุโส ปาราชิกํ อชฺฌาปนฺโน”ติ. ตํ สํโฆ ปาราชิเกน กาเรติ, ธมฺมิกํ ปฏิญฺญาต\-กรณํ.}

\prose{ภิกฺขุ สํฆาทิเสสํ ฯเปฯ ถุลฺลจฺจยํ ฯเปฯ ปาจิตฺติยํ ฯเปฯ ปาฏิเทสนียํ ฯเปฯ ทุกฺกฏํ ฯเปฯ ทุพฺภาสิตํ อชฺฌาปนฺโน โหติ, ตเมนํ โจเทติ สํโฆ วา สมฺพหุลา วา เอกปุคฺคโล วา “ทุพฺภาสิตํ อายสฺมา อชฺฌาปนฺโน”ติ. โส เอวํ วเทติ “อาม อาวุโส ทุพฺภาสิตํ อชฺฌาปนฺโน”ติ. ตํ สํโฆ ทุพฺภาสิเตน กาเรติ, ธมฺมิกํ ปฏิญฺญาต\-กรณํ. เอวํ โข ภิกฺขเว ธมฺมิกํ โหติ ปฏิญฺญาต\-กรณนฺติ.}

\csromanmatikaanchor{mtk.108}
\csromantocmarkref{section}{5. เยภุยฺยสิกา}{mtk.108}
\bookmark[dest={mtk.108},level=1]{5. เยภุยฺยสิกา}
\csromanheader{1}{5. \textbf{เยภุยฺยสิกา}}
\proseitem{202}{เตน โข ปน สมเยน ภิกฺขู สํฆมชฺเฌ ภณฺฑนชาตา กลหชาตา วิวาทาปนฺนา อญฺญมญฺญํ มุขสตฺตีหิ วิตุทนฺตา วิหรนฺติ, น สกฺโกนฺติ ตํ อธิกรณํ วูปสเมตุํ. ภควโต เอตมตฺถํ อาโรเจสุํ. อนุชานามิ ภิกฺขเว เอวรูปํ อธิกรณํ เยภุยฺยสิกาย วูปสเมตุํ. ปญฺจหงฺเคหิ สมนฺนาคโต ภิกฺขุ สลาก\-คฺคาหาปโก สมฺมนฺนิตพฺโพ, โย น ฉนฺทาคติํ คจฺเฉยฺย, น โทสาคติํ คจฺเฉยฺย, น โมหาคติํ คจฺเฉยฺย, น ภยาคติํ คจฺเฉยฺย, คติตาคหิตญฺจ ชาเนยฺย. เอวญฺจ ปน ภิกฺขเว สมฺมนฺนิตพฺโพ, ปฐมํ ภิกฺขุ ยาจิตพฺโพ, ยาจิตฺวา พฺยตฺเตน ภิกฺขุนา ปฏิพเลน สํโฆ ญาเปตพฺโพ\csromandash{}}

\proseitem{203}{“สุณาตุ เม ภนฺเต สํโฆ, ยทิ สํฆสฺส ปตฺตกลฺลํ, สํโฆ อิตฺถนฺนามํ ภิกฺขุํ สลาก\-คฺคาหาปกํ สมฺมนฺเนยฺย, เอสา ญตฺติ.}

\prose{สุณาตุ เม ภนฺเต สํโฆ, สํโฆ อิตฺถนฺนามํ ภิกฺขุํ สลาก\-คฺคาหาปกํ สมฺมนฺนติ, ยสฺสายสฺมโต ขมติ อิตฺถนฺนามสฺส ภิกฺขุโน สลาก\-คฺคาหาปกสฺส สมฺมุติ, โส ตุณฺหสฺส, ยสฺส นกฺขมติ, โส ภาเสยฺย.}

\prose{สมฺมโต สํเฆน อิตฺถนฺนาโม ภิกฺขุ สลาก\-คฺคาหาปโก, ขมติ สํฆสฺส, ตสฺมา ตุณฺหี, เอวเมตํ ธารยามี”ติ.}

\proseitem{204}{\csromanfolio{205}ทสยิเม ภิกฺขเว อธมฺมิกา สลากคฺคาหา, ทส ธมฺมิกา\csromansharedfootnote{103f397d08e069d0}{ทส ธมฺมิกา สลากคฺคาหา (ก)}. กตเม ทส อธมฺมิกา สลากคฺคาหา. โอรมตฺตกญฺจ อธิกรณํ โหติ, น จ คติคตํ โหติ, น จ สริตสาริตํ โหติ, ชานาติ “อธมฺมวาที พหุตรา”ติ, อปฺเปว นาม อธมฺมวาที พหุตรา อสฺสูติ, ชานาติ “สํโฆ ภิชฺชิสฺสตี”ติ, อปฺเปว นาม สํโฆ ภิชฺเชยฺยาติ, อธมฺเมน คณฺหนฺติ, วคฺคา คณฺหนฺติ, น จ ยถาทิฏฺฐิยา คณฺหนฺติ. อิเม ทส อธมฺมิกา สลากคฺคาหา.}

\prose{กตเม ทส ธมฺมิกา สลากคฺคาหา. น จ โอรมตฺตกํ อธิกรณํ โหติ, คติคตญฺจ โหติ, สริตสาริตญฺจ โหติ, ชานาติ “ธมฺมวาที พหุตรา”ติ, อปฺเปว นาม ธมฺมวาที พหุตรา อสฺสูติ, ชานาติ “สํโฆ น ภิชฺชิสฺสตี”ติ, อปฺเปว นาม สํโฆ น ภิชฺเชยฺยาติ, ธมฺเมน คณฺหนฺติ, สมคฺคา คณฺหนฺติ, ยถาทิฏฺฐิยา จ คณฺหนฺติ. อิเม ทส ธมฺมิกา สลากคฺคาหาติ.}

\csromanmatikaanchor{mtk.109}
\csromantocmarkref{section}{6. ตสฺสปาปิยสิกา}{mtk.109}
\bookmark[dest={mtk.109},level=1]{6. ตสฺสปาปิยสิกา}
\csromanheader{1}{6. \textbf{ตสฺสปาปิยสิกา}}
\proseitem{205}{เตน โข ปน สมเยน อุปวาโฬ ภิกฺขุ สํฆมชฺเฌ อาปตฺติยา อนุยุญฺชิยมาโน อวชานิตฺวา ปฏิชานาติ, ปฏิชานิตฺวา อวชานาติ, อญฺเญนญฺญํ ปฏิจรติ, สมฺปชานมุสา ภาสติ. เย เต ภิกฺขู อปฺปิจฺฉา ฯเปฯ เต อุชฺฌายนฺติ ขิยฺยนฺติ วิปาเจนฺติ “กถํ หิ นาม อุปวาโฬ ภิกฺขุ สํฆมชฺเฌ อาปตฺติยา อนุยุญฺชิยมาโน อวชานิตฺวา ปฏิชานิสฺสติ, ปฏิชานิตฺวา อวชานิสฺสติ, อญฺเญนญฺญํ ปฏิจริสฺสติ, สมฺปชานมุสา ภาสิสฺสตี”ติ. อถ โข เต ภิกฺขู ภควโต เอตมตฺถํ อาโรเจสุํ ฯเปฯ. สจฺจํ กิร ภิกฺขเว ฯเปฯ. สจฺจํ ภควาติ ฯเปฯ วิครหิตฺวา ฯเปฯ ธมฺมิํ กถํ กตฺวา ภิกฺขู อามนฺเตสิ\csromandash{} เตน หิ ภิกฺขเว สํโฆ อุปวาฬสฺส ภิกฺขุโน ตสฺสปาปิยสิกา\-กมฺมํ กโรตุ. เอวญฺจ ปน ภิกฺขเว กาตพฺพํ, ปฐมํ อุปวาโฬ ภิกฺขุ โจเทตพฺโพ, โจเทตฺวา สาเรตพฺโพ, สาเรตฺวา อาปตฺติํ อาโรเปตพฺโพ, อาปตฺติํ อาโรเปตฺวา พฺยตฺเตน ภิกฺขุนา ปฏิพเลน สํโฆ ญาเปตพฺโพ\csromandash{}}

\proseitem{206}{“สุณาตุ เม ภนฺเต สํโฆ, อยํ อุปวาโฬ ภิกฺขุ สํฆมชฺเฌ อาปตฺติยา อนุยุญฺชิยมาโน อวชานิตฺวา ปฏิชานาติ, ปฏิชานิตฺวา อวชานาติ, อญฺเญนญฺญํ ปฏิจรติ, สมฺปชานมุสา ภาสติ, \csromanfolio{206}ยทิ สํฆสฺส ปตฺตกลฺลํ, สํโฆ อุปวาฬสฺส ภิกฺขุโน ตสฺสปาปิยสิกา\-กมฺมํ กเรยฺย, เอสา ญตฺติ.}

\prose{สุณาตุ เม ภนฺเต สํโฆ, อยํ อุปวาโฬ ภิกฺขุ สํฆมชฺเฌ อาปตฺติยา อนุยุญฺชิยมาโน อวชานิตฺวา ปฏิชานาติ, ปฏิชานิตฺวา อวชานาติ, อญฺเญนญฺญํ ปฏิจรติ, สมฺปชานมุสา ภาสติ, สํโฆ อุปวาฬสฺส ภิกฺขุโน ตสฺสปาปิยสิกา\-กมฺมํ กโรติ, ยสฺสายสฺมโต ขมติ อุปวาฬสฺส ภิกฺขุโน ตสฺสปาปิยสิกา\-กมฺมสฺส กรณํ, โส ตุณฺหสฺส, ยสฺส นกฺขมติ, โส ภาเสยฺย.}

\prose{ทุติยมฺปิ เอตมตฺถํ วทามิ ฯเปฯ. ตติยมฺปิ เอตมตฺถํ วทามิ ฯเปฯ.}

\prose{กตํ สํเฆน อุปวาฬสฺส ภิกฺขุโน ตสฺสปาปิยสิกา\-กมฺมํ, ขมติ สํฆสฺส, ตสฺมา ตุณฺหี, เอวเมตํ ธารยามี”ติ.}

\proseitem{207}{ปญฺจิมานิ ภิกฺขเว ธมฺมิกานิ ตสฺสปาปิยสิกา\-กมฺมสฺส กรณานิ. อสุจิ จ โหติ, อลชฺชี จ, สานุวาโท จ, ตสฺส สํโฆ ตสฺสปาปิยสิกา\-กมฺมํ กโรติ ธมฺเมน สมคฺเคน. อิมานิ โข ภิกฺขเว ปญฺจ ธมฺมิกานิ ตสฺสปาปิยสิกา\-กมฺมสฺส กรณานิ.}

\csromanmatikaanchor{mtk.110}
\csromantocmarkref{subsection}{อธมฺมกมฺมทฺวาทสก}{mtk.110}
\bookmark[dest={mtk.110},level=2]{อธมฺมกมฺมทฺวาทสก}
\csromanheader{2}{อธมฺมกมฺม\-ทฺวาทสก}
\proseitem{208}{ตีหิ ภิกฺขเว องฺเคหิ สมนฺนาคตํ ตสฺสปาปิยสิกา\-กมฺมํ อธมฺมกมฺมญฺจ โหติ อวินยกมฺมญฺจ ทุวูปสนฺตญฺจ. อสมฺมุขา กตํ โหติ, อปฺปฏิปุจฺฉากตํ โหติ, อปฺปฏิญฺญาย กตํ โหติ ฯเปฯ. อธมฺเมน กตํ โหติ, วคฺเคน กตํ โหติ. อิเมหิ โข ภิกฺขเว ตีหงฺเคหิ สมนฺนาคตํ ตสฺสปาปิยสิกา\-กมฺมํ อธมฺมกมฺมญฺจ โหติ อวินยกมฺมญฺจ ทุวูปสนฺตญฺจ.}

\csromanmatikaanchor{mtk.111}
\csromantocmarkref{subsection}{ธมฺมกมฺมทฺวาทสก}{mtk.111}
\bookmark[dest={mtk.111},level=2]{ธมฺมกมฺมทฺวาทสก}
\csromanheader{2}{ธมฺมกมฺม\-ทฺวาทสก}
\proseitem{209}{ตีหิ ภิกฺขเว องฺเคหิ สมนฺนาคตํ ตสฺสปาปิยสิกา\-กมฺมํ ธมฺมกมฺมญฺจ โหติ วินยกมฺมญฺจ สุวูปสนฺตญฺจ. สมฺมุขา กตํ โหติ, ปฏิปุจฺฉากตํ โหติ, ปฏิญฺญาย กตํ โหติ ฯเปฯ. ธมฺเมน กตํ โหติ, สมคฺเคน กตํ โหติ. อิเมหิ โข ภิกฺขเว ตีหงฺเคหิ สมนฺนาคตํ ตสฺสปาปิยสิกา\-กมฺมํ ธมฺมกมฺมญฺจ โหติ วินยกมฺมญฺจ สุวูปสนฺตญฺจ.}

\csromanmatikaanchor{mtk.112}
\csromantocmarkref{subsection}{อากงฺขมานฉกฺก}{mtk.112}
\bookmark[dest={mtk.112},level=2]{อากงฺขมานฉกฺก}
\proseitem{210}{\csromanfolio{207}ตีหิ ภิกฺขเว องฺเคหิ สมนฺนาคตสฺส ภิกฺขุโน อากงฺขมาโน สํโฆ ตสฺสปาปิยสิกา\-กมฺมํ กเรยฺย. ภณฺฑนการโก โหติ กลหการโก วิวาทการโก ภสฺสการโก สํเฆ อธิกรณ\-การโก, พาโล โหติ อพฺยตฺโต อาปตฺติพหุโล อนปทาโน, คิหิสํสฏฺโฐ วิหรติ อนนุโลมิเกหิ คิหิสํสคฺเคหิ. อิเมหิ โข ภิกฺขเว ตีหงฺเคหิ สมนฺนาคตสฺส ภิกฺขุโน อากงฺขมาโน สํโฆ ตสฺสปาปิยสิกา\-กมฺมํ กเรยฺย. (1)}

\prose{อปเรหิปิ ภิกฺขเว ตีหงฺเคหิ สมนฺนาคตสฺส ภิกฺขุโน อากงฺขมาโน สํโฆ ตสฺสปาปิยสิกา\-กมฺมํ กเรยฺย. อธิสีเล สีลวิปนฺโน โหติ, อชฺฌาจาเร อาจารวิปนฺโน โหติ, อติทิฏฺฐิยา ทิฏฺฐิวิปนฺโน โหติ. อิเมหิ โข ภิกฺขเว ฯเปฯ. (2)}

\prose{อปเรหิปิ ภิกฺขเว ฯเปฯ. พุทฺธสฺส อวณฺณํ ภาสติ, ธมฺมสฺส อวณฺณํ ภาสติ, สํฆสฺส อวณฺณํ ภาสติ. อิเมหิ โข ภิกฺขเว ตีหงฺเคหิ สมนฺนาคตสฺส ภิกฺขุโน อากงฺขมาโน สํโฆ ตสฺสปาปิยสิกา\-กมฺมํ กเรยฺย. (3)}

\prose{ติณฺณํ ภิกฺขเว ภิกฺขูนํ อากงฺขมาโน สํโฆ ตสฺสปาปิยสิกา\-กมฺมํ กเรยฺย. เอโก ภณฺฑนการโก โหติ ฯเปฯ สํเฆ อธิกรณ\-การโก, เอโก พาโล โหติ อพฺยตฺโต อาปตฺติพหุโล อนปทาโน, เอโก คิหิสํสฏฺโฐ วิหรติ อนนุโลมิเกหิ คิหิสํสคฺเคหิ. อิเมสํ โข ภิกฺขเว ติณฺณํ ภิกฺขูนํ อากงฺขมาโน สํโฆ ตสฺสปาปิยสิกา\-กมฺมํ กเรยฺย. (4)}

\prose{อปเรสมฺปิ ภิกฺขเว ติณฺณํ ภิกฺขูนํ อากงฺขมาโน สํโฆ ตสฺสปาปิยสิกา\-กมฺมํ กเรยฺย. เอโก อธิสีเล สีลวิปนฺโน โหติ, เอโก อชฺฌาจาเร อาจารวิปนฺโน โหติ, เอโก อติทิฏฺฐิยา ทิฏฺฐิวิปนฺโน โหติ. อิเมสํ โข ภิกฺขเว ฯเปฯ. (5)}

\prose{อปเรสมฺปิ ภิกฺขเว ฯเปฯ. เอโก พุทฺธสฺส อวณฺณํ ภาสติ, เอโก ธมฺมสฺส อวณฺณํ ภาสติ, เอโก สํฆสฺส อวณฺณํ ภาสติ. อิเมสํ โข ภิกฺขเว ติณฺณํ ภิกฺขูนํ อากงฺขมาโน สํโฆ ตสฺสปาปิยสิกา\-กมฺมํ กเรยฺย. (6)}

\csromanmatikaanchor{mtk.113}
\csromantocmarkref{subsection}{อฏฺฐารสวตฺต}{mtk.113}
\bookmark[dest={mtk.113},level=2]{อฏฺฐารสวตฺต}
\csromanheader{2}{\textbf{อฏฺฐารสวตฺต}}
\proseitem{211}{\csromanfolio{208}ตสฺสปาปิยสิกากมฺม\-กเตน ภิกฺขเว ภิกฺขุนา สมฺมา วตฺติตพฺพํ. ตตฺรายํ สมฺมาวตฺตนา\csromandash{}}

\prose{น อุปสมฺ\-ปาเท\-ตพฺพํ, น นิสฺสโย ทาตพฺโพ, น สามเณโร อุปฏฺฐาเปตพฺโพ, น ภิกฺขุโนวาทก\-สมฺมุติ สาทิตพฺพา, สมฺมเตนปิ ภิกฺขุนิโย น โอวทิตพฺพา ฯเปฯ น ภิกฺขูหิ สมฺป\-โยเช\-ตพฺพนฺติ. อถ โข สํโฆ อุปวาฬสฺส ภิกฺขุโน ตสฺสปาปิยสิกา\-กมฺมํ อกาสิ.}

\csromanmatikaanchor{mtk.114}
\csromantocmarkref{section}{7. ติณวตฺถารก}{mtk.114}
\bookmark[dest={mtk.114},level=1]{7. ติณวตฺถารก}
\csromanheader{1}{7. \textbf{ติณวตฺถารก}}
\proseitem{212}{เตน โข ปน สมเยน ภิกฺขูนํ ภณฺฑน\-ชาตานํ กลหชาตานํ วิวาทาปนฺนานํ วิหรตํ พหุํ อสฺสามณกํ อชฺฌาจิณฺณํ โหติ ภาสิต\-ปริกฺกนฺตํ. อถ โข เตสํ ภิกฺขูนํ เอตทโหสิ “อมฺหากํ โข ภณฺฑน\-ชาตานํ กลหชาตานํ วิเวทาปนฺนานํ วิหรตํ พหุํ อสฺสามณกํ อชฺฌาจิณฺณํ ภาสิต\-ปริกฺกนฺตํ, สเจ มยํ อิมาหิ อาปตฺตีหิ อญฺญมญฺญํ กาเรสฺสาม, สิยาปิ ตํ อธิกรณํ กกฺขฬตฺตาย วาฬตฺตาย\csromansharedfootnote{b776a3c9024c52f7}{กกฺขฬตาย วาฬตาย (สฺยา, กํ)} เภทาย สํวตฺเตยฺย, กถํ นุ โข อเมฺหหิ ปฏิปชฺชิตพฺพนฺ”ติ. ภควโต เอตมตฺถํ อาโรเจสุํ.}

\prose{อิธ ปน ภิกฺขเว ภิกฺขูนํ ภณฺฑน\-ชาตานํ กลหชาตานํ วิวาทาปนฺนานํ วิหรตํ พหุํ อสฺสามณกํ อชฺฌาจิณฺณํ โหติ ภาสิต\-ปริกฺกนฺตํ, ตตฺร เจ ภิกฺขูนํ\csromansharedfootnote{2ba13bad601ee147}{ตตฺร เจ ภิกฺขเว ภิกฺขูนํ (สฺยา)} เอวํ โหติ “อมฺหากํ โข ภณฺฑน\-ชาตานํ กลหชาตานํ วิวาทาปนฺนานํ วิหรตํ พหุํ อสฺสามณกํ อชฺฌาจิณฺณํ ภาสิต\-ปริกฺกนฺตํ, สเจ มยํ อิมาหิ อาปตฺตีหิ อญฺญมญฺญํ กาเรสฺสาม, สิยาปิ ตํ อธิกรณํ กกฺขฬตฺตาย วาฬตฺตาย เภทาย สํวตฺเตยฺยา”ติ. อนุชานามิ ภิกฺขเว เอวรูปํ อธิกรณํ ติณวตฺถารเกน วูปสเมตุํ. เอวญฺจ ปน ภิกฺขเว วูปสเมตพฺพํ, สพฺเพเหว เอกชฺฌํ สนฺนิปติตพฺพํ, สนฺนิปติตฺวา พฺยตฺเตน ภิกฺขุนา ปฏิพเลน สํโฆ ญาเปตพฺโพ\csromandash{}}

\prose{“สุณาตุ เม ภนฺเต สํโฆ, อมฺหากํ ภณฺฑน\-ชาตานํ กลหชาตานํ วิวาทาปนฺนานํ วิหรตํ พหุํ อสฺสามณกํ อชฺฌาจิณฺณํ ภาสิต\-ปริกฺกนฺตํ, สเจ มยํ อิมาหิ อาปตฺตีหิ อญฺญมญฺญํ กาเรสฺสาม, สิยาปิ \csromanfolio{209}ตํ อธิกรณํ กกฺขฬตฺตาย วาฬตฺตาย เภทาย สํวตฺเตยฺย, ยทิ สํฆสฺส ปตฺตกลฺลํ, สํโฆ อิมํ อธิกรณํ ติณวตฺถารเกน วูปสเมยฺย ฐเปตฺวา ถุลฺลวชฺชํ ฐเปตฺวา คิหิปฺปฏิสํยุตฺตนฺ”ติ. เอกโต\-ปกฺขิกานํ ภิกฺขูนํ พฺยตฺเตน ภิกฺขุนา ปฏิพเลน สโก ปกฺโข ญาเปตพฺโพ\csromandash{}}

\prose{“สุณนฺตุ เม อายสฺมนฺตา, อมฺหากํ ภณฺฑน\-ชาตานํ กลหชาตานํ วิวาทาปนฺนานํ วิหรตํ พหุํ อสฺสามณกํ อชฺฌาจิณฺณํ ภาสิต\-ปริกฺกนฺตํ, สเจ มยํ อิมาหิ อาปตฺตีหิ อญฺญมญฺญํ กาเรสฺสาม, สิยาปิ ตํ อธิกรณํ กกฺขฬตฺตาย วาฬตฺตาย เภทาย สํวตฺเตยฺย, ยทา\-ยสฺมนฺตานํ ปตฺตกลฺลํ, อหํ ยา เจว อายสฺมนฺตานํ อาปตฺติ, ยา จ อตฺตโน อาปตฺติ, อายสฺมนฺตาน\-ญฺเจว อตฺถาย อตฺตโน จ อตฺถาย สํฆมชฺเฌ ติณวตฺถารเกน เทเสยฺยํ ฐเปตฺวา ถุลฺลวชฺชํ, ฐเปตฺวา คิหิปฺปฏิสํยุตฺตนฺ”ติ.}

\prose{อถาปเรสํ เอกโต\-ปกฺขิกานํ ภิกฺขูนํ พฺยตฺเตน ภิกฺขุนา ปฏิพเลน สโก ปกฺโข ญาเปตพฺโพ\csromandash{}}

\prose{“สุณนฺตุ เม อายสฺมนฺตา, อมฺหากํ ภณฺฑน\-ชาตานํ กลหชาตานํ วิวาทาปนฺนานํ วิหรตํ พหุํ อสฺสามณกํ อชฺฌาจิณฺณํ ภาสิต\-ปริกฺกนฺตํ, สเจ มยํ อิมาหิ อาปตฺตีหิ อญฺญมญฺญํ กาเรสฺสาม, สิยาปิ ตํ อธิกรณํ กกฺขฬตฺตาย วาฬตฺตาย เภทาย สํวตฺเตยฺย, ยทา\-ยสฺมนฺตานํ ปตฺตกลฺลํ อหํ ยา เจว อายสฺมนฺตานํ อาปตฺติ, ยา จ อตฺตโน อาปตฺติ, อายสฺมนฺตาน\-ญฺเจว อตฺถาย อตฺตโน จ อตฺถาย สํฆมชฺเฌ ติณวตฺถารเกน เทเสยฺยํ ฐเปตฺวา ถุลฺลวชฺชํ ฐเปตฺวา คิหิปฺปฏิสํยุตฺตนฺ”ติ.}

\proseitem{213}{อถาปเรสํ เอกโต\-ปกฺขิกานํ ภิกฺขูนํ พฺยตฺเตน ภิกฺขุนา ปฏิพเลน สํโฆ ญาเปตพฺโพ\csromandash{}}

\prose{“สุณาตุ เม ภนฺเต สํโฆ, อมฺหากํ ภณฺฑน\-ชาตานํ กลหชาตานํ วิวาทาปนฺนานํ วิหรตํ พหุํ อสฺสามณกํ อชฺฌาจิณฺณํ ภาสิต\-ปริกฺกนฺตํ, สเจ มยํ อิมาหิ อาปตฺตีหิ อญฺญมญฺญํ กาเรสฺสาม, สิยาปิ ตํ อธิกรณํ กกฺขฬตฺตาย วาฬตฺตาย เภทาย สํวตฺเตยฺย, ยทิ สํฆสฺส ปตฺตกลฺลํ อหํ ยา เจว อิเมสํ อายสฺมนฺตานํ อาปตฺติ, ยา จ อตฺตโน อาปตฺติ, อิเมสญฺเจว อายสฺมนฺตานํ อตฺถาย อตฺตโน \csromanfolio{210}จ อตฺถาย สํฆมชฺเฌ ติณวตฺถารเกน เทเสยฺยํ ฐเปตฺวา ถุลฺลวชฺชํ, ฐเปตฺวา คิหิปฺปฏิสํยุตฺตํ, เอสา ญตฺติ.}

\prose{สุณาตุ เม ภนฺเต สํโฆ, อมฺหากํ ภณฺฑน\-ชาตานํ กลหชาตานํ วิวาทาปนฺนานํ วิหรตํ พหุํ อสฺสามณกํ อชฺฌาจิณฺณํ ภาสิต\-ปริกฺกนฺตํ, สเจ มยํ อิมาหิ อาปตฺตีหิ อญฺญมญฺญํ กาเรสฺสาม, สิยาปิ ตํ อธิกรณํ กกฺขฬตฺตาย วาฬตฺตาย เภทาย สํวตฺเตยฺย, อหํ ยา เจว อิเมสํ อายสฺมนฺตานํ อาปตฺติ, ยา จ อตฺตโน อาปตฺติ, อิเมสญฺเจว อายสฺมนฺตานํ อตฺถาย อตฺตโน จ อตฺถาย สํฆมชฺเฌ ติณวตฺถารเกน เทเสมิ ฐเปตฺวา ถุลฺลวชฺชํ, ฐเปตฺวา คิหิปฺปฏิสํยุตฺตํ, ยสฺสายสฺมโต ขมติ อมฺหากํ อิมาสํ อาปตฺตีนํ สํฆมชฺเฌ ติณวตฺถารเกน เทสนา ฐเปตฺวา ถุลฺลวชฺชํ, ฐเปตฺวา คิหิปฺปฏิสํยุตฺตํ, โส ตุณฺหสฺส, ยสฺส นกฺขมติ, โส ภาเสยฺย.}

\prose{เทสิตา อมฺหากํ อิมา อาปตฺติโย สํฆมชฺเฌ ติณวตฺถารเกน ฐเปตฺวา ถุลฺลวชฺชํ, ฐเปตฺวา คิหิปฺปฏิสํยุตฺตํ, ขมติ สํฆสฺส, ตสฺมา ตุณฺหี, เอวเมตํ ธารยามี”ติ.}

\proseitem{214}{อถาปเรสํ เอกโต\-ปกฺขิกานํ ภิกฺขูนํ พฺยตฺเตน ภิกฺขุนา ปฏิพเลน สํโฆ ญาเปตพฺโพ\csromandash{}}

\prose{“สุณาตุ เม ภนฺเต สํโฆ, อมฺหากํ ภณฺฑน\-ชาตานํ กลหชาตานํ วิวาทาปนฺนานํ วิหรตํ พหุํ อสฺสามณกํ อชฺฌาจิณฺณํ ภาสิต\-ปริกฺกนฺตํ, สเจ มยํ อิมาหิ อาปตฺตีหิ อญฺญมญฺญํ กาเรสฺสาม, สิยาปิ ตํ อธิกรณํ กกฺขฬตฺตาย วาฬตฺตาย เภทาย สํวตฺเตยฺย, ยทิ สํฆสฺส ปตฺตกลฺลํ, อหํ ยา เจว อิเมสํ อายสฺมนฺตานํ อาปตฺติ, ยา จ อตฺตโน อาปตฺติ, อิเมสญฺเจว อายสฺมนฺตานํ อตฺถาย อตฺตโน จ อตฺถาย สํฆมชฺเฌ ติณวตฺถารเกน เทเสยฺยํ ฐเปตฺวา ถุลฺลวชฺชํ, ฐเปตฺวา คิหิปฺปฏิสํยุตฺตํ, เอสา ญตฺติ.}

\prose{สุณาตุ เม ภนฺเต สํโฆ, อมฺหากํ ภณฺฑน\-ชาตานํ กลหชาตานํ วิวาทาปนฺนานํ วิหรตํ พหุํ อสฺสามณกํ อชฺฌาจิณฺณํ ภาสิต\-ปริกฺกนฺตํ, สเจ มยํ อิมาหิ อาปตฺตีหิ อญฺญมญฺญํ กาเรสฺสาม, สิยาปิ ตํ อธิกรณํ กกฺขฬตฺตาย วาฬตฺตาย เภทาย สํวตฺเตยฺย, อหํ ยา เจว อิเมสํ อายสฺมนฺตานํ อาปตฺติ, ยา จ อตฺตโน อาปตฺติ, อิเมสญฺเจว \csromanfolio{211}อายสฺมนฺตานํ อตฺถาย อตฺตโน จ อตฺถาย สํฆมชฺเฌ ติณวตฺถารเกน เทเสมิ ฐเปตฺวา ถุลฺลวชฺชํ, ฐเปตฺวา คิหิปฺปฏิสํยุตฺตํ, ยสฺสายสฺมโต ขมติ อมฺหากํ อิมาสํ อาปตฺตีนํ สํฆมชฺเฌ ติณวตฺถารเกน เทสนา ฐเปตฺวา ถุลฺลวชฺชํ, ฐเปตฺวา คิหิปฺปฏิสํยุตฺตํ, โส ตุณฺหสฺส, ยสฺส นกฺขมติ, โส ภาเสยฺย.}

\prose{เทสิตา อมฺหากํ อิมา อาปตฺติโย สํฆมชฺเฌ ติณวตฺถารเกน ฐเปตฺวา ถุลฺลวชฺชํ, ฐเปตฺวา คิหิปฺปฏิสํยุตฺตํ, ขมติ สํฆสฺส, ตสฺมา ตุณฺหี, เอวเมตํ ธารยามี”ติ.}

\prose{เอวญฺจ ปน ภิกฺขเว เต ภิกฺขู ตาหิ อาปตฺตีหิ วุฏฺฐิตา โหนฺติ ฐเปตฺวา ถุลฺลวชฺชํ, ฐเปตฺวา คิหิปฺปฏิสํยุตฺตํ ฐเปตฺวา ทิฏฺฐาวิกมฺมํ ฐเปตฺวา เย น ตตฺถ โหนฺตีติ.}

\csromanmatikaanchor{mtk.115}
\csromantocmarkref{section}{8. อธิกรณ}{mtk.115}
\bookmark[dest={mtk.115},level=1]{8. อธิกรณ}
\csromanheader{1}{8. \textbf{อธิกรณ}}
\proseitem{215}{เตน โข ปน สมเยน (ภิกฺขูปิ ภิกฺขูหิ วิวทนฺติ)\csromansharedfootnote{48ea464a4f1db7c7}{( ) นตฺถิ (สี, สฺยา, กํ)} ภิกฺขูปิ ภิกฺขุนีหิ วิวทนฺติ, ภิกฺขุนิโยปิ ภิกฺขูหิ วิวทนฺติ, ฉนฺโนปิ ภิกฺขุ ภิกฺขุนีนํ อนุปขชฺช ภิกฺขูหิ สทฺธิํ วิวทติ, ภิกฺขุนีนํ ปกฺขํ คาเหติ. เย เต ภิกฺขู อปฺปิจฺฉา ฯเปฯ เต อุชฺฌายนฺติ ขิยฺยนฺติ วิปาเจนฺติ “กถํ หิ นาม ฉนฺโน ภิกฺขุ ภิกฺขุนีนํ อนุปขชฺช ภิกฺขูหิ สทฺธิํ วิวทิสฺสติ, ภิกฺขุนีนํ ปกฺขํ คาเหสฺสตี”ติ. อถ โข เต ภิกฺขู ภควโต เอตมตฺถํ อาโรเจสุํ ฯเปฯ. สจฺจํ กิร ภิกฺขเว ฯเปฯ. สจฺจํ ภควาติ ฯเปฯ วิครหิตฺวา ฯเปฯ ธมฺมิํ กถํ กตฺวา ภิกฺขู อามนฺเตสิ\csromandash{}}

\prose{\csromansymbolfootnote{*}{วิ 5. 169 ปิฏฺเฐปิ. 2. กิจฺจาธิกรณญฺจ (ก)}จตฺตาริมานิ ภิกฺขเว อธิกรณานิ วิวาทา\-ธิกรณํ อนุวาทา\-ธิกรณํ อาปตฺตา\-ธิกรณํ กิจฺจา\-ธิกรณํ2.}

\prose{ตตฺถ กตมํ วิวาทา\-ธิกรณํ.\csromansymbolfootnote{+}{วิ 5. 202 ปิฏฺเฐปิ.}อิธ ปน ภิกฺขเว ภิกฺขู\csromansharedfootnote{0f5a368cb025ebb5}{ภิกฺขู ภิกฺขุํ (ก)} วิวทนฺติ ธมฺโมติ วา อธมฺโมติ วา, วินโยติ วา อวินโยติ วา, ภาสิตํ ลปิตํ ตถาคเตนาติ วา อภาสิตํ อลปิตํ ตถาคเตนาติ วา, อาจิณฺณํ ตถาคเตนาติ วา อนาจิณฺณํ ตถาคเตนาติ วา, ปญฺญตฺตํ ตถาคเตนาติ วา อปฺปญฺญตฺตํ ตถาคเตนาติ วา, อาปตฺตีติ วา อนาปตฺตีติ \csromanfolio{212}วา, ลหุกา อาปตฺตีติ วา ครุกา อาปตฺตีติ วา, สาวเสสา อาปตฺตีติ วา อนวเสสา อาปตฺตีติ วา, ทุฏฺฐุลฺลา อาปตฺตีติ วา อทุฏฺฐุลฺลา อาปตฺตีติ วา. ยํ ตตฺถ ภณฺฑนํ กลโห วิคฺคโห วิวาโท นานาวาโท อญฺญถาวาโท วิปจฺจตาย โวหาโร เมธคํ, อิทํ วุจฺจติ วิวาทา\-ธิกรณํ.}

\prose{ตตฺถ กตมํ อนุวาทา\-ธิกรณํ.\csromansymbolfootnote{*}{วิ 5. 203 ปิฏฺเฐปิ.}อิธ ปน ภิกฺขเว ภิกฺขู ภิกฺขุํ อนุวทนฺติ สีลวิปตฺติยา วา อาจารวิปตฺติยา วา ทิฏฺฐิ\-วิปตฺติยา วา อาชีววิปตฺติยา วา. โย ตตฺถ อนุวาโท อนุวทนา อนุลฺลปนา อนุภณนา อนุสมฺปวงฺกตา อพฺภุสฺสหนตา อนุพล\-ปฺปทานํ, อิทํ วุจฺจติ อนุวาทา\-ธิกรณํ.}

\prose{ตตฺถ กตมํ อาปตฺตา\-ธิกรณํ. ปญฺจปิ อาปตฺติกฺขนฺธา อาปตฺตา\-ธิกรณํ, สตฺตปิ อาปตฺติกฺขนฺธา อาปตฺตา\-ธิกรณํ, อิทํ วุจฺจติ อาปตฺตา\-ธิกรณํ.}

\prose{ตตฺถ กตมํ กิจฺจา\-ธิกรณํ.\csromansymbolfootnote{+}{วิ 5. 204 ปิฏฺเฐปิ. ++ วิ 5. 165; อํ 2. 294; ม 3. 33 ปิฏฺเฐสุปิ.}ยา สํฆสฺส กิจฺจยตา กรณียตา อปโลกน\-กมฺมํ ญตฺติกมฺมํ ญตฺติ\-ทุติย\-กมฺมํ ญตฺติ\-จตุตฺถ\-กมฺมํ, อิทํ วุจฺจติ กิจฺจา\-ธิกรณํ.}

\proseitem{216}{วิวาทา\-ธิกรณสฺส กิํ มูลํ. ฉ วิวาทมูลานิ วิวาทาหิกรณสฺส มูลํ, ตีณิปิ อกุสลมูลนิ, วิวาทา\-ธิกรณสฺส มูลํ, ตีณิปิ กุสลมูลานิ วิวาทา\-ธิกรณสฺส มูลํ. ++ กตมานิ ฉ วิวาทมูลานิ วิวาทา\-ธิกรณสฺส มูลํ. อิธ ปน ภิกฺขเว ภิกฺขุ โกธโน โหติ อุปนาหี, โย โส ภิกฺขเว ภิกฺขุ โกธโน โหติ อุปนาหี, โส สตฺถริปิ อคารโว วิหรติ อปฺปติสฺโส, ธมฺเมปิ อคารโว วิหรติ อปฺปติสฺโส, สํเฆปิ อคารโว วิหรติ อปฺปติสฺโส, สิกฺขายปิ น ปริปูรการี โหติ, โย โส ภิกฺขเว ภิกฺขุ สตฺถริปิ อคารโว วิหรติ อปฺปติสฺโส, ธมฺเมปิ ฯเปฯ สํเฆปิ ฯเปฯ สิกฺขายปิ น ปริปูรการี, โส สํเฆ วิวาทํ ชเนติ, โย โหติ วิวาโท พหุชนาหิตาย พหุชนา\-สุขาย พหุโน ชนสฺส อนตฺถาย อหิตาย ทุกฺขาย เทวมนุสฺสานํ, เอวรูปญฺเจ ตุเมฺห ภิกฺขเว วิวาทมูลํ อชฺฌตฺตํ วา พหิทฺธา วา สมนุปสฺเสยฺยาถ, ตตฺร ตุเมฺห ภิกฺขเว ตสฺเสว ปาปกสฺส วิวาทมูลสฺส ปหานาย วายเมยฺยาถ, เอวรูปญฺเจ ตุเมฺห ภิกฺขเว วิวาทมูลํ อชฺฌตฺตํ วา พหิทฺธา วา น สมนุปสฺเสยฺยาถ, ตตฺร ตุเมฺห ภิกฺขเว ตสฺเสว ปาปกสฺส \csromanfolio{213}วิวาทมูลสฺส อายติํ อนวสฺสวาย ปฏิปชฺเชยฺยาถ, เอวเมตสฺส ปาปกสฺส วิวาทมูลสฺส ปหานํ โหติ, เอวเมตสฺส ปาปกสฺส วิวาทมูลสฺส อายติํ อนวสฺสโว โหติ.}

\prose{\csromansymbolfootnote{*}{วิ 5. 166 ปิฏฺเฐปิ.}ปุน จปรํ ภิกฺขเว ภิกฺขุ มกฺขี โหติ ปฬาสี ฯเปฯ อิสฺสุกี โหติ มจฺฉรี. สโฐ โหติ มายาวี. ปาปิจฺโฉ โหติ มิจฺฉาทิฏฺฐี. สนฺทิฏฺฐิ\-ปรามาสี โหติ อาธานคฺคาหี ทุปฺปฏินิสฺสคฺคี, โย โส ภิกฺขเว ภิกฺขุ สนฺทิฏฺฐิ\-ปรามาสี โหติ อาธานคฺคาหี ทุปฺปฏินิสฺสคฺคี, โส สตฺถริปิ อคารโว วิหรติ อปฺปติสฺโส, ธมฺเมปิ อคารโว วิหรติ อปฺปติสฺโส, สํเฆปิ อคารโว วิหรติ อปฺปติสฺโส, สิกฺขายปิ น ปริปูรการี โหติ, โย โส ภิกฺขเว ภิกฺขุ สตฺถริปิ อคารโว วิหรติ อปฺปติสฺโส, ธมฺเมปิ ฯเปฯ สํเฆปิ ฯเปฯ สิกฺขายปิ น ปริปูรการี, โส สํเฆ วิวาทํ ชเนติ, โย โหติ วิวาโท พหุชนาหิตาย พหุชนา\-สุขาย พหุโน ชนสฺส อนตฺถาย อหิตาย ทุกฺขาย เทวมนุสฺสานํ, เอวรูปญฺเจ ตุเมฺห ภิกฺขเว วิวาทมูลํ อชฺฌตฺตํ วา พหิทฺธา วา สมนุปสฺเสยฺยาถ, ตตฺร ตุเมฺห ภิกฺขเว ตสฺเสว ปาปกสฺส วิวาทมูลสฺส ปหานาย วายเมยฺยาถ, เอวรูปญฺเจ ตุเมฺห ภิกฺขเว วิวาทมูลํ อชฺฌตฺตํ วา พหิทฺธา วา น สมนุปสฺเสยฺยาถ, ตตฺร ตุเมฺห ภิกฺขเว ตสฺเสว ปาปกสฺส วิวาทมูลสฺส อายติํ อนวสฺสวาย ปฏิปชฺเชยฺยาถ, เอวเมตสฺส ปาปกสฺส วิวาทมูลสฺส ปหานํ โหติ, เอวเมตสฺส ปาปกสฺส วิวาทมูลสฺส อายติํ อนวสฺสโว โหติ. อิมานิ ฉ วิวาทมูลานิ วิวาทา\-ธิกรณสฺส มูลํ.}

\prose{กตมานิ ตีณิ อกุสลมูลานิ วิวาทา\-ธิกรณสฺส มูลํ. อิธ ปน ภิกฺขเว ภิกฺขู ลุทฺธจิตฺตา วิวทนฺติ. ทุฏฺฐจิตฺตา วิวทนฺติ. มูฬฺหจิตฺตา วิวทนฺติ ธมฺโมติ วา อธมฺโมติ วา, วินโยติ วา อวินโยติ วา, ภาสิตํ ลปิตํ ตถาคเตนาติ วา อภาสิตํ อลปิตํ ตถาคเตนาติ วา, อาจิณฺณํ ตถาคเตนาติ วา อนาจิณฺณํ ตถาคเตนาติ วา, ปญฺญตฺตํ ตถาคเตนาติ วา อปฺปญฺญตฺตํ ตถาคเตนาติ วา, อาปตฺตีติ วา อนาปตฺตีติ วา, ลหุกา อาปตฺตีติ วา ครุกา อาปตฺตีติ วา, สาวเสสา อาปตฺตีติ วา อนวเสสา อาปตฺตีติ วา, ทุฏฺฐุลฺลา อาปตฺตีติ วา อทุฏฺฐุลฺลา อาปตฺตีติ วา. อิมานิ ตีณิ อกุสลมูลานิ วิวาทา\-ธิกรณสฺส มูลํ.}

\prose{\csromanfolio{214}กตมานิ ตีณิ กุสลมูลานิ วิวาทา\-ธิกรณสฺส มูลํ. อิธ ปน ภิกฺขเว ภิกฺขู อลุทฺธจิตฺตา วิวทนฺติ. อทุฏฺฐจิตฺตา วิวทนฺติ. อมูฬฺหจิตฺตา วิวทนฺติ ธมฺโมติ วา อธมฺโมติ วา ฯเปฯ ทุฏฺฐุลฺลา อาปตฺตีติ วา อทุฏฺฐุลฺลา อาปตฺตีติ วา. อิมานิ ตีณี กุสลมูลานิ วิวาทาธิกรสฺส มูลํ.}

\proseitem{217}{\csromansymbolfootnote{*}{วิ 5. 166 ปิฏฺเฐปิ.}อนุวาทา\-ธิกรณสฺส กิํ มูลํ. ฉ อนุวาทมูลานิ อนุวาทา\-ธิกรณสฺส มูลํ, ตีณิปิ อกุสลมูลานิ อนุวาทา\-ธิกรณสฺส มูลํ, ตีณิปิ กุสลมูลานิ อนุวาทา\-ธิกรณสฺส มูลํ, กาโยปิ อนุวาทา\-ธิกรณสฺส มูลํ, วาจาปิ อนุวาทา\-ธิกรณสฺส มูลํ. กตมานิ ฉ อนุวาทมูลานิ อนุวาทา\-ธิกรณสฺส มูลํ. อิธ ปน ภิกฺขเว ภิกฺขุ โกธโน โหติ อุปนาหี, โย โส ภิกฺขเว ภิกฺขุ โกธโน โหติ อุปนาหี, โส สตฺถริปิ อคารโว วิหรติ อปฺปติสฺโส, ธมฺเมปิ อคารโว วิหรติ อปฺปติสฺโส, สํเฆปิ อคารโว วิหรติ อปฺปติสฺโส, สิกฺขายปิ น ปริปูรการี โหติ, โย โส ภิกฺขเว ภิกฺขุ สตฺถริปิ อคารโว วิหรติ อปฺปติสฺโส, ธมฺเมปิ ฯเปฯ สํเฆปิ ฯเปฯ สิกฺขายปิ น ปริปูรการี, โส สํเฆ อนุวาทํ ชเนติ, โย โหติ อนุวาโท พหุชนาหิตาย พหุชนา\-สุขาย พหุโน ชนสฺส อนตฺถาย อหิตาย ทุกฺขาย เทวมนุสฺสานํ, เอวรูปญฺเจ ตุเมฺห ภิกฺขเว อนุวาทมูลํ อชฺฌตฺตํ วา พหิทฺธา วา สมนุปสฺเสยฺยาถ, ตตฺร ตุเมฺห ภิกฺขเว ตสฺเสว ปาปกสฺส อนุวาทมูลสฺส ปหานาย วายเมยฺยาถ, เอวรูปญฺเจ ตุเมฺห ภิกฺขเว อนุวาทมูลํ อชฺฌตฺตํ วา พหิทฺธา วา น สมนุปสฺเสยฺยาถ, ตตฺร ตุเมฺห ภิกฺขเว ตสฺเสว ปาปกสฺส อนุวาทมูลสฺส อายติํ อนวสฺสวาย ปฏิปชฺเชยฺยาถ, เอวเมตสฺส ปาปกสฺส อนุวาทมูลสฺส ปหานํ โหติ, เอวเมตสฺส ปาปกสฺส อนุวาทมูลสฺส อายติํ อนวสฺสโว โหติ.}

\prose{ปุน จปรํ ภิกฺขเว ภิกฺขุ มกฺขี โหติ ปฬาสี ฯเปฯ อิสฺสุกี โหติ มจฺฉรี. สโฐ โหติ มายาวี. ปาปิจฺโฉ โหติ มิจฺฉาทิฏฺฐี. สนฺทิฏฺฐิ\-ปรามาสี โหติ อาธานคฺคาหี ทุปฺปฏินิสฺสคฺคี, โย โส ภิกฺขเว ภิกฺขุ สนฺทิฏฺฐิ\-ปรามาสี โหติ อาธานคฺคาหี ทุปฺปฏินิสฺสคฺคี, โส สตฺถริปิ อคารโว วิหรติ อปฺปติสฺโส, ธมฺเมปิ อคารโว วิหรติ อปฺปติสฺโส, สํเฆปิ อคารโว วิหรติ อปฺปติสฺโส, สิกฺขายปิ น ปริปูรการี โหติ, โย โส ภิกฺขเว ภิกฺขุ สตฺถริปิ อคารโว วิหรติ \csromanfolio{215}อปฺปติสฺโส, ธมฺเมปิ ฯเปฯ สํเฆปิ ฯเปฯ สิกฺขายปิ น ปริปูรการี, โส สํเฆ อนุวาทํ ชเนติ, โย โหติ อนุวาโท พหุชนาหิตาย พหุชนา\-สุขาย พหุโน ชนสฺส อนตฺถาย อหิตาย ทุกฺขาย เทวมนุสฺสานํ. เอวรูปญฺเจ ตุเมฺห ภิกฺขเว อนุวาทมูลํ อชฺฌตฺตํ วา พหิทฺธา วา สมนุปสฺเสยฺยาถ, ตตฺร ตุเมฺห ภิกฺขเว ตสฺเสว ปาปกสฺส อนุวาทมูลสฺส ปหานาย วายเมยฺยาถ, เอวรุปญฺเจ ตุเมฺห ภิกฺขเว อนุวาทมูลํ อชฺฌตฺตํ วา พหิทฺธา วา น สมนุปสฺเสยฺยาถ, ตตฺร ตุเมฺห ภิกฺขเว ตสฺเสว ปาปกสฺส อนุวาทมูลสฺส อายติํ อนวสฺสวาย ปฏิปชฺเชยฺยาถ, เอวเมตสฺส ปาปกสฺส อนุวาทมูลสฺส ปหานํ โหติ, เอวเมตสฺส ปาปกสฺส อนุวาทมูลสฺส อายติํ อนวสฺสโว โหติ. อิมานิ ฉ อนุวาทมูลานิ อนุวาทา\-ธิกรณสฺส มูลํ.}

\prose{กตมานิ ตีณิ อกุสลมูลานิ อนุวาทา\-ธิกรณสฺส มูลํ. อิธ ปน ภิกฺขเว ภิกฺขู ภิกฺขุํ ลุทฺธจิตฺตา อนุวทนฺติ. ทุฏฺฐจิตฺตา อนุวทนฺติ. มูฬฺหจิตฺตา อนุวทนฺติ สีลวิปตฺติยา วา อาจารวิปตฺติยา วา ทิฏฺฐิ\-วิปตฺติยา วา อาชีววิปตฺติยา วา. อิมานิ ตีณิ อกุสลมูลานิ อนุวาทา\-ธิกรณสฺส มูลํ.}

\prose{กตมานิ ตีณิ กุสลมูลานิ อนุวาทา\-ธิกรณสฺส มูลํ. อิธ ปน ภิกฺขเว ภิกฺขู ภิกฺขุํ อลุทฺธจิตฺตา อนุวทนฺติ. อทุฏฺฐจิตฺตา อนุวทนฺติ. อมูฬฺหจิตฺตา อนุวทนฺติ สีลวิปตฺติยา วา อาจารวิปตฺติยา วา ทิฏฺฐิ\-วิปตฺติยา วา อาชีววิปตฺติยา วา. อิมานิ ตีณิ กุสลมูลานิ อนุวาทา\-ธิกรณสฺส มูลํ.}

\prose{กตโม กาโย\csromansharedfootnote{1b863bcb5944e7ad}{กตโม จ กาโย (สฺยา, กํ)} อนุวาทา\-ธิกรณสฺส มูลํ. อิเธกจฺโจ ทุพฺพณฺโณ โหติ ทุทฺทสฺสิโก โอโกฏิมโก พหฺวาพาโธ กาโณ วา กุณี วา ขญฺโช วา ปกฺขหโต วา, เยน นํ อนุวทนฺติ, อยํ กาโย อนุวาทา\-ธิกรณสฺส มูลํ.}

\prose{กตมา วาจา อนุวาทา\-ธิกรณสฺส มูลํ. อิเธกจฺโจ ทุพฺพโจ โหติ มมฺมโน เอฬคลวาโจ, ยาย นํ อนุวทนฺติ, อยํ วาจา อนุวาทา\-ธิกรณสฺส มูลํ.}

\proseitem{218}{อาปตฺตา\-ธิกรณสฺส กิํ มูลํ. ฉ อาปตฺติ\-สมุฏฺฐานา อาปตฺตา\-ธิกรณสฺส มูลํ. อตฺถาปตฺติ กายโต สมุฏฺฐาติ น วาจโต \csromanfolio{216}น จิตฺตโต, อตฺถาปตฺติ วาจโต สมุฏฺฐาติ น กายโต น จิตฺตโต, อตฺถาปตฺติ กายโต จ วาจโต จ สมุฏฺฐาติ น จิตฺตโต, อตฺถาปตฺติ กายโต จ จิตฺตโต จ สมุฏฺฐาติ น วาจโต, อตฺถาปตฺติ วาจโต จ จิตฺตโต จ สมุฏฺฐาติ น กายโต, อตฺถาปตฺติ กายโต จ วาจโต จ จิตฺตโต จ สมุฏฺฐาติ. อิเม ฉ อาปตฺติ\-สมุฏฺฐานา อาปตฺตา\-ธิกรณสฺส มูลํ.}

\proseitem{219}{กิจฺจา\-ธิกรณสฺส กิํ มูลํ. กิจฺจา\-ธิกรณสฺส เอกํ มูลํ สํโฆ.}

\proseitem{220}{วิวาทา\-ธิกรณํ กุสลํ อกุสลํ อพฺยากตํ? วิวาทา\-ธิกรณํ สิยา กุสลํ สิยา อกุสลํ สิยา อพฺยากตํ. ตตฺถ กตมํ วิวาทา\-ธิกรณํ กุสลํ. อิธ ปน ภิกฺขเว ภิกฺขู กุสลจิตฺตา วิวทนฺติ ธมฺโมติ วา อธมฺโมติ วา ฯเปฯ ทุฏฺฐุลฺลา อาปตฺตีติ วา อทุฏฺฐุลฺลา อาปตฺตีติ วา. ยํ ตตฺถ ภณฺฑนํ กลโห วิคฺคโห วิวาโท นานาวาโท อญฺญถาวาโท วิปจฺจตาย โวหาโร เมธคํ, อิทํ วุจฺจติ วิวาทา\-ธิกรณํ กุสลํ.}

\prose{ตตฺถ กตมํ วิวาทา\-ธิกรณํ อกุสลํ. อิธ ปน ภิกฺขเว ภิกฺขู อกุสลจิตฺตา วิวทนฺติ ธมฺโมติ วา อธมฺโมติ วา ฯเปฯ ทุฏฺฐุลฺลา อาปตฺตีติ วา อทุฏฺฐุลฺลา อาปตฺตีติ วา. ยํ ตตฺถ ภณฺฑนํ กลโห วิคฺคโห วิวาโท นานาวาโท อญฺญถาวาโท วิปจฺจตาย โวหาโร เมธคํ, อิทํ วุจฺจติ วิวาทา\-ธิกรณํ อกุสลํ.}

\prose{ตตฺถ กตมํ วิวาทา\-ธิกรณํ อพฺยากตํ. อิธ ปน ภิกฺขเว ภิกฺขู อพฺยากตจิตฺตา วิวทนฺติ ธมฺโมติ วา อธมฺโมติ วา ฯเปฯ ทุฏฺฐุลฺลา อาปตฺตีติ วา อทุฏฺฐุลฺลา อาปตฺตีติ วา. ยํ ตตฺถ ภณฺฑนํ กลโห วิคฺคโห วิวาโท นานาวาโท อญฺญถาวาโท วิปจฺจตาย โวหาโร เมธคํ, อิทํ วุจฺจติ วิวาทา\-ธิกรณํ อพฺยากตํ.}

\proseitem{221}{อนุวาทา\-ธิกรณํ กุสลํ อกุสลํ อพฺยากตํ? อนุวาทา\-ธิกรณํ สิยา กุสลํ สิยา อกุสลํ สิยา อพฺยากตํ. ตตฺถ กตมํ อนุวาทา\-ธิกรณํ กุสลํ. อิธ ปน ภิกฺขเว ภิกฺขู ภิกฺขุํ กุสลจิตฺตา อนุวทนฺติ สีลวิปตฺติยา วา อาจารวิปตฺติยา วา ทิฏฺฐิ\-วิปตฺติยา วา อาชีววิปตฺติยา วา. โย ตตฺถ อนุวาโท อนุวทนา อนุลฺลปนา \csromanfolio{217}อนุภณนา อนุสมฺปวงฺกตา อพฺภุสฺสหนตา อนุพล\-ปฺปทานํ, อิทํ วุจฺจติ อนุวาทา\-ธิกรณํ กุสลํ.}

\prose{ตตฺถ กตมํ อนุวาทา\-ธิกรณํ อกุสลํ. อิธ ปน ภิกฺขเว ภิกฺขู ภิกฺขุํ อกุสลจิตฺตา อนุวทนฺติ สีลวิปตฺติยา วา อาจารวิปตฺติยา วา ทิฏฺฐิ\-วิปตฺติยา วา อาชีววิปตฺติยา วา. โย ตตฺถ อนุวาโท อนุวทนา อนุลฺลปนา อนุภณนา อนุสมฺปวงฺกตา อพฺภุสฺสหนตา อนุพล\-ปฺปทานํ, อิทํ วุจฺจติ อนุวาทา\-ธิกรณํ อกุสลํ.}

\prose{ตตฺถ กตมํ อนุวาทา\-ธิกรณํ อพฺยากตํ. อิธ ปน ภิกฺขเว ภิกฺขู ภิกฺขุํ อพฺยากตจิตฺตา อนุวทนฺติ สีลวิปตฺติยา วา อาจารวิปตฺติยา วา ทิฏฺฐิ\-วิปตฺติยา วา อาชีววิปตฺติยา วา. โย ตตฺถ อนุวาโท อนุวทนา อนุลฺลปนา อนุภณนา อนุสมฺปวงฺกตา อพฺภุสฺสหนตา อนุพล\-ปฺปทานํ, อิทํ วุจฺจติ อนุวาทา\-ธิกรณํ อพฺยากตํ.}

\proseitem{222}{อาปตฺตา\-ธิกรณํ กุสลํ\csromansharedfootnote{1572d5858345b76b}{อิทํ ปทํ เกสุจิ นตฺถิ.} อกุสลํ อพฺยากตํ? อาปตฺตา\-ธิกรณํ สิยา อกุสลํ สิยา อพฺยากตํ, นตฺถิ อาปตฺตา\-ธิกรณํ กุสลํ. ตตฺถ กตมํ อาปตฺตา\-ธิกรณํ อกุสลํ. ยํ ชานนฺโต สญฺชานนฺโต เจจฺจ อภิวิตริตฺวา วิติกฺกโม, อิทํ วุจฺจติ อาปตฺตา\-ธิกรณํ อกุสลํ.}

\prose{ตตฺถ กตมํ อาปตฺตา\-ธิกรณํ อพฺยากตํ. ยํ อชานนฺโต อสญฺชานนฺโต อเจจฺจ อนภิวิตริตฺวา วีติกฺกโม, อิทํ วุจฺจติ อาปตฺตา\-ธิกรณํ อพฺยากตํ.}

\proseitem{223}{กิจฺจา\-ธิกรณํ กุสลํ อกุสลํ อพฺยากตํ? กิจฺจา\-ธิกรณํ สิยา กุสลํ สิยา อกุสลํ สิยา อพฺยากตํ. ตตฺถ กตมํ กิจฺจา\-ธิกรณํ กุสลํ. ยํ สํโฆ กุสลจิตฺโต กมฺมํ กโรติ อปโลกน\-กมฺมํ ญตฺติกมฺมํ ญตฺติ\-ทุติย\-กมฺมํ ญตฺติ\-จตุตฺถ\-กมฺมํ, อิทํ วุจฺจติ กิจฺจา\-ธิกรณํ กุสลํ.}

\prose{ตตฺถ กตมํ กิจฺจา\-ธิกรณํ อกุสลํ. ยํ สํโฆ อกุสลจิตฺโต กมฺมํ กโรติ อปโลกน\-กมฺมํ ญตฺติกมฺมํ ญตฺติ\-ทุติย\-กมฺมํ ญตฺติ\-จตุตฺถ\-กมฺมํ, อิทํ วุจฺจติ กิจฺจา\-ธิกรณํ อกุสลํ.}

\prose{\csromanfolio{218}ตตฺถ กตมํ กิจฺจา\-ธิกรณํ อพฺยากตํ. ยํ สํโฆ อพฺยากตจิตฺโต กมฺมํ กโรติ อปโลกน\-กมฺมํ ญตฺติกมฺมํ ญตฺติ\-ทุติย\-กมฺมํ ญตฺติ\-จตุตฺถ\-กมฺมํ, อิทํ วุจฺจติ กิจฺจา\-ธิกรณํ อพฺยากตํ.}

\proseitem{224}{\csromansymbolfootnote{*}{วิ 5. 275 ปิฏฺเฐปิ.}วิวาโท วิวาทา\-ธิกรณํ, วิวาโท โน อธิกรณํ, อธิกรณํ โน วิวาโท, อธิกรณ\-ญฺเจว วิวาโท จ. สิยา วิวาโท วิวาทา\-ธิกรณํ, สิยา วิวาโท โน อธิกรณํ, สิยา อธิกรณํ โน วิวาโท, สิยา อธิกรณ\-ญฺเจว วิวาโท จ.}

\prose{ตตฺถ กตโม วิวาโท วิวาทา\-ธิกรณํ. อิธ ปน ภิกฺขเว ภิกฺขู วิวทนฺติ ธมฺโมติ วา อธมฺโมติ วา ฯเปฯ ทุฏฺฐุลฺลา อาปตฺตีติ วา อทุฏฺฐุลฺลา อาปตฺตีติ วา. ยํ ตตฺถ ภณฺฑนํ กลโห วิคฺคโห วิวาโท นานาวาโท อญฺญถาวาโท วิปจฺจตาย โวหาโร เมธคํ, อยํ วิวาโท วิวาทามิกรณํ.}

\prose{ตตฺถ กตโม วิวาโท โน อธิกรณํ. มาตาปิ ปุตฺเตน วิวทติ, ปุตฺโตปิ มาตรา วิวทติ, ปิตาปิ ปุตฺเตน วิวทติ, ปุตฺโตปิ ปิตรา วิวทติ, ภาตาปิ ภาตรา วิวทติ, ภาตาปิ ภคินิยา วิวทติ, ภคินีปิ ภาตรา วิวทติ, สหาโยปิ สหาเยน วิวทติ, อยํ วิวาโท โน อธิกรณํ.}

\prose{ตตฺถ กตมํ อธิกรณํ โน วิวาโท. อนุวาทา\-ธิกรณํ อาปตฺตา\-ธิกรณํ กิจฺจา\-ธิกรณํ, อิทํ อธิกรณํ โน วิวาโท.}

\prose{ตตฺถ กตมํ อธิกรณ\-ญฺเจว วิวาโท จ. วิวาทา\-ธิกรณํ อธิกรณ\-ญฺเจว วิวาโท จ.}

\proseitem{225}{\csromansymbolfootnote{+}{วิ 5. 276 ปิฏฺเฐปิ.}อนุวาโท อนุวาทา\-ธิกรณํ, อนุวาโท โน อธิกรณํ, อธิกรณํ โน อนุวาโท, อธิกรณ\-ญฺเจว อนุวาโท จ. สิยา อนุวาโท อนุวาทา\-ธิกรณํ, สิยา อนุวาโท โน อธิกรณํ, สิยา อธิกรณํ โน อนุวาโท, สิยา อธิกรณ\-ญฺเจว อนุวาโท จ.}

\prose{ตตฺถ กตโม อนุวาโท อนุวาทา\-ธิกรณํ. อิธ ปน ภิกฺขเว ภิกฺขู ภิกฺขุํ อนุวทนฺติ สีลวิปตฺติยา วา อาจารวิปตฺติยา วา ทิฏฺฐิ\-วิปตฺติยา วา อาชีววิปตฺติยา วา. โย ตตฺถ อนุวาโท อนุวทนา อนุลฺลปนา \csromanfolio{219}อนุภณนา อนุสมฺปวงฺกตา อพฺภุสฺสหนตา อนุพล\-ปฺปทานํ, อยํ อนุวาโท อนุวาทา\-ธิกรณํ.}

\prose{ตตฺถ กตโม อนุวาโท โน อธิกรณํ. มาตาปิ ปุตฺตํ อนุวทติ, ปุตฺโตปิ มาตรํ อนุวทติ, ปิตาปิ ปุตฺตํ อนุวทติ, ปุตฺโตปิ ปิตรํ อนุวทติ, ภาตาปิ ภาตรํ อนุวทติ, ภาตาปิ ภคินิํ อนุวทติ, ภคินีปิ ภาตรํ อนุวทติ, สหาโยปิ สหายํ อนุวทติ, อยํ อนุวาโท โน อธิกรณํ.}

\prose{ตตฺถ กตมํ อธิกรณํ โน อนุวาโท. อาปตฺตา\-ธิกรณํ กิจฺจา\-ธิกรณํ วิเวทาธิกรณํ, อิทํ อธิกรณํ โน อนุวาโท.}

\prose{ตตฺถ กตมํ อธิกรณ\-ญฺเจว อนุวาโท จ. อนุวาทา\-ธิกรณํ อธิกรณ\-ญฺเจว อนุวาโท จ.}

\proseitem{226}{\csromansymbolfootnote{*}{วิ 5. 277 ปิฏฺเฐปิ.}อาปตฺติ อาปตฺตา\-ธิกรณํ, อาปตฺติ โน อธิกรณํ, อธิกรณํ โน อาปตฺติ, อธิกรณ\-ญฺเจว อาปตฺติ จ. สิยา อาปตฺติ อาปตฺตา\-ธิกรณํ, สิยา อาปตฺติ โน อธิกรณํ, สิยา อธิกรณํ โน อาปตฺติ, สิยา อธิกรณ\-ญฺเจว อาปตฺติ จ. ตตฺถ กตมํ อาปตฺติ อาปตฺตา\-ธิกรณํ. ปญฺจปิ อาปตฺติกฺขนฺธา อาปตฺตา\-ธิกรณํ, สตฺตปิ อาปตฺถิกฺขนฺธา อาปตฺตา\-ธิกรณํ, อยํ อาปตฺติ อาปตฺตา\-ธิกรณํ.}

\prose{ตตฺถ กตมํ อาปตฺติ โน อธิกรณํ. โสตาปตฺติ สมาปตฺติ, อยํ อาปตฺติ โน อธิกรณํ.}

\prose{ตตฺถ กตมํ อธิกรณํ โน อาปตฺติ. กิจฺจา\-ธิกรณํ วิวาทา\-ธิกรณํ อนุวาทา\-ธิกรณํ, อิทํ อธิกรณํ โน อาปตฺติ.}

\prose{ตตฺถ กตมํ อธิกรณ\-ญฺเจว อาปตฺติ จ. อาปตฺตา\-ธิกรณํ อธิกรณ\-ญฺเจว อาปตฺติ จ.}

\proseitem{227}{\csromansymbolmark{*}กิจฺจํ กิจฺจา\-ธิกรณํ, กิจฺจํ โน อธิกรณํ, อธิกรณํ โน กิจฺจํ, อธิกรณ\-ญฺเจว กิจฺจญฺจ. สิยา กิจฺจํ กิจฺจา\-ธิกรณํ, สิยา กิจฺจํ โน อธิกรณํ, สิยา อธิกรณํ โน กิจฺจํ, สิยา อธิกรณ\-ญฺเจว กิจฺจญฺจ.}

\prose{\csromanfolio{220}ตตฺถ กตมํ กิจฺจํ กิจฺจา\-ธิกรณํ. ยา สํฆสฺส กิจฺจยตา กรณียตา อปโลกน\-กมฺมํ ญตฺติกมฺมํ ญตฺติ\-ทุติย\-กมฺมํ ญตฺติ\-จตุตฺถ\-กมฺมํ, อิทํ กิจฺจํ กิจฺจา\-ธิกรณํ.}

\prose{ตตฺถ กตมํ กิจฺจํ โน อธิกรณํ. อาจริยกิจฺจํ อุปชฺฌาย\-กิจฺจํ สมานุ\-ปชฺฌาย\-กิจฺจํ สมานา\-จริย\-กิจฺจํ, อิทํ กิจฺจํ โน อธิกรณํ.}

\prose{ตตฺถ กตมํ อธิกรณํ โน กิจฺจํ, วิวาทา\-ธิกรณํ อนุวาทา\-ธิกรณํ อาปตฺตา\-ธิกรณํ, อิทํ อธิกรณํ โน กิจฺจํ.}

\prose{ตตฺถ กตมํ อธิกรณ\-ญฺเจว กิจฺจญฺจ. กิจฺจา\-ธิกรณํ อธิกรณ\-ญฺเจว กิจฺจญฺจ.}

\csromanmatikaanchor{mtk.116}
\csromantocmarkref{section}{9. อธิกรณวูปสมนสมถ}{mtk.116}
\bookmark[dest={mtk.116},level=1]{9. อธิกรณวูปสมนสมถ}
\csromanheader{1}{9. อธิกรณ\-วูปสมน\-สมถ}
\csromanmatikaanchor{mtk.117}
\csromantocmarkref{subsection}{สมฺมุขาวินย}{mtk.117}
\bookmark[dest={mtk.117},level=2]{สมฺมุขาวินย}
\csromanheader{2}{\textbf{สมฺมุขาวินย}}
\proseitem{228}{\csromansymbolfootnote{*}{วิ 5. 183, 195 ปิฏฺฐาทีสุปิ.}วิวาทา\-ธิกรณํ กติหิ สมเถหิ สมฺมติ. วิวาทา\-ธิกรณํ ทฺวีหิ สมเถหิ สมฺมติ สมฺมุขา\-วินเยน จ เยภุยฺยสิกาย จ. สิยา วิวาทา\-ธิกรณํ เอกํ สมถํ อนาคมฺม เยภุยฺยสิกํ เอเกน สมเถน สเมยฺย สมฺมุขา\-วินเยนาติ? สิยาติสฺส วจนียํ. ยถา กถํ วิย. อิธ ปน ภิกฺขเว ภิกฺขู วิวทนฺติ ธมฺโมติ วา อธมฺโมติ วา, วินโยติ วา อวินโยติ วา, ภาสิตํ ลปิตํ ตถาคเตนาติ วา อภาสิตํ อลปิตํ ตถาคเตนาติ วา, อาจิณฺณํ ตถาคเตนาติ วา อนาจิณฺณํ ตถาคเตนาติ วา, ปญฺญตฺตํ ตถาคเตนาติ วา อปญฺญตฺตํ ตถาคเตนาติ วา, อาปตฺตีติ วา อนาปตฺตีติ วา, ลหุกา อาปตฺตีติ วา ครุกา อาปตฺตีติ วา, สาวเสสา อาปตฺตีติ วา อนวเสสา อาปตฺตีติ วา, ทุฏฺฐุลฺลา อาปตฺตีติ วา อทุฏฺฐุลฺลา อาปตฺตีติ วา.}

\prose{เต เจ ภิกฺขเว ภิกฺขู สกฺโกนฺติ ตํ อธิกรณํ วูปสเมตุํ, อิทํ วุจฺจติ ภิกฺขเว อธิกรณํ วุปสนฺตํ. เกน วูปสนฺตํ, สมฺมุขา\-วินเยน. กิญฺจ ตตฺถ สมฺมุขา\-วินยสฺมิํ, สํฆสมฺมุขตา ธมฺม\-สมฺมุขตา วินย\-สมฺมุขตา ปุคฺคล\-สมฺมุขตา. กา จ ตตฺถ สํฆสมฺมุขตา, ยาวติกา ภิกฺขู กมฺมปฺปตฺตา, เต อาคตา โหนฺติ, ฉนฺทารหานํ ฉนฺโท อาหโฏ โหติ, สมฺมุขีภูตา น ปฏิกฺโกสนฺติ, อยํ ตตฺถ สํฆสมฺมุขตา. กา \csromanfolio{221}จ ตตฺถ ธมฺม\-สมฺมุขตา วินย\-สมฺมุขตา, เยน ธมฺเมน เยน วินเยน เยน สตฺถุสาสเนน ตํ อธิกรณํ วูปสมฺมติ, อยํ ตตฺถ ธมฺม\-สมฺมุขตา วินย\-สมฺมุขตา. กา จ ตตฺถ ปุคฺคล\-สมฺมุขตา, โย จ วิวทติ, เยน จ วิวทติ, อุโภ อตฺถ\-ปจฺจตฺถิกา สมฺมุขีภูตา โหนฺติ, อยํ ตตฺถ ปุคฺคล\-สมฺมุขตา. เอวํ วูปสนฺตํ เจ ภิกฺขเว อธิกรณํ การโก อุกฺโกเฏติ, อุกฺโกฏนกํ ปาจิตฺติยํ. ฉนฺททายโก ขียติ, ขียนกํ ปาจิตฺติยํ.}

\proseitem{229}{เต เจ ภิกฺขเว ภิกฺขู น สกฺโกนฺติ ตํ อธิกรณํ ตสฺมิํ อาวาเส วูปสเมตุํ, เตหิ ภิกฺขเว ภิกฺขูหิ ยสฺมิํ อาวาเส สมฺพหุลา\csromansharedfootnote{6524e48ba8b710f5}{พหุตรา (สี, สฺยา)} ภิกฺขู, โส อาวาโส คนฺตพฺโพ. เต เจ ภิกฺขเว ภิกฺขู ตํ อาวาสํ คจฺฉนฺตา อนฺตรามคฺเค สกฺโกนฺติ ตํ อธิกรณํ วูปสเมตุํ. อิทํ วุจฺจติ ภิกฺขเว อธิกรณํ วูปสนฺตํ. เกน วูปสนฺตํ, สมฺมุขา\-วินเยน. กิญฺจ ตตฺถ สมฺมุขา\-วินยสฺมิํ, สํฆสมฺมุขตา ธมฺม\-สมฺมุขตา วินย\-สมฺมุขตา ปุคฺคล\-สมฺมุขตา. กา จ ตตฺถ สํฆสมฺมุขตา, ยาวติกา ภิกฺขู กมฺมปฺปตฺตา, เต อาคตา โหนฺติ, ฉนฺทารหานํ ฉนฺโท อาหโฏ โหติ, สมฺมุขีภูตา น ปฏิกฺโกสนฺติ, อยํ ตตฺถ สํฆสมฺมุขตา. กา จ ตตฺถ ธมฺม\-สมฺมุขตา วินย\-สมฺมุขตา, เยน ธมฺเมน เยน วินเยน เยน สตฺถุสาสเนน ตํ อธิกรณํ วูปสมฺมติ, อยํ ตตฺถ ธมฺม\-สมฺมุขตา วินย\-สมฺมุขตา. กา จ ตตฺถ ปุคฺคล\-สมฺมุขตา, โย จ วิวทติ, เยน จ วิวทติ, อุโภ อตฺถ\-ปจฺจตฺถิกา สมฺมุขีภูตา โหนฺติ, อยํ ตตฺถ ปุคฺคล\-สมฺมุขตา. เอวํ วูปสนฺตํ เจ ภิกฺขเว อธิกรณํ การโก อุกฺโกเฏติ, อุกฺโกฏนกํ ปาจิตฺติยํ. ฉนฺททายโก ขียติ, ขียนกํ ปาจิตฺติยํ.}

\proseitem{230}{เต เจ ภิกฺขเว ภิกฺขู ตํ อาวาสํ อาคจฺฉนฺตา อนฺตรามคฺเค น สกฺโกนฺติ ตํ อธิกรณํ วูปสเมตุํ, เตหิ ภิกฺขเว ภิกฺขูหิ ตํ อาวาสํ คนฺตฺวา อาวาสิกา ภิกฺขู เอวมสฺสุ วจนียา “อิทํ โข อาวุโส อธิกรณํ เอวํ ชาตํ เอวํ สมุปฺปนฺนํ, สาธายสฺมนฺตา อิมํ อธิกรณํ วูปสเมนฺตุ ธมฺเมน วินเยน สตฺถุสาสเนน, ยถยิทํ อธิกรณํ สุวูปสนฺตํ อสฺสา”ติ.}

\prose{\csromanfolio{222}สเจ ภิกฺขเว อาวาสิกา ภิกฺขู วุฑฺฒตรา โหนฺติ, อาคนฺตุกา ภิกฺขู นวกตรา, เตหิ ภิกฺขเว อาวาสิเกหิ ภิกฺขูหิ อาคนฺตุกา ภิกฺขู เอวมสฺสุ วจนียา “อิงฺฆ ตุเมฺห อายสฺมนฺโต มุหุตฺตํ เอกมนฺตํ โหถ, ยาว มยํ มนฺเตมา”ติ. สเจ ปน ภิกฺขเว อาวาสิกา ภิกฺขู นวกตรา โหนฺติ, อาคนฺตุกา ภิกฺขู วุฑฺฒุตรา, เตหิ ภิกฺขเว อาวาสิเกหิ ภิกฺขูหิ อาคนฺตุกา ภิกฺขู เอวมสฺสุ วจนียา “เตน หิ ตุเมฺห อายสฺมนฺโต มุหุตฺตํ อิเธว ตาว โหถ, ยาว มยํ มนฺเตมา”ติ.}

\prose{สเจ ปน ภิกฺขเว อาวาสิกานํ ภิกฺขูนํ มนฺตยมานานํ เอวํ โหติ “น มยํ สกฺโกม อิมํ อธิกรณํ วูปสเมตุํ ธมฺเมน วินเยน สตฺถุสาสเนนา”ติ, น ตํ อธิกรณํ อาวาสิเกหิ ภิกฺขูหิ สมฺปฏิจฺฉิตพฺพํ. สเจ ปน ภิกฺขเว อาวาสิกานํ ภิกฺขูนํ มนฺตยมานานํ เอวํ โหติ “สกฺโกม มยํ อิมํ อธิกรณํ วูปสเมตุํ ธมฺเมน วินเยน สตฺถุสาสเนนา”ติ, เตหิ ภิกฺขเว อาวาสิเกหิ ภิกฺขูหิ อาคนฺตุกา ภิกฺขู เอวมสฺสุ วจนียา “สเจ ตุเมฺห อายสฺมนฺโต อมฺหากํ อิมํ อธิกรณํ ยถาชาตํ ยถา\-สมุปฺปนฺนํ อาโรเจสฺสถ, ยถา จ มยํ อิมํ อธิกรณํ วูปสเมสฺสาม ธมฺเมน วินเยน สตฺถุสาสเนน, ตถาสุ วูปสนฺตํ ภวิสฺสติ, เอวํ มยํ อิมํ อธิกรณํ สมฺปฏิจฺฉิสฺสาม. โน เจ ตุเมฺห อายสฺมนฺโต อมฺหากํ อิมํ อธิกรณํ ยถาชาตํ ยถา\-สมุปฺปนฺนํ อาโรเจสฺสถ, ยถา จ มยํ อิมํ อธิกรณํ วูปสเมสฺสาม ธมฺเมน วินเยน สตฺถุสาสเนน, ตถา น สุวูปสนฺตํ ภวิสฺสติ, น มยํ อิมํ อธิกรณํ สมฺปฏิจฺฉิสฺสามา”ติ. เอวํ สุปริคฺคหิตํ โข ภิกฺขเว กตฺวา อาวาสิเกหิ ภิกฺขูหิ ตํ อธิกรณํ สมฺปฏิจฺฉิตพฺพํ.}

\prose{เตหิ ภิกฺขเว อาคนฺตุเกหิ ภิกฺขูหิ อาวาสิกา ภิกฺขู เอวมสฺสุ วจนียา “ยถาชาตํ ยถา\-สมุปฺปนฺนํ มยํ อิมํ อธิกรณํ อายสฺมนฺตานํ อาโรเจสฺสาม. สเจ อายสฺมนฺตา สกฺโกนฺติ เอตฺตเกน วา เอตฺตเกน วา อนฺตเรน อิมํ อธิกรณํ วูปสเมตุํ ธมฺเมน วินเยน สตฺถุสาสเนน, ตถา สุวูปสนฺตํ ภวิสฺสติ, เอวํ มยํ อิมํ อธิกรณํ อายสฺมนฺตานํ นิยฺยาเทสฺสาม. โน เจ อายสฺมนฺตา สกฺโกนฺติ เอตฺตเกน วา เอตฺตเกน วา อนฺตเรน อิมํ อธิกรณํ วูปสเมตุํ ธมฺเมน วินเยน สตฺถุสาสเนน, ตถา น สุวูปสนฺตํ ภวิสฺสติ, น มยํ อิมํ อธิกรณํ อายสฺมนฺตานํ \csromanfolio{223}นิยฺยาเทสฺสาม, มยเมว อิมสฺส อธิกรณสฺส สามิโน ภวิสฺสามา”ติ. เอวํ สุปริคฺคหิตํ โข ภิกฺขเว กตฺวา อาคนฺตุเกหิ ภิกฺขูหิ ตํ อธิกรณํ อาวาสิกานํ ภิกฺขูนํ นิยฺยาเทตพฺพํ.}

\prose{เต เจ ภิกฺขเว ภิกฺขู สกฺโกนฺติ ตํ อธิกรณํ วูปสเมตุํ, อิทํ วุจฺจติ ภิกฺขเว อธิกรณํ วูปสนฺตํ. เกน วูปสนฺตํ, สมฺมุขา\-วินเยน. กิญฺจ ตตฺถ สมฺมุขา\-วินยสฺมิํ, สํฆสมฺมุขตา ธมฺม\-สมฺมุขตา วินย\-สมฺมุขตา ปุคฺคล\-สมฺมุขตา ฯเปฯ. เอวํ วูปสนฺตํ เจ ภิกฺขเว อธิกรณํ การโก อุกฺโกเฏติ, อุกฺโกฏนกํ ปาจิตฺติยํ. ฉนฺททายโก ขียติ, ขียนกํ ปาจิตฺติยํ.}

\csromanmatikaanchor{mtk.118}
\csromantocmarkref{subsection}{อุพฺพาหิกายวูปสมน}{mtk.118}
\bookmark[dest={mtk.118},level=2]{อุพฺพาหิกายวูปสมน}
\csromanheader{2}{อุพฺพาหิกาย\-วูปสมน}
\proseitem{231}{เตหิ เจ ภิกฺขเว ภิกฺขูหิ ตสฺมิํ อธิกรเณ วินิจฺฉิยมาเน อนนฺตานิ\csromansharedfootnote{554e35fd0356bcfb}{อนคฺคานิ (สี)} เจว ภสฺสานิ ชายนฺติ, น เจกสฺส\csromansharedfootnote{a69860b6d4e07d95}{น เจตสฺส (ก)} ภาสิตสฺส อตฺโถ วิญฺญายติ, อนุชานามิ ภิกฺขเว เอวรูปํ อธิกรณํ อุพฺพาหิกาย วูปสเมตุํ. ทสหงฺเคหิ สมนฺนาคโต ภิกฺขุ อุพฺพาหิกาย สมฺมนฺนิตพฺโพ. สีลวา โหติ ปาติโมกฺข\-สํวร\-สํวุโต วิหรติ อาจารโคจร\-สมฺปนฺโน อณุมตฺเตสุ วชฺเชสุ ภยทสฺสาวี สมาทาย สิกฺขติ สิกฺขาปเทสุ, พหุสฺสุโต โหติ สุตธโร สุตสนฺนิจโย, เย เต ธมฺมา อาทิกลฺยาณา มชฺเฌกลฺยาณา ปริโยสาน\-กลฺยาณา สาตฺถํ สพฺยญฺชนํ เกวล\-ปริปุณฺณํ ปริสุทฺธํ พฺรหฺมจริยํ อภิวทนฺติ, ตถารูปสฺส ธมฺมา พหุสฺสุตา โหนฺติ ธาตา\csromansharedfootnote{1c25e3e705596ff0}{ธตา (สี, สฺยา)} วจสา ปริจิตา มนสา\-นุเปกฺขิตา ทิฏฺฐิยา สุปฺปฏิวิทฺธา, อุภยานิ โข ปนสฺส ปาติโมกฺขานิ วิตฺถาเรน สฺวาคตานิ โหนฺติ สุวิภตฺตานิ สุปฺปวตฺตีนิ สุวินิจฺฉิตานิ สุตฺตโส อนุพฺยญฺชนโส, วินเย โข ปน ฐิโต\csromansharedfootnote{a5e68d375bde21cb}{เฉโก (ก)} โหติ อสํหีโร, ปฏิพโล โหติ อุโภ อตฺถ\-ปจฺจตฺถิเก อสฺสาเสตุํ สญฺญาเปตุํ นิชฺฌาเปตุํ เปกฺเขตุํ ปสาเทตุํ, อธิกรณ\-สมุปฺปาท\-วูปสมน\-กุสโล โหติ, อธิกรณํ ชานาติ, อธิกรณ\-สมุทยํ ชานาติ, อธิกรณ\-นิโรธํ ชานาติ, อธิกรณนิโรธคามินิปฏิปทํ ชานาติ. อนุชานามิ ภิกฺขเว อิเมหิ ทสหงฺเคหิ สมนฺนาคตํ ภิกฺขุํ อุพฺพาหิกาย สมฺมนฺนิตุํ. \csromanfolio{224}เอวญฺจ ปน ภิกฺขเว สมฺมนฺนิตพฺโพ, ปฐมํ ภิกฺขุ ยาจิตพฺโพ, ยาจิตฺวา พฺยตฺเตน ภิกฺขุนา ปฏิพเลน สํโฆ ญาเปตพฺโพ\csromandash{}}

\proseitem{232}{“สุณาตุ เม ภนฺเต สํโฆ, อมฺหากํ อิมสฺมิํ อธิกรเณ วินิจฺฉิยมาเน อนนฺตานิ เจว ภสฺสานิ ชายนฺติ, น เจกสฺส ภาสิตสฺส อตฺโถ วิญฺญายติ, ยทิ สํฆสฺส ปตฺตกลฺลํ, สํโฆ อิตฺถนฺนามญฺจ อิตฺถนฺนามญฺจ ภิกฺขุํ สมฺมนฺเนยฺย อุพฺพาหิกาย อิมํ อธิกรณํ วูปสเมตุํ, เอสา ญตฺติ.}

\prose{สุณาตุ เม ภนฺเต สํโฆ, อมฺหากํ อิมสฺมิํ อธิกรเณ วินิจฺฉิยมาเน อนนฺตานิ เจว ภสฺสานิ ชายนฺติ, น เจกสฺส ภาสิตสฺส อตฺโถ วิญฺญายติ, สํโฆ อิตฺถนฺนามญฺจ อิตฺถนฺนามญฺจ ภิกฺขุํ สมฺมนฺนติ อุพฺพาหิกาย อิมํ อธิกรณํ วูปสเมตุํ, ยสฺสายสฺมโต ขมติ อิตฺถนฺนามสฺส จ อิตฺถนฺนามสฺส จ ภิกฺขูโน สมฺมุติ อุพฺพาหิกาย อิมํ อธิกรณํ วูปสเมตุํ, โส ตุณฺหสฺส, ยสฺส นกฺขมติ, โส ภาเสยฺย.}

\prose{สมฺมโต สํเฆน อิตฺถนฺนาโม จ อิตฺถนฺนาโม จ ภิกฺขุ อุพฺพาหิกาย อิมํ อธิกรณํ วูปสเมตุํ, ขมติ สํฆสฺส, ตสฺมา ตุณฺหี, เอวเมตํ ธารยามี”ติ.}

\prose{เต เจ ภิกฺขเว ภิกฺขู สกฺโกนฺติ ตํ อธิกรณํ อุพฺพาหิกาย วูปสเมตุํ, อิทํ วุจฺจติ ภิกฺขเว อธิกรณํ วูปสนฺตํ. เกน วูปสนฺตํ, สมฺมุขา\-วินเยน. กิญฺจ ตตฺถ สมฺมุขา\-วินยสฺมิํ, ธมฺม\-สมฺมุขตา วินย\-สมฺมุขตา ปุคฺคล\-สมฺมุขตา ฯเปฯ. เอวํ วูปสนฺตํ เจ ภิกฺขเว อธิกรณํ การโก อุกฺโกเฏติ, อุกฺโกฏนกํ ปาจิตฺติยํ.}

\proseitem{233}{เตหิ เจ ภิกฺขเว ภิกฺขูหิ ตสฺมิํ อธิกรเณ วินิจฺฉิยมาเน ตตฺราสฺส ภิกฺขุ ธมฺมกถิโก, ตสฺส เนว สุตฺตํ อาคตํ โหติ, โน สุตฺตวิภงฺโค, โส อตฺถํ อสลฺลกฺเขนฺโต พฺยญฺชน\-จฺฉายาย อตฺถํ ปฏิพาหติ. พฺยตฺเตน ภิกฺขุนา ปฏิพเลน เต ภิกฺขู ญาเปตพฺพา\csromandash{}}

\prose{“สุณนฺตุ เม อายสฺมนฺตา, อยํ อิตฺถนฺนาโม ภิกฺขู ธมฺมกถิโก, อิมสฺส เนว สุตฺตํ อาคตํ โหติ, โน สุตฺตวิภงฺโค, โส อตฺถํ อสลฺลกฺเขนฺโต พฺยญฺชน\-จฺฉายาย อตฺถํ ปฏิพาหติ, ยทา\-ยสฺมนฺตานํ \csromanfolio{225}ปตฺตกลฺลํ, อิตฺถนฺนามํ ภิกฺขุํ วุฏฺฐาเปตฺวา อวเสสา อิมํ อธิกรณํ วูปสเมยฺยามา”ติ.}

\prose{เต เจ ภิกฺขเว ภิกฺขู ตํ ภิกฺขุํ วุฏฺฐาเปตฺวา สกฺโกนฺติ ตํ อธิกรณํ วูปสเมตุํ, อิทํ วุจฺจติ ภิกฺขเว อธิกรณํ วูปสนฺตํ. เกน วูปสนฺตํ, สมฺมุขา\-วินเยน. กิญฺจ ตตฺถ สมฺมุขา\-วินยสฺมิํ, ธมฺม\-สมฺมุขตา วินย\-สมฺมุขตา ปุคฺคล\-สมฺมุขตา ฯเปฯ. เอวํ วูปสนฺตํ เจ ภิกฺขเว อธิกรณํ กรโก อุกฺโกเฏติ, อุกฺโกฏนกํ ปาจิตฺติยํ.}

\prose{เตหิ เจ ภิกฺขเว ภิกฺขูหิ ตสฺมิํ อธิกรเณ วินิจฺฉิยมาเน ตตฺราสฺส ภิกฺขุ ธมฺมกถิโก, ตสฺส สุตฺตํ หิ โข อาคตํ โหติ, โน สุตฺตวิภงฺโค, โส อตฺถํ อสลฺลกฺเขนฺโต พฺยญฺชน\-จฺฉายาย อตฺถํ ปฏิพาหติ. พฺยตฺเตน ภิกฺขูนา ปฏิพเลน เต ภิกฺขู ญาเปตพฺพา\csromandash{}}

\prose{“สุณนฺตุ เม อายสฺมนฺตา, อยํ อิตฺถนฺนาโม ภิกฺขุ ธมฺมกถิโก, อิมสฺส สุตฺตํ หิ โข อาคตํ โหติ, โน สุตฺตวิภงฺโค, โส อตฺถํ อสลฺลกฺเขนฺโต พฺยญฺชน\-จฺฉายาย อตฺถํ ปฏิพาหติ, ยทา\-ยสฺมนฺตานํ ปตฺตกลฺลํ, อิตฺถนฺนามํ ภิกฺขุํ วุฏฺฐาเปตฺวา อวเสสา อิมํ อธิกรณํ วูปสเมยฺยามา”ติ.}

\prose{เต เจ ภิกฺขเว ภิกฺขู ตํ ภิกฺขุํ วุฏฺฐาเปตฺวา สกฺโกนฺติ ตํ อธิกรณํ วูปสเมตุํ, อิทํ วุจฺจติ ภิกฺขเว อธิกรณํ วูปสนฺตํ. เกน วูปสนฺตํ, สมฺมุขา\-วินเยน. กิญฺจ ตตฺถ สมฺมุขา\-วินยสฺมิํ, ธมฺม\-สมฺมุขตา วินย\-สมฺมุขตา ปุคฺคล\-สมฺมุขตา ฯเปฯ. เอวํ วูปสนฺตํ เจ ภิกฺขเว อธิกรณํ การโก อุกฺโกเฏติ, อุกฺโกฏนกํ ปาจิตฺติยํ.}

\csromanmatikaanchor{mtk.119}
\csromantocmarkref{subsection}{เยภุยฺยสิกาวินย}{mtk.119}
\bookmark[dest={mtk.119},level=2]{เยภุยฺยสิกาวินย}
\csromanheader{2}{เยภุยฺยสิกา\-วินย}
\proseitem{234}{เต เจ ภิกฺขเว ภิกฺขู น สกฺโกนฺติ ตํ อธิกรณํ อุพฺพาหิกาย วูปสเมตุํ. เตหิ ภิกฺขเว ภิกฺขูหิ ตํ อธิกรณํ สํฆสฺส นิยฺยาเทตพฺพํ “น มยํ\csromansharedfootnote{4a92d711efccdbfc}{น จ มยํ (ก)} ภนฺเต สกฺโกม อิมํ อธิกรณํ อุพฺพาหิกาย วูปสเมตุํ, สํโฆว อิมํ อธิกรณํ วูปสเมตู”ติ. อนุชานามิ ภิกฺขเว เอวรูปํ อธิกรณํ เยภุยฺยสิกาย วูปสเมตุํ. ปญฺจหงฺเคหิ สมนฺนาคโต ภิกฺขุ สลาก\-คฺคาหาปโก สมฺมนฺนิตพฺโพ, โย น ฉนฺทาคติํ คจฺเฉยฺย. \csromanfolio{226}น โทสาคติํ คจฺเฉยฺย, น โมหาคติํ คจฺเฉยฺย, น ภยาคติํ คจฺเฉยฺย, คหิตาคหิตญฺจ ชาเนยฺย ฯเปฯ. เอวญฺจ ปน ภิกฺขเว สมฺมนฺนิตพฺโพ, ปฐมํ ภิกฺขุ ยาจิตพฺโพ, ยาจิตฺวา พฺยตฺเตน ภิกฺขุนา ปฏิพเลน สํโฆ ญาเปตพฺโพ\csromandash{}}

\prose{“สุณาตุ เม ภนฺเต สํโฆ, ยทิ สํฆสฺส ปตฺตกลฺลํ, สํโฆ อิตฺถนฺนามํ ภิกฺขุํ สลาก\-คฺคาหาปกํ สมฺมนฺเนยฺย, เอสา ญตฺติ.}

\prose{สุณาตุ เม ภนฺเต สํโฆ, สํโฆ อิตฺถนฺนามํ ภิกฺขุํ สลาก\-คฺคาหาปกํ สมฺมนฺนติ, ยสฺสายสฺมโต ขมติ อิตฺถนฺนามสฺส ภิกฺขุโน สลาก\-คฺคาหาปกสฺส สมฺมุติ, โส ตุณฺหสฺส, ยสฺส นกฺขมติ, โส ภาเสยฺย.}

\prose{สมฺมโต สํเฆน อิตฺถนฺนาโม ภิกฺขุ สลาก\-คฺคาหาปโก, ขมติ สํฆสฺส, ตสฺมา ตุณฺหี, เอวเมตํ ธารยามี”ติ.}

\setlength{\csromangathapagewidth}{0pt}
\setlength{\csromangathawakwidth}{0pt}
\csromangathameasuregroup{เกน วูปสนฺตํ, \\ ฉนฺททายโก ขียติ,}{\csromangathabat{เกน วูปสนฺตํ,}{สมฺมุขาวินเยน จ เยภุยฺยสิกาย จ.} \\ \csromangathabat{ฉนฺททายโก ขียติ,}{ขียนกํ ปาจิตฺติยนฺติ.}}
\csromangathasetleft{เกน วูปสนฺตํ, \\ ฉนฺททายโก ขียติ,}
\csromangathagroup{\csromangathastanza{\csromangathaitembat{234}{เกน วูปสนฺตํ,}{สมฺมุขา\-วินเยน จ เยภุยฺยสิกาย จ.}}\csromangathastanza{\csromangathaitembat{234}{ฉนฺททายโก ขียติ,}{ขียนกํ ปาจิตฺติยนฺติ.}}}

\prose{เตน สลาก\-คฺคาหาปเกน ภิกฺขุนา สลากา คาเหตพฺพา, ยถา พหุตรา ภิกฺขู ธมฺมวาทิโน วทนฺติ, ตถา ตํ อธิกรณํ วูปสเมตพฺพํ, อิทํ วุจฺจติ ภิกฺขเว อธิกรณํ วูปสนฺตํ.}

\prose{กิญฺจ ตตฺถ สมฺมุขา\-วินยสฺมิํ, สํฆสมฺมุขตา ธมฺม\-สมฺมุขตา วินย\-สมฺมุขตา ปุคฺคล\-สมฺมุขตา.}

\prose{กา จ ตตฺถ สํฆสมฺมุขตา, ยาวติกา ภิกฺขู กมฺมปฺปตฺตา, เต อาคตา โหนฺติ, ฉนฺทารหานํ ฉนฺโท อาหโฏ โหติ, สมฺมุขีภูตา น ปฏิกฺโกสนฺติ, อยํ ตตฺถ สํฆสมฺมุขตา.}

\prose{กา จ ตตฺถ ธมฺม\-สมฺมุขตา วินย\-สมฺมุขตา, เยน ธมฺเมน เยน วินเยน เยน สตฺถุสาสเนน ตํ อธิกรณํ วูปสมฺมติ, อยํ ตตฺถ ธมฺม\-สมฺมุขตา วินย\-สมฺมุขตา.}

\prose{กา จ ตตฺถ ปุคฺคล\-สมฺมุขตา, โย จ วิวทติ, เยน จ วิวทติ, อุโภ อตฺถ\-ปจฺจตฺถิกา สมฺมุขีภูตา โหนฺติ, อยํ ตตฺถ ปุคฺคล\-สมฺมุขตา.}

\prose{กา จ ตตฺถ เยภุยฺยสิกาย, ยา เยภุยฺยสิกา\-กมฺมสฺส กิริยา กรณํ อุปคมนํ อชฺฌุปคมนํ อธิวาสนา อปฺปฏิกฺโกสนา, อยํ ตตฺถ เยภุยฺยสิกาย.}

\prose{เอวํ วูปสนฺตํ เจ ภิกฺขเว อธิกรณํ การโก อุกฺโกเฏติ, อุกฺโกฏนกํ ปาจิตฺติยํ.}

\csromanmatikaanchor{mtk.120}
\csromantocmarkref{subsection}{ติวิธสลากคฺคาห}{mtk.120}
\bookmark[dest={mtk.120},level=2]{ติวิธสลากคฺคาห}
\csromanheader{2}{\textbf{ติวิธสลากคฺคาห}}
\proseitem{253}{\csromanfolio{227}เตน โข ปน สมเยน สาวตฺถิยา เอวํ ชาตํ เอวํ สมุปฺปนฺนํ อธิกรณํ โหติ. อถ โข เต ภิกฺขู อสนฺตุฏฺฐา สาวตฺถิยา สํฆสฺส อธิกรณ\-วูปสมเนน. อโสสุํ โข “อมุกสฺมิํ กิร อาวาเส สมฺพหุลา เถรา วิหรนฺติ พหุสฺสุตา อาคตาคมา ธมฺมธรา วินยธรา มาติกาธรา ปณฺฑิตา วิยตฺตา เมธาวิโน ลชฺชิโน กุกฺกุจฺจกา สิกฺขากามา, เต เจ เถรา อิมํ อธิกรณํ วูปสเมยฺยุํ ธมฺเมน วินเยน สตฺถุสาสเนน, เอวมิทํ อธิกรณํ สุวูปสนฺตํ อสฺสา”ติ. อถ โข เต ภิกฺขู ตํ อาวาสํ คนฺตฺวา เต เถเร เอตทโวจุํ “อิทํ ภนฺเต อธิกรณํ เอวํ ชาตํ เอวํ สมุปฺปนฺนํ, สาธุ ภนฺเต เถรา อิมํ อธิกรณํ วูปสเมนฺตุ ธมฺเมน วินเยน สตฺถุสาสเนน, ยถยิทํ อธิกรณํ สุวูปสนฺตํ อสฺสา”ติ. อถ โข เต เถรา “ยถา สาวตฺถิยา สํเฆน อธิกรณํ วูปสมิตํ, ตถา สุวูปสนฺตนฺ”ติ\csromansharedfootnote{5c5010a572700401}{ยถา สุวูปสนฺตํ (สี, สฺยา)} ตถา ตํ อธิกรณํ วูปสเมสุํ.}

\prose{อถ โข เต ภิกฺขู อสนฺตุฏฺฐา สาวตฺถิยา สํฆสฺส อธิกรณ\-วูปสมเนน อสนฺตุฏฺฐา สมฺพหุลานํ เถรานํ อธิกรณ\-วูปสมเนน. อสฺโสสุํ โข “อมุกสฺมิํ กิร อาวาเส ตโย เถรา วิหรนฺติ ฯเปฯ เทฺว เถรา วิหรนฺติ ฯเปฯ เอโก เถโร วิหรติ พหุสฺสุโต อาคตาคโม ธมฺมธโร วินยธโร มาติกาธโร ปณฺฑิโต วิยตฺโต เมธาวี ลชฺชี กุกฺกุจฺจโก สิกฺขากาโม, โส เจ เถโร อิมํ อธิกรณํ วูปสเมยฺย ธมฺเมน วินเยน สตฺถุสาสเนน, เอวมิทํ อธิกรณํ สุวูปสนฺตํ อสฺสา”ติ. อถ โข เต ภิกฺขู ตํ อาวาสํ คนฺตฺวา ตํ เถรํ เอตทโวจุํ “อิทํ ภนฺเต อธิกรณํ เอวํ ชาตํ เอวํ สมุปฺปนฺนํ, สาธุ ภนฺเต เถโร อิมํ อธิกรณํ วูปสเมตุ ธมฺเมน วินเยน สตฺถุสาสเนน, ยถยิทํ อธิกรณํ สุวูปสนฺตํ อสฺสา”ติ. อถ โข โส เถโร “ยถา สาวตฺถิยา สํเฆน อธิกรณํ วูปสมิตํ ยถา สมฺพหุเลหิ เถเรหิ อธิกรณํ วูปสมิตํ ยถา ตีหิ เถเรหิ อธิกรณํ วูปสมิตํ ยถา ทฺวีหิ เถเรหิ อธิกรณํ วูปสมิตํ, ตถา สุวูปสนฺตนฺ”ติ ตถา ตํ อธิกรณํ วูปสเมสิ.}

\prose{\csromanfolio{228}อถ โข เต ภิกฺขู อสนฺตุฏฺฐา สาวตฺถิยา สํฆสฺส อธิกรณ\-วูปสมเนน อสนฺตุฏฺฐา สมฺพหุลานํ เถรานํ อธิกรณ\-วูปสมเนน อสนฺตุฏฺฐา ติณฺณํ เถเรนํ อธิกรณ\-วูปสมเนน อสนฺตุฏฺฐา ทฺวินฺนํ เถรานํ อธิกรณ\-วูปสมเนน อสนฺตุฏฺฐา เอกสฺส เถรสฺส อธิกรณ\-วูปสมเนน เยน ภควา เตนุ\-ปสงฺกมิํสุ, อุปสงฺกมิตฺวา ภควโต เอตมตฺถํ อาโรเจสุํ ฯเปฯ. นิหตเมตํ ภิกฺขเว อธิกรณํ สนฺตํ วูปสนฺตํ สุวูปสนฺตํ. อนุชานามิ ภิกฺขเว เตสํ ภิกฺขูนํ สญฺญตฺติยา ตโย สลากคฺคาเห คุฬฺหกํ สกณฺณ\-ชปฺปกํ วิวฏกํ.}

\prose{กถญฺจ ภิกฺขเว คุฬฺหโก สลากคฺคาโห โหติ. เตน สลาก\-คฺคาหาปเกน ภิกฺขุนา สลากาโย วณฺณาวณฺณาโย กตฺวา เอกเมโก ภิกฺขุ อุปสงฺกมิตฺวา เอวมสฺส วจนีโย “อยํ เอวํวาทิสฺส สลากา, อยํ เอวํวาทิสฺส สลากา, ยํ อิจฺฉสิ, ตํ คณฺหาหี”ติ. คหิเต วตฺตพฺโพ “มา จ กสฺสจิ ทสฺเสหี”ติ. สเจ ชานาติ “อธมฺมวาที พหุตรา”ติ, “ทุคฺคโห”ติ ปจฺจุกฺกฑฺฒิตพฺพํ. สเจ ชานาติ “ธมฺมวาที พหุตรา”ติ, “สุคฺคโห”ติ สาเวตพฺพํ. เอวํ โข ภิกฺขเว คูฬฺหโก สลากคฺคาโห โหติ.}

\prose{กถญฺจ ภิกฺขเว สกณฺณชปฺปโก สลากคฺคาโห โหติ. เตน สลาก\-คฺคาหาปเกน ภิกฺขุนา เอกเมกสฺส ภิกฺขุโน อุปกณฺณเก อาโรเจตพฺพํ “อยํ เอวํวาทิสฺส สลากา, อยํ เอวํวาทิสฺส สลากา, ยํ อิจฺฉสิ, ตํ คณฺหาหี”ติ. คหิเต วตฺตพฺโพ “มา จ กสฺสจิ อาโรเจหี”ติ. สเจ ชานาติ “อธมฺมวาที พหุตรา”ติ, “ทุคฺคโห”ติ ปจฺจุกฺกฑฺฒิตพฺพํ. สเจ ชานาติ “ธมฺมวาที พหุตรา”ติ, “สุคฺคโห”ติ สาเวตพฺพํ. เอวํ โข ภิกฺขเว สกณฺณชปฺปโก สลากคฺคาโห โหติ.}

\prose{กถญฺจ ภิกฺขเว วิวฏโก สลากคฺคาโห โหติ. สเจ ชานาติ “ธมฺมวาที พหุตรา”ติ, วิสฺสฏฺเฐเนว วิวเฏน คาเหตพฺโพ. เอวํ โข ภิกฺขเว วิวฏโก สลากคฺคาโห โหติ. อิเม โข ภิกฺขเว ตโย สลากคฺคาหาติ.}

\csromanmatikaanchor{mtk.121}
\csromantocmarkref{subsection}{สติวินย}{mtk.121}
\bookmark[dest={mtk.121},level=2]{สติวินย}
\csromanheader{2}{\textbf{สติวินย}}
\proseitem{236}{\csromanfolio{229}\csromansymbolfootnote{*}{วิ 5. 184, 195 ปิฏฺฐาทีสุปิ.}อนุวาทา\-ธิกรณํ กติหิ สมเถหิ สมฺมติ. อนุวาทา\-ธิกรณํ จตูหิ สมเถหิ สมฺมติ สมฺมุขา\-วินเยน จ สติวินเยน จ อมูฬฺหวินเยน จ ตสฺสปาปิยสิกาย จ. สิยา อนุวาทา\-ธิกรณํ เทฺว สมเถ อนาคมฺม อมูฬฺหวินยญฺจ ตสฺสปาปิยสิกญฺจ ทฺวีหิ สมเถหิ สมฺเมยฺย สมฺมุขา\-วินเยน จ สติวินเยน จาติ? สิยาติสฺส วจนียํ. ยถา กถํ วิย. อิธ ปน ภิกฺขเว ภิกฺขู ภิกฺขุํ อมูลิกาย สีลวิปตฺติยา อนุทฺธํเสนฺติ, ตสฺส โข ภิกฺขเว\csromansharedfootnote{7b942117c801da66}{ตสฺส โข ตํ ภิกฺขเว (สฺยา, ก)} ภิกฺขุโน สติเวปุลฺล\-ปฺปตฺตสฺส สติวินโย ทาตพฺโพ. เอวญฺจ ปน ภิกฺขเว ทาตพฺโพ,}

\prose{เตน ภิกฺขเว ภิกฺขุนา สํฆํ อุปสงฺกมิตฺวา เอกํสํ อุตฺตราสงฺคํ กริตฺวา ฯเปฯ เอวมสฺส วจนีโย “มํ ภนฺเต ภิกฺขู อมูลิกาย สีลวิปตฺติยา อนุทฺธํเสนฺติ, โสหํ ภนฺเต สติเวปุลฺล\-ปฺปตฺโต สํฆํ สติวินยํ ยาจามี”ติ. ทุติยมฺปิ ยาจิตพฺโพ. ตติยมฺปิ ยาจิตพฺโพ. พฺยตฺเตน ภิกฺขุนา ปฏิพเลน สํโฆ ญาเปตพฺโพ\csromandash{}}

\prose{“สุณาตุ เม ภนฺเต สํโฆ, ภิกฺขู อิตฺถนฺนามํ ภิกฺขุํ อมูลิกาย สีลวิปตฺติยา อนุทฺธํเสนฺติ, โส สติเวปุลฺล\-ปฺปตฺโต สํฆํ สติวินยํ ยาจติ, ยทิ สํฆสฺส ปตฺตกลฺลํ, สํโฆ อิตฺถนฺนามสฺส ภิกฺขุโน สติเวปุลฺล\-ปฺปตฺตสฺส สติวินยํ ทเทยฺย, เอสา ญตฺติ.}

\prose{สุณาตุ เม ภนฺเต สํโฆ, ภิกฺขู อิตฺถนฺนามํ ภิกฺขุํ อมูลิกาย สีลวิปตฺติยา อนุทฺธํเสนฺติ, โส สติเวปุลฺล\-ปฺปตฺโต สํฆํ สติวินยํ ยาจติ, สํโฆ อิตฺถนฺนามสฺส ภิกฺขุโน สติเวปุลฺล\-ปฺปตฺตสฺส สติวินยํ เทติ, ยสฺสายสฺมโต ขมติ อิตฺถนฺนามสฺส ภิกฺขุโน สติเวปุลฺล\-ปฺปตฺตสฺส สติวินยสฺส ทานํ, โส ตุณฺหสฺส, ยสฺส นกฺขมติ, โส ภาเสยฺย.}

\prose{ทุติยมฺปิ เอตมตฺถํ วทามิ ฯเปฯ. ตติยมฺปิ เอตมตฺถํ วทามิ ฯเปฯ.}

\prose{ทินฺโน สํเฆน อิตฺถนฺนามสฺส ภิกฺขุโน สติเวปุลฺล\-ปฺปตฺตสฺส สติวินโย, ขมติ สํฆสฺส, ตสฺมา ตุณฺหี, เอวเมตํ ธารยามี”ติ.}

\prose{อิทํ วุจฺจติ ภิกฺขเว อธิกรณํ วูปสนฺตํ. เกน วูปสนฺตํ, สมฺมุขา\-วินเยน จ สติวินเยน จ. กิญฺจ ตตฺถ สมฺมุขา\-วินยสฺมิํ, สํฆสมฺมุขตา \csromanfolio{230}ธมฺม\-สมฺมุขตา วินย\-สมฺมุขตา ปุคฺคล\-สมฺมุขตา ฯเปฯ. กา จ ตตฺถ ปุคฺคล\-สมฺมุขตา, โย จ อนุวทติ, ยญฺจ อนุวทติ, อุโภ สมฺมุขีภูตา โหนฺติ, อยํ ตตฺถ ปุคฺคล\-สมฺมุขตา. กิญฺจ ตตฺถ สติวินยสฺมิํ, ยา สติวินยสฺส กมฺมสฺส กิริยา กรณํ อุปคมนํ อชฺฌุปคมนํ อธิวาสนา อปฺปฏิกฺโกสนา, อิทํ ตตฺถ สติวินยสฺมิํ. เอวํ วูปสนฺตํ เจ ภิกฺขเว อธิกรณํ การโก อุกฺโกเฏติ, อุกฺโกฏนกํ ปาจิตฺติยํ. ฉนฺททายโก ขียติ, ขียนกํ ปาจิตฺติยํ.}

\csromanmatikaanchor{mtk.122}
\csromantocmarkref{subsection}{อมูฬฺหวินย}{mtk.122}
\bookmark[dest={mtk.122},level=2]{อมูฬฺหวินย}
\csromanheader{2}{\textbf{อมูฬฺหวินย}}
\proseitem{237}{สิยา อนุวาทา\-ธิกรณํ เทฺว สมเถ อนาคมฺม สติวินยญฺจ ตสฺสปาปิยสิกญฺจ ทฺวีหิ สมเถหิ สมฺเมยฺย สมฺมุขา\-วินเยน จ อมูฬวินเยน จาติ? สิยาติสฺส วจนียํ. ยถา กถํ วิย. อิธ ปน ภิกฺขเว ภิกฺขุ อุมฺมตฺตโก โหติ จิตฺตวิปริยาสกโต, เตน อุมฺมตฺตเกน จิตฺตวิปริยาสกเตน พหุํ อสฺสามณกํ อชฺฌาจิณฺณํ โหติ ภาสิต\-ปริกฺกนฺตํ, ตํ ภิกฺขู อุมฺมตฺตเกน จิตฺตวิปริยาสกเตน อชฺฌาจิณฺเณน อาปตฺติยา โจเทนฺติ “สรตายสฺมา เอวรูปิํ อาปตฺติํ อาปชฺชิตา”ติ. โส เอวํ วเทติ “อหํ โข อาวุโส อุมฺมตฺตโก อโหสิํ จิตฺตวิปริยาสกโต, เตน เม อุมฺมตฺตเกน จิตฺตวิปริยาสกเตน พหุํ อสฺสามณกํ อชฺฌาจิณฺณํ ภาสิต\-ปริกฺกนฺตํ, นาหํ ตํ สรามิ, มูเฬฺหน เม เอตํ กตนฺ”ติ. เอวมฺปิ นํ วุจฺจมานา โจเทนฺเตว “สรตายสฺมา เอวรูปิํ อาปตฺติํ อาปชฺชิตา”ติ. ตสฺส โข ภิกฺขเว ภิกฺขุโน อมูฬฺหสฺส อมูฬฺหวินโย ทาตพฺโพ. เอวญฺจ ปน ภิกฺขเว ทาตพฺโพ.}

\prose{เตน ภิกฺขเว ภิกฺขุนา สํฆํ อุปสงฺกมิตฺวา เอกํสํ อุตฺตราสงฺคํ กริตฺวา ฯเปฯ เอวมสฺส วจนีโย “อหํ ภนฺเต อุมฺมตฺตโก อโหสิํ จิตฺตวิปริยาสกโต, เตน เม อุมฺมตฺตเกน จิตฺตวิปริยาสกเตน พหุํ อสฺสามณกํ อชฺฌาจิณฺณํ ภาสิต\-ปริกฺกนฺตํ, มํ ภิกฺขู อุมฺมตฺตเกน จิตฺตวิปริยาสกเตน อชฺฌาจิณฺเณน อาปตฺติยา โจเทนฺติ “สรตายสฺมา เอวรูปิํ อาปตฺติํ อาปชฺชิตา”ติ, ตฺยาหํ เอวํ วทามิ ‘อหํ โข อาวุโส อุมฺมตฺตโก อโหสิํ จิตฺตวิปริยาสกโต, เตน เม อุมฺมตฺตเกน จิตฺตวิปริยาสกเตน พหุํ อสฺสามณกํ อชฺฌาจิณฺณํ ภาสิต\-ปริกฺกนฺตํ, นาหํ ตํ สรามิ, มูเฬฺหน เม เอตํ กตนฺ’ติ, เอวมฺปิ มํ วุจฺจมาโน โจเทนฺเตว \csromanfolio{231}“สรตายสฺมา เอวรูปิํ อาปตฺติํ อาปชฺชิตา’ติ, โสหํ ภนฺเต อมูโฬฺห สํฆํ อมูฬฺหวินยํ ยาจามี”ติ. ทุติยมฺปิ ยาจิตพฺโพ. ตติยมฺปิ ยาจิตพฺโพ. พฺยตฺเตน ภิกฺขุนา ปฏิพเลน สํโฆ ญาเปตพฺโพ\csromandash{}}

\prose{“สุณาตุ เม ภนฺเต สํโฆ, อยํ อิตฺถนฺนาโม ภิกฺขุ อุมฺมตฺตโก อโหสิ จิตฺตวิปริยาสกโต, เตน อุมฺมตฺตเกน จิตฺตวิปริยาสกเตน พหุํ อสฺสามณกํ อชฺฌาจิณฺณํ ภาสิต\-ปริกฺกนฺตํ, ตํ ภิกฺขู อุมฺมตฺตเกน จิตฺตวิปริยาสกเตน อชฺฌาจิณฺเณน อาปตฺติยา โจเทนฺติ ‘สรตายสฺมา เอวรูปิํ อาปตฺติํ อาปชฺชิตา’ติ, โส เอวํ วเทติ ‘อหํ โข อาวุโส อุมฺมตฺตโก อโหสิํ จิตฺตวิปริยาสกโต, เตน เม อุมฺมตฺตเกน จิตฺตวิปริยาสกเตน พหุํ อสฺสามณกํ อชฺฌาจิณฺณํ ภาสิต\-ปริกฺกนฺตํ, นาหํ ตํ สรามิ, มูเฬฺหน เม เอตํ กตนฺ’ติ เอวมฺปิ นํ วุจฺจมานา โจเทนฺเตว ‘สรตายสฺมา เอวรูปิํ อาปตฺติํ อาปชฺชิตา’ติ. โส อมูโฬฺห สํฆํ อมูฬฺหวินยํ ยาจติ, ยทิ สํฆสฺส ปตฺตกลฺลํ, สํโฆ อิตฺถนฺนามสฺส ภิกฺขุโน อมูฬฺหสฺส อมูฬฺหวินยํ ทเทยฺย, เอสา ญตฺติ.}

\prose{สุณาตุ เม ภนฺเต สํโฆ, อยํ อิตฺถนฺนาโม ภิกฺขุ อุมฺมตฺตโก อโหสิ จิตฺตวิปริยาสกโต, เตน อุมฺมตฺตเกน จิตฺตวิปริยาสกเตน พหุํ อสฺสามณกํ อชฺฌาจิณฺณํ ภาสิต\-ปริกฺกนฺตํ, ตํ ภิกฺขู อุมฺมตฺตเกน จิตฺตวิปริยาสกเตน อชฺฌาจิณฺเณน อาปตฺติยา โจเทนฺติ ‘สรตายสฺมา เอวรูปิํ อาปตฺติํ อาปชฺชิตา’ติ, โส เอวํ วเทติ ‘อหํ โข อาวุโส อุมฺมตฺตโก อโหสิํ จิตฺตวิปริยาสกโต, เตน เม อุมฺมตฺตเกน จิตฺตวิปริยาสกเตน พหุํ อสฺสามณกํ อชฺฌาจิณฺณํ ภาสิต\-ปริกฺกนฺตํ, นาหํ ตํ สรามิ, มูเฬฺหน เม เอตํ กตนฺ’ติ, เอวมฺปิ นํ วุจฺจมานา โจเทนฺเตว ‘สรตายสฺมา เอวรูปิํ อาปตฺติํ อาปชฺชิตา’ติ, โส อมูโฬฺห สํฆํ อมูฬฺหวินยํ ยาจติ, สํโฆ อิตฺถนฺนามสฺส ภิกฺขุโน อมูฬฺหสฺส อมูฬฺหวินยํ เทติ, ยสฺสายสฺมโต ขมติ อิตฺถนฺนามสฺส ภิกฺขุโน อมูฬฺหสฺส อมูฬฺห\-วินยสฺส ทานํ, โส ตุณฺหสฺส, ยสฺส นกฺขมติ, โส ภาเสยฺย.}

\prose{ทุติยมฺปิ เอตมตฺถํ วทามิ ฯเปฯ. ตติยมฺปิ เอตมตฺถํ วทามิ ฯเปฯ.}

\prose{ทินฺโน สํเฆน อิตฺถนฺนามสฺส ภิกฺขุโน อมูฬฺหสฺส อมูฬฺหวินโย, ขมติ สํฆสฺส, ตสฺมา ตุณฺหี, เอวเมตํ ธารยามี”ติ.}

\prose{\csromanfolio{232}อิทํ วุจฺจติ ภิกฺขเว อธิกรณํ วูปสนฺตํ. เกน วูปสนฺตํ, สมฺมุขา\-วินเยน จ อมูฬฺหวินเยน จ. กิญฺจ ตตฺถ สมฺมุขา\-วินยสฺมิํ, สํฆสมฺมุขตา ธมฺม\-สมฺมุขตา วินย\-สมฺมุขตา ปุคฺคล\-สมฺมุขตา ฯเปฯ. กิญฺจ ตตฺถ อมูฬฺห\-วินยสฺมิํ, ยา อมูฬฺห\-วินยสฺส กมฺมสฺส กิริยา กรณํ อุปคมนํ อชฺฌุปคมนํ อธิวาสนา อปฺปฏิกฺโกสนา, อิทํ ตตฺถ อมูฬฺห\-วินยสฺมิํ. เอวํ วูปสนฺตํ เจ ภิกฺขเว อธิกรณํ การโก อุกฺโกเฏติ, อุกฺโกฏนกํ ปาจิตฺติยํ. ฉนฺททายโก ขียติ, ขียนกํ ปาจิตฺติยํ.}

\csromanmatikaanchor{mtk.123}
\csromantocmarkref{subsection}{ตสฺสปาปิยสิกาวินย}{mtk.123}
\bookmark[dest={mtk.123},level=2]{ตสฺสปาปิยสิกาวินย}
\csromanheader{2}{ตสฺสปาปิยสิกา\-วินย}
\proseitem{238}{สิยา อนุวาทา\-ธิกรณํ เทฺว สมเถ อนาคมฺม สติวินยญฺจ อมูฬฺหวินยญฺจ ทฺวีหิ สมเถหิ สมฺเมยฺย สมฺมุขา\-วินเยน จ ตสฺสปาปิยสิกาย จาติ? สิยาติสฺส วจนียํ. ยถา กถํ วิย. อิธ ปน ภิกฺขเว ภิกฺขุ ภิกฺขุํ สํฆมชฺเฌ ครุกาย อาปตฺติยา โจเทติ “สรตายสฺมา เอวรูปิํ ครุกํ อาปตฺติํ อาปชฺชิตา ปาราชิกํ วา ปาราชิก\-สามนฺตํ วา”ติ. โส เอวํ วเทติ “น โข อหํ อาวุโส สรามิ เอวรูปิํ ครุกํ อาปตฺติํ อาปชฺชิตา ปาราชิกํ วา ปาราชิก\-สามนฺตํ วา”ติ. ตเมนํ โส นิพฺเพเฐนฺตํ อติเวเฐติ “อิงฺฆายสฺมา สาธุกเมว ชานาติ, ยทิ สรสิ เอวรูปิํ ครุกํ อาปตฺติํ อาปชฺชิตา ปาราชิกํ วา ปาราชิก\-สามนฺตํ วา”ติ. โส เอวํ วเทติ “น โข อหํ อาวุโส สรามิ เอวรูปิํ ครุกํ อาปตฺติํ อาปชฺชิตา ปาราชิกํ วา ปาราชิก\-สามนฺตํ วา, สรามิ จ โข อหํ อาวุโส เอวรูปิํ อปฺปมตฺติกํ อาปตฺติํ อาปชฺชิตา”ติ. ตเมนํ โส นิพฺเพเฐนฺตํ อติเวเฐติ “อิงฺฆายสฺมา สาธุกเมว ชานาหิ, ยทิ สรสิ เอวรูปิํ ครุกํ อาปตฺติํ อาปชฺชิตา ปาราชิกํ วา ปาราชิก\-สามนฺตํ วา”ติ. โส เอวํ วเทติ “อิมํ หิ นามาหํ อาวุโส อปฺปมตฺติกํ อาปตฺติํ อาปชฺชิตฺวา อปุฏฺโฐ ปฏิชานิสฺสามิ, กิํ ปนาหํ เอวรูปิํ ครุกํ อาปตฺติํ อาปชฺชีตฺวา ปาราชิกํ วา ปาราชิก\-สามนฺตํ วา ปุฏฺโฐ น ปฏิชานิสฺสามี”ติ. โส เอวํ วเทติ “อิมํ หิ นาม ตฺวํ อาวุโส อปฺปมตฺติกํ อาปตฺติํ อาปชฺชิตฺวา อปุฏฺโฐ น ปฏิชานิสฺสสิ, กิํ ปน ตฺวํ เอวรูปิํ ครุกํ อาปตฺติํ อาปชฺชิตฺวา ปาราชิกํ วา ปาราชิก\-สามนฺตํ วา อปุฏฺโฐ ปฏิชานิสฺสสิ, อิงฺฆายสฺมา สาธุกเมว ชานาหิ, ยทิ สรสิ เอวรูปิํ ครุกํ อาปตฺติํ อาปชฺชิตา ปาราชิกํ วา ปาราชิก\-สามนฺตํ วา”ติ. โส เอวํ วเทสิ “สรามิ โข อหํ อาวุโส \csromanfolio{233}เอวรูปิํ ครุกํ อาปตฺติํ อาปชฺชิตา ปาราชิกํ วา ปาราชิก\-สามนฺตํ วา, ทวา เม เอตํ วุตฺตํ, รวา เม เอตํ วุตฺตํ, นาหํ ตํ สรามิ เอวรูปิํ ครุกํ อาปตฺติํ อาปชฺชิตา ปาราชิกํ วา ปาราชิก\-สามนฺตํ วา”ติ. ตสฺส โข ภิกฺขเว\csromansharedfootnote{c556731f51dd52fc}{ตสฺส โข ตํ ภิกฺขเว (ก), ตสฺส เขฺวตํ ภิกฺขเว (สฺยา)} ภิกฺขุโน ตสฺสปาปิยสิกา\-กมฺมํ กาตพฺพํ. เอวญฺจ ปน ภิกฺขเว กาตพฺพํ, พฺยตฺเตน ภิกฺขุนา ปฏิพเลน สํโฆ ญาเปตพฺโพ\csromandash{}}

\prose{“สุณาตุ เม ภนฺเต สํโฆ, อยํ อิตฺถนฺนาโม ภิกฺขุ สํฆมชฺเฌ ครุกาย อาปตฺติยา อนุยุญฺชิยมาโน อวชานิตฺวา ปฏิชานาติ, ปฏิชานิตฺวา อวชานาติ, อญฺเญนาญฺญํ ปฏิจรติ, สมฺปชานมุสา ภาสติ, ยทิ สํฆสฺส ปตฺตกลฺลํ, สํโฆ อิตฺถนฺนามสฺส ภิกฺขุโน ตสฺสปาปิยสิกา\-กมฺมํ กเรยฺย, เอสา ญตฺติ.}

\prose{สุณาตุ เม ภนฺเต สํโฆ, อยํ อิตฺถนฺนาโม ภิกฺขุ สํฆมชฺเฌ ครุกาย อาปตฺติยา อนุยุญฺชิยมาโน อวชานิตฺวา ปฏิชานาติ ปฏิชานิตฺวา อวชานาติ, อญฺเญนาญฺญํ ปฏิจรติ, สมฺปชานมุสา ภาสติ. สํโฆ อิตฺถนฺนามสฺส ภิกฺขุโน ตสฺสปาปิยสิกา\-กมฺมํ กโรติ, ยสฺสายสฺมโต ขมติ อิตฺถนฺนามสฺส ภิกฺขุโน ตสฺสปาปิยสิกา\-กมฺมสฺส กรณํ, โส ตุณฺหสฺส, ยสฺส นกฺขมติ, โส ภาเสยฺย.}

\prose{ทุติยมฺปิ เอตมตฺถํ วทามิ ฯเปฯ. ตติยมฺปิ เอตมตฺถํ วทามิ ฯเปฯ.}

\prose{กตํ สํเฆน อิตฺถนฺนามสฺส ภิกฺขุโน ตสฺสปาปิยสิกา\-กมฺมํ, ขมติ สํฆสฺส, ตสฺมา ตุณฺหี, เอวเมตํ ธารยามี”ติ.}

\prose{อิทํ วุจฺจติ ภิกฺขเว อธิกรณํ วูปสนฺตํ. เกน วูปสนฺตํ, สมฺมุขา\-วินเยน จ ตสฺสปาปิยสิกาย จ. กิญฺจ ตตฺถ สมฺมุขา\-วินยสฺมิํ, สํฆสมฺมุขตา ธมฺม\-สมฺมุขตา วินย\-สมฺมุขตา ฯเปฯ. กา จ ตตฺถ ตสฺสปาปิยสิกาย, ยา ตสฺสปาปิยสิกา\-กมฺมสฺส กิริยา กรณํ อุปคมนํ อชฺฌุปคมนํ อธิวาสนา อปฺปฏิกฺโกสนา, อยํ ตตฺถ ตสฺสปาปิยสิกาย. เอวํ วูปสนฺตํ เจ ภิกฺขเว อธิกรณํ การโก อุกฺโกเฏติ, อุกฺโกฏนกํ ปาจิตฺติยํ. ฉนฺททายโก ขียติ, ขียนกํ ปาจิตฺติยํ.}

\csromanmatikaanchor{mtk.124}
\csromantocmarkref{subsection}{ปฏิญฺญาตกรณ}{mtk.124}
\bookmark[dest={mtk.124},level=2]{ปฏิญฺญาตกรณ}
\csromanheader{2}{ปฏิญฺญาต\-กรณ}
\proseitem{239}{\csromansymbolfootnote{*}{วิ 5. 184, 195, 198 ปิฏฺฐาทีสุปิ.}อาปตฺตา\-ธิกรณํ กติหิ สมเถหิ สมฺมติ. อาปตฺตา\-ธิกรณํ ตีหิ สมเถหิ สมฺมติ สมฺมุขา\-วินเยน จ ปฏิญฺญาต\-กรเณน จ \csromanfolio{234}ติณวตฺถารเกน จ. สิยา อาปตฺตา\-ธิกรณํ เอกํ สมถํ อนาคมฺม ติณวตฺถารกํ ทฺวีหิ สมเถหิ สมฺเมยฺย สมฺมุขา\-วินเยน จ ปฏิญฺญาต\-กรเณน จาติ? สิยาติสฺส วจนียํ. ยถา กถํ วิย. อิธ ปน ภิกฺขเว ภิกฺขุ ลหุกํ อาปตฺติํ อาปนฺโน โหติ, เตน ภิกฺขเว ภิกฺขุโน เอกํ ภิกฺขุํ อุปสงฺกมิตฺวา เอกํสํ อุตฺตราสงฺคํ กริตฺวา อุกฺกุฏิกํ นิสีทิตฺวา อญฺชลิํ ปคฺคเหตฺวา เอวมสฺส วจนีโย “อหํ อาวุโส อิตฺถนฺนามํ อาปตฺติํ อาปนฺโน, ตํ ปฏิเทเสมี”ติ. เตน วตฺตพฺโพ “ปสฺสสี”ติ. “อาม ปสฺสามี”ติ. “อายติํ สํวเรยฺยาสี”ติ.}

\prose{อิทํ วุจฺจติ ภิกฺขเว อธิกรณํ วูปสนฺตํ. เกน วูปสนฺตํ, สมฺมุขา\-วินเยน จ ปฏิญฺญาต\-กรเณน จ. กิญฺจ ตตฺถ สมฺมุขา\-วินยสฺมิํ, ธมฺม\-สมฺมุขตา วินย\-สมฺมุขตา ปุคฺคล\-สมฺมุขตา ฯเปฯ. กา จ ตตฺถ ปุคฺคล\-สมฺมุขตา, โย จ เทเสติ, ยสฺส จ เทเสติ, อุโภ สมฺมุขีภูตา โหนฺติ. อยํ ตตฺถ ปุคฺคล\-สมฺมุขตา. กิญฺจ ตตฺถ ปฏิญฺญาต\-กรณสฺมิํ, ยา ปฏิญฺญาต\-กรณสฺส กมฺมสฺส กิริยา กรณํ อุปคมนํ อชฺฌุปคมนํ อธิวาสนา อปฺปฏิกฺโกสนา, อิทํ ตตฺถ ปฏิญฺญาต\-กรณสฺมิํ. เอวํ วูปสนฺตํ เจ ภิกฺขเว อธิกรณํ ปฏิคฺคาหโก อุกฺโกเฏติ, อุกฺโกฏนกํ ปาจิตฺติยํ. เอวญฺเจตํ ลเภถ, อิจฺเจตํ กุสลํ. โน เจ ลเภถ, เตน ภิกฺขเว ภิกฺขุนา สมฺพหุเล ภิกฺขู อุปสงฺกมิตฺวา เอกํสํ อุตฺตราสงฺคํ กริตฺวา วุฑฺฒานํ ภิกฺขูนํ ปาเท วนฺทิตฺวา อุกฺกุฏิกํ นิสีทิตฺวา อญฺชลิํ ปคฺคเหตฺวา เอวมสฺสุ วจนียา “อหํ ภนฺเต อิตฺถนฺนามํ อาปตฺติํ อาปนฺโน,ตํ ปฏิเทเสมี”ติ. พฺยตฺเตน ภิกฺขุนา ปฏิพเลน เต ภิกฺขู ญาเปตพฺพา\csromandash{}}

\prose{“สุณนฺตุ เม อายสฺมนฺตา, อยํ อิตฺถนฺนาโม ภิกฺขุ อาปตฺติํ สรติ วิวรติ อุตฺตานิํ กโรติ เทเสติ, ยทา\-ยสฺมนฺตานํ ปตฺตกลฺลํ, อหํ อิตฺถนฺนามสฺส ภิกฺขุโน อาปตฺติํ ปฏิคฺคเณฺหยฺยนฺ”ติ. เตน วตฺตพฺโพ “ปสฺสสี”ติ. “อาม ปสฺสามี”ติ. “อายติํ สํวเรยฺยาสี”ติ.}

\prose{อิทํ วุจฺจติ ภิกฺขเว อธิกรณํ วูปสนฺตํ. เกน วูปสนฺตํ, สมฺมุขา\-วินเยน จ ปฏิญฺญาต\-กรเณน จ. กิญฺจ ตตฺถ สมฺมุขา\-วินยสฺมิํ, ธมฺม\-สมฺมุขตา วินย\-สมฺมุขตา ปุคฺคล\-สมฺมุขตา ฯเปฯ. กา จ ตตฺถ ปุคฺคล\-สมฺมุขตา, โย จ เทเสติ, ยสฺส จ เทเสติ, อุโภ สมฺมุขีภูตา \csromanfolio{235}โหนฺติ. อยํ ตตฺถ ปุคฺคล\-สมฺมุขตา. กิญฺจ ตตฺถ ปฏิญฺญาต\-กรณสฺมิํ, ยา ปฏิญฺญาต\-กรณสฺส กมฺมสฺส กิริยา กรณํ อุปคมนํ อชฺฌุปคมนํ อธิวาสนา อปฺปฏิกฺโกสนา, อิทํ ตตฺถ ปฏิญฺญาต\-กรณสฺมิํ. เอวํ วูปสนฺตํ เจ ภิกฺขเว อธิกรณํ ปฏิคฺคาหโก อุกฺโกเฏติ, อุกฺโกฏนกํ ปาจิตฺติยํ. เอวญฺเจตํ ลเภถ, อิจฺเจตํ กุสลํ. โน เจ ลเภถ, เตน ภิกฺขเว ภิกฺขุนา สํฆํ อุปสงฺกมิตฺวา เอกํสํ อุตฺตราสงฺคํ กริตฺวา วุฑฺฒานํ ภิกฺขูนํ ปาเท วนฺทิตฺวา อุกฺกุฏิกํ นิสีทิตฺวา อญฺชลิํ ปคฺคเหตฺวา เอวมสฺส วจนีโย “อหํ ภนฺเต อิตฺถนฺนามํ อาปตฺติํ อาปนฺโน, ตํ ปฏิเทเสมี”ติ. พฺยตฺเตน ภิกฺขุนา ปฏิพเลน สํโฆ ญาเปตพฺโพ\csromandash{}}

\prose{“สุณาตุ เม ภนฺเต สํโฆ, อยํ อิตฺถนฺนาโม ภิกฺขุ อาปตฺติํ สรติ วิวรติ อุตฺตานิํ กโรติ เทเสติ, ยทิ สํฆสฺส ปตฺตกลฺลํ, อหํ อิตฺถนฺนามสฺส ภิกฺขุโน อาปตฺติํ ปฏิคฺคเณฺหยฺยนฺ”ติ. เตน วตฺตพฺโพ “ปสฺสสี”ติ. “อาม ปสฺสามี”ติ. “อายติํ สํวเรยฺยาสี”ติ.}

\prose{อิทํ วุจฺจติ ภิกฺขเว อธิกรณํ วูปสนฺตํ. เกน วูปสนฺตํ, สมฺมุขา\-วินเยน จ ปฏิญฺญาต\-กรเณน จ. กิญฺจ ตตฺถ สมฺมุขา\-วินยสฺมิํ, สํฆสมฺมุขตา ธมฺม\-สมฺมุขตา วินย\-สมฺมุขตา ปุคฺคล\-สมฺมุขตา ฯเปฯ. เอวํ วูปสนฺตํ เจ ภิกฺขเว อธิกรณํ ปฏิคฺคาหโก อุกฺโกเฏติ, อุกฺโกฏนกํ ปาจิตฺติยํ. ฉนฺททายโก ขียติ, ขียนกํ ปาจิตฺติยํ.}

\csromanmatikaanchor{mtk.125}
\csromantocmarkref{subsection}{ติณวตฺถารก}{mtk.125}
\bookmark[dest={mtk.125},level=2]{ติณวตฺถารก}
\csromanheader{2}{\textbf{ติณวตฺถารก}}
\proseitem{240}{สิยา อาปตฺตา\-ธิกรณํ เอกํ สมถํ อนาคมฺม ปฏิญฺญาต\-กรณํ ทฺวีหิ สมเถหิ สมฺเมยฺย สมฺมุขา\-วินเยน จ ติณวตฺถารเกน จาติ? สิยาติสฺส วจนียํ. ยถา กถํ วิย. อิธ ปน ภิกฺขเว ภิกฺขูนํ ภณฺฑน\-ชาตานํ กลหชาตานํ วิวาทาปนฺนานํ วิหรตํ พหุํ อสฺสามณกํ อชฺฌาจิณฺณํ โหติ ภาสิต\-ปริกฺกนฺตํ. ตตฺร เจ ภิกฺขูนํ เอวํ โหติ “อมฺหากํ โข ภณฺฑน\-ชาตานํ กลหชาตานํ วิเวทาปนฺนานํ วิหรตํ พหุํ อสฺสามณกํ อชฺฌาจิณฺณํ ภาสิต\-ปริกฺกนฺตํ, สเจ มยํ อิมาหิ อาปตฺตีหิ อญฺญมญฺญํ กาเรสฺสาม, สิยาปิ ตํ อธิกรณํ กกฺขฬตฺตาย วาฬตฺตาย เภทาย สํวตฺเตยฺยา”ติ. อนุชานามิ ภิกฺขเว เอวรูปํ อธิกรณํ ติณวตฺถารเกน วูปสเมตุํ. เอวญฺจ ปน ภิกฺขเว วูปสเมตพฺพํ, \csromanfolio{236}สพฺเพเหว เอกชฺฌํ สนฺนิปติตพฺพํ, สนฺนิปติตฺวา พฺยตฺเตน ภิกฺขุนา ปฏิพเลน สํโฆ ญาเปตพฺโพ\csromandash{}}

\prose{“สุณาตุ เม ภนฺเต สํโฆ, อมฺหากํ ภณฺฑน\-ชาตานํ กลหชาตานํ วิวาทาปนฺนานํ วิหรตํ พหุํ อสฺสามณกํ อชฺฌาจิณฺณํ ภาสิต\-ปริกฺกนฺตํ, สเจ มยํ อิมาหิ อาปตฺตีหิ อญฺญมญฺญํ กาเรสฺสาม, สิยาปิ ตํ อธิกรณํ กกฺขฬตฺตาย วาฬตฺตาย เภทาย สํวตฺเตยฺย, ยทิ สํฆสฺส ปตฺตกลฺลํ, สํโฆ อิมํ อธิกรณํ ติณวตฺถารเกน วูปสเมยฺย ฐเปตฺวา ถุลฺลวชฺชํ, ฐเปตฺวา คิหิปฺปฏิสํยุตฺตนฺ”ติ. เอกโต\-ปกฺขิกานํ ภิกฺขูนํ พฺยตฺเตน ภิกฺขุนา ปฏิพเลน สโก ปกฺโข ญาเปตพฺโพ\csromandash{}}

\prose{“สุณนฺตุ เม อายสฺมนฺตา, อมฺหากํ ภณฺฑน\-ชาตานํ กลหชาตานํ วิวาทาปนฺนานํ วิหรตํ พหุํ อสฺสามณกํ อชฺฌาจิณฺณํ ภาสิต\-ปริกฺกนฺตํ, สเจ มยํ อิมาหิ อาปตฺตีหิ อญฺญมญฺญํ กาเรสฺสาม, สิยาปิ ตํ อธิกรณํ กกฺขฬตฺตาย วาฬตฺตาย เภทาย สํวตฺเตยฺย, ยทา\-ยสฺมนฺตานํ ปตฺตกลฺลํ, อหํ ยา เจว อายสฺมนฺตานํ อาปตฺติ, ยา จ อตฺตโน อาปตฺติ, อายสฺมนฺตาน\-ญฺเจว อตฺถาย อตฺตโน จ อตฺถาย สํฆมชฺเฌ ติณวตฺถารเกน เทเสยฺยํ ฐเปตฺวา ถุลฺลวชฺชํ, ฐเปตฺวา คิหิปฺปฏิสํยุตฺตนฺ”ติ.}

\proseitem{241}{อถาปเรสํ เอกโต\-ปกฺขิกานํ ภิกฺขูนํ พฺยตฺเตน ภิกฺขุนา ปฏิพเลน สโก ปกฺโข ญาเปตพฺโพ\csromandash{}}

\prose{“สุณนฺตุ เม อายสฺมนฺตา, อมฺหากํ ภณฺฑน\-ชาตานํ กลหชาตานํ วิวาทาปนฺนานํ วิหรตํ พหุํ อสฺสามณกํ อชฺฌาจิณฺณํ ภาสิต\-ปริกฺกนฺตํ, สเจ มยํ อิมาหิ อาปตฺตีหิ อญฺญมญฺญํ กาเรสฺสาม, สิยาปิ ตํ อธิกรณํ กกฺขฬตฺตาย วาฬตฺตาย เภทาย สํวตฺเตยฺย, ยทา\-ยสฺมนฺตานํ ปตฺตกลฺลํ, อหํ ยา เจว อายสฺมนฺตานํ อาปตฺติ, ยา จ อตฺตโน อาปตฺติ, อายสฺมนฺตาน\-ญฺเจว อตฺถาย อตฺตโน จ อตฺถาย สํฆมชฺเฌ ติณวตฺถารเกน เทเสยฺยํ ฐเปตฺวา ถุลฺลวชฺชํ, ฐเปตฺวา คิหิปฺปฏิสํยุตฺตนฺ”ติ.}

\proseitem{242}{อถาปเรสํ เอกโต\-ปกฺขิกานํ ภิกฺขูนํ พฺยตฺเตน ภิกฺขุนา ปฏิพเลน สํโฆ ญาเปตพฺโพ\csromandash{}}

\prose{\csromanfolio{237}“สุณาตุ เม ภนฺเต สํโฆ, อมฺหากํ ภณฺฑน\-ชาตานํ กลหชาตานํ วิวาทาปนฺนานํ วิหรตํ พหุํ อสฺสามณกํ อชฺฌาจิณฺณํ ภาสิต\-ปริกฺกนฺตํ, สเจ มยํ อิมาหิ อาปตฺตีหิ อญฺญมญฺญํ กาเรสฺสาม, สิยาปิ ตํ อธิกรณํ กกฺขฬตฺตาย วาฬตฺตาย เภทาย สํวตฺเตยฺย, ยทิ สํฆสฺส ปตฺตกลฺลํ, อหํ ยา เจว อิเมสํ อายสฺมนฺตานํ อาปตฺติ, ยา จ อตฺตโน อาปตฺติ, อิเมสญฺเจว อายสฺมนฺตานํ อตฺถาย อตฺตโน จ อตฺถาย สํฆมชฺเฌ ติณวตฺถารเกน เทเสยฺยํ ฐเปตฺวา ถุลฺลวชฺชํ, ฐเปตฺวา คิหิปฺปฏิสํยุตฺตํ, เอสา ญตฺติ.}

\prose{สุณาตุ เม ภนฺเต สํโฆ, อมฺหากํ ภณฺฑน\-ชาตานํ กลหชาตานํ วิวาทาปนฺนานํ วิหรตํ พหุํ อสฺสามณกํ อชฺฌาจิณฺณํ ภาสิต\-ปริกฺกนฺตํ, สเจ มยํ อิมาหิ อาปตฺตีหิ อญฺญมญฺญํ กาเรสฺสาม, สิยาปิ ตํ อธิกรณํ กกฺขฬตฺตาย วาฬตฺตาย เภทาย สํวตฺเตยฺย, อหํ ยา เจว อิเมสํ อายสฺมนฺตานํ อาปตฺติ, ยา จ อตฺตโน อาปตฺติ, อิเมสญฺเจจ อายสฺมนฺตานํ อตฺถาย อตฺตโน จ อตฺถาย สํฆมชฺเฌ ติณวตฺถารเกน เทเสมิ ฐเปตฺวา ถุลฺลวชฺชํ, ฐเปตฺวา คิหิปฺปฏิสํยุตฺตํ, ยสฺสายสฺมโต ขมติ อมฺหากํ อิมาสํ อาปตฺตีนํ สํฆมชฺเฌ ติณวตฺถารเกน เทสนา ฐเปตฺวา ถุลฺลวชฺชํ, ฐเปตฺวา คิหิปฺปฏิสํยุตฺตํ, โส ตุณฺหสฺส, ยสฺส นกฺขมติ, โส ภาเสยฺย.}

\prose{เทสิตา อมฺหากํ อิมา อาปตฺติโย สํฆมชฺเฌ ติณวตฺถารเกน ฐเปตฺวา ถุลฺลวชฺชํ, ฐเปตฺวา คิหิปฺปฏิสํยุตฺตํ, ขมติ สํฆสฺส, ตสฺมา ตุณฺหี, เอวเมตํ ธารยามี”ติ.}

\prose{อถาปเรสํ ฯเปฯ เอวเมตํ ธารยามีติ.}

\prose{อิทํ วุจฺจติ ภิกฺขเว อธิกรณํ วูปสนฺตํ. เกน วูปสนฺตํ, สมฺมุขา\-วินเยน จ ติณวตฺถารเกน จ. กิญฺจ ตตฺถ สมฺมุขา\-วินยสฺมิํ, สํฆสมฺมุขตา ธมฺม\-สมฺมุขตา วินย\-สมฺมุขตา ปุคฺคล\-สมฺมุขตา.}

\prose{กา จ ตตฺถ สํฆสมฺมุขตา, ยาวติกา ภิกฺขู กมฺมปฺปตฺตา, เต อาคตา โหนฺติ, ฉนฺทารหานํ ฉนฺโท อาหโฏ โหติ สมฺมุขีภูตา น ปฏิกฺโกสนฺติ, อยํ ตตฺถ สํฆสมฺมุขตา.}

\prose{\csromanfolio{238}กา จ ตตฺถ ธมฺม\-สมฺมุขตา วินย\-สมฺมุขตา, เยน ธมฺเมน เยน วินเยน เยน สตฺถุสาสเนน ตํ อธิกรณํ วูปสมฺมติ, อยํ ตตฺถ ธมฺม\-สมฺมุขตา วินย\-สมฺมุขตา. กา จ ตตฺถ ปุคฺคล\-สมฺมุขตา, โย จ เทเสติ, ยสฺส จ เทเสติ, อุโภ สมฺมุขีภูตา โหติ, อยํ ตตฺถ ปุคฺคล\-สมฺมุขตา.}

\prose{กิญฺจ ตตฺถ ติณวตฺถารกสฺมิํ, ยา ติณวตฺถารกสฺส กมฺมสฺส กิริยา กรณํ อุปคมนํ อชฺฌุปคมนํ อธิวาสนา อปฺปฏิกฺโกสนา, อิทํ ตตฺถ ติณวตฺถารกสฺมิํ. เอวํ วูปสนฺตํ เจ ภิกฺขเว อธิกรณํ ปฏิคฺคาหโก อุกฺโกเฏติ, อุกฺโกฏนกํ ปาจิตฺติยํ. ฉนฺททายโก ขียติ, ขียนกํ ปาจิตฺติยํ.}

\prose{\csromansymbolfootnote{*}{วิ 5. 185, 195, 198 ปิฏฺฐาทีสุปิ.}กิจฺจา\-ธิกรณํ กติหิ สมเถหิ สมฺมติ. กิจฺจา\-ธิกรณํ เอเกน สมเถน สมฺมติ สมฺมุขาวินเยนาติ.}

\vspace{\tipitakaparskip}
\csromancenter{สมถ\-กฺขนฺธโก นิฏฺฐิโต จตุตฺโถ.}
\csromanmatikaanchor{mtk.126}
\csromantocmarknopage{chapter}{5. ขุทฺทกวตฺถุกฺขนฺธก}{mtk.126}
\bookmark[dest={mtk.126},level=0]{5. ขุทฺทกวตฺถุกฺขนฺธก}
\chapterheadpageread{5. ขุทฺทก\-วตฺถุ\-กฺขนฺธก}
\csromanmatikaanchor{mtk.127}
\csromantocmarkref{section}{ขุทฺทกวตฺถูนิ}{mtk.127}
\bookmark[dest={mtk.127},level=1]{ขุทฺทกวตฺถูนิ}
\csromanheader{1}{ขุทฺทก\-วตฺถูนิ}
\proseitem{243}{\csromanfolio{239}เตน สมเยน พุทฺโธ ภควา ราชคเห วิหรติ เวฬุวเน กลนฺทก\-นิวาเป. เตน โข ปน สมเยน ฉพฺพคฺคิยา ภิกฺขู นหายมานา รุกฺเข กายํ อุคฺฆํเสนฺติ, อูรุมฺปิ พาหุมฺปิ อุรมฺปิ ปิฏฺฐิมฺปิ. มนุสฺสา อุชฺฌายนฺติ ขิยฺยนฺติ วิปาเจนฺติ “กถํ หิ นาม สมณา สกฺยปุตฺติยา นหายมานา รุกฺเข กายํ อุคฺฆํเสสฺสนฺติ, อูรุมฺปิ พาหุมฺปิ อุรมฺปิ ปิฏฺฐิมฺปิ เสยฺยถาปิ มลฺลมุฏฺฐิกา คามโมทฺทวา”ติ\csromansharedfootnote{47f0b8e2fd16ff32}{คามโปทฺทวา (สี), คามปูตวา (สฺยา)}. อสฺโสสุํ โข ภิกฺขู เตสํ มนุสฺสานํ อุชฺฌายนฺตานํ ขิยฺยนฺตานํ วิปาเจนฺตานํ. อถ โข เต ภิกฺขู ภควโต เอตมตฺถํ อาโรเจสุํ. อถ โข ภควา เอตสฺมิํ นิทาเน เอตสฺมิํ ปกรเณ ภิกฺขุสํฆํ สนฺนิปาตาเปตฺวา ภิกฺขู ปฏิปุจฺฉิ “สจฺจํ กิร ภิกฺขเว ฉพฺพคฺคิยา ภิกฺขู นหายมานา รุกฺเข กายํ อุคฺฆํเสนฺติ, อูรุํปิ พาหุํปิ อุรํปิ ปิฏฺฐิํปี”ติ. สจฺจํ ภควาติ. วิครหิ พุทฺโธ ภควา อนนุจฺฉวิกํ ภิกฺขเว เตสํ โมฆปุริสานํ อนนุโลมิกํ อปฺปติรูปํ อสฺสามณกํ อกปฺปิยํ อกรณียํ, กถํ หิ นาม เต ภิกฺขเว โมฆปุริสา นหายมานา รุกฺเข กายํ อุคฺฆํเสสฺสนฺติ, อูรุมฺปิ พาหุมฺปิ อุรมฺปิ ปิฏฺฐิมฺปิ. เนตํ ภิกฺขเว อปฺปสนฺนานํ วา ปสาทาย ฯเปฯ วิครหิตฺวา ฯเปฯ ธมฺมิํ กถํ กตฺวา ภิกฺขู อามนฺเตสิ “น ภิกฺขเว นหายมาเนน ภิกฺขุนา รุกฺเข กาโย อุคฺฆํเสตพฺโพ, โย อุคฺฆํเสยฺย, อาปตฺติ ทุกฺกฏสฺสา”ติ.}

\prose{เตน โข ปน สมเยน ฉพฺพคฺคิยา ภิกฺขู นหายมานา ถมฺเภ กายํ อุคฺฆํเสนฺติ, อูรุมฺปิ พาหุมฺปิ อุรมฺปิ ปิฏฺฐิมฺปิ. มนุสฺสา อุชฺฌายนฺติ ขิยฺยนฺติ วิปาเจนฺติ “กถํ หิ นาม สมณา สกฺยปุตฺติยา นหายมานา ถมฺเภ กายํ อุคฺฆํเสสฺสนฺติ, อูรุมฺปิ พาหุมฺปิ อุรมฺปิ ปิฏฺฐิมฺปิ เสยฺยถาปิ มลฺลมุฏฺฐิกา คามโมทฺทวา”ติ. อสฺโสสุํ โข ภิกฺขู เตสํ มนุสฺสานํ อุชฺฌายนฺตานํ ขิยฺยนฺตานํ วิปาเจนฺตานํ. อถ โข เต ภิกฺขู ภควโต เอตมตฺถํ อาโรเจสุํ ฯเปฯ. สจฺจํ ภควาติ ฯเปฯ วิครหิตฺวา ฯเปฯ ธมฺมิํ กถํ กตฺวา ภิกฺขู อามนฺเตสิ “น ภิกฺขเว นหายมาเนน ภิกฺขุนา ถมฺเภ กาโย อุคฺฆํเสตพฺโพ, โย อุคฺฆํเสยฺย, อาปตฺติ ทุกฺกฏสฺสา”ติ.}

\prose{\csromanfolio{240}เตน โข ปน สมเยน ฉพฺพคฺคิยา ภิกฺขู นหายมานา กุฏฺเฏ\csromansharedfootnote{3a39bb80825f35ec}{กุฑฺเฑ (สี, สฺยา)} กายํ อุคฺฆํเสนฺติ, อูรุมฺปิ พาหุมฺปิ อุรมฺปิ ปิฏฺฐิมฺปิ. มนุสฺสา อุชฺฌายนฺติ ขิยฺยนฺติ วิปาเจนฺติ “กถํ หิ นาม สมณา สกฺยปุตฺติยา นหายมานา กุฏฺเฏ กายํ อุคฺฆํเสสฺสนฺติ, อูรุมฺปิพาหุมฺปิ อุรมฺปิ ปิฏฺฐิมฺปิ เสยฺยถาปิ มลฺลมุฏฺฐิกา คามโมทฺทวา”ติ ฯเปฯ น ภิกฺขเว นหายมาเนน ภิกฺขุนา กุฏฺเฏ กาโย อุคฺฆํเสตพฺโพ, โย อุคฺฆํเสยฺย, อาปตฺติ ทุกฺกฏสฺสาติ.}

\prose{เตน โข ปน สมเยน ฉพฺพคฺคิยา ภิกฺขู อฏฺฏาเน\csromansharedfootnote{18d68139d0c96a04}{อฏฺฐาเน (สี, สฺยา)} นหายนฺติ. มนุสฺสา อุชฺฌายนฺติ ขิยฺยนฺติ วิปาเจนฺติ “ฯเปฯ เสยฺยถาปิ คิหี กามโภคิโน”ติ. อสฺโสสุํ โข ภิกฺขู เตสํ มนุสฺสานํ อุชฺฌายนฺตานํ ขิยฺยนฺตานํ วิปาเจนฺตานํ. อถ โข เต ภิกฺขู ภควโต เอตมตฺถํ อาโรเจสุํ ฯเปฯ. สจฺจํ ภควาติ ฯเปฯ วิครหิตฺวา ฯเปฯ ธมฺมิํ กถํ กตฺวา ภิกฺขู อามนฺเตสิ “น ภิกฺขเว อฏฺฏาเน นหายิตพฺพํ, โย นหาเยยฺย, อาปตฺติ ทุกฺกฏสฺสา”ติ.}

\prose{เตน โข ปน สมเยน ฉพฺพคฺคิยา ภิกฺขู คนฺธพฺพหตฺถเกน นหายนฺติ. มนุสฺสา อุชฺฌายนฺติ ขิยฺยนฺติ วิปาเจนฺติ “ฯเปฯ เสยฺยถาปิ คิหี กามโภคิโน”ติ. อสฺโสสุํ โข ภิกฺขู เตสํ มนุสฺสานํ อุชฺฌายนฺตานํ ขิยฺยนฺตานํ วิปาเจนฺตานํ. อถ โข เต ภิกฺขู ภควโต เอตมตฺถํ อาโรเจสุํ ฯเปฯ น ภิกฺขเว คนฺธพฺพหตฺถเกน นหายิตพฺพํ, โย นหาเยยฺย, อาปตฺติ ทุกฺกฏสฺสาติ.}

\prose{เตน โข ปน สมเยน ฉพฺพคฺคิยา ภิกฺขู กุรุวินฺทก\-สุตฺติยา นหายนฺติ. มนุสฺสา อุชฺฌายนฺติ ขิยฺยนฺติ วิปาเจนฺติ “ฯเปฯ เสยฺยถาปิ คิหี กามโภคิโน”ติ. ภควโต เอตมตฺถํ อาโรเจสุํ. น ภิกฺขเว กุรุวินฺทก\-สุตฺติยา นหายิตพฺพํ, โย นหาเยยฺย, อาปตฺติ ทุกฺกฏสฺสาติ.}

\prose{เตน โข ปน สมเยน ฉพฺพคฺคิยา ภิกฺขู วิคฺคยฺห ปริกมฺมํ การาเปนฺติ. มนุสฺสา อุชฺฌายนฺติ ขิยฺยนฺติ วิปาเจนฺติ “ฯเปฯ เสยฺยถาปิ คิหี กามโภคิโน”ติ. ภควโต เอตมตฺถํ อาโรเจสุํ. น ภิกฺขเว วิคฺคยฺห ปริกมฺมํ การาเปตพฺพํ, โย การาเปยฺย, อาปตฺติ ทุกฺกฏสฺสาติ.}

\prose{\csromanfolio{241}เตน โข ปน สมเยน ฉพฺพคฺคิยา ภิกฺขู มลฺลเกน นหายนฺติ. มนุสฺสา อุชฺฌายนฺติ ขิยฺยนฺติ วิปาเจนฺติ “ฯเปฯ เสยฺยถาปิ คิหี กามโภคิโน”ติ. ภควโต เอตมตฺถํ อาโรเจสุํ. น ภิกฺขเว มลฺลเกน นหายิตพฺพํ, โย นหาเยยฺย, อาปตฺติ ทุกฺกฏสฺสาติ.}

\proseitem{244}{เตน โข ปน สมเยน อญฺญตรสฺส ภิกฺขุโน กจฺฉุ\-โรคา\-พาโธ โหติ, น ตสฺส วินา มลฺลเกน ผาสุ โหติ. ภควโต เอตมตฺถํ อาโรเจสุํ. อนุชานามิ ภิกฺขเว คิลานสฺส อกตมลฺลกนฺติ.}

\prose{เตน โข ปน สมเยน อญฺญตโร ภิกฺขุ ชราทุพฺพโล นหายมาโน น สกฺโกติ อตฺตโน กายํ อุคฺฆํเสตุํ. ภควโต เอตมตฺถํ อาโรเจสุํ. อนุชานามิ ภิกฺขเว อุกฺกาสิกนฺติ.}

\prose{เตน โข ปน สมเยน ภิกฺขู ปิฏฺฐิ\-ปริกมฺมํ กาตุํ กุกฺกุจฺจายนฺติ. ภควโต เอตมตฺถํ อาโรเจสุํ. อนุชานามิ ภิกฺขเว ปุถุปาณิกนฺติ.}

\proseitem{245}{เตน โข ปน สมเยน ฉพฺพคฺคิยา ภิกฺขู วลฺลิกํ ธาเรนฺติ ฯเปฯ ปามงฺคํ ธาเรนฺติ ฯเปฯ กณฺฐสุตฺตกํ ธาเรนฺติ ฯเปฯ กฏิสุตฺตกํ ธาเรนฺติ ฯเปฯ โอวฏฺฏิกํ ธาเรนฺติ ฯเปฯ กายุรํ ธาเรนฺติ ฯเปฯ หตฺถาภรณํ ธาเรนฺติ ฯเปฯ องฺคุลิมุทฺทิกํ ธาเรนฺติ. มนุสฺสา อุชฺฌายนฺติ ขิยฺยนฺติ วิปาเจนฺติ “เสยฺยถาปิ คิหี กามโภคิโน”ติ. อสฺโสสุํ โข ภิกฺขู เตสํ มนุสฺสานํ อุชฺฌายนฺตานํ ขิยฺยนฺตานํ วิปาเจนฺตานํ. อถ โข เต ภิกฺขู ภควโต เอตมตฺถํ อาโรเจสุํ. สจฺจํ กิร ภิกฺขเว ฉพฺพคฺคิยา ภิกฺขู วลฺลิกํ ธาเรนฺติ ฯเปฯ ปามงฺคํ ธาเรนฺติ. กณฺฐสุตฺตกํ ธาเรนฺติ. กฏิสุตฺตกํ ธาเรนฺติ. โอวฏฺฏิกํ ธาเรนฺติ. กายุรํ ธาเรนฺติ. หตฺถาภรณํ ธาเรนฺติ. องฺคุลิมุทฺทิกํ ธาเรนฺตีติ. สจฺจํ ภควาติ ฯเปฯ วิครหิตฺวา ฯเปฯ ธมฺมิํ กถํ กตฺวา ภิกฺขู อามนฺเตสิ “น ภิกฺขเว วลฺลิกา ธาเรตพฺพา ฯเปฯ น ปามงฺโค ธาเรตพฺโพ. น กณฺฐสุตฺตกํ ธาเรตพฺพํ. น กฏิสุตฺตกํ ธาเรตพฺพํ. น โอวฏฺฏิกํ ธาเรตพฺพํ. น กายุรํ ธาเรตพฺพํ. น หตฺถาภรณํ ธาเรตพฺพํ. น องฺคุลิมุทฺทิกา ธาเรตพฺพา, โย ธาเรยฺย, อาปตฺติ ทุกฺกฏสฺสา”ติ.}

\proseitem{246}{เตน โข ปน สมเยน ฉพฺพคฺคิยา ภิกฺขู ทีเฆ เกเส ธาเรนฺติ. มนุสฺสา อุชฺฌายนฺติ ขิยฺยนฺติ วิปาเจนฺติ “เสยฺยถาปิ คิหี กามโภคิโน”ติ. ภควโต เอตมตฺถํ อาโรเจสุํ. น ภิกฺขเว ทีฆา \csromanfolio{242}เกสา ธาเรตพฺพา, โย ธาเรยฺย, อาปตฺติ ทุกฺกฏสฺส. อนุชานามิ ภิกฺขเว ทุมาสิกํ วา ทุวงฺคุลํ วาติ.}

\prose{เตน โข ปน สมเยน ฉพฺพคฺคิยา ภิกฺขู โกจฺเฉน เกเส โอสณฺเฐนฺติ\csromansharedfootnote{0f2282a74d5427d5}{โอสเณฺหนฺติ (สี, สฺยา)} ฯเปฯ ผณเกน เกเส โอสณฺเฐนฺติ. หตฺถ\-ผณเกน เกเส โอสณฺฐนฺติ. สิตฺถเตลเกน เกเส โอสณฺเฐนฺติ. อุทกเตลเกน เกเส โอสณฺเฐนฺติ. มนุสฺสา อุชฺฌายนฺติ ขิยฺยนฺติ วิปาเจนฺติ “เสยฺยถาปิ คิหี กามโภคิโน”ติ. ภควโต เอตมตฺถํ อาโรเจสุํ. น ภิกฺขเว โกจฺเฉน เกสา โอสณฺเฐตพฺพา ฯเปฯ. น สิตฺถเตลเกน เกสา โอสณฺเฐตพฺพา. น อุทกเตลเกน เกสา โอสณฺเฐตพฺพา, โย โอสณฺเฐยฺย, อาปตฺติ ทุกฺกฏสฺสาติ.}

\proseitem{247}{เตน โข ปน สมเยน ฉพฺพคฺคิยา ภิกฺขู อาทาเสปิ อุทกปตฺเตปิ มุขนิมิตฺตํ โอโลเกนฺติ. มนุสฺสา อุชฺฌายนฺติ ขิยฺยนฺติ วิปาเจนฺติ “เสยฺยถาปิ คิหี กามโภคิโน”ติ. ภควโต เอตมตฺถํ อาโรเจสุํ. น ภิกฺขเว อาทาเส วา อุทกปตฺเต วา มุขนิมิตฺตํ โอโลเกตพฺพํ, โย โอโลเกยฺย, อาปตฺติ ทุกฺกฏสฺสาติ.}

\prose{เตน โข ปน สมเยน อญฺญตรสฺส ภิกฺขุโน มุเข วโณ โหติ, โส ภิกฺขู ปุจฺฉิ “กีทิโส เม อาวุโส วโณ”ติ. ภิกฺขู เอวมาหํสุ “เอทิโส เต อาวุโส วโณ”ติ. โส น สทฺทหติ. ภควโต เอตมตฺถํ อาโรเจสุํ. อนุชานามี ภิกฺขเว อาพาธปฺปจฺจยา อาทาเส วา อุทกปตฺเต วา มุขนิมิตฺตํ โอโลเกตุนฺติ.}

\prose{เตน โข ปน สมเยน ฉพฺพคฺคิยา ภิกฺขู มุขํ อาลิมฺปนฺติ ฯเปฯ มุขํ อุมฺมทฺเทนฺติ. มุขํ จุณฺเณนฺติ. มโนสิลิกาย มุขํ ลญฺเฉนฺติ. องฺคราคํ กโรนฺติ. มุขราคํ กโรนฺติ. องฺคราค\-มุขราคํ กโรนฺติ. มนุสฺสา อุชฺฌายนฺติ ขิยฺยนฺติ วิปาเจนฺติ “เสยฺยถาปิ คิหี กามโภคิโน”ติ. ภควโต เอตมตฺถํ อาโรเจสุํ. น ภิกฺขเว มุขํ อาลิมฺปิตพฺพํ ฯเปฯ น มุขํ อุมฺมทฺทิตพฺพํ. น มุขํ จุณฺเณตพฺพํ. น มโนสิลิกาย มุขํ ลญฺเฉ ตพฺพํ. น องฺคราโค กาตพฺโพ. น มุขราโค กาตพฺโพ. น องฺคราค\-มุขราโค กาตพฺโพ, โย กเรยฺย, อาปตฺติ ทุกฺกฏสฺสาติ.}

\prose{\csromanfolio{243}เตน โข ปน สมเยน อญฺญตรสฺส ภิกฺขุโน จกฺขุ\-โรคา\-พาโธ โหติ. ภควโต เอตมตฺถํ อาโรเจสุํ. อนุชานามิ ภิกฺขเว อาพาธปฺปจฺจยา มุขํ อาลิมฺปิตุนฺติ.}

\proseitem{248}{เตน โข ปน สมเยน ราชคเห คิรคฺคสมชฺโช\csromansharedfootnote{c79192bea71c3268}{สมชฺชา (อภิธานคนฺเถสุ)} โหติ, ฉพฺพคฺคิยา ภิกฺขู คิรคฺค\-สมชฺชํ ทสฺสนาย อคมํสุ. มนุสฺสา อุชฺฌายนฺติ ขิยฺยนฺติ วิปาเจนฺติ “กถํ หิ นาม สมณา สกฺยปุตฺติยา นจฺจมฺปิ คีตมฺปิ วาทิตมฺปิ ทสฺสนาย คจฺฉิสฺสนฺติ เสยฺยถาปิ คิหี กามโภคิโน”ติ. ภควโต เอตมตฺถํ อาโรเจสุํ. น ภิกฺขเว นจฺจํ วา คีตํ วา วาทิตํ วา ทสฺสนาย คนฺตพฺพํ, โย คจฺเฉยฺย, อาปตฺติ ทุกฺกฏสฺสาติ.}

\proseitem{249}{เตน โข ปน สมเยน ฉพฺพคฺคิยา ภิกฺขู อายตเกน คีตสฺสเรน ธมฺมํ คายนฺติ. มนุสฺสา อุชฺฌายนฺติ ขิยฺยนฺติ วิปาเจนฺติ “ยเถว\csromansharedfootnote{97dd1823d991399d}{ยถา จ (ก)} มยํ คายาม, เอวเมวิเม สมณา สกฺยปุตฺติยา อายตเกน คีตสฺสเรน ธมฺมํ คายนฺตี”ติ. อสฺโสสุํ โข ภิกฺขู เตสํ มนุสฺสานํ อุชฺฌายนฺตานํ ขิยฺยนฺตานํ วิปาเจนฺตานํ. เย เต ภิกฺขู อปฺปิจฺฉา ฯเปฯ เต อุชฺฌายนฺติ ขิยฺยนฺติ วิปาเจนฺติ “กถํ หิ นาม ฉพฺพคฺคิยา ภิกฺขู อายตเกน คีตสฺสเรน ธมฺมํ คายิสฺสนฺตี”ติ. อถ โข เต ภิกฺขู ภควโต เอตมตฺถํ อาโรเจสุํ. สจฺจํ กิร ภิกฺขเว ฯเปฯ. สจฺจํ ภควาติ ฯเปฯ ธมฺมิํ กถํ กตฺวา ภิกฺขู อามนฺเตสิ “ปญฺจิเม ภิกฺขเว อาทีนวา อายตเกน คีตสฺสเรน ธมฺมํ คายนฺตสฺส. อตฺตนาปิ ตสฺมิํ สเร สารชฺชติ, ปเรปิ ตสฺมิํ สเร สารชฺชนฺติ, คหปติกาปิ อุชฺฌายนฺติ, สรกุตฺติมฺปิ นิกามย\-มานสฺส สมาธิสฺส ภงฺโค โหติ, ปจฺฉิมา ชนตา ทิฏฺฐานุคติํ อาปชฺชติ. อิเม โข ภิกฺขเว ปญฺจ อาทีนวา อายตเกน คีตสฺสเรน ธมฺมํ คายนฺตสฺส. น ภิกฺขเว อายตเกน คีตสฺสเรน ธมฺโม คายิตพฺโพ, โย คาเยยฺย, อาปตฺติ ทุกฺกฏสฺสา”ติ.}

\prose{เตน โข ปน สมเยน ภิกฺขู สรภญฺเญ กุกฺกุจฺจายนฺติ. ภควโต เอตมตฺถํ อาโรเจสุํ. อนุชานามิ ภิกฺขเว สรภญฺญนฺติ.}

\prose{เตน โข ปน สมเยน ฉพฺพคฺคิยา ภิกฺขู พาหิรโลมิํ\csromansharedfootnote{7e575948f14e0068}{พาหิยโลมิํ (ก)} อุณฺณิํ ธาเรนฺติ. มนุสฺสา อุชฺฌายนฺติ ขิยฺยนฺติ วิปาเจนฺติ “เสยฺยถาปิ คิหี \csromanfolio{244}กามโภคิโน”ติ. ภควโต เอตมตฺถํ อาโรเจสุํ. น ภิกฺขเว พาหิยโลมิ อุณฺณิ ธาเรตพฺพา, โย ธาเรยฺย, อาปตฺติ ทุกฺกฏสฺสาติ.}

\proseitem{250}{เตน โข ปน สมเยน รญฺโญ มาคธสฺส เสนิยสฺส พิมฺพิสารสฺส อาราเม อมฺพา ผลิโน โหนฺติ. รญฺโญ มาคเธน เสนิเยน พิมฺพิสาเรน อนุญฺญาตํ โหติ “ยถาสุขํ อยฺยา อมฺพํ ปริภุญฺชนฺตู”ติ. ฉพฺพคฺคิยา ภิกฺขู ตรุณญฺเญว อมฺพํ ปาตาเปตฺวา ปริภุญฺชิํสุ. รญฺโญ จ มาคธสฺส เสนิยสฺส พิมฺพิสารสฺส อมฺเพน อตฺโถ โหติ. อถ โข ราชา มาคโธ เสนิโย พิมฺพิสาโร มนุสฺเส อาณาเปสิ “คจฺฉถ ภเณ, อารามํ คนฺตฺวา อมฺพํ อาหรถา”ติ. “เอวํ เทวา”ติ โข เต มนุสฺสา รญฺโญ มาคธสฺส เสนิยสฺส พิมฺพิสารสฺส ปฏิสฺสุตฺวา อารามํ คนฺตฺวา อารามปาลํ เอตทโวจุํ “เทวสฺส ภเณ อมฺเพน อตฺโถ, อมฺพํ เทถา”ติ. นตฺถายฺยา อมฺพํ, ตรุณญฺเญว อมฺพํ ปาตาเปตฺวา ภิกฺขู ปริภุญฺชิํสูติ. อถ โข เต มนุสฺสา รญฺโญ มาคธสฺส เสนิยสฺส พิมฺพิสารสฺส เอตมตฺถํ อาโรเจสุํ. สุปริภุตฺตํ ภเณ อยฺเยหิ อมฺพํ, อปิ จ ภควตา มตฺตา วณฺณิตาติ. มนุสฺสา อุชฺฌายนฺติ ขิยฺยนฺติ วิปาเจนฺติ “กถํ หิ นาม สมณา สกฺยปุตฺติยา น มตฺตํ ชานิตฺวา รญฺโญ อมฺพํ ปริภุญฺชิสฺสนฺตี”ติ. อสฺโสสุํ โข ภิกฺขู เตสํ มนุสฺสานํ อุชฺฌายนฺตานํ ขิยฺยนฺตานํ วิปาเจนฺตานํ. อถ โข เต ภิกฺขู ภควโต เอตมตฺถํ อาโรเจสุํ. น ภิกฺขเว อมฺพํ ปริภุญฺชิ\-ตพฺพํ, โย ปริภุญฺเชยฺย, อาปตฺติ ทุกฺกฏสฺสาติ.}

\prose{เตน โข ปน สมเยน อญฺญตรสฺส ปูคสฺส สํฆภตฺตํ โหติ, สูเป อมฺพเปสิกาโย ปกฺขิตฺตา โหนฺติ, ภิกฺขู กุกฺกุจฺจายนฺตา นปฺปฏิคฺคณฺหนฺติ - ป-. ปฏิคฺคณฺหถ ภิกฺขเว ปริภุญฺชถ, อนุชานามิ ภิกฺขเว อมฺพเปสิกนฺติ.}

\prose{เตน โข ปน สมเยน อญฺญตรสฺส ปูคสฺส สํฆภตฺตํ โหติ, เต น ปริยาปุณิํสุ อมฺพเปสิกํ กาตุํ, ภตฺตคฺเค สกเลเหว อมฺเพหิ เทนฺติ, ภิกฺขู กุกฺกุจฺจายนฺตา นปฺปฏิคฺคณฺหนฺติ. ปฏิคฺคณฺหถ ภิกฺขเว ปริภุญฺชถ, อนุชานามิ ภิกฺขเว ปญฺจหิ สมณกปฺเปหิ ผลํ ปริภุญฺชิตุํ, \csromanfolio{245}อคฺคิปริจิตํ สตฺถ\-ปริจิตํ นขปริวิตํ อพีชํ นิพฺพตฺต\-พีชญฺเญว\csromansharedfootnote{7a9384f708842f4f}{นิพฺพฏฺฏพีชํ (สี, สฺยา)} ปญฺจมํ. อนาชานามิ ภิกฺขเว อิเมหิ ปญฺจหิ สมณกปฺเปหิ ผลํ ปริภุญฺชิ\-ตุนฺติ.}

\proseitem{251}{เตน โข ปน สมเยน อญฺญตโร ภิกฺขุ อหินา ทฏฺโฐ กาลงฺกโต โหติ. ภควโต เอตมตฺถํ อาโรเจสุํ. น หิ นูน โส ภิกฺขเว ภิกฺขุ อิมานิ จตฺตาริ อหิราชกุลานิ เมตฺเตน จิตฺเตน ผริ. สเจ หิ โส ภิกฺขเว ภิกฺขุ อิมานิ จตฺตาริ อหิราชกุลานิ เมตฺเตน จิตฺเตน ผเรยฺย, น หิ โส ภิกฺขเว ภิกฺขุ อหินา ทฏฺโฐ กาลํ กเรยฺย. กตมานิ จตฺตาริ อหิราชกุลานิ, วิรูปกฺขํ อหิราชกุลํ, เอราปถํ อหิราชกุลํ, ฉพฺยาปุตฺตํ อหิราชกุลํ, กณฺหา โคตมํ อหิราชกุลํ. น หิ นูน โส ภิกฺขเว ภิกฺขุ อิมานิ จตฺตาริ อหิราชกุลานิ เมตฺเตน จิตฺเตน ผริ. สเจ หิ โส ภิกฺขเว ภิกฺขุ อิมานิ จตฺตาริ อหิราชกุลานิ เมตฺเตน จิตฺเตน ผเรยฺย, น หิ โส ภิกฺขเว ภิกฺขุ อหินา ทฏฺโฐ กาลํ กเรยฺย. อนุชานามิ ภิกฺขเว อิมานิ จตฺตาริ อหิราชกุลานิ เมตฺเตน จิตฺเตน ผริตุํ อตฺตคุตฺติยา อตฺตรกฺขาย อตฺตปริตฺตํ กาตุํ. เอวญฺจ ปน ภิกฺขเว กาตพฺพํ\csromandash{}}

\prose{\csromansymbolfootnote{*}{ขุ 5. 53; อํ 1. 384 ปิฏฺเฐสุปิ อิทํ วตฺถุ อาคตํ. 2. สริํสปานิ (สี, สฺยา)}“วิรูปกฺเขหิ เม เมตฺตํ, เมตฺตํ เอราปเถหิ เม.}

\prose{ฉพฺยาปุตฺเตหิ เม เมตฺตํ, เมตฺตํ กณฺหา\-โคตม\-เกหิ จ.}

\prose{อปาทเกหิ เม เมตฺตํ, เมตฺตํ ทฺวิปาทเกหิ เม.}

\prose{จตุปฺปเทหิ เม เมตฺตํ, เมตฺตํ พหุปฺปเทหิ เม.}

\prose{มา มํ อปาทโก หิํสิ, มา มํ หิํสิ ทฺวิปาทโก.}

\prose{มา มํ จตุปฺปโท หิํสิ, มา มํ หิํสิ พหุปฺปโท.}

\prose{สพฺเพ สตฺตา สพฺเพ ปาณา, สพฺเพ ภูตา จ เกวลา.}

\prose{สพฺเพ ภทฺรานิ ปสฺสนฺตุ, มา กญฺจิ ปาปมาคมา.}

\prose{อปฺปมาโณ พุทฺโธ, อปฺปมาโณ ธมฺโม.}

\prose{อปฺปมาโณ สํโฆ, ปมาณวนฺตานิ สรีสปานิ2.}

\prose{อหิ วิจฺฉิกา สตปที อุณฺณนาภิ สรพู มูสิกา.}

\prose{กตา เม รกฺขา กตํ เม ปริตฺตํ, ปฏิกฺกมนฺตุ ภูตานิ.}

\prose{โสหํ นโม ภควโต, นโม สตฺตนฺนํ สมฺมาสมฺพุทฺธานนฺ”ติ.}

\prose{\csromanfolio{246}เตน โข ปน สมเยน อญฺญตโร ภิกฺขุ อนภิรติยา ปีฬิโต อตฺตโน องฺคชาตํ ฉินฺทิ. ภควโต เอตมตฺถํ อาโรเจสุํ. อญฺญมฺหิ โส ภิกฺขเว โมฆปุริโส เฉตพฺพมฺหิ อญฺญํ ฉินฺทิ, น ภิกฺขเว อตฺตโน องฺคชาตํ เฉตพฺพํ, โย ฉินฺเทยฺย, อาปตฺติ ถุลฺลจฺจยสฺสาติ.}

\proseitem{252}{เตน โข ปน สมเยน ราชคหกสฺส เสฏฺฐิสฺส มหคฺฆสฺส จนฺทนสฺส\csromansharedfootnote{f67a035c27284a24}{จนฺทนสารสฺส (สี, สฺยา)} จนฺทนคณฺฐิ อุปฺปนฺนา โหติ. อถ โข ราชคหกสฺส เสฏฺฐิสฺส เอตทโหสิ “ยํนูนาหํ อิมาย จนฺทนคณฺฐิยา ปตฺตํ เลขาเปยฺยํ, เลขญฺจ เม ปริโภคํ ภวิสฺสติ, ปตฺตญฺจ ทานํ ทสฺสามี”ติ. อถ โข ราชคหโก เสฏฺฐิ ตาย จนฺทนคณฺฐิยา ปตฺตํ เลขาเปตฺวา สิกฺกาย อุฑฺฑิตฺวา\csromansharedfootnote{bea76841c5ecf04d}{วาหิตฺวา (สี)} เวฬคฺเค อาลคฺเคตฺวา เวฬุ\-ปรมฺปราย พนฺธิตฺวา\csromansharedfootnote{34397b475f037ae6}{วาหิตฺวา (สฺยา)} เอวมาห “โย สมโณ วา พฺราหฺมโณ วา อรหา เจว อิทฺธิมา จ, ทินฺนํเยว ปตฺตํ, โอหรตู”ติ. อถ โข ปูรโณ กสฺสโป เยน ราชคหโก เสฏฺฐิ เตนุปสงฺกมิ, อุปสงฺกมิตฺวา ราชคหกํ เสฏฺฐิํ เอตทโวจ “อหํ หิ คหปติ อรหา เจว อิทฺธิมา จ, เทหิ เม ปตฺตนฺ”ติ. สเจ ภนฺเต อายสฺมา อรหา เจว อิทฺธิมา จ, ทินฺนํเยว ปตฺตํ, โอหรตูติ. อถ โข มกฺขลิ โคสาโล. อชิโต เกสกมฺพโล. ปกุโธ กจฺจายโน. สญฺจโย เพลฏฺฐปุตฺโต\csromansharedfootnote{d2a9d77abdc4f762}{สญฺชโย เพลฺลฏฺฐิปุตฺโต (สี)}. นิคณฺโฐ นาฏปุตฺโต\csromansharedfootnote{6f66c12c39288f07}{นาถปุตฺโต (สี)} เยน ราชคหโก เสฏฺฐิ เตนุปสงฺกมิ, อุปสงฺกมิตฺวา ราชคหกํ เสฏฺฐิํ เอตทโวจ “อหํ หิ คหปติ อรหา เจว อิทฺธิมา จ, เทหิ เม ปตฺตนฺ”ติ. สเจ ภนฺเต อายสฺมา อรหา เจว อิทฺธิมา จ, ทินฺนํเยว ปตฺตํ โอหรตูติ.}

\prose{เตน โข ปน สมเยน อายสฺมา จ มหาโมคฺคลฺลาโน อายสฺมา จ ปิณฺโฑล\-ภารทฺวาโช ปุพฺพณฺหสมยํ นิวาเสตฺวา ปตฺต\-จีวรมา\-ทาย ราชคหํ ปิณฺฑายา ปวิสิํสุ. อถ โข อายสฺมา ปิณฺโฑล\-ภารทฺวาโช อายสฺมนฺตํ มหา\-โมคฺคลฺลานํ เอตทโวจ “อายสฺมา โข มหาโมคฺคลฺลาโน อรหา เจว อิทฺธิมา จ, คจฺฉาวุโส โมคฺคลฺลาน, เอตํ ปตฺตํ โอหร, ตุเยฺหโส ปตฺโต”ติ. อายสฺมา โข ภารทฺวาโช อรหา เจว อิทฺธิมา จ, คจฺฉาวุโส ภารทฺวาช เอตํ ปตฺตํ โอหร, ตุเยฺหโส ปตฺโตติ. อถ โข อายสฺมา ปิณฺโฑล\-ภารทฺวาโช เวหาสํ อพฺภุคฺคนฺตฺวา ตํ ปตฺตํ คเหตฺวา ติกฺขตฺตุํ ราชคหํ อนุปริยายิ.}

\prose{\csromanfolio{247}เตน โข ปน สมเยน ราชคหโก เสฏฺฐิ สปุตฺตทาโร สเก นิเวสเน ฐิโต โหติ ปญฺชลิโก นมสฺสมาโน “อิเธว ภนฺเต อยฺโย ภารทฺวาโช อมฺหากํ นิเวสเน ปติฏฺฐาตู”ติ. อถ โข อายสฺมา ปิณฺโฑล\-ภารทฺวาโช ราชคหกสฺส เสฏฺฐิสฺส นิเวสเน ปติฏฺฐาสิ. อถ โข ราชคหโก เสฏฺฐิ อายสฺมโต ปิณฺโฑล\-ภารทฺวาชสฺส หตฺถโต ปตฺตํ คเหตฺวา มหคฺฆสฺส ขาทนียสฺส ปูเรตฺวา อายสฺมโต ปิณฺโฑล\-ภารทฺวาชสฺส อทาสิ. อถ โข อายสฺมา ปิณฺโฑล\-ภารทฺวาโช ตํ ปตฺตํ คเหตฺวา อารามํ อคมาสิ. อสฺโสสุํ โข มนุสฺสา “อยฺเยน กิร ปิณฺโฑล\-ภารทฺวาเชน ราชคหกสฺส เสฏฺฐิสฺส ปตฺโต โอหาริโต”ติ. เต จ มนุสฺสา อุจฺจาสทฺทา มหาสทฺทา อายสฺมนฺตํ ปิณฺโฑล\-ภารทฺวาชํ ปิฏฺฐิโต ปิฏฺฐิโต อนุพนฺธิํสุ.}

\prose{อสฺโสสิ โข ภควา อุจฺจาสทฺทํ มหาสทฺทํ, สุตฺวาน อายสฺมนฺตํ อานนฺทํ อามนฺเตสิ “กิํ นุ โข โส อานนฺท อุจฺจาสทฺโท มหาสทฺโท”ติ. อายสฺมตา ภนฺเต ปิณฺโฑล\-ภารทฺวาเชน ราชคหกสฺส เสฏฺฐิสฺส ปตฺโต โอหาริโต, อสฺโสสุํ โข ภนฺเต มนุสฺสา อยฺเยน กิร ปิณฺโฑล\-ภารทฺวาเชน ราชคหกสฺส เสฏฺฐิสฺส ปตฺโต โอหาริโตติ, เต จ ภนฺเต มนุสฺสา อุจฺจาสทฺทา มหาสทฺทา อายสฺมนฺตํ ปิณฺโฑล\-ภารทฺวาชํ ปิฏฺฐิโต ปิฏฺฐิโต อนุพนฺธา, โส เอโส ภนฺเต ภควา อุจฺจาสทฺโท มหาสทฺโทติ. อถ โข ภควา เอตสฺมิํ นิทาเน เอตสฺมิํ ปกรเณ ภิกฺขุสํฆํ สนฺนิปาตาเปตฺวา อายสฺมนฺตํ ปิณฺโฑล\-ภารทฺวาชํ ปฏิปุจฺฉิ “สจฺจํ กิร ตยา ภารทฺวาช ราชคหกสฺส เสฏฺฐิสฺส ปตฺโต โอหาริโต”ติ. สจฺจํ ภควาติ. วิครหิ พุทฺโธ ภควา อนนุจฺฉวิกํ ภารทฺวาช อนนุโลมิกํ อปฺปติรูปํ อสฺสามณกํ อกปฺปิยํ อกรณียํ, กถํ หิ นาม ตฺวํ ภารทฺวาช ฉวสฺส ทารุปกฺกสฺส การณา คิหีนํ อุตฺตริ\-มนุสฺสธมฺมํ อิทฺธิ\-ปาฏิหาริยํ ทสฺเสสฺสสิ, เสยฺยถาปิ ภารทฺวาช มาตุคาโม ฉวสฺส มาสกรูปสฺส การณา โกปินํ ทสฺเสติ, เอวเมว โข ตยา ภารทฺวาช ฉวสฺส ทารุปตฺตสฺส การณา คิหีนํ อุตฺตริ\-มนุสฺสธมฺมํ อิทฺธิ\-ปาฏิหาริยํ ทสฺสิตํ. เนตํ ภารทฺวาช อปฺปสนฺนานํ วา ปสาทาย - ป- วิครหิตฺวา ฯเปฯ ธมฺมิํ กถํ กตฺวา ภิกฺขู อามนฺเตสิ “น ภิกฺขเว คิหีนํ อุตฺตริ\-มนุสฺสธมฺมํ อิทฺธิ\-ปาฏิหาริยํ ทสฺเสตพฺพํ, โย ทสฺเสยฺย, อาปตฺติ ทุกฺกฏสฺส. ภินฺทเถตํ ภิกฺขเว ทารุปตฺตํ, สกลิกํ สกลิกํ \csromanfolio{248}กตฺวา ภิกฺขูนํ อญฺชนุ\-ปปิสนํ เทถ. น จ ภิกฺขเว ทารุปตฺโต ธาเรตพฺโพ, โย ธาเรยฺย, อาปตฺติ ทุกฺกฏสฺสา”ติ.}

\prose{เตน โข ปน สมเยน ฉพฺพคฺคิยา ภิกฺขู อุจฺจาวเจ ปตฺเต ธาเรนฺติ โสวณฺณมยํ รูปิยมยํ. มนุสฺสา อุชฺฌายนฺติ ขิยฺยนฺติ วิปาเจนฺติ “เสยฺยถาปิ คิหี กามโภคิโน”ติ. ภควโต เอตมตฺถํ อาโรเจสุํ. น ภิกฺขเว โสวณฺณมโย ปตฺโต ธาเรตพฺโพ ฯเปฯ. น รูปิยมโย ปตฺโต ธาเรตพฺโพ. น มณิมโย ปตฺโต ธาเรตพฺโพ. น เวฬุริยมโย ปตฺโต ธาเรตพฺโพ. น ผลิกมโย ปตฺโต ธาเรตพฺโพ. น กํสมโย ปตฺโต ธาเรตพฺโพ. น กาจมโย ปตฺโต ธาเรตพฺโพ. น ติปุมโย ปตฺโต ธาเรตพฺโพ. น สีสมโย ปตฺโต ธาเรตพฺโพ. น ตมฺพโลหมโย ปตฺโต ธาเรตพฺโพ, โย ธาเรยฺย, อาปตฺติ ทุกฺกฏสฺส. อนุชานามิ ภิกฺขเว เทฺว ปตฺเต อโยปตฺตํ มตฺติกา\-ปตฺตนฺติ.}

\proseitem{253}{เตน โข ปน สมเยน ปตฺตมูลํ ฆํสิยติ. ภควโต เอตมตฺถํ อาโรเจสุํ. อนุชานามิ ภิกฺขเว ปตฺต\-มณฺฑลนฺติ.}

\prose{เตน โข ปน สมเยน ฉพฺพคฺคิยา ภิกฺขู อุจฺจาวจานิ ปตฺตมณฺฑลานิ ธาเรนฺติ โสวณฺณมยํ รูปิยมยํ. มนุสฺสา อุชฺฌายนฺติ ขิยฺยนฺติ วิปาเจนฺติ “เสยฺยถาปิ คิหี กามโภคิโน”ติ, ภควโต เอตมตฺถํ อาโรเจสุํ. น ภิกฺขเว อุจฺจาวจานิ ปตฺตมณฺฑลานิ ธาเรตพฺพานิ, โย ธาเรยฺย, อาปตฺติ ทุกฺกฏสฺส. อนุชานามิ ภิกฺขเว เทฺว ปตฺตมณฺฑลานิ ติปุมยํ สีสมยนฺติ. พหลานิ มณฺฑลานิ น อจฺฉุปิยนฺติ. ภควโต เอตมตฺถํ อาโรเจสุํ. อนุชานามิ ภิกฺขเว ลิขิตุนฺติ. วลี\csromansharedfootnote{cba67ab7a0c116f1}{จลี (ก)} โหนฺติ. ภควโต เอตมตฺถํ อาโรเจสุํ. อนุชานามิ ภิกฺขเว มกรทนฺตกํ ฉินฺทิตุนฺติ.}

\prose{เตน โข ปน สมเยน ฉพฺพคฺคิยา ภิกฺขู จิตฺรานิ ปตฺตมณฺฑลานิ ธาเรนฺติ รูปกากิณฺณานิ ภิตฺติกมฺม\-กตานิ, ตานิ รถิกายปิ ทสฺเสนฺตา อาหิณฺฑนฺติ. มนุสฺสา อุชฺฌายนฺติ ขิยฺยนฺติ วิปาเจนฺติ “เสยฺยถาปิ คิหี กามโภคิโน”ติ. ภควโต เอตมตฺถํ อาโรเจสุํ. น ภิกฺขเว จิตฺรานิ ปตฺตมณฺฑลานิ ธาเรตพฺพานิ รูปกากิณฺณานิ ภิตฺติกมฺม\-กตานิ, โย ธาเรยฺย, อาปตฺติ ทุกฺกฏสฺส. อนุชานามิ ภิกฺขเว ปกติ\-มณฺฑลนฺติ.}

\proseitem{254}{\csromanfolio{249}เตน โข ปน สมเยน ภิกฺขู โสทกํ ปตฺตํ ปฏิสาเมนฺติ, ปตฺโต ทุสฺสติ. ภควโต เอตมตฺถํ อาโรเจสุํ. น ภิกฺขเว โสทโก ปตฺโต ปฏิสาเมตพฺโพ, โย ปฏิสาเมยฺย, อาปตฺติ ทุกฺกฏสฺส. อนุชานามิ ภิกฺขเว โอตาเปตฺวา ปตฺตํ ปฏิสาเมตุนฺติ.}

\prose{เตน โข ปน สมเยน ภิกฺขู โสทกํ\csromansharedfootnote{bfed45c2270cd5df}{สอุทกํ (ก)} ปตฺตํ โอตาเปนฺติ, ปตฺโต ทุคฺคนฺโธ โหติ. ภควโต เอตมตฺถํ อาโรเจสุํ. น ภิกฺขเว โสทโก ปตฺโต โอตาเปตพฺโพ, โย โอตาเปยฺย, อาปตฺติ ทุกฺกฏสฺส. อนุชานามิ ภิกฺขเว โวทกํ กตฺวา โอตาเปตฺวา ปตฺตํ ปฏิสาเมตุนฺติ.}

\prose{เตน โข ปน สมเยน ภิกฺขู อุเณฺห ปตฺตํ นิทหนฺติ, ปตฺตสฺส วณฺโณ ทุสฺสติ. ภควโต เอตมตฺถํ อาโรเจสุํ. น ภิกฺขเว อุเณฺห ปตฺโต นิทหิตพฺโพ, โย นิทเหยฺย, อาปตฺติ ทุกฺกฏสฺส. อนุชานามิ ภิกฺขเว มุหุตฺตํ อุเณฺห โอตาเปตฺวา ปตฺตํ ปฏิสาเมตุนฺติ.}

\prose{เตน โข ปน สมเยน สมฺพหุลา ปตฺตา อชฺโฌกาเส อนาธารา นิกฺขิตฺตา โหนฺติ, วาต\-มณฺฑลิกาย อาวฏฺเฏตฺวา ปตฺตา ภิชฺชิํสุ. ภควโต เอตมตฺถํ อาโรเจสุํ. อนุชานามิ ภิกฺขเว ปตฺตา\-ธารกนฺติ.}

\prose{เตน โข ปน สมเยน ภิกฺขู มิฑฺฒนฺเต ปตฺตํ นิกฺขิปนฺติ, ปริปติตฺวา\csromansharedfootnote{c9f72f2c087b0c00}{ปริวฏฺฏิตฺวา (สฺยา)} ปตฺโต ภิชฺชติ. ภควโต เอตมตฺถํ อาโรเจสุํ. น ภิกฺขเว มิฑฺฒนฺเต ปตฺโต นิกฺขิปิตพฺโพ, โย นิกฺขิเปยฺย, อาปตฺติ ทุกฺกฏสฺสาติ.}

\prose{เตน โข ปน สมเยน ภิกฺขู ปริภณฺฑนฺเต ปตฺตํ นิกฺขิปนฺติ, ปริปติตฺวา ปตฺโต ภิชฺชติ. ภควโต เอตมตฺถํ อาโรเจสุํ. น ภิกฺขเว ปริภณฺฑนฺเต ปตฺโต นิกฺขิปิตพฺโพ, โย นิกฺขิเปยฺย, อาปตฺติ ทุกฺกฏสฺสาติ.}

\prose{เตน โข ปน สมเยน ภิกฺขู ฉมาย ปตฺตํ นิกฺกุชฺชนฺติ, โอฏฺโฐ ฆํสิยติ. ภควโต เอตมตฺถํ อาโรเจสุํ. อนุชานามิ ภิกฺขเว ติณสนฺถารกนฺติ. ติณสนฺถารโก อุปจิกาหิ ขชฺชติ. ภควโต เอตมตฺตํ อาโรเจสุํ. อนุชานามิ ภิกฺขเว โจฬกนฺติ. โจฬกํ อุปจิกาหิ ขชฺชติ. ภควโต เอตมตฺถํ อาโรเจสุํ. อนุชานามิ ภิกฺขเว ปตฺตมาฬกนฺติ. ปตฺตมาฬโก ปริปติตฺวา ปตฺโต ภิชฺชติ. ภควโต \csromanfolio{250}เอตมตฺถํ อาโรเจสุํ. อนาชานามิ ภิกฺขเว ปตฺต\-กุณฺโฑลิกนฺติ. ปตฺต\-กุณฺโฑลิกาย ปตฺโต ฆํสิยติ. ภควโต เอตมตฺถํ อาโรเจสุํ. อนุชานามิ ภิกฺขเว ปตฺต\-ถวิกนฺติ. อํสพทฺธโก น โหติ. ภควโต เอตมตฺถํ อาโรเจสุํ. อนุชานามิ ภิกฺขเว อํสพทฺธกํ พนฺธน\-สุตฺตกนฺติ.}

\prose{เตน โข ปน สมเยน ภิกฺขู ภิตฺติขิเลปิ นาคทนฺตเกปิ ปตฺตํ ลคฺเคนฺติ, ปริปติตฺวา ปตฺโต ภิชฺชติ. ภควโต เอตมตฺถํ อาโรเจสุํ. น ภิกฺขเว ปตฺโต ลคฺเคตพฺโพ, โย ลคฺเคยฺย, อาปตฺติ ทุกฺกฏสฺสาติ.}

\prose{เตน โข ปน สมเยน ภิกฺขู มญฺจ ปตฺตํ นิกฺขิปนฺติ, สติสมฺโมสา นิสีทนฺตา โอตฺถริตฺวา ปตฺตํ ภินฺเทนฺติ. ภควโต เอตมตฺถํ อาโรเจสุํ. น ภิกฺขเว มญฺเจ ปตฺโต นิกฺขิปิตพฺโพ, โย นิกฺขิเปยฺย, อาปตฺติ ทุกฺกฏสฺสาติ.}

\prose{เตน โข ปน สมเยน ภิกฺขู ปีเฐ ปตฺตํ นิกฺขิปนฺติ, สติสมฺโมสา นิสีทนฺตา โอตฺถริตฺวา ปตฺตํ ภินฺเทนฺติ. ภควโต เอตมตฺถํ อาโรเจสุํ. น ภิกฺขเว ปีเฐ ปตฺโต นิกฺขิปิตพฺโพ, โย นิกฺขิเปยฺย, อาปตฺติ ทุกฺกฏสฺสาติ.}

\prose{เตน โข ปน สมเยน ภิกฺขู องฺเก ปตฺตํ นิกฺขิปนฺติ, สติสมฺโมสา อุฏฺฐหนฺติ, ปริปติตฺวา ปตฺโต ภิชฺชติ. ภควโต เอตมตฺถํ อาโรเจสุํ. น ภิกฺขเว องฺเก ปตฺโต นิกฺขิปิตพฺโพ, โย นิกฺขิเปยฺย, อาปตฺติ ทุกฺกฏสฺสาติ.}

\prose{เตน โข ปน สมเยน ภิกฺขู ฉตฺเต ปตฺตํ นิกฺขิปนฺติ, วาต\-มณฺฑลิกาย ฉตฺตํ อุกฺขิปิยติ, ปริปติตฺวา ปตฺโต ภิชฺชติ. ภควโต เอตมตฺถํ อาโรเจสุํ. น ภิกฺขเว ฉตฺเต ปตฺโต นิกฺขิปิตพฺโพ, โย นิกฺขิเปยฺย, อาปตฺติ ทุกฺกฏสฺสาติ.}

\proseitem{255}{เตน โข ปน สมเยน ภิกฺขู ปตฺตหตฺถา กวาฏํ ปณาเมนฺติ, กวาโฏ อาวฏฺฏิตฺวา ปตฺโต ภิชฺชติ. ภควโต เอตมตฺถํ อาโรเจสุํ. น ภิกฺขเว ปตฺตหตฺเถน กวาฏํ ปณาเมตพฺพํ\csromansharedfootnote{3bf3d0168887b41d}{กวาโฏ ปณาเมตพฺโพ (ก)}, โย ปณาเมยฺย, อาปตฺติ ทุกฺกฏสฺสาติ.}

\prose{\csromanfolio{251}เตน โข ปน สมเยน ภิกฺขู ตุมฺพกฏาเห ปิณฺฑาย จรนฺติ. มนุสฺสา อุชฺฌายนฺติ ขิยฺยนฺติ วิปาเจนฺติ “เสยฺยถาปิ ติตฺถิยา”ติ. ภควโต เอกมตฺถํ อาโรเจสุํ. น ภิกฺขเว ตุมฺพกฏาเห ปิณฺฑาย จริตพฺพํ, โย จเรยฺย, อาปตฺติ ทุกฺกฏสฺสาติ.}

\prose{เตน โข ปน สมเยน ภิกฺขู ฆฏิกฏาเห\csromansharedfootnote{c55dabf2004d700c}{ฆฏิกฏาเหน (สฺยา)} ปิณฺฑาย จรนฺติ. มนุสฺสา อุชฺฌายนฺติ ขิยฺยนฺติ วิปาเจนฺติ “เสยฺยถาปิ ติตฺถิยา”ติ. ภควโต เอตมตฺถํ อาโรเจสุํ. น ภิกฺขเว ฆฏิกฏาเห ปิณฺฑาย จริตพฺพํ, โย จเรยฺย, อาปตฺติ ทุกฺกฏสฺสาติ.}

\prose{เตน โข ปน สมเยน อญฺญตโร ภิกฺขุ สพฺพ\-ปํสุกูลิโก โหติ, โส ฉวสีสสฺส ปตฺตํ ธาเรติ, อญฺญตรา อิตฺถี ปสฺสิตฺวา ภีตา วิสฺสรมกาสิ “อภุํ เม ปิสาโจ วตายนฺ”ติ. มนุสฺสา อุชฺฌายนฺติ ขิยฺยนฺติ วิปาเจนฺติ “กยํ หิ นาม สมณา สกฺยปุตฺติยา ฉวสีสสฺส ปตฺตํ ธาเรสฺสนฺติ, เสยฺยถาปิ ปิสาจิลฺลิกา”ติ. ภควโต เอตมตฺถํ อาโรเจสุํ. น ภิกฺขเว ฉวสีสสฺส ปตฺโต ธาเรตพฺโพ, โย ธาเรยฺย, อาปตฺติ ทุกฺกฏสฺส. น จ ภิกฺขเว สพฺพ\-ปํสุกูลิเกน สวิตพฺพํ, โย ภเวยฺย, อาปตฺติ ทุกฺกฏสฺสาติ.}

\prose{เตน โข ปน สมเยน ภิกฺขู จลกานิปิ อฏฺฐิกานิปิ อุจฺฉิฏฺโฐ\-ทกมฺปิ ปตฺเตน นีหรนฺติ. มนุสฺสา อุชฺฌายนฺติ ขิยฺยนฺติ วิปาเจนฺติ “ยสฺมิํ เยวิเม สมณา สกฺยปุตฺติยา ภุญฺชนฺติ, โสว เนสํ ปฏิคฺคโห”ติ. ภควโต เอตมตฺถํ อาโรเจสุํ. น ภิกฺขเว จลกานิ วา อฏฺฐิกานิ วา อุจฺฉิฏฺโฐทกํ วา ปตฺเตน นีหริตพฺพํ, โย นีหเรยฺย, อาปตฺติ ทุกฺกฏสฺส. อนุชานามิ ภิกฺขเว ปฏิคฺคหนฺติ.}

\proseitem{256}{เตน โข ปน สมเยน ภิกฺขู หตฺเถน วิปฺผาเฬตฺวา จีวรํ สิพฺเพนฺติ, จีวรํ วิโลมิกํ โหติ. ภควโต เอตมตฺถํ อาโรเจสุํ. อนุชานามิ ภิกฺขเว สตฺถกํ นมตกนฺติ.}

\prose{เตน โข ปน สมเยน สํฆสฺส ทณฺฑสตฺถกํ อุปฺปนฺนํ โหติ. ภควโต เอตมตฺถํ อาโรเจสุํ. อนุชานามิ ภิกฺขเว ทณฺฑสตฺถกนฺติ.}

\prose{\csromanfolio{252}เตน โข ปน สมเยน ฉพฺพคฺคิยา ภิกฺขู อุจฺจาวเจ สตฺถกทณฺเฑ ธาเรนฺติ โสวณฺณมยํ รูปิยมยํ. มนุสฺสา อุชฺฌายนฺติ ขิยฺยนฺติ วิปาเจนฺติ “เสยฺยถาปิ คิหี กามโภคิโน”ติ. ภควโต เอตมตฺถํ อาโรเจสุํ. น ภิกฺขเว อุจฺจาวจา สตฺถกทณฺฑา ธาเรตพฺพา, โย ธาเรยฺย, อาปตฺติ ทุกฺกฏสฺส. อนุชานามิ ภิกฺขเว อฏฺฐิมยํ ทนฺตมยํ วิสาณมยํ นฬมยํ เวฬุมยํ กฏฺฐมยํ ชตุมยํ ผลมยํ โลหมยํ สงฺข\-นาภิ\-มยนฺติ.}

\prose{เตน โข ปน สมเยน ภิกฺขู กุกฺกุฏปตฺเตนปิ เวฬุเปสิกายปิ จีวรํ สิพฺเพนฺติ, จีวรํ ทุสฺสิพฺพิตํ โหติ. ภควโต เอตมตฺถํ อาโรเจสุํ. อนุชานามิ ภิกฺขเว สูจินฺติ. สูจิโย กณฺณกิตาโย โหนฺติ ฯเปฯ. อนุชานามิ ภิกฺขเว สูจินาฬิกนฺติ. สูจินาฬิกายปิ กณฺณกิตาโย โหนฺติ ฯเปฯ. อนุชานามิ ภิกฺขเว กิณฺเณน ปุเรตุนฺติ. กิณฺเณปิ กณฺณกิตาโย โหนฺติ ฯเปฯ. อนุชานามิ ภิกฺขเว สตฺตุยา ปุเรตุนฺติ. สตฺตุยาปิ กณฺณกิตาโย โหนฺติ ฯเปฯ. อนุชานามิ ภิกฺขเว สริตกนฺติ. สริตเกปิ กณฺณกิตาโย โหนฺติ ฯเปฯ. อนุชานามิ ภิกฺขเว มธุสิตฺถ\-เกน สาเรตุนฺติ. สริตกํ ปริภิชฺชติ ฯเปฯ. อนุชานามิ ภิกฺขเว สริตก\-สิปาฏิกนฺติ.}

\prose{เตน โข ปน สมเยน ภิกฺขู ตตฺถ ตตฺถ ขิลํ นิกฺขณิตฺวา สมฺพนฺธิตฺวา จีวรํ สิพฺเพนฺติ, จีวรํ วิกณฺณํ โหติ. ภควโต เอตมตฺถํ อาโรเจสุํ. อนุชานามิ ภิกฺขเว กถินํ กถินรชฺชุํ\csromansharedfootnote{1649318cd8548796}{กฐินํ กฐินรชฺชุํ (สี, สฺยา)}, ตตฺถ ตตฺถ โอพนฺธิตฺวา จีวรํ สิพฺเพตุนฺติ. วิสเม กถินํ ปตฺถรนฺติ, กถินํ ปริภิชฺชติ ฯเปฯ. น ภิกฺขเว วิสเม กถินํ ปตฺถริตพฺพํ, โย ปตฺถเรยฺย, อาปตฺติ ทุกฺกฏสฺสาติ.}

\prose{ฉมาย กถินํ ปตฺถรนฺติ, กถินํ ปํสุกิตํ โหติ ฯเปฯ. อนุชานามิ ภิกฺขเว ติณสนฺถารกนฺติ. กถินสฺส อนฺโต ชีรติ ฯเปฯ. อนุชานามิ ภิกฺขเว อนุวาตํ ปริภณฺฑํ อาโรเปตุนฺติ. กถินํ นปฺปโหติ ฯเปฯ. อนุชานามิ ภิกฺขเว ทณฺฑกถินํ พิทลกํ สลากํ วินนฺธน\-รชฺชุํ วินนฺธน\-สุตฺตํ วินนฺธิตฺวา จีวรํ สิพฺเพตุนฺติ. สุตฺตนฺตริกาโย วิสมา โหนฺติ ฯเปฯ. \csromanfolio{253}อนุชานามิ ภิกฺขเว กฬิมฺภกนฺติ. สุตฺตา วงฺกา โหนฺติ ฯเปฯ. อนุชานามิ ภิกฺขเว โมฆ\-สุตฺตกนฺติ.}

\prose{เตน โข ปน สมเยน ภิกฺขู อโธเตหิ ปาเทหิ กถินํ อกฺกมนฺติ กถินํ ทุสฺสติ. ภควโต เอตมตฺถํ อาโรเจสุํ. น ภิกฺขเว อโธเตหิ ปาเทหิ กถินํ อกฺกมิตพฺพํ, โย อกฺกเมยฺย อาปตฺติ ทุกฺกฏสฺสาติ.}

\prose{เตน โข ปน สมเยน ภิกฺขู อลฺเลหิ ปาเทหิ กถินํ อกฺกมนฺติ, กถินํ ทุสฺสติ. ภควโต เอตมตฺถํ อาโรเจสุํ. น ภิกฺขเว อลฺเลหิ ปาเทหิ กถินํ อกฺกมิตพฺพํ, โย อกฺกเมยฺย, อาปตฺติ ทุกฺกฏสฺสาติ.}

\prose{เตน โข ปน สมเยน ภิกฺขู สอุปาหนา กถินํ อกฺกมนฺติ, กถินํ ทุสฺสติ. ภควโต เอตมตฺถํ อาโรเจสุํ. น ภิกฺขเว สอุปาหเนน กถินํ อกฺกมิตพฺพํ, โย อกฺกเมยฺย, อาปตฺติ ทุกฺกฏสฺสาติ.}

\prose{เตน โข ปน สมเยน ภิกฺขู จีวรํ สิพฺพนฺตา องฺคุลิยา ปฏิคฺคณฺหนฺติ, องฺคุลิโย ทุกฺขา โหนฺติ. ภควโต เอตมตฺถํ อาโรเจสุํ. อนุชานามิ ภิกฺขเว ปฏิคฺคหนฺติ.}

\proseitem{257}{เตน โข ปน สมเยน ฉพฺพคฺคิยา ภิกฺขู อุจฺจาวเจ ปฏิคฺคเห ธาเรนฺติ โสวณฺณมยํ รูปิยมยํ. มนุสฺสา อุชฺฌายนฺติ ขิยฺยนฺติ วิปาเจนฺติ “เสยฺยถาปิ คิหี กามโภคิโน”ติ. ภควโต เอตมตฺถํ อาโรเจสุํ. น ภิกฺขเว อุจฺจาวจา ปฏิคฺคหา ธาเรตพฺพา, โย ธาเรยฺย, อาปตฺติ ทุกฺกฏสฺส. อนุชามานิ ภิกฺขเว อฏฺฐิมยํ ฯเปฯ สงฺข\-นาภิ\-มยนฺติ.}

\prose{เตน โข ปน สมเยน สูจิโยปิ สตฺถกาปิ ปฏิคฺคหาปิ นสฺสนฺติ. ภควโต เอตมตฺถํ อาโรเจสุํ. อนุชานามิ ภิกฺขเว อาเวสนวิตฺถกนฺติ. อาเวสนวิตฺถเก สเมกุลา โหนฺติ. ภควโต เอตมตฺถํ อาโรเจสุํ. อนุชานามิ ภิกฺขเว ปฏิคฺคหถ\-วิกนฺติ. อํสพทฺธโก น โหติ. ภควโต เอตมตฺถํ อาโรเจสุํ. อนุชานามิ ภิกฺขเว อํสพทฺธกํ พนฺธน\-สุตฺตกนฺติ.}

\prose{เตน โข ปน สมเยน ภิกฺขู อพฺโภกาเส จีวรํ สิพฺพนฺตา สีเตนปิ อุเณฺหนปิ กิลมนฺติ. ภควโต เอตมตฺถํ อาโรเจสุํ. อนุชานามิ ภิกฺขเว กถินสาลํ กถิน\-มณฺฑปนฺติ. กถินสาลา \csromanfolio{254}นีจวตฺถุกา โหติ, อุทเกน โอตฺถริยฺยติ. ภควโต เอตมตฺถํ อาโรเจสุํ. อนุชานามิ ภิกฺขเว อุจฺจวตฺถุกํ กาตุนฺติ. จโย ปริปตติ ฯเปฯ. อนุชานามิ ภิกฺขเว จินิตุํ ตโย จเย อิฏฺฐกาจยํ สิลาจยํ ทารุจยนฺติ. อาโรหนฺตา วิหญฺญนฺติ ฯเปฯ. อนุชานามิ ภิกฺขเว ตโย โสปาเน อิฏฺฐกาโสปานํ สิลาโสปานํ ทารุโสปานนฺติ. อาโรหนฺตา ปริปตนฺติ ฯเปฯ. อนุชานามิ ภิกฺขเว อาลมฺพนพาหนฺติ. กถินสาลาย ติณจุณฺณํ ปริปตติ ฯเปฯ. อนุชานามิ ภิกฺขเว โอคุมฺเผตฺวา\csromansharedfootnote{564b07207693bac1}{โอคุมฺเพตฺวา (สี, สฺยา)} อุลฺลิตฺตา\-วลิตฺตํ กาตุํ, เสตวณฺณํ กาฬวณฺณํ เครุก\-ปริกมฺมํ มาลากมฺมํ ลตากมฺมํ มกรทนฺตกํ ปญฺจปฏิกํ จีวรวํสํ จีวร\-รชฺชุกนฺติ.}

\prose{เตน โข ปน สมเยน ภิกฺขู จีวรํ สิพฺเพตฺวา ตตฺเถว กถินํ อุชฺฌิตฺวา ปกฺกมนฺติ, อุนฺทูเรหิปิ อุปจิกาหิปิ ขชฺชติ. ภควโต เอตมตฺถํ อาโรเจสุํ. อนุชานามิ ภิกฺขเว กถินํ สงฺฆริตุนฺติ. กถินํ ปริภิชฺชติ ฯเปฯ. อนุชานามิ ภิกฺขเว โคฆํสิกาย กถินํ สงฺฆริตุนฺติ. กถินํ วินิเวฐิยติ ฯเปฯ. อนุชานามิ ภิกฺขเว พนฺธน\-รชฺชุนฺติ.}

\prose{เตน โข ปน สมเยน ภิกฺขู กุฏฺเฏปิ ถมฺเภปิ กถินํ อุสฺสาเปตฺวา ปกฺกมนฺติ, ปริปติตฺวา กถินํ ภิชฺชติ. ภควโต เอตมตฺถํ อาโรเจสุํ. อนุชานามิ ภิกฺขเว ภิตฺติขิเล วา นาคทนฺเต วา ลคฺเคตุนฺติ.}

\proseitem{258}{อถ โข ภควา ราชคเห ยถาภิรนฺตํ วิหริตฺวา เยน เวสาลี เตน จาริกํ ปกฺกามิ. เตน โข ปน สมเยน ภิกฺขู สูจิกมฺปิ สตฺถกมฺปิ เภสชฺชมฺปิ ปตฺเตน อาทาย คจฺฉนฺติ. ภควโต เอตมตฺถํ อาโรเจสุํ. อนุชานามิ ภิกฺขเว เภสชฺช\-ตฺถวิกนฺติ. อํสพทฺธโก น โหติ ฯเปฯ. อนุชานามิ ภิกฺขเว อํสพทฺธกํ พนฺธน\-สุตฺตกนฺติ.}

\prose{เตน โข ปน สมเยน อญฺญตโร ภิกฺขุ อุปาหนาโย กายพนฺธเนน พนฺธิตฺวา คามํ ปิณฺฑาย ปาวิสิ, อญฺญตโร อุปาสโก ตํ ภิกฺขุํ อภิวาเทนฺโต อุปาหนาโย สีเสน ฆฏฺเฏติ, โส ภิกฺขุ มงฺกุ อโหสิ. อถ โข โส ภิกฺขุ อารามํ คนฺตฺวา ภิกฺขูนํ เอตมตฺถํ อาโรเจสิ. ภิกฺขู ภควโต เอตมตฺถํ อาโรเจสุํ. อนุชานามิ ภิกฺขเว อุปาหน\-ตฺถวิกนฺติ. อํสพทฺธโก น โหติ ฯเปฯ. อนุชานามิ ภิกฺขเว อํสพทฺธกํ พนฺธน\-สุตฺตกนฺติ.}

\prose{\csromanfolio{255}เตน โข ปน สมเยน อนฺตรามคฺเค อุทกํ อกปฺปิยํ โหติ, ปริสฺสาวนํ น โหติ. ภควโต เอตมตฺถํ อาโรเจสุํ. อนุชานามิ ภิกฺขเว ปริสฺสาวนนฺติ. โจฬกํ นปฺปโหติ ฯเปฯ. อนุชานามิ ภิกฺขเว กฏจฺฉุปริสฺสาวนนฺติ. โจฬกํ นปฺปโหติ ฯเปฯ. อนุชานามิ ภิกฺขเว ธมฺมกรณนฺติ\csromansharedfootnote{4c8b587411f2400f}{ธมฺมกรกํ (สี, สฺยา), ธมกรณํ (ก)}.}

\proseitem{259}{เตน โข ปน สมเยน เทฺว ภิกฺขู โกสเลสุ ชนปเท อทฺธาน\-มคฺค\-ปฺปฏิปนฺนา โหนฺติ, เอโก ภิกฺขุํ อนาจารํ อาจรติ, ทุติโย ภิกฺขุ ตํ ภิกฺขุํ เอตทโวจ “มา อาวุโส เอวรูปํ อกาสิ, เนตํ กปฺปตี”ติ. โส ตสฺมิํ อุปนนฺธิ. อถ โข โส ภิกฺขุ ปิปาสาย ปีฬิโต อุปนทฺธํ ภิกฺขุํ เอตทโวจ “เทหิ เม อาวุโส ปริสฺสาวนํ, ปานียํ ปิวิสฺสามี”ติ. อุปนทฺโธ ภิกฺขุ น อทาสิ, โส ภิกฺขุ ปิปาสาย ปีฬิโต กาลมกาสิ. อถ โข โส ภิกฺขุ อารามํ คนฺตฺวา ภิกฺขูนํ เอตมตฺถํ อาโรเจสิ. กิํ ปน ตฺวํ อาวุโส ปริสฺสาวนํ ยาจิยมาโน น อทาสีติ. เอวมาวุโสติ. เย เต ภิกฺขู อปฺปิจฺฉา ฯเปฯ เต อุชฺฌายนฺติ ขิยฺยนฺติ วิปาเจนฺติ “กถํ หิ นาม ภิกฺขุ ปริสฺสาวนํ ยาจิยมาโน น ทสฺสตี”ติ. อถ โข เต ภิกฺขู ภควโต เอตมตฺถํ อาโรเจสุํ. อถ โข ภควา เอตสฺมิํ นิทาเน เอตสฺมิํ ปกรเณ ภิกฺขุสํฆํ สนฺนิปาตาเปตฺวา ตํ ภิกฺขุํ ปฏิปุจฺฉิ “สจฺจํ กิร ตฺวํ ภิกฺขู ปริสฺสาวนํ ยาจิยมาโน น อทาสี”ติ. สจฺจํ ภควาติ. วิครหิ พุทฺโธ ภควา อนนุจฺฉวิกํ โมฆปุริส อนนุโลมิกํ อปฺปติรูปํ อสฺสามณกํ อกปฺปิยํ อกรณียํ, กถํ หิ นาม ตฺวํ โมฆปุริส ปริสฺสาวนํ ยาจิยมาโน น ทสฺสสิ. เนตํ โมฆปุริส อปฺปสนฺนานํ วา ปสาทาย ฯเปฯ วิครหิตฺวา ฯเปฯ ธมฺมิํ กถํ กตฺวา ภิกฺขู อามนฺเตสิ “น ภิกฺขเว อทฺธาน\-มคฺค\-ปฺปฏิปนฺเนน ภิกฺขุนา ปริสฺสาวนํ ยาจิยมาเนน น ทาตพฺพํ, โย น ทเทยฺย, อาปตฺติ ทุกฺกฏสฺส. น จ ภิกฺขเว อปริสฺสาวนเกน อทฺธาโน ปฏิปชฺชิตพฺโพ, โย ปฏิปชฺเชยฺย, อาปตฺติ ทุกฺกฏสฺส. สเจ น โหติ ปริสฺสาวนํ วา ธมฺมกรโณ วา, สํฆาฏิกณฺโณปิ อธิฏฺฐาตพฺโพ อิมินา ปริสฺสาเวตฺวา ปิวิสฺสามี”ติ. อถ โข ภควา อนุปุพฺเพน จาริกํ จรมาโน เยน เวสาลี ตทวสริ, ตตฺร สุทํ ภควา เวสาลิยํ วิหรติ มหาวเน \csromanfolio{256}กูฏาคาร\-สาลายํ. เตน โข ปน สมเยน ภิกฺขู นวกมฺมํ กโรนฺติ, ปริสฺสาวนํ น สมฺมติ. ภควโต เอตมตฺถํ อาโรเจสุํ. อนุชานามิ ภิกฺขเว ทณฺฑปริสฺสาวนนฺติ. ทณฺฑปริสฺสาวนํ น สมฺมติ. ภควโต เอตมตฺถํ อาโรเจสุํ. อนุชานามิ ภิกฺขเว โอตฺถรกนฺติ.}

\prose{เตน โข ปน สมเยน ภิกฺขู มกเสหิ อุพฺพาฬฺหา โหนฺติ. ภควโต เอตมตฺถํ อาโรเจสุํ. อนุชานามิ ภิกฺขเว มกส\-กุฏิกนฺติ.}

\proseitem{260}{เตน โข ปน สมเยน เวสาลิยํ ปณีตานํ ภตฺตานํ สตฺตปฏิปาฏิ อฏฺฐิตา โหติ. ภิกฺขู ปณีตานิ โภชนานิ ภุญฺชิตฺวา อภิสนฺนกายา โหนฺติ พหฺวาพาธา\csromansharedfootnote{562793cab0345a7f}{พวฺหาพาธา (สี)}. อถ โข ชีวโก โกมารภจฺโจ เวสาลิํ อคมาสิ เกนจิเทว กรณีเยน. อทฺทสา โข ชีวโก โกมารภจฺโจ ภิกฺขู อภิสนฺนกาเย พหฺวาพาเธ, ทิสฺวาน เยน ภควา เตนุปสงฺกมิ, อุปสงฺกมิตฺวา ภควนฺตํ อภิวาเทตฺวา เอกมนฺตํ นิสีทิ, เอกมนฺตํ นิสินฺโน โข ชีวโก โกมารภจฺโจ ภควนฺตํ เอตทโวจ “เอตรหิ ภนฺเต ภิกฺขู อภิสนฺนกายา พหฺวาพาธา, สาธุ ภนฺเต ภควา ภิกฺขูนํ จงฺกมญฺจ ชนฺตาฆรญฺจ อนุชานาตุ, เอวํ ภิกฺขู อปฺปาพาธา ภวิสฺสนฺตี”ติ. อถ โข ภควา ชีวกํ โกมารภจฺจํ ธมฺมิยา กถาย สนฺทสฺเสสิ สมาทเปสิ สมุตฺเตเชสิ สมฺปหํเสสิ. อถ โข ชีวโก โกมารภจฺโจ ภควตา ธมฺมิกา กถาย สนฺทสฺสิโต สมาทปิโต สมุตฺเตชิโต สมฺปหํสิโต อุฏฺฐายาสนา ภควนฺตํ อภิวาเทตฺวา ปทกฺขิณํ กตฺวา ปกฺกามิ. อถ โข ภควา เอตสฺมิํ นิทาเน เอตสฺมิํ ปกรเณ ธมฺมิํ กถํ กตฺวา ภิกฺขู อามนฺเตสิ “อนุชานามิ ภิกฺขเว จงฺกมญฺจ ชนฺตาฆรญฺจา”ติ.}

\prose{เตน โข ปน สมเยน ภิกฺขู วิสเม จงฺกเม จงฺกมนฺติ, ปาทา ทุกฺขา โหนฺติ. ภควโต เอตมตฺถํ อาโรเจสุํ. อนุชานามิ ภิกฺขเว สมํ กาตุนฺติ. จงฺกโม นีจวตฺถุโก โหติ, อุทเกน โอตฺถริยฺยติ ฯเปฯ. อนุชานามิ ภิกฺขเว อุจฺจวตฺถุตํ กาตุนฺติ. จโย ปริปตติ ฯเปฯ. อนุชานามิ ภิกฺขเว จินิตุํ ตโย จเย อิฏฺฐกาจยํ สิลาจยํ ทารุจยนฺติ. อาโรหนฺตา วิหญฺญนฺติ ฯเปฯ. อนุชานามิ ภิกฺขเว ตโย โสปาเน อิฏฺฐกาโสปานํ สิลาโสปานํ ทารุโสปานนฺติ. อาโร หนฺตา ปริปตนฺติ ฯเปฯ. อนุชานามิ ภิกฺขเว อาลมฺพนพาหนฺติ.}

\prose{\csromanfolio{257}เตน โข ปน สมเยน ภิกฺขู จงฺกเม จงฺกมนฺตา ปริปตนฺติ. ภควโต เอตมตฺถํ อาโรเจสุํ. อนุชานามิ ภิกฺขเว จงฺกมน\-เวทิกนฺติ.}

\prose{เตน โข ปน สมเยน ภิกฺขู อชฺโฌกาเส จงฺกมนฺตา สีเตนปิ อุเณฺหนปิ กิลมนฺติ. ภควโต เอตมตฺถํ อาโรเจสุํ. อนุชานามิ ภิกฺขเว จงฺกมนสาลนฺติ. จงฺกมนสาลายํ ติณจุณฺณํ ปริปตติ ฯเปฯ. อนุชานามิ ภิกฺขเว โอคุมฺเผตฺวา อุลฺลิตฺตา\-วลิตฺตํ กาตุํ เสตวณฺณํ กาฬวณฺณํ เครุก\-ปริกมฺมํ มาลากมฺมํ ลตากมฺมํ มกรทนฺตกํ ปญฺจปฏิกํ จีวรวํสํ จีวรรชฺชุนฺติ. ชนฺตาฆรํ นีจวตฺถุกํ โหติ, อุทเกน โอตฺถริยฺยติ ฯเปฯ. อนุชานามิ ภิกฺขเว อุจฺจวตฺถุกํ กาตุนฺติ. จโย ปริปตติ ฯเปฯ. อนุชานามิ ภิกฺขเว จินิตุํ ตโย จเย อิฏฺฐกาจยํ สิลาจยํ ทารุจยนฺติ. อาโรหนฺตา วิหญฺญนฺติ ฯเปฯ. อนุชานามิ ภิกฺขเว ตโย โสปาเน อิฏฺฐกาโสปานํ สิลาโสปานํ ทารุโสปานนฺติ. อาโรหนฺตา ปริปตนฺติ ฯเปฯ. อนุชานามิ ภิกฺขเว อาลมฺพนพาหนฺติ. ชนฺตาฆรสฺส กวาฏํ น โหติ ฯเปฯ. อนุชานามิ ภิกฺขเว กวาฏํ ปิฏฺฐสํฆาฏํ\csromansharedfootnote{8a2500f39eae48ed}{ปีฐสํฆาฏํ (ก)} อุทุกฺขลิกํ อุตฺตรปาสกํ อคฺคฬวฏฺฏิํ กปิสีสกํ สูจิกํ ฆฏิกํ ตาฬจฺฉิทฺทํ อาวิญฺฉน\-ฉิทฺทํ อาวิญฺฉน\-รชฺชุนฺติ. ชนฺตาฆรสฺส กุฏฺฏปาโท ชีรติ ฯเปฯ. อนุชานามิ ภิกฺขเว มณฺฑลิกํ กาตุนฺติ. ชนฺตาฆรสฺส ธูมเนตฺตํ น โหติ ฯเปฯ. อนุชานามิ ภิกฺขเว ธูมเนตฺตนฺติ.}

\prose{เตน โข ปน สมเยน ภิกฺขู ขุทฺทเก ชนฺตาฆเร มชฺเฌ อคฺคิฏฺฐานํ กโรนฺติ, อุปจาโร น โหติ. ภควโต เอตมตฺถํ อาโรเจสุํ. อนุชานามิ ภิกฺขเว ขุทฺทเก ชนฺตาฆเร เอกมนฺตํ อคฺคิฏฺฐานํ กาตุํ มหลฺลเก มชฺเฌติ. ชนฺตาฆเร อคฺคิ มุขํ ฑหติ ฯเปฯ. อนุชานามิ ภิกฺขเว มุขมตฺติกนฺติ. หตฺเถ มตฺติกํ เตเมนฺติ ฯเปฯ. อนุชานามิ ภิกฺขเว มตฺติกา\-โทณิกนฺติ. มตฺติกา ทุคฺคนฺธา โหติ ฯเปฯ. อนุชานามิ ภิกฺขเว วาเสตุนฺติ. ชนฺตาฆเร อคฺคิ กายํ ฑหติ ฯเปฯ. อนุชานามิ ภิกฺขเว อุทกํ อติหริตุนฺติ. ปาติยาปิ ปตฺเตนปิ อุทกํ อติหรนฺติ ฯเปฯ. อนุชานามิ ภิกฺขเว อุทกฏฺฐานํ อุทก\-สราวกนฺติ. ชนฺตาฆรํ ติณจฺฉทนํ น เสเทติ\csromansharedfootnote{8ae960d6ef2aef88}{ติณจฺฉาทเนน ฉาเทติ (ก)} ฯเปฯ. อนุชานามิ ภิกฺขเว โอคุมฺเผตฺวา อุลฺลิตฺตา\-วลิตฺตํ กาตุนฺติ. \csromanfolio{258}ชนฺตาฆรํ จิกฺขลฺลํ โหติ ฯเปฯ. อนุชานามิ ภิกฺขเว สนฺถริตุํ ตโย สนฺถเร อิฏฺฐกา\-สนฺถรํ สิลาสนฺถรํ ทารุ\-สนฺถรนฺติ. จิกฺขลฺลํเยว โหติ ฯเปฯ. อนุชานามิ ภิกฺขเว โธวิตุนฺติ. อุทกํ สนฺติฏฺฐติ ฯเปฯ. อนุชานามิ ภิกฺขเว อุทกนิทฺธมนนฺติ.}

\prose{เตน โข ปน สมเยน ภิกฺขู ชนฺตาฆเร ฉมาย นิสีทนฺติ, คตฺตานิ กณฺฑูวนฺติ. ภควโต เอตมตฺถํ อาโรเจสุํ. อนุชานามิ ภิกฺขเว ชนฺตาฆร\-ปีฐนฺติ. เตน โข ปน สมเยน ชนฺตาฆรํ อปริกฺขิตฺตํ โหติ ฯเปฯ. อนุชานามิ ภิกฺขเว ปริกฺขิปิตุํ ตโย ปากาเร อิฏฺฐกาปาการํ สิลาปาการํ ทารุปาการนฺติ. โกฏฺฐโก น โหติ ฯเปฯ. อนุชานามิ ภิกฺขเว โกฏฺฐกนฺติ. โกฏฺฐโก นีจวตฺถุโก โหติ, อุทเกน โอตฺถริยฺยติ ฯเปฯ. อนุชานามิ ภิกฺขเว อุจฺจวตฺถุกํ กาตุนฺติ. จโย ปริปตติ ฯเปฯ. อนุชานามิ ภิกฺขเว จินิตุํ ตโย จเย อิฏฺฐกาจยํ สิลาจยํ ทารุจยนฺติ. อาโรหนฺตา วิหญฺญนฺติ ฯเปฯ. อนุชานามิ ภิกฺขเว ตโย โสปาเน อิฏฺฐกาโสปานํ สิลาโสปานํ ทารุโสปานนฺติ. อาโรหนฺตา ปริปตนฺติ ฯเปฯ. อนุชานามิ ภิกฺขเว อาลมฺพนพาหนฺติ. โกฏฺฐกสฺส กวาฏํ น โหติ ฯเปฯ. อนุชานามิ ภิกฺขเว กวาฏํ ปิฏฺฐสํฆาฏํ อุทุกฺขลิกํ อุตฺตรปาสกํ อคฺคฬวฏฺฏิํ กปิสีสกํ สูจิกํ ฆฏิกํ ตาฬจฺฉิทฺทํ อาวิญฺฉนฉทฺทํ อาวิญฺฉน\-รชฺชุนฺติ. โกฏฺฐเก ติณจุณฺณํ ปริปตติ ฯเปฯ. อนุชานามิ ภิกฺขเว โอคุมฺเผตฺวา อุลฺลิตฺตา\-วลิตฺตํ กาตุํ กาฬวณฺณํ เครุก\-ปริกมฺมํ มาลากมฺมํ ลเตกมฺมํ มกรทนฺตกํ ปญฺจปฏิกนฺติ. ปริเวณํ จิกฺขลฺลํ โหติ ฯเปฯ. อนุชานามิ ภิกฺขเว มรุมฺพํ อุปกิริตุนฺติ. น ปริยาปุณนฺติ ฯเปฯ. อนุชานามิ ภิกฺขเว ปทรสิลํ นิกฺขิปิตุนฺติ. อุทกํ สนฺติฏฺฐติ ฯเปฯ. อนุชานามิ ภิกฺขเว อุทกนิทฺธมนนฺติ.}

\proseitem{261}{เตน โข ปน สมเยน ภิกฺขู นคฺคา นคฺคํ อภิวาเทนฺติ ฯเปฯ. นคฺคา นคฺคํ อภิวาทาเปนฺติ. นคฺคา นคฺคสฺส ปริกมฺมํ กโรนฺติ. นคฺคา นคฺคสฺส ปริกมฺมํ การาเปนฺติ. นคฺคา นคฺคสฺส เทนฺติ. นคฺคา ปฏิคฺคณฺหนฺติ. นคฺคา ขาทนฺติ. นคฺคา ภุญฺชนฺติ. นคฺคา สายนฺติ. นคฺคา ปิวนฺติ. ภควโต เอตมตฺตํ อาโรเจสุํ. น ภิกฺขเว นคฺเคน\csromansharedfootnote{7b0a0adb284b89f9}{อิทํ ปทํ กณฺฑจิ นตฺถิ.} นคฺโค อภิวาเทตพฺโพ ฯเปฯ. น นคฺเคน อภิวาเทตพฺพํ. น นคฺเคน\csromansharedfootnote{7b0a0adb284b89f9}{อิทํ ปทํ กณฺฑจิ นตฺถิ.} นคฺโค อภิวาทาเป\-ตพฺโพ. น นคฺเคน อภิวาทาเป\-ตพฺพํ. น นคฺเคน นคฺคสฺส \csromanfolio{259}ปริกมฺมํ กาตพฺพํ. น นคฺเคน นคฺคสฺส ปริกมฺมํ การาเปตพฺพํ. น นคฺเคน นคฺคสฺส ทาตพฺพํ. น นคฺเคน ปฏิคฺคเหตพฺพํ. น นคฺเคน ขาทิตพฺพํ. น นคฺเคน ภุญฺชิตพฺพํ. น นคฺเคน สายิตพฺพํ. น นคฺเคน ปาตพฺพํ, โย ปิเวยฺย, อาปตฺติ ทุกฺกฏสฺสาติ.}

\prose{เตน โข ปน สมเยน ภิกฺขู ชนฺตาฆเร ฉมาย จีวรํ นิกฺขิปนฺติ, จีวรํ ปํสุกิตํ โหติ. ภควโต เอตมตฺถํ อาโรเจสุํ. อนุชานามิ ภิกฺขเว จีวรวํสํ จีวรรชฺชุนฺติ. เทเว วสฺสนฺเต จีวรํ โอวสฺสติ ฯเปฯ. อนุชานามิ ภิกฺขเว ชนฺตา\-ฆร\-สาลนฺติ. ชนฺตาฆรสาลา นีจวตฺถุกา โหติ, อุทเกน โอตฺถริยฺยติ ฯเปฯ. อนุชานามิ ภิกฺขเว อุจฺจวตฺถุกํ กาตุนฺติ. จโย ปริปตติ ฯเปฯ. อนุชานามิ ภิกฺขเว จินิตุํ ฯเปฯ อาโรหนฺตา วิหญฺญนฺติ ฯเปฯ อาโรหนฺตา ปริปตนฺติ ฯเปฯ. อนุชานามิ ภิกฺขเว อาลมฺพนพาหนฺติ. ชนฺตา\-ฆร\-สาลาย ติณจุณฺณํ ปริปตติ ฯเปฯ. อนุชานามิ ภิกฺขเว โอคุมฺเผตฺวา อุลฺลิตฺตา\-วลิตฺตํ กาตุํ ฯเปฯ จีวรวํสํ จีวรรชฺชุนฺติ.}

\prose{เตน โข ปน สมเยน ภิกฺขู ชนฺตาฆเรปิ อุทเกปิ ปริกมฺมํ กาตุํ กุกฺกุจฺจายนฺติ. ภควโต เอตมตฺถํ อาโรเจสุํ. อนุชานามิ ภิกฺขเว ติสฺโส ปฏิจฺฉาทิโย ชนฺตาฆร\-ปฏิจฺฉาทิํ อุทก\-ปฏิจฺฉาทิํ วตฺถ\-ปฏิจฺฉาทินฺติ.}

\prose{เตน โข ปน สมเยน ชนฺตาฆเร อุทกํ น โหติ. ภควโต เอตมตฺถํ อาโรเจสุํ. อนุชานามิ ภิกฺขเว อุทปานนฺติ. อุทปานสฺส กูลํ ลุชฺชติ ฯเปฯ. อนุชานามิ ภิกฺขเว จินิตุํ ตโย จเย อิฏฺฐกาจยํ สิลาจยํ ทารุจยนฺติ. อุทปาโน นีจวตฺถุโก โหติ, อุทเกน โอตฺถริยฺยติ ฯเปฯ. อนุชานามิ ภิกฺขเว อุจฺจวตฺถุกํ กาตุนฺติ. จโย ปริปตติ ฯเปฯ อาโรหนฺตา วิหญฺญนฺติ ฯเปฯ อาโรหนฺตา ปริปตนฺติ ฯเปฯ. อนุชานามิ ภิกฺขเว อาลมฺพนพาหนฺติ.}

\proseitem{262}{เตน โข ปน สมเยน ภิกฺขู วลฺลิกายปิ กายพนฺธเนนปิ อุทกํ วาเหนฺติ\csromansharedfootnote{970a45ecc8a10d7e}{วาหนฺติ (ก)}. ภควโต เอตมตฺถํ อาโรเจสุํ. อนุชานามิ ภิกฺขเว อุทกวาหนรชฺชุนฺติ. หตฺถา ทุกฺขา โหนฺติ ฯเปฯ. อนุชานามิ ภิกฺขเว ตุลํ กรกฏกํ จกฺกวฏฺฏกนฺติ. ภาชนา พหู ภิชฺชนฺติ ฯเปฯ. อนุชานามิ ภิกฺขเว ตโย วารเก โลหวารกํ ทารุวารกํ จมฺม\-กฺขณฺฑนฺติ.}

\prose{\csromanfolio{260}เตน โข ปน สมเยน ภิกฺขู อชฺโฌกาเส อุทกํ วาเหนฺตา สีเตนปิ อุเณฺหนปิ กิลมนฺติ. ภควโต เอตมตฺถํ อาโรเจสุํ. อนุชานามิ ภิกฺขเว อุทปานสาลนฺติ. อุทปานสาลาย ติณจุณฺณํ ปริปตติ ฯเปฯ. อนุชานามิ ภิกฺขเว โอคุมฺเผตฺวา อุลฺลิตฺตา\-วลิตฺตํ กาตุํ เสตวณฺณํ กาฬวณฺณํ เครุก\-ปริกมฺมํ มาลากมฺมํ ลตากมฺมํ มกรทนฺตกํ ปญฺจปฏิกํ จีวรวํสํ จีวรรชฺชุนฺติ. อุทปาโน อปารุโต โหติ, ติณจุณฺเณหิปิ ปํสุเกหิปิ โอกิริยฺยติ ฯเปฯ. อนุชานามิ ภิกฺขเว อปิธานนฺติ. อุทกภาชนํ น สํวิชฺชติ ฯเปฯ. อนุชานามิ ภิกฺขเว อุทกโทณิํ อุทกกฏาหนฺติ.}

\prose{เตน โข ปน สมเยน ภิกฺขู อาราเม ตหํ ตหํ นหายนฺติ, อาราโม จิกฺขลฺโล โหติ. ภควโต เอตมตฺถํ อาโรเจสุํ. อนุชานามิ ภิกฺขเว จนฺทนิกนฺติ. จนฺทนิกา ปากฏา โหติ. ภิกฺขู หิริยนฺติ นหายิตุํ ฯเปฯ. อนุชานามิ ภิกฺขเว ปริกฺขิปิตุํ ตโย ปากาเร อิฏฺฐกาปาการํ สิลาปาการํ ทารุปาการนฺติ. จนฺทนิกา จิกฺขลฺลา โหติ ฯเปฯ. อนุชานามิ ภิกฺขเว สนฺตริตุํ ตโย สนฺถเร อิฏฺฐกา\-สนฺถรํ สิลาสนฺถรํ ทารุ\-สนฺถรนฺติ. อุทกํ สนฺติฏฺฐติ ฯเปฯ. อนุชานามิ ภิกฺขเว อุกกนิทฺธมนนฺติ.}

\prose{เตน โข ปน สมเยน ภิกฺขูนํ คตฺตานิ สีติคตานิ โหนฺติ, ภควโต เอตมตฺถํ อาโรเจสุํ ฯเปฯ. อนุชานามิ ภิกฺขเว อุทกปุญฺฉนิํ โจฬเกนปิ ปจฺจุทฺธริตุนฺติ.}

\proseitem{263}{เตน โข ปน สมเยน อญฺญตโร อุปาสโก สํฆสฺส อตฺถาย โปกฺขรณิํ กาเรตุกาโม โหติ ภควโต เอกมตฺถํ อาโรเจสุํ ฯเปฯ. อนุชานามิ ภิกฺขเว โปกฺขรณินฺติ. โปกฺขรณิยา กูลํ ลุชฺชติ ฯเปฯ. อนุชานามิ ภิกฺขเว จินิตุํ ตโย จเย อิฏฺฐกาจยํ สิลาจยํ ทารุจยนฺติ. อาโรหนฺตา วิหญฺญนฺติ ฯเปฯ. อนุชานามิ ภิกฺขเว ตโย โสปาเน อิฏฺฐกาโสปานํ สิลาโสปานํ ทารุโสปานนฺติ. อาโรหนฺตา ปริปตนฺติ ฯเปฯ. อนุชานามิ ภิกฺขเว อาลมฺพนพาหนฺติ. โปกฺขรณิยา อุทกํ ปุราณํ โหติ ฯเปฯ. อนุชานามิ ภิกฺขเว อุทกมาติกํ อุทกนิทฺธมนนฺติ.}

\prose{\csromanfolio{261}เตน โข ปน สมเยน อญฺญตโร ภิกฺขุ สํฆสฺส อตฺถาย นิลฺเลขํ ชนฺตาฆรํ กตฺตุกาโม โหติ. ภควโต เอตมตฺถํ อาโรเจสุํ ฯเปฯ. อนุชานามิ ภิกฺขเว นิลฺเลขํ ชนฺตาฆรนฺติ.}

\prose{เตน โข ปน สมเยน ฉพฺพคฺคิยา ภิกฺขู จาตุมาสํ นิสีทเนน วิปฺปวสนฺติ. ภควโต เอตมตฺถํ อาโรเจสุํ ฯเปฯ. น ภิกฺขเว จาตุมาสํ นิสีทเนน วิปฺปวสิตพฺพํ. โย วิปฺปวเสยฺย, อาปตฺติ ทุกฺกฏสฺสาติ.}

\proseitem{264}{เตน โข ปน สมเยน ฉพฺพคฺคิยา ภิกฺขู ปุปฺผา\-ภิกิณฺเณสุ สยเนสุ สยนฺติ, มนุสฺสา วิหารจาริกํ อาหิณฺฑนฺตา ปสฺสิตฺวา อุชฺฌายนฺติ ขิยฺยนฺติ วิปาเจนฺติ “เสยฺยถาปิ คิหี กามโภคิโน”ติ. ภควโต เอตมตฺถํ อาโรเจสุํ. น ภิกฺขเว ปุปฺผา\-ภิกิณฺเณสุ สยเนสุ สยิตพฺพํ, โย สเยยฺย, อาปตฺติ ทุกฺกฏสฺสาติ.}

\prose{เตน โข ปน สมเยน มนุสฺสา คนฺธมฺปิ มาลมฺปิ อาทาย อารามํ อาคจฺฉนฺติ, ภิกฺขู กุกฺกุจฺจายนฺตา น ปฏิคฺคณฺหนฺติ. ภควโต เอตมตฺถํ อาโรเจสุํ. อนุชานามิ ภิกฺขเว คนฺธํ คเหตฺวา กวาเฏ ปญฺจงฺคุลิกํ ทาตุํ, ปุปฺผํ คเหตฺวา วิหาเร เอกมนฺตํ นิกฺขิปิตุนฺติ.}

\prose{เตน โข ปน สมเยน สํฆสฺส นมตกํ อุปฺปนฺนํ โหติ. ภควโต เอตมตฺถํ อาโรเจสุํ. อนุชานามิ ภิกฺขเว นมตกนฺติ. อถ โข ภิกฺขูนํ เอตทโหสิ “นมตกํ อธิฏฺฐาตพฺพํ นุ โข อุทาหุ วิกปฺเปตพฺพนฺ”ติ ฯเปฯ. น ภิกฺขเว นมตกํ อธิฏฺฐาตพฺพํ, น วิกปฺเป\-ตพฺพนฺติ.}

\prose{เตน โข ปน สมเยน ฉพฺพคฺคิยา ภิกฺขู อาสิตฺตกู\-ปธาเน ภุญฺชนฺติ. มนุสฺสา อุชฺฌายนฺติ ขิยฺยนฺติ วิปาเจนฺติ “เสยฺยถาปิ คิหี กามโภคิโน”ติ. ภควโต เอตมตฺถํ อาโรเจสุํ ฯเปฯ. น ภิกฺขเว อาสิตฺตกู\-ปธาเน ภุญฺชิตพฺพํ, โย ภุญฺเชยฺย, อาปตฺติ ทุกฺกฏสฺสาติ.}

\prose{เตน โข ปน สมเยน อญฺญตโร ภิกฺขุ คิลาโน โหติ, โส ภุญฺชมาโน น สกฺโกติ หตฺเถน ปตฺตํ สนฺธาเรตุํ. ภควโต เอตมตฺถํ อาโรเจสุํ. อนุชานามิ ภิกฺขเว มโฬริกนฺติ.}

\prose{เตน โข ปน สมเยน ฉพฺพคฺคิยา ภิกฺขู เอกภาชเนปิ ภุญฺชนฺติ ฯเปฯ เอกถาลเกปิ ปิวนฺติ. เอกมญฺเจปิ ตุวฏฺเฏนฺติ. เอกตฺถรณาปิ ตุวฏฺเฏนฺติ. เอกปาวุรณาปิ ตุวฏฺเฏนฺติ. เอก\-ตฺถรณ\-ปาวุรณาปิ ตุวฏฺเฏนฺติ. มนุสฺสา \csromanfolio{262}อุชฺฌายนฺติ ขิยฺยนฺติ วิปาเจนฺติ “เสยฺยถาปิ คิหี กามโภคิโน”ติ. ภควโต เอตมตฺถํ อาโรเจสุํ. น ภิกฺขเว เอกภาชเน ภุญฺชิตพฺพํ ฯเปฯ น เอกถาลเก ปาตพฺพํ. น เอกมญฺเจ ตุวฏฺฏิตพฺพํ. น เอกตฺถรณา\csromansharedfootnote{9608a3fc99a67c36}{น เอกตฺถรเณ.} ตุวฏฺฏิตพฺพํ. น เอกปาวุรณา\csromansharedfootnote{eed9051104ccbe10}{น เอกปาวุรเณ.} ตุวฏฺฏิตพฺพํ. น เอก\-ตฺถรณ\-ปาวุรณา\csromansharedfootnote{637ab091b1272b29}{น เอกตฺถรณปาวุรเณ (สฺยา)} ตุวฏฺฏิตพฺพํ, โย ตุวฏฺเฏยฺย, อาปตฺติ ทุกฺกฏสฺสาติ.}

\proseitem{265}{เตน โข ปน สมเยน วฑฺโฒ ลิจฺฉวี เมตฺติย\-ภูมชกานํ ภิกฺขูนํ สหาโย โหติ. อถ โข วฑฺโฒ ลิจฺฉวี เยน เมตฺติย\-ภูมชกา ภิกฺขู เตนุปสงฺกมิ, อุปสงฺกมิตฺวา เมตฺติย\-ภูมชเก ภิกฺขู เอตทโวจ “วนฺทามิ อยฺยา”ติ. เอวํ วุตฺเต เมตฺติย\-ภูมชกา ภิกฺขู นาลปิํสุ. ทุติยมฺปิ โข วฑฺโฒ ลิจฺฉวี เมตฺติย\-ภูมชเก ภิกฺขู เอตทโวจ “วนฺทามิ อยฺยา”ติ. ทุติยมฺปิ โข เมตฺติย\-ภูมชกา ภิกฺขู นาลปิํสุ. ตติยมฺปิ โข วฑฺโฒ ลิจฺฉวี เมตฺติย\-ภูมชเก ภิกฺขู เอตทโวจ “วนฺทามิ อยฺยา”ติ. ตติยมฺปิ โข เมตฺติย\-ภูมชกา ภิกฺขู นาลปิํสุ. กฺยาหํ อยฺยานํ อปรชฺฌามิ, กิสฺส มํ อยฺยา นาลปนฺตีติ. ตถา หิ ปน ตฺวํ อาวุโส วฑฺฒ อเมฺห ทพฺเพน มลฺลปุตฺเตน วิเหฐิยมาเน อชฺฌุเปกฺขสีติ. กฺยาหํ อยฺยา กโรมีติ. สเจ โข ตฺวํ อาวุโส วฑฺฒ อิจฺเฉยฺยาสิ, อชฺเชว ภควา อายสฺมนฺตํ ทพฺพํ มลฺลปุตฺตํ นาสาเปยฺยาติ. กฺยาหํ อยฺยา กโรมิ, กิํ มยา สกฺกา กาตุนฺติ. เอหิ ตฺวํ อาวุโส วฑฺฒ เยน ภควา เตนุปสงฺกม, อุปสงฺกมิตฺวา ภควนฺตํ เอวํ วเทหิ “อิทํ ภนฺเต นจฺฉนฺนํ นปฺปติรูปํ, ยายํ ภนฺเต ทิสา อภยา อนีติกา อนุปทฺทวา, สายํ ทิสา สภยา สอีติกา สอุปทฺทวา, ยโต นิวาตํ, ตโต สวาตํ, อุทกํ มญฺเญ อาทิตฺตํ, อยฺเยน เม ทพฺเพน มลฺลปุตฺเตน ปชาปติ ทูสิตา”ติ. “เอวํ อยฺยา”ติ โข วฑฺโฒ ลิจฺฉวี เมตฺติย\-ภูมชกานํ ภิกฺขูนํ ปฏิสฺสุตฺวา เยน ภควา เตนุปสงฺกมิ, อุปสงฺกมิตฺวา ภควนฺตํ อภิวาเทตฺวา เอกมนฺตํ นิสีทิ, เอกมนฺตํ นิสินฺโน โข วฑฺโฒ ลิจฺฉวี ภควนฺตํ เอตทโวจ “อิทํ ภนฺเต นจฺฉนฺนํ นปฺปติรูปํ, ยายํ ภนฺเต ทิสา อภยา อนีติกา อนุปทฺทวา, สายํ ทิสา สภยา สอีติกา สอุปทฺทวา, ยโต นิวาตํ, ตโต สวาตํ, อุทกํ มญฺเญ อาทิตฺตํ, อยฺเยน เม ทพฺเพน มลฺลปุตฺเตน ปชาปติ ทูสิตา”ติ.}

\prose{\csromanfolio{263}อถ โข ภควา เอตสฺมิํ นิทาเน เอตสฺมิํ ปกรเณ ภิกฺขุสํฆํ สนฺนิปาตาเปตฺวา อายสฺมนฺตํ ทพฺพํ มลฺลปุตฺตํ ปฏิปุจฺฉิ “สรสิ ตฺวํ ทพฺพ เอวรูปํ กตฺตา ยถายํ วฑฺโฒ อาหา”ติ ฯเปฯ. ยถา มํ ภนฺเต ภควา ชานาตีติ. ทุติยมฺปิ โข ภควา ฯเปฯ. ตติยมฺปิ โข ภควา อายสฺมนฺตํ ทพฺพํ มลฺลปุตฺตํ เอตทโวจ “สรสิ ตฺวํ ทพฺพ เอวรูปํ กตฺตา ยถายํ วฑฺโฒ อาหา”ติ. ยถา มํ ภนฺเต ภควา ชานาตี’ติ. น โข ทพฺพ ทพฺพา เอวํ นิพฺเพเฐนฺติ, สเจ ตยา กตํ กตนฺติ วเทหิ, สเจ อกตํ อกตนฺติ วเทหีติ. ยโต อหํ ภนฺเต ชาโต นาภิชานามิ สุปินนฺเตนปิ เมถุนํ ธมฺมํ ปฏิเสวิตา, ปเคว ชาคโรติ. อถ โข ภควา ภิกฺขู อามนฺเตสิ\csromandash{} เตน หิ ภิกฺขเว สํโฆ วฑฺฒสฺส ลิจฺฉวิสฺส ปตฺตํ นิกฺกุชฺชตุ, อสมฺโภคํ สํเฆน กโรตุ.}

\prose{อฏฺฐหิ ภิกฺขเว องฺเคหิ สมนฺนาคตสฺส อุปาสกสฺส ปตฺโต นิกฺกุชฺชิตพฺโพ. ภิกฺขูนํ อลาภาย ปริสกฺกติ, ภิกฺขูนํ อนตฺถาย ปริสกฺกติ, ภิกฺขูนํ อวาสาย\csromansharedfootnote{fb64d2bf42bb9c0a}{อนาวาสาย (สฺยา)} ปริสกฺกติ, ภิกฺขู อกฺโกสติ ปริภาสติ, ภิกฺขู ภิกฺขูหิ เภเทติ, พุทฺธสฺส อวณฺณํ ภาสติ, ธมฺมสฺส อวณฺณํ ภาสติ, สํฆสฺส อวณฺณํ ภาสติ. อนุชานามิ ภิกฺขเว อิเมหิ อฏฺฐหงฺเคหิ สมนฺนาคตสฺส อุปาสกสฺส ปตฺตํ นิกฺกุชฺชิตุํ. เอวญฺจ ปน ภิกฺขเว นิกฺกุชฺชิตพฺโพ, พฺยตฺเตน ภิกฺขุนา ปฏิพเลน สํโฆ ญาเปตพฺโพ\csromandash{}}

\proseitem{266}{“สุณาตุ เม ภนฺเต สํโฆ, วฑฺโฒ ลิจฺฉวี อายสฺมนฺตํ ทพฺพํ มลฺลปุตฺตํ อมูลิกาย สีลวิปตฺติยา อนุทฺธํเสติ, ยทิ สํฆสฺส ปตฺตกลฺลํ, สํโฆ วฑฺฒสฺส ลิจฺฉวิสฺส ปตฺตํ นิกฺกุชฺเชยฺย, อสมฺโภคํ สํเฆน กเรยฺย, เอสา ญตฺติ.}

\prose{สุณาตุ เม ภนฺเต สํโฆ, วฑฺโฒ ลิจฺฉวี อายสฺมนฺตํ ทพฺพํ มลฺลปุตฺตํ อมูลิกาย สีลวิปตฺติยา อนุทฺธํเสติ, สํโฆ วฑฺฒสฺส ลิจฺฉวิสฺส ปตฺตํ นิกฺกุชฺชติ, อสมฺโภคํ สํเฆน กโรติ, ยสฺสายสฺมโต ขมติ วฑฺฒสฺส ลิจฺฉวิสฺส ปตฺตสฺส นิกฺกุชฺชนา, อสมฺโภคํ สํเฆน กรณํ, โส ตุณฺหสฺส, ยสฺส นกฺขมติ, โส ภาเสยฺย.}

\prose{นิกฺกุชฺชิโต สํเฆน วฑฺฒสฺส ลิจฺฉวิสฺส ปตฺโต อสมฺโภโค สํเฆน, ขมติ สํฆสฺส, ตสฺมา หุณฺหี, เอวเมตํ ธารยามี”ติ.}

\prose{\csromanfolio{264}อถ โข อายสฺมา อานนฺโท ปุพฺพณฺหสมยํ นิวาเสตฺวา ปตฺตจีวร มาทาย เยน วฑฺฒสฺส ลิจฺฉวิสฺส นิเวสนํ เตนุปสงฺกมิ, อุปสงฺกมิตฺวา วฑฺฒํ ลิจฺฉวิํ เอตทโวจ “สํเฆน เต อาวุโส วฑฺฒ ปตฺโต นิกฺกุชฺชิโต อสมฺโภโคสิ สํเฆนา”ติ. อถ โข วฑฺโฒ ลิจฺฉวี “สํเฆน กิร เม ปตฺโต นิกฺกุชฺชิโต อสมฺโภโคมฺหิ กิร สํเฆนา”ติ ตตฺเถว มุจฺฉิโต ปปโต. อถ โข วฑฺฒสฺส ลิจฺฉวิสฺส มิตฺตามจฺจา ญาติสาโลหิตา วฑฺฒํ ลิจฺฉวิํ เอตทโวจุํ “อลํ อาวุโส วฑฺฒ, มา โสจิ มา ปริเทวิ, มยํ ภควนฺตํ ปสาเทสฺสาม ภิกฺขุสํฆญฺจา”ติ.}

\prose{อถ โข วฑฺโฒ ลิจฺฉวี สปุตฺตทาโร สมิตฺตามจฺโจ สญาติสาโลหิโต อลฺลวตฺโถ อลฺลเกโส เยน ภควา เตนุปสงฺกมิ, อุปสงฺกมิตฺวา ภควโต ปาเทสุ สิรสา นิปติตฺวา ภควนฺตํ เอตทโวจ “อจฺจโย มํ ภนฺเต อจฺจคมา ยถาพาลํ ยถามูฬฺหํ ยถาอกุสลํ, โยหํ อยฺยํ ทพฺพํ มลฺลปุตฺตํ อมูลิกาย สีลวิปตฺติยา อนุทฺธํเสสิํ, ตสฺส เม ภนฺเต ภควา อจฺจยํ อจฺจยโต ปฏิคฺคณฺหาตุ อายติํ สํวรายา”ติ. ตคฺฆ ตฺวํ อาวุโส วฑฺฒ อจฺจโย อจฺจคมา ยถาพาลํ ยถามูฬฺหํ ยถาอกุสลํ, ยํ ตฺวํ ทพฺพํ มลฺลปุตฺตํ อมูลิกาย สีลวิปตฺติยา อนุทฺธํเสสิ, ยโต จ โข ตฺวํ อาวุโส วฑฺฒ อจฺจยํ อจฺจยโต ทิสฺวา ยถาธมฺมํ ปฏิกโรสิ, ตํ เต มยํ ปฏิคฺคณฺหาม, วุฑฺฒิเหสา อาวุโส วฑฺฒ อริยสฺส วินเย โย อจฺจยํ อจฺจยโต ทิสฺวา ยถาธมฺมํ ปฏิกโรติ, อายติํ สํวรํ อาปชฺชตีติ. อถ โข ภควา ภิกฺขู อามนฺเตสิ\csromandash{} เตน หิ ภิกฺขเว สํโฆ วฑฺฒสฺส ลิจฺฉวิสฺส ปตฺตํ อุกฺกุชฺชตุ, สมฺโภคํ สํเฆน กโรตุ.}

\prose{อฏฺฐหิ ภิกฺขเว องฺเคหิ สมนฺนาคตสฺส อุปาสกสฺส ปตฺโต อุกฺกุชฺชิตพฺโพ. น ภิกฺขูนํ อลาภาย ปริสกฺกติ, น ภิกฺขูนํ อนตฺถาย ปริสกฺกติ, น ภิกฺขูนํ อวาสาย ปริสกฺกติ, น ภิกฺขู อกฺโกสติ ปริภาสติ, น ภิกฺขู ภิกฺขูหิ เภเทติ, น พุทฺธสฺส อวณฺณํ ภาสติ, น ธมฺมสฺส อวณฺณํ ภาสติ, น สํฆสฺส อวณฺณํ ภาสติ. อนุชานามิ ภิกฺขเว อิเมหิ อฏฺฐหงฺเคหิ สมนฺนาคตสฺส อุปาสกสฺส ปตฺตํ อุกฺกุชฺชิตุํ. เอวญฺจ ปน ภิกฺขเว อุกฺกุชฺชิตพฺโพ. เตน ภิกฺขเว วฑฺเฒน ลิจฺฉวินา สํฆํ อุปสงฺกมิตฺวา เอกํสํ อุตฺตราสงฺคํ กริตฺวา ภิกฺขูนํ ปาเท วนฺทิตฺวา อุกฺกุฏิกํ นิสีทิตฺวา อญฺชลิํ ปคฺคเหตฺวา เอวมสฺส วจนีโย “สํเฆน เม ภนฺเต \csromanfolio{265}ปตฺโต นิกฺกุชฺชิโต, อสมฺโภโคมฺหิ สํเฆน, โสหํ ภนฺเต สมฺมา วตฺตามิ, โลมํ ปาเตมิ, เนตฺถารํ วตฺตามิ, สํฆํ ปตฺตุกฺกุชฺชนํ ยาจามี”ติ. ทุติยมฺปิ ยาจิตพฺโพ. ตติยมฺปิ ยาจิตพฺโพ. พฺยตฺเตน ภิกฺขุนา ปฏิพเลน สํโฆ ญาเปตพฺโพ\csromandash{}}

\proseitem{267}{“สุณาตุ เม ภนฺเต สํโฆ, สํเฆน วฑฺฒสฺส ลิจฺฉวิสฺส ปตฺโต นิกฺกุชฺชิโต, อสมฺโภโค สํเฆน, โส สมฺมา วตฺตติ, โลมํ ปาเตติ, เนตฺถารํ วตฺตติ, สํฆํ ปตฺตุกฺกุชฺชนํ ยาจติ, ยทิ สํฆสฺส ปตฺตกลฺลํ สํโฆ วฑฺฒสฺส ลิจฺฉวิสฺส ปตฺตํ อุกฺกุชฺเชยฺย, สมฺโภคํ สํเฆน กเรยฺย, เอสา ญตฺติ.}

\prose{สุณาตุ เม ภนฺเต สํโฆ, สํเฆน วฑฺฒสฺส ลิจฺฉวิสฺส ปตฺโต นิกฺกุชฺชิโต, อสมฺโภโค สํเฆน, โส สมฺมา วตฺตติ, โลมํ ปาเตติ, เนตฺถารํ วตฺตติ, สํฆํ ปตฺตุกฺกุชฺชนํ ยาจติ, สํโฆ วฑฺฒสฺส ลิจฺฉวิสฺส ปตฺตํ อุกฺกุชฺชติ, สมฺโภคํ สํเฆน กโรติ, ยสฺสายสฺมโต ขมติ วฑฺฒสฺส ลิจฺฉวิสฺส ปตฺตสฺส อุกฺกุชฺชนา, สมฺโภคํ สํเฆน กรณํ, โส ตุณฺหสฺส, ยสฺส นกฺขมติ, โส ภาเสยฺย.}

\prose{อุกฺกุชฺชิโต สํเฆน วฑฺฒสฺส ลิจฺฉวิสฺส ปตฺโต สมฺโภโค สํเฆน, ขมติ สํฆสฺส, ตสฺมา ตุณฺหี, เอวเมตํ ธารยามี”ติ.}

\proseitem{268}{อถ โข ภควา เวสาลิยํ ยถาภิรนฺตํ วิหริตฺวา เยน ภคฺคา เตน จาริกํ ปกฺกามิ, อนุปุพฺเพน จาริกํ จรมาโน เยน ภคฺคา ตทวสริ, ตตฺร สุทํ ภควา ภคฺเคสุ วิหรติ สุสุมารคิเร\csromansharedfootnote{e4b6aaabfd556bad}{สุํสูมารคิเร (สี, สฺยา), สํสุมารคิเร (ก)} เภสกฬวเน มิคทาเย. เตน โข ปน สมเยน โพธิสฺส ราชกุมารสฺส โกกนโท\csromansharedfootnote{6204b5cd431ed457}{โกกนุโท (ก)} นาม ปาสาโท อจิรการิโต โหติ อนชฺฌาวุตฺโถ สมเณน วา พฺราหฺมเณน วา เกนจิ วา มนุสฺสภูเตน. อถ โข โพธิ ราช กุมาโร สญฺจิกาปุตฺตํ มาณวํ อามนฺเตสิ\csromandash{}เอหิ ตฺวํ สมฺม สญฺจิกาปุตฺต, เยน ภควา เตนุปสงฺกม, อุปสงฺกมิตฺวา มม วจเนน ภควโต ปาเท สิรสา วนฺท, อปฺปาพาธํ อปฺปาตงฺกํ ลหุฏฺฐานํ พลํ ผาสุวิหารํ ปุจฺฉ “โพธิ ภนฺเต ราชกุมาโร ภควโต ปาเท สิรสา วนฺทติ, อปฺปาพาธํ \csromanfolio{266}อปฺปาตงฺกํ ลหุฏฺฐานํ พลํ ผาสุวิหารํ ปุจฺฉตี”ติ. เอวญฺจ วเทหิ “อธิวาเสตุ กิร ภนฺเต ภควา โพธิสฺส ราชกุมารสฺส สฺวาตนาย ภตฺตํ สทฺธิํ ภิกฺขุสํเฆนา”ติ. “เอวํ โภ”ติ โข สญฺจิกาปุตฺโต มาณโว โพธิสฺส ราชกุมารสฺส ปฏิสฺสุตฺวา เยน ภควา เตนุปสงฺกมิ, อุปสงฺกมิตฺวา ภควตา สทฺธิํ สมฺโมทิ, สมฺโมทนียํ กถํ สารณียํ วีติสาเรตฺวา เอกมนฺตํ นิสีทิ, เอกมนฺตํ นิสินฺโน โข สญฺจิกาปุตฺโต มาณโว ภควนฺตํ เอตทโวจ\csromandash{}โพธิ โข ราชกุมาโร โภโต โคตมสฺส ปาเท สิรสา วนฺทติ, อปฺปาพาธํ อปฺปาตงฺกํ ลหุฏฺฐานํ พลํ ผาสุวิหารํ ปุจฺฉติ, เอวญฺจ วเทติ “อธิวาเสตุ กิร ภวํ โคตโม โพธิสฺส ราชกุมารสฺส สฺวาตนาย ภตฺตํ สทฺธิํ ภิกฺขุสํเฆนา”ติ. อธิวาเสสิ ภควา ตุณฺหีภาเวน. อถ โข สญฺจิกาปุตฺโต มาณโว ภควโต อธิวาสนํ วิทิตฺวา อุฏฺฐายาสนา เยน โพธิ ราชกุมาโร เตนุปสงฺกมิ, อุปสงฺกมิตฺวา โพธิํ ราชกุมารํ เอตทโวจ “อโวจุมฺห โข มยํ โภโต วจเนน ตํ ภวนฺตํ โคตมํ ‘โพธิ โข ราชกุมาโร โภโต โคตมสฺส ปาเท สิรสา วนฺทติ, อปฺปาพาธํ อปฺปาตงฺกํ ลหุฏฺฐานํ พลํ ผาสุวิหารํ ปุจฺฉติ, เอวญฺจ วเทติ อธิวาเสตุ กิร ภวํ โคตโม โพธิสฺส ราชกุมารสฺส สฺวาตนาย ภตฺตํ สทฺธิํ ภิกฺขุสํเฆนา’ติ, อธิวุตฺถญฺจ ปน สมเณน โคตเมนา”ติ.}

\prose{อถ โข โพธิ ราชกุมาโร ตสฺสา รตฺติยา อจฺจเยน ปณีตํ ขาทนียํ โภชนียํ ปฏิยาทาเปตฺวา โกกนทญฺจ ปาสาทํ โอทาเตหิ ทุสฺเสหิ สนฺตราเปตฺวา ยาว ปจฺฉิม\-โสปาน\-กเฬวรา สญฺจิกาปุตฺตํ มาณวํ อามนฺเตสิ\csromandash{}เอหิ ตฺวํ สมฺม สญฺจิกาปุตฺต, เยน ภควา เตนุปสงฺกม, อุปสงฺกมิตฺวา ภควโต กาลํ อาโรเจหิ “กาโล ภนฺเต, นิฏฺฐิตํ ภตฺตนฺ”ติ. “เอวํ โภ”ติ โข สญฺจิกาปุตฺโต มาณโว โพธิสฺส ราชกุมารสฺส ปฏิสฺสุตฺวา เยน ภควา เตนุปสงฺกมิ, อุปสงฺกมิตฺวา ภควโต กาลํ อาโรเจสิ “กาโล โภ โคตม, นิฏฺฐิตํ ภตฺตนฺ”ติ.}

\prose{อถ โข ภควา ปุพฺพณฺหสมยํ นิวาเสตฺวา ปตฺตจีวรทาย เยน โพธิสฺส ราชกุมารสฺส นิเวสนํ เตนุปสงฺกมิ. เตน โข ปน สมเยน โพธิ ราชกุมาโร พหิทฺวาร\-โกฏฺฐเก ฐิโต โหติ ภควนฺตํ อาคมยมาโน. อทฺทสา โข โพธิ ราชกุมาโร ภควนฺตํ \csromanfolio{267}ทูรโตว อาคจฺฉนฺตํ, ทิสฺวาน ตโต ปจฺจุคฺคนฺตฺวา ภควนฺตํ อภิวาเทตฺวา ปุเรกฺขตฺวา เยน โกกนโท ปาสาโท เตนุปสงฺกมิ. อถ โข ภควา ปจฺฉิม\-โสปาน\-กเฬวรํ นิสฺสาย อฏฺฐาสิ. อถ โข โพธิ ราชกุมาโร ภควนฺตํ เอตทโวจ “อกฺกมตุ ภนฺเต ภควา ทุสฺสานิ, อกฺกมตุ สุคโต ทุสฺสานิ, ยํ มม อสฺส ทีฆรตฺตํ หิตาย สุขายา”ติ. เอวํ วุตฺเต ภควา ตุณฺหี อโหสิ. ทุติยมฺปิ โข ฯเปฯ. ตติยมฺปิ โข โพธิ ราชกุมาโร ภควนฺตํ เอตทโวจ “อกฺกมตุ ภนฺเต ภควา ทุสฺสานิ, อกฺกมตุ สุคโต ทุสฺสานิ, ยํ มม อสฺส ทีฆรตฺตํ หิตาย สุขายา”ติ. อถ โข ภควา อายสฺมนฺตํ อานนฺทํ อปโลเกสิ. อถ โข อายสฺมา อานนฺโท โพธิํ ราชกุมารํ เอตทโวจ “สํหรนฺตุ ราชกุมาร ทุสฺสานิ, น ภควา เจลปฏิกํ\csromansharedfootnote{4bf9b8385cc8f2c6}{เจลปตฺติกํ (สี)} อกฺกมิสฺสติ, ปจฺฉิมํ ชนตํ ตถาคโต อนุกมฺปตี”ติ.}

\prose{อถ โข โพธิ ราชกุมาโร ทุสฺสานิ สํหราเปตฺวา อุปริโกกนเท ปาสาเท อาสนํ ปญฺญเปสิ. อถ โข ภควา โกกนทํ ปาสาทํ อภิรุหิตฺวา ปญฺญตฺเต อาสเน นิสีทิ สทฺธิํ ภิกฺขุ\-สํเฆน. อถ โข โพธิ ราชกุมาโร พุทฺธ\-ปฺปมุขํ ภิกฺขุสํฆํ ปณีเตน ขาทนีเยน โภชนีเยน สหตฺถา สนฺตปฺเปตฺวา สมฺปวาเรตฺวา ภควนฺตํ ภุตฺตาวิํ โอนีต\-ปตฺต\-ปาณิํ เอกมนฺตํ นิสีทิ, เอกมนฺตํ นิสินฺนํ โข โพธิํ ราชกุมารํ ภควา ธมฺมิยา กถาย สนฺทสฺเสตฺวา สมาทเปตฺวา สมุตฺเตเชตฺวา สมฺปหํเสตฺวา อุฏฺฐายาสนา ปกฺกามิ. อถ โข ภควา เอตสฺมิํ นิทาเน เอตสฺมิํ ปกรเณ ธมฺมิํ กถํ กตฺวา ภิกฺขู อามนฺเตสิ “น ภิกฺขเว เจลปฺปฏิกา อกฺกมิตพฺพา, โย อกฺกเมยฺย, อาปตฺติ ทุกฺกฏสฺสา”ติ.}

\prose{เตน โข ปน สมเยน อญฺญตรา อิตฺตี อปคตคพฺภา ภิกฺขู นิมนฺเตตฺวา ทุสฺสํ ปญฺญเปตฺวา เอตทโวจ “อกฺกมถ ภนฺเต ทุสฺสนฺ”ติ. ภิกฺขู กุกฺกุจฺจายนฺตา น อกฺกมนฺติ. อกฺกมถ ภนฺเต ทุสฺสํ มงฺคล\-ตฺถา\-ยาติ. ภิกฺขู กุกฺกุจฺจายนฺตา น อกฺกมิํสุ. อถ โข สา อิตฺถี อุชฺฌายติ ขิยฺยติ วิปาเจติ “กถํ หิ นาม อยฺยา มงฺคลตฺถาย ยาจิยมานา เจลปฺปฏิกํ น อกฺกมิสฺสนฺตี”ติ. อโสสุํ โข ภิกฺขู ตสฺสา อิตฺถิยา อุชฺฌายนฺติยา ขิยฺยนฺติยา วิปาเจนฺติยา. อถ โข เต ภิกฺขู ภควโต \csromanfolio{268}เอตมตฺถํ อาโรเจสุํ. คิหี ภิกฺขเว มงฺคลิกา. อนุชานามิ ภิกฺขเว คิหีนํ มงฺคลตฺถาย ยาจิยมาเนน เจลปฺปฏิกํ อกฺกมิตุนฺติ.}

\prosewithcloser{เตน โข ปน สมเยน ภิกฺขู โธตปาทกํ อกฺกมิตุํ กุกฺกุจฺจายนฺติ. ภควโต เอตมตฺถํ อาโรเจสุํ. อนุชานามิ ภิกฺขเว โธตปาทกํ อกฺกมิตุนฺติ.}{ทุติย\-ภาณวาโร นิฏฺฐิโต.}
\proseitem{269}{อถ โข ภควา ภคฺเคสุ ยถาภิรนฺตํ วิหริตฺวา เยน สาวตฺถิ เตน จาริกํ ปกฺกามิ, อนุปุพฺเพน จาริกํ จรมาโน เยน สาวตฺถิ ตทวสริ, ตตฺร สุทํ ภควา สาวตฺถิยํ วิหรติ เชตวเน อนาถ\-ปิณฺฑิกสฺส อาราเม. อถ โข วิสาขา มิคารมาตา ฆฏกญฺจ กตกญฺจ สมฺมชฺชนิญฺจ อาทาย เยน ภควา เตนุปสงฺกมิ, อุปสงฺกมิตฺวา ภควนฺตํ อภิวาเทตฺวา เอกมนฺตํ นิสีทิ, เอกมนฺตํ นิสินฺนา โข วิสาขา มิคารมาตา ภควนฺตํ เอตทโวจ “ปฏิคฺคณฺหาตุ เม ภนฺเต ภควา ฆฏกญฺจ กตกญฺจ สมฺมชฺชนิญฺจ, ยํ มม อสฺส ทีฆรตฺตํ หิตาย สุขายา”ติ. ปฏิคฺคเหสิ ภควา ฆฏกญฺจ สมฺมชฺชนิญฺจ, น ภควา กตกํ ปฏิคฺคเหสิ. อถ โข ภควา วิสาขํ มิคารมาตรํ ธมฺมิยา กถาย สนฺทสฺเสสิ สมาทเปสิ สมุตฺเตเชสิ สมฺปหํเสสิ. อถ โข วิสาขา มิคารมาตา ภควตา ธมฺมิยา กถาย สนฺทสฺสิตา สมาทปิตา สมุตฺเตชิตา สมฺปหํสิตา อุฏฺฐายาสนา ภควนฺตํ อภิวาเทตฺวา ปทกฺขิณํ กตฺวา ปกฺกามิ.}

\prose{อถ โข ภควา เอตสฺมิํ นิทาเน เอตสฺมิํ ปกรเณ ธมฺมิํ กถํ กตฺวา ภิกฺขู อามนฺเตสิ “อนุชานามิ ภิกฺขเว ฆฏกญฺจ สมฺมชฺชนิญฺจ. น ภิกฺขเว กตกํ ปริภุญฺชิ\-ตพฺพํ, โย ปริภุญฺเชยฺย, อาปตฺติ ทุกฺกฏสฺส. อนุชานามิ ภิกฺขเว ติสฺโส ปาทฆํสนิโย สกฺขรํ กถลํ สมุทฺทเผณกนฺ”ติ.}

\prose{อถ โข วิสาขา มิคารมาตา วิธูปนญฺจ ตาลวณฺฏญฺจ\csromansharedfootnote{fd84ab9873cf176f}{ตาลวณฺฑญฺจ (ก)} อาทาย เยน ภควา เตนุปสงฺกมิ, อุปสงฺกมิตฺวา ภควนฺตํ อภิวาเทตฺวา เอกมนฺตํ นิสีทิ, เอกมนฺตํ นิสินฺนา โข วิสาขา มิคารมาตา ภควนฺตํ นิสินฺนา โข วิสาขา มิคารมาตา ภควนฺตํ เอตทโวจ “ปฏิคฺคณฺหาตุ เม ภนฺเต ภควา วิธูปนญฺจ ตาลวณฺฏญฺจ, ยํ มม อสฺส ทีฆรตฺตํ หิตาย สุขายา”ติ. ปฏิคฺคเหสิ ภควา วิธูปนญฺจ ตาลวณฺฏญฺจ.}

\prose{\csromanfolio{269}อถ โข ภควา วิสาขํ มิคารมาตรํ ธมฺมิยา กถาย สนฺทสฺเสสิ สมาทเปสิ สมุตฺเตเชสิ สมฺปหํเสสิ ฯเปฯ ปทกฺขิณํ กตฺวา ปกฺกามิ. อถ โข ภควา เอตสฺมิํ นิทาเน เอตสฺมิํ ปกรเณ ธมฺมิํ กถํ กตฺวา ภิกฺขู อามนฺเตสิ “อนุชานามิ ภิกฺขเว วิธูปนญฺจ ตาลวณฺฏญฺจา”ติ.}

\prose{เตน โข ปน สมเยน สํฆสฺส มกสพีชนี อุปฺปนฺนา โหติ. ภควโต เอตมตฺถํ อาโรเจสุํ. อนุชานามิ ภิกฺขเว มกส\-พีชนินฺติ. จามริพีชนี อุปฺปนฺนา โหติ. ภควโต เอตมตฺถํ อาโรเจสุํ. น ภิกฺขเว จามริพีชนี ธาเรตพฺพา, โย ธาเรยฺย, อาปตฺติ ทุกฺกฏสฺส. อนุชานามิ ภิกฺขเว ติสฺโส พีชนิโย วากมยํ อุสีรมยํ โมรปิญฺฉ\-มยนฺติ.}

\proseitem{270}{เตน โข ปน สมเยน สํฆสฺส ฉตฺตํ อุปฺปนฺนํ โหติ. ภควโต เอตมตฺถํ อาโรเจสุํ. อนุชานามิ ภิกฺขเว ฉตฺตนฺติ.}

\prose{เตน โข ปน สมเยน ฉพฺพคฺคิยา ภิกฺขู ฉตฺต\-ปฺปคฺคหิตา\csromansharedfootnote{d6b9667eb322e5d3}{ฉตฺตํ มคฺคเหตฺวา (ก)} อาหิณฺฑนฺติ. เตน โข ปน สมเยน อญฺญตโร อุปาสโก สมฺพหุเลหิ อาชีวก\-สาวเกหิ สทฺธิํ อุยฺยานํ อคมาสิ. อทฺทสาสุํ โข เต อาชีวกสาวกา ฉพฺพคฺคิเย ภิกฺขู ทูรโตว ฉตฺต\-ปฺปคฺคหิเต อาคจฺฉนฺเต, ทิสฺวาน ตํ อุปาสกํ เอตทโวจุํ “เอเต โข อยฺยา\csromansharedfootnote{d8149a5003b73eef}{อยฺโย (ก)} ตุมฺหากํ ภทนฺตา ฉตฺต\-ปฺปคฺคหิตา อาคจฺฉนฺติ เสยฺยถาปิ คณก\-มหามตฺตา”ติ. นายฺยา เอเต ภิกฺขู ปริพฺพาชกาติ. “ภิกฺขู น ภิกฺขู”ติ อพฺภุตํ อกํสุ. อถ โข โส อุปาสโก อุปคเต สญฺชานิตฺวา อุชฺฌายติ ขิยฺยติ วิปาเจติ “กถํ หิ นาม ภทนฺตา ฉตฺต\-ปฺปคฺคหิตา อาหิณฺฑิสฺสนฺตี”ติ. อสฺโสสุํ โข ภิกฺขู ตสฺส อุปาสกสฺส อุชฺฌายนฺตสฺส ขิยฺยนฺตสฺส วิปาเจนฺตสฺส. อถ โข เต ภิกฺขู ภควโต เอตมตฺถํ อาโรเจสุํ. สจฺจํ กิร ภิกฺขเว ฯเปฯ. สจฺจํ ภควาติ ฯเปฯ วิครหิตฺวา ฯเปฯ ธมฺมิํ กถํ กตฺวา ภิกฺขู อามนฺเตสิ “น ภิกฺขเว ฉตฺตํ ธาเรตพฺพํ, โย ธาเรยฺย, อาปตฺติ ทุกฺกฏสฺสา”ติ.}

\prose{เตน โข ปน สมเยน อญฺญตโร ภิกฺขุ คิลาโน โหติ, ตสฺส ภิกฺขุโน วินา ฉตฺตํ น ผาสุ โหติ. ภควโต เอตมตฺถํ อาโรเจสุํ. อนุชานามิ ภิกฺขเว คิลานสฺส ฉตฺตํ ธาเรตุนฺติ.}

\prose{\csromanfolio{270}เตน โข ปน สมเยน ภิกฺขู “คิลานสฺเสว ภควตา ฉตฺตํ อนุญฺญาตํ โน อคิลานสฺสา”ติ อาราเม อารามูปจาเร ฉตฺตํ ธาเรตุํ กุกฺกุจฺจายนฺติ. ภควโต เอตมตฺถํ อาโรเจสุํ. อนุชานามิ ภิกฺขเว อคิลาเนนปิ อาราเม อารามูปจาเร ฉตฺตํ ธาเรตุนฺติ.}

\prose{เตน โข ปน สมเยน อญฺญตโร ภิกฺขุ สิกฺกาย ปตฺตํ อุฏฺฏิตฺวา ทณฺเฑ อาลคฺคิตฺวา วิกาเล อญฺญตเรน คามทฺวาเรน อติกฺกมติ. มนุสฺสา “เอสยฺโย โจโร คจฺฉติ, อสิสฺส วิชฺโชตลตี”ติ อนุปติตฺวา คเหตฺวา สญฺชานิตฺวา มุญฺจิํสุ. อถ โข โส ภิกฺขุ อารามํ คนฺตฺวา ภิกฺขูนํ เอตมตฺถํ อาโรเจสิ. กิํ ปน ตฺวํ อาวุโส ทณฺฑสิกฺกํ ธาเรสีติ. เอวมาวุโสติ. เย เต ภิกฺขู อปฺปิจฺฉา ฯเปฯ เต อุชฺฌายนฺติ ขิยฺยนฺติ วิปาเจนฺติ “กถํ หิ นาม ภิกฺขุ ทณฺฑสิกฺกํ ธาเรสฺสสี”ติ. อถ โข เต ภิกฺขู ภควโต เอตมตฺถํ อาโรเจสุํ ฯเปฯ. สจฺจํ ภควาติ ฯเปฯ วิครหิตฺวา ฯเปฯ ธมฺมิํ กถํ กตฺวา ภิกฺขู อามนฺเตสิ “น ภิกฺขเว ทณฺฑสิกฺกา ธาเรตพฺพา, โย ธาเรยฺย, อาปตฺติ ทุกฺกฏสฺสา”ติ.}

\prose{เตน โข ปน สมเยน อญฺญตโร ภิกฺขุ คิลาโน โหติ, น สกฺโกติ วินา ทณฺเฑน อาหิณฺฑิตุํ. ภควโต เอตมตฺถํ อาโรเจสุํ. อนุชานามิ ภิกฺขเว คิลานสฺส ภิกฺขุโน ทณฺฑสมฺมุติํ ทาตุํ. เอวญฺจ ปน ภิกฺขเว ทาตพฺพา, เตน คิลาเนน ภิกฺขุนา สํฆํ อุปสงฺกมิตฺวา เอกํสํ อุตฺตราสงฺคํ กริตฺวา วุฑฺฒานํ ภิกฺขูนํ ปาเท วนฺทิตฺวา อุกฺกุฏิกํ นิสีทิตฺวา อญฺชลิํ ปคฺคเหตฺวา เอวมสฺส วจนีโย “อหํ ภนฺเต คิลาโน น สกฺโกมิ วินา ทณฺเฑน อาหิณฺฑิตุํ, โสหํ ภนฺเต สํฆํ ทณฺฑสมฺมุติํ ยาจมี”ติ. ทุติยมฺปิ ยาจิตพฺพา. ตติยมฺปิ ยาจิตพฺพา. พฺยตฺเตน ภิกฺขุนา ปฏิพเลน สํโฆ ญาเปตพฺโพ\csromandash{}}

\prose{“สุณาตุ เม ภนฺเต สํโฆ, อยํ อิตฺถนฺนาโม ภิกฺขุ คิลาโน น สกฺโกติ วินา ทณฺเฑน อาหิณฺฑิตุํ, โส สํฆํ ทณฺฑสมฺมุติํ ยาจติ, ยทิ สํฆสฺส ปตฺตกลฺลํ, สํโฆ อิตฺถนฺนามสฺส ภิกฺขุโน ทณฺฑสมฺมุติํ ทเทยฺย, เอสา ญตฺติ.}

\prose{สุณาตุ เม ภนฺเต สํโฆ, อยํ อิตฺถนฺนาโม ภิกฺขุ คิลาโน น สกฺโกติ วิเน ทณฺเฑน อาหิณฺฑิตุํ, โส สํฆํ ทณฺฑสมฺมุติํ ยาจติ, \csromanfolio{271}สํโฆ อิตฺถนฺนามสฺส ภิกฺขุโน ทณฺฑสมฺมุติํ เทติ, ยสฺสายสฺมโต ขมติ อิตฺถนฺนามสฺส ภิกฺขุโน ทณฺฑสมฺมุติยา ทานํ, โส ตุณฺหสฺส, ยสฺส นกฺขมติ, โส ภาเสยฺย.}

\prose{ทินฺนา สํเฆน อิตฺถนฺนามสฺส ภิกฺขุโน ทณฺฑสมฺมุติ, ขมติ สํฆสฺส, ตสฺมา ตุณฺหี, เอวเมตํ ธารยามี”ติ.}

\proseitem{271}{เตน โข ปน สมเยน อญฺญตโร ภิกฺขุ คิลาโน โหติ, น สกฺโกติ วินา สิกฺกาย ปตฺตํ ปริหริตุํ. ภควโต เอตมตฺถํ อาโรเจสุํ. อนุชานามิ ภิกฺขเว คิลานสฺส ภิกฺขุโน สิกฺกาสมฺมุติํ ทาตุํ. เอวญฺจ ปน ภิกฺขเว ทาตพฺพา, เตน คิลาเนน ภิกฺขุนา สํฆํ อุปสงฺกมิตฺวา เอกํสํ อุตฺตราสงฺคํ กริตฺวา วุฑฺฒานํ ภิกฺขูนํ ปาเท วนฺทิตฺวา อุกฺกุฏิกํ นิสีทิตฺวา อญฺชลิํ ปคฺคเหตฺวา เอวมสฺส วจนีโย “อหํ ภนฺเต คิลาโน น สกฺโกมิ วินา สิกฺกาย ปตฺตํ ปริหริตุํ, โสหํ ภนฺเต สํฆํ สิกฺกาสมฺมุติํ ยาจามี”ติ. ทุติยมฺปิ ยาจิตพฺพา. ตติยมฺปิ ยาจิตพฺพา. พฺยตฺเตน ภิกฺขุนา ปฏิพเลน สํโฆ ญาเปตพฺโพ\csromandash{}}

\prose{“สุณาตุ เม ภนฺเต สํโฆ, อยํ อิตฺถนฺนาโม ภิกฺขุ คิลาโน น สกฺโกติ วินา สิกฺกาย ปตฺตํ ปริหริตุํ, โส สํฆํ สิกฺกาสมฺมุติํ ยาจติ, ยทิ สํฆสฺส ปตฺตกลฺลํ, สํโฆ อิตฺถนฺนามสฺส ภิกฺขุโน สิกฺกาสมฺมุติํ ทเทยฺย, เอสา ญตฺติ.}

\prose{สุณาตุ เม ภนฺเต สํโฆ, อยํ อิตฺถนฺนาโม ภิกฺขุ คิลาโน น สกฺโกติ วินา สิกฺกาย ปตฺตํ ปริหริตุํ, โส สํฆํ สิกฺกาสมฺมุติํ ยาจติ, สํโฆ อิตฺถนฺนามสฺส ภิกฺขุโน สิกฺกาสมฺมุติํ เทติ, ยสฺสายสฺมโต ขมติ อิตฺถนฺนามสฺส ภิกฺขุโน สิกฺกาสมฺมุติยา ทานํ, โส ตุณฺหสฺส, ยสฺส นกฺขมติ, โส ภาเสยฺย.}

\prose{ทินฺนา สํเฆน อิตฺถนฺนามสฺส ภิกฺขุโน สิกฺกาสมฺมุติ, ขมติ สํฆสฺส, ตสฺมา ตุณฺหี, เอวเมตํ ธารยามี”ติ.}

\proseitem{272}{เตน โข ปน สมเยน อญฺญตโร ภิกฺขุ คิลาโน โหติ, น สกฺโกติ วินา ทณฺเฑน อาหิณฺฑิตุํ, น สกฺโกติ วินา สิกฺกาย ปตฺตํ ปริหริตุํ. ภควโต เอตมตฺถํ อาโรเจสุํ. อนุชานามิ ภิกฺขเว คิลานสฺส ภิกฺขุโน ทณฺฑ\-สิกฺกา\-สมฺมุติํ ทาตุํ. เอวญฺจ ปน ภิกฺขเว \csromanfolio{272}ทาตพฺพา, เตน คิลาเนน ภิกฺขุนา สํฆํ อุปสงฺกมิตฺวา เอกํสํ อุตฺตราสงฺคํ กริตฺวา วุฑฺฒานํ ภิกฺขูนํ ปาเท วนฺทิตฺวา อุกฺกุฏิกํ นิสีทิตฺวา อญฺจลิํ ปคฺคเหตฺวา เอวมสฺส วจนีโย “อหํ ภนฺเต คิลาโน น สกฺโกมิ วินา ทณฺเฑน อาหิณฺฑิตุํ, น สกฺโกมิ วินา สิกฺกาย ปตฺตํ ปริหริตุํ, โสหํ ภนฺเต สํฆํ ทณฺฑ\-สิกฺกา\-สมฺมุติํ ยาจามี”ติ. ทุติยมฺปิ ยาจิตพฺพา. ตติยมฺปิ ยาจิตพฺพา. พฺยตฺเตน ภิกฺขุนา ปฏิพเลน สํโฆ ญาเปตพฺโพ\csromandash{}}

\prose{“สุณาตุ เม ภนฺเต สํโฆ, อยํ อิตฺถนฺนาโม ภิกฺขุ คิลาโน น สกฺโกติ วินา ทณฺเฑน อาหิณฺฑิตุํ, น สกฺโกติ วินา สิกฺกาย ปตฺตํ ปริหริตุํ, โส สํฆํ ทณฺฑ\-สิกฺกา\-สมฺมุติํ ยาจติ, ยทิ สํฆสฺส ปตฺตกลฺลํ, สํโฆ อิตฺถนฺนามสฺส ภิกฺขุโน ทณฺฑ\-สิกฺกา\-สมฺมุติํ ทเทยฺย, เอสา ญตฺติ.}

\prose{สุณาตุ เม ภนฺเต สํโฆ, อยํ อิตฺถนฺนาโม ภิกฺขุ คิลาโน น สกฺโกติ วินา ทณฺเฑน อาหิณฺฑิตุํ, น สกฺโกติ วินา สิกฺกาย ปตฺตํ ปริหริตุํ, โส สํฆํ ทณฺฑ\-สิกฺกา\-สมฺมุติํ ยาจติ, สํโฆ อิตฺถนฺนามสฺส ภิกฺขุโน ทณฺฑ\-สิกฺกา\-สมฺมุติํ เทติ, ยสฺสายสฺมโต ขมติ อิตฺถนฺนามสฺส ภิกฺขุโน ทณฺฑ\-สิกฺกา\-สมฺมุติยา ทานํ, โส ตุณฺหสฺส, ยสฺส นกฺขมติ, โส ภาเสยฺย.}

\prose{ทินฺนา สํเฆน อิตฺถนฺนามสฺส ภิกฺขุโน ทณฺฑ\-สิกฺกา\-สมฺมุติ, ขมติ สํฆสฺส, ตสฺมา ตุณฺหี, เอวเมตํ ธารยามี”ติ.}

\proseitem{273}{เตน โข ปน สมเยน อญฺญตโร ภิกฺขุ โรมนฺถโก\csromansharedfootnote{b36b94bffa2a32bc}{เอรามฏฺฐโก (ก)} โหติ, โส โรมนฺถิตฺวา โรมนฺถิตฺวา อชฺโฌหรติ. ภิกฺขู อุชฺฌายนฺติ ขิยฺยนฺติ วิปาเจนฺติ “วิกาลายํ\csromansharedfootnote{d3cf567d1884f979}{กถํ หิ นาม วิกาลายํ (ก)} ภิกฺขุ โภชนํ ภุญฺชตี”ติ. ภควโต เอตมตฺถํ อาโรเจสุํ. เอโส ภิกฺขเว ภิกฺขุ อจิรํโคโยนิยา จุโต. อนุชานามิ ภิกฺขเว โรมนฺถกสฺส โรมนฺถนํ. น จ ภิกฺขเว พหิมุขทฺวารํ นีหริตฺวา อชฺโฌหริตพฺพํ, โย อชฺโฌหเรยฺย, ยถาธมฺโม กาเรตพฺโพติ.}

\prose{\csromanfolio{273}เตน โข ปน สมเยน อญฺญตรสฺส ปูคสฺส สํฆภตฺตํ โหติ, ภตฺตคฺเค พหุสิตฺถานิ ปกิริยิํสุ\csromansharedfootnote{711c1ed127ed5522}{วิปฺปกิรียิํสุ (สี), ปริกิริํสุ (สฺยา)}. มนุสฺสา อุชฺฌายนฺติ ขิยฺยนฺติ วิปาเจนฺติ “กถํ หิ นาม สมณา สกฺยปุตฺติยา โอทเน ทิยฺยมาเน น สกฺกจฺจํ ปฏิคฺคเหสฺสนฺติ, เอกเมกํ สิตฺถํ กมฺมสเตน นิฏฺฐายตี”ติ. อสฺโสสุํ โข ภิกฺขู เตสํ มนุสฺสานํ อุชฺฌายนฺตานํ ขิยฺยนฺตานํ วิปาเจนฺตานํ. อถ โข เต ภิกฺขู ภควโต เอตมตฺถํ อาโรเจสุํ. อนุชานามิ ภิกฺขเว ยํ ทิยฺยมานํ ปตติ, ตํ สามํ คเหตฺวา ปริภุญฺชิตุํ, ปริจฺจตฺตํ ตํ ภิกฺขเว ทายเกหีติ.}

\proseitem{274}{เตน โข ปน สมเยน อญฺญตโร ภิกฺขุ ทีเฆหิ นเขหิ ปิณฺฑาย จรติ, อญฺญตรา อิตฺถี ปสฺสิตฺวา ตํ ภิกฺขุํ เอตทโวจ “เอหิ ภนฺเต เมถุนํ ธมฺมํ ปฏิเสวา”ติ. อลํ ภคินิ, เนตํ กปฺปตีติ. สเจ โข ตฺวํ ภนฺเต นปฺปฏิเสวิสฺสสิ, อิทานาหํ อตฺตโน นเขหิ คตฺตานิ วิลิขิตฺวา กุปฺปํ กริสฺสามิ “อยํ มํ ภิกฺขุ วิปฺปกโรตี”ติ. ปชานาหิ ตฺวํ ภคินีติ. อถ โข สา อิตฺถี อตฺตโน นเขหิ คตฺตานิ วิลิขิตฺวา กุปฺปํ อกาสิ “อยํ มํ ภิกฺขุ วิปฺปกโรตี”ติ. มนุสฺสา อุปธาวิตฺวา ตํ ภิกฺขุํ อคฺคเหสุํ. อทฺทสาสุํ โข เต มนุสฺสา ตสฺสา อิตฺถิยา นเข ฉวิมฺปิ โลหิตมฺปิ. ทิสฺวาน “อิมิสฺสาเยว อิตฺถิยา อิทํ กมฺมํ, อการโก ภิกฺขู”ติ ตํ ภิกฺขุํ มุญฺจิํสุ. อถ โข โส ภิกฺขุ อารามํ คนฺตฺวา ภิกฺขูนํ เอตมตฺถํ อาโรเจสิ. กิํ ปน ตฺวํ อาวุโส ทีเฆ นเข ธาเรสีติ. เอวมาวุโสติ. เย เต ภิกฺขู อปฺปิจฺฉา ฯเปฯ เต อุชฺฌายนฺติ ขิยฺยนฺติ วิปาเจนฺติ “กถํ หิ นาม ภิกฺขุ ทีเฆ นเข ธาเรสฺสสี”ติ. อถ โข เต ภิกฺขู ภควโต เอตมตฺถํ อาโรเจสุํ. น ภิกฺขเว ทีฆา นขา ธาเรตพฺพา, โย ธาเรยฺย, อาปตฺติ ทุกฺกฏสฺสาติ.}

\prose{เตน โข ปน สมเยน ภิกฺขู นเขนปิ นขํ ฉินฺทนฺติ, มุเขนปิ นขํ ฉินฺทนฺติ, กุฏฺเฏปิ ฆํสนฺติ, องฺคุลิโย ทุกฺขา โหนฺติ. ภควโต เอตมตฺถํ อาโรเจสุํ. อนุชานามิ ภิกฺขเว นขจฺเฉทนนฺติ. สโลหิตํ นขํ ฉินฺทนฺติ, องฺคุลิโย ทุกฺขา โหนฺติ ฯเปฯ. อนุชานามิ ภิกฺขเว มํสปฺปมาเณน นขํ ฉินฺทิตุนฺติ.}

\prose{\csromanfolio{274}เตน โข ปน สมเยน ฉพฺพคฺคิยา ภิกฺขู วีสติมฏฺฐํ\csromansharedfootnote{96e5a1088cdca8f1}{วีสติมฏฺฏํ (สี)} การาเปนฺติ. มนุสฺสา อุชฺฌายนฺติ ขิยฺยนฺติ วิปาเจนฺติ “เสยฺยถาปิ คิหี กาม โภคิโน”ติ. ภควโต เอตมตฺถํ อาโรเจสุํ. น ภิกฺขเว วีสติมฏฺฐํ การาเปตพฺพํ, โย การาเปยฺย, อาปตฺติ ทุกฺกฏสฺส. อนุชานามิ ภิกฺขเว มลมตฺตํ อปกฑฺฒิตุนฺติ.}

\proseitem{257}{เตน โข ปน สมเยน ภิกฺขูนํ เกสา ทีฆา โหนฺติ. ภควโต เอตมตฺถํ อาโรเจสุํ. อุสฺสหนฺติ ปน ภิกฺขเว ภิกฺขู อญฺญมญฺญํ เกเส โอโรเปตุนฺติ. อุสฺสหนฺติ ภควาติ. อถ โข ภควา เอตสฺมิํ นิทาเน เอตสฺมิํ ปกรเณ ภิกฺขุสํฆํ สนฺนิปาตาเปตฺวา ฯเปฯ ภิกฺขู อามนฺเตสิ “อนุชานามิ ภิกฺขเว ขุรํ ขุรสิลํ ขุรสิปาฏิกํ นมตกํ สพฺพํ ขุรภณฺฑนฺ”ติ.}

\prose{เตน โข ปน สมเยน ฉพฺพคฺคิยา ภิกฺขู มสฺสุํ กปฺปาเปนฺติ ฯเปฯ มสฺสุํ วฑฺฒาเปนฺติ. โคโลมิกํ การาเปนฺติ. จตุรสฺสกํ การาเปนฺติ. ปริมุขํ การาเปนฺติ. อฑฺฒทุกํ\csromansharedfootnote{015470f29cb38d89}{อฑฺฒุรกํ (สี), อฑฺฒรุกํ (สฺยา)} การาเปนฺติ. ทาฐิกํ ฐเปนฺติ. สมฺพาเธ โลมํ สํหราเปนฺติ. มนุสฺสา อุชฺฌายนฺติ ขิยฺยนฺติ วิปาเจนฺติ “เสยฺยถาปิ คิหี กามโภคิโน”ติ. ภควโต เอตมตฺถํ อาโรเจสุํ. น ภิกฺขเว มสฺสุ กปฺปาเปตพฺพํ ฯเปฯ น มสฺสุ วฑฺฒาเปตพฺพํ. น โคโลมิกํ การาเปตพฺพํ. น จตุรสฺสกํ การาเปตพฺพํ. น ปริมุขํ การาเปตพฺพํ. น อฑฺฒทุกํ การาเปตพฺพํ. น ทาฐิกา ฐเปตพฺพา. น สมฺพาเธ โลมํ สํหราเป\-ตพฺพํ, โย สํหราเปยฺย, อาปตฺติ ทุกฺกฏสฺสาติ.}

\prose{เตน โข ปน สมเยน อญฺญตรสฺส ภิกฺขุโน สมฺพาเธ วโณ โหติ, เภสชฺชํ น สนฺติฏฺฐติ. ภควโต เอตมตฺถํ อาโรเจสุํ. อนุชานามิ ภิกฺขเว อาพาธปฺปจฺจยา สมฺพาเธ โลมํ สํหราเปตุนฺติ.}

\prose{เตน โข ปน สมเยน ฉพฺพคฺคิยา ภิกฺขู กตฺตริกาย เกเส เฉทาเปนฺติ. มนุสฺสา อุชฺฌายนฺติ ขิยฺยนฺติ วิปาเจนฺติ “เสยฺยถาปิ คิหี กามโภคิโน”ติ. ภควโต เอตมตฺถํ อาโรเจสุํ. น ภิกฺขเว กตฺตริกาย เกสา เฉทาเปตพฺพา, โย เฉทาเปยฺย, อาปตฺติ ทุกฺกฏสฺสาติ.}

\prose{\csromanfolio{275}เตน โข ปน สมเยน อญฺญตรสฺส ภิกฺขุโน สีเส วโณ โหติ, น สกฺโกติ ขุเรน เกเส โอโรเปตุํ. ภควโต เอตมตฺถํ อาโรเจสุํ. อนุชานามิ ภิกฺขเว อาพาธปฺปจฺจยา กตฺตริกาย เกเส เฉทาเปตุนฺติ.}

\prose{เตน โข ปน สมเยน ภิกฺขู ทีฆานิ นาสิกาโลมานิ ธาเรนฺติ. มนุสฺสา อุชฺฌายนฺติ ขิยฺยนฺติ วิปาเจนฺติ “เสยฺยถาปิ ปิสาจิลฺลิกา”ติ. ภควโต เอตมตฺถํ อาโรเจสุํ. น ภิกฺขเว ทีฆํ นาสิกาโลมํ ธาเรตพฺพํ, โย ธาเรยฺย, อาปตฺติ ทุกฺกฏสฺสาติ.}

\prose{เตน โข ปน สมเยน ภิกฺขู สกฺขริกายปิ มธุสิตฺถ\-เกนปิ นาสิกาโลมํ คาหาเปนฺติ, นาสิกา ทุกฺขา โหนฺติ. ภควโต เอตมตฺถํ อาโรเจสุํ. อนุชานามิ ภิกฺขเว สณฺฑาสนฺติ.}

\prose{เตน โข ปน สมเยน ฉพฺพคฺคิยา ภิกฺขู ปลิตํ คาหาเปนฺติ. มนุสฺสา อุชฺฌายนฺติ ขิยฺยนฺติ วิปาเจนฺติ “เสยฺยถาปิ คิหี กามโภคิโน”ติ. ภควโต เอตมตฺถํ อาโรเจสุํ. น ภิกฺขเว ปลิตํ คาหาเปตพฺพํ, โย คาหาเปยฺย, อาปตฺติ ทุกฺกฏสฺส.}

\proseitem{276}{เตน โข ปน สมเยน อญฺญตรสฺส ภิกฺขุโน กณฺณคูถเกหิ กณฺณา ถกิตา โหนฺติ. ภควโต เอตมตฺถํ อาโรเจสุํ. อนุชานามิ ภิกฺขเว กณฺณมล\-หรณินฺติ.}

\prose{เตน โข ปน สมเยน ฉพฺพคฺคิยา ภิกฺขู อุจฺจาวจา กณฺณมล\-หรณิโย ธาเรนฺติ โสวณฺณมยํ รูปิยมยํ. มนุสฺสา อุชฺฌายนฺติ ขิยฺยนฺติ วิปาเจนฺติ “เสยฺยถาปิ คิหี กามโภคิโน”ติ. ภควโต เอตมตฺถํ อาโรเจสุํ. น ภิกฺขเว อุจฺจาวจา กณฺณมล\-หรณิโย ธาเรตพฺพา, โย ธาเรยฺย, อาปตฺติ ทุกฺกฏสฺสาติ. อนุชานามิ ภิกฺขเว อฏฺฐิมยํ ทนฺตมยํ วิสาณมยํ นฬมยํ เวฬุมยํ กฏฺฐมยํ ชตุมยํ ผลมยํ โลหมยํ สงฺข\-นาภิ\-มยนฺติ.}

\proseitem{277}{เตน โข ปน สมเยน ฉพฺพคฺคิยา ภิกฺขู พหุํ โลหภณฺฑํ กํสภณฺฑํ สนฺนิจยํ กโรนฺติ. มนุสฺสา วิหารจาริกํ อาหิณฺฑนฺตา ปสฺสิตฺวา อุชฺฌายนฺติ ขิยฺยนฺติ วิปาเจนฺติ “กถํ หิ นาม สมณา สกฺยปุตฺติยา พหุํ โลหภณฺฑํ กํสภณฺฑํ สนฺนิจยํ กริสฺสนฺติ \csromanfolio{276}เสยฺยถาปิ กํสปตฺถริกา”ติ. ภควโต เอตมตฺถํ อาโรเจสุํ. น ภิกฺขเว พหุํ โลหภณฺฑํ กํสภณฺฑํ สนฺนิจโย กาตพฺโพ, โย กเรยฺย, อาปตฺติ ทุกฺกฏสฺสาติ.}

\prose{เตน โข ปน สมเยน ภิกฺขู อญฺชนิมฺปิ อญฺชนิ\-สลากมฺปิ กณฺณมล\-หรณิมฺปิ พนฺธน\-มตฺตมฺปิ กุกฺกุจฺจายนฺติ. ภควโต เอตมตฺถํ อาโรเจสุํ. อนุชานามิ ภิกฺขเว อญฺชนิํ อญฺชนิสลากํ กณฺณมล\-หรณิํ พนฺธน\-มตฺตนฺติ.}

\prose{เตน โข ปน สมเยน ฉพฺพคฺคิยา ภิกฺขู สํฆาฏิปลฺลตฺถิกาย นิสีทนฺติ, สํฆาฏิยา ปตฺตา ลุชฺชนฺติ. ภควโต เอตมตฺถํ อาโรเจสุํ. น ภิกฺขเว สํฆาฏิปลฺลตฺถิกาย นิสีทิตพฺพํ, โย นิสีเทยฺย, อาปตฺติ ทุกฺกฏสฺสาติ.}

\prose{เตน โข ปน สมเยน อญฺญตโร ภิกฺขุ คิลาโน โหติ, ตสฺส วินา อาโยเคน\csromansharedfootnote{13b29f02c9b61063}{อาโยคา (สี, สฺยา)} น ผาสุ โหติ. ภควโต เอตมตฺถํ อาโรเจสุํ. อนุชานามิ ภิกฺขเว อาโยคนฺติ. อถ โข ภิกฺขูนํ เอตทโหสิ “กถํ นุ โข อาโยโค กาตพฺโพ”ติ. ภควโต เอตมตฺถํ อาโรเจสุํ. อนุชานามิ ภิกฺขเว ตนฺตกํ เวมํ กวฏํ\csromansharedfootnote{713f79ac14757923}{เวมกํ วฏฺฏํ (สี), เวมกํ วฏํ (สฺยา)} สลากํ สพฺพํ ตนฺต\-ภณฺฑกนฺติ.}

\proseitem{278}{เตน โข ปน สมเยน อญฺญตโร ภิกฺขุ อกายพนฺธโน คามํ ปิณฺฑาย ปาวิสิ, ตสฺส รถิกาย อนฺตรวาสโก ปภสฺสิตฺถ. มนุสฺสา อุกฺกุฏฺฐิมกํสุ. โส ภิกฺขุ มงฺกุ อโหสิ. อถ โข โส ภิกฺขุ อารามํ คนฺตา ภิกฺขูนํ เอตมตฺถํ อาโรเจสิ. ภิกฺขู ภควโต เอตมตฺถํ อาโรเจสุํ. น ภิกฺขเว อกายพนฺธเนน คาโม ปวิสิตพฺโพ, โย ปวิเสยฺย, อาปตฺติ ทุกฺกฏสฺส. อนุชานามิ ภิกฺขเว กาย\-พนฺธนนฺติ.}

\prose{เตน โข ปน สมเยน ฉพฺพคฺคิยา ภิกฺขู อุจฺจาวจานิ กายพนฺธนานิ ธาเรนฺติ กลาพุกํ เทฑฺฑุภกํ มุรชํ มทฺทวีณํ. มนุสฺสา อุชฺฌายนฺติ ขิยฺยนฺติ วิปาเจนฺติ “เสยฺยถาปิ คิหี กามโภคิโน”ติ. ภควโต เอตมตฺถํ อาโรเจสุํ. น ภิกฺขเว อุจฺจาวจานิ กายพนฺธนานิ ธาเรตพฺพานิ กลาพุกํ เทฑฺฑุภกํ มุรชํ มทฺทวีณํ, โย ธาเรยฺย, อาปตฺติ ทุกฺกฏสฺส. อนุชานามิ ภิกฺขเว เทฺว กายพนฺธนานิ ปฏฺฏิกํ สูกรนฺตกนฺติ. กาย\-พนฺธนสฺส \csromanfolio{277}ทสา ชีรนฺติ ฯเปฯ. อนุชานามิ ภิกฺขเว มุรชํ มทฺทวีณนฺติ. กาย\-พนฺธนสฺส อนฺโต ชีรติ ฯเปฯ. อนุชานามิ ภิกฺขเว โสภณํ คุณกนฺติ. กาย\-พนฺธนสฺส ปวนนฺโต ชีรติ ฯเปฯ. อนุชานามิ ภิกฺขเว วิธนฺติ\csromansharedfootnote{b43d3431b8ebdf56}{วีถนฺติ (สี, สฺยา)}.}

\prose{เตน โข ปน สมเยน ฉพฺพคฺคิยา ภิกฺขู อุจฺจาวเจ วิเธ ธาเรนฺติ โสวณฺณมยํ รูปิยมยํ. มนุสฺสา อุชฺฌายนฺติ ขิยฺยนฺติ วิปาเจนฺติ “เสยฺยถาปิ คิหี กามโภคิโน”ติ. ภควโต เอตมตฺถํ อาโรเจสุํ. น ภิกฺขเว อุจฺจาวจา วิธา ธาเรตพฺพา, โย ธาเรยฺย, อาปตฺติ ทุกฺกฏสฺส. อนุชานามิ ภิกฺขเว อฏฺฐิมยํ ฯเปฯ สงฺข\-นาภิมยํ สุตฺตมยนฺติ.}

\proseitem{279}{เตน โข ปน สมเยน อายสฺมา อานนฺโท ลหุกา สํฆาฏิโย ปารุปิตฺวา คามํ ปิณฺฑาย ปาวิสิ, วาต\-มณฺฑลิกาย สํฆาฏิโย อุกฺขิปิยิํสุ. อถ โข อายสฺมา อานนฺโท อารามํ คนฺตา ภิกฺขูนํ เอตมตฺถํ อาโรเจสิ. ภิกฺขู ภควโต เอตมตฺถํ อาโรเจสุํ. อนุชานามิ ภิกฺขเว คณฺฐิกํ ปาสกนฺติ.}

\prose{เตน โข ปน สมเยน ฉพฺพคฺคิยา ภิกฺขู อุจฺจาวจา คณฺฐิกาโย ธาเรนฺติ โสวณฺณมยํ รูปิยมยํ. มนุสฺสา อุชฺฌายนฺติ ขิยฺยนฺติ วิปาเจนฺติ “เสยฺยถาปิ คิหี กามโภคิโน”ติ. ภควโต เอตมตฺถํ อาโรเจสุํ. น ภิกฺขเว อุจฺจาวจา คณฺฐิกา ธาเรตพฺพา, โย ธาเรยฺย, อาปตฺติ ทุกฺกฏสฺส. อนุชานามิ ภิกฺขเว อฏฺฐิมยํ ทนฺตมยํ วิสาณมยํ นฬมยํ เวฬุมยํ กฏฺฐมยํ ชตุมยํ ผลมยํ โลหมยํ สงฺข\-นาภิมยํ สุตฺตมยนฺติ.}

\prose{เตน โข ปน สมเยน ภิกฺขู คณฺฐิกมฺปิ ปาสกมฺปิ จีวเร อปฺเปนฺติ. จีวรํ ชีรติ. ภควโต เอตมตฺถํ อาโรเจสุํ. อนุชานามิ ภิกฺขเว คณฺฐิก\-ผลกํ ปาสก\-ผลกนฺติ. คณฺฐิก\-ผลกมฺปิ ปาสก\-ผลกมฺปิ อนฺเต อปฺเปนฺติ, โกฏฺโฏ\csromansharedfootnote{2d1873491cb0fb2e}{โกโณ (สี, สฺยา)} วิวริยติ. ภควโต เอตมตฺถํ อาโรเจสุํ. อนุชานามิ ภิกฺขเว คณฺฐิก\-ผลกํ อนฺเต อปฺเปตุํ, ปาสกผลกํ สตฺตงฺคุลํ วา อฏฺฐงฺคุลํ วา โอคาเหตฺวา อปฺเปตุนฺติ.}

\proseitem{280}{เตน โข ปน สมเยน ฉพฺพคฺคิยา ภิกฺขู คิหินิวตฺถํ นิวาเสนฺติ หตฺถิโสณฺฑกํ มจฺฉวาฬกํ จตุกณฺณกํ ตาลวณฺฏกํ สตวลิกํ. มนุสฺสา \csromanfolio{278}อุชฺฌายนฺติ ขิยฺยนฺติ วิปาเจนฺติ “เสยฺยถาปิ คิหี กามโภคิโน”ติ. ภควโต เอตมตฺถํ อาโรเจสุํ. น ภิกฺขเว คิหินิวตฺถํ นิวาเสตพฺพํ หตฺถิโสณฺฑกํ มจฺฉวาฬกํ จตุกณฺณกํ ตาลวณฺฏกํ สตวลิกํ, โย นิวาเสยฺย, อาปตฺติ ทุกฺกฏสฺสาติ.}

\prose{เตน โข ปน สมเยน ฉพฺพคฺคิยา ภิกฺขู คิหิปารุตํ ปารุปนฺติ. มนุสฺสา อุชฺฌายนฺติ ขิยฺยนฺติ วิปาเจนฺติ “เสยฺยถาปิ คิหี กามโภคิโน”ติ. ภควโต เอตมตฺถํ อาโรเจสุํ. น ภิกฺขเว คิหิปารุตํ ปารุปิตพฺพํ, โย ปารุเปยฺย, อาปตฺติ ทุกฺกฏสฺสาติ.}

\prose{เตน โข ปน สมเยน ฉพฺพคฺคิยา ภิกฺขู สํเวลฺลิยํ นิวาเสนฺติ. มนุสฺสา อุชฺฌายนฺติ ขิยฺยนฺติ วิปาเจนฺติ “เสยฺยถาปิ รญฺโญ มุณฺฑวฏฺฏี”ติ. ภควโต เอตมตฺถํ อาโรเจสุํ. น ภิกฺขเว สํเวลฺลิยํ นิวาเสตพฺพํ, โย นิวาเสยฺย, อาปตฺติ ทุกฺกฏสฺสาติ.}

\proseitem{281}{เตน โข ปน สมเยน ฉพฺพคฺคิยา ภิกฺขู อุภโตกาชํ หรนฺติ. มนุสฺสา อุชฺฌายนฺติ ขิยฺยนฺติ วิปาเจนฺติ “เสยฺยถาปิ รญฺโญ มุณฺฑวฏฺฏี”ติ. ภควโต เอตมตฺถํ อาโรเจสุํ. น ภิกฺขเว อุภโตกาชํ หริตพฺพํ, โย หเรยฺย, อาปตฺติ ทุกฺกฏสฺส. อนุชานามิ ภิกฺขเว เอกโตกาชํ อนฺตรากาชํ สีสภารํ ขนฺธภารํ กฏิภารํ โอลมฺพกนฺติ.}

\proseitem{282}{เตน โข ปน สมเยน ภิกฺขู ทนฺตกฏฺฐํ น ขาทนฺติ, มุขํ ทุคฺคนฺธํ โหติ. ภควโต เอตมตฺถํ อาโรเจสุํ.}

\prose{ปญฺจิเม ภิกฺขเว อาทีนวา ทนฺตกฏฺฐสฺส อขาทเน. อจกฺขุสฺสํ, มุขํ ทุคฺคนฺธํ โหติ, รสหรณิโย น วิสุชฺฌนฺติ, ปิตฺตํ เสมฺหํ ภตฺตํ ปริโยนนฺธติ, ภตฺตมสฺส นจฺฉาเทติ. อิเม โข ภิกฺขเว ปญฺจ อาทีนวา ทนฺตกฏฺฐสฺส อขาทเน.}

\prose{ปญฺจิเม ภิกฺขเว อานิสํสา ทนฺตกฏฺฐสฺส ขาทเน. จกฺขุสฺสํ, มุขํ น ทุคฺคนฺธํ โหติ, รสหรณิโย วิสุชฺฌนฺติ, ปิตฺตํ เสมฺหํ ภตฺตํ น ปริโยนนฺธติ, ภตฺตมสฺส ฉาเทติ. อิเม โข ภิกฺขเว ปญฺจ อานิสํสา ทนฺตกฏฺฐสฺส ขาทเน. อนุชานามิ ภิกฺขเว ทนฺตกฏฺฐนฺติ.}

\prose{เตน โข ปน สมเยน ฉพฺพคฺคิยา ภิกฺขู ทีฆานิ ทนฺตกฏฺฐานิ ขาทนฺติ, เตเหว สามเณรํ อาโกเฏนฺติ. ภควโต เอตมตฺถํ อาโรเจสุํ. \csromanfolio{279}น ภิกฺขเว ทีฆํ ทนฺตกฏฺฐํ ขาทิตพฺพํ, โย ขาเทยฺย, อาปตฺติ ทุกฺกฏสฺส. อนุชานามิ ภิกฺขเว อฏฺฐงฺคุล\-ปรมํ ทนฺตกฏฺฐํ, น จ เตน สามเณโร อาโกเฏตพฺโพ, โย อาโกเฏยฺย, อาปตฺติ ทุกฺกฏสฺสาติ.}

\prose{เตน โข ปน สมเยน อญฺญตรสฺส ภิกฺขุโน อติมฏาหกํ ทนฺตกฏฺฐํ ขาทนฺตสฺส กณฺเฐ วิลคฺคํ โหติ. ภควโต เอตมตฺถํ อาโรเจสุํ. น ภิกฺขเว อติมฏาหกํ ทนฺตกฏฺฐํ ขาทิตพฺพํ, โย ขาเทยฺย, อาปตฺติ ทุกฺกฏสฺส. อนุชานามิ ภิกฺขเว จตุรงฺคุล\-ปจฺฉิมํ\csromansharedfootnote{79fb646633057195}{จตุรงฺคุลํ ปจฺฉิมํ (ก)} ทนฺตกฏฺฐนฺติ.}

\proseitem{283}{เตน โข ปน สมเยน ฉพฺพคฺคิยา ภิกฺขู ทายํ อาลิมฺเปนฺติ. มนุสฺสา อุชฺฌายนฺติ ขิยฺยนฺติ วิปาเจนฺติ “เสยฺยถาปิ ทวฑาหกา”ติ. ภควโต เอตมตฺถํ อาโรเจสุํ. น ภิกฺขเว ทาโย อาลิมฺปิตพฺโพ, โย อาลิมฺเปยฺย, อาปตฺติ ทุกฺกฏสฺสาติ.}

\prose{เตน โข ปน สมเยน วิหารา ติณคหนา โหนฺติ, ทวฑาเห ฑยฺหมาเน วิหารา ฑยฺหนฺติ, ภิกฺขู กุกฺกุจฺจายนฺติ ปฏคฺคิํ ทาตุํ ปริตฺตํ กาตุํ. ภควโต เอตมตฺถํ อาโรเจสุํ. อนุชานามิ ภิกฺขเว ทวฑาเห ฑยฺหมาเน ปฏคฺคิํ ทาตุํ ปริตฺตํ กาตุนฺติ.}

\proseitem{284}{เตน โข ปน สมเยน ฉพฺพคฺคิยา ภิกฺขู รุกฺขํ อภิรุหนฺติ, รุกฺขา รุกฺขํ สงฺกมนฺติ. มนุสฺสา อุชฺฌายนฺติ ขิยฺยนฺติ วิปาเจนฺติ “เสยฺยถาปิ มกฺกฏา”ติ. ภควโต เอตมตฺถํ อาโรเจสุํ. น ภิกฺขเว รุกฺโข อภิรุหิตพฺโพ, โย อภิรุเหยฺย, อาปตฺติ ทุกฺกฏสฺสาติ.}

\prose{เตน โข ปน สมเยน อญฺญตรสฺส ภิกฺขุโน โกสเลสุ ชนปเท สาวตฺถิํ คจฺฉนฺตสฺส อนฺตรามคฺเค หตฺถี ปริยุฏฺฐาติ. อถ โข โส ภิกฺขุ รุกฺขมูลํ อุปธาวิตฺวา กุกฺกุจฺจายนฺโต รุกฺขํ น อภิรุหิ, โส หตฺถี อญฺเญน อคมาสิ. อถ โข โส ภิกฺขุ สาวตฺถิํ คนฺตฺวา ภิกฺขูนํ เอตมตฺถํ อาโรเจสิ. ภิกฺขู ภควโต เอตมตฺถํ อาโรเจสุํ. อนุชานามิ ภิกฺขเว สติ กรณีเย โปริสํ รุกฺขํ อภิรุหิตุํ, อาปทาสุ ยาวทตฺถนฺติ.}

\proseitem{285}{\csromanfolio{280}เตน โข ปน สมเยน ยเมฬเกกุฏา นาม\csromansharedfootnote{7904bb4790a117e8}{ยเมฬุเตกุลา นาม (สี), เมฏฺฐโกกุฏฺฐา นาม (สฺยา)} ภิกฺขู เทฺว ภาติกา โหนฺติ พฺราหฺมณชาติกา กลฺยาณวาจา กลฺยาณ\-วากฺกรณา, เต เยน ภควา เตนุ\-ปสงฺกมิํสุ, อุปสงฺกมิตฺวา ภควนฺตํ อภิวาเทตฺวา เอกมนฺตํ นิสีทิํสุ, เอกมนฺตํ นิสินฺนา โข เต ภิกฺขู ภควนฺตํ เอตทโวจุํ “เอตรหิ ภนฺเต ภิกฺขู นานานามา นานาโคตฺตา นานาชจฺจา นานากุลา ปพฺพชิตา, เต สกาย นิรุตฺติยา พุทฺธวจนํ ทูเสนฺติ, หนฺท มยํ ภนฺเต พุทฺธวจนํ ฉนฺทโส อาโรเปมา”ติ. วิครหิ พุทฺโธ ภควา ฯเปฯ. กถํ หิ นาม ตุเมฺห โมฆปุริสา เอวํ วกฺขถ “หนฺท มยํ ภนฺเต พุทฺธวจนํ ฉนฺทโส อาโรเปมา”ติ. เนตํ โมฆปุริสา อปฺปสนฺนานํ วา ปสาทย ฯเปฯ วิครหิตฺวา ฯเปฯ ธมฺมิํ กถํ กตฺวา ภิกฺขู อามนฺเตสิ “น ภิกฺขเว พุทฺธวจนํ ฉนฺทโส อาโรเปตพฺพํ, โย อาโรเปยฺย, อาปตฺติ ทุกฺกฏสฺส. อนุชานามิ ภิกฺขเว สกาย นิรุตฺติยา พุทฺธวจนํ ปริยาปุณิตุนฺ”ติ.}

\proseitem{286}{เตน โข ปน สมเยน ฉพฺพคฺคิยา ภิกฺขู โลกายตํ ปริยาปุณนฺติ. มนุสฺสา อุชฺฌายนฺติ ขิยฺยนฺติ วิปาเจนฺติ “เสยฺยถาปิ คิหี กามโภคิโน”ติ. อสฺโสสุํ โข ภิกฺขู เตสํ มนุสฺสานํ อุชฺฌายนฺตานํ ขิยฺยนฺตานํ วิปาเจนฺตานํ. อถ โข เต ภิกฺขู ภควโต เอตมตฺถํ อาโรเจสุํ. อปิ นุ โข ภิกฺขเว โลกายเต สารทสฺสาวี อิมสฺมิํ ธมฺมวินเย วุทฺธิํ วิรุฬฺหิํ เวปุลฺลํ อาปชฺเชยฺยาติ. โนเหตํ ภนฺเต. อิมสฺมิํ วา ปน ธมฺมวินเย สารทสฺสาวี โลกายตํ ปริยาปุเณยฺยาติ. โนเหตํ ภนฺเต. น ภิกฺขเว โลกายตํ ปริยาปุณิตพฺพํ, โย ปริยาปุเณยฺย, อาปตฺติ ทุกฺกฏสฺสาติ.}

\prose{เตน โข ปน สมเยน ฉพฺพคฺคิยา ภิกฺขู โลกายตํ วาเจนฺติ. มนุสฺสา อุชฺฌายนฺติ ขิยฺยนฺติ วิปาเจนฺติ “เสยฺยถาปิ คิหี กามโภคิโน”ติ. ภควโต เอตมตฺถํ อาโรเจสุํ. น ภิกฺขเว โลกายตํ วาเจตพฺพํ, โย วาเจยฺย, อาปตฺติ ทุกฺกฏสฺสาติ.}

\proseitem{287}{เตน โข ปน สมเยน ฉพฺพคฺคิยา ภิกฺขู ติรจฺฉาน\-วิชฺชํ ปริยาปุณนฺติ. มนุสฺสา อุชฺฌายนฺติ ขิยฺยนฺติ วิปาเจนฺติ “เสยฺยถาปิ คิหี \csromanfolio{281}กามโภคิโน”ติ. ภควโต เอตมตฺถํ อาโรเจสุํ. น ภิกฺขเว ติรจฺฉาน\-วิชฺชา ปริยาปุณิตพฺพา, โย ปริยาปุเณยฺย, อาปตฺติ ทุกฺกฏสฺสาติ.}

\prose{เตน โข ปน สมเยน ฉพฺพคฺคิยา ภิกฺขู ติรจฺฉาน\-วิชฺชํ วาเจนฺติ. มนุสฺสา อุชฺฌายนฺติ ขิยฺยนฺติ วิปาเจนฺติ “เสยฺยถาปิ คิหี กามโภคิโน”ติ. ภควโต เอตมตฺถํ อาโรเจสุํ. น ภิกฺขเว ติรจฺฉาน\-วิชฺชา วาเจตพฺพา, โย วาเจยฺย, อาปตฺติ ทุกฺกฏสฺสาติ.}

\proseitem{288}{เตน โข ปน สมเยน ภควา มหติยา ปริสาย ปริวุโต ธมฺมํ เทเสนฺโต ขิปิ. ภิกฺขู “ชีวตุ ภนฺเต ภควา, ชีวตุ สุคโต”ติ อุจฺจาสทฺทํ มหาสทฺทํ อกํสุ, เตน สทฺเทน ธมฺมกถา อนฺตรา อโหสิ. อถ โข ภควา ภิกฺขู อามนฺเตสิ “อปิ นุ โข ภิกฺขเว ขิปิเต ‘ชีวา’ติ วุตฺโต\csromansharedfootnote{9920ff0b2d94a109}{วุตฺเต (ก)} ตปฺปจฺจยา ชีเวยฺย วา มเรยฺย วา”ติ. โนเหตํ ภนฺเต. น ภิกฺขเว ขิปิเต “ชีวา”ติ วตฺตพฺโพ, โย วเทยฺย, อาปตฺติ ทุกฺกฏสฺสาติ.}

\prose{เตน โข ปน สมเยน มนุสฺสา ภิกฺขูนํ ขิปิเต “ชีวถ ภนฺเต”ติ วทนฺติ, ภิกฺขู กุกฺกุจฺจายนฺตา นาลปนฺติ. มนุสฺสา อุชฺฌายนฺติ ขิยฺยนฺติ วิปาเจนฺติ “กถํ หิ นาม สมณา สกฺยปุตฺติยา ‘ชีวถ ภนฺเต’ติ วุจฺจมานา นาลปิสฺสนฺตี”ติ. ภควโต เอตมตฺถํ อาโรเจสุํ. คิหี ภิกฺขเว มงฺคลิกา. อนุชานามิ ภิกฺขเว คิหีนํ “ชีวถ ภนฺเต”ติ วุจฺจมาเนน “จิรํ ชีวา”ติ วตฺตุนฺติ.}

\proseitem{289}{เตน โข ปน สมเยน ภควา มหติยา ปริสาย ปริวุโต ธมฺมํ เทเสนฺโต นิสินฺโน โหติ, อญฺญตเรน ภิกฺขุนา ลสุณํ ขายิตํ โหติ. โส “มา ภิกฺขู พฺยาพาธิํสู”ติ เอกมนฺตํ นิสีทิ. อทฺทสา โข ภควา ตํ ภิกฺขุํ เอกมนฺตํ นิสินฺนํ, ทิสฺวาน ภิกฺขู อามนฺเตสิ “กิํ นุ โข โส ภิกฺขเว ภิกฺขุ เอกมนฺตํ นิสินฺโน”ติ. เอเตน ภนฺเต ภิกฺขุนา ลสุณํ ขายิตํ, โส “มา ภิกฺขู พฺยาพาธิํสู”ติ เอกมนฺตํ นิสินฺโนติ. อปิ นุ โข ภิกฺขเว\csromansharedfootnote{9f3954669b4798f2}{ภิกฺขเว ภิกฺขุนา (สฺยา)} ตํ ขาทิตพฺพํ, ยํ ขาทิตฺวา เอวรูปาย ธมฺมกถาย ปริพาหิโย อสฺสาติ. โนเหตํ ภนฺเต. น ภิกฺขเว ลสุณํ ขาทิตพฺพํ, โย ขาเทยฺย, อาปตฺติ ทุกฺกฏสฺสาติ.}

\prose{\csromanfolio{282}เตน โข ปน สมเยน อายสฺมโต สาริปุตฺตสฺส อุทรวาตาพาโธ โหติ. อถ โข อายสฺมา มหาโมคฺคลฺลาโน เยนายสฺมา สาริปุตฺโต เตนุปสงฺกมิ, อุปสงฺกมิตฺวา อายสฺมนฺตํ สาริปุตฺตํ เอตทโวจ “ปุพฺเพ เต อาวุโส สาริปุตฺต อุทรวาตาพาโธ เกน ผาสุ โหตี”ติ. ลสุเณน เม อาวุโสติ\csromansharedfootnote{6529da0c43ceec59}{อถ โข อายสฺมา สาริปุตฺโต อายสฺมนฺตํ มหาโมคฺคลฺลานํ เอตทโวจ “ปุพฺเพ โข เม อาวุโส โมคฺคลฺลาน อุทรวาตาพาโธ ลสุเณสุ ผาสุ โหตี”ติ (ก)}. ภควโต เอตมตฺถํ อาโรเจสุํ. อนุชานามิ ภิกฺขเว อาพาธปฺปจฺจยา ลสุณํ ขาทิตุนฺติ.}

\proseitem{290}{เตน โข ปน สมเยน ภิกฺขู อาราเม ตหํ ตหํ ปสฺสาวํ กโรนฺติ, อาราโม ทุสฺสติ. ภควโต เอตมตฺถํ อาโรเจสุํ. อนุชานามิ ภิกฺขเว เอกมนฺตํ ปสฺสาวํ กาตุนฺติ. อาราโม ทุคฺคนฺโธ โหติ ฯเปฯ. อนุชานามิ ภิกฺขเว ปสฺสาว\-กุมฺภินฺติ. ทุกฺขํ นิสินฺนา ปสฺสาวํ กโรนฺติ ฯเปฯ. อนุชานามิ ภิกฺขเว ปสฺสาว\-ปาทุกนฺติ. ปสฺสาวปาทุกา ปากฏา โหนฺติ, ภิกฺขู หิริยนฺติ ปสฺสาวํ กาตุํ ฯเปฯ. อนุชานามิ ภิกฺขเว ปริกฺขิปิตุํ ตโย ปากาเร อิฏฺฐกาปาการํ สิลาปาการํ ทารุปาการนฺติ. ปสฺสาวกุมฺภี อปารุตา ทุคฺคนฺธา โหติ ฯเปฯ. อนุชานามิ ภิกฺขเว อปิธานนฺติ.}

\proseitem{291}{เตน โข ปน สมเยน ภิกฺขู อาราเม ตหํ ตหํ วจฺจํ กโรนฺติ, อาราโม ทุสฺสติ. ภควโต เอตมตฺถํ อาโรเจสุํ. อนุชานามิ ภิกฺขเว เอกมนฺตํ วจฺจํ กาตุนฺติ. อาราโม ทุคฺคนฺโธ โหติ ฯเปฯ. อนุชานามิ ภิกฺขเว วจฺจกูปนฺติ. วจฺจกูปสฺส กูลํ ลุชฺชติ ฯเปฯ. อนุชานามิ ภิกฺขเว จินิตุํ ตโย จเย อิฏฺฐกาจยํ สิลาจยํ ทารุจยนฺติ. วจฺจกูโป นีจวตฺถุโก โหติ, อุทเกน โอตฺถริยฺยติ ฯเปฯ. อนุชานามิ ภิกฺขเว อุจฺจวตฺถุกํ กาตุนฺติ. จโย ปริปตติ ฯเปฯ. อนุชานามิ ภิกฺขเว จินิตุํ ตโย จเย อิฏฺฐกาจยํ สิลาจยํ ทารุจยนฺติ. อาโรหนฺตา วิหญฺญนฺติ ฯเปฯ. อนุชานามิ ภิกฺขเว ตโย โสปาเน อิฏฺฐกาโสปานํ สิลาโสปานํ ทารุโสปานนฺติ. อาโรหนฺตา ปริปตนฺติ ฯเปฯ. อนุชานามิ ภิกฺขเว อาลมฺพนพาหนฺติ. อนฺเต นิสินฺนา วจฺจํ กโรนฺตา ปริปตนฺติ ฯเปฯ. อนุชานามิ ภิกฺขเว สนฺถริตฺวา มชฺเฌ ฉิทฺทํ กตฺวา วจฺจํ กาตุนฺติ. ทุกฺขํ นิสินฺนา วจฺจํ กโรนฺติ ฯเปฯ. อนุชานามิ ภิกฺขเว วจฺจปาทุกนฺติ. พหิทฺธา \csromanfolio{283}ปสฺสาวํ กโรนฺติ ฯเปฯ. อนุชานามิ ภิกฺขเว ปสฺสาว\-โทณิกนฺติ. อวเลขน\-กฏฺฐํ น โหติ ฯเปฯ. อนุชานามิ ภิกฺขเว อวเลขน\-กฏฺฐนฺติ. อวเลขน\-ปิฐโร น โหติ ฯเปฯ. อนุชานามิ ภิกฺขเว อวเลขน\-ปิฐรนฺติ. วจฺจกูโป อปารุโต ทุคฺคนฺโธ โหติ ฯเปฯ. อนุชานามิ ภิกฺขเว อปิธานนฺติ. อชฺโฌกาเส วจฺจํ กโรนฺตา สีเตนปิ อุเณฺหนปิ กิลมนฺติ ฯเปฯ. อนุชานามิ ภิกฺขเว วจฺจกุฏินฺติ. วจฺจกุฏิยา กวาฏํ น โหติ ฯเปฯ. อนุชานามิ ภิกฺขเว กวาฏํ ปิฏฺฐสํฆาฏํ อุทุกฺขลิกํ อุตฺตรปาสกํ อคฺคฬวฏฺฏิํ กปิสีสกํ สูจิกํ ฆฏิกํ ตาฬจฺฉิทฺทํ อาวิญฺฉน\-จฺฉิทฺทํ อาวิญฺฉน\-รชฺชุนฺติ. วจฺจกุฏิยา ติณจุณฺณํ ปริปตติ ฯเปฯ. อนุชานามิ ภิกฺขเว โอคุมฺเผตฺวา อุลฺลิตฺตา\-วลิตฺตํ กาตุํ เสตวณฺณํ กาฬวณฺณํ เครุก\-ปริกมฺมํ มาลากมฺมํ ลตากมฺมํ มกรทนฺตกํ ปญฺจปฏิกํ จีวรวํสํ จีวรรชฺชุนฺติ.}

\proseitem{292}{เตน โข ปน สมเยน อญฺญตโร ภิกฺขุ ชราทุพฺพโล วจฺจํ กตฺวา วุฏฺฐหนฺโต ปริปตติ. ภควโต เอตมตฺถํ อาโรเจสุํ. อนุชานามิ ภิกฺขเว โอลมฺพกนฺติ. วจฺจกุฏิ อปริกฺขิตฺตา โหติ ฯเปฯ. อนุชานามิ ภิกฺขเว ปริกฺขิปิตุํ ตโย ปากาเร อิฏฺฐกาปาการํ สิลาปาการํ ทารุปาการนฺติ. โกฏฺฐโก น โหติ ฯเปฯ. (อนุชานามิ ภิกฺขเว โกฏฺฐกนฺติ. โกฏฺฐกสฺส กวาฏํ น โหติ.)\csromansharedfootnote{bcbb607e11ce8c8c}{(อนุชานามิ ภิกฺขเว โกฏฺฐกนฺติ. โกฏฺฐโก นีจวตฺถุโก โหติ ฯเปฯ. อนุชานามิ ภิกฺขเว อุจฺจวตฺถุกํ กาตุนฺติ. จโย ปริปตติ ฯเปฯ. อนุชานามิ ภิกฺขเว จินิตุํ ตโย จเย อิฏฺฐกาจยํ สิลาจยํ ทารุจยนฺติ. อาโรหนฺตา วิหญฺญนฺติ ฯเปฯ. อนุชานามิ ภิกฺขเว ตโย โสปาเน อิฏฺฐกาโสปานํ สิลาโสปานํ ทารุโสปานนฺติ. อาโรหนฺตา ปริปตนฺติ ฯเปฯ. อนุชานามิ ภิกฺขเว อาลมฺพนพาหนฺติ. โกฏฺฐกสฺส กวาฏํ น โหติ.) (สฺยา, กํ)} ฯเปฯ. อนุชานามิ ภิกฺขเว กวาฏํ ปิฏฺฐสํฆาฏํ อุทุกฺขลิกํ อุตฺตรปาสกํ อคฺคฬวฏฺฏิํ กปิสีสกํ สูจิกํ ฆฏิกํ ตาฬจฺฉิทฺทํ อาวิญฺฉน\-จฺฉิทฺทํ อาวิญฺฉน\-รชฺชุนฺติ. โกฏฺฐเก ติณจุณฺณํ ปริปตติ ฯเปฯ. อนุชานามิ ภิกฺขเว โอคุมฺเผตฺวา อุลฺลิตฺตา\-วลิตฺตํ กาตุํ เสตวณฺณํ กาฬวณฺณํ เครุก\-ปริกมฺมํ มาลากมฺมํ ลตากมฺมํ มกรทนฺตกํ ปญฺจปฏิกนฺติ. ปริเวณํ จิกฺขลฺลํ โหติ ฯเปฯ. อนุชานามิ ภิกฺขเว มรุมฺพํ ปกิริตุนฺติ. น ปริยาปุณนฺติ ฯเปฯ. อนุชานามิ ภิกฺขเว ปทรสิลํ\csromansharedfootnote{2a46b1632770b6bb}{ปฏฺฏสิลํ (ก)} นิกฺขิปิตุนฺติ. อุทกํ สนฺติฏฺฐติ ฯเปฯ. อนุชานามิ ภิกฺขเว อุทกนิทฺธมนนฺติ. อาจมนกุมฺภี น โหติ ฯเปฯ. อนุชานามิ ภิกฺขเว อาจมนกุมฺภินฺติ. \csromanfolio{284}อาจมนสราวโก น โหติ ฯเปฯ. อนุชานามิ ภิกฺขเว อาจมนสราวกนฺติ. ทุกฺขํ นิสินฺนา อาจเมนฺติ ฯเปฯ. อนุชานามิ ภิกฺขเว อาจมน ปาทุกนฺติ. อาจมนปาทุกา ปากฏา โหนฺติ, ภิกฺขู หิริยนฺติ อาจเมตุํ ฯเปฯ. อนุชานามิ ภิกฺขเว ปริกฺขิปิตุํ ตโย ปากาเร อิฏฺฐกาปาการํ สิลาปาการํ ทารุปาการนฺติ. อาจมนกุมฺภี อปารุตา โหติ, ติณจุณฺเณหิปิ ปํสุเกหิปิ โอกิริยฺยติ ฯเปฯ. อนุชานามิ ภิกฺขเว อปิธานนฺติ.}

\csromanmatikaanchor{mtk.128}
\csromantocmarkref{section}{อุทฺทานคาถา}{mtk.128}
\bookmark[dest={mtk.128},level=1]{อุทฺทานคาถา}
\proseitem{293}{เตน โข ปน สมเยน ฉพฺพคฺคิยา ภิกฺขู เอวรูปํ อนาจารํ อาจรนฺติ, มาลาวจฺฉํ โรเปนฺติปิ โรปาเปนฺติปิ, สิญฺจนฺติปิ สิญฺจาเปนฺติปิ, โอจินนฺติปิ โอจินาเปนฺติปิ, คนฺเถนฺติปิ คนฺถาเปนฺติปิ, เอกโตวณฺฏิกมาลํ กโรนฺติปิ การาเปนฺติปิ, อุภโตวณฺฏิกมาลํ กโรนฺติปิ การาเปนฺติปิ, มญฺชริกํ กโรนฺติปิ การาเปนฺติปิ, วิธูติกํ กโรนฺติปิ การาเปนฺติปิ, วฏํสกํ กโรนฺติปิ การาเปนฺติปิ, อาเวฬํ กโรนฺติปิ การาเปนฺติปิ, อุรจฺฉทํ กโรนฺติปิ การาเปนฺติปิ. เต กุลิตฺถีนํ กุลธีตานํ กุลกุมารีนํ กุลสุณฺหานํ กุลทาสีนํ เอกโตวณฺฏิกมาลํ หรนฺติปิ หราเปนฺติปิ, อุภโตวณฺฏิกมาลํ หรนฺติปิ หราเปนฺติปิ, มญฺชริกํ หรนฺติปิ หราเปนฺติปิ, วิธูติกํ หรนฺติปิ หราเปนฺติปิ, วฏํสกํ หรนฺติปิ หราเปนฺติปิ, อาเวฬํ หรนฺติปิ หาราเปนฺติปิ, อุรจฺฉทํ หรนฺติปิ หราเปนฺติปิ. เต กุลิตฺถีหิ กุลธีตาหิ กุลกุมารีหิ กุลสุณฺหาหิ กุลทาสีหิ สทฺธิํ เอกภาชเนปิ ภุญฺชนฺติ, เอกถาลเกปิ ปิวนฺติ, เอกาสเนปิ นิสีทนฺติ, เอกมญฺเจปิ ตุวฏฺเฏนฺติ, เอกตฺถรณาปิ ตุวฏฺเฏนฺติ, เอกปาวุรณาปิ ตุวฏฺเฏนฺติ, เอก\-ตฺถรณ\-ปาวุรณาปิ ตุวฏฺเฏนฺติ, วิกาเลปิ ภุญฺชนฺติ, มชฺชมฺปิ ปิวนฺติ, มาลา\-คนฺธ\-วิเลปนมฺปิ ธาเรนฺติ, นจฺจนฺติปิ คายนฺติปิ วาเทนฺติปิ ลาเสนฺติปิ, นจฺจนฺติยาปิ นจฺจนฺติ นจฺจนฺติยาปิ คายนฺติ นจฺจนฺติยาปิ วาเทนฺติ นจฺจนฺติยาปิ ลาเสนฺติ ฯเปฯ. ลาเสนฺติยาปิ นจฺจนฺติ ลาเสนฺติยาปิ คายนฺติ ลาเสนฺติยาปิ วาเทนฺติ ลาเสนฺติยาปิ ลาเสนฺติ, อฏฺฐปเทปิ กีฬนฺติ, ทสปเทปิ กีฬนฺติ, อากาเสปิ กีฬนฺติ, ปริหาร\-ปเถปิ กีฬนฺติ, สนฺติกายปิ กีฬนฺติ, ขลิกายปิ กีฬนฺติ, ฆฏิกายปิ กีฬนฺติ, สลากหตฺเถนปิ กีฬนฺติ, อกฺเขนปิ กีฬนฺติ, ปงฺคจีเรนปิ กีฬนฺติ, วงฺกเกนปิ กีฬนฺติ, โมกฺขจิกายปิ กีฬนฺติ, จิงฺคุลเกนปิ กีฬนฺติ, ปตฺตาฬฺหเกนปิ กีฬนฺติ, รถเกนปิ กีฬนฺติ, ธนุเกนปิ กีฬนฺติ, อกฺขริกายปิ กีฬนฺติ, มเนสิกายปิ กีฬนฺติ, ยถาวชฺเชนปิ กีฬนฺติ, หตฺถิสฺมิมฺปิ สิกฺขนฺติ, อสฺสสฺมิมฺปิ สิกฺขนฺติ, รถสฺมิมฺปิ สิกฺขนฺติ, \csromanfolio{285}ธนุสฺมิมฺปิ สิกฺขนฺติ, ถรุสฺมิมฺปิ สิกฺขนฺติ, หตฺถิสฺสปิ ปุรโต ธาวนฺติ, อสฺสสฺสปิ ปุรโต ธาวนฺติ, รถสฺสปิ ปุรโต ธาวนฺติปิ อาธาวนฺติปิ, อุสฺเสเฬนฺติปิ, อปฺโผเฏนฺติปิ, นิพฺพุชฺฌนฺติปิ, มุฏฺฐีหิปิ ยุชฺฌนฺติ, รงฺคมชฺเฌปิ สํฆาฏิํ ปตฺถริตฺวา นจฺจกิํ เอวํ วทนฺติ “อิธ ภคินิ นจฺจสฺสู”ติ, นลาฏิกมฺปิ เทนฺติ, วิวิธมฺปิ อนาจารํ อาจรนฺติ. ภควโต เอตมตฺถํ อาโรเจสุํ ฯเปฯ. น ภิกฺขเว วิวิธํ อนาจารํ อาจริตพฺพํ, โย อาจเรยฺย, ยถาธมฺโม กาเรตพฺโพติ.}

\prose{เตน โข ปน สมเยน อายสฺมนฺเต อุรุเวลกสฺสเป ปพฺพชิเต สํฆสฺส พหุํ โลหภณฺฑํ ทารุภณฺฑํ มตฺติกาภณฺฑํ อุปฺปนฺนํ โหติ. อถ โข ภิกฺขูนํ เอตทโหสิ “กิํ นุ โข ภควตา โลหภณฺฑํ อนุญฺญาตํ, กิํ อนนุญฺญาตํ, กิํ ทารุภณฺฑํ อนุญฺญาตํ, กิํ อนนุญฺญาตํ, กิํ มตฺติกาภณฺฑํ อนุญฺญาตํ, กิํ อนนุญฺญาตนฺ”ติ. ภควโต เอตมตฺถํ อาโรเจสุํ. อถ โข ภควา เอตสฺมิํ นิทาเน เอตสฺมิํ ปกรเณ ธมฺมิํ กถํ กตฺวา ภิกฺขู อามนฺเตสิ “อนุชานามิ ภิกฺขเว ฐเปตฺวา ปหรณิํ สพฺพํ โลหภณฺฑํ, ฐเปตฺวา อาสนฺทิํ ปลฺลงฺกํ ทารุปตฺตํ ทารุปาทุกํ สพฺพํ ทารุภณฺฑํ, ฐเปตฺวา กตกญฺจ กุมฺภการิกญฺจ สพฺพํ มตฺติกาภณฺฑนฺ”ติ.}

\vspace{\tipitakaparskip}
\csromancenter{ขุทฺทก\-วตฺถุ\-กฺขนฺธโก ปญฺจโม.}
\csromancenter{\textbf{ตสฺสุทฺทานํ}}
\setlength{\csromangathapagewidth}{0pt}
\setlength{\csromangathawakwidth}{0pt}
\csromangathameasuregroup{รุกฺเข ถมฺเภ จ กุฏฺเฏ จ, \\ วิคยฺห มลฺลโก กจฺฉุ, \\ วลฺลิกาปิ จ ปามงฺโค, \\ กฏิ โอวฏฺฏิ กายุรํ, \\ ทีเฆ โกจฺเฉ ผเณ หตฺเถ, \\ อาทาสุทปตฺตวณา, \\ ลญฺเฉนฺติ องฺคราคญฺจ, \\ จกฺขุโรคํ คิรคฺคญฺจ, \\ อมฺพเปสิสกเลหิ, \\ อุจฺจาวจา ปตฺตมูลา, \\ จิตฺรา ทุสฺสติ ทุคฺคนฺโธ, \\ ปริภณฺฑํ ติณํ โจฬํ, \\ ถวิกา อํสพทฺธญฺจ, \\ ขิเล มญฺเจ จ ปีเฐ จ, \\ ตุมฺพฆฏิฉวสีสํ, \\ วิปฺผาลิทณฺฑโสวณฺณํ, \\ กิณฺณสตฺตุ สริตญฺจ, \\ วิกณฺณํ พนฺธิวิสมํ, \\ กฬิมฺภํ โมฆสุตฺตญฺจ, \\ องฺคุเล ปฏิคฺคหญฺจ, \\ อชฺโฌกาเส นีจวตฺถุ, \\ ปริปติ ติณจุณฺณํ, \\ เสตํ กาฬกวณฺณญฺจ, \\ มาลากมฺมํ ลตากมฺมํ, \\ จีวรวํสํ รชฺชุญฺจ, \\ อุชฺฌิตฺวา ปกฺกมนฺติ จ, \\ วินิเวฐิยติ กุฏฺเฏปิ, \\ ถวิกา พนฺธสุตฺตญฺจ, \\ อุปาหนตฺถวิกญฺจ, \\ อุทกากปฺปิยํ มคฺเค, \\ ธมฺมกรณํ เทฺว ภิกฺขู, \\ ทณฺฑํ โอตฺถรกํ ตตฺถ, \\ มกเสหิ ปณีเตน, \\ จงฺกมนชนฺตาฆรํ, \\ ตโย จเย วิหญฺญนฺติ, \\ อชฺโฌกาเส ติณจุณฺณํ, \\ เสตกํ กาฬวณฺณญฺจ, \\ มาลากมฺมํ ลตากมฺมํ, \\ วํสํ จีวรรชฺชุญฺจ, \\ จโย โสปานพาหญฺจ, \\ อุทุกฺขลุตฺตรปาสกํ, \\ สูจิฆฏิตาฬจฺฉิทฺทํ, \\ มณฺฑลํ ธูมเนตฺตญฺจ, \\ โทณิ ทุคฺคนฺธา ทหติ, \\ น เสเทติ จ จิกฺขลฺลํ, \\ ปีฐญฺจ โกฏฺฐเก กมฺมํ, \\ นคฺคา ฉมาย วสฺสนฺเต, \\ อุทปานํ ลุชฺชติ นีจํ, \\ ตุลํ กฏกฏํ จกฺกํ, \\ โลหทารุจมฺมขณฺฑํ, \\ โทณิจนฺทนิ ปาการํ, \\ สีติคตํ โปกฺขรณิํ, \\ จาตุมาสํ สยนฺติ จ, \\ อาสิตฺตกํ มโฬริกํ, \\ วฑฺโฒ โพธิ น อกฺกมิ, \\ สกฺขรํ กถลญฺเจว, \\ วิธูปนํ ตาลวณฺฏํ, \\ ฉตฺตํ วินา จ อาราเม, \\ โรมสิตฺถา นขา ทีฆา, \\ สโลหิตํ ปมาณญฺจ, \\ ขุรํ สิลํ สิปาฏิกํ, \\ มสฺสุํ กปฺเปนฺติ วฑฺเฒนฺติ, \\ ปริมุขํ อฑฺฒทุกญฺจ, \\ อาพาธา กตฺตริวโณ, \\ ปลิตํ ถกิตํ อุจฺจา, \\ ปลฺลตฺถิกญฺจ อาโยโค, \\ กลาพุกํ เทฑฺฑุภกํ, \\ ปฏฺฏิกา สูกรนฺตญฺจ, \\ อนฺโต โสภํ คุณกญฺจ, \\ คณฺฐิกา อุจฺจาวจญฺจ, \\ คิหิวตฺถํ หตฺถิโสณฺฑํ, \\ ตาลวณฺฏํ สตวลิ, \\ สํเวลฺลิ อุภโตกาชํ, \\ กณฺเฐ วิลคฺคํ ทายญฺจ, \\ ยเมเฬ โลกายตกํ, \\ ติรจฺฉานกถา วิชฺฌา, \\ วาตาพาโธ ทุสฺสติ จ, \\ หิริยนฺติ ปารุทุคฺคนฺโธ, \\ ทุคฺคนฺโธ กูปํ ลุชฺชนฺติ, \\ โสปานาลมฺพนพาหา, \\ พหิทฺธา โทณิ กฏฺฐญฺจ, \\ วจฺจกุฏิํ กวาฏญฺจ, \\ อุทุกฺขลุตฺตรปาโส, \\ สูจิฆฏิตาฬจฺฉิทฺทํ, \\ รชฺชุ อุลฺลิตฺตาวลิตฺตํ, \\ มาลากมฺมํ ลตากมฺมํ, \\ จีวรวํสํ รชฺชุญฺจ, \\ โกฏฺฐเก จาปิ ตเถว, \\ สนฺติฏฺฐติ นิทฺธมนํ, \\ ทุกฺขํ หิริ อปิธานํ, \\ โลหภณฺฑํ อนุญฺญาสิ, \\ ฐปยิตฺวาสนฺทิปลฺลงฺกํ, \\ สพฺพํ ทารุมยํ ภณฺฑํ, \\ กตกํ กุมฺภการญฺจ, \\ สพฺพมฺปิ มตฺติกาภณฺฑํ, \\ ยสฺส วตฺถุสฺส นิทฺเทโส, \\ ตมฺปิ สํขิตฺตํ อุทฺทาเน, \\ เอวํ ทสสตา วตฺถู, \\ สทฺธมฺมฏฺฐิติโก เจว, \\ สุสิกฺขิโต วินยธโร, \\ ปทีปกรโณ ธีโร,}{\csromangathabat{รุกฺเข ถมฺเภ จ กุฏฺเฏ จ,}{อฏฺฏาเน คนฺธสุตฺติยา.} \\ \csromangathabat{วิคยฺห มลฺลโก กจฺฉุ,}{ชรา จ ปุถุปาณิกา.} \\ \csromangathabat{วลฺลิกาปิ จ ปามงฺโค,}{กณฺฐสุตฺตํ นธารเย.} \\ \csromangathabat{กฏิ โอวฏฺฏิ กายุรํ,}{หตฺถาภรณมุทฺทิกา.} \\ \csromangathabat{ทีเฆ โกจฺเฉ ผเณ หตฺเถ,}{สิตฺถา อุทกเตลเก.} \\ \csromangathabat{อาทาสุทปตฺตวณา,}{อาเลโปมฺมทฺทจุณฺณนา.} \\ \csromangathabat{ลญฺเฉนฺติ องฺคราคญฺจ,}{มุขราคํ ตทูภยํ.} \\ \csromangathabat{จกฺขุโรคํ คิรคฺคญฺจ,}{อายตํ สรพาหิรํ.} \\ \csromangathabat{อมฺพเปสิสกเลหิ,}{อหิจฺฉินฺทิ จ จนฺทนํ.} \\ \csromangathabat{อุจฺจาวจา ปตฺตมูลา,}{สุวณฺโณ พหลา วลี.} \\ \csromangathabat{จิตฺรา ทุสฺสติ ทุคฺคนฺโธ,}{อุเณฺห ภิชฺชิํสุ มิฑฺฒิยา.} \\ \csromangathabat{ปริภณฺฑํ ติณํ โจฬํ,}{มาฬํ กุณฺโฑลิกาย จ.} \\ \csromangathabat{ถวิกา อํสพทฺธญฺจ,}{ตถา พนฺธนสุตฺตกา.} \\ \csromangathabat{ขิเล มญฺเจ จ ปีเฐ จ,}{องฺเก ฉตฺเต ปณามนา.} \\ \csromangathabat{ตุมฺพฆฏิฉวสีสํ,}{จลกานิ ปฏิคฺคโห.} \\ \csromangathabat{วิปฺผาลิทณฺฑโสวณฺณํ,}{ปตฺเต เปสิ จ นาฬิกา.} \\ \csromangathabat{กิณฺณสตฺตุ สริตญฺจ,}{มธุสิตฺถํ สิปาฏิกํ.} \\ \csromangathabat{วิกณฺณํ พนฺธิวิสมํ,}{ฉมาชิรปโหติ จ.} \\ \csromangathabat{กฬิมฺภํ โมฆสุตฺตญฺจ,}{อโธตลฺลํ อุปาหนา.} \\ \csromangathabat{องฺคุเล ปฏิคฺคหญฺจ,}{วิตฺถกํ ถวิกพทฺธกา.} \\ \csromangathabat{อชฺโฌกาเส นีจวตฺถุ,}{จโย จาปิ วิหญฺญเร.} \\ \csromangathabat{ปริปติ ติณจุณฺณํ,}{อุลฺลิตฺตอวลิตฺตกํ.} \\ \csromangathabat{เสตํ กาฬกวณฺณญฺจ,}{ปริกมฺมญฺจ เครุกํ.} \\ \csromangathabat{มาลากมฺมํ ลตากมฺมํ,}{มกรทนฺตกปาฏิกํ.} \\ \csromangathabat{จีวรวํสํ รชฺชุญฺจ,}{อนุญฺญาสิ วินายโก.} \\ \csromangathabat{อุชฺฌิตฺวา ปกฺกมนฺติ จ,}{กถินํ ปริภิชฺชติ.} \\ \csromangathabat{วินิเวฐิยติ กุฏฺเฏปิ,}{ปตฺเตนาทาย คจฺฉเร.} \\ \csromangathabat{ถวิกา พนฺธสุตฺตญฺจ,}{พนฺธิตฺวา จ อุปาหนา.} \\ \csromangathabat{อุปาหนตฺถวิกญฺจ,}{อํสพทฺธญฺจ สุตฺตกํ.} \\ \csromangathabat{อุทกากปฺปิยํ มคฺเค,}{ปริสฺสาวนโจฬกํ.} \\ \csromangathabat{ธมฺมกรณํ เทฺว ภิกฺขู,}{เวสาลิํ อคมา มุนิ.} \\ \csromangathabat{ทณฺฑํ โอตฺถรกํ ตตฺถ,}{อนุญฺญาสิ ปริสฺสาวนํ.} \\ \csromangathabat{มกเสหิ ปณีเตน,}{พหฺวาพาธา จ ชีวโก.} \\ \csromangathabat{จงฺกมนชนฺตาฆรํ,}{วิสเม นีจวตฺถุโก.} \\ \csromangathabat{ตโย จเย วิหญฺญนฺติ,}{โสปานาลมฺพเวทิตํ.} \\ \csromangathabat{อชฺโฌกาเส ติณจุณฺณํ,}{อุลฺลิตฺตอวลิตฺตกํ.} \\ \csromangathabat{เสตกํ กาฬวณฺณญฺจ,}{ปริกมฺมญฺจ เครุกํ.} \\ \csromangathabat{มาลากมฺมํ ลตากมฺมํ,}{มกรทนฺตกปาฏิกํ.} \\ \csromangathabat{วํสํ จีวรรชฺชุญฺจ,}{อุจฺจํ จ วตฺถุกํ กเร.} \\ \csromangathabat{จโย โสปานพาหญฺจ,}{กวาฏํ ปิฏฺฐสํฆาฏํ.} \\ \csromangathabat{อุทุกฺขลุตฺตรปาสกํ,}{วฏฺฏิญฺจ กปิสีสกํ.} \\ \csromangathabat{สูจิฆฏิตาฬจฺฉิทฺทํ,}{อาวิญฺฉนญฺจ รชฺชุกํ.} \\ \csromangathabat{มณฺฑลํ ธูมเนตฺตญฺจ,}{มชฺเฌ จ มุขมตฺติกา.} \\ \csromangathabat{โทณิ ทุคฺคนฺธา ทหติ,}{อุทกฏฺฐานํ สราวกํ.} \\ \csromangathabat{น เสเทติ จ จิกฺขลฺลํ,}{โธวิ นิทฺธมนํ กเร.} \\ \csromangathabat{ปีฐญฺจ โกฏฺฐเก กมฺมํ,}{มรุมฺพา สิลา นิทฺธมนํ.} \\ \csromangathabat{นคฺคา ฉมาย วสฺสนฺเต,}{ปฏิจฺฉาที ตโย ตหิํ.} \\ \csromangathabat{อุทปานํ ลุชฺชติ นีจํ,}{วลฺลิยา กายพนฺธเน.} \\ \csromangathabat{ตุลํ กฏกฏํ จกฺกํ,}{พหู ภิชฺชนฺติ ภาชนา.} \\ \csromangathabat{โลหทารุจมฺมขณฺฑํ,}{สาลาติณาปิธานิ จ.} \\ \csromangathabat{โทณิจนฺทนิ ปาการํ,}{จิกฺขลฺลํ นิทฺธมเนน จ.} \\ \csromangathabat{สีติคตํ โปกฺขรณิํ,}{ปุราณญฺจ นิลฺเลขณํ.} \\ \csromangathabat{จาตุมาสํ สยนฺติ จ,}{นมตกญฺจ\textsuperscript{1} นธิฏฺฐเห.} \\ \csromangathabat{อาสิตฺตกํ มโฬริกํ,}{ภุญฺชนฺเตกํ ตุวฏฺเฏยฺยุํ.} \\ \csromangathabat{วฑฺโฒ โพธิ น อกฺกมิ,}{ฆฏํ กตกสมฺมชฺชนิ.} \\ \csromangathabat{สกฺขรํ กถลญฺเจว,}{เผณกํ ปาทฆํสนี.} \\ \csromangathabat{วิธูปนํ ตาลวณฺฏํ,}{มกสญฺจาปิ จามรี.} \\ \csromangathabat{ฉตฺตํ วินา จ อาราเม,}{ตโย สิกฺกาย สมฺมุติ.} \\ \csromangathabat{โรมสิตฺถา นขา ทีฆา,}{ฉินฺทนฺตงฺคุลิกา ทุกฺขา.} \\ \csromangathabat{สโลหิตํ ปมาณญฺจ,}{วีสติ ทีฆเกสตา.} \\ \csromangathabat{ขุรํ สิลํ สิปาฏิกํ,}{นมตกํ ขุรภณฺฑกํ.} \\ \csromangathabat{มสฺสุํ กปฺเปนฺติ วฑฺเฒนฺติ,}{โคโลมิจตุรสฺสกํ.} \\ \csromangathabat{ปริมุขํ อฑฺฒทุกญฺจ,}{ทาฐิสมฺพาธสํหเร.} \\ \csromangathabat{อาพาธา กตฺตริวโณ,}{ทีฆํ สกฺขริกาย จ.} \\ \csromangathabat{ปลิตํ ถกิตํ อุจฺจา,}{โลหภณฺฑญฺชนี สห.} \\ \csromangathabat{ปลฺลตฺถิกญฺจ อาโยโค,}{วฏํ สลากพนฺธนํ.} \\ \csromangathabat{กลาพุกํ เทฑฺฑุภกํ,}{มุรชํ มทฺทวีณกํ.} \\ \csromangathabat{ปฏฺฏิกา สูกรนฺตญฺจ,}{ทสา มุรชเวณิกา.} \\ \csromangathabat{อนฺโต โสภํ คุณกญฺจ,}{ปวนนฺโตปิ ชีรติ.} \\ \csromangathabat{คณฺฐิกา อุจฺจาวจญฺจ,}{ผลกนฺเตปิ โอคาเห.} \\ \csromangathabat{คิหิวตฺถํ หตฺถิโสณฺฑํ,}{มจฺฉกํ จตุกณฺณกํ.} \\ \csromangathabat{ตาลวณฺฏํ สตวลิ,}{คิหิปารุตปารุปํ.} \\ \csromangathabat{สํเวลฺลิ อุภโตกาชํ,}{ทนฺตกฏฺฐํ อาโกฏเน.} \\ \csromangathabat{กณฺเฐ วิลคฺคํ ทายญฺจ,}{ปฏคฺคิ รุกฺขหตฺถินา.} \\ \csromangathabat{ยเมเฬ โลกายตกํ,}{ปริยาปุณิํสุ วาจยุํ.} \\ \csromangathabat{ติรจฺฉานกถา วิชฺฌา,}{ขิปิ มงฺคลํ ขาทิ จ.} \\ \csromangathabat{วาตาพาโธ ทุสฺสติ จ,}{ทุคฺคนฺโธ ทุกฺขปาทุกา.} \\ \csromangathabat{หิริยนฺติ ปารุทุคฺคนฺโธ,}{ตหํ ตหํ กโรนฺติ จ.} \\ \csromangathabat{ทุคฺคนฺโธ กูปํ ลุชฺชนฺติ,}{อุจฺจวตฺถุ จเยน จ.} \\ \csromangathabat{โสปานาลมฺพนพาหา,}{อนฺเต ทุกฺขญฺจ ปาทุกา.} \\ \csromangathabat{พหิทฺธา โทณิ กฏฺฐญฺจ,}{ปิฐโร จ อปารุโต.} \\ \csromangathabat{วจฺจกุฏิํ กวาฏญฺจ,}{ปิฏฺฐสํฆาฏเมว จ.} \\ \csromangathabat{อุทุกฺขลุตฺตรปาโส,}{วฏฺฏิญฺจ กปิสีสกํ.} \\ \csromangathabat{สูจิฆฏิตาฬจฺฉิทฺทํ,}{อาวิญฺฉนจฺฉิทฺทเมว จ.} \\ \csromangathabat{รชฺชุ อุลฺลิตฺตาวลิตฺตํ,}{เสตวณฺณญฺจ กาฬกํ.} \\ \csromangathabat{มาลากมฺมํ ลตากมฺมํ,}{มกรํ ปญฺจปฏิกํ.} \\ \csromangathabat{จีวรวํสํ รชฺชุญฺจ,}{ชราทุพฺพลปาการํ.} \\ \csromangathabat{โกฏฺฐเก จาปิ ตเถว,}{มรุมฺพํ ปทรสิลา.} \\ \csromangathabat{สนฺติฏฺฐติ นิทฺธมนํ,}{กุมฺภิญฺจาปิ สราวกํ.} \\ \csromangathabat{ทุกฺขํ หิริ อปิธานํ,}{อนาจารญฺจ อาจรุํ.} \\ \csromangathabat{โลหภณฺฑํ อนุญฺญาสิ,}{ฐปยิตฺวา ปหรณิํ.} \\ \csromangathabat{ฐปยิตฺวาสนฺทิปลฺลงฺกํ,}{ทารุปตฺตญฺจ ปาทุกํ.} \\ \csromangathabat{สพฺพํ ทารุมยํ ภณฺฑํ,}{อนุญฺญาสิ มหามุนิ.} \\ \csromangathabat{กตกํ กุมฺภการญฺจ,}{ฐปยิตฺวา ตถาคโต.} \\ \csromangathabat{สพฺพมฺปิ มตฺติกาภณฺฑํ,}{อนุญฺญาสิ อนุกมฺปโก.} \\ \csromangathabat{ยสฺส วตฺถุสฺส นิทฺเทโส,}{ปุริเมน ยทิ สมํ.} \\ \csromangathabat{ตมฺปิ สํขิตฺตํ อุทฺทาเน,}{นยโต ตํ วิชานิยา.} \\ \csromangathabat{เอวํ ทสสตา วตฺถู,}{วินเย ขุทฺทกวตฺถุเก.} \\ \csromangathabat{สทฺธมฺมฏฺฐิติโก เจว,}{เปสลานญฺจนุคฺคโห.} \\ \csromangathabat{สุสิกฺขิโต วินยธโร,}{หิตจิตฺโต สุเปสโล.} \\ \csromangathabat{ปทีปกรโณ ธีโร,}{ปูชารโห พหุสฺสุโตติ.}}
\csromangathasetleft{รุกฺเข ถมฺเภ จ กุฏฺเฏ จ, \\ วิคยฺห มลฺลโก กจฺฉุ, \\ วลฺลิกาปิ จ ปามงฺโค, \\ กฏิ โอวฏฺฏิ กายุรํ, \\ ทีเฆ โกจฺเฉ ผเณ หตฺเถ, \\ อาทาสุทปตฺตวณา, \\ ลญฺเฉนฺติ องฺคราคญฺจ, \\ จกฺขุโรคํ คิรคฺคญฺจ, \\ อมฺพเปสิสกเลหิ, \\ อุจฺจาวจา ปตฺตมูลา, \\ จิตฺรา ทุสฺสติ ทุคฺคนฺโธ, \\ ปริภณฺฑํ ติณํ โจฬํ, \\ ถวิกา อํสพทฺธญฺจ, \\ ขิเล มญฺเจ จ ปีเฐ จ, \\ ตุมฺพฆฏิฉวสีสํ, \\ วิปฺผาลิทณฺฑโสวณฺณํ, \\ กิณฺณสตฺตุ สริตญฺจ, \\ วิกณฺณํ พนฺธิวิสมํ, \\ กฬิมฺภํ โมฆสุตฺตญฺจ, \\ องฺคุเล ปฏิคฺคหญฺจ, \\ อชฺโฌกาเส นีจวตฺถุ, \\ ปริปติ ติณจุณฺณํ, \\ เสตํ กาฬกวณฺณญฺจ, \\ มาลากมฺมํ ลตากมฺมํ, \\ จีวรวํสํ รชฺชุญฺจ, \\ อุชฺฌิตฺวา ปกฺกมนฺติ จ, \\ วินิเวฐิยติ กุฏฺเฏปิ, \\ ถวิกา พนฺธสุตฺตญฺจ, \\ อุปาหนตฺถวิกญฺจ, \\ อุทกากปฺปิยํ มคฺเค, \\ ธมฺมกรณํ เทฺว ภิกฺขู, \\ ทณฺฑํ โอตฺถรกํ ตตฺถ, \\ มกเสหิ ปณีเตน, \\ จงฺกมนชนฺตาฆรํ, \\ ตโย จเย วิหญฺญนฺติ, \\ อชฺโฌกาเส ติณจุณฺณํ, \\ เสตกํ กาฬวณฺณญฺจ, \\ มาลากมฺมํ ลตากมฺมํ, \\ วํสํ จีวรรชฺชุญฺจ, \\ จโย โสปานพาหญฺจ, \\ อุทุกฺขลุตฺตรปาสกํ, \\ สูจิฆฏิตาฬจฺฉิทฺทํ, \\ มณฺฑลํ ธูมเนตฺตญฺจ, \\ โทณิ ทุคฺคนฺธา ทหติ, \\ น เสเทติ จ จิกฺขลฺลํ, \\ ปีฐญฺจ โกฏฺฐเก กมฺมํ, \\ นคฺคา ฉมาย วสฺสนฺเต, \\ อุทปานํ ลุชฺชติ นีจํ, \\ ตุลํ กฏกฏํ จกฺกํ, \\ โลหทารุจมฺมขณฺฑํ, \\ โทณิจนฺทนิ ปาการํ, \\ สีติคตํ โปกฺขรณิํ, \\ จาตุมาสํ สยนฺติ จ, \\ อาสิตฺตกํ มโฬริกํ, \\ วฑฺโฒ โพธิ น อกฺกมิ, \\ สกฺขรํ กถลญฺเจว, \\ วิธูปนํ ตาลวณฺฏํ, \\ ฉตฺตํ วินา จ อาราเม, \\ โรมสิตฺถา นขา ทีฆา, \\ สโลหิตํ ปมาณญฺจ, \\ ขุรํ สิลํ สิปาฏิกํ, \\ มสฺสุํ กปฺเปนฺติ วฑฺเฒนฺติ, \\ ปริมุขํ อฑฺฒทุกญฺจ, \\ อาพาธา กตฺตริวโณ, \\ ปลิตํ ถกิตํ อุจฺจา, \\ ปลฺลตฺถิกญฺจ อาโยโค, \\ กลาพุกํ เทฑฺฑุภกํ, \\ ปฏฺฏิกา สูกรนฺตญฺจ, \\ อนฺโต โสภํ คุณกญฺจ, \\ คณฺฐิกา อุจฺจาวจญฺจ, \\ คิหิวตฺถํ หตฺถิโสณฺฑํ, \\ ตาลวณฺฏํ สตวลิ, \\ สํเวลฺลิ อุภโตกาชํ, \\ กณฺเฐ วิลคฺคํ ทายญฺจ, \\ ยเมเฬ โลกายตกํ, \\ ติรจฺฉานกถา วิชฺฌา, \\ วาตาพาโธ ทุสฺสติ จ, \\ หิริยนฺติ ปารุทุคฺคนฺโธ, \\ ทุคฺคนฺโธ กูปํ ลุชฺชนฺติ, \\ โสปานาลมฺพนพาหา, \\ พหิทฺธา โทณิ กฏฺฐญฺจ, \\ วจฺจกุฏิํ กวาฏญฺจ, \\ อุทุกฺขลุตฺตรปาโส, \\ สูจิฆฏิตาฬจฺฉิทฺทํ, \\ รชฺชุ อุลฺลิตฺตาวลิตฺตํ, \\ มาลากมฺมํ ลตากมฺมํ, \\ จีวรวํสํ รชฺชุญฺจ, \\ โกฏฺฐเก จาปิ ตเถว, \\ สนฺติฏฺฐติ นิทฺธมนํ, \\ ทุกฺขํ หิริ อปิธานํ, \\ โลหภณฺฑํ อนุญฺญาสิ, \\ ฐปยิตฺวาสนฺทิปลฺลงฺกํ, \\ สพฺพํ ทารุมยํ ภณฺฑํ, \\ กตกํ กุมฺภการญฺจ, \\ สพฺพมฺปิ มตฺติกาภณฺฑํ, \\ ยสฺส วตฺถุสฺส นิทฺเทโส, \\ ตมฺปิ สํขิตฺตํ อุทฺทาเน, \\ เอวํ ทสสตา วตฺถู, \\ สทฺธมฺมฏฺฐิติโก เจว, \\ สุสิกฺขิโต วินยธโร, \\ ปทีปกรโณ ธีโร,}
\csromangathagroupwithclosermajor{\csromangathastanza{\csromangathabat{รุกฺเข ถมฺเภ จ กุฏฺเฏ จ,}{อฏฺฏาเน คนฺธสุตฺติยา.} \\ \csromangathabat{วิคยฺห มลฺลโก กจฺฉุ,}{ชรา จ ปุถุปาณิกา.}}\csromangathastanza{\csromangathabat{วลฺลิกาปิ จ ปามงฺโค,}{กณฺฐสุตฺตํ นธารเย.} \\ \csromangathabat{กฏิ โอวฏฺฏิ กายุรํ,}{หตฺถาภรณมุทฺทิกา.}}\csromangathastanza{\csromangathabat{ทีเฆ โกจฺเฉ ผเณ หตฺเถ,}{สิตฺถา อุทกเตลเก.} \\ \csromangathabat{อาทาสุทปตฺตวณา,}{อาเลโปมฺมทฺทจุณฺณนา.}}\csromangathastanza{\csromangathabat{ลญฺเฉนฺติ องฺคราคญฺจ,}{มุขราคํ ตทูภยํ.} \\ \csromangathabat{จกฺขุโรคํ คิรคฺคญฺจ,}{อายตํ สรพาหิรํ.}}\csromangathastanza{\csromangathabat{อมฺพเปสิสกเลหิ,}{อหิจฺฉินฺทิ จ จนฺทนํ.} \\ \csromangathabat{อุจฺจาวจา ปตฺตมูลา,}{สุวณฺโณ พหลา วลี.}}\csromangathastanza{\csromanfolio{286}\csromangathabat{จิตฺรา ทุสฺสติ ทุคฺคนฺโธ,}{อุเณฺห ภิชฺชิํสุ มิฑฺฒิยา.} \\ \csromangathabat{ปริภณฺฑํ ติณํ โจฬํ,}{มาฬํ กุณฺโฑลิกาย จ.}}\csromangathastanza{\csromangathabat{ถวิกา อํสพทฺธญฺจ,}{ตถา พนฺธนสุตฺตกา.} \\ \csromangathabat{ขิเล มญฺเจ จ ปีเฐ จ,}{องฺเก ฉตฺเต ปณามนา.}}\csromangathastanza{\csromangathabat{ตุมฺพฆฏิฉวสีสํ,}{จลกานิ ปฏิคฺคโห.} \\ \csromangathabat{วิปฺผาลิทณฺฑโสวณฺณํ,}{ปตฺเต เปสิ จ นาฬิกา.}}\csromangathastanza{\csromangathabat{กิณฺณสตฺตุ สริตญฺจ,}{มธุสิตฺถํ สิปาฏิกํ.} \\ \csromangathabat{วิกณฺณํ พนฺธิวิสมํ,}{ฉมาชิรปโหติ จ.}}\csromangathastanza{\csromangathabat{กฬิมฺภํ โมฆสุตฺตญฺจ,}{อโธตลฺลํ อุปาหนา.} \\ \csromangathabat{องฺคุเล ปฏิคฺคหญฺจ,}{วิตฺถกํ ถวิกพทฺธกา.}}\csromangathastanza{\csromangathabat{อชฺโฌกาเส นีจวตฺถุ,}{จโย จาปิ วิหญฺญเร.} \\ \csromangathabat{ปริปติ ติณจุณฺณํ,}{อุลฺลิตฺตอวลิตฺตกํ.}}\csromangathastanza{\csromangathabat{เสตํ กาฬกวณฺณญฺจ,}{ปริกมฺมญฺจ เครุกํ.} \\ \csromangathabat{มาลากมฺมํ ลตากมฺมํ,}{มกรทนฺตกปาฏิกํ.}}\csromangathastanza{\csromangathabat{จีวรวํสํ รชฺชุญฺจ,}{อนุญฺญาสิ วินายโก.} \\ \csromangathabat{อุชฺฌิตฺวา ปกฺกมนฺติ จ,}{กถินํ ปริภิชฺชติ.}}\csromangathastanza{\csromangathabat{วินิเวฐิยติ กุฏฺเฏปิ,}{ปตฺเตนาทาย คจฺฉเร.} \\ \csromangathabat{ถวิกา พนฺธสุตฺตญฺจ,}{พนฺธิตฺวา จ อุปาหนา.}}\csromangathastanza{\csromangathabat{อุปาหนตฺถวิกญฺจ,}{อํสพทฺธญฺจ สุตฺตกํ.} \\ \csromangathabat{อุทกากปฺปิยํ มคฺเค,}{ปริสฺสาวนโจฬกํ.}}\csromangathastanza{\csromangathabat{ธมฺมกรณํ เทฺว ภิกฺขู,}{เวสาลิํ อคมา มุนิ.} \\ \csromangathabat{ทณฺฑํ โอตฺถรกํ ตตฺถ,}{อนุญฺญาสิ ปริสฺสาวนํ.}}\csromangathastanza{\csromangathabat{มกเสหิ ปณีเตน,}{พหฺวาพาธา จ ชีวโก.} \\ \csromangathabat{จงฺกมนชนฺตาฆรํ,}{วิสเม นีจวตฺถุโก.}}\csromangathastanza{\csromangathabat{ตโย จเย วิหญฺญนฺติ,}{โสปานาลมฺพเวทิตํ.} \\ \csromangathabat{อชฺโฌกาเส ติณจุณฺณํ,}{อุลฺลิตฺตอวลิตฺตกํ.}}\csromangathastanza{\csromanfolio{287}\csromangathabat{เสตกํ กาฬวณฺณญฺจ,}{ปริกมฺมญฺจ เครุกํ.} \\ \csromangathabat{มาลากมฺมํ ลตากมฺมํ,}{มกรทนฺตกปาฏิกํ.}}\csromangathastanza{\csromangathabat{วํสํ จีวรรชฺชุญฺจ,}{อุจฺจํ จ วตฺถุกํ กเร.} \\ \csromangathabat{จโย โสปานพาหญฺจ,}{กวาฏํ ปิฏฺฐสํฆาฏํ.}}\csromangathastanza{\csromangathabat{อุทุกฺขลุตฺตรปาสกํ,}{วฏฺฏิญฺจ กปิสีสกํ.} \\ \csromangathabat{สูจิฆฏิตาฬจฺฉิทฺทํ,}{อาวิญฺฉนญฺจ รชฺชุกํ.}}\csromangathastanza{\csromangathabat{มณฺฑลํ ธูมเนตฺตญฺจ,}{มชฺเฌ จ มุขมตฺติกา.} \\ \csromangathabat{โทณิ ทุคฺคนฺธา ทหติ,}{อุทกฏฺฐานํ สราวกํ.}}\csromangathastanza{\csromangathabat{น เสเทติ จ จิกฺขลฺลํ,}{โธวิ นิทฺธมนํ กเร.} \\ \csromangathabat{ปีฐญฺจ โกฏฺฐเก กมฺมํ,}{มรุมฺพา สิลา นิทฺธมนํ.}}\csromangathastanza{\csromangathabat{นคฺคา ฉมาย วสฺสนฺเต,}{ปฏิจฺฉาที ตโย ตหิํ.} \\ \csromangathabat{อุทปานํ ลุชฺชติ นีจํ,}{วลฺลิยา กายพนฺธเน.}}\csromangathastanza{\csromangathabat{ตุลํ กฏกฏํ จกฺกํ,}{พหู ภิชฺชนฺติ ภาชนา.} \\ \csromangathabat{โลหทารุจมฺมขณฺฑํ,}{สาลาติณาปิธานิ จ.}}\csromangathastanza{\csromangathabat{โทณิจนฺทนิ ปาการํ,}{จิกฺขลฺลํ นิทฺธมเนน จ.} \\ \csromangathabat{สีติคตํ โปกฺขรณิํ,}{ปุราณญฺจ นิลฺเลขณํ.}}\csromangathastanza{\csromangathabat{จาตุมาสํ สยนฺติ จ,}{นมตกญฺจ\csromansharedfootnote{9ffaab2f626dc5f9}{คนฺธปุปฺผํ (สฺยา)} นธิฏฺฐเห.} \\ \csromangathabat{อาสิตฺตกํ มโฬริกํ,}{ภุญฺชนฺเตกํ ตุวฏฺเฏยฺยุํ.}}\csromangathastanza{\csromangathabat{วฑฺโฒ โพธิ น อกฺกมิ,}{ฆฏํ กตกสมฺมชฺชนิ.} \\ \csromangathabat{สกฺขรํ กถลญฺเจว,}{เผณกํ ปาทฆํสนี.}}\csromangathastanza{\csromangathabat{วิธูปนํ ตาลวณฺฏํ,}{มกสญฺจาปิ จามรี.} \\ \csromangathabat{ฉตฺตํ วินา จ อาราเม,}{ตโย สิกฺกาย สมฺมุติ.}}\csromangathastanza{\csromangathabat{โรมสิตฺถา นขา ทีฆา,}{ฉินฺทนฺตงฺคุลิกา ทุกฺขา.} \\ \csromangathabat{สโลหิตํ ปมาณญฺจ,}{วีสติ ทีฆเกสตา.}}\csromangathastanza{\csromangathabat{ขุรํ สิลํ สิปาฏิกํ,}{นมตกํ ขุรภณฺฑกํ.} \\ \csromangathabat{มสฺสุํ กปฺเปนฺติ วฑฺเฒนฺติ,}{โคโลมิจตุรสฺสกํ.}}\csromangathastanza{\csromanfolio{288}\csromangathabat{ปริมุขํ อฑฺฒทุกญฺจ,}{ทาฐิสมฺพาธสํหเร.} \\ \csromangathabat{อาพาธา กตฺตริวโณ,}{ทีฆํ สกฺขริกาย จ.}}\csromangathastanza{\csromangathabat{ปลิตํ ถกิตํ อุจฺจา,}{โลหภณฺฑญฺชนี สห.} \\ \csromangathabat{ปลฺลตฺถิกญฺจ อาโยโค,}{วฏํ สลากพนฺธนํ.}}\csromangathastanza{\csromangathabat{กลาพุกํ เทฑฺฑุภกํ,}{มุรชํ มทฺทวีณกํ.} \\ \csromangathabat{ปฏฺฏิกา สูกรนฺตญฺจ,}{ทสา มุรชเวณิกา.}}\csromangathastanza{\csromangathabat{อนฺโต โสภํ คุณกญฺจ,}{ปวนนฺโตปิ ชีรติ.} \\ \csromangathabat{คณฺฐิกา อุจฺจาวจญฺจ,}{ผลกนฺเตปิ โอคาเห.}}\csromangathastanza{\csromangathabat{คิหิวตฺถํ หตฺถิโสณฺฑํ,}{มจฺฉกํ จตุกณฺณกํ.} \\ \csromangathabat{ตาลวณฺฏํ สตวลิ,}{คิหิปารุตปารุปํ.}}\csromangathastanza{\csromangathabat{สํเวลฺลิ อุภโตกาชํ,}{ทนฺตกฏฺฐํ อาโกฏเน.} \\ \csromangathabat{กณฺเฐ วิลคฺคํ ทายญฺจ,}{ปฏคฺคิ รุกฺขหตฺถินา.}}\csromangathastanza{\csromangathabat{ยเมเฬ โลกายตกํ,}{ปริยาปุณิํสุ วาจยุํ.} \\ \csromangathabat{ติรจฺฉาน\-กถา วิชฺฌา,}{ขิปิ มงฺคลํ ขาทิ จ.}}\csromangathastanza{\csromangathabat{วาตาพาโธ ทุสฺสติ จ,}{ทุคฺคนฺโธ ทุกฺขปาทุกา.} \\ \csromangathabat{หิริยนฺติ ปารุทุคฺคนฺโธ,}{ตหํ ตหํ กโรนฺติ จ.}}\csromangathastanza{\csromangathabat{ทุคฺคนฺโธ กูปํ ลุชฺชนฺติ,}{อุจฺจวตฺถุ จเยน จ.} \\ \csromangathabat{โสปานาลมฺพนพาหา,}{อนฺเต ทุกฺขญฺจ ปาทุกา.}}\csromangathastanza{\csromangathabat{พหิทฺธา โทณิ กฏฺฐญฺจ,}{ปิฐโร จ อปารุโต.} \\ \csromangathabat{วจฺจกุฏิํ กวาฏญฺจ,}{ปิฏฺฐสํฆาฏเมว จ.}}\csromangathastanza{\csromangathabat{อุทุกฺขลุตฺตรปาโส,}{วฏฺฏิญฺจ กปิสีสกํ.} \\ \csromangathabat{สูจิฆฏิตาฬจฺฉิทฺทํ,}{อาวิญฺฉนจฺฉิทฺทเมว จ.}}\csromangathastanza{\csromangathabat{รชฺชุ อุลฺลิตฺตา\-วลิตฺตํ,}{เสตวณฺณญฺจ กาฬกํ.} \\ \csromangathabat{มาลากมฺมํ ลตากมฺมํ,}{มกรํ ปญฺจปฏิกํ.}}\csromangathastanza{\csromangathabat{จีวรวํสํ รชฺชุญฺจ,}{ชราทุพฺพลปาการํ.} \\ \csromangathabat{โกฏฺฐเก จาปิ ตเถว,}{มรุมฺพํ ปทรสิลา.}}\csromangathastanza{\csromangathabat{สนฺติฏฺฐติ นิทฺธมนํ,}{กุมฺภิญฺจาปิ สราวกํ.} \\ \csromangathabat{ทุกฺขํ หิริ อปิธานํ,}{อนาจารญฺจ อาจรุํ.}}\csromangathastanza{\csromanfolio{289}\csromangathabat{โลหภณฺฑํ อนุญฺญาสิ,}{ฐปยิตฺวา ปหรณิํ.} \\ \csromangathabat{ฐปยิตฺวาสนฺทิปลฺลงฺกํ,}{ทารุปตฺตญฺจ ปาทุกํ.}}\csromangathastanza{\csromangathabat{สพฺพํ ทารุมยํ ภณฺฑํ,}{อนุญฺญาสิ มหามุนิ.} \\ \csromangathabat{กตกํ กุมฺภการญฺจ,}{ฐปยิตฺวา ตถาคโต.}}\csromangathastanza{\csromangathabat{สพฺพมฺปิ มตฺติกาภณฺฑํ,}{อนุญฺญาสิ อนุกมฺปโก.} \\ \csromangathabat{ยสฺส วตฺถุสฺส นิทฺเทโส,}{ปุริเมน ยทิ สมํ.}}\csromangathastanza{\csromangathabat{ตมฺปิ สํขิตฺตํ อุทฺทาเน,}{นยโต ตํ วิชานิยา.} \\ \csromangathabat{เอวํ ทสสตา วตฺถู,}{วินเย ขุทฺทกวตฺถุเก.}}\csromangathastanza{\csromangathabat{สทฺธมฺมฏฺฐิติโก เจว,}{เปสลานญฺจนุคฺคโห.} \\ \csromangathabat{สุสิกฺขิโต วินยธโร,}{หิตจิตฺโต สุเปสโล.} \\ \csromangathabat{ปทีปกรโณ ธีโร,}{ปูชารโห พหุสฺสุโตติ.}}}{ขุทฺทก\-วตฺถุ\-กฺขนฺธโก นิฏฺฐิโต.}
\csromanmatikaanchor{mtk.129}
\csromantocmarknopage{chapter}{6. เสนาสนกฺขนฺธก}{mtk.129}
\bookmark[dest={mtk.129},level=0]{6. เสนาสนกฺขนฺธก}
\chapterheadpageread{6. เสนาสน\-กฺขนฺธก}
\csromanmatikaanchor{mtk.130}
\csromantocmarkref{section}{1. ปฐมภาณวาร}{mtk.130}
\bookmark[dest={mtk.130},level=1]{1. ปฐมภาณวาร}
\csromanheader{1}{1. ปฐม\-ภาณวาร}
\csromanmatikaanchor{mtk.131}
\csromantocmarkref{subsection}{วิหารานุชานน}{mtk.131}
\bookmark[dest={mtk.131},level=2]{วิหารานุชานน}
\csromanheader{2}{\textbf{วิหารานุชานน}}
\proseitem{294}{\csromanfolio{290}เตน สมเยน พุทฺโธ ภควา ราชคเห วิหรติ เวฬุวเน กลนฺทก\-นิวาเป. เตน โข ปน สมเยน ภควตา ภิกฺขูนํ เสนาสนํ อปญฺญตฺตํ โหติ. เต จ ภิกฺขู ตหํ ตหํ วิหรนฺติ อรญฺเญ รุกฺขมูเล ปพฺพเต กนฺทรายํ คิริคุหายํ สุสาเน วนปตฺเถ อชฺโฌกาเส ปลาลปุญฺเช, เต กาลสฺเสว ตโต ตโต อุปนิกฺขมนฺติ อรญฺญา รุกฺขมูลา ปพฺพตา กนฺทรา คิริคุหา สุสานา วนปตฺถา อชฺโฌกาสา ปลาลปุญฺชา ปาสาทิเกน อภิกฺกนฺเตน ปฏิกฺกนฺเตน อาโลกิเตน วิโลกิเตน สมิญฺชิเตน ปสาริเตน โอกฺขิตฺตจกฺขู อิริยาปถ\-สมฺปนฺนา. เตน โข ปน สมเยน ราชคหโก เสฏฺฐี\csromansharedfootnote{d4d6c9ddb07cc09d}{เสฏฺฐิ (ก)} กาลสฺเสว อุยฺยานํ อคมาสิ. อทฺทสา โข ราชคหโก เสฏฺฐี เต ภิกฺขู กาลสฺเสว ตโต ตโต อุปนิกฺขมนฺเต อรญฺญา รุกฺขมูลา ปพฺพตา กนฺทรา คิริคุหา สุสานา วนปตฺถา อชฺโฌกาสา ปลาลปุญฺชา ปาสาทิเกน อภิกฺกนฺเตน ปฏิกฺกนฺเตน อาโลกิเตน วิโลกิเตน สมิญฺชิเตน ปสาริเตน โอกฺขิตฺตจกฺขู อิริยาปถ\-สมฺปนฺเน, ทิสฺวานสฺส จิตฺตํ ปสีทิ. อถ โข ราชคหโก เสฏฺฐี เยน เต ภิกฺขู เตนุปสงฺกมิ, อุปสงฺกมิตฺวา เต ภิกฺขู เอตทโวจ “สจาหํ ภนฺเต วิหาเร การาเปยฺยํ, วเสยฺยาถ เม วิหาเรสู”ติ. น โข คหปติ ภควตา วิหารา อนุญฺญาตาติ. เตน หิ ภนฺเต ภควนฺตํ ปฏิปุจฺฉิตฺวา มม อาโรเจยฺยาถาติ. “เอวํ คหปตี”ติ โข เต ภิกฺขู ราชคหกสฺส เสฏฺฐิสฺส ปฏิสฺสุตฺวา เยน ภควา เตนุ\-ปสงฺกมิํสุ, อุปสงฺกมิตฺวา ภควนฺตํ อภิวาเทตฺวา เอกมนฺตํ นิสีทิํสุ. เอกมนฺตํ นิสินฺนา โข เต ภิกฺขู ภควนฺตํ เอตทโวจุํ “ราชคหโก ภนฺเต เสฏฺฐี วิหาเร การาเปตุกาโม, กถํ นุ โข ภนฺเต อเมฺหหิ\csromansharedfootnote{b5ecd50093aa9023}{ภนฺเต (สี, ก)} ปฏิปชฺชิตพฺพนฺ”ติ. อถ โข ภควา เอตสฺมิํ นิทาเน เอตสฺมิํ ปกรเณ ธมฺมิํ กถํ กตฺวา ภิกฺขู อามนฺเตสิ “อนุชานามิ ภิกฺขเว ปญฺจ เลณานิ\csromansharedfootnote{aeeea879ed441d15}{ปญฺจ เสนาสนานิ (สฺยา)} วิหารํ อฑฺฒโยคํ ปาสาทํ หมฺมิยํ คุหนฺ”ติ.}

\prose{\csromanfolio{291}อถ โข เต ภิกฺขู เยน ราชคหโก เสฏฺฐี เตนุ\-ปสงฺกมิํสุ, อุปสงฺกมิตฺวา ราชคหกํ เสฏฺฐิํ เอตทโวจุํ “อนุญฺญาตา โข คหปติ ภควตา วิหารา, ยสฺสทานิ กาลํ มญฺญสี”ติ. อถ โข ราชคหโก เสฏฺฐี เอกาเหเนว สฏฺฐิวิหาเร ปติฏฺฐาเปสิ. อถ โข ราชคหโก เสฏฺฐี เต สฏฺฐิวิหาเร ปริโยสาเปตฺวา เยน ภควา เตนุปสงฺกมิ, อุปสงฺกมิตฺวา ภควนฺตํ อภิวาเทตฺวา เอกมนฺตํ นิสีทิ. เอกมนฺตํ นิสินฺโน โข ราชคหโก เสฏฺฐี ภควนฺตํ เอตทโวจ “อธิวาเสตุ เม ภนฺเต ภควา สฺวาตนาย ภตฺตํ สทฺธิํ ภิกฺขุสํเฆนา”ติ. อธิวาเสสิ ภควา ตุณฺหีภาเวน. อถ โข ราชคหโก เสฏฺฐี ภควโต อธิวาสนํ วิทิตฺวา อุฏฺฐายาสนา ภควนฺตํ อภิวาเทตฺวา ปทกฺขิณํ กตฺวา ปกฺกามิ.}

\prose{อถ โข ราชคหโก เสฏฺฐี ตสฺสา รตฺติยา อจฺจเยน ปณีตํ ขาทนียํ โภชนียํ ปฏิยาทาเปตฺวา ภควโต กาลํ อาโรจาเปสิ “กาโล ภนฺเต, นิฏฺฐิตํ ภตฺตนฺ”ติ. อถ โข ภควา ปุพฺพณฺหสมยํ นิวาเสตฺวา ปตฺต\-จีวรมา\-ทาย เยน ราชคหกสฺส เสฏฺฐิสฺส นิเวสนํ เตนุปสงฺกมิ, อุปสงฺกมิตฺวา ปญฺญตฺเต อาสเน นิสีทิ สทฺธิํ ภิกฺขุ\-สํเฆน. อถ โข ราชคหโก เสฏฺฐี พุทฺธ\-ปฺปมุขํ ภิกฺขุสํฆํ ปณีเตน ขาทนีเยน โภชนีเยน สหตฺถา สนฺตปฺเปตฺวา สมฺปวาเรตฺวา ภควนฺตํ ภุตฺตาวิํ โอนีต\-ปตฺต\-ปาณิํ เอกมนฺตํ นิสีทิ. เอกมนฺตํ นิสินฺโน โข ราชคหโก เสฏฺฐี ภควนฺตํ เอตทโวจ “เอเต เม ภนฺเต สฏฺฐิวิหารา ปุญฺญตฺถิเกน สคฺคตฺถิเกน การาปิตา, กถาหํ ภนฺเต เตสุ วิหาเรสุ ปฏิปชฺชามี”ติ. เตน หิ ตฺวํ คหปติ เต สฏฺฐิวิหาเร อาคตานาคตสฺส จาตุทฺทิสสฺส สํฆสฺส ปติฏฺฐาเปหีติ. “เอวํ ภนฺเต”ติ โข ราชคหโก เสฏฺฐี ภควโต ปฏิสฺสุตฺวา เต สฏฺฐิวิหาเร อาคตานาคตสฺส จาตุทฺทิสสฺส สํฆสฺส ปติฏฺฐาเปสิ.}

\proseitem{295}{อถ โข ภควา ราชคหกํ เสฏฺฐิํ อิมาหิ คาถาหิ อนุโมทิ\csromandash{}}

\setlength{\csromangathapagewidth}{0pt}
\setlength{\csromangathawakwidth}{0pt}
\csromangathameasuregroup{“สีตํ อุณฺหํ ปฏิหนฺติ\textsuperscript{1}, \\ สรีสเป จ มกเส, \\ ตโต วาตาตโป โฆโร, \\ เลณตฺถญฺจ สุขตฺถญฺจ, \\ วิหารทานํ สํฆสฺส, \\ ตสฺมา หิ ปณฺฑิโต โปโส, \\ วิหาเร การเย รมฺเม, \\ เตสํ อนฺนญฺจ ปานญฺจ, \\ ทเทยฺย อุชุภูเตสุ, \\ เต ตสฺส ธมฺมํ เทเสนฺติ, \\ ยํ โส ธมฺมํ อิธญฺญาย,}{\csromangathabat{“สีตํ อุณฺหํ ปฏิหนฺติ\textsuperscript{1},}{ตโต วาฬมิคานิ จ.} \\ \csromangathabat{สรีสเป จ มกเส,}{สิสิเร จาปิ วุฏฺฐิโย.} \\ \csromangathabat{ตโต วาตาตโป โฆโร,}{สญฺชาโต\textsuperscript{1} ปฏิหญฺญติ.} \\ \csromangathabat{เลณตฺถญฺจ สุขตฺถญฺจ,}{ฌายิตุญฺจ วิปสฺสิตุํ.} \\ \csromangathabat{วิหารทานํ สํฆสฺส,}{อคฺคํ พุทฺเธน\textsuperscript{1} วณฺณิตํ.} \\ \csromangathabat{ตสฺมา หิ ปณฺฑิโต โปโส,}{สมฺปสฺสํ อตฺถมตฺตโน.} \\ \csromangathabat{วิหาเร การเย รมฺเม,}{วาสเยตฺถ พหุสฺสุเต.} \\ \csromangathabat{เตสํ อนฺนญฺจ ปานญฺจ,}{วตฺถเสนาสนานิ จ.} \\ \csromangathabat{ทเทยฺย อุชุภูเตสุ,}{วิปฺปสนฺเนน เจตสา.} \\ \csromangathabat{เต ตสฺส ธมฺมํ เทเสนฺติ,}{สพฺพทุกฺขาปนูทนํ.} \\ \csromangathabat{ยํ โส ธมฺมํ อิธญฺญาย,}{ปรินิพฺพาติ อนาสโว”ติ.}}
\csromangathasetleft{“สีตํ อุณฺหํ ปฏิหนฺติ\textsuperscript{1}, \\ สรีสเป จ มกเส, \\ ตโต วาตาตโป โฆโร, \\ เลณตฺถญฺจ สุขตฺถญฺจ, \\ วิหารทานํ สํฆสฺส, \\ ตสฺมา หิ ปณฺฑิโต โปโส, \\ วิหาเร การเย รมฺเม, \\ เตสํ อนฺนญฺจ ปานญฺจ, \\ ทเทยฺย อุชุภูเตสุ, \\ เต ตสฺส ธมฺมํ เทเสนฺติ, \\ ยํ โส ธมฺมํ อิธญฺญาย,}
\csromangathagroup{\csromangathastanza{\csromangathabat{“สีตํ อุณฺหํ ปฏิหนฺติ\csromansharedfootnote{19f997f30f4c0aba}{ปฏิหนติ (ก)},}{ตโต วาฬมิคานิ จ.} \\ \csromangathabat{สรีสเป จ มกเส,}{สิสิเร จาปิ วุฏฺฐิโย.}}\csromangathastanza{\csromanfolio{292}\csromangathabat{ตโต วาตาตโป โฆโร,}{สญฺชาโต\csromansharedfootnote{732e47f95dbd69fe}{วาตาตเป โฆเร, สญฺชเต (ก, สทฺทนีติ)} ปฏิหญฺญติ.} \\ \csromangathabat{เลณตฺถญฺจ สุขตฺถญฺจ,}{ฌายิตุญฺจ วิปสฺสิตุํ.}}\csromangathastanza{\csromangathabat{วิหารทานํ สํฆสฺส,}{อคฺคํ พุทฺเธน\csromansharedfootnote{8eb094223f30d5be}{พุทฺเธหิ (สฺยา)} วณฺณิตํ.} \\ \csromangathabat{ตสฺมา หิ ปณฺฑิโต โปโส,}{สมฺปสฺสํ อตฺถมตฺตโน.}}\csromangathastanza{\csromangathabat{วิหาเร การเย รมฺเม,}{วาสเยตฺถ พหุสฺสุเต.} \\ \csromangathabat{เตสํ อนฺนญฺจ ปานญฺจ,}{วตฺถ\-เสนาสนานิ จ.}}\csromangathastanza{\csromangathabat{ทเทยฺย อุชุภูเตสุ,}{วิปฺปสนฺเนน เจตสา.} \\ \csromangathabat{เต ตสฺส ธมฺมํ เทเสนฺติ,}{สพฺพทุกฺขาปนูทนํ.} \\ \csromangathabat{ยํ โส ธมฺมํ อิธญฺญาย,}{ปรินิพฺพาติ อนาสโว”ติ.}}}

\prose{อถ โข ภควา ราชคหกํ เสฏฺฐิํ อิมาหิ คาถาหิ อนุโมทิตฺวา อุฏฺฐายาสนา ปกฺกามิ.}

\proseitem{296}{อสฺโสสุํ โข มนุสฺสา “ภควโต กิร วิหารา อนุญฺญาตา”ติ สกฺกจฺจํ\csromansharedfootnote{0266fd0e79b724f6}{เต สกฺกจฺจํ (สฺยา, กํ)} วิหาเร การาเปนฺติ, เต วิหารา อกวาฏกา โหนฺติ, อหีปิ วิจฺฉิกาปิ สตปทิโยปิ ปวิสนฺติ. ภควโต เอตมตฺถํ อาโรเจสุํ. อนุชานามิ ภิกฺขเว กวาฏนฺติ. ภิตฺติฉิทฺทํ กริตฺวา วลฺลิยาปิ รชฺชุยาปิ กวาฏํ พนฺธนฺติ, อุนฺทูเรหิปิ อุปจิกาหิปิ ขชฺชนฺติ, ขยิต\-พนฺธนานิ กวาฏานิ ปตนฺติ. ภควโต เอตมตฺถํ อาโรเจสุํ. อนุชานามิ ภิกฺขเว ปิฏฺฐสํฆาฏํ อุทุกฺขลิกํ อุตฺตรปาสกนฺติ. กวาฏา น ผุสียนฺติ. ภควโต เอตมตฺถํ อาโรเจสุํ. อนุชานามิ ภิกฺขเว อาวิญฺฉน\-จฺฉิทฺทํ อาวิญฺฉน\-รชฺชุนฺติ. กวาฏา น ถกิยนฺติ. ภควโต เอตมตฺถํ อาโรเจสุํ. อนุชานามิ ภิกฺขเว อคฺคฬวฏฺฏิํ กปิสีสกํ สูจิกํ ฆฏิกนฺติ.}

\prose{เตน โข ปน สมเยน ภิกฺขู น สกฺโกนฺติ กวาฏํ อปาปุริตุํ. ภควโต เอตมตฺถํ อาโรเจสุํ. อนุชานามิ ภิกฺขเว ตาฬจฺฉิทฺทํ, ตีณี ตาฬานิ โลหตาฬํ กฏฺฐตาฬํ วิสาณตาฬนฺติ, เยหิ\csromansharedfootnote{3d40054fbd27ba41}{เยปิ (สี, สฺยา)} เต อุคฺฆาเฏตฺวา ปวิสนฺติ\csromansharedfootnote{01609121a334be58}{วิสาณตาฬํ, เยหิ เต อุคฺฆาเฏตฺวา ปวิสนฺตีติ (ก)}. วิหารา อคุตฺตา โหนฺติ. ภควโต เอตมตฺถํ อาโรเจสุํ. อนุชานามิ ภิกฺขเว ยนฺตกํ สูจิกนฺติ.}

\prose{\csromanfolio{293}เตน โข ปน สมเยน วิหารา ติณจฺฉทนา โหนฺติ, สีตกาเล สีตา, อุณฺหกาเล อุณฺหา. ภควโต เอตมตฺถํ อาโรเจสุํ. อนุชานามิ ภิกฺขเว โอคุมฺเผตฺวา อุลฺลิตฺตา\-วลิตฺตํ กาตุนฺติ.}

\prose{เตน โข ปน สมเยน วิหารา อวาตปานกา โหนฺติ อจกฺขุสฺสา ทุคฺคนฺธา. ภควโต เอตมตฺถํ อาโรเจสุํ. อนุชานามิ ภิกฺขเว ตีณิ วาตปานานิ เวทิกา\-วาตปานํ ชาลวาตปานํ สลาก\-วาตปานนฺติ. วาตปาน\-นฺตริกาย กาฬกาปิ วคฺคุลิโยปิ ปวิสนฺติ. ภควโต เอตม\-ตฺถํ\-อาโรเจสุํ. อนุชานามิ ภิกฺขเว วาตปาน\-จกฺกลิกนฺติ. จกฺกลิกนฺตริกายปิ กาฬกาปิ วคฺคุลิโยปิ ปวิสนฺติ. ภควโต เอตมตฺถํ อาโรเจสุํ. อนุชานามิ ภิกฺขเว วาตปาน\-กวาฏกํ วาตปาน\-ภิสิกนฺติ.}

\prose{เตน โข ปนสมเยน ภิกฺขู ฉมายํ สยนฺติ, คตฺตานิปิ จีวรานิปิ ปํสุกิตานิ โหนฺติ. ภควโต เอตมตฺถํ อาโรเจสุํ. อนุชานามิ ภิกฺขเว ติณสนฺถารกนฺติ. ติณสนฺถารโก อุนฺทูเรหิปิ อุปจิกาหิปิ ขชฺชติ. ภควโต เอตมตฺถํ อาโรเจสุํ. อนุชานามิ ภิกฺขเว มิฑฺฒินฺติ\csromansharedfootnote{988fe9de5dfb77bc}{มีฒนฺติ (สี), มิฒินฺติ (สฺยา)}. มิฑฺฒิยา คตฺตานิ ทุกฺขา โหนฺติ. ภควโต เอตมตฺถํ อาโรเจสุํ. อนุชานามิ ภิกฺขเว พิทลมญฺจกนฺติ.}

\csromanmatikaanchor{mtk.132}
\csromantocmarkref{subsection}{มญฺจปีฐาทิอนุชานน}{mtk.132}
\bookmark[dest={mtk.132},level=2]{มญฺจปีฐาทิอนุชานน}
\csromanheader{2}{มญฺจปีฐา\-ทิ\-อนุชานน}
\proseitem{297}{เตน โข ปน สมเยน สํฆสฺส โสสานิโก มสารโก มญฺโจ อุปฺปนฺโน โหติ. ภควโต เอตมตฺถํ อาโรเจสุํ. อนุชานามิ ภิกฺขเว มสารกํ มญฺจนฺติ. มสารกํ ปีฐํ อุปฺปนฺนํ โหติ. ภควโต เอตมตฺถํ อาโรเจสุํ. อนุชานามิ ภิกฺขเว มสารกํ ปีฐนฺติ.}

\prose{เตน โข ปน สมเยน สํฆสฺส โสสานิโก พุนฺทิกาพทฺโธ มญฺโจ อุปฺปนฺโน โหติ. ภควโต เอตมตฺถํ อาโรเจสุํ. อนุชานามิ ภิกฺขเว พุนฺทิกาพทฺธํ มญฺจนฺติ. พุนฺทิกาพทฺธํ ปีฐํ อุปฺปนฺนํ โหติ. ภควโต เอตมตฺถํ อาโรเจสุํ. อนุชานามิ ภิกฺขเว พุนฺทิกาพทฺธํ ปิฐนฺติ.}

\prose{เตน โข ปน สมเยน สํฆสฺส โสสานิโก กุฬีรปาทโก มญฺโจ อุปฺปนฺโน โหติ. ภควโต เอตมตฺถํ อาโรเจสุํ. อนุชานามิ \csromanfolio{294}ภิกฺขเว กุฬีรปาทกํ มญฺจนฺติ. กุฬีรปาทกํ ปีฐํ อุปฺปนฺนํ โหติ. ภควโต เอตมตฺถํ อาโรเจสุํ. อนุชานามิ ภิกฺขเว กุฬีรปาทกํ ปีฐนฺติ.}

\prose{เตน โข ปน สมเยน สํฆสฺส โสสานิโก อาหจฺจปาทโก มญฺโจ อุปฺปนฺโน โหติ. ภควโต เอตมตฺถํ อาโรเจสุํ. อนุชานามิ ภิกฺขเว อาหจฺจปาทกํ มญฺจนฺติ. อาหาจฺจปาทกํ ปีฐํ อุปฺปนฺนํ โหติ. ภควโต เอตมตฺถํ อาโรเจสุํ. อนุชานามิ ภิกฺขเว อาหจฺจปาทกํ ปีฐนฺติ.}

\prose{เตน โข ปน สมเยน สํฆสฺส อาสนฺทิโก อุปฺปนฺโน โหติ. ภควโต เอตมตฺถํ อาโรเจสุํ. อนุชานามิ ภิกฺขเว อาสนฺทิกนฺติ. อุจฺจโก อาสนฺทิโก อุปฺปนฺโน โหติ. ภควโต เอตมตฺถํ อาโรเจสุํ. อนุชานามิ ภิกฺขเว อุจฺจกมฺปิ อาสนฺทิกนฺติ. สตฺตงฺโค อุปฺปนฺโน โหติ. ภควโต เอตมตฺถํ อาโรเจสุํ. อนุชานามิ ภิกฺขเว สตฺตงฺคนฺติ. อุจฺจโก สตฺตงฺโค อุปฺปนฺโน โหติ. ภควโต เอตมตฺถํ อาโรเจสุํ. อนุชานามิ ภิกฺขเว อุจฺจกมฺปิ สตฺตงฺคนฺติ. ภทฺทปีฐํ อุปฺปนฺนํ โหติ. ภควโต เอตมตฺถํ อาโรเจสุํ. อนุชานามิ ภิกฺขเว ภทฺทปีฐนฺติ. ปีฐิกา อุปฺปนฺนา โหติ. ภควโต เอตมตฺถํ อาโรเจสุํ. อนุชานามิ ภิกฺขเว ปีฐิกนฺติ. เอฬกปาทกํ ปีฐํ อุปฺปนฺนํ โหติ. ภควโต เอตมตฺถํ อาโรเจสุํ. อนุชานามิ ภิกฺขเว เอฬกปาทกํ ปีฐนฺติ. อามลก\-วฏฺฏิกํ\csromansharedfootnote{ef3cd11ea0257cbd}{อามลกวณฺฏิกํ (สฺยา)} ปีฐํ อุปฺปนฺนํ โหติ. ภควโต เอตมตฺถํ อาโรเจสุํ. อนุชานามิ ภิกฺขเว อามลก\-วฏฺฏิกํ ปีฐนฺติ. ผลกํ อุปฺปนฺนํ โหติ. ภควโต เอตมตฺถํ อาโรเจสุํ. อนุชานามิ ภิกฺขเว ผลกนฺติ. โกจฺฉํ อุปฺปนฺนํ โหติ. ภควโต เอตมตฺถํ อาโรเจสุํ. อนุชานามิ ภิกฺขเว โกจฺฉนฺติ. ปลาลปีฐํ อุปฺปนฺนํ โหติ. ภควโต เอตมตฺถํ อาโรเจสุํ. อนุชานามิ ภิกฺขเว ปลาลปีฐนฺติ.}

\prose{เตน โข ปน สมเยน ฉพฺพคฺคิยา ภิกฺขู อุจฺเจ มญฺเจ สยนฺติ. มนุสฺสา วิหารจาริกํ อาหิณฺฑนฺตา ปสฺสิตฺวา อุชฺฌายนฺติ ขิยฺยนฺติ วิปาเจนฺติ “เสยฺยถาปิ คิหี กามโภคิโน”ติ. ภควโต เอตมตฺถํ อาโรเจสุํ. น ภิกฺขเว อุจฺเจ มญฺเจ สยิตพฺพํ, โย สเยยฺย, อาปตฺติ ทุกฺกฏสฺสาติ.}

\prose{\csromanfolio{295}เตน โข ปน สมเยน อญฺญตโร ภิกฺขุ นีเจ มญฺเจ สยนฺโต อหินา ทฏฺโฐ โหติ. ภควโต เอตมตฺถํ อาโรเจสุํ. อนุชานามิ ภิกฺขเว มญฺจ\-ปฏิปาทกนฺติ.}

\prose{เตน โข ปน สมเยน ฉพฺพคฺคิยา ภิกฺขู อุจฺเจ มญฺจ\-ปฏิปาทเก ธาเรนฺติ, สห มญฺจ\-ปฏิปาทเกหิ ปเวเธนฺติ. มนุสฺสา วิหารจาริกํ อาหิณฺฑนฺตา ปสฺสิตฺวา อุชฺฌายนฺติ ขิยฺยนฺติ วิปาเจนฺติ “เสยฺยถาปิ คิหี กามโภคิโน”ติ. ภควโต เอตมตฺถํ อาโรเจสุํ. น ภิกฺขเว อุจฺจา มญฺจ\-ปฏิปาทกา ธาเรตพฺพา, โย ธาเรยฺย, อาปตฺติ ทุกฺกฏสฺส. อนุชานามิ ภิกฺขเว อฏฺฐงฺคุล\-ปรมํ มญฺจ\-ปฏิปาทกนฺติ.}

\prose{เตน โข ปน สมเยน สํฆสฺส สุตฺตํ อุปฺปนฺนํ โหติ. ภควโต เอตมตฺถํ อาโรเจสุํ. อนุชานามิ ภิกฺขเว สุตฺตํ มญฺจํ เวเฐตุนฺติ\csromansharedfootnote{e2a20af975cfa09f}{เวตุนฺติ (สี)}. องฺคานิ พหุสุตฺตํ ปริยาทิยนฺติ. ภควโต เอตมตฺถํ อาโรเจสุํ. อนุชานามิ ภิกฺขเว องฺเค วิชฺฌิตฺวา อฏฺฐปทกํ เวเฐตุนฺติ. โจฬกํ อุปฺปนฺนํ โหติ. ภควโต เอตมตฺถํ อาโรเจสุํ. อนุชานามิ ภิกฺขเว จิมิลิกํ กาตุนฺติ. ตูลิกา อุปฺปนฺนาโหติ. ภควโต เอตมตฺถํ อาโรเจสุํ. อนุชานามิ ภิกฺขเว วิชเฏตฺวา พิพฺโพหนํ\csromansharedfootnote{5507a21645e35f7e}{พิมฺโพหนํ (สี, สฺยา, พิมฺพ + โอหนํ)} กาตุํ, ตีณิ ตูลานิ รุกฺขตูลํ ลตาตูลํ โปฏกิตูลนฺติ.}

\prose{เตน โข ปน สมเยน ฉพฺพคฺคิยา ภิกฺขู อทฺธกายิกานิ\csromansharedfootnote{9806e74a681bad6c}{อฑฺฒกายิกานิ (ก)} พิพฺโพหนานิ ธาเรนฺติ. มนุสฺสา วิหารจาริกํ อาหิณฺฑนฺตา ปสฺสิตฺวา อุชฺฌายนฺติ ขิยฺยนฺติ วิปาเจนฺติ “เสยฺยถาปิ คิหี กามโภคิโน”ติ. ภควโต เอตมตฺถํ อาโรเจสุํ. น ภิกฺขเว อฑฺฒกายิกานิ พิพฺโพหนานิ ธาเรตพฺพานิ, โย ธาเรยฺย, อาปตฺติ ทุกฺกฏสฺส. อนุชานามิ ภิกฺขเว สีสปฺปมาณํ พิพฺโพหนํ กาตุนฺติ.}

\prose{เตน โข ปน สมเยน ราชคเห คิรคฺคสมชฺโช โหติ, มนุสฺสา มหามตฺตานํ อตฺถาย ภิสิโย ปฏิยาเทนฺติ อุณฺณภิสิํ โจฬภิสิํ วากภิสิํ ติณภิสิํ ปณฺณภิสิํ, เต วีติวตฺเต สมชฺเช ฉวิํ อุปฺปาเฏตฺวา หรนฺติ. อทฺทสาสุํ โข ภิกฺขู สมชฺชฏฺฐาเน พหุํ อุณฺณมฺปิ โจฬกมฺปิ วากมฺปิ ติณมฺปิ ปณฺณมฺปิ ฉฏฺฏิตํ, ทิสฺวาน ภควโต เอตมตฺถํ \csromanfolio{296}อาโรเจสุํ. อนุชานามิ ภิกฺขเว ปญฺจ ภิสิโย อุณฺณภิสิํ โจฬภิสิํ วากภิสิํ ติณภิสิํ ปณฺณภิสินฺติ.}

\prose{เตน โข ปน สมเยน สํฆสฺส เสนาสน\-ปริกฺขาริกํ ทุสฺสํ อุปฺปนฺนํ โหติ. ภควโต เอตมตฺถํ อาโรเจสุํ. อนุชานามิ ภิกฺขเว ภิสิํ โอนนฺธิตุนฺติ\csromansharedfootnote{1797714b0bd2340c}{โอนทฺธิตุํ (สฺยา)}.}

\prose{เตน โข ปน สมเยน ภิกฺขู มญฺจภิสิํ ปีเฐ สนฺถรนฺติ, ปีฐภิสิํ มญฺเจ สนฺถรนฺติ, ภิสิโย ปริภิชฺชนฺติ. ภควโต เอตมตฺถํ อาโรเจสุํ. อนุชานามิ ภิกฺขเว โอนทฺธมญฺจํ\csromansharedfootnote{de82735017a197bc}{โอนนฺธมญฺจ (ก) เอวมุปริปิ.} โอนทฺธ\-ปีฐนฺติ. อุลฺโลกํ อกริตฺวา สนฺถรนฺติ, เหฏฺฐโต นิปตนฺติ\csromansharedfootnote{116c860929b6fe11}{นิปฏนฺติ (ก), นิปฺผฏนฺติ (สี), นิปฺปาเฏนฺติ (สฺยา)} ฯเปฯ. อนุชานามิ ภิกฺขเว อุลฺโลกํ กริตฺวา สนฺถริตฺวา ภิสิํ โอนนฺธิตุนฺติ. ฉวิํ อุปฺปาเฏตฺวา หรนฺติ ฯเปฯ. อนุชานามิ ภิกฺขเว โผสิตุนฺติ. หรนฺติเยว ฯเปฯ. อนุชานามิ ภิกฺขเว ภตฺติกมฺมนฺติ. หรนฺติเยว ฯเปฯ. อนุชานามิ ภิกฺขเว หตฺถ\-ภตฺติกมฺมนฺติ. หรนฺติเยว ฯเปฯ. อนุชานามิ ภิกฺขเว หตฺถ\-ภตฺตินฺติ.}

\csromanmatikaanchor{mtk.133}
\csromantocmarkref{subsection}{เสตวณฺณาทิอนุชานน}{mtk.133}
\bookmark[dest={mtk.133},level=2]{เสตวณฺณาทิอนุชานน}
\csromanheader{2}{เสตวณฺณาทิอนุชานน}
\proseitem{298}{เตน โข ปน สมเยน ติตฺถิยานํ เสยฺยาโย เสตวณฺณา โหนฺติ, กาฬวณฺณกตา ภูมิ, เครุก\-ปริกมฺม\-กตา ภิตฺติ, พหู มนุสฺสา เสยฺยาเปกฺขกา คจฺฉนฺติ. ภควโต เอตมตฺถํ อาโรเจสุํ. อนุชานามิ ภิกฺขเว วิหาเร เสตวณฺณํ กาฬวณฺณํ เครุก\-ปริกมฺมนฺติ.}

\prose{เตน โข ปน สมเยน ผรุสาย ภิตฺติยา เสตวณฺโณ น นิปตติ. ภควโต เอตมตฺถํ อาโรเจสุํ. อนุชานามิ ภิกฺขเว ถุสปิณฺฑํ ทตฺวา ปาณิกาย ปฏิพาเหตฺวา เสตวณฺณํ นิปาเตตุนฺติ. เสตวณฺโณ อนิพนฺธนีโย โหติ ฯเปฯ. อนุชานามิ ภิกฺขเว สณฺหมตฺติกํ ทตฺวา ปาณิกาย ปฏิพาเหตฺวา เสตวณฺณํ นิปาเตตุนฺติ. เสตวณฺโณ อนิพนฺธนีโย โหติ ฯเปฯ. อนุชานามิ ภิกฺขเว อิกฺกาสํ ปิฏฺฐมทฺทนฺติ.}

\prose{เตน โข ปน สมเยน ผรุสาย ภิตฺติยา เครุกา น นิปตติ. ภควโต เอตมตฺถํ อาโรเจสุํ. อนุชานามิ ภิกฺขเว ถุสปิณฺฑํ ทตฺวา ปาณิกาย ปฏิพาเหตฺวา เครุกํ นิปาเตตุนฺติ. เครุกา อนิพนฺธนียา \csromanfolio{297}โหติ ฯเปฯ. อนุชานามิ ภิกฺขเว กุณฺฑกมตฺติกํ ทตฺวา ปาณิกาย ปฏิพาเหตฺวา เครุกํ นิปาเตตุนฺติ. เครุกา อนิพนฺธนียา โหติ ฯเปฯ. อนุชานามิ ภิกฺขเว สาสปกุฏฺฏํ สิตฺถ\-เตลกนฺติ. อจฺจุสฺสนฺนํ โหติ ฯเปฯ. อนุชานามิ ภิกฺขเว โจฬเกน ปจฺจุทฺธริตุนฺติ.}

\prose{เตน โข ปน สมเยน ผรุสาย ภูมิยา กาฬวณฺโณ น นิปตติ. ภควโต เอตมตฺถํ อาโรเจสุํ. อนุชานามิ ภิกฺขเว ถุสปิณฺฑํ ทตฺวา ปาณิกาย ปฏิพาเหตฺวา กาฬวณฺณํ นิปาเตตุนฺติ. กาฬวณฺโณ อนิพนฺธนีโย โหติ ฯเปฯ. อนุชานามิ ภิกฺขเว คณฺฑุมตฺติกํ ทตฺวา ปาณิกาย ปฏิพาเหตฺวา กาฬวณฺณํ นิปาเตตุนฺติ. กาฬวณฺโณ อนิพนฺธนีโย โหติ ฯเปฯ. อนุชานามิ ภิกฺขเว อิกฺกาสํ กสาวนฺติ.}

\csromanmatikaanchor{mtk.134}
\csromantocmarkref{subsection}{ปฏิภานจิตฺตปฏิกฺเขป}{mtk.134}
\bookmark[dest={mtk.134},level=2]{ปฏิภานจิตฺตปฏิกฺเขป}
\csromanheader{2}{ปฏิภานจิตฺต\-ปฏิกฺเขป}
\proseitem{299}{เตน โข ปน สมเยน ฉพฺพคฺคิยา ภิกฺขู วิหาเร ปฏิภาน\-จิตฺตํ การาเปนฺติ อิตฺถิรูปกํ ปุริสรูปกํ. มนุสฺสา วิหารจาริกํ อาหิณฺฑนฺตา ปสฺสิตฺวา อุชฺฌายนฺติ ขิยฺยนฺติ วิปาเจนฺติ “เสยฺยถาปิ คิหี กามโภคิโน”ติ. ภควโต เอตมตฺถํ อาโรเจสุํ. น ภิกฺขเว ปฏิภาน\-จิตฺตํ การาเปตพฺพํ อิตฺถิรูปกํ ปุริสรูปกํ, โย การาเปยฺย, อาปตฺติ ทุกฺกฏสฺส. อนุชานามิ ภิกฺขเว มาลากมฺมํ ลตากมฺมํ มกรทนฺตกํ ปญฺจปฏิกนฺติ.}

\csromanmatikaanchor{mtk.135}
\csromantocmarkref{subsection}{อิฏฺฐกาจยาทิอนุชานน}{mtk.135}
\bookmark[dest={mtk.135},level=2]{อิฏฺฐกาจยาทิอนุชานน}
\csromanheader{2}{อิฏฺฐกาจยา\-ทิ\-อนุชานน}
\proseitem{300}{เตน โข ปน สมเยน วิหารา นีจวตฺถุกา โหนฺติ, อุทเกน โอตฺถริยฺยนฺติ. ภควโต เอตมตฺถํ อาโรเจสุํ. อนุชานามิ ภิกฺขเว อุจฺจวตฺถุกํ กาตุนฺติ. จโย ปริปตติ ฯเปฯ. อนุชานามิ ภิกฺขเว จินิตุํ ตโย จเย อิฏฺฐกาจยํ สิลาจยํ ทารุจยนฺติ. อาโรหนฺตา วิหญฺญนฺติ ฯเปฯ. อนุชานามิ ภิกฺขเว ตโย โสปาเน อิฏฺฐกาโสปานํ สิลาโสปานํ ทารุโสปานนฺติ. อาโรหนฺตา ปริปตนฺติ ฯเปฯ. อนุชานามิ ภิกฺขเว อาลมฺพนพาหนฺติ.}

\prose{เตน โข ปน สมเยน วิหารา อาฬกมนฺทา โหนฺติ, ภิกฺขู หิริยนฺติ นิปชฺชิตุํ. ภควโต เอตมตฺถํ อาโรเจสุํ. อนุชานามิ ภิกฺขเว ติโรกรณินฺติ. ติโรกรณิํ อุกฺขิปิตฺวา โอโลเกนฺติ ฯเปฯ. อนุชานามิ \csromanfolio{298}ภิกฺขเว อฑฺฒ\-กุฏฺฏกนฺติ, อฑฺฒกุฏฺฏกา อุปริโต โอโลเกนฺติ ฯเปฯ. อนุชานามิ ภิกฺขเว ตโย คพฺเภ สิวิกาคพฺภํ นาฬิกาคพฺภํ หมฺมิย\-คพฺภนฺติ.}

\prose{เตน โข ปน สมเยน ภิกฺขู ขุทฺทเก วิหาเร มชฺเฌ คพฺภํ กโรนฺติ, อุปจาโร น โหติ. ภควโต เอตมตฺถํ อาโรเจสุํ. อนุชานามิ ภิกฺขเว ขุทฺทเก วิหาเร เอกมนฺตํ คพฺภํ กาตุํ, มหลฺลเก มชฺเฌติ.}

\prose{เตน โข ปน สมเยน วิหารสฺส กุฏฺฏปาโท ชีรติ. ภควโต เอตมตฺถํ อาโรเจสุํ. อนุชานามิ ภิกฺขเว กุลงฺกปาทกนฺติ\csromansharedfootnote{97bc192b97efb57c}{กุฬุงฺกปาทกนฺติ (สี)}. วิหารสฺส กุฏฺโฏ โอวสฺสติ ฯเปฯ. อนุชานามิ ภิกฺขเว ปริตฺตาณ\-กิฏิกํ อุทฺทสุธนฺติ\csromansharedfootnote{c3ce05a2ac290218}{อุทฺธาสุธนฺติ (สฺยา)}.}

\prose{เตน โข ปน สมเยน อญฺญตรสฺส ภิกฺขุโน ติณจฺฉทนา อหิ ขนฺเธ ปตติ, โส ภีโต วิสฺสรมกาสิ. ภิกฺขู อุปธาวิตฺวา ตํ ภิกฺขุํ เอตทโวจุํ “กิสฺส ตฺวํ อาวุโส วิสฺสรมกาสี”ติ. อถ โข โส ภิกฺขูนํ เอตมตฺถํ อาโรเจสิ. ภิกฺขู ภควโต เอตมตฺถํ อาโรเจสุํ. อนุชานามิ ภิกฺขเว วิตานนฺติ.}

\prose{เตน โข ปน สมเยน ภิกฺขู มญฺจปาเทปิ ปีฐปาเทปิ ถวิกาโย ลคฺเคนฺติ, อุนฺทูเรหิปิ อุปจิกาหิปิ ขชฺชนฺติ. ภควโต เอตมตฺถํ อาโรเจสุํ. อนุชานามิ ภิกฺขเว ภิตฺติขิลํ นาคทนฺตกนฺติ.}

\prose{เตน โข ปน สมเยน ภิกฺขู มญฺเจปิ ปีเฐปิ จีวรํ นิกฺขิปนฺติ, จีวรํ ปริภิชฺชิ. ภควโต เอตมตฺถํ อาโรเจสุํ. อนุชานามิ ภิกฺขเว วิหาเร จีวรวํสํ จีวรรชฺชุนฺติ.}

\prose{เตน โข ปน สมเยน วิหารา อนาฬินฺทกา โหนฺติ อปฺปฏิสฺสรณา. ภควโต เอตมตฺถํ อาโรเจสุํ. อนุชานามิ ภิกฺขเว อาฬินฺทํ ปฆนํ ปกุฏฺฏํ\csromansharedfootnote{1cb80b8011dc6ecf}{ปกุฑฺฑํ (สี)} โอสารกนฺติ. อาฬินฺทา ปากฏา โหนฺติ, ภิกฺขู หิริยนฺติ นิปชฺชิตุํ ฯเปฯ. อนุชานามิ ภิกฺขเว สํสรณ\-กิฏิกํ อุคฺฆาฏนกิฏิกนฺติ.}

\csromanmatikaanchor{mtk.136}
\csromantocmarkref{subsection}{อุปฏฺฐานสาลาอนุชานน}{mtk.136}
\bookmark[dest={mtk.136},level=2]{อุปฏฺฐานสาลาอนุชานน}
\csromanheader{2}{อุปฏฺฐานสาลาอนุชานน}
\proseitem{301}{เตน โข ปน สมเยน ภิกฺขู อชฺโฌกาเส ภตฺตวิสฺสคฺคํ กโรนฺตา สีเตนปิ อุเณฺหนปิ กิลมนฺติ. ภควโต เอตมตฺถํ อาโรเจสุํ. อนุชานามิ ภิกฺขเว อุปฏฺฐานสาลา นีจวตฺถุกา \csromanfolio{299}โหติ, อุทเกน โอตฺถริยฺยติ ฯเปฯ. อนุชานามิ ภิกฺขเว อุจฺจวตฺถุกํ กาตุนฺติ. จโย ปริปตติ ฯเปฯ. อนุชานามิ ภิกฺขเว จินิตุํ ตโย จเย อิฏฺฐกาจยํ สิลาจยํ ทารุจยนฺติ. อาโรหนฺตา วิหญฺญนฺติ ฯเปฯ. อนุชานามิ ภิกฺขเว ตโย โสปาเน อิฏฺฐกาโสปานํ สิลาโสปานํ ทารุโสปานนฺติ. อาโรหนฺตา ปริปตนฺติ ฯเปฯ. อนุชานามิ ภิกฺขเว อาลมฺพนพาหนฺติ. อุปฏฺฐาน\-สาลาย ติณจุณฺณํ ปริปตติ ฯเปฯ. อนุชานามิ ภิกฺขเว โอคุมฺเผตฺวา อุลฺลิตฺตา\-วลิตฺตํ กาตุํ เสตวณฺณํ กาฬวณฺณํ เครุก\-ปริกมฺมํ มาลากมฺมํ ลตากมฺมํ มกรทนฺตกํ ปญฺจปฏิกํ จีวรวํสํ จีวรรชฺชุนฺติ.}

\prose{เตน โข ปน สมเยน ภิกฺขู อชฺโฌกาเส ฉมาย จีวรํ ปตฺถรนฺติ, จีวรํ ปํสุกิตํ โหติ. ภควโต เอตมตฺถํ อาโรเจสุํ. อนุชานามิ ภิกฺขเว อชฺโฌกาเส จีวรวํสํ จีวรรชฺชุนฺติ. ปานียํ โอตปฺปติ ฯเปฯ. อนุชานามิ ภิกฺขเว ปานิยสาลํ ปานีย\-มณฺฑปนฺติ. ปานียสาลา นีจวตฺถุกา โหติ, อุทเกน โอตฺถริยฺยติ ฯเปฯ. อนุชานามิ ภิกฺขเว อุจฺจวตฺถุกํ กาตุนฺติ. จโย ปริปตติ ฯเปฯ. อนุชานามิ ภิกฺขเว จินิตุํ ตโย จเย อิฏฺฐกาจยํ สิลาจยํ ทารุจยนฺติ. อาโรหนฺตา วิหญฺญนฺติ ฯเปฯ. อนุชานามิ ภิกฺขเว ตโย โสปาเน อิฏฺฐกาโสปานํ สิลาโสปานํ ทารุโสปานนฺติ. อาโรหนฺตา ปริปตนฺติ ฯเปฯ. อนุชานามิ ภิกฺขเว อาลมฺพนพาหนฺติ. ปานียสาลาย ติณจุณฺณํ ปริปตติ ฯเปฯ. อนุชานามิ ภิกฺขเว โอคุมฺเผตฺวา อุลฺลิตฺตา\-วลิตฺตํ กาตุํ เสตวณฺณํ กาฬวณฺณํ เครุก\-ปริกมฺมํ มาลากมฺมํ ลตากมฺมํ มกรทนฺตกํ ปญฺจปฏิกํ จีวรวํสํ จีวรรชฺชุนฺติ. ปานียภาชนํ น สํวิชฺชติ ฯเปฯ. อนุชานามิ ภิกฺขเว ปานิยสงฺขํ ปานิยสราวกนฺติ.}

\csromanmatikaanchor{mtk.137}
\csromantocmarkref{subsection}{ปาการาทิอนุชานน}{mtk.137}
\bookmark[dest={mtk.137},level=2]{ปาการาทิอนุชานน}
\csromanheader{2}{ปาการาทิอนุชานน}
\proseitem{302}{เตน โข ปน สมเยน วิหารา อปริกฺขิตฺตา หนฺติ. ภควโต เอตมตฺถํ อาโรเจสุํ. อนุชานามิ ภิกฺขเว ปริกฺขิปิตุํ ตโย ปากาเร อิฏฺฐกาปาการํ สิลาปาการํ ทารุปาการนฺติ. โกฏฺฐโก น โหติ ฯเปฯ. อนุชานามิ ภิกฺขเว โกฏฺฐกนฺติ. โกฏฺฐโก นีจวตฺถุโก โหติ, อุทเกน โอตฺถริยฺยติ ฯเปฯ. อนุชานามิ ภิกฺขเว อุจฺจวตฺถุกํ กาตุนฺติ. โกฏฺฐกสฺส กวาฏํ น โหติ ฯเปฯ. อนุชานามิ ภิกฺขเว กวาฏํ ปิฏฺฐสํฆาฏํ อุทุกฺขลิกํ อุตฺตรปาสกํ อคฺคฬวฏฺฏิํ กปิสีสกํ สูจิกํ ฆฏิกํ \csromanfolio{300}ตาฬจฺฉิทฺทํ อาวิญฺฉนจฺฉิทํ อาวิญฺฉน\-รชฺชุนฺติ. โกฏฺฐเก ติณจุณฺณํ ปริปตติ ฯเปฯ. อนุชานามิ ภิกฺขเว โอคุมฺเผตฺวา อุลฺลิตฺตา\-วลิตฺตํ กาตุํ เสตวณฺณํ กาฬวณฺณํ เครุก\-ปริกมฺมํ มาลากมฺมํ ลตากมฺมํ มกรทนฺตกํ ปญฺจปฏิกนฺติ.}

\prose{เตน โข ปน สมเยน ปริเวณํ จิกฺขลฺลํ โหติ. ภควโต เอตมตฺถํ อาโรเจสุํ. อนุชานามิ ภิกฺขเว มรุมฺพํ อุปกิริตุนฺติ. น ปริยาปุณนฺติ ฯเปฯ. อนุชานามิ ภิกฺขเว ปทรสิลํ นิกฺขิปิตุนฺติ. อุทกํ สนฺติฏฺฐติ ฯเปฯ. อนุชานามิ ภิกฺขเว อุทกนิทฺธมนนฺติ.}

\prose{เตน โข ปน สมเยน ภิกฺขู ปริเวเณ ตหํ ตหํ อคฺคิฏฺฐานํ กโรนฺติ, ปริเวณํ อุกฺลาปํ โหติ. ภควโต เอตมตฺถํ อาโรเจสุํ. อนุชานามิ ภิกฺขเว เอกมนฺตํ อคฺคิสาลํ กาตุนฺติ. อคฺคิสาลา นีจวตฺถุกา โหติ, อุทเกน โอตฺถริยฺยติ ฯเปฯ. อนุชานามิ ภิกฺขเว อุจฺจวตฺถุกํ กาตุนฺติ. จโย ปริปตติ ฯเปฯ. อนุชานามิ ภิกฺขเว จินิตุํ ตโย จเย อิฏฺฐกาจยํ สิลาจยํ ทารุจยนฺติ. อาโรหนฺตา วิหญฺญนฺติ ฯเปฯ. อนุชานามิ ภิกฺขเว ตโย โสปาเน อิฏฺฐกาโสปานํ สิลาโสปานํ ทารุโสปานนฺติ. อาโรหนฺตา ปริปตนฺติ ฯเปฯ. อนุชานามิ ภิกฺขเว อาลมฺพนพาหนฺติ. อคฺคิสาลาย กวาฏํ น โหติ ฯเปฯ. อนุชานามิ ภิกฺขเว กวาฏํ ปิฏฺฐสํฆาฏํ อุทุกฺขลิกํ อุตฺตรปาสกํ อคฺคฬวฏฺฏิํ กปิสีสกํ สูจิกํ ฆฏิกํ ตาฬจฺฉิทฺทํ อาวิญฺฉน\-จฺฉิทฺทํ อาวิญฺฉน\-รชฺชุนฺติ. อคฺคิสาลาย ติณจุณฺณํ ปริปตติ ฯเปฯ. อนุชานามิ ภิกฺขเว โอคุมฺเผตฺวา อุลฺลิตฺตา\-วลิตฺตํ กาตุํ เสตวณฺณํ กาฬวณฺณํ เครุก\-ปริกมฺมํ มาลากมฺมํ ลตากมฺมํ มกรทนฺตกํ ปญฺจปฏิกํ ปญฺจปฏิกํ จีวรวํสํ จีวรรชฺชุนฺติ.}

\csromanmatikaanchor{mtk.138}
\csromantocmarkref{subsection}{อารามปริกฺเขปอนุชานน}{mtk.138}
\bookmark[dest={mtk.138},level=2]{อารามปริกฺเขปอนุชานน}
\csromanheader{2}{อารามปริกฺเขปอนุชานน}
\proseitem{303}{เตน โข ปน สมเยน อาราโม อปริกฺขิตฺโต โหติ, อชกาปิ ปสุกาปิ อุปโรเป วิเหเฐนฺติ. ภควโต เอตมตฺถํ อาโรเจสุํ. อนุชานามิ ภิกฺขเว ปริกฺขิปิตุํ ตโย วาเฏ เวฬุวาฏํ กณฺฑกวาฏํ\csromansharedfootnote{77d3b64232cb6f62}{วเฏ เวฬุวฏํ กณฺฑกวฏํ (สฺยา)} ปริกฺขนฺติ. โกฏฺฐโก น โหติ. ตเถว อชกาปิ ปสุกาปิ อุปโรเป วิเหเฐนฺติ ฯเปฯ. อนุชานามิ ภิกฺขเว โกฏฺฐกํ อเปสิํ ยมกกวาฏํ โตรณํ ปลิฆนฺติ. โกฏฺฐเก ติณจุณฺณํ ปริปตติ ฯเปฯ. \csromanfolio{301}อนุชานามิ ภิกฺขเว โอคุมฺเผตฺวา อุลฺลิตฺตา\-วลิตฺตํ กาตุํ เสตวณฺณํ กาฬวณฺณํ เครุก\-ปริกมฺมํ มาลากมฺมํ ลตากมฺมํ มกรทนฺตกํ ปญฺจปฏิกนฺติ. อาราโม จิกฺขลฺโล โหติ ฯเปฯ. อนุชานามิ ภิกฺขเว มรุมฺพํ อุปกิริตุนฺติ. น ปริยาปุณนฺติ ฯเปฯ. อนุชานามิ ภิกฺขเว ปทรสิลํ นิกฺขิปิตุนฺติ. อุทกํ สนฺติฏฺฐติ ฯเปฯ. อนุชานามิ ภิกฺขเว อุทกนิทฺธมนนฺติ.}

\prosewithcloser{เตน โข ปน สมเยน ราชา มาคโธ เสนิโย พิมฺพิสาโร สํฆสฺส อตฺถาย สุธา\-มตฺติกา\-เลปนํ ปาสาทํ กาเรตุกาโม โหติ. อถ โข ภิกฺขูนํ เอตทโหสิ “กิํ นุ โข ภควตา ฉทนํ อนุญฺญาตํ, กิํ อนนุญฺญาตนฺ”ติ. ภควโต เอตมตฺถํ อาโรเจสุํ. อนุชานามิ ภิกฺขเว ปญฺจ ฉทนานิ อิฏฺฐกา\-ฉทนํ สิลาฉทนํ สุธาฉทนํ ติณจฺฉทนํ ปณฺณ\-จฺฉทนนฺติ.}{ปฐม\-ภาณวาโร นิฏฺฐิโต.}
\csromanmatikaanchor{mtk.139}
\csromantocmarkref{section}{2. ทุติยภาณวาร}{mtk.139}
\bookmark[dest={mtk.139},level=1]{2. ทุติยภาณวาร}
\csromanheader{1}{2. ทุติย\-ภาณวาร}
\csromanmatikaanchor{mtk.140}
\csromantocmarkref{section}{อนาถปิณฺฑิกวตฺถุ}{mtk.140}
\bookmark[dest={mtk.140},level=1]{อนาถปิณฺฑิกวตฺถุ}
\csromanheader{1}{อนาถปิณฺฑิก\-วตฺถุ}
\proseitem{304}{เตน โข ปน สมเยน อนาถปิณฺฑิโก คหปติ ราชคหกสฺส เสฏฺฐิสฺส ภคินิปติโก โหติ. อถ โข อนาถปิณฺฑิโก คหปติ ราชคหํ อคมาสิ เกนจิเทว กรณีเยน. เตน โข ปน สมเยน ราชคหเกน เสฏฺฐินา สฺวาตนาย พุทฺธปฺปมุโขสํโฆ นิมนฺติโต โหติ. อถ โข ราชคหโก เสฏฺฐี ทาเส จ กมฺมกาเร\csromansharedfootnote{a1ff5f71d6961de6}{กมฺมกเร (สี, สฺยา)} จ อาณาเปสิ “เตน หิ ภเณ กาลสฺเสว อุฏฺฐาย ยาคุโย ปจถ, ภตฺตานิ ปจถ, สูปานิ สมฺปาเทถ, อุตฺตริภงฺคานิ สมฺปาเทถา”ติ. อถ โข อนาถ\-ปิณฺฑิกสฺส คหปติสฺส เอตทโหสิ “ปุพฺเพ ขฺวายํ คหปติ มยิ อาคเต สพฺพกิจฺจานิ นิกฺขิปิตฺวา มมญฺเญว สทฺธิํ ปฏิสมฺโมทติ, โสทานายํ วิกฺขิตฺตรูโป ทาเส จ กมฺมกาเร จ อาณาเปสิ ‘เตน หิ ภเณ กาลสฺเสว อุฏฺฐาย ยาคุโย ปจถ, ภตฺตานิ ปจถ, สูปานิ สมฺปาเทถ, อุตฺตริภงฺคานิ สมฺปาเทถา’ติ. กิํ นุ โข อิมสฺส คหปติสฺส อาวาโห วา ภวิสฺสกิ, วิวาโห วา ภวิสฺสติ, มหายญฺโญ วา \csromanfolio{302}ปจฺจุปฏฺฐิโต, ราชา วา มาคโธ เสนิโย พิมฺพิสาโร นิมนฺติโต สฺวาตนาย สทฺธิํ พลกาเยนา”ติ.}

\prose{อถ โข ราชคหโก เสฏฺฐี ทาเส จ กมฺมกาเร จ อาณาเปตฺวา เยน อนาถปิณฺฑิโก คหปติ เตนุปสงฺกมิ, อุปสงฺกมิตฺวา อนาถ\-ปิณฺฑิเกน คหปตินา สทฺธิํ ปฏิสมฺโมทิตฺวา เอกมนฺตํ นิสีทิ. เอกมนฺตํ นิสินฺนํ โข ราชคหกํ เสฏฺฐิํ อนาถปิณฺฑิโก คหปติ เอตทโวจ “ปุพฺเพ โข ตฺวํ คหปติ มยิ อาคเต สพฺพกิจฺจานิ นิกฺขิปิตฺวา มมญฺเญว สทฺธิํ ปฏิสมฺโมทสิ, โสทานิ ตฺวํ วิกฺขิตฺตรูโป ทาเส จ กมฺมกาเร จ อาณาเปสิ ‘เตน หิ ภเณ กาลสฺเสว อุฏฺฐาย ยาคุโย ปจถ, ภตฺตานิ ปจถ, สูปานิ สมฺปาเทถ, อุตฺตริภงฺคานิ สมฺปาเทถา’ติ. กิํ นุ โข เต คหปติ อาวาโห วา ภวิสฺสติ, วิวาโห วา ภวิสฺสติ, มหายญฺโญ วา ปจฺจุปฏฺฐิโต, ราชา วา มาคโธ เสนิโย พิมฺพิสาโร นิมนฺติโต สฺวาตนาย สทฺธิํ พลกาเยนา”ติ. น เม คหปติ อาวาโห วา ภวิสฺสติ, นาปิ วิวาโห วา ภวิสฺสติ. นาปิ ราชา วา มาคโธ เสนิโย พิมฺพิสาโร นิมนฺติโต สฺวาตนาย สทฺธิํ พลกาเยน. อปิ จ เม มหายญฺโญ ปจฺจุปฏฺฐิโต สฺวาตนาย พุทฺธปฺปมุโข สํโฆ นิมนฺติโตติ. “พุทฺโธ”ติ ตฺวํ คหปติ วเทสีติ. “พุทฺโธ”ตฺยาหํ คหปติ วทามีติ. “พุทฺโธ”ติ ตฺวํ คหปติ วเทสีติ. “พุทฺโธ”ตฺยาหํ คหปติ วทามีติ. “พุทฺโธ”ติ ตฺวํ คหปติ วเทสีติ. “พุทฺโธ”ตฺยาหํ คหปติ วทามีติ. โฆโสปิ โข เอโส คหปติ ทุลฺลโภ โลกสฺมิํ, ยทิทํ “พุทฺโธ พุทฺโธ”ติ. สกฺกา นุ โข คหปติ อิมํ กาลํ ตํ ภควนฺตํ ทสฺสนาย อุปสงฺกมิตุํ อรหนฺตํ สมฺมา\-สมฺพุทฺธนฺติ. อกาโล โข คหปติ อิมํ กาลํ ตํ ภควนฺตํ ทสฺสนาย อุปสงฺกมิตุํ อรหนฺตํ สมฺมา\-สมฺพุทฺธํ, เสฺวทานิ ตฺวํ กาเลน ตํ ภควนฺตํ ทสฺสนาย อุปสงฺกมิสฺสสิ อรหนฺตํ สมฺมา\-สมฺพุทฺธนฺติ. อถ โข อนาถปิณฺฑิโก คหปติ “เสฺวทานาหํ กาเลน ตํ ภควนฺตํ ทสฺสนาย อุปสงฺกมิสฺสามิ อรหนฺตํ สมฺมาสมฺพุทฺธนฺ”ติ พุทฺธคตาย สติยา นิปชฺชิตฺวา รตฺติยา สุทฺทํ ติกฺขตฺตุํ วุฏฺฐาสิ ปภาตํ มญฺญมาโน.}

\proseitem{305}{\csromansymbolfootnote{*}{สํ 1. 213 ปิฏฺเฐปิ.}อถ โข อนาถปิณฺฑิโก คหปติ เยน สิวกทฺวารํ\csromansharedfootnote{44c926a446a16441}{สีวทฺวารํ (สี), สีตวนทฺวารํ (สฺยา)} เตนุปสงฺกมิ. อมนุสฺสา ทฺวารํ วิวริํสุ. อถ โข อนาถ\-ปิณฺฑิกสฺส \csromanfolio{303}คหปติสฺส นครมฺหา นิกฺขนฺตสฺส อาโลโก อนฺตรธายิ, อนฺธกาโร ปาตุรโหสิ, ภยํ ฉมฺภิตตฺตํ โลมหํโส อุทปาทิ, ตโตว ปุน นิวตฺติตุกาโม อโหสิ. อถ โข สิวโก\csromansharedfootnote{e64746f0e825f1db}{สีวโก (สี, สฺยา)} ยกฺโข อนฺตรหิโต สทฺทมนุสฺสเวสิ\csromandash{}}

\setlength{\csromangathapagewidth}{0pt}
\setlength{\csromangathawakwidth}{0pt}
\csromangathameasuregroup{“สตํ หตฺถี สตํ อสฺสา, \\ สตํ กญฺญาสหสฺสานิ, \\ เอกสฺส ปทวีติหารสฺส,}{\csromangathabat{“สตํ หตฺถี สตํ อสฺสา,}{สตํ อสฺสตรีรถา.} \\ \csromangathabat{สตํ กญฺญาสหสฺสานิ,}{อามุกฺกมณิกุณฺฑลา.} \\ \csromangathabat{เอกสฺส ปทวีติหารสฺส,}{กลํ นาคฺฆนฺติ โสฬสิํ\textsuperscript{1}.} \\ อภิกฺกม คหปติ อภิกฺกม คหปติ, \\ อภิกฺกนฺตํ เต เสยฺโย โน ปฏิกฺกนฺตนฺ”ติ.}
\csromangathasetleft{“สตํ หตฺถี สตํ อสฺสา, \\ สตํ กญฺญาสหสฺสานิ, \\ เอกสฺส ปทวีติหารสฺส,}
\csromangathagroup{\csromangathastanza{\csromangathabat{“สตํ หตฺถี สตํ อสฺสา,}{สตํ อสฺสตรีรถา.} \\ \csromangathabat{สตํ กญฺญาสหสฺสานิ,}{อามุกฺกมณิกุณฺฑลา.}}\csromangathastanza{\csromangathabat{เอกสฺส ปทวีติหารสฺส,}{กลํ นาคฺฆนฺติ โสฬสิํ\csromansharedfootnote{32ab573c8365650b}{โสฬสินฺติ (สี, ก)}.} \\ อภิกฺกม คหปติ อภิกฺกม คหปติ, \\ อภิกฺกนฺตํ เต เสยฺโย โน ปฏิกฺกนฺตนฺ”ติ.}}

\prose{อถ โข อนาถ\-ปิณฺฑิกสฺส คหปติสฺส อนฺธกาโร อนฺตรธายิ, อาโลโก ปาตุรโหสิ, ยํ อโหสิ ภยํ ฉมฺภิตตฺตํ โลมหํโส, โส ปฏิปฺปสฺสมฺภิ. ทุติยมฺปิ โข ฯเปฯ. ตติยมฺปิ โข อนาถ\-ปิณฺฑิกสฺส คหปติสฺส อาโลโก อนฺตรธายิ, อนฺธกาโร ปาตุรโหสิ, ภยํ ฉมฺภิตตฺตํ โลมหํโส อุทปาทิ, ตโตว ปุน นิวตฺติตุกาโม อโหสิ. ตติยมฺปิ โข สิวโก ยกฺโข อนฺตรหิโต สทฺทม\-นุสฺสาเวสิ\csromandash{}}

\setlength{\csromangathapagewidth}{0pt}
\setlength{\csromangathawakwidth}{0pt}
\csromangathameasuregroup{“สตํ หตฺถี สตํ อสฺสา, \\ สตํ กญฺญาสหสฺสานิ, \\ เอกสฺส ปทวีติหารสฺส,}{\csromangathabat{“สตํ หตฺถี สตํ อสฺสา,}{สตํ อสฺสตรีรถา.} \\ \csromangathabat{สตํ กญฺญาสหสฺสานิ,}{อามุกฺกมณิกุณฺฑลา.} \\ \csromangathabat{เอกสฺส ปทวีติหารสฺส,}{กลํ นาคฺฆนฺติ โสฬสิํ.} \\ อภิกฺกม คหปติ อภิกฺกม คหปติ, \\ อภิกฺกนฺตํ เต เสยฺโย โน ปฏิกฺกนฺตนฺ”ติ.}
\csromangathasetleft{“สตํ หตฺถี สตํ อสฺสา, \\ สตํ กญฺญาสหสฺสานิ, \\ เอกสฺส ปทวีติหารสฺส,}
\csromangathagroup{\csromangathastanza{\csromangathabat{“สตํ หตฺถี สตํ อสฺสา,}{สตํ อสฺสตรีรถา.} \\ \csromangathabat{สตํ กญฺญาสหสฺสานิ,}{อามุกฺกมณิกุณฺฑลา.}}\csromangathastanza{\csromangathabat{เอกสฺส ปทวีติหารสฺส,}{กลํ นาคฺฆนฺติ โสฬสิํ.} \\ อภิกฺกม คหปติ อภิกฺกม คหปติ, \\ อภิกฺกนฺตํ เต เสยฺโย โน ปฏิกฺกนฺตนฺ”ติ.}}

\prose{ตติยมฺปิ โข อนาถ\-ปิณฺฑิกสฺส คหปติสฺส อนฺธกาโร อนฺตรธายิ, อาโลโก ปาตุรโหสิ, ยํ อโหสิ ภยํ ฉมฺภิตตฺตํ โลมหํโส, โส ปฏิปฺปสฺสมฺภิ. อถ โข อนาถปิณฺฑิโก คหปติ เยน สีตวนํ เตนุปสงฺกมิ. เตน โข ปน สมเยน ภควา รตฺติยา ปจฺจูสสมยํ ปจฺจุฏฺฐาย อชฺโฌกาเส จงฺกมติ. อทฺทสา โข ภควา อนาถปิณฺฑิกํ คหปติํ ทูรโตว อาคจฺฉนฺตํ, ทิสฺวาน จงฺกมา โอโรหิตฺวา ปญฺญตฺเต อาสเน นิสีทิ, นิสชฺช โข ภควา อนาถปิณฺฑิกํ \csromanfolio{304}คหปติํ เอตทโวจ “เอหิ สุทตฺตา”ติ. อถ โข อนาถปิณฺฑิโก คหปติ “นาเมน มํ ภควา อาลปตี”ติ หฏฺโฐ อุทคฺโค เยน ภควา เตนุปสงฺกมิ, อุปสงฺกมิตฺวา ภควโต ปาเทสุ สิรสา นิปติตฺวา ภควนฺตํ เอตทโวจ ‘กจฺจิ ภนฺเต ภควา สุขํ สยิตฺถา”ติ.}

\setlength{\csromangathapagewidth}{0pt}
\setlength{\csromangathawakwidth}{0pt}
\csromangathameasuregroup{\textsuperscript{*}สพฺพทา เว สุขํ เสติ, \\ โย น ลิมฺปติ กาเมสุ, \\ สพฺพา อาสตฺติโย เฉตฺวา, \\ อุปสนฺโต สุขํ เสติ,}{\csromangathabat{\textsuperscript{*}สพฺพทา เว สุขํ เสติ,}{พฺราหฺมโณ ปรินิพฺพุโต.} \\ \csromangathabat{โย น ลิมฺปติ กาเมสุ,}{สีติภูโต นิรูปธิ.} \\ \csromangathabat{สพฺพา อาสตฺติโย เฉตฺวา,}{วิเนยฺย หทเย ทรํ.} \\ \csromangathabat{อุปสนฺโต สุขํ เสติ,}{สนฺติํ ปปฺปุยฺย เจตสาติ\textsuperscript{1}.}}
\csromangathasetleft{\textsuperscript{*}สพฺพทา เว สุขํ เสติ, \\ โย น ลิมฺปติ กาเมสุ, \\ สพฺพา อาสตฺติโย เฉตฺวา, \\ อุปสนฺโต สุขํ เสติ,}
\csromangathagroup{\csromangathastanza{\csromangathabat{\csromansymbolfootnote{*}{สํ 1. 214 ปิฏฺเฐปิ.}สพฺพทา เว สุขํ เสติ,}{พฺราหฺมโณ ปรินิพฺพุโต.} \\ \csromangathabat{โย น ลิมฺปติ กาเมสุ,}{สีติภูโต นิรูปธิ.}}\csromangathastanza{\csromangathabat{สพฺพา อาสตฺติโย เฉตฺวา,}{วิเนยฺย หทเย ทรํ.} \\ \csromangathabat{อุปสนฺโต สุขํ เสติ,}{สนฺติํ ปปฺปุยฺย เจตสาติ\csromansharedfootnote{0d3e6fef1705f05b}{เจตโสติ (สี, สฺยา)}.}}}

\prose{อถ โข ภควา อนาถ\-ปิณฺฑิกสฺส คหปติสฺส อนุปุพฺพิํ กถํ\csromansharedfootnote{ecd5bf4b9bb6b992}{อานุปุพฺพิกถํ (สี)} กเถสิ, เสยฺยถิทํ, ทานกถํ สีลกถํ สคฺคกถํ กามานํ อาทีนวํ โอการํ สํกิเลสํ เนกฺขมฺเม อานิสํสํ ปกาเสสิ. ยทา ภควา อญฺญาสิ อนาถปิณฺฑิกํ คหปติํ กลฺลจิตฺตํ มุทุจิตฺตํ วินีวรณ\-จิตฺตํ อุทคฺคจิตฺตํ ปสนฺนจิตฺตํ, อถ ยา พุทฺธานํ สามุกฺกํสิกา ธมฺมเทสนา, ตํ ปกาเสสิ ทุกฺขํ สมุทยํ นิโรธํ มคฺคํ, เสยฺยถาปิ นาม สุทฺธํ วตฺถํ อปคตกาฬกํ สมฺมเทว รชนํ ปฏิคฺคเณฺหยฺย, เอวเมว อนาถ\-ปิณฺฑิกสฺส คหปติสฺส ตสฺมิํเยว อาสเน วิรชํ วีตมลํ ธมฺมจกฺขุํ อุทปาทิ “ยํ กิญฺจิ สมุทย\-ธมฺมํ สพฺพํ ตํ นิโรธธมฺมนฺ”ติ. อถ โข อนาถปิณฺฑิโก คหปติ ทิฏฺฐธมฺโม ปตฺตธมฺโม วิทิตธมฺโม ปริโยคาฬฺห\-ธมฺโม ติณฺณ\-วิจิกิจฺโฉ วิคต\-กถํกโถ เวสารชฺช\-ปฺปตฺโต อปรปฺปจฺจโย สตฺถุสาสเน ภควนฺตํ เอตทโวจ “อภิกฺกนฺตํ ภนฺเต, อภิกฺกนฺตํ ภนฺเต, เสยฺยถาปิ ภนฺเต นิกฺกุชฺชิตํ วา อุกฺกุชฺเชยฺย, ปฏิจฺฉนฺนํ วา วิวเรยฺย, มูฬฺหสฺส วา มคฺคํ อาจิกฺเขยฺย, อนฺธกาเร วา เตลปชฺโชตํ ธาเรยฺย ‘จกฺขุมนฺโต รูปานิ ทกฺขนฺตี’ติ, เอวเมวํ ภควตา อเนก\-ปริยาเยน ธมฺโม ปกาสิโต. เอสาหํ ภนฺเต ภควนฺตํ สรณํ คจฺฉามิ ธมฺมญฺจ ภิกฺขุสํฆญฺจ, อุปาสกํ มํ ภควา ธาเรตุ อชฺชตคฺเค ปาณุเปตํ สรณํ คตํ, อธิวาเสตุ จ เม ภนฺเต ภควา สฺวาตนาย ภตฺตํ สทฺธิํ ภิกฺขุสํเฆนาติ. อธิวาเสสิ ภควา ตุณฺหีภาเวน.}

\prose{อถ โข อนาถปิณฺฑิโก คหปติ ภควโต อธิวาสนํ วิทิตฺวา อุฏฺฐายาสนา ภควนฺตํ อภิวาเทตฺวา ปทกฺขิณํ กตฺวา ปกฺกามิ. อสฺโสสิ \csromanfolio{305}โข ราชคหโก เสฏฺฐี “อนาถ\-ปิณฺฑิเกน กิร คหปตินา สฺวาตนาย พุทฺธปฺปมุโข สํโฆ นิมนฺติโต”ติ.}

\proseitem{306}{อถ โข ราชคหโก เสฏฺฐี อนาถปิณฺฑิกํคหปติํ เอตทโวจ “ตยา กิร คหปติ สฺวาตนาย พุทฺธปฺปมุโข สํโฆ นิมนฺติโต, ตฺวญฺจาสิ อาคนฺตุโก, เทมิ เต คหปติ, เวยฺยายิกํ, เยน ตฺวํ พุทฺธ\-ปฺปมุขสฺส สํฆสฺส ภตฺตํ กเรยฺยาสี”ติ. อลํ คหปติ, อตฺถิ เม เวยฺยายิกํ, เยนาหํ พุทฺธ\-ปฺปมุขสฺส สํฆสฺส ภตฺตํ กริสฺสามีติ.}

\prose{อสฺโสสิ โข ราชคหโก เนคโม “อนาถ\-ปิณฺฑิเกน กิร คหปตินา สฺวาตนาย พุทฺธปฺปมุโข สํโฆ นิมนฺติโต”ติ. อถ โข ราชคหโก เนคโม อนาถปิณฺฑิกํ คหปติํ เอตทโวจ “ตยา กิร คหปติ สฺวาตนาย พุทฺธปฺปมุโข สํโฆ นิมนฺติโต, ตฺวญฺจาสิ อาคนฺตุโก, เทมิ เก คหปติ เวยฺยายิกํ, เยน ตฺวํ พุทฺธ\-ปฺปมุขสฺส สํฆสฺส ภตฺตํ กเรยฺยาสี”ติ. อลํ อยฺย, อตฺถิ เม เวยฺยายิกํ, เยนาหํ พุทฺธ\-ปฺปมุขสฺส สํฆสฺส ภตฺตํ กริสฺสามีติ.}

\prose{อสฺโสสิ โขราชา มาคโธ เสนิโย พิมฺพิสาโร “อนาถ\-ปิณฺฑิเกน กิร คหปตินา สฺวาตนาย พุทฺธปฺปมุโข สํโฆ นิมนฺติโต”ติ. อถ โข ราชา มาคโธ เสนิโย พิมฺพิสาโร อนาถปิณฺฑิกํ คหปติํ เอตทโวจ “ตยา กิร คหปติ สฺวาตนาย พุทฺธปฺปมุโข สํโฆ นิมนฺติโต, ตฺวญฺจาสิ อาคนฺตุโก, เทมิ เต คหปติ เวยฺยายิกํ, เยน ตฺวํ พุทฺธ\-ปฺปมุขสฺส สํฆสฺส ภตฺตํ กเรยฺยาสี”ติ. อลํ เทว, อตฺถิ เม เวยฺยายิกํ, เยนาหํ พุทฺธ\-ปฺปมุขสฺส สํฆสฺส ภตฺตํ กริสฺสมีติ.}

\prose{อถ โข อนาถปิณฺฑิโก คหปติ ตสฺสา รตฺติยา อจฺจเยน ราชคหกสฺส เสฏฺฐิสฺส นิเวสเน ปณีตํ ขาทนียํ โภชนิยํ ปฏิยาทาเปตฺวา ภควโต กาลํ อาโรจาเปสิ “กาโล ภนฺเต นิฏฺฐิตํ ภตฺตนฺ”ติ. อถ โข ภควา ปุพฺพณฺหสมยํ นิวาเสตฺวา ปตฺต\-จีวรมา\-ทาย เยน ราชคหกสฺส เสฏฺฐิสฺส นิเวสนํ เตนุปสงฺกมิ, อุปสงฺกมิตฺวา ปญฺญตฺเต อาสเน นิสีทิ สทฺธิํ ภิกฺขุ\-สํเฆน. อถ โข อนาถปิณฺฑิโก คหปติ พุทฺธ\-ปฺปมุขํ ภิกฺขุสํฆํ ปณีเตน ขาทนีเยน โภชนีเยน สหตฺถา สนฺตปฺเปตฺวา สมฺปวาเรตฺวา ภควนฺตํ ภุตฺตาวิํ โอนีตปตฺตวาณิํ เอกมนฺตํ \csromanfolio{306}นิสีทิ, เอกมนฺตํ นิสินฺโน โข อนาถปิณฺฑิโก คหปติ ภควนฺตํ เอตทโวจ “อธิวาเสตุ เม ภนฺเต ภควา สาวตฺถิยํ วสฺสาวาสํ สทฺธิํ ภิกฺขุสํเฆนา”ติ. สุญฺญาคาเร โข คหปติ ตถาคตา อภิรมนฺตีติ. อญฺญาตํ ภควา, อญฺญาตํ สุคตาติ. อถ โข ภควา อนาถปิณฺฑิกํ คหปติํ ธมฺมิยา กถาย สนฺทสฺเสตฺวา สมาทเปตฺวา สมุตฺเตเชตฺวา สมฺปหํเสตฺวา อุฏฺฐายาสนา ปกฺกามิ.}

\proseitem{307}{เตน โข ปน สมเยน อนาถปิณฺฑิโก คหปติ พหุมิตฺโต โหติ พหุสหาโย อาเทยฺยวาโจ. อถ โข อนาถปิณฺฑิโก คหปติ ราชคเห ตํ กรณียํ ตีเรตฺวา เยน สาวตฺถิ เตน ปกฺกามิ. อถ โข อนาถปิณฺฑิโก คหปติ อนฺตรามคฺเค มนุสฺเส อาณาเปสิ “อาราเม อยฺยา กโรถ, วิหาเร ปติฏฺฐาเปถ, ทานานิ ปฏฺฐเปถ, พุทฺโธ โลเก อุปฺปนฺโน, โส จ มยา ภควา นิมนฺติโต อิมินา มคฺเคน อาคจฺฉิสฺสตี”ติ. อถ โข เต มนุสฺสา อนาถ\-ปิณฺฑิเกน คหปตินา อุยฺโยชิตา อาราเม อกํสุ, วิหาเร ปติฏฺฐาเปสุํ, ทานานิ ปฏฺฐเปสุํ.}

\prose{อถ โข นาถปิณฺฑิโก คหปติ สาวตฺถิํ คนฺตา สมนฺตา สาวตฺถิํ อนุวิโลเกสิ “กตฺถ นุ โข ภควา วิหเรยฺย, ยํ อสฺส คามโต เนว อติทูเร น อจฺจาสนฺเน คมนาคมน\-สมฺปนฺนํ อตฺถิกานํ อตฺถิกานํ มนุสฺสานํ อภิกฺกมนียํ ทิวา อปฺปากิณฺณํ รตฺติํ อปฺปสทฺทํ อปฺปนิคฺโฆสํ วิชนวาตํ มนุสฺส\-ราหสฺเสยฺยกํ ปฏิสลฺลานสารุปฺปนฺ”ติ.}

\prose{อทฺทสา โข อนาถปิณฺฑิโก คหปติ เชตสฺส กุมารสฺส\csromansharedfootnote{c334dfe167a7784e}{ราชกุมารสฺส (สี, สฺยา, กํ)} อุยฺยานํ คามโต เนว อติทูเร น อจฺจาสนฺเน คมนาคมน\-สมฺปนฺนํ อตฺถิกานํ อตฺถิกานํ มนุสฺสานํ อภิกฺกมนียํ ทิวา อปฺปากิณฺณํ รตฺติํ อปฺปสทฺทํ อปฺปนิคฺโฆสํ วิชนวาตํ มนุสฺส\-ราหสฺเสยฺยกํ ปฏิสลฺลาน\-สารุปฺปํ, ทิสฺวาน เยน เชโต กุมาโร เตนุปสงฺกมิ, อุปสงฺกมิตฺวา เชตํ กุมารํ เอตทโวจ “เทหิ เม อยฺยปุตฺต อุยฺยานํ อารามํ กาตุนฺ”ติ\csromansharedfootnote{faddcaf487b19903}{เกตุํ (วชีรพุทฺธิฏีกายํ)}. อเทยฺโย คหปติ อาราโม อปิ โกฏิสนฺถเรนาติ. คหิโต อยฺยปุตฺต อาราโมติ. น คหปติ คหิโต อาราโมติ. “คหิโต น คหิโต”ติ โวหาริเก มหามตฺเต ปุจฺฉิํสุ. มหามตฺตา \csromanfolio{307}เอวมาหํสุ “ยโต ตยา อยฺยปุตฺต อคฺโฆ กโต, คหิโต อาราโม”ติ. อถ โข อนาถปิณฺฑิโก คหปติ สกเฏหิ หิรญฺญํ นิพฺพาหาเปตฺวา เชตวนํ โกฏิสนฺถรํ สนฺถราเปสิ, สกิํ นีหฏํ หิรญฺญํ โถกสฺส โอกาสสฺส โกฏฺฐก\-สามนฺตา นปฺปโหติ. อถ โข อนาถปิณฺฑิโก คหปติ มนุสฺเส อาณาเปสิ “คจฺฉถ ภเณ, หิรญฺญํ อาหรถ, อิมํ โอกาสํ สนฺถริสฺสามา”ติ.}

\prose{อถ โข เชตสฺส กุมารสฺส เอตทโหสิ “น โข อิทํ โอรกํ ภวิสฺสติ, ยถายํ คหปติ ตาว พหุํ หิรญฺญํ ปริจฺจชตี”ติ อนาถปิณฺฑิกํ คหปติํ เอตทโวจ “อลํ คหปติ, มา ตํ โอกาสํ สนฺถราเปสิ, เทหิ เม เอตํ โอกาสํ, มเมตํ ทานํ ภวิสฺสตี”ติ. อถ โข อนาถปิณฺฑิโก คหปติ “อยํ โข เชโต กุมาโร อภิญฺญาโก ญาตมนุสฺโส, มหตฺถิโก โข ปน เอวรูปานํ ญาตมนุสฺสานํ อิมสฺมิํ ธมฺมวินเย ปสาโท”ติ ตํ โอกาสํ เชตสฺส กุมารสฺส ปาทาสิ\csromansharedfootnote{5b0c729655da8e57}{อทาสิ (สฺยา)}. อถ โข เชโต กุมาโร ตสฺมิํ โอกาเส โกฏฺฐกํ มาเปสิ.}

\prose{อถ โข อนาถปิณฺฑิโก คหปติ เชตวเน วิหาเร การาเปสิ, ปริเวณานิ การาเปสิ, โกฏฺฐเก การาเปสิ, อุปฏฺฐาน\-สาลาโย การาเปสิ, อคฺคิสาลาโย การาเปสิ, กปฺปิยกุฏิโย การาเปสิ, วจฺจกุฏิโย การาเปสิ, จงฺกเม การาเปสิ, จงฺกมนสาลาโย การาเปสิ, อุทปาเน การาเปสิ, อุทปานสาลาโย การาเปสิ, ชนฺตาฆเร การาเปสิ, ชนฺตา\-ฆร\-สาลาโย การาเปสิ, โปกฺขรณิโย การาเปสิ, มณฺฑเป การาเปสิ.}

\csromanmatikaanchor{mtk.141}
\csromantocmarkref{subsection}{นวกมฺมทาน}{mtk.141}
\bookmark[dest={mtk.141},level=2]{นวกมฺมทาน}
\csromanheader{2}{\textbf{นวกมฺมทาน}}
\proseitem{308}{อถ โข ภควา ราชคเห ยถาภิรนฺตํ วิหริตฺวา เยน เวสาลี เตน จาริกํ ปกฺกามิ. อนุปุพฺเพน จาริกํ จรมาโน เยน เวสาลี ตทวสริ, ตตฺร สุทํ ภควา เวสาลิยํ วิหรติ มหาวเน กูฏาคาร\-สาลายํ. เตน โข ปน สมเยน มนุสฺสา สกฺกจฺจํ นวกมฺมํ กโรนฺติ, เยปิ ภิกฺขู นวกมฺมํ อธิฏฺเฐนฺติ, เตปิ สกฺกจฺจํ อุปฏฺเฐนฺติ. \csromanfolio{308}จีวร\-ปิณฺฑปาต\-เสนาสน\-คิลานปฺปจฺจย\-เภสชฺช\-ปริกฺขาเรน. อถ โข อญฺญตรสฺส ทลิทฺทสฺส ตุนฺนวายสฺส เอตทโหสิ “น โข อิทํ โอรกํ ภวิสฺสติ, ยถยิเม มนุสฺสา สกฺกจฺจํ นวกมฺมํ กโรนฺติ, ยํนูนาหมฺปิ นวกมฺมํ กเรยฺยนฺ”ติ. อถ โข โส ทลิทฺโท ตุนฺนวาโย สามํ จิกฺขลฺลํ มทฺทิตฺวา อิฏฺฐกาโย จินิตฺวา กุฏฺฏํ อุฏฺฐาเปสิ, เตน อกุสลเกน จิตา วงฺกา ภิตฺติ ปริปติ. ทุติยมฺปิ โข ฯเปฯ. ตติยมฺปิ โข โส ทลิทฺโท ตุนฺนวาโย สามํ จิกฺขลฺลํ มทฺทิตฺวา อิฏฺฐกาโย จินิตฺวา กุฏฺฏํ อุฏฺฐาเปสิ, เตน อกุสลเกน จิตา วงฺกา ภิตฺติ ปริปติ. อถ โข โส ทลิทฺโท ตุนฺนวาโย อุชฺฌายติ ขิยฺยติ วิปาเจติ “เย อิเมสํ สมณานํ สกฺย\-ปุตฺติ\-ยานํ เทนฺติ จีวร\-ปิณฺฑปาต\-เสนาสน\-คิลานปฺปจฺจย\-เภสชฺช\-ปริกฺขารํ, เต อิเม โอวทนฺติ อนุสาสนฺติ, เตสญฺจ นวกมฺมํ อธิฏฺเฐนฺติ, อหํ ปนมฺหิ ทลิทฺโท, น มํ โกจิ โอวทติ วา อนุสาสติ วา, นวกมฺมํ วา อธิฏฺเฐตี”ติ. อสฺโสสุํ โข ภิกฺขู ตสฺส ทลิทฺทสฺส ตุนฺนวายสฺส อุชฺฌายนฺตสฺส ขิยฺยนฺตสฺส วิปาเจนฺตสฺส. อถ โข เต ภิกฺขู ภควโต เอตมตฺถํ อาโรเจสุํ. อถ โข ภควา เอตสฺมิํ นิทาเน เอตสฺมิํ ปกรเณ ธมฺมิํ กถํ กตฺวา ภิกฺขู อามนฺเตสิ\csromandash{}อนุชานามิ ภิกฺขเว นวกมฺมํ ทาตุํ, นวกมฺมิโก ภิกฺขเว ภิกฺขุ อุสฺสุกฺกํ อาปชฺชิสฺสติ “กินฺติ นุ โข วิหาโร ขิปฺปํ ปริโยสานํ คจฺเฉยฺยา”ติ, ขณฺฑํ ผุลฺลํ ปฏิสงฺขริสฺสติ. เอวญฺจ ปน ภิกฺขเว ทาตพฺพํ, ปฐมํ ภิกฺขุ ยาจิตพฺโพ, ยาจิตฺวา พฺยตฺเตน ภิกฺขุนา ปฏิพเลน สํโฆ ญาเปตพฺโพ\csromandash{}}

\proseitem{309}{“สุณาตุ เม ภนฺเต สํโฆ, ยทิ สํฆสฺส ปตฺตกลฺลํ, สํโฆ อิตฺถนฺนามสฺส คหปติโน วิหารํ อิตฺถนฺนามสฺส ภิกฺขุโน นวกมฺมํ ทเทยฺย, เอสา ญตฺติ.}

\prose{สุณาตุ เม ภนฺเต สํโฆ, สํโฆ อิตฺถนฺนามสฺส คหปติโน วิหารํ อิตฺถนฺนามสฺส ภิกฺขุโน นวกมฺมํ เทติ, ยสฺสายสฺมโต ขมติ อิตฺถนฺนามสฺส คหปติโน วิหารํ อิตฺถนฺนามสฺส ภิกฺขุโน นวกมฺมสฺส ทานํ, โส ตุณฺหสฺส, ยสฺส นกฺขมติ, โส ภาเสยฺย.}

\prose{ทินฺโน สํเฆน อิตฺถนฺนามสฺส คหปติโน วิหาโร อิตฺถนฺนามสฺส ภิกฺขุโน นวกมฺมํ, ขมติ สํฆสฺส, ตสฺมา ตุณฺหี, เอวเมตํ ธารยามี”ติ.}

\csromanmatikaanchor{mtk.142}
\csromantocmarkref{subsection}{อคฺคาสนาทิอนุชานน}{mtk.142}
\bookmark[dest={mtk.142},level=2]{อคฺคาสนาทิอนุชานน}
\csromanheader{2}{อคฺคาสนา\-ทิ\-อนุชานน}
\proseitem{310}{\csromanfolio{309}อถ โข ภควา เวสาลิยํ ยถาภิรนฺตํ วิหริตฺวา เยน สาวตฺถิ เตน จาริกํ ปกฺกามิ. เตน โข ปน สมเยน ฉพฺพคฺคิยานํ ภิกฺขูนํ อนฺเตวาสิกา ภิกฺขู พุทฺธ\-ปฺปมุขสฺส สํฆสฺส ปุรโต ปุรโต คนฺตฺวา วิหาเร ปริคฺคณฺหนฺติ เสยฺยาโย ปริคฺคณฺหนฺติ “อิทํ อมฺหากํ อุปชฺฌายานํ ภวิสฺสติ, อิทํ อมฺหากํ อาจริยานํ ภวิสฺสติ, อิทํ อมฺหากํ ภวิสฺสตี”ติ.}

\prose{อถ โข อายสฺมา สาริปุตฺโต พุทฺธ\-ปฺปมุขสฺส สํฆสฺส ปิฏฺฐิโต ปิฏฺฐิโต คนฺตฺวา วิหาเรสุ ปริคฺคหิเตสุ เสยฺยาสุ ปริคฺคหิตาสุ เสยฺยํ อลภมาโน อญฺญตรสฺมิํ รุกฺขมูเล นิสีทิ. อถ โข ภควา รตฺติยา ปจฺจูสสมยํ ปจฺจุฏฺฐาย อุกฺกาสิ. อายสฺมาปิ สาริปุตฺโต อุกฺกาสิ. โก เอตฺถาติ. อหํ ภควา สาริปุตฺโตติ. กิสฺส ตฺวํ สาริปุตฺต อิธ นิสินฺโนติ. อถ โข อายสฺมา สาริปุตฺโต ภควโต เอตมตฺถํ อาโรเจสิ. อถ โข ภควา เอตสฺมิํ นิทาเน เอตสฺมิํ ปกรเณ ภิกฺขุสํฆํ สนฺนิปาตาเปตฺวา ภิกฺขู ปฏิปุจฺฉิ\csromandash{}สจฺจํ กิร ภิกฺขเว ฉพฺพคฺคิยานํ ภิกฺขูนํ อนฺเตวาสิกา ภิกฺขู พุทฺธ\-ปฺปมุขสฺส สํฆสฺส ปุรโต ปุรโต คนฺตฺวา วิหาเร ปริคฺคณฺหนฺติ เสยฺยาโย ปริคฺคณฺหนฺติ “อิทํ อมฺหากํ อุปชฺฌายานํ ภวิสฺสติ, อิทํ อมฺหากํ อาจริยานํ ภวิสฺสติ, อิทํ อมฺหากํ ภวิสฺสตี”ติ. สจฺจํ ภควาติ. วิครหิ พุทฺโธ ภควา ฯเปฯ. กถํ หิ นาม เต ภิกฺขเว โมฆปุริสา พุทฺธ\-ปฺปมุขสฺส สํฆสฺส ปุรโต ปุรโต คนฺตฺวา วิหาเร ปริคฺคเหสฺสนฺติ เสยฺยาโย ปริคฺคเหสฺสนฺติ “อิทํ อมฺหากํ อุปชฺฌายานํ ภวิสฺสติ, อิทํ อมฺหากํ อาจริยานํ ภวิสฺสติ, อิทํ อมฺหากํ ภวิสฺสตี”ติ. เนตํ ภิกฺขเว อปฺปสนฺนานํ วา ปสาทาย ฯเปฯ วิครหิตฺวา ฯเปฯ ธมฺมิํ กถํ กตฺวา ภิกฺขู อามนฺเตสิ “โก ภิกฺขเว อรหติ อคฺคาสนํ อคฺโคทกํ อคฺคปิณฺฑนฺ”ติ.}

\prose{เอกจฺเจ ภิกฺขู เอวมาหํสุ “โย ภควา ขตฺติยกุลา ปพฺพชิโต, โส อรหติ อคฺคาสนํ อคฺโคทกํ อคฺคปิณฺฑนฺ”ติ. เอกจฺเจ ภิกฺขู เอวมาหํสุ “โย ภควา พฺราหฺมณกุลา ปพฺพชิโต, โส อรหติ อคฺคาสนํ อคฺโคทกํ อคฺคปิณฺฑนฺ”ติ. เอกจฺเจ ภิกฺขู เอวมาหํสุ “โย ภควา คหปติกุลา ปพฺพชิโต, โส อรหติ อคฺคาสนํ อคฺโคทกํ อคฺคปิณฺฑนฺ”ติ. \csromanfolio{310}เอกจฺเจ ภิกฺขู เอวมาหํสุ “โย ภควา สุตฺตนฺติโก, โส อรหติ อคฺคาสนํ อคฺโคทกํ อคฺคปิณฺฑนฺ”ติ. เอกจฺเจ ภิกฺขู เอวมาหํสุ “โย ภควา วินยธโร, โส อรหติ อคฺคาสนํ อคฺโคทกํ อคฺคปิณฺฑนฺ”ติ. เอกจฺเจ ภิกฺขู เอวมาหํสุ “โย ภควา ธมฺมกถิโก, โส อรหติ อคฺคาสนํ อคฺโคทกํ อคฺคปิณฺฑนฺ”ติ. เอกจฺเจ ภิกฺขู เอวมาหํสุ “โย ภควา ปฐมสฺส ฌานสฺส ลาภี, โส อรหติ อคฺคาสนํ อคฺโคทกํ อคฺคปิณฺฑนฺ”ติ. เอกจฺเจ ภิกฺขู เอวมาหํสุ “โย ภควา ทุติยสฺส ฌานสฺส ลาภี, โส อรหติ อคฺคาสนํ อคฺโคทกํ อคฺคปิณฺฑนฺ”ติ. เอกจฺเจ ภิกฺขู เอวมาหํสุ “โย ภควา ตติยสฺส ฌานสฺส ลาภี, โส อรหติ อคฺคาสนํ อคฺโคทกํ อคฺคปิณฺฑนฺ”ติ. เอกจฺเจ ภิกฺขู เอวมาหํสุ “โย ภควา จตุตฺถสฺส ฌานสฺส ลาภี, โส อรหติ อคฺคาสนํ อคฺโคทกํ อคฺคปิณฺฑนฺ”ติ. เอกจฺเจ ภิกฺขู เอวมาหํสุ “โย ภควา โสตาปนฺโน, โส อรหติ อคฺคาสนํ อคฺโคทกํ อคฺคปิณฺฑนฺ”ติ. เอกจฺเจ ภิกฺขู เอวมาหํสุ “โย ภควา สกทาคามี ฯเปฯ. โย ภควา อนาคามี ฯเปฯ. โย ภควา อรหา, โส อรหติ อคฺคาสนํ อคฺโคทกํ อคฺคปิณฺฑนฺ”ติ. เอกจฺเจ ภิกฺขู เอวมาหํสุ “โย ภควา เตวิชฺโช, โส อรหติ อคฺคาสนํ อคฺโคทกํ อคฺคปิณฺฑนฺ”ติ. เอกจฺเจ ภิกฺขู เอวมาหํสุ “โย ภควา ฉฬภิญฺโญ, โส อรหติ อคฺคาสนํ อคฺโคทกํ อคฺคปิณฺฑนฺ”ติ.}

\proseitem{311}{อถ โข ภควา ภิกฺขู อามนฺเตสิ “ภูตปุพฺพํ ภิกฺขเว หิมวนฺต\-ปเทเส\csromansharedfootnote{6426805b6f363f30}{หิมวนฺตสฺส (สี, สฺยา)} มหานิโคฺรโธ อโหสิ, ตํ ตโย สหายา อุปนิสฺสาย วิหริํสุ ติตฺติโร จ มกฺกโฏ จ หตฺถินาโค จ, เต อญฺญมญฺญํ อคารวา อปฺปติสฺสา อสภาค\-วุตฺติกา วิหรนฺติ. อถ โข ภิกฺขเว เตสํ สหายานํ เอตทโหสิ “อโห นูน มยํ ชาเนยฺยาม, ยํ อมฺหากํ ชาติยา มหนฺตตรํ, ตํ มยํ สกฺกเรยฺยาม ครุํ กเรยฺยาม มาเนยฺยาม ปูเชยฺยาม, ตสฺส จ มยํ โอวาเท ติฏฺเฐยฺยามา”ติ.}

\prose{อถ โข ภิกฺขเว ติตฺติโร จ มกฺกโฏ จ หตฺถินาคํ ปุจฺฉิํสุ “ตฺวํ สมฺม กิํ โปราณํ สรสี”ติ. ยทาหํ สมฺมา โปโต โหมิ, อิมํ นิโคฺรธํ \csromanfolio{311}อนฺตรา สตฺถีนํ\csromansharedfootnote{5fbc8cebfc3803a3}{อนฺตราสตฺถิกํ (สี)} กริตฺวา อติกฺกมามิ, อคฺคงฺกุรกํ เม อุทรํ ฉุปติ, อิมาหํ สมฺมา โปราณํ สรามีติ.}

\prose{อถ โข ภิกฺขเว ติตฺติโร จ หตฺถินาโค จ มกฺกฏํ ปุจฺฉิํสุ “ตฺวํ สมฺม กิํ โปราณํ สรสี”ติ. ยทาหํ สมฺมา ฉาโป โหมิ, ฉมายํ นิสีทิตฺวา อิมสฺส นิโคฺรธสฺส อคฺคงฺกุรกํ ขาทามิ, อิมาหํ สมฺมา โปราณํ สรามีติ.}

\prose{อถ โข ภิกฺขเว มกฺกโฏ จ หตฺถินาโค จ ติตฺติรํ ปุจฺฉิํสุ “ตฺวํ สมฺม กิํ โปราณํ สรสี”ติ. อมุกสฺมิํ สมฺมา โอกาเส มหานิโคฺรโธ อโหสิ, ตโต อหํ ผลํ ภกฺขิตฺวา อิมสฺมิํ โอกาเส วจฺจํ อกาสิํ, ตสฺสายํ นิโคฺรโธ ชาโต, ตทาหํ สมฺมา ชาติยา มหนฺตตโรติ.}

\prose{อถ โข ภิกฺขเว มกฺกโฏ จ หตฺถิ นาโค จ ติตฺติรํ เอตทโวจุํ “ตฺวํ สมฺม อมฺหากํ ชาติยา มหนฺตตโร, ตํ มยํ สกฺกริสฺสาม ครุํ กริสฺสาม มาเนสฺสาม ปูเชสฺสาม, ตุยฺหญฺจ มยํ โอวาเท ปติฏฺฐิสฺสามา”ติ. อถ โข ภิกฺขเว ติตฺติโร มกฺกฏญฺจ หตฺถินาคญฺจ ปญฺจสุ สีเลสุ สมาทเปสิ, อตฺตนา จ ปญฺจสุ สีเลสุ สมาทาย วตฺตติ, เต อญฺญมญฺญํ สคารวา สปฺปติสฺสา สภาควุตฺติกา วิหริตฺวา กายสฺส เภทา ปรํ มรณา สุคติํ สคฺคํ โลกํ อุปปชฺชิํสุ. เอวํ โข ตํ ภิกฺขเว ติตฺติริยํ นาม พฺรหฺมจริยํ อโหสิ.}

\setlength{\csromangathapagewidth}{0pt}
\setlength{\csromangathawakwidth}{0pt}
\csromangathameasuregroup{\textsuperscript{*}เย วุฑฺฒมปจายนฺติ, \\ ทิฏฺเฐ ธมฺเม จ ปาสํสา,}{\csromangathabat{\textsuperscript{*}เย วุฑฺฒมปจายนฺติ,}{นรา ธมฺมสฺส โกวิทา.} \\ \csromangathabat{ทิฏฺเฐ ธมฺเม จ ปาสํสา,}{สมฺปราเย จ สุคฺคตีติ.}}
\csromangathasetleft{\textsuperscript{*}เย วุฑฺฒมปจายนฺติ, \\ ทิฏฺเฐ ธมฺเม จ ปาสํสา,}
\csromangathagroup{\csromangathastanza{\csromangathabat{\csromansymbolfootnote{*}{ขุ 5. 9 ปิฏฺเฐ ชาตเก.}เย วุฑฺฒมปจายนฺติ,}{นรา ธมฺมสฺส โกวิทา.} \\ \csromangathabat{ทิฏฺเฐ ธมฺเม จ ปาสํสา,}{สมฺปราเย จ สุคฺคตีติ.}}}

\prose{เต หิ นาม ภิกฺขเว ติรจฺฉานคตา ปาณา อญฺญมญฺญํ สคารวา สปฺปติสฺสา สภาควุตฺติกา วิหริสฺสนฺติ. อิธ โข ตํ ภิกฺขเว โสเภถ, ยํ ตุเมฺห เอวํ สฺวากฺขาเต ธมฺมวินเย ปพฺพชิตา สมานา อญฺญมญฺญํ อคารวา อปฺปติสฺสา อสภาค\-วุตฺติกา วิหเรยฺยาถ. เนตํ ภิกฺขเว อปฺปสนฺนานํ วา ปสาทาย ฯเปฯ วิครหิตฺวา ฯเปฯ ธมฺมิํ กถํ กตฺวา ภิกฺขู อามนฺเตสิ “อนุชานามิ ภิกฺขเว ยถาวุฑฺฒํ อภิวาทนํ ปจฺจุฏฺฐานํ อญฺชลิกมฺมํ สามีจิกมฺมํ อคฺคาสนํ อคฺโคทกํ อคฺคปิณฺฑํ, น จ ภิกฺขเว สํฆิกํ ยถาวุฑฺฒํ ปฏิพาหิตพฺพํ, โย ปฏิพาเหยฺย,อาปตฺติ ทุกฺกฏสฺสา”ติ.}

\csromanmatikaanchor{mtk.143}
\csromantocmarkref{subsection}{อวนฺทิยาทิปุคฺคล}{mtk.143}
\bookmark[dest={mtk.143},level=2]{อวนฺทิยาทิปุคฺคล}
\csromanheader{2}{อวนฺทิยา\-ทิ\-ปุคฺคล}
\proseitem{312}{\csromanfolio{312}ทสยิเม ภิกฺขเว อวนฺทิยา. ปุเร อุปสมฺปนฺเนน ปจฺฉา อุปสมฺปนฺโน อวนฺทิโย, อนุปสมฺปนฺโน อวนฺทิโย, นานาสํวาสโก วุฑฺฒตโร อธมฺมวาที อวนฺทิโย, มาตุคาโม อวนฺทิโย, ปณฺฑโก อวนฺทิโย, ปาริวาสิโก อวนฺทิโย, มูลาย\-ปฏิกสฺสนา\-รโห อวนฺทิโย, มานตฺตารโห อวนฺทิโย, มานตฺตจาริโก อวนฺทิโย, อพฺภานารโห อวนฺทิโย. อิเม โข ภิกฺขเว ทส อวนฺทิยา.}

\prose{ตโยเม ภิกฺขเว วนฺทิยา. ปจฺฉา\-อุปสมฺปนฺเนน ปุเร\-อุปสมฺปนฺโน วนฺทิโย, นานาสํวาสโก วุฑฺฒตโร ธมฺมวาที วนฺทิโย, สเทวเก ภิกฺขเว โลเก สมารเก สพฺรหฺมเก สสฺสมณ\-พฺราหฺมณิยา ปชาย สเทว\-มนุสฺสาย ตถาคโต อรหํ สมฺมาสมฺพุทฺโธ วนฺทิโย. อิเม โข ภิกฺขเว ตโย วนฺทิยาติ.}

\csromanmatikaanchor{mtk.144}
\csromantocmarkref{subsection}{อาสนปฺปฏิพาหนปฏิกฺเขป}{mtk.144}
\bookmark[dest={mtk.144},level=2]{อาสนปฺปฏิพาหนปฏิกฺเขป}
\csromanheader{2}{อาสน\-ปฺปฏิพาหน\-ปฏิกฺเขป}
\proseitem{313}{เตน โข ปน สมเยน มนุสฺสา สํฆํ อุทฺทิสฺส มณฺฑเป ปฏิยาเทนฺติ, สนฺถเร ปฏิยาเทนฺติ, โอกาเส ปฏิยาเทนฺติ, ฉพฺพคฺคิยานํ ภิกฺขูนํ อนฺเตวาสิกา ภิกฺขู “สํฆิกญฺเญว ภควตา ยถาวุฑฺฒํ อนุญฺญาตํ, โน อุทฺทิสฺสกตนฺ”ติ พุทฺธ\-ปฺปมุขสฺส สํฆสฺส ปุรโต ปุรโต คนฺตฺวา มณฺฑเปปิ ปริคฺคณฺหนฺติ สนฺตเรปิ ปริคฺคณฺหนฺติ โอกาเสปิ ปริคฺคณฺหนฺติ “อิทํ อมฺหากํ อุปชฺฌายานํ ภวิสฺสติ, อิทํ อมฺหากํ อาจริยานํ ภวิสฺสติ, อิทํ อมฺหากํ ภวิสฺสตี”ติ. อถ โข อายสฺมา สาริปุตฺโต พุทฺธ\-ปฺปมุขสฺส สํฆสฺส ปิฏฺฐิโต ปิฏฺฐิโต คนฺตฺวา มณฺฑเปสุ ปริคฺคหิเตสุ สนฺถเรสุ ปริคฺคหิเตสุ โอกาเสสุ ปริคฺคหิเตสุ โอกาสํ อลภมาโน อญฺญตรสฺมิํ รุกฺขมูเล นิสีทิ. อถ โข ภควา รตฺติยา ปจฺจูสสมยํ ปจฺจุฏฺฐาย อุกฺกาสิ. อายสฺมาปิ สาริปุตฺโต อุกฺกาสิ. โก เอตฺถาติ. อหํ ภควา สาริปุตฺโตติ. กิสฺส ตฺวํ สาริปุตฺต อิธ นิสินฺโนติ. อถ โข อายสฺมา สาริปุตฺโต ภควโต เอตมตฺถํ อาโรเจสิ. อถ โข ภควา เอตสฺมิํ นิทาเน เอตสฺมิํ ปกรเณ ภิกฺขุสํฆํ สนฺนิปาตาเปตฺวา ภิกฺขู ปฏิปุจฺฉิ “สจฺจํ กิร ภิกฺขเว ฉพฺพคฺคิยานํ ภิกฺขูนํ อนฺเตวาสิกา ภิกฺขู ‘สํฆิกญฺเญว ภควตา ยถาวุฑฺฒํ อนุญฺญาตํ, โน อุทฺทิสฺสกตนฺ’ติ พุทฺธ\-ปฺปมุขสฺส สํฆสฺส ปุรโต ปุรโต \csromanfolio{313}คนฺตฺวา มณฺฑเป ปริคฺคณฺหนฺติ, สนฺถเร ปริคฺคณฺหนฺติ, โอกาเส ปริคฺคณฺหนฺติ, อิทํ อมฺหากํ อุปชฺฌายานํ ภวิสฺสติ, อิทํ อมฺหากํ อาจริยานํ ภวิสฺสติ, อิทํ อมฺหากํ ภวิสฺสตี”ติ. สจฺจํ ภควาติ ฯเปฯ วิครหิตฺวา ฯเปฯ ธมฺมิํ กถํ กตฺวา ภิกฺขู อามนฺเตสิ “น ภิกฺขเว อุทฺทิสฺสกตมฺปิ ยถาวุฑฺฒํ ปฏิพาหิตพฺพํ, โย ปฏิพาเหยฺย, อาปตฺติ ทุกฺกฏสฺสา”ติ.}

\csromanmatikaanchor{mtk.145}
\csromantocmarkref{subsection}{คิหิวิกตอนุชานน}{mtk.145}
\bookmark[dest={mtk.145},level=2]{คิหิวิกตอนุชานน}
\csromanheader{2}{คิหิวิกตอนุชานน}
\proseitem{314}{เตน โข ปน สมเยน มนุสฺสา ภตฺตคฺเค อนฺตรฆเร อุจฺจาสยน\-มหาสยนานิ ปญฺญเปนฺติ. เสยฺยถิทํ, อาสนฺทิํ ปลฺลงฺกํ โคนกํ จิตฺตกํ ปฏิกํ ปฏลิกํ ตูลิกํ วิกติกํ อุทฺทโลมิํ เอกนฺตโลมิํ กฏฺฏิสฺสํ โกเสยฺยํ\csromansharedfootnote{b1ba8a5e2970c347}{โกเสยฺยํ กมฺพลํ (สี, สฺยา)} กุตฺตกํ หตฺถตฺถรํ อสฺสตฺถรํ รถตฺถรํ อชินปเวณิํ กทลิ\-มิค\-ปวร\-ปจฺจตฺถรณํ ส- อุตฺตรจฺฉทํ อุภโต\-โลหิตกู\-ปธานํ. ภิกฺขู กุกฺกุจฺจายนฺตา นาภินิสีทนฺติ. ภควโต เอตมตฺถํ อาโรเจสุํ. อนุชานามิ ภิกฺขเว ฐเปตฺวา ตีณิ อาสนฺทิํ ปลฺลงฺกํ ตูลิกํ คิหิวิกตํ\csromansharedfootnote{4d85591305645ff4}{คิหิวิกฏํ (สี, ก), อวเสสํ คิหิวิกฏํ (สฺยา)} อภินิสีทิตุํ, นเตฺวว อภินิ\-ปชฺชิ\-ตุนฺติ.}

\prose{เตน โข ปน สมเยน มนุสฺสา ภตฺตคฺเค อนฺตรฆเร ตูโลนทฺธํ มญฺจํปิ ปีฐํปิ ปญฺญเปนฺติ. ภิกฺขู กุกฺกุจฺจายนฺตา นาภินิสีทนฺติ. ภควโต เอตมตฺถํ อาโรเจสุํ. อนุชานามิ ภิกฺขเว คิหิวิกตํ อภินิสีทิตุํ, นเตฺวว อภินิ\-ปชฺชิ\-ตุนฺติ.}

\csromanmatikaanchor{mtk.146}
\csromantocmarkref{subsection}{เชตวนวิหารานุโมทนา}{mtk.146}
\bookmark[dest={mtk.146},level=2]{เชตวนวิหารานุโมทนา}
\csromanheader{2}{เชตวน\-วิหารา\-นุโมทนา}
\proseitem{315}{อถ โข ภควา อนุปุพฺเพน จาริกํ จรมาโน เยน สาวตฺถิ ตทวสริ, ตตฺร สุทํ ภควา สาวตฺถิยํ วิหรติ เชตวเน อนาถ\-ปิณฺฑิกสฺส อาราเม. อถ โข อนาถปิณฺฑิโก คหปติ เยน ภควา เตนุปสงฺกมิ, อุปสงฺกมิตฺวา ภควนฺตํ อภิวาเทตฺวา เอกมนฺตํ นิสีทิ, เอกมนฺตํ นิสินฺโน โข อนาถปิณฺฑิโก คหปติ ภควนฺตํ เอตทโวจ “อธิวาเสตุ เม ภนฺเต ภควา สฺวาตนาย ภตฺตํ สทฺธิํ ภิกฺขุสํเฆนา”ติ. อธิวาเสสิ ภควา ตุณฺหีภาเวน. อถ โข อนาถปิณฺฑิโก คหปติ ภควโต อธิวาสนํ วิทิตฺวา อุฏฺฐายาสนา \csromanfolio{314}ภควนฺตํ อภิวาเทตฺวา ปทกฺขิณํ กตฺวา ปกฺกามิ. อถ โข อนาถปิณฺฑิโก คหปติ ตสฺสา รตฺติยา อจฺจเยน ปณีตํ ขาทนียํ โภชนียํ ปฏิยาทาเปตฺวา ภควโต กาลํ อาโรจาเปสิ “กาโล ภนฺเต, นิฏฺฐิตํ ภตฺตนฺ”ติ. อถ โข ภควา ปุพฺพณฺหสมยํ นิวาเสตฺวา ปตฺต\-จีวรมา\-ทาย เยน อนาถ\-ปิณฺฑิกสฺส คหปติสฺส นิเวสนํ เตนุปสงฺกมิ, อุปสงฺกมิตฺวา ปญฺญตฺเต อาสเน นิสีทิ สทฺธิํ ภิกฺขุ\-สํเฆน. อถ โข อนาถปิณฺฑิโก คหปติ พุทฺธ\-ปฺปมุขํ ภิกฺขุสํฆํ ปณีเตน ขาทนีเยน โภชนีเยน สหตฺถา สนฺตปฺเปตฺวา สมฺปวาเรตฺวา ภควนฺตํ ภุตฺตาวิํ โอนีต\-ปตฺต\-ปาณิํ เอกมนฺตํ นิสีทิ, เอกมนฺตํ นิสินฺโน โข อนาถปิณฺฑิโก คหปติ ภควนฺตํ เอตทโวจ “กถาหํ ภนฺเต เชตวเน ปฏิปชฺชามี”ติ. เตน หิ ตฺวํ คหปติ เชตวนํ อาคตานาคตสฺส จาตุทฺทิสสฺส สํฆสฺส ปติฏฺฐเปหีติ. “เอวํ ภนฺเต”ติ โข อนาถปิณฺฑิโก คหปติ ภควโต ปฏิสฺสุตฺวา เชตวนํ อาคตานาคตสฺส จาตุทฺทิสสฺส สํฆสฺส ปติฏฺฐาเปสิ.}

\prose{อถ โข ภควา อนาถปิณฺฑิกํ คหปติํ อิมาหิ คาถาหิ อนุโมทิ\csromandash{}}

\setlength{\csromangathapagewidth}{0pt}
\setlength{\csromangathawakwidth}{0pt}
\csromangathameasuregroup{“สีตํ อุณฺหํ ปฏิหนฺติ, \\ สริสเป จ มกเส, \\ ตโต วาตาตโป โฆโร, \\ เลณตฺถญฺจ สุขตฺถญฺจ, \\ วิหารทานํ สํฆสฺส, \\ ตสฺมา หิ ปณฺฑิโต โปโส, \\ วิหาเร การเย รมฺเม, \\ เตสํ อนฺนญฺจ ปานญฺจ, \\ ทเทยฺย อุชุภูเตสุ, \\ เต ตสฺส ธมฺมํ เทเสนฺติ, \\ ยํ โส ธมฺมํ อิธญฺญาย,}{\csromangathabat{“สีตํ อุณฺหํ ปฏิหนฺติ,}{ตโต วาฬมิคานิ จ.} \\ \csromangathabat{สริสเป จ มกเส,}{สิสิเร จาปิ วุฏฺฐิโย.} \\ \csromangathabat{ตโต วาตาตโป โฆโร,}{สญฺชาโต ปฏิหญฺญติ.} \\ \csromangathabat{เลณตฺถญฺจ สุขตฺถญฺจ,}{ฌายิตุญฺจ วิปสฺสิตุํ.} \\ \csromangathabat{วิหารทานํ สํฆสฺส,}{อคฺคํ พุทฺเธน วณฺณิตํ.} \\ \csromangathabat{ตสฺมา หิ ปณฺฑิโต โปโส,}{สมฺปสฺสํ อตฺถมตฺตโน.} \\ \csromangathabat{วิหาเร การเย รมฺเม,}{วาสเยตฺถ พหุสฺสุเต.} \\ \csromangathabat{เตสํ อนฺนญฺจ ปานญฺจ,}{วตฺถเสนาสนานิ จ.} \\ \csromangathabat{ทเทยฺย อุชุภูเตสุ,}{วิปฺปสนฺเนน เจตสา.} \\ \csromangathabat{เต ตสฺส ธมฺมํ เทเสนฺติ,}{สพฺพทุกฺขาปนูทนํ.} \\ \csromangathabat{ยํ โส ธมฺมํ อิธญฺญาย,}{ปรินิพฺพาติ อนาสโว”ติ.}}
\csromangathasetleft{“สีตํ อุณฺหํ ปฏิหนฺติ, \\ สริสเป จ มกเส, \\ ตโต วาตาตโป โฆโร, \\ เลณตฺถญฺจ สุขตฺถญฺจ, \\ วิหารทานํ สํฆสฺส, \\ ตสฺมา หิ ปณฺฑิโต โปโส, \\ วิหาเร การเย รมฺเม, \\ เตสํ อนฺนญฺจ ปานญฺจ, \\ ทเทยฺย อุชุภูเตสุ, \\ เต ตสฺส ธมฺมํ เทเสนฺติ, \\ ยํ โส ธมฺมํ อิธญฺญาย,}
\csromangathagroup{\csromangathastanza{\csromangathabat{“สีตํ อุณฺหํ ปฏิหนฺติ,}{ตโต วาฬมิคานิ จ.} \\ \csromangathabat{สริสเป จ มกเส,}{สิสิเร จาปิ วุฏฺฐิโย.}}\csromangathastanza{\csromangathabat{ตโต วาตาตโป โฆโร,}{สญฺชาโต ปฏิหญฺญติ.} \\ \csromangathabat{เลณตฺถญฺจ สุขตฺถญฺจ,}{ฌายิตุญฺจ วิปสฺสิตุํ.}}\csromangathastanza{\csromangathabat{วิหารทานํ สํฆสฺส,}{อคฺคํ พุทฺเธน วณฺณิตํ.} \\ \csromangathabat{ตสฺมา หิ ปณฺฑิโต โปโส,}{สมฺปสฺสํ อตฺถมตฺตโน.}}\csromangathastanza{\csromangathabat{วิหาเร การเย รมฺเม,}{วาสเยตฺถ พหุสฺสุเต.} \\ \csromangathabat{เตสํ อนฺนญฺจ ปานญฺจ,}{วตฺถ\-เสนาสนานิ จ.}}\csromangathastanza{\csromangathabat{ทเทยฺย อุชุภูเตสุ,}{วิปฺปสนฺเนน เจตสา.} \\ \csromangathabat{เต ตสฺส ธมฺมํ เทเสนฺติ,}{สพฺพทุกฺขาปนูทนํ.} \\ \csromangathabat{ยํ โส ธมฺมํ อิธญฺญาย,}{ปรินิพฺพาติ อนาสโว”ติ.}}}

\prose{อถ โข ภควา อนาถปิณฺฑิกํ คหปติํ อิมาหิ คาถาหิ อนุโมทิตฺวา อุฏฺฐายาสนา ปกฺกามิ.}

\csromanmatikaanchor{mtk.147}
\csromantocmarkref{subsection}{อาสนปฺปฏิพาหนาทิ}{mtk.147}
\bookmark[dest={mtk.147},level=2]{อาสนปฺปฏิพาหนาทิ}
\csromanheader{2}{อาสน\-ปฺปฏิพาหนาทิ}
\proseitem{316}{\csromanfolio{315}เตน โข ปน สมเยน อญฺญตรสฺส อาชีวก\-สาวกสฺส มหามตฺตสฺส สํฆภตฺตํ โหติ. อายสฺมา อุปนนฺโท สกฺยปุตฺโต ปจฺฉา อาคนฺตฺวา วิปฺปกต\-โภชนํ อานนฺตริกํ ภิกฺขุํ วุฏฺฐาเปสิ, ภตฺตคฺคํ โกลาหลํ อโหสิ. อถ โข โส มหามตฺโต อุชฺฌายติ ขิยฺยติ วิปาเจติ “กถํ หิ นาม สมณา สกฺยปุตฺติยา ปจฺฉา อาคนฺตฺวา วิปฺปกต\-โภชนํ อานนฺตริกํ ภิกฺขุํ วุฏฺฐาเปสฺสนฺติ, ภตฺตคฺคํ โกลาหลํ อโหสิ, นนุ นาม ลพฺภา อญฺญตฺราปิ นิสินฺเนน ยาวทตฺถํ ภุญฺชิตุนฺ”ติ. อสฺโสสุํ โข ภิกฺขู ตสฺส มหามตฺตสฺส อุชฺฌายนฺตสฺส ขิยฺยนฺตสฺส วิปาเจนฺตสฺส. เย เต ภิกฺขู อปฺปิจฺฉา ฯเปฯ เต อุชฺฌายนฺติ ขิยฺยนฺติ วิปาเจนฺติ “กถํ หิ นาม อายสฺมา อุปนนฺโท สกฺยปุตฺโต ปจฺฉา อาคนฺตฺวา วิปฺปกต\-โภชนํ อานนฺตริกํ ภิกฺขุํ วุฏฺฐาเปสฺสติ, ภตฺตคฺคํ โกลาหลํ อโหสี”ติ. อถ โข เต ภิกฺขู ภควโต เอตมตฺถํ อาโรเจสุํ. สจฺจํ กิร ตฺวํ อุปนนฺท ปจฺฉา อาคนฺตฺวา วิปฺปกต\-โภชนํ อานนฺตริกํ ภิกฺขุํ วุฏฺฐาเปสิ, ภตฺตคฺคํ โกลาหลํ อโหสีติ. สจฺจํ ภควาติ. วิครหิ พุทฺโธ ภควา ฯเปฯ. กถํ หิ นาม ตฺวํ โมฆปุริส ปจฺฉา อาคนฺตฺวา วิปฺปกต\-โภชนํ อานนฺตริกํ ภิกฺขุํ วุฏฺฐาเปสฺสสิ, ภตฺตคฺคํ โกลาหลํ อโหสิ. เนตํ โมฆปุริส อปฺปสนฺนานํ วา ปสาทาย ฯเปฯ วิครหิตฺวา ฯเปฯ ธมฺมิํ กถํ กตฺวา ภิกฺขู อามนฺเตสิ “น ภิกฺขเว วิปฺปกต\-โภชโน\csromansharedfootnote{fc2a95ba6d54a207}{โภชโน อานนฺตริโก (สฺยา)} ภิกฺขู วุฏฺฐาเปตพฺโพ, โย วุฏฺฐาเปยฺย, อาปตฺติ ทุกฺกฏสฺส. สเจ วุฏฺฐาเปติ, ปวาริโต จ โหติ, ‘คจฺฉ อุทกํ อาหรา’ติ วตฺตพฺโพ. เอวญฺเจตํ ลเภถ, อิจฺเจตํ กุสลํ. โน เจ ลเภถ, สาธุกํ สิตฺถานิ คิลิตฺวา วุฑฺฒตรสฺส ภิกฺขุโน อาสนํ ทาตพฺพํ. น เตฺววาหํ ภิกฺขเว ‘เกนจิ ปริยาเยน วุฑฺฒตรสฺส ภิกฺขุโน อาสนํ ปฏิพาหิตพฺพนฺ’ติ วทามิ. โย ปฏิพาเหยฺย, อาปตฺติ ทุกฺกฏสฺสา”ติ.}

\prose{เตน โข ปน สมเยน ฉพฺพคฺคิยา ภิกฺขู คิลาเน ภิกฺขู วุฏฺฐาเปนฺติ. คิลานา เอวํ วเทนฺติ “น มยํ อาวุโส สกฺโกม วุฏฺฐาตุํ, คิลานามฺหา”ติ. “มยํ อายสฺมนฺเต วุฏฺฐาเปสฺสามา”ติ ปริคฺคเหตฺวา วุฏฺฐาเปตฺวา ฐิตเก มุญฺจนฺติ, คิลานา มุจฺฉิตา ปปตนฺติ. ภควโต เอตมตฺถํ \csromanfolio{316}อาโรเจสุํ. น ภิกฺขเว คิลาโน วุฏฺฐาเปตพฺโพ, โย วุฏฺฐาเปยฺย, อาปตฺติ ทุกฺกฏสฺสาติ.}

\prose{เตน โข ปน สมเยน ฉพฺพคฺคิยา ภิกฺขู “คิลานา มยมฺหา อวุฏฺฐาปนียา”ติ วรเสยฺยาโย ปลิพุนฺเธนฺติ\csromansharedfootnote{b4531de9d1411b3d}{ปลิพุนฺธนฺติ (ก)}. ภควโต เอตมตฺถํ อาโรเจสุํ. อนุชานามิ ภิกฺขเว คิลานสฺส ปติรูปํ เสยฺยํ ทาตุนฺติ.}

\prose{เตน โข ปน สมเยน ฉพฺพคฺคิยา ภิกฺขู เลสกปฺเปน เสนาสนํ ปฏิพาหนฺติ. ภควโต เอตมตฺถํ อาโรเจสุํ. น ภิกฺขเว เลสกปฺเปน เสนาสนํ ปฏิพาหิตพฺพํ, โย ปฏิพาเหยฺย, อาปตฺติ ทุกฺกฏสฺสาติ.}

\prose{\csromansymbolfootnote{*}{วิ 2. 63 ปิฏฺเฐปิ.}เตน โข ปน สมเยน สตฺตรส\-วคฺคิยา ภิกฺขู อญฺญตรํ ปจฺจนฺติมํ มหาวิหารํ ปฏิสงฺขโรนฺติ “อิธ มยํ วสฺสํ วสิสฺสามา”ติ. อทฺทสํสุ\csromansharedfootnote{22771c1fe10b17c8}{อทฺทสาสุํ (ก)} โข ฉพฺพคฺคิยา ภิกฺขู สตฺตรส\-วคฺคิเย ภิกฺขู วิหารํ ปฏิสงฺขโรนฺเต, ทิสฺวาน เอวมาหํสุ “อิเม อาวุโส สตฺตรส\-วคฺคิยา ภิกฺขู วิหารํ\csromansharedfootnote{2beb6482bc696a9f}{อญฺญตรํ วิหารํ (ก)} ปฏิสงฺขโรนฺติ, หนฺท เน วุฏฺฐาเปสฺสามา”ติ. เอกจฺเจ เอวมาหํสุ “อาคเมถาวุโส ยาว ปฏิสงฺขโรนฺติ, ปฏิสงฺขเต วุฏฺฐาเปสฺสามา”ติ. อถ โข ฉพฺพคฺคิยา ภิกฺขู สตฺตรส\-วคฺคิเย ภิกฺขู เอตทโวจุํ “อุฏฺเฐถาวุโส อมฺหากํ วิหาโร ปาปุณาตี”ติ. นนุ อาวุโส ปฏิกจฺเจว อาจิกฺขิตพฺพํ, มยญฺจญฺญํ ปฏิสงฺขเรยฺยามาติ. นนุ อาวุโส สํฆิโก วิหาโรติ. อามาวุโส สํฆิโก วิหาโรติ. อุฏฺเฐถาวุโส, อมฺหากํ วิหาโร ปาปุณาตีติ. มหลฺลโก อาวุโส วิหาโร, ตุเมฺหปิ วสถ, มยมฺปิ วสิสฺสามาติ. “อุฏฺเฐถาวุโส, อมฺหากํ วิหาโร ปาปุณาตี”ติ กุปิตา อนตฺตมนา คีวายํ คเหตฺวา นิกฺกฑฺฒนฺติ. เต นิกฺกฑฺฒิย\-มานา โรทนฺติ. ภิกฺขู เอวมาหํสุ “กิสฺส ตุเมฺห อาวุโส โรทถา”ติ. อิเม อาวุโส ฉพฺพคฺคิยา ภิกฺขู กุปิตา อนตฺตมนา อเมฺห สํฆิกา วิหารา นิกฺกฑฺฒนฺตีติ. เย เต ภิกฺขู อปฺปิจฺฉา ฯเปฯ เต อุชฺฌายนฺติ ขิยฺยนฺติ วิปาเจนฺติ “กถํ หิ นาม ฉพฺพคฺคิยา ภิกฺขู กุปิตา อนตฺตมนา ภิกฺขู สํฆิกา วิหารา นิกฺกฑฺฒิสฺสนฺตี”ติ. อถ โข เต ภิกฺขู ภควโต เอตมตฺถํ อาโรเจสุํ ฯเปฯ. สจฺจํ กิร ตุเมฺห ภิกฺขเว กุปิตา อนตฺตมนา สํฆิกา วิหารา ภิกฺขู นิกฺกฑฺฒถาติ. สจฺจํ ภควาติ ฯเปฯ วิครหิตฺวา ฯเปฯ ธมฺมิํ กถํ กตฺวา ภิกฺขู อามนฺเตสิ “น ภิกฺขเว \csromanfolio{317}กุปิเตน อนตฺตมเนน ภิกฺขู สํฆิกา วิหารา นิกฺกฑฺฒิตพฺโพ, โย นิกฺกฑฺเฒยฺย, ยถาธมฺโม กเรตพฺโพ. อนุชานามิ ภิกฺขเว เสนาสนํ คาเหตุนฺ”ติ.}

\csromanmatikaanchor{mtk.148}
\csromantocmarkref{subsection}{เสนาสนคฺคาหาปกสมฺมุติ}{mtk.148}
\bookmark[dest={mtk.148},level=2]{เสนาสนคฺคาหาปกสมฺมุติ}
\csromanheader{2}{เสนาสนคฺคาหาปก\-สมฺมุติ}
\proseitem{317}{อถ โข ภิกฺขูนํ เอตทโหสิ “เกน นุ โข เสนาสนํ คาเหตพฺพนฺ”ติ. ภควโต เอตมตฺถํ อาโรเจสุํ ฯเปฯ. อนุชานามิ ภิกฺขเว ปญฺจหงฺเคหิ สมนฺนาคตํ ภิกฺขุํ เสนาสน\-คฺคาหาปกํ สมฺมนฺนิตุํ. โย น ฉนฺทาคติํ คจฺเฉยฺย, น โทสาคติํ คจฺเฉยฺย, น โมหาคติํ คจฺเฉยฺย, น ภยาคติํ คจฺเฉยฺย, คหิตาคหิตญฺจ ชาเนยฺย. เอวญฺจ ปน ภิกฺขเว สมฺมนฺนิตพฺโพ.}

\prose{ปฐมํ ภิกฺขู ยาจิตพฺโพ, ยาจิตฺวา พฺยตฺเตน ภิกฺขุนา ปฏิพเลน สํโฆ ญาเปตพฺโพ\csromandash{}}

\prose{“สุณาตุ เม ภนฺเต สํโฆ, ยทิ สํฆสฺส ปตฺตกลฺลํ, สํโฆ อิตฺถนฺนามํ ภิกฺขุํ เสนาสน\-คฺคาหาปกํ สมฺมนฺเนยฺย, เอสา ญตฺติ.}

\prose{สุณาตุ เม ภนฺเต สํโฆ, สํโฆ อิตฺถนฺนามํ ภิกฺขุํ เสนาสน\-คฺคาหาปกํ สมฺมนฺนติ, ยสฺสายสฺมโต ขมติ อิตฺถนฺนามสฺส ภิกฺขุโน เสนาสน\-คฺคาหาปกสฺส สมฺมุติ, โส ตุณฺหสฺส, ยสฺส นกฺขมติ, โส ภาเสยฺย.}

\prose{สมฺมโต สํเฆน อิตฺถนฺนาโม ภิกฺขุ เสนาสน\-คฺคาหาปโก, ขมติ สํฆสฺส, ตสฺมา ตุณฺหี, เอวเมตํ ธารยามี”ติ.}

\proseitem{318}{อถ โข เสนาสน\-คฺคาหาปกานํ ภิกฺขุนํ เอตทโหสิ “กถํ นุ โข เสนาสนํ คาเหตพฺพนฺ”ติ. ภควโต เอตมตฺถํ อาโรเจสุํ. อนุชานามิ ภิกฺขเว ปฐมํ ภิกฺขู คเณตุํ, ภิกฺขู คเณตฺวา เสยฺยา คเณตุํ, เสยฺยา คเณตฺวา เสยฺยคฺเคน คาเหตุนฺติ. เสยฺยคฺเคน คาเหนฺตา เสยฺยา อุสฺสารยิํสุ ฯเปฯ. อนุชานามิ ภิกฺขเว วิหารคฺเคน คาเหตุนฺติ. วิหารคฺเคน คาเหนฺตา วิหารา อุสฺสารยิํสุ ฯเปฯ. อนุชานามิ ภิกฺขเว ปริเวณคฺเคน คาเหตุนฺติ. ปริเวณคฺเคน คาเหนฺตา ปริเวณา อุสฺสารยิํสุ ฯเปฯ. อนุชานามิ ภิกฺขเว อนุภาคมฺปิ ทาตุํ, คหิเต อนุภาเค อญฺโญ ภิกฺขุ อาคจฺฉติ, อกามา น ทาตพฺโพติ.}

\prose{\csromanfolio{318}เตน โข ปน สมเยน ภิกฺขู นิสฺสีเม ฐิตสฺส เสนาสนํ คาเหนฺติ. ภควโต เอตมตฺถํ อาโรเจสุํ. น ภิกฺขเว นิสฺสีเม ฐิตสฺสํ เสนาสนํ คาเหตพฺพํ, โย คาเหยฺย, อาปตฺติ ทุกฺกฏสฺสาติ.}

\prose{เตน โข ปน สมเยน ภิกฺขู เสนาสนํ คเหตฺวา สพฺพกาลํ ปฏิพาหนฺติ. ภควโต เอตมตฺถํ อาโรเจสุํ. น ภิกฺขเว เสนาสนํ คเหตฺวา สพฺพกาลํ ปฏิพาหิตพฺพํ, โย ปฏิพาเหยฺย, อาปตฺติ ทุกฺกฏสฺส. อนุชานามิ ภิกฺขเว วสฺสานํ เตมาสํ ปฏิพาหิตุํ, อุตุกาลํ ปน น ปฏิพาหิตุนฺติ.}

\prosewithcloser{อถ โข ภิกฺขูนํ เอตทโหสิ “กติ นุ โข เสนาสนคฺคาหา”ติ. ภควโต เอตมตฺถํ อาโรเจสุํ. ตโย เม ภิกฺขเว เสนาสนคฺคาหา, ปุริมโก ปจฺฉิมโก อนฺตรามุตฺตโก. อปรชฺชุคตาย อาสาฬฺหิยา ปุริมโก คาเหตพฺโพ, มาสคตาย อาสาฬฺหิยา ปจฺฉิมโก คาเหตพฺโพ, อปรชฺชุคตาย ปวารณาย อายติํ วสฺสาวาส\-ตฺถาย อนฺตรามุตฺตโก คาเหตพฺโพ. อิเม โข ภิกฺขเว ตโย เสนาสนคฺคาหาติ.}{ทุติย\-ภาณวาโร นิฏฺฐิโต.}
\csromanmatikaanchor{mtk.149}
\csromantocmarkref{section}{3. ตติยภาณวาร}{mtk.149}
\bookmark[dest={mtk.149},level=1]{3. ตติยภาณวาร}
\csromanheader{1}{3. ตติย\-ภาณวาร}
\proseitem{319}{เตน โข ปน สมเยน อายสฺมา อุปนนฺโท สกฺยปุตฺโต สาวตฺถิยํ เสนาสนํ คเหตฺวา อญฺญตรํ คามกาวาสํ อคมาสิ, ตตฺถปิ เสนาสนํ อคฺคเหสิ. อถ โข เตสํ ภิกฺขูนํ เอตทโหสิ “อยํ อาวุโส อายสฺมา อุปนนฺโท สกฺยปุตฺโต ภณฺฑนการโก กลหการโก วิวาทการโก ภสฺสการโก สํเฆ อธิกรณ\-การโก, สจายํ อิธ วสฺสํ วสิสฺสติ, สพฺเพว มยํ น ผาสุ ภวิสฺสาม, หนฺท นํ ปุจฺฉามา”ติ. อถ โข เต ภิกฺขู อายสฺมนฺตํ อุปนนฺทํ สกฺยปุตฺตํ เอตทโวจุํ “นนุ ตยา อาวุโส อุปนนฺท สาวตฺถิยํ เสนาสนํ คหิตนฺ”ติ. เอวมาวุโสติ. กิํ ปน ตฺวํ อาวุโส อุปนนฺท เอโก เทฺว ปฏิพาหสีติ. อิธทานาหํ อาวุโส มุญฺจามิ, ตตฺถ คณฺหามีติ. เย เต ภิกฺขู อปฺปิจฺฉา ฯเปฯ เต อุชฺฌายนฺติ ขิยฺยนฺติ วิปาเจนฺติ “กถํ หิ นาม อายสฺมา อุปนนฺโท \csromanfolio{319}สกฺยปุตฺโต เอโก เทฺว ปฏิพาเหสฺสตี”ติ. ภควโต เอตมตฺถํ อาโรเจสุํ. สจฺจํ กิร ตฺวํ อุปนนฺท เอโก เทฺว ปฏิพาหสีติ. สจฺจํ ภควาติ. วิครหิ พุทฺโธ ภควา ฯเปฯ. กถํ หิ นาม ตฺวํ โมฆปุริส เอโก เทฺว ปฏิพาหิสฺสสิ, ตตฺถ ตยา โมฆปุริส คหิตํ อิธ มุตฺตํ, อิธ ตยา คหิตํ ตตฺร มุตฺตํ, เอวํ โข ตฺวํ โมฆปุริส อุภยตฺถ ปริพาหิโร. เนตํ โมฆปุริส อปฺปสนฺนานํ วา ปสาทาย ฯเปฯ วิครหิตฺวา ฯเปฯ ธมฺมิํ กถํ กตฺวา ภิกฺขู อามนฺเตสิ “น ภิกฺขเว เอเกน เทฺว ปฏิพาหิตพฺพา, โย ปฏิพาเหยฺย, อาปตฺติ ทุกฺกฏสฺสา”ติ.}

\proseitem{320}{\csromansymbolfootnote{*}{วิ 2. 187 ปิฏฺเฐปิ.}เตน โข ปน สมเยน ภควา ภิกฺขูนํ อเนก\-ปริยาเยน วินยกถํ กเถสิ, วินยสฺส วณฺณํ ภาสติ, วินย\-ปริยตฺติยา วณฺณํ ภาสติ, อาทิสฺส อาทิสฺส อายสฺมโต อุปาลิสฺส วณฺณํ ภาสติ. ภิกฺขูนํ เอตทโหสิ “ภควา โข อเนก\-ปริยาเยน วินยกถํ กเถติ, วินยสฺส วณฺณํ ภาสติ, วินย\-ปริยตฺติยา วณฺณํ ภาสติ, อาทิสฺส อาทิสฺส อายสฺมโต อุปาลิสฺส วณฺณํ ภาสติ, หนฺท มยํ อาวุโส อายสฺมโต อุปาลิสฺส สนฺติเก วินยํ ปริยาปุณามา”ติ. เต'ธ\csromansharedfootnote{5ff24605c9dcba0e}{เต จ (สฺยา, ก)} พหู ภิกฺขู เถรา จ นวา จ มชฺฌิมา จ อายสฺมโต อุปาลิสฺส สนฺติเก วินยํ ปริยาปุณนฺติ, อายสฺมา อุปาลิ ฐิตโกว อุทฺทิสติ เถรานํ ภิกฺขูนํ คารเวน, เถราปิ ภิกฺขู ฐิตกาว อุทฺทิสาเปนฺติ ธมฺมคารเวน, ตตฺถ เถรา เจว ภิกฺขู กิลมนฺติ, อายสฺมา จ อุปาลิ กิลมติ. ภควโต เอตมตฺถํ อาโรเจสุํ. อนุชานามิ ภิกฺขเว นวเกน ภิกฺขุนา อุทฺทิสนฺเตน สมเก วา อาสเน นิสีทิตุํ อุจฺจตเร วา ธมฺมคารเวน, เถเรน ภิกฺขุนา อุทฺทิสาเปนฺเตน สมเก วา อาสเน นิสีทิตุํ นีจตเร วา ธมฺมคารเวนาติ.}

\prose{เตน โข ปน สมเยน พหู ภิกฺขู อายสฺมโต อุปาลิสฺส สนฺติเก ฐิตกา อุทฺเทสํ ปฏิมาเนนฺตา กิลมนฺติ. ภควโต เอตมตฺถํ อาโรเจสุํ. อนุชานามิ ภิกฺขเว สมานาสนิเกหิ สห นิสีทิตุนฺติ. อถ โข ภิกฺขูนํ เอตทโหสิ “กิตฺตาวตา นุ โข สมานาสนิโก โหตี”ติ. ภควโต เอตมตฺถํ อาโรเจสุํ. อนุชานามิ ภิกฺขเว ติวสฺสนฺตเรน สห นิสีทิตุนฺติ.}

\prose{\csromanfolio{320}เตน โข ปน สมเยน สมฺพหุลา ภิกฺขู สมานาสนิกา มญฺเจ\csromansharedfootnote{be60bcf12d40b1f4}{เอกมญฺเจ (สฺยา)} นิสีทิตฺวา มญฺจํ ภินฺทิํสุ, ปีเฐ\csromansharedfootnote{e1b4241766cded70}{เอกปีเฐ (สฺยา)} นิสีทิตฺวา ปีฐํ ภินฺทิํสุ. ภควโต เอตมตฺถํ อาโรเจสุํ. อนุชานามิ ภิกฺขเว ติวคฺคสฺส มญฺจํ ติวคฺคสฺส ปีฐนฺติ. ติวคฺโคปิ มญฺเจ นิสีทิตฺวา มญฺจํ ภินฺทิ. ปีเฐ นิสีทิตฺวา ปีฐํ ภินฺทิ ฯเปฯ. อนุชานามิ ภิกฺขเว ทุวคฺคสฺส มญฺจํ ทุวคฺคสฺส ปีฐนฺติ.}

\prose{เตน โข ปน สมเยน ภิกฺขู อสมานาสนิเกหิ สห ทีฆาสเน นิสีทิตุํ กุกฺกุจฺจายนฺติ. ภควโต เอตมตฺถํ อาโรเจสุํ. อนุชานามิ ภิกฺขเว ฐเปตฺวา ปณฺฑกํ มาตุคามํ อุภโต\-พฺยญฺชนกํ อสมานาสนิเกหิ สห ทีฆาสเน นิสีทิตุนฺติ. อถ โข ภิกฺขูนํ เอตทโหสิ “กิตฺตกํ ปจฺฉิมํ นุ โข ทีฆาสนํ โหตี”ติ. ภควโต เอตมตฺถํ อาโรเจสุํ. อนุชานามิ ภิกฺขเว ยํ ติณฺณํ ปโหติ, เอตฺตกํ ปจฺฉิมํ ทีฆาสนนฺติ.}

\prose{เตน โข ปน สมเย วิสาขา มิคารมาตา สํฆสฺส อตฺถาย สาฬินฺทํ ปาสาทํ การาเปตุกามา โหติ หตฺถินขกํ. อถ โข ภิกฺขูนํ เอตทโหสิ “กิํ นุ โข ภควตา ปาสาท\-ปริโภโค อนุญฺญาโต, กิํ อนนุญฺญาโต”ติ. ภควโต เอตมตฺถํ อาโรเจสุํ. อนุชานามิ ภิกฺขเว สพฺพํ ปาสาท\-ปริโภคนฺติ.}

\prose{เตน โข ปน สมเยน รญฺโญ ปเสนทิสฺส โกสลสฺส อยฺยิกา กาลงฺกตา โหติ, ตสฺส กาลงฺกิริยาย สํฆสฺส พหุํ อกปฺปิย\-ภณฺฑํ อุปนฺนํ โหติ, เสยฺยถิทํ, อาสนฺทิ ปลฺลงฺโก โคนโก จิตฺตโก ปฏิกา ปฏลิกา ตูลิกา วิกติกา อุทฺทโลมี เอกนฺตลิมี กฏฺฏิสฺสํ โกเสยฺยํ กุตฺตกํ หตฺถตฺถรํ อสฺสตฺถรํ รถตฺถรํ อชินปฺปเวณิ กทลิมิค\-ปฺปวร\-ปจฺจตฺถรณํ สอุตฺตรจฺฉทํ อุภโต\-โลหิตกู\-ปธานํ. ภควโต เอตมตฺถํ อาโรเจสุํ. อนุชานามิ ภิกฺขเว อาสนฺทิยา ปาเท ฉินฺทิตฺวา ปริภุญฺชิตุํ. ปลฺลงฺกสฺส วาเฬ ภินฺทิตฺวา ปริภุญฺชิตุํ. ตูลิกํ วิชเฏตฺวา พิพฺโพหนํ กาตุํ. อวเสสํ ภูมตฺถรณํ\csromansharedfootnote{69556c887454ad23}{ภุมฺมตฺถรณํ (สี, สฺยา)} กาตุนฺติ.}

\csromanmatikaanchor{mtk.150}
\csromantocmarkref{subsection}{อวิสฺสชฺชิยวตฺถุ}{mtk.150}
\bookmark[dest={mtk.150},level=2]{อวิสฺสชฺชิยวตฺถุ}
\csromanheader{2}{อวิสฺสชฺชิย\-วตฺถุ}
\proseitem{321}{\csromanfolio{321}เตน โข ปน สมเยน สาวตฺถิยา อวิทูเร อญฺญตรสฺมิํ คามกาวาเส อาวาสิกา ภิกฺขู อุปทฺทุตา โหนฺติ อาคนฺตุก\-คมิกานํ ภิกฺขูนํ เสนาสนํ ปญฺญเปนฺตา. อถ โข เตสํ ภิกฺขูนํ เอตทโหสิ “เอตรหิ โข มยํ อาวุโส อุปทฺทุตา อาคนฺตุก\-คมิกานํ ภิกฺขูนํ เสนาสนํ ปญฺญเปนฺตา, หนฺท มยํ อาวุโส สพฺพํ สํฆิกํ เสนาสนํ เอกสฺส เทม, ตสฺส สนฺตกํ ปริภุญฺชิสฺสามา”ติ. เต สพฺพํ สํฆิกํ เสนาสนํ เอกสฺส อทํสุ. อาคนฺตุกา ภิกฺขู เต ภิกฺขู เอตทโวจุํ “อมฺหากํ อาวุโส เสนาสนํ ปญฺญเปถา”ติ. นตฺถาวุโส สํฆิกํ เสนาสนํ, สพฺพํ อมฺหาหิ เอกสฺส ทินฺนนฺติ. กิํ ปน ตุเมฺห อาวุโส สํฆิกํ เสนาสนํ วิสฺสชฺเชถาติ. เอวมาวุโสติ. เย เต ภิกฺขู อปฺปิจฺฉา ฯเปฯ เต อุชฺฌายนฺติ ขิยฺยนฺติ วิปาเจนฺติ “กถํ หิ นาม ภิกฺขู สํฆิกํ เสนาสนํ วิสฺสชฺเชสฺสนฺตี”ติ\csromansharedfootnote{e1d597037766355e}{วิสฺสชฺชิสฺสนฺตีติ (ก)}. อถ โข เต ภิกฺขู ภควโต เอตมตฺถํ อาโรเจสุํ ฯเปฯ. สจฺจํ กิร ภิกฺขเว ภิกฺขู สํฆิกํ เสนาสนํ วิสฺสชฺเชนฺตีติ. สจฺจํ ภควาติ. วิครหิ พุทฺโธ ภควา ฯเปฯ. กถํ หิ นาม เต ภิกฺขเว โมฆปุริสา สํฆิกํ เสนาสนํ วิสฺสชฺเชสฺสนฺติ. เนตํ ภิกฺขเว อปฺปสนฺนานํ วา ปสาทย ฯเปฯ วิครหิตฺวา ฯเปฯ ธมฺมิํ กถํ กตฺวา ภิกฺขู อามนฺเตสิ\csromandash{}}

\prose{ปญฺจิมานิ ภิกฺขเว อวิสฺสชฺชิยานิ, น วิสฺสชฺเช\-ตพฺพานิ\csromansharedfootnote{190a54ebf62c6599}{น วิสฺสชฺชิตพฺพานิ (ก)} สํเฆน วา คเณน วา ปุคฺคเลน วา, วิสฺสชฺชิตานิปิ อวิสฺสชฺชิตานิ โหนฺติ, โย วิสฺสชฺเชยฺย, อาปตฺติ ถุลฺลจฺจยสฺส. กตมานิ ปญฺจ. อาราโม, อารามวตฺถุ, อิทํ ปฐมํ อวิสฺสชฺชิยํ, น วิสฺสชฺเชตพฺพํ สํเฆน วา คเณน วา ปุคฺคเลน วา วิสฺสชฺชิตมฺปิ อวิสฺสชฺชิตํ โหติ, โย วิสฺสชฺเชยฺย, อาปตฺติ ถุลฺลจฺจยสฺส.}

\prose{วิหาโร, วิหารวตฺถุ, อิทํ ทุติยํ อวิสฺสชฺชิยํ, น วิสฺสชฺเชตพฺพํ, สํเฆน วา คเณน วา ปุคฺคเลน วา, วิสฺสชฺชิตมฺปิ อวิสฺสชฺชิตํ โหติ, โย วิสฺสชฺเชยฺย, อาปตฺติ ถุลฺลจฺจยสฺส.}

\prose{มญฺโจ, ปีฐํ, ภิสิ, พิพฺโพหนํ, อิทํ ตติยํ อวิสฺสชฺชิยํ, น วิสฺสชฺเชตพฺพํ สํเฆน วา คเณน วา ปุคฺคเลน วา, วิสฺสชฺชิตมฺปิ อวิสฺสชฺชิตํ โหติ, โย วิสฺสชฺเชยฺย, อาปตฺติ ถุลฺลจฺจยสฺส.}

\prose{\csromanfolio{322}โลหกุมฺภี, โลหภาณกํ, โลหวารโก, โลหกฏาหํ, วาสิ, ปรสุ\csromansharedfootnote{ab7fda305bddc2c2}{ผรสุ (สี, สฺยา, ก)}, กุฐารี\csromansharedfootnote{7434db6a7e44588b}{กุธารี (ก)}, กุทาโล, นิขาทนํ, อิทํ จตุตฺถํ อวิสฺสชฺชิยํ, น วิสฺสชฺเชตพฺพํ สํเฆน วา คเณน วา ปุคฺคเลน วา, วิสฺสชฺชิตมฺปิ อวิสฺสชฺชิตํ โหติ, โย วิสฺสชฺเชยฺย, อาปตฺติ ถุลฺลจฺจยสฺส.}

\prose{วลฺลิ, เวฬุ, มุญฺชํ, ปพฺพชํ\csromansharedfootnote{97cf7925dc26e945}{พพฺพชํ (สี)}, มตฺติกา, ทารุภณฺฑํ, มตฺติกาภณฺฑํ, อิทํ ปญฺจมํ อวิสฺสชฺชิยํ, น วิสฺสชฺเชตพฺพํ สํเฆน วา คเณน วา ปุคฺคเลน วา, วิสฺสชฺชิตมฺปิ อวิสฺสชฺชิตํ โหติ, โย วิสฺสชฺเชยฺย, อาปตฺติ ถุลฺลจฺจยสฺส. อิมานิ โข ภิกฺขเว ปญฺจ อวิสฺสชฺชิยานิ, น วิสฺสชฺเช\-ตพฺพานิ สํเฆน วา คเณน วา ปุคฺคเลน วา, วิสฺสชฺชิตานิปิ อวิสฺสชฺชิตานิ โหนฺติ, โย วิสฺสชฺเชยฺย, อาปตฺติ ถุลฺลจฺจยสฺสาติ.}

\csromanmatikaanchor{mtk.151}
\csromantocmarkref{subsection}{อเวภงฺคิยวตฺถุ}{mtk.151}
\bookmark[dest={mtk.151},level=2]{อเวภงฺคิยวตฺถุ}
\csromanheader{2}{อเวภงฺคิย\-วตฺถุ}
\proseitem{322}{อถ โข ภควา สาวตฺถิยํ ยถาภิรนฺตํ วิหริตฺวา เยน กีฏาคิริ เตน จาริกํ ปกฺกามิ มหตา ภิกฺขุ\-สํเฆน สทฺธิํ ปญฺจมตฺเตหิ ภิกฺขุสเตหิ สาริปุตฺต\-โมคฺคลฺลาเนหิ จ. อสฺโสสุํ โข อสฺสชิ\-ปุนพฺพสุกา ภิกฺขู “ภควา กิร กีฏาคิริํ อาคจฺฉติ มหตา ภิกฺขุ\-สํเฆน สทฺธิํ ปญฺจมตฺเตหิ ภิกฺขุสเตหิ สาริปุตฺต\-โมคฺคลฺลาเนหิ จ. หนฺท มยํ อาวุโส สพฺพํ สํฆิกํ เสนาสนํ ภาเชม, ปาปิจฺฉา สาริปุตฺต\-โมคฺคลฺลานา ปาปิกานํ อิจฺฉานํ วสํ คตา, น มยํ เตสํ เสนาสนํ ปญฺญเปสฺสามา”ติ. เต สพฺพํ สํฆิกํ เสนาสนํ ภาเชสุํ. อถ โข ภควา อนุปุพฺเพน จาริกํ จรมาโน เยน กีฏาคิริ ตทวสริ. อถ โข ภควา สมฺพหุเล ภิกฺขู อามนฺเตสิ\csromandash{}คจฺฉถ ตุเมฺห ภิกฺขเว, อสฺสชิปุนพฺพสุเก ภิกฺขู อุปสงฺกมิตฺวา เอวํ วเทถ “ภควา อาวุโส อาคจฺฉติ มหตา ภิกฺขุ\-สํเฆน สทฺธิํ ปญฺจมตฺเตหิ ภิกฺขุสเตหิ สาริปุตฺต\-โมคฺคลฺลาเนหิ จ, ภควโต จ อาวุโส เสนาสนํ ปญฺญเปถ ภิกฺขุ\-สํฆสฺส จ สาริปุตฺต\-โมคฺคลฺลานานญฺจา”ติ. “เอวํ ภนฺเต”ติ โข เต ภิกฺขู ภควโต ปฏิสฺสุตฺวา เยน อสฺสชิ\-ปุนพฺพสุกา ภิกฺขู เตนุ\-ปสงฺกมิํสุ, อุปสงฺกมิตฺวา อสฺสชิปุนพฺพสุเก ภิกฺขู เอตทโวจุํ “ภควา อาวุโส อาคจฺฉติ มหตา ภิกฺขุ\-สํเฆน สทฺธิํ ปญฺจมตฺเตหิ ภิกฺขุสเตหิ สาริปุตฺต\-โมคฺคลฺลาเนหิ จ, ภควโต จ อาวุโส \csromanfolio{323}เสนาสนํ ปญฺญเปถ ภิกฺขุ\-สํฆสฺส จ สาริปุตฺต\-โมคฺคลฺลานานญฺจา”ติ. นตฺถาวุโส สํฆิกํ เสนาสนํ, สพฺพํ อเมฺหหิ ภาชิตํ, สฺวาคตํ อาวุโส ภควโต, ยสฺมิํ วิหาเร ภควา อิจฺฉิสฺสติ, ตสฺมิํ วิหาเร วสิสฺสติ, ปาปิจฺฉา สาริปุตฺต\-โมคฺคลฺลานา ปาปิกานํ อิจฺฉานํ วสํ คตา, น มยํ เตสํ เสนสนํ ปญฺญเปสฺสามาติ. กิํ ปน ตุเมฺห อาวุโส สํฆิกํ เสนาสนํ ภาชิตฺถาติ. เอวมาวุโสติ. เย เต ภิกฺขู อปฺปิจฺฉา ฯเปฯ เต อุชฺฌายนฺติ ขิยฺยนฺติ วิปาเจนฺติ “กถํ หิ นาม อสฺสชิ\-ปุนพฺพสุกา ภิกฺขู สํฆิกํ เสนาสนํ ภาเชสฺสนฺตี”ติ. อถ โข เต ภิกฺขู ภควโต เอตมตฺถํ อาโรเจสุํ ฯเปฯ. สจฺจํ กิร ภิกฺขเว ฯเปฯ. สจฺจํ ภควาติ. วิครหิ พุทฺโธ ภควา ฯเปฯ. กถํ หิ นาม เต ภิกฺขเว โมฆปุริสา สํฆิกํ เสนาสนํ ภาเชสฺสนฺติ. เนตํ ภิกฺขเว อปฺปสนฺนานํ วา ปสาทาย ฯเปฯ วิครหิตฺวา ฯเปฯ ธมฺมิํ กถํ กตฺวา ภิกฺขู อามนฺเตสิ\csromandash{}}

\prose{ปญฺจิมานิ ภิกฺขเว อเวภงฺคิยานิ\csromansharedfootnote{db74d5d9d8cf2bc7}{อเวภงฺคิกานิ (ก)}, น วิภชิตพฺพานิ สํเฆน วา คเณน วา ปุคฺคเลน วา, วิภตฺตานิปิ อวิภตฺตานิ โหนฺติ, โย วิภเชยฺย, อาปตฺติ ถุลฺลจฺจยสฺส. กตมานิ ปญฺจ. อาราโม, อารามวตฺถุ, อิทํ ปฐมํ อเวภงฺคิยํ, น วิภชิตพฺพํ สํเฆน วา คเณน วา ปุคฺคเลน วา, วิภตฺตมฺปิ อวิภตฺตํ โหติ, โย วิภเชยฺย, อาปตฺติ ถุลฺลจฺจยสฺส.}

\prose{วิหาโร, วิหารวตฺถุ, อิทํ ทุติยํ อเวภงฺคิยํ, น วิภชิตพฺพํ สํเฆน วา คเณน วา ปุคฺคเลน วา, วิภตฺตมฺปิ อวิภตฺตํ โหติ, โย วิภเชยฺย, อาปตฺติ ถุลฺลจฺจยสฺส.}

\prose{มญฺโจ, ปีฐํ, ภิสิ, พิพฺโพหนํ, อิทํ ตติยํ อเวงฺคิยํ, น วิภชิตพฺพํ สํเฆน วา คเณน วา ปุคฺคเลน วา, วิภตฺตมฺปิ อวิภตฺตํ โหติ, โย วิภเชยฺย, อาปตฺติ ถุลฺลจฺจยสฺส.}

\prose{โลหกุมฺภี, โลหภาณกํ, โลหวารโก, โลหกฏาหํ, วาสี, ปรสุ, กุฐารี, กุทาโล, นิขาทนํ, อิทํ จตุตฺถํ อเวภงฺคิยํ, น วิภชิตพฺพํ สํเฆน วา คเณน วา ปุคฺคเลน วา, วิภตฺตมฺปิ อวิภตฺตํ โหติ, โย วิภเชยฺย, อาปตฺติ ถุลฺลจฺจยสฺส.}

\prose{วลฺลี, เวฬุ, มุญฺชํ, ปพฺพชํ, ติณํ, มตฺติกา, ทารุภณฺฑํ, มตฺติกาภณฺฑํ, อิทํ ปญฺจมํ อเวภงฺคิยํ, น วิภชิตพฺพํ สํเฆน วา คเณน วา ปุคฺคเลน วา, \csromanfolio{324}วิภตฺตํปิ อวิภตฺตํ โหติ, โย วิภเชยฺย, อาปตฺติ ถุลฺลจฺจยสฺส. อิมานิ โข ภิกฺขเว ปญฺจ ภิกฺขเว ปญฺจ อเวภงฺคิยานิ, น วิภชิตพฺพานิ สํเฆน วา คเณน วา ปุคฺคเลน วา, วิภตฺตานิปิ อวิภตฺตานิ โหนฺติ, โย วิภเชยฺย, อาปตฺติ ถุลฺลจฺจยสฺสาติ.}

\csromanmatikaanchor{mtk.152}
\csromantocmarkref{subsection}{นวกมฺมทานกถา}{mtk.152}
\bookmark[dest={mtk.152},level=2]{นวกมฺมทานกถา}
\csromanheader{2}{นวกมฺม\-ทาน\-กถา}
\proseitem{323}{อถ โข ภควา กีฏาคิริสฺมิํ ยถาภิรนฺตํ วิหริตฺวา เยน อาฬวี เตน จาริกํ ปกฺกามิ. อนุปุพฺเพน จาริกํ จรมาโน เยน อาฬวี ตทวสริ, ตตฺร สุทํ ภควา อาฬวิยํ วิหรติ อคฺคาฬเว เจติเย. เตน โข ปน สมเยน อาฬวกา\csromansharedfootnote{abe21c3da7cb608d}{อาฬวิกา (สฺยา, ก)} ภิกฺขู เอวรุปานิ นวกมฺมานิ เทนฺติ, ปิณฺฑ\-นิกฺเขปน\-มตฺเตนปิ นวกมฺมํ เทนฺติ, กุฏฺฏ\-เลปน\-มตฺเตนปิ นวกมฺมํ เทนฺติ, ทฺวารฏฺฐปน\-มตฺเตนปิ นวกมฺมํ เทนฺติ, อคฺคฬวฏฺฏิ\-กรณ\-มตฺเตนปิ นวกมฺมํ เทนฺติ, อาโลกสนฺธิ\-กรณ\-มตฺเตนปิ นวกมฺมํ เทนฺติ, เสตวณฺณ\-กรณ\-มตฺเตนปิ นวกมฺมํ เทนฺติ, กาฬวณฺณ\-กรณ\-มตฺเตนปิ นวกมฺมํ เทนฺติ, เครุกปริกมฺมกรณ- มตฺเตนปิ นวกมฺมํ เทนฺติ, ฉาทนมตฺเตนปิ นวกมฺมํ เทนฺติ, พนฺธน\-มตฺเตนปิ นวกมฺมํ เทนฺติ, ภณฺฑิกา\-ฏฺฐปน\-มตฺเตนปิ นวกมฺมํ เทนฺติ, ขณฺฑผุลฺล\-ปฏิสงฺขรณ\-มตฺเตนปิ นวกมฺมํ เทนฺติ, ปริภณฺฑ\-กรณ\-มตฺเตนปิ นวกมฺมํ เทนฺติ, วีสติ\-วสฺสิกมฺปิ นวกมฺมํ เทนฺติ, ติํสวสฺสิกมฺปิ นวกมฺมํ เทนฺติ, ยาวชีวิกมฺปิ นวกมฺมํ เทนฺติ, ธูมกาลิกมฺปิ ปริโยสิตํ วิหารํ นวกมฺมํ เทนฺติ. เย เต ภิกฺขู อปฺปิจฺฉา ฯเปฯ เต อุชฺฌายนฺติ ขิยฺยนฺติ วิปาเจนฺติ “กถํ หิ นาม อาฬวกา ภิกฺขู เอวรูปานิ นวกมฺมานิ ทสฺสนฺติ, ปิณฺฑ\-นิกฺเขปน\-มตฺเตนปิ นวกมฺมํ ทสฺสนฺติ, กุฏฺฏ\-เลปน\-มตฺเตนปิ. ทฺวารฏฺฐปน\-มตฺเตนปิ. อคฺคฬวฏฺฏิ\-กรณ\-มตฺเตนปิ. อาโลกสนฺธิกรณ- มตฺเตนปิ. เสตวณฺณ\-กรณ\-มตฺเตนปิ. กาฬวณฺณ\-กรณ\-มตฺเตนปิ. เครุกปริกมฺม\-กรณ\-มตฺเตนปิ. ฉาทนมตฺเตนปิ. พนฺธน\-มตฺเตนปิ. ภณฺฑิกา\-ฏฺฐปน\-มตฺเตนปิ. ขณฺฑผุลฺล\-ปฏิสงฺขรณ\-มตฺเตนปิ. ปริภณฺฑ\-กรณ\-มตฺเตนปิ. วีสติ\-วสฺสิกมฺปิ. ติํสวสฺสิกมฺปิ. ยาวชีวิกมฺปิ. ธูมกาลิ กมฺปิ ปริโยสิตํ วิหารํ นวกมฺมํ ทสฺสนฺตี”ติ. ภควโต เอตมตฺถํ อาโรเจสุํ ฯเปฯ. สจฺจํ กิร ภิกฺขเว ฯเปฯ. สจฺจํ ภควาติ ฯเปฯ วิครหิตฺวา ฯเปฯ ธมฺมิํ กถํ กตฺวา ภิกฺขู อามนฺเตสิ “น ภิกฺขเว ปิณฺฑ\-นิกฺเขปน\-มตฺเตน นวกมฺมํ ทาตพฺพํ ฯเปฯ. น กุฏฺฏ\-เลปน\-มตฺเตน นวกมฺมํ ทาตพฺพํ. \csromanfolio{325}น ทฺวารฏฺฐปน\-มตฺเตน นวกมฺมํ ทาตพฺพํ. น อคฺคฬวฏฺฏิกรณ- มตฺเตน นวกมฺมํ ทาตพฺพํ. น อาโลกสนฺธิ\-กรณ\-มตฺเตน นวกมฺมํ ทาตพฺพํ. น เสตวณฺณ\-กรณ\-มตฺเตน นวกมฺมํ ทาตพฺพํ. น กาฬวณฺณ\-กรณ\-มตฺเตน นวกมฺมํ ทาตพฺพํ. น เครุกปริกมฺม\-กรณ\-มตฺเตน นวกมฺมํ ทาตพฺพํ. น ฉาทนมตฺเตน นวกมฺมํ ทาตพฺพํ. น พนฺธน\-มตฺเตน นวกมฺมํ ทาตพฺพํ. น ภณฺฑิกา\-ฏฺฐปน\-มตฺเตน นวกมฺมํ ทาตพฺพํ. น ขณฺฑผุลฺล\-ปฏิสงฺขรณ\-มตฺเตน นวกมฺมํ ทาตพฺพํ. น ปริภณฺฑ\-กรณ\-มตฺเตน นวกมฺมํ ทาตพฺพํ. น วีสติวสฺสิกํ นวกมฺมํ ทาตพฺพํ. น ติํสวสฺสิกํ นวกมฺมํ ทาตพฺพํ. น ยาวชีวิกํ นวกมฺมํ ทาตพฺพํ. น ธูมกาลิกมฺปิ ปริโยสิตํ วิหารํ นวกมฺมํ ทาตพฺพํ, โย ทเทยฺย, อาปตฺติ ทุกฺกฏสฺส. อนุชานามิ ภิกฺขเว อกตํ วา วิปฺปกตํ วา นวกมฺมํ ทาตุํ, ขุทฺทเก วิหาเร กมฺมํ โอโลเกตฺวา ฉปฺปญฺจ\-วสฺสิกํ นวกมฺมํ ทาตุํ, อฑฺฒโยเค กมฺมํ โอโลเกตฺวา สตฺต\-ฏฺฐ\-วสฺสิกํ นวกมฺมํ ทาตุํ, มหลฺลเก วิหาเร ปาสาเท วา กมฺมํ โอโลเกตฺวา ทสทฺวาทส\-วสฺสิกํ นวกมฺมํ ทาตุนฺ”ติ.}

\prose{เตน โข ปน สมเยน ภิกฺขู สพฺเพ วิหาเร นวกมฺมํ เทนฺติ. ภควโต เอตมตฺถํ อาโรเจสุํ. น ภิกฺขเว สพฺเพ วิหาเร นวกมฺมํ ทาตพฺพํ, โย ทเทยฺย, อาปตฺติ ทุกฺกฏสฺสาติ.}

\prose{เตน โข ปน สมเยน ภิกฺขู เอกสฺส เทฺว เทนฺติ. ภควโต เอตมตฺถํ อาโรเจสุํ. น ภิกฺขเว เอกสฺส เทฺว ทาตพฺพา, โย ทเทยฺย, อาปตฺติ ทุกฺกฏสฺสาติ.}

\prose{เตน โข ปน สมเยน ภิกฺขู นวกมฺมํ คเหตฺวา อญฺญํ วาเสนฺติ. ภควโต เอตมตฺถํ อาโรเจสุํ. น ภิกฺขเว นวกมฺมํ คเหตฺวา อญฺโญ วาเสตพฺโพ, โย วาเสยฺย, อาปตฺติ ทุกฺกฏสฺสาติ.}

\prose{เตน โข ปน สมเยน ภิกฺขู นวกมฺมํ คเหตฺวา สํฆิกํ ปฏิพาหนฺติ. ภควโต เอตมตฺถํ อาโรเจสุํ. น ภิกฺขเว นวกมฺมํ คเหตฺวา สํฆิกํ ปฏิพาหิตพฺพํ, โย ปฏิพาเหยฺย, อาปตฺติ ทุกฺกฏสฺส. อนุชานามิ ภิกฺขเว เอกํ วรเสยฺยํ คเหตุนฺติ.}

\prose{เตน โข ปน สมเยน ภิกฺขู นิสฺสีเม ฐิตสฺส นวกมฺมํ เทนฺติ. ภควโต เอตมตฺถํ อาโรเจสุํ. น ภิกฺขเว นิสฺสีเม ฐิตสฺส นวกมฺมํ ทาตพฺพํ, โย ทเทยฺย อาปตฺติ ทุกฺกฏสฺสาติ.}

\prose{\csromanfolio{326}เตน โข ปน สมเยน ภิกฺขู นวกมฺมํ คเหตฺวา สพฺพกาลํ ปฏิพาหนฺติ. ภควโต เอตมตฺถํ อาโรเจสุํ. น ภิกฺขเว นวกมฺมํ คเหตฺวา สพฺพกาลํ ปฏิพาหิตพฺพํ, โย ปฏิพาเหยฺย, อาปตฺติ ทุกฺกฏสฺส. อนุชานามิ ภิกฺขเว วสฺสานํ เตมาสํ ปฏิพาหิตุํ. อุตุกาลํ ปน น ปฏิพาหิตุนฺติ.}

\prose{เตน โข ปน สมเยน ภิกฺขู นวกมฺมํ คเหตฺวา ปกฺกมนฺติปิ วิพฺภมนฺติปิ, กาลํปิ กโรนฺติ, สามเณราปิ ปฏิชานนฺติ, สิกฺขํ ปจฺจกฺขาต\-กาปิ ปฏิชานนฺติ, อนฺติมวตฺถุํ อชฺฌาปนฺนกาปิ ปฏิชานนฺติ, อุมฺมตฺตกาปิ ปฏิชานนฺติ, ขิตฺตจิตฺตาปิ ปฏิชานนฺติ, เวทนาฏฺฏาปิ ปฏิชานนฺติ, อาปตฺติยา อทสฺสเน อุกฺขิตฺตกาปิ ปฏิชานนฺติ, อาปตฺติยา อปฺปฏิกมฺเม อุกฺขิตฺตกาปิ ปฏิชานนฺติ, ปาปิกาย ทิฏฺฐิยา อปฺปฏินิสฺสคฺเค อุกฺขิตฺตกาปิ ปฏิชานนฺติ, ปณฺฑกาปิ ปฏิชานนฺติ, เถยฺยสํวาส\-กาปิ ปฏิชานนฺติ, ติตฺถิยปกฺกนฺตกาปิ ปฏิชานนฺติ, ติรจฺฉาน\-คตาปิ ปฏิชานนฺติ, มาตุฆาตกาปิ ปฏิชานนฺติ, ปิตุฆาตกาปิ ปฏิชานนฺติ, อรหนฺต\-ฆาตกาปิ ปฏิชานนฺติ, ภิกฺขุนิทูสกาปิ ปฏิชานนฺติ, สํฆเภทกาปิ ปฏิชานนฺติ, โลหิตุปฺปาท\-กาปิ ปฏิชานนฺติ, อุภโตพฺยญฺชน\-กาปิ ปฏิชานนฺติ. ภควโต เอตมตฺถํ อาโรเจสุํ.}

\prose{อิธ ปน ภิกฺขเว ภิกฺขุ นวกมฺมํ คเหตฺวา ปกฺกมติ, “มา สํฆสฺส หายี”ติ อญฺญสฺส ทาตพฺพํ.}

\prose{อิธ ปน ภิกฺขเว ภิกฺขุ นวกมฺมํ คเหตฺวา วิพฺภมติ ฯเปฯ กาลํ กโรติ. สามเณโร ปฏิชานาติ. สิกฺขํ ปจฺจกฺขาตโก ปฏิชานาติ. อนฺติมวตฺถุํ อชฺฌาปนฺนโก ปฏิชานาติ. อุมฺมตฺตโก ปฏิชานาติ. ขิตฺตจิตฺโต ปฏิชานาติ. เวทนาฏฺโฏ ปฏิชานาติ. อาปตฺติยา อทสฺสเน อุกฺขิตฺตโก ปฏิชานาติ. อาปตฺติยา อปฺปฏิกมฺเม อุกฺขิตฺตโก ปฏิชานาติ. ปาปิกาย ทิฏฺฐิยา อปฺปฏินิสฺสคฺเค อุกฺขิตฺตโก ปฏิชานาติ. ปณฺฑโก ปฏิชานาติ. เถยฺย\-สํวาสโก ปฏิชานาติ. ติตฺถิย\-ปกฺกนฺตโก ปฏิชานาติ. ติรจฺฉานคโต ปฏิชานาติ. มาตุฆาตโก ปฏิชานาติ. ปิตุฆาตโก ปฏิชานาติ. อรหนฺต\-ฆาตโก ปฏิชานาติ. ภิกฺขุนิ\-ทูสโก ปฏิชานาติ. สํฆเภทโก ปฏิชานาติ. โลหิตุปฺปาทโก ปฏิชานาติ. อุภโต\-พฺยญฺชนโก ปฏิชานาติ, “มา สํฆสฺส หายี”ติ อญฺญสฺส ทาตพฺพํ.}

\prose{\csromanfolio{327}อิธ ปน ภิกฺขเว ภิกฺขุ นวกมฺมํ คเหตฺวา วิปฺปกเต ปกฺกมติ, “มา สํฆสฺส หายี”ติ อญฺญสฺส ทาตพฺพํ.}

\prose{อิธ ปน ภิกฺขเว ภิกฺขุ นวกมฺมํ คเหตฺวา วิปฺปกเต วิพฺภมติ ฯเปฯ. อุภโต\-พฺยญฺชนโก ปฏิชานาติ, “มา สํฆสฺส หายี”ติ อญฺญสฺส ทาตพฺพํ.}

\prose{อิธ ปน ภิกฺขเว ภิกฺขุ นวกมฺมํ คเหตฺวา ปริโยสิเต ปกฺกมติ, ตสฺเสเวตํ.}

\prose{อิธ ปน ภิกฺขเว ภิกฺขุ นวกมฺมํ คเหตฺวา ปริโยสิเต วิพฺภมติ ฯเปฯ กาลํ กโรติ. สามเณโร ปฏิชานาติ. สิกฺขํ ปจฺจกฺขาตโก ปฏิชานาติ. อนฺติมวตฺถุํ อชฺฌาปนฺนโก ปฏิชานาติ. สํโฆ สามี.}

\prose{อิธ ปน ภิกฺขเว ภิกฺขุ นวกมฺมํ คเหตฺวา ปริโยสิเต อุมฺมตฺตโก ปฏิชานาติ. ขิตฺตจิตฺโต ปฏิชานาติ. เวทนาฏฺโฏ ปฏิชานาติ. อาปตฺติยา อทสฺสเน อุกฺขิตฺตโก ปฏิชานาติ. อาปตฺติยา อปฺปฏิกมฺเม อุกฺขิตฺตโก ปฏิชานาติ. ปาปิกาย ทิฏฺฐิยา อปฺปฏินิสฺสคฺเค อุกฺขิตฺตโก ปฏิชานาติ, ตสฺเสเวตํ.}

\prose{อิธ ปน ภิกฺขเว ภิกฺขุ นวกมฺมํ คเหตฺวา ปริโยสิเต ปณฺฑโก ปฏิชานาติ. เถยฺย\-สํวาสโก ปฏิชานาติ. ติตฺถิย\-ปกฺกนฺตโก ปฏิชานาติ. ติรจฺฉานคโต ปฏิชานาติ. มาตุฆาตโก ปฏิชานาติ. ปิตุฆาตโก ปฏิชานาติ. อรหนฺต\-ฆาตโก ปฏิชานาติ. ภิกฺขุนี\-ทูสโก ปฏิชานาติ. สํฆเภทโก ปฏิชานาติ. โลหิตุปฺปาทโก ปฏิชานาติ. อุภโต\-พฺยญฺชนโก ปฏิชานาติ, สํโฆ สามี.}

\csromanmatikaanchor{mtk.153}
\csromantocmarkref{subsection}{อญฺญตฺรปริโภคปฏิกฺเขปาทิ}{mtk.153}
\bookmark[dest={mtk.153},level=2]{อญฺญตฺรปริโภคปฏิกฺเขปาทิ}
\csromanheader{2}{อญฺญตฺร\-ปริโภค\-ปฏิกฺเขปาทิ}
\proseitem{324}{\csromansymbolfootnote{*}{วิ 1. 83 ปิฏฺเฐปิ.}เตน โข ปน สมเยน ภิกฺขู อญฺญตรสฺส อุปาสกสฺส วิหาร\-ปริโภคํ เสนาสนํ อญฺญตฺร ปริภุญฺชนฺติ. อถ โข โส อุปาสโก อุชฺฌายติ ขิยฺยติ วิปาเจติ “กถํ หิ นาม ภทนฺตา อญฺญตฺร ปริโภคํ อญฺญตฺร ปริภุญฺชิสฺสนฺตี”ติ. ภควโต เอตมตฺถํ อาโรเจสุํ. น ภิกฺขเว อญฺญตร ปริโภโค อญฺญตฺร ปริภุญฺชิ\-ตพฺโพ, โย ปริภุญฺเชยฺย, อาปตฺติ ทุกฺกฏสฺสาติ.}

\prose{\csromanfolio{328}\csromansymbolfootnote{*}{วิ 1. 83 ปิฏฺเฐปิ.}เตน โข ปน สมเยน ภิกฺขู อุโปสถคฺคมฺปิ สนฺนิสชฺชมฺปิ หริตุํ กุกฺกุจฺจายนฺตา ฉมาย นิสีทนฺติ, คตฺตานิปิ จีวรานิปิ ปํสุกิตานิ โหนฺติ. ภควโต เอตมตฺถํ อาโรเจสุํ. อนุชานามิ ภิกฺขเว ตาวกาลิกํ หริตุนฺติ.}

\prose{เตน โข ปน สมเยน สํฆสฺส มหาวิหาโร อุนฺทฺริยติ, ภิกฺขู กุกฺกุจฺจายนฺตา เสนาสนํ นาติหรนฺติ\csromansharedfootnote{6bb6c01f87a0db42}{นาภิหรนฺติ (ก)}. ภควโต เอตมตฺถํ อาโรเจสุํ. อนุชานามิ ภิกฺขเว คุตฺตตฺถาย หริตุนฺติ.}

\prose{เตน โข ปน สมเยน สํฆสฺส เสนาสน\-ปริกฺขาริโก มหคฺโฆ กมฺพโล อุปฺปนฺโน โหติ. ภควโต เอตมตฺถํ อาโรเจสุํ. อนุชานามิ ภิกฺขเว ผาติกมฺม\-ตฺถาย ปริวตฺเตตุนฺติ.}

\prose{เตน โข ปน สมเยน สํฆสฺส เสนาสน\-ปริกฺขาริกํ มฆคฺฆํ ทุสฺสํ อุปฺปนฺนํ โหติ. ภควโต เอตมตฺถํ อาโรเจสุํ. อนุชานามิ ภิกฺขเว ผาติกมฺม\-ตฺถาย ปริวตฺเตตุนฺติ.}

\prose{เตน โข ปน สมเยน สํฆสฺส อจฺฉจมฺมํ อุปฺปนฺนํ โหติ. ภควโต เอตมตฺถํ อาโรเจสุํ. อนุชานามิ ภิกฺขเว ปาทปุจฺฉนิํ กาตุนฺติ.}

\prose{เตน โข ปน สมเยน สํฆสฺส จกฺกลิกํ อุปฺปนฺนํ โหติ. ภควโต เอตมตฺถํ อาโรเจสุํ. อนุชานามิ ภิกฺขเว ปาทปุญฺฉนิํ กาตุนฺติ.}

\prose{เตน โข ปน สมเยน สํฆสฺส โจฬกํ อุปฺปนฺนํ โหติ. ภควโต เอตมตฺถํ อาโรเจสุํ. อนุชานามิ ภิกฺขเว ปาทปุญฺฉนิํ กาตุนฺติ.}

\prose{เตน โข ปน สมเยน ภิกฺขู อโธเตหิ ปาเทหิ เสนาสนํ อกฺกมนฺติ, เสนาสนํ ทุสฺสติ. ภควโต เอตมตฺถํ อาโรเจสุํ. น ภิกฺขเว อโธเตหิ ปาเทหิ เสนาสนํ อกฺกมิตพฺพํ, โย อกฺกเมยฺย, อาปตฺติ ทุกฺกฏสฺสาติ.}

\prose{เตน โข ปน สมเยน ภิกฺขู อลฺเลหิ ปาเทหิ เสนาสนํ อกฺกมนฺติ, เสนาสนํ ทุสฺสติ. ภควโต เอตมตฺถํ อาโรเจสุํ. น ภิกฺขเว อลฺเลหิ ปาเทหิ เสนาสนํ อกฺกมิตพฺพํ, โย อกฺกเมยฺย, อาปตฺติ ทุกฺกฏสฺสาติ.}

\prose{\csromanfolio{329}เตน โข ปน สมเยน ภิกฺขู สอุปาหนา เสนาสนํ อกฺกมนฺติ, เสนาสนํ ทุสฺสติ. ภควโต เอตมตฺถํ อาโรเจสุํ. น ภิกฺขเว สอุปาหเนน เสนาสนํ อกฺกมิตพฺพํ, โย อกฺกเมยฺย, อาปตฺติ ทุกฺกฏสฺสาติ.}

\prose{เตน โข ปน สมเยน ภิกฺขู ปริกมฺม\-กตาย ภูมิยา นิฏฺฐุภนฺติ, วณฺโณ ทุสฺสติ. ภควโต เอตมตฺถํ อาโรเจสุํ. น ภิกฺขเว ปริกมฺม\-กตาย ภูมิยา นิฏฺฐุภิตพฺพํ, โย นิฏฺฐุเภยฺย, อาปตฺติ ทุกฺกฏสฺส. อนุชานามิ ภิกฺขเว เขฬมลฺลกนฺติ.}

\prose{เตน โข ปน สมเยน มญฺจปาทาปิ ปีฐปาทาปิ ปริกมฺมกตํ ภูมิํ วิลิขนฺติ. ภควโต เอตมตฺถํ อาโรเจสุํ. อนุชานามิ ภิกฺขเว โจฬเกน ปลิเวเฐตุนฺติ.}

\prose{เตน โข ปน สมเยน ภิกฺขู ปริกมฺมกตํ ภิตฺติํ อปสฺเสนฺติ, วณฺโณ ทุสฺสติ. ภควโต เอตมตฺถํ อาโรเจสุํ. น ภิกฺขเว ปริกมฺมกตา ภิตฺติ อปสฺเสตพฺพา, โย อปสฺเสยฺย, อาปตฺติ ทุกฺกฏสฺส. อนุชานามิ ภิกฺขเว อปสฺเสนผลกนฺติ. อปสฺเสนผลกํ เหฏฺฐโต ภูมิํ วิลิขติ อุปริโต ภิตฺติํ จ. ภควโต เอตมตฺถํ อาโรเจสุํ. อนุชานามิ ภิกฺขเว เหฏฺฐโต จ อุปริโต จ โจฬเกน ปลิเวเฐตุนฺติ.}

\prose{เตน โข ปน สมเยน ภิกฺขู โธตปาทกา นิปชฺชิตุํ กุกฺกุจฺจายนฺติ. ภควโต เอตมตฺถํ อาโรเจสุํ. อนุชานามิ ภิกฺขเว ปจฺจตฺถริตฺวา นิปชฺชิตุนฺติ.}

\csromanmatikaanchor{mtk.154}
\csromantocmarkref{subsection}{สํฆภตฺตาทิอนุชานน}{mtk.154}
\bookmark[dest={mtk.154},level=2]{สํฆภตฺตาทิอนุชานน}
\csromanheader{2}{สํฆภตฺตา\-ทิ\-อนุชานน}
\proseitem{325}{อถ โข ภควา อาฬวิยํ ยถาภิรนฺตํ วิหริตฺวา เยน ราชคหํ เตน จาริกํ ปกฺกามิ. อนุปุพฺเพน จาริกํ จรมาโน เยน ราชคหํ ตทวสริ, ตตฺร สุทํ ภควา ราชคเห วิหรติ เวฬุวเน กลนฺทก\-นิวาเป. เตน โข ปน สมเยน ราชคหํ ทุพฺภิกฺขํ โหติ, มนุสฺสา น สกฺโกนฺติ สํฆภตฺตํ กาตุํ, อิจฺฉนฺติ อุทฺเทสภตฺตํ นิมนฺตนํ สลากภตฺตํ ปกฺขิกํ อุโปสถิกํ ปาฏิปทิกํ กาตุํ. ภควโต เอตมตฺถํ อาโรเจสุํ. อนุชานามิ ภิกฺขเว สํฆภตฺตํ อุทฺเทสภตฺตํ นิมนฺตนํ สลากภตฺตํ ปกฺขิกํ อุโปสถิกํ ปาฏิปทิกนฺติ.}

\csromanmatikaanchor{mtk.155}
\csromantocmarkref{subsection}{ภตฺตุทฺเทสกสมฺมุติ}{mtk.155}
\bookmark[dest={mtk.155},level=2]{ภตฺตุทฺเทสกสมฺมุติ}
\csromanheader{2}{ภตฺตุทฺเทสก\-สมฺมุติ}
\csromanmatikaanchor{mtk.156}
\csromantocmarkref{subsection}{เสนาสนปญฺญาปกาทิสมฺมุติ}{mtk.156}
\bookmark[dest={mtk.156},level=2]{เสนาสนปญฺญาปกาทิสมฺมุติ}
\proseitem{326}{\csromanfolio{330}เตน โข ปน สมเยน ฉพฺพคฺคิยา ภิกฺขู อตฺตโน วรภตฺตานิ คเหตฺวา ลามกานิ ภตฺตานิ ภิกฺขูนํ เทนฺติ. ภควโต เอตมตฺถํ อาโรเจสุํ. อนุชานามิ ภิกฺขเว ปญฺจหงฺเคหิ สมนฺนาคตํ ภิกฺขุํ ภตฺตุทฺเทสกํ สมฺมนฺนิตุํ, โย น ฉนฺทาคติํ คจฺเฉยฺย, น โทสาคติํ คจฺเฉยฺย, น โมหาคติํ คจฺเฉยฺย, น ภยาคติํ คจฺเฉยฺย, อุทฺทิฏฺฐานุทฺทิฏฺฐญฺจ ชาเนยฺย. เอวญฺจ ปน ภิกฺขเว สมฺมนฺนิตพฺโพ, ปฐมํ ภิกฺขุ ยาจิตพฺโพ, ยาจิตฺวา พฺยตฺเตน ภิกฺขุนา ปฏิพเลน สํโฆ ญาเปตพฺโพ\csromandash{}}

\prose{“สุณาตุ เม ภนฺเต สํโฆ, ยทิ สํฆสฺส ปตฺตกลฺลํ, สํโฆ อิตฺถนฺนามํ ภิกฺขุํ ภตฺตุทฺเทสกํ สมฺมนฺเนยฺย, เอสา ญตฺติ.}

\prose{สุณาตุ เม ภนฺเต สํโฆ, สํโฆ อิตฺถนฺนามํ ภิกฺขุํ ภตฺตุทฺเทสกํ สมฺมนฺนติ, ยสฺสายสฺมโต ขมติ อิตฺถนฺนามสฺส ภิกฺขุโน ภตฺตุ\-ทฺเทสกสฺส สมฺมุติ, โส ตุณฺหสฺส, ยสฺส นกฺขมติ, โส ภาเสยฺย.}

\prose{สมฺมโต สํเฆน อิตฺถนฺนาโม ภิกฺขุ ภตฺตุทฺเทสโก, ขมติ สํฆสฺส, ตสฺมา ตุณฺหี, เอวเมตํ ธารยามี”ติ.}

\prose{อถ โข ภตฺตุ\-ทฺเทสกานํ ภิกฺขูนํ เอตทโหสิ “กถํ นุ โข ภตฺตํ อุทฺทิสิตพฺพนฺ”ติ. ภควโต เอตมตฺถํ อาโรเจสุํ. อนุชานามิ ภิกฺขเว สลากาย วา ปฏฺฏิกาย วา\csromansharedfootnote{323c3d4441e6e4ce}{ปฏิกาย วา (สฺยา)} อุปนิพนฺธิตฺวา โอปุญฺชิตฺวา ภตฺตํ อุทฺทิสิตุนฺติ.}

\csromanheader{2}{\textbf{เสนาสน ปญฺญาปกาทิสมฺมุติ}}
\proseitem{327}{เตน โข ปน สมเยน สํฆสฺส เสนาสน\-ปญฺญาปโก น โหติ ฯเปฯ ภณฺฑาคาริโก น โหติ ฯเปฯ จีวร\-ปฺปฏิคฺคาหโก น โหติ ฯเปฯ จีวรภาชโก น โหติ ฯเปฯ ยาคุภาชโก น โหติ ฯเปฯ ผลภาชโก น โหติ ฯเปฯ ขชฺชกภาชโก น โหติ, ขชฺชกํ อภาชิยมานํ นสฺสติ. ภควโต เอตมตฺถํ อาโรเจสุํ. อนุชานามิ ภิกฺขเว ปญฺจหงฺเคหิ สมนฺนาคตํ ภิกฺขุํ ขชฺชกภาชกํ สมฺมนฺนิตุํ, โย น ฉนฺทาคติํ คจฺเฉยฺย, น โทสาคติํ คจฺเฉยฺย, น โมหาคติํ คจฺเฉยฺย, น ภยาคติํ \csromanfolio{331}คจฺเฉยฺย, ภาชิตาภาชิตญฺจ ชาเนยฺย. เอวญฺจ ปน ภิกฺขเว สมฺมนฺนิตพฺโพ, ปฐมํ ภิกฺขุ ยาจิตพฺโพ, ยาจิตฺวา พฺยตฺเตน ภิกฺขุนา ปฏิพเลน สํโฆ ญาเปตพฺโพ\csromandash{}}

\prose{“สุณาตุ เม ภนฺเต สํโฆ, ยทิ สํฆสฺส ปตฺตกลฺลํ, สํโฆ อิตฺถนฺนามํ ภิกฺขุํ ขชฺชกภาชกํ สมฺมนฺเนยฺย, เอสา ญตฺติ.}

\prose{สุณาตุ เม ภนฺเต สํโฆ, สํโฆ อิตฺถนฺนามํ ภิกฺกุํ ขชฺชกภาชกํ สมฺมนฺนติ, ยสฺสายสฺมโต ขมติ อิตฺถนฺนามสฺส ภิกฺขุโน ขชฺชกภาชกสฺส สมฺมุติ, โส ตุณฺหสฺส, ยสฺส นกฺขมติ, โส ภาเสยฺย.}

\prose{สมฺมโต สํเฆน อิตฺถนฺนาโม ภิกฺขุ ขชฺชกภาชโก, ขมติ สํฆสฺส, ตสฺมา ตุณฺหี, เอวเมตํ ธารยามี”ติ.}

\csromanmatikaanchor{mtk.157}
\csromantocmarkref{subsection}{อปฺปมตฺตกวิสฺสชฺชกสมฺมุติ}{mtk.157}
\bookmark[dest={mtk.157},level=2]{อปฺปมตฺตกวิสฺสชฺชกสมฺมุติ}
\csromanheader{2}{อปฺปมตฺตกวิสฺสชฺชก\-สมฺมุติ}
\proseitem{328}{เตน โข ปน สมเยน สํฆสฺส ภณฺฑาคาเร อปฺปมตฺตโก ปริกฺขาโร อุปฺปนฺโน\csromansharedfootnote{bdc4b54f6c286e3b}{อุสฺสนฺโน (สฺยา)} โหติ. ภควโต เอตมตฺถํ อาโรเจสุํ. อนุชานามิ ภิกฺขเว ปญฺจหงฺเคหิ สมนฺนาคตํ ภิกฺขุํ อปฺปมตฺตก\-วิสฺสชฺชกํ สมฺมนฺนิตุํ, โย น ฉนฺทาคติํ คจฺเฉยฺย, น โทสาคติํ คจฺเฉยฺย, น โมหาคติํ คจฺเฉยฺย, น ภยาคติํ คจฺเฉยฺย, วิสฺสชฺชิตาวิสฺสชฺชิตญฺจ ชาเนยฺย. เอวญฺจ ปน ภิกฺขเว สมฺมนฺนิตพฺโพ, ปฐมํ ภิกฺขุ ยาจิตพฺโพ, ยาจิตฺวา พฺยตฺเตน ภิกฺขุนา ปฏิพเลน สํโฆ ญาเปตพฺโพ\csromandash{}}

\prose{“สุณาตุ เม ภนฺเต สํโฆ, ยทิ สํฆสฺส ปตฺตกลฺลํ, สํโฆ อิตฺถนฺนามํ ภิกฺขุํ อปฺปมตฺตก\-วิสฺสชฺชกํ สมฺมนฺเนยฺย, เอสา ญตฺติ.}

\prose{สุณาตุ เม ภนฺเต สํโฆ, สํโฆ อิตฺถนฺนามํ ภิกฺขุํ อปฺปมตฺตก\-วิสฺสชฺชกํ สมฺมนฺนติ, ยสฺสายสฺมโต ขมติ อิตฺถนฺนามสฺส ภิกฺขุโน อปฺปมตฺตก\-วิสฺสชฺชกสฺส สมฺมุติ, โส ตุณฺหสฺส, ยสฺส นกฺขมติ, โส ภาเสยฺย.}

\prose{สมฺมโต สํเฆน อิตฺถนฺนาโม ภิกฺขุ อปฺปมตฺตก\-วิสฺสชฺชโก, ขมติ สํฆสฺส, ตสฺมา ตุณฺหี, เอวเมตํ ธารยามี”ติ.}

\prose{\csromanfolio{332}เตน อปฺปมตฺตก\-วิสฺสชฺชเกน ภิกฺขุนา เอกา\csromansharedfootnote{98af2c26c70ae412}{เอเตกา (สี)} สูจิ ทาตพฺพา, สตฺถกํ ทาตพฺพํ, อุปาหนา ทาตพฺพา, กายพนฺธนํ ทาตพฺพํ, อํสพทฺธโก ทาตพฺโพ, ปริสฺสาวนํ ทาตพฺพํ, ธมฺมกรโณ ทาตพฺโพ, กุสิ ทาตพฺพา, อฑฺฒกุสิ ทาตพฺพา, มณฺฑลํ ทาตพฺพํ, อฑฺฒมณฺฑลํ ทาตพฺพํ, อนุวาโต ทาตพฺโพ, ปริภณฺฑํ ทาตพฺพํ. สเจ โหติ สํฆสฺส สปฺปิ วา เตลํ วา มธุ วา ผาณิตํ วา, สกิํ ปฏิสายิตุํ ทาตพฺพํ. สเจ ปุนปิ อตฺโถ โหติ, ปุนปิ ทาตพฺพํ.}

\csromanmatikaanchor{mtk.158}
\csromantocmarkref{subsection}{สาฏิยคฺคาหาปกาทิสมฺมุติ}{mtk.158}
\bookmark[dest={mtk.158},level=2]{สาฏิยคฺคาหาปกาทิสมฺมุติ}
\csromanheader{2}{สาฏิยคฺคาหาปกา\-ทิ\-สมฺมุติ}
\proseitem{329}{เตน โข ปน สมเยน สํฆสฺส สาฏิย\-คฺคาหาปโก น โหติ ฯเปฯ ปตฺตคฺคาหาปโก น โหติ ฯเปฯ อารามิกเปสโก น โหติ ฯเปฯ สามเณรเปสโก น โหติ, สามเณรา อเปสิยมานา กมฺมํ น กโรนฺติ. ภควโต เอตมตฺถํ อาโรเจสุํ. อนุชานามิ ภิกฺขเว ปญฺจหงฺเคหิ สมนฺนาคตํ ภิกฺขุํ สามเณร\-เปสกํ สมฺมนฺนิตุํ, โย น ฉนฺทาคติํ คจฺเฉยฺย, น โทสาคติํ คจฺเฉยฺย, น โมหาคติํ คจฺเฉยฺย, น ภยาคติํ คจฺเฉยฺย, เปสิตาเปสิตญฺจ ชาเนยฺย. เอวญฺจ ปน ภิกฺขเว สมฺมนฺนิตพฺโพ, ปฐมํ ภิกฺขุ ยาจิตพฺโพ, ยาจิตฺวา พฺยตฺเตน ภิกฺขุนา ปฏิพเลน สํโฆ ญาเปตพฺโพ\csromandash{}}

\prose{“สุณาตุ เม ภนฺเต สํโฆ, ยทิ สํฆสฺส ปตฺตกลฺลํ, สํโฆ อิตฺถนฺนามํ ภิกฺขุํ สามเณร\-เปสกํ สมฺมนฺเนยฺย, เอสา ญตฺติ.}

\prose{สุณาตุ เม ภนฺเต สํโฆ, สํโฆ อิตฺถนฺนามํ ภิกฺขุํ สามเณร\-เปสกํ สมฺมนฺนติ, ยสฺสายสฺมโต ขมติ อิตฺถนฺนามสฺส ภิกฺขุโน สามเณร\-เปสกสฺส สมฺมุติ, โส ตุณฺหสฺส, ยสฺส นกฺขมติ, โส ภาเสยฺย.}

\prosewithcloser{สมฺมโต สํเฆน อิตฺถนฺนาโม ภิกฺขุ สามเณรเปสโก, ขมติ สํฆสฺส, ตสฺมา ตุณฺหี, เอวเมตํ ธารยามี”ติ.}{ตติย\-ภาณวาโร นิฏฺฐิโต.}
\csromancenter{เสนาสน\-กฺขนฺธโก ฉฏฺโฐ.}
\csromancenter{\textbf{ตสฺสุทฺทานํ}}
\csromanmatikaanchor{mtk.159}
\csromantocmarkref{section}{อุทฺทานคาถา}{mtk.159}
\bookmark[dest={mtk.159},level=1]{อุทฺทานคาถา}
\setlength{\csromangathapagewidth}{0pt}
\setlength{\csromangathawakwidth}{0pt}
\csromangathameasuregroup{วิหารํ พุทฺธเสฏฺเฐน, \\ ตหํ ตหํ นิกฺขมนฺติ, \\ เสฏฺฐี คหปติ ทิสฺวา, \\ การาเปยฺยํ วเสยฺยาถ, \\ วิหารํ อฑฺฒโยคญฺจ, \\ ปญฺจเลณํ อนุญฺญาสิ, \\ ชโน วิหารํ กาเรติ, \\ กวาฏํ ปิฏฺฐสํฆาฏํ, \\ อาวิญฺฉนจฺฉิทฺทํ รชฺชุํ, \\ สูจิฆฏิตาฬจฺฉิทฺทํ, \\ ยนฺตกํ สูจิกญฺเจว, \\ เวทิชาลสลากญฺจ, \\ มิฑฺฒิ พิทลมญฺจญฺจ, \\ พุนฺทิกุฬิรปาทญฺจ, \\ สตฺตงฺโค จ ภทฺทปีฐํ, \\ อามลาผลกา โกจฺฉา, \\ อุจฺจาหิปฏิปาทกา, \\ สุตฺตํ อฏฺฐปทํ โจฬํ, \\ คิรคฺโค ภิสิโย จาปิ, \\ โอนทฺธํ เหฏฺฐา ปตติ, \\ ภตฺติญฺจ หตฺถภตฺติญฺจ, \\ ติตฺถิยา วิหาเร จาปิ, \\ อิกฺกาสํ ปาณิกํ กุณฺฑํ, \\ อุสฺสนฺเน ปจฺจุทฺธริตุํ, \\ อิกฺกาสํ ปฏิภานญฺจ, \\ ปริปตนฺติ อาฬกา, \\ ขุทฺทเก กุฏฺฏปาโท จ, \\ จีวรวํสํ รชฺชุญฺจ, \\ อาลมฺพนํ ติณจุณฺณํ, \\ อชฺโฌกาเส โอตปฺปติ, \\ วิหาโร โกฏฺฐโก เจว, \\ อาราเม จ ปุน โกฏฺเฐ, \\ สุธํ อนาถปิณฺฑิ จ, \\ ทิฏฺฐธมฺโม นิมนฺเตสิ, \\ อาณาเปสนฺตรามคฺเค, \\ เวสาลิยํ นวกมฺมํ, \\ โก อรหติ ภตฺตคฺเค, \\ ปริคฺคหิตนฺตรฆรา, \\ ปติฏฺฐาเปสิ อารามํ, \\ คิลานา วรเสยฺยา จ, \\ เกน นุ โข กถํ นุ โข, \\ ปริเวณํ อนุภาคญฺจ, \\ นิสฺสีมํ สพฺพกาลญฺจ, \\ อุปนนฺโท จ วณฺเณสิ, \\ สมานาสนิกา ภินฺทิํสุ, \\ อสมานาสนิกา ทีฆํ, \\ อยฺยิกา จ อวิทูเร, \\ อาฬวี ปิณฺฑกกุฏฺเฏหิ, \\ อาโลกเสตกาฬญฺจ, \\ ภณฺฑิขณฺฑปริภณฺฑํ, \\ โอสิเต อกตํ วิปฺปํ, \\ อฑฺฒโยเค จ สตฺตฏฺฐ, \\ สพฺพํ วิหารํ เอกสฺส, \\ นิสฺสีมํ สพฺพกาลญฺจ, \\ กาลญฺจ สามเณรญฺจ, \\ อุมฺมตฺตขิตฺตจิตฺตา จ, \\ อปฺปฏิกมฺมทิฏฺฐิยา, \\ ติรจฺฉานมาตุปิตุ, \\ เภทกา โลหิตุปฺปาทา, \\ มา สํฆสฺส ปริหายิ, \\ วิปฺปกเต จ อญฺญสฺส, \\ วิพฺภมติ กาลงฺกโต, \\ ปจฺจกฺขาโต จ สิกฺขาย, \\ สํโฆว สามิโก โหติ, \\ อทสฺสนาปฺปฏิกมฺเม, \\ ปณฺฑโก เถยฺยติตฺถี จ, \\ ฆาตโก ทูสโก จาปิ, \\ ปฏิชานาติ ยทิ โส, \\ หรนฺตญฺญตฺร กุกฺกุจฺจํ, \\ ทุสฺสํ จ จมฺมจกฺกลี, \\ อลฺลา อุปาหนานิฏฺฐุ, \\ อปสฺเสนํ ลิขเตว, \\ ราชคเห น สกฺโกนฺติ, \\ กถํ นุ โข ปญฺญาปกํ, \\ ปฏิคฺคาหภาชโก จาปิ, \\ ขชฺชกภาชโก เจว, \\ สาฏิยคฺคาหาปโก เจว, \\ อารามิกสามเณร, \\ สพฺพาภิภู โลกวิทู,}{\csromangathabat{วิหารํ พุทฺธเสฏฺเฐน,}{อปญฺญตฺตํ ตทา อหุ.} \\ \csromangathabat{ตหํ ตหํ นิกฺขมนฺติ,}{วาสา เต ชินสาวกา.} \\ \csromangathabat{เสฏฺฐี คหปติ ทิสฺวา,}{ภิกฺขูนํ อิทมพฺรวิ.} \\ \csromangathabat{การาเปยฺยํ วเสยฺยาถ,}{ปฏิปุจฺฉิํสุ นายกํ.} \\ \csromangathabat{วิหารํ อฑฺฒโยคญฺจ,}{ปาสาทํ หมฺมิยํ คุหํ.} \\ \csromangathabat{ปญฺจเลณํ อนุญฺญาสิ,}{วิหาเร เสฏฺฐิ การยิ.} \\ \csromangathabat{ชโน วิหารํ กาเรติ,}{อกวาฏํ อสํวุตํ.} \\ \csromangathabat{กวาฏํ ปิฏฺฐสํฆาฏํ,}{อุทุกฺขลญฺจ อุตฺตริ.} \\ \csromangathabat{อาวิญฺฉนจฺฉิทฺทํ รชฺชุํ,}{วฏฺฏิญฺจ กปิสีสกํ.} \\ \csromangathabat{สูจิฆฏิตาฬจฺฉิทฺทํ,}{โลหกฏฺฐวิสาณกํ.} \\ \csromangathabat{ยนฺตกํ สูจิกญฺเจว,}{ฉทนํ อุลฺลิตฺตาวลิตฺตํ.} \\ \csromangathabat{เวทิชาลสลากญฺจ,}{จกฺกลิ สนฺถเรน จ.} \\ \csromangathabat{มิฑฺฒิ พิทลมญฺจญฺจ,}{โสสานิกมสารโก.} \\ \csromangathabat{พุนฺทิกุฬิรปาทญฺจ,}{อาหจฺจาสนฺทิ อุจฺจเก.} \\ \csromangathabat{สตฺตงฺโค จ ภทฺทปีฐํ,}{ปีฐเกฬกปาทกํ.} \\ \csromangathabat{อามลาผลกา โกจฺฉา,}{ปลาลปีฐเมว จ.} \\ \csromangathabat{อุจฺจาหิปฏิปาทกา,}{อฏฺฐงฺคุลิ จ ปาทกา.} \\ \csromangathabat{สุตฺตํ อฏฺฐปทํ โจฬํ,}{ตูลิกํ อฑฺฒกายิกํ.} \\ \csromangathabat{คิรคฺโค ภิสิโย จาปิ,}{ทุสฺสํ เสนาสนํปิ จ.} \\ \csromangathabat{โอนทฺธํ เหฏฺฐา ปตติ,}{อุปฺปาเฏตฺวา หรนฺติ จ.} \\ \csromangathabat{ภตฺติญฺจ หตฺถภตฺติญฺจ,}{อนุญฺญาสิ ตถาคโต.} \\ \csromangathabat{ติตฺถิยา วิหาเร จาปิ,}{ถุสํ สณฺหญฺจ มตฺติกา.} \\ \csromangathabat{อิกฺกาสํ ปาณิกํ กุณฺฑํ,}{สาสปํ สิตฺถเตลกํ.} \\ \csromangathabat{อุสฺสนฺเน ปจฺจุทฺธริตุํ,}{ปรุสํ คณฺฑุมตฺติกํ.} \\ \csromangathabat{อิกฺกาสํ ปฏิภานญฺจ,}{นีจา จโย จ อารุหํ.} \\ \csromangathabat{ปริปตนฺติ อาฬกา,}{อฑฺฒกุฏฺฏํ ตโย ปุน.} \\ \csromangathabat{ขุทฺทเก กุฏฺฏปาโท จ,}{โอวสฺสติ สรํ ขิลํ.} \\ \csromangathabat{จีวรวํสํ รชฺชุญฺจ,}{อาฬินฺทํ กิฏิเกน จ.} \\ \csromangathabat{อาลมฺพนํ ติณจุณฺณํ,}{เหฏฺฐามคฺเค นยํ กเร.} \\ \csromangathabat{อชฺโฌกาเส โอตปฺปติ,}{สาลํ เหฏฺฐา จ ภาชนํ.} \\ \csromangathabat{วิหาโร โกฏฺฐโก เจว,}{ปริเวณคฺคิสาลกํ.} \\ \csromangathabat{อาราเม จ ปุน โกฏฺเฐ,}{เหฏฺฐญฺเญว นยํ กเร.} \\ \csromangathabat{สุธํ อนาถปิณฺฑิ จ,}{สทฺโธ สีตวนํ อคา.} \\ \csromangathabat{ทิฏฺฐธมฺโม นิมนฺเตสิ,}{สห สํเฆน นายกํ.} \\ \csromangathabat{อาณาเปสนฺตรามคฺเค,}{อารามํ การยี คโณ.} \\ \csromangathabat{เวสาลิยํ นวกมฺมํ,}{ปุรโต จ ปริคฺคหิ.} \\ \csromangathabat{โก อรหติ ภตฺตคฺเค,}{ติตฺติรญฺจ อวนฺทิยา.} \\ \csromangathabat{ปริคฺคหิตนฺตรฆรา,}{ตูโล สาวตฺถิ โอสริ.} \\ \csromangathabat{ปติฏฺฐาเปสิ อารามํ,}{ภตฺตคฺเค จ โกลาหลํ.} \\ \csromangathabat{คิลานา วรเสยฺยา จ,}{เลสา สตฺตรสา ตหิํ.} \\ \csromangathabat{เกน นุ โข กถํ นุ โข,}{วิหารคฺเคน ภาชยิ.} \\ \csromangathabat{ปริเวณํ อนุภาคญฺจ,}{อกามา ภาคํ โน ทเท.} \\ \csromangathabat{นิสฺสีมํ สพฺพกาลญฺจ,}{คาหา เสนาสเน ตโย.} \\ \csromangathabat{อุปนนฺโท จ วณฺเณสิ,}{ฐิตกา สมกาสนา.} \\ \csromangathabat{สมานาสนิกา ภินฺทิํสุ,}{ติวคฺคา จ ทุวคฺคิกํ.} \\ \csromangathabat{อสมานาสนิกา ทีฆํ,}{สาฬินฺทํ ปริภุญฺชิตุํ.} \\ \csromangathabat{อยฺยิกา จ อวิทูเร,}{ภาชิตญฺจ กีฏาคิเร.} \\ \csromangathabat{อาฬวี ปิณฺฑกกุฏฺเฏหิ,}{ทฺวารอคฺคฬวฏฺฏิกา.} \\ \csromangathabat{อาโลกเสตกาฬญฺจ,}{เครุฉาทนพนฺธนา.} \\ \csromangathabat{ภณฺฑิขณฺฑปริภณฺฑํ,}{วีส ติํสา จ กาลิกา.} \\ \csromangathabat{โอสิเต อกตํ วิปฺปํ,}{ขุทฺเท ฉปฺปญฺจวสฺสิกํ.} \\ \csromangathabat{อฑฺฒโยเค จ สตฺตฏฺฐ,}{มหลฺเล ทส ทฺวาทส.} \\ \csromangathabat{สพฺพํ วิหารํ เอกสฺส,}{อญฺญํ วาเสนฺติ สํฆิกํ.} \\ \csromangathabat{นิสฺสีมํ สพฺพกาลญฺจ,}{ปกฺกมิ วิพฺภมนฺติ จ.} \\ \csromangathabat{กาลญฺจ สามเณรญฺจ,}{สิกฺขาปจฺจกฺขอนฺติมํ.} \\ \csromangathabat{อุมฺมตฺตขิตฺตจิตฺตา จ,}{เวทนาปตฺติทสฺสนา.} \\ \csromangathabat{อปฺปฏิกมฺมทิฏฺฐิยา,}{ปณฺฑกา เถยฺยติตฺถิยา.} \\ \csromangathabat{ติรจฺฉานมาตุปิตุ,}{อรหนฺตา จ ทูสกา.} \\ \csromangathabat{เภทกา โลหิตุปฺปาทา,}{อุภโต จาปิ พฺยญฺชนกา.} \\ \csromangathabat{มา สํฆสฺส ปริหายิ,}{กมฺมํ อญฺญสฺส ทาตเว.} \\ \csromangathabat{วิปฺปกเต จ อญฺญสฺส,}{กเต ตสฺเสว ปกฺกเม.} \\ \csromangathabat{วิพฺภมติ กาลงฺกโต,}{สามเณโร จ ชายติ.} \\ \csromangathabat{ปจฺจกฺขาโต จ สิกฺขาย,}{อนฺติมชฺฌาปนฺนโก ยทิ.} \\ \csromangathabat{สํโฆว สามิโก โหติ,}{อุมฺมตฺตขิตฺตเวทนา.} \\ \csromangathabat{อทสฺสนาปฺปฏิกมฺเม,}{ทิฏฺฐิ ตสฺเสว โหติ ตํ.} \\ \csromangathabat{ปณฺฑโก เถยฺยติตฺถี จ,}{ติรจฺฉานมาตุเปตฺติกํ.} \\ \csromangathabat{ฆาตโก ทูสโก จาปิ,}{เภทโลหิตพฺยญฺชนา.} \\ \csromangathabat{ปฏิชานาติ ยทิ โส,}{สํโฆว โหติ สามิโก.} \\ \csromangathabat{หรนฺตญฺญตฺร กุกฺกุจฺจํ,}{อุนฺทฺริยติ จ กมฺพลํ.} \\ \csromangathabat{ทุสฺสํ จ จมฺมจกฺกลี,}{โจฬกํ อกฺกมนฺติ จ.} \\ \csromangathabat{อลฺลา อุปาหนานิฏฺฐุ,}{ลิขนฺติ อปสฺเสนฺติ จ.} \\ \csromangathabat{อปสฺเสนํ ลิขเตว,}{โธตปจฺจตฺถเรน จ.} \\ \csromangathabat{ราชคเห น สกฺโกนฺติ,}{ลามกํ ภตฺตุทฺเทสกํ.} \\ \csromangathabat{กถํ นุ โข ปญฺญาปกํ,}{ภณฺฑาคาริกสมฺมุติ.} \\ \csromangathabat{ปฏิคฺคาหภาชโก จาปิ,}{ยาคุ จ ผลภาชโก.} \\ \csromangathabat{ขชฺชกภาชโก เจว,}{อปฺปมตฺตกวิสฺสชฺเช.} \\ \csromangathabat{สาฏิยคฺคาหาปโก เจว,}{ตเถว ปตฺตคฺคาหโก.} \\ \csromangathabat{อารามิกสามเณร,}{เปสกสฺส จ สมฺมุติ.} \\ \csromangathabat{สพฺพาภิภู โลกวิทู,}{หิตจิตฺโต วินายโก.}}
\csromangathasetleft{วิหารํ พุทฺธเสฏฺเฐน, \\ ตหํ ตหํ นิกฺขมนฺติ, \\ เสฏฺฐี คหปติ ทิสฺวา, \\ การาเปยฺยํ วเสยฺยาถ, \\ วิหารํ อฑฺฒโยคญฺจ, \\ ปญฺจเลณํ อนุญฺญาสิ, \\ ชโน วิหารํ กาเรติ, \\ กวาฏํ ปิฏฺฐสํฆาฏํ, \\ อาวิญฺฉนจฺฉิทฺทํ รชฺชุํ, \\ สูจิฆฏิตาฬจฺฉิทฺทํ, \\ ยนฺตกํ สูจิกญฺเจว, \\ เวทิชาลสลากญฺจ, \\ มิฑฺฒิ พิทลมญฺจญฺจ, \\ พุนฺทิกุฬิรปาทญฺจ, \\ สตฺตงฺโค จ ภทฺทปีฐํ, \\ อามลาผลกา โกจฺฉา, \\ อุจฺจาหิปฏิปาทกา, \\ สุตฺตํ อฏฺฐปทํ โจฬํ, \\ คิรคฺโค ภิสิโย จาปิ, \\ โอนทฺธํ เหฏฺฐา ปตติ, \\ ภตฺติญฺจ หตฺถภตฺติญฺจ, \\ ติตฺถิยา วิหาเร จาปิ, \\ อิกฺกาสํ ปาณิกํ กุณฺฑํ, \\ อุสฺสนฺเน ปจฺจุทฺธริตุํ, \\ อิกฺกาสํ ปฏิภานญฺจ, \\ ปริปตนฺติ อาฬกา, \\ ขุทฺทเก กุฏฺฏปาโท จ, \\ จีวรวํสํ รชฺชุญฺจ, \\ อาลมฺพนํ ติณจุณฺณํ, \\ อชฺโฌกาเส โอตปฺปติ, \\ วิหาโร โกฏฺฐโก เจว, \\ อาราเม จ ปุน โกฏฺเฐ, \\ สุธํ อนาถปิณฺฑิ จ, \\ ทิฏฺฐธมฺโม นิมนฺเตสิ, \\ อาณาเปสนฺตรามคฺเค, \\ เวสาลิยํ นวกมฺมํ, \\ โก อรหติ ภตฺตคฺเค, \\ ปริคฺคหิตนฺตรฆรา, \\ ปติฏฺฐาเปสิ อารามํ, \\ คิลานา วรเสยฺยา จ, \\ เกน นุ โข กถํ นุ โข, \\ ปริเวณํ อนุภาคญฺจ, \\ นิสฺสีมํ สพฺพกาลญฺจ, \\ อุปนนฺโท จ วณฺเณสิ, \\ สมานาสนิกา ภินฺทิํสุ, \\ อสมานาสนิกา ทีฆํ, \\ อยฺยิกา จ อวิทูเร, \\ อาฬวี ปิณฺฑกกุฏฺเฏหิ, \\ อาโลกเสตกาฬญฺจ, \\ ภณฺฑิขณฺฑปริภณฺฑํ, \\ โอสิเต อกตํ วิปฺปํ, \\ อฑฺฒโยเค จ สตฺตฏฺฐ, \\ สพฺพํ วิหารํ เอกสฺส, \\ นิสฺสีมํ สพฺพกาลญฺจ, \\ กาลญฺจ สามเณรญฺจ, \\ อุมฺมตฺตขิตฺตจิตฺตา จ, \\ อปฺปฏิกมฺมทิฏฺฐิยา, \\ ติรจฺฉานมาตุปิตุ, \\ เภทกา โลหิตุปฺปาทา, \\ มา สํฆสฺส ปริหายิ, \\ วิปฺปกเต จ อญฺญสฺส, \\ วิพฺภมติ กาลงฺกโต, \\ ปจฺจกฺขาโต จ สิกฺขาย, \\ สํโฆว สามิโก โหติ, \\ อทสฺสนาปฺปฏิกมฺเม, \\ ปณฺฑโก เถยฺยติตฺถี จ, \\ ฆาตโก ทูสโก จาปิ, \\ ปฏิชานาติ ยทิ โส, \\ หรนฺตญฺญตฺร กุกฺกุจฺจํ, \\ ทุสฺสํ จ จมฺมจกฺกลี, \\ อลฺลา อุปาหนานิฏฺฐุ, \\ อปสฺเสนํ ลิขเตว, \\ ราชคเห น สกฺโกนฺติ, \\ กถํ นุ โข ปญฺญาปกํ, \\ ปฏิคฺคาหภาชโก จาปิ, \\ ขชฺชกภาชโก เจว, \\ สาฏิยคฺคาหาปโก เจว, \\ อารามิกสามเณร, \\ สพฺพาภิภู โลกวิทู,}
\csromangathagroupwithclosermajor{\csromangathastanza{\csromanfolio{333}\csromangathabat{วิหารํ พุทฺธเสฏฺเฐน,}{อปญฺญตฺตํ ตทา อหุ.} \\ \csromangathabat{ตหํ ตหํ นิกฺขมนฺติ,}{วาสา เต ชินสาวกา.}}\csromangathastanza{\csromangathabat{เสฏฺฐี คหปติ ทิสฺวา,}{ภิกฺขูนํ อิทมพฺรวิ.} \\ \csromangathabat{การาเปยฺยํ วเสยฺยาถ,}{ปฏิปุจฺฉิํสุ นายกํ.}}\csromangathastanza{\csromangathabat{วิหารํ อฑฺฒโยคญฺจ,}{ปาสาทํ หมฺมิยํ คุหํ.} \\ \csromangathabat{ปญฺจเลณํ อนุญฺญาสิ,}{วิหาเร เสฏฺฐิ การยิ.}}\csromangathastanza{\csromangathabat{ชโน วิหารํ กาเรติ,}{อกวาฏํ อสํวุตํ.} \\ \csromangathabat{กวาฏํ ปิฏฺฐสํฆาฏํ,}{อุทุกฺขลญฺจ อุตฺตริ.}}\csromangathastanza{\csromangathabat{อาวิญฺฉน\-จฺฉิทฺทํ รชฺชุํ,}{วฏฺฏิญฺจ กปิสีสกํ.} \\ \csromangathabat{สูจิฆฏิตาฬจฺฉิทฺทํ,}{โลหกฏฺฐวิสาณกํ.}}\csromangathastanza{\csromangathabat{ยนฺตกํ สูจิกญฺเจว,}{ฉทนํ อุลฺลิตฺตาวลิตฺตํ.} \\ \csromangathabat{เวทิชาลสลากญฺจ,}{จกฺกลิ สนฺถเรน จ.}}\csromangathastanza{\csromangathabat{มิฑฺฒิ พิทลมญฺจญฺจ,}{โสสานิกมสารโก.} \\ \csromangathabat{พุนฺทิกุฬิรปาทญฺจ,}{อาหจฺจาสนฺทิ อุจฺจเก.}}\csromangathastanza{\csromangathabat{สตฺตงฺโค จ ภทฺทปีฐํ,}{ปีฐเกฬกปาทกํ.} \\ \csromangathabat{อามลาผลกา โกจฺฉา,}{ปลาลปีฐเมว จ.}}\csromangathastanza{\csromangathabat{อุจฺจาหิปฏิปาทกา,}{อฏฺฐงฺคุลิ จ ปาทกา.} \\ \csromangathabat{สุตฺตํ อฏฺฐปทํ โจฬํ,}{ตูลิกํ อฑฺฒกายิกํ.}}\csromangathastanza{\csromangathabat{คิรคฺโค ภิสิโย จาปิ,}{ทุสฺสํ เสนาสนํปิ จ.} \\ \csromangathabat{โอนทฺธํ เหฏฺฐา ปตติ,}{อุปฺปาเฏตฺวา หรนฺติ จ.}}\csromangathastanza{\csromangathabat{ภตฺติญฺจ หตฺถภตฺติญฺจ,}{อนุญฺญาสิ ตถาคโต.} \\ \csromangathabat{ติตฺถิยา วิหาเร จาปิ,}{ถุสํ สณฺหญฺจ มตฺติกา.}}\csromangathastanza{\csromangathabat{อิกฺกาสํ ปาณิกํ กุณฺฑํ,}{สาสปํ สิตฺถเตลกํ.} \\ \csromangathabat{อุสฺสนฺเน ปจฺจุทฺธริตุํ,}{ปรุสํ คณฺฑุมตฺติกํ.}}\csromangathastanza{\csromangathabat{อิกฺกาสํ ปฏิภานญฺจ,}{นีจา จโย จ อารุหํ.} \\ \csromangathabat{ปริปตนฺติ อาฬกา,}{อฑฺฒกุฏฺฏํ ตโย ปุน.}}\csromangathastanza{\csromanfolio{334}\csromangathabat{ขุทฺทเก กุฏฺฏปาโท จ,}{โอวสฺสติ สรํ ขิลํ.} \\ \csromangathabat{จีวรวํสํ รชฺชุญฺจ,}{อาฬินฺทํ กิฏิเกน จ.}}\csromangathastanza{\csromangathabat{อาลมฺพนํ ติณจุณฺณํ,}{เหฏฺฐามคฺเค นยํ กเร.} \\ \csromangathabat{อชฺโฌกาเส โอตปฺปติ,}{สาลํ เหฏฺฐา จ ภาชนํ.}}\csromangathastanza{\csromangathabat{วิหาโร โกฏฺฐโก เจว,}{ปริเวณคฺคิสาลกํ.} \\ \csromangathabat{อาราเม จ ปุน โกฏฺเฐ,}{เหฏฺฐญฺเญว นยํ กเร.}}\csromangathastanza{\csromangathabat{สุธํ อนาถปิณฺฑิ จ,}{สทฺโธ สีตวนํ อคา.} \\ \csromangathabat{ทิฏฺฐธมฺโม นิมนฺเตสิ,}{สห สํเฆน นายกํ.}}\csromangathastanza{\csromangathabat{อาณาเปสนฺตรามคฺเค,}{อารามํ การยี คโณ.} \\ \csromangathabat{เวสาลิยํ นวกมฺมํ,}{ปุรโต จ ปริคฺคหิ.}}\csromangathastanza{\csromangathabat{โก อรหติ ภตฺตคฺเค,}{ติตฺติรญฺจ อวนฺทิยา.} \\ \csromangathabat{ปริคฺคหิตนฺตรฆรา,}{ตูโล สาวตฺถิ โอสริ.}}\csromangathastanza{\csromangathabat{ปติฏฺฐาเปสิ อารามํ,}{ภตฺตคฺเค จ โกลาหลํ.} \\ \csromangathabat{คิลานา วรเสยฺยา จ,}{เลสา สตฺตรสา ตหิํ.}}\csromangathastanza{\csromangathabat{เกน นุ โข กถํ นุ โข,}{วิหารคฺเคน ภาชยิ.} \\ \csromangathabat{ปริเวณํ อนุภาคญฺจ,}{อกามา ภาคํ โน ทเท.}}\csromangathastanza{\csromangathabat{นิสฺสีมํ สพฺพกาลญฺจ,}{คาหา เสนาสเน ตโย.} \\ \csromangathabat{อุปนนฺโท จ วณฺเณสิ,}{ฐิตกา สมกาสนา.}}\csromangathastanza{\csromangathabat{สมานาสนิกา ภินฺทิํสุ,}{ติวคฺคา จ ทุวคฺคิกํ.} \\ \csromangathabat{อสมานาสนิกา ทีฆํ,}{สาฬินฺทํ ปริภุญฺชิตุํ.}}\csromangathastanza{\csromangathabat{อยฺยิกา จ อวิทูเร,}{ภาชิตญฺจ กีฏาคิเร.} \\ \csromangathabat{อาฬวี ปิณฺฑกกุฏฺเฏหิ,}{ทฺวารอคฺคฬวฏฺฏิกา.}}\csromangathastanza{\csromangathabat{อาโลกเสตกาฬญฺจ,}{เครุฉาทนพนฺธนา.} \\ \csromangathabat{ภณฺฑิขณฺฑปริภณฺฑํ,}{วีส ติํสา จ กาลิกา.}}\csromangathastanza{\csromangathabat{โอสิเต อกตํ วิปฺปํ,}{ขุทฺเท ฉปฺปญฺจ\-วสฺสิกํ.} \\ \csromangathabat{อฑฺฒโยเค จ สตฺตฏฺฐ,}{มหลฺเล ทส ทฺวาทส.}}\csromangathastanza{\csromangathabat{สพฺพํ วิหารํ เอกสฺส,}{อญฺญํ วาเสนฺติ สํฆิกํ.} \\ \csromangathabat{นิสฺสีมํ สพฺพกาลญฺจ,}{ปกฺกมิ วิพฺภมนฺติ จ.}}\csromangathastanza{\csromanfolio{335}\csromangathabat{กาลญฺจ สามเณรญฺจ,}{สิกฺขาปจฺจกฺขอนฺติมํ.} \\ \csromangathabat{อุมฺมตฺตขิตฺตจิตฺตา จ,}{เวทนาปตฺติทสฺสนา.}}\csromangathastanza{\csromangathabat{อปฺปฏิกมฺมทิฏฺฐิยา,}{ปณฺฑกา เถยฺยติตฺถิยา.} \\ \csromangathabat{ติรจฺฉานมาตุปิตุ,}{อรหนฺตา จ ทูสกา.}}\csromangathastanza{\csromangathabat{เภทกา โลหิตุปฺปาทา,}{อุภโต จาปิ พฺยญฺชนกา.} \\ \csromangathabat{มา สํฆสฺส ปริหายิ,}{กมฺมํ อญฺญสฺส ทาตเว.}}\csromangathastanza{\csromangathabat{วิปฺปกเต จ อญฺญสฺส,}{กเต ตสฺเสว ปกฺกเม.} \\ \csromangathabat{วิพฺภมติ กาลงฺกโต,}{สามเณโร จ ชายติ.}}\csromangathastanza{\csromangathabat{ปจฺจกฺขาโต จ สิกฺขาย,}{อนฺติมชฺฌาปนฺนโก ยทิ.} \\ \csromangathabat{สํโฆว สามิโก โหติ,}{อุมฺมตฺตขิตฺตเวทนา.}}\csromangathastanza{\csromangathabat{อทสฺสนาปฺปฏิกมฺเม,}{ทิฏฺฐิ ตสฺเสว โหติ ตํ.} \\ \csromangathabat{ปณฺฑโก เถยฺยติตฺถี จ,}{ติรจฺฉานมาตุเปตฺติกํ.}}\csromangathastanza{\csromangathabat{ฆาตโก ทูสโก จาปิ,}{เภทโลหิตพฺยญฺชนา.} \\ \csromangathabat{ปฏิชานาติ ยทิ โส,}{สํโฆว โหติ สามิโก.}}\csromangathastanza{\csromangathabat{หรนฺตญฺญตฺร กุกฺกุจฺจํ,}{อุนฺทฺริยติ จ กมฺพลํ.} \\ \csromangathabat{ทุสฺสํ จ จมฺมจกฺกลี,}{โจฬกํ อกฺกมนฺติ จ.}}\csromangathastanza{\csromangathabat{อลฺลา อุปาหนานิฏฺฐุ,}{ลิขนฺติ อปสฺเสนฺติ จ.} \\ \csromangathabat{อปสฺเสนํ ลิขเตว,}{โธตปจฺจตฺถเรน จ.}}\csromangathastanza{\csromangathabat{ราชคเห น สกฺโกนฺติ,}{ลามกํ ภตฺตุทฺเทสกํ.} \\ \csromangathabat{กถํ นุ โข ปญฺญาปกํ,}{ภณฺฑาคาริกสมฺมุติ.}}\csromangathastanza{\csromangathabat{ปฏิคฺคาหภาชโก จาปิ,}{ยาคุ จ ผลภาชโก.} \\ \csromangathabat{ขชฺชกภาชโก เจว,}{อปฺปมตฺตกวิสฺสชฺเช.}}\csromangathastanza{\csromangathabat{สาฏิย\-คฺคาหาปโก เจว,}{ตเถว ปตฺตคฺคาหโก.} \\ \csromangathabat{อารามิกสามเณร,}{เปสกสฺส จ สมฺมุติ.} \\ \csromangathabat{สพฺพาภิภู โลกวิทู,}{หิตจิตฺโต วินายโก.}}}{เลณตฺถญฺจ สุขตฺถญฺจ, ฌายิตุญฺจ วิปสฺสิตุนฺติ. เสนาสน\-กฺขนฺธโก นิฏฺฐิโต.}
\csromanmatikaanchor{mtk.160}
\csromantocmarknopage{chapter}{7. สํฆเภทกกฺขนฺธก}{mtk.160}
\bookmark[dest={mtk.160},level=0]{7. สํฆเภทกกฺขนฺธก}
\chapterheadpageread{7. สํฆเภทก\-กฺขนฺธก}
\csromanmatikaanchor{mtk.161}
\csromantocmarkref{section}{1. ปฐมภาณวาร}{mtk.161}
\bookmark[dest={mtk.161},level=1]{1. ปฐมภาณวาร}
\csromanheader{1}{1. ปฐม\-ภาณวาร}
\csromanmatikaanchor{mtk.162}
\csromantocmarkref{subsection}{ฉสกฺยปพฺพชฺชากถา}{mtk.162}
\bookmark[dest={mtk.162},level=2]{ฉสกฺยปพฺพชฺชากถา}
\csromanheader{2}{ฉสกฺย\-ปพฺพชฺชา\-กถา}
\proseitem{330}{\csromanfolio{336}เตน สมเยน พุทฺโธ ภควา อนุปิยายํ วิหรติ อนุปิยํ นาม มลฺลานํ นิคโม. เตน โข ปน สมเยน อภิญฺญาตา อภิญฺญาตา สกฺยกุมารา ภควนฺตํ ปพฺพชิตํ อนุปพฺพชนฺติ. เตน โข ปน สมเยน มหานาโม จ สกฺโก อนุรุทฺโธ จ สกฺโก เทฺวภาติกา โหนฺติ. อนุรุทฺโธ สกฺโก สุขุมาโล โหติ, ตสฺส ตโย ปาสาทา โหนฺติ เอโก เหมนฺติโก เอโก คิมฺหิโก เอโก วสฺสิโก, โส วสฺสิเก ปาสาเท จตฺตาโร มาเส\csromansharedfootnote{4959b46ba7ed28d4}{วสฺสิเก ปาสาเท วสฺสิเก จตฺตาโร มาเส (สี)} นิปฺปุริเสหิ ตูริเยหิ ปริจารยมาโน\csromansharedfootnote{a46ebbc38925421b}{ปริจาริยมาโน (ก)} น เหฏฺฐาปาสาทํ โอโรหติ. อถ โข มหานามสฺส สกฺกสฺส เอตทโหสิ “เอตรหิ โข อภิญฺญาตา อภิญฺญาตา สกฺยกุมารา ภควนฺตํ ปพฺพชิตํ อนุปพฺพชนฺติ, อมฺหากญฺจ ปน กุลา นตฺถิ โกจิ อคารสฺมา อนคาริยํ ปพฺพชิโต, ยํนูนาหํ วา ปพฺพเชยฺยํ, อนุรุทฺโธ วา”ติ. อถ โข มหานาโม สกฺโก เยน อนุรุทฺโธ สกฺโก เตนุปสงฺกมิ, อุปสงฺกมิตฺวา อนุรุทฺธํ สกฺกํ เอตทโวจ “เอตรหิ ตาต อนุรุทฺธ อภิญฺญาตา อภิญฺญาตา สกฺยกุมารา ภควนฺตํ ปพฺพชิตํ อนุปพฺพชนฺติ, อมฺหากญฺจ ปน กุลา นตฺถิ โกจิ อคารสฺมา อนคาริยํ ปพฺพชิโต, เตน หิ ตฺวํ วา ปพฺพช, อหํ วา ปพฺพชิสฺสามี”ติ, อหํ โข สุขุมาโล, นาหํ สกฺโกมิ อคารสฺมา อนคาริยํ ปพฺพชิตุํ, ตฺวํ ปพฺพชาหีติ. เอหิ โข เต ตาต อนุรุทฺธ, ฆราวาสตฺถํ อนุสาสิสฺสามิ. ปฐมํ เขตฺตํ กสาเปตพฺพํ, กสาเปตฺวา วปาเปตพฺพํ, วปาเปตฺวา อุทกํ อภิเนตพฺพํ, อุทกํ อภิเนตฺวา อุทกํ นินฺเนตพฺพํ, อุทกํ นินฺเนตฺวา นิทฺธาเปตพฺพํ, นิทฺธาเปตฺวา\csromansharedfootnote{39733c4eb3a459a3}{นิฑฺฑเหตพฺพํ, นิฑฺฑตฺวา (สี)} ลวาเปตพฺพํ, ลวาเปตฺวา อุพฺพาหาเปตพฺพํ, อุพฺพาหาเปตฺวา ปุญฺชํ การาเปตพฺพํ, ปุญฺชํ การาเปตฺวา มทฺทาเปตพฺพํ, มทฺทาเปตฺวา ปลาลานิ อุทฺธราเป\-ตพฺพานิ, ปลาลานิ อุทฺธราเปตฺวา ภุสิกา อุทฺธราเปตพฺพา, ภุสิกํ อุทฺธราเปตฺวา โอปุนาเปตพฺพํ\csromansharedfootnote{b390a83088e4c128}{โอผุนาเปตพฺพํ (สฺยา, ก), โอผุณาเปตพฺพํ (โยชนา)}, โอปุนาเปตฺวา อติหราเป\-ตพฺพํ, อติหราเปตฺวา \csromanfolio{337}อายติมฺปิ วสฺสํ เอวเมว กาตพฺพํ, อายติมฺปิ วสฺสํ เอวเมว กาตพฺพนฺติ. น กมฺมา ขียนฺติ, น กมฺมานํ อนฺโต ปญฺญายติ, กทา กมฺมา ขียิสฺสนฺติ, กทา กมฺมานํ อนฺโต ปญฺญายิสฺสติ, กทา มยํ อปฺโปสฺสุกฺกา ปญฺจหิ กามคุเณหิ สมปฺปิตา สมงฺคีภูตา ปริจาเรสฺสามาติ. น หิ ตาต อนุรุทฺธ กมฺมา ขียนฺติ, น กมฺมานํ อนฺโต ปญฺญายติ, อขีเณว กมฺเม ปิตโร จ ปิตามหา จ กาลงฺกตาติ. เตน หิ ตฺวญฺเญว ฆราวาส\-ตฺเถน อุปชานาหิ, อหํ อคารสฺมา อนคาริยํ ปพฺพชิสฺสามีติ.}

\prose{อถ โข อนุรุทฺโธ สกฺโก เยน มาตา เตนุปสงฺกมิ, อุปสงฺกมิตฺวา มาตรํ เอตทโวจ “อิจฺฉามหํ อมฺม อคารสฺมา อนคาริยํ ปพฺพชิตุํ, อนุชานาหิ มํ อคารสฺมา อนคาริยํ ปพฺพชฺชายา”ติ. เอวํ วุตฺเต อนุรุทฺธสฺส สกฺกสฺส มาตา อนุรุทฺธํ สกฺกํ เอตทโวจ “ตุเมฺห โข เม ตาต อนุรุทฺธ เทฺว ปุตฺตา ปิยา มนาปา อปฺปฏิกูลา, มรเณนปิ โว อกามกา วินา ภวิสฺสามิ, กิํ ปนาหํ ตุเมฺห ชีวนฺเต อนุชานิสฺสามิ อคารสฺมา อนคาริยํ ปพฺพชฺชายา”ติ. ทุติยมฺปิ โข ฯเปฯ. ตติยมฺปิ โข อนุรุทฺโธ สกฺโก มาตรํ เอตทโวจ “อิจฺฉามหํ อมฺม อคารสฺมา อนคาริยํ ปพฺพชิตุํ, อนุชานาหิ มํ อคารสฺมา อนคาริยํ ปพฺพชฺชายา”ติ\csromansharedfootnote{a9271ba5a8269c2b}{อิมสฺส อนนฺตรํ เกสุจิ โปตฺถเกสุ เอวมฺปิ ปาโฐ ทิสฺสติ\csromandash{} “เอวํ วุตฺเต อนุรุทฺธสฺส สกฺกสฺส มตา เอวมาห ‘สเจ ตาต อนุรุทฺธ ภทฺทิโย สกฺยราชา สกฺยานํ รชฺชํ กาเรติ อคารสฺมา อนคาริยํ ปพฺพชติ, เอวํ ตฺวมฺปิ ปพฺพชาหี’ติ”.}. เตน โข ปน สมเยน ภทฺทิโย สกฺยราชา สกฺยานํ รชฺชํ กาเรสิ, โส จ อนุรุทฺธสฺส สกฺกสฺส สหาโย โหติ. อถ โข อนุรุทฺธสฺส สกฺกสฺส มาตา “อยํ โข ภทฺทิโย สกฺยราชา สกฺยานํ รชฺชํ กาเรติ อนุรุทฺธสฺส สกฺกสฺส สหาโย, โส น อุสฺสหติ อคารสฺมา อนคาริยํ ปพฺพชิตุนฺ”ติ อนุรุทฺธํ สกฺกํ เอตทโวจ “สเจ ตาต อนุรุทฺธ ภทฺทิโย สกฺยราชา อคารสฺมา อนคาริยํ ปพฺพชติ, เอวํ ตฺวมฺปิ ปพฺพชาหี”ติ. อถ โข อนุรุทฺโธ สกฺโก เยน ภทฺทิโย สกฺยราชา เตนุปสงฺกมิ, อุปสงฺกมิตฺวา ภทฺทิยํ สกฺยราชานํ เอตทโวจ “มม โข สมฺม ปพฺพชฺชา ตว ปฏิพทฺธา”ติ. สเจ เต สมฺม ปพฺพชฺชา มม ปฏิพทฺธา วา อปฺปฏิพทฺธา วา สา โหตุ, อหํ ตยา, ยถา สุขํ ปพฺพชาหีติ. เอหิ สมฺม อุโภ อคารสฺมา อนคาริยํ ปพฺพชิสฺสามาติ. นาหํ สมฺม สกฺโกมิ อคารสฺมา อนคาริยํ ปพฺพชิตุนฺติ. ยํ เต สกฺกา อญฺญํ \csromanfolio{338}มยา กาตุํ, กฺยาหํ\csromansharedfootnote{d4de712fded8fc12}{ตฺยาหํ (สี, สฺยา)} กริสฺสามิ, ตฺวํ ปพฺพชาหีติ. มาตา โข มํ สมฺม เอวมาห “สเจ ตาต อนุรุทฺธ ภทฺทิโย สกฺยราชา อคารสฺมา อนคาริยํ ปพฺพชติ, เอวํ ตฺวมฺปิ ปพฺพชาหี”ติ. ภาสิตา โข ปน เต สมฺม เอสา วาจา “สเจ เต สมฺม ปพฺพชฺชา มม ปฏิพทฺธา วา อปฺปฏิพทฺธา วา สา โหตุ, อหํ ตยา, ยถา สุขํ ปพฺพชาหี”ติ. เอหิ สมฺม, อุโภ อคารสฺมา อนคาริยํ ปพฺพชิสฺสามาติ.}

\prose{เตน โข ปน สมเยน มนุสฺสา สจฺจวาทิโน โหนฺติ สจฺจปฏิญฺญา. อถ โข ภทฺทิโย สกฺยราชา อนุรุทฺธํ สกฺกํ เอตทโวจ “อาคเมหิ สมฺม สตฺตวสฺสานิ, สตฺตนฺนํ วสฺสานํ อจฺจเยน อุโภ อคารสฺมา อนคาริยํ ปพฺพชิสฺสามา”ติ. อติจิรํ สมฺม สตฺตวสฺสานิ, นาหํ สกฺโกมิ สตฺตวสฺสานิ อาคเมตุนฺติ. อาคเมหิ สมฺม ฉวสฺสานิ ฯเปฯ ปญฺจวสฺสานิ. จตฺตาริ วสฺสานิ. ตีณิ วสฺสานิ. เทฺว วสฺสานิ. เอกํ วสฺสํ, เอกสฺส วสฺสสฺส อจฺจเยน อุโภ อคารสฺมา อนคาริยํ ปพฺพชิสฺสามาติ. อติจิรํ สมฺม เอกวสฺสํ, นาหํ สกฺโกมิ เอกํ วสฺสํ อาคเมตุนฺติ. อาคเมหิ สมฺม สตฺตมาเส, สตฺตนฺนํ มาสานํ อจฺจเยน อุโภปิ อคารสฺมา อนคาริยํ ปพฺพชิสฺสามาติ. อติจิรํ สมฺม สตฺตมาสา, นาหํ สกฺโกมิ สตฺตมาเส อาคเมตุนฺติ. อาคเมหิ สมฺม ฉ มาเส ฯเปฯ ปญฺจ มาเส. จตฺตาโร มาเส. ตโย มาเส. เทฺว มาเส. เอกํ มาสํ. อฑฺฒมาสํ, อฑฺฒมาสสฺส อจฺจเยน อุโภปิ อคารสฺมา อนคาริยํ ปพฺพชิสฺสามาติ. อติจิรํ สมฺม อฑฺฒมาโส, นาหํ สกฺโกมิ อฑฺฒมาสํ อาคเมตุนฺติ. อาคเมหิ สมฺม สตฺตาหํ, ยาวาหํ ปุตฺเต จ ภาตโร จ รชฺชํ นิยฺยาเทมีติ. น จิรํ สมฺม สตฺตาโห อาคเมสฺสามีติ.}

\proseitem{331}{อถ โข ภทฺทิโย จ สกฺยราชา อนุรุทฺโธ จ อานนฺโท จ ภคุ จ กิมิโล จ เทวทตฺโต จ อุปาลิกปฺปเกน สตฺตมา ยถา ปุเร จตุรงฺคินิยา เสนาย อุยฺยานภูมิํ นิยฺยนฺติ, เอวเมว จตุรงฺคินิยา เสนาย นิยฺยิํสุ. เต ทูรํ คนฺตฺวา เสนํ นิวตฺตาเปตฺวา ปรวิสยํ โอกฺกมิตฺวา อาภรณํ โอมุญฺจิตฺวา อุตฺตราสงฺเคน ภณฺฑิกํ พนฺธิตฺวา อุปาลิํ กปฺปกํ เอตทโวจุํ “หนฺท ภเณ อุปาลิ นิวตฺตสฺสุ, อลํ เต เอตฺตกํ ชีวิกายา”ติ. อถ โข อุปาลิสฺส กปฺปกสฺส นิวตฺตนฺตสฺส เอตทโหสิ “จณฺฑา โข สากิยา, ‘อิมินา กุมารา นิปฺปาติตา’ติ ฆาตาเปยฺยุมฺปิ มํ. \csromanfolio{339}อิเม หิ นาม สกฺยกุมารา อคารสฺมา อนคาริยํ ปพฺพชิสฺสนฺติ, กิมงฺคํ\csromansharedfootnote{000f45a8a6c5a7bf}{กิมงฺค (สี)} ปนาหนฺ’ติ ภณฺฑิกํ มุญฺจิตฺวา ตํ ภณฺฑํ รุกฺเข อาลคฺเคตฺวา “โย ปสฺสติ, ทินฺนํเยว หรตู”ติ วตฺวา เยน เต สกฺยกุมารา เตนุปสงฺกมิ, อทฺทสาสุํ โข เต สกฺยกุมารา อุปาลิํ กปฺปกํ ทูรโตว อาคจฺฉนฺตํ, ทิสฺวาน อุปาลิํ กปฺปกํ เอตทโวจุํ “กิสฺส ภเณ อุปาลิ นิวตฺเตสี”ติ. อิธ เม อยฺยปุตฺตา นิวตฺตนฺตสฺส เอตทโหสิ “จณฺฑา โข สากิยา, ‘อิมินา กุมารา นิปฺปาติตา’ติ ฆาตาเปยฺยุมฺปิ มํ. อิเม หิ นาม สกฺยกุมารา อคารสฺมา อนคาริยํ ปพฺพชิสฺสนฺติ, กิมงฺคํ ปนาหนฺ”ติ. โส โข อหํ อยฺยปุตฺตา ภณฺฑิกํ มุญฺจิตฺวา ตํ ภณฺฑํ รุกฺเข อาลคฺเคตฺวา “โย ปสฺสติ, ทินฺนํเยว หรตู”ติ วตฺวา ตโตมฺหิ ปฏินิวตฺโตติ. สุฏฺฐุ ภเณ อุปาลิ อกาสิ ยมฺปิ น นิวตฺโต\csromansharedfootnote{24090f378eb08039}{ยํ นิวตฺโต (สี), ยํ ปน นิวตฺโต (สฺยา)}. จณฺฑา โข สากิยา “อิมินา กุมารา นิปฺปาติตา”ติ ฆาตาเปยฺยุมฺปิ ตนฺติ.}

\prose{อถ โข สกฺยกุมารา อุปาลิํ กปฺปกํ อาทาย เยน ภควา เตนุ\-ปสงฺกมิํสุ, อุปสงฺกมิตฺวา ภควนฺตํ อภิวาเทตฺวา เอกมนฺตํ นิสีทิํสุ. เอกมนฺตํ นิสินฺนา โข เต สกฺยกุมารา ภควนฺตํ เอตทโวจุํ “มยํ ภนฺเต สากิยา นาม มานสฺสิโน, อยํ ภนฺเต อุปาลิ กปฺปโก อมฺหากํ ทีฆรตฺตํ ปริจารโก, อิมํ ภควา ปฐมํ ปพฺพาเชตุ, อิมสฺส มยํ อภิวาทน\-ปจฺจุฏฺฐาน\-อญฺชลิกมฺม\-สามีจิกมฺมํ กริสฺสาม, เอวํ อมฺหากํ สากิยานํ สากิยมาโน นิมฺมานายิสฺสตี”ติ\csromansharedfootnote{ffdc73d60520d2e8}{นิมฺมาทยิสฺสตีติ (สี), นิมฺมนิยิสฺสตีติ (สฺยา)}.}

\prose{อถ โข ภควา อุปาลิํ กปฺปกํ ปฐมํ ปพฺพาเชสิ, ปจฺฉา เต สกฺยกุมาเร. อถ โข อายสฺมา ภทฺทิโย เตเนว อนฺตรวสฺเสน ติสฺโส วิชฺชา สจฺฉากาสิ, อายสฺมา อนุรุทฺโธ ทิพฺพจกฺขุํ อุปฺปาเทสิ, อายสฺมา อานนฺโท โสตาปตฺติผลํ สจฺฉากาสิ, เทวทตฺโต โปถุชฺชนิกํ อภินิปฺผาเทสิ.}

\proseitem{332}{เตน โข ปน สมเยน อายสฺมา ภทฺทิโย อรญฺญคโตปิ รุกฺขมูล\-คโตปิ สุญฺญาคาร\-คโตปิ อภิกฺขณํ อุทานํ อุทาเนสิ “อโห สุขํ, อโห สุขนฺ”ติ. อถ โข สมฺพหุลา ภิกฺขู เยน ภควา เตนุ\-ปสงฺกมิํสุ, อุปสงฺกมิตฺวา ภควนฺตํ อภิวาเทตฺวา เอกมนฺตํ นิสีทิํสุ. \csromanfolio{340}เอกมนฺตํ นิสินฺนา โข เต ภิกฺขู ภควนฺตํ เอตทโวจุํ “อายสฺมา ภนฺเต ภทฺทิโย อรญฺญคโตปิ รุกฺขมูล\-คโตปิ สุญฺญาคาร\-คโตปิ อภิกฺขณํ อุทานํ อุทาเนสิ ‘อโห สุขํ, อโห สุขนฺ’ติ, นิสฺสํสยํ โข ภนฺเต อายสฺมา ภทฺทิโย อนภิรโตว พฺรหฺมจริยํ จรติ, ตํเยว วา ปุริมํ รชฺชสุขํ สมนุสฺสรนฺโต อรญฺญคโตปิ รุกฺขมูล\-คโตปิ สุญฺญาคาร\-คโตปิ อภิกฺขณํ อุทานํ อุทาเนสิ ‘อโห สุขํ, อโห สุขนฺ’ติ”.}

\prose{อถ โข ภควา อญฺญตรํ ภิกฺขุํ อามนฺเตสิ\csromandash{}เอหิ ตฺวํ ภิกฺขุ มม วจเนน ภทฺทิยํ ภิกฺขุํ อามนฺเตหิ “สตฺถา ตํ อาวุโส ภทฺทิย อามนฺเตตี”ติ. “เอวํ ภนฺเต”ติ โข โส ภิกฺขุ ภควโต ปฏิสฺสุตฺวา เยนายสฺมา ภทฺทิโย เตนุปสงฺกมิ, อุปสงฺกมิตฺวา อายสฺมนฺตํ ภทฺทิยํ เอตทโวจ “สตฺถา ตํ อาวุโส ภทฺทิย อามนฺเตตี”ติ. “เอวมาวุโส”ติ โข อายสฺมา ภทฺทิโย ตสฺส ภิกฺขุโน ปฏิสฺสุตฺวา เยน ภควา เตนุปสงฺกมิ, อุปสงฺกมิตฺวา ภควนฺตํ อภิวาเทตฺวา เอกมนฺตํ นิสีทิ. เอกมนฺตํ นิสินฺนํ โข อายสฺมนฺตํ ภทฺทิยํ ภควา เอตทโวจ “สจฺจํ กิร ตฺวํ ภทฺทิย อรญฺญคโตปิ รุกฺขมูล\-คโตปิ สุญฺญาคาร\-คโตปิ อภิกฺขณํ อุทานํ อุทาเนสิ “อโห สุขํ อโห สุขนฺ”ติ. เอวํ ภนฺเตติ. กิํ ปน ตฺวํ ภทฺทิย อตฺถวสํ สมฺปสฺสมาโน อรญฺญคโตปิ รุกฺขมูล\-คโตปิ สุญฺญาคาร\-คโตปิ อภิกฺขณํ อุทานํ อุทาเนสิ “อโห สุขํ อโห สุขนฺ”ติ. ปุพฺเพ เม ภนฺเต รญฺโญ สโตปิ อนฺโตปิ อนฺเตปุเร รกฺขา สุสํวิหิตา โหติ, พหิปิ อนฺเตปุเร รกฺขา สุสํวิหิตา โหติ, อนฺโตปิ นคเร รกฺขา สุสํวิหิตา โหติ, พหิปิ นคเร รกฺขา สุสํวิหิตา โหติ, อนฺโตปิ ชนปเท รกฺขา สุสํวิหิตา โหติ, พหิปิ ชนปเท รกฺขา สุสํวิหิตา โหติ. โส โข อหํ ภนฺเต เอวํ รกฺขิโตปิ โคปิโตปิ สนฺโต ภีโต อุพฺพิคฺโค อุสฺสงฺกี อุตฺรสฺโต วิหรามิ. เอตรหิ โข ปน อหํ เอโก ภนฺเต อรญฺญคโตปิ รุกฺขมูล\-คโตปิ สุญฺญาคาร\-คโตปิ อภีโต อนุพฺพิคฺโค อนุสฺสงฺกี อนุตฺรสฺโต อปฺโปสฺสุกฺโก ปนฺนโลโม ปรทตฺตวุตฺโต มิคภูเตน เจตสา วิหรามีติ อิมํ โข อหํ ภนฺเต อตฺถวสํ สมฺปสฺสมาโน อรญฺญคโตปิ รุกฺขมูล\-คโตปิ สุญฺญาคาร\-คโตปิ อภิกฺขณํ อุทานํ อุทาเนมิ “อโห สุขํ อโห สุขนฺ”ติ. อถ โข ภควา เอตมตฺถํ วิทิตฺวา ตายํ เวลายํ อิมํ อุทานํ อุทาเนสิ\csromandash{}}

\setlength{\csromangathapagewidth}{0pt}
\setlength{\csromangathawakwidth}{0pt}
\csromangathameasuregroup{}{\textsuperscript{*}“ยสฺสนฺตรโต น สนฺติ โกปา, \\ อิติ ภวาภวตญฺจ วีติวตฺโต. \\ ตํ วิคตภยํ สุขิํ อโสกํ, \\ เทวา นานุภวนฺติ ทสฺสนายา”ติ.}
\csromangathameasurewak{\textsuperscript{*}“ยสฺสนฺตรโต น สนฺติ โกปา, \\ อิติ ภวาภวตญฺจ วีติวตฺโต. \\ ตํ วิคตภยํ สุขิํ อโสกํ, \\ เทวา นานุภวนฺติ ทสฺสนายา”ติ.}
\setlength{\csromangathaleftcol}{0pt}
\csromangathagroup{\csromangathastanza{\csromangathawak{\csromanfolio{341}\csromansymbolfootnote{*}{ขุ 1. 100 ปิฏฺเฐ อุทาเนปิ.}“ยสฺสนฺตรโต น สนฺติ โกปา, \\ อิติ ภวาภวตญฺจ วีติวตฺโต. \\ ตํ วิคตภยํ สุขิํ อโสกํ, \\ เทวา นานุภวนฺติ ทสฺสนายา”ติ.}}}

\csromanmatikaanchor{mtk.163}
\csromantocmarkref{subsection}{เทวทตฺตวตฺถุ}{mtk.163}
\bookmark[dest={mtk.163},level=2]{เทวทตฺตวตฺถุ}
\csromanheader{2}{เทวทตฺต\-วตฺถุ}
\proseitem{33}{อถ โข ภควา อนุปิยายํ ยถาภิรนฺตํ วิหริตฺวา เยน โกสมฺพี เตน จาริกํ ปกฺกามิ, อนุปุพฺเพน จาริกํ จรมาโน เยน โกสมฺพี ตทวสริ, ตตฺร สุทํ ภควา โกสมฺพิยํ วิหรติ โฆสิตาราเม. อถ โข เทวทตฺตสฺส รโหคตสฺส ปฏิสลฺลีนสฺส เอวํ เจตโส ปริวิตกฺโก อุทปาทิ “กํ นุ โข อหํ ปสาเทยฺยํ, ยสฺมิํ เม ปสนฺเน พหุลาภสกฺกาโร อุปฺปชฺเชยฺยา”ติ. อถ โข เทวทตฺตสฺส เอตทโหสิ “อยํ โข อชาตสตฺตุ กุมาโร ตรุโณ เจว อายติํ ภทฺโท จ, ยํนุนาหํ อชาตสตฺตุํ กุมารํ ปสาเทยฺยํ, ตสฺมิํ เม ปสนฺเน พหุลาภสกฺกาโร อุปฺปชฺชิสฺสตี”ติ.}

\prose{อถ โข เทวทตฺโต เสนาสนํ สํสาเมตฺวา ปตฺต\-จีวรมา\-ทาย เยน ราชคหํ เตน ปกฺกามิ, อนุปุพฺเพน เยน ราชคหํ ตทวสริ. อถ โข เทวทตฺโต สกวณฺณํ ปฏิสํหริตฺวา กุมารกวณฺณํ อภินิมฺมินิตฺวา อหิเมขลิกาย อชาตสตฺตุสฺส กุมารสฺส อุจฺฉงฺเค\csromansharedfootnote{14c1408ebe8f2606}{อุจฺจงฺเก (สฺยา)} ปาตุรโหสิ. อถ โข อชาตสตฺตุ กุมาโร ภีโต อโหสิ อุพฺพิคฺโค อุสฺสงฺกี อุตฺรสฺโต. อถ โข เทวทตฺโต อชาตสตฺตุํ กุมารํ เอตทโวจ “ภายสิ มํ ตฺวํ กุมารา”ติ. อาม ภายามิ, โกสิ ตฺวนฺติ. อหํ เทวทตฺโตติ. สเจ โข ตฺวํ ภนฺเต อยฺโย เทวทตฺโต อิงฺฆ สเกเนว วณฺเณน ปาตุภวสฺสูติ. อถ โข เทวทตฺโต กุมารกวณฺณํ ปฏิสํหริตฺวา สํฆาฏิปตฺตจีวรธโร อชาตสตฺตุสฺส กุมารสฺส ปุรโต อฏฺฐาสิ. อถ โข อชาตสตฺตุ กุมาโร เทวทตฺตสฺส อิมินา อิทฺธิ\-ปาฏิหาริเยน อภิปฺปสนฺโน ปญฺจหิ รถสเตหิ สายํ ปาตํ อุปฏฺฐานํ คจฺฉติ, ปญฺจ จ ถาลิปาก\-สตานิ ภตฺตาภิหาโร อภิหรียติ. อถ โข เทวทตฺตสฺส ลาภ\-สกฺการ\-สิโลเกน อภิภูตสฺส ปริยาทินฺน\-จิตฺตสฺส เอวรูปํ \csromanfolio{342}อิจฺฉาคตํ อุปฺปชฺชิ “อหํ ภิกฺขุสํฆํ ปริหริสฺสามี”ติ. สห จิตฺตุปฺปาทาว เทวทตฺโต ตสฺสา อิทฺธิยา ปริหายิ.}

\prose{\csromansymbolfootnote{*}{อํ 2. 107 ปิฏฺเฐปิ.}เตน โข ปน สมเยน กกุโธ นาม โกฬิยปุตฺโต อายสฺมโต มหา\-โมคฺคลฺลานสฺส อุปฏฺฐาโก อธุนา กาลงฺกโต อญฺญตรํ มโนมยํ กายํ อุปปนฺโน. ตสฺส เอวรูโป อตฺตภาว\-ปฺปฏิลาโภ โหติ เสยฺยถาปิ นาม เทฺว วา ตีณิ วา มาคธกานิ\csromansharedfootnote{21df51e7d072614f}{มาคธิกานิ (สฺยา)} คามกฺเขตฺตานิ. โส เตน อตฺตภาว\-ปฺปฏิลาเภน เนว อตฺตานํ น ปรํ พฺยาพาเธติ. อถ โข กกุโธ เทวปุตฺโต เยนายสฺมา มหาโมคฺคลฺลาโน เตนุปสงฺกมิ, อุปสงฺกมิตฺวา อายสฺมนฺตํ มหา\-โมคฺคลฺลานํ อภิวาเทตฺวา เอกมนฺตํ อฏฺฐาสิ. เอกมนฺตํ ฐิโต โข กกุโธ เทวปุตฺโต อายสฺมนฺตํ มหา\-โมคฺคลฺลานํ เอตทโวจ “เทวทตฺตสฺส ภนฺเต ลาภ\-สกฺการ\-สิโลเกน อภิภูตสฺส ปริยาทินฺน\-จิตฺตสฺส\csromansharedfootnote{b31e5978c0545de7}{ปริยาทิณฺณจิตฺตสฺส (ก)} เอวรูปํ อิจฺฉาคตํ อุปฺปชฺชิ ‘อหํ ภิกฺขุสํฆํ ปริหริสฺสามี”ติ. สห จิตฺตุปฺปาทาว ภนฺเต เทวทตฺโต ตสฺสา อิทฺธิยา ปริหีโน”ติ. อิทมโวจ กกุโธ เทวปุตฺโต, อิทํ วตฺวา อายสฺมนฺตํ มหา\-โมคฺคลฺลานํ อภิวาเทตฺวา ปทกฺขิณํ กตฺวา ตตฺเถว อนฺตรธายิ.}

\prose{อถ โข อายสฺมา มหาโมคฺคลฺลาโน เยน ภควา เตนุปสงฺกมิ, อุปสงฺกมิตฺวา ภควนฺตํ อภิวาเทตฺวา เอกมนฺตํ นิสีทิ. เอกมนฺตํ นิสินฺโน โข อายสฺมา มหาโมคฺคลฺลาโน ภควนฺตํ เอตทโวจ “กกุโธ นาม ภนฺเต โกฬิยปุตฺโต มม อุปฏฺฐาโก อธุนา กาลงฺกโต อญฺญตรํ มโนมยํ กายํ อุปปนฺโน, ตสฺส เอวรูโป อตฺตภาว\-ปฺปฏิลาโภ เสยฺยถาปิ นาม เทฺว วา ตีณิ วา มาคธกานิ คามกฺเขตฺตานิ. โส เตน อตฺตภาว\-ปฺปฏิลาเภน เนว อตฺตานํ น ปรํ พฺยาพาเธติ. อถ โข ภนฺเต กกุโธ เทวปุตฺโต เยนาหํ เตนุปสงฺกมิ, อุปสงฺกมิตฺวา มํ อภิวาเทตฺวา เอกมนฺตํ อฏฺฐาสิ. เอกมนฺตํ ฐิโต โข ภนฺเต กกุโธ เทวปุตฺโต มํ เอตทโวจ “เทวทตฺตสฺส ภนฺเต ลาภ\-สกฺการ\-สิโลเกน อภิภูตสฺส ปริยาทินฺน\-จิตฺตสฺส เอวรูปํ อิจฺฉาคตํ อุปฺปชฺชิ ‘อหํ ภิกฺขุสํฆํ ปริหริสฺสามี’ติ, สห จิตฺตุปฺปาทาว ภนฺเต เทวทตฺโต ตสฺสา อิทฺธิยา ปริหีโนติ. อิทมโวจ ภนฺเต กกุโธ เทวปุตฺโต, อิทํ วตฺวา มํ อภิวาเทตฺวา ปทกฺขิณํ กตฺวา ตตฺเถว อนฺตรธายี”ติ.}

\prose{\csromanfolio{343}กิํ ปน เต โมคฺคลฺลาน กกุโธ เทวปุตฺโต เจตสา เจโต ปริจฺจ วิทิโต, ยํ กิญฺจิ กกุโธ เทวปุตฺโต ภาสติ, สพฺพํ ตํ ตเถว โหติ, โน อญฺญถาติ. เจตสา เจโต ปริจฺจ วิทิโต จ เม ภนฺเต กกุโธ เทวปุตฺโต, ยํ กิญฺจิ กกุโธ เทวปุตฺโต ภาสติ, สพฺพํ ตํ ตเถว โหติ, โน อญฺญถาติ. รกฺขสฺเสตํ โมคฺคลฺลาน วาจํ, รกฺขสฺเสตํ โมคฺคลฺลาน วาจํ, อิทานิ โส โมฆปุริโส อตฺตนาว อตฺตานํ ปาตุกริสฺสติ.}

\csromanmatikaanchor{mtk.164}
\csromantocmarkref{subsection}{ปญฺจสตฺถุกถา}{mtk.164}
\bookmark[dest={mtk.164},level=2]{ปญฺจสตฺถุกถา}
\csromanheader{2}{ปญฺจ\-สตฺถุ\-กถา}
\proseitem{334}{\csromansymbolfootnote{*}{อํ 2. 108 ปิฏฺเฐปิ.}ปญฺจิเม โมคฺคลฺลาน สตฺถาโร สนฺโต สํวิชฺชมานา โลกสฺมิํ. กตเม ปญฺจ. อิธ โมคฺคลฺลาน เอกจฺโจ สตฺถา อปริสุทฺธ\-สีโล สมาโน “ปริสุทฺธสีโลมฺหี”ติ ปฏิชานาติ “ปริสุทฺธํ เม สีลํ ปริโยทาตํ อสํกิลิฏฺฐนฺ”ติ จ. ตเมนํ สาวกา เอวํ ชานนฺติ “อยํ โข ภวํ สตฺถา อปริสุทฺธ\-สีโล สมาโน ‘ปริสุทฺธสีโลมฺหี’ติ ปฏิชานาติ ‘ปริสุทฺธํ เม สีลํ ปริโยทาตํ อสํกิลิฏฺฐนฺ’ติ จ. มยญฺเจว โข ปน คิหีนํ อาโรเจยฺยาม ‘นาสฺสสฺส มนาปํ, ยํ โข ปนสฺส อมนาปํ กถํ นํ มยํ เตน สมุทาจเรยฺยาม. สมฺมนฺนติ โข ปน จีวร\-ปิณฺฑปาต\-เสนาสน\-คิลานปฺปจฺจย\-เภสชฺช\-ปริกฺขาเรน. ยํ ตุโม กริสฺสติ, ตุโมว เตน ปญฺญายิสฺสตี”ติ. เอวรูปํ โข โมคฺคลฺลาน สตฺถารํ สาวกา สีลโต รกฺขนฺติ, เอวรูโป จ ปน สตฺถา สาวเกหิ สีลโต รกฺขํ ปจฺจาสีสติ\csromansharedfootnote{4c7995c09b7d4ebc}{ปจฺจาสิํสติ (สี, สฺยา)}.}

\prose{ปุน จปรํ โมคฺคลฺลาน อิเธกจฺโจ สตฺถา อปริสุทฺธา\-ชีโว สมาโน “ปริสุทฺธาชีโวมฺหี”ติ ปฏิชานาติ “ปริสุทฺโธ เม อาชีโว ปริโยทาโต อสํกิลิฏฺโฐ”ติ จ. ตเมนํ สาวกา เอวํ ชานนฺติ “อยํ โข ภวํ สตฺถา อปริสุทฺธา\-ชีโว สมาโน ‘ปริสุทฺธาชีโวมฺหี’ติ ปฏิชานาติ ‘ปริสุทฺโธ เม อาชีโว ปริโยทาโต อสํกิลิฏฺโฐ’ติ จ. มยญฺเจว โข ปน คิหีนํ อาโรเจยฺยาม ‘นาสฺสสฺส มนาปํ, ยํ โข ปนสฺส อมนาปํ, กถํ นํ มยํ เตน สมุทาจเรยฺยาม. สมฺมนฺนติ โข ปน จีวร\-ปิณฺฑปาต\-เสนาสน\-คิลานปฺปจฺจย\-เภสชฺช\-ปริกฺขาเรน. ยํ \csromanfolio{344}ตุโม กริสฺสติ, ตุโมว เตน ปญฺญายิสฺสตี”ติ. เอวรูปํ โข โมคฺคลฺลาน สตฺถารํ สาวกา อาชีวโต รกฺขนฺติ, เอวรูโป จ ปน สตฺถา สาวเกหิ อาชีวโต รกฺขํ ปจฺจาสีสติ.}

\prose{ปุน จปรํ โมคฺคลฺลาน อิเธกจฺโจ สตฺถา อปริสุทฺธธมฺมเทสโน สมาโน “ปริสุทฺธธมฺมเทสโนมฺหี”ติ ปฏิชานาติ “ปริสุทฺธา เม ธมฺมเทสนา ปริโยทาตา อสํกิลิฏฺฐา”ติ จ. ตเมนํ สาวกา เอวํ ชานนฺติ “อยํ โข ภวํ สตฺถา อปริสุทฺธธมฺมเทสโน สมาโน ‘ปริสุทฺธธมฺมเทสโนมฺหี’ติ ปฏิชานาติ ‘ปริสุทฺธา เม ธมฺมเทสนา ปริโยทาตา อสํกิลิฏฺฐา’ติ จ. มยญฺเจว โข ปน คิหีนํ อาโรเจยฺยาม ‘นาสฺสสฺส มนาปํ, ยํ โข ปนสฺส อมนาปํ. กถํ นํ มยํ เตน สมุทาจเรยฺยาม. สมฺมนฺนติ โข ปน จีวร\-ปิณฺฑปาต\-เสนาสน\-คิลานปฺปจฺจย\-เภสชฺช\-ปริกฺขาเรน. ยํ ตุโม กริสฺสติ, ตุโมว เตน ปญฺญายิสฺสาตี”ติ. เอวรูปํ โข โมคฺคลฺลาน สตฺถารํ สาวกา ธมฺมเทสนโต รกฺขนฺติ, เอวรูโป จ ปน สตฺถา สาวเกหิ ธมฺมเทสนโต รกฺขํ ปจฺจาสีสติ.}

\prose{ปุน จปรํ โมคฺคลฺลาน อิเธกจฺโจ สตฺถา อปริสุทฺธ\-เวยฺยากรโณ สมาโน “ปริสุทฺธเวยฺยากรโณมฺหี”ติ ปฏิชานาติ “ปริสุทฺธํ เม เวยฺยากรณํ ปริโยทาตํ อสํกิลิฏฺฐนฺ”ติ จ. ตเมนํ สาวกา เอวํ ชานนฺติ “อยํ โข ภวํ สตฺถา อปริสุทฺธ\-เวยฺยากรโณ สมาโน ‘ปริสุทฺธเวยฺยากรโณมฺหี’ติ ปฏิชานาติ ‘ปริสุทฺธํ เม เวยฺยากรณํ ปริโยทาตํ อสํกิลิฏฺฐนฺ’ติ จ. มยญฺเจว โข ปน คิหีนํ อาโรเจยฺยาม ‘นาสฺสสฺส มนาปํ, ยํ โข ปนสฺส อมนาปํ, กถํ นํ มยํ เตน สมุทาจเรยฺยาม. สมฺมนฺนติ โข ปน จีวร\-ปิณฺฑปาต\-เสนาสน\-คิลานปฺปจฺจย\-เภสชฺช\-ปริกฺขาเรน. ยํ ตุโม กริสฺสติ, ตุโมว เตน ปญฺญายิสฺสตี”ติ. เอวรูปํ โข โมคฺคลฺลาน สตฺถารํ สาวกา เวยฺยากรณโต รกฺขนฺติ, เอวรูโป จ ปน สตฺถา สาวเกหิ เวยฺยากรณโต รกฺขํ ปจฺจาสีสติ.}

\prose{ปุน จปรํ โมคฺคลฺลาน อิเธกจฺโจ สตฺถา อปริสุทฺธ\-ญาณทสฺสโน สมาโน “ปริสุทฺธ\-ญาณทสฺสโนมฺหี”ติ ปฏิชานาติ “ปริสุทฺธํ เม ญาณทสฺสนํ ปริโยทาตํ อสํกิลิฏฺฐนฺ”ติ จ. ตเมนํ สาวกา เอวํ ชานนฺติ “อยํ โข ภวํ สตฺถา อปริสุทฺธ\-ญาณทสฺสโน สมาโน \csromanfolio{345}‘ปริสุทฺธ\-ญาณทสฺสโนมฺหี’ติ ปฏิชานาติ ‘ปริสุทฺธํ เม ญาณทสฺสนํ ปริโยทาตํ อสํกิลิฏฺฐนฺ’ติ จ. มยญฺเจว โข ปน คิหีนํ อาโรเจยฺยาม ‘นาสฺสสฺส มนาปํ, ยํ โข ปนสฺส อมนาปํ, กถํ นํ มยํ เตน สมุทาจเรยฺยาม. สมฺมนฺนติ โข ปน จีวร\-ปิณฺฑปาต\-เสนาสน\-คิลานปฺปจฺจย\-เภสชฺช\-ปริกฺขาเรน. ยํ ตุโม กริสฺสติ, ตุโมว เตน ปญฺญายิสฺสตี”ติ. เอวรูปํ โข โมคฺคลฺลาน สตฺถารํ สาวกา ญาณทสฺสนโต รกฺขนฺติ, เอวรูโป จ ปน สตฺถา สาวเกหิ ญาณทสฺสนโต รกฺขํ ปจฺจาสีสตีติ. อิเม โข โมคฺคลฺลาน ปญฺจ สตฺถาโร สนฺโต สํวิชฺชมานา โลกสฺมิํ.}

\prose{อหํ โข ปน โมคฺคลฺลาน ปริสุทฺธสีโล สมาโน “ปริสุทฺธสีโลมฺหี”ติ ปฏิชานามิ “ปริสุทฺธํ เม สีลํ ปริโยทาตํ อสํกิลิฏฺฐนฺ”ติ จ, น จ มํ สาวกา สีลโต รกฺขนฺติ, น จาหํ สาวเกหิ สีลโต รกฺขํ ปจฺจาสีสามิ. ปริสุทฺธาชีโว สมาโน ฯเปฯ. ปริสุทฺธธมฺมเทสโน สมาโน ฯเปฯ. ปริสุทฺธ\-เวยฺยากรโณ สมาโน ฯเปฯ. ปริสุทฺธ\-ญาณทสฺสโน สมาโน “ปริสุทฺธ\-ญาณทสฺสโนมฺหี”ติ ปฏิชานามิ “ปริสุทฺธํ เม ญาณทสฺสนํ ปริโยทาตํ อสํกิลิฏฺฐนฺ”ติ จ. น จ มํ สาวกา ญาณทสฺสนโต รกฺขนฺติ, น จาหํ สาวเกหิ ญาณทสฺสนโต รกฺขํ ปจฺจาสีสามีติ.}

\proseitem{335}{อถ โข ภควา โกสมฺพิยํ ยถาภิรนฺตํ วิหริตฺวา เยน ราชคหํ เตน จาริกํ ปกฺกามิ, อนุปุพฺเพน จาริกํ จรมาโน เยน ราชคหํ ตทวสริ, ตตฺร สุทํ ภควา ราชคเห วิหรติ เวฬุวเน กลนฺทก\-นิวาเป. อถ โข สมฺพหุลา ภิกฺขู เยน ภควา เตนุ\-ปสงฺกมิํสุ, อุปสงฺกมิตฺวา ภควนฺตํ อภิวาเทตฺวา เอกมนฺตํ นิสีทิํสุ. เอกมนฺตํ นิสินฺนา โข เต ภิกฺขู ภควนฺตํ เอตทโวจุํ “เทวทตฺตสฺส ภนฺเต อชาตสตฺตุ กุมาโร ปญฺจหิ รถสเตหิ สายํ ปาตํ อุปฏฺฐานํ คจฺฉติ, ปญฺจ จ ถาลิปาก\-สตานิ ภตฺตาภิหาโร อภิหรียตี”ติ. มา ภิกฺขเว เทวทตฺตสฺส ลาภ\-สกฺการ\-สิโลกํ ปิหยิตฺถ, ยาวกีวญฺจ ภิกฺขเว เทวทตฺตสฺส อชาตสตฺตุ กุมาโร ปญฺจหิ รถสเตหิ สายํ ปาตํ อุปฏฺฐานํ คมิสฺสติ, ปญฺจ จ ถาลิปาก\-สตานิ ภตฺตาภิหาโร อภิหรียิสฺสติ, หานิเยว ภิกฺขเว เทวทตฺตสฺส ปาฏิกงฺขา กุสเลสุ ธมฺเมสุ, โน วุฑฺฒิ.}

\prose{\csromanfolio{346}เสยฺยถาปิ ภิกฺขเว จณฺฑสฺส กุกฺกุรสฺส นาสาย ปิตฺตํ ภินฺเทยฺยุํ, เอวญฺหิ โส ภิกฺขเว กุกฺกุโร ภิยฺโยโส\-มตฺตาย จณฺฑตโร อสฺส, เอวเมว โข ภิกฺขเว ยาวกีวญฺจ เทวทตฺตสฺส อชาตสตฺตุ กุมาโร ปญฺจหิ รถสเตหิ สายํ ปาตํ อุปฏฺฐานํ คมิสฺสติ, ปญฺจ จ ถาลิปาก\-สตานิ ภตฺตาภิหาโร อภิหรียิสฺสติ. หานิเยว ภิกฺขเว เทวทตฺตสฺส ปาฏิกงฺขา กุสเลสุ ธมฺเมสุ, โน วุฑฺฒิ.}

\prose{\csromansymbolfootnote{*}{สํ 1. 436; อํ 1. 385 ปิฏฺเฐสุปิ.}อตฺตวธาย ภิกฺขเว เทวทตฺตส ลาภ\-สกฺการ\-สิโลโก อุทปาทิ, ปราภวาย เทวทตฺตสฺส ลาภ\-สกฺการ\-สิโลโก อุทปาทิ.}

\prose{เสยฺยถาปิ ภิกฺขเว กทลี อตฺตวธาย ผลํ เทติ, ปราภวาย ผลํ เทติ, เอวเมว โข ภิกฺขเว อตฺตวธาย เทวทตฺตสฺส ลาภ\-สกฺการ\-สิโลโก อุทปาทิ, ปราภวาย เทวทตฺตสฺส ลาภ\-สกฺการ\-สิโลโก อุทปาทิ.}

\prose{เสยฺยถาปิ ภิกฺขเว เวฬุ อตฺตวธาย ผลํ เทติ, ปราภวาย ผลํ เทติ, เอวเมว โข ภิกฺขเว อตฺตวธาย เทวทตฺตสฺส ลาภ\-สกฺการ\-สิโลโก อุทปาทิ, ปราภวาย เทวทตฺตสฺส ลาภ\-สกฺการ\-สิโลโก อุทปาทิ.}

\prose{เสยฺยถาปิ ภิกฺขเว นโฬ อตฺถวธาย ผลํ เทติ, ปราภวาย ผลํ เทติ, เอวเมว โข ภิกฺขเว อตฺตวธาย เทวทตฺตสฺส ลาภ\-สกฺการ\-สิโลโก อุทปาทิ, ปราภวาย เทวทตฺตสฺส ลาภ\-สกฺการ\-สิโลโก อุทปาทิ.}

\prose{เสยฺยถาปิ ภิกฺขเว อสฺสตรี อตฺตวธาย คพฺภํ คณฺหาติ, ปราภวาย คพฺภํ คณฺหาติ, เอวเมว โข ภิกฺขเว อตฺตวธาย เทวทตฺตสฺส ลาภ\-สกฺการ\-สิโลโก อุทปาทิ, ปราภวาย เทวทตฺตสฺส ลาภ\-สกฺการ\-สิโลโก อุทปาทีติ.}

\setlength{\csromangathapagewidth}{0pt}
\setlength{\csromangathawakwidth}{0pt}
\csromangathameasuregroup{\textsuperscript{+}ผลํ เว กทลิํ หนฺติ, \\ สกฺกาโร กาปุริสํ หนฺติ,}{\csromangathabat{\textsuperscript{+}ผลํ เว กทลิํ หนฺติ,}{ผลํ เวฬุํ ผลํ นฬํ.} \\ \csromangathabat{สกฺกาโร กาปุริสํ หนฺติ,}{คพฺโภ อสฺสตริํ ยถาติ.}}
\csromangathasetleft{\textsuperscript{+}ผลํ เว กทลิํ หนฺติ, \\ สกฺกาโร กาปุริสํ หนฺติ,}
\csromangathagroupwithcloser{\csromangathastanza{\csromangathabat{\csromansymbolfootnote{+}{สํ 1. 156; ขุ 10. 109 ปิฏฺเฐสุปิ.}ผลํ เว กทลิํ หนฺติ,}{ผลํ เวฬุํ ผลํ นฬํ.} \\ \csromangathabat{สกฺกาโร กาปุริสํ หนฺติ,}{คพฺโภ อสฺสตริํ ยถาติ.}}}{ปฐม\-ภาณวาโร นิฏฺฐิโต.}
\csromanmatikaanchor{mtk.165}
\csromantocmarkref{section}{2. ทุติยภาณวาร}{mtk.165}
\bookmark[dest={mtk.165},level=1]{2. ทุติยภาณวาร}
\csromanheader{1}{2. ทุติย\-ภาณวาร}
\csromanmatikaanchor{mtk.166}
\csromantocmarkref{subsection}{ปกาสนียกมฺม}{mtk.166}
\bookmark[dest={mtk.166},level=2]{ปกาสนียกมฺม}
\csromanheader{2}{ปกาสนีย\-กมฺม}
\proseitem{336}{\csromanfolio{347}เตน โข ปน สมเยน ภควา มหติยา ปริสาย ปริวุโต ธมฺมํ เทเสนฺโต นิสินฺโน โหติ สราชิกาย ปริสาย. อถ โข เทวทตฺโต อุฏฺฐายาสนา เอกํสํ อุตฺตราสงฺคํ กริตฺวา เยน ภควา เตนญฺชลิํ ปณาเมตฺวา ภควนฺตํ เอตทโวจ “ชิณฺโณ ทานิ ภนฺเต ภควา วุฑฺโฒ มหลฺลโก อทฺธคโต วโยอนุปฺปตฺโต, อปฺโปสฺสุกฺโก ทานิ ภนฺเต ภควา ทิฏฺฐ\-ธมฺม\-สุข\-วิหารํ อนุยุตฺโต วิหรตุ, มมํ ภิกฺขุสํฆํ นิสฺสชฺชตุ, อหํ ภิกฺขุสํฆํ ปริหริสฺสามี”ติ. อลํ เทวทตฺต, มา เต รุจฺจิ ภิกฺขุสํฆํ ปริหริตุนฺติ. ทุติยมฺปิ โข เทวทตฺโต ฯเปฯ. ตติยมฺปิ โข เทวทตฺโต ภควนฺตํ เอตทโวจ “ชิณฺโณ ทานิ ภนฺเต ภควา วุฑฺโฒ มหลฺลโก อทฺธคโต วโยอนุปฺปตฺโต, อปฺโปสฺสุกฺโก ทานิ ภนฺเต ภควา ทิฏฺฐ\-ธมฺม\-สุข\-วิหารํ อนุยุตฺโต วิหรตุ, มมํ ภิกฺขุสํฆํ นิสฺสชฺชตุ, อหํ ภิกฺขุสํฆํ ปริหริสฺสามี”ติ. สาริปุตฺตโมคฺคลฺลนานมฺปิ โข อหํ เทวทตฺต ภิกฺขุสํฆํ น นิสฺสชฺเชยฺยํ, กิํ ปน ตุยฺหํ ฉวสฺส เขฬาสกสฺสาติ. อถ โข เทวทตฺโต “สราชิกาย มํ ภควา ปริสาย เขฬาสก\-วาเทน อปสาเทติ, สาริปุตฺต\-โมคฺคลฺลาเนว อุกฺกํสตี”ติ กุปิโต อนตฺตมโน ภควนฺตํ อภิวาเทตฺวา ปทกฺขิณํ กตฺวา ปกฺกามิ. อยญฺจรหิ เทวทตฺตสฺส ภควติ ปฐโม อาฆาโต อโหสิ.}

\prose{อถ โข ภควา ภิกฺขู อามนฺเตสิ\csromandash{} เตน หิ ภิกฺขเว สํโฆ เทวทตฺตสฺส ราชคเห ปกาสนียํ กมฺมํ กโรตุ “ปุพฺเพ เทวทตฺตสฺส อญฺญา ปกติ อโหสิ, อิทานิ อญฺญา ปกติ, ยํ เทวทตฺโต กเรยฺย กาเยน วาจาย, น เตน พุทฺโธ วา ธมฺโม วา สํโฆ วา ทฏฺฐพฺโพ, เทวทตฺโตว เตน ทฏฺฐพฺโพ”ติ. เอวญฺจ ปน ภิกฺขเว กาตพฺพํ, พฺยตฺเตน ภิกฺขุนา ปฏิพเลน สํโฆ ญาเปตพฺโพ\csromandash{}}

\proseitem{337}{“สุณาตุ เม ภนฺเต สํโฆ, ยทิ สํฆสฺส ปตฺตกลฺลํ, สํโฆ เทวทตฺตสฺส ราชคเห ปกาสนียํ กมฺมํ กเรยฺย ‘ปุพฺเพ เทวทตฺถสฺส อญฺญา ปกติ อโหสิ, อิทานิ อญฺญา ปกติ, ยํ เทวทตฺโต กเรยฺย \csromanfolio{348}กาเยน วาจายาน เตน พุทฺโธ วา ธมฺโม วา สํโฆ วา ทฏฺฐพฺโพ, เทวทตฺโตว เตน ทฏฺฐพฺโพ’ติ, เอสา ญตฺติ.}

\prose{สุณาตุ เม ภนฺเต สํโฆ, สํโฆ เทวทตฺตสฺส ราชคเห ปกาสนียํ กมฺมํ กโรติ ‘ปุพฺเพ เทวทตฺตสฺส อญฺญา ปกติ อโหสิ, อิทานิ อญฺญา ปกติ, ยํ เทวทตฺโต กเรยฺย กาเยน วาจาย, น เตน พุทฺโธ วา ธมฺโม วา สํโฆ วา ทฏฺฐพฺโพ, เทวทตฺโตว เตน ทฏฺฐพฺโพ’ติ, ยสฺสายสฺมโต ขมติ เทวทตฺตสฺส ราชคเห ปกาสนียสฺส กมฺมสฺส กรณํ ‘ปุพฺเพ เทวทตฺตสฺส อญฺญา ปกติ อโหสิ, อิทานิ อญฺญา ปกติ, ยํ เทวทตฺโต กเรยฺย กาเยน วาจาย, น เตน พุทฺโธ วา ธมฺโม วา สํโฆ วา ทฏฺฐพฺโพ, เทวทตฺโตว เตน ทฏฺฐพฺโพ’ติ, โส ตุณฺหสฺส, ยสฺส นกฺขมติ, โส ภาเสยฺย.}

\prose{กถํ สํเฆน เทวทตฺตสฺส ราชคเห ปกาสนียํ กมฺมํ ‘ปุพฺเพ เทวทตฺตสฺส อญฺญา ปกติ อโหสิ, อิทานิ อญฺญา ปกติ, ยํ เทวทตฺโต กเรยฺย กาเยน วาจาย, น เตน พุทฺโธ วา ธมฺโม วา สํโฆ วา ทฏฺฐพฺโพ, เทวทตฺโตว เตน ทฏฺฐพฺโพ’ติ, ขมติ สํฆสฺส, ตสฺมา ตุณฺหี, เอวเมตํ ธารยามี”ติ.}

\proseitem{338}{อถ โข ภควา อายสฺมนฺตํ สาริปุตฺตํ อามนฺเตสิ “เตน หิ ตฺวํ สาริปุตฺต เทวทตฺตํ ราชคเห ปกาเสหี”ติ. ปุพฺเพ มยา ภนฺเต เทวทตฺตสฺส ราชคเห วณฺโณ ภาสิโต “มหิทฺธิโก โคธิปุตฺโต, มหานุภาโว โคธิปุตฺโต”ติ. กถาหํ ภนฺเต เทวทตฺตํ ราชคเห ปกาเสมีติ. นนุ ตยา สาริปุตฺต ภูโตเยว เทวทตฺตสฺส ราชคเห วณฺโณ ภาสิโต “มหิทฺธิโก โคธิปุตฺโต, มหานุภาโว โคธิปุตฺโต”ติ. เอวํ ภนฺเตติ. เอวเมว โข ตฺวํ สาริปุตฺต ภูตํเยว เทวทตฺตํ ราชคเห ปกาเสหีติ. “เอวํ ภนฺเต”ติ โข อายสฺมา สาริปุตฺโต ภควโต ปจฺจสฺโสสิ.}

\prose{อถ โข ภควา ภิกฺขู อามนฺเตสิ “เตน หิ ภิกฺขเว สํโฆ สาริปุตฺตํ สมฺมนฺนตุ เทวทตฺตํ ราชคเห ปกาเสตุํ ‘ปุพฺเพ เทวทตฺตสฺส อญฺญา ปกติ อโหสิ, อิทานิ อญฺญา ปกติ, ยํ เทวทตฺโต กเรยฺย กาเยน วาจาย, น เตน พุทฺโธ วา ธมฺโม วา สํโฆ วา ทฏฺฐพฺโพ, เทวทตฺโตว เตน ทฏฺฐพฺโพ’ติ”. เอวญฺจ ปน ภิกฺขเว สมฺมนฺนิตพฺโพ, ปฐมํ \csromanfolio{349}สาริปุตฺโต ยาจิตพฺโพ, ยาจิตฺวา พฺยตฺเตน ภิกฺขุนา ปฏิพเลน สํโฆ ญาเปตพฺโพ\csromandash{}}

\prose{“สุณาตุ เม ภนฺเต สํโฆ, ยทิ สํฆสฺส ปตฺตกลฺลํ, สํโฆ อายสฺมนฺตํ สาริปุตฺตํ สมฺมนฺเนยฺย เทวทตฺตํ ราชคเห ปกาเสตุํ ‘ปุพฺเพ เทวทตฺตสฺส อญฺญา ปกติ อโหสิ, อิทานิ อญฺญา ปกติ, ยํ เทวทตฺโต กเรยฺย กาเยน วาจาย, น เตน พุทฺโธ วา ธมฺโม วา สํโฆ วา ทฏฺฐพฺโพ, เทวทตฺโตว เตน ทฏฺฐพฺโพ’ติ, เอสา ญตฺติ.}

\prose{สุณาตุ เม ภนฺเต สํโฆ, สํโฆ อายสฺมนฺตํ สาริปุตฺตํ สมฺมนฺนติ เทวทตฺตํ ราชคเห ปกาเสตุํ ‘ปุพฺเพ เทวทตฺตส อญฺญา ปกติ อโหสิ, อิทานิ อญฺญา ปกติ, ยํ เทวทตฺโต กเรยฺย กาเยน วาจาย, น เตน พุทฺโธ วา ธมฺโม วา สํโฆ วา ทฏฺฐพฺโพ, เทวทตฺโตว เตน ทฏฺฐพฺโพ’ติ, ยสฺสายสฺมโต ขมติ อายสฺมโต สาริปุตฺตสฺส สมฺมุติ เทวทตฺตํ ราชคเห ปกาเสตุํ ‘ปุพฺเพ เทวทตฺตสฺส อญฺญา ปกติ อโหสิ, อิทานิ อญฺญา ปกติ, ยํ เทวทตฺโต กเรยฺย กาเยน วาจาย, น เตน พุทฺโธ วา ธมฺโม วา สํโฆ วา ทฏฺฐพฺโพ, เทวทตฺโตว เตน ทฏฺฐพฺโพ’ติ. โส ตุณฺหสฺส, ยสฺส นกฺขมติ, โส ภาเสยฺย.}

\prose{สมฺมโต สํเฆน อายสฺมา สาริปุตฺโต เทวทตฺตํ ราชคเห ปกาเสตุํ ‘ปุพฺเพ เทวทตฺตสฺส อญฺญา ปกติ อโหสิ, อิทานิ อญฺญา ปกติ, ยา เทวทตฺโต กเรยฺย กาเยน วาจาย, น เตน พุทฺโธ วา ธมฺโม วา สํโฆ วา ทฏฺฐพฺโพ, เทวทตฺโตว เตน ทฏฺฐพฺโพ’ติ, ขมติ สํฆสฺส, ตสฺมา ตุณฺหี, เอวเมตํ ธารยามี”ติ.}

\prose{สมฺมโต จ อายสฺมา สาริปุตฺโต สมฺพหุเลหิ ภิกฺขูหิ สทฺธิํ ราชคหํ ปวิสิตฺวา เทวทตฺตํ ราชคเห ปกาเสสิ “ปุพฺเพ เทวทตฺตสฺส อญฺญา ปกติ อโหสิ, อิทานิ อญฺญา ปกติ, ยํ เทวทตฺโต กเรยฺย กาเยน วาจาย, น เตน พุทฺโธ วา ธมฺโม วา สํโฆ วา ทฏฺฐพฺโพ, เทวทตฺโตว เตน ทฏฺฐพฺโพ”ติ. ตตฺถ เย เต มนุสฺสา อสฺสทฺธา อปฺปสนฺนา ทุพฺพุทฺธิโน, เต เอวมาหํสุ “อุสูยกา อิเม สมณา สกฺยปุตฺติยา เทวทตฺตสฺส ลาภสกฺการํ อุสูยนฺตี”ติ. เย ปน เต มนุสฺสา สทฺธา ปสนฺนา ปณฺฑิตา พฺยตฺตา พุทฺธิมนฺโต, เต เอวมาหํสุ “น โข อิทํ โอรกํ ภวิสฺสติ, ยถา ภควา เทวทตฺตํ ราชคเห ปกาสาเปตี”ติ.}

\csromanmatikaanchor{mtk.167}
\csromantocmarkref{subsection}{อชาตสตฺตุกุมารวตฺถุ}{mtk.167}
\bookmark[dest={mtk.167},level=2]{อชาตสตฺตุกุมารวตฺถุ}
\csromanheader{2}{อชาตสตฺตุกุมาร\-วตฺถุ}
\proseitem{339}{\csromanfolio{350}อถ โข เทวทตฺโต เยน อชาตสตฺตุ กุมาโร เตนุปสงฺกมิ, อุปสงฺกมิตฺวา อชาตสตฺตุํ กุมารํ เอตทโวจ “ปุพฺเพ โข กุมาร มนุสฺสา ทีฆายุกา, เอตรหิ อปฺปายุกา. ฐานํ โข ปเนตํ วิชฺชติ, ยํ ตฺวํ กุมาโรว สมาโน กาลํ กเรยฺยาสิ, เตน หิ ตฺวํ กุมาร ปิตรํ หนฺตฺวา ราชา โหติ, อหํ ภควนฺตํ หนฺตฺวา พุทฺโธ ภวิสฺสามี”ติ.}

\prose{อถ โข อชาตสตฺตุ กุมาโร “อยฺโย โข เทวทตฺโต มหิทฺธิโก มหานุภาโว, ชาเนยฺยาสิ อยฺโย เทวทตฺโต”ติ อูรุยา โปตฺถนิกํ พนฺธิตฺวา ทิวา ทิวสฺส\csromansharedfootnote{ca96bb623b17e9b8}{ทิวา ทิวสสฺส (ก)} ภีโต อุพฺพิคฺโค อุสฺสงฺกี อุตฺรสฺโต สหสา อนฺเตปุรํ ปาวิสิ. อทฺทสาสุํ โข อนฺเตปุเร อุปจารกา มหามตฺตา อชาตสตฺตุํ กุมารํ ทิวา ทิวสฺส ภีตํ อุพฺพิคฺคํ อุสฺสงฺกิํ อุตฺรสฺตํ สหสา อนฺเตปุรํ ปวิสนฺตํ, ทิสฺวาน อคฺคเหสุํ. เต วิจินนฺตา อูรุยา โปตฺถนิกํ พทฺธํ\csromansharedfootnote{7b3663c35a2ca658}{พนฺธํ (ก)} ทิสฺวาน อชาตสตฺตุํ กุมารํ เอตทโวจุํ “กิํ ตฺวํ กุมาร กตฺตุกาโมสี”ติ. ปิตรมฺหี หนฺตุกาโมติ. เกนาสิ อุสฺสาหิโตติ. อยฺเยน เทวทตฺเตนาติ. เอกจฺเจ มหามตฺตา เอวํ มติํ อกํสุ “กุมาโร จ หนฺตพฺโพ เทวทตฺโต จ, สพฺเพ จ ภิกฺขู หนฺตพฺพา”ติ. เอกจฺเจ มหามตฺตา เอวํ มติํ อกํสุ “น ภิกฺขู หนฺตพฺพา, น ภิกฺขู กิญฺจิ อปรชฺฌนฺติ, กุมาโร จ หนฺตพฺโพ เทวทตฺโต จา”ติ. เอกจฺเจ มหามตฺตา เอวํ มติํ อกํสุ “น กุมาโร จ หนฺตพฺโพ น เทวทตฺโต, น ภิกฺขู หนฺตพฺพา, รญฺโญ อาโรเจตพฺพํ, ยถา ราชา วกฺขติ ตถา กริสฺสามา”ติ.}

\prose{อถ โข เต มหามตฺตา อชาตสตฺตุํ กุมารํ อาทาย เยน ราชา มาคโธ เสนิโย พิมฺพิสาโร เตนุ\-ปสงฺกมิํสุ, อุปสงฺกมิตฺวา รญฺโญ มาคธสฺส เสนิยสฺส พิมฺพิสารสฺส เอตมตฺถํ อาโรเจสุํ. กถํ ภเณ มหามตฺเตหิ มติ กตาติ. เอกจฺเจ เทว มหามตฺตา เอวํ มติํ อกํสุ “กุมาโร จ หนฺตพฺโพ เทวทตฺโต จ, สพฺเพ จ ภิกฺขู หนฺตพฺพา”ติ. เอกจฺเจ มหามตฺตา เอวํ มติํ อกํสุ “น ภิกฺขู หนฺตพฺพา, น ภิกฺขู กิญฺจิ อปรชฺฌนฺติ, กุมาโร จ หนฺตพฺโพ เทวทตฺโต จา”ติ. เอกจฺเจ มหามตฺตา เอวํ มติํ อกํสุ “น กุมาโร จ หนฺตพฺโพ, น เทวทตฺโต, น ภิกฺขู หนฺตพฺพา, \csromanfolio{351}รญฺโญ อาโรเจตพฺพํ, ยถา ราชา วกฺขติ ตถา กริสฺสามา”ติ. กิํ ภเณ กริสฺสติ พุทฺโธ วา ธมฺโม วา สํโฆ วา, นนุ ภควตา ปฏิกจฺเจว เทวทตฺโต ราชคเห ปกาสาปิโต “ปุพฺเพ เทวทตฺตสฺส อญฺญา ปกติ อโหสิ, อิทานิ อญฺญา ปกติ, ยํ เทวทตฺโต กเรยฺย กาเยน วาจาย, น เตน พุทฺโธ วา ธมฺโม วา สํโฆ วา ทฏฺฐพฺโพ, เทวทตฺโตว เตน ทฏฺฐพฺโพ”ติ. ตตฺถ เย เต มหามตฺตา เอวํ มติํ อกํสุ “กุมาโร จ หนฺตพฺโพ เทวทตฺโต จ, สพฺเพ จ ภิกฺขู หนฺตพฺพา”ติ. เต อฏฺฐาเน อกาสิ. เย เต มหามตฺตา เอวํ มติํ อกํสุ “น ภิกฺขู หนฺตพฺพา, น ภิกฺขู กิญฺจิ อปรชฺฌนฺติ, กุมาโร จ หนฺตพฺโพ เทวทตฺโต จา”ติ. เต นีเจ ฐาเน ฐเปสิ. เย เต มหามตฺตา เอวํ มติํ อกํสุ “น กุมาโร จ หนฺตพฺโพ น เทวทตฺโต, น ภิกฺขู หนฺตพฺพา, รญฺโญ อาโรเจตพฺพํ, ยถา ราชา วกฺขติ ตถา กริสฺสามา”ติ. เต อุจฺเจ ฐาเน ฐเปสิ. อถ โข ราชา มาคโธ เสนิโย พิมฺพิสาโร อชาตสตฺตุํ กุมารํ เอตทโวจ “กิสฺส มํ ตฺวํ กุมาร หนฺตุกาโมสี”ติ. รชฺเชนามฺหิ เทว อตฺถิโกติ. สเจ โข ตฺวํ กุมาร รชฺเชน อตฺถิโก, “เอตํ เต รชฺชนฺ”ติ อชาตสตฺตุสฺส กุมารสฺส รชฺชํ นิยฺยาเทสิ.}

\csromanmatikaanchor{mtk.168}
\csromantocmarkref{subsection}{อภิมารเปสน}{mtk.168}
\bookmark[dest={mtk.168},level=2]{อภิมารเปสน}
\csromanheader{2}{\textbf{อภิมารเปสน}}
\proseitem{340}{อถ โข เทวทตฺโต เยน อชาตสตฺตุ กุมาโร เตนุปสงฺกมิ, อุปสงฺกมิตฺวา อชาตสตฺตุํ กามารํ เอตทโวจ “ปุริเส มหาราช อาณาเปหิ, เย สมณํ โคตมํ ชีวิตา โวโรเปสฺสนฺตี”ติ. อถ โข อชาตสตฺตุ กุมาโร มนุสฺเส อาณาเปสิ “ยถา ภเณ อยฺโย เทวทตฺโต อาห, ตถา กเรยฺยาถา”ติ. อถ โข เทวทตฺโต เอกํ ปุริสํ อาณาเปสิ “คจฺฉาวุโส, อมุกสฺมิํ โอกาเส สมโณ โคตโม วิหรติ, ตํ ชีวิตา โวโรเปตฺวา อิมินา มคฺเคน อาคจฺฉา”ติ. ตสฺมิํ มคฺเค เทฺว ปุริเส ฐเปสิ “โย อิมินา มคฺเคน เอโก ปุริโส อาคจฺฉติ, ตํ ชีวิตา โวโรเปตฺวา อิมินา มคฺเคน อาคจฺฉถา”ติ. ตสฺมิํ มคฺเค จตฺตาโร ปุริเส ฐเปสิ “เย อิมินา มคฺเคน เทฺว ปุริสา อาคจฺฉนฺติ, เต ชีวิตา โวโรเปตฺวา อิมินา มคฺเคน อาคจฺฉถา”ติ. ตสฺมิํ มคฺเค อฏฺฐ ปุริเส ฐเปสิ “เย อิมินา มคฺเคน จตฺตาโร ปุริสา อาคจฺฉนฺติ, เต ชีวิตา โวโรเปตฺวา อิมินา \csromanfolio{352}มคฺเคน อาคจฺฉถา”ติ. ตสฺมิํ มคฺเค โสฬส ปุริเส ฐเปสิ “เย อิมินา มคฺเคน อฏฺฐ ปุริสา อาคจฺฉนฺติ, เต ชีวิตา โวโรเปตฺวา อาคจฺฉถา”ติ.}

\prose{อถ โข โส เอโก ปุริโส อสิจมฺมํ คเหตฺวา ธนุกลาปํ สนฺนยฺหิตฺวา เยน ภควา เตนุปสงฺกมิ, อุปสงฺกมิตฺวา ภควโต อวิทูเร ภีโต อุพฺพิคฺโค อุสฺสงฺกี อุตฺรสฺโต ปตฺถทฺเธน กาเยน อฏฺฐาสิ. อทฺทสา โข ภควา ตํ ปุริสํ ภีตํ อุพฺพิคฺคํ อุสฺสงฺกิํ อุตฺรสฺตํ ปตฺถทฺเธน กาเยน ฐิตํ, ทิสฺวาน ตํ ปุริสํ เอตทโวจ “เอหาวุโส, มา ภายี”ติ. อถ โข โส ปุริโส อสิจมฺมํ เอกมนฺตํ กริตฺวา ธนุกลาปํ นิกฺขิปิตฺวา เยน ภควา เตนุปสงฺกมิ, อุปสงฺกมิตฺวา ภควโต ปาเทสุ สิรสา นิปติตฺวา ภควนฺตํ เอตทโวจ “อจฺจโย มํ ภนฺเต อจฺจคมา ยถาพาลํ ยถามูฬฺหํ ยถาอกุสลํ, โยหํ ทุฏฺฐจิตฺโต วธกจิตฺโต อิธูปสงฺกนฺโต, ตสฺส เม ภนฺเต ภควา อจฺจยํ อจฺจยโต ปฏิคฺคณฺหาตุ อายติํ สํวรายา”ติ. ตคฺฆ ตฺวํ อาวุโส อจฺจโย อจฺจคมา ยถาพาลํ ยถามูฬฺหํ ยถาอกุสลํ, ยํ ตฺวํ ทุฏฺฐจิตฺโต วธกจิตฺโต อิธูปสงฺกนฺโต, ยโต จ โข ตฺวํ อาวุโส อจฺจยํ อจฺจยโต ทิสฺวา ยถาธมฺมํ ปฏิกโรสิ, ตํ เต มยํ ปฏิคฺคณฺหาม, วุทฺธิ เหสา อาวุโส อริยสฺส วินเย, โย อจฺจยํ อจฺจยโต ทิสฺวา ยถาธมฺมํ ปฏิกโรติ, อายติํ สํวรํ อาปชฺชตีติ.}

\prose{อถ โข ภควา ตสฺส ปุริสสฺส อนุปุพฺพิํ กถํ กถสิ, เสยฺยถิทํ ทานกถํ สีลกถํ สคฺคกถํ กามานํ อาทีนวํ โอการํ สํกิเลสํ เนกฺขมฺเม อานิสํสํ ปกาเสสิ. ยทา ตํ ภควา อญฺญาสิ กลฺลจิตฺตํ มุทุจิตฺตํ วินีวรณ\-จิตฺตํ อุทคฺคจิตฺตํ ปสนฺนจิตฺตํ, อถ ยา พุทฺธานํ สามุกฺกํสิกา ธมฺมเทสนา, ตํ ปกาเสสิ ทุกฺขํ สมุทยํ นิโรธํ มคฺคํ, เสยฺยถาปิ นาม สุทฺธํ วตฺถํ อปคตกาฬกํ สมฺมเทว รชนํ ปฏิคฺคเหยฺย, เอวเมว ตสฺส ปุริสสฺส ตสฺมิํเยว อาสเน วิรชํ วีตมลํ ธมฺมจกฺขุํ อุทปาทิ “ยํ กิญฺจิ สมุทย\-ธมฺมํ, สพฺพํ ตํ นิโรธธมฺมนฺ”ติ. อถ โข โส ปุริโส ทิฏฺฐธมฺโม ปตฺตธมฺโม วิทิตธมฺโม ปริโยคาฬฺห\-ธมฺโม ติณฺณ\-วิจิกิจฺโฉ วิคต\-กถํกโถ เวสารชฺช\-ปฺปตฺโต อปรปฺปจฺจโย สตฺถุสาสเน ภควนฺตํ เอตทโวจ “อภิกฺกนฺตํ ภนฺเต, อภิกฺกนฺตํ ภนฺเต, เสยฺยถาปิ ภนฺเต นิกฺกุชฺชิตํ วา อุกฺกุชฺเชยฺย, ปฏิจฺฉนฺนํ วา วิวเรยฺย, \csromanfolio{353}มูฬสฺส วา มคฺคํ อาจิกฺเขยฺย, อนฺธกาเร วา เตลปชฺโชตํ ธาเรยฺย ‘จกฺขุมนฺโต รูปานิ ทกฺขนฺตี’ติ. เอวเมวํ ภควตา อเนก\-ปริยาเยน ธมฺโม ปกาสิโต. เอสาหํ ภนฺเต ภควนฺตํ สรณํ คจฺฉามิ ธมฺมญฺจ ภิกฺขุสํฆญฺจ, อุปาสกํ มํ ภควา ธาเรตุ อชฺชตคฺเค ปาณุเปตํ สรณํ คตนฺ”ติ. อถ โข ภควา ตํ ปุริสํ เอตทโวจ “มา โข ตฺวํ อาวุโส อิมินา มคฺเคน คจฺฉ, อิมินา มคฺเคน คจฺฉาหี”ติ, อญฺเญน มคฺเคน อุยฺโยเชสิ.}

\prose{อถ โข เต เทฺว ปุริสา “กิํ นุ โข โส เอโก ปุริโส จิเรน อาคจฺฉตี”ติ ปฏิปถํ คจฺฉนฺตา อทฺทสาสุํ ภควนฺตํ อญฺญตรสฺมิํ รุกฺขมูเล นิสินฺนํ, ทิสฺวาน เยน ภควา เตนุสงฺกมิํสุ, อุปสงฺกมิตฺวา ภควนฺตํ อภิวาเทตฺวา เอกมนฺตํ นิสีทิํสุ. เตสํ ภควา อนุปุพฺพิํ กถํ กเถสิ ฯเปฯ อปรปฺปจฺจยา สตฺถุสาสเน ภควนฺตํ เอตทโวจุํ “อภิกฺกนฺตํ ภนฺเต ฯเปฯ อุปาสเก โน ภควา ธาเรตุ อชฺชตคฺเค ปาณุเปตํ สรณํ คเต”ติ. อถ โข ภควา เต ปุริเส เอตทโวจ “มา โข ตุเมฺห อาวุโส อิมินา มคฺเคน คจฺฉิตฺถ, อิมินา มคฺเคน คจฺฉถา, อญฺเญน มคฺเคน อุยฺโยเชสิ.}

\prose{อถ โข เต จตฺตาโร ปุริสา ฯเปฯ. อถ โข เต อฏฺฐปุริสา ฯเปฯ. อถ โข เต โสฬส ปุริสา “กิํ นุ โข เต อฏฺฐ ปุริสา จิเรน อาคจฺฉนฺตี”ติ ปฏิปถํ คจฺฉนฺตา อทฺทสาสุํ ภควนฺตํ อญฺญตรสฺมิํ รุกฺขมูเล นิสินฺนํ, ทิสฺวาน เยน ภควา เตนุ\-ปสงฺกมิํสุ, อุปสงฺกมิตฺวา ภควนฺตํ อภิวาเทตฺวา เอกมนฺตํ นิสีทิํสุ. เตสํ ภควา อนุปุพฺพิํ กถํ กเถสิ, เสยฺยถิทํ, ทานกถํ ฯเปฯ อปรปฺปจฺจยา สตฺถุสาสเน ภควนฺตํ เอตทโวจุํ “อภิกฺกนฺตํ ภนฺเต ฯเปฯ อุปาสเก โน ภควา ธาเรตุ อชฺชตคฺเค ปาณุเปตํ สรณํ คเต”ติ. อถ โข ภควา เต ปุริเส เอตทโวจ “มา โข ตุเมฺห อาวุโส อิมินา มคฺเคน คจฺฉิตฺถ, อิมินา มคฺเคน คจฺฉถา”ติ, อญฺเญน มคฺเคน อุยฺโยเชสิ.}

\prose{อถ โข โส เอโก ปุริโส เยน เทวทตฺโต เตนุปสงฺกมิ, อุปสงฺกมิตฺวา เทวทตฺตํ เอตทโวจ “นาหํ ภนฺเต สกฺโกมิ ตํ ภควนฺตํ ชีวิตา โวโรเปตุํ, มหิทฺธิโก โส ภควา มหานุภาโว”ติ. อลํ อาวุโส, มา ตฺวํ สมณํ โคตมํ ชีวิตา โวโรเปสิ, อหเมว สมณํ โคตมํ ชีวิตา โวโรเปสฺสามีติ.}

\csromanmatikaanchor{mtk.169}
\csromantocmarkref{subsection}{โลหิตุปฺปาทกกมฺม}{mtk.169}
\bookmark[dest={mtk.169},level=2]{โลหิตุปฺปาทกกมฺม}
\csromanheader{2}{โลหิตุปฺปาทก\-กมฺม}
\proseitem{341}{\csromanfolio{354}เตน โข ปนสมเยน ภควา คิชฺฌกูฏสฺส ปพฺพตสฺส ฉายายํ จงฺกมติ. อถ โข เทวทตฺโต คิชฺฌกูฏํ ปพฺพตํ อารุหิตฺวา มหติํ สิลํ ปวิชฺฌิ “อิมาย สมณํ โคตมํ ชีวิตา โวโรเปสฺสามี”ติ. เทฺว ปพฺพตกูฏา สมาคนฺตฺวา ตํ สิลํ สมฺปฏิจฺฉิํสุ, ตโต ปปติกา อุปฺปติตฺวา ภควโต ปาเท รุหิรํ อุปฺปาเทสิ. อถ โข ภควา อุทฺธํ อุลฺโลเกตฺวา เทวทตฺตํ เอตทโวจ “พหุํ ตยา โมฆปุริส อปุญฺญํ ปสุตํ, ยํ ตฺวํ ทุฏฺฐจิตฺโต วธกจิตฺโต ตถาคตสฺส รุหิรํ อุปฺปาเทสี”ติ. อถ โข ภควา ภิกฺขู อามนฺเตสิ “อิทํ ภิกฺขเว เทวทตฺเตน ปฐมํ อานนฺตริยํ กมฺมํ อุปจิตํ, ยํทุฏฺฐจิตฺเตน วธกจิตฺเตน ตถาคตสฺส รุหิรํ อุปฺปาทิตนฺ”ติ.}

\prose{อสฺโสสุํ โข ภิกฺขู “เทวทตฺเตน กิร ภควโต วโธ ปยุตฺโต”ติ. เต จ ภิกฺขู ภควโต วิหารสฺส ปริโต ปริโต จงฺกมนฺติ อุจฺจาสทฺทา มหาสทฺทา สชฺฌายํ กโรนฺตา ภควโต รกฺขา\-วรณ\-คุตฺติยา. อสฺโสสิ โข ภควา อุจฺจาสทฺทํ มหาสทฺทํ สชฺฌายสทฺทํ, สุตฺวาน อายสฺมนฺตํ อานนฺทํ อามนฺเตสิ “กิํ นุ โข โส อานนฺท อุจฺจาสทฺโท มหาสทฺโท สชฺฌายสทฺโท”ติ. อสฺโสสุํ โข ภนฺเต ภิกฺขู “เทวทตฺเตน กิร ภควโต วโธ ปยุตฺโต”ติ. เต จ\csromansharedfootnote{c17595623c81b551}{เตธ (สี)} ภนฺเต ภิกฺขู ภควโต วิหารสฺส ปริโต ปริโต จงฺกมนฺติ อุจฺจาสทฺทา มหาสทฺทา สชฺฌายํ กโรนฺตา ภควโต รกฺขา\-วรณ\-คุตฺติยา, โส เอโส ภควา อุจฺจาสทฺโท มหาสทฺโท สชฺฌาย\-สทฺโทติ. เตน หานนฺท มม วจเนน เต ภิกฺขู อามนฺเตหิ “สตฺถา อายสฺมนฺเต อามนฺเตตี”ติ. “เอวํ ภนฺเต”ติ โข อายสฺมา อานนฺโท ภควโต ปฏิสฺสุตฺวา เยน เต ภิกฺขู เตนุปสงฺกมิ, อุปสงฺกมิตฺวา เต ภิกฺขู เอตทโวจ “สตฺถา อายสฺมนฺเต อามนฺเตตี”ติ. “เอวมาวุโส”ติ โข เต ภิกฺขู อายสฺมโต อานนฺทสฺส ปฏิสฺสุตฺวา เยน ภควา เตนุ\-ปสงฺกมิํสุ, อุปสงฺกมิตฺวา ภควนฺตํ อภิวาเทตฺวา เอกมนฺตํ นิสีทิํสุ. เอกมนฺตํ นิสินฺเน โข เต ภิกฺขู ภควา เอตทโวจ\csromandash{}}

\prose{\csromanfolio{355}อฏฺฐานเมตํ ภิกฺขเว อนวกาโส, ยํ ปรูปกฺกเมน ตถาคตํ ชีวิตา โวโรเปยฺย, อนุปกฺกเมน ภิกฺขเว ตถาคตา ปรินิพฺพายนฺติ.}

\prose{\csromansymbolfootnote{*}{วิ 4. 343; อํ 2. 108 ปิฏฺเฐสุปิ.}ปญฺจิเม ภิกฺขเว สตฺถาโร สนฺโต สํวิชฺชมานา โลกสฺมิํ, กตเม ปญฺจ, อิธ ภิกฺขเว เอกจฺโจ สตฺถา อปริสุทฺธ\-สีโล สมาโน “ปริสุทฺธสีโลมฺหี”ติ ปฏิชานาติ “ปริสุทฺธํ เม สีลํ ปริโยทาตํ อสํกิลิฏฺฐนฺ”ติ จ. ตเมนํ สาวกา เอวํ ชานนฺติ “อยํ โข ภวํ สตฺถา อปริสุทฺธ\-สีโล สมาโน “ปริสุทฺธสีโลมฺหี”ติ ปฏิชานาติ “ปริสุทฺธํ เม สีลํ ปริโยทาตํ อสํกิลิฏฺฐนฺ”ติ จ. มยญฺเจว โข ปน คิหีนํ อาโรเจยฺยาม “นาสฺสสฺส มนาปํ, ยํ โข ปนสฺส อมนาปํ, กถํ นํ มยํ เตน สมุทาจเรยฺยาม. สมฺมนฺนติ โข ปน จีวร\-ปิณฺฑปาต\-เสนาสน\-คิลานปฺปจฺจย\-เภสชฺช\-ปริกฺขาเรน. ยํ ตุโม กริสฺสติ, ตุโมว เตน ปญฺญายิสฺสตี”ติ. เอวรูปํ โข ภิกฺขเว สตฺถารํ สาวกา สีลโต รกฺขนฺติ, เอวรูโป จ ปน สตฺถา สาวเกหิ สีลโต รกฺขํ ปจฺจาสีสติ.}

\prose{ปุน จปรํ ภิกฺขเว อิเธกจฺโจ สตฺถา อปริสุทฺธอาชีโว สมาโน ฯเปฯ อปริสุทฺธธมฺมเทสโน สมาโน ฯเปฯ อปริสุทฺธ\-เวยฺยากรโณ สมาโน ฯเปฯ อปริสุทฺธ\-ญาณทสฺสโน สมาโน “ปริสุทฺธ\-ญาณทสฺสโนมฺหี”ติ ปฏิชานาติ “ปริสุทฺธํ เม ญาณทสฺสนํ ปริโยทาตํ อสํกิลิฏฺฐนฺ”ติ จ. ตเมนํ สาวกา เอวํ ชานนฺติ อยํ โข ภวํ สตฺถา อปริสุทฺธ\-ญาณทสฺสโน สมโน “ปริสุทฺธ\-ญาณทสฺสโนมฺหี”ติ ปฏิชานาติ “ปริสุทฺธํ เม ญาณทสฺสนํ ปริโยทาตํ อสํกิลิฏฺฐนฺ”ติ จ. มยญฺเจว โข ปน คิหีนํ อาโรเจยฺยาม “นาสฺสสฺส มนาปํ, ยํ โข ปนสฺส อมนาปํ, กถํ นํ มยํ เตน สมุทาจเรยฺยาม, สมฺมนฺนติ โข ปน จีวร\-ปิณฺฑปาต\-เสนาสน\-คิลานปฺปจฺจย\-เภสชฺช\-ปริกฺขาเรน. ยํ ตุโม กริสฺสติ, ตุโมว เตน ปญฺญายิสฺสตี”ติ. เอวรูปํ โข ภิกฺขเว สตฺถารํ สาวกา ญาณทสฺสนโต รกฺขนฺติ, เอวรูโป จ ปน สตฺถา สาวเกหิ ญาณทสฺสนโต รกฺขํ ปจฺจาสีสติ. อิเม โข ภิกฺขเว ปญฺจ สตฺถาโร สนฺโต สํวิชฺชมานา โลกสฺมิํ. อหํ โข ปน ภิกฺขเว ปริสุทฺธสีโล สมาโน “ปริสุทฺธสีโลมฺหี”ติ ปฏิชานามิ “ปริสุทฺธํ เม สีลํ ปริโยทาตํ อสํกิลิฏฺฐนฺ”ติ จ, น จ มํ สาวกา สีลโต รกฺขนฺติ, น จาหํ สาวเกหิ สีลโต รกฺขํ ปจฺจาสีสามิ. อหํ โข \csromanfolio{356}ปน ภิกฺขเว ปริสุทฺธาชีโว สมาโน ฯเปฯ ปริสุทฺธธมฺมเทสโน สมาโน ฯเปฯ ปริสุทฺธ\-เวยฺยากรโณ สมาโน ฯเปฯ ปริสุทฺธ\-ญาณทสฺสโน สมาโน “ปริสุทฺธ\-ญาณทสฺสโนมฺหี”ติ ปฏิชานามิ “ปริสุทฺธํ เม ญาณทสฺสนํ ปริโยทาตํ อสํกิลิฏฺฐนฺ”ติ จ, น จ มํ สาวกา ญาณทสฺสนโต รกฺขนฺติ, น จาหํ สาวเกหิ ญาณทสฺสนโต รกฺขํ ปจฺจาสีสามิ. อฏฺฐานเมตํ ภิกฺขเว อนวกาโส, ยํ ปรูปกฺกเมน ตถาคตํ ชีวิตา โวโรเปยฺย, อนุปกฺกเมน ภิกฺขเว ตถาคตา ปรินิพฺพายนฺติ. คจฺฉถ ตุเมฺห ภิกฺขเว ยถาวิหารํ, อรกฺขิยา ภิกฺขเว ตถาคตาติ.}

\csromanmatikaanchor{mtk.170}
\csromantocmarkref{subsection}{นาฬาคิริเปสน}{mtk.170}
\bookmark[dest={mtk.170},level=2]{นาฬาคิริเปสน}
\csromanheader{2}{\textbf{นาฬาคิริเปสน}}
\proseitem{342}{เตน โข ปน สมเยน ราชคเห นาฬาคิริ นาม หตฺถี จณฺโฑ โหติ มนุสฺสฆาตโก. อถ โข เทวทตฺโต ราชคหํ ปวิสิตฺวา หตฺถิสาลํ คนฺตฺวา หตฺถิภณฺเฑ เอตทโวจ “มยํ โข ภเณ ราชญาตกา นาม ปฏิพลา นีจฏฺฐานิยํ อุจฺจฏฺฐาเน ฐเปตุํ, ภตฺตมฺปิ เวตนมฺปิ วฑฺฒาเปตุํ. เตน หิ ภเณ ยทา สมโณ โคตโม อิมํ รจฺฉํ ปฏิปนฺโน โหติ, ตทา อิมํ นาฬาคิริํ หตฺถิํ มุญฺเจตฺวา อิมํ รจฺฉํ ปฏิปาเทถา”ติ. “เอวํ ภนฺเต”ติ โข เต หตฺถิภณฺฑา เทวทตฺตสฺส ปจฺจสฺโสสุํ. อถ โข ภควา ปุพฺพณฺหสมยํ นิวาเสตฺวา ปตฺต\-จีวรมา\-ทาย สมฺพหุเลหิ ภิกฺขูหิ สทฺธิํ รชคหํ ปิณฺฑาย ปาวิสิ. อถ โข ภควา ตํ รจฺฉํ ปฏิปชฺชิ. อทฺทสาสุํ โข เต หตฺถิภณฺฑา ภควนฺตํ ตํ รจฺฉํ ปฏิปนฺนํ, ทิสฺวาน นาฬาคิริํ หตฺถิํ มุญฺจิตฺวา ตํ รจฺฉํ ปฏิปาเทสุํ. อทฺทสา โข นาฬาคิริ หตฺถี ภควนฺตํ ทูรโตว อาคจฺฉนฺตํ, ทิสฺวาน โสณฺฑํ อุสฺสาเปตฺวา ปหฏฺฐ\-กณฺณ\-วาโล เยน ภควา เตน อภิธาวิ. อทฺทสาสุํ โข เต ภิกฺขู นาฬาคิริํ หตฺถิํ ทูรโตว อาคจฺฉนฺตํ, ทิสฺวาน ภควนฺตํ เอตทโวจุํ “อยํ ภนฺเต นาฬาคิริ หตฺถี จณฺโฑ มนุสฺสฆาตโก อิมํ รจฺฉํ ปฏิปนฺโน, ปฏิกฺกมตุ ภนฺเต ภควา, ปฏิกฺกมตุ สุคโต”ติ. อาคจฺฉถ ภิกฺขเว, มา ภายิตฺถ. อฏฺฐานเมตํ ภิกฺขเว อนวกาโส, ยํ ปรูปกฺกเมน ตถาคตํ ชีวิตา โวโรเปยฺย, อนุปกฺกเมน ภิกฺขเว ตถาคตา ปรินิพฺพายนฺตีติ. ทุติยมฺปิ โข เต ภิกฺขู ฯเปฯ. ตติยมฺปิ โข เต ภิกฺขู ภควนฺตํ เอตทโวจุํ “อยํ ภนฺเต นาฬาคิริ หตฺถี จณฺโฑ มนุสฺสฆาตโก อิมํ รจฺฉํ ปฏิปนฺโน, ปฏิกฺกมตุ ภนฺเต ภควา, ปฏิกฺกมตุ สุคโต”ติ. อาคจฺฉถ ภิกฺขเว, มา ภายิตฺถ. \csromanfolio{357}อฏฺฐานเมตํ ภิกฺขเว อนวกาโส, ยํ ปรูปกฺกเมน ตถาคตํ ชีวิตา โวโรเปยฺย, อนุปกฺกเมน ภิกฺขเว ตถาคตา ปรินิพฺพายนฺตีติ.}

\prose{เตน โข ปน สมเยน มนุสฺสา ปาสาเทสุปิ หมฺมิเยสุปิ ฉทเนสุปิ อารุฬฺหา อจฺฉนฺติ. ตตฺถ เย เต มนุสฺสา อสฺสทฺธา อปฺปสนฺนา ทุพฺพุทฺธิโน, เต เอวมาหํสุ “อภิรูโป วต โภ\csromansharedfootnote{7f77e5c45963db79}{อภิรูโป วต โภ โคตโม (สฺยา, กํ)} มหาสมโณ นาเคน วิเหฐียิสฺสตี”ติ. เย ปน เต มนุสฺสาสทฺธา ปสนฺนา ปณฺฑิตา พฺยตฺตา พุทฺธิมนฺโต, เต เอวมาหํสุ “นจิรสฺสํ วต โภ นาโค นาเคน สงฺคาเมสฺสตี”ติ. อถ โข ภควา นาฬาคิริํ หตฺถิํ เมตฺเตน จิตฺเตน ผริ. อถ โข นาฬาคิริ หตฺถี ภควโต\csromansharedfootnote{a5b66bfa43168243}{ภควตา (สี)} เมตฺเตน จิตฺเตน ผุฏฺโฐ\csromansharedfootnote{15eae50defa7a0c3}{ผุโฏ (ก)} โสณฺฑํ โอโรเปตฺวา เยน ภควา เตนุปสงฺกมิ, อุปสงฺกมิตฺวา ภควโต ปุรโต อฏฺฐาสิ. อถ โข ภควา ทกฺขิเณน หตฺเถน นาฬาคิริสฺส หตฺถิสฺส กุมฺภํ ปรามสนฺโต นาฬาคิริํ หตฺถิํ อิมาหิ คาถาหิ อชฺฌภาสิ\csromandash{}}

\setlength{\csromangathapagewidth}{0pt}
\setlength{\csromangathawakwidth}{0pt}
\csromangathameasuregroup{“มา กุญฺชร นาคมาสโท, \\ น หิ นาคหตสฺส กุญฺชร, \\ มา จ มโท มา จ ปมาโท, \\ ตฺวญฺเญว ตถา กริสฺสสิ,}{\csromangathabat{“มา กุญฺชร นาคมาสโท,}{ทุกฺขํ หิ กุญฺชร นาคมาสโท.} \\ \csromangathabat{น หิ นาคหตสฺส กุญฺชร,}{สุคติ โหติ อิโต ปรํ ยโต.} \\ \csromangathabat{มา จ มโท มา จ ปมาโท,}{น หิ ปมตฺตา สุคติํ วชนฺติ เต.} \\ \csromangathabat{ตฺวญฺเญว ตถา กริสฺสสิ,}{เยน ตฺวํ สุคติํ คมิสฺสสี”ติ.}}
\csromangathasetleft{“มา กุญฺชร นาคมาสโท, \\ น หิ นาคหตสฺส กุญฺชร, \\ มา จ มโท มา จ ปมาโท, \\ ตฺวญฺเญว ตถา กริสฺสสิ,}
\csromangathagroup{\csromangathastanza{\csromangathabat{“มา กุญฺชร นาคมาสโท,}{ทุกฺขํ หิ กุญฺชร นาคมาสโท.} \\ \csromangathabat{น หิ นาคหตสฺส กุญฺชร,}{สุคติ โหติ อิโต ปรํ ยโต.}}\csromangathastanza{\csromangathabat{มา จ มโท มา จ ปมาโท,}{น หิ ปมตฺตา สุคติํ วชนฺติ เต.} \\ \csromangathabat{ตฺวญฺเญว ตถา กริสฺสสิ,}{เยน ตฺวํ สุคติํ คมิสฺสสี”ติ.}}}

\prose{อถ โข นาฬาคิริ หตฺถี โสณฺฑาย ภควโต ปาทปํสูนิ คเหตฺวา อุปริมุทฺธนิ อากิริตฺวา ปฏิกุฏิโยว\csromansharedfootnote{101bdefc839626d6}{ปฏิกุฏิโต ปฏิสกฺกิ. (สี, สฺยา)} โอสกฺกิ. ยาว ภควนฺตํ อทฺทกฺขิ. อถ โข นาฬาคิริ หตฺถี หตฺถิสาลํ คนฺตฺวา สเก ฐาเน อฏฺฐาสิ, ตถา ทนฺโต จ ปน นาฬาคิริ หตฺถี อโหสิ. เตน โข ปน สมเยน มนุสฺสา อิมํ คาถํ คายนฺติ\csromandash{}}

\setlength{\csromangathapagewidth}{0pt}
\setlength{\csromangathawakwidth}{0pt}
\csromangathameasuregroup{\textsuperscript{*}“ทณฺเฑเนเก ทมยนฺติ, \\ อทณฺเฑน อสตฺเถน,}{\csromangathabat{\textsuperscript{*}“ทณฺเฑเนเก ทมยนฺติ,}{องฺกุเสหิ กสาหิ จ.} \\ \csromangathabat{อทณฺเฑน อสตฺเถน,}{นาโค ทนฺโต มเหสินา”ติ.}}
\csromangathasetleft{\textsuperscript{*}“ทณฺเฑเนเก ทมยนฺติ, \\ อทณฺเฑน อสตฺเถน,}
\csromangathagroup{\csromangathastanza{\csromangathaitembat{342}{\csromansymbolfootnote{*}{ม 2. 308; ขุ 2. 335 ปิฏฺเฐสุปิ.}“ทณฺเฑเนเก ทมยนฺติ,}{องฺกุเสหิ กสาหิ จ.} \\ \csromangathacontbat{อทณฺเฑน อสตฺเถน,}{นาโค ทนฺโต มเหสินา”ติ.}}}

\prose{มนุสฺสา อุชฺฌายนฺติ ขิยฺยนฺติ วิปาเจนฺติ “ยาว ปาโป อยํ เทวทตฺโต อลกฺขิโก, ยตฺร หิ นาม สมณสฺส โคตมสฺส เอวํ\-มหิทฺธิกสฺส เอวํ\-มหา\-นุภาวสฺส วธาย ปรกฺกมิสฺสตี”ติ. เทวทตฺตสฺส ลาภสกฺกาโร ปริหายิ, ภควโต จ ลาภสกฺกาโร อภิวฑฺฒิ.}

\csromanmatikaanchor{mtk.171}
\csromantocmarkref{subsection}{ปญฺจวตฺถุยาจนกถา}{mtk.171}
\bookmark[dest={mtk.171},level=2]{ปญฺจวตฺถุยาจนกถา}
\csromanheader{2}{ปญฺจ\-วตฺถุ\-ยาจน\-กถา}
\proseitem{343}{\csromanfolio{358}\csromansymbolfootnote{*}{วิ 2. 97 ปิฏฺเฐปิ.}เตน โข ปน สมเยน เทวทตฺโต ปริหีน\-ลาภสกฺกาโร สปริโส กุเลสุ วิญฺญาเปตฺวา ภุญฺชติ. มนุสฺสา อุชฺฌายนฺติ ขิยฺยนฺติ วิปาจินฺติ “กถํ หิ นาม สมณา สกฺยปุตฺติยา กุเลสุ วิญฺญาเปตฺวา ภุญฺชิสฺสนฺติ, กสฺส สมฺปนฺนํ น มนาปํ, กสฺส สาทุํ น รุจฺจตี”ติ. อสฺโสสุํ โข ภิกฺขู เตสํ มนุสฺสานํ อุชฺฌายนฺตานํ ขิยฺยนฺตานํ วิปาเจนฺตานํ. เย เต ภิกฺขู อปฺปิจฺฉา ฯเปฯ เต อุชฺฌายนฺติ ขิยฺยนฺติ วิปาเจนฺติ “กถํ หิ นาม เทวทตฺโต สปริโส กุเลสุ วิญฺญาเปตฺวา วิญฺญาเปตฺวา ภุญฺชิสฺสตี”ติ. ภควโต เอตมตฺถํ อาโรเจสุํ ฯเปฯ. สจฺจํ กิร ตฺวํ เทวทตฺต สปริโส กุเลสุ วิญฺญาเปตฺวา วิญฺญาเปตฺวา ภุญฺชสีติ. สจฺจํ ภควาติ ฯเปฯ วิครหิตฺวา ฯเปฯ ธมฺมิํ กถํ กตฺวา ภิกฺขู อามนฺเตสิ “เตน หิ ภิกฺขเว ภิกฺขูนํ กุเลสุ ติกโภชนํ ปญฺญเปสฺสามิ ตโย อตฺถวเส ปฏิจฺจ, ทุมฺมงฺกูนํ ปุคฺคลานํ นิคฺคหาย, เปสลานํ ภิกฺขูนํ ผาสุวิหาราย ‘มา ปาปิจฺฉา ปกฺขํ นิสฺสาย สํฆํ ภินฺเทยฺยุนฺ’ติ, กุลานุทฺทยาย\csromansharedfootnote{5592a0dfb94b620d}{กุลานุทยตาย (สี, สฺยา)} จ. คณโภชเน ยถาธมฺโม กาเรตพฺโพ”ติ.}

\prose{\csromansymbolfootnote{+}{วิ 1. 263 ปิฏฺเฐปิ.}อถ โข เทวทตฺโต เยน โกกาลิโก กฏโมทกติสฺสโก\csromansharedfootnote{bdcfbc07a8c5d26f}{กุฏโมรกติสฺสโก (สี, สฺยา)} ขณฺฑเทวิยา ปุตฺโต สมุทฺททตฺโต เตนุปสงฺกมิ, อุปสงฺกมิตฺวา โกกาลิกํ กฏโมทกติสฺสกํ ขณฺฑเทวิยา ปุตฺตํ สมุทฺททตฺตํ เอตทโวจ “เอถ มยํ อาวุโส สมณสฺส โคตมสฺส สํฆเภทํ กริสฺสาม จกฺกเภทนฺ”ติ. เอวํ วุตฺเต โกกาลิโก เทวทตฺตํ เอตทโวจ “สมโณ โข อาวุโส โคตโม มหิทฺธิโก มหานุภาโว, กถํ มยํ สมณสฺส โคตมสฺส สํฆเภทํ กริสฺสาม จกฺกเภทนฺ”ติ. เอถ มยํ อาวุโส สมณํ โคตมํ อุปสงฺกมิตฺวา ปญฺจ วตฺถูนิ ยาจิสฺสาม “ภควา ภนฺเต อเนก\-ปริยาเยน อปฺปิจฺฉสฺส สนฺตุฏฺฐสฺส สลฺเลขสฺส ธุตสฺส ปาสาทิกสฺส อปจยสฺส วีริยา\-รมฺภสฺส วณฺณวาที, อิมานิ ภนฺเต ปญฺจ วตฺถูนิ อเนก\-ปริยาเยน อปฺปิจฺฉตาย สนฺตุฏฺฐิยา สลฺเลขาย ธุตตาย ปาสาทิกตาย อปจยาย วีริยารมฺภาย สํวตฺตนฺติ, สาธุ ภนฺเต ภิกฺขู ยาวชีวํ อารญฺญิกา อสฺสุ, โย \csromanfolio{359}คามนฺตํ โอสเรยฺย, วชฺชํ นํ ผุเสยฺย. ยาวชีวํ ปิณฺฑปาติกา อสฺสุ, โย นิมนฺตนํ สาทิเยยฺย, วชฺชํ นํ ผุเสยฺย. ยาวชีวํ ปํสุกูลิกา อสฺสุ, โย คหปติ\-จีวรํ สาทิเยยฺย, วชฺชํ นํ ผุเสยฺย. ยาวชีวํ รุกฺขมูลิกา อสฺสุ, โย ฉนฺนํ อุปคจฺเฉยฺย, วชฺชํ นํ ผุเสยฺย. ยาวชีวํ มจฺฉมํสํ น ขาเทยฺยุํ, โย มจฺฉมํสํ ขาเทยฺย, วชฺชํ นํ ผุเสยฺยา”ติ. อิมานิ ปญฺจ วตฺถูนิ สมโณ โคตโม นานุชานิสฺสติ. เต มยํ อิเมหิ ปญฺจหิ วตฺถูหิ ชนํ สญฺญาเปสฺสามาติ. สกฺกา โข อาวุโส อิเมหิ ปญฺจหิ วตฺถูหิ สมณสฺส โคตมสฺส สํฆเภโท กาตุํ จกฺกเภโท, ลูขปฺปสนฺนา หิ อาวุโส มนุสฺสาติ.}

\prose{อถ โข เทวทตฺโต สปริโส เยน ภควา เตนุปสงฺกมิ, อุปสงฺกมิตฺวา ภควนฺตํ อภิวาเทตฺวา เอกมนฺตํ นิสีทิ, เอกมนฺตํ นิสินฺโน โข เทวทตฺโต ภควนฺตํ เอตทโวจ “ภควา ภนฺเต อเนก\-ปริยาเยน อปฺปิจฺฉสฺส สนฺตุฏฺฐสฺส สลฺเลขสฺส ธุตสฺส ปาสาทิกสฺส อปจยสฺส วีริยา\-รมฺภสฺส วณฺณวาที, อิมานิ ภนฺเต ปญฺจ วตฺถูนิ อเนก\-ปริยาเยน อปฺปิจฺฉตาย สนฺตุฏฺฐิยา สลฺเลขาย ธุตตาย ปาสาทิกตาย อปจยาย วีริยารมฺภาย สํวตฺตนฺติ, สาธุ ภนฺเต ภิกฺขู ยาวชีวํ อารญฺญิกา อสฺสุ, โย คามนฺตํ โอสเรยฺย, วชฺชํ นํ ผุเสยฺย. ยาวชีวํ ปิณฺฑปาติกา อสฺสุ, โย นิมนฺตนํ สาทิเยยฺย, วชฺชํ นํ ผุเสยฺย. ยาวชีวํ ปํสุกูลิกา อสฺสุ, โย คหปติ\-จีวรํ สาทิเยยฺย, วชฺชํ นํ ผุเสยฺย. ยาวชีวํ รุกฺขมูลิกา อสฺสุ, โย ฉนฺนํ อุปคจฺเฉยฺย, วชฺชํ นํ ผุเสยฺย. ยาวชีวํ มจฺฉมํสํ น ขาเทยฺยุํ, โย มจฺฉมํสํ ขาเทยฺย, วชฺชํ นํ ผุเสยฺยา”ติ. อลํ เทวทตฺต, โย อิจฺฉติ อารญฺญิโก โหตุ, โย อิจฺฉติ คามนฺเต วิหรตุ, โย อิจฺฉติ ปิณฺฑปาติโก โหตุ, โย อิจฺฉติ นิมนฺตนํ สาทิยตุ, โย อิจฺฉติ ปํสุกูลิโก โหตุ, โย อิจฺฉติ คหปติ\-จีวรํ สาทิยตุ, อฏฺฐมาเส โข มยา เทวทตฺต รุกฺขมูล\-เสนาสนํ อนุญฺญาตํ, ติโกฏิ\-ปริสุทฺธํ มจฺฉมํสํ อทิฏฺฐํ อสฺสุตํ อปริสงฺกิตนฺติ. อถ โข เทวทตฺโต “น ภควา อิมานิ ปญฺจ วตฺถูนิ อนุชานาตี”ติ หฏฺโฐ อุทคฺโค สปริโส อุฏฺฐายาสนา ภควนฺตํ อภิวาเทตฺวา ปทกฺขิณํ กตฺวา ปกฺกามิ.}

\prose{อถ โข เทวทตฺโต สปริโส ราชคหํ ปวิสิตฺวา ปญฺจหิ วตฺถูหิ ชนํ สญฺญาเปสิ “มยํ อาวุโส สมณํ โคตมํ อุปสงฺกมิตฺวา ปญฺจ \csromanfolio{360}วตฺถูนิ ยาจิมฺหา ‘ภควา ภนฺเต อเนก\-ปริยาเยน อปฺปิจฺฉสฺส ฯเปฯ วีริยา\-รมฺภสฺส วณฺณวาที, อิมานิ ภนฺเต ปญฺจ วตฺถูนิ อเนก\-ปริยาเยน อปฺปิจฺฉตาย ฯเปฯ วีริยารมฺภาย สํวตฺตนฺติ, สาธุ ภนฺเต ภิกฺขู ยาวชีวํ อารญฺญิกา อสฺสุ, โย คามนฺตํ โอสเรยฺย, วชฺชํ นํ ผุเสยฺย ฯเปฯ. ยาวชีวํ มจฺฉมํสํ น ขาเทยฺยุํ, โย มจฺฉมํสํ ขาเทยฺย, วชฺชํ นํ ผุเสยฺยา’ติ. อิมานิ ปญฺจ วตฺถูนิ สมโณ โคตโม นานุชานาติ, เต มยํ อิเมหิ ปญฺจหิ วตฺถูหิ สมาทาย วตฺตามา”ติ.}

\prose{ตตฺถ เย เต มนุสฺสา อสฺสทฺธา อปฺปสนฺนา ทุพฺพุทฺธิโน, เต เอวมาหํสุ “อิเม โข สมณา สกฺยปุตฺติยา ธุตา สลฺเลข\-วุตฺติโน, สมโณ ปน โคตโม พาหุลฺลิโก พาหุลฺลาย เจเตตี”ติ. เย ปน เต มนุสฺสา สทฺธา ปสนฺนา ปณฺฑิตา พฺยตฺตา พุทฺธิมนฺโต, เต อุชฺฌายนฺติ ขิยฺยนฺติ วิปาเจนฺติ “กถํ หิ นาม เทวทตฺโต ภควโต สํฆเภทาย ปรกฺกมิสฺสติ จกฺกเภทายา”ติ. อสฺโสสุํ โข ภิกฺขู เตสํ มนุสฺสานํ อุชฺฌายนฺตานํ ขิยฺยนฺตานํ วิปาเจนฺตานํ. เย เต ภิกฺขู อปฺปิจฺฉา ฯเปฯ เต อุชฺฌายนฺติ ขิยฺยนฺติ วิปาเจนฺติ “กถํ หิ นาม เทวทตฺโต สํฆเภทาย ปรกฺกมิสฺสติ จกฺกเภทายา”ติ. อถ โข เต ภิกฺขู ภควโต เอตมตฺถํ อาโรเจสุํ ฯเปฯ. สจฺจํ กิร ตฺวํ เทวทตฺต สํฆเภทาย ปรกฺกมสิ จกฺกเภทายาติ. สจฺจํ ภควาติ. อลํ เทวทตฺต, มา เต รุจฺจิ สํฆเภโท, ครุโก โข เทวทตฺต สํฆเภโท, โย โข เทวทตฺต สมคฺคํ สํฆํ ภินฺทติ, กปฺปฏฺฐิกํ\csromansharedfootnote{697d06f218d7c4f7}{กปฺปฏฺฐิติกํ (สฺยา)} กิพฺพิสํ ปสวติ, กปฺปํ นิรยมฺหิ ปจฺจติ. โย จ โข เทวทตฺต ภินฺนํ สํฆํ สมคฺคํ กโรติ, พฺรหฺมํ ปุญฺญํ ปสวติ, กปฺปํ สคฺคมฺหิ โมทติ, อลํ เทวทตฺต, มา เต รุจฺจิ สํฆเภโท, ครุโก โข เทวทตฺต สํฆเภโทติ.}

\prose{\csromansymbolfootnote{*}{ขุ 1. 149 ปิฏฺเฐ อุทาเนปิ.}อถ โข อายสฺมา อานนฺโท ปุพฺพณฺหสมยํ นิวาเสตฺวา ปตฺต\-จีวรมา\-ทาย ราชคหํ ปิณฺฑาย ปาวิสิ. อทฺทสา โข เทวทตฺโต อายสฺมนฺตํ อานนฺทํ ราชคเห ปิณฺฑาย จรนฺตํ, ทิสฺวาน เยนายสฺมา อานนฺโท เตนุปสงฺกมิ, อุปสงฺกมิตฺวา อายสฺมนฺตํ อานนฺทํ เอตทโวจ “อชฺชตคฺเค\-ทานาหํ อาวุโส อานนฺท อญฺญเตฺรว ภควตา, อญฺญเตฺรว ภิกฺขุสํฆา อุโปสถํ กริสฺสามิ, สํฆกมฺมํ กริสฺสามี”ติ.}

\csromanmatikaanchor{mtk.172}
\csromanmatikaanchor{mtk.173}
\csromantocmarkref{section}{3. ตติยภาณวาร}{mtk.172}
\bookmark[dest={mtk.172},level=1]{3. ตติยภาณวาร}
\csromantocmarkref{subsection}{สํฆเภทกถา}{mtk.173}
\bookmark[dest={mtk.173},level=2]{สํฆเภทกถา}
\prose{\csromanfolio{361}อถ โข อายสฺมา อานนฺโท ราชคเห ปิณฺฑาย จริตฺวา ปจฺฉาภตฺตํ ปิณฺฑปาต\-ปฺปฏิกฺกนฺโต เยน ภควา เตนุปสงฺกมิ, อุปสงฺกมิตฺวา ภควนฺตํ อภิวาเทตฺวา เอกมนฺตํ นิสีทิ, เอกมนฺตํ นิสินฺโน โข อายสฺมา อานนฺโท ภควนฺตํ เอตทโวจ “อิธาหํ ภนฺเต ปุพฺพณฺหสมยํ นิวาเสตฺวา ปตฺต\-จีวรมา\-ทาย ราชคหํ ปิณฺฑาย ปาวิสิํ, อทฺทสา โข มํ ภนฺเต เทวทตฺโต ราชคเห ปิณฺฑาย จรนฺตํ, ทิสฺวาน เยนาหํ เตนุปสงฺกมิ, อุปสงฺกมิตฺวา มํ เอตทโวจ ‘อชฺชตคฺเค\-ทานาหํ อาวุโส อานนฺท อญฺญเตฺรว ภควตา, อญฺญเตฺรว ภิกฺขุสํฆา อุโปสถํ กริสฺสามิ, สํฆกมฺมํ กริสฺสามี”ติ, อชฺชตคฺเค ภนฺเต เทวทตฺโต สํฆํ ภินฺทิสฺสตี”ติ. อถ โข ภควา เอตมตฺถํ วิทิตฺวา ตายํ เวลายํ อิมํ อุทานํ อุทาเนสิ\csromandash{}}

\setlength{\csromangathapagewidth}{0pt}
\setlength{\csromangathawakwidth}{0pt}
\csromangathameasuregroup{\textsuperscript{*}“สุกรํ สาธุนา สาธุํ, \\ ปาปํ ปาเปน สุกรํ,}{\csromangathabat{\textsuperscript{*}“สุกรํ สาธุนา สาธุํ,}{สาธุํ ปาเปน ทุกฺกรํ.} \\ \csromangathabat{ปาปํ ปาเปน สุกรํ,}{ปาปมริเยหิ ทุกฺกรนฺ”ติ.}}
\csromangathasetleft{\textsuperscript{*}“สุกรํ สาธุนา สาธุํ, \\ ปาปํ ปาเปน สุกรํ,}
\csromangathagroupwithcloser{\csromangathastanza{\csromangathabat{\csromansymbolfootnote{*}{ขุ 1. 150 ปิฏฺเฐ อุทาเนปิ.}“สุกรํ สาธุนา สาธุํ,}{สาธุํ ปาเปน ทุกฺกรํ.} \\ \csromangathabat{ปาปํ ปาเปน สุกรํ,}{ปาปมริเยหิ ทุกฺกรนฺ”ติ.}}}{ทุติย\-ภาณวาโร นิฏฺฐิโต.}
\csromanheader{1}{3. ตติย\-ภาณวาร}
\csromanheader{2}{\textbf{สํฆเภทถา}}
\proseitem{344}{อถ โข เทวทตฺโต ตทหุโปสเถ อุฏฺฐายาสนา สลากํ คาเหสิ “มยํ อาวุโส สมณํ โคตมํ อุปสงฺกมิตฺวา ปญฺจ วตฺถูนิ ยาจิมฺหา ‘ภควา ภนฺเต อเนก\-ปริยาเยน อปฺปิจฺฉสฺส ฯเปฯ วีริยา\-รมฺภสฺส วณฺณวาที, อิมานิ ภนฺเต ปญฺจ วตฺถูนิ อเนก\-ปริยาเยน อปฺปิจฺฉตาย ฯเปฯ วีริยารมฺภาย สํวตฺตนฺติ, สาธุ ภนฺเต ภิกฺขู ยาวชีวํ อารญฺญิกา อสฺสุ, โย คามนฺตํ โอสเรยฺย, วชฺชํ นํ ผุเสยฺย ฯเปฯ. ยาวชีวํ มจฺฉมํสํ น ขาเทยฺยุํ, โย มจฺฉมํสํ ขาเทยฺย, วชฺชํ นํ ผุเสยฺยา’ติ. อิมานิ ปญฺจ วตฺถูนิ สมโณ โคตโม นานุชานาติ, เต มยํ อิเมหิ ปญฺจหิ วตฺถูหิ สมาทาย วตฺตาม, ยสฺสายสฺมโต อิมานิ ปญฺจ วตฺถูนิ ขมนฺติ, โส สลากํ คณฺหาตู”ติ.}

\prose{เตน โข ปน สมเยน เวสาลิกา วชฺชิปุตฺตกา ปญฺจมตฺตานิ ภิกฺขุสตานิ นวกา เจว โหนฺติ อปฺปกตญฺญุโน จ, เต “อยํ \csromanfolio{362}ธมฺโม, อยํ วินโย, อิทํ สตฺถุสาสนนฺ”ติ สลากํ คณฺหิํสุ. อถ โข เทวทตฺโต สํฆํ ภินฺทิตฺวา ปญฺจมตฺตานิ ภิกฺขุสตานิ อาทาย เยน คยาสีสํ เตน ปกฺกามิ. อถ โข สาริปุตฺต\-โมคฺคลฺลานา เยน ภควา เตนุ\-ปสงฺกมิํสุ, อุปสงฺกมิตฺวา ภควนฺตํ อภิวาเทตฺวา เอกมนฺตํ นิสีทิํสุ. เอกมนฺตํ นิสินฺโน โข อายสฺมา สาริปุตฺโต ภควนฺตํ เอตทโวจ “เทวทตฺโต ภนฺเต สํฆํ ภินฺทิตฺวา ปญฺจมตฺตานิ ภิกฺขุสตานิ อาทาย เยน คยาสีสํ, เตน ปกฺกนฺโต”ติ. น หิ นาม ตุมฺหากํ สาริปุตฺตา เตสุ นวเกสุ ภิกฺขูสุ การุญฺญมฺปิ ภวิสฺสติ, คจฺฉถ ตุเมฺห สาริปุตฺตา, ปุรา เต ภิกฺขู อนยพฺยสนํ อาปชฺชนฺตีติ. “เอวํ ภนฺเต”ติ โข สาริปุตฺต\-โมคฺคลฺลานา ภควโต ปฏิสฺสุตฺวา อุฏฺฐายาสนา ภควนฺตํ อภิวาเทตฺวา ปทกฺขิณํ กตฺวา เยน คยาสีสํ เตนุ\-ปสงฺกมิํสุ.}

\prose{เตน โข ปน สมเยน อญฺญตโร ภิกฺขุ ภควโต อวิทูเร โรทมาโน ฐิโต โหติ. อถ โข ภควา ตํ ภิกฺขุํ เอตทโวจ “กิสฺส ตฺวํ ภิกฺขุ โรทสี”ติ. เยปิ เต ภนฺเต ภควโต อคฺคสาวกา สาริปุตฺตโมคลฺลานา, เตปิ เทวทตฺตสฺส สนฺติเก คจฺฉนฺติ เทวทตฺตสฺส ธมฺมํ โรเจนฺตาติ. อฏฺฐานเมตํ ภิกฺขุ อนวกาโส, ยํ สาริปุตฺต\-โมคฺคลฺลานา เทวทตฺตสฺส ธมฺมํ โรเจยฺยุํ. อปิ จ เต คตา ภิกฺขูนํ สญฺญตฺติยาติ\csromansharedfootnote{c36269873a0d74a2}{ภิกฺขุสญฺญตฺติยาติ (สี, สฺยา), ภิกฺขู สญฺญตฺติยา (ก)}.}

\proseitem{345}{เตน โข ปน สมเยน เทวทตฺโต มหติยา ปริสาย ปริวุโต ธมฺมํ เทเสนฺโต นิสินฺโน โหติ. อทฺทสา โข เทวทตฺโต สาริปุตฺต\-โมคฺคลฺลาเน ทูรโตว อาคจฺฉนฺเต, ทิสฺวาน ภิกฺขู อามนฺเตสิ “ปสฺสถ ภิกฺขเว, ยาว สฺวากฺขาโต มยํ ธมฺโม, เยปิ เต สมณสฺส โคตมสฺส อคฺคสาวกา สาริปุตฺต\-โมคฺคลฺลานา, เตปิ มม สนฺติเก อาคจฺฉนฺติ มม ธมฺมํ โรเจนฺตา”ติ. เอวํ วุตฺเต โกกาลิโก เทวทตฺตํ เอตทโวจ “มา อาวุโส เทวทตฺต สาริปุตฺต\-โมคฺคลฺลาเน วิสฺสาสิ, ปาปิจฺฉา สาริปุตฺต\-โมคฺคลฺลานา ปาปิกานํ อิจฺฉานํ วสํ คตา”ติ. อลํ อาวุโส, สฺวาคตํ เตสํ, ยโต เม ธมฺมํ โรเจนฺตีติ.}

\prose{อถ โข เทวทตฺโต อายสฺมนฺตํ สาริปุตฺตํ อุปฑฺฒาสเนน นิมนฺเตสิ “เอหาวุโส สาริปุตฺต, อิธ นิสีทาหี”ติ. “อลํ อาวุโส”ติ โข \csromanfolio{363}อายสฺมา สาริปุตฺโต อญฺญตรํ อาสนํ คเหตฺวา เอกมนฺตํ นิสีทิ อายสฺมาปิ โข มหาโมคฺคลฺลาโน อญฺญตรํ อาสนํ คเหตฺวา เอกมนฺตํ นิสีทิ. อถ โข เทวทตฺโต พหุเทว รตฺติํ ภิกฺขู ธมฺมิยา กถาย สนฺทสฺเสตฺวา สมาทเปตฺวา สมุตฺเตเชตฺวา สมฺปหํเสตฺวา อายสฺมนฺตํ สาริปุตฺตํ อชฺเฌสิ “วิคต\-ถินมิทฺโธ โข อาวุโส สาริปุตฺต ภิกฺขุสํโฆ, ปฏิภาตุ ตํ อาวุโส สาริปุตฺต ภิกฺขูนํ ธมฺมี กถา, ปิฏฺฐิ เม อาคิลายติ, ตมหํ อายมิสฺสามี”ติ. “เอวมาวุโส”ติ โข อายสฺมา สาริปุตฺโต เทวทตฺตสฺส ปจฺจสฺโสสิ. อถ โข เทวทตฺโต จตุคฺคุณํ สํฆาฏิํ ปญฺญเปตฺวา ทกฺขิเณน ปสฺเสน เสยฺยํ กปฺเปสิ. ตสฺส กิลมนฺตสฺส มุฏฺฐสฺสติสฺส อสมฺปชานสฺส มุหุตฺตเกเนว นิทฺทา โอกฺกมิ.}

\prose{อถ โข อายสฺมา สาริปุตฺโต อาเทสนาปาฏิหาริยา\-นุสาสนิยา ภิกฺขู ธมฺมิยา กถาย โอวทิ อนุสาสิ, อายสฺมา มหาโมคฺคลฺลาโน อิทฺธิปาฏิหาริยา\-นุสาสนิยา ภิกฺขู ธมฺมิยา กถาย โอวทิ อนุสาสิ. อถ โข เตสํ ภิกฺขูนํ อายสฺมตา สาริปุตฺเตน อาเทสนาปาฏิหาริยา\-นุสาสนิยา, อายสฺมตา จ มหา\-โมคฺคลฺลาเนน อิทฺธิปาฏิหาริยา\-นุสาสนิยา โอวทิยมานานํ อนุสาสิย\-มานานํ วิรชํ วีตมลํ ธมฺมจกฺขุํ อุทปาทิ “ยํ กิญฺจิ สมุทย\-ธมฺมํ, สพฺพํ ตํ นิโรธธมฺมนฺ”ติ.}

\prose{อถ โข อายสฺมา สาริปุตฺโต ภิกฺขู อามนฺเตสิ “คจฺฉาม มยํ อาวุโส ภควโต สนฺติเก, โย ตสฺส ภควโต ธมฺมํ โรเจสิ, โส อาคจฺฉตู”ติ. อถ โข สาริปุตฺต\-โมคฺคลฺลานา ตานิ ปญฺจ ภิกฺขุสตานิ อาทาย เยน เวฬุวนํ เตนุ\-ปสงฺกมิํสุ. อถ โข โกกาลิโก เทวทตฺตํ อุฏฺฐาเปสิ, อุฏฺเฐหิ อาวุโส เทวทตฺต, นีตา เต ภิกฺขู สาริปุตฺต\-โมคฺคลฺลาเนหิ, นนุ ตฺวํ อาวุโส เทวทตฺต มยา วุตฺโต “มา อาวุโส เทวทตฺต สาริปุตฺต\-โมคฺคลฺลาเน วิสฺสาสิ, ปาปิจฺฉา สาริปุตฺต\-โมคฺคลฺลานา ปาปิกานํ อิจฺฉานํ วสํ คตา”ติ. อถ โข เทวทตฺตสฺส ตตฺเถว อุณฺหํ โลหิตํ มุขโต อุคฺคญฺฉิ.}

\prose{อถ โข สาริปุตฺต\-โมคฺคลฺลานา เยน ภควา เตนุ\-ปสงฺกมิํสุ, อุปสงฺกมิตฺวา ภควนฺตํ อภิวาเทตฺวา เอกมนฺตํ นิสีทิํสุ. เอกมนฺตํ นิสินฺโน โข อายสฺมา สาริปุตฺโต ภควนฺตํ เอตทโวจ “สาธุ \csromanfolio{364}ภนฺเต เภทกา\-นุวตฺตกา ภิกฺขู ปุน อุปสมฺปชฺเชยฺยุนฺ”ติ. อลํ สาริปุตฺต, มา เต รุจฺจิ เภทกา\-นุวตฺตกานํ ภิกฺขูนํ ปุน อุปสมฺปทา. เตน หิ ตฺวํ สาริปุตฺต เภทกา\-นุวตฺตเก ภิกฺขู ถุลฺลจฺจยํ เทสาเปหิ. กถํ ปน เต สาริปุตฺต เทวทตฺโต ปฏิปชฺชีติ. ยเถว ภนฺเต ภควา พหุเทว รตฺติํ ภิกฺขู ธมฺมิยา กถาย สนฺทสฺเสตฺวา สมาทเปตฺวา สมุตฺเตเชตฺวา สมฺปหํเสตฺวา มํ อชฺเฌสติ “วิคต\-ถินมิทฺโธ โข สาริปุตฺต ภิกฺขุสํโฆ, ปฏิภาตุ ตํ สาริปุตฺต ภิกฺขูนํ ธมฺมี กถา, ปิฏฺฐิ เม อาคิลายติ, ตมหํ อายมิสฺสามี”ติ. เอวเมว โข ภนฺเต เทวทตฺโต ปฏิปชฺชีติ.}

\proseitem{346}{อถ โข ภควา ภิกฺขู อามนฺเตสิ “ภูตปุพฺพํ ภิกฺขเว อรญฺญายตเน มหาสรสี, ตํ นาคา อุปนิสฺสาย วิหริํสุ, เต ตํ สรสิํ โอคาเหตฺวา โสณฺฑาย ภิสมุฬาลํ อพฺพุหิตฺวา สุวิกฺขาลิตํ วิกฺขาเลตฺวา อกทฺทมํ สํขาทิตฺวา อชฺโฌหรนฺติ. เตสํ ตํ วณฺณาย เจว โหติ พลาย จ, น จ ตโตนิทานํ มรณํ วา นิคจฺฉนฺติ, มรณมตฺตํ วา ทุกฺขํ. เตสํเยว โข ปน ภิกฺขเว มหานาคานํ อนุสิกฺขมานา ตรุณา ภิงฺกจฺฉาปา เต ตํ สรสิํ โอคาเหตฺวา โสณฺฑาย ภิสมุฬาลํ อพฺพุหิตฺวา น สุวิกฺขาลิตํ วิกฺขาลิตฺวา สกทฺทมํ สํขาทิตฺวา อชฺโฌหรนฺติ. เตสํ ตํ เนว วณฺณาย โหติ, น พลาย, ตโตนิทานญฺจ มรณํ วา นิคจฺฉนฺติ, มรณมตฺตํ วา ทุกฺขํ. เอวเมว โข ภิกฺขเว เทวทตฺโต มมานุกฺรุพฺพํ\csromansharedfootnote{94fdffeafd7d35ad}{มมานุกุพฺพํ (สี, สฺยา)} กปโณ มริสฺสตี”ติ.}

\setlength{\csromangathapagewidth}{0pt}
\setlength{\csromangathawakwidth}{0pt}
\csromangathameasuregroup{}{มหาวราหสฺส มหิํ วิกฺรุพฺพโต\textsuperscript{1}, \\ ภิสํ ฆสานสฺส\textsuperscript{1} นทีสุ ชคฺคโต. \\ ภิงฺโกว ปงฺกํ อภิภกฺขยิตฺวา, \\ มมานุกฺรุพฺพํ กปโณ มริสฺสตีติ.}
\csromangathameasurewak{มหาวราหสฺส มหิํ วิกฺรุพฺพโต\textsuperscript{1}, \\ ภิสํ ฆสานสฺส\textsuperscript{1} นทีสุ ชคฺคโต. \\ ภิงฺโกว ปงฺกํ อภิภกฺขยิตฺวา, \\ มมานุกฺรุพฺพํ กปโณ มริสฺสตีติ.}
\setlength{\csromangathaleftcol}{0pt}
\csromangathagroup{\csromangathastanza{\csromangathawak{มหาวราหสฺส มหิํ วิกฺรุพฺพโต\csromansharedfootnote{f3da5e30322c556f}{วิกุพฺพโต (สี, สฺยา)}, \\ ภิสํ ฆสานสฺส\csromansharedfootnote{d5ececd8c7a32163}{ฆสมานสฺส (ก)} นทีสุ ชคฺคโต. \\ ภิงฺโกว ปงฺกํ อภิภกฺขยิตฺวา, \\ มมานุกฺรุพฺพํ กปโณ มริสฺสตีติ.}}}

\proseitem{347}{\csromansymbolfootnote{*}{อํ 3. 37 ปิฏฺเฐปิ.}อฏฺฐหิ ภิกฺขเว องฺเคหิ สมนฺนาคโต ภิกฺขุ ทูเตยฺยํ คนฺตุมรหติ. กตเมหิ อฏฺฐหิ. อิธ ภิกฺขเว ภิกฺขุ โสตา จ โหติ, สาเวตา จ, อุคฺคเหตา จ, ธาเรตา จ วิญฺญาตา จ, วิญฺญาเปตา จ, กุสโล จ สหิตา\-สหิตสฺส, โน จ กลหการโก. อิเมหิ โข ภิกฺขเว อฏฺฐหงฺเคหิ สมนฺนาคโต ภิกฺขุ ทูเตยฺยํ คนฺตุมรหติ.}

\prose{\csromanfolio{365}\csromansymbolfootnote{*}{อํ 3. 37 ปิฏฺเฐปิ. 1. พฺยาธติ (สี, สฺยา) 2. อขาตา (ก) + อํ 3. 10 ปิฏฺเฐปิ.}อฏฺฐหิ ภิกฺขเว องฺเคหิ สมนฺนาคโต สาริปุตฺโต ทูเตยฺยํ คนฺตุมรหติ. กตเมหิ อฏฺฐหิ. อิธ ภิกฺขเว สาริปุตฺโต โสตา จ โหติ, สาเวตา จ, อุคฺคเหตา จ, ธาเรตา จ, วิญฺญาตา จ, วิญฺญาเปตา จ, กุสโล จ สหิตา\-สหิตสฺส, โน จ กลหการโก. อิเมหิ โข ภิกฺขเว อฏฺฐหงฺเคหิ สมนฺนาคโต สาริปุตฺโต ทูเตยฺยํ คนฺตุมรหตีติ.}

\setlength{\csromangathapagewidth}{0pt}
\setlength{\csromangathawakwidth}{0pt}
\csromangathameasuregroup{\textsuperscript{*}โย เว น พฺยถติ1 ปตฺวา, \\ น จ หาเปติ วจนํ, \\ อสนฺทิทฺโธ จ อกฺขาติ2, \\ ส เว ตาทิสโก ภิกฺขุ,}{\csromangathabat{\textsuperscript{*}โย เว น พฺยถติ1 ปตฺวา,}{ปริสํ อุคฺควาทินิํ.} \\ \csromangathabat{น จ หาเปติ วจนํ,}{น จ ฉาเทติ สาสนํ.} \\ \csromangathabat{อสนฺทิทฺโธ จ อกฺขาติ2,}{ปุจฺฉิโต จ น กุปฺปติ.} \\ \csromangathabat{ส เว ตาทิสโก ภิกฺขุ,}{ทูเตยฺยํ คนฺตุมรหตีติ.}}
\csromangathasetleft{\textsuperscript{*}โย เว น พฺยถติ1 ปตฺวา, \\ น จ หาเปติ วจนํ, \\ อสนฺทิทฺโธ จ อกฺขาติ2, \\ ส เว ตาทิสโก ภิกฺขุ,}
\csromangathagroup{\csromangathastanza{\csromangathabat{\csromansymbolmark{*}โย เว น พฺยถติ1 ปตฺวา,}{ปริสํ อุคฺควาทินิํ.} \\ \csromangathabat{น จ หาเปติ วจนํ,}{น จ ฉาเทติ สาสนํ.}}\csromangathastanza{\csromangathabat{อสนฺทิทฺโธ จ อกฺขาติ2,}{ปุจฺฉิโต จ น กุปฺปติ.} \\ \csromangathabat{ส เว ตาทิสโก ภิกฺขุ,}{ทูเตยฺยํ คนฺตุมรหตีติ.}}}

\proseitem{348}{\csromansymbolmark{+}อฏฺฐหิ ภิกฺขเว อสทฺธมฺเมหิ อภิภูโต ปริยาทินฺน\-จิตฺโต เทวทตฺโต อาปายิโก เนรยิโก กปฺปฏฺโฐ อเตกิจฺโฉ. กตเมหิ อฏฺฐหิ. ลาเภน ภิกฺขเว อภิภูโต ปริยาทินฺน\-จิตฺโต เทวทตฺโต อาปายิโก เนรยิโก กปฺปฏฺโฐ อเตกิจฺโฉ. อลาเภน ภิกฺขเว ฯเปฯ. ยเสน ภิกฺขเว ฯเปฯ. อยเสน ภิกฺขเว ฯเปฯ. สกฺกาเรน ภิกฺขเว ฯเปฯ. อสกฺกาเรน ภิกฺขเว ฯเปฯ. ปาปิจฺฉตาย ภิกฺขเว ฯเปฯ. ปาปมิตฺตตาย ภิกฺขเว อภิภูโต ปริยาทินฺน\-จิตฺโต เทวทตฺโต อาปายิโก เนรยิโก กปฺปฏฺโฐ อเตกิจฺโฉ. อิเมหิ โข ภิกฺขเว อฏฺฐหิ อสทฺธมฺเมหิ อภิภูโต ปริยาทินฺน\-จิตฺโต เทวทตฺโต อาปายิโก เนรยิโก กปฺปฏฺโฐ อเตกิจฺโฉ.}

\proseitem{349}{สาธุ ภิกฺขเว ภิกฺขุ อุปฺปนฺนํ ลาภํ อภิภุยฺย อภิภุยฺย วิหเรยฺย. อุปฺปนฺนํ อลาภํ ฯเปฯ. อุปฺปนฺนํ ยสํ. อุปฺปนฺนํ อยสํ. อุปฺปนฺนํ สกฺการํ. อุปฺปนฺนํ อสกฺการํ. อุปฺปนฺนํ ปาปิจฺฉตํ. อุปฺปนฺนํ ปาปมิตฺตตํ อภิภุยฺย อภิภุยฺย วิหเรยฺย. กถญฺจ ภิกฺขเว ภิกฺขุ อตฺถวสํ ปฏิจฺจ อุปฺปนฺนํ ลาภํ อภิภุยฺย อภิภุยฺย วิหเรยฺย. อุปฺปนฺนํ อลาภํ ฯเปฯ. อุปฺปนฺนํ ยสํ. อุปฺปนฺนํ อยสํ. อุปฺปนฺนํ สกฺการํ. อุปฺปนฺนํ อสกฺการํ. อุปฺปนฺนํ ปาปิจฺฉตํ. อุปฺปนฺนํ ปาปมิตฺตตํ อภิภุยฺย อภิภุยฺย วิหเรยฺย. ยํ หิสฺส ภิกฺขเว อุปฺปนฺนํ ลาภํ อนภิภุยฺย วิหรโต อุปฺปชฺเชยฺยุํ อาสวา วิฆาต\-ปริฬาหา, อุปฺปนฺนํ ลาภํ อภิภุยฺย \csromanfolio{366}อภิภุยฺย วิหรโต เอวํสเต อาสวา วิฆาต\-ปริฬาหา น โหนฺติ. ยํ หิสฺส ภิกฺขเว อุปฺปนฺนํ อลาภํ ฯเปฯ อุปฺปนฺนํ ยสํ. อุปฺปนฺนํ อยสํ. อุปฺปนฺนํ สกฺการํ. อุปฺปนฺนํ อสกฺการํ. อุปฺปนฺนํ ปาปิจฺฉตํ. อุปฺปนฺนํ ปาปมิตฺตตํ อนภิภุยฺย วิหรโต อุปฺปชฺเชยฺยุํ อาสวา วิฆาต\-ปริฬาหา, อุปฺปนฺนํ ปาปมิตฺตตํ อภิภุยฺย อภิภุยฺย วิหรโต เอวํสเต อาสวา วิฆาต\-ปริฬาหา น โหนฺติ. อิทํ โข ภิกฺขเว อตฺถวสํ ปฏิจฺจ อุปฺปนฺนํ ลาภํ อภิภุยฺย อภิภุยฺย วิหเรยฺย. อุปฺปนฺนํ อลาภํ ฯเปฯ อุปฺปนฺนํ ยสํ. อุปฺปนฺนํ อยสํ. อุปฺปนฺนํ สกฺการํ. อุปฺปนฺนํ อสกฺการํ. อุปฺปนฺนํ ปาปิจฺฉตํ. อุปฺปนฺนํ ปาปมิตฺตตํ อภิภุยฺย อภิภุยฺย วิหเรยฺย. ตสฺมา ติห ภิกฺขเว อุปฺปนฺนํ ลาภํ อภิภุยฺย อภิภุยฺย วิหริสฺสาม. อุปฺปนฺนํ อลาภํ ฯเปฯ อุปฺปนฺนํ ยสํ. อุปฺปนฺนํ อยสํ. อุปฺปนฺนํ สกฺการํ. อุปฺปนฺนํ อสกฺการํ. อุปฺปนฺนํ ปาปิจฺฉตํ. อุปฺปนฺนํ ปาปมิตฺตตํ อภิภุยฺย อภิภุยฺย วิหริสฺสามาติ. เอวญฺหิ โว ภิกฺขเว สิกฺขิตพฺพนฺติ.}

\proseitem{350}{ตีหิ ภิกฺขเว อสทฺธมฺเมหิ อภิภูโต ปริยาทินฺน\-จิตฺโต เทวทตฺโต อาปายิโก เนรยิโก กปฺปฏฺโฐ อเตกิจฺโฉ. กตเมหิ ตีหิ. ปาปิจฺฉตา, ปาปมิตฺตตา, โอรมตฺตเกน วิเสสา\-ธิคเมน อนฺตรา โวสานํ อาปาทิ. อิเมหิ โข ภิกฺขเว ตีหิ อสทฺธมฺเมหิ อภิภูโต ปริยาทินฺน\-จิตฺโต เทวทตฺโต อาปายิโก เนรยิโก กปฺปฏฺโฐ อเตกิจฺโฉติ.}

\setlength{\csromangathapagewidth}{0pt}
\setlength{\csromangathawakwidth}{0pt}
\csromangathameasuregroup{มา ชาตุ โกจิ โลกสฺมิํ, \\ ตทมินาปิ ชานาถ, \\ ปณฺฑิโตติ สมญฺญาโต, \\ ชลํว ยสสา อฏฺฐา, \\ โส ปมาทํ อนุจิณฺโณ, \\ อวีจินิรยํ ปตฺโต, \\ อทุฏฺฐสฺส หิ โย ทุพฺเภ, \\ ตเมว ปาปํ ผุสติ, \\ สมุทฺทํ วิสกุมฺเภน, \\ น โส เตน ปทูเสยฺย, \\ เอวเมว ตถาคตํ, \\ สมคฺคตํ\textsuperscript{1} สนฺตจิตฺตํ, \\ ตาทิสํ มิตฺตํ กฺรุพฺเพถ\textsuperscript{1}, \\ ยสฺส มคฺคานุโค ภิกฺขุ,}{\csromangathabat{มา ชาตุ โกจิ โลกสฺมิํ,}{ปาปิจฺโฉ อุทปชฺชถ.} \\ \csromangathabat{ตทมินาปิ ชานาถ,}{ปาปิจฺฉานํ ยถาคติ.} \\ \csromangathabat{ปณฺฑิโตติ สมญฺญาโต,}{ภาวิตตฺโตติ สมฺมโต.} \\ \csromangathabat{ชลํว ยสสา อฏฺฐา,}{เทวทตฺโตติ เม สุตํ.} \\ \csromangathabat{โส ปมาทํ อนุจิณฺโณ,}{อาสชฺช นํ ตถาคตํ.} \\ \csromangathabat{อวีจินิรยํ ปตฺโต,}{จตุทฺวารํ ภยานกํ.} \\ \csromangathabat{อทุฏฺฐสฺส หิ โย ทุพฺเภ,}{ปาปกมฺมํ อกฺรุพฺพโต.} \\ \csromangathabat{ตเมว ปาปํ ผุสติ,}{ทุฏฺฐจิตฺตํ อนาทรํ.} \\ \csromangathabat{สมุทฺทํ วิสกุมฺเภน,}{โย มญฺเญยฺย ปทูสิตุํ\textsuperscript{1}.} \\ \csromangathabat{น โส เตน ปทูเสยฺย,}{เภสฺมา หิ อุทธี มหา.} \\ \csromangathabat{เอวเมว ตถาคตํ,}{โย วาเทนุปหิํสติ.} \\ \csromangathabat{สมคฺคตํ\textsuperscript{1} สนฺตจิตฺตํ,}{วาโท ตมฺหิ น รูหติ.} \\ \csromangathabat{ตาทิสํ มิตฺตํ กฺรุพฺเพถ\textsuperscript{1},}{ตญฺจ เสเวถ ปณฺฑิโต.} \\ \csromangathabat{ยสฺส มคฺคานุโค ภิกฺขุ,}{ขยํ ทุกฺขสฺส ปาปุเณติ.}}
\csromangathasetleft{มา ชาตุ โกจิ โลกสฺมิํ, \\ ตทมินาปิ ชานาถ, \\ ปณฺฑิโตติ สมญฺญาโต, \\ ชลํว ยสสา อฏฺฐา, \\ โส ปมาทํ อนุจิณฺโณ, \\ อวีจินิรยํ ปตฺโต, \\ อทุฏฺฐสฺส หิ โย ทุพฺเภ, \\ ตเมว ปาปํ ผุสติ, \\ สมุทฺทํ วิสกุมฺเภน, \\ น โส เตน ปทูเสยฺย, \\ เอวเมว ตถาคตํ, \\ สมคฺคตํ\textsuperscript{1} สนฺตจิตฺตํ, \\ ตาทิสํ มิตฺตํ กฺรุพฺเพถ\textsuperscript{1}, \\ ยสฺส มคฺคานุโค ภิกฺขุ,}
\csromangathagroup{\csromangathastanza{\csromangathabat{มา ชาตุ โกจิ โลกสฺมิํ,}{ปาปิจฺโฉ อุทปชฺชถ.} \\ \csromangathabat{ตทมินาปิ ชานาถ,}{ปาปิจฺฉานํ ยถาคติ.}}\csromangathastanza{\csromangathabat{ปณฺฑิโตติ สมญฺญาโต,}{ภาวิตตฺโตติ สมฺมโต.} \\ \csromangathabat{ชลํว ยสสา อฏฺฐา,}{เทวทตฺโตติ เม สุตํ.}}\csromangathastanza{\csromangathabat{โส ปมาทํ อนุจิณฺโณ,}{อาสชฺช นํ ตถาคตํ.} \\ \csromangathabat{อวีจินิรยํ ปตฺโต,}{จตุทฺวารํ ภยานกํ.}}\csromangathastanza{\csromangathabat{อทุฏฺฐสฺส หิ โย ทุพฺเภ,}{ปาปกมฺมํ อกฺรุพฺพโต.} \\ \csromangathabat{ตเมว ปาปํ ผุสติ,}{ทุฏฺฐจิตฺตํ อนาทรํ.}}\csromangathastanza{\csromangathabat{สมุทฺทํ วิสกุมฺเภน,}{โย มญฺเญยฺย ปทูสิตุํ\csromansharedfootnote{25784d358fd6e13b}{ปทุสฺสิตุํ (ก)}.} \\ \csromangathabat{น โส เตน ปทูเสยฺย,}{เภสฺมา หิ อุทธี มหา.}}\csromangathastanza{\csromanfolio{367}\csromangathabat{เอวเมว ตถาคตํ,}{โย วาเทนุปหิํสติ.} \\ \csromangathabat{สมคฺคตํ\csromansharedfootnote{5658e3a8976323ad}{สมฺมาคตํ (สี), สมคตํ (สฺยา)} สนฺตจิตฺตํ,}{วาโท ตมฺหิ น รูหติ.}}\csromangathastanza{\csromangathabat{ตาทิสํ มิตฺตํ กฺรุพฺเพถ\csromansharedfootnote{340171264ace26f8}{กุพฺเพถ (สี, สฺยา)},}{ตญฺจ เสเวถ ปณฺฑิโต.} \\ \csromangathabat{ยสฺส มคฺคานุโค ภิกฺขุ,}{ขยํ ทุกฺขสฺส ปาปุเณติ.}}}

\csromanmatikaanchor{mtk.174}
\csromantocmarkref{subsection}{อุปาลิปญฺหา}{mtk.174}
\bookmark[dest={mtk.174},level=2]{อุปาลิปญฺหา}
\csromanheader{2}{\textbf{อุปาลิปญฺหา}}
\proseitem{351}{อถ โข อายสฺมา อุปาลิ เยน ภควา เตนุปสงฺกมิ, อุปสงฺกมิตฺวา ภควนฺตํ อภิวาเทตฺวา เอกมนฺตํ นิสีทิ, เอกมนฺตํ นิสินฺโน โข อายสฺมา อุปาลิ ภควนฺตํ เอตทโวจ “สํฆราชิ สํฆราชี”ติ ภนฺเต วุจฺจติ, กิตฺตาวตา นุ โข ภนฺเต สํฆราชิ โหติ, โน จ สํฆเภโท. กิตฺตาวตา จ ปน สํฆราชิ เจว โหติ สํฆเภโท จา”ติ.}

\prose{เอกโต อุปาลิ เอโก โหติ, เอกโต เทฺว, จตุตฺโถ อนุสฺสาเวติ, สลากํ คาเหติ “อยํ ธมฺโม, อยํ วินโย, อิทํ สตฺถุสาสนํ, อิมํ คณฺหถ, อิมํ โรเจถา”ติ. เอวมฺปิ โข อุปาลิ สํฆราชิ โหติ, โน จ สํฆเภโท. เอกโต อุปาลิ เทฺว โหนฺติ, เอกโต เทฺว. ปญฺจโม อนุสฺสาเวติ, สลากํ คาเหติ “อยํ ธมฺโม, อยํ วินโย, อิทํ สตฺถุสาสนํ, อิมํ คณฺหถ, อิมํ โรเจถา”ติ. เอวมฺปิ โข อุปาลิ สํฆราชิ โหติ, โน จ สํฆเภโท. เอกโต อุปาลิ เทฺว โหนฺติ. เอกโต ตโย, ฉฏฺโฐ อนุสฺสาเวติ, สลากํ คาเหติ “อยํ ธมฺโม, อยํ วินโย, อิทํ สตฺถุสาสนํ, อิมํ คณฺหถ, อิมํ โรเจถา”ติ. เอวมฺปิ โข อุปาลิ สํฆราชิ โหติ, โน จ สํฆเภโท. เอกโต อุปาลิ ตโย โหนฺติ. เอกโต ตโย, สตฺตโม อนุสฺสาเวติ, สลากํ คาเหติ “อยํ ธมฺโม, อยํ วินโย, อิทํ สตฺถุสาสนํ, อิมํ คณฺหถ, อิมํ โรเจถา”ติ. เอวมฺปิ โข อุปาลิ สํฆราชิ โหติ, โน จ สํฆเภโท. เอกโต อุปาลิ ตโย โหนฺติ. เอกโต จตฺตาโร, อฏฺฐโม อนุสฺสาเวติ, สลากํ คาเหติ “อยํ ธมฺโม, อยํ วินโย, อิทํ สตฺถุสาสนํ, อิมํ คณฺหถ, อิมํ โรเจถา”ติ. เอวมฺปิ โข อุปาลิ สํฆราชิ โหติ, โน จ สํฆเภโท. เอกโต อุปาลิ จตฺตาโร โหนฺติ, เอกโต จตฺตาโร, นวโม อนุสฺสาเวติ, สลากํ คาเหติ “อยํ ธมฺโม, อยํ วินโย, อิทํ สตฺถุสาสนํ, อิมํ คณฺหถ, อิมํ โรเจถา”ติ. \csromanfolio{368}เอวํ โข อุปาลิ สํฆราชิ เจว โหติ สํฆเภโท จ. นวนฺนํ วา อุปาลิ อติเรก\-นวนฺนํ วา สํฆราชิ เจว โหติ สํฆเภโท จ. น โข อุปาลิ ภิกฺขุนี สํฆํ ภินฺทติ, อปิ จ เภทาย ปรกฺกมติ. น สิกฺขมานา สํฆํ ภินฺทติ ฯเปฯ. น สามเณโร สํฆํ ภินฺทติ. น สามเณรี สํฆํ ภินฺทติ. น อุปาสโก สํฆํ ภินฺทติ. น อุปาสิกา สํฆํ ภินฺทติ, อปิ จ เภทาย ปรกฺกมติ. ภิกฺขุ โข อุปาลิ ปกตตฺโต สมาน\-สํวาสโก สมานสีมายํ ฐิโต สํฆํ ภินฺทตีติ.}

\proseitem{352}{\csromansymbolfootnote{*}{อํ 3. 314 ปิฏฺเฐทีสุปิ.}“สํฆเภโท สํฆเภโท”ติ ภนฺเต วุจฺจติ, กิตฺตาวตา นุ โข ภนฺเต สํโฆ ภินฺโน โหตีติ.}

\prose{อิธุปาลิ ภิกฺขู อธมฺมํ “ธมฺโม”ติ ทีเปนฺติ, ธมฺมํ “อธมฺโม”ติ ทีเปนฺติ, อวินยํ “วินโย”ติ ทีเปนฺติ, วินยํ “อวินโย”ติ ทีเปนฺติ, อภาสิตํ อลปิตํ ตถาคเตน “ภาสิตํ ลปิตํ ตถาคเตนา”ติ ทีเปนฺติ, ภาสิตํ ลปิตํ ตถาคเตน “อภาสิตํ อลปิตํ ตถาคเตนา”ติ ทีเปนฺติ, อนาจิณฺณํ ตถาคเตน “อาจิณฺณํ ตถาคเตนา”ติ ทีเปนฺติ, อาจิณฺณํ ตถาคเตน “อนาจิณฺณํ ตถาคเตนา”ติ ทีเปนฺติ, อปญฺญตฺตํ ตถาคเตน “ปญฺญตฺตํ ตถาคเตนา”ติ ทีเปนฺติ, ปญฺญตฺตํ ตถาคเตน “อปญฺญตฺตํ ตถาคเตนา”ติ ทีเปนฺติ, อนาปตฺติํ “อาปตฺตี”ติ ทีเปนฺติ, อาปตฺติํ “อานาปตฺตี”ติ ทีเปนฺติ, ลหุกํ อาปตฺติํ “ครุกา อาปตฺตี”ติ ทีเปนฺติ, ครุกํ อาปตฺติํ “ลหุกา อาปตฺตี”ติ ทีเปนฺติ, สาวเสสํ อาปตฺติํ “อนวเสสา อาปตฺตี”ติ ทีเปนฺติ, อนวเสสํ อาปตฺติํ “สาวเสสา อาปตฺตี”ติ ทีเปนฺติ, ทุฏฺฐุลฺลํ อาปตฺติํ “อทุฏฺฐุลฺลา อาปตฺตี”ติ ทีเปนฺติ, อทุฏฺฐุลฺลํ อาปตฺติํ “ทุฏฺฐุลฺลา อาปตฺตี”ติ ทีเปนฺติ. เต อิเมหิ อฏฺฐารสหิ วตฺถูหิ อปกสฺสนฺติ, อวปกสฺสนฺติ, อาเวนิํ\csromansharedfootnote{8b93ddb9622a98e2}{อาเวณิ (สี), อาเวณิกํ (สฺยา)} อุโปสถํ กโรนฺติ, อาเวนิํ ปวารณํ กโรนฺติ, อาเวนิํ สํฆกมฺมํ กโรนฺติ. เอตฺตาวตา โข อุปาลิ สํโฆ ภินฺโน โหตีติ.}

\proseitem{353}{\csromansymbolmark{*}“สํฆสามคฺคี สํฆสามคฺคี”ติ ภนฺเต วุจฺจติ, กิตฺตาวตา นุ โข ภนฺเต สํโฆ สมคฺโค โหตีติ. อิธุปาลิ ภิกฺขู อธมฺมํ “อธมฺโม”ติ ทีเปนฺติ, ธมฺมํ “ธมฺโม”ติ ทีเปนฺติ, อวินยํ “อวินโย”ติ ทีเปนฺติ, วินยํ \csromanfolio{369}“วินโย”ติ ทีเปนฺติ, อภาสิตํ อลปิตํ ตถาคเตน “อภาสิตํ อลปิตํ ตถาคเตนา”ติ ทีเปนฺติ, ภาสิตํ ลปิตํ ตถาคเตน “ภาสิตํ ลปิตํ ตถาคเตนา”ติ ทีเปนฺติ, อนาจิณฺณํ ตถาคเตน “อนาจิณฺณํ ตถาคเตนา”ติ ทีเปนฺติ, อาจิณฺณํ ตถาคเตน “อาจิณฺณํ ตถาคเตนา”ติ ทีเปนฺติ, อปญฺญตฺตํ ตถาคเตน “อปญฺญตฺตํ ตถาคเตนา”ติ ทีเปนฺติ, ปญฺญตฺตํ ตถาคเตน “ปญฺญตฺตํ ตถาคเตนา”ติ ทีเปนฺติ, อนาปตฺติํ “อนาปตฺตี”ติ ทีเปนฺติ, อาปตฺติํ “อาปตฺตี”ติ ทีเปนฺติ, ลหุกํ อาปตฺติํ “ลหุกา อาปตฺตี”ติ ทีเปนฺติ, ครุกํ อาปตฺติํ “ครุกา อาปตฺตี”ติ ทีเปนฺติ, สาวเสสํ อาปตฺติํ “สาวเสสา อาปตฺตี”ติ ทีเปนฺติ, อนวเสสํ อาปตฺติํ “อนวเสสา อาปตฺตี”ติ ทีเปนฺติ, ทุฏฺฐุลฺลํ อาปตฺติํ “ทุฏฺฐุลฺลา อาปตฺตี”ติ ทีเปนฺติ, อทุฏฺฐุลฺลํ อาปตฺติํ “อทุฏฺฐุลฺลา อาปตฺตี”ติ ทีเปนฺติ. เต อิเมหิ อฏฺฐารสหิ วตฺถูหิ น อปกสฺสนฺติ, อวปกสฺสนฺติ, น อาเวนิํ อุโปสถํ กโรนฺติ, น อาเวนิํ ปวารณํ กโรนฺติ, น อาเวนิํ สํฆกมฺมํ กโรนฺติ. เอตฺตาวตา โข อุปาลิ สํโฆ สมคฺโค โหตีติ.}

\proseitem{345}{สมคฺคํ ปน ภนฺเต สํฆํ ภนฺทิตฺวา กิํ โส ปสวตีติ. สมคฺคํ โข อุปาลิ สํฆํ ภินฺทิตฺวา กปฺปฏฺฐิกํ กิพฺพิสํ ปสวติ, กปฺปํ นิรยมฺหิ ปจฺจตีติ.}

\prose{\csromansymbolfootnote{*}{ขุ 1. 203; อํ 3. 315, 316 ปิฏฺเฐสุปิ.}อาปายิโก เนรยิโก, กปฺปฏฺโฐ สํฆเภทโก.}

\setlength{\csromangathapagewidth}{0pt}
\setlength{\csromangathawakwidth}{0pt}
\csromangathameasuregroup{วคฺครโต อธมฺมฏฺโฐ, \\ สํฆํ สมคฺคํ ภินฺทิตฺวา,}{\csromangathabat{วคฺครโต อธมฺมฏฺโฐ,}{โยคกฺเขมา ปธํสติ.} \\ \csromangathabat{สํฆํ สมคฺคํ ภินฺทิตฺวา,}{กปฺปํ นิรยมฺหิ ปจฺจตีติ.}}
\csromangathasetleft{วคฺครโต อธมฺมฏฺโฐ, \\ สํฆํ สมคฺคํ ภินฺทิตฺวา,}
\csromangathagroup{\csromangathastanza{\csromangathabat{วคฺครโต อธมฺมฏฺโฐ,}{โยคกฺเขมา ปธํสติ.} \\ \csromangathabat{สํฆํ สมคฺคํ ภินฺทิตฺวา,}{กปฺปํ นิรยมฺหิ ปจฺจตีติ.}}}

\prose{ภินฺนํ ปน ภนฺเต สํฆํ สมคฺคํ กตฺวา กิํ โส ปสวตีติ. ภินฺนํ โข อุปาลิ สํฆํ สมคฺคํ กตฺวา พฺรหฺมํ ปุญฺญํ ปสวติ, กปฺปํ สคฺคมฺหิ โมทตีติ.}

\prose{\csromansymbolmark{*}สุขา สํฆสฺส สามคฺคี, สมคฺคานญฺจ อนุคฺคโห.}

\setlength{\csromangathapagewidth}{0pt}
\setlength{\csromangathawakwidth}{0pt}
\csromangathameasuregroup{สมคฺครโต ธมฺมฏฺโฐ, \\ สํฆํ สมคฺคํ กตฺวาน,}{\csromangathabat{สมคฺครโต ธมฺมฏฺโฐ,}{โยคกฺเขมา น ธํสติ.} \\ \csromangathabat{สํฆํ สมคฺคํ กตฺวาน,}{กปฺปํ สคฺคมฺหิ โมทตีติ.}}
\csromangathasetleft{สมคฺครโต ธมฺมฏฺโฐ, \\ สํฆํ สมคฺคํ กตฺวาน,}
\csromangathagroup{\csromangathastanza{\csromangathabat{สมคฺครโต ธมฺมฏฺโฐ,}{โยคกฺเขมา น ธํสติ.} \\ \csromangathabat{สํฆํ สมคฺคํ กตฺวาน,}{กปฺปํ สคฺคมฺหิ โมทตีติ.}}}

\proseitem{355}{\csromansymbolfootnote{+}{วิ 5. 350 ปิฏฺเฐปิ.}สิยา นุ โข ภนฺเต สํฆเภทโก อาปายิโก เนรยิโก กปฺปฏฺโฐ อเตกิจฺโฉติ. สิยา อุปาลิ สํฆเภทโก อาปายิโก เนรยิโก กปฺปฏฺโฐ อเตกิจฺโฉติ.}

\prose{\csromanfolio{370}\csromansymbolfootnote{*}{วิ 5. 350 ปิฏฺเฐปิ.}สิยา\csromansharedfootnote{977edc8700f78998}{สิยา นุ โข (สฺยา, กํ)} ปน ภนฺเต สํฆเภทโก น อาปายิโก น เนรยิโก น กปฺปฏฺโฐ น อเตกิจฺโฉติ. สิยา อุปาลิ สํฆเภทโก น อาปายิกา น เนรยิโก น กปฺปฏฺโฐ น อเตกิจฺโฉติ.}

\prose{กตโม ปน ภนฺเต สํฆเภทโก อาปายิโก เนรยิโก กปฺปฏฺโฐ อเตกิจฺโฉติ. อิธุปาลิ ภิกฺขุ อธมฺมํ “ธมฺโม”ติ ทีเปติ. ตสฺมิํ อธมฺมทิฏฺฐิ เภเท อธมฺมทิฏฺฐิ วินิธาย ทิฏฺฐิํ วินิธาย ขนฺติํ วินิธาย รุจิํ วินิธาย ภาวํ อนุสฺสาเวติ, สลากํ คาเหติ “อยํ ธมฺโม, อยํ วินโย, อิทํ สตฺถุสาสนํ, อิมํ คณฺหถ, อิมํ โรเจตา”ติ. อยมฺปิ โข อุปาลิ สํฆเภทโก อาปายิโก เนรยิโก กปฺปฏฺโฐ อเตกิจฺโฉ.}

\prose{ปุน จปรํ อุปาลิ ภิกฺขุ อธมฺมํ “ธมฺโม”ติ ทีเปติ, ตสฺมิํ อธมฺมทิฏฺฐิ เภเท ธมฺมทิฏฺฐิ วินิธาย ทิฏฺฐิํ วินิธาย ขนฺติํ วินิธาย รุจิํ วินิธาย ภาวํ อนุสฺสาเวติ, สลากํ คาเหติ “อยํ ธมฺโม, อยํ วินโย, อิทํ สตฺถุสาสนํ, อิมํ คณฺหถ, อิมํ โรเจถา”ติ. อยมฺปิ โข อุปาลิ สํฆเภทโก อาปายิโก เนรยิโก กปฺปฏฺโฐ อเตกิจฺโฉ.}

\prose{ปุน จปรํ อุปาลิ ภิกฺขุ อธมฺมํ “ธมฺโม”ติ ทีเปติ, ตสฺมิํ อธมฺมทิฏฺฐิ เภเท เวมติโก วินิธาย ทิฏฺฐิํ วินิธาย ขนฺติํ วินิธาย รุจิํ วินิธาย ภาวํ อนุสฺสาเวติ, สลากํ คาเหติ “อยํ ธมฺโม, อยํ วินโย, อิทํ สตฺถุสาสนํ, อิมํ คณฺหถ, อิมํ โรเจถา”ติ. อยมฺปิ โข อุปาลิ สํฆเภทโก อาปายิโก เนรยิโก กปฺปฏฺโฐ อเตกิจฺโฉ.}

\prose{ปุน จปรํ อุปาลิ ภิกฺขุ อธมฺมํ “ธมฺโม”ติ ทีเปติ, ตสฺมิํ ธมฺมทิฏฺฐิ เภเท อธมฺมทิฏฺฐิ ฯเปฯ (ตสฺมิํ ธมฺมทิฏฺฐิ เภเท ธมฺมทิฏฺฐิ)\csromansharedfootnote{e932c7f13211565f}{( ) สฺยามโปตฺถเก นตฺถิ, วิมติวิโนทนีฏีกาย สเมติ.}. ตสฺมิํ ธมฺมทิฏฺฐิ เภเท เวมติโก. ตสฺมิํ เวมติโก เภเท อธมฺมทิฏฺฐิ. ตสฺมิํ เวมติโก เภเท ธมฺมทิฏฺฐิ. ตสฺมิํ เวมติโก เภเท เวมติโก วินิธาย ทิฏฺฐิํ วินิธาย ขนฺติํ วินิธาย รุจิํ วินิธาย ภาวํ อนุสฺสาเวติ, สลากํ คาเหติ “อยํ ธมฺโม, อยํ วินโย, อิทํ สตฺถุสาสนํ, อิมํ คณฺหถ, อิมํ โรเจถา”ติ. อยมฺปิ โข อุปาลิ สํฆเภทโก อาปายิโก เนรยิโก กปฺปฏฺโฐ อเตกิจฺโฉ.}

\prose{\csromanfolio{371}ปุน จปรํ อุปาลิ ภิกฺขุ ธมฺมํ “อธมฺโม”ติ ทีเปติ ฯเปฯ อวินยํ “วินโย”ติ ทีเปติ. วินยํ “อวินโย”ติ ทีเปติ. อภายิตํ อลปิตํ ตถาคเตน “ภาสิตํ ลปิตํ ตถาคเตนา”ติ ทีเปติ. ภาสิตํ ลปิตํ ตถาคเตน “อภาสิตํ อลปิตํ ตถาคเตนา”ติ ทีเปติ. อนาจิณฺณํ ตถาคเตน “อาจิณฺณํ ตถาคเตนา”ติ ทีเปติ. อาจิณฺณํ ตถาคเตน “อนาจิณฺณํ ตถาคเตนา”ติ ทีเปติ. อปญฺญตฺตํ ตถาคเตน “ปญฺญตฺตํ ตถาคเตนา”ติ ทีเปติ. ปญฺญตฺตํ ตถาคเตน “อปญฺญตฺตํ ตถาคเตนา”ติ ทีเปติ. อนาปตฺติํ “อาปตฺตี”ติ ทีเปติ. อาปตฺติํ “อนาปตฺตี”ติ ทีเปติ. ลหุกํ อาปตฺติํ “ครุกา อาปตฺตี”ติ ทีเปติ. ครุกํ อาปตฺติํ “ลหุกา อาปตฺตี”ติ ทีเปติ. สาวเสสํ อาปตฺติํ “อนวเสสา อาปตฺตี”ติ ทีเปติ. อนวเสสํ อาปตฺติํ “สาวเสสา อาปตฺตี”ติ ทีเปติ. ทุฏฺฐุลฺลํ อาปตฺติํ “อทุฏฺฐุลฺลา อาปตฺตี”ติ ทีเปติ. อทุฏฺฐุลฺลํ อาปตฺติํ “ทุฏฺฐุลฺลา อาปตฺตี”ติ ทีเปติ. ตสฺมิํ อธมฺมทิฏฺฐิ เภเท อธมฺมทิฏฺฐิ ฯเปฯ. ตสฺมิํ อธมฺมทิฏฺฐิ เภเท ธมฺมทิฏฺฐิ. ตสฺมิํ อธมฺมทิฏฺฐิ เภเท เวมติโก. ตสฺมิํ ธมฺมทิฏฺฐิ เภเท อธมฺมทิฏฺฐิ. (ตสฺมิํ ธมฺมทิฏฺฐิ เภเท ธมฺมทิฏฺฐิ)\csromansharedfootnote{ded81c193db1fcf4}{( ) สฺยามโปตฺถเก นตฺถิ.}. ตสฺมิํ ธมฺมทิฏฺฐิ เภเท เวมติโก. ตสฺมิํ เวมติโก เภเท อธมฺมทิฏฺฐิ. ตสฺมิํ เวมติโก เภเท ธมฺมทิฏฺฐิ. ตสฺมิํ เวมติโก เภเท เวมติโก วินิธาย ทิฏฺฐิํ วินิธาย ขนฺติํ วินิธาย รุจิํ วินิธาย ภาวํ อนุสฺสาเวติ, สลากํ คาเหติ “อยํ ธมฺโม, อยํ วินโย, อิทํ สตฺถุสาสนํ, อิมํ คณฺหถ, อิมํ โรเจถา”ติ. อยมฺปิ โข อุปาลิ สํฆเภทโก อาปายิโก เนรยิโก กปฺปฏฺโฐ อเตกิจฺโฉติ.}

\prose{กตโม ปน ภนฺเต สํฆเภทโก น อาปายิโก น เนรยิโก น กปฺปฏฺโฐ น อเตกิจฺโฉติ. อิธุปาลิ ภิกฺขุ อธมฺมํ “ธมฺโม”ติ ทีเปติ, ตสฺมิํ ธมฺมทิฏฺฐิ เภเท ธมฺมทิฏฺฐิ อวินิธาย ทิฏฺฐิํ อวินิธาย ขนฺติํ อวินิธาย รุจิํ อวินิธาย ภาวํ อนุสฺสาเวติ, สลากํ คาเหติ “อยํ ธมฺโม, อยํ วินโย, อิทํ สตฺถุสาสนํ, อิมํ คณฺหถ, อิมํ โรเจถา”ติ. อยมฺปิ โข อุปาลิ สํฆเภทโก น อาปายิโก น เนรยิโก น กปฺปฏฺโฐ น อเตกิจฺโฉ.}

\csromanmatikaanchor{mtk.175}
\csromantocmarkref{section}{อุทฺทานคาถา}{mtk.175}
\bookmark[dest={mtk.175},level=1]{อุทฺทานคาถา}
\prosewithcloser{\csromanfolio{372}ปุน จปรํ อุปาลิ ภิกฺขุ ธมฺมํ “อธมฺโม”ติ ทีเปติ ฯเปฯ ทุฏฺฐุลฺลํ อาปตฺติํ “อทุฏฺฐุลฺลา อาปตฺตี”ติ ทีเปติ, ตสฺมิํ ธมฺมทิฏฺฐิ เภเท ธมฺมทิฏฺฐิ อวินิธาย ทิฏฺฐิํ อวินิธาย ขนฺติํ อวินิธาย รุจิํ อวินิธาย ภาวํ อนุสฺสาเวติ, สลากํ คาเหติ “อยํ ธมฺโม, อยํ วินโย, อิทํ สตฺถุสาสนํ, อิมํ คณฺหถ, อิมํ โรเจถา”ติ. อยมฺปิ โข อุปาลิ สํฆเภทโก น อาปายิโก น เนรยิโก น กปฺปฏฺโฐ น อเตกิจฺโฉติ.}{ตติย\-ภาณวาโร นิฏฺฐิโต.}
\csromancenter{สํฆเภทก\-กฺขนฺธโก สตฺตโม.}
\csromancenter{\textbf{ตสฺสุทฺทานํ}}
\setlength{\csromangathapagewidth}{0pt}
\setlength{\csromangathawakwidth}{0pt}
\csromangathameasuregroup{อนุปิเย อภิญฺญาตา, \\ กสา วปา อภิ นินฺเน, \\ ปุญฺชมทฺทปลาลญฺจ, \\ อายติมฺปิ น ขียนฺติ, \\ ภทฺทิโย อนุรุทฺโธ จ, \\ สกฺยมาโน จ โกสมฺพิํ, \\ ปกาเสสิ ปิตุโน จ, \\ ติกปญฺจครุโก โข, \\ ตโย อฏฺฐ ปุน ตีณิ,}{\csromangathabat{อนุปิเย อภิญฺญาตา,}{สุขุมาโล น อิจฺฉติ.} \\ \csromangathabat{กสา วปา อภิ นินฺเน,}{นิทฺธา ลาเว จ อุพฺพเห.} \\ \csromangathabat{ปุญฺชมทฺทปลาลญฺจ,}{ภุสโอผุณนีหเร.} \\ \csromangathabat{อายติมฺปิ น ขียนฺติ,}{ปิตโร จ ปิตามหา.} \\ \csromangathabat{ภทฺทิโย อนุรุทฺโธ จ,}{อานนฺโท ภคุ กิมิโล.} \\ \csromangathabat{สกฺยมาโน จ โกสมฺพิํ,}{ปริหายิ กกุเธน จ.} \\ \csromangathabat{ปกาเสสิ ปิตุโน จ,}{ปุริเส สิลํ นาฬาคิริํ.} \\ \csromangathabat{ติกปญฺจครุโก โข,}{ภินฺทิ ถุลฺลจฺจเยน จ.} \\ \csromangathabat{ตโย อฏฺฐ ปุน ตีณิ,}{ราชิ เภทา สิยา นุ โขติ.}}
\csromangathasetleft{อนุปิเย อภิญฺญาตา, \\ กสา วปา อภิ นินฺเน, \\ ปุญฺชมทฺทปลาลญฺจ, \\ อายติมฺปิ น ขียนฺติ, \\ ภทฺทิโย อนุรุทฺโธ จ, \\ สกฺยมาโน จ โกสมฺพิํ, \\ ปกาเสสิ ปิตุโน จ, \\ ติกปญฺจครุโก โข, \\ ตโย อฏฺฐ ปุน ตีณิ,}
\csromangathagroupwithclosermajor{\csromangathastanza{\csromangathabat{อนุปิเย อภิญฺญาตา,}{สุขุมาโล น อิจฺฉติ.} \\ \csromangathabat{กสา วปา อภิ นินฺเน,}{นิทฺธา ลาเว จ อุพฺพเห.}}\csromangathastanza{\csromangathabat{ปุญฺชมทฺทปลาลญฺจ,}{ภุสโอผุณนีหเร.} \\ \csromangathabat{อายติมฺปิ น ขียนฺติ,}{ปิตโร จ ปิตามหา.}}\csromangathastanza{\csromangathabat{ภทฺทิโย อนุรุทฺโธ จ,}{อานนฺโท ภคุ กิมิโล.} \\ \csromangathabat{สกฺยมาโน จ โกสมฺพิํ,}{ปริหายิ กกุเธน จ.}}\csromangathastanza{\csromangathabat{ปกาเสสิ ปิตุโน จ,}{ปุริเส สิลํ นาฬาคิริํ.} \\ \csromangathabat{ติกปญฺจครุโก โข,}{ภินฺทิ ถุลฺลจฺจเยน จ.} \\ \csromangathabat{ตโย อฏฺฐ ปุน ตีณิ,}{ราชิ เภทา สิยา นุ โขติ.}}}{สํฆเภทก\-กฺขนฺธโก นิฏฺฐิโต.}
\csromanmatikaanchor{mtk.176}
\csromantocmarknopage{chapter}{8. วตฺตกฺขนฺธก}{mtk.176}
\bookmark[dest={mtk.176},level=0]{8. วตฺตกฺขนฺธก}
\chapterheadpageread{8. วตฺต\-กฺขนฺธก}
\csromanmatikaanchor{mtk.177}
\csromantocmarkref{section}{1. อาคนฺตุกวตฺตกถา}{mtk.177}
\bookmark[dest={mtk.177},level=1]{1. อาคนฺตุกวตฺตกถา}
\csromanheader{1}{1. อาคนฺตุกวตฺต\-กถา}
\proseitem{356}{\csromanfolio{373}เตน สมเยน พุทฺโธ ภควา สาวตฺถิยํ วิหรติ เชตวเน อนาถ\-ปิณฺฑิกสฺส อาราเม. เตน โข ปน สมเยน อาคนฺตุกา ภิกฺขู สอุปาหนาปิ อารามํ ปวิสนฺติ, ฉตฺตปคฺคหิตาปิ อารามํ ปวิสนฺติ, โอคุณฺฐิตาปิ อารามํ ปวิสนฺติ, สีเสปิ จีวรํ กริตฺวา อารามํ ปวิสนฺติ, ปานีเยนปิ ปาเท โธวนฺติ, วุฑฺฒตเรปิ อาวาสิเก ภิกฺขู น อภิวาเทนฺติ, นปิ เสนาสนํ ปุจฺฉนฺติ, อญฺญตโรปิ อาคนฺตุโก ภิกฺขุ อนชฺฌาวุฏฺฐํ วิหารํ ฆฏิกํ อุคฺฆาเฏตฺวา กวาฏํ ปณาเมตฺวา สหสา ปาวิสิ, ตสฺส อุปริปิฏฺฐิโต\csromansharedfootnote{67296bb262e100ba}{อุปริปิฏฺฐโต (?)} อหิ ขนฺเธ ปปติ, โส ภีโต วิสฺสรมกาสิ, ภิกฺขู อุปธาวิตฺวา ตํ ภิกฺขุํ เอตทโวจุํ “กิสฺส ตฺวํ อาวุโส วิสฺสรมกาสี”ติ. อถ โข โส ภิกฺขุ ภิกฺขูนํ เอตมตฺถํ อาโรเจสิ. เย เต ภิกฺขู อปฺปิจฺฉา ฯเปฯ เต อุชฺฌายนฺติ ขิยฺยนฺติ วิปาเจนฺติ “กถํ หิ นาม อาคนฺตุกา ภิกฺขู สอุปาหนาปิ อารามํ ปวิสิสฺสนฺติ, ฉตฺตปคฺคหิตาปิ อารามํ ปวิสิสฺสนฺติ, โอคุณฺฐิตาปิ อารามํ ปวิสิสฺสนฺติ, สีเสปิ จีวรํ กริตฺวา อารามํ ปวิสิสฺสนฺติ, ปานีเยนปิ ปาเท โธวิสฺสนฺติ, วุฑฺฒตเรปิ อาวาสิเก ภิกฺขู น อภิวาเทสฺสนฺติ, นปิ เสนาสนํ ปุจฺฉิสฺสนฺตี”ติ. อถ โข เต ภิกฺขู ภควโต เอตมตฺถํ อาโรเจสุํ ฯเปฯ. สจฺจํ กิร ภิกฺขเว อาคนฺตุกา ภิกฺขู สอุปาหนาปิ อารามํ ปวิสนฺติ, ฉตฺตปคฺคหิตาปิ อารามํ ปวิสนฺติ, โอคุณฺฐิตาปิ อารามํ ปวิสนฺติ, สีเสปิ จีวรํ กริตฺวา อารามํ ปวิสนฺติ, ปานีเยนปิ ปาเท โธวนฺติ, วุฑฺฒตเรปิ อาวาสิเก ภิกฺขู น อภิวาเทนฺติ, นปิ เสนาสนํ ปุจฺฉนฺตีติ. สจฺจํ ภควาติ. วิครหิ พุทฺโธ ภควา ฯเปฯ. กถํ หิ นาม ภิกฺขเว อาคนฺตุกา ภิกฺขู สอุปาหนาปิ อารามํ ปวิสิสฺสนฺติ, ฉตฺตปคฺคหิตาปิ อารามํ ปวิสิสฺสนฺติ, โอคุณฺฐิตาปิ อารามํ ปวิสิสฺสนฺติ, สีเสปิ จีวรํ กริตฺวา อารามํ ปวิสิสฺสนฺติ, ปานีเยนปิ ปาเท โธวิสฺสนฺติ, วุฑฺฒตเรปิ อาวาสิเก ภิกฺขู น อภิวาเทสฺสนฺติ, นปิ เสนาสนํ ปุจฺฉิสฺสนฺติ, เนตํ ภิกฺขเว อปฺปสนฺนานํ วา ปสาทาย ฯเปฯ วิครหิตฺวา ฯเปฯ ธมฺมิํ กถํ กตฺวา ภิกฺขู อามนฺเตสิ\csromandash{}}

\proseitem{357}{\csromanfolio{374}เตน หิ ภิกฺขเว อาคนฺตุกานํ ภิกฺขูนํ วตฺตํ ปญฺญเปสฺสามิ, ยถา อาคนฺตุเกหิ ภิกฺขูหิ สมฺมา วตฺติตพฺพํ. อาคนฺตุเกน ภิกฺขเว ภิกฺขุนา “อิทานิ อารามํ ปวิสิสฺสามี”ติ อุปาหนา โอมุญฺจิตฺวา นีจํ กตฺว ปปฺโผเฏตฺวา คเหตฺวา ฉตฺตํ อปนาเมตฺวา สีสํ วิวริตฺวา สีเส จีวรํ\csromansharedfootnote{68a61d26b2ad9c71}{วิวริตฺวา จีวรํ (ก)} ขนฺเธ กตฺวา สาธุกํ อตรมาเนน อาราโม ปวิสิตพฺโพ, อารามํ ปวิสนฺเตน สลฺลกฺเขตพฺพํ “กตฺถ อาวาสิกา ภิกฺขู ปฏิกฺกมนฺตี”ติ. ยตฺถ อาวาสิกา ภิกฺขู ปฏิกฺกมนฺติ อุปฏฺฐาน\-สาลาย วา มณฺฑเป วา รุกฺขมูเล วา, ตตฺถ คนฺตฺวา เอกมนฺตํ ปตฺโต นิกฺขิปิตพฺโพ, เอกมนฺตํ จีวรํ นิกฺขิปิตพฺพํ, ปติรูปํ อาสนํ คเหตฺวา นิสีทิตพฺพํ, ปานียํ ปุจฺฉิตพฺพํ, ปริโภชนียํ ปุจฺฉิตพฺพํ “กตมํ ปานียํ, กตมํ ปริโภชนียนฺ”ติ. สเจ ปานีเยน อตฺโถ โหติ, ปานียํ คเหตฺวา ปาตพฺพํ, สเจ ปริโภชนีเยน อตฺโถ โหติ, ปริโภชนียํ คเหตฺวา ปาทา โธวิตพฺพา, ปาเท โธวนฺเตน เอเกน หตฺเถน อุทกํ อาสิญฺจิตพฺพํ, เอเกน หตฺเถน ปาทา โธวิตพฺพา, น เตเนว หตฺเถน อุทกํ อาสิญฺจิตพฺพํ\csromansharedfootnote{7f6c90984801d9ef}{เยน หตฺเถน อุทกํ อาสิญฺจิตพฺพํ (สฺยา)}, น เตเนว หตฺเถน ปาทา โธวิตพฺพา, อุปาหนา\-ปุญฺฉนโจฬกํ ปุจฺฉิตฺวา อุปาหนา ปุญฺฉิตพฺพา, อุปาหนา ปุญฺฉนฺเตน ปฐมํ สุกฺเขน โจฬเกน ปุญฺฉิตพฺพา, ปจฺฉา อลฺเลน, อุปาหนา\-ปุญฺฉนโจฬกํ โธวิตฺวา\csromansharedfootnote{520529c761a09bdc}{ปีเฬตฺวา (สฺยา)} เอกมนฺตํ วิสฺสชฺเชตพฺพํ.}

\prose{สเจ อาวาสิโก ภิกฺขุ วุฑฺโฒ โหติ, อภิวาเทตพฺโพ, สเจ นวโก โหติ, อภิวาทาเป\-ตพฺโพ. เสนาสนํ ปุจฺฉิตพฺพํ “กตมํ เม เสนาสนํ ปาปุณาตี”ติ, อชฺฌาวุฏฺฐํ วา อนชฺฌาวุฏฺฐํ วา ปุจฺฉิตพฺพํ, โคจโร ปุจฺฉิตพฺโพ, อโคจโร ปุจฺฉิตพฺโพ, เสกฺข\-สมฺมตานิ\csromansharedfootnote{682db2bc1e1d1cc0}{เสขสมฺมตานิ (ก)} กุลานิ ปุจฺฉิตพฺพานิ, วจฺจฏฺฐานํ ปุจฺฉิตพฺพํ, ปสฺสาวฏฺฐานํ ปุจฺฉิตพฺพํ, ปานียํ ปุจฺฉิตพฺพํ, ปริโภชนียํ ปุจฺฉิตพฺพํ, กตฺตรทณฺโฑ ปุจฺฉิตพฺโพ, สํฆสฺส กติก\-สณฺฐานํ ปุจฺฉิตพฺพํ “กํ กาลํ ปวิสิตพฺพํ, กํ กาลํ นิกฺขมิตพฺพนฺ”ติ. สเจ วิหาโร อนชฺฌาวุฏฺโฐ โหติ, กวาฏํ อาโกเฏตฺวา มุหุตฺตํ อาคเมตฺวา ฆฏิกํ อุคฺฆาเฏตฺวา กวาฏํ ปณาเมตฺวา พหิ ฐิเตน นิลฺโลเกตพฺโพ.}

\prose{\csromanfolio{375}สเจ วิหาโร อุกฺลาโป โหติ, มญฺเจ วา มญฺโจ อาโรปิโต โหติ, ปีเฐ วา ปีฐํ อาโรปิตํ โหติ, เสนาสนํ อุปริ ปุญฺชีกตํ\csromansharedfootnote{6fc89402dcb0b6bc}{ปุญฺชกิตํ (ก)} โหติ, สเจ อุสฺสหติ โสเธตพฺโพ,\csromansymbolfootnote{*}{วิ 3. 60, 65 ปิฏฺฐาทีสุปิ (โถกํ วิสทิสํ)}วิหารํ โสเธนฺเตน ปฐมํ ภูมตฺถรณํ นีหริตฺวา เอกมนฺตํ นิกฺขิปิตพฺพํ, มญฺจ\-ปฏิปาทกา นีหริตฺวา เอกมนฺตํ นิกฺขิปิตพฺพา, ภิสิพิพฺโพหนํ นีหริตฺวา เอกมนฺตํ นิกฺขิปิตพฺพํ, นิสีทน\-ปจฺจตฺถรณํ นีหริตฺวา เอกมนฺตํ นิกฺขิปิตพฺพํ, มญฺโจ นีจํ กตฺวา สาธุกํ อปฺปฏิฆํสนฺเตน อสํฆฏฺเฏนฺเตน กวาฏปิฏฺฐํ นีหริตฺวา เอกมนฺตํ นิกฺขิปิตพฺโพ, ปีฐํ นีจํ กตฺวา สาธุกํ อปฺปฏิฆํสนฺเตน อสํฆฏฺเฏนฺเตน กวาฏปิฏฺฐํ นีหริตฺวา เอกมนฺตํ นิกฺขิปิตพฺพํ, เขฬมลฺลโก นีหริตฺวา เอกมนฺตํ นิกฺขิปิตพฺโพ, อปสฺเสนผลกํ นีหริตฺวา เอกมนฺตํ นิกฺขิปิตพฺพํ. สเจ วิหาเร สนฺตานกํ โหติ, อุลฺโลกา ปฐมํ โอหาเรตพฺพํ, อาโลกสนฺธิ\-กณฺณภาคา ปมชฺชิตพฺพา. สเจ เครุก\-ปริกมฺม\-กตา ภิตฺติ กณฺณกิตา โหติ, โจฬกํ เตเมตฺวา ปีเฬตฺวา ปมชฺชิตพฺพา. สเจ กาฬวณฺณกตา ภูมิ กณฺณกิตา โหติ, โจฬกํ เตเมตฺวา ปีเฬตฺวา ปมชฺชิตพฺพา. สเจ อกตา โหติ ภูมิ, อุทเกน ปริปฺโผสิตฺวา สมฺมชฺชิตพฺพา “มา วิหาโร รเชน อุหญฺญี”ติ\csromansharedfootnote{8b2f7d14ebd3850c}{อูหญฺญีติ (สี, สฺยา)}, สงฺการํ วิจินิตฺวา เอกมนฺตํ ฉฑฺเฑตพฺพํ.}

\prose{\csromansymbolmark{*}ภูมตฺถรณํ โอตาเปตฺวา โสเธตฺวา ปปฺโผเฏตฺวา อติหริตฺวา ยถาฐาเน\csromansharedfootnote{0e3824ade5202327}{ยถาปญฺญตฺตํ (สี, สฺยา), ยถาภาคํ (ก)} ปญฺญเปตพฺพํ, มญฺจ\-ปฏิปาทกา โอตาเปตฺวา ปมชฺชิตฺวา อติหริตฺวา ยถาฐาเน\csromansharedfootnote{a431dafbe9e2d779}{ยถาภาคํ (สฺยา, ก)} ฐเปตพฺพา, มญฺโจ โอตาเปตฺวา โสเธตฺวา ปปฺโผเฏตฺวา นีจํ กตฺวา สาธุกํ อปฺปฏิฆํสนฺเตน อสํฆฏฺเฏนฺเตน กวาฏปิฏฺฐํ อติหริตฺวา ยถาฐาเน\csromansharedfootnote{a431dafbe9e2d779}{ยถาภาคํ (สฺยา, ก)} ปญฺญเปตพฺโพ, ปีฐํ โอตาเปตฺวา โสเธตฺวา ปปฺโผเฏตฺวา นีจํ กตฺวา สาธุกํ อปฺปฏิฆํสนฺเตน อสํฆฏฺเฏนฺเตน กวาฏปิฏฺฐํ อติหริตฺวา ยถาฐาเน\csromansharedfootnote{a431dafbe9e2d779}{ยถาภาคํ (สฺยา, ก)} ปญฺญเปตพฺพํ, ภิสิพิพฺโพหนํ โอตาเปตฺวา โสเธตฺวา ปปฺโผเฏตฺวา อติหริตฺวา ยถาภาคํ ปญฺญเปตพฺพํ, นิสีทน\-ปจฺจตฺถรณํ โอตาเปตฺวา โสเธตฺวา ปปฺโผเฏตฺวา อติหริตฺวา ยถาภาคํ ปญฺญเปตพฺพํ, เขฬมลฺลโก โอตาเปตฺวา ปมชฺชิตฺวา อติหริตฺวา ยถาภาคํ ฐเปตพฺโพ, อปสฺเสนผลกํ โอตาเปตฺวา ปมชฺชิตฺวา อติหริตฺวา ยถาภาคํ ฐเปตพฺพํ. ปตฺตจีวรํ นิกฺขิปิตพฺพํ, ปตฺตํ \csromanfolio{376}นิกฺขิปนฺเตน เอเกน หตฺเถน ปตฺตํ คเตตฺวา เอเกน หตฺเถน เหฏฺฐามญฺจํ วา เหฏฺฐาปีฐํ วา ปรามสิตฺวา ปตฺโต นิกฺขิปิตพฺโพ, น จ อนนฺตรหิตาย ภูมิยา ปตฺโต นิกฺขิปิตพฺโพ. จีวรํ นิกฺขิปนฺเตน เอเกน หตฺเถน จีวรํ คเหตฺวา เอเกน หตฺเถน จีวรวํสํ วา จีวรรชฺชุํ วา ปมชฺชิตฺวา ปารโต อนฺตํ โอรโต โภคํ กตฺวา จีวรํ นิกฺขิปิตพฺพํ.}

\prose{สเจ ปุรตฺถิมา สรชา วาตา วายนฺติ, ปุรตฺถิมา วาตปานา ถเกตพฺพา. สเจ ปจฺฉิมา สรชา วาตา วายนฺติ, ปจฺฉิมา วาตปานา ตเกตพฺพา. สเจ อุตฺตรา สรชา วาตา วายนฺติ, อุตฺตรา วาตปานา ถเกตพฺพา. สเจ ทกฺขิณา สรชา วาตา วายนฺติ, ทกฺขิณา วาตปานา ถเกตพฺพา. สเจ สีตกาโล โหติ, ทิวา วาตปานา วิวริตพฺพา, รตฺติํ ถเกตพฺพา. สเจ อุณฺหกาโล โหติ, ทิวา วาตปานา ถเกตพฺพา, รตฺติํ วิวริตพฺพา.}

\prose{สเจ ปริเวณํ อุกฺลาปํ โหติ, ปริเวณํ สมฺมชฺชิตพฺพํ. สเจ โกฏฺฐโก อุกฺลาโป โหติ, โกฏฺฐโก สมฺมชฺชิตพฺโพ. สเจ อุปฏฺฐานสาลา อุกฺลาปา โหติ, อุปฏฺฐานสาลา สมฺมชฺชิตพฺพา. สเจ อคฺคิสาลา อุกฺลาปา โหติ, อคฺคิสาลา สมฺมชฺชิตพฺพา. สเจ วจฺจกุฏิ อุกฺลาปา โหติ, วจฺจกุฏิ สมฺมชฺชิตพฺพา. สเจ ปานียํ น โหติ, ปานียํ อุปฏฺฐาเป\-ตพฺพํ. สเจ ปริโภชนียํ น โหติ, ปริโภชนียํ อุปฏฺฐาเป\-ตพฺพํ. สเจ อาจมนกุมฺภิยา อุทกํ น โหติ, อาจมนกุมฺภิยา อุทกํ อาสิญฺจิตพฺพํ. อิทํ โข ภิกฺขเว อาคนฺตุกานํ ภิกฺขูนํ วตฺตํ, ยถา อาคนฺตุเกหิ ภิกฺขูหิ สมฺมา วตฺติตพฺพนฺติ.}

\csromanmatikaanchor{mtk.178}
\csromantocmarkref{section}{2. อาวาสิกวตฺตกถา}{mtk.178}
\bookmark[dest={mtk.178},level=1]{2. อาวาสิกวตฺตกถา}
\csromanheader{1}{2. อาวาสิกวตฺต\-กถา}
\proseitem{358}{เตน โข ปน สมเยน อาวาสิกา ภิกฺขู อาคนฺตุเก ภิกฺขู ทิสฺวา เนว อาสนํ ปญฺญเปนฺติ, น ปาโททกํ ปาทปีฐํ ปาทกถลิกํ อุปนีกฺขิปนฺติ, น ปจฺจุคฺคนฺตฺวา ปตฺตจีวรํ ปฏิคฺคณฺหนฺติ, น ปานีเยน ปุจฺฉนฺติ\csromansharedfootnote{9257d98c2d070a27}{น ปานีเยน ปุจฺฉนฺติ, น ปริโภชนีเยน ปุจฺฉนฺติ (สฺยา, กํ)}, น วุฑฺฒตเรปิ อาคนฺตุเก ภิกฺขู อภิวาเทนฺติ, น เสนาสนํ ปญฺญเปนฺติ, เย เต ภิกฺขู อปฺปิจฺฉา ฯเปฯ เต อุชฺฌายนฺติ ขิยฺยนฺติ วิปาเจนฺติ “กถํ หิ นาม อาวาสิกา ภิกฺขู อาคนฺตุเก ภิกฺขู ทิสฺวา เนว อาสนํ ปญฺญเปสฺสนฺติ, น ปาโททกํ ปาทปีฐํ ปาทกถลิกํ อุปนิกฺขิปิสฺสนฺติ, \csromanfolio{377}น ปจฺจุคฺคนฺตฺวา ปตฺตจีวรํ ปฏิคฺคหิสฺสนฺติ, น ปานีเยน ปุจฺฉิสฺสนฺติ. วุฑฺฒตเรปิ อาคนฺตุเก ภิกฺขู น อภิวาเทสฺสนฺติ, น เสนาสนํ ปญฺญเปสฺสนฺตี”ติ. อถ โข เต ภิกฺขู ภควโต เอตมตฺถํ อาโรเจสุํ ฯเปฯ. สจฺจํ กิร ภิกฺขเว ฯเปฯ. สจฺจํ ภควาติ ฯเปฯ วิครหิตฺวา ฯเปฯ ธมฺมิํ กถํ กตฺวา ภิกฺขู อามนฺเตสิ\csromandash{}}

\proseitem{359}{เตน หิ ภิกฺขเว อาวาสิกานํ ภิกฺขูนํ วตฺตํ ปญฺญเปสฺสามิ, ยถา อาวาสิเกหิ ภิกฺขูหิ สมฺมา วตฺติตพฺพํ. อาวาสิเกน ภิกฺขเว ภิกฺขุนา อาคนฺตุกํ ภิกฺขุํ วุฑฺฒตรํ ทิสฺวา อาสนํ ปญฺญเปตพฺพํ, ปาโททกํ ปาทปีฐํ ปาทกถลิกํ อุปนิกฺขิปิตพฺพํ, ปจฺจุคฺคนฺตฺวา ปตฺตจีวรํ ปฏิคฺคเหตพฺพํ, ปานีเยน ปุจฺฉิตพฺโพ\csromansharedfootnote{708651b746680c91}{ปานีเยน ปุจฺฉิตพฺโพ, ปริโภชนีเยน ปุจฺฉิตพฺโพ (สฺยา)}, สเจ อุสฺสหติ, อุปาหนา ปุญฺฉิตพฺพา, อุปาหนา ปุญฺฉนฺเตน ปฐมํ สุกฺเขน โจฬเกน ปุญฺฉิตพฺพา, ปจฺฉา อลฺเลน, อุปาหนา\-ปุญฺฉนโจฬกํ โธวิตฺวา\csromansharedfootnote{182545861982840a}{โธวิตฺวา ปีเฬตฺวา (สฺยา)} เอกมนฺตํ วิสฺสชฺเชตพฺพํ.}

\prose{อาคนฺตุโก ภิกฺขุ วุฑฺฒตโร อภิวาเทตพฺโพ, เสนาสนํ ปญฺญเปตพฺพํ “เอตํ เต เสนาสนํ ปาปุณาตี”ติ. อชฺฌาวุฏฺฐํ วา อนชฺฌาวุฏฺฐํ วา อาจิกฺขิตพฺพํ, โคจโร อาจิกฺขิตพฺโพ, อาโคจโร อาจิกฺขิตพฺโพ, เสกฺข\-สมฺมตานิ กุลานิ อาจิกฺขิตพฺพานิ, วจฺจฏฺฐานํ อาจิกฺขิตพฺพํ, ปสฺสาวฏฺฐานํ อาจิกฺขิตพฺพํ, ปานียํ อาจิกฺขิตพฺพํ, ปริโภชนียํ อาจิกฺขิตพฺพํ, กตฺตรทณฺโฑ อาจิกฺขิตพฺโพ, สํฆสฺส กติก\-สณฺฐานํ อาจิกฺขิตพฺพํ “อิมํ กาลํ ปวิสิตพฺพํ, อิมํ กาลํ นิกฺขมิตพฺพนฺ”ติ.}

\prose{สเจ นวโก โหติ, นิสินฺนเกเนว อาจิกฺขิตพฺพํ “อตฺร ปตฺตํ นิกฺขิปาหิ, อตฺร จีวรํ นิกฺขิปาหิ, อิทํ อาสนํ นิสีทาหี”ติ, ปานียํ อาจิกฺขิตพฺพํ, ปริโภชนียํ อาจิกฺขิตพฺพํ, อุปาหนา\-ปุญฺฉนโจฬกํ อาจิกฺขิตพฺพํ, อาคนฺตุโก ภิกฺขุ นวโก อภิวาทาเป\-ตพฺโพ, เสนาสนํ อาจิกฺขิตพฺพํ “เอตํ เต เสนาสนํ ปาปุณาตี”ติ, อชฺฌาวุฏฺฐํ วา อนชฺฌาวุฏฺฐํ วา อาจิกฺขิตพฺพํ, โคจโร อาจิกฺขิตพฺโพ, อโคจโร อาจิกฺขิตพฺโพ, เสกฺข\-สมฺมตานิ กุลานิ อาจิกฺขิตพฺพานิ, วจฺจฏฺฐานํ อาจิกฺขิตพฺพํ, ปสฺสาวฏฺฐานํ อาจิกฺขิตพฺพํ, ปานียํ อาจิกฺขิตพฺพํ, ปริโภชนียํ อาจิกฺขิตพฺพํ, กตฺตรทณฺโฑ อาจิกฺขิตพฺโพ, สํฆสฺส กติก\-สณฺฐานํ อาจิกฺขิตพฺพํ “อิมํ กาลํ ปวิสิตพฺพํ, อิมํ กาลํ นิกฺขมิตพฺพนฺ”ติ. อิทํ โข \csromanfolio{378}ภิกฺขเว อาวาสิกานํ ภิกฺขูนํ วตฺตํ, ยถา อาวาสิเกหิ ภิกฺขูหิ สมฺมา วตฺติตพฺพนฺติ.}

\csromanmatikaanchor{mtk.179}
\csromantocmarkref{section}{3. คมิกวตฺตกถา}{mtk.179}
\bookmark[dest={mtk.179},level=1]{3. คมิกวตฺตกถา}
\csromanheader{1}{3. คมิกวตฺต\-กถา}
\proseitem{360}{เตน โข ปน สมเยน คมิกา ภิกฺขู ทารุภณฺฑํ มตฺติกาภณฺฑํ อปฺปฏิสาเมตฺวา ทฺวารวาตปานํ วิวริตฺวา เสนาสนํ อนาปุจฺฉา ปกฺกมนฺติ, ทารุภณฺฑํ มตฺติกาภณฺฑํ นสฺสติ, เสนาสนํ อคุตฺตํ โหติ. เย เต ภิกฺขู อปฺปิจฺฉา ฯเปฯ เต อุชฺฌายนฺติ ขิยฺยนฺติ วิปาเจนฺติ “กถํ หิ นาม คมิกา ภิกฺขู ทารุภณฺฑํ มตฺติกาภณฺฑํ อปฺปฏิสาเมตฺวา ทฺวารวาตปานํ วิวริตฺวา เสนาสนํ อนาปุจฺฉา ปกฺกมิสฺสนฺติ, ทารุภณฺฑํ มตฺติกาภณฺฑํ นสฺสติ, เสนาสนํ อคุตฺตํ โหตี”ติ. อถ โข เต ภิกฺขู ภควโต เอตมตฺถํ อาโรเจสุํ ฯเปฯ. สจฺจํ กิร ภิกฺขเว ฯเปฯ. สจฺจํ ภควาติ ฯเปฯ วิครหิตฺวา ฯเปฯ ธมฺมิํ กถํ กตฺวา ภิกฺขู อามนฺเตสิ\csromandash{}}

\proseitem{361}{เตน หิ ภิกฺขเว คมิกานํ ภิกฺขูนํ วตฺตํ ปญฺญเปสฺสามิ, ยถา คมิเกหิ ภิกฺขูหิ สมฺมา วตฺติตพฺพํ. คมิเกน ภิกฺขเว ภิกฺขุนา ทารุภณฺฑํ มตฺติกาภณฺฑํ ปฏิสาเมตฺวา ทฺวารวาตปานํ ถเกตฺวา เสนาสนํ อาปุจฺฉา ปกฺกมิตพฺพํ\csromansharedfootnote{01909a8817ee8eec}{อาปุจฺฉิตพฺพํ (สฺยา)}. สเจ ภิกฺขุ น โหติ, สามเณโร อาปุจฺฉิตพฺโพ. สเจ สามเณโร น โหติ, อารามิโก อาปุจฺฉิตพฺโพ. สเจ อารามิโก น โหติ, อุปาสโก อาปุจฺฉิตพฺโพ. สเจ น โหติ ภิกฺขุ วา สามเณโร วา อารามิโก วา อุปาสโก วา, จตูสุ ปาสาเณสุ มญฺจํ ปญฺญเปตฺวา มญฺเจ มญฺจํ อาโรเปตฺวา ปีเฐ ปีฐํ อาโรเปตฺวา เสนาสนํ อุปาริ ปุญฺชํ กริตฺวา ทารุภณฺฑํ มตฺติกาภณฺฑํ ปฏิสาเมตฺวา ทฺวารวาตปานํ ถเกตฺวา ปกฺกมิตพฺพํ. สเจ วิหาโร โอวสฺสติ, สเจ อุสฺสหติ, ฉาเทตพฺโพ, อุสฺสุกฺกํ วา กาตพฺพํ “กินฺติ นุ โข วิหาโร ฉาทิเยถา”ติ, เอวญฺเจตํ ลเภถ, อิจฺเจตํ กุสลํ. โน เจ ลเภถ, โย เทโส อโนวสฺสโก โหติ, ตตฺถ จตูสุ ปาสาเณสุ มญฺจํ ปญฺญเปตฺวา มญฺเจ มญฺจํ อาโรเปตฺวา ปีเฐ ปีฐํ อาโรเปตฺวา เสนาสนํ อุปริ ปุญฺชํ กริตฺวา ทารุภณฺฑํ มตฺติกาภณฺฑํ ปฏิสาเมตฺวา ทฺวารวาตปานํ ถเกตฺวา ปกฺกมิตพฺพํ. สเจ สพฺโพ วิหาโร โอวสฺสติ, สเจ อุสฺสหติ, เสนาสนํ คามํ อติหริตพฺพํ, อุสฺสุกฺกํ วา \csromanfolio{379}กาตพฺพํ “กินฺติ นุ โข เสนาสนํ คามํ อติหริเยถา”ติ, เอวญฺเจตํ ลเภถ, อิจฺเจตํ กุสลํ. โน เจ ลเภถ, อชฺโฌกาเส จตูสุ ปาสาเณสุ มญฺจํ ปญฺญเปตฺวา มญฺเจ มญฺจํ อาโรเปตฺวา ปีเฐ ปีฐํ อาโรเปตฺวา เสนาสนํ อุปริ ปุญฺชํ กริตฺวา ทารุภณฺฑํ มตฺติกาภณฺฑํ ปฏิสาเมตฺวา ติเณน วา ปณฺเณน วา ปฏิจฺฉาเทตฺวา ปกฺกมิตพฺพํ “อปฺเปว นาม องฺคานิปิ เสเสยฺยุนฺ”ติ. อิทํ โข ภิกฺขเว คมิกานํ ภิกฺขูนํ วตฺตํ, ยถา คมิเกหิ ภิกฺขูหิ สมฺมา วตฺติตพฺพนฺติ.}

\csromanmatikaanchor{mtk.180}
\csromantocmarkref{section}{4. อนุโมทนวตฺตกถา}{mtk.180}
\bookmark[dest={mtk.180},level=1]{4. อนุโมทนวตฺตกถา}
\csromanheader{1}{4. อนุโมทนวตฺต\-กถา}
\proseitem{362}{เตน โข ปน สมเยน ภิกฺขู ภตฺตคฺเค น อนุโมทนฺติ. มนุสฺสา อุชฺฌายนฺติ ขิยฺยนฺติ วิปาเจนฺติ “กถํ หิ นาม สมณา สกฺยปุตฺติยา ภตฺตคฺเค น อนุโมทิสฺสนฺตี”ติ. อสฺโสสุํ โข ภิกฺขู เตสํ มนุสฺสานํ อุชฺฌายนฺตานํ ขิยฺยนฺตานํ วิปาเจนฺตานํ. อถ โข เต ภิกฺขู ภควโต เอตมตฺถํ อาโรเจสุํ. อถ โข ภควา เอตสฺมิํ นิทาเน เอตสฺมิํ ปกรเณ ธมฺมิํ กถํ กตฺวา ภิกฺขู อามนฺเตสิ “อนุชานามิ ภิกฺขเว ภตฺตคฺเค อนุโมทิตุนฺ”ติ. อถ โข เตสํ ภิกฺขูนํ เอตทโหสิ “เกน นุ โข ภตฺตคฺเค อนุโมทิตพฺพนฺ”ติ. ภควโต เอตมตฺถํ อาโรเจสุํ. อถ โข ภควา เอตสฺมิํ นิทาเน เอตสฺมิํ ปกรเณ ธมฺมิํ กถํ กตฺวา ภิกฺขู อามนฺเตสิ “อนุชานามิ ภิกฺขเว เถเรน ภิกฺขุนา ภตฺตคฺเค อนุโมทิตุนฺ”ติ.}

\prose{เตน โข ปน สมเยน อญฺญตรสฺส ปูคสฺส สํฆภตฺตํ โหติ, อายสฺมา สาริปุตฺโต สํฆตฺเถโร โหติ, ภิกฺขู “ภควตา อนุญฺญาตํ เถเรน ภิกฺขุนา ภตฺตคฺเค อนุโมทิตุนฺ”ติ อายสฺมนฺตํ สาริปุตฺตํ เอกกํ โอหาย ปกฺกมิํสุ. อถ โข อายสฺมา สาริปุตฺโต เต มนุสฺเส ปฏิสมฺโมทิตฺวา ปจฺฉา เอกโก อคมาสิ. อทฺทสา โข ภควา อายสฺมนฺตํ สาริปุตฺตํ ทูรโตว เอกกํ อาคจฺฉนฺตํ, ทิสฺวาน อายสฺมนฺตํ สาริปุตฺตํ เอตทโวจ “กจฺจิ สาริปุตฺต ภตฺตํ อิทฺธํ อโหสี”ติ. อิทฺธํ โข ภนฺเต ภตฺตํ อโหสิ, อปิ จ มํ ภิกฺขู เอกกํ โอหาย ปกฺกนฺตาติ. อถ โข ภควา เอตสฺมิํ นีทาเน เอตสฺมิํ ปกรเณ ธมฺมิํ กถํ กตฺวา ภิกฺขู อามนฺเตสิ “อนุชานามิ ภิกฺขเว ภตฺตคฺเค จตูหิ ปญฺจหิ เถรานุเถเรหิ ภิกฺขูหิ อาคเมตุนฺ”ติ.}

\prose{\csromanfolio{380}เตน โข ปน สมเยน อญฺญตโร เถโร ภตฺตคฺเค วจฺจิโต อาคเมสิ, โส วจฺจํ สทฺธาเรตุํ อสกฺโกนฺโต มุจฺฉิโต ปปติ. ภควโต เอตมตฺถํ อาโรเจสุํ. อนุชานามิ ภิกฺขเว สติ กรณีเย อานนฺตริกํ ภิกฺขุํ อาปุจฺฉิตฺวา คนฺตุนฺติ.}

\csromanmatikaanchor{mtk.181}
\csromantocmarkref{section}{5. ภตฺตคฺควตฺตกถา}{mtk.181}
\bookmark[dest={mtk.181},level=1]{5. ภตฺตคฺควตฺตกถา}
\csromanheader{1}{5. ภตฺตคฺควตฺต\-กถา}
\proseitem{363}{เตน โข ปน สมเยน ฉพฺพคฺคิยา ภิกฺขู ทุนฺนิวตฺถา ทุปฺปารุตา อนากปฺป\-สมฺปนฺนา ภตฺตคฺคํ คจฺฉนฺติ, โวกฺกมฺมปิ เถรานํ ภิกฺขูนํ ปุรโต ปุรโต คจฺฉนฺติ, เถเรปิ ภิกฺขู อนุปขชฺช นิสีทนฺติ, นเวปิ ภิกฺขู อาสเนน ปฏิพาหนฺติ, สํฆาฏิมฺปิ โอตฺถริตฺวา อนฺตรฆเร นิสีทนฺติ. เย เต ภิกฺขู อปฺปิจฺฉา ฯเปฯ เต อุชฺฌายนฺติ ขิยฺยนฺติ วิปาเจนฺติ “กถํ หิ นาม ฉพฺพคฺคิยา ภิกฺขู ทุนฺนิวตฺถา ทุปฺปารุตา อนากปฺป\-สมฺปนฺนา ภตฺตคฺคํ คจฺฉิสฺสนฺติ, โวกฺกมฺมปิ เถรานํ ภิกฺขูนํ ปุรโต ปุรโต คจฺฉิสฺสนฺติ, เถเรปิ ภิกฺขู อนุปขชฺช นิสีทิสฺสนฺติ, นเวปิ ภิกฺขู อาสเนนปิ ปฏิพาหิสฺสนฺติ, สํฆาฏิมฺปิ โอตฺถริตฺวา อนฺตรฆเร นิสีทิสฺสนฺตี”ติ. อถ โข เต ภิกฺขู ภควโต เอตมตฺถํ อาโรเจสุํ ฯเปฯ. สจฺจํ กิร ภิกฺขเว ฉพฺพคฺคิยา ภิกฺขู ทุนฺนิวตฺถา ทุปฺปารุตา อนากปฺป\-สมฺปนฺนา ภตฺตคฺคํ คจฺฉนฺติ, โวกฺกมฺมปิ เถรานํ ภิกฺขูนํ ปุรโต ปุรโต คจฺฉนฺติ, เถเรปิ ภิกฺขู อนุปขชฺช นิสีทนฺติ, นเวปิ ภิกฺขู อาสเนน ปฏิพาหนฺติ, สํฆาฏิมฺปิ โอตฺถริตฺวา อนฺตรฆเร นิสีทนฺตีติ. สจฺจํ ภควาติ ฯเปฯ วิครหิตฺวา ฯเปฯ ธมฺมิํ กถํ กตฺวา ภิกฺขู อามนฺเตสิ\csromandash{}}

\proseitem{364}{เตน หิ ภิกฺขเว ภิกฺขูนํ ภตฺตคฺค\-วตฺตํ ปญฺญเปสฺสามิ, ยถา ภิกฺขูหิ ภตฺตคฺเค สมฺมา วตฺติตพฺพํ. สเจ อาราเม กาโล อาโรจิโต โหติ, ติมณฺฑลํ ปฏิจฺฉาเทนฺเตน ปริมณฺฑลํ นิวาเสตฺวา กายพนฺธนํ พนฺธิตฺวา สคุณํ กตฺวา สํฆาฏิโย ปารุปิตฺวา คณฺฐิกํ ปฏิมุญฺจิตฺวา โธวิตฺวา ปตฺตํ คเหตฺวา สาธุกํ อตรมาเนน คาโม ปวิสิตพฺโพ.}

\prose{น โวกฺกมฺม เถรานํ ภิกฺขูนํ ปุรโต ปุรโต คนฺตพฺพํ, สุปฺปฏิจฺฉนฺเนน อนฺตรฆเร คนฺตพฺพํ, สุสํวุเตน อนฺตรฆเร คนฺตพฺพํ, โอกฺขิตฺต\-จกฺขุนา อนฺตรฆเร คนฺตพฺพํ, น อุกฺขิตฺตกาย อนฺตรฆเร คนฺตพฺพํ, น อุชฺชคฺฆิกาย อนฺตรฆเร คนฺตพฺพํ, อปฺปสทฺเทน อนฺตรฆเร คนฺตพฺพํ, น กายปฺปจาลกํ อนฺตรฆเร \csromanfolio{381}คนฺตพฺพํ, น พาหุปฺปจาลกํ อนฺตรฆเร คนฺตพฺพํ, น สีสปฺปจาลกํ อนฺตรฆเร คนฺตพฺพํ, น ขมฺภกเตน อนฺตรฆเร คนฺตพฺพํ, น โอคุณฺฐิเตน อนฺตรฆเร คนฺตพฺพํ, น อุกฺกุฏิกาย อนฺตรฆเร คนฺตพฺพํ.}

\prose{สุปฺปฏิจฺฉนฺเนน อนฺตรฆเร นิสีทิตพฺพํ, สุสํวุเตน อนฺตรฆเร นิสีทิตพฺพํ โอกฺขิตฺต\-จกฺขุนา อนฺตรฆเร นิสีทิตพฺพํ, น อุกฺขิตฺตกาย อนฺตรฆเร นิสีทิตพฺพํ, น อุชฺชคฺฆิกาย อนฺตรฆเร นิสีทิตพฺพํ, อปฺปสทฺเทน อนฺตรฆเร นิสีทิตพฺพํ, น กายปฺปจาลกํ อนฺตรฆเร นิสีทิตพฺพํ, น พาหุปฺปจาลกํ อนฺตรฆเร นิสีทิตพฺพํ, น สีสปฺปจาลกํ อนฺตรฆเร นิสีทิตพฺพํ, น ขมฺภกเตน อนฺตรฆเร นิสีทิตพฺพํ, น โอคุณฺฐิเตน อนฺตรฆเร นิสีทิตพฺพํ, น ปลฺลตฺถิกาย อนฺตรฆเร นิสีทิตพฺพํ, น เถเร ภิกฺขู อนุปขชฺช นิสีทิตพฺพํ, น นวา ภิกฺขู อาสเนน ปฏิพาหิตพฺพา, น สํฆาฏิํ โอตฺถริตฺวา อนฺตรฆเร นิสีทิตพฺพํ.}

\prose{อุทเก ทิยฺยมาเน อุโภหิ หตฺเถหิ ปตฺตํ ปฏิคฺคเหตฺวา อุทกํ ปฏิคฺคเหตพฺพํ, นีจํ กตฺวา สาธุกํ อปฺปฏิฆํสนฺเตน ปตฺโต โธวิตพฺโพ. สเจ อุทก\-ปฺปฏิคฺคาหโก โหติ, นีจํ กตฺวา อุทก\-ปฺปฏิคฺคเห อุทกํ อาสิญฺจิตพฺพํ “มา อุทก\-ปฺปฏิคฺคาหโก อุทเกน โอสิญฺจิ\csromansharedfootnote{299b899c20bd787e}{โอสิญฺจิยฺยี (ก)}, มา สามนฺตา ภิกฺขู อุทเกน โอสิญฺจิํสุ\csromansharedfootnote{299ceccc49b91213}{โอสิญฺจิยฺยิํสุ (ก)}, มา สํฆาฏิ อุทเกน โอสิญฺจี”ติ. สเจ อุทก\-ปฺปฏิคฺคาหโก น โหติ, นีจํ กตฺวา ฉมาย อุทกํ อาสิญฺจิตพฺพํ “มา สามนฺตา ภิกฺขู อุทเกน โอสิญฺจิํสุ, มา สํฆาฏิ อุทเกน โอสิญฺจี”ติ.}

\prose{โอทเน ทิยฺยมาเน อุโภหิ หตฺเถหิ ปตฺตํ ปฏิคฺคเหตฺวา โอทโน ปฏิคฺคเหตพฺโพ, สูปสฺส โอกาโส กาตพฺโพ. สเจ โหติ สปฺปิ วา เตลํ วา อุตฺตริภงฺคํ วา, เถเรน วตฺตพฺโพ “สพฺเพสํ สมกํ สมฺปาเทหี”ติ. สกฺกจฺจํ ปิณฺฑปาโต ปฏิคฺคเหตพฺโพ, ปตฺตสญฺญินา ปิณฺฑปาโต ปฏิคฺคเหตพฺโพ, สมสูปโก ปิณฺฑปาโต ปฏิคฺคเหตพฺโพ, สมติตฺติโก ปิณฺฑปาโต ปฏิคฺคเหตพฺโพ.}

\prose{น ตาว เถเรน ภุญฺชิตพฺพํ, ยาว น สพฺเพสํ โอทโน สมฺปตฺโต โหติ, สกฺกจฺจํ ปิณฺฑปาโต ภุญฺชิตพฺโพ, ปตฺตสญฺญินา ปิณฺฑปาโต ภุญฺชิตพฺโพ, สปทานํ ปิณฺฑปาโต ภุญฺชิตพฺโพ, สมสูปโก ปิณฺฑปาโต \csromanfolio{382}ภุญฺชิตพฺโพ, น ถูปกโต โอมทฺทิตฺวา ปิณฺฑปาโต ภุญฺชิตพฺโพ, น สูปํ วา พฺยญฺชนํ วา โอทเนน ปฏิจฺฉาเทตพฺพํ ภิยฺโยกมฺยตํ อุปาทาย, น สูปํ วา โอทนํ วา อคิลาเนน อตฺตโน อตฺถาย วิญฺญาเปตฺวา ภุญฺชิตพฺพํ, น อุชฺฌานสญฺญินา ปเรสํ ปตฺโต โอโลเกตพฺโพ, นาติมหนฺโต กพโฬ กาตพฺโพ, ปฏิมณฺฑโล อาโลโป กาตพฺโพ, น อนาหเฏ กพเฬ มุขทฺวารํ วิวริตพฺพํ, น ภุญฺชมาเนน สพฺโพ หตฺโถ มุเข ปกฺขิปิตพฺโพ, น สกพเฬน มุเขน พฺยาหริตพฺพํ, น ปิณฺฑุกฺเขปกํ ภุญฺชิตพฺพํ, น กพฬาวจฺเฉทกํ ภุญฺชิตพฺพํ, น อวคณฺฑการกํ ภุญฺชิตพฺพํ, น หตฺถนิทฺธุนกํ ภุญฺชิตพฺพํ, น สิตฺถาวการกํ ภุญฺชิตพฺพํ, น ชีวฺหานิจฺฉารกํ ภุญฺชิตพฺพํ, น จปุจปุการกํ ภุญฺชิตพฺพํ, น สุรุสุรุการกํ ภุญฺชิตพฺพํ, น หตฺถนิลฺเลหกํ ภุญฺชิติพฺพํ, น ปตฺตนิลฺเลหกํ ภุญฺชิตพฺพํ, น โอฏฺฐนิลฺเลหกํ ภุญฺชิตพฺพํ.}

\prose{น สามิเสน หตฺเถน ปานียถาลโก ปฏิคฺคเหตพฺโพ, น ตาว เถเรน อุทกํ ปฏิคฺคเหตพฺพํ, ยาว น สพฺเพว ภุตฺตาวิโน โหนฺติ, อุทเก ทิยฺยมาเน อุโภหิ หตฺเถหิ ปตฺตํ ปฏิคฺคเหตฺวา อุทกํ ปฏิคฺคเหตพฺพํ, นีจํ กตฺวา สาธุกํ อปฺปฏิฆํสนฺเตน ปตฺโต โธวิตพฺโพ. สเจ อุทก\-ปฺปฏิคฺคาหโก โหติ, นีจํ กตฺวา อุทก\-ปฺปฏิคฺคเห อุทกํ อาสิญฺจิตพฺพํ “มา อุทก\-ปฺปฏิคฺคาหโก อุทเกน โอสิญฺจิ, มา สามนฺตา ภิกฺขู อุทเกน โอสิญฺจิํสุ, มา สํฆาฏิ อุทเกน โอสิญฺจี”ติ. สเจ อุทก\-ปฺปฏิคฺคาหโก น โหติ, นีจํ กตฺวา ฉมาย อุทกํ อาสิญฺจิตพฺพํ “มา สามนฺตา ภิกฺขู อุทเกน โอสิญฺจิํสุ, มา สํฆาฏิ อุทเกน โอสิญฺจี”ติ, น สสิตฺถกํ ปตฺตโธวนํ อนฺตรฆเร ฉฑฺเฑตพฺพํ.}

\prosewithcloser{นิวตฺตนฺเตน นวเกหิ ภิกฺขูหิ ปฐมตรํ นิวตฺติตพฺพํ, ปจฺฉา เถเรหิ, สุปฺปฏิจฺฉนฺเนน อนฺตรฆเร คนฺตพฺพํ, สุสํวุเตน อนฺตรฆเร คนฺตพฺพํ, โอกฺขิตฺต\-จกฺขุนา อนฺตรฆเร คนฺตพฺพํ, น อุกฺขิตฺตกาย อนฺตรฆเร คนฺตพฺพํ, น อุชฺชคฺฆิกาย อนฺตรฆเร คนฺตพฺพํ, อปฺปสทฺเทน อนฺตรฆเร คนฺตพฺพํ, น กายปฺปจาลกํ อนฺตรฆเร คนฺตพฺพํ, น พาหุปฺปจาลกํ อนฺตรฆเร คนฺตพฺพํ, น สีสปฺปจาลกํ อนฺตรฆเร คนฺตพฺพํ, น ขมฺภกเตน อนฺตรฆเร คนฺตพฺพํ, น โอคุณฺฐิเตน อนฺตรฆเร คนฺตพฺพํ, น อุกฺกุฏิกาย อนฺตรฆเร คนฺตพฺพํ. อิทํ โข ภิกฺขเว ภิกฺขูนํ ภตฺตคฺค\-วตฺตํ, ยถา ภิกฺขูหิ ภตฺตคฺเค สมฺมา วตฺติตพฺพนฺติ.}{ปฐม\-ภาณวาโร นิฏฺฐิโต.}
\csromanmatikaanchor{mtk.182}
\csromantocmarkref{section}{6. ปิณฺฑจาริกวตฺตกถา}{mtk.182}
\bookmark[dest={mtk.182},level=1]{6. ปิณฺฑจาริกวตฺตกถา}
\csromanheader{1}{6. ปิณฺฑจาริกวตฺต\-กถา}
\proseitem{365}{\csromanfolio{383}เตน โข ปน สมเยน ปิณฺฑจาริกา ภิกฺขู ทุนฺนิวตฺถา ทุปฺปารุตา อนากปฺป\-สมฺปนฺนา ปิณฺฑาย จรนฺติ, อสลฺลกฺเขตฺวาปิ นิเวสนํ ปวิสนฺติ, อสลฺลกฺเขตฺวาปิ นิกฺขมนฺติ, อติสหสาปิ ปวิสนฺติ, อติสหสาปิ นิกฺขมนฺติ, อติทูเรปิ ติฏฺฐนฺติ, อจฺจาสนฺเนปิ ติฏฺฐนฺติ, อติจิรมฺปิ ติฏฺฐนฺติ, อติลหุมฺปิ นิวตฺตนฺติ, อญฺญตโรปิ ปิณฺฑจาริโก ภิกฺขุ อสลฺลกฺเขตฺวา นิเวสนํ ปาวิสิ, โส จ ทฺวารํ มญฺญมาโน อญฺญตรํ โอวรกํ ปาวิสิ, ตสฺมิํปิ โอวรเก อิตฺถี นคฺคา อุตฺตานา นิปนฺนา โหติ. อทฺทสา โข โส ภิกฺขุ ตํ อิตฺถิํ นคฺคํ อุตฺตานํ นิปนฺนํ, ทิสฺวาน “นยิทํ ทฺวารํ, โอวรกํ อิทนฺ”ติ ตมฺหา โอวรกา นิกฺขมิ. อทฺทสา โข ตสฺสา อิตฺถิยา สามิโก ตํ อิตฺถิํ นคฺคํ อุตฺตานํ นิปนฺนํ, ทิสฺวาน “อิมินา เม ภิกฺขุนา ปชาปติ ทูสิตา”ติ ตํ ภิกฺขุํ คเหตฺวา อาโกเฏสิ. อถ โข สา อิตฺถี เตน สทฺเทน ปฏิพุชฺฌิตฺวา ตํ ปุริสํ เอตทโวจ “กิสฺส ตฺวํ อยฺย อิมํ ภิกฺขุํ อาโกเฏสี”ติ. อิมินาสิ ตฺวํ ภิกฺขุนา ทูสิตาติ. “นาหํ อยฺย อิมินา ภิกฺขุนา ทูสิตา, อการโก โส ภิกฺขู”ติ ตํ ภิกฺขุํ มุญฺจาเปสิ. อถ โข โส ภิกฺขุ อารามํ คนฺตฺวา ภิกฺขูนํ เอตมตฺถํ อาโรเจสิ. เย เต ภิกฺกู อปฺปิจฺฉา ฯเปฯ เต อุชฺฌายนฺติ ขิยฺยนฺติ วิปาเจนฺติ “กถํ หิ นาม ปิณฺฑจาริกา ภิกฺขู ทุนฺนิวตฺถา ทุปฺปารุตา อนากปฺป\-สมฺปนฺนา ปิณฺฑาย จริสฺสนฺติ, อสลฺลกฺเขตฺวาปิ นิเวสนํ ปวิสิสฺสนฺติ, อสลฺลกฺเขตฺวาปิ นิกฺขมิสฺสนฺติ, อติสหสาปิ ปวิสิสฺสนฺติ, อติสหสาปิ นิกฺขมิสฺสนฺติ, อติทูเรปิ ติฏฺฐิสฺสนฺติ, อจฺจาสนฺเนปิ ติฏฺฐิสฺสนฺติ, อติจิรมฺปิ ติฏฺฐิสฺสนฺติ, อติลหุมฺปิ นิวตฺติสฺสนฺตี”ติ. อถ โข เต ภิกฺขุ ภควโต เอตมตฺถํ อาโรเจสุํ ฯเปฯ. สจฺจํ กิร ภิกฺขเว ฯเปฯ. สจฺจํ ภควาติ ฯเปฯ วิครหิตฺวา ฯเปฯ ธมฺมิํ กถํ กตฺวา ภิกฺขู อามนฺเตสิ\csromandash{}}

\proseitem{366}{เตน หิ ภิกฺขเว ปิณฺฑจาริกานํ ภิกฺขูนํ วตฺตํ ปญฺญเปสฺสามิ, ยถา ปิณฺฑจาริเกหิ ภิกฺขูหิ สมฺมา วตฺติตพฺพํ. ปิณฺฑจาริเกน ภิกฺขเว ภิกฺขุนา “อิทานิ คามํ ปวิสิสฺสามี”ติ ติมณฺฑลํ ปฏิจฺฉาเทนฺเตน ปริมณฺฑลํ นิวาเสตฺวา กายพนฺธนํ พนฺธิตฺวา สคุณํ กตฺวา สํฆาฏิโย ปารุปิตฺวา คณฺฐิตํ ปฏิมุญฺจิตฺวา โธวิตฺวา ปตฺตํ คเหตฺวา สาธุกํ อตรมาเนน คาโม ปวิสิตพฺโพ.}

\prose{\csromanfolio{384}สุปฺปฏิจฺฉนฺเนน อนฺตรฆเร คนฺตพฺพํ, สุสํวุเตน อนฺตเรฆเร คนฺตพฺพํ, โอกฺขิตฺต\-จกฺขุนา อนฺตรฆเร คนฺตพฺพํ, น อุกฺขิตฺตกาย อนฺตรฆเร คนฺตพฺพํ, น อุชฺชคฺฆิกาย อนฺตรฆเร คนฺตพฺพํ, อปฺปสทฺเทน อนฺตรฆเร คนฺตพฺพํ, น กายปฺปจาลกํ อนฺตรฆเร คนฺตพฺพํ, น พาหุปฺปจาลกํ อนฺตรฆเร คนฺตพฺพํ, น สีสปฺปจาลกํ อนฺตรฆเร คนฺตพฺพํ, น ขมฺภกเตน อนฺตรฆเร คนฺตพฺพํ, น โอคุณฺฐิเตน อนฺตรฆเร คนฺตพฺพํ, น อุกฺกุฏิกาย อนฺตรฆเร คนฺตพฺพํ.}

\prose{นิเวสนํ ปวิสนฺเตน สลฺเลกฺเข\-ตพฺพํ “อิมินา ปวิสิสฺสมิ, อิมินา นิกฺขมิสฺสามี”ติ, นาติสหสา ปวิสิตพฺพํ, นาติสหสา นิกฺขมิตพฺพํ, นาติทูเร ฐาตพฺพํ, นาจฺจาสนฺเน ฐาตพฺพํ, นาติจิรํ ฐาตพฺพํ, นาติลหุํ นิวตฺติตพฺพํ, ฐิตเกน สลฺลกฺเขตพฺพํ “ภิกฺขํ ทาตุกามา วา อทาตุกามา วา”ติ. สเจ กมฺมํ วา นิกฺขิปติ, อาสนา วา วุฏฺฐาติ, กฏจฺฉุํ วา ปรามสติ, ภาชนํ วา ปรามสติ, ฐเปติ\csromansharedfootnote{d9ce8fe5fe0f60e3}{ฐาเปติ (ก)} วา, “ทาตุกามสฺสา”ติ\csromansharedfootnote{384cf62e0c6085d5}{ทาตุกามิยาติ (สฺยา), ทาตุกามา วิยาติ (สี)} ฐาตพฺพํ, ภิกฺขาย ทิยฺยมานาย วาเมน หตฺเถน สํฆาฏิํ อุจฺจาเรตฺวา ทกฺขิเณน หตฺเถน ปตฺตํ ปณาเมตฺวา อุโภหิ หตฺเถหิ ปตฺตํ ปฏิคฺคเหตฺวา ภิกฺขา ปฏิคฺคเหตพฺพา, น จ ภิกฺขา\-ทายิกาย มุขํ อุลฺโลเกตพฺพํ\csromansharedfootnote{43fca1fa36f2c5f6}{โอโลเกตพฺพํ (สฺยา)}, สลฺลกฺเขตพฺพํ “สูปํ ทาตุกามา วา อทาตุกามา วา”ติ. สเจ กฏจฺฉุํ วา ปรามสติ, ภาชนํ วา ปรามสติ, ฐเปติ วา, “ทาตุกามสฺสา”ติ ฐาตพฺพํ, ภิกฺขาย ทินฺนาย สํฆาฏิยา ปตฺตํ ปฏิจฺฉาเทตฺวา สาธุกํ อตรมาเนน นิวตฺติตพฺพํ.}

\prose{สุปฺปฏิจฺฉนฺเนน อนฺตรฆเร คนฺตพฺพํ, สุสํวุเตน อนฺตรฆเร คนฺตพฺพํ, โอกฺขิตฺต\-จกฺขุนา อนฺตรฆเร คนฺตพฺพํ, น อุกฺขิตฺตกาย อนฺตรฆเร คนฺตพฺพํ, น อุชฺชคฺฆิกาย อนฺตรฆเร คนฺตพฺพํ, อปฺปสทฺเทน อนฺตรฆเร คนฺตพฺพํ, น กายปฺปจาลกํ อนฺตรฆเร คนฺตพฺพํ, น พาหุปฺปจาลกํ อนฺตรฆเร คนฺตพฺพํ, น สีสปฺปจาลกํ อนฺตรฆเร คนฺตพฺพํ, น ขมฺภกเตน อนฺตรฆเร คนฺตพฺพํ, น โอคุณฺฐิเตน อนฺตรฆเร คนฺตพฺพํ, น อุกฺกุฏิกาย อนฺตรฆเร คนฺตพฺพํ, โย ปฐมํ คามโต ปิณฺฑาย ปฏิกฺกมติ, เตน อาสนํ ปญฺญเปตพฺพํ, ปาโททกํ ปาทปีฐํ ปาทกถลิกํ อุปนิกฺขิปิตพฺพํ, อวกฺการปาติ โธวิตฺวา อุปฏฺฐาเปตพฺพา, ปานียํ ปริโภชนียํ อุปฏฺฐาเป\-ตพฺพํ, โย ปจฺฉา คามโต ปิณฺฑาย ปฏิกฺกมติ, สเจ โหติ ภุตฺตาวเสโส, สเจ อากงฺขติ, \csromanfolio{385}ภุญฺชิตพฺพํ, โน เจ อากงฺขติ, อปฺปหริเต วา ฉฑฺเฑตพฺพํ, อปฺปาณเก วา อุทเก โอปิลาเปตพฺพํ, เตน อาสนํ อุทฺธริตพฺพํ, ปาโททกํ ปาทปีฐํ ปาทกถลิกํ ปฏิสาเมตพฺพํ, อวกฺการปาติ โธวิตฺวา ปฏิสาเมตพฺพา, ปานียํ ปริโภชนียํ ปฏิสาเมตพฺพํ, ภตฺตคฺคํ สมฺมชฺชิตพฺพํ, โย ปสฺสติ ปานียฆฏํ วา ปริโภชนียฆฏํ วา วจฺจฆฏํ วา ริตฺตํ ตุจฺฉํ, เตน อุปฏฺฐาเป\-ตพฺพํ. สจสฺส โหติ, อวิสยฺหํ, หตฺถวิกาเรน ทุติยํ อามนฺเตตฺวา หตฺถ\-วิลงฺฆเกน อุปฏฺฐาเป\-ตพฺพํ, น จ ตปฺปจฺจยา วาจา ภินฺทิตพฺพา. อิทํ โข ภิกฺขเว ปิณฺฑจาริกานํ ภิกฺขูนํ วตฺตํ, ยถา ปิณฺฑจาริเกหิ ภิกฺขูหิ สมฺมา วตฺติตพฺพนฺติ.}

\csromanmatikaanchor{mtk.183}
\csromantocmarkref{section}{7. อารญฺญิกวตฺตกถา}{mtk.183}
\bookmark[dest={mtk.183},level=1]{7. อารญฺญิกวตฺตกถา}
\csromanheader{1}{7. อารญฺญิกวตฺต\-กถา}
\proseitem{367}{เตน โข ปน สมเยน สมฺพหุลา ภิกฺขู อรญฺเญ วิหรนฺติ, เต เนว ปานียํ อุปฏฺฐาเปนฺติ, น ปริโภชนียํ อุปฏฺฐาเปนฺติ, น อคฺคิํ อุปฏฺฐาเปนฺติ, น อรณิสหิตํ อุปฏฺฐาเปนฺติ, น นกฺขตฺตปทานิ ชานนฺติ, น ทิสาภาคํ ชานนฺติ. โจรา ตตฺถ คนฺตฺวา เต ภิกฺขู เอตทโวจุํ “อตฺถิ ภนฺเต ปานียนฺ”ติ. นตฺถาวุโสติ. อตฺถิ ภนฺเต ปริโภชนียนฺติ. นตฺตาวุโสติ. อตฺถิ ภนฺเต อคฺคีติ. นตฺถาวุโสติ. อตฺถิ ภนฺเต อรณิสหิตนฺติ. นตฺถาวุโสติ. ( )\csromansharedfootnote{c719e5be3e52b144}{(อตฺถิ ภนฺเต นกฺขตฺตปทานีติ, น ชานาม อาวุโสติ, อตฺถิ ภนฺเต ทิสาภาคนฺติ, น ชานาม อาวุโสติ.) สี, วิมติฏีกาย ปน สเมติ.} เกนชฺช ภนฺเต ยุตฺตนฺติ. น โข มยํ อาวุโส ชานามาติ. กตมายํ ภนฺเต ทิสาติ. น โข มยํ อาวุโส ชานามาติ. อถ โข เต โจรา “เนวิเมสํ ปานียํ อตฺถิ, น ปริโภชนียํ อตฺถิ, น อคฺคิ อตฺถิ, น อรณิสหิตํ อตฺถิ, น นกฺขตฺตปทานิ ชานนฺติ, น ทิสาภาคํ ชานนฺติ, โจรา อิเม, นยิเม ภิกฺขู”ติ อาโกเฏตฺวา ปกฺกมิํสุ. อถ โข เต ภิกฺขู ภิกฺขูนํ เอตมตฺถํ อาโรเจสุํ. ภิกฺขู ภควโต เอตมตฺถํ อาโรเจสุํ. อถ โข ภควา เอตสฺมิํ นิทาเน เอตสฺมิํ ปกรเณ ธมฺมิํ กถํ กตฺวา ภิกฺขู อามนฺเตสิ\csromandash{}}

\proseitem{368}{เตน หิ ภิกฺขเว อารญฺญิกานํ ภิกฺขูนํ วตฺตํ ปญฺญเปสฺสามิ, ยถา อารญฺญิเกหิ ภิกฺขูหิ สมฺมา วตฺติตพฺพํ. อารญฺญิเกน ภิกฺขเว ภิกฺขุนา กาลสฺเสว อุฏฺฐาย ปตฺตํ ถวิกาย ปกฺขิปิตฺวา อํเส \csromanfolio{386}อาลคฺเคตฺวา จีวรํ ขนฺเธ กริตฺวา อุปาหนา อาโรหิตฺวา ทารุภณฺฑํ มตฺติกาภณฺฑํ ปฏิสาเมตฺวา ทฺวารวาตปานํ ถเกตฺวา เสนาสนา โอตริตพฺพํ. “อิทานิ คามํ ปวิสิสฺสามี”ติ อุปาหนา โอมุญฺจิตฺวา นีจํ กตฺวา ปปฺโผเฏตฺวา ถวิกาย ปกฺขิปิตฺวา อํเส อาลคฺเคตฺวา ติมณฺฑลํ ปฏิจฺฉาเทนฺเตน ปริมณฺฑลํ นิวาเสตฺวา กายพนฺธนํ พนฺธิตฺวา สคุณํ กตฺวา สํฆาฏิโย ปารุปิตฺวา คณฺฐิกํ ปฏิมุญฺจิตฺวา โธวิตฺวา ปตฺตํ คเหตฺวา สาธุกํ อตรมาเนน คาโม ปวิสิตพฺโพ, สุปฺปฏิจฺฉนฺเนน อนฺตรฆเร คนฺตพฺพํ ฯเปฯ น ขมฺภกเตน อนฺตรฆเร คนฺตพฺพํ, น โอคุณฺฐิเตน อนฺตรฆเร คนฺตพฺพํ, น อุกฺกุฏิกาย อนฺตรฆเร คนฺตพฺพํ.}

\prose{นิเวสนํ ปวิสนฺเตน สลฺลกฺเขตพฺพํ “อิมินา ปวิสิสฺสามิ, อิมินา นิกฺขมิสฺสามี”ติ. นาติสหสา ปวิสิตพฺพํ, นาติสหสา นิกฺขมิตพฺพํ, นาติทูเร ฐาตพฺพํ, นาจฺจาสนฺเน ฐาตพฺพํ, นาติจิรํ ฐาตพฺพํ, นาติลหุํ นิวตฺติตพฺพํ, ฐิตเกน สลฺลกฺเขตพฺพํ “ภิกฺขํ ทาตุกามา วา อทาตุกามาวา”ติ. สเจ กมฺมํ วา นิกฺขิปติ, อาสนา วา วุฏฺฐาติ, กฏจฺฉุํ วา ปรามสติ, ภาชนํ วา ปรามสติ, ฐเปติ วา, “ทาตุกามสฺสา”ติ ฐาตพฺพํ, ภิกฺขาย ทิยฺยมานาย วาเมน หตฺเถน สํฆาฏิํ อุจฺจาเรตฺวา ทกฺขิเณน หตฺเถน ปตฺตํ ปณาเมตฺวา อุโภหิ หตฺเถหิ ปตฺตํ ปฏิคฺคเหตฺวา ภิกฺขา ปฏิคฺคเหตพฺพา, น จ ภิกฺขา\-ทายิกาย มุขํ อุลฺโลเกตพฺพํ, สลฺลกฺเขตพฺพํ “สูปํ ทาตุกามา วา อทาตุกามา วา”ติ. สเจ กฏจฺฉุ วา ปรามสติ, ภาชนํ วา ปรามสติ, ฐเปติ วา, “ทาตุกามสฺสา”ติ ฐาตพฺพํ. ภิกฺขาย ทินฺนาย สํฆาฏิยา ปตฺตํ ปฏิจฺฉาเทตฺวา สาธุกํ อตรมาเนน นิวตฺติตพฺพํ.}

\prose{สุปฺปฏิจฺฉนฺเนน อนฺตรฆเร คนฺตพฺพํ ฯเปฯ น อุกฺกุฏิกาย อนฺตรฆเร คนฺตพฺพํ. คามโต นิกฺขมิตฺวา ปตฺตํ ถวิกาย ปกฺขิปิตฺวา อํเส อาลคฺเคตฺวา จีวรํ สงฺฆริตฺวา สีเส กริตฺวา อุปาหนา อาโรหิตฺวา คนฺตพฺพํ.}

\prose{อารญฺญิเกน ภิกฺขเว ภิกฺขุนา ปานียํ อุปฏฺฐาเป\-ตพฺพํ, ปริโภชนียํ อุปฏฺฐาเป\-ตพฺพํ, อคฺคิ อุปฏฺฐาเปตพฺโพ, อรณิสหิตํ อุปฏฺฐาเป\-ตพฺพํ, กตฺตรทณฺโฑ อุปฏฺฐาเปตพฺโพ, นกฺขตฺตปทานิ อุคฺคเหตพฺพานิ สกลานิ วา เอกเทสานิ วา, ทิสากุสเลน ภวิตพฺพํ. อิทํ โข ภิกฺขเว อารญฺญิกานํ ภิกฺขูนํ วตฺตํ, ยถา อารญฺญิเกหิ ภิกฺขูหิ สมฺมา วตฺติตพฺพนฺติ.}

\csromanmatikaanchor{mtk.184}
\csromantocmarkref{section}{8. เสนาสนวตฺตกถา}{mtk.184}
\bookmark[dest={mtk.184},level=1]{8. เสนาสนวตฺตกถา}
\csromanheader{1}{8. เสนาสนวตฺต\-กถา}
\proseitem{369}{\csromanfolio{387}เตน โข ปน สมเยน สมฺพหุลา ภิกฺขู อชฺโฌกาเส จีวรกมฺมํ กโรนฺติ. ฉพฺพคฺคิยา ภิกฺขู ปฏิวาเต องฺคเณ\csromansharedfootnote{93cef606c70de5e0}{ปงฺคเณ (สี, สฺยา)} เสนาสนํ ปปฺโผเฏสุํ, ภิกฺขู รเชน โอกิริํสุ. เย เต ภิกฺขู อปฺปิจฺฉา ฯเปฯ เต อุชฺฌายนฺติ ขิยฺยนฺติ วิปาเจนฺติ “กถํ หิ นาม ฉพฺพคฺคิยา ภิกฺขู ปฏิวาเต องฺคเณ เสนาสนํ ปปฺโผเฏสฺสนฺติ, ภิกฺขู รเชน โอกิริํสู”ติ. อถ โข เต ภิกฺขู ภควโต เอตมตฺถํ อาโรเจสุํ ฯเปฯ. สจฺจํ กิร ภิกฺขเว ฉพฺพคฺคิยา ภิกฺขู ปฏิวาเต องฺคเณ เสนาสนํ ปปฺโผเฏนฺติ, ภิกฺขู รเชน โอกิริํสูติ. สจฺจํ ภควาติ ฯเปฯ วิครหิตฺวา ฯเปฯ ธมฺมิํ กตฺวา ภิกฺขู อามนฺเตสิ\csromandash{}}

\proseitem{370}{เตน หิ ภิกฺขเว ภิกฺขูนํ เสนาสนวตฺตํ ปญฺญเปสฺสามิ, ยถา ภิกฺขูหิ เสนาสเน สมฺมา วตฺติตพฺพํ. ยสฺมิํ วิหาเร วิหรติ, สเจ โส วิหาโร อุกฺลาโป โหติ, สเจ อุสฺสหติ โสเธตพฺโพ,\csromansymbolfootnote{*}{วิ 3. 60, 65; วิ 4. 375 ปิฏฺฐาทีสุปิ.}วิหารํ โสเธนฺเตน ปฐมํ ปตฺตจีวรํ นีหริตฺวา เอกมนฺตํ นิกฺขิปิตพฺพํ, นิสีทน\-ปจฺจตฺถรณํ นีหริตฺวา เอกมนฺตํ นิกฺขิปิตพฺพํ, ภิสิพิพฺโพหนํ นีหริตฺวา เอกมนฺตํ นิกฺขิปิตพฺพํ, มญฺโจ นีจํ กตฺวา สาธุกํ อปฺปฏิฆสนฺเตน อสํฆฏฺเฏนฺเตน กวาฏปิฏฺฐํ นีหริตฺวา เอกมนฺตํ นิกฺขิปิตพฺโพ, ปีฐํ นีจํ กตฺวา สาธุกํ อปฺปฏิฆํสนฺเตน อสํฆฏฺเฏนฺเตน กวาฏปิฏฺฐํ นีหริตฺวา เอกมนฺตํ นิกฺขิปิตพฺพํ, มญฺจ\-ปฏิปาทกา นีหริตฺวา เอกมนฺตํ นิกฺขิปิตพฺพา, เขฬมลฺลโก นีหริตฺวา เอกมนฺตํ นิกฺขิปิตพฺโพ, อปสฺเสนผลกํ นีหริตฺวา เอกมนฺตํ นิกฺขิปิตพฺพํ, ภูมตฺถรณํ ยถาปญฺญตฺตํ สลฺลกฺเขตฺวา นีหริตฺวา เอกมนฺตํ นิกฺขิปิตพฺพํ. สเจ วิหาเร สนฺตานกํ โหติ, อุลฺโลกา ปฐมํ โอหาเรตพฺพํ, อาโลกสนฺธิ\-กณฺณภาคา ปมชฺชิตพฺพา. สเจ เครุก\-ปริกมฺม\-กตา ภิตฺติ กณฺณกิตา โหติ, โจฬกํ เตเมตฺวา ปีเฬตฺวา ปมชฺชิตพฺพา. สเจ กาฬวณฺณกตา ภูมิ กณฺณกิตา โหติ, โจฬกํ เตเมตฺวา ปีเฬตฺวา ปมชฺชิตพฺพา. สเจ อกตา โหติ ภูมิ, อุทเกน ปริปฺโผสิตฺวา ปริปฺโผสิตฺวา สมฺมชฺชิตพฺพา “มา วิหาโร รเชน อุหญฺญี”ติ, สงฺการํ วิจินิตฺวา เอกมนฺตํ ฉฑฺเฑตพฺพํ.}

\prose{\csromanfolio{388}น ภิกฺขุสามนฺตา เสนาสนํ ปปฺโผเฏตพฺพํ, น วิหารสามนฺตา เสนาสนํ ปปฺโผเฏตพฺพํ, น ปานียสามนฺตา เสนาสนํ ปปฺโผเฏตพฺพํ, น ปริโภชนีย\-สามนฺตา เสนาสนํ ปปฺโผเฏตพฺพํ, น ปฏิวาเต องฺคเณ เสนาสนํ ปปฺโผเฏตพฺพํ, อโธวาเต เสนาสนํ ปปฺโผเฏตพฺพํ.}

\prose{ภูมตฺถรณํ เอกมนฺตํ โอตาเปตฺวา โสเธตฺวา ปปฺโผเฏตฺวา อติหริตฺวา ยถาปญฺญตฺตํ, ปญฺญเปตพฺพํ, มญฺจ\-ปฏิปาทกา เอกมนฺตํ โอตาเปตฺวา ปมชฺชิตฺวา อติหริตฺวา ยถาฐาเน ฐเปตพฺพา, มญฺโจ เอกมนฺตํ โอตาเปตฺวา โสเธตฺวา ปปฺโผเฏตฺวา นีจํ กตฺวา สาธุกํ อปฺปฏิฆํสนฺเตน อสํฆฏฺเฏนฺเตน กวาฏปิฏฺฐํ อติหริตฺวา ยถาปญฺญตฺตํ ปญฺญเปตพฺโพ, ปีฐํ เอกมนฺตํ โอตาเปตฺวา โสเธตฺวา ปปฺโผเฏตฺวา นีจํ กตฺวา สาธุกํ อปฺปฏิฆํสนฺเตน อสํฆฏฺเฏนฺเตน กวาฏปิฏฺฐํ อติหริตฺวา ยถาปญฺญตฺตํ ปญฺญเปตพฺพํ, ภิสิพิพฺโพหนํ เอกมนฺตํ โอตาเปตฺวา โสเธตฺวา ปปฺโผเฏตฺวา อติหริตฺวา ยถาปญฺญตฺตํ ปญฺญเปตพฺพํ, นิสีทน\-ปจฺจตฺถรณํ เอกมนฺตํ โอตาเปตฺวา โสเธตฺวา ปปฺโผเฏตฺวา อติหริตฺวา ยถาปญฺญตฺตํ ปญฺญเปตพฺพํ, เขฬมลฺลโก เอกมนฺตํ โอตาเปตฺวา ปมชฺชิตฺวา อติหริตฺวา ยถาฐาเน ฐเปตพฺโพ, อปสฺเสนผลกํ เอกมนฺตํ โอตาเปตฺวา ปมชฺชิตฺวา อติหริตฺวา ยถาฐาเน ฐเปตพฺพํ. ปตฺตจีวรํ นิกฺขิปิตพฺพํ, ปตฺตํ นิกฺขิปนฺเตน เอเกน หตฺเถน ปตฺตํ คเหตฺวา เอเกน หตฺเถน เหฏฺฐามญฺจํ วา เหฏฺฐาปีฐํ วา ปรามสิตฺวา ปตฺโต นิกฺขิปิตพฺโพ, น จ อนนฺตรหิตาย ภูมิยา ปตฺโต นิกฺขิปิตพฺโพ. จีวรํ นิกฺขิปนฺเตน เอเกน หตฺเถน จีวรํ คเหตฺวา เอเกน หตฺเถน จีวรวํสํ วา จีวรรชฺชุํ วา ปมชฺชิตฺวา ปารโต อนฺตํ โอรโต โภคํ กตฺวา จีวรํ นิกฺขิปิตพฺพํ.}

\prose{สเจ ปุรตฺถิมา สรชา วาตา วายนฺติ, ปุรตฺถิมา วาตปานา ถเกตพฺพา. สเจ ปจฺฉิมา สรชา วาตา วายนฺติ, ปจฺฉิมา วาตปานา ถเกตพฺพา. สเจ อุตฺตรา สรชา วาตา วายนฺติ, อุตฺตรา วาตปานา ถเกตพฺพา. สเจ ทกฺขิณา สรชา วาตา วายนฺติ, ทกฺขิณา วาตปานา ถเกตพฺพา. สเจ สีตกาโล โหติ, ทิวา วาตปานา วิวริตพฺพา รตฺติํ ถเกตพฺพา. สเจ อุณฺหกาโล โหติ, ทิวา วาตปานา ถเกตพฺพา, รตฺติํ วิวริตพฺพา.}

\prose{\csromanfolio{389}สเจ ปริเวณํ อุกฺลาปํ โหติ, ปริเวณํ สมฺมชฺชิตพฺพํ. สเจ โกฏฺฐโก อุกฺลาโป โหติ, โกฏฺฐโก สมฺมชฺชิตพฺโพ. สเจ อุปฏฺฐานสาลา อุกฺลาปา โหติ, อุปฏฺฐานสาลา สมฺมชฺชิตพฺโพ. สเจ อคฺคิสาลา อุกฺลาปา โหติ, อคฺคิสาลา สมฺมชฺชิตพฺพา. สเจ วจฺจกุฏิ อุกฺลาปา โหติ, วจฺจกุฏิ สมฺมชฺชิตพฺพา. สเจ ปานียํ น โหติ, ปานียํ อุปฏฺฐาเป\-ตพฺพํ. สเจ ปริโภชนียํ น โหติ, ปริโภชนียํ อุปฏฺฐาเป\-ตพฺพํ. สเจ อาจมนกุมฺภิยา อุทกํ น โหติ, อาจมนกุมฺภิยา อุทกํ อาสิญฺจิตพฺพํ.}

\prose{สเจ วุฑฺเฒน สทฺธิํ เอกวิหาเร วิหรติ, น วุฑฺฒํ อนาปุจฺฉา อุทฺเทโส ทาตพฺโพ, น ปริปุจฺฉา ทาตพฺพา, น สชฺฌาโย กาตพฺโพ, น ธมฺโม ภาสิตพฺโพ, น ปทีโป กาตพฺโพ, น ปทีโป วิชฺฌาเปตพฺโพ, น วาตปานา วิวริตพฺพา, น วาตปานา ถเกตพฺพา. สเจ วุฑฺเฒน สทฺธิํ เอกจงฺกเม จงฺกมติ, เยน วุฑฺโฒ เตน ปริวตฺติตพฺพํ, น จ วุฑฺโฒ สํฆาฏิกณฺเณน ฆฏฺเฏตพฺโพ. อิทํ โข ภิกฺขเว ภิกฺขูนํ เสนาสนวตฺตํ, ยถา ภิกฺขูหิ เสนาสเน สมฺมา วตฺติตพฺพนฺติ.}

\csromanmatikaanchor{mtk.185}
\csromantocmarkref{section}{9. ชนฺตาฆรวตฺตกถา}{mtk.185}
\bookmark[dest={mtk.185},level=1]{9. ชนฺตาฆรวตฺตกถา}
\csromanheader{1}{9. ชนฺตาฆรวตฺต\-กถา}
\csromanmatikaanchor{mtk.186}
\csromantocmarkref{section}{10. วจฺจกุฏิวตฺตกถา}{mtk.186}
\bookmark[dest={mtk.186},level=1]{10. วจฺจกุฏิวตฺตกถา}
\proseitem{371}{เตน โข ปน สมเยน ฉพฺพคฺคิยา ภิกฺขู ชนฺตาฆเร เถเรหิ ภิกฺขูหิ นิวาริยมานา อนาทริยํ ปฏิจฺจ ปหูตํ กฏฺฐํ อาโรเปตฺวา อคฺคิํ ทตฺวา ทฺวารํ ถเกตฺวา ทฺวาเร นิสีทนฺติ, ภิกฺขู\csromansharedfootnote{3f3437b44c337cd6}{เถรา จ ภิกฺขู (สฺยา, กํ)} อุณฺหาภิตตฺตา ทฺวารํ อลภมานา มุจฺฉิตา ปปตนฺติ. เย เต ภิกฺขู อปฺปิจฺฉา ฯเปฯ เต อุชฺฌายนฺติ ขิยฺยนฺติ วิปาเจนฺติ “กถํ หิ นาม ฉพฺพคฺคิยา ภิกฺขู ชนฺตาฆเร เถเรหิ ภิกฺขูหิ นิวาริยมานา อนาทริยํ ปฏิจฺจ ปหูตํ กฏฺฐํ อาโรเปตฺวา อคฺคิํ ทตฺวา ทฺวารํ ถเกตฺวา ทฺวาเร นิสีทิสฺสนฺติ, ภิกฺขู อุณฺหาภิตตฺตา ทฺวารํ อลภมานา มุจฺฉิตา ปปตนฺตี”ติ. อถ โข เต ภิกฺขู ภควโต เอตมตฺถํ อาโรเจสุํ ฯเปฯ. สจฺจํ กิร ภิกฺขเว ฉพฺพคฺคิยา ภิกฺขู ชนฺตาฆเร เถเรหิ ภิกฺขูหิ นิวาริยมานา อนาทริยํ ปฏิจฺจ ปหูตํ กฏฺฐํ อาโรเปตฺวา อคฺคิํ ทตฺวา ทฺวารํ ถเกตฺวา ทฺวาเร นิสีทนฺติ, ภิกฺขู อุณฺหาภิตตฺตา ทฺวารํ อลภมานา มุจฺฉิตา ปปตนฺตีติ. สจฺจํ ภควาติ ฯเปฯ วิครหิตฺวา ฯเปฯ ธมฺมิํ กถํ กตฺวา ภิกฺขู อามนฺเตสิ\csromandash{} \csromanfolio{390}น ภิกฺขเว ชนฺตาฆเร เถเรน ภิกฺขุนา นิวาริยมาเนน อนาทริยํ ปฏิจฺจ ปหูตํ กฏฺฐํ อาโรเปตฺวา อคฺคิ ทาตพฺโพ, โย ทเทยฺย, อาปตฺติ ทุกฺกฏสฺส. น ภิกฺขเว ทฺวารํ ถเกตฺวา ทฺวาเร นิสีทิตพฺพํ, โย นิสีเทยฺย, อาปตฺติ ทุกฺกฏสฺส.}

\proseitem{372}{เตน หิ ภิกฺขเว ภิกฺขูนํ ชนฺตาฆร\-วตฺตํ ปญฺญเปสฺสามิ, ยถา ภิกฺขูหิ ชนฺตาฆเร สมฺมา วตฺติตพฺพํ. โย ปฐมํ ชนฺตาฆรํ คจฺฉติ, สเจ ฉาริกา อุสฺสนฺนา โหติ, ฉาริกา ฉฑฺเฑตพฺพา. สเจ ชนฺตาฆรํ อุกฺลาปํ โหติ, ชนฺตาฆรํ สมฺมชฺชิตพฺพํ. สเจ ปริภณฺฑํ อุกฺลาปํ โหติ, ปริภณฺฑํ สมฺมชฺชิตพฺพํ. สเจ ปริเวณํ อุกฺลาปํ โหติ, ปริเวณํ สมฺมชฺชิตพฺพํ. สเจ โกฏฺฐโก อุกฺลาโป โหติ, โกฏฺฐโก สมฺมชฺชิตพฺโพ. สเจ ชนฺตาฆรสาลา อุกฺลาปา โหติ, ชนฺตาฆรสาลา สมฺมชฺชิตพฺพา.}

\prose{จุณฺณํ สนฺเนตพฺพํ, มตฺติกา เตเมตพฺพา, อุทกโทณิกาย อุทกํ อาสิญฺจิตพฺพํ, ชนฺตาฆรํ ปวิสนฺเตน มตฺติกาย มุขํ มกฺเขตฺวา ปุรโต จ ปจฺฉโต จ ปฏิจฺฉาเทตฺวา ชนฺตาฆรํ ปวิสิตพฺพํ, น เถเร ภิกฺขู อนุปขชฺช นิสีทิตพฺพํ, น นวา ภิกฺขู อาสเนน ปฏิพาหิตพฺพา. สเจ อุสฺสหติ, ชนฺตาฆเร เถรานํ ภิกฺขุนํ ปริกมฺมํ กาตพฺพํ, ชนฺตาฆรา นิกฺขมนฺเตน ชนฺตาฆรปีฐํ อาทาย ปุรโต จ ปจฺฉโต จ ปฏิจฺฉาเทตฺวา ชนฺตาฆรา นิกฺขมิตพฺพํ. สเจ อุสฺสหติ, อุทเกปิ เถรานํ ภิกฺขูนํ ปริกมฺมํ กาตพฺพํ, น เถรานํ ภิกฺขูนํ ปุรโตปิ นหายิตพฺพํ, น อุปริโตปิ นหายิตพฺพํ, นหาเตน อุตฺตรนฺเตน โอตรนฺตานํ มคฺโค ทาตพฺโพ. โย ปจฺฉา ชนฺตาฆรา นิกฺขมติ, สเจ ชนฺตาฆรํ จิกฺขลฺลํ โหติ, โธวิตพฺพํ, มตฺติกาโทณิกํ โธวิตฺวา ชนฺตาฆรปีฐํ ปฏิสาเมตฺวา อคฺคิํ วิชฺฌาเปตฺวา ทฺวารํ ถเกตฺวา ปกฺกมิตพฺพํ. อิทํ โข ติกฺขเว ภิกฺขูนํ ชนฺตาฆร\-วตฺตํ, ยถา ภิกฺขูหิ ชนฺตาฆเร สมฺมา วตฺติตพฺพนฺติ.}

\csromanheader{2}{10. \textbf{วจฺจกุฏิวตฺตตถา}}
\proseitem{373}{เตน โข ปน สมเยน อญฺญตโร ภิกฺขุ พฺราหฺมณชาติโก วจฺจํ กตฺวา น อิจฺฉติ อาจเมตุํ “โก อิมํ วสลํ ทุคฺคนฺธํ \csromanfolio{391}อามสิสฺสตี”ติ\csromansharedfootnote{41b0dfc74d2ae7af}{อาจมิสฺสตีติ (ก)}. ตสฺส วจฺจมคฺเค กิมิ สณฺฐาติ. อถ โข โส ภิกฺขุ ภิกฺขูนํ เอตมตฺถํ อาโรเจสิ. กิํ ปน ตฺวํ อาวุโส วจฺจํ กตฺวา น อาจเมสีติ. เอวมาวุโสติ. เย เต ภิกฺขู อปฺปิจฺฉา ฯเปฯ เต อุชฺฌายนฺติ ขิยฺยนฺติ วิปาเจนฺติ “กถํ หิ นาม ภิกฺขุ วจฺจํ กตฺวา น อาจเมสฺสตี”ติ. อถ โข เต ภิกฺขู ภควโต เอตมตฺถํ อาโรเจสุํ ฯเปฯ. สจฺจํ กิร ตฺวํ ภิกฺขุ วจฺจํ กตฺวา น อาจเมสีติ. สจฺจํ ภควาติ ฯเปฯ วิครหิตฺวา ฯเปฯ ธมฺมิํ กถํ กตฺวา ภิกฺขู อามนฺเตสิ “น ภิกฺขเว วจฺจํ กตฺวา สติ อุทเก นาจเมตพฺพํ, โย นาจเมยฺย, อาปตฺติ ทุกฺกฏสฺสา”ติ.}

\prose{เตน โข ปน สมเยน ภิกฺขู วจฺจกุฏิยา ยถาวุฑฺฒํ วจฺจํ กโรนฺติ, นวกา ภิกฺขู ปฐมตรํ อาคนฺตฺวา วจฺจิตา อาคเมนฺติ, เต วจฺจํ สนฺธาเรนฺตา มุจฺฉิตา ปปตนฺติ. ภควโต เอตมตฺถํ อาโรเจสุํ ฯเปฯ. สจฺจํ กิร ภิกฺขเว ฯเปฯ. สจฺจํ ภควาติ ฯเปฯ. น ภิกฺขเว วจฺจกุฏิยา ยถาวุฑฺฒํ วจฺโจ กาตพฺโพ, โย กเรยฺย อาปตฺติ ทุกฺกฏสฺส. อนุชานามิ ภิกฺขเว อาคต\-ปฏิปาฏิยา วจฺจํ กาตุนฺติ.}

\prose{เตน โข ปน สมเยน ฉพฺพคฺคิยา ภิกฺขู อติสหสาปิ วจฺจกุฏิํ ปวิสนฺติ, อุพฺภชิตฺวาปิ\csromansharedfootnote{b2033fbf8b64f599}{อุพฺภุชฺชิตฺวาปิ (สี), อุพฺภุชิตฺวา (สฺยา)} ปวิสนฺติ, นิตฺถุนนฺตาปิ วจฺจํ กโรนฺติ, ทนฺตกฏฺฐํ ขาทนฺตาปิ วจฺจํ กโรนฺติ, พหิทฺธาปิ วจฺจโทณิกาย วจฺจํ กโรนฺติ, พหิทฺธาปิ ปสฺสาว\-โทณิกาย ปสฺสาวํ กโรนฺติ, ปสฺสาว\-โทณิกายปิ เขฬํ กโรนฺติ, ผรุเสนปิ กฏฺเฐน อวเลขนฺติ, อวเลขน\-กฏฺฐมฺปิ วจฺจกูปมฺหิ ปาเตนฺติ, อติสหสาปิ นิกฺขมนฺติ, อุพฺภชิตฺวาปิ นิกฺขมนฺติ, จปุจปุ\-การกมฺปิ อาจเมนฺติ, อาจมนสราวเกปิ อุทกํ เสเสนฺติ. เย เต ภิกฺขู อปฺปิจฺฉา ฯเปฯ เต อุชฺฌายนฺติ ขิยฺยนฺติ วิปาเจนฺติ “กถํ หิ นาม ฉพฺพคฺคิยา ภิกฺขู อติสหสาปิ วจฺจกุฏิํ ปวิสิสฺสนฺติ, อุพฺภชิตฺวาปิ ปวิสิสฺสนฺติ, นิตฺถุนนฺตาปิ วจฺจํ กริสฺสนฺติ, ทนฺตกฏฺฐํ ขาทนฺตาปิ วจฺจํ กริสฺสนฺติ, พหิทฺธาวิ วจฺจโทณิกาย วจฺจํ กริสฺสนฺติ, พหิทฺธาปิ ปสฺสาว\-โทณิกาย ปสฺสาวํ กริสฺสนฺติ, ปสฺสาว\-โทณิกายปิ เขฬํ กริสฺสนฺติ, ผรุเสนปิ กฏฺเฐน อวเลขิสฺสนฺติ, อวเลขน\-กฏฺฐมฺปิ วจฺจกูปมฺหิ ปาเตสฺสนฺติ, อติสหสาปิ นิกฺขมิสฺสนฺติ, อุพฺภชิตฺวาปิ นิกฺขมิสฺสนฺติ, จปุจปุ\-การกมฺปิ อาจเมสฺสนฺติ, อาจมนสราวเกปิ อุทกํ เสเสสฺสนฺตี”ติ. อถ โข เต ภิกฺขู ภควโต \csromanfolio{392}เอตมตฺถํ อาโรเจสุํ ฯเปฯ. สจฺจํ กิร ภิกฺขเว ฯเปฯ. สจฺจํ ภควาติ ฯเปฯ วิครหิตฺวา ฯเปฯ ธมฺมิํ กถํ กตฺวา ภิกฺขู อามนฺเตสิ\csromandash{}}

\proseitem{374}{เตน หิ ภิกฺขเว ภิกฺขูนํ วจฺจกุฏิ\-วตฺตํ ปญฺญเปสฺสามิ, ยถา ภิกฺขูหิ วจฺจกุฏิยา สมฺมา วตฺติตพฺพํ. โย วจฺจกุฏิํ คจฺฉติ, เตน พหิ ฐิเตน\csromansharedfootnote{2a1530c25a3d7640}{พหิ ฐิเตน (สี, ก)} อุกฺกาสิตพฺพํ, อนฺโต นิสินฺเนนปิ อุกฺกาสิตพฺพํ, จีวรวํเส วา จีวรรชฺชุยา วา จีวรํ นิกฺขิปิตฺวา สาธุกํ อตรมาเนน วจฺจกุฏิ ปวิสิตพฺพา, นาติสหสา ปวิสิตพฺพา, น อุพฺภชิตฺวา ปวิสิตพฺพา, วจฺจปาทุกาย ฐิเตน อุพฺภชิตพฺพํ, น นิตฺถุนนฺเตน วจฺโจ กาตพฺโพ, น ทนฺตกฏฺฐํ ขาทนฺเตน วจฺโจ กาตพฺโพ, น พหิทฺธา วจฺจโทณิกาย วจฺโจ กาตพฺโพ, น พหิทฺธา ปสฺสาว\-โทณิกาย ปสฺสาโว กาตพฺโพ, น ปสฺสาว\-โทณิกาย เขโฬ กาตพฺโพ, น ผรุเสน กฏฺเฐน อวเลขิตพฺพํ, น อวเลขน\-กฏฺฐํ วจฺจกูปมฺหิ ปาเตตพฺพํ, วจฺจปาทุกาย ฐิเตน ปฏิจฺฉาเทตพฺพํ, นาติสหสา นิกฺขมิตพฺพํ, น อุพฺภชิตฺวา นิกฺขมิตพฺพํ, อาจมนปาทุกาย ฐิเตน อุพฺภชิตพฺพํ, น จปุจปุการกํ อาจเมตพฺพํ, น อาจมนสราวเก อุทกํ เสเสตพฺพํ, อาจมนปาทุกาย ฐิเตน ปฏิจฺฉาเทตพฺพํ.}

\prose{สเจ วจฺจกุฏิ อุหตา\csromansharedfootnote{fa794d48e5de5de9}{อูหตา (สี, สฺยา)} โหติ, โธวิตพฺพา. สเจ อวเลขน\-ปิธโร ปูโร โหติ, อวเลขน\-กฏฺฐํ ฉฑฺเฑตพฺพํ. สเจ วจฺจกุฏิ อุกฺลาปา โหติ, วจฺจกุฏิ สมฺมชฺชิตพฺพา. สเจ ปริภณฺฑํ อุกฺลาปํ โหติ, ปริภณฺฑํ สมฺมชฺชิตพฺพํ. สเจ ปริเวณํ อุกฺลาปํ โหติ, ปริเวณํ สมฺมชฺชิตพฺพํ. สเจ โกฏฺฐโก อุกฺลาโป โหติ, โกฏฺฐโก สมฺมชฺชิตพฺโพ. สเจ อาจมนกุมฺภิยา อุทกํ น โหติ, อาจมนกุมฺภิยา อุทกํ อาสิญฺจิตพฺพํ. อิทํ โข ภิกฺขเว ภิกฺขูนํ วจฺจกุฏิ\-วตฺตํ, ยถา ภิกฺขูหิ วจฺจกุฏิยา สมฺมา วตฺติตพฺพนฺติ.}

\csromanmatikaanchor{mtk.187}
\csromantocmarkref{section}{11. อุปชฺฌายวตฺตกถา}{mtk.187}
\bookmark[dest={mtk.187},level=1]{11. อุปชฺฌายวตฺตกถา}
\csromanheader{1}{11. อุปชฺฌายวตฺต\-กถา}
\proseitem{375}{เตน โข ปน สมเยน สทฺธิวิหาริกา อุปชฺฌาเยสุ น สมฺมา วตฺตนฺติ. เย เต ภิกฺขู อปฺปิจฺฉา ฯเปฯ เต อุชฺฌายนฺติ ขิยฺยนฺติ วิปาเจนฺติ “กถํ หิ นาม สทฺธิวิหาริกา อุปชฺฌาเยสุ น สมฺมา วตฺติสฺสนฺตี”ติ. อถ โข เต ภิกฺขู ภควโต เอตมตฺถํ อาโรเจสุํ ฯเปฯ. สจฺจํ กิร ภิกฺขเว สทฺธิวิหาริกา อุปชฺฌาเยสุ น สมฺมา วตฺตนฺตีติ. สจฺจํ ภควาติ. วิครหิ \csromanfolio{393}พุทฺโธ ภควา ฯเปฯ. กถํ หิ นาม ภิกฺขเว สทฺธิวิหาริกา อุปชฺฌาเยสุ น สมฺมา วตฺติสฺสนฺติ. เนตํ ภิกฺขเว อปฺปสนฺนานํ วา ปสาทาย ฯเปฯ วิครหิตฺวา ฯเปฯ ธมฺมิํ กตํ กตฺวา ภิกฺขู อามนฺเตสิ\csromandash{}}

\proseitem{376}{เตน หิ ภิกฺขเว สทฺธิ\-วิหาริกานํ อุปชฺฌาเยสุ วตฺตํ ปญฺญเปสฺสามิ, ยถา สทฺธิ\-วิหาริเกหิ อุปชฺฌาเยสุ สมฺมา วตฺติตพฺพํ.\csromansymbolfootnote{*}{วิ 3. 58 ปิฏฺเฐปิ.}สทฺธิ\-วิหาริเกน ภิกฺขเว อุปชฺฌายมฺหิ สมฺมา วตฺติตพฺพํ. ตตฺรายํ สมฺมาวตฺตนา\csromandash{}}

\prose{กาลสฺเสว อุฏฺฐาย อุปาหนา โอมุญฺจิตฺวา เอกํสํ อุตฺตราสงฺคํ กริตฺวา ทนฺตกฏฺฐํ ทาตพฺพํ, มุโขทกํ ทาตพฺพํ อาสนํ ปญฺญเปตพฺพํ. สเจ ยาคุ โหติ, ภาชนํ โธวิตฺวา ยาคุ อุปนาเมตพฺพา, ยาคุํ ปีตสฺส อุทกํ ทตฺวา ภาชนํ ปฏิคฺคเหตฺวา นีจํ กตฺวา สาธุกํ อปฺปฏิฆํสนฺเตน โธวิตฺวา ปฏิสาเมตพฺพํ, อุปชฺฌายมฺหิ วุฏฺฐิเต อาสนํ อุทฺธริตพฺพํ. สเจ โส เทโส อุกฺลาโป โหติ, โส เทโส สมฺมชฺชิตพฺโพ.}

\prose{สเจ อุปชฺฌาโย คามํ ปวิสิตุกาโม โหติ, นิวาสนํ ทาตพฺพํ, ปฏินิวาสนํ ปฏิคฺคเหตพฺพํ, กายพนฺธนํ ทาตพฺพํ, สคุณํ กตฺวา สํฆาฏิโย ทาตพฺพา, โธวิตฺวา ปตฺโต โสทโก\csromansharedfootnote{6a149cae96967c2d}{สอุทโก (ก)} ทาตพฺโพ. สเจ อุปชฺฌาโย ปจฺฉาสมณํ อากงฺขติ, ติมณฺฑลํ ปฏิจฺฉาเทนฺเตน ปริมณฺฑลํ นิวาเสตฺวา กายพนฺธนํ พนฺธิตฺวา สคุณํ กตฺวา สํฆาฏิโย ปารุปิตฺวา คณฺฐิกํ ปฏิมุญฺจิตฺวา โธวิตฺวา ปตฺตํ คเหตฺวา อุปชฺฌายสฺส ปจฺฉาสมเณน โหตพฺพํ, นาติทูเร คนฺตพฺพํ, นาจฺจาสนฺเน คนฺตพฺพํ, ปตฺต\-ปริยาปนฺนํ ปฏิคฺคเหตพฺพํ, น อุปชฺฌายสฺส ภณมานสฺส อนฺตรนฺตรา กถา โอปาเตตพฺพา, อุปชฺฌาโย อาปตฺติสามนฺตา ภณมาโน นิวาเรตพฺโพ.}

\prose{นิวตฺตนฺเตน ปฐมตรํ อาคนฺตฺวา อาสนํ ปญฺญเปตพฺพํ, ปาโททกํ ปาทปีฐํ ปาทกถลิกํ อุปนิกฺขิปิตพฺพํ, ปจฺจุคฺคนฺตฺวา ปตฺตจีวรํ ปฏิคฺคเหตพฺพํ ปฏินิวาสนํ ทาตพฺพํ, นิวาสนํ ปฏิคฺคเหตพฺพํ. สเจ จีวรํ สินฺนํ โหติ. มุหุตฺตํ อุเณฺห โอตาเปตพฺพํ, น จ อุเณฺห จีวรํ นิทหิตพฺพํ, จีวรํ สงฺฆริตพฺพํ, จีวรํ สงฺฆรนฺเตน จตุรงฺคุลํ กณฺณํ อุสฺสาเรตฺวา จีวรํ สงฺฆริตพฺพํ “มา มชฺเฌ ภงฺโค อโหสี”ติ, โอโภเค กายพนฺธนํ กาตพฺพํ.}

\prose{\csromanfolio{394}สเจ ปิณฺฑปาโต โหติ, อุปชฺฌาโย จ ภุญฺชิตุกาโม โหติ, อุทกํ ทตฺวา ปิณฺฑปาโต อุปนาเมตพฺโพ, อุปชฺฌาโย ปานีเยน ปุจฺฉิตพฺโพ, ภุตฺตาวิสฺส อุทกํ ทตฺวา ปตฺตํ ปฏิคฺคเหตฺวา นีจํ กตฺวา สาธุกํ อปฺปฏิฆํสนฺเตน โธวิตฺวา โวทกํ กตฺวา มุหุตฺตํ อุเณฺห โอตาเปตพฺโพ, น จ อุเณฺห ปตฺโต นิทหิตพฺโพ. ปตฺตจีวรํ นิกฺขิปิตพฺพํ ปตฺตํ นิกฺขิปนฺเตน เอเกน หตฺเถน ปตฺตํ คเหตฺวา เอเกน หตฺเถน เหฏฺฐามญฺจํ วา เหฏฺฐาปีฐํ วา ปรามสิตฺวา ปตฺโต นิกฺขิปิตพฺโพ, น จ อนนฺตรหิตาย ภูมิยา ปตฺโต นิกฺขิปิตพฺโพ. จีวรํ นิกฺขิปนฺเตน เอเกน หตฺเถน จีวรํ คเหตฺวา เอเกน หตฺเถน จีวรวํสํ วา จีวรรชฺชุํ วา ปมชฺชิตฺวา ปารโต อนฺตํ โอรโต โภคํ กตฺวา จีวรํ นิกฺขิปิตพฺพํ. อุปชฺฌายมฺหิ วุฏฺฐิเต อาสนํ อุทฺธริตพฺพํ, ปาโททกํ ปาทปีฐํ ปาทกถลิกํ ปฏิสาเมตพฺพํ. สเจ โส เทโส อุกฺลาโป โหติ, โสเทโส สมฺมชฺชิตพฺโพ.}

\prose{สเจ อุปชฺฌาโย นหายิตุกาโม โหติ, นหานํ ปฏิยาเทตพฺพํ. สเจ สีเตน อตฺโถ โหติ, สีตํ ปฏิยาเทตพฺพํ. สเจ อุเณฺหน อตฺโถ โหติ, อุณฺหํ ปฏิยาเทตพฺพํ.}

\prose{สเจ อุปชฺฌาโย ชนฺตาฆรํ ปวิสิตุกาโม โหติ, จุณฺณํ สนฺเนตพฺพํ, มตฺติกา เตเมตพฺพา, ชนฺตาฆรปีฐํ อาทาย อุปชฺฌายสฺส ปิฏฺฐิโต ปิฏฺฐิโต คนฺตฺวา ชนฺตาฆรปีฐํ ทตฺวา จีวรํ ปฏิคฺคเหตฺวา เอกมนฺตํ นิกฺขิปิตพฺพํ, จุณฺณํ ทาตพฺพํ, มตฺติกา ทาตพฺพา. สเจ อุสฺสหติ, ชนฺตาฆรํ ปวิสิตพฺพํ, ชนฺตาฆรํ ปวิสนฺเตน มตฺติกาย มุขํ มกฺเขตฺวา ปุรโต จ ปจฺฉโต จ ปฏิจฺฉาเทตฺวา ชนฺตาฆรํ ปวิสิตพฺพํ, น เถเร ภิกฺขู อนุปขชฺช นิสีทิตพฺพํ, น นวา ภิกฺขู อาสเนน ปฏิพาหิตพฺพา, ชนฺตาฆเร อุปชฺฌายสฺส ปริกมฺมํ กาตพฺพํ, ชนฺตาฆรา นิกฺขมนฺเตน ชนฺตาฆรปีฐํ อาทาย ปุรโต จ ปจฺฉโต จ ปฏิจฺฉาเทตฺวา ชนฺตาฆรา นิกฺขมิตพฺพํ.}

\prose{อุทเกปิ อุปชฺฌายสฺส ปริกมฺมํ กาตพฺพํ, นหาเตน ปฐมตรํ อุตฺตริตฺวา อตฺตโน คตฺตํ โวทกํ กตฺวา นิวาเสตฺวา อุปชฺฌายสฺส คตฺตโต อุทกํ ปมชฺชิตพฺพํ, นิวาสนํ ทาตพฺพํ, สํฆาฏิ ทาตพฺโพ, ชนฺตาฆรปีฐํ อาทาย ปฐมตรํ อาคนฺตฺวา อาสนํ ปญฺญเปตพฺพํ, ปาโททกํ ปาทปิฐํ ปาทกถลิกํ อุปนิกฺขิปิตพฺพํ, อุปชฺฌาโย ปานีเยน ปุจฺฉิตพฺโพ. สเจ อุทฺทิสาเปตุกาโม โหติ อุทฺทิสิตพฺโพ. สเจ ปริปุจฺฉิตุ\-กาโม โหติ ปริปุจฺฉิตพฺโพ.}

\prose{\csromanfolio{395}ยสฺมิํ วิหาเร อุปชฺฌาโย วิหรติ, สเจ โส วิหาโร อุกฺลาโป โหติ, สเจ อุสฺสหติ, โสเธตพฺโพ, วิหารํ โสเธนฺเตน ปฐมํ ปตฺตจีวรํ นีหริตฺวา เอกมนฺตํ นิกฺขิปิตพฺพํ, นิสีทน\-ปจฺจตฺถรณํ นีหริตฺวา เอกมนฺตํ นิกฺขิปิตพฺพํ, ภิสิพิพฺโพหนํ นีหริตฺวา เอกมนฺตํ นิกฺขิปิตพฺพํ, มญฺโจ นีจํ กตฺวา สาธุกํ อปฺปฏิฆํสนฺเตน อสํฆฏฺเฏนฺเตน กวาฏปิฏฺฐํ นีหริตฺวา เอกมนฺตํ นิกฺขิปิตพฺโพ, ปีฐํ นีจํ กตฺวา สาธุกํ อปฺปฏิฆํสนฺเตน อสํฆฏฺเฏนฺเตน กวาฏปิฏฺฐํ นีหริตฺวา เอกมนฺตํ นิกฺขิปิตพฺพํ, มญฺจ\-ปฏิปาทกา นีหริตฺวา เอกมนฺตํ นิกฺขิปิตพฺพา, เขฬมลฺลโก นีหริตฺวา เอกมนฺตํ นิกฺขิปิตพฺโพ, อปสฺเสนผลกํ นีหริตฺวา เอกมนฺตํ นิกฺขิปิตพฺพํ ภูมตฺถรณํ ยถาปญฺญตฺตํ สลฺลกฺเขตฺวา นีหริตฺวา เอกมนฺตํ นิกฺขิปิตพฺพํ. สเจ วิหาเร สนฺตานกํ โหติ, อุลฺโลกา ปฐมํ โอหาเรตพฺพํ, อาโลกสนฺธิ\-กณฺณภาคา ปมชฺชิตพฺพา. สเจ เครุก\-ปริกมฺม\-กตา ภิตฺติ กณฺณกิตา โหติ, โจฬกํ เตเมตฺวา ปีเฬตฺวา ปมชฺชิตพฺพา. สเจ กาฬวณฺณกตา ภูมิ กณฺณกิตา โหติ, โจฬกํ เตเมตฺวา ปีเฬตฺวา ปมชฺชิตพฺพา. สเจ อกตา โหติ ภูมิ, อุทเกน ปริปฺโผสิตฺวา ปริปฺโผสิตฺวา สมฺมชฺชิตพฺพา “มา วิหาโร รเชน อุหญฺญี”ติ, สงฺการํ วิจินิตฺวา เอกมนฺตํ ฉฑฺเฑตพฺพํ.}

\prose{ภูมตฺถรณํ โอตาเปตฺวา โสเธตฺวา ปปฺโผเฏตฺวา อติหริตฺวา ยถาปญฺญตฺตํ ปญฺญเปตพฺพํ, มญฺจ\-ปฏิปาทกา โอตาเปตฺวา ปมชฺชิตฺวา อติหริตฺวา ยถาฐาเน ฐเปตพฺพา, มญฺโจ โอตาเปตฺวา โสเธตฺวา ปปฺโผเฏตฺวา นีจํ กตฺวา สาธุกํ อปฺปฏิฆํสนฺเตน อสํฆฏฺเฏนฺเตน กวาฏปิฏฺฐํ อติหริตฺวา ยถาปญฺญตฺตํ ปญฺญเปตพฺโพ, ปีฐํ โอตาเปตฺวา โสเธตฺวา ปปฺโผเฏตฺวา นีจํ กตฺวา สาธุกํ อปฺปฏิฆํสนฺเตน อสํฆฏฺเฏนฺเตน กวาฏปิฏฺฐํ อติหริตฺวา ยถาปญฺญตฺตํ ปญฺญเปตพฺพํ, ภิสิพิพฺโพหนํ โอตาเปตฺวา โสเธตฺวา ปปฺโผเฏตฺวา อติหริตฺวา ยถาปญฺญตฺตํ ปญฺญเปตพฺพํ, นิสีทน\-ปจฺจตฺถรณํ โอตาเปตฺวา โสเธตฺวา ปปฺโผเฏตฺวา อติหริตฺวา ยถาปญฺญตฺตํ ปญฺญเปตพฺพํ, เขฬมลฺลโก โอตาเปตฺวา ปมชฺชิตฺวา อติหริตฺวา ยถาฐาเน ฐเปตพฺโพ, อปสฺเสนผลกํ โอตาเปตฺวา ปมชฺชิตฺวา อติหริตฺวา ยถาฐาเน ฐเปตพฺพํ, ปตฺตจีวรํ นิกฺขิปิตพฺพํ, ปตฺตํ นิกฺขิปนฺเตน เอเกน หตฺเถน ปตฺตํ คเหตฺวา เอเกน หตฺเถน เหฏฺฐามญฺจํ วา เหฏฺฐาปีฐํ วา ปรามสิตฺวา \csromanfolio{396}ปตฺโต นิกฺขิปิตพฺโพ, น จ อนนฺตรหิตาย ภูมิยา ปตฺโต นิกฺขิปิตพฺโพ, จีวรํ นิกฺขิปนฺเตน เอเกน หเถน จีวรํ คเหตฺวา เอเกน หตฺเถน จีวรํวํสํ วา จีวรรชฺชุํ วา ปมชฺชิตฺวา ปารโต อนฺตํ โอรโต โภคํ กตฺวา จีวรํ นิกฺขิปิตพฺพํ.}

\prose{สเจ ปุรตฺถิมา สรชา วาตา วายนฺติ, ปุรตฺถิมา วาตปานา ถเกตพฺพา. สเจ ปจฺฉิมา สรชา วาตา วายนฺติ, ปจฺฉิมา วาตปานา ถเกตพฺพา. สเจ อุตฺตรา สรชา วาตา วายนฺติ, อุตฺตรา วาตปานา ถเกตพฺพา. สเจ ทกฺขิณา สรชา วาตา วายนฺติ, ทกฺขิณา วาตปานา ถเกตพฺพา. สเจ สีตกาโล โหติ, ทิวา วาตปานา วิวริตพฺพา, รตฺติํ ถเกตพฺพา. สเจ อุณฺหกาโล โหติ, ทิวา วา ตปานา ถเกตพฺพา, รตฺติํ วิวริตพฺพา.}

\prose{สเจ ปริเวณํ อุกฺลาปํ โหติ, ปริเวณํ สมฺมชฺชิตพฺพํ. สเจ โกฏฺฐโก อุกฺลาโป โหติ, โกฏฺฐโก สมฺมชฺชิตพฺโพ. สเจ อุปฏฺฐานสาลา อุกฺลาปา โหติ, อุปฏฺฐานสาลา สมฺมชฺชิตพฺพา. สเจ อคฺคิสาลา อุกฺลาปา โหติ, อคฺคิสาลา สมฺมชฺชิตพฺพา. สเจ วจฺจกุฏิ อุกฺลาปา โหติ, วจฺจกุฏิ สมฺมชฺชิตพฺพา. สเจ ปานียํ น โหติ, ปานียํ อุปฏฺฐาเป\-ตพฺพํ. สเจ ปริโภชนียํ น โหติ, ปริโภชนียํ อุปฏฺฐาเป\-ตพฺพํ. สเจ อาจมนกุมฺภิยา อุทกํ น โหติ, อาจมนกุมฺภิยา อุทกํ อาสิญฺจิตพฺพํ.}

\prose{สเจ อุปชฺฌายสฺส อนภิรติ อุปฺปนฺนา โหติ, สทฺธิ\-วิหาริเกน วูปกาเสตพฺโพ วูปกาสาเป\-ตพฺโพ ธมฺมกถา วาสฺส กาตพฺพา. สเจ อุปชฺฌายสฺส กุกฺกุจฺจํ อุปฺปนฺนํ โหติ, สทฺธิ\-วิหาริเกน วิโนเทตพฺพํ วิโนทาเปตพฺพํ ธมฺมกถา วาสฺส กาตพฺพา. สเจ อุปชฺฌายสฺส ทิฏฺฐิคตํ อุปฺปนฺนํ โหติ, สทฺธิ\-วิหาริเกน วิเวเจตพฺพํ วิเวจาเปตพฺพํ ธมฺมกถา วาสฺส กาตพฺพา. สเจ อุปชฺฌาโย ครุธมฺมํ อชฺฌาปนฺโน โหติ ปริวาสารโห, สทฺธิ\-วิหาริเกน อุสฺสุกฺกํ กาตพฺพํ “กินฺติ นุ โข สํโฆ อุปชฺฌายสฺส ปริวาสํ ทเทยฺยา”ติ. สเจ อุปชฺฌาโย มูลาย\-ปฏิกสฺสนา\-รโห โหติ, สทฺธิ\-วิหาริเกน อุสฺสุกฺกํ กาตพฺพํ “กินฺติ นุ โข สํโฆ อุปชฺฌายํ มูลาย ปฏิกสฺเสยฺยา”ติ. สเจ อุปชฺฌาโย มานตฺตารโห โหติ, สทฺธิ\-วิหาริเกน อุสฺสุกฺกํ กาตพฺพํ “กินฺติ นุ โข สํโฆ อุปชฺฌายสฺส มานตฺตํ ทเทยฺยา”ติ. สเจ อุปชฺฌาโย อพฺภานารโห โหติ, สทฺธิ\-วิหาริเกน อุสฺสุกฺกํ กาตพฺพํ “กินฺติ นุ \csromanfolio{397}โข สํโฆ อุปชฺฌายํ อพฺเภยฺยา”ติ. สเจ สํโฆ อุปชฺฌายสฺส กมฺมํ กตฺตุกาโม โหติ ตชฺชนียํ วา นิยสฺสํ วา ปพฺพาชนียํ วา ปฏิสารณียํ วา อุกฺเขปนียํ วา, สทฺธิ\-วิหาริเกน อุสฺสุกฺกํ กาตพฺพํ “กินฺติ นุ โข สํโฆ อุปชฺฌายสฺส กมฺมํ น กเรยฺย, ลหุกาย วา ปริณาเมยฺยา”ติ. กตํ วา ปนสฺส โหติ สํเฆน กมฺมํ ตชฺชนียํ วา นิยสฺสํ วา ปพฺพาชนียํ วา ปฏิสารณียํ วา อุกฺเขปนียํ วา, สทฺธิ\-วิหาริเกน อุสฺสุกฺกํ กาตพฺพํ “กินฺติ นุ โข อุปชฺฌาโย สมฺมา วตฺเตยฺย, โลมํ ปาเตยฺย, เนตฺถารํ วตฺเตยฺย, สํโฆ ตํ กมฺมํ ปฏิปฺปสฺสมฺเภยฺยา”ติ.}

\prose{สเจ อุปชฺฌายสฺส จีวรํ โธวิตพฺพํ โหติ, สทฺธิ\-วิหาริเกน โธวิตพฺพํ, อุสฺสุกฺกํ วา กาตพฺพํ “กินฺติ นุ โข อุปชฺฌายสฺส จีวรํ โธวิเยถา”ติ. สเจ อุปชฺฌายสฺส จีวรํ กาตพฺพํ โหติ, สทฺธิ\-วิหาริเกน กาตพฺพํ, อุสฺสุกฺกํ วา กาตพฺพํ “กินฺติ นุ โข อุปชฺฌายสฺส จีวรํ กริเยถา”ติ. สเจ อุปชฺฌายสฺส รชนา ปจิตพฺพา โหติ, สทฺธิ\-วิหาริเกน ปจิตพฺพา, อุสฺสุกฺกํ วา กาตพฺพํ “กินฺติ นุ โข อุปชฺฌายสฺส รชนํ ปจิเยถา”ติ. สเจ อุปชฺฌายสฺส จีวรํ รชิตพฺพํ\csromansharedfootnote{d0f4cbe6aeda3cb2}{รเชตพฺพํ (สฺยา)} โหติ, สทฺธิ\-วิหาริเกน รชิตพฺพํ, อุสฺสุกฺกํ วา กาตพฺพํ “กินฺติ นุ โข อุปชฺฌายสฺส จีวรํ รชิเยถา”ติ. จีวรํ รชนฺเตน\csromansharedfootnote{2f39cbba58fcb43c}{รเชนฺเตน (สฺยา)} สาธุกํ สมฺปริวตฺตกํ สมฺปริวตฺตกํ รชิตพฺพํ, น จ อจฺฉินฺเน เถเว ปกฺกมิตพฺพํ.}

\prose{น อุปชฺฌายํ อนาปุจฺฉา เอกจฺจสฺส ปตฺโต ทาตพฺโพ, น เอกจฺจสฺส ปตฺโต ปฏิคฺคเหตพฺโพ, น เอกจฺจสฺส จีวรํ ทาตพฺพํ, น เอกจฺจสฺส จีวรํ ปฏิคฺคเหตพฺพํ, น เอกจฺจสฺส ปริกฺขาโร ทาตพฺโพ, น เอกจฺจสฺส ปริกฺขาโร ปฏิคฺคเหตพฺโพ, น เอกจฺจสฺส เกสา เฉเทตพฺพา\csromansharedfootnote{4c72805cf12f1b7a}{เฉตฺตพฺพา (สี), เฉทิตพฺพา (ก)}, น เอกจฺเจน เกสา เฉทาเปตพฺพา, น เอกจฺจสฺส ปริกมฺมํ กาตพฺพํ, น เอกจฺเจน ปริกมฺมํ การาเปตพฺพํ, น เอกจฺจสฺส เวยฺยาวจฺโจ\csromansharedfootnote{2cee68ae19c794e5}{เวยฺยาวจฺจํ (ก)} กาตพฺโพ, น เอกจฺเจน เวยฺยาวจฺโจ การาเปตพฺโพ, น เอกจฺจสฺส ปจฺฉาสมเณน โหตพฺพํ, น เอกจฺโจ ปจฺฉาสมโณ อาทาตพฺโพ, น เอกจฺจสฺส ปิณฺฑปาโต นีหริตพฺโพ, น เอกจฺเจน ปิณฺฑปาโต นีหราเปตพฺโพ, น อุปชฺฌายํ อนาปุจฺฉา คาโม ปวิสิตพฺโพ, น สุสานํ คนฺตพฺพํ, น ทิสา ปกฺกมิตพฺพา. สเจ อุปชฺฌาโย คิลาโน โหติ, ยาวชีวํ \csromanfolio{398}อุปฏฺฐาตพฺโพ, วุฏฺฐานมสฺส อาคเมตพฺพํ. อิทํ โข ภิกฺขเว สทฺธิ\-วิหาริกานํ อุปชฺฌาเยสุ วตฺตํ, ยถา สทฺธิ\-วิหาริเกหิ อุปชฺฌาเยสุ สมฺมา วตฺติตพฺพนฺติ.}

\csromanmatikaanchor{mtk.188}
\csromantocmarkref{section}{12. สทฺธิวิหาริกวตฺตกถา}{mtk.188}
\bookmark[dest={mtk.188},level=1]{12. สทฺธิวิหาริกวตฺตกถา}
\csromanheader{1}{12. สทฺธิวิหาริกวตฺต\-กถา}
\proseitem{377}{เตน โข ปน สมเยน อุปชฺฌายา สทฺธิ\-วิหาริเกสุ น สมฺมา วตฺตนฺติ. เย เต ภิกฺขู อปฺปิจฺฉา ฯเปฯ เต อุชฺฌายนฺติ ขิยฺยนฺติ วิปาเจนฺติ “กถํ หิ นาม อุปชฺฌายา สทฺธิ\-วิหาริเกสุ น สมฺมา วตฺติสฺสนฺตี”ติ. อถ โข เต ภิกฺขู ภควโต เอตมตฺถํ อาโรเจสุํ ฯเปฯ. สจฺจํ กิร ภิกฺขเว อุปชฺฌายา สทฺธิ\-วิหาริเกสุ น สมฺมา วตฺตนฺตีติ. สจฺจํ ภควาติ ฯเปฯ วิครหิตฺวา ฯเปฯ ธมฺมิํ กถํ กตฺวา ภิกฺขู อามนฺเตสิ\csromandash{}}

\proseitem{378}{เตน หิ ภิกฺขเว อุปชฺฌายานํ สทฺธิ\-วิหาริเกสุ วตฺตํ ปญฺญเปสฺสามิ, ยถา อุปชฺฌาเยหิ สทฺธิ\-วิหาริเกสุ สมฺมา วตฺติตพฺพํ.\csromansymbolfootnote{*}{วิ 3. 63 ปิฏฺเฐปิ.}อุปชฺฌาเยน ภิกฺขเว สทฺธิ\-วิหาริกมฺหิ สมฺมา วตฺติตพฺพํ, ตตฺรายํ สมฺมาวตฺตนา\csromandash{}}

\prose{อุปชฺฌาเยน ภิกฺขเว สทฺธิวิหาริโก สงฺคเหตพฺโพ อนุคฺคเหตพฺโพ อุทฺเทเสน ปริปุจฺฉาย โอวาเทน อนุสาสนิยา. สเจ อุปชฺฌายสฺส ปตฺโต โหติ สทฺธิ\-วิหาริกสฺส ปตฺโต น โหติ, อุปชฺฌาเยน สทฺธิ\-วิหาริกสฺส ปตฺโต ทาตพฺโพ, อุสฺสุกฺกํ วา กาตพฺพํ “กินฺติ นุ โข สทฺธิ\-วิหาริกสฺส ปตฺโต อุปฺปชฺชิเยถา”ติ. สเจ อุปชฺฌายสฺส จีวรํ โหติ สทฺธิ\-วิหาริกสฺส จีวรํ น โหติ, อุปชฺฌาเยน สทฺธิ\-วิหาริกสฺส จีวรํ ทาตพฺพํ, อุสฺสุกฺกํ วา กาตพฺพํ “กินฺติ นุ โข สทฺธิ\-วิหาริกสฺส จีวรํ อุปฺปชฺชิเยถา”ติ. สเจ อุปชฺฌายสฺส ปริกฺขาโร โหติ สทฺธิ\-วิหาริกสฺส ปริกฺขาโร น โหติ, อุปชฺฌาเยน สทฺธิ\-วิหาริกสฺส ปริกฺขาโร ทาตพฺโพ, อุสฺสุกฺกํ วา กาตพฺพํ “กินฺติ นุ โข สทฺธิ\-วิหาริกสฺส ปริกฺขาโร อุปฺปชฺชิเยถา”ติ.}

\prose{สเจ สทฺธิวิหาริโก คิลาโน โหติ, กาลสฺเสว อุฏฺฐาย ทนฺตกฏฺฐํ ทาตพฺพํ, มุโขทกํ ทาตพฺพํ, อาสนํ ปญฺญเปตพฺพํ. สเจ ยาคุ โหติ, ภาชนํ โธวิตฺวา ยาคุ อุปนาเมตพฺพา, ยาคุํ ปีตสฺส \csromanfolio{399}อุทกํ ทตฺวา ภาชนํ ปฏิคฺคเหตฺวา นีจํ กตฺวา สาธุกํ อปฺปฏิฆํสนฺเตน โธวิตฺวา ปฏิสาเมตพฺพํ สทฺธิ\-วิหาริกมฺหิ วุฏฺฐิเต อาสนํ อุทฺธริตพฺพํ. สเจ โส เทโส อุกฺลาโป โหติ, โส เทโส สมฺมชฺชิตพฺโพ.}

\prose{สเจ สทฺธิวิหาริโก คามํ ปวิสิตุกาโม โหติ, นิวาสนํ ทาตพฺพํ, ปฏินิวาสนํ ปฏิคฺคเหตพฺพํ, กายพนฺธนํ ทาตพฺพํ, สคุณํ กตฺวา สํฆาฏิโย ทาตพฺพา, โธวิตฺวา ปตฺโต โสทโก ทาตพฺโพ.}

\prose{“เอตฺตาวตา นิวตฺติสฺสตี”ติ อาสนํ ปญฺญเปตพฺพํ, ปาโททกํ ปาทปีฐํ ปาทกถลิกํ อุปนิกฺขิปิตพฺพํ. ปจฺจุคฺคนฺตฺวา ปตฺตจีวรํ ปฏิคฺคเหตพฺพํ ปฏินิวาสนํ ทาตพฺพํ, นิวาสนํ ปฏิคฺคเหตพฺพํ. สเจ จีวรํ สินฺนํ โหติ, มุหุตฺตํ อุเณฺห โอตาเปตพฺพํ, น จ อุเณฺห จีวรํ นิทหิตพฺพํ, จีวรํ สงฺฆริตพฺพํ, จีวรํ สงฺฆรนฺเตน จตุรงฺคุลํ กณฺณํ อุสฺสาเรตฺวา จีวรํ สงฺฆริตพฺพํ “มา มชฺเฌ ภงฺโค อโหสี”ติ, โอโภเค กายพนฺธนํ กาตพฺพํ.}

\prose{สเจ ปิณฺฑปาโต โหติ สทฺธิวิหาริโก จ ภุญฺชิตุกาโม โหติ, อุทกํ ทตฺวา ปิณฺฑปาโต อุปนาเมตพฺโพ, สทฺธิวิหาริโกปานีเยน ปุจฺฉิตพฺโพ, ภุตฺตาวิสฺส อุทกํ ทตฺวา ปตฺตํ ปฏิคฺคเหตฺวา นีจํ กตฺวา สาธุกํ อปฺปฏิฆํสนฺเตน โธวิตฺวา โวทกํ กตฺวา มุหุตฺตํ อุเณฺห โอตาเปตพฺโพ, น จ อุเณฺห ปตฺโต นิทหิตพฺโพ. ปตฺตจีวรํ นิกฺขิปิตพฺพํ, ปตฺตํ นิกฺขิปนฺเตน เอเกน หตฺเถน ปตฺตํ คเหตฺวา เอเกน หตฺเถน เหฏฺฐามญฺจํ วา เหฏฺฐาปีฐํ วา ปรามสิตฺวา ปตฺโต นิกฺขิปิตพฺโพ, น จ อนนฺตรหิตาย ภูมิยา ปตฺโต นิกฺขิปิตพฺโพ. จีวรํ นิกฺขิปนฺเตน เอเกน หตฺเถน จีวรํ คเหตฺวา เอเกน หตฺเถน จีวรวํสํ วา จีวรรชฺชุํ วา ปมชฺชิตฺวา ปารโต อนฺตํ โอรโต โภคํ กตฺวา จีวรํ นิกฺขิปิตพฺพํ. สทฺธิ\-วิหาริกมฺหิ วุฏฺฐิเต อาสนํ อุทฺธริตพฺพํ, ปาโททกํ ปาทปีฐํ ปาทกถลิกํ ปฏิสาเมตพฺพํ. สเจ โส เทโส อุกฺลาโป โหติ, โส เทโส สมฺมชฺชิตพฺโพ.}

\prose{สเจ สทฺธิวิหาริโก นหายิตุกาโม โหติ, นหานํ ปฏิยาเทตพฺพํ. สเจ สีเตน อตฺโถ โหติ, สีตํ ปฏิยาเทตพฺพํ. สเจ อุเณฺหน อตฺโถ โหติ, อุณฺหํ ปฏิยาเทตพฺพํ.}

\prose{\csromanfolio{400}สเจ สทฺธิวิหาริโก ชนฺตาฆรํ ปวิสิตุกาโม โหติ, จุณฺณํ สนฺเนตพฺพํ, มตฺติกา เตเมตพฺพา, ชนฺตาฆรปีฐํ อาทาย\csromansharedfootnote{c7bc02be51ad6982}{อาทาย สทฺธิวิหาริกสฺส ปิฏฺฐิโต ปิฏฺฐิโต (ก)} คนฺตฺวา ชนฺตาฆรปีฐํ ทตฺวา จีวรํ ปฏิคฺคเหตฺวา เอกมนฺตํ นิกฺขิปิตพฺพํ, จุณฺณํ ทาตพฺพํ, มตฺติกา ทาตพฺพา. สเจ อุสฺสหติ, ชนฺตาฆรํ ปวิสิตพฺพํ, ชนฺตาฆรํ ปวิสนฺเตน มตฺติกาย มุขํ มกฺเขตฺวา ปุรโต จ ปจฺฉโต จ ปฏิจฺฉาเทตฺวา ชนฺตาฆรํ ปวิสิตพฺพํ, น เถเร ภิกฺขู อนุปขชฺช นิสีทิตพฺพํ, น นวา ภิกฺขู อาสเนน ปฏิพาหิตพฺพา, ชนฺตาฆเร สทฺธิ\-วิหาริกสฺส ปริกมฺมํ กาตพฺพํ, ชนฺตาฆรา นิกฺขมนฺเตน ชนฺตาฆรปีฐํ อาทาย ปุรโต จ ปจฺฉโต จ ปฏิจฺฉาเทตฺวา ชนฺตาฆรา นิกฺขมิตพฺพํ.}

\prose{อุทเกปิ สทฺธิ\-วิหาริกสฺส ปริกมฺมํ กาตพฺพํ, นหาเตน ปฐมตรํ อุตฺตริตฺวา อตฺตโน คตฺตํ โวทกํ กตฺวา นิวาเสตฺวา สทฺธิ\-วิหาริกสฺส คตฺตโต อุทกํ ปมชฺชิตพฺพํ, นิวาสนํ ทาตพฺพํ, สํฆาฏิ ทาตพฺพา, ชนฺตาฆรปีฐํ อาทาย ปฐมตรํ อาคนฺตฺวา อาสนํ ปญฺญเปตพฺพํ, ปาโททกํ ปาทปีฐํ ปาทกถลิกํ อุปนิกฺขิปิตพฺพํ, สทฺธิวิหาริโก ปานีเยน ปุจฺฉิตพฺโพ.}

\prose{ยสฺมิํ วิหาเร สทฺธิวิหาริโก วิหรติ, สเจ โส วิหาโร อุกฺลาโป โหติ, สเจ อุสฺสหติ, โสเธตพฺโพ, วิหารํ โสเธนฺเตน ปฐมํ ปตฺตจีวรํ นีหริตฺวา เอกมนฺตํ นิกฺขิปิตพฺพํ ฯเปฯ. สเจ อาจมนกุมฺภิยา อุทกํ น โหติ, อาจมนกุมฺภิยา อุทกํ อาสิญฺจิตพฺพํ.}

\prose{สเจ สทฺธิ\-วิหาริกสฺส อนภิรติ อุปฺปนฺนา โหติ, อุปชฺฌาเยน วูปกาเสตพฺโพ วูปกาสาเป\-ตพฺโพ ธมฺมกถา วาสฺส กาตพฺพา. สเจ สทฺธิ\-วิหาริกสฺส กุกฺกุจฺจํ อุปฺปนฺนํ โหติ, อุปชฺฌาเยน วิโนเทตพฺพํ วิโนทาเปตพฺพํ, ธมฺมกถา วาสฺส กาตพฺพา. สเจ สทฺธิ\-วิหาริกสฺส ทิฏฺฐิคตํ อุปฺปนฺนํ โหติ, อุปชฺฌาเยน วิเวเจตพฺพํ วิเวจาเปตพฺพํ ธมฺมกถา วาสฺส กาตพฺพา. สเจ สทฺธิวิหาริโก ครุธมฺมํ อชฺฌาปนฺโน โหติ ปริวาสารโห, อุปชฺฌาเยน อุสฺสุกฺกํ กาตพฺพํ “กินฺติ นุ โข สํโฆ สทฺธิ\-วิหาริกสฺส ปริวาสํ ทเทยฺยาติ. สเจ สทฺธิวิหาริโก มูลาย\-ปฏิกสฺสนา\-รโห โหติ, อุปชฺฌาเยน อุสฺสุกฺกํ กาตพฺพํ “กินฺติ นุ โข สํโฆ สทฺธิ\-วิหาริกํ มูลาย ปฏิกสฺเสยฺยา”ติ. สเจ สทฺธิวิหาริโก มานตฺตารโห โหติ, อุปชฺฌาเยน อุสฺสุกฺกํ กาตพฺพํ “กินฺติ นุ โข \csromanfolio{401}สํโฆ สทฺธิ\-วิหาริกสฺส มานตฺตํ ทเทยฺยา”ติ. สเจ สทฺธิวิหาริโก อพฺภานารโห โหติ, อุปชฺฌาเยน อุสฺสุกฺกํ กาตพฺพํ “กินฺติ นุ โข สํโฆ สทฺธิ\-วิหาริกํ อพฺเภยฺยา”ติ. สเจ สํโฆ สทฺธิ\-วิหาริกสฺส กมฺมํ กตฺตุกาโม โหติ ตชฺชนียํ วา นิยสฺสํ วา ปพฺพาชนียํ วา ปฏิสารณียํ วา อุกฺเขปนียํ วา, อุปชฺฌาเยน อุสฺสุกฺกํ กาตพฺพํ “กินฺติ นุ โข สํโฆ สทฺธิ\-วิหาริกสฺส กมฺมํ น กเรยฺย, ลหุกาย วา ปริณาเมยฺยา”ติ. กตํ วา ปนสฺส โหติ สํเฆน กมฺมํ ตชฺชนียํ วา นิยสฺสํ วา ปพฺพาชนียํ วา ปฏิสารณียํ วา อุกฺเขปนียํ วา, อุปชฺฌาเยน อุสฺสุกฺกํ กาตพฺพํ “กินฺติ นุ โข สทฺธิวิหาริโก สมฺมา วตฺเตยฺย, โลมํ ปาเตยฺย, เนตฺถารํ วตฺเตยฺย, สํโฆ ตํ กมฺมํ ปฏิปฺปสฺสมฺเภยา”ติ.}

\prosewithcloser{สเจ สทฺธิ\-วิหาริกสฺส จีวรํ โธวิตพฺพํ โหติ, อุปชฺฌาเยน อาจิกฺขิตพฺพํ “เอวํ โธเวยฺยาสี”ติ, อุสฺสุกฺกํ วา กาตพฺพํ “กินฺติ นุ โข สทฺธิ\-วิหาริกสฺส จีวรํ โธวิเยถา”ติ. สเจ สทฺธิ\-วิหาริกสฺส จีวรํ กาตพฺพํ โหติ, อุปชฺฌาเยน อาจิกฺขิตพฺพํ “เอวํ กเรยฺยาสี”ติ, อุสฺสุกฺกํ วา กาตพฺพํ “กินฺติ นุ โข สทฺธิ\-วิหาริกสฺส จีวรํ กริเยถา”ติ. สเจ สทฺธิ\-วิหาริกสฺส รชนํ ปจิตพฺพํ โหติ, อุปชฺฌาเยน อาจิกฺขิตพฺพํ “เอวํ ปเจยฺยาสี”ติ, อุสฺสุกฺกํ วา กาตพฺพํ “กินฺติ นุ โข สทฺธิ\-วิหาริกสฺส รชนํ ปจิเยถา”ติ. สเจ สทฺธิ\-วิหาริกสฺส จีวรํ รชิตพฺพํ โหติ, อุปชฺฌาเยน อาจิกฺขิตพฺพํ “เอวํ รเชยฺยาสี”ติ, อุสฺสุกฺกํ วา กาตพฺพํ “กินฺติ นุ โข สทฺธิ\-วิหาริกสฺส จีวรํ รชิเยถา”ติ. จีวรํ รชนฺเตน สาธุกํ สมฺปริวตฺตกํ สมฺปริวตฺตกํ รชิตพฺพํ, น จ อจฺฉินฺเน เถเว ปกฺกมิตพฺพํ, สเจ สทฺธิวิหาริโก คิลาโน โหติ, ยาวชีวํ อุปฏฺฐาตพฺโพ, วุฏฺฐานมสฺส อาคเมตพฺพํ. อิทํ โข ภิกฺขเว อุปชฺฌายานํ สทฺธิ\-วิหาริเกสุ วตฺตํ, ยถา อุปชฺฌาเยหิ สทฺธิ\-วิหาริเกสุ สมฺมา วตฺติตพฺพนฺติ.}{ทุติย\-ภาณวาโร นิฏฺฐิโต.}
\csromanmatikaanchor{mtk.189}
\csromantocmarkref{section}{13. อาจริยวตฺตกถา}{mtk.189}
\bookmark[dest={mtk.189},level=1]{13. อาจริยวตฺตกถา}
\csromanheader{1}{13. อาจริยวตฺต\-กถา}
\proseitem{379}{เตน โข ปน สมเยน อนฺเตวาสิกา อาจริเยสุ น สมฺมา วตฺตนฺติ. เย เต ภิกฺขู อปฺปิจฺฉา ฯเปฯ เต อุชฺฌายนฺติ ขิยฺยนฺติ วิปาเจนฺติ “กถํ หิ นาม อนฺเตวาสิกา อาจริเยสุ น สมฺมา วตฺติสฺสนฺตี”ติ. \csromanfolio{402}อถ โข เต ภิกฺขู ภควโต เอตมตฺถํ อาโรเจสุํ ฯเปฯ. สจฺจํ กิร ภิกฺขเว อนฺเตวาสิกา อาจริเยสุ น สมฺมา วตฺตนฺตีติ. สจฺจํ ภควาติ ฯเปฯ วิครหิตฺวา ฯเปฯ ธมฺมิํ กถํ กตฺวา ภิกฺขู อามนฺเตสิ\csromandash{}}

\proseitem{380}{เตน หิ ภิกฺขเว อนฺเตวาสิกานํ อาจริเยสุ วตฺตํ ปญฺญเปสฺสามิ, ยถา อนฺเตวาสิเกหิ อาจริเยสุ สมฺมา วตฺติตพฺพํ.\csromansymbolfootnote{*}{วิ 3. 77 ปิฏฺเฐปิ.}อนฺเตวาสิเกน ภิกฺขเว อาจริยมฺหิ สมฺมา วตฺติตพฺพํ, ตตฺรายํ สมฺมาวตฺตนา\csromandash{}}

\prose{กาลสฺเสว อุฏฺฐาย อุปาหนา โอมุญฺจิตฺวา เอกํสํ อุตฺตราสงฺคํ กริตฺวา ทนฺตกฏฺฐํ ทาตพฺพํ, มุโขทกํ ทาตพฺพํ, อาสนํ ปญฺญเปตพฺพํ. สเจ ยาคุ โหติ, ภาชนํ โธวิตฺวา ยาคุ อุปนาเมตพฺพา, ยาคุํ ปีตสฺส อุทกํ ทตฺวา ภาชนํ ปฏิคฺคเหตฺวา นีจํ กตฺวา สาธุกํ อปฺปฏิฆํสนฺเตน โธวิตฺวา ปฏิสาเมตพฺพํ, อาจริยมฺหิ วุฏฺฐิเต อาสนํ อุทฺธริตพฺพํ. สเจ โส เทโส อุกฺลาโป โหติ, โส เทโส สมฺมชฺชิตพฺโพ.}

\prose{สเจ อาจริโย คามํ ปวิสิตุกาโม โหติ, นิวาสนํ ทาตพฺพํ, ปฏินิวาสนํ ปฏิคฺคเหตพฺพํ, กายพนฺธนํ ทาตพฺพํ, สคุณํ กตฺวา สํฆาฏิโย ทาตพฺพา, โธวิตฺวา ปตฺโต โสทโก ทาตพฺโพ. สเจ อาจริโย ปจฺฉาสมณํ อากงฺขติ, ติมณฺฑลํ ปฏิจฺฉาเทนฺเตน ปริมณฺฑลํ นิวาเสตฺวา กายพนฺธนํ พนฺธิตฺวา สคุณํ กตฺวา สํฆาฏิโย ปารุปิตฺวา คณฺฐิกํ ปฏิมุญฺจิตฺวา โธวิตฺวา ปตฺตํ คเหตฺวา อาจริยสฺส ปจฺฉาสมเณน โหตพฺพํ, นาติทูเร คนฺตพฺพํ, นาจฺจาสนฺเน คนฺตพฺพํ, ปตฺต\-ปริยาปนฺนํ ปฏิคฺคเหตพฺพํ, น อาจริยสฺส ภณมานสฺส อนฺตรนฺตรา กถา โอปาเตตพฺพา, อาจริโย อาปตฺติสามนฺตา ภณมาโน นิวาเรตพฺโพ.}

\prose{นิวตฺตนฺเตน ปฐมตรํ อาคนฺตฺวา อาสนํ ปญฺญเปตพฺพํ, ปาโททกํ ปาทปีฐํ ปาทกถลิกํ อุปนิกฺขิปิตพฺพํ, ปจฺจุคฺคนฺตฺวา ปตฺตจีวรํ ปฏิคฺคเหตพฺพํ, ปฏินิวาสนํ ทาตพฺพํ, นิวาสนํ ปฏิคฺคเหตพฺพํ. สเจ จีวรํ สินฺนํ โหติ, มุหุตฺตํ อุเณฺห โอตาเปตพฺพํ, น จ อุเณฺห จีวรํ นิทหิตพฺพํ, จีวรํ สงฺฆริตพฺพํ, จีวรํ สงฺฆรนฺเตน จตุรงฺคุลํ กณฺณํ อุสฺสาเรตฺวา จีวรํ สงฺฆริตพฺพํ “มา มชฺเฌ ภงฺโค อโหสี”ติ, โอโภเค กายพนฺธนํ กาตพฺพํ.}

\prose{สเจ ปิณฺฑปาโต โหติ, อาจริโย จ ภุญฺชิตุกาโม โหติ, อุทกํ ทตฺวา ปิณฺฑปาโต อุปนาเมตพฺโพ, อาจริโย ปานีเยน \csromanfolio{403}ปุจฺฉิตพฺโพ, ภุตฺตาวิสฺส อุทกํ ทตฺวา ปตฺตํ ปฏิคฺคเหตฺวา นีจํ กตฺวา สาธุกํ อปฺปฏิฆํสนฺเตน โธวิตฺวา โวทกํ กตฺวา มุหุตฺตํ อุเณฺห โอตาเปตพฺโพ, น จ อุเณฺห ปตฺโต นิทหิตพฺโพ. ปตฺตจีวรํ นิกฺขิปิตพฺพํ, ปตฺตํ นิกฺขิปนฺเตน เอเกน หตฺเถน ปตฺตํ คเหตฺวา เอเกน หตฺเถน เหฏฺฐามญฺจํ วา เหฏฺฐาปีฐํ วา ปรามสิตฺวา ปตฺโต นิกฺขิปิตพฺโพ, น จ อนนฺตรหิตาย ภูมิยา ปตฺโต นิกฺขิปิตพฺโพ. จีวรํ นิกฺขิปนฺเตน เอเกน หตฺเถน จีวรํ คเหตฺวา เอเกน หตฺเถน จีวรวํสํ วา จีวรรชฺชุํ วา ปมชิตฺวา ปารโต อนฺตํ โอรโต โภคํ กตฺวา จีวรํ นิกฺขิปิตพฺพํ. อาจริยมฺหิ วุฏฺฐิเต อาสนํ อุทฺธริตพฺพํ, ปาโททกํ ปาทปีฐํ ปาทกถลิกํ ปฏิสาเมตพฺพํ. สเจ โส เทโส อุกฺลาโป โหติ, โส เทโส สมฺมชฺชิตพฺโพ.}

\prose{สเจ อาจริโย นหายิตุกาโม โหติ, นหานํ ปฏิยาเทตพฺพํ. สเจ สีเตน อตฺโถ โหติ, สีตํ ปฏิยาเทตพฺพํ. สเจ อุเณฺหน อตฺโถ โหติ, อุณฺหํ ปฏิยาเทตพฺพํ.}

\prose{สเจ อาจริโย ชนฺตาฆรํ ปวิสิตุกาโม โหติ, จุณฺณํ สนฺเนตพฺพํ, มตฺติกา เตเมตพฺพา, ชนฺตาฆรปีฐํ อาทาย อาจริยสฺส ปิฏฺฐิโต ปิฏฺฐิโต คนฺตฺวา ชนฺตาฆรปีฐํ ทตฺวา จีวรํ ปฏิคฺคเหตฺวา เอกมนฺตํ นิกฺขิปิตพฺพํ, จุณฺณํ ทาตพฺพํ, มตฺติกา ทาตพฺพา. สเจ อุสฺสหติ, ชนฺตาฆรํ ปวิสิตพฺพํ, ชนฺตาฆรํ ปวิสนฺเตน มตฺติกาย มุขํ มกฺเขตฺวา ปุรโต จ ปจฺฉโต จ ปฏิจฺฉาเทตฺวา ชนฺตาฆรํ ปวิสิตพฺพํ, น เถเร ภิกฺขู อนุปขชฺช นิสีทิตพฺพํ, น นวา ภิกฺขู อาสเนน ปฏิพาหิตพฺพา, ชนฺตาฆเร อาจริยสฺส ปริกมฺมํ กาตพฺพํ, ชนฺตาฆรา นิกฺขมนฺเตน ชนฺตาฆรปีฐํ อาทาย ปุรโต จ ปจฺฉโต จ ปฏิจฺฉาเทตฺวา ชนฺตาฆรา นิกฺขมิตพฺพํ.}

\prose{อุทเกปิ อาจริยสฺส ปริกมฺมํ กาตพฺพํ, นหาเตน ปฐมตรํ อุตฺตริตฺวา อตฺตโน คตฺตํ โวทกํ กตฺวา นิวาเสตฺวา อาจริยสฺส คตฺตโต อุทกํ ปมชฺชิตพฺพํ, นิวาสนํ ทาตพฺพํ, สํฆาฏิ ทาตพฺพา, ชนฺตาฆรปีฐํ อาทาย ปฐมตรํ อาคนฺตฺวา อาสนํ ปญฺญเปตพฺพํ, ปาโททกํ ปาทปีฐํ ปาทกถลิกํ อุปนิกฺขิปิตพฺพํ, อาจริโย ปานีเยน ปุจฺฉิตพฺโพ. สเจ อุทฺทิสาเปตุกาโม โหติ อุทฺทิสิตพฺโพ, สเจ ปริปุจฺฉิตุ\-กาโม โหติ ปริปุจฺฉิตพฺโพ.}

\prose{\csromanfolio{404}ยสฺมิํ วิหาเร อาจริโย วิหรติ, สเจ โส วิหาโร อุกฺลาโป โหติ, สเจ อุสฺสหติ, โสเธตพฺโพ. วิหารํ โสเธนฺเตน ปฐมํ ปตฺตจีวรํ นีหริตฺวา เอกมนฺตํ นิกฺขิปิตพฺพํ, นิสีทน\-ปจฺจตฺถรณํ นีหริตฺวา เอกมนฺตํ นิกฺขิปิตพฺพํ, ภิสิพิพฺโพหนํ นีหริตฺวา เอกมนฺตํ นิกฺขิปิตพฺพํ, มญฺโจ นีจํ กตฺวา สาธุกํ อปฺปฏิฆํสนฺเตน อสํฆฏฺเฏนฺเตน กวาฏปิฏฺฐํ นีหริตฺวา เอกมนฺตํ นิกฺขิปิตพฺโพ, ปีฐํ นีจํ กตฺวา สาธุกํ อปฺปฏิฆํสนฺเตน อสํฆฏฺเฏนฺเตน กวาฏปิฏฺฐํ นีหริตฺวา เอกมนฺตํ นิกฺขิปิตพฺพํ, มญฺจ\-ปฏิปาทกา นีหริตฺวา เอกมนฺตํ นิกฺขิปิตพฺพา, เขฬมลฺลโก นีหริตฺวา เอกมนฺตํ นิกฺขิปิตพฺโพ, อปสฺเสนผลกํ นีหริตฺวา เอกมนฺตํ นิกฺขิปิตพฺพํ, ภูมตฺถรณํ ยถาปญฺญตฺตํ สลฺลกฺเขตฺวา นีหริตฺวา เอกมนฺตํ นิกฺขิปิตพฺพํ. สเจ วิหาเร สนฺตานกํ โหติ, อุลฺโลกา ปฐมํ โอหาเรตพฺพํ, อาโลกสนฺธิ\-กณฺณภาคา ปมชฺชิตพฺพา. สเจ เครุก\-ปริกมฺม\-กตา ภิตฺติ กณฺณกิตา โหติ, โจฬกํ เตเมตฺวา ปีเฬตฺวา ปมชฺชิตพฺพา. สเจ กาฬวณฺณกตา ภูมิ กณฺณกิตา โหติ, โจฬกํ เตเมตฺวา ปีเฬตฺวา ปมชฺชิตพฺพา. สเจ อกตา โหติ ภูมิ, อุทเกน ปริปฺโผสิตฺวา ปริปฺโผสิตฺวา สมฺมชฺชิตพฺพา “มา วิหาโร รเชน อุหญฺญี”ติ. สงฺการํ วิจินิตฺวา เอกมนฺตํ ฉฑฺเฑตพฺพํ.}

\prose{ภูมตฺถรณํ โอตาเปตฺวา โสเธตฺวา ปปฺโผเฏตฺวา อติหริตฺวา ยถาปญฺญตฺตํ ปญฺญเปตพฺพํ, มญฺจ\-ปฏิปาทกา โอตาเปตฺวา ปมชฺชิตฺวา อติหริตฺวา ยถาฐาเน ฐเปตพฺพา, มญฺโจ โอตาเปตฺวา ปปฺโผเฏตฺวา นีจํ กตฺวา สาธุกํ อปฺปฏิฆํสนฺเตน อสํฆฏฺเฏนฺเตน กวาฏปิฏฺฐํ อติหริตฺวา ยถาปญฺญตฺตํ ปญฺญเปตพฺโพ, ปีฐํ โอตาเปตฺวา โสเธตฺวา ปปฺโผเฏตฺวา นีจํ กตฺวา สาธุกํ อปฺปฏิฆํสนฺเตน อสํฆฏฺเฏนฺเตน กวาฏปิฏฺฐํ อติหริตฺวา ยถาปญฺญตฺตํ ปญฺญเปตพฺพํ, ภิสิพิพฺโพหนํ โอตาเปตฺวา โสเธตฺวา ปปฺโผเฏตฺวา อติหริตฺวา ยถาปญฺญตฺตํ ปญฺญเปตพฺพํ, นิสีทน\-ปจฺจตฺถรณํ โอตาเปตฺวา โสเธตฺวา ปปฺโผเฏตฺวา อติหริตฺวา ยถาปญฺญตฺตํ ปญฺญเปตพฺพํ, เขฬมลฺลโก โอตาเปตฺวา ปมชฺชิตฺวา อติหริตฺวา ยถาฐาเน ฐเปตพฺโพ, อปสฺเสนผลกํ โอตาเปตฺวา ปมชฺชิตฺวา อติหริตฺวา ยถาฐาเน ฐเปตพฺพํ. ปตฺตจีวรํ นิกฺขิปิตพฺพํ, ปตฺตํ นิกฺขิปนฺเตน เอเกน หตฺเถน ปตฺตํ คเหตฺวา เอเกน หตฺเถน เหฏฺฐามญฺจํ วา เหฏฺฐาปีฐํ วา ปรามสิตฺวา ปตฺโต นิกฺขิปิตพฺโพ, น จ อนนฺตรหิตาย \csromanfolio{405}ภูมิยา ปตฺโต นิกฺขิปิตพฺโพ. จีวรํ นิกฺขิปนฺเตน เอเกน หตฺเถน จีวรํ คเหตฺวา เอเกน หตฺเถน จีวรวํสํ วา จีวรรชฺชุํ วา ปมชฺชิตฺวา ปารโต อนฺตํ โอรโต โภคํ กตฺวา จีวรํ นิกฺขิปิตพฺพํ.}

\prose{สเจ ปุรตฺถิมา สรชา วาตา วายนฺติ, ปุรตฺถิมา วาตปานา ถเกตพฺพา. สเจ ปจฺฉิมา สรชา วาตา วายนฺติ, ปจฺฉิมา วาตปานา ถเกตพฺพา. สเจ อุตฺตรา สรชา วาตา วายนฺติ, อุตฺตรา วาตปานา ถเกตพฺพา. สเจ ทกฺขิณา สรชา วาตา วายนฺติ, ทกฺขิณา วาตปานา ถเกตพฺพา. สเจ สีตกาโล โหติ, ทิวา วาตปานา วิวริตพฺพา, รตฺติํ ถเกตพฺพา. สเจ อุณฺหกาโล โหติ, ทิวา วาตปานา ถเกตพฺพา, รตฺติํ วิวริตพฺพา.}

\prose{สเจ ปริเวณํ อุกฺลาปํ โหติ, ปริเวณํ สมฺมชฺชิตพฺพํ. สเจ โกฏฺฐโก อุกฺลาโป โหติ, โกฏฺฐโก สมฺมชฺชิตพฺโพ. สเจ อุปฏฺฐานสาลา อุกฺลาปา โหติ, อุปฏฺฐานสาลา สมฺมชฺชิตพฺพา. สเจ อคฺคิสาลา อุกฺลาปา โหติ, อคฺคิสาลา สมฺมชฺชิตพฺพา. สเจ วจฺจกุฏิ อุกฺลาปา โหติ, วจฺจกุฏิ สมฺมชฺชิตพฺพา. สเจ ปานียํ น โหติ, ปานียํ อุปฏฺฐาเป\-ตพฺพํ, สเจ ปริโภชนียํ น โหติ, ปริโภชนียํ อุปฏฺฐาเป\-ตพฺพํ. สเจ อาจมนกุมฺภิยา อุทกํ น โหติ, อาจมนกุมฺภิยา อุทกํ อาสิญฺจิตพฺพํ.}

\prose{สเจ อาจริยสฺส อนภิรติ อุปฺปนฺนา โหติ, อนฺเตวาสิเกน วูปกาเสตพฺโพ วูปกาสาเป\-ตพฺโพ ธมฺมกถา วาสฺส กาตพฺพา. สเจ อาจริยสฺส กุกฺกุจฺจํ อุปฺปนฺนํ โหติ, อนฺเตวาสิเกน วิโนเทตพฺพํ วิโนทาเปตพฺพํ ธมฺมกถา วาสฺส กาตพฺพา. สเจ อาจริยสฺส ทิฏฺฐิคตํ อุปฺปนฺนํ โหติ, อนฺเตวาสิเกน วิเวเจตพฺพํ วิเวจาเปตพฺพํ ธมฺมกถา วาสฺส กาตพฺพา. สเจ อาจริโย ครุธมฺมํ อชฺฌาปนฺโน โหติ ปริวาสารโห, อนฺเตวาสิเกน อุสฺสุกฺกํ กาตพฺพํ “กินฺติ นุ โข สํโฆ อาจริยสฺส ปริวาสํ ทเทยฺยา”ติ. สเจ อาจริโย มูลาย\-ปฏิกสฺสนา\-รโห โหติ, อนฺเตวาสิเกน อุสฺสุกฺกํ กาตพฺพํ “กินฺติ นุ โข สํโฆ อาจริยํ มูลาย ปฏิกสฺเสยฺยา”ติ. สเจ อาจริโย มานตฺตารโห โหติ, อนฺเตวาสิเกน อุสฺสุกฺกํ กาตพฺพํ “กินฺติ นุ โข สํโฆ อาจริยสฺส มานตฺตํ ทเทยฺยา”ติ. สเจ อาจริโย อพฺภานารโห โหติ, อนฺเตวาสิเกน อุสฺสุกฺกํ กาตพฺพํ “กินฺติ นุ โข สํโฆ อาจริยํ อพฺเภยฺยา”ติ. สเจ สํโฆ \csromanfolio{406}อาจริยสฺส กมฺมํ กตฺตุกาโม โหติ ตชฺชนียํ วา นิยสฺสํ วา ปพฺพาชนียํ วา ปฏิสารณียํ วา อุกฺเขปนียํ วา, อนฺเตวาสิเกน อุสฺสุกฺกํ กาตพฺพํ “กินฺติ นุ โข สํโฆ อาจริยสฺส กมฺมํ น กเรยฺย ลหุกาย วา ปริณาเมยฺยา”ติ. กตํ วา ปนสฺส โหติ สํเฆน กมฺมํ ตชฺชนียํ วา นิยสฺสํ วา ปพฺพาชนียํ วา ปฏิสารณียํ อุกฺเขปนียํ วา, อนฺเตวาสิเกน อุสฺสุกฺกํ กาตพฺพํ “กินฺติ นุ โข อาจริโย สมฺมา วตฺเตยฺย, โลมํ ปาเตยฺย, เนตฺถารํ วตฺเตยฺย, สํโฆ ตํ กมฺมํ ปฏิปฺปสฺสมฺเภยฺยา”ติ.}

\prose{สเจ อาจริยสฺส จีวรํ โธวิตพฺพํ โหติ, อนฺเตวาสิเกน โธวิตพฺพํ, อุสฺสุกฺกํ วา กาตพฺพํ “กินฺติ นุ โข อาจริยสฺส จีวรํ โธวิเยถา”ติ. สเจ อาจริยสฺส จีวรํ กาตพฺพํ โหติ, อนฺเตวาสิเกน กาตพฺพํ, อุสฺสุกฺกํ วา กาตพฺพํ “กินฺติ นุ โข อาจริยสฺส จีวรํ กริเยถา”ติ. สเจ อาจริยสฺส รชนํ ปจิตพฺพํ โหติ, อนฺเตวาสิเกน ปจิตพฺพํ, อุสฺสุกฺกํ วา กาตพฺพํ “กินฺติ นุ โข อาจริยสฺส รชนํ ปจิเยถา”ติ. สเจ อาจริยสฺส จีวรํ รชิตพฺพํ โหติ, อนฺเตวาสิเกน รชิตพฺพํ, อุสฺสุกฺกํ วา กาตพฺพํ “กินฺติ นุ โข อาจริยสฺส จีวรํ รชิเยถา”ติ. จีวรํ รชนฺเตน สาธุกํ สมฺปริวตฺตกํ สมฺปริวตฺตกํ รชิตพฺพํ, น จ อจฺฉินฺเน เถว ปกฺกมิตพฺพํ.}

\prose{น อาจริยํ อนาปุจฺฉา เอกจฺจสฺส ปตฺโต ทาตพฺโพ, น เอกจฺจสฺส ปตฺโต ปฏิคฺคเหตพฺโพ, น เอกจฺจสฺส จีวรํ ทาตพฺพํ, น เอกจฺจสฺส จีวรํ ปฏิคฺคเหตพฺพํ, น เอกจฺจสฺส ปริกฺขาโร ทาตพฺโพ, น เอกจฺจสฺส ปริกฺขาโร ปฏิคฺคเหตพฺโพ, น เอกจฺจสฺส เกสา เฉทิตพฺพา, น เอกจฺเจน เกสา เฉทาเปตพฺพา, น เอกจฺจสฺส ปริกมฺมํ กาตพฺพํ, น เอกจฺเจน ปริกมฺมํ การาเปตพฺพํ, น เอกจฺจสฺส เวยฺยาวจฺโจ กาตพฺโพ, น เอกจฺเจน เวยฺยาวจฺโจ การาเปตพฺโพ, น เอกจฺจสฺส ปจฺฉาสมเณน โหตพฺพํ, น เอกจฺโจ ปจฺฉาสมโณ อาทาตพฺโพ, น เอกจฺจสฺส ปิณฺฑปาโต นีหริตพฺโพ, น เอกจฺเจน ปิณฺฑปาโต นีหราเปตพฺโพ, น อาจริยํ อนาปุจฺฉา คาโม ปวิสิตพฺโพ, น สุสานํ คนฺตพฺพํ, น ทิสา ปกฺกมิตพฺพา, สเจ อาจริโย คิลาโน โหติ, ยาวชีวํ \csromanfolio{407}อุปฏฺฐาตพฺโพ, วุฏฺฐานมสฺส อาคเมตพฺพํ. อิทํ โข ภิกฺขเว อนฺเตวาสิกานํ อาจริเยสุ วตฺตํ, ยถา อนฺเตวาสิเกหิ อาจริเยสุ สมฺมา วตฺติตพฺพนฺติ.}

\csromanmatikaanchor{mtk.190}
\csromantocmarkref{section}{14. อนฺเตวาสิกวตฺตกถา}{mtk.190}
\bookmark[dest={mtk.190},level=1]{14. อนฺเตวาสิกวตฺตกถา}
\csromanheader{1}{14. อนฺเตวาสิกวตฺต\-กถา}
\proseitem{381}{เตน โข ปน สมเยน อาจริยา อนฺเตวาสิเกสุ น สมฺมา วตฺตนฺติ. เย เต ภิกฺขู อปฺปิจฺฉา ฯเปฯ เต อุชฺฌายนฺติ ขิยฺยนฺติ วิปาเจนฺติ “กถํ หิ นาม อาจริยา อนฺเตวาสิเกสุ น สมฺมา วตฺติสฺสนฺตี”ติ. อถ โข เต ภิกฺขู ภควโต เอตมตฺถํ อาโรเจสุํ. อถ โข ภควา เอตสฺมิํ นิทาเน เอตสฺมิํ ปกรเณ ภิกฺขุสํฆํ สนฺนิปาตาเปตฺวา ภิกฺขู ปฏิปุจฺฉิ “สจฺจํ กิร ภิกฺขเว อาจริยา อนฺเตวาสิเกสุ น สมฺมา วตฺตนฺตี”ติ. สจฺจํ ภควาติ ฯเปฯ วิครหิตฺวา ฯเปฯ ธมฺมิํ กถํ กตฺวา ภิกฺขู อามนฺเตสิ\csromandash{}}

\proseitem{382}{เตน หิ ภิกฺขเว อาจริยานํ อนฺเตวาสิเกสุ วตฺตํ ปญฺญเปสฺสามิ, ยถา อาจาริเยหิ อนฺเตวาสิเกสุ สมฺมา วตฺติตพฺพํ.\csromansymbolfootnote{*}{วิ 3. 81 ปิฏฺเฐปิ.}อาจริเยน ภิกฺขเว อนฺเตวาสิกมฺหิ สมฺมา วตฺติตพฺพํ. ตตฺรายํ สมฺมาวตฺตนา\csromandash{}}

\prose{อาจริเยน ภิกฺขเว อนฺเตวาสิโก สงฺคเหตพฺโพ อนุคฺคเหตพฺโพ อุทฺเทเสน ปริปุจฺฉาย โอวาเทน อนุสาสนิยา, สเจ อาจริยสฺส ปตฺโต โหติ อนฺเตวาสิกสฺส ปตฺโต น โหติ, อาจริเยน อนฺเตวาสิกสฺส ปตฺโต ทาตพฺโพ, อุสฺสุกฺกํ วา กาตพฺพํ “กินฺติ นุ โข อนฺเตวาสิกสฺส ปตฺโต อุปฺปชฺชิเยถา”ติ. สเจ อาจริยสฺส จีวรํ โหติ อนฺเตวาสิกสฺส จีวรํ น โหติ, อาจริเยน อนฺเตวาสิกสฺส จีวรํ ทาตพฺพํ, อุสฺสุกฺกํ วา กาตพฺพํ “กินฺติ นุ โข อนฺเตวาสิกสฺส จีวรํ อุปฺปชฺชิเยถา”ติ. สเจ อาจริยสฺส ปริกฺขาโร โหติ อนฺเตวาสิกสฺส ปริกฺขาโร น โหติ, อาจริเยน อนฺเตวาสิกสฺส ปริกฺขาโร ทาตพฺโพ, อุสฺสุกฺกํ วา กาตพฺพํ “กินฺติ นุ โข อนฺเตวาสิกสฺส ปริกฺขาโร อุปฺปชฺชิเยถา”ติ.}

\prose{สเจ อนฺเตวาสิโก คิลาโน โหติ, กลาสฺเสว อุฏฺฐาย ทนฺตกฏฺฐํ ทาตพฺพํ, มุโขทกํ ทาตพฺพํ, อาสนํ ปญฺญเปตพฺพํ, สเจ ยาคุ \csromanfolio{408}โหติ, ภาชนํ โธวิตฺวา ยาคุ อุปนาเมตพฺพา, ยาคุํ ปีตสฺส อุทกํ ทตฺวา ภาชนํ ปฏิคฺคเหตฺวา นีจํ กตฺวา สาธุกํ อปฺปฏิฆํสนฺเตน โธวิตฺวา ปฏิสาเมตพฺพํ, อนฺเตวาสิกมฺหิ วุฏฺฐิเต อาสนํ อุทฺธริตพฺพํ, สเจ โส เทโส อุกฺลาโป โหติ, โส เทโส สมฺมชฺชิตพฺโพ.}

\prose{สเจ อนฺเตวาสิโก คามํ ปวิสิตุกาโม โหติ, นิวาสนํ ทาตพฺพํ, ปฏินิวาสนํ ปฏิคฺคเหตพฺพํ, กายพนฺธนํ ทาตพฺพํ, สคุณํ กตฺวา สํฆาฏิโย ทาตพฺพา, โธวิตฺวา ปตฺโต โสทโก ทาตพฺโพ.}

\prose{“เอตฺตาวตา นิวตฺติสฺสตี”ติ อาสนํ ปญฺญเปตพฺพํ, ปาโททกํ ปาทปีฐํ ปาทกถลิกํ อุปนิกฺขิปิตพฺพํ, ปจฺจุคฺคนฺตฺวา ปตฺตจีวรํ ปฏิคฺคเหตพฺพํ, ปฏินิวาสนํ ทาตพฺพํ, นิวาสนํ ปฏิคฺคเหตพฺพํ, สเจ จีวรํ สินฺนํ โหติ, มุหุตฺตํ อุเณฺห โอตาเปตพฺพํ, น จ อุเณฺห จีวรํ นิทหิตพฺพํ, จีวรํ สงฺฆริตพฺพํ, จีวรํ สงฺฆรนฺเตน จตุรงฺคุลํ กณฺณํ อุสฺสาเรตฺวา จีวรํ สงฺฆริตพฺพํ “มา มชฺเฌ ภงฺโค อโหสี”ติ. โอโภเค กายพนฺธนํ กาตพฺพํ.}

\prose{สเจ ปิณฺฑปาโต โหติ อนฺเตวาสิโก จ ภุญฺชิตุกาโม โหติ, อุทกํ ทตฺวา ปิณฺฑปาโต อุปนาเมตพฺโพ, อนฺเตวาสิโก ปานีเยน ปุจฺฉิตพฺโพ, ภุตฺตาวิสฺส อุทกํ ทตฺวา ปตฺตํ ปฏิคฺคเหตฺวา นีจํ กตฺวา สาธุกํ อปฺปฏิฆํสนฺเตน โธวิตฺวา โวทกํ กตฺวา มุหุตฺตํ อุเณฺห โอตาเปตพฺโพ, น จ อุเณฺห ปตฺโต นิทหิตพฺโพ. ปตฺตจีวรํ นิกฺขิปิตพฺพํ, ปตฺตํ นิกฺขิปนฺเตน ฯเปฯ จีวรํ นิกฺขิปนฺเตน ฯเปฯ ปารโต อนฺตํ โอรโต โภคํ กตฺวา จีวรํ นิกฺขิปิตพฺพํ. อนฺเตวาสิกมฺหิ อุฏฺฐิเต อาสนํ อุทฺธริตพฺพํ, ปาโททกํ ปาทปีฐํ ปาทกถลิกํ ปฏิสาเมตพฺพํ, สเจ โส เทโส อุกฺลาโป โหติ, โส เทโส สมฺมชฺชิตพฺโพ.}

\prose{สเจ อนฺเตวาสิโก นหายิตุกาโม โหติ, นหานํ ปฏิยาเทตพฺพํ, สเจ สีเตน อตฺโถ โหติ, สีตํ ปฏิยาเทตพฺพํ, สเจ อุเณฺหน อตฺโถ โหติ, อุณฺหํ ปฏิยาเทตพฺพํ.}

\prose{สเจ อนฺเตวาสิโก ชนฺตาฆรํ ปวิสิตุกาโม โหติ, จุณฺณํ สนฺเนตพฺพํ, มตฺติกา เตเมตพฺพา, ชนฺตาฆรปีฐํ อาทาย\csromansharedfootnote{8f2facc40cbc7307}{อาทาย อนฺเตวาสิกสฺส ปิฏฺฐิโต ปิฏฺฐิโต (ก)} คนฺตฺวา ชนฺตาฆรปีฐํ ทตฺวา จีวรํ ปฏิคฺคเหตฺวา เอกมนฺตํ นิกฺขิปิตพฺพํ, จุณฺณํ ทาตพฺพํ, มตฺติกา ทาตพฺพา, สเจ อุสฺสหติ, ชนฺตาฆรํ ปวิสิตพฺพํ, ชนฺตาฆรํ \csromanfolio{409}ปวิสนฺเตน มตฺติกาย มุขํ มกฺเขตฺวา ปุรโต จ ปจฺฉโต จ ปฏิจฺฉาเทตฺวา ชนฺตาฆรํ ปวิสิตพฺพํ, น เถเร ภิกฺขู อนุปขชฺช นิสีทิตพฺพํ, น นวา ภิกฺขู อาสเนน ปฏิพาหิตพฺพา, ชนฺตาฆเร อนฺเตวาสิกสฺส ปริกมฺมํ กาตพฺพํ, ชนฺตาฆรา นิกฺขมนฺเตน ชนฺตาฆรปีฐํ อาทาย ปุรโต จ ปจฺฉโต จ ปฏิจฺฉาเทตฺวา ชนฺตาฆรา นิกฺขมิตพฺพํ.}

\prose{อุทเกปิ อนฺเตวาสิกสฺส ปริกมฺมํ กาตพฺพํ, นหาเตน ปฐมตรํ อุตฺตริตฺวา อตฺตโน คตฺตํ โวทกํ กตฺวา นิวาเสตฺวา อนฺเตวาสิกสฺส คตฺตโต อุทกํ ปมชฺชิตพฺพํ, นิวาสนํ ทาตพฺพํ, สํฆาฏิ ทาตพฺพา, ชนฺตาฆรปีฐํ อาทาย ปฐมตรํ อาคนฺตฺวา อาสนํ ปญฺญเปตพฺพํ, ปาโททกํ ปาทปีฐํ ปาทกถลิกํ อุปนิกฺขิปิตพฺพํ, อนฺเตวาสิโก ปานีเยน ปุจฺฉิตพฺโพ.}

\prose{ยสฺมิํ วิหาเร อนฺเตวาสิโก วิหรติ, สเจ โส วิหาโร อุกฺลาโป โหติ, สเจ อุสฺสหติ, โสเธตพฺโพ. วิหารํ โสเธนฺเตน ปฐมํ ปตฺตจีวรํ นีหริตฺวา เอกมนฺตํ นิกฺขิปิตพฺพํ ฯเปฯ อาจมนกุมฺภิยา อุทกํ น โหติ, อาจมนกุมฺภิยา อุทกํ อาสิญฺจิตพฺพํ.}

\csromanmatikaanchor{mtk.191}
\csromantocmarkref{section}{อุทฺทานคาถา}{mtk.191}
\bookmark[dest={mtk.191},level=1]{อุทฺทานคาถา}
\prose{สเจ อนฺเตวาสิกสฺส อนภิรติ อุปฺปนฺนา โหติ, อาจริเยน วูปกาเสตพฺโพ วูปกาสาเป\-ตพฺโพ ธมฺมกถา วาสฺส กาตพฺพา. สเจ อนฺเตวาสิกสฺส กุกฺกุจฺจํ อุปฺปนฺนํ โหติ, อาจริเยน วิโนเทตพฺพํ วิโนทาเปตพฺพํ ธมฺมกถา วาสฺส กาตพฺพา. สเจ อนฺเตวาสิกสฺส ทิฏฺฐิคตํ อุปฺปนฺนํ โหติ, อาจริเยน วิเวเจตพฺพํ วิเวจาเปตพฺพํ ธมฺมกถา วาสฺส กาตพฺพา. สเจ อนฺเตวาสิโก ครุธมฺมํ อชฺฌาปนฺโน โหติ ปริวาสารโห, อาจริเยน อุสฺสุกฺกํ กาตพฺพํ “กินฺติ นุ โข สํโฆ อนฺเตวาสิกสฺส ปริวาสํ ทเทยฺยา”ติ. สเจ อนฺเตวาสิโก มูลาย\-ปฏิกสฺสนา\-รโห โหติ, อาจริเยน อุสฺสุกฺกํ กาตพฺพํ “กินฺติ นุ โข สํโฆ อนฺเตวาสิกํ มูลาย ปฏิกสฺเสยฺยา”ติ. สเจ อนฺเตวาสิโก มานตฺตารโห โหติ, อาจริเยน อุสฺสุกฺกํ กาตพฺพํ “กินฺติ นุ โข สํโฆ อนฺเตวาสิกสฺส มานตฺตํ ทเทยฺยา”ติ. สเจ อนฺเตวาสิโก อพฺภานารโห โหติ, อาจริเยน อุสฺสุกฺกํ กาตพฺพํ “กินฺติ นุ โข สํโฆ อนฺเตวาสิกํ อพฺเภยฺยา”ติ. สเจ สํโฆ อนฺเตวาสิกสฺส กมฺมํ กตฺตุกาโม โหติ ตชฺชนียํ วา นิยสฺสํ วา ปพฺพาชนียํ วา ปฏิสารณียํ วา อุกฺเขปนียํ วา, อาจริเยน อุสฺสุกฺกํ กาตพฺพํ “กินฺติ นุ โข สํโฆ อนฺเตวาสิกสฺส กมฺมํ น \csromanfolio{410}กเรยฺย ลหุกาย วา ปริณาเมยฺยา”ติ. กตํ วา ปนสฺส โหติ สํเฆน กมฺมํ ตชฺชนียํ วา นิยสฺสํ วา ปพฺพาชนียํ วา ปฏิสารณียํ วา อุกฺเขปนียํ วา, อาจริเยน อุสฺสุกฺกํ กาตพฺพํ “กินฺติ นุ โข อนฺเตวาสิโก สมฺมา วตฺเตยฺย, โลมํ ปาเตยฺย, เนตฺถารํ วตฺเตยฺย, สํโฆ ตํ กมฺมํ ปฏิปฺปสฺสมฺเภยฺยา”ติ.}

\prose{สเจ อนฺเตวาสิกสฺส จีวรํ โธวิตพฺพํ โหติ, อาจริเยน อาจิกฺขิตพฺพํ “เอวํ โธเวยฺยาสี”ติ, อุสฺสุกฺกํ วา กาตพฺพํ “กินฺติ นุ โข อนฺเตวาสิกสฺส จีวรํ โธวิเยถา”ติ. สเจ อนฺเตวาสิกสฺส จีวรํ กาตพฺพํ โหติ, อาจริเยน อาจิกฺขิตพฺพํ “เอวํ กเรยฺยาสี”ติ, อุสฺสุกฺกํ วา กาตพฺพํ “กินฺติ นุ โข อนฺเตวาสิกสฺส จีวรํ กริเยถา”ติ. สเจ อนฺเตวาสิกสฺส รชนํ ปจิตพฺพํ โหติ, อาจริเยน อาจิกฺขิตพฺพํ “กฺวํ ปเจยฺยาสี”ติ, อุสฺสุกฺกํ วา กาตพฺพํ “กินฺติ นุ โข อนฺเตวาสิกสฺส รชนํ ปจิเยถา”ติ. สเจ อนฺเตวาสิกสฺส จีวรํ รชิตพฺพํ โหติ, อาจริเยน อาจิกฺขิตพฺพํ “เอวํ รเชยฺยาสี”ติ, อุสฺสุกฺกํ วา กาตพฺพํ “กินฺติ นุ โข อนฺเตวาสิกสฺส จีวรํ รชิเยถา”ติ. จีวรํ รชนฺเตน สาธุกํ สมฺปริวตฺตกํ สมฺปริวตฺตกํ รชิตพฺพํ, น จ อจฺฉินฺเน เถเว ปกฺกมิตพฺพํ. สเจ อนฺเตวาสิโก คิลาโน โหติ, ยาวชีวํ อุปฏฺฐาตพฺโพ, วุฏฺฐานมสฺส อาคเมตพฺพํ. อิทํ โข ภิกฺขเว อาจริยานํ อนฺเตวาสิเกสุ วตฺตํ, ยถา อาจริเยหิ อนฺเตวาสิเกสุ สมฺมา วตฺติตพฺพนฺติ.}

\vspace{\tipitakaparskip}
\csromancenter{วตฺต\-กฺขนฺธโก อฏฺฐโม.}
\setlength{\csromangathapagewidth}{0pt}
\setlength{\csromangathawakwidth}{0pt}
\csromangathameasuregroup{อิมมฺหิ ขนฺธเก วตฺถู เอกูนวีสติ,}{\csromangathabat{อิมมฺหิ ขนฺธเก วตฺถู เอกูนวีสติ,}{วตฺตา จุทฺทส.}}
\csromangathasetleft{อิมมฺหิ ขนฺธเก วตฺถู เอกูนวีสติ,}
\csromangathagroup{\csromangathastanza{\csromangathabat{อิมมฺหิ ขนฺธเก วตฺถู เอกูนวีสติ,}{วตฺตา จุทฺทส.}}}

\vspace{\tipitakaparskip}
\csromancenter{\textbf{ตสฺสุทฺทานํ}}
\setlength{\csromangathapagewidth}{0pt}
\setlength{\csromangathawakwidth}{0pt}
\csromangathameasuregroup{}{สอุปาหนา ฉตฺตา จ โอคุณฺฐิ สีสํ ปานียํ. \\ นาภิวาเท น ปุจฺฉนฺติ, \\ อหิ อุชฺฌนฺติ เปสลา. \\ โอมุญฺจิ ฉตฺตํ ขนฺเธ จ, \\ อตรญฺจ ปฏิกฺกมํ. \\ ปตฺตจีวรํ นิกฺขิปา, \\ ปติรูปญฺจ ปุจฺฉิตา. \\ อาสิญฺเจยฺย โธวิเตน, \\ สุกฺเขนลฺเลนุปาหนา. \\ วุฑฺโฒ นวโก ปุจฺเฉยฺย, \\ อชฺฌาวุฏฺฐญฺจ โคจรา. \\ เสกฺขา วจฺจา ปานี ปริ, \\ กตฺตรํ กติกํ ตโต. \\ กาลํ มุหุตฺตํ อุกฺลาโป, \\ ภูมตฺถรณํ นีหเร. \\ ปฏิปาโท ภิสิพิพฺโพ, \\ มญฺจปีฐญฺจ มลฺลกํ. \\ อปสฺเสนุลฺโลกกณฺณา, \\ เครุกา กาฬ อกตา. \\ สงฺการญฺจ ภูมตฺถรณํ, \\ ปฏิปาทกํ มญฺจปีฐํ. \\ ภิสิ นิสีทนมฺปิ, \\ มลฺลกํ อปสฺเสน จ. \\ ปตฺตจีวรํ ภูมิ จ, \\ ปารนฺตํ โอรโต โภคํ. \\ ปุรตฺถิมา ปจฺฉิมา จ, \\ อุตฺตรา อถ ทกฺขิณา. \\ สีตุเณฺห จ ทิวารตฺติํ, \\ ปริเวณญฺจ โกฏฺฐโก. \\ อุปฏฺฐานคฺคิ สาลา จ, \\ วตฺตํ วจฺจกุฏีสุ จ. \\ ปานี ปริโภชนิยา, \\ กุมฺภิ อาจมเนสุ จ. \\ อโนปเมน ปญฺญตฺตํ, \\ วตฺตํ อาคนฺตุเกหิเม\textsuperscript{1}. \\ เนวาสนํ น อุทกํ, \\ น ปจฺจุ น จ ปานิยํ. \\ นาภิวาเท นปญฺญเป, \\ อุชฺฌายนฺติ จ เปสลา. \\ วุฑฺฒาสนญฺจ อุทกํ, \\ ปจฺจุคฺคนฺตฺวา จ ปานิยํ. \\ อุปาหเน เอกมนฺตํ, \\ อภิวาเท จ ปญฺญเป. \\ วุตฺถํ โคจรเสกฺโข จ, \\ ฐานํ ปานียโภชนํ. \\ กตฺตรํ กติกํ กาลํ, \\ นวกสฺส นิสินฺนเก. \\ อภิวาทเย อาจิกฺเข, \\ ยถา เหฏฺฐา ตถา นเย. \\ นิทฺทิฏฺฐํ สตฺถวาเหน, \\ วตฺตํ อาวาสิเกหิเม. \\ คมิกา ทารุมตฺติ จ, \\ วิวริตฺวา น ปุจฺฉิย. \\ นสฺสนฺติ จ อคุตฺตญฺจ, \\ อุชฺฌายนฺติ จ เปสลา. \\ ปฏิสาเมตฺวา ถเกตฺวา, \\ อาปุจฺฉิตฺวาว ปกฺกเม. \\ ภิกฺขุ วา สามเณโร วา, \\ อารามิโก อุปาสโก. \\ ปาสาณเกสุ จ ปุญฺชํ, \\ ปฏิสาเมถเกยฺย จ. \\ สเจ อุสฺสหติ อุสฺสุกฺกํ, \\ อโนวสฺเส ตเถว จ. \\ สพฺโพ โอวสฺสติ คามํ, \\ อชฺโฌกาเส ตเถว จ. \\ อปฺเปวงฺคานิ เสเสยฺยุํ, \\ วตฺตํ คมิกภิกฺขุนา. \\ นานุโมทนฺติ เถเรน, \\ โอหาย จตุปญฺจหิ. \\ วจฺจิโต มุจฺฉิโต อาสิ, \\ วตฺตานุโมทเนสุเม. \\ ฉพฺพคฺคิยา ทุนฺนิวตฺถา, \\ อโถปิ จ ทุปฺปารุกา. \\ อนากปฺปา จ โวกฺกมฺม, \\ เถเร จ อนุปขชฺชเน. \\ นเว ภิกฺขู จ สํฆาฏิ, \\ อุชฺฌายนฺติ จ เปสลา. \\ ติมณฺฑลํ นิวาเสตฺวา, \\ กายสคุณคณฺฐิกา. \\ น โวกฺกมฺม ปฏิจฺฉนฺนํ, \\ สุสํวุโตกฺขิตฺตจกฺขุ. \\ อุกฺขิตฺโตชฺชคฺฆิกาสทฺโท, \\ ตโย เจว ปจาลนา. \\ ขมฺโภคุณฺฐิอุกฺกุฏิกา, \\ ปฏิจฺฉนฺนํ สุสํวุโต. \\ โอกฺขิตฺตุกฺขิตฺตอุชฺชคฺฆิ, \\ อปฺปสทฺโท ตโย จลา. \\ ขมฺโภคุณฺฐิปลฺลตฺถิ จ, \\ อนุปขชฺช นาสเน. \\ โอตฺถริตฺวาน อุทเก, \\ นีจํ กตฺวาน สิญฺจิยา. \\ ปฏิ สามนฺตา สํฆาฏิ, \\ โอทเน จ ปฏิคฺคเห. \\ สูปํ อุตฺตริภงฺเคน, \\ สพฺเพสํ สมติตฺถิ จ. \\ สกฺกจฺจํ, \\ ปตฺตสญฺญี จ, สปทานญฺจ สูปกํ. \\ น ถูปโต ปฏิจฺฉาเท, \\ วิญฺญตฺตุชฺฌานสญฺญินา. \\ มหนฺตมณฺฑลทฺวารํ, \\ สพฺพหตฺโถ น พฺยาหเร. \\ อุกฺเขโป เฉทนาคณฺฑ, \\ ธุนํ สิตฺถาวการกํ. \\ ชิวฺหานิจฺฉารกญฺเจว, \\ จปุจปุ สุรุสุรุ. \\ หตฺถปตฺโตฏฺฐนิลฺเลหํ, \\ สามิเสน ปฏิคฺคเห. \\ ยาว น สพฺเพ อุทเก, \\ นีจํ กตฺวาน สิญฺจิยํ. \\ ปฏิ สามนฺตา สํฆาฏิ, \\ นีจํ กตฺวา ฉมาย จ. \\ สสิตฺถกํ นิวตฺตนฺเต, \\ สุปฺปฏิจฺฉนฺนมุกฺกุฏิ. \\ ธมฺมราเชน ปญฺญตฺตํ, \\ อิทํ ภตฺตคฺควตฺตนํ. \\ ทุนฺนิวตฺถา อนากปฺปา, \\ อสลฺลกฺเขตฺวา จ สหสา. \\ ทูเร อจฺจ จิรํ ลหุํ, \\ ตเถว ปิณฺฑจาริโก. \\ ปฏิจฺฉนฺโนว คจฺเฉยฺย, \\ สุสํวุโตกฺขิตฺตจกฺขุ. \\ อุกฺขิตฺโตชฺชคฺฆิกาสทฺโท, \\ ตโย เจว ปจาลนา. \\ ขมฺโภคุณฺฐิอุกฺกุฏิกา, \\ สลฺลกฺเขตฺวา จ สหสา. \\ ทูเร อจฺจ จิรํ ลหุํ, \\ อาสนกํ กฏจฺฉุกา. \\ ภาชนํ วา ฐเปติ จ, \\ อุจฺจาเรตฺวา ปณาเมตฺวา. \\ ปฏิคฺคเห น อุลฺโลเก, \\ สูเปสุปิ ตเถว ตํ. \\ ภิกฺขุ สํฆาฏิยา ฉาเท, \\ ปฏิจฺฉนฺเนว คจฺฉิยํ. \\ สํวุโตกฺขิตฺตจกฺขุ จ, \\ อุกฺขิตฺโตชฺชคฺฆิกาย จ. \\ อปฺปสทฺโท ตโย จาลา, \\ ขมฺโภคุณฺฐิกอุกฺกุฏิ. \\ ปฐมาสนวกฺการ, \\ ปานิยํ ปริโภชนี. \\ ปจฺฉากงฺขติ ภุญฺเชยฺย, \\ โอปิลาเปยฺย อุทฺธเร. \\ ปฏิสาเมยฺย สมฺมชฺเช, \\ ริตฺตํ ตุจฺฉํ อุปฏฺฐเป. \\ หตฺถวิกาเร ภินฺเทยฺย, \\ วตฺติทํ ปิณฺฑจาริเก. \\ ปานี ปริ อคฺคิรณิ, \\ นกฺขตฺตทิสโจรา จ.}
\csromangathameasurewak{สอุปาหนา ฉตฺตา จ โอคุณฺฐิ สีสํ ปานียํ. \\ นาภิวาเท น ปุจฺฉนฺติ, \\ อหิ อุชฺฌนฺติ เปสลา. \\ โอมุญฺจิ ฉตฺตํ ขนฺเธ จ, \\ อตรญฺจ ปฏิกฺกมํ. \\ ปตฺตจีวรํ นิกฺขิปา, \\ ปติรูปญฺจ ปุจฺฉิตา. \\ อาสิญฺเจยฺย โธวิเตน, \\ สุกฺเขนลฺเลนุปาหนา. \\ วุฑฺโฒ นวโก ปุจฺเฉยฺย, \\ อชฺฌาวุฏฺฐญฺจ โคจรา. \\ เสกฺขา วจฺจา ปานี ปริ, \\ กตฺตรํ กติกํ ตโต. \\ กาลํ มุหุตฺตํ อุกฺลาโป, \\ ภูมตฺถรณํ นีหเร. \\ ปฏิปาโท ภิสิพิพฺโพ, \\ มญฺจปีฐญฺจ มลฺลกํ. \\ อปสฺเสนุลฺโลกกณฺณา, \\ เครุกา กาฬ อกตา. \\ สงฺการญฺจ ภูมตฺถรณํ, \\ ปฏิปาทกํ มญฺจปีฐํ. \\ ภิสิ นิสีทนมฺปิ, \\ มลฺลกํ อปสฺเสน จ. \\ ปตฺตจีวรํ ภูมิ จ, \\ ปารนฺตํ โอรโต โภคํ. \\ ปุรตฺถิมา ปจฺฉิมา จ, \\ อุตฺตรา อถ ทกฺขิณา. \\ สีตุเณฺห จ ทิวารตฺติํ, \\ ปริเวณญฺจ โกฏฺฐโก. \\ อุปฏฺฐานคฺคิ สาลา จ, \\ วตฺตํ วจฺจกุฏีสุ จ. \\ ปานี ปริโภชนิยา, \\ กุมฺภิ อาจมเนสุ จ. \\ อโนปเมน ปญฺญตฺตํ, \\ วตฺตํ อาคนฺตุเกหิเม\textsuperscript{1}. \\ เนวาสนํ น อุทกํ, \\ น ปจฺจุ น จ ปานิยํ. \\ นาภิวาเท นปญฺญเป, \\ อุชฺฌายนฺติ จ เปสลา. \\ วุฑฺฒาสนญฺจ อุทกํ, \\ ปจฺจุคฺคนฺตฺวา จ ปานิยํ. \\ อุปาหเน เอกมนฺตํ, \\ อภิวาเท จ ปญฺญเป. \\ วุตฺถํ โคจรเสกฺโข จ, \\ ฐานํ ปานียโภชนํ. \\ กตฺตรํ กติกํ กาลํ, \\ นวกสฺส นิสินฺนเก. \\ อภิวาทเย อาจิกฺเข, \\ ยถา เหฏฺฐา ตถา นเย. \\ นิทฺทิฏฺฐํ สตฺถวาเหน, \\ วตฺตํ อาวาสิเกหิเม. \\ คมิกา ทารุมตฺติ จ, \\ วิวริตฺวา น ปุจฺฉิย. \\ นสฺสนฺติ จ อคุตฺตญฺจ, \\ อุชฺฌายนฺติ จ เปสลา. \\ ปฏิสาเมตฺวา ถเกตฺวา, \\ อาปุจฺฉิตฺวาว ปกฺกเม. \\ ภิกฺขุ วา สามเณโร วา, \\ อารามิโก อุปาสโก. \\ ปาสาณเกสุ จ ปุญฺชํ, \\ ปฏิสาเมถเกยฺย จ. \\ สเจ อุสฺสหติ อุสฺสุกฺกํ, \\ อโนวสฺเส ตเถว จ. \\ สพฺโพ โอวสฺสติ คามํ, \\ อชฺโฌกาเส ตเถว จ. \\ อปฺเปวงฺคานิ เสเสยฺยุํ, \\ วตฺตํ คมิกภิกฺขุนา. \\ นานุโมทนฺติ เถเรน, \\ โอหาย จตุปญฺจหิ. \\ วจฺจิโต มุจฺฉิโต อาสิ, \\ วตฺตานุโมทเนสุเม. \\ ฉพฺพคฺคิยา ทุนฺนิวตฺถา, \\ อโถปิ จ ทุปฺปารุกา. \\ อนากปฺปา จ โวกฺกมฺม, \\ เถเร จ อนุปขชฺชเน. \\ นเว ภิกฺขู จ สํฆาฏิ, \\ อุชฺฌายนฺติ จ เปสลา. \\ ติมณฺฑลํ นิวาเสตฺวา, \\ กายสคุณคณฺฐิกา. \\ น โวกฺกมฺม ปฏิจฺฉนฺนํ, \\ สุสํวุโตกฺขิตฺตจกฺขุ. \\ อุกฺขิตฺโตชฺชคฺฆิกาสทฺโท, \\ ตโย เจว ปจาลนา. \\ ขมฺโภคุณฺฐิอุกฺกุฏิกา, \\ ปฏิจฺฉนฺนํ สุสํวุโต. \\ โอกฺขิตฺตุกฺขิตฺตอุชฺชคฺฆิ, \\ อปฺปสทฺโท ตโย จลา. \\ ขมฺโภคุณฺฐิปลฺลตฺถิ จ, \\ อนุปขชฺช นาสเน. \\ โอตฺถริตฺวาน อุทเก, \\ นีจํ กตฺวาน สิญฺจิยา. \\ ปฏิ สามนฺตา สํฆาฏิ, \\ โอทเน จ ปฏิคฺคเห. \\ สูปํ อุตฺตริภงฺเคน, \\ สพฺเพสํ สมติตฺถิ จ. \\ สกฺกจฺจํ, \\ ปตฺตสญฺญี จ, สปทานญฺจ สูปกํ. \\ น ถูปโต ปฏิจฺฉาเท, \\ วิญฺญตฺตุชฺฌานสญฺญินา. \\ มหนฺตมณฺฑลทฺวารํ, \\ สพฺพหตฺโถ น พฺยาหเร. \\ อุกฺเขโป เฉทนาคณฺฑ, \\ ธุนํ สิตฺถาวการกํ. \\ ชิวฺหานิจฺฉารกญฺเจว, \\ จปุจปุ สุรุสุรุ. \\ หตฺถปตฺโตฏฺฐนิลฺเลหํ, \\ สามิเสน ปฏิคฺคเห. \\ ยาว น สพฺเพ อุทเก, \\ นีจํ กตฺวาน สิญฺจิยํ. \\ ปฏิ สามนฺตา สํฆาฏิ, \\ นีจํ กตฺวา ฉมาย จ. \\ สสิตฺถกํ นิวตฺตนฺเต, \\ สุปฺปฏิจฺฉนฺนมุกฺกุฏิ. \\ ธมฺมราเชน ปญฺญตฺตํ, \\ อิทํ ภตฺตคฺควตฺตนํ. \\ ทุนฺนิวตฺถา อนากปฺปา, \\ อสลฺลกฺเขตฺวา จ สหสา. \\ ทูเร อจฺจ จิรํ ลหุํ, \\ ตเถว ปิณฺฑจาริโก. \\ ปฏิจฺฉนฺโนว คจฺเฉยฺย, \\ สุสํวุโตกฺขิตฺตจกฺขุ. \\ อุกฺขิตฺโตชฺชคฺฆิกาสทฺโท, \\ ตโย เจว ปจาลนา. \\ ขมฺโภคุณฺฐิอุกฺกุฏิกา, \\ สลฺลกฺเขตฺวา จ สหสา. \\ ทูเร อจฺจ จิรํ ลหุํ, \\ อาสนกํ กฏจฺฉุกา. \\ ภาชนํ วา ฐเปติ จ, \\ อุจฺจาเรตฺวา ปณาเมตฺวา. \\ ปฏิคฺคเห น อุลฺโลเก, \\ สูเปสุปิ ตเถว ตํ. \\ ภิกฺขุ สํฆาฏิยา ฉาเท, \\ ปฏิจฺฉนฺเนว คจฺฉิยํ. \\ สํวุโตกฺขิตฺตจกฺขุ จ, \\ อุกฺขิตฺโตชฺชคฺฆิกาย จ. \\ อปฺปสทฺโท ตโย จาลา, \\ ขมฺโภคุณฺฐิกอุกฺกุฏิ. \\ ปฐมาสนวกฺการ, \\ ปานิยํ ปริโภชนี. \\ ปจฺฉากงฺขติ ภุญฺเชยฺย, \\ โอปิลาเปยฺย อุทฺธเร. \\ ปฏิสาเมยฺย สมฺมชฺเช, \\ ริตฺตํ ตุจฺฉํ อุปฏฺฐเป. \\ หตฺถวิกาเร ภินฺเทยฺย, \\ วตฺติทํ ปิณฺฑจาริเก. \\ ปานี ปริ อคฺคิรณิ, \\ นกฺขตฺตทิสโจรา จ.}
\setlength{\csromangathaleftcol}{0pt}
\csromangathagroup{\csromangathastanza{\csromangathawak{สอุปาหนา ฉตฺตา จ โอคุณฺฐิ สีสํ ปานียํ. \\ นาภิวาเท น ปุจฺฉนฺติ, \\ อหิ อุชฺฌนฺติ เปสลา. \\ โอมุญฺจิ ฉตฺตํ ขนฺเธ จ,}}\csromangathastanza{\csromangathawak{อตรญฺจ ปฏิกฺกมํ. \\ ปตฺตจีวรํ นิกฺขิปา, \\ ปติรูปญฺจ ปุจฺฉิตา. \\ อาสิญฺเจยฺย โธวิเตน,}}\csromangathastanza{\csromangathawak{สุกฺเขนลฺเลนุปาหนา. \\ วุฑฺโฒ นวโก ปุจฺเฉยฺย, \\ อชฺฌาวุฏฺฐญฺจ โคจรา. \\ เสกฺขา วจฺจา ปานี ปริ,}}\csromangathastanza{\csromangathawak{\csromanfolio{411}กตฺตรํ กติกํ ตโต. \\ กาลํ มุหุตฺตํ อุกฺลาโป, \\ ภูมตฺถรณํ นีหเร. \\ ปฏิปาโท ภิสิพิพฺโพ,}}\csromangathastanza{\csromangathawak{มญฺจปีฐญฺจ มลฺลกํ. \\ อปสฺเสนุลฺโลกกณฺณา, \\ เครุกา กาฬ อกตา. \\ สงฺการญฺจ ภูมตฺถรณํ,}}\csromangathastanza{\csromangathawak{ปฏิปาทกํ มญฺจปีฐํ. \\ ภิสิ นิสีทนมฺปิ, \\ มลฺลกํ อปสฺเสน จ. \\ ปตฺตจีวรํ ภูมิ จ,}}\csromangathastanza{\csromangathawak{ปารนฺตํ โอรโต โภคํ. \\ ปุรตฺถิมา ปจฺฉิมา จ, \\ อุตฺตรา อถ ทกฺขิณา. \\ สีตุเณฺห จ ทิวารตฺติํ,}}\csromangathastanza{\csromangathawak{ปริเวณญฺจ โกฏฺฐโก. \\ อุปฏฺฐานคฺคิ สาลา จ, \\ วตฺตํ วจฺจกุฏีสุ จ. \\ ปานี ปริโภชนิยา,}}\csromangathastanza{\csromangathawak{กุมฺภิ อาจมเนสุ จ. \\ อโนปเมน ปญฺญตฺตํ, \\ วตฺตํ อาคนฺตุเกหิเม\csromansharedfootnote{51283409dc78d62e}{เว (ก, เอวมุปริปิ)}. \\ เนวาสนํ น อุทกํ,}}\csromangathastanza{\csromangathawak{น ปจฺจุ น จ ปานิยํ. \\ นาภิวาเท นปญฺญเป, \\ อุชฺฌายนฺติ จ เปสลา. \\ วุฑฺฒาสนญฺจ อุทกํ,}}\csromangathastanza{\csromangathawak{ปจฺจุคฺคนฺตฺวา จ ปานิยํ. \\ อุปาหเน เอกมนฺตํ, \\ อภิวาเท จ ปญฺญเป. \\ วุตฺถํ โคจรเสกฺโข จ,}}\csromangathastanza{\csromangathawak{ฐานํ ปานียโภชนํ. \\ กตฺตรํ กติกํ กาลํ, \\ นวกสฺส นิสินฺนเก. \\ อภิวาทเย อาจิกฺเข,}}\csromangathastanza{\csromangathawak{ยถา เหฏฺฐา ตถา นเย. \\ นิทฺทิฏฺฐํ สตฺถวาเหน, \\ วตฺตํ อาวาสิเกหิเม. \\ คมิกา ทารุมตฺติ จ,}}\csromangathastanza{\csromangathawak{วิวริตฺวา น ปุจฺฉิย. \\ นสฺสนฺติ จ อคุตฺตญฺจ, \\ อุชฺฌายนฺติ จ เปสลา. \\ ปฏิสาเมตฺวา ถเกตฺวา,}}\csromangathastanza{\csromangathawak{อาปุจฺฉิตฺวาว ปกฺกเม. \\ ภิกฺขุ วา สามเณโร วา, \\ อารามิโก อุปาสโก. \\ ปาสาณเกสุ จ ปุญฺชํ,}}\csromangathastanza{\csromangathawak{ปฏิสาเมถเกยฺย จ. \\ สเจ อุสฺสหติ อุสฺสุกฺกํ, \\ อโนวสฺเส ตเถว จ. \\ สพฺโพ โอวสฺสติ คามํ,}}\csromangathastanza{\csromangathawak{\csromanfolio{412}อชฺโฌกาเส ตเถว จ. \\ อปฺเปวงฺคานิ เสเสยฺยุํ, \\ วตฺตํ คมิกภิกฺขุนา. \\ นานุโมทนฺติ เถเรน,}}\csromangathastanza{\csromangathawak{โอหาย จตุปญฺจหิ. \\ วจฺจิโต มุจฺฉิโต อาสิ, \\ วตฺตานุโมทเนสุเม. \\ ฉพฺพคฺคิยา ทุนฺนิวตฺถา,}}\csromangathastanza{\csromangathawak{อโถปิ จ ทุปฺปารุกา. \\ อนากปฺปา จ โวกฺกมฺม, \\ เถเร จ อนุปขชฺชเน. \\ นเว ภิกฺขู จ สํฆาฏิ,}}\csromangathastanza{\csromangathawak{อุชฺฌายนฺติ จ เปสลา. \\ ติมณฺฑลํ นิวาเสตฺวา, \\ กายสคุณคณฺฐิกา. \\ น โวกฺกมฺม ปฏิจฺฉนฺนํ,}}\csromangathastanza{\csromangathawak{สุสํวุโตกฺขิตฺตจกฺขุ. \\ อุกฺขิตฺโตชฺชคฺฆิกาสทฺโท, \\ ตโย เจว ปจาลนา. \\ ขมฺโภคุณฺฐิอุกฺกุฏิกา,}}\csromangathastanza{\csromangathawak{ปฏิจฺฉนฺนํ สุสํวุโต. \\ โอกฺขิตฺตุกฺขิตฺตอุชฺชคฺฆิ, \\ อปฺปสทฺโท ตโย จลา. \\ ขมฺโภคุณฺฐิปลฺลตฺถิ จ,}}\csromangathastanza{\csromangathawak{อนุปขชฺช นาสเน. \\ โอตฺถริตฺวาน อุทเก, \\ นีจํ กตฺวาน สิญฺจิยา. \\ ปฏิ สามนฺตา สํฆาฏิ,}}\csromangathastanza{\csromangathawak{โอทเน จ ปฏิคฺคเห. \\ สูปํ อุตฺตริภงฺเคน, \\ สพฺเพสํ สมติตฺถิ จ. \\ สกฺกจฺจํ,}}\csromangathastanza{\csromangathawak{ปตฺตสญฺญี จ, สปทานญฺจ สูปกํ. \\ น ถูปโต ปฏิจฺฉาเท, \\ วิญฺญตฺตุชฺฌานสญฺญินา. \\ มหนฺตมณฺฑลทฺวารํ,}}\csromangathastanza{\csromangathawak{สพฺพหตฺโถ น พฺยาหเร. \\ อุกฺเขโป เฉทนาคณฺฑ, \\ ธุนํ สิตฺถาวการกํ. \\ ชิวฺหานิจฺฉารกญฺเจว,}}\csromangathastanza{\csromangathawak{จปุจปุ สุรุสุรุ. \\ หตฺถปตฺโตฏฺฐนิลฺเลหํ, \\ สามิเสน ปฏิคฺคเห. \\ ยาว น สพฺเพ อุทเก,}}\csromangathastanza{\csromangathawak{นีจํ กตฺวาน สิญฺจิยํ. \\ ปฏิ สามนฺตา สํฆาฏิ, \\ นีจํ กตฺวา ฉมาย จ. \\ สสิตฺถกํ นิวตฺตนฺเต,}}\csromangathastanza{\csromangathawak{สุปฺปฏิจฺฉนฺนมุกฺกุฏิ. \\ ธมฺมราเชน ปญฺญตฺตํ, \\ อิทํ ภตฺตคฺควตฺตนํ. \\ ทุนฺนิวตฺถา อนากปฺปา,}}\csromangathastanza{\csromangathawak{\csromanfolio{413}อสลฺลกฺเขตฺวา จ สหสา. \\ ทูเร อจฺจ จิรํ ลหุํ, \\ ตเถว ปิณฺฑจาริโก. \\ ปฏิจฺฉนฺโนว คจฺเฉยฺย,}}\csromangathastanza{\csromangathawak{สุสํวุโตกฺขิตฺตจกฺขุ. \\ อุกฺขิตฺโตชฺชคฺฆิกาสทฺโท, \\ ตโย เจว ปจาลนา. \\ ขมฺโภคุณฺฐิอุกฺกุฏิกา,}}\csromangathastanza{\csromangathawak{สลฺลกฺเขตฺวา จ สหสา. \\ ทูเร อจฺจ จิรํ ลหุํ, \\ อาสนกํ กฏจฺฉุกา. \\ ภาชนํ วา ฐเปติ จ,}}\csromangathastanza{\csromangathawak{อุจฺจาเรตฺวา ปณาเมตฺวา. \\ ปฏิคฺคเห น อุลฺโลเก, \\ สูเปสุปิ ตเถว ตํ. \\ ภิกฺขุ สํฆาฏิยา ฉาเท,}}\csromangathastanza{\csromangathawak{ปฏิจฺฉนฺเนว คจฺฉิยํ. \\ สํวุโตกฺขิตฺตจกฺขุ จ, \\ อุกฺขิตฺโตชฺชคฺฆิกาย จ. \\ อปฺปสทฺโท ตโย จาลา,}}\csromangathastanza{\csromangathawak{ขมฺโภคุณฺฐิกอุกฺกุฏิ. \\ ปฐมาสนวกฺการ, \\ ปานิยํ ปริโภชนี. \\ ปจฺฉากงฺขติ ภุญฺเชยฺย,}}\csromangathastanza{\csromangathawak{โอปิลาเปยฺย อุทฺธเร. \\ ปฏิสาเมยฺย สมฺมชฺเช, \\ ริตฺตํ ตุจฺฉํ อุปฏฺฐเป. \\ หตฺถวิกาเร ภินฺเทยฺย,}}\csromangathastanza{\csromangathawak{วตฺติทํ ปิณฺฑจาริเก. \\ ปานี ปริ อคฺคิรณิ, \\ นกฺขตฺตทิสโจรา จ.}}}

\vspace{\tipitakaparskip}
\csromancenter{“สพฺพํ นตฺถี”ติ โกฏฺเฏตฺวา, ปตฺตํเส จีวรํ ตโต.}
\setlength{\csromangathapagewidth}{0pt}
\setlength{\csromangathawakwidth}{0pt}
\csromangathameasuregroup{อิทานิ อํเส ลคฺเคตฺวา, \\ ยถา ปิณฺฑจาริวตฺตํ, \\ ปตฺตํเส จีวรํ สีเส, \\ ปริโภชนิยํ อคฺคิ, \\ นกฺขตฺตํ สปฺปเทสํ วา, \\ สตฺตุตฺตเมน ปญฺญตฺตํ, \\ อชฺโฌกาเส โอกิริํสุ, \\ สเจ วิหาโร อุกฺลาโป, \\ ภิสิพิพฺโพหนํ มญฺจํ, \\ อปสฺเสนุลฺโลกกณฺณา, \\ สงฺการํ ภิกฺขุสามนฺตา, \\ ปริโภชนสามนฺตา, \\ อโธวาเต อตฺถรณํ, \\ ปีฐํ ภิสิ นิสีทนํ, \\ ปตฺตจีวรํ ภูมิ จ, \\ ปุรตฺถิมา จ ปจฺฉิมา, \\ สีตุเณฺห จ ทิวา รตฺติํ, \\ อุปฏฺฐานคฺคิสาลา จ, \\ อาจมกุมฺภิ วุฑฺเฒ จ, \\ ธมฺโม ปทีปํ วิชฺฌาเป, \\ เยน วุฑฺโฒ ปริวตฺติ, \\ ปญฺญเปสิ มหาวีโร, \\ นิวาริยมานา ทฺวารํ, \\ ฉาริกํ ฉฑฺฑเย ชนฺตา, \\ ปริเวณํ โกฏฺฐโก สาลา, \\ มุขํ ปุรโต น เถเร, \\ ปุรโต อุปริมคฺโค, \\ วิชฺฌาเปตฺวา ถเกตฺวา จ, \\ นาจเมติ ยถาวุฑฺฒุํ, \\ อุพฺภชิ นิตฺถุโน กฏฺฐํ, \\ ผรุสา กูป สหสา, \\ พหิ อนฺโต จ อุกฺกาเส, \\ สหสา อุพฺภชิ ฐิเต, \\ ปสฺสาว เขฬ ผรุสา, \\ นติสหสา อุพฺภชิ, \\ น เสสเย ปฏิจฺฉาเท, \\ วจฺจกุฏี ปริภณฺฑํ, \\ อาจมเน จ อุทกํ, \\ อุปาหนา ทนฺตกฏฺฐํ, \\ ยาคุ อุทกํ โธวิตฺวา, \\ นิวาสนา กายพนฺธา, \\ ปจฺฉา ติมณฺฑโล เจว, \\ สคุณํ โธวิตฺวา ปจฺฉา, \\ ภณมานสฺส อาปตฺติ, \\ อุทกํ ปีฐกถลิ, \\ โอตาเป นิทหิ ภงฺโค, \\ ปานียํ อุทกํ นีจํ, \\ ปตฺตจีวรํ ภูมิ จ, \\ อุทฺธเร ปฏิสาเม จ, \\ สีตํ อุณฺหํ ชนฺตาฆรํ, \\ ปีฐญฺจ จีวรํ จุณฺณํ, \\ ปุรโต เถเร นเว จ, \\ ปุรโต อุทเก นฺหาเต, \\ นิวาสนญฺจ สํฆาฏิ, \\ ปาโท ปีฐํ กถลิญฺจ, \\ อุกฺลาปํ สุโสเธยฺย, \\ นิสีทนปจฺจตฺถรณํ, \\ มญฺโจ ปีฐํ ปฏิปาทํ, \\ ภูม สนฺตาน อาโลก, \\ ภูมตฺถรปฏิปาทา, \\ นิสีทถรณํ เขฬ, \\ ปุรตฺถิมา ปจฺฉิมา จ, \\ สีตุณฺหญฺจ ทิวา รตฺติํ, \\ อุปฏฺฐานคฺคิสาลา จ, \\ อาจมํ อนภิรติ, \\ มูลมานตฺตอพฺภานํ, \\ ปพฺพาช ปฏิสารณี, \\ โธเว กาตพฺพํ รชญฺจ, \\ ปตฺตญฺจ จีวรญฺจาปิ, \\ ปริกมฺมํ เวยฺยาวจฺจํ, \\ น สุสานํ ทิสา เจว, \\ สทฺธิวิหาริเกเนตํ, \\ โอวาทสาสนุทฺเทสา, \\ ปริกฺขาโร คิลาโน จ, \\ อุปชฺฌาเยสุ เย วตฺตา, \\ สทฺธิวิหาริเก วตฺตา, \\ อาคนฺตุเกสุ เย วตฺตา, \\ คมิกานุโมทนิกา, \\ อารญฺญเกสุ ยํ วตฺตํ, \\ ชนฺตาฆเร วจฺจกุฏี, \\ อาจริเยสุ ยํ วตฺตํ, \\ เอกูนวีสติ วตฺถู, \\ วตฺตํ อปริปูเรนฺโต, \\ อสุทฺธสีโล ทุปฺปญฺโญ, \\ วิกฺขิตฺตจิตฺโตเนกคฺโค, \\ อปสฺสมาโน สทฺธมฺมํ, \\ ยํ วตฺตํ ปริปูเรนฺโต, \\ วิสุทฺธสีโล สปฺปญฺโญ, \\ อวิกฺขิตฺตจิตฺโต เอกคฺโค, \\ สมฺปสฺสมาโน สทฺธมฺมํ, \\ ตสฺมา หิ วตฺตํ ปูเรยฺย, \\ โอวาทํ พุทฺธเสฏฺฐสฺส,}{\csromangathabat{อิทานิ อํเส ลคฺเคตฺวา,}{ติมณฺฑลํ ปริมณฺฑลํ.} \\ \csromangathabat{ยถา ปิณฺฑจาริวตฺตํ,}{นเย อารญฺญเกสุปิ.} \\ \csromangathabat{ปตฺตํเส จีวรํ สีเส,}{อาโรหิตฺวา จ ปานิยํ.} \\ \csromangathabat{ปริโภชนิยํ อคฺคิ,}{อรณี จาปิ กตฺตรํ.} \\ \csromangathabat{นกฺขตฺตํ สปฺปเทสํ วา,}{ทิสาปิ กุสโล ภเว.} \\ \csromangathabat{สตฺตุตฺตเมน ปญฺญตฺตํ,}{วตฺตํ อรญฺญเกสุเม.} \\ \csromangathabat{อชฺโฌกาเส โอกิริํสุ,}{อุชฺฌายนฺติ จ เปสลา.} \\ \csromangathabat{สเจ วิหาโร อุกฺลาโป,}{ปฐมํ ปตฺตจีวรํ.} \\ \csromangathabat{ภิสิพิพฺโพหนํ มญฺจํ,}{ปีฐญฺจ เขฬมลฺลกํ.} \\ \csromangathabat{อปสฺเสนุลฺโลกกณฺณา,}{เครุกา กาฬ อกตา.} \\ \csromangathabat{สงฺการํ ภิกฺขุสามนฺตา,}{เสนาวิหารปานิยํ.} \\ \csromangathabat{ปริโภชนสามนฺตา,}{ปฏิวาเต จ องฺคเณ.} \\ \csromangathabat{อโธวาเต อตฺถรณํ,}{ปฏิปาทกมญฺโจ จ.} \\ \csromangathabat{ปีฐํ ภิสิ นิสีทนํ,}{มลฺลกํ อปสฺเสน จ.} \\ \csromangathabat{ปตฺตจีวรํ ภูมิ จ,}{ปารนฺตํ โอรโต โภคํ.} \\ \csromangathabat{ปุรตฺถิมา จ ปจฺฉิมา,}{อุตฺตรา อถ ทกฺขิณา.} \\ \csromangathabat{สีตุเณฺห จ ทิวา รตฺติํ,}{ปริเวณญฺจ โกฏฺฐโก.} \\ \csromangathabat{อุปฏฺฐานคฺคิสาลา จ,}{วจฺจกุฏี จ ปานิยํ.} \\ \csromangathabat{อาจมกุมฺภิ วุฑฺเฒ จ,}{อุทฺเทสปุจฺฉนา สชฺฌา.} \\ \csromangathabat{ธมฺโม ปทีปํ วิชฺฌาเป,}{น วิวเร นปิ ถเก.} \\ \csromangathabat{เยน วุฑฺโฒ ปริวตฺติ,}{กณฺเณนปิ น ฆฏฺเฏเย.} \\ \csromangathabat{ปญฺญเปสิ มหาวีโร,}{วตฺตํ เสนาสเนสุ ตํ.} \\ \csromangathabat{นิวาริยมานา ทฺวารํ,}{มุจฺฉิตุชฺฌนฺติ เปสลา.} \\ \csromangathabat{ฉาริกํ ฉฑฺฑเย ชนฺตา,}{ปริภณฺฑํ ตเถว จ.} \\ \csromangathabat{ปริเวณํ โกฏฺฐโก สาลา,}{จุณฺณมตฺติกโทณิกา.} \\ \csromangathabat{มุขํ ปุรโต น เถเร,}{น นเว อุสฺสหติ สเจ.} \\ \csromangathabat{ปุรโต อุปริมคฺโค,}{จิกฺขลฺลํ มตฺติ ปีฐกํ.} \\ \csromangathabat{วิชฺฌาเปตฺวา ถเกตฺวา จ,}{วตฺตํ ชนฺตาฆเรสุเม.} \\ \csromangathabat{นาจเมติ ยถาวุฑฺฒุํ,}{ปฏิปาฏิ จ สหสา.} \\ \csromangathabat{อุพฺภชิ นิตฺถุโน กฏฺฐํ,}{วจฺจํ ปสฺสาว เขฬกํ.} \\ \csromangathabat{ผรุสา กูป สหสา,}{อุพฺภชิ จปุ เสเสน.} \\ \csromangathabat{พหิ อนฺโต จ อุกฺกาเส,}{รชฺชุ อตรมานญฺจ.} \\ \csromangathabat{สหสา อุพฺภชิ ฐิเต,}{นิตฺถุเน กฏฺฐ วจฺจญฺจ.} \\ \csromangathabat{ปสฺสาว เขฬ ผรุสา,}{กูปญฺจ วจฺจปาทุเก.} \\ \csromangathabat{นติสหสา อุพฺภชิ,}{ปาทุกาย จปุจปุ.} \\ \csromangathabat{น เสสเย ปฏิจฺฉาเท,}{อุหตปิธเรน จ.} \\ \csromangathabat{วจฺจกุฏี ปริภณฺฑํ,}{ปริเวณญฺจ โกฏฺฐโก.} \\ \csromangathabat{อาจมเน จ อุทกํ,}{วตฺตํ วจฺจกุฏีสุเม.} \\ \csromangathabat{อุปาหนา ทนฺตกฏฺฐํ,}{มุโขทกญฺจ อาสนํ.} \\ \csromangathabat{ยาคุ อุทกํ โธวิตฺวา,}{อุทฺธารุกฺลาป คาม จ.} \\ \csromangathabat{นิวาสนา กายพนฺธา,}{สคุณํ ปตฺตโสทกํ.} \\ \csromangathabat{ปจฺฉา ติมณฺฑโล เจว,}{ปริมณฺฑล พนฺธนํ.} \\ \csromangathabat{สคุณํ โธวิตฺวา ปจฺฉา,}{นาติทูเร ปฏิคฺคเห.} \\ \csromangathabat{ภณมานสฺส อาปตฺติ,}{ปฐมาคนฺตฺวาน อาสนํ.} \\ \csromangathabat{อุทกํ ปีฐกถลิ,}{ปจฺจุคฺคนฺตฺวา นิวาสนํ.} \\ \csromangathabat{โอตาเป นิทหิ ภงฺโค,}{โอโภเค ภุญฺชิตุ นเม.} \\ \csromangathabat{ปานียํ อุทกํ นีจํ,}{มุหุตฺตํ น จ นิทเห.} \\ \csromangathabat{ปตฺตจีวรํ ภูมิ จ,}{ปารนฺตํ โอรโต โภคํ.} \\ \csromangathabat{อุทฺธเร ปฏิสาเม จ,}{อุกฺลาโป จ นหายิตุํ.} \\ \csromangathabat{สีตํ อุณฺหํ ชนฺตาฆรํ,}{จุณฺณํ มตฺติก ปิฏฺฐิโต.} \\ \csromangathabat{ปีฐญฺจ จีวรํ จุณฺณํ,}{มตฺติกุสฺสหติ มุขํ.} \\ \csromangathabat{ปุรโต เถเร นเว จ,}{ปริกมฺมญฺจ นิกฺขเม.} \\ \csromangathabat{ปุรโต อุทเก นฺหาเต,}{นิวาเสตฺวา อุปชฺฌายํ.} \\ \csromangathabat{นิวาสนญฺจ สํฆาฏิ,}{ปีฐกํ อาสเนน จ.} \\ \csromangathabat{ปาโท ปีฐํ กถลิญฺจ,}{ปานียุทฺเทสปุจฺฉนา.} \\ \csromangathabat{อุกฺลาปํ สุโสเธยฺย,}{ปฐมํ ปตฺตจีวรํ.} \\ \csromangathabat{นิสีทนปจฺจตฺถรณํ,}{ภิสิ พิพฺโพหนานิ จ.} \\ \csromangathabat{มญฺโจ ปีฐํ ปฏิปาทํ,}{มลฺลกํ อปสฺเสน จ.} \\ \csromangathabat{ภูม สนฺตาน อาโลก,}{เครุกา กาฬ อกตา.} \\ \csromangathabat{ภูมตฺถรปฏิปาทา,}{มญฺโจ ปีฐํ พิพฺโพหนํ.} \\ \csromangathabat{นิสีทถรณํ เขฬ,}{อปสฺเส ปตฺตจีวรํ.} \\ \csromangathabat{ปุรตฺถิมา ปจฺฉิมา จ,}{อุตฺตรา อถ ทกฺขิณา.} \\ \csromangathabat{สีตุณฺหญฺจ ทิวา รตฺติํ,}{ปริเวณญฺจ โกฏฺฐโก.} \\ \csromangathabat{อุปฏฺฐานคฺคิสาลา จ,}{วจฺจปานิยโภชนี.} \\ \csromangathabat{อาจมํ อนภิรติ,}{กุกฺกุจฺจํ ทิฏฺฐิ จ ครุ.} \\ \csromangathabat{มูลมานตฺตอพฺภานํ,}{ตชฺชนียํ นิยสฺสกํ.} \\ \csromangathabat{ปพฺพาช ปฏิสารณี,}{อุกฺเขปญฺจ กตํ ยทิ.} \\ \csromangathabat{โธเว กาตพฺพํ รชญฺจ,}{รเช สมฺปริวตฺตกํ.} \\ \csromangathabat{ปตฺตญฺจ จีวรญฺจาปิ,}{ปริกฺขารญฺจ เฉทนํ.} \\ \csromangathabat{ปริกมฺมํ เวยฺยาวจฺจํ,}{ปจฺฉา ปิณฺฑํ ปวิสนํ.} \\ \csromangathabat{น สุสานํ ทิสา เจว,}{ยาวชีวํ อุปฏฺฐเห.} \\ \csromangathabat{สทฺธิวิหาริเกเนตํ,}{วตฺตุปชฺฌายเกสุเม.} \\ \csromangathabat{โอวาทสาสนุทฺเทสา,}{ปุจฺฉา ปตฺตญฺจ จีวรํ.} \\ \csromangathabat{ปริกฺขาโร คิลาโน จ,}{น ปจฺฉาสมโณ ภเว.} \\ \csromangathabat{อุปชฺฌาเยสุ เย วตฺตา,}{เอวํ อาจริเยสุปิ.} \\ \csromangathabat{สทฺธิวิหาริเก วตฺตา,}{ตเถว อนฺเตวาสิเก.} \\ \csromangathabat{อาคนฺตุเกสุ เย วตฺตา,}{ปุน อาวาสิเกสุ จ.} \\ \csromangathabat{คมิกานุโมทนิกา,}{ภตฺตคฺเค ปิณฺฑจาริเก.} \\ \csromangathabat{อารญฺญเกสุ ยํ วตฺตํ,}{ยญฺจ เสนาสเนสุปิ.} \\ \csromangathabat{ชนฺตาฆเร วจฺจกุฏี,}{อุปชฺฌา สทฺธิวิหาริเก.} \\ \csromangathabat{อาจริเยสุ ยํ วตฺตํ,}{ตเถว อนฺเตวาสิเก.} \\ \csromangathabat{เอกูนวีสติ วตฺถู,}{วตฺตา จุทฺทส ขนฺธเก.} \\ \csromangathabat{วตฺตํ อปริปูเรนฺโต,}{น สีลํ ปริปูรติ.} \\ \csromangathabat{อสุทฺธสีโล ทุปฺปญฺโญ,}{จิตฺเตกคฺคํ น วินฺทติ.} \\ \csromangathabat{วิกฺขิตฺตจิตฺโตเนกคฺโค,}{สมฺมา ธมฺมํ น ปสฺสติ.} \\ \csromangathabat{อปสฺสมาโน สทฺธมฺมํ,}{ทุกฺขา น ปริมุจฺจติ.} \\ \csromangathabat{ยํ วตฺตํ ปริปูเรนฺโต,}{สีลมฺปิ ปริปูรติ.} \\ \csromangathabat{วิสุทฺธสีโล สปฺปญฺโญ,}{จิตฺเตกคฺคมฺปิ วินฺทติ.} \\ \csromangathabat{อวิกฺขิตฺตจิตฺโต เอกคฺโค,}{สมฺมา ธมฺมํ วิปสฺสติ.} \\ \csromangathabat{สมฺปสฺสมาโน สทฺธมฺมํ,}{ทุกฺขา โส ปริมุจฺจติ.} \\ \csromangathabat{ตสฺมา หิ วตฺตํ ปูเรยฺย,}{ชินปุตฺโต วิจฺจกฺขโณ.} \\ \csromangathabat{โอวาทํ พุทฺธเสฏฺฐสฺส,}{ตโต นิพฺพานเมหิตีติ.}}
\csromangathasetleft{อิทานิ อํเส ลคฺเคตฺวา, \\ ยถา ปิณฺฑจาริวตฺตํ, \\ ปตฺตํเส จีวรํ สีเส, \\ ปริโภชนิยํ อคฺคิ, \\ นกฺขตฺตํ สปฺปเทสํ วา, \\ สตฺตุตฺตเมน ปญฺญตฺตํ, \\ อชฺโฌกาเส โอกิริํสุ, \\ สเจ วิหาโร อุกฺลาโป, \\ ภิสิพิพฺโพหนํ มญฺจํ, \\ อปสฺเสนุลฺโลกกณฺณา, \\ สงฺการํ ภิกฺขุสามนฺตา, \\ ปริโภชนสามนฺตา, \\ อโธวาเต อตฺถรณํ, \\ ปีฐํ ภิสิ นิสีทนํ, \\ ปตฺตจีวรํ ภูมิ จ, \\ ปุรตฺถิมา จ ปจฺฉิมา, \\ สีตุเณฺห จ ทิวา รตฺติํ, \\ อุปฏฺฐานคฺคิสาลา จ, \\ อาจมกุมฺภิ วุฑฺเฒ จ, \\ ธมฺโม ปทีปํ วิชฺฌาเป, \\ เยน วุฑฺโฒ ปริวตฺติ, \\ ปญฺญเปสิ มหาวีโร, \\ นิวาริยมานา ทฺวารํ, \\ ฉาริกํ ฉฑฺฑเย ชนฺตา, \\ ปริเวณํ โกฏฺฐโก สาลา, \\ มุขํ ปุรโต น เถเร, \\ ปุรโต อุปริมคฺโค, \\ วิชฺฌาเปตฺวา ถเกตฺวา จ, \\ นาจเมติ ยถาวุฑฺฒุํ, \\ อุพฺภชิ นิตฺถุโน กฏฺฐํ, \\ ผรุสา กูป สหสา, \\ พหิ อนฺโต จ อุกฺกาเส, \\ สหสา อุพฺภชิ ฐิเต, \\ ปสฺสาว เขฬ ผรุสา, \\ นติสหสา อุพฺภชิ, \\ น เสสเย ปฏิจฺฉาเท, \\ วจฺจกุฏี ปริภณฺฑํ, \\ อาจมเน จ อุทกํ, \\ อุปาหนา ทนฺตกฏฺฐํ, \\ ยาคุ อุทกํ โธวิตฺวา, \\ นิวาสนา กายพนฺธา, \\ ปจฺฉา ติมณฺฑโล เจว, \\ สคุณํ โธวิตฺวา ปจฺฉา, \\ ภณมานสฺส อาปตฺติ, \\ อุทกํ ปีฐกถลิ, \\ โอตาเป นิทหิ ภงฺโค, \\ ปานียํ อุทกํ นีจํ, \\ ปตฺตจีวรํ ภูมิ จ, \\ อุทฺธเร ปฏิสาเม จ, \\ สีตํ อุณฺหํ ชนฺตาฆรํ, \\ ปีฐญฺจ จีวรํ จุณฺณํ, \\ ปุรโต เถเร นเว จ, \\ ปุรโต อุทเก นฺหาเต, \\ นิวาสนญฺจ สํฆาฏิ, \\ ปาโท ปีฐํ กถลิญฺจ, \\ อุกฺลาปํ สุโสเธยฺย, \\ นิสีทนปจฺจตฺถรณํ, \\ มญฺโจ ปีฐํ ปฏิปาทํ, \\ ภูม สนฺตาน อาโลก, \\ ภูมตฺถรปฏิปาทา, \\ นิสีทถรณํ เขฬ, \\ ปุรตฺถิมา ปจฺฉิมา จ, \\ สีตุณฺหญฺจ ทิวา รตฺติํ, \\ อุปฏฺฐานคฺคิสาลา จ, \\ อาจมํ อนภิรติ, \\ มูลมานตฺตอพฺภานํ, \\ ปพฺพาช ปฏิสารณี, \\ โธเว กาตพฺพํ รชญฺจ, \\ ปตฺตญฺจ จีวรญฺจาปิ, \\ ปริกมฺมํ เวยฺยาวจฺจํ, \\ น สุสานํ ทิสา เจว, \\ สทฺธิวิหาริเกเนตํ, \\ โอวาทสาสนุทฺเทสา, \\ ปริกฺขาโร คิลาโน จ, \\ อุปชฺฌาเยสุ เย วตฺตา, \\ สทฺธิวิหาริเก วตฺตา, \\ อาคนฺตุเกสุ เย วตฺตา, \\ คมิกานุโมทนิกา, \\ อารญฺญเกสุ ยํ วตฺตํ, \\ ชนฺตาฆเร วจฺจกุฏี, \\ อาจริเยสุ ยํ วตฺตํ, \\ เอกูนวีสติ วตฺถู, \\ วตฺตํ อปริปูเรนฺโต, \\ อสุทฺธสีโล ทุปฺปญฺโญ, \\ วิกฺขิตฺตจิตฺโตเนกคฺโค, \\ อปสฺสมาโน สทฺธมฺมํ, \\ ยํ วตฺตํ ปริปูเรนฺโต, \\ วิสุทฺธสีโล สปฺปญฺโญ, \\ อวิกฺขิตฺตจิตฺโต เอกคฺโค, \\ สมฺปสฺสมาโน สทฺธมฺมํ, \\ ตสฺมา หิ วตฺตํ ปูเรยฺย, \\ โอวาทํ พุทฺธเสฏฺฐสฺส,}
\csromangathagroupwithclosermajor{\csromangathastanza{\csromangathabat{อิทานิ อํเส ลคฺเคตฺวา,}{ติมณฺฑลํ ปริมณฺฑลํ.} \\ \csromangathabat{ยถา ปิณฺฑจาริวตฺตํ,}{นเย อารญฺญเกสุปิ.}}\csromangathastanza{\csromangathabat{ปตฺตํเส จีวรํ สีเส,}{อาโรหิตฺวา จ ปานิยํ.} \\ \csromangathabat{ปริโภชนิยํ อคฺคิ,}{อรณี จาปิ กตฺตรํ.}}\csromangathastanza{\csromangathabat{นกฺขตฺตํ สปฺปเทสํ วา,}{ทิสาปิ กุสโล ภเว.} \\ \csromangathabat{สตฺตุตฺตเมน ปญฺญตฺตํ,}{วตฺตํ อรญฺญเกสุเม.}}\csromangathastanza{\csromangathabat{อชฺโฌกาเส โอกิริํสุ,}{อุชฺฌายนฺติ จ เปสลา.} \\ \csromangathabat{สเจ วิหาโร อุกฺลาโป,}{ปฐมํ ปตฺตจีวรํ.}}\csromangathastanza{\csromangathabat{ภิสิพิพฺโพหนํ มญฺจํ,}{ปีฐญฺจ เขฬมลฺลกํ.} \\ \csromangathabat{อปสฺเสนุลฺโลกกณฺณา,}{เครุกา กาฬ อกตา.}}\csromangathastanza{\csromanfolio{414}\csromangathabat{สงฺการํ ภิกฺขุสามนฺตา,}{เสนาวิหารปานิยํ.} \\ \csromangathabat{ปริโภชนสามนฺตา,}{ปฏิวาเต จ องฺคเณ.}}\csromangathastanza{\csromangathabat{อโธวาเต อตฺถรณํ,}{ปฏิปาทกมญฺโจ จ.} \\ \csromangathabat{ปีฐํ ภิสิ นิสีทนํ,}{มลฺลกํ อปสฺเสน จ.}}\csromangathastanza{\csromangathabat{ปตฺตจีวรํ ภูมิ จ,}{ปารนฺตํ โอรโต โภคํ.} \\ \csromangathabat{ปุรตฺถิมา จ ปจฺฉิมา,}{อุตฺตรา อถ ทกฺขิณา.}}\csromangathastanza{\csromangathabat{สีตุเณฺห จ ทิวา รตฺติํ,}{ปริเวณญฺจ โกฏฺฐโก.} \\ \csromangathabat{อุปฏฺฐานคฺคิสาลา จ,}{วจฺจกุฏี จ ปานิยํ.}}\csromangathastanza{\csromangathabat{อาจมกุมฺภิ วุฑฺเฒ จ,}{อุทฺเทสปุจฺฉนา สชฺฌา.} \\ \csromangathabat{ธมฺโม ปทีปํ วิชฺฌาเป,}{น วิวเร นปิ ถเก.}}\csromangathastanza{\csromangathabat{เยน วุฑฺโฒ ปริวตฺติ,}{กณฺเณนปิ น ฆฏฺเฏเย.} \\ \csromangathabat{ปญฺญเปสิ มหาวีโร,}{วตฺตํ เสนาสเนสุ ตํ.}}\csromangathastanza{\csromangathabat{นิวาริยมานา ทฺวารํ,}{มุจฺฉิตุชฺฌนฺติ เปสลา.} \\ \csromangathabat{ฉาริกํ ฉฑฺฑเย ชนฺตา,}{ปริภณฺฑํ ตเถว จ.}}\csromangathastanza{\csromangathabat{ปริเวณํ โกฏฺฐโก สาลา,}{จุณฺณมตฺติกโทณิกา.} \\ \csromangathabat{มุขํ ปุรโต น เถเร,}{น นเว อุสฺสหติ สเจ.}}\csromangathastanza{\csromangathabat{ปุรโต อุปริมคฺโค,}{จิกฺขลฺลํ มตฺติ ปีฐกํ.} \\ \csromangathabat{วิชฺฌาเปตฺวา ถเกตฺวา จ,}{วตฺตํ ชนฺตาฆเรสุเม.}}\csromangathastanza{\csromangathabat{นาจเมติ ยถาวุฑฺฒุํ,}{ปฏิปาฏิ จ สหสา.} \\ \csromangathabat{อุพฺภชิ นิตฺถุโน กฏฺฐํ,}{วจฺจํ ปสฺสาว เขฬกํ.}}\csromangathastanza{\csromangathabat{ผรุสา กูป สหสา,}{อุพฺภชิ จปุ เสเสน.} \\ \csromangathabat{พหิ อนฺโต จ อุกฺกาเส,}{รชฺชุ อตรมานญฺจ.}}\csromangathastanza{\csromangathabat{สหสา อุพฺภชิ ฐิเต,}{นิตฺถุเน กฏฺฐ วจฺจญฺจ.} \\ \csromangathabat{ปสฺสาว เขฬ ผรุสา,}{กูปญฺจ วจฺจปาทุเก.}}\csromangathastanza{\csromangathabat{นติสหสา อุพฺภชิ,}{ปาทุกาย จปุจปุ.} \\ \csromangathabat{น เสสเย ปฏิจฺฉาเท,}{อุหตปิธเรน จ.}}\csromangathastanza{\csromanfolio{415}\csromangathabat{วจฺจกุฏี ปริภณฺฑํ,}{ปริเวณญฺจ โกฏฺฐโก.} \\ \csromangathabat{อาจมเน จ อุทกํ,}{วตฺตํ วจฺจกุฏีสุเม.}}\csromangathastanza{\csromangathabat{อุปาหนา ทนฺตกฏฺฐํ,}{มุโขทกญฺจ อาสนํ.} \\ \csromangathabat{ยาคุ อุทกํ โธวิตฺวา,}{อุทฺธารุกฺลาป คาม จ.}}\csromangathastanza{\csromangathabat{นิวาสนา กายพนฺธา,}{สคุณํ ปตฺตโสทกํ.} \\ \csromangathabat{ปจฺฉา ติมณฺฑโล เจว,}{ปริมณฺฑล พนฺธนํ.}}\csromangathastanza{\csromangathabat{สคุณํ โธวิตฺวา ปจฺฉา,}{นาติทูเร ปฏิคฺคเห.} \\ \csromangathabat{ภณมานสฺส อาปตฺติ,}{ปฐมาคนฺตฺวาน อาสนํ.}}\csromangathastanza{\csromangathabat{อุทกํ ปีฐกถลิ,}{ปจฺจุคฺคนฺตฺวา นิวาสนํ.} \\ \csromangathabat{โอตาเป นิทหิ ภงฺโค,}{โอโภเค ภุญฺชิตุ นเม.}}\csromangathastanza{\csromangathabat{ปานียํ อุทกํ นีจํ,}{มุหุตฺตํ น จ นิทเห.} \\ \csromangathabat{ปตฺตจีวรํ ภูมิ จ,}{ปารนฺตํ โอรโต โภคํ.}}\csromangathastanza{\csromangathabat{อุทฺธเร ปฏิสาเม จ,}{อุกฺลาโป จ นหายิตุํ.} \\ \csromangathabat{สีตํ อุณฺหํ ชนฺตาฆรํ,}{จุณฺณํ มตฺติก ปิฏฺฐิโต.}}\csromangathastanza{\csromangathabat{ปีฐญฺจ จีวรํ จุณฺณํ,}{มตฺติกุสฺสหติ มุขํ.} \\ \csromangathabat{ปุรโต เถเร นเว จ,}{ปริกมฺมญฺจ นิกฺขเม.}}\csromangathastanza{\csromangathabat{ปุรโต อุทเก นฺหาเต,}{นิวาเสตฺวา อุปชฺฌายํ.} \\ \csromangathabat{นิวาสนญฺจ สํฆาฏิ,}{ปีฐกํ อาสเนน จ.}}\csromangathastanza{\csromangathabat{ปาโท ปีฐํ กถลิญฺจ,}{ปานียุทฺเทสปุจฺฉนา.} \\ \csromangathabat{อุกฺลาปํ สุโสเธยฺย,}{ปฐมํ ปตฺตจีวรํ.}}\csromangathastanza{\csromangathabat{นิสีทน\-ปจฺจตฺถรณํ,}{ภิสิ พิพฺโพหนานิ จ.} \\ \csromangathabat{มญฺโจ ปีฐํ ปฏิปาทํ,}{มลฺลกํ อปสฺเสน จ.}}\csromangathastanza{\csromangathabat{ภูม สนฺตาน อาโลก,}{เครุกา กาฬ อกตา.} \\ \csromangathabat{ภูมตฺถรปฏิปาทา,}{มญฺโจ ปีฐํ พิพฺโพหนํ.}}\csromangathastanza{\csromangathabat{นิสีทถรณํ เขฬ,}{อปสฺเส ปตฺตจีวรํ.} \\ \csromangathabat{ปุรตฺถิมา ปจฺฉิมา จ,}{อุตฺตรา อถ ทกฺขิณา.}}\csromangathastanza{\csromanfolio{416}\csromangathabat{สีตุณฺหญฺจ ทิวา รตฺติํ,}{ปริเวณญฺจ โกฏฺฐโก.} \\ \csromangathabat{อุปฏฺฐานคฺคิสาลา จ,}{วจฺจปานิยโภชนี.}}\csromangathastanza{\csromangathabat{อาจมํ อนภิรติ,}{กุกฺกุจฺจํ ทิฏฺฐิ จ ครุ.} \\ \csromangathabat{มูลมานตฺตอพฺภานํ,}{ตชฺชนียํ นิยสฺสกํ.}}\csromangathastanza{\csromangathabat{ปพฺพาช ปฏิสารณี,}{อุกฺเขปญฺจ กตํ ยทิ.} \\ \csromangathabat{โธเว กาตพฺพํ รชญฺจ,}{รเช สมฺปริวตฺตกํ.}}\csromangathastanza{\csromangathabat{ปตฺตญฺจ จีวรญฺจาปิ,}{ปริกฺขารญฺจ เฉทนํ.} \\ \csromangathabat{ปริกมฺมํ เวยฺยาวจฺจํ,}{ปจฺฉา ปิณฺฑํ ปวิสนํ.}}\csromangathastanza{\csromangathabat{น สุสานํ ทิสา เจว,}{ยาวชีวํ อุปฏฺฐเห.} \\ \csromangathabat{สทฺธิวิหาริเกเนตํ,}{วตฺตุปชฺฌายเกสุเม.}}\csromangathastanza{\csromangathabat{โอวาทสาสนุทฺเทสา,}{ปุจฺฉา ปตฺตญฺจ จีวรํ.} \\ \csromangathabat{ปริกฺขาโร คิลาโน จ,}{น ปจฺฉาสมโณ ภเว.}}\csromangathastanza{\csromangathabat{อุปชฺฌาเยสุ เย วตฺตา,}{เอวํ อาจริเยสุปิ.} \\ \csromangathabat{สทฺธิวิหาริเก วตฺตา,}{ตเถว อนฺเตวาสิเก.}}\csromangathastanza{\csromangathabat{อาคนฺตุเกสุ เย วตฺตา,}{ปุน อาวาสิเกสุ จ.} \\ \csromangathabat{คมิกานุโมทนิกา,}{ภตฺตคฺเค ปิณฺฑจาริเก.}}\csromangathastanza{\csromangathabat{อารญฺญเกสุ ยํ วตฺตํ,}{ยญฺจ เสนาสเนสุปิ.} \\ \csromangathabat{ชนฺตาฆเร วจฺจกุฏี,}{อุปชฺฌา สทฺธิวิหาริเก.}}\csromangathastanza{\csromangathabat{อาจริเยสุ ยํ วตฺตํ,}{ตเถว อนฺเตวาสิเก.} \\ \csromangathabat{เอกูนวีสติ วตฺถู,}{วตฺตา จุทฺทส ขนฺธเก.}}\csromangathastanza{\csromangathabat{วตฺตํ อปริปูเรนฺโต,}{น สีลํ ปริปูรติ.} \\ \csromangathabat{อสุทฺธสีโล ทุปฺปญฺโญ,}{จิตฺเตกคฺคํ น วินฺทติ.}}\csromangathastanza{\csromangathabat{วิกฺขิตฺตจิตฺโตเนกคฺโค,}{สมฺมา ธมฺมํ น ปสฺสติ.} \\ \csromangathabat{อปสฺสมาโน สทฺธมฺมํ,}{ทุกฺขา น ปริมุจฺจติ.}}\csromangathastanza{\csromanfolio{417}\csromangathabat{ยํ วตฺตํ ปริปูเรนฺโต,}{สีลมฺปิ ปริปูรติ.} \\ \csromangathabat{วิสุทฺธสีโล สปฺปญฺโญ,}{จิตฺเตกคฺคมฺปิ วินฺทติ.}}\csromangathastanza{\csromangathabat{อวิกฺขิตฺตจิตฺโต เอกคฺโค,}{สมฺมา ธมฺมํ วิปสฺสติ.} \\ \csromangathabat{สมฺปสฺสมาโน สทฺธมฺมํ,}{ทุกฺขา โส ปริมุจฺจติ.}}\csromangathastanza{\csromangathabat{ตสฺมา หิ วตฺตํ ปูเรยฺย,}{ชินปุตฺโต วิจฺจกฺขโณ.} \\ \csromangathabat{โอวาทํ พุทฺธ\-เสฏฺฐสฺส,}{ตโต นิพฺพานเมหิตีติ.}}}{วตฺต\-กฺขนฺธโก นิฏฺฐิโต.}
\csromanmatikaanchor{mtk.192}
\csromantocmarknopage{chapter}{9. ปาติโมกฺขฏฺฐปนกฺขนฺธก}{mtk.192}
\bookmark[dest={mtk.192},level=0]{9. ปาติโมกฺขฏฺฐปนกฺขนฺธก}
\chapterheadpageread{9. ปาติโมกฺขฏฺฐปน\-กฺขนฺธก}
\csromanmatikaanchor{mtk.193}
\csromantocmarkref{section}{1. ปาติโมกฺขุทฺเทสยาจน}{mtk.193}
\bookmark[dest={mtk.193},level=1]{1. ปาติโมกฺขุทฺเทสยาจน}
\csromanheader{1}{1. ปาติโมกฺขุทฺเทส\-ยาจน}
\proseitem{383}{\csromanfolio{418}\csromansymbolfootnote{*}{ขุ 1. 138; อํ 3. 44 ปิฏฺเฐสุปิ.}เตน สมเยน พุทฺโธ ภควา สาวตฺถิยํ วิหรติ ปุพฺพาราเม มิคารมาตุ ปาสาเท. เตน โข ปน สมเยน ภควา ตทหุโปสเถ ภิกฺขุ\-สํฆ\-ปริวุโต นิสินฺโน โหติ. อถ โข อายสฺมา อานนฺโท อภิกฺกนฺตาย รตฺติยา นิกฺขนฺเต ปฐเม ยาเม อุฏฺฐายาสนา เอกํสํ อุตฺตราสงฺคํ กริตฺวา เยน ภควา เตนญฺชลิํ ปณาเมตฺวา ภควนฺตํ เอตทโวจ “อภิกฺกนฺตา ภนฺเต รตฺติ, นิกฺขนฺโต ปฐโม ยาโม, จิรนิสินฺโน ภิกฺขุสํโฆ, อุทฺทิสตุ ภนฺเต ภควา ภิกฺขูนํ ปาติโมกฺขนฺ”ติ. เอวํ วุตฺเต ภควา ตุณฺหี อโหสิ. ทุติยมฺปิ โข อายสฺมา อานนฺโท อภิกฺกนฺตาย รตฺติยา นิกฺขนฺเต มชฺฌิเม ยาเม อุฏฺฐายาสนา เอกํสํ อุตฺตราสงฺคํ กริตฺวา เยน ภควา เตนญฺชลิํ ปณาเมตฺวา ภควนฺตํ เอตทโวจ “อภิกฺกนฺตา ภนฺเต รตฺติ, นิกฺขนฺโต มชฺฌิโม ยาโม, จิรนิสินฺโน ภิกฺขุสํโฆ, อุทฺทิสตุ ภนฺเต ภควา ภิกฺขูนํ ปาติโมกฺขนฺ”ติ. ทุติยมฺปิ โข ภควา ตุณฺหี อโหสิ. ตติยมฺปิ โข อายสฺมา อานนฺโท อภิกฺกนฺตาย รตฺติยา นิกฺขนฺเต ปจฺฉิเม ยาเม อุทฺธสฺเต อรุเณ นนฺทิมุขิยา รตฺติยา อุฏฺฐายาสนา เอกํสํ อุตฺตราสงฺคํ กริตฺวา เยน ภควา เตนญฺชลิํ ปณาเมตฺวา ภควนฺตํ เอตทโวจ “อภิกฺกนฺตา ภนฺเต รตฺติ, นิกฺขนฺโต ปจฺฉิโม ยาโม, อุทฺธสฺตํ อรุณํ, นนฺทิมุขิ รตฺติ, จิรนิสินฺโน ภิกฺขุสํโฆ, อุทฺทิสตุ ภนฺเต ภควา ภิกฺขูนํ ปาติโมกฺขนฺ”ติ. อปริสุทฺธา อานนฺท ปริสาติ.}

\prose{อถ โข อายสฺมโต มหา\-โมคฺคลฺลานสฺส เอตทโหสิ “กํ นุ โข ภควา ปุคฺคลํ สนฺธาย เอวมาห ‘อปริสุทฺธา อานนฺท ปริสา’ติ”. อถ โข อายสฺมา มหาโมคฺคลฺลาโน สพฺพาวนฺตํ ภิกฺขุสํฆํ เจตสา เจโต ปริจฺจ มนสากาสิ. อทฺทสา โข อายสฺมา มหาโมคฺคลฺลาโน ตํ ปุคฺคลํ ทุสฺสีลํ ปาปธมฺมํ อสุจิ\-สงฺกสฺสร\-สมาจารํ ปฏิจฺฉนฺน\-กมฺมนฺตํ อสฺสมณํ สมณปฏิญฺญํ อพฺรหฺมจาริํ พฺรหฺมจาริ\-ปฏิญฺญํ อนฺโตปูติํ อวสฺสุตํ กสมฺพุชาตํ\csromansharedfootnote{8e986869aa2f8cdd}{กสมฺพุกชาตํ (ก)} มชฺเฌ ภิกฺขุ\-สํฆสฺส นิสินฺนํ, ทิสฺวาน เยน โส \csromanfolio{419}ปุคฺคโล เตนุปสงฺกมิ, อุปสงฺกมิตฺวา ตํ ปุคฺคลํ เอตทโวจ “อุฏฺเฐหิ อาวุโส, ทิฏฺโฐสิ ภควตา, นตฺถิ เต ภิกฺขูหิ สทฺธิํ สํวาโส”ติ. เอวํ วุตฺเต โส ปุคฺคโล ตุณฺหี อโหสิ. ทุติยมฺปิ โข อายสฺมา มหาโมคฺคลฺลาโน ตํ ปุคฺคลํ เอตทโวจ “อุฏฺเฐหิ อาวุโส, ทิฏฺโฐสิ ภควตา, นตฺถิ เต ภิกฺขูหิ สทฺธิํ สํวาโส”ติ. ทุติยมฺปิ โข โส ปุคฺคโล ตุณฺหี อโหสิ. ตติยมฺปิ โข อายสฺมา มหาโมคฺคลฺลาโน ตํ ปุคฺคลํ เอตทโวจ “อุฏฺเฐหิ อาวุโส, ทิฏฺโฐสิ ภควตา, นตฺถิ เต ภิกฺขูหิ สทฺธิํ สํวาโส”ติ. ตติยมฺปิ โข โส ปุคฺคโล ตุณฺหี อโหสิ. อถ โข อายสฺมา มหาโมคฺคลฺลาโน ตํ ปุคฺคลํ พาหายํ คเหตฺวา พหิทฺวาร\-โกฏฺฐกา นิกฺขาเมตฺวา สูจิฆฏิกํ ทตฺวา เยน ภควา เตนุปสงฺกมิ, อุปสงฺกมิตฺวา ภควนฺตํ เอตทโวจ “นิกฺขามิโต โส ภนฺเต ปุคฺคโล มยา, สุทฺธา ปริสา อุทฺทิสตุ ภนฺเต ภควา ภิกฺขูนํ ปาติโมกฺขนฺ”ติ.}

\prose{อจฺฉริยํ โมคฺคลฺลาน, อพฺภุตํ โมคฺคลฺลาน, ยาว พาหาคหณาปิ นาม โส โมฆปุริโส อาคเมสฺสตีติ. อถ โข ภควา ภิกฺขู อามนฺเตสิ\csromandash{}}

\csromanmatikaanchor{mtk.194}
\csromantocmarkref{section}{2. มหาสมุทฺเท อฏฺฐจฺฉริย}{mtk.194}
\bookmark[dest={mtk.194},level=1]{2. มหาสมุทฺเท อฏฺฐจฺฉริย}
\csromanheader{1}{2. มหาสมุทฺเทอฏฺฐจฺฉริย}
\proseitem{384}{\csromansymbolfootnote{*}{ขุ 1. 139; อํ 3. 39 ปิฏฺเฐสุปิ.}อฏฺฐิเม ภิกฺขเว มหาสมุทฺเท อจฺฉริยา อพฺภุตา ธมฺมา, เย ทิสฺวา ทิสฺวา อสุรา มหาสมุทฺเท อภิรมนฺติ. กตเม อฏฺฐ.}

\prose{มหาสมุทฺโท ภิกฺขเว อนุปุพฺพนินฺโน อนุปุพฺพโปโณ อนุปุพฺพ\-ปพฺภาโร น อายตเกเนว ปปาโต. ยมฺปิ ภิกฺขเว มหาสมุทฺโท อนุปุพฺพนินฺโน อนุปุพฺพโปโณ อนุปุพฺพ\-ปพฺภาโร น อายตเกเนว ปปาโต, อยํ ภิกฺขเว มหาสมุทฺเท ปฐโม อจฺฉริโย อพฺภุโต ธมฺโม, ยํ ทิสฺวา ทิสฺวา อสุรา มหาสมุทฺเท อภิรมนฺติ.}

\prose{ปุน จปรํ ภิกฺขเว มหาสมุทฺโท ฐิตธมฺโม เวลํ นาติวตฺตติ, ยมฺปิ ภิกฺขเว มหาสมุทฺโท ฐิตธมฺโม เวลํ นาติวตฺตติ. อยํ\csromansharedfootnote{7543d4bec678b728}{อยมฺปิ (สฺยา)} ภิกฺขเว มหาสมุทฺเท ทุติโย อจฺฉริโย อพฺภุโต ธมฺโม, ยํ ทิสฺวา ทิสฺวา อสุรา มหาสมุทฺเท อภิรมนฺติ.}

\prose{\csromanfolio{420}ปุน จปรํ ภิกฺขเว มหาสมุทฺโท น มเตน กุณเปน สํวสติ. ยํ โหติ มหาสมุทฺเท มตํ กุณปํ, ตํ ขิปฺปญฺเญว ตีรํ วาเหติ, ถลํ อุสฺสาเรติ. ยมฺปิ ภิกฺขเว มหาสมุทฺโท น มเตน กุณเปน สํวสติ. ยํ โหติ มหาสมุทฺเท มตํ กุณปํ, ตํ ขิปฺปญฺเญว ตีรํ วาเหติ, ถลํ อุสฺสาเรติ. อยํ ภิกฺขเว มหาสมุทฺเท ตติโย อจฺฉริโย อพฺภุโต ธมฺโม, ยํ ทิสฺวา ทิสฺวา อสุรา มหาสมุทฺเท อภิรมนฺติ.}

\prose{ปุน จปรํ ภิกฺขเว ยา กาจิ มหานทิโย, เสยฺยถิทํ, คงฺคา ยมุนา อจิรวตี สรภู มหี, ตา มหาสมุทฺทํ ปตฺตา ชหนฺติ ปุริมานิ นามโคตฺตานิ, “มหาสมุทฺโท”เตฺวว สงฺขํ คจฺฉนฺติ. ยมฺปิ ภิกฺขเว ยากาจิ มหานทิโย, เสยฺยถิทํ, คงฺคา ยมุนา อจิรวตี สรภู มหี, ตา มหาสมุทฺทํ ปตฺตา ชหนฺติ ปุริมานิ นามโคตฺตานิ, “มหาสมุทฺโท”เตฺวว สงฺขํ คจฺฉนฺติ, อยํ ภิกฺขเว มหาสมุทฺเท จตุตฺโถ อจฺฉริโย อพฺภุโต ธมฺโม, ยํ ทิสฺวา ทิสฺวา อสุรา มหาสมุทฺเท อภิรมนฺติ.}

\prose{ปุน จปรํ ภิกฺขเว ยา จ\csromansharedfootnote{8627fccd20806815}{ยากาจิ (สฺยา)} โลเก สวนฺติโย มหาสมุทฺทํ อปฺเปนฺติ, ยา จ อนฺตลิกฺขา ธารา ปปตนฺติ, น เตน มหาสมุทฺทสฺส อูนตฺตํ วา ปูรตฺตํ วา ปญฺญยติ. ยมฺปิ ภิกฺขเว ยา จ โลเก สวนฺติโย มหาสมุทฺทํ อปฺเปนฺติ, ยา จ อนฺตลิกฺขา ธารา ปปตนฺติ, น เตน มหาสมุทฺทสฺส อูนตฺตํ วา ปูรตฺตํ วา ปญฺญายติ, อยํ ภิกฺขเว มหาสมุทฺเท ปญฺจโม อจฺฉริโย อพฺภุโต ธมฺโม, ยํ ทิสฺวา ทิสฺวา อสุรา มหาสมุทฺเท อภิรมนฺติ.}

\prose{ปุน จปรํ ภิกฺขเว มหาสมุทฺโท เอกรโส โลณรโส. ยมฺปิ ภิกฺขเว มหาสมุทฺโท เอกรโส โลณรโส, อยํ ภิกฺขเว มหาสมุทฺเท ฉฏฺโฐ อจฺฉริโย อพฺภุโต ธมฺโม, ยํ ทิสฺวา ทิสฺวา อสุรา มหาสมุทฺเท อภิรมนฺติ.}

\prose{ปุน จปรํ ภิกฺขเว มหาสมุทฺโท พหุรตโน\csromansharedfootnote{768df3bdf564dd17}{ปหูตรตโน (ก)} อเนกรตโน, ตตฺริมานิ รตนานิ, เสยฺยถิทํ, มุตฺตา มณิ เวฬุริโย สงฺโข สิลา ปวาฬํ รชตํ ชาตรูปํ โลหิตโก มสารคลฺลํ. ยมฺปิ ภิกฺขเว มหาสมุทฺโท พหุรตโน อเนกรตโน, ตตฺริมานิ รตนานิ เสยฺยถิทํ, มุตฺตา มณิ เวฬุริโย สงฺโข สิลา ปวาฬํ รชตํ ชตรูปํ \csromanfolio{421}โลหิตโก มสารคลฺลํ, อยํ ภิกฺขเว มหาสมุทฺเท สตฺตโม อจฺฉริโย อพฺภุโต ธมฺโม, ยํ ทิสฺวา ทิสฺวา อสุรา มหาสมุทฺเท อภิรมนฺติ.}

\prose{ปุน จปรํ ภิกฺขเว มหาสมุทฺโท มหตํ ภูตานํ อาวาโส, ตตฺริเม ภูตา, ติมิ ติมิงฺคโล ติมิติมิงฺคโล\csromansharedfootnote{4424468a9aeb4e26}{ติมิรปิงฺคโล (สี), ติมิติมิงฺคโล มหาติมิงฺคโล (สฺยา, กํ)} อสุรา นาคา คนฺธพฺพา, สนฺติ มหาสมุทฺเท โยชนสติกาปิ อตฺตภาวา, ทฺวิโยชน\-สติกาปิ อตฺตภาวา, ติโยชน\-สติกาปิ อตฺตภาวา, จตุโยชน\-สติกาปิ อตฺตภาวา, ปญฺจ\-โยชน\-สติกาปิ อตฺตภาวา. ยมฺปิ ภิกฺขเว มหาสมุทฺโท มหตํ ภูตานํ อาวาโส, ตตฺริเม ภูตา, ติมิ ติมิงฺคโล ติมิติมิงฺคโล อสุรา นาคา คนฺธพฺพา, สนฺติ มหาสมุทฺเท โยชนสติกาปิ อตฺตภาวา ฯเปฯ ปญฺจ\-โยชน\-สติกาปิ อตฺตภาวา, อยํ ภิกฺขเว มหาสมุทฺเท อฏฺฐโม อจฺฉริโย อพฺภุโต ธมฺโม, ยํ ทิสฺวา ทิสฺวา อสุรา มหาสมุทฺเท อภิรมนฺติ. อิเม โข ภิกฺขเว มหาสมุทฺเท อฏฺฐ อจฺฉริยา อพฺภุตา ธมฺมา, เย ทิสฺวา ทิสฺวา อสุรา มหาสมุทฺเท อภิรมนฺติ.}

\csromanmatikaanchor{mtk.195}
\csromantocmarkref{section}{3. อิมสฺมิํ ธมฺมวินเย อฏฺฐจฺฉริย}{mtk.195}
\bookmark[dest={mtk.195},level=1]{3. อิมสฺมิํ ธมฺมวินเย อฏฺฐจฺฉริย}
\csromanheader{1}{3. \textbf{อิมสฺมิํ ธมฺมวินเย อฏฺฐจฺฉริย}}
\proseitem{385}{เอวเมว โข ภิกฺขเว อิมสฺมิํ ธมฺมวินเย อฏฺฐ อจฺฉริยา อพฺภุตา ธมฺมา, เย ทิสฺวา ทิสฺวา ภิกฺขู อิมสฺมิํ ธมฺมวินเย อภิรมนฺติ. กตเม อฏฺฐ.}

\prose{เสยฺยถาปิ ภิกฺขเว มหาสมุทฺโท อนุปุพฺพนินฺโน อนุปุพฺพโปโณ อนุปุพฺพ\-ปพฺภาโร น อายตเกเนว ปปาโต, เอวเมว โข อภิกฺขเว อิมสฺมิํ ธมฺมวินเย อนุปุพฺพสิกฺขา อนุปุพฺพกิริยา อนุปุพฺพ\-ปฏิปทา น อายตเกเนว อญฺญาปฏิเวโธ. ยมฺปิ ภิกฺขเว อิมสฺมิํ ธมฺมวินเย อนุปุพฺพสิกฺขา อนุปุพฺพกิริยา อนุปุพฺพ\-ปฏิปทา น อายตเกเนว อญฺญาปฏิเวโธ, อยํ ภิกฺขเว อิมสฺมิํ ธมฺมวินเย ปฐโม อจฺฉริโย อพฺภุโต ธมฺโม, ยํ ทิสฺวา ทิสฺวา ภิกฺขู อิมสฺมิํ ธมฺมวินเย อภิรมนฺติ.}

\prose{เสยฺยถาปิ ภิกฺขเว มหาสมุทฺโท ฐิตธมฺโม เวลํ นาติวตฺตติ เอวเมว โข ภิกฺขเว ยํ มยา มม สาวกานํ สิกฺขาปทํ ปญฺญตฺตํ, ตํ มม สาวกา ชีวิตเหตุปิ นาติกฺกมนฺติ. ยมฺปิ ภิกฺขเว มยา มม สาวกานํ \csromanfolio{422}สิกฺขาปทํ ปญฺญตฺตํ, ตํ มม สาวกา ชีวิตเหตุปิ นาติกฺกมนฺติ, อยํ ภิกฺขเว อิมสฺมิํ ธมฺมวินเย ทุติโย อจฺฉริโย อพฺภุโต ธมฺโม, ยํ ทิสฺวา ทิสฺวา ภิกฺขู อิมสฺมิํ ธมฺมวินเย อภิรมนฺติ.}

\prose{เสยฺยถาปิ ภิกฺขเว มหาสมุทฺโท น มเตน กุณเปน สํวสติ, ยํ โหติ มหาสมุทฺเท มตํ กุณปํ, ตํ ขิปฺปเมว ตีรํ วาเหติ, ถลํ อุสฺสาเรติ, เอวเมว โข ภิกฺขเว โย โส ปุคฺคโล ทุสฺสีโล ปาปธมฺโม อสุจิ\-สงฺกสฺสร\-สมาจาโร ปฏิจฺฉนฺน\-กมฺมนฺโต อสฺสมโณ สมณปฏิญฺโญ อพฺรหฺมจารี พฺรหฺมจาริ\-ปฏิญฺโญ อนฺโตปูติ อวสฺสุโต กสมฺพุชาโต, น เตน สํโฆ สํวสติ, ขิปฺปเมว นํ สนฺนิปติตฺวา อุกฺขิปติ, กิญฺจาปิ โข โส โหติ มชฺเฌ ภิกฺขุ\-สํฆสฺส นิสินฺโน, อถ โข โส อารกาว สํฆมฺหา สํโฆ จ เตน. ยมฺปิ ภิกฺขเว โย โส ปุคฺคโล ทุสฺสีโล ปาปธมฺโม อสุจิ\-สงฺกสฺสร\-สมาจาโร ปฏิจฺฉนฺน\-กมฺมนฺโต อสฺสมโณ สมณปฏิญฺโญ อพฺรหฺมจารี พฺรหฺมจาริ\-ปฏิญฺโญ อนฺโตปูติ อวสฺสุโต กสมฺพุชาโต, น เตน สํโฆ สํวสติ, ขิปฺปเมว นํ สนฺนิปติตฺวา อุกฺขิปติ, กิญฺจาปิ โข โส โหติ มชฺเฌ ภิกฺขุ\-สํฆสฺส นิสินฺโน, อถ โข โส อารกาว สํฆมฺหา สํโฆ จ เตน, อยํ ภิกฺขเว อิมสฺมิํ ธมฺมวินเย ตติโย อจฺฉริโย อพฺภุโต ธมฺโม, ยํ ทิสฺวา ทิสฺวา ภิกฺขู อิมสฺมิํ ธมฺมวินเย อภิรมนฺติ.}

\prose{เสยฺยถาปิ ภิกฺขเว ยา กาจิ มหานทิโย, เสยฺยถิทํ, คงฺคา ยมุนา อจิรวตี สรภู มหี, ตา มหาสมุทฺทํ ปตฺตา ชหนฺติ ปุริมานิ นามโคตฺตานิ, “มหาสมุทฺโท”เตฺวว สงฺขํ คจฺฉนฺติ, เอวเมว โข ภิกฺขเว จตฺตาโรเม วณฺณา ขตฺติยา พฺราหฺมณา เวสฺสา สุทฺทา, เต ตถาคต\-ปฺปเวทิเต ธมฺมวินเย อคารสฺมา อนคาริยํ ปพฺพชิตฺวา\csromansharedfootnote{05082cd777dad9f9}{ปพฺพชิตา (สี)} ชหนฺติ ปุริมานิ นามโคตฺตานิ, “สมณา สกฺยปุตฺติยา”เตฺวว สงฺขํ คจฺฉนฺติ. ยมฺปิ ภิกฺขเว จตฺตาโรเม วณฺณา ขตฺติยา พฺราหฺมณา เวสฺสา สุทฺทา, เต ตถาคต\-ปฺปเวทิเต ธมฺมวินเย อคารสฺมา อนคาริยํ ปพฺพชิตฺวา ชหนฺติ ปุริมานิ นามโคตฺตานิ, “สมณา สกฺยปุตฺติยา”เตฺวว สงฺขํ คจฺฉนฺติ, อยํ ภิกฺขเว อิมสฺมิํ ธมฺมวินเย จตุตฺโถ อจฺฉริโย อพฺภุโต ธมฺโม, ยํ ทิสฺวา ทิสฺวา ภิกฺขู อิมสฺมิํ ธมฺมวินเย อภิรมนฺติ.}

\prose{\csromanfolio{423}เสยฺยถาปิ ภิกฺขเว ยา จ โลเก สวนฺติโย มหาสมุทฺทํ อปฺเปนฺติ, ยา จ อนฺตลิกฺขา ธารา ปปตนฺติ, น เตน มหาสมุทฺทสฺส อูนตฺตํ วา ปูรตฺตํ วา ปญฺญายติ, เอวเมว โข ภิกฺขเว พหู เจปิ ภิกฺขู อนุปาทิเสสาย นิพฺพานธาตุยา ปรินิพฺพายนฺติ, น เตน นิพฺพานธาตุยา อูนตฺตํ วา ปูรตฺตํ วา ปญฺญายติ. ยมฺปิ ภิกฺขเว พหู เจปิ ภิกฺขู อนุปาทิเสสาย นิพฺพานธาตุยา ปรินิพฺพายนฺติ, น เตน นิพฺพานธาตุยา อูนตฺตํ วา ปูรตฺตํ วา ปญฺญายติ, อยํ ภิกฺขเว อิมสฺมิํ ธมฺมวินเย ปญฺจโม อจฺฉริโย อพฺภุโต ธมฺโม, ยํ ทิสฺวา ทิสฺวา ภิกฺขู อิมสฺมิํ ธมฺมวินเย อภิรมนฺติ.}

\prose{เสยฺยถาปิ ภิกฺขเว มหาสมุทฺโท เอกรโส โลณรโส, เอวเมว โข ภิกฺขเว อยํ ธมฺมวินโย เอกรโส วิมุตฺติรโส. ยมฺปิ ภิกฺขเว อยํ ธมฺมวินโย เอกรโส วิมุตฺติรโส, อยํ ภิกฺขเว อิมสฺมิํ ธมฺมวินเย ฉฏฺโฐ อจฺฉริโย อพฺภุโต ธมฺโม, ยํ ทิสฺวา ทิสฺวา ภิกฺขู อิมสฺมิํ ธมฺมวินเย อภิรมนฺติ.}

\prose{เสยฺยถาปิ ภิกฺขเว มหาสมุทฺโท พหุรตโน อเนกรตโน, ตตฺริมานิ รตนานิ, เสยฺยถิทํ, มุตฺตา มณิ เวฬุริโย สงฺโข สิลา ปวาฬํ รชตํ ชาตรูปํ โลหิตโก มสารคลฺลํ, เอวเมว โข ภิกฺขเว อยํ ธมฺมวินโย พหุรตโน อเนกรตโน, ตตฺริมานิ รตนานิ, เสยฺยถิทํ, จตฺตาโร สติปฏฺฐานา จตฺตาโร สมฺมปฺปธานา จตฺตาโร อิทฺธิปาทา ปญฺจินฺทฺริยานิ ปญฺจ พลานิ สตฺต โพชฺฌงฺคา อริโย อฏฺฐงฺคิโก มคฺโค. ยมฺปิ ภิกฺขเว อยํ ธมฺมวินโย พหุรตโน อเนกรตโน, ตตฺริมานิ รตนานิ, เสยฺยถิทํ, จตฺตาโร สติปฏฺฐานา ฯเปฯ อริโย อฏฺฐงฺคิโก มคฺโค, อยํ ภิกฺขเว อิมสฺมิํ ธมฺมวินเย สตฺตโม อจฺฉริโย อพฺภุโต ธมฺโม, ยํ ทิสฺวา ทิสฺวา ภิกฺขู อิมสฺมิํ ธมฺมวินเย อภิรมนฺติ.}

\prose{เสยฺยถาปิ ภิกฺขเว มหาสมุทฺโท มหตํ ภูตานํ อาวาโส, ตตฺริเม ภูตา, ติมิ ติมิงฺคโล ติมิติมิงฺคโล อสุรา นาคา คนฺธพฺพา, สนฺติ มหาสมุทฺเท โยชนสติกาปิ อตฺตภาวา, ทฺวิโยชน\-สติกาปิ อตฺตภาวา, ติโยชน\-สติกาปิ อตฺตภาวา, จตุโยชน\-สติกาปิ อตฺตภาวา, ปญฺจ\-โยชน\-สติกาปิ อตฺตภาวา, เอวเมว โข ภิกฺขเว อยํ ธมฺมวินโย มหตํ ภูตานํ อาวาโส, ตตฺริเม ภูตา, โสตาปนฺโน, โสตาปตฺติผล\-สจฺฉิกิริยาย ปฏิปนฺโน, สกทาคามี, \csromanfolio{424}สกทาคามิ\-ผล\-สจฺฉิกิริยาย ปฏิปนฺโน, อนาคามี, อนาคามิ\-ผล\-สจฺฉิกิริยาย ปฏิปนฺโน อรหา, อรหตฺตผล\-สจฺฉิกิริยาย ปฏิปนฺโน. ยมฺปิ ภิกฺขเว อยํ ธมฺมวินโย มหตํ ภูตานํ อาวาโส, ตตฺริเม ภูตา, โสตาปนฺโน, โสตาปตฺติผล\-สจฺฉิกิริยาย ปฏิปนฺโน ฯเปฯ อรหา, อรหตฺตผล\-สจฺฉิกิริยาย ปฏิปนฺโน, อยํ ภิกฺขเว อิมสฺมิํ ธมฺมวินเย อฏฺฐโม อจฺฉริโย อพฺภุโต ธมฺโม, ยํ ทิสฺวา ทิสฺวา ภิกฺขู อิมสฺมิํ ธมฺมวินเย อภิรมนฺติ. อิเม โข ภิกฺขเว อิมสฺมิํ ธมฺมวินเย อฏฺฐ อจฺฉริยา อพฺภุตา ธมฺมา, เย ทิสฺวา ทิสฺวา ภิกฺขู อิมสฺมิํ ธมฺมวินเย อภิรมนฺตีติ.}

\prose{อถ โข ภควา เอตมตฺถํ วิทิตฺวา ตายํ เวลายํ อิมํ อุทานํ อุทาเนสิ\csromandash{}}

\setlength{\csromangathapagewidth}{0pt}
\setlength{\csromangathawakwidth}{0pt}
\csromangathameasuregroup{\textsuperscript{*}“ฉนฺนมติวสฺสติ\textsuperscript{1}, \\ ตสฺมา ฉนฺนํ วิวเรถ,}{\csromangathabat{\textsuperscript{*}“ฉนฺนมติวสฺสติ\textsuperscript{1},}{วิวฏํ นาติวสฺสติ.} \\ \csromangathabat{ตสฺมา ฉนฺนํ วิวเรถ,}{เอวํ ตํ นาติวสฺสตี”ติ.}}
\csromangathasetleft{\textsuperscript{*}“ฉนฺนมติวสฺสติ\textsuperscript{1}, \\ ตสฺมา ฉนฺนํ วิวเรถ,}
\csromangathagroup{\csromangathastanza{\csromangathabat{\csromansymbolfootnote{*}{วิ 1. 44 ปิฏฺเฐ อุทาเนปิ.}“ฉนฺนม\-ติวสฺสติ\csromansharedfootnote{8874de77f6161e98}{สุฉนฺทมติวสฺสติ (ก)},}{วิวฏํ นาติวสฺสติ.} \\ \csromangathabat{ตสฺมา ฉนฺนํ วิวเรถ,}{เอวํ ตํ นาติวสฺสตี”ติ.}}}

\csromanmatikaanchor{mtk.196}
\csromantocmarkref{section}{4. ปาติโมกฺขสวนารห}{mtk.196}
\bookmark[dest={mtk.196},level=1]{4. ปาติโมกฺขสวนารห}
\csromanheader{1}{4. ปาติโมกฺข\-สวนารห}
\proseitem{386}{อถ โข ภควา ภิกฺขู อามนฺเตสิ\csromandash{}นทานาหํ ภิกฺขเว อิโต ปรํ อุโปสถํ กริสฺสามิ, ปาติโมกฺขํ อุทฺทิสิสฺสามิ, ตุเมฺหวทานิ ภิกฺขเว อิโต ปรํ อุโปสถํ กเรยฺยาถ, ปาติโมกฺขํ อุทฺทิเสยฺยาถ. อฏฺฐานเมตํ ภิกฺขเว อนวกาโส, ยํ ตถาคโต อปริสุทฺธาย ปริสาย อุโปสถํ กเรยฺย, ปาติโมกฺขํ อุทฺทิเสยฺย. น จ ภิกฺขเว\csromansharedfootnote{8eb6dff44e3880eb}{น จ ภิกฺขเว ภิกฺขุนา (สฺยา, กํ)} สาปตฺติเกน ปาติโมกฺขํ โสตพฺพํ, โย สุเณยฺย, อาปตฺติ ทุกฺกฏสฺส. อนุชานามิ ภิกฺขเว โย สาปตฺติโก ปาติโมกฺขํ สุณาติ, ตสฺส ปาติโมกฺขํ ฐเปตุํ. เอวญฺจ ปน ภิกฺขเว ฐเปตพฺพํ, ตทหุโปสเถ จาตุทฺทเส วา ปนฺนรเส วา ตสฺมิํ ปุคฺคเล สมฺมุขีภูเต สํฆมชฺเฌ อุทาหริตพฺพํ\csromandash{}}

\prose{“สุณาตุ เม ภนฺเต สํโฆ, อิตฺถนฺนาโม ปุคฺคโล สาปตฺถิโก, ตสฺส ปาติโมกฺขํ ฐเปมิ ‘น ตสฺมิํ สมฺมุขีภูเต ปาติโมกฺขํ อุทฺทิสิตพฺพนฺ’ติ”. ฐปิตํ โหติ ปาติโมกฺขนฺติ.}

\prose{\csromanfolio{425}เตน โข ปน สมเยน ฉพฺพคฺคิยา ภิกฺขู “นาเมฺห โกจิ ชานาตี”ติ สาปตฺติกาว ปาติโมกฺขํ สุณนฺติ, เถรา ภิกฺขู ปรจิตฺต\-วิทุโน ภิกฺขูนํ อาโรเจนฺติ “อิตฺถนฺนาโม จ อิตฺถนฺนาโม จ อาวุโส ฉพฺพคฺคิยา ภิกฺขู ‘นาเมฺห โกจิ ชานาตี’ติ สาปตฺติกาว ปาติโมกฺขํ สุณนฺตี”ติ. อสฺโสสุํ โข ฉพฺพคฺคิยา ภิกฺขู เถรา กิร ภิกฺขู ปรจิตฺต\-วิทุโน อเมฺห ภิกฺขูนํ อาโรเจนฺติ “อิตฺถนฺนาโม จ อิตฺถนฺนาโม จ อาวุโส ฉพฺพคฺคิยา ภิกฺขู ‘นาเมฺห โกจิ ชานาตี’ติ สาปตฺติกาว ปาติโมกฺขํ สุณนฺตี”ติ.}

\prose{เต “ปุรมฺหากํ เปสลา ภิกฺขู ปาติโมกฺขํ ฐเปนฺตี”ติ ปฏิกจฺเจว สุทฺธานํ ภิกฺขูนํ อนาปตฺติกานํ อวตฺถุสฺมิํ อการเณ ปาติโมกฺขํ ฐเปนฺติ. เย เต ภิกฺขู อปฺปิจฺฉา ฯเปฯ เต อุชฺฌายนฺติ ขิยฺยนฺติ วิปาเจนฺติ “กถํ หิ นาม ฉพฺพคฺคิยา ภิกฺขู สุทฺธานํ ภิกฺขูนํ อนาปตฺติกานํ อวตฺถุสฺมิํ อการเณ ปาติโมกฺขํ ฐเปสฺสนฺตี”ติ. อถ โข เต ภิกฺขู ภควโต เอตมตฺถํ อาโรเจสุํ ฯเปฯ. สจฺจํ กิร ภิกฺขเว ฉพฺพคฺคิยา ภิกฺขู สุทฺธานํ ภิกฺขูนํ อนาปตฺติกานํ อวตฺถุสฺมิํ อการเณ ปาติโมกฺขํ ฐเปนฺตีติ. สจฺจํ ภควาติ ฯเปฯ วิครหิตฺวา ฯเปฯ ธมฺมิํ กถํ กตฺวา ภิกฺขู อามนฺเตสิ\csromandash{}น ภิกฺขเว สุทฺธานํ ภิกฺขูนํ อนาปตฺติกานํ อวตฺถุสฺมิํ อการเณ ปาติโมกฺขํ ฐเปตพฺพํ, โย ฐเปยฺย, อาปตฺติ ทุกฺกฏสฺส.}

\csromanmatikaanchor{mtk.197}
\csromantocmarkref{section}{5. ธมฺมิกาธมฺมิกปาติโมกฺขฏฺฐปน}{mtk.197}
\bookmark[dest={mtk.197},level=1]{5. ธมฺมิกาธมฺมิกปาติโมกฺขฏฺฐปน}
\csromanheader{1}{5. ธมฺมิกาธมฺมิก\-ปาติโมกฺข\-ฏฺฐปน}
\proseitem{387}{เอกํ ภิกฺขเว อธมฺมิกํ ปาติโมกฺข\-ฏฺฐปนํ, เอกํ ธมฺมิกํ ปาติโมกฺข\-ฏฺฐปนํ\csromansharedfootnote{50f0d3e8c60bac19}{เอกํ ธมฺมิกํ (สี, สฺยา)}. เทฺว อธมฺมิกานิ ปาติโมกฺข\-ฏฺฐปนานิ, เทฺว ธมฺมิกานิ. ตีณิ อธมฺมิกานิ ปาติโมกฺข\-ฏฺฐปนานิ, ตีณิ ธมฺมิกานิ. จตฺตาริ อธมฺมิกานิ ปาติโมกฺข\-ฏฺฐปนานิ, จตฺตาริ ธมฺมิกานิ. ปญฺจ อธมฺมิกานิ ปาติโมกฺข\-ฏฺฐปนานิ, ปญฺจ ธมฺมิกานิ. ฉ อธมฺมิกานิ ปาติโมกฺข\-ฏฺฐปนานิ, ฉ ธมฺมิกานิ. สตฺต อธมฺมิกานิ ปาติโมกฺข\-ฏฺฐปนานิ, สตฺต ธมฺมิกานิ. อฏฺฐ อธมฺมิกานิ ปาติโมกฺข\-ฏฺฐปนานิ, อฏฺฐ ธมฺมิกานิ. นว อธมฺมิกานิ ปาติโมกฺข\-ฏฺฐปนานิ, นว ธมฺมิกานิ. ทส อธมฺมิกานิ ปาติโมกฺข\-ฏฺฐปนานิ, ทส ธมฺมิกานิ.}

\prose{\csromanfolio{426}กตมํ เอกํ อธมฺมิกํ ปาติโมกฺข\-ฏฺฐปนํ. อมูลิกาย สีลวิปตฺติยา ปาติโมกฺขํ ฐเปติ. อิทํ เอกํ อธมฺมิกํ ปาติโมกฺข\-ฏฺฐปนํ. กตมํ เอกํ ธมฺมิกํ ปาติโมกฺข\-ฏฺฐปนํ. สมูลิกาย สีลวิปตฺติยา ปาติโมกฺขํ ฐเปติ. อิทํ เอกํ ธมฺมิกํ ปาติโมกฺข\-ฏฺฐปนํ.}

\prose{กตมานิ เทฺว อธมฺมิกานิ ปาติโมกฺข\-ฏฺฐปนานิ. อมูลิกาย สีลวิปตฺติยา ปาติโมกฺขํ ฐเปติ, อมูลิกาย อาจารวิปตฺติยา ปาติโมกฺขํ ฐเปติ. อิมานิ เทฺว อธมฺมิกานิ ปาติโมกฺข\-ฏฺฐปนานิ. กตมานิ เทฺว ธมฺมิกานิ ปาติโมกฺข\-ฏฺฐปนานิ. สมูลิกาย สีลวิปตฺติยา ปาติโมกฺขํ ฐเปติ, สมูลิกาย อาจารวิปตฺติยา ปาติโมกฺขํ ฐเปติ. อิมานิ เทฺว ธมฺมิกานิ ปาติโมกฺข\-ฏฺฐปนานิ.}

\prose{กตมานิ ตีณิ อธมฺมิกานิ ปาติโมกฺข\-ฏฺฐปนานิ. อมูลิกาย สีลวิปตฺติยา ปาติโมกฺขํ ฐเปติ, อมูลิกาย อาจารวิปตฺติยา ปาติโมกฺขํ ฐเปติ, อมูลิกาย ทิฏฺฐิ\-วิปตฺติยา ปาติโมกฺขํ ฐเปติ. อิมานิ ตีณิ อธมฺมิกานิ ปาติโมกฺข\-ฏฺฐปนานิ. กตมานิ ตีณิ ธมฺมิกานิ ปาติโมกฺข\-ฏฺฐปนานิ. สมูลิกาย สีลวิปตฺติยา ปาติโมกฺขํ ฐเปติ, สมูลิกาย อาจารวิปตฺติยา ปาติโมกฺขํ ฐเปติ, สมูลิกาย ทิฏฺฐิ\-วิปตฺติยา ปาติโมกฺขํ ฐเปติ. อิมานิ ตีณิ ธมฺมิกานิ ปาติโมกฺข\-ฏฺฐปนานิ.}

\prose{กตมานิ จตฺตาริ อธมฺมิกานิ ปาติโมกฺข\-ฏฺฐปนานิ. อมูลิกาย สีลวิปตฺติยา ปาติโมกฺขํ ฐเปติ, อมูลิกาย อาจารวิปตฺติยา ปาติโมกฺขํ ฐเปติ, อมูลิกาย ทิฏฺฐิ\-วิปตฺติยา ปาติโมกฺขํ ฐเปติ, อมูลิกาย อาชีววิปตฺติยา ปาติโมกฺขํ ฐเปติ. อิมานิ จตฺตาริ อธมฺมิกานิ ปาติโมกฺข\-ฏฺฐปนานิ. กตมานิ จตฺตาริ ธมฺมิกานิ ปาติโมกฺข\-ฏฺฐปนานิ. สมูลิกาย สีลวิปตฺติยา ปาติโมกฺขํ ฐเปติ, สมูลิกาย อาจารวิปตฺติยา ปาติโมกฺขํ ฐเปติ, สมูลิกาย ทิฏฺฐิ\-วิปตฺติยา ปาติโมกฺขํ ฐเปติ, สมูลิกาย อาชีววิปตฺติยา, ปาติโมกฺขํ ฐเปติ. อิมานิ จตฺตาริ ธมฺมิกานิ ปาติโมกฺข\-ฏฺฐปนานิ.}

\prose{กตมานิ ปญฺจ อธมฺมิกานิ ปาติโมกฺข\-ฏฺฐปนานิ. อมูลเกน ปาราชิเกน ปาติโมกฺขํ ฐเปติ, อมูลเกน สํฆาทิเสเสน ฯเปฯ อมูลเกน ปาจิตฺติเยน. อมูลเกน ปาฏิเทสนีเยน. อมูลเกน ทุกฺกเฏน \csromanfolio{427}ปาติโมกฺขํ ฐเปติ. อิมานิ ปญฺจ อธมฺมิกานิ ปาติโมกฺข\-ฏฺฐปนานิ. กตมานิ ปญฺจ ธมฺมิกานิ ปาติโมกฺข\-ฏฺฐปนานิ. สมูลเกน ปาราชิเกน ปาติโมกฺขํ ฐเปติ, สมูลเกน สํฆาทิเสเสน ฯเปฯ ปาจิตฺติเยน. ปาฏิเทสนีเยน. ทุกฺกเฏน ปาติโมกฺขํ ฐเปติ. อิมานิ ปญฺจ ธมฺมิกานิ ปาติโมกฺข\-ฏฺฐปนานิ.}

\prose{กตมานิ ฉ อธมฺมิกานิ ปาติโมกฺข\-ฏฺฐปนานิ. อมูลิกาย สีลวิปตฺติยา ปาติโมกฺขํ ฐเปติ อกตาย, อมูลิกาย สีลวิปตฺติยา ปาติโมกฺขํ ฐเปติ กตาย, อมูลิกาย อาจารวิปตฺติยา ปาติโมกฺขํ ฐเปติ อกตาย, อมูลิกาย อาจารวิปตฺติยา ปาติโมกฺขํ ฐเปติ กตาย, อมูลิกาย ทิฏฺฐิ\-วิปตฺติยา ปาติโมกฺขํ ฐเปติ อกตาย, อมูลิกาย ทิฏฺฐิ\-วิปตฺติยา ปาติโมกฺขํ ฐเปติ กตาย. อิมานิ ฉ อธมฺมิกานิ ปาติโมกฺข\-ฏฺฐปนานิ. กตมานิ ฉ ธมฺมิกานิ ปาติโมกฺข\-ฏฺฐปนานิ. สมูลิกาย สีลวิปตฺติยา ปาติโมกฺขํ ฐเปติ อกตาย, สมุลิกาย สีลวิปตฺติยา ปาติโมกฺขํ ฐเปติ กตาย, สมูลิกาย อาจารวิปตฺติยา ฯเปฯ ทิฏฺฐิ\-วิปตฺติยา ปาติโมกฺขํ ฐเปติ อกตาย ฯเปฯ กตาย. อิมานิ ฉ ธมฺมิกานิ ปาติโมกฺข\-ฏฺฐปนานิ.}

\prose{กตมานิ สตฺต อธมฺมิกานิ ปาติโมกฺข\-ฏฺฐปนานิ. อมูลเกน ปาราชิเกน ปาติโมกฺขํ ฐเปติ, อมูลเกน สํฆาทิเสเสน ฯเปฯ ถุลฺลจฺจเยน. ปาจิตฺติเยน. ปาฏิเทสนีเยน. ทุกฺกเฏน. ทุพฺภาสิเตน ปาติโมกฺขํ ฐเปติ. อิมานิ สตฺต อธมฺมิกานิ ปาติโมกฺข\-ฏฺฐปนานิ. กตมนิ สตฺต ธมฺมิกานิ ปาติโมกฺข\-ฏฺฐปนานิ. สมูลเกน ปาราชิเกน ปาติโมกฺขํ ฐเปติ, สมูลเกน สํฆาทิเสเสน ฯเปฯ ถุลฺลจฺจเยน. ปาจิตฺติเยน. ปาฏิเทสนีเยน. ทุกฺกเฏน. ทุพฺภาสิเตน ปาติโมกฺขํ ฐเปติ. อิมานิ สตฺต ธมฺมิกานิ ปาติโมกฺข\-ฏฺฐปนานิ.}

\prose{กตมานิ อฏฺฐ อธมฺมิกานิ ปาติโมกฺข\-ฏฺฐปนานิ. อมูลิกาย สีลวิปตฺติยา ปาติโมกฺขํ ฐเปติ อกตาย, อมูลิกาย สีลวิปตฺติยา ปาติโมกฺขํ ฐเปติ กตาย, อมูลิกาย อาจารวิปตฺติยา ฯเปฯ. ทิฏฺฐิ\-วิปตฺติยา. อาชีววิปตฺติยา ปาติโมกฺขํ ฐเปติ อกตาย ฯเปฯ กตาย. อิมานิ อฏฺฐ อธมฺมิกานิ ปาติโมกฺข\-ฏฺฐปนานิ. กตมานิ อฏฺฐ ธมฺมิกานิ ปาติโมกฺข\-ฏฺฐปนานิ. สมูลิกาย สีลวิปตฺติยา ปาติโมกฺขํ ฐเปติ อกตาย, สมูลิกาย สีลวิปตฺติยา ปาติโมกฺขํ \csromanfolio{428}ฐเปติ กตาย, สมูลิกาย อาจารวิปตฺติยา ฯเปฯ ทิฏฺฐิ\-วิปตฺติยา. อาชีววิปตฺติยา ปาติโมกฺขํ ฐเปติ อกตาย ฯเปฯ กตาย. อิมานิ อฏฺฐ ธมฺมิกานิ ปาติโมกฺข\-ฏฺฐปนานิ.}

\prose{กตมานิ นว อธมฺมิกานิ ปาติโมกฺข\-ฏฺฐปนานิ. อมูลิกาย สีลวิปตฺติยา ปาติโมกฺขํ ฐเปติ อกตาย, อมูลิกาย สีลวิปตฺติยา ปาติโมกฺขํ ฐเปติ กตาย, อมูลิกาย สีลวิปตฺติยา ปาติโมกฺขํ ฐเปติ กตากตาย, อมูลิกาย อาจารวิปตฺติยา ฯเปฯ ทิฏฺฐิ\-วิปตฺติยา ปาติโมกฺขํ ฐเปติ อกตาย ฯเปฯ กตาย ฯเปฯ กตากตาย. อิมานิ นว อธมฺมิกานิ ปาติโมกฺข\-ฏฺฐปนานิ. กตามานิ นว ธมฺมิกานิ ปาติโมกฺข\-ฏฺฐปนานิ. สมูลิกาย สีลวิปตฺติยา ปาติโมกฺขํ ฐเปติ อกตาย, สมูลิกาย สีลวิปตฺติยา ปาติโมกฺขํ ฐเปติ กตาย, สมูลิกาย, สีลวิปตฺติยา ปาติโมกฺขํ ฐเปติ กตากตาย, สมูลิกาย, อาจารวิปตฺติยา ฯเปฯ ทิฏฺฐิ\-วิปตฺติยา ปาติโมกฺขํ ฐเปติ อกตาย ฯเปฯ กตาย ฯเปฯ กตากตาย. อิมานิ นว ธมฺมิกานิ ปาติโมกฺข\-ฏฺฐปนานิ.}

\prose{กตมานิ ทส อธมฺมิกานิ ปาติโมกฺข\-ฏฺฐปนานิ. น ปาราชิโก ตสฺสํ ปริสายํ นิสินฺโน โหติ, น ปาราชิกกถา วิปฺปกตา โหติ, น สิกฺขํ ปจฺจกฺขาตโก ตสฺสํ ปริสายํ นิสินฺโน โหติ, น สิกฺขํ ปจฺจกฺขาต\-กถา วิปฺปกตา โหติ, ธมฺมิกํ สามคฺคิํ อุเปติ, น ธมฺมิกํ สามคฺคิํ ปจฺจาทิยติ, น ธมฺมิกาย สามคฺคิยา ปจฺจาทานกถา วิปฺปกตา โหติ, น สีลวิปตฺติยา ทิฏฺฐ\-สุต\-ปริสงฺกิโต โหติ, น อาจารวิปตฺติยา ทิฏฺฐ\-สุต\-ปริสงฺกิโต โหติ, น ทิฏฺฐิ\-วิปตฺติยา ทิฏฺฐ\-สุต\-ปริสงฺกิโต โหติ. อิมานิ ทส อธมฺมิกานิ ปาติโมกฺข\-ฏฺฐปนานิ. กตมานิ ทส ธมฺมิกานิ ปาติโมกฺข\-ฏฺฐปนานิ. ปาราชิโก ตสฺสํ ปริสายํ นิสินฺโน โหติ, ปาราชิกกถา วิปฺปกตา โหติ, สิกฺขํ ปจฺจกฺขาตโก ตสฺสํ ปริสายํ นิสินฺโน โหติ, สิกฺขํ ปจฺจกฺขาต\-กถา วิปฺปกตา โหติ, ธมฺมิกํ สามคฺคิํ น อุเปติ, ธมฺมิกํ สามคฺคิํ ปจฺจาทิยติ, ธมฺมิกาย สามคฺคิยา ปจฺจาทานกถา วิปฺปกตา โหติ, สีลวิปตฺติยา ทิฏฺฐ\-สุต\-ปริสงฺกิโต โหติ, อาจารวิปตฺติยา ทิฏฺฐ\-สุต\-ปริสงฺกิโต โหติ, ทิฏฺฐิ\-วิปตฺติยา ทิฏฺฐ\-สุต\-ปริสงฺกิโต โหติ. อิมานิ ทส ธมฺมิกานิ ปาติโมกฺข\-ฏฺฐปนานิ.}

\csromanmatikaanchor{mtk.198}
\csromantocmarkref{section}{6. ธมฺมิกปาติโมกฺขฏฺฐปน}{mtk.198}
\bookmark[dest={mtk.198},level=1]{6. ธมฺมิกปาติโมกฺขฏฺฐปน}
\csromanheader{1}{6. ธมฺมิก\-ปาติโมกฺข\-ฏฺฐปน}
\proseitem{388}{\csromanfolio{429}กถํ ปาราชิโก ตสฺสํ ปริสายํ นิสินฺโน โหติ. อิธ ภิกฺขเว เยหิ อากาเรหิ เยหิ ลิงฺเคหิ เยหิ นิมิตฺเตหิ ปาราชิกสฺส ธมฺมสฺส อชฺฌาปตฺติ โหติ, เตหิ อากาเรหิ เตหิ ลิงฺเคหิ เตหิ นิมิตฺเตหิ ภิกฺขู ภิกฺขุํ ปสฺสติ ปาราชิกํ ธมฺมํ อชฺฌาปชฺชนฺตํ. น เหว โข ภิกฺขุ ภิกฺขุํ ปสฺสติ ปาราชิกํ ธมฺมํ อชฺฌาปชฺชนฺตํ, อปิ จ อญฺโญ ภิกฺขุ ภิกฺขุสฺส อาโรเจติ “อิตฺถนฺนาโม อาวุโส ภิกฺขุ ปาราชิกํ ธมฺมํ อชฺฌาปนฺโน”ติ. น เหว โข ภิกฺขุ ภิกฺขุํ ปสฺสติ ปาราชิกํ ธมฺมํ อชฺฌาปชฺชนฺตํ, นาปิ\csromansharedfootnote{abcf37e7cdeb701b}{นาปิ จ (ก)} อญฺโญ ภิกฺขุ ภิกฺขุสฺส อาโรเจติ “อิตฺถนฺนาโม อาวุโส ภิกฺขุ ปาราชิกํ ธมฺมํ อชฺฌาปนฺโน”ติ, อปิ จ โสว ภิกฺขุ ภิกฺขุสฺส อาโรเจติ “อหํ อาวุโส ปาราชิกํ ธมฺมํ อชฺฌาปนฺโน”ติ. อากงฺขมาโน ภิกฺขเว ภิกฺขุ เตน ทิฏฺเฐน เตน สุเตน ตาย ปริสงฺกาย ตทหุโปสเถ จาตุทฺทเส วา ปนฺนรเส วา ตสฺมิํ ปุคฺคเล สมฺมุขีภูเต สํฆมชฺเฌ อุทาหเรยฺย\csromandash{}}

\prose{“สุณาตุ เม ภนฺเต สํโฆ, อิตฺถนฺนาโม ปุคฺคโล ปาราชิกํ ธมฺมํ อชฺฌาปนฺโน, ตสฺส ปาติโมกฺขํ ฐเปมิ ‘น ตสฺมิํ สมฺมุขีภูเต ปาติโมกฺขํ อุทฺทิสิตพฺพนฺ’ติ”. ธมฺมิกํ ปาติโมกฺข\-ฏฺฐปนํ.}

\proseitem{389}{ภิกฺขุสฺส ปาติโมกฺเข ฐปิเต ปริสา วุฏฺฐาติ ทสนฺนํ อนฺตรายานํ อญฺญตเรน\csromansharedfootnote{063f0c1fbba49687}{อญฺญตเรน อนฺตราเยน (สฺยา, กํ)} ราชนฺตราเยน วา โจรนฺตราเยน วา อคฺยนฺตราเยน วา อุทกนฺตราเยน วา มนุสฺสนฺตราเยน วา อมนุสฺสนฺตราเยน วา วาฬนฺตราเยน วา สรีสปนฺตราเยน วา ชีวิต\-นฺตราเยน วา พฺรหฺมจริย\-นฺตราเยน วา, อากงฺขมาโน ภิกฺขเว ภิกฺขุ ตสฺมิํ วา อาวาเส อญฺญสฺมิํ วา อาวาเส ตสฺมิํ ปุคฺคเล สมฺมุขีภูเต สํฆมชฺเฌ อุทาหเรยฺย\csromandash{}}

\prose{“สุณาตุ เม ภนฺเต สํโฆ, อิตฺถนฺนามสฺส ปุคฺคลสฺส ปาราชิกกถา วิปฺปกตา, ตํ วตฺถุ อวินิจฺฉิตํ, ยทิ สํฆสฺส ปตฺตกลฺลํ, สํโฆ ตํ วตฺถุํ วินิจฺฉิเนยฺยา”ติ.}

\prose{เอวญฺเจตํ ลเภถ, อิจฺเจตํ กุสลํ. โน เจ ลเภถ, ตทหุโปสเถ จาตุทฺทเส วา ปนฺนรเส วา ตสฺมิํ ปุคฺคเล สมฺมุขีภูเต สํฆมชฺเฌ อุทาหริตพฺพํ\csromandash{}}

\prose{\csromanfolio{430}“สุณาตุ เม ภนฺเต สํโฆ, อิตฺถนฺนามสฺส ปุคฺคลสฺส ปาราชิกกถา วิปฺปกตา, ตํ วตฺถุ อวินิจฺฉิตํ, ตสฺส ปาติโมกฺขํ ฐเปมิ ‘น ตสฺมิํ สมฺมุขีภูเต ปาติโมกฺขํ อุทฺทิสิตพฺพนฺ’ติ”. ธมฺมิกํ ปาติโมกฺข\-ฏฺฐปนํ.}

\proseitem{390}{กถํ สิกฺขํ ปจฺจกฺขาตโก ตสฺสํ ปริสายํ นิสินฺโน โหติ. อิธ ปน ภิกฺขเว ภิกฺขุ เยหิ อากาเรหิ เยหิ ลิงฺเคหิ เยหิ นิมิตฺเตหิ สิกฺขา ปจฺจกฺขาตา โหติ, เตหิ อากาเรหิ เตหิ ลิงฺเคหิ เตหิ นิมิตฺเตหิ ภิกฺขุ ภิกฺขุํ ปสฺสติ สิกฺขํ ปจฺจกฺขนฺตํ. น เหว โข ภิกฺขุ ภิกฺขุํ ปสฺสติ สิกฺขํ ปจฺจกฺขนฺตํ, อปิ จ อญฺโญ ภิกฺขุ ภิกฺขุสฺส อาโรเจติ “อิตฺถนฺนาเมน อาวุโส ภิกฺขุนา สิกฺขา ปจฺจกฺขาตา”ติ. น เหว โข ภิกฺขุ ภิกฺขุํ ปสฺสติ สิกฺขํ ปจฺจกฺขนฺตํ, นาปิ อญฺโญ ภิกฺขุ ภิกฺขุสฺส อาโรเจติ “อิตฺถนฺนาเมน อาวุโส ภิกฺขุนา สิกฺขา ปจฺจกฺขาตา”ติ. อปิ จ โสว ภิกฺขุ ภิกฺขุสฺส อาโรเจติ “มยา อาวุโส สิกฺขา ปจฺจกฺขาตา”ติ. อากงฺขมาโน ภิกฺขเว ภิกฺขุ เตน ทิฏฺเฐน เตน สุเตน ตาย ปริสงฺกาย ตทหุโปสเถ จาตุทฺทเส วา ปนฺนรเส วา ตสฺมิํ ปุคฺคเล สมฺมุขีภูเต สํฆมชฺเฌ อุทาหเรยฺย\csromandash{}}

\prose{“สุณาตุ เม ภนฺเต สํโฆ, อิตฺถนฺนาเมน ปุคฺคเลน สิกฺขา ปจฺจกฺขาตา, ตสฺส ปาติโมกฺขํ ฐเปมิ ‘น ตสฺมิํ สมฺมุขีภูเต ปาติโมกฺขํ อุทฺทิสิตพฺพนฺ’ติ”. ธมฺมิกํ ปาติโมกฺข\-ฏฺฐปนํ.}

\proseitem{391}{ภิกฺขุสฺส ปาติโมกฺเข ฐปิเต ปริสา วุฏฺฐาติ ทสนฺนํ อนฺตรายานํ อญฺญตเรน ราชนฺตราเยน วา ฯเปฯ พฺรหฺมจริย\-นฺตราเยน วา. อากงฺขมาโน ภิกฺขเว ภิกฺขุ ตสฺมิํ วา อาวาเส อญฺญสฺมิํ วา อาวาเส ตสฺมิํ ปุคฺคเล สมฺมุขีภูเต สํฆมชฺเฌ อุทาหเรยฺย\csromandash{}}

\prose{“สุณาตุ เม ภนฺเต สํโฆ, อิตฺถนฺนามสฺส ปุคฺคลสฺส สิกฺขํ ปจฺจกฺขาต\-กถา\csromansharedfootnote{abe5156492d6ace1}{สิกฺขาปจฺจกฺขาตกถา (สี)} วิปฺปกตา, ตํ วตฺถุ อวินิจฺฉิตํ, ยทิ สํฆสฺส ปตฺตกลฺลํ, สํโฆ ตํ วตฺถุํ วินิจฺฉิเนยฺยา”ติ.}

\prose{เอวญฺเจตํ ลเภถ, อิจฺเจตํ กุสลํ. โน เจ ลเภถ, ตทหุโปสเถ จาตุทฺทเส วา ปนฺนรเส วา ตสฺมิํ ปุคฺคเล สมฺมุขีภูเต สํฆมชฺเฌ อุทาหริตพฺพํ\csromandash{}}

\prose{\csromanfolio{431}“สุณาตุ เม ภนฺเต สํโฆ, อิตฺถนฺนามสฺส ปุคฺคลสฺส สิกฺขํ ปจฺจกฺขาต\-กถา วิปฺปกตา, ตํ วตฺถุ อวินิจฺฉิตํ, ตสฺส ปาติโมกฺขํ ฐเปมิ ‘น ตสฺมิํ สมฺมุขีภูเต ปาติโมกฺขํ อุทฺทิสิตพฺพนฺ’ติ”. ธมฺมิกํ ปาติโมกฺข\-ฏฺฐปนํ.}

\proseitem{392}{กถํ ธมฺมิกํ สามคฺคิํ น อุเปติ. อิธ ปน ภิกฺขเว ภิกฺขุ เยหิ อากาเรหิ เยหิ ลิงฺเคหิ เยหิ นิมิตฺเตหิ ธมฺมิกาย สามคฺคิยา\-นุปคมนํ โหติ, เตหิ อากาเรหิ เตหิ ลิงฺเคหิ เตหิ นิมิตฺเตหิ ภิกฺขุ ภิกฺขุํ ปสฺสติ ธมฺมิกํ สามคฺคิํ น อุเปนฺตํ. น เหว โข ภิกฺขุ ภิกฺขุํ ปสฺสติ ธมฺมิกํ สามคฺคิํ น อุเปนฺตํ, อปิ จ อญฺโญ ภิกฺขุ ภิกฺขุสฺส อาโรเจติ “อิตฺถนฺนาโม อาวุโส ภิกฺขุ ธมฺมิกํ สามคฺคิํ น อุเปตี”ติ. น เหว โข ภิกฺขุ ภิกฺขุํ ปสฺสติ ธมฺมิกํ สามคฺคิํ น อุเปนฺตํ, นาปิ อญฺโญ ภิกฺขุ ภิกฺขุสฺส อาโรเจติ “อิตฺถนฺนาโม อาวุโส ภิกฺขุ ธมฺมิกํ สามคฺคิํ น อุเปตี”ติ, อปิ จ โสว ภิกฺขุ ภิกฺขุสฺส อาโรเจติ “อหํ อาวุโส ธมฺมิกํ สามคฺคิํ น อุเปมี”ติ. อากงฺขมาโน ภิกฺขเว ภิกฺขุ เตน ทิฏฺเฐน เตน สุเตน ตาย ปริสงฺกาย ตทหุโปสเถ จาตุทฺทเส วา ปนฺนรเส วา ตสฺมิํ ปุคฺคเล สมฺมุขีภูเต สํฆมชฺเฌ อุทาหเรยฺย\csromandash{}}

\prose{“สุณาตุ เม ภนฺเต สํโฆ, อิตฺถนฺนาโม ปุคฺคโล ธมฺมิกํ สามคฺคิํ น อุเปติ, ตสฺส ปาติโมกฺขํ ฐเปมิ ‘น ตสฺมิํ สมฺมุขีภูเต ปาติโมกฺขํ อุทฺทิสิตพฺพนฺ’ติ”. ธมฺมิกํ ปาติโมกฺข\-ฏฺฐปนํ.}

\proseitem{393}{กถํ ธมฺมิกํ สามคฺคิํ ปจฺจาทิยติ. อิธ ภิกฺขเว ภิกฺขุ เยหิ อากาเรหิ เยหิ ลิงฺเคหิ เยหิ นิมิตฺเตหิ ธมฺมิกาย สามคฺคิยา ปจฺจาทานํ โหติ, เตหิ อากาเรหิ เตหิ ลิงฺเคหิ เตหิ นิมิตฺเตหิ ภิกฺขุ ภิกฺขุํ ปสฺสติ ธมฺมิกํ สามคฺคิํ ปจฺจาทิยนฺตํ. น เหว โข ภิกฺขุ ภิกฺขุํ ปสฺสติ ธมฺมิกํ สามคฺคิํ ปจฺจาทิยนฺตํ, อปิ จ อญฺโญ ภิกฺขุ ภิกฺขุสฺส อาโรเจติ “อิตฺถนฺนาโม อาวุโส ภิกฺขุ ธมฺมิกํ สามคฺคิํ ปจฺจาทิยตี”ติ. น เหว โข ภิกฺขุ ภิกฺขุํ ปสฺสติ ธมฺมิกํ สามคฺคิํ ปจฺจาทิยนฺตํ, นาปิ อญฺโญ ภิกฺขุ ภิกฺขุสฺส อาโรเจติ “อิตฺถนฺนาโม อาวุโส ภิกฺขุ ธมฺมิกํ สามคฺคิํ ปจฺจาทิยตี”ติ, อปิ จ โสว ภิกฺขุ ภิกฺขุสฺส อาโรเจติ “อหํ อาวุโส ธมฺมิกํ สามคฺคิํ \csromanfolio{432}ปจฺจาทิยามี”ติ. อากงฺขมาโน ภิกฺขเว ภิกฺขุ เตน ทิฏฺเฐน เตน สุเตน ตาย ปริสงฺกาย ตทหุโปสเถ จาตุทฺทเส วา ปนฺนรเส วา ตสฺมิํ ปุคฺคเล สมฺมุขีภูเต สํฆมชฺเฌ อุทาหเรยฺย\csromandash{}}

\prose{“สุณาตุ เม ภนฺเต สํโฆ, อิตฺถนฺนโม ปุคฺคโล ธมฺมิกํ สามคฺคิํ ปจฺจาทิยติ, ตสฺส ปาติโมกฺขํ ฐเปมิ ‘น ตสฺมิํ สมฺมุขีภูเต ปาติโมกฺขํ อุทฺทิสิตพฺพนฺ’ติ”. ธมฺมิกํ ปาติโมกฺข\-ฏฺฐปนํ.}

\proseitem{394}{ภิกฺขุสฺส ปาติโมกฺเข ฐปิเต ปริสา วุฏฺฐาติ ทสนฺนํ อนฺตรายานํ อญฺญตเรน ราชนฺตราเยน วา ฯเปฯ พฺรหฺมจริย\-นฺตราเยน วา. อากงฺขมาโน ภิกฺขเว ภิกฺขุ ตสฺมิํ วา อาวาเส อญฺญสฺมิํ วา อาวาเส ตสฺมิํ ปุคฺคเล สมฺมุขีภูเต สํฆมชฺเฌ อุทาหเรยฺย\csromandash{}}

\prose{“สุณาตุ เม ภนฺเต สํโฆ, อิตฺถนฺนามสฺส ปุคฺคลสฺส ธมฺมิกาย สามคฺคิยา ปจฺจาทานกถา วิปฺปกตา, ตํ วตฺถุ อวินิจฺฉิตํ, ยทิ สํฆสฺส ปตฺตกลฺลํ, สํโฆ ตํ วตฺถุํ วินิจฺฉิเนยฺยา”ติ.}

\prose{เอวญฺเจตํ ลเภถ, อิจฺเจตํ กุสลํ. โน เจ ลเภถ, ตทหุโปสเถ จาตุทฺทเส วา ปนฺนรเส วา ตสฺมิํ ปุคฺคเล สมฺมุขีภูเต สํฆมชฺเฌ อุทาหริตพฺพํ\csromandash{}}

\prose{“สุณาตุ เม ภนฺเต สํโฆ, อิตฺถนฺนามสฺส ปุคฺคลสฺส ธมฺมิกาย สามคฺคิยา ปจฺจาทานกถา วิปฺปกตา, ตํ วตฺถุ อวินิจฺฉิตํ, ตสฺส ปาติโมกฺขํ ฐเปมิ ‘น ตสฺมิํ สมฺมุขีภูเต ปาติโมกฺขํ อุทฺทิสิตพฺพนฺ’ติ”. ธมฺมิกํ ปาติโมกฺข\-ฏฺฐปนํ.}

\proseitem{395}{กตํ สีลวิปตฺติยา ทิฏฺฐ\-สุต\-ปริสงฺกิโต โหติ. อิธ ปน ภิกฺขเว ปน ภิกฺขุ เยหิ อากาเรหิ เยหิ ลิงฺเคหิ เยหิ นิมิตฺเตหิ สีลวิปตฺติยา ทิฏฺฐ\-สุต\-ปริสงฺกิโต โหติ, เตหิ อากาเรหิ เตหิ ลิงฺเคหิ เตหิ นิมิตฺเตหิ ภิกฺขุ ภิกฺขุํ ปสฺสติ สีลวิปตฺติยา ทิฏฺฐ\-สุต\-ปริสงฺกิตํ. น เหว โข ภิกฺขุ ภิกฺขุํ ปสฺสติ สีลวิปตฺติยา ทิฏฺฐ\-สุต\-ปริสงฺกิตํ, อปิจ อญฺโญ ภิกฺขุ ภิกฺขุสฺส อาโรเจติ “อิตฺถนฺนาโม อาวุโส ภิกฺขุ สีลวิปตฺติยา ทิฏฺฐ\-สุต\-ปริสงฺกิโต”ติ. น เหว โข ภิกฺขุ ภิกฺขุํ ปสฺสติ สีลวิปตฺติยา ทิฏฺฐ\-สุต\-ปริสงฺกิตํ, นาปิ อญฺโญ ภิกฺขุ ภิกฺขุสฺส อาโรเจติ “อิตฺถนฺนาโม อาวุโส \csromanfolio{433}ภิกฺขุ สีลวิปตฺติยา ทิฏฺฐ\-สุต\-ปริสงฺกิโต”ติ. อปิ จ โสว ภิกฺขุ ภิกฺขุสฺส อาโรเจติ “อหํ อาวุโส สีลวิปตฺติยา ทิฏฺฐ\-สุต\-ปริสงฺกิโตมฺหี”ติ. อากงฺขมาโน ภิกฺขเว ภิกฺขุ เตน ทิฏฺเฐน เตน สุเตน ตาย ปริสงฺกาย ตทหุโปสเถ จาตุทฺทเส วา ปนฺนรเส วา ตสฺมิํ ปุคฺคเล สมฺมุขีภูเต สํฆมชฺเฌ อุทาหเรยฺย\csromandash{}}

\prose{“สุณาตุ เม ภนฺเต สํโฆ, อิตฺถนฺนาโม ปุคฺคโล สีลวิปตฺติยา ทิฏฺฐ\-สุต\-ปริสงฺกิโต\csromansharedfootnote{fffeff91fa6bcdbf}{ทิฏฺฐสุตปริสงฺกิโต โหติ (สฺยา, กํ)}, ตสฺส ปาติโมกฺขํ ฐเปมิ น ตสฺมิํ สมฺมุขีภูเต ปาติโมกฺขํ อุทฺทิสิตพฺพนฺ”ติ. ธมฺมิกํ ปาติโมกฺข\-ฏฺฐปนํ.}

\proseitem{396}{กถํ อาจารวิปตฺติยา ทิฏฺฐ\-สุต\-ปริสงฺกิโต โหติ. อิธ ภิกฺขเว ภิกฺขุ เยหิ อากาเรหิ เยหิ ลิงฺเคหิ เยหิ นิมิตฺเตหิ อาจารวิปตฺติยา ทิฏฺฐ\-สุต\-ปริสงฺกิโต โหติ, เตหิ อากาเรหิ เตหิ ลิงฺเคหิ เตหิ นิมิตฺเตหิ ภิกฺขุ ภิกฺขุํ ปสฺสติ อาจารวิปตฺติยา ทิฏฺฐ\-สุต\-ปริสงฺกิตํ. น เหว โข ภิกฺขุ ภิกฺขุํ ปสฺสติ อาจารวิปตฺติยา ทิฏฺฐ\-สุต\-ปริสงฺกิตํ, อปิ จ อญฺโญ ภิกฺขุ ภิกฺขุสฺส อาโรเจติ “อิตฺถนฺนาโม อาวุโส ภิกฺขุ อาจารวิปตฺติยา ทิฏฺฐ\-สุต\-ปริสงฺกิโต”ติ. น เหว โข ภิกฺขุ ภิกฺขุํ ปสฺสติ อาจารวิปตฺติยา ทิฏฺฐ\-สุต\-ปริสงฺกิตํ, นาปิ อญฺโญ ภิกฺขุ ภิกฺขุสฺส อาโรเจติ “อิตฺถนฺนาโม อาวุโส ภิกฺขุ อาจารวิปตฺติยา ทิฏฺฐ\-สุต\-ปริสงฺกิโต”ติ. อปิ จ โสว ภิกฺขุ ภิกฺขุสฺส อาโรเจติ “อหํ อาวุโส อาจารวิปตฺติยา ทิฏฺฐ\-สุต\-ปริสงฺกิโตมฺหี”ติ. อากงฺขมาโน ภิกฺขเว ภิกฺขุ เตน ทิฏฺเฐน เตน สุเตน ตาย ปริสงฺกาย ตทหุโปสเถ จาตุทฺทเส วา ปนฺนรเส วา ตสฺมิํ ปุคฺคเล สมฺมุขีภูเต สํฆมชฺเฌ อุทาหเรยฺย\csromandash{}}

\prose{“สุณาตุ เม ภนฺเต สํโฆ, อิตฺถนฺนาโม ปุคฺคโล อาจารวิปตฺติยา ทิฏฺฐ\-สุต\-ปริสงฺกิโต, ตสฺส ปาติโมกฺขํ ฐเปมิ ‘น ตสฺมิํ สมฺมุขีภูเต ปาติโมกฺขํ อุทฺทิสิตพฺพนฺ’ติ”. ธมฺมิกํ ปาติโมกฺข\-ฏฺฐปนํ.}

\proseitem{397}{กถํ ทิฏฺฐิ\-วิปตฺติยา ทิฏฺฐ\-สุต\-ปริสงฺกิโต โหติ. อิธ ปน ภิกฺขเว ภิกฺขุ เยหิ อากาเรหิ เยหิ ลิงฺเคหิ เยหิ นิมิตฺเตหิ \csromanfolio{434}ทิฏฺฐิ\-วิปตฺติยา ทิฏฺฐ\-สุต\-ปริสงฺกิโต โหติ, เตหิ อากาเรหิ เตหิ ลิงฺเคหิ เตหิ นิมิตฺเตหิ ภิกฺขุ ภิกฺขุํ ปสฺสติ ทิฏฺฐิ\-วิปตฺติยา ทิฏฺฐ\-สุต\-ปริสงฺกิตํ. น เหว โข ภิกฺขุ ภิกฺขุํ ปสฺสติ ทิฏฺฐิ\-วิปตฺติยา ทิฏฺฐ\-สุต\-ปริสงฺกิตํ, อปิ จ อญฺโญ ภิกฺขุ ภิกฺขุสฺส อาโรเจติ “อิตฺถนฺนาโม อาวุโส ภิกฺขุ ทิฏฺฐิ\-วิปตฺติยา ทิฏฺฐ\-สุต\-ปริสงฺกิโต”ติ. น เหว โข ภิกฺขุ ภิกฺขุํ ปสฺสติ ทิฏฺฐิ\-วิปตฺติยา ทิฏฺฐ\-สุต\-ปริสงฺกิตํ, นาปิ อญฺโญ ภิกฺขุ ภิกฺขุสฺส อาโรเจติ “อิตฺถนฺนาโม อาวุโส ภิกฺขุ ทิฏฺฐิ\-วิปตฺติยา, ทิฏฺฐ\-สุต\-ปริสงฺกิโต”ติ, อปิ จ โสว ภิกฺขุ ภิกฺขุสฺส อาโรเจติ “อหํ อาวุโส ทิฏฺฐิ\-วิปตฺติยา ทิฏฺฐ\-สุต\-ปริสงฺกิโตมฺหี”ติ. อากงฺขมาโน ภิกฺขเว ภิกฺขุ เตน ทิฏฺเฐน เตน สุเตน ตาย ปริสงฺกาย ตทหุโปสเถ จาตุทฺทเส วา ปนฺนรเส วา ตสฺมิํ ปุคฺคเล สมฺมุขีภูเต สํฆมชฺเฌ อุทาหเรยฺย\csromandash{}}

\prose{“สุณาตุ เม ภนฺเต สํโฆ, อิตฺถนฺนาโม ปุคฺคโล ทิฏฺฐิ\-วิปตฺติยา ทิฏฺฐ\-สุต\-ปริสงฺกิโต, ตสฺส ปาติโมกฺขํ ฐเปมิ ‘น ตสฺมิํ สมฺมุขีภูเต ปาติโมกฺขํ อุทฺทิสิตพฺพนฺ’ติ”. ธมฺมิกํ ปาติโมกฺข\-ฏฺฐปนํ.}

\vspace{\tipitakaparskip}
\csromancenter{อิมานิ ทส ธมฺมิกานิ ปาติโมกฺขฏฺฐปนานีติ.}
\nitthitam{ปฐม\-ภาณวาโร นิฏฺฐิโต.}
\csromanmatikaanchor{mtk.199}
\csromantocmarkref{section}{7. อตฺตาทานองฺค}{mtk.199}
\bookmark[dest={mtk.199},level=1]{7. อตฺตาทานองฺค}
\csromanheader{1}{7. อตฺตาทานองฺค}
\proseitem{398}{อถ โข อายสฺมา อุปาลิ เยน ภควา เตนุปสงฺกมิ, อุปสงฺกมิตฺวา ภควนฺตํ อภิวาเทตฺวา เอกมนฺตํ นิสีทิ, เอกมนฺตํ เอกมนฺตํ นิสินฺโน โข อายสฺมา อุปาลิ ภควนฺตํ เอตทโวจ “อตฺตาทานํ อาทาตุกาเมน ภนฺเต ภิกฺขุนา กตม\-งฺค\-สมนฺนาคตํ\csromansharedfootnote{54b89a07df83fd57}{กตงฺคสมนฺนาคตํ (ก)} อตฺตาทานํ อาทาตพฺพนฺ”ติ.}

\prose{\csromansymbolfootnote{*}{วิ 5. 331 ปิฏฺเฐปิ.}อตฺตาทานํ อาทาตุกาเมน อุปาลิ ภิกฺขุนา ปญฺจ\-งฺค\-สมนฺนาคตํ อตฺตาทานํ อาทาตพฺพํ, อตฺตาทานํ อาทาตุกาเมน อุปาลิ ภิกฺขุนา เอวํ ปจฺจเวกฺขิตพฺพํ “ยํ โข อหํ อิมํ อตฺตาทานํ อาทาตุกาโม, กาโล นุ โข อิมํ \csromanfolio{435}อตฺตาทานํ อาทาตุํ, อุทาหุ โน”ติ. สเจ อุปาลิ ภิกฺขุ ปจฺจเวกฺขมาโน เอวํ ชานาติ “อกาโล อิมํ อตฺตาทานํ อาทาตุํ, โน กาโล”ติ. น ตํ อุปาลิ อตฺตาทานํ อาทาตพฺพํ.}

\prose{สเจ ปนุปาลิ ภิกฺขุ ปจฺจเวกฺขมาโน เอวํ ชานาติ “กาโล อิมํ อตฺตาทานํ อาทาตุํ, โน อกาโล”ติ. เตนุปาลิ ภิกฺขุนา อุตฺตริ ปจฺจเวกฺขิตพฺพํ “ยํ โข อหํ อิมํ อตฺตาทานํ อาทาตุกาโม, ภูตํ นุ โข อิทํ อตฺตาทานํ, อุทาหุ โน”ติ. สเจ อุปาลิ ภิกฺขุ ปจฺจเวกฺขมาโน เอวํ ชานาติ “อภูตํ อิทํ อตฺตาทานํ, โน ภูตนฺ”ติ, น ตํ อุปาลิ อตฺตาทานํ อาทาตพฺพํ.}

\prose{สเจ ปนุปาลิ ภิกฺขุ ปจฺจเวกฺขมาโน เอวํ ชานาติ “ภูตํ อิทํ อตฺตาทานํ, โน อภูตนฺ”ติ, เตนุปาลิ ภิกฺขุนา อุตฺตริ ปจฺจเวกฺขิตพฺพํ “ยํ โข อหํ อิมํ อตฺตาทานํ อาทาตุกาโม, อตฺถสญฺหิตํ นุ โข อิทํ อตฺตาทานํ, อุทาหุ โน”ติ. สเจ อุปาลิ ภิกฺขุ ปจฺจเวกฺขมาโน เอวํ ชานาติ “อนตฺถสญฺหิตํ อิทํ อตฺตาทานํ, โน อตฺถสญฺหิตนฺ”ติ, น ตํ อุปาลิ อตฺตาทานํ อาทาตพฺพํ.}

\prose{สเจ ปนุปาลิ ภิกฺขุ ปจฺจเวกฺขมาโน เอวํ ชานาติ “อตฺถสญฺหิตํ อิทํ อตฺตาทานํ, โน อนตฺถสญฺหิตนฺ”ติ, เตนุปาลิ ภิกฺขุนา อุตฺตริ ปจฺจเวกฺขิตพฺพํ “อิมํ โข อหํ อตฺตาทานํ อาทิยมาโน ลภิสฺสามิ สนฺทิฏฺเฐ สมฺภตฺเต ภิกฺขู ธมฺมโต วินยโต ปกฺเข, อุทาหุ โน”ติ. สเจ อุปาลิ ภิกฺขุ ปจฺจเวกฺขมาโน เอวํ ชานาติ “อิมํ โข อหํ อตฺตาทานํ อาทิยมาโน น ลภิสฺสามิ สนฺทิฏฺเฐ สมฺเภตฺเต ภิกฺขู ธมฺมโต วินยโต ปกฺเข”ติ, น ตํ อุปาลิ อตฺตาทานํ อาทาตพฺพํ.}

\prose{สเจ ปนุปาลิ ภิกฺขุ ปจฺจเวกฺขมาโน เอวํ ชานาติ “อิมํ โข อหํ อตฺตาทานํ อาทิยมาโน ลภิสฺสามิ สนฺทิฏฺเฐ สมฺภตฺเต ภิกฺขู ธมฺมโต วินยโต ปกฺเข”ติ, เตนุปาลิ ภิกฺขุนา อุตฺตริ ปจฺจเวกฺขิตพฺพํ “อิมํ โข เม อตฺตาทานํ อาทิยโต ภวิสฺสติ สํฆสฺส ตโตนิทานํ ภณฺฑนํ กลโห วิคฺคโห วิวาโท สํฆเภโท สํฆราชิ สํฆววตฺถานํ สํฆนานากรณํ, อุทาหุ โน”ติ. สเจ อุปาลิ ภิกฺขุ ปจฺจเวกฺขมาโน เอวํ ชานาติ “อิมํ โข เม อตฺตาทานํ อาทิยโต ภวิสฺสติ สํฆสฺส ตโตนิทานํ ภณฺฑํ กลโห วิคฺคโห วิวาโท สํฆเภโท สํฆราชิ \csromanfolio{436}สํฆววตฺถานํ สํฆนานากรณนฺ”ติ, น ตํ อุปาลิ อตฺตาทานํ อาทาตพฺพํ. สเจ ปนุปาลิ ภิกฺขุ ปจฺจเวกฺขมาโน เอวํ ชานาติ “อิมํ โข เม อตฺตาทานํ อาทิยโต น ภวิสฺสติ สํฆสฺส ตโตนิทานํ ภณฺฑนํ กลโห วิคฺคโห วิวาโท สํฆเภโท สํฆราชิ สํฆววตฺถานํ สํฆนานากรณนฺ”ติ, อาทาตพฺพํ ตํ อุปาลิ อตฺตาทานํ. เอวํ ปญฺจ\-งฺค\-สมนฺนาคตํ โข อุปาลิ อตฺตาทานํ อาทินฺนํ ปจฺฉาปิ อวิปฺปฏิสารการกรํ ภวิสฺสตีติ.}

\csromanmatikaanchor{mtk.200}
\csromantocmarkref{section}{8. โจทเกนปจฺจเวกฺขิตพฺพธมฺม}{mtk.200}
\bookmark[dest={mtk.200},level=1]{8. โจทเกนปจฺจเวกฺขิตพฺพธมฺม}
\csromanheader{1}{8. \textbf{โจทเกนปจฺจเวกฺขิตพฺพธมฺม}}
\proseitem{399}{\csromansymbolfootnote{*}{วิ 5. 328 ปิฏฺเฐปิ. (อํ 3. 318 ปิฏฺเฐ ปน โถกํ วิสทิสํ)}โจทเกน ภนฺเต ภิกฺขุนา ปรํ โจเทตุกาเมน กติ ธมฺเม อชฺฌตฺตํ ปจฺจเวกฺขิตฺวา ปโร โจเทตพฺโพติ. โจทเกน อุปาลิ ภิกฺขุนา ปรํ โจเทตุกาเมน ปญฺจ ธมฺเม อชฺฌตฺตํ ปจฺจเวกฺขิตฺวา ปโร โจเทตพฺโพ.}

\prose{โจทเกน อุปาลิ ภิกฺขุนา ปรํ โจเทตุกาเมน เอวํ ปจฺจเวกฺขิตพฺพํ “ปริสุทฺธ\-กาย\-สมาจาโร นุ โขมฺหิ, ปริสุทฺเธนมฺหิ กายสมาจาเรน สมนฺนาคโต อจฺฉิทฺเทน อปฺปฏิมํเสน, สํวิชฺชติ นุ โข เม เอโส ธมฺโม, อุทาหุ โน”ติ. โน เจ อุปาลิ ภิกฺขุ ปริสุทฺธ\-กาย\-สมาจาโร โหติ ปริสุทฺเธน กายสมาจาเรน สมนฺนาคโต อจฺฉิทฺเทน อปฺปฏิมํเสน, ตสฺส ภวนฺติ วตฺตาโร “อิงฺฆ ตาว อายสฺมา กายิกํ สิกฺขสฺสู”ติ, อิติสฺส ภวนฺติ วตฺตาโร.}

\prose{ปุน จปรํ อุปาลิ โจทเกน ภิกฺขุนา ปรํ โจเทตุกาเมน เอวํ ปจฺจเวกฺขิตพฺพํ “ปริสุทฺธ\-วจี\-สมาจาโร นุ โขมฺหิ, ปริสุทฺเธนมฺหิ วจีสมาจาเรน สมนฺนาคโต อจฺฉิทฺเทน อปฺปฏิมํเสน, สํวิชฺชติ นุ โข เม เอโส ธมฺโม, อุทาหุ โน”ติ. โน เจ อุปาลิ ภิกฺขุ ปริสุทฺธ\-วจี\-สมาจาโร โหติ ปริสุทฺเธน วจีสมาจาเรน สมนฺนาคโต อจฺฉิทฺเทน อปฺปฏิมํเสน, ตสฺส ภวนฺติ วตฺตาโร “อิงฺฆ ตาว อายสฺมา วาจสิกํ สิกฺขสฺสู”ติ, อิติสฺส ภวนฺติ วตฺตาโร.}

\prose{ปุน จปรํ อุปาลิ โจทเกน ภิกฺขุนา ปรํ โจเทตุกาเมน เอวํ ปจฺจเวกฺขิตพฺพํ “เมตฺตํ นุ โข เม จิตฺตํ ปจฺจุปฏฺฐิตํ สพฺรหฺมจารีสุ อนาฆาตํ, สํวิชฺชติ นุ โข เม เอโส ธมฺโม, อุทาหุ โน”ติ. โน เจ อุปาลิ ภิกฺขุโน เมตฺตจิตฺตํ ปจฺจุปฏฺฐิตํ โหติ สพฺรหฺมจารีสุ อนาฆาตํ, \csromanfolio{437}ตสฺส ภวนฺติ วตฺตาโร “อิงฺฆ ตาว อายสฺมา สพฺรหฺมจารีสุ เมตฺตจิตฺตํ อุปฏฺฐาเปหี”ติ, อิติสฺส ภวนฺติ วตฺตาโร.}

\prose{ปุน จปรํ อุปาลิ โจทเกน ภิกฺขุนา ปรํ โจเทตุกาเมน เอวํ ปจฺจเวกฺขิตพฺพํ “พหุสฺสุโต นุ โขมฺหิ สุตธโร สุตสนฺนิจฺจโย, เย เต ธมฺมา อาทิกลฺยาณา มชฺเฌกลฺยาณา ปริโยสาน\-กลฺยาณา สาตฺถํ สพฺยญฺชนํ เกวล\-ปริปุณฺณํ ปริสุทฺธํ พฺรหฺมจริยํ อภิวทนฺติ, ตถารูปา เม ธมฺมา พหุสฺสุตา โหนฺติ ธาตา วจสา ปริจิตา มนสา\-นุเปกฺขิตา ทิฏฺฐิยา สุปฺปฏิวิทฺธา, สํวิชฺชติ นุ โข เม เอโส ธมฺโม, อุทาหุ โน”ติ. โน เจ อุปาลิ ภิกฺขุ พหุสฺสุโต โหติ สุตธโร สุตสนฺนิจฺจโย, เย เต ธมฺมา อาทิกลฺยาณา มชฺเฌกลฺยาณา ปริโยสาน\-กลฺยาณา สาตฺถํ สพฺยญฺชนํ เกวล\-ปริปุณฺณํ ปริสุทฺธํ พฺรหฺมจริยํ อภิวทนฺติ, ตถารูปสฺส ธมฺมา น พหุสฺสุตา โหนฺติ ธาตา วจสา ปริจิตา มนสา\-นุเปกฺขิตา ทิฏฺฐิยา สุปฺปฏิวิทฺธา, ตสฺส ภวนฺติ วตฺตาโร “อิงฺฆ ตาว อายสฺมา อาคมํ ปริยาปุณสฺสู”ติ, อิติสฺส ภวนฺติ วตฺตาโร.}

\prose{ปุน จปรํ อุปาลิ โจทเกน ภิกฺขุนา ปรํ โจเทตุกาเมน เอวํ ปจฺจเวกฺขิตพฺพํ “อุภยานิ โข เม ปาติโมกฺขานิ วิตฺถาเรน สฺวาคตานิ โหนฺติ สุวิภตฺตานิ สุปฺปวตฺตีนิ สุวินิจฺฉิตานิ สุตฺตโส อนุพฺยญฺชนโส. สํวิชฺชติ นุ โข เม เอโส ธมฺโม, อุทาหุ โน”ติ. โน เจ อุปาลิ ภิกฺขุโน อุภยานิ ปาติโมกฺขานิ วิตฺถาเรน สฺวาคตานิ โหนฺติ สุวิภตฺตานิ สุปฺปวตฺตีนิ สุวินิจฺฉิตานิ สุตฺตโส อนุพฺยญฺชนโส, “อิทํ ปนาวุโส กตฺถ วุตฺตํ ภควตา”ติ อิติ ปุฏฺโฐ น สมฺปายติ\csromansharedfootnote{9ac72db7e319927e}{น สมฺปาทยติ (ก)}, ตสฺส ภวนฺติ วตฺตาโร “อิงฺฆ ตาว อายสฺมา วินยํ ปริยาปุณสฺสู”ติ, อิติสฺส ภวนฺติ วตฺตาโร. โจทเกนุปาลิ ภิกฺขุนา ปรํ โจเทตุกาเมน อิเม ปญฺจ ธมฺเม อชฺฌตฺตํ ปจฺจเวกฺขิตฺวา ปโร โจเทตพฺโพติ.}

\csromanmatikaanchor{mtk.201}
\csromantocmarkref{section}{9. โจทเกนอุปฏฺฐาเปตพฺพธมฺม}{mtk.201}
\bookmark[dest={mtk.201},level=1]{9. โจทเกนอุปฏฺฐาเปตพฺพธมฺม}
\csromanheader{1}{9. โจทเกนอุปฏฺฐาเปตพฺพธมฺม}
\proseitem{400}{\csromansymbolfootnote{*}{วิ 5. 330 ปิฏฺเฐปิ.}โจทเกน ภนฺเต ภิกฺขุนา ปรํ โจเทตุกาเมน กติ ธมฺเม อชฺฌตฺตํ อุปฏฺฐาเปตฺวา ปโร โจเทตพฺโพติ. โจทเกนุปาลิ ภิกฺขุนา ปรํ โจเทตุกาเมน ปญฺจ ธมฺเม อชฺฌตฺตํ อุปฏฺฐาเปตฺวา ปโร โจเทตพฺโพ “กาเลน วกฺขามิ, โน อกาเลน, ภูเตน วกฺขามิ, โน อภูเตน, \csromanfolio{438}สเณฺหน วกฺขามิ, โน ผรุเสน, อตฺถสํหิเตน วกฺขามิ, โน อนตฺถ\-สํหิเตน, เมตฺตจิตฺโต วกฺขามิ, โน โทสนฺตโร”ติ. โจทเกนุปาลิ ภิกฺขุนา ปรํ โจเทตุกาเมน อิเม ปญฺจ ธมฺเม อชฺฌตฺตํ อุปฏฺฐาเปตฺวา ปโร โจเทตพฺโพติ.}

\csromanmatikaanchor{mtk.202}
\csromantocmarkref{section}{10. โจทกจุทิตกปฏิสํยุตฺตกถา}{mtk.202}
\bookmark[dest={mtk.202},level=1]{10. โจทกจุทิตกปฏิสํยุตฺตกถา}
\csromanheader{1}{10. โจทก\-จุทิตก\-ปฏิสํยุตฺต\-กถา}
\proseitem{401}{อธมฺม\-โจทกสฺส ภนฺเต ภิกฺขุโน กติหากาเรหิ วิปฺปฏิสาโร อุปทหาตพฺโพติ. อธมฺม\-โจทกสฺส อุปาลิ ภิกฺขุโน ปญฺจหากาเรหิ วิปฺปฏิสาโร อุปทหาตพฺโพ “อกาเลนายสฺมา โจเทสิ, โน กาเลน, อลํ เต วิปฺปฏิสาราย. อภูเตนายสฺมา โจเทสิ, โน ภูเตน, อลํ เต วิปฺปฏิสาราย. ผรุเสนายสฺมา โจเทสิ, โน สเณฺหน, อลํ เต วิปฺปฏิสาราย. อนตฺถสํหิเตนา\-ยสฺมา โจเทสิ, โน อตฺถสํหิเตน, อลํ เต วิปฺปฏิสาราย. โทสนฺตโร อายสฺมา โจเทสิ, โน เมตฺตจิตฺโต, อลํ เต วิปฺปฏิสารายา”ติ. อธมฺม\-โจทกสฺส อุปาลิ ภิกฺขุโน อิเมหิ ปญฺจหากาเรหิ วิปฺปฏิสาโร อุปทหาตพฺโพ, ตํ กิสฺส เหตุ, ยถา น อญฺโญปิ ภิกฺขุ อภูเตน โจเทตพฺพํ มญฺเญยฺยาติ.}

\prose{อธมฺม\-จุทิตสฺส ปน ภนฺเต ภิกฺขุโน กติหากาเรหิ อวิปฺปฏิสาโร อุปทหาตพฺโพติ. อธมฺม\-จุทิตสฺส อุปาลิ ภิกฺขุโน ปญฺจหากาเรหิ อวิปฺปฏิสาโร อุปทหาตพฺโพ “อกาเลนายสฺมา จุทิโต, โน กาเลน, อลํ เต อวิปฺปฏิสาราย. อภูเตนายสฺมา จุทิโต, โน ภูเตน, อลํ เต อวิปฺปฏิสาราย. ผรุเสนายสฺมา จุทิโต, โน สเณฺหน, อลํ เต อวิปฺปฏิสาราย. อนตฺถสํหิเตนา\-ยสฺมา จุทิโต, โน อตฺถสํหิเตน, อลํ เต อวิปฺปฏิสาราย. โทสนฺตเรนา\-ยสฺมา จุทิโต, โน เมตฺตจิตฺเตน, อลํ เต อวิปฺปฏิสารายา”ติ. อธมฺมจุทิตส อุปาลิ ภิกฺขุโน อิเมหิ ปญฺจหากาเรหิ อวิปฺปฏิสาโร อุปทหาตพฺโพติ.}

\prose{ธมฺมโจทกสฺส ภนฺเต ภิกฺขุโน กติหากาเรหิ อวิปฺปฏิสาโร อุปทหาตพฺโพติ. ธมฺมโจทกสฺส อุปาลิ ภิกฺขุโน ปญฺจหากาเรหิ อวิปฺปฏิสาโร อุปทหาตพฺโพ “กาเลนายสฺมา โจเทสิ, โน อกาเลน, อลํ เต อวิปฺปฏิสาราย. ภูเตนายสฺมา โจเทสิ, โน \csromanfolio{439}อภูเตน, อลํ เต อวิปฺปฏิสาราย. สเณฺหนายสฺมา โจเทสิ, โน ผรุเสน, อลํ เต อวิปฺปฏิสาราย. อตฺถสํหิเตนา\-ยสฺมา โจเทสิ, โน อนตฺถ\-สํหิเตน, อลํ เต อวิปฺปฏิสาราย. เมตฺตจิตฺโต อายสฺมา โจเทสิ, โน โทสนฺตโร, อลํ เต อวิปฺปฏิสารายา”ติ. ธมฺมโจทกสฺส อุปาลิ ภิกฺขุโน อิเมหิ ปญฺจหากาเรหิ อวิปฺปฏิสาโร อุปทหาตพฺโพ. ตํ กิสฺส เหตุ, ยถา อญฺโญปิ ภิกฺขุ ภูเตน โจเทตพฺพํ มญฺเญยฺยาติ.}

\prose{ธมฺม\-จุทิตสฺส ปน ภนฺเต ภิกฺขุโน กติหากาเรหิ วิปฺปฏิสาโร อุปทหาตพฺโพติ. ธมฺม\-จุทิตสฺส อุปาลิ ภิกฺขุโน ปญฺจหากาเรหิ วิปฺปฏิสาโร อุปทหาตพฺโพ “กาเลนายสฺมา จุทิโต, โน อกาเลน, อลํ เต วิปฺปฏิสาราย. ภูเตนายสฺมา จุทิโต, โน อภูเตน, อลํ เต วิปฺปฏิสาราย. สเณฺหนายสฺมา จุทิโต, โน ผรุเสน, อลํ เต วิปฺปฏิสาราย. อตฺถสํหิเตนา\-ยสฺมา จุทิโต, โน อนตฺถ\-สํหิเตน, อลํ เต วิปฺปฏิสาราย. เมตฺตจิตฺเตนา\-ยสฺมา จุทิโต, โน โทสนฺตเรน, อลํ เต วิปฺปฏิสารายา”ติ. ธมฺม\-จุทิตสฺส อุปาลิ ภิกฺขุโน อิเมหิ ปญฺจหากาเรหิ วิปฺปฏิสาโร อุปทหาตพฺโพติ.}

\prose{โจทเกน ภนฺเต ภิกฺขุนา ปรํ โจเทตุกาเมน กติ ธมฺเม อชฺฌตฺตํ มนสิ กริตฺวา ปโร โจเทตพฺโพติ. โจทเกนุปาลิ ภิกฺขุนา ปรํ โจเทตุกาเมน ปญฺจ ธมฺเม อชฺฌตฺตํ มนสิ กริตฺวา ปโร โจเทตพฺโพ. การุญฺญตา หิเตสิตา อนุกมฺปิตา\csromansharedfootnote{d7b780e041ea56e3}{อนุกมฺปตา (ก)} อาปตฺติวุฏฺฐานตา วินยปุเรกฺขารตาติ. โจทเกนุปาลิ ภิกฺขุนา ปรํ โจเทตุกาเมน อิเม ปญฺจ ธมฺเม อชฺฌตฺตํ มนสิ กริตฺวา ปโร โจเทตพฺโพติ.}

\prosewithcloser{จุทิเตน ปน ภนฺเต ภิกฺขุนา กติสุ ธมฺเมสุ ปติฏฺฐาตพฺพนฺติ. จุทิเตนุปาลิ ภิกฺขุนา ทฺวีสุ ธมฺเมสุ ปติฏฺฐาตพฺพํ สจฺเจ จ อกุปฺเป จาติ.}{ทุติย\-ภาณวาโร นิฏฺฐิโต.}
\csromancenter{ปาติโมกฺขฏฺฐปน\-กฺขนฺธโก นวโม.}
\csromancenter{อิมมฺหิ ขนฺธเก วตฺถู ติํส.}
\csromancenter{\textbf{ตสฺสุทฺทานํ}}
\csromanmatikaanchor{mtk.203}
\csromantocmarkref{section}{อุทฺทานคาถา}{mtk.203}
\bookmark[dest={mtk.203},level=1]{อุทฺทานคาถา}
\setlength{\csromangathapagewidth}{0pt}
\setlength{\csromangathawakwidth}{0pt}
\csromangathameasuregroup{อุโปสเถ ยาวติกํ, \\ โมคฺคลฺลาเนน นิจฺฉุทฺโธ, \\ นินฺโนนุปุพฺพสิกฺขา จ, \\ กุณปุกฺขิปติ สํโฆ, \\ สวนฺติ ปรินิพฺพนฺติ, \\ พหุ ธมฺมวินโยปิ, \\ สมุทฺทํ อุปมํ กตฺวา, \\ อุโปสเถ ปาติโมกฺขํ, \\ ปฏิกจฺเจว อุชฺฌนฺติ, \\ ปญฺจ ฉ สตฺต อฏฺฐานิ, \\ สีล อาจาร ทิฏฺฐิ จ, \\ ปาราชิกญฺจ สํฆาทิ, \\ ทุกฺกฏํ ปญฺจภาเคสุ, \\ อกตาย กตาย จ, \\ ปาราชิกญฺจ สํฆาทิ, \\ ปาฏิเทสนิยญฺเจว, \\ สีลาจารวิปตฺติ จ, \\ ยา จ อฏฺฐา กตากเต, \\ อกตาย กตายาปิ, \\ เอวํ นววิธา วุตฺตา, \\ ปาราชิโก วิปฺปกตา, \\ อุเปติ ปจฺจาทิยติ, \\ สีลาจารวิปตฺติ จ, \\ ทิฏฺฐสุตปริสงฺกิตํ, \\ ภิกฺขุ วิปสฺสติ ภิกฺขุํ, \\ โส เยว ตสฺส อกฺขาติ\textsuperscript{1}, \\ วุฏฺฐาติ อนฺตราเยน, \\ มนุสฺสอมนุสฺสา จ, \\ ทสนฺนมญฺญตเรน, \\ ธมฺมิกาธมฺมิกา เจว, \\ กาลภูตตฺถสํหิตํ, \\ กายวาจสิกา เมตฺตา, \\ กาลภูเตน สเณฺหน, \\ วิปฺปฏิสารธมฺเมน, \\ ธมฺมโจทจุทิตสฺส, \\ กรุณา หิตานุกมฺปิ, \\ โจทกสฺส ปฏิปตฺติ, \\ สจฺเจ เจว อกุปฺเป จ,}{\csromangathabat{อุโปสเถ ยาวติกํ,}{ปาปภิกฺขุ น นิกฺขมิ.} \\ \csromangathabat{โมคฺคลฺลาเนน นิจฺฉุทฺโธ,}{อจฺเฉรา ชินสาสเน.} \\ \csromangathabat{นินฺโนนุปุพฺพสิกฺขา จ,}{ฐิตธมฺโม นาติกฺกมฺม.} \\ \csromangathabat{กุณปุกฺขิปติ สํโฆ,}{สวนฺติโย ชหนฺติ จ.} \\ \csromangathabat{สวนฺติ ปรินิพฺพนฺติ,}{เอกรส วิมุตฺติ จ.} \\ \csromangathabat{พหุ ธมฺมวินโยปิ,}{ภูตฏฺฐาริยปุคฺคลา.} \\ \csromangathabat{สมุทฺทํ อุปมํ กตฺวา,}{วาเจสิ\textsuperscript{1} สาสเน คุณํ.} \\ \csromangathabat{อุโปสเถ ปาติโมกฺขํ,}{น อเมฺห โกจิ ชานาติ.} \\ \csromangathabat{ปฏิกจฺเจว อุชฺฌนฺติ,}{เอโก เทฺว ตีณิ จตฺตาริ.} \\ \csromangathabat{ปญฺจ ฉ สตฺต อฏฺฐานิ,}{นวา จ ทสมานิ จ.} \\ \csromangathabat{สีล อาจาร ทิฏฺฐิ จ,}{อาชีวํ จตุภาคิเก.} \\ \csromangathabat{ปาราชิกญฺจ สํฆาทิ,}{ปาจิตฺติ ปาฏิเทสนิ.} \\ \csromangathabat{ทุกฺกฏํ ปญฺจภาเคสุ,}{สีลาจารวิปตฺติ จ.} \\ \csromangathabat{อกตาย กตาย จ,}{ฉภาเคสุ ยถาวิธิ.} \\ \csromangathabat{ปาราชิกญฺจ สํฆาทิ,}{ถุลฺลํ ปาจิตฺติเยน จ.} \\ \csromangathabat{ปาฏิเทสนิยญฺเจว,}{ทุกฺกฏญฺจ ทุพฺภาสิตํ.} \\ \csromangathabat{สีลาจารวิปตฺติ จ,}{ทิฏฺฐิอาชีววิปตฺติ.} \\ \csromangathabat{ยา จ อฏฺฐา กตากเต,}{เตเนตา สิลาจารทิฏฺฐิยา.} \\ \csromangathabat{อกตาย กตายาปิ,}{กตากตายเมว จ.} \\ \csromangathabat{เอวํ นววิธา วุตฺตา,}{ยถาภูเตน ญายโต.} \\ \csromangathabat{ปาราชิโก วิปฺปกตา,}{ปจฺจกฺขาโต ตเถว จ.} \\ \csromangathabat{อุเปติ ปจฺจาทิยติ,}{ปจฺจาทานกถา จ ยา.} \\ \csromangathabat{สีลาจารวิปตฺติ จ,}{ตถา ทิฏฺฐิวิปตฺติยา.} \\ \csromangathabat{ทิฏฺฐสุตปริสงฺกิตํ,}{ทสธา ตํ วิชานาถ.} \\ \csromangathabat{ภิกฺขุ วิปสฺสติ ภิกฺขุํ,}{อญฺโญ จาโรจยาติ ตํ.} \\ \csromangathabat{โส เยว ตสฺส อกฺขาติ\textsuperscript{1},}{ปาติโมกฺขํ ฐเปติ โส.} \\ \csromangathabat{วุฏฺฐาติ อนฺตราเยน,}{ราชโจรคฺคุทกา จ.} \\ \csromangathabat{มนุสฺสอมนุสฺสา จ,}{วาฬสรีสปา ชีวิพฺรหฺมํ.} \\ \csromangathabat{ทสนฺนมญฺญตเรน,}{ตสฺมิํ อญฺญตเรสุ วา.} \\ \csromangathabat{ธมฺมิกาธมฺมิกา เจว,}{ยถา มคฺเคน ชานาถ.} \\ \csromangathabat{กาลภูตตฺถสํหิตํ,}{ลภิสฺสามิ ภวิสฺสติ.} \\ \csromangathabat{กายวาจสิกา เมตฺตา,}{พาหุสจฺจํ อุภยานิ.} \\ \csromangathabat{กาลภูเตน สเณฺหน,}{อตฺถเมตฺเตน โจทเย.} \\ \csromangathabat{วิปฺปฏิสารธมฺเมน,}{ตถา วาจา\textsuperscript{1} วิโนทเย.} \\ \csromangathabat{ธมฺมโจทจุทิตสฺส,}{วิโนเทติ วิปฺปฏิสาโร.} \\ \csromangathabat{กรุณา หิตานุกมฺปิ,}{วุฏฺฐานปุเรกฺขารโต.} \\ \csromangathabat{โจทกสฺส ปฏิปตฺติ,}{สมฺพุทฺเธน ปกาสิตา.} \\ \csromangathabat{สจฺเจ เจว อกุปฺเป จ,}{จุทิตสฺเสว ธมฺมตาติ.}}
\csromangathasetleft{อุโปสเถ ยาวติกํ, \\ โมคฺคลฺลาเนน นิจฺฉุทฺโธ, \\ นินฺโนนุปุพฺพสิกฺขา จ, \\ กุณปุกฺขิปติ สํโฆ, \\ สวนฺติ ปรินิพฺพนฺติ, \\ พหุ ธมฺมวินโยปิ, \\ สมุทฺทํ อุปมํ กตฺวา, \\ อุโปสเถ ปาติโมกฺขํ, \\ ปฏิกจฺเจว อุชฺฌนฺติ, \\ ปญฺจ ฉ สตฺต อฏฺฐานิ, \\ สีล อาจาร ทิฏฺฐิ จ, \\ ปาราชิกญฺจ สํฆาทิ, \\ ทุกฺกฏํ ปญฺจภาเคสุ, \\ อกตาย กตาย จ, \\ ปาราชิกญฺจ สํฆาทิ, \\ ปาฏิเทสนิยญฺเจว, \\ สีลาจารวิปตฺติ จ, \\ ยา จ อฏฺฐา กตากเต, \\ อกตาย กตายาปิ, \\ เอวํ นววิธา วุตฺตา, \\ ปาราชิโก วิปฺปกตา, \\ อุเปติ ปจฺจาทิยติ, \\ สีลาจารวิปตฺติ จ, \\ ทิฏฺฐสุตปริสงฺกิตํ, \\ ภิกฺขุ วิปสฺสติ ภิกฺขุํ, \\ โส เยว ตสฺส อกฺขาติ\textsuperscript{1}, \\ วุฏฺฐาติ อนฺตราเยน, \\ มนุสฺสอมนุสฺสา จ, \\ ทสนฺนมญฺญตเรน, \\ ธมฺมิกาธมฺมิกา เจว, \\ กาลภูตตฺถสํหิตํ, \\ กายวาจสิกา เมตฺตา, \\ กาลภูเตน สเณฺหน, \\ วิปฺปฏิสารธมฺเมน, \\ ธมฺมโจทจุทิตสฺส, \\ กรุณา หิตานุกมฺปิ, \\ โจทกสฺส ปฏิปตฺติ, \\ สจฺเจ เจว อกุปฺเป จ,}
\csromangathagroupwithclosermajor{\csromangathastanza{\csromanfolio{440}\csromangathabat{อุโปสเถ ยาวติกํ,}{ปาปภิกฺขุ น นิกฺขมิ.} \\ \csromangathabat{โมคฺคลฺลาเนน นิจฺฉุทฺโธ,}{อจฺเฉรา ชินสาสเน.}}\csromangathastanza{\csromangathabat{นินฺโนนุปุพฺพสิกฺขา จ,}{ฐิตธมฺโม นาติกฺกมฺม.} \\ \csromangathabat{กุณปุกฺขิปติ สํโฆ,}{สวนฺติโย ชหนฺติ จ.}}\csromangathastanza{\csromangathabat{สวนฺติ ปรินิพฺพนฺติ,}{เอกรส วิมุตฺติ จ.} \\ \csromangathabat{พหุ ธมฺมวินโยปิ,}{ภูตฏฺฐาริยปุคฺคลา.}}\csromangathastanza{\csromangathabat{สมุทฺทํ อุปมํ กตฺวา,}{วาเจสิ\csromansharedfootnote{0f3a0869b226166a}{ฐาเปสิ (ก)} สาสเน คุณํ.} \\ \csromangathabat{อุโปสเถ ปาติโมกฺขํ,}{น อเมฺห โกจิ ชานาติ.}}\csromangathastanza{\csromangathabat{ปฏิกจฺเจว อุชฺฌนฺติ,}{เอโก เทฺว ตีณิ จตฺตาริ.} \\ \csromangathabat{ปญฺจ ฉ สตฺต อฏฺฐานิ,}{นวา จ ทสมานิ จ.}}\csromangathastanza{\csromangathabat{สีล อาจาร ทิฏฺฐิ จ,}{อาชีวํ จตุภาคิเก.} \\ \csromangathabat{ปาราชิกญฺจ สํฆาทิ,}{ปาจิตฺติ ปาฏิเทสนิ.}}\csromangathastanza{\csromangathabat{ทุกฺกฏํ ปญฺจภาเคสุ,}{สีลาจารวิปตฺติ จ.} \\ \csromangathabat{อกตาย กตาย จ,}{ฉภาเคสุ ยถาวิธิ.}}\csromangathastanza{\csromangathabat{ปาราชิกญฺจ สํฆาทิ,}{ถุลฺลํ ปาจิตฺติเยน จ.} \\ \csromangathabat{ปาฏิเทสนิยญฺเจว,}{ทุกฺกฏญฺจ ทุพฺภาสิตํ.}}\csromangathastanza{\csromangathabat{สีลาจารวิปตฺติ จ,}{ทิฏฺฐิอาชีววิปตฺติ.} \\ \csromangathabat{ยา จ อฏฺฐา กตากเต,}{เตเนตา สิลาจารทิฏฺฐิยา.}}\csromangathastanza{\csromangathabat{อกตาย กตายาปิ,}{กตากตายเมว จ.} \\ \csromangathabat{เอวํ นววิธา วุตฺตา,}{ยถาภูเตน ญายโต.}}\csromangathastanza{\csromangathabat{ปาราชิโก วิปฺปกตา,}{ปจฺจกฺขาโต ตเถว จ.} \\ \csromangathabat{อุเปติ ปจฺจาทิยติ,}{ปจฺจาทานกถา จ ยา.}}\csromangathastanza{\csromangathabat{สีลาจารวิปตฺติ จ,}{ตถา ทิฏฺฐิ\-วิปตฺติยา.} \\ \csromangathabat{ทิฏฺฐ\-สุต\-ปริสงฺกิตํ,}{ทสธา ตํ วิชานาถ.}}\csromangathastanza{\csromanfolio{441}\csromangathabat{ภิกฺขุ วิปสฺสติ ภิกฺขุํ,}{อญฺโญ จาโรจยาติ ตํ.} \\ \csromangathabat{โส เยว ตสฺส อกฺขาติ\csromansharedfootnote{979131f797b8f1c5}{วิปสฺสญฺโญ จาโรจติ. ตํ สุทฺเธว ตสฺส อกฺขาติ (ก)},}{ปาติโมกฺขํ ฐเปติ โส.}}\csromangathastanza{\csromangathabat{วุฏฺฐาติ อนฺตราเยน,}{ราชโจรคฺคุทกา จ.} \\ \csromangathabat{มนุสฺสอมนุสฺสา จ,}{วาฬสรีสปา ชีวิพฺรหฺมํ.}}\csromangathastanza{\csromangathabat{ทสนฺนมญฺญตเรน,}{ตสฺมิํ อญฺญตเรสุ วา.} \\ \csromangathabat{ธมฺมิกาธมฺมิกา เจว,}{ยถา มคฺเคน ชานาถ.}}\csromangathastanza{\csromangathabat{กาลภูตตฺถสํหิตํ,}{ลภิสฺสามิ ภวิสฺสติ.} \\ \csromangathabat{กายวาจสิกา เมตฺตา,}{พาหุสจฺจํ อุภยานิ.}}\csromangathastanza{\csromangathabat{กาลภูเตน สเณฺหน,}{อตฺถเมตฺเตน โจทเย.} \\ \csromangathabat{วิปฺปฏิสารธมฺเมน,}{ตถา วาจา\csromansharedfootnote{868957fbc3e6f527}{ตเถวาปิ (สฺยา)} วิโนทเย.}}\csromangathastanza{\csromangathabat{ธมฺมโจทจุทิตสฺส,}{วิโนเทติ วิปฺปฏิสาโร.} \\ \csromangathabat{กรุณา หิตานุกมฺปิ,}{วุฏฺฐานปุเรกฺขารโต.}}\csromangathastanza{\csromangathabat{โจทกสฺส ปฏิปตฺติ,}{สมฺพุทฺเธน ปกาสิตา.} \\ \csromangathabat{สจฺเจ เจว อกุปฺเป จ,}{จุทิตสฺเสว ธมฺมตาติ.}}}{ปาติโมกฺขฏฺฐปน\-กฺขนฺธโก นิฏฺฐิโต.}
\csromanmatikaanchor{mtk.204}
\csromantocmarknopage{chapter}{10. ภิกฺขุนิกฺขนฺธก}{mtk.204}
\bookmark[dest={mtk.204},level=0]{10. ภิกฺขุนิกฺขนฺธก}
\chapterheadpageread{10. ภิกฺขุนิ\-กฺขนฺธก}
\csromanmatikaanchor{mtk.205}
\csromantocmarkref{section}{1. ปฐมภาณวาร}{mtk.205}
\bookmark[dest={mtk.205},level=1]{1. ปฐมภาณวาร}
\csromanheader{1}{1. ปฐม\-ภาณวาร}
\csromanheader{2}{\textbf{มหาปชาปติ โคตมีวตฺถุ}}
\csromanmatikaanchor{mtk.206}
\csromantocmarkref{subsection}{มหาปชาปติโคตมีวตฺถุ}{mtk.206}
\bookmark[dest={mtk.206},level=2]{มหาปชาปติโคตมีวตฺถุ}
\proseitem{402}{\csromanfolio{442}\csromansymbolfootnote{*}{อํ 3. 101 ปิฏฺฐาทีสุปิ อิทํ วตฺถุ อาคตํ.}เตน สมเยน พุทฺโธ ภควา สกฺเกสุ วิหรติ กปิล\-วตฺถุสฺมิํ นิโคฺรธาราเม. อถ โข มหาปชาปติ\csromansharedfootnote{3d7073d5ee5a87cc}{มหาปชาปตี (สี, สฺยา)}โคตมี เยน ภควา เตนุปสงฺกมิ, อุปสงฺกมิตฺวา ภควนฺตํ อภิวาเทตฺวา เอกมนฺตํ อฏฺฐาสิ, เอกมนฺตํ ฐิตา โข มหาปชาปติ โคตมี ภควนฺตํ เอตทโวจ “สาธุ ภนฺเต, ลเภยฺย มาตุคาโม ตถาคต\-ปฺปเวทิเต ธมฺมวินเย อคารสฺมา อนคาริยํ ปพฺพชฺชนฺ”ติ. อลํ โคตมิ, มา เต รุจฺจิ มาตุคามสฺส ตถาคต\-ปฺปเวทิเต ธมฺมวินเย อคารสฺมา อนคาริยํ ปพฺพชฺชาติ. ทุติยมฺปิ โข ฯเปฯ. ตติยมฺปิ โข มหาปชาปติ โคตมี ภควนฺตํ เอตทโวจ “สาธุ ภนฺเต, ลเภยฺย มาตุคาโม ตถาคต\-ปฺปเวทิเต ธมฺมวินเย อคารสฺมา อนคาริยํ ปพฺพชฺชนฺ”ติ. อลํ โคตมิ, มา เต รุจฺจิ มาตุคามสฺส ตถาคต\-ปฺปเวทิเต ธมฺมวินเย อคารสฺมา อนคาริยํ ปพฺพชฺชาติ. อถ โข มหาปชาปติ โคตมี “น ภควา อนุชานาติ มาตุคามสฺส ตถาคต\-ปฺปเวทิเต ธมฺมวินเย อคารสฺมา อนคาริยํ ปพฺพชฺชนฺ”ติ ทุกฺขี ทุมฺมนา อสฺสุมุขี รุทมานา ภควนฺตํ อภิวาเทตฺวา ปทกฺขิณํ กตฺวา ปกฺกามิ.}

\prose{อถ โข ภควา กปิล\-วตฺถุสฺมิํ ยถาภิรนฺตํ วิหริตฺวา เยน เวสาลี เตน จาริกํ ปกฺกามิ, อนุปุพฺเพน จาริกํ จรมาโน เยน เวสาลี ตทวสริ, ตตฺร สุทํ ภควา เวสาลิยํ วิหรติ มหาวเน กูฏาคาร\-สาลายํ. อถ โข มหาปชาปติ โคตมี เกเส เฉทาเปตฺวา กาสายานิ วตฺถานิ อจฺฉาเทตฺวา สมฺพหุลาหิ สากิยานีหิ สทฺธิํ เยน เวสาลี เตน ปกฺกามิ. อนุปุพฺเพน เยน เวสาลี มหาวนํ กูฏาคารสาลา เตนุปสงฺกมิ. อถ โข มหาปชาปติ โคตมี สูเนหิ ปาเทหิ รโชกิณฺเณน คตฺเตน ทุกฺขี ทุมฺมนา อสฺสุมุขี รุทมานา พหิทฺวาร\-โกฏฺฐเก อฏฺฐาสิ, อทฺทสา โข อายสฺมา อานนฺโท มหาปชาปติํ โคตมิํ สูเนหิ ปาเทหิ รโชกิณฺเณน คตฺเตน ทุกฺขิํ ทุมฺมนํ อสฺสุมุขิํ รุทมานํ พหิทฺวาร\-โกฏฺฐเก ฐิตํ, ทิสฺวาน มหาปชาปติํ \csromanfolio{443}โคตมิํ เอตทโวจ “กิสฺส ตฺวํ โคตมิ สูเนหิ ปาเทหิ รโชกิณฺเณน คตฺเตน ทุกฺขี ทุมฺมนา อสฺสุมุขี รุทมานา พหิทฺวาร\-โกฏฺฐเก ฐิตา”ติ. ตถา หิ ปน ภนฺเต อานนฺท น ภควา อนุชานาติ มาตุคามสฺส ตถาคต\-ปฺปเวทิเต ธมฺมวินเย อคารสฺมา อนคาริยํ ปพฺพชฺชนฺติ. เตน หิ ตฺวํ โคตมิ มุหุตฺตํ อิเธว ตาว โหติ, ยาวาหํ ภควนฺตํ ยาจามิ มาตุคามสฺส ตถาคต\-ปฺปเวทิเต ธมฺมวินเย อคารสฺมา อนคาริยํ ปพฺพชฺชนฺติ.}

\prose{อถ โข อายสฺมา อานนฺโท เยน ภควา เตนุปสงฺกมิ, อุปสงฺกมิตฺวา ภควนฺตํ อภิวาเทตฺวา เอกมนฺตํ นิสีทิ, เอกมนฺตํ นิสินฺโน โข อายสฺมา อานนฺโท ภควนฺตํ เอตทโวจ “เอสา ภนฺเต มหาปชาปติ โคตมี สูเนหิ ปาเทหิ รโชกิณฺเณน คตฺเตน ทุกฺขี ทุมฺมนา อสฺสุมุขี รุทมานา พหิทฺวาร\-โกฏฺฐเก ฐิตา ‘น ภควา อนุชานาติ มาตุคามสฺส ตถาคต\-ปฺปเวทิเต ธมฺมวินเย อคารสฺมา อนคาริยํ ปพฺพชฺชนฺ’ติ, สาธุ ภนฺเต, ลเภยฺย มาตุคาโม ตถาคต\-ปฺปเวทิเต ธมฺมวินเย อคารสฺมา อนคาริยํ ปพฺพชฺชนฺ”ติ. อลํ อานนฺท, มา เต รุจฺจิ มาตุคามสฺส ตถาคต\-ปฺปเวทิเต ธมฺมวินเย อคารสฺมา อนคาริยํ ปพฺพชฺชาติ. ทุติยมฺปิ โข ฯเปฯ. ตติยมฺปิ โข อายสฺมา อานนฺโท ภควนฺตํ เอตทโวจ “สาธุ ภนฺเต, ลเภยฺย มาตุคาโม ตถาคต\-ปฺปเวทิเต ธมฺมวินเย อคารสฺมา อนคาริยํ ปพฺพชฺชนฺ”ติ. อลํ อานนฺท, มา เต รุจฺจิ มาตุคามสฺส ตถาคต\-ปฺปเวทิเต ธมฺมวินเย อคารสฺมา อนคาริยํ ปพฺพชฺชาติ.}

\prose{อถ โข อายสฺมา อานนฺโท\csromansharedfootnote{9ea42d48cc417d34}{อํ 3. 103 ปิฏฺเฐ ปน “อถ โข อายสฺมโต อานนฺทสฺส เอตทโหสี”ติ อาคตํ.} “น ภควา อนุชานาติ มาตุคามสฺส ตถาคต\-ปฺปเวทิเต ธมฺมวินเย อคารสฺมา อนคาริยํ ปพฺพชฺชํ, ยํนูนาหํ อญฺเญนปิ ปริยาเยน ภควนฺตํ ยาเจยฺยํ มาตุคามสฺส ตถาคต\-ปฺปเวทิเต ธมฺมวินเย อคารสฺมา อนคาริยํ ปพฺพชฺชนฺ”ติ. อถ โข อายสฺมา อานนฺโท ภควนฺตํ เอตทโวจ “ภพฺโพ นุ โข ภนฺเต มาตุคาโม ตถาคต\-ปฺปเวทิเต ธมฺมวินเย อคารสฺมา อนคาริยํ ปพฺพชิตฺวา โสตาปตฺติผลํ วา สกทาคามิผลํ\csromansharedfootnote{6481ffeb77302739}{สกิทาคามิผลํ (สฺยา)} วา อนาคามิผลํ วา อรหตฺตผลํ วา สจฺฉิกาตุนฺ”ติ. ภพฺโพ อานนฺท มาตุคาโม ตถาคต\-ปฺปเวทิเต ธมฺมวินเย อคารสฺมา อนคาริยํ ปพฺพชิตฺวา โสตาปตฺติ\-ผลมฺปิ สกทาคามิ\-ผลมฺปิ \csromanfolio{444}อนาคามิ\-ผลมฺปิ อรหตฺต\-ผลมฺปิ สจฺฉิกาตุนฺติ. สเจ ภนฺเต ภพฺโพ มาตุคาโม ตถาคต\-ปฺปเวทิเต ธมฺมวินเย อคารสฺมา อนคาริยํ ปพฺพชิตฺวา โสตาปตฺติ\-ผลมฺปิ สกทาคามิ\-ผลมฺปิ อนาคามิ\-ผลมฺปิ อรหตฺต\-ผลมฺปิ สจฺฉิกาตุํ, พหูปการา ภนฺเต มหาปชาปติ โคตมี ภควโต มาตุจฺฉา อาปาทิกา โปสิกา ขีรสฺส ทายิกา ภควนฺตํ ชเนตฺติยา กาลงฺกตาย ถญฺญํ ปาเยสิ, สาธุ ภนฺเต, ลเภยฺย มาตุคาโม ตถาคต\-ปฺปเวทิเต ธมฺมวินเย อคารสฺมา อนคาริยํ ปพฺพชฺชนฺติ.}

\csromanmatikaanchor{mtk.207}
\csromantocmarkref{subsection}{อฏฺฐครุธมฺม}{mtk.207}
\bookmark[dest={mtk.207},level=2]{อฏฺฐครุธมฺม}
\csromanheader{2}{อฏฺฐ\-ครุธมฺม}
\proseitem{403}{สเจ อานนฺท มหาปชาปติ โคตมี อฏฺฐ ครุธมฺเม ปฏิคฺคณฺหาติ, สาวสฺสา โหตุ อุปสมฺปทา.}

\prose{\csromansymbolfootnote{*}{วิ 2. 74 ปิฏฺเฐปิ.}วสฺสสตู\-ปสมฺปนฺนาย ภิกฺขุนิยา ตทหุ\-ปสมฺปนฺนสฺส ภิกฺขุโน อภิวาทนํ ปจฺจุฏฺฐานํ อญฺชลิกมฺมํ สามีจิกมฺมํ กาตพฺพํ, อยมฺปิ ธมฺโม สกฺกตฺวา ครุกตฺวา\csromansharedfootnote{d33b3b4b14afb7ee}{ครุํกตฺวา (ก)} มาเนตฺวา ปูเชตฺวา ยาวชีวํ อนติกฺกมนีโย. (1)}

\prose{น ภิกฺขุนิยา อภิกฺขุเก อาวาเส วสฺสํ วสิตพฺพํ, อยมฺปิ ธมฺโม สกฺกตฺวา ครุกตฺวา มาเนตฺวา ปูเชตฺวา ยาวชีวํ อนติกฺกมนีโย. (2)}

\prose{อนฺวทฺธมาสํ ภิกฺขุนิยา ภิกฺขุ\-สํฆโต เทฺว ธมฺมา ปจฺจาสีสิตพฺพา อุโปสถปุจฺฉกญฺจ โอวาทูปสงฺกมนญฺจ, อยมฺปิ ธมฺโม สกฺกตฺวา ครุกตฺวา มาเนตฺวา ปูเชตฺวา ยาวชีวํ อนติกฺกมนีโย. (3)}

\prose{วสฺสํวุฏฺฐาย ภิกฺขุนิยา อุภโตสํเฆ ตีหิ ฐาเนหิ ปวาเรตพฺพํ ทิฏฺเฐน วา สุเตน วา ปริสงฺกาย วา, อยมฺปิ ธมฺโม สกฺกตฺวา ครุกตฺวา มาเนตฺวา ปูเชตฺวา ยาวชีวํ อนติกฺกมนีโย. (4)}

\prose{ครุธมฺมํ อชฺฌาปนฺนาย ภิกฺขุนิยา อุภโตสํเฆ ปกฺขมานตฺตํ จริตพฺพํ, อยมฺปิ ธมฺโม สกฺกตฺวา ครุกตฺวา มาเนตฺวา ปูเชตฺวา ยาวชีวํ อนติกฺกมนีโย. (5)}

\prose{เทฺว วสฺสานิ ฉสุ ธมฺเมสุ สิกฺขิต\-สิกฺขาย สิกฺขมานาย อุภโตสํเฆ อุปสมฺปทา ปริเยสิตพฺพา, อยมฺปิ ธมฺโม สกฺกตฺวา ครุกตฺวา มาเนตฺวา ปูเชตฺวา ยาวชีวํ อนติกฺกมนีโย. (6)}

\prose{\csromanfolio{445}น ภิกฺขุนิยา เกนจิ ปริยาเยน ภิกฺขุ อกฺโกสิตพฺโพ ปริภาสิตพฺโพ, อยมฺปิ ธมฺโม สกฺกตฺวา ครุกตฺวา มาเนตฺวา ปูเชตฺวา ยาวชีวํ อนติกฺกมนีโย. (7)}

\prose{อชฺชตคฺเค โอวโฏ ภิกฺขุนีนํ ภิกฺขูสุ วจนปโถ, อโนวโฏ ภิกฺขูนํ ภิกฺขุนีสุ วจนปโถ, อยมฺปิ ธมฺโม สกฺกตฺวา ครุกตฺวา มาเนตฺวา ปูเชตฺวา ยาวชีวํ อนติกฺกมนีโย. (8)}

\prose{สเจ อานนฺท มหาปชาปติ โคตมี อิเม อฏฺฐ ครุธมฺเม ปฏิคฺคณฺหาติ, สาวสฺสา โหตุ อุปสมฺปทาติ.}

\prose{อถ โข อายสฺมา อนนฺโท ภควโต สนฺติเก อฏฺฐ ครุธมฺเม อุคฺคเหตฺวา เยน มหาปชาปติ โคตมี เตนุปสงฺกมิ, อุปสงฺกมิตฺวา มหาปชาปติํ โคตมิํ เอตทโวจ “สเจ โข ตฺวํ โคตมิ อฏฺฐ ครุธมฺเม ปฏิคฺคเณฺหยฺยาสิ, สาว เต ภวิสฺสติ อุปสมฺปทา.}

\prose{วสฺสสตู\-ปสมฺปนฺนาย ภิกฺขุนิยา ตทหุ\-ปสมฺปนฺนสฺส ภิกฺขุโน อภิวาทนํ ปจฺจุฏฺฐานํ อญฺชลิกมฺมํ สามีจิกมฺมํ กาตพฺพํ, อยมฺปิ ธมฺโม สกฺกตฺวา ครุกตฺวา มาเนตฺวา ปูเชตฺวา ยาวชีวํ อนติกฺกมนีโย.}

\prose{น ภิกฺขุนิยา อภิกฺขุเก อาวาเส วสฺสํ วสิตพฺพํ, อยมฺปิ ธมฺโม สกฺกตฺวา ครุกตฺวา มาเนตฺวา ปูเชตฺวา ยาวชีวํ อนติกฺกมนีโย.}

\prose{อนฺวทฺธมาสํ ภิกฺขุนิยา ภิกฺขุ\-สํฆโต เทฺว ธมฺมา ปจฺจาสีสิตพฺพา อุโปสถปุจฺฉกญฺจ โอวาทูปสงฺกมนญฺจ, อยมฺปิ ธมฺโม สกฺกตฺวา ครุกตฺวา มาเนตฺวา ปูเชตฺวา ยาวชีวํ อนติกฺกมนีโย.}

\prose{วสฺสํวุฏฺฐาย ภิกฺขุนิยา อุภโตสํเฆ ตีหิ ฐาเนหิ ปวาเรตพฺพํ ทิฏฺเฐน วา สุเตน วา ปริสงฺกาย วา, อยมฺปิ ธมฺโม สกฺกตฺวา ครุกตฺวา มาเนตฺวา ปูเชตฺวา ยาวชีวํ อนติกฺกมนีโย.}

\prose{ครุธมฺมํ อชฺฌาปนฺนาย ภิกฺขุนิยา อุภโตสํเฆ ปกฺขมานตฺตํ จริตพฺพํ, อยมฺปิ ธมฺโม สกฺกตฺวา ครุกตฺวา มาเนตฺวา ปูเชตฺวา ยาวชีวํ อนติกฺกมนีโย.}

\prose{เทฺว วสฺสานิ ฉสุ ธมฺเมสุ สิกฺขิต\-สิกฺขาย สิกฺขมานาส อุภโตสํเฆ อุปสมฺปทา ปริเยสิตพฺพา, อยมฺปิ ธมฺโม สกฺกตฺวา ครุกตฺวา มาเนตฺวา ปูเชตฺวา ยาวชีวํ อนติกฺกมนีโย.}

\prose{\csromanfolio{446}น ภิกฺขุนิยา เกนจิ ปริยาเยน ภิกฺขุ อกฺโกสิตพฺโพ ปริภาสิตพฺโพ, อยมฺปิ ธมฺโม สกฺกตฺวา ครุกตฺวา มาเนตฺวา ปูเชตฺวา ยาวชีวํ อนติกฺกมนิโย.}

\prose{อชฺชตคฺเค โอวโฏ ภิกฺขุนีนํ ภิกฺขูสุ วจนปโถ, อโนวโฏ ภิกฺขูนํ ภิกฺขุนีสุ วจนปโถ, อยมฺปิ ธมฺโม สกฺกตฺวา ครุกตฺวา มาเนตฺวา ปูเชตฺวา ยาวชีวํ อนติกฺกมนีโย.}

\prose{สเจ โข ตฺวํ โคตมิ อิเม อฏฺฐ ครุธมฺเม ปฏิคฺคเณฺหยฺยาสิ, สาว เต ภวิสฺสติ อุปสมฺปทา”ติ.}

\prose{เสยฺยถาปิ ภนฺเต อานนฺท อิตฺถี วา ปุริโส วา ทหโร ยุวา มณฺฑนก\-ชาติโก สีสํนหาโต อุปฺปลมาลํ วา วสฺสิกมาลํ วา อติมุตฺตก\-มาลํ\csromansharedfootnote{9f434821ea5fcd2f}{อธิมตฺตกมาลํ (สฺยา)} วา ลภิตฺวา อุโภหิ หตฺเถหิ ปฏิคฺคเหตฺวา อุตฺตมงฺเค สิรสฺมิํ ปติฏฺฐาเปยฺย, เอวเมว โข อหํ ภนฺเต อานนฺท อิเม อฏฺฐ ครุธมฺเม ปฏิคฺคณฺหามิ ยาวชีวํ อนติกฺกมนีเยติ.}

\prose{อถ โข อายสฺมา อานนฺโท เยน ภควา เตนุปสงฺกมิ, อุปสงฺกมิตฺวา ภควนฺตํ อภิวาเทตฺวา เอกมนฺตํ นิสีทิ, เอกมนฺตํ นิสินฺโน โข อายสฺมา อานนฺโท ภควนฺตํ เอตทโวจ “ปฏิคฺคหิตา ภนฺเต มหาปชาปติยา โคตมิยา อฏฺฐ ครุธมฺมา, อุปสมฺปนฺนา ภควโต มาตุจฺฉา”ติ.}

\prose{สเจ อานนฺท นาลภิสฺส มาตุคาโม ตถาคต\-ปฺปเวทิเต ธมฺมวินเย อคารสฺมา อนคาริยํ ปพฺพชฺชํ, จิรฏฺฐิติกํ อานนฺท พฺรหฺมจริยํ อภวิสฺส, วสฺสสหสฺสํ สทฺธมฺโม ติฏฺเฐยฺย. ยโต จ โข อานนฺท มาตุคาโม ตถาคต\-ปฺปเวทิเต ธมฺมวินเย อคารสฺมา อนคาริยํ ปพฺพชิโต, น ทานิ อานนฺท พฺรหฺมจริยํ จิรฏฺฐิติกํ ภวิสฺสติ, ปญฺเจว ทานิ อานนฺท วสฺสสตานิ สทฺธมฺโม ฐสฺสติ.}

\prose{เสยฺยถาปิ อานนฺท ยานิ กานิจิ กุลานิ พหุตฺถิกานิ\csromansharedfootnote{1b9a10bf236aceea}{พหุอิตฺถิกานิ (สี, สฺยา)} อปฺปปุริสกานิ, ตานิ สุปฺปธํสิยานิ โหนฺติ โจเรหิ กุมฺภเถนเกหิ, เอวเมว โข อานนฺท ยสฺมิํ ธมฺมวินเย ลภติ มาตุคาโม อคารสฺมา อนคาริยํ ปพฺพชฺชํ, น ตํ พฺรหฺมจริยํ จิรฏฺฐิติกํ โหติ.}

\prose{\csromanfolio{447}เสยฺยถาปิ อานนฺท สมฺปนฺเน สาลิกฺเขตฺเต เสตฏฺฏิกา นาม โรคชาติ นิปตติ, เอวํ ตํ สาลิกฺเขตฺตํ น จิรฏฺฐิติกํ โหติ, เอวเมว โข อานนฺท ยสฺมิํ ธมฺมวินเย ลภติ มาตุคาโม อคารสฺมา อนคาริยํ ปพฺพชฺชํ, น ตํ พฺรหฺมจริยํ จิรฏฺฐิติกํ โหติ.}

\prose{เสยฺยถาปิ อานนฺท สมฺปนฺเน อุจฺฉุกฺเขตฺเต มญฺชิฏฺฐิกา\csromansharedfootnote{318f14f13d7966c5}{มญฺเชฏฺฐิกา (สี, สฺยา)} นาม โรคชาติ นิปตติ, เอวํ ตํ อุจฺฉุกฺเขตฺตํ น จิรฏฺฐิติกํ โหติ, เอวเมว โข อานนฺท ยสฺมิํ ธมฺมวินเย ลภติ มาตุคาโม อคารสฺมา อนคาริยํ ปพฺพชฺชํ, น ตํ พฺรหฺมจริยํ จิรฏฺฐิติกํ โหติ.}

\prosewithcloser{เสยฺยถาปิ อานนฺท ปุริโส มหโต ตฬากสฺส ปฏิกจฺเจว อาฬิํ พนฺเธยฺย ยาวเทว อุทกสฺส อนติกฺกมนาย, เอวเมว โข อานนฺท มยา ปฏิกจฺเจว ภิกฺขุนีนํ อฏฺฐ ครุธมฺมา ปญฺญตฺตา ยาวชีวํ อนติกฺกมนียาติ.}{ภิกฺขุนีนํ อฏฺฐ ครุธมฺมา นิฏฺฐิตา.}
\csromanmatikaanchor{mtk.208}
\csromantocmarkref{subsection}{ภิกฺขุนีอุปสมฺปทานุชานน}{mtk.208}
\bookmark[dest={mtk.208},level=2]{ภิกฺขุนีอุปสมฺปทานุชานน}
\csromanheader{2}{ภิกฺขุนี\-อุปสมฺปทา\-นุชานน}
\proseitem{404}{อถ โข มหาปชาปติ โคตมี เยน ภควา เตนุปสงฺกมิ, อุปสงฺกมิตฺวา ภควนฺตํ อภิวาเทตฺวา เอกมนฺตํ อฏฺฐาสิ, เอกมนฺตํ ฐิตา โข มหาปชาปติ โคตมี ภควนฺตํ เอตทโวจ “กถาหํ ภนฺเต อิมาสุ สากิยานีสุ ปฏิปชฺชามี”ติ. อถ โข ภควา มหาปชาปติํ โคตมิํ ธมฺมิยา กถาย สนฺทสฺเสสิ สมาทเปสิ สมุตฺเตเชสิ สมฺปหํเสสิ. อถ โข มหาปชาปติ โคตมี ภควตา ธมฺมิยา กถาย สนฺทสฺสิตา สมาทปิตา สมุตฺเตชิตา สมฺปหํสิตา ภควนฺตํ อภิวาเทตฺวา ปทกฺขิณํ กตฺวา ปกฺกามิ. อถ โข ภควา เอตสฺมิํ นิทาเน เอตสฺมิํ ปกรเณ ธมฺมิํ กถํ กตฺวา ภิกฺขู อามนฺเตสิ “อนุชานามิ ภิกฺขเว ภิกฺขูหิ ภิกฺขุนิโย อุปสมฺปาเทตุนฺ”ติ.}

\prose{อถ โข ตา ภิกฺขุนิโย มหาปชาปติํ โคตมิํ เอตทโวจุํ “อยฺยา อนุปสมฺปนฺนา, มยญฺจมฺหา อุปสมฺปนฺนา, เอวํ หิ ภควตา ปญฺญตฺตํ, ‘ภิกฺขูหิ ภิกฺขุนิโย อุปสมฺ\-ปาเท\-ตพฺพา’ติ”. อถ โข มหาปชาปติ \csromanfolio{448}โคตมี เยนายสฺมา อานนฺโท เตนุปสงฺกมิ, อุปสงฺกมิตฺวา อายสฺมนฺตํ อานนฺทํ อภิวาเทตฺวา เอกมนฺตํ อฏฺฐาสิ, เอกมนฺตํ ฐิตา โข มหาปชาปติ โคตมี อายสฺมนฺตํ อานนฺทํ เอตทโวจ “อิมา มํ ภนฺเต อานนฺท ภิกฺขุนิโย เอวมาหํสุ ‘อยฺยา อนุปสมฺปนฺนา, มยญฺจมฺหา อุปสมฺปนฺนา, เอวํ หิ ภควตา ปญฺญตฺตํ ภิกฺขูหิ ภิกฺขุนิโย อุปสมฺ\-ปาเท\-ตพฺพา’ติ”.}

\prose{อถ โข อายสฺมา อานนฺโท เยน ภควา เตนุปสงฺกมิ, อุปสงฺกมิตฺวา ภควนฺตํ อภิวาเทตฺวา เอกมนฺตํ นิสีทิ, เอกมนฺตํ นิสินฺโน โข อายสฺมา อานนฺโท ภควนฺตํ เอตทโวจ “มหาปชาปติ ภนฺเต โคตมี เอวมาห ‘อิมา มํ ภนฺเต อานนฺท ภิกฺขุนิโย เอวมาหํสุ อยฺยา อนุปสมฺปนฺนา, มยญฺจมฺหา อุปสมฺปนฺนา, เอวํ หิ ภควตา ปญฺญตฺตํ ภิกฺขูหิ ภิกฺขุนิโย อุปสมฺ\-ปาเท\-ตพฺพา’ติ”.}

\prose{ยทคฺเคน อานนฺท มหาปชาปติยา โคตมิยา อฏฺฐ ครุธมฺมา ปฏิคฺคหิตา, ตเทว สา\csromansharedfootnote{608bd2a709a4b611}{ตเทวสฺสา (ก)} อุปสมฺปนฺนาติ.}

\proseitem{405}{อถ โข มหาปชาปติ โคตมี เยนายสฺมา อานนฺโท เตนุปสงฺกมิ, อุปสงฺกมิตฺวา อายสฺมนฺตํ อานนฺทํ อภิวาเทตฺวา เอกมนฺตํ อฏฺฐาสิ, เอกมนฺตํ ฐิตา โข มหาปชาปติ โคตมี อายสฺมนฺตํ อานนฺทํ เอตทโวจ “เอกาหํ ภนฺเต อานนฺท ภควนฺตํ วรํ ยาจามิ, สาธุ ภนฺเต, ภควา อนุชาเนยฺย ภิกฺขูนญฺจ ภิกฺขุนีนญฺจ ยถาวุฑฺฒํ อภิวาทนํ ปจฺจุฏฺฐานํ อญฺชลิกมฺมํ สามีจิกมฺมนฺ”ติ.}

\prose{อถ โข อายสฺมา อานนฺโท เยน ภควา เตนุปสงฺกมิ, อุปสงฺกมิตฺวา ภควนฺตํ อภิวาเทตฺวา เอกมนฺตํ นิสีทิ, เอกมนฺตํ นิสินฺโน โข อายสฺมา อานนฺโท ภควนฺตํ เอตทโวจ “มหาปชาปติ ภนฺเต โคตมี เอวมาห ‘เอกาหํ ภนฺเต อานนฺท ภควนฺตํ วรํ ยาจามิ, สาธุ ภนฺเต ภควา อนุชาเนยฺย ภิกฺขูนญฺจ ภิกฺขุนีนญฺจ ยถาวุฑฺฒํ อภิวาทนํ ปจฺจุฏฺฐานํ อญฺชลิกมฺมํ สามีจิกมฺมนฺ’ติ”.}

\prose{อฏฺฐานเมตํ อานนฺท อนวกาโส, ยํ ตถาคโต อนุชาเนยฺย มาตุคามสฺส อภิวาทนํ ปจฺจุฏฺฐานํ อญฺชลิกมฺมํ สามีจิกมฺมํ, อิเมหิ นาม อานนฺท อญฺญติตฺถิยา ทุรกฺขาต\-ธมฺมา มาตุคามสฺส อภิวาทนํ ปจฺจุฏฺฐานํ \csromanfolio{449}อญฺชลิกมฺมํ สามีจิกมฺมํ น กริสฺสนฺติ, กิมงฺคํ ปน ตถาคโต อนุชานิสฺสติ มาตุคามสฺส อภิวาทนํ ปจฺจุฏฺฐานํ อญฺชลิกมฺมํ สามีจิกมฺมนฺติ.}

\prose{อถ โข ภควา เอตสฺมิํ นิทาเน เอตสฺมิํ ปกรเณ ธมฺมิํ กถํ กตฺวา ภิกฺขู อามนฺเตสิ “น ภิกฺขเว มาตุคามสฺส อภิวาทนํ ปจฺจุฏฺฐานํ อญฺชลิกมฺมํ สามีจิกมฺมํ กาตพฺพํ, โย กเรยฺย, อาปตฺติ ทุกฺกฏสฺสา”ติ.}

\prose{อถ โข มหาปชาปติ โคตมี เยน ภควา เตนุปสงฺกมิ อุปสงฺกมิตฺวา ภควนฺตํ อภิวาเทตฺวา เอกมนฺตํ อฏฺฐาสิ, เอกมนฺตํ ฐิตา โข มหาปชาปติ โคตมี ภควนฺตํ เอตทโวจ “ยานิ ตานิ ภนฺเต ภิกฺขุนีนํ สิกฺขาปทานิ ภิกฺขูหิ สาธารณานิ, กถํ มยํ ภนฺเต เตสุ สิกฺขาปเทสุ ปฏิปชฺชามา”ติ. ยานิ ตานิ โคตมิ ภิกฺขุนีนํ สิกฺขาปทานิ ภิกฺขูหิ สาธารณานิ, ยถา ภิกฺขู สิกฺขนฺติ ตถา เตสุ สิกฺขาปเทสุ สิกฺขถาติ. ยานิ ปน ตานิ ภนฺเต ภิกฺขุนีนํ สิกฺขาปทานิ ภิกฺขูหิ อสาธารณานิ, กถํ มยํ ภนฺเต เตสุ สิกฺขาปเทสุ ปฏิปชฺชามาติ. ยานิ ตานิ โคตมิ ภิกฺขุนีนํ สิกฺขาปทานิ ภิกฺขูหิ อสาธรณานิ, ยถา\-ปญฺญตฺเตสุ สิกฺขาปเทสุ สิกฺขถาติ.}

\proseitem{406}{\csromansymbolfootnote{*}{อํ 3. 106 ปิฏฺเฐปิ.}อถ โข มหาปชาปติ โคตมี เยน ภควา เตนุปสงฺกมิ, อุปสงฺกมิตฺวา ภควนฺตํ อภิวาเทตฺวา เอกมนฺตํ อฏฺฐาสิ, เอกมนฺตํ ฐิตา โข มหาปชาปติ โคตมี ภควนฺตํ เอตทโวจ “สาธุ เม ภนฺเต ภควา สํขิตฺเตน ธมฺมํ เทเสตุ, ยมหํ ภควโต ธมฺมํ สุตฺวา เอกา วูปกฏฺฐา อปฺปมตฺตา อาตาปินี ปหิตตฺตา วิหเรยฺยนฺ”ติ. เย โข ตฺวํ โคตมิ ธมฺเม ชาเนยฺยาสิ, อิเม ธมฺมา สราคาย สํวตฺตนฺติ, โน วิราคาย, สญฺโญคาย สํวตฺตนฺติ, โน วิสญฺโญคาย, อาจยาย สํวตฺตนฺติ, โน อปจยาย, มหิจฺฉตาย สํวตฺตนฺติ, โน อปฺปิจฺฉตาย, อสนฺตุฏฺฐิยา สํวตฺตนฺติ, โน สนฺตุฏฺฐิยา, สงฺคณิกาย สํวตฺตนฺติ, โน ปวิเวกาย, โกสชฺชาย สํวตฺตนิ, โน วีริยารมฺภาย, ทุพฺภรตาย สํวตฺตนฺติ, โน สุภรตาย, เอกํเสน โคตมิ ธาเรยฺยาสิ “เนโส ธมฺโม, เนโส วินโย, เนตํ สตฺถุสาสนนฺ”ติ. เย จ โข ตฺวํ โคตมิ ธมฺเม ชาเนยฺยาสิ, อิเม ธมฺมา วิราคาย สํวตฺตนฺติ, โน สราคาย, วิสญฺโญคาย สํวตฺตนฺติ, โน สญฺโญคาย, อปจยาย สํวตฺตนฺติ, โน อาจยาย, อปฺปิจฺฉตาย สํวตฺตนฺติ, \csromanfolio{450}โน มหิจฺฉตาย, สนฺตุฏฺฐิยา สํวตฺตนฺติ, โน อสนฺตุฏฺฐิยา, ปวิเวกาย สํวตฺตนฺติ, โน สงฺคณิกาย, วีริยารมฺภาย สํวตฺตนฺติ, โน โกสชฺชาย, สุภรตาย สํวตฺตนฺติ, โน ทุพฺภรตาย, เอกํเสน โคตมิ ธาเรยฺยาสิ “เอโส ธมฺโม, เอโส วินโย, เอตํ สตฺถุสาสนนฺ”ติ.}

\proseitem{407}{เตน โข ปน สมเยน ภิกฺขุนีนํ ปาติโมกฺขํ น อุทฺทิสียติ. ภควโต เอตมตฺถํ อาโรเจสุํ ฯเปฯ. อนุชานามิ ภิกฺขเว ภิกฺขุนีนํ ปาติโมกฺขํ อุทฺทิสิตุนฺติ. อถ โข ภิกฺขูนํ เอตทโหสิ “เกน นุ โข ภิกฺขุนีนํ ปาติโมกฺขํ อุทฺทิสิตพฺพนฺ”ติ. ภควโต เอตมตฺถํ อาโรเจสุํ. อนุชานามิ ภิกฺขเว ภิกฺขูหิ ภิกฺขุนีนํ ปาติโมกฺขํ อุทฺทิสิตุนฺติ.}

\prose{เตน โข ปน สมเยน ภิกฺขู ภิกฺขุนุ\-ปสฺสยํ อุปสงฺกมิตฺวา ภิกฺขุนีนํ ปาติโมกฺขํ อุทฺทิสนฺติ. มนุสฺสา อุชฺฌายนฺติ ขิยฺยนฺติ วิปาเจนฺติ “ชายาโย อิมา อิเมสํ, ชาริโย อิมา อิเมสํ, อิทานิ อิเม อิมาหิ สทฺธิํ อภิรมิสฺสนฺตี”ติ. ภควโต เอตมตฺถํ อาโรเจสุํ. น ภิกฺขเว ภิกฺขูหิ ภิกฺขุนีนํ ปาติโมกฺขํ อุทฺทิสิตพฺพํ, โย อุทฺทิเสยฺย, อาปตฺติ ทุกฺกฏสฺส. อนุชานามิ ภิกฺขเว ภิกฺขุนีหิ ภิกฺขุนีนํ ปาติโมกฺขํ อุทฺทิสิตุนฺติ. ภิกฺขุนิโย น ชานนฺติ “เอวํ ปาติโมกฺขํ อุทฺทิสิตพฺพนฺ”ติ. ภควโต เอตมตฺถํ อาโรเจสุํ. อนุชานามิ ภิกฺขเว ภิกฺขูหิ ภิกฺขุนีนํ อาจิกฺขิตุํ “เอวํ ปาติโมกฺขํ อุทฺทิเสยฺยาถา”ติ.}

\proseitem{408}{เตน โข ปน สมเยน ภิกฺขุนิโย อาปตฺติํ น ปฏิกโรนฺติ. ภควโต เอตมตฺถํ อาโรเจสุํ. น ภิกฺขเว ภิกฺขุนิยา อาปตฺติ น ปฏิกาตพฺพา, ยา น ปฏิกเรยฺย, อาปตฺติ ทุกฺกฏสฺสาติ. ภิกฺขุนิโย น ชานนฺติ “เอวํ อาปตฺติ ปฏิกาตพฺพา”ติ. ภควโต เอตมตฺถํ อาโรเจสุํ. อนุชานามิ ภิกฺขเว ภิกฺขูหิ ภิกฺขุนีนํ อาจิกฺขิตุํ “เอวํ อาปตฺติํ ปฏิกเรยฺยาถา”ติ. อถ โข ภิกฺขูนํ เอตทโหสิ “เกน นุ โข ภิกฺขุนีนํ อาปตฺติ ปฏิคฺคเหตพฺพา”ติ. ภควโต เอตมตฺถํ อาโรเจสุํ. อนุชานามิ ภิกฺขเว ภิกฺขูหิ ภิกฺขุนีนํ อาปตฺติํ ปฏิคฺคเหตุนฺติ.}

\prose{เตน โข ปน สมเยน ภิกฺขุนิโย รถิกายปิ พฺยูเหปิ สิงฺฆาฏเกปิ ภิกฺขุํ ปสฺสิตฺวา ปตฺตํ ภูมิยํ นิกฺขิปิตฺวา เอกํสํ อุตฺตราสงฺคํ กริตฺวา \csromanfolio{451}อุกฺกุฏิกํ นิสีทิตฺวา อญฺชลิํ ปคฺคเหตฺวา อาปตฺติํ ปฏิกโรนฺติ. มนุสฺสา อุชฺฌายนฺติ ขิยฺยนฺติ วิปาเจนฺติ “ชายาโย อิมา อิเมสํ, ชาริโย อิมา อิเมสํ, รตฺติํ วิมาเนตฺวา อิทานิ ขมาเปนฺตี”ติ. ภควโต เอตมตฺถํ อาโรเจสุํ. น ภิกฺขเว ภิกฺขูหิ ภิกฺขุนีนํ อาปตฺติ ปฏิคฺคเหตพฺพา, โย ปฏิคฺคเณฺหยฺย, อาปตฺติ ทุกฺกฏสฺส. อนุชานามิ ภิกฺขเว ภิกฺขุนีหิ ภิกฺขุนีนํ อาปตฺติํ ปฏิคฺคเหตุนฺติ. ภิกฺขุนิโย น ชานนฺติ “เอวํ อาปตฺติ ปฏิคฺคเหตพฺพา”ติ. ภควโต เอตมตฺถํ อาโรเจสุํ. อนุชานามิ ภิกฺขเว ภิกฺขูหิ ภิกฺขุนีนํ อาจิกฺขิตุํ “เอวํ อาปตฺติํ ปฏิคฺคเณฺหยฺยาถา”ติ.}

\proseitem{409}{เตน โข ปน สมเยน ภิกฺขุนีนํ กมฺมํ น กริยติ. ภควโต เอตมตฺถํ อาโรเจสุํ. อนุชานามิ ภิกฺขเว ภิกฺขุนีนํ กมฺมํ กาตุนฺติ. อถ โข ภิกฺขูนํ เอตทโหสิ “เกน นุ โข ภิกฺขุนีนํ กมฺมํ กาตพฺพนฺ”ติ. ภควโต เอตมตฺถํ อาโรเจสุํ. อนุชานามิ ภิกฺขเว ภิกฺขูหิ ภิกฺขุนีนํ กมฺมํ กาตุนฺติ.}

\prose{เตน โข ปน สมเยน กตกมฺมา ภิกฺขุนิโย รถิกายปิ พฺยูเหปิ สิงฺฆาฏเกปิ ภิกฺขุํ ปสฺสิตฺวา ปตฺตํ ภูมิยํ นิกฺขิปิตฺวา เอกํสํ อุตฺตราสงฺคํ กริตฺวา อุกฺกุฏิกํ นิสีทิตฺวา อญฺชลิํ ปคฺคเหตฺวา ขมาเปนฺติ “เอวํ นูน กาตพฺพนฺ”ติ มญฺญมานา. มนุสฺสา ตเถว อุชฺฌายนฺติ ขิยฺยนฺติ วิปาเจนฺติ “ชายาโย อิมา อิเมสํ, ชาริโย อิมา อิเมสํ, รตฺติํ วิมาเนตฺวา อิทานิ ขมาเปนฺตี”ติ. ภควโต เอตมตฺถํ อาโรเจสุํ. น ภิกฺขเว ภิกฺขูหิ ภิกฺขุนีนํ กมฺมํ กาตพฺพํ, โย กเรยฺย, อาปตฺติ ทุกฺกฏสฺส. อนุชานามิ ภิกฺขเว ภิกฺขุนีหิ ภิกฺขุนีนํ กมฺมํ กาตุนฺติ. ภิกฺขุนิโย น ชานนฺติ “เอวํ กมฺมํ กาตพฺพนฺ”ติ. ภควโต เอตมตฺถํ อาโรเจสุํ. อนุชานามิ ภิกฺขเว ภิกฺขูหิ ภิกฺขุนีนํ อาจิกฺขิตุํ “เอวํ กมฺมํ กเรยฺยาถา”ติ.}

\proseitem{410}{เตน โข ปน สมเยน ภิกฺขุนิโย สํฆมชฺเฌ ภณฺฑนชาตา กลหชาตา วิวาทาปนฺนา อญฺญมญฺญํ มุขสตฺตีหิ วิตุทนฺตา วิหรนฺติ, น สกฺโกนฺติ ตํ อธิกรณํ วูปสเมตุํ. ภควโต เอตมตฺถํ อาโรเจสุํ. อนุชานามิ ภิกฺขเว ภิกฺขูหิ ภิกฺขุนีนํ อธิกรณํ วูปสเมตุนฺติ.}

\prose{\csromanfolio{452}เตน โข ปน สมเยน ภิกฺขู ภิกฺขุนีนํ อธิกรณํ วูปสเมนฺติ, ตสฺมิํ โข ปน อธิกรเณ วินิจฺฉิยมาเน ทิสฺสนฺติ ภิกฺขุนิโย กมฺม\-ปฺปตฺตาโยปิ อาปตฺติคามินิโยปิ, ภิกฺขุนิโย เอวมาหํสุ “สาธุ ภนฺเต, อยฺยาว ภิกฺขุนีนํ กมฺมํ กโรนฺตุ, อยฺยาว ภิกฺขุนีนํ อาปตฺติํ ปฏิคฺคณฺหนฺตุ, เอวํ หิ ภควตา ปญฺญตฺตํ ภิกฺขูหิ ภิกฺขุนีนํ อธิกรณํ วูปสเมตพฺพนฺ”ติ. ภควโต เอตมตฺถํ อาโรเจสุํ. อนุชานามิ ภิกฺขเว ภิกฺขูหิ ภิกฺขุนีนํ กมฺมํ อาโรเปตฺวา ภิกฺขุนีนํ นิยฺยาเทตุํ, ภิกฺขุนีหิ ภิกฺขุนีนํ กมฺมํ กาตุํ, ภิกฺขูหิ ภิกฺขุนีนํ อาปตฺติํ อาโรเปตฺวา ภิกฺขุนีนํ นิยฺยาเทตุํ, ภิกฺขุนีนํ อาปตฺติํ ปฏิคฺคเหตุนฺติ.}

\prosewithcloser{เตน โข ปน สมเยน อุปฺปลวณฺณาย ภิกฺขุนิยา อนฺเตวาสินี ภิกฺขุนี สตฺต วสฺสานิ ภควนฺตํ อนุพนฺธา โหติ วินยํ ปริยาปุณนฺตี, ตสฺสา มุฏฺฐ\-สฺสติ\-นิยา คหิโต คหิโต มุสฺสติ. อสฺโสสิ โข สา ภิกฺขุนี “ภควา กิร สาวตฺถิํ คนฺตุกาโม”ติ. อถ โข ตสฺสา ภิกฺขุนิยา เอตทโหสี “อหํ โข สตฺต วสฺสานิ ภควนฺตํ อนุพนฺธิํ วินยํ ปริยาปุณนฺตี, ตสฺสา เม มุฏฺฐ\-สฺสติ\-นิยา คหิโต คหิโต มุสฺสติ, ทุกฺกรํ โข ปน มาตุคาเมน ยาวชีวํ สตฺถารํ อนุพนฺธิตุํ, กถํ นุ โข มยา ปฏิปชฺชิตพฺพนฺ”ติ. อถ โข สา ภิกฺขุนี ภิกฺขุนีนํ เอตมตฺถํ อาโรเจสิ, ภิกฺขุนิโย ภิกฺขูนํ เอตมตฺถํ อาโรเจสุํ, ภิกฺขู ภควโต เอตมตฺถํ อาโรเจสุํ. อนุชานามิ ภิกฺขเว ภิกฺขูหิ ภิกฺขุนีนํ วินยํ วาเจตุนฺติ.}{ปฐม\-ภาณวาโร นิฏฺฐิโต.}
\csromanmatikaanchor{mtk.209}
\csromantocmarkref{section}{2. ทุติยภาณวาร}{mtk.209}
\bookmark[dest={mtk.209},level=1]{2. ทุติยภาณวาร}
\csromanheader{1}{2. ทุติย\-ภาณวาร}
\proseitem{411}{อถ โข ภควา เวสาลิยํ ยถาภิรนฺตํ วิหริตฺวา เยน สาวตฺถิ เตน จาริกํ ปกฺกามิ, อนุปุพฺเพน จาริกํ จรมาโน เยน สาวตฺถิ ตทวสริ, ตตฺร สุทํ ภควา สาวตฺถิยํ วิหรติ เชตวเน อนาถ\-ปิณฺฑิกสฺส อาราเม. เตน โข ปน สมเยน ฉพฺพคฺคิยา ภิกฺขู ภิกฺขุนิโย กทฺทโมทเกน โอสิญฺจนฺติ “อปฺเปว นาม อเมฺหสุ สารชฺเชยฺยุนฺ”ติ. ภควโต เอตมตฺถํ อาโรเจสุํ. น ภิกฺขเว \csromanfolio{453}ภิกฺขุนา ภิกฺขุนิโย กทฺทโมทเกน โอสิญฺจิตพฺพา, โย โอสิญฺเจยฺย, อาปตฺติ ทุกฺกฏสฺส. อนุชานามิ ภิกฺขเว ตสฺส ภิกฺขุโน ทณฺฑกมฺมํ กาตุนฺติ. อถ โข ภิกฺขูนํ เอตทโหสิ “กิํ นุ โข ทณฺฑกมฺมํ กาตพฺพนฺ”ติ. ภควโต เอตมตฺถํ อาโรเจสุํ. อวนฺทิโย โส ภิกฺขเว ภิกฺขุ ภิกฺขุนิ\-สํเฆน กาตพฺโพติ.}

\prose{เตน โข ปน สมเยน ฉพฺพคฺคิยา ภิกฺขู กายํ วิวริตฺวา ภิกฺขุนีนํ ทสฺเสนฺติ ฯเปฯ. อูรุํ วิวริตฺวา ภิกฺขุนีนํ ทสฺเสนฺติ. องฺคชาตํ วิวริตฺวา ภิกฺขุนีนํ ทสฺเสนฺติ. ภิกฺขุนิโย โอภาเสนฺติ. ภิกฺขุนีหิ สทฺธิํ สมฺปโยเชนฺติ “อปฺเปว นาม อเมฺหสุ สารชฺเชยฺยุนฺ”ติ. ภควโต เอตมตฺถํ อาโรเจสุํ. น ภิกฺขเว ภิกฺขุนา กาโย วิวริตฺวา ภิกฺขุนีนํ ทสฺเสตพฺโพ ฯเปฯ. น อูรุ วิวริตฺวา ภิกฺขุนีนํ ทสฺเสตพฺโพ. น องฺคชาตํ วิวริตฺวา ภิกฺขุนีนํ ทสฺเสตพฺพํ. น ภิกฺขุนิโย โอภาสิตพฺพา. น ภิกฺขุนีหิ สทฺธิํ สมฺป\-โยเช\-ตพฺพํ, โย สมฺปโยเชยฺย, อาปตฺติ ทุกฺกฏสฺส. อนุชานามิ ภิกฺขเว ตสฺส ภิกฺขุโน ทณฺฑกมฺมํ กาตุนฺติ. อถ โข ภิกฺขูนํ เอตทโหสิ “กิํ นุ โข ทณฺฑกมฺมํ กาตพฺพนฺ”ติ. ภควโต เอตมตฺถํ อาโรเจสุํ. อวนฺทิโย โส ภิกฺขเว ภิกฺขุ ภิกฺขุนิ\-สํเฆน กาตพฺโพติ.}

\prose{เตน โข ปน สมเยน ฉพฺพคฺคิยา ภิกฺขุนิโย ภิกฺขุํ กทฺทโมทเกน โอสิญฺจนฺติ “อปฺเปว นาม อเมฺหสุ สารชฺเชยฺยุนฺ”ติ. ภควโต เอตมตฺถํ อาโรเจสุํ. น ภิกฺขเว ภิกฺขุนิยา ภิกฺขุ กทฺทโมทเกน โอสิญฺจิตพฺโพ, ยา โอสิญฺเจยฺย, อาปตฺติ ทุกฺกฏสฺส. อนุชานามิ ภิกฺขเว ตสฺสา ภิกฺขุนิยา ทณฺฑกมฺมํ กาตุนฺติ. อถ โข ภิกฺขูนํ เอตทโหสิ “กิํ นุ โข ทณฺฑกมฺมํ กาตพฺพนฺ”ติ. ภควโต เอตมตฺถํ อาโรเจสุํ. อนุชานามิ ภิกฺขเว อาวรณํ กาตุนฺติ. อาวรเณ กเต น อาทิยนฺติ. ภควโต เอตมตฺถํ อาโรเจสุํ. อนุชานามิ ภิกฺขเว โอวาทํ ฐเปตุนฺติ.}

\prose{เตน โข ปน สมเยน ฉพฺพคฺคิยา ภิกฺขุนิโย กายํ วิวริตฺวา ภิกฺขูนํ ทสฺเสนฺติ, ถนํ วิวริตฺวา ภิกฺขูนํ ทสฺเสนฺติ, อูรุํ วิวริตฺวา ภิกฺขูนํ ทสฺเสนฺติ, องฺคชาตํ วิวริตฺวา ภิกฺขูนํ ทสฺเสนฺติ, ภิกฺขู โอภาเสนฺติ, ภิกฺขูหิ สทฺธิํ สมฺปโยเชนฺติ “อปฺเปว นาม อเมฺหสุ สารชฺเชยฺยุนฺ”ติ. ภควโต เอตมตฺถํ อาโรเจสุํ. น ภิกฺขเว ภิกฺขุนิยา กาโย วิวริตฺวา ภิกฺขูนํ ทสฺเสตพฺโพ ฯเปฯ. น ถโน \csromanfolio{454}วิวริตฺวา ภิกฺขูนํ ทสฺเสตพฺโพ. น อูรุ วิวริตฺวา ภิกฺขูนํ ทสฺเสตพฺโพ. น องฺคชาตํ วิวริตฺวา ภิกฺขูนํ ทสฺเสตพฺพํ. น ภิกฺขู โอภาสิตพฺพา. น ภิกฺขูหิ สทฺธิํ สมฺป\-โยเช\-ตพฺพํ, ยา สมฺปโยเชยฺย, อาปตฺติ ทุกฺกฏสฺส. อนุชานามิ ภิกฺขเว ตสฺสา ภิกฺขุนิยา ทณฺฑกมฺมํ กาตุนฺติ. อถ โข ภิกฺขูนํ เอตทโหสิ “กิํ นุ โข ทณฺฑกมฺมํ กาตพฺพนฺ”ติ. ภควโต เอตมตฺถํ อาโรเจสุํ. อนุชานามิ ภิกฺขเว อาวรณํ กาตุนฺติ. อาวรเณ กเต น อาทิยนฺติ. ภควโต เอตมตฺถํ อาโรเจสุํ. อนุชานามิ ภิกฺขเว โอวาทํ ฐเปตุนฺติ. อถ โข ภิกฺขูนํ เอตทโหสิ “กปฺปติ นุ โข โอวาท\-ฏฺฐปิตาย\csromansharedfootnote{f8957b58a9f0f9c6}{โอวาทณฺฐปิตาย (สฺยา), โอวาทํฐปิตาย (ก)} ภิกฺขุนิยา สทฺธิํ อุโปสถํ กาตุํ, น นุ โข กปฺปตี”ติ. ภควโต เอตมตฺถํ อาโรเจสุํ. น ภิกฺขเว โอวาท\-ฏฺฐปิตาย ภิกฺขุนิยา สทฺธิํ อุโปสโถ กาตพฺโพ, ยาว น ตํ อธิกรณํ วูปสนฺตํ โหตีติ.}

\proseitem{412}{เตน โข ปน สมเยน อายสฺมา อุทายี โอวาทํ ฐเปตฺวา จาริกํ ปกฺกามิ. ภิกฺขุนิโย อุชฺฌายนฺติ ขิยฺยนฺติ วิปาเจนฺติ “กถํ หิ นาม อยฺโย อุทายี โอวาทํ ฐเปตฺวา จาริกํ ปกฺกมิสฺสตี”ติ. ภควโต เอตมตฺถํ อาโรเจสุํ. น ภิกฺขเว โอวาทํ ฐเปตฺวา จาริกา ปกฺกมิตพฺพา, โย ปกฺกเมยฺย, อาปตฺติ ทุกฺกฏสฺสาติ.}

\prose{เตน โข ปน สมเยน\csromansharedfootnote{449a23f609c826ce}{เตน โข ปน สมเยน ภิกฺขู (สฺยา, กํ)} พาลา อพฺยตฺตา โอวาทํ ฐเปนฺติ. ภควโต เอตมตฺถํ อาโรเจสุํ. น ภิกฺขเว พาเลน อพฺยตฺเตน โอวาโท ฐเปตพฺโพ, โย ฐเปยฺย, อาปตฺติ ทุกฺกฏสฺสาติ.}

\prose{เตน โข ปน สมเยน ภิกฺขู อวตฺถุสฺมิํ อการเณ โอวาทํ ฐเปนฺติ. ภควโต เอตมตฺถํ อาโรเจสุํ. น ภิกฺขเว อวตฺถุสฺมิํ อการเณ โอวาโท ฐเปตพฺโพ, โย ฐเปยฺย, อาปตฺติ ทุกฺกฏสฺสาติ.}

\prose{เตน โข ปน สมเยน ภิกฺขู โอวาทํ ฐเปตฺวา วินิจฺฉยํ น เทนฺติ. ภควโต เอตมตฺถํ อาโรเจสุํ. น ภิกฺขเว โอวาทํ ฐเปตฺวา วินิจฺฉโย น ทาตพฺโพ, โย น ทเทยฺย, อาปตฺติ ทุกฺกฏสฺสาติ.}

\proseitem{413}{\csromanfolio{455}\csromansymbolfootnote{*}{วิ 2. 413 ปิฏฺเฐปิ.}เตน โข ปน สมเยน ภิกฺขุนิโย โอวาทํ น คจฺฉนฺติ. ภควโต เอตมตฺถํ อาโรเจสุํ. น ภิกฺขเว ภิกฺขุนิยา โอวาโท น คนฺตพฺโพ, ยา น คจฺเฉยฺย, ยถาธมฺโม กาเรตพฺโพติ.}

\prose{เตน โข ปน สมเยน สพฺโพ ภิกฺขุนิ\-สํโฆ โอวาทํ คจฺฉติ. มนุสฺสา อุชฺฌายนฺติ ขิยฺยนฺติ วิปาเจนฺติ “ชายาโย อิมา อิเมสํ, ชาริโย อิมา อิเมสํ, อิทานิ อิเม อิมาหิ สทฺธิํ อภิรมิสฺสนฺตี”ติ. ภควโต เอตมตฺถํ อาโรเจสุํ. น ภิกฺขเว สพฺเพน ภิกฺขุนิ\-สํเฆน โอวาโท คนฺตพฺโพ, คจฺเฉยฺย เจ, อาปตฺติ ทุกฺกฏสฺส. อนุชานามิ ภิกฺขเว จตูหิ ปญฺจหิ ภิกฺขุนีหิ โอวาทํ คนฺตุนฺติ.}

\prose{เตน โข ปน สมเยน จตสฺโส ปญฺจ ภิกฺขุนิโย โอวาทํ คจฺฉนฺติ. มนุสฺสา ตเถว อุชฺฌายนฺติ ขิยฺยนฺติ วิปาเจนฺติ “ชายาโย อิมา อิเมสํ, ชาริโย อิมา อิเมสํ, อิทานิ อิเม อิมาหิ สทฺธิํ อภิรมิสฺสนฺตี”ติ. ภควโต เอตมตฺถํ อาโรเจสุํ. น ภิกฺขเว จตูหิ ปญฺจหิ ภิกฺขุนีหิ โอวาโท คนฺตพฺโพ, คจฺเฉยฺยุํ เจ, อาปตฺติ ทุกฺกฏสฺส. อนุชานามิ ภิกฺขเว เทฺว ติสฺโส ภิกฺขุนิโย\csromansharedfootnote{c9a5dfc9f430050d}{ภิกฺขุนีหิ (ก)} โอวาทํ คนฺตุํ, เอกํ ภิกฺขุํ อุปสงฺกมิตฺวา เอกํสํ อุตฺตราสงฺคํ กริตฺวา ปาเท วนฺทิตฺวา อุกฺกุฏิกํ นิสีทิตฺวา อญฺชลิํ ปคฺคเหตฺวา เอวมสฺส วจนีโย “ภิกฺขุนิ\-สํโฆ อยฺย ภิกฺขุ\-สํฆสฺส ปาเท วนฺทติ, โอวทูปสงฺกมนญฺจ ยาจติ, ลภตุ กิร อยฺย ภิกฺขุนิ\-สํโฆ โอวาทูปสงฺกมนนฺ”ติ. เตน ภิกฺขุนา ปาติโมกฺขุ\-ทฺเทสโก อุปสงฺกมิตฺวา เอวมสฺส วจนีโย “ภิกฺขุนิ\-สํโฆ ภนฺเต ภิกฺขุ\-สํฆสฺส ปาเท วนฺทติ, โอวาทูปสงฺกมนญฺจ ยาจติ, ลภตุ กิร ภนฺเต ภิกฺขุนิ\-สํโฆ โอวาทูปสงฺกมนนฺ”ติ. ปาติโมกฺขุ\-ทฺเทสเกน วตฺตพฺโพ “อตฺถิ โกจิ ภิกฺขุ ภิกฺขุโน\-วาทโก สมฺมโต”ติ. สเจ โหติ โกจิ ภิกฺขุ ภิกฺขุโน\-วาทโก สมฺมโต, ปาติโมกฺขุ\-ทฺเทสเกน วตฺตพฺโพ “อิตฺถนฺนาโม ภิกฺขุ ภิกฺขุโน\-วาทโก สมฺมโต, ตํ ภิกฺขุนิ\-สํโฆ อุปสงฺกมตู”ติ. สเจ น โหติ โกจิ ภิกฺขุ ภิกฺขุโน\-วาทโก สมฺมโต, ปาติโมกฺขุ\-ทฺเทสเกน วตฺตพฺโพ “โก อายสฺมา อุสฺสหติ ภิกฺขุนิโย โอวทิตุนฺ”ติ. สเจ โกจิ อุสฺสหติ ภิกฺขุนิโย โอวทิตุํ, โส จ โหติ อฏฺฐหงฺเคหิ สมนฺนาคโต, สมฺมนฺนิตฺวา วตฺตพฺโพ “อิตฺถนฺนาโม \csromanfolio{456}ภิกฺขุ ภิกฺขุโน\-วาทโก สมฺมโต, ตํ ภิกฺขุนิ\-สํโฆ อุปสงฺกมตู”ติ. สเจ น โกจิ อุสฺสหภิ ภิกฺขุนิโย โอวทิตุํ, ปาติโมกฺขุ\-ทฺเทสเกน วตฺตพฺโพ “นตฺถิ โกจิ ภิกฺขุ ภิกฺขุโน\-วาทโก สมฺมโต, ปาสาทิเกน ภิกฺขุนิ\-สํโฆ สมฺปาเทตู”ติ.}

\proseitem{414}{เตน โข ปน สมเยน ภิกฺขู โอวาทํ น คณฺหนฺติ. ภควโต เอตมตฺถํ อาโรเจสุํ. น ภิกฺขเว โอวาโท น คเหตพฺโพ, โย น คเณฺหยฺย, อาปตฺติ ทุกฺกฏสฺสาติ.}

\prose{เตน โข ปน สมเยน อญฺญตโร ภิกฺขุ พาโล โหติ, ตํ ภิกฺขุนิโย อุปสงฺกมิตฺวา เอตทโวจุํ “โอวาทํ อยฺย คณฺหาหี”ติ, อหํ หิ ภคินี พาโล, กถาหํ โอวาทํ คณฺหามีติ. คณฺหาหยฺย โอวาทํ, เอวํ หิ ภควตา ปญฺญตฺตํ ภิกฺขูหิ ภิกฺขุนีนํ โอวาโท คเหตพฺโพ”ติ. ภควโต เอตมตฺถํ อาโรเจสุํ. อนุชานามิ ภิกฺขเว ฐเปตฺวา พาลํ อวเสเสหิ โอวาทํ คเหตุนฺติ.}

\prose{เตน โข ปน สมเยน อญฺญตโร ภิกฺขุ คิลาโน โหติ, ตํ ภิกฺขุนิโย อุปสงฺกมิตฺวา เอตทโวจุํ “โอวาทํ อยฺย คณฺหาหี”ติ. อหํ หิ ภคินี คิลาโน, กถาหํ โอวาทํ คณฺหามีติ. คณฺหาหยฺย โอวาทํ, เอวํ หิ ภควตา ปญฺญตฺตํ “ฐเปตฺวา พาลํ อวเสเสหิ โอวาโท คเหตพฺโพ”ติ. ภควโต เอตมตฺถํ อาโรเจสุํ. อนุชานามิ ภิกฺขเว ฐเปตฺวา พาลํ, ฐเปตฺวา คิลานํ อวเสเสหิ โอวาทํ คเหตุนฺติ.}

\prose{เตน โข ปน สมเยน อญฺญตโร ภิกฺขุ คมิโก โหติ, ตํ ภิกฺขุนิโย อุปสงฺกมิตฺวา เอตทโวจุํ “โอวาทํ อยฺย คณฺหาหี”ติ. อหํ หิ ภคินี คมิโก, กถาหํ โอวาทํ คณฺหามีติ. คณฺหาหยฺย โอวาทํ, เอวํ หิ ภควตา ปญฺญตฺตํ “ฐเปตฺวา พาลํ, ฐเปตฺวา คิลานํ อวเสเสหิ โอวาโท คเหตพฺโพ”ติ. ภควโต เอตมตฺถํ อาโรเจสุํ. อนุชานามิ ภิกฺขเว ฐเปตฺวา พาลํ, ฐเปตฺวา คิลานํ, ฐเปตฺวา คมิกํ อวเสเสหิ โอวาทํ คเหตุนฺติ.}

\prose{เตน โข ปน สมเยน อญฺญตโร ภิกฺขุ อรญฺเญ วิหรติ, ตํ ภิกฺขุนิโย อุปสงฺกมิตฺวา เอตทโวจุํ “โอวาทํ อยฺย คณฺหาหี”ติ. อหํ หิ ภคินี อรญฺเญ วิหรามิ, กถาหํ โอวาทํ คณฺหามีติ. \csromanfolio{457}คณฺหาหยฺย โอวาทํ, เอวํ หิ ภควตา ปญฺญตฺตํ “ฐเปตฺวา พาลํ, ฐเปตฺวา คิลานํ, ฐเปตฺวา คมิกํ อวเสเสหิ โอวาโท คเหตพฺโพ”ติ. ภควโต เอตมตฺถํ อาโรเจสุํ. อนุชานามิ ภิกฺขเว อรญฺญิเกน ภิกฺขุนา โอวาทํ คเหตุํ, สงฺเกตญฺจ กาตุํ “อตฺร ปติหริสฺสามี”ติ.}

\proseitem{415}{เตน โข ปน สมเยน ภิกฺขู โอวาทํ คเหตฺวา น อาโรเจนฺติ. ภควโต เอตมตฺถํ อาโรเจสุํ. น ภิกฺขเว โอวาโท น อาโรเจตพฺโพ, โย น อาโรเจยฺย, อาปตฺติ ทุกฺกฏสฺสาติ.}

\prose{เตน โข ปน สมเยน ภิกฺขู โอวาทํ คเหตฺวา น ปจฺจาหรนฺติ. ภควโต เอตมตฺถํ อาโรเจสุํ. น ภิกฺขเว โอวาโท น ปจฺจาหริตพฺโพ, โย น ปจฺจาหเรยฺย, อาปตฺติ ทุกฺกฏสฺสาติ.}

\prose{เตน โข ปน สมเยน ภิกฺขุนิโย สงฺเกตํ น คจฺฉนฺติ. ภควโต เอตมตฺถํ อาโรเจสุํ. น ภิกฺขเว ภิกฺขุนิยา สงฺเกตํ น คนฺตพฺพํ, ยา น คจฺเฉยฺย, อาปตฺติ ทุกฺกฏสฺสาติ.}

\proseitem{416}{เตน โข ปน สมเยน ภิกฺขุนิโย ทีฆานิ กายพนฺธนานิ ธาเรนฺติ, เตเหว ผาสุกา นาเมนฺติ. มนุสฺสา อุชฺฌายนฺติ ขิยฺยนฺติ วิปาเจนฺติ “ฯเปฯ เสยฺยถาปิ คิหินี กามโภคินิโย”ติ. ภควโต เอตมตฺถํ อาโรเจสุํ. น ภิกฺขเว ภิกฺขุนิยา ทีฆํ กายพนฺธนํ ธาเรตพฺพํ, ยา ธาเรยฺย, อาปตฺติ ทุกฺกฏสฺส. อนุชานามิ ภิกฺขเว ภิกฺขุนิยา เอกปริยากตํ\csromansharedfootnote{8e4bf34d7e8a1b2e}{เอกปริยายกตํ (สฺยา)} กายพนฺธนํ, น จ เตน ผาสุกา นาเมตพฺพา, ยา นาเมยฺย, อาปตฺติ ทุกฺกฏสฺสาติ.}

\prose{เตน โข ปน สมเยน ภิกฺขุนิโย วิลีเวน\csromansharedfootnote{d13bfafe50c5689e}{วิลิเวน(ก)} ปฏฺเฏน ผาสุกา นาเมนฺติ ฯเปฯ จมฺมปฏฺเฏน ผาสุกา นาเมนฺติ. ทุสฺสปฏฺเฏน ผาสุกา นาเมนฺติ. ทุสฺสเวณิยา ผาสุกา นาเมนฺติ. ทุสฺสวฏฺฏิยา ผาสุกา นาเมนฺติ. โจฬปฏฺเฏน ผาสุกา นาเมนฺติ. โจฬเวณิยา ผาสุกา นาเมนฺติ. โจฬวฏฺฏิยา ผาสุกา นาเมนฺติ. สุตฺตเวณิยา ผาสุกา นาเมนฺติ. สุตฺตวฏฺฏิยา ผาสุกา นาเมนฺติ. มนุสฺสา อุชฺฌายนฺติ ขิยฺยนฺติ วิปาเจนฺติ “ฯเปฯ เสยฺยถาปิ คิหินี กามโภคินิโย”ติ. ภควโต เอตมตฺถํ \csromanfolio{458}อาโรเจสุํ. น ภิกฺขเว ภิกฺขุนิยา วิลีเวน ปฏฺเฏน ผาสุกา นาเมตพฺพา ฯเปฯ. น สุตฺตวฏฺฏิยา ผาสุกา นาเมตพฺพา, ยา นาเมยฺย อาปตฺติ ทุกฺกฏสฺสาติ.}

\prose{เตน โข ปน สมเยน ภิกฺขุนิโย อฏฺฐิลฺเลน ชฆนํฆํสาเปนฺติ ฯเปฯ โคหนุเกน ชฆนํ โกฏฺฏาเปนฺติ. หตฺถํ โกฏฺฏาเปนฺติ. หตฺถโกจฺฉํ โกฏฺฏาเปนฺติ, ปาทํ โกฏฺฏาเปนฺติ. ปาทโกจฺฉํ โกฏฺฏาเปนฺติ. อูรุํ โกฏฺฏาเปนฺติ. มุขํ โกฏฺฏาเปนฺติ. ทนฺตมํสํ โกฏฺฏาเปนฺติ. มนุสฺสา อุชฺฌายนฺติ ขิยฺยนฺติ วิปาเจนฺติ “ฯเปฯ เสยฺยถาปิ คิหินี กามโภคินิโย”ติ. ภควโต เอตมตฺถํ อาโรเจสุํ. น ภิกฺขเว ภิกฺขุนิยา อฏฺฐิลฺเลน ชฆนํ ฆํสาเปตพฺพํ ฯเปฯ. น ทนฺตมํสํ โกฏฺฏาเปตพฺพํ, ยา โกฏฺฏาเปยฺย, อาปตฺติ ทุกฺกฏสฺสาติ.}

\proseitem{417}{เตน โข ปน สมเยน ฉพฺพคฺคิยา ภิกฺขุนิโย มุขํ อาลิมฺปนฺติ ฯเปฯ. มุขํ อุมฺมทฺเทนฺติ. มุขํ จุณฺเณนฺติ. มโนสิลิกาย มุขํ ลญฺเฉนฺติ. องฺคราคํ กโรนฺติ. มุขราคํ กโรนฺติ. องฺคราค\-มุขราคํ กโรนฺติ. มนุสฺสา อุชฺฌายนฺติ ขิยฺยนฺติ วิปาเจนฺติ “ฯเปฯ เสยฺยถาปิคิหินี กามโภคินิโย”ติ. ภควโต เอตมตฺถํ อาโรเจสุํ. น ภิกฺขเว ภิกฺขุนิยา มุขํ อาลิมฺปิตพฺพํ ฯเปฯ. น มุขํ อุมฺมทฺทิตพฺพํ. น มุขํ จุณฺเณตพฺพํ. น มโนสิลิกาย มุขํ ลญฺฉิตพฺพํ. น องฺคราโค กาตพฺโพ. น มุขราโค กาตพฺโพ. น องฺคราค\-มุขราโค กาตพฺโพ, ยา กเรยฺย, อาปตฺติ ทุกฺกฏสฺสาติ.}

\prose{เตน โข ปน สมเยน ฉพฺพคฺคิยา ภิกฺขุนิโย อวงฺคํ\csromansharedfootnote{9689a4e5e0682c98}{อปางฺคํ (?)} กโรนฺติ ฯเปฯ. วิเสสกํ กโรนฺติ. โอโลกนเกน โอโลเกนฺติ. สาโลเก ติฏฺฐนฺติ. นจฺจํ\csromansharedfootnote{5723850b21ed9da8}{สนจฺจํ (สี, สฺยา), สมชฺชํ (ก)} การาเปนฺติ. เวสิํ วุฏฺฐาเปนฺติ. ปานาคารํ ฐเปนฺติ. สูนํ ฐเปนฺติ. อาปณํ ปสาเรนฺติ. วฑฺฒิํ ปโยเชนฺติ. วณิชฺชํ ปโยเชนฺติ. ทาสํ อุปฏฺฐาเปนฺติ. ทาสิํ อุปฏฺฐาเปนฺติ. กมฺมการํ อุปฏฺฐาเปนฺติ. กมฺมการิํ อุปฏฺฐาเปนฺติ. ติรจฺฉานคตํ อุปฏฺฐาเปนฺติ. หรีตก\-ปตฺติกํ\csromansharedfootnote{a79945f1eb187ee5}{หรีตกปณฺณิกํ (ก)} ปกิณนฺติ. นมตกํ ธาเรนฺติ. มนุสฺสา อุชฺฌายนฺติ ขิยฺยนฺติ วิปาเจนฺติ “ฯเปฯ เสยฺยถาปิ คิหินี กามโภคินิโย”ติ. ภควโต เอตมตฺถํ อาโรเจสุํ. น ภิกฺขเว ภิกฺขุนิยา อวงฺคํ กาตพฺพํ ฯเปฯ. น วิเสสกํ กาตพฺพํ. น \csromanfolio{459}โอโลกนเกน โอโลเกตพฺพํ. น สาโลเก ฐาตพฺพํ. น นจฺจํ การาเปตพฺพํ. น เวสี วุฏฺฐาเปตพฺพา. น ปานาคารํ ฐเปตพฺพํ. น สูนา ฐเปตพฺพา. น อาปโณ ปสาเรตพฺโพ. น วฑฺฒิ ปโยเชตพฺพา. น วณิชฺชา ปโยเชตพฺพา. น ทาโส อุปฏฺฐาเปตพฺโพ. น ทาสี อุปฏฺฐาเปตพฺพา. น กมฺมกาโร อุปฏฺฐาเปตพฺโพ. น กมฺมการี อุปฏฺฐาเปตพฺพา. น ติรจฺฉานคโต อุปฏฺฐาเปตพฺโพ. น หรีตก\-ปกฺกิกํ ปกิณิตพฺพํ. น นมตกํ ธาเรตพฺพํ, ยา ธาเรยฺย, อาปตฺติ ทุกฺกฏสฺสาติ.}

\proseitem{418}{เตน โข ปน สมเยน ฉพฺพคฺคิยา ภิกฺขุนิโย สพฺพนีลกานิ จีวรานิ ธาเรนฺติ ฯเปฯ. สพฺพปีตกานิ จีวรานิ ธาเรนฺติ. สพฺพ\-โลหิตกานิ จีวรานิ ธาเรนฺติ. สพฺพ\-มญฺชิฏฺฐิ\-กานิ จีวรานิ ธาเรนฺติ. สพฺพกณฺหานิ จีวรานิ ธาเรนฺติ. สพฺพ\-มหารงฺค\-รตฺตานิ จีวรานิ ธาเรนฺติ. สพฺพ\-มหา\-นาม\-รตฺตานิ จีวรานิ ธาเรนฺติ. อจฺฉินฺนทสานิ จีวรานิ ธาเรนฺติ. ทีฆทสานิ จีวรานิ ธาเรนฺติ. ปุปฺผทสานิ จีวรานิ ธาเรนฺติ. ผลทสานิ จีวรานิ ธาเรนฺติ. กญฺจุกํ ธาเรนฺติ. ติรีฏกํ ธาเรนฺติ. มนุสฺสา อุชฺฌายนฺติ ขิยฺยนฺติ วิปาเจนฺติ “ฯเปฯ เสยฺยถาปิ คิหินี กามโภคินิโย”ติ. ภควโต เอตมตฺถํ อาโรเจสุํ. น ภิกฺขเว ภิกฺขุนิยา สพฺพนีลกานิ จีวรานิ ธาเรตพฺพานิ ฯเปฯ. น ติรีฏกํ ธาเรตพฺพํ, ยา ธาเรยฺย, อาปตฺติ ทุกฺกฏสฺสาติ.}

\proseitem{419}{เตน โข ปน สมเยน อญฺญตรา ภิกฺขุนี กาลํ กโรนฺตี เอวมาห “มมจฺจเยน มยฺหํ ปริกฺขาโร สํฆสฺส โหตู”ติ. ตตฺถ ภิกฺขู จ ภิกฺขุนิโย จ วิวทนฺติ “อมฺหากํ โหติ, อมฺหากํ โหตี”ติ. ภควโต เอตมตฺถํ อาโรเจสุํ. ภิกฺขุนี เจ ภิกฺขเว กาลํ กโรนฺตี เอวํ วเทยฺย “มมจฺจเยน มยฺหํ ปริกฺขาโร สํฆสฺส โหตู”ติ, อนิสฺสโร ตตฺถ ภิกฺขุสํโฆ, ภิกฺขุนิสํฆสฺเสเวตํ. สิกฺขมานา เจ ภิกฺขเว ฯเปฯ. สามเณรี เจ ภิกฺขเว กาลํ กโรนฺตี เอวํ วเทยฺย “มมจฺจเยน มยฺหํ ปริกฺขาโร สํฆสฺส โหตู”ติ, อนิสฺสโร ตตฺถ ภิกฺขุสํโฆ, ภิกฺขุนิสํฆสฺเสเวตํ. ภิกฺขุ เจ ภิกฺขเว กาลํ กโรนฺโต เอวํ วเทยฺย “มมจฺจเยน มยฺหํ ปริกฺขาโร สํฆสฺส โหตู”ติ, อนิสฺสโร ตตฺถ ภิกฺขุนิ\-สํโฆ, ภิกฺขุสํฆสฺเสเวตํ. สามเณโร เจ ภิกฺขเว ฯเปฯ. อุปาสโก เจ ภิกฺขเว ฯเปฯ. อุปาสิกา เจ ภิกฺขเว ฯเปฯ. อญฺโญ เจ ภิกฺขเว โกจิ กาลํ กโรนฺโต เอวํ วเทยฺย “มมจฺจเยน มยฺหํ \csromanfolio{460}ปริกฺขาโร สํฆสฺส โหตู”ติ, อนิสฺสโร ตตฺถ ภิกฺขุนิ\-สํโฆ, ภิกฺขุสํฆสฺเสเวตนฺติ.}

\proseitem{420}{เตน โข ปน สมเยน อญฺญตรา อิตฺถี ปุราณมลฺลี ภิกฺขุนีสุ ปพฺพชิตา โหติ, สา รถิกาย ทุพฺพลกํ ภิกฺขุํ ปสฺสิตฺวา อํสกูเฏน ปหารํ ทตฺวา ปาเตสิ\csromansharedfootnote{3e702a8df02febc3}{ปวฏฺเฏสิ (สี)}. ภิกฺขู อุชฺฌายนฺติ ขิยฺยนฺติ วิปาเจนฺติ “กถํ หิ นาม ภิกฺขุนี ภิกฺขุสฺส ปหารํ ทสฺสตี”ติ. ภควโต เอตมตฺถํ อาโรเจสุํ. น ภิกฺขเว ภิกฺขุนิยา ภิกฺขุสฺส ปหาโร ทาตพฺโพ, ยา ทเทยฺย, อาปตฺติ ทุกฺกฏสฺส. อนุชานามิ ภิกฺขเว ภิกฺขุนิยา ภิกฺขุํ ปสฺสิตฺวา ทูรโตว โอกฺกมิตฺวา มคฺคํ ทาตุนฺติ.}

\prose{เตน โข ปน สมเยน อญฺญตรา อิตฺถี ปวุตฺถปติกา ชาเรน คพฺภินี โหติ, สา คพฺภํ ปาเตตฺวา กุลูปิกํ ภิกฺขุนิํ เอตทโวจ “หนฺทยฺเย อิมํ คพฺภํ ปตฺเตน นีหรา”ติ. อถ โข สา ภิกฺขุนี ตํ คพฺภํ ปตฺเต ปกฺขิปิตฺวา สํฆาฏิยา ปฏิจฺฉาเทตฺวา อคมาสิ. เตน โข ปน สมเยน อญฺญตเรน ปิณฺฑจาริเกน ภิกฺขุนา สมาทานํ กตํ โหติ “ยาหํ ปฐมํ ภิกฺขํ ลภิสฺสามิ, น ตํ อทตฺวา ภิกฺขุสฺส วา ภิกฺขุนิยา วา ปริภุญฺชิสฺสามี”ติ. อถ โข โส ภิกฺขุ ตํ ภิกฺขุนิํ ปสฺสิตฺวา เอตทโวจ “หนฺท ภคินิ ภิกฺขํ ปฏิคฺคณฺหา”ติ. อลํ อยฺยาติ. ทุติยมฺปิ โข ฯเปฯ. ตติยมฺปิ โข โส ภิกฺขุ ตํ ภิกฺขุนิํ เอตทโวจ “หนฺท ภคินิ ภิกฺขํ ปฏิคฺคณฺหา”ติ. อลํ อยฺยาติ. มยา โข ภคินิ สมาทานํ กตํ “ยาหํ ปฐมํ ภิกฺขํ ลภิสฺสามิ, น ตํ อทตฺวา ภิกฺขุสฺส วา ภิกฺขุนิยา วา ปริภุญฺชิสฺสามี”ติ. หนฺท ภคินิ ภิกฺขํ ปฏิคฺคณฺหาติ. อถ โข สา ภิกฺขุนี เตน ภิกฺขุนา นิปฺปีฬิยมานา นีหริตฺวา ปตฺตํ ทสฺเสสิ “ปสฺส อยฺย ปตฺเต คพฺภํ, มา จ กสฺสจิ อาโรเจสี”ติ. อถ โข โส ภิกฺขุ อุชฺฌายติ ขิยฺยติ วิปาเจติ “กถํ หิ นาม ภิกฺขุนี ปตฺเตน คพฺภํ นีหริสฺสตี”ติ. อถ โข โส ภิกฺขุ ภิกฺขูนํ เอตมตฺถํ อาโรเจสิ. เย เต ภิกฺขู อปฺปิจฺฉา ฯเปฯ เต อุชฺฌายนฺติ ขิยฺยนฺติ วิปาเจนฺติ “กถํ หิ นาม ภิกฺขุนี ปตฺเตน คพฺภํ นีหริสฺสตี”ติ. ภควโต เอตมตฺถํ อาโรเจสุํ. น ภิกฺขเว ภิกฺขุนิยา ปตฺเตน คพฺโภ นีหริตพฺโพ, \csromanfolio{461}ยา นีหเรยฺย, อาปตฺติ ทุกฺกฏสฺส. อนุชานามิ ภิกฺขเว ภิกฺขุนิยา ภิกฺขุํ ปสฺสิตฺวา นีหริตฺวา ปตฺตํ ทสฺเสตุนฺติ.}

\prose{เตน โข ปน สมเยน ฉพฺพคฺคิยา ภิกฺขุนิโย ภิกฺขุํ ปสฺสิตฺวา ปริวตฺเตตฺวา ปตฺตมูลํ ทสฺเสนฺติ. ภิกฺขู อุชฺฌายนฺติ ขิยฺยนฺติ วิปาเจนฺติ “กถํ หิ นาม ฉพฺพคฺคิยา ภิกฺขุนิโย ภิกฺขุํ ปสฺสิตฺวา ปริวตฺเตตฺวา ปตฺตมูลํ ทสฺเสสฺสนฺตี”ติ. อถ โข เต ภิกฺขู ภควโต เอตมตฺถํ อาโรเจสุํ. น ภิกฺขเว ภิกฺขุนิยา ภิกฺขุํ ปสฺสิตฺวา ปริวตฺเตตฺวา ปตฺตมูลํ ทสฺเสตพฺพํ, ยา ทสฺเสยฺย, อาปตฺติ ทุกฺกฏสฺส. อนุชานามิ ภิกฺขเว ภิกฺขุนิยา ภิกฺขุํ ปสฺสิตฺวา อุกฺกุชฺชิตฺวา ปตฺตํ ทสฺเสตุํ, ยญฺจ ปตฺเต อามิสํ โหติ, เตน จ ภิกฺขุ นิมนฺเตตพฺโพติ.}

\prose{เตน โข ปน สมเยน สาวตฺถิยํ รถิกาย ปุริส\-พฺยญฺชนํ ฉฑฺฑิตํ โหติ, ตํ ภิกฺขุนิโย สกฺกจฺจํ อุปนิชฺฌายิํสุ, มนุสฺสา อุกฺกุฏฺฐิํ อกํสุ, ตา ภิกฺขุนิโย มงฺกู อเหสุํ. อถ โข ตา ภิกฺขุนิโย อุปสฺสยํ คนฺตฺวา ภิกฺขุนีนํ เอตมตฺถํ อาโรเจสุํ. ยา ตา ภิกฺขุนิโย อปฺปิจฺฉา ฯเปฯ ตา อุชฺฌายนฺติ ขิยฺยนฺติ วิปาเจนฺติ “กถํ หิ นาม ภิกฺขุนิโย ปุริส\-พฺยญฺชนํ อุปนิชฺฌายิสฺสนฺตี”ติ. อถ โข ตา ภิกฺขุนิโย ภิกฺขูนํ เอตมตฺถํ อาโรเจสุํ. ภิกฺขู ภควโต เอตมตฺถํ อาโรเจสุํ. น ภิกฺขเว ภิกฺขุนิยา ปุริส\-พฺยญฺชนํ อุปนิชฺฌายิตพฺพํ, ยา อุปนิชฺฌาเยยฺย, อาปตฺติ ทุกฺกฏสฺสาติ.}

\proseitem{421}{เตน โข ปน สมเยน มนุสฺสา ภิกฺขูนํ อามิสํ เทนฺติ, ภิกฺขู ภิกฺขุนีนํ เทนฺติ. มนุสฺสา อุชฺฌายนฺติ ขิยฺยนฺติ วิปาเจนฺติ “กถํ หิ นาม ภทนฺตา อตฺตโน ปริโภค\-ตฺถาย ทินฺนํ อญฺเญสํ ทสฺสนฺติ, มยมฺปิ น ชานาม ทานํ ทาตุนฺ”ติ. ภควโต เอตมตฺถํ อาโรเจสุํ. น ภิกฺขเว อตฺตโน ปริโภค\-ตฺถาย ทินฺนํ อญฺเญสํ ทาตพฺพํ, โย ทเทยฺย, อาปตฺติ ทุกฺกฏสฺสาติ.}

\prose{เตน โข ปน สมเยน ภิกฺขูนํ อามิสํ อุสฺสนฺนํ โหติ. ภควโต เอตมตฺถํ อาโรเจสุํ. อนุชานามิ ภิกฺขเว สํฆสฺส ทาตุนฺติ. พาฬฺหตรํ อุสฺสนฺนํ โหติ. ภควโต เอตมตฺถํ อาโรเจสุํ. อนุชานามิ ภิกฺขเว ปุคฺคลิกมฺปิ ทาตุนฺติ.}

\prose{\csromanfolio{462}เตน โข ปน สมเยน ภิกฺขูนํ สนฺนิธิกตํ อามิสํ อุสฺสนฺนํ โหติ. ภควโต เอตมตฺถํ อาโรเจสุํ. อนุชานามิ ภิกฺขเว ภิกฺขูนํ สนฺนิธิํ ภิกฺขุนีหิ\csromansharedfootnote{18256b268bfbeffa}{ภิกฺขุนีหิ ภิกฺขูหิ (สี)} ปฏิคฺคาหาเปตฺวา\csromansharedfootnote{1cac5903a0fe9977}{ปฏิคฺคหาเปตฺวา (ก)} ปริภุญฺชิ\-ตุนฺติ.}

\prose{เตน โข ปน สมเยน มนุสฺสา ภิกฺขุนีนํ อามิสํ เทนฺติ, ภิกฺขุนิโย ภิกฺขูนํ เทนฺติ. มนุสฺสา อุชฺฌายนฺติ ขิยฺยนฺติ วิปาเจนฺติ “กถํ หิ นาม ภิกฺขุนิโย อตฺตโน ปริโภค\-ตฺถาย ทินฺนํ อญฺเญสํ ทสฺสนฺติ, มยมฺปิ น ชานาม ทานํ ทาตุนฺ”ติ. ภควโต เอตมตฺถํ อาโรเจสุํ. น ภิกฺขเว ภิกฺขุนิยา อตฺตโน ปริโภค\-ตฺถาย ทินฺนํ อญฺเญสํ ทาตพฺพํ, ยา ทเทยฺย, อาปตฺติ ทุกฺกฏสฺสาติ.}

\prose{เตน โข ปน สมเยน ภิกฺขุนีนํ อามิสํ อุสฺสนฺนํ โหติ. ภควโต เอตมตฺถํ อาโรเจสุํ. อนุชานามิ ภิกฺขเว สํฆสฺส ทาตุนฺติ. พาฬฺหตรํ อุสฺสนฺนํ โหติ. ภควโต เอตมตฺถํ อาโรเจสุํ. อนุชานามิ ภิกฺขเว ปุคฺคลิกมฺปิ ทาตุนฺติ.}

\prose{เตน โข ปน สมเยน ภิกฺขุนีนํ สนฺนิธิกตํ อามิสํ อุสฺสนฺนํ โหติ. ภควโต เอตมตฺถํ อาโรเจสุํ. อนุชานามิ ภิกฺขเว ภิกฺขุนีนํ สนฺนิธิํ ภิกฺขูหิ\csromansharedfootnote{964506ffe179e630}{ภิกฺขูหิ ภิกฺขุนีหิ (สี)} ปฏิคฺคาหาเปตฺวา ปริภุญฺชิ\-ตุนฺติ.}

\proseitem{442}{เตน โข ปน สมเยน ภิกฺขูนํ เสนาสนํ อุสฺสนฺนํ โหติ, ภิกฺขุนีนํ\csromansharedfootnote{7e30edec3e4d1aed}{ภิกฺขุนีนํ เสนาสนํ (สฺยา, กํ)} น โหติ. ภิกฺขุนิโย ภิกฺขูนํ สนฺติเก ทูตํ ปาเหสุํ “สาธุ ภนฺเต, อยฺยา อมฺหากํ เสนาสนํ เทนฺตุ ตาวกาลิกนฺ”ติ. ภควโต เอตมตฺถํ อาโรเจสุํ. อนุชานามิ ภิกฺขเว ภิกฺขุนีนํ เสนาสนํ ทาตุํ ตาวกาลิกนฺติ.}

\prose{เตน โข ปน สมเยน อุตุนิโย ภิกฺขุนิโย โอนทฺธมญฺจํ โอนทฺธปีฐํ อภินิ\-สีท\-นฺติปิ อภินิ\-ปชฺช\-นฺติปิ, เสนาสนํ โลหิเตน มกฺขิยฺยติ. ภควโต เอตมตฺถํ อาโรเจสุํ. น ภิกฺขเว ภิกฺขุนิยา โอนทฺธมญฺจํ โอนทฺธปีฐํ อภินิ\-สีทิ\-ตพฺพํ อภินิ\-ปชฺชิ\-ตพฺพํ, ยา อภินิสีเทยฺย วา อภินิปชฺเชยฺย วา, อาปตฺติ ทุกฺกฏสฺส. อนุชานามิ ภิกฺขเว อาวสถจีวรนฺติ. อาวสถจีวรํ โลหิเตน มกฺขิยฺยติ. ภควโต เอตมตฺถํ อาโรเจสุํ. \csromanfolio{463}อนุชานามิ ภิกฺขเว อาณิโจฬกนฺติ. โจฬกํ นิปตติ. ภควโต เอตมตฺถํ อาโรเจสุํ. อนุชานามิ ภิกฺขเว สุตฺตเกน พนฺธิตฺวา อูรุยา พนฺธิตุนฺติ. สุตฺตํ ฉิชฺชติ. ภควโต เอตมตฺถํ อาโรเจสุํ. อนุชานามิ ภิกฺขเว สํเวลฺลิยํ กฏิสุตฺตกนฺติ.}

\prosewithcloser{เตน โข ปน สมเยน ฉพฺพคฺคิยา ภิกฺขุนิโย สพฺพกาลํ กฏิสุตฺตกํ ธาเรนฺติ. มนุสฺสา อุชฺฌายนฺติ ขิยฺยนฺติ วิปาเจนฺติ “เสยฺยถาปิ คิหินี กามโภคินิโย”ติ. ภควโต เอตมตฺถํ อาโรเจสุํ. น ภิกฺขเว ภิกฺขุนิยา สพฺพกาลํ กฏิสุตฺตกํ ธาเรตพฺพํ, ยา ธาเรยฺย, อาปตฺติ ทุกฺกฏสฺส. อนุชานามิ ภิกฺขเว อุตุนิยา กฏิสุตฺตกนฺติ.}{ทุติย\-ภาณวาโร นิฏฺฐิโต.}
\csromanmatikaanchor{mtk.210}
\csromantocmarkref{section}{3. ตติยภาณวาร}{mtk.210}
\bookmark[dest={mtk.210},level=1]{3. ตติยภาณวาร}
\csromanheader{1}{3. ตติย\-ภาณวาร}
\proseitem{423}{เตน โข ปน สมเยน อุปสมฺปนฺนาโย ทิสฺสนฺติ อนิมิตฺตาปิ นิมิตฺตมตฺตาปิ อโลหิตาปิ ธุวโลหิตาปิ ธุวโจฬาปิ ปคฺฆรนฺตีปิ สิขรณีปิ อิตฺถิปณฺฑกาปิ เวปุริสิกาปิ สมฺภินฺนาปิ อุภโตพฺยญฺชนาปิ. ภควโต เอตมตฺถํ อาโรเจสุํ. อนุชานามิ ภิกฺขเว อุปสมฺปาเทนฺติยา จตุวีสติ อนฺตรายิเก ธมฺเม ปุจฺฉิตุํ. เอวญฺจ ปน ภิกฺขเว ปุจฺฉิตพฺพา “นสิ อนิมิตฺตา, นสิ นิมิตฺตมตฺตา, นสิ อโลหิตา, นสิ ธุวโลหิตา, นสิ ธุวโจฬา, นสิ ปคฺฆรนฺตี, นสิ สิขรณี\csromansharedfootnote{4878e5f990f3b135}{สิขริณี (สี, สฺยา)}, นสิ อิตฺถิปณฺฑกา, นสิ เวปุริสิกา, นสิ สมฺภินฺนา, นสิ อุภโตพฺยญฺชนา. สนฺติ เต เอวรูปา อาพาธา กุฏฺฐํ คณฺโฑ กิลาโส โสโส อปมาโร, มนุสฺสาสิ, อิตฺถีสิ, ภุชิสฺสาสิ, อณณาสิ, นสิ ราชภฏี, อนุญฺญาตาสิ มาตาปิตูหิ สามิเกน, ปริปุณฺณ\-วีสติวสฺสาสิ, ปริปุณฺณํ เต ปตฺตจีวรํ, กินฺนามาสิ, กานามา เต ปวตฺตินี”ติ.}

\prose{เตน โข ปน สมเยน ภิกฺขู ภิกฺขุนีนํ อนฺตรายิเก ธมฺเม ปุจฺฉนฺติ, อุปสมฺปทา\-เปกฺขาโย วิตฺถายนฺติ, มงฺกู โหนฺติ, น สกฺโกนฺติ วิสฺสชฺเชตุํ. ภควโต เอตมตฺถํ อาโรเจสุํ. อนุชานามิ ภิกฺขเว เอกโต\-อุปสมฺปนฺนาย ภิกฺขุนิ\-สํเฆ วิสุทฺธาย ภิกฺขุสํเฆ อุปสมฺปาเทตุนฺติ.}

\prose{\csromanfolio{464}เตน โข ปน สมเยน ภิกฺขุนิโย อนนุสิฏฺฐา อุปสมฺปทา\-เปกฺขาโย อนฺตรายิเก ธมฺเม ปุจฺฉนฺติ, อุปสมฺปทา\-เปกฺขาโย วิตฺถายนฺติ, มงฺกู โหนฺติ, น สกฺโกนฺติ วิสฺสชฺเชตุํ. ภควโต เอตมตฺถํ อาโรเจสุํ. อนุชานามิ ภิกฺขเว ปฐมํ อนุสาสิตฺวา ปจฺฉา อนฺตรายิเก ธมฺเม ปุจฺฉิตุนฺติ.}

\prose{ตตฺเถว สํฆมชฺเฌ อนุสาสนฺติ, อุปสมฺปทา\-เปกฺขาโย ตเถว วิตฺถายนฺติ, มงฺกู โหนฺติ, น สกฺโกนฺติ วิสฺสชฺเชตุํ. ภควโต เอตมตฺถํ อาโรเจสุํ. อนุชานามิ ภิกฺขเว เอกมนฺตํ อนุสาสิตฺวา สํฆมชฺเฌ อนฺตรายิเก ธมฺเม ปุจฺฉิตุํ. เอวญฺจ ปน ภิกฺขเว อนุสาสิตพฺพา\csromandash{}}

\proseitem{424}{ปฐมํ อุปชฺฌํ คาหาเปตพฺพา, อุปชฺฌํ คาหาเปตฺวา ปตฺตจีวรํ อาจิกฺขิตพฺพํ, อยํ เต ปตฺโต, อยํ สํฆาฏิ, อยํ อุตฺตราสงฺโค, อยํ อนฺตรวาสโก, อิทํ สํกจฺจิกํ\csromansharedfootnote{4bf5dd4614a8e5c0}{สงฺกจฺฉิกํ (สฺยา)}, อยํ อุทกสาฏิกา. คจฺฉ อมุมฺหิ โอกาเส ติฏฺฐาหีติ.}

\prose{พาลา อพฺยตฺตา อนุสาสนฺติ, ทุรนุสิฏฺฐา อุปสมฺปทา\-เปกฺขาโย วิตฺถายนฺติ, มงฺกู โหนฺติ, น สกฺโกนฺติ วิสฺสชฺเชตุํ. ภควโต เอตมตฺถํ อาโรเจสุํ. น ภิกฺขเว พาลาย อพฺยตฺตาย อนุสาสิตพฺพา, ยา อนุสาเสยฺย, อาปตฺติ ทุกฺกฏสฺส. อนุชานามิ ภิกฺขเว พฺยตฺตาย ภิกฺขุนิยา ปฏิพลาย อนุสาสิตุนฺติ.}

\prose{อสมฺมตา อนุสาสนฺติ. ภควโต เอตมตฺถํ อาโรเจสุํ. น ภิกฺขเว อสมฺมตาย อนุสาสิตพฺพา, ยา อนุสาเสยฺย, อาปตฺติ ทุกฺกฏสฺสา. อนุชานามิ ภิกฺขเว สมฺมตาย อนุสาสิตุํ. เอวญฺจ ปน ภิกฺขเว สมฺมนฺนิตพฺพา, อตฺตนา วา อตฺตานํ สมฺมนฺนิตพฺพํ, ปราย วา ปรา สมฺมนฺนิตพฺพา.}

\prose{กถญฺจ อตฺตนาว\csromansharedfootnote{586bbcce12b9911d}{อตฺตนา วา (ก)} อตฺตานํ สมฺมนฺนิตพฺพํ, พฺยตฺตาย ภิกฺขุนิยา ปฏิพลาย สํโฆ ญาเปตพฺโพ “สุณาตุ เม อยฺเย สํโฆ, อิตฺถนฺนามา อิตฺถนฺนามาย อยฺยาย อุปสมฺปทา\-เปกฺขา, ยทิ สํฆสฺส ปตฺตกลฺลํ, อหํ อิตฺถนฺนามา อิตฺถนฺนามํ อนุสาเสยฺยนฺ”ติ, เอวํ อตฺตนาว อตฺตานํ สมฺมนฺนิตพฺพํ.}

\prose{กถญฺจ ปราย\csromansharedfootnote{46e1f5df37e2dd13}{ปราย วา (ก)} ปรา สมฺมนฺนิตพฺพา, พฺยตฺตาย ภิกฺขุนิยา ปฏิพลาย สํโฆ ญาเปตพฺโพ “สุณาตุ เม อยฺเย สํโฆ, อิตฺถนฺนามา \csromanfolio{465}อิตฺถนฺนามาย อยฺยาย อุปสมฺปทา\-เปกฺขา, ยทิ สํฆสฺส ปตฺตกลฺลํ, อิตฺถนฺนามา อิตฺถนฺนามํ อนุสาเสยฺยา”ติ, เอวํ ปราย ปรา สมฺมนฺนิตพฺพา.}

\prose{ตาย สมฺมตาย ภิกฺขุนิยา อุปสมฺปทา\-เปกฺขํ อุปสงฺกมิตฺวา เอวมสฺส วจนียา “สุณสิ อิตฺถนฺนาเม อยํ เต สจฺจกาโล ภูตกาโล, ยํ ชาตํ ตํ สํฆมชฺเฌ ปุจฺฉนฺเต สนฺตํ ‘อตฺถี’ติ วตฺตพฺพํ, อสนฺตํ ‘นตฺถี’ติ วตฺตพฺพํ, มา โข วิตฺถายิ\csromansharedfootnote{cd3352a64e5098ce}{วิตฺถาสิ (ก)}, มา โข มงฺกุ อโหสิ, เอวํ ตํ ปุจฺฉิสฺสนฺติ\csromandash{}นสิ อนิมิตฺตา, นสิ นิมิตฺตมตฺตา, นสิ อโลหิตา, นสิ ธุวโลหิตา, นสิ ธุวโจฬา, นสิ ปคฺฆรนฺตี, นสิ สิขรณี, นสิ อิตฺถิปณฺฑกา, นสิ เวปุริสิกา, นสิ สมฺภินฺนา, นสิ อุภโตพฺยญฺชนา. สนฺติ เต เอวรูปา อาพาธา กุฏฺฐํ คณฺโฑ กิลาโส โสโส อปมาโร, มนุสฺสาสิ, อิตฺถีสิ, ภุชิสฺสาสิ, อณณาสิ, นสิ ราชภฏี, อนุญฺญาตาสิ มาตาปิตูหิ สามิเกน, ปริปุณฺณ\-วีสติวสฺสาสิ, ปริปุณฺณํ เต ปตฺตจีวรํ กินฺนามาสิ, กานามา เต ปวตฺตินี”ติ.}

\prose{เอกโต อาคจฺฉนฺติ. น เอกโต อาคนฺตพฺพํ, อนุสาสิกาย ปฐมตรํ อาคนฺตฺวา สํโฆ ญาเปตพฺโพ “สุณาตุ เม อยฺเย สํโฆ, อิตฺถนฺนามา อิตฺถนฺนามาย อยฺยาย อุปสมฺปทา\-เปกฺขา, อนุสิฏฺฐา สา มยา, ยทิ สํฆสฺส ปตฺตกลฺลํ, อิตฺถนฺนามา อาคจฺเฉยฺยา”ติ. "อาคจฺฉาหี”ติ วตฺตพฺพา.}

\prose{เอกํสํ อุตฺตราสงฺคํ การาเปตฺวา ภิกฺขุนีนํ ปาเท วนฺทาเปตฺวา อุกฺกุฏิกํ นิสีทาเปตฺวา อญฺชลิํ ปคฺคณฺหาเปตฺวา อุปสมฺปทํ ยาจาเปตพฺพา “สํฆํ อยฺเย อุปสมฺปทํ ยาจามิ, อุลฺลุมฺปตุ มํ อยฺเย สํโฆ อนุกมฺปํ อุปาทาย. ทุติยมฺปิ อยฺเย สํฆํ อุปสมฺปทํ ยาจามิ, อุลฺลุมฺปตุ มํ อยฺเย สํโฆ อนุกมฺปํ อุปาทาย. ตติยมฺปิ อยฺเย สํฆํ อุปสมฺปทํ ยาจามิ, อุลฺลุมฺปตุ มํ อยฺเย สํโฆ อนุกมฺปํ อุปาทายา”ติ. พฺยตฺตาย ภิกฺขุนิยา ปฏิพลาย สํโฆ ญาเปตพฺโพ\csromandash{}}

\prose{“สุณาตุ เม อยฺเย สํโฆ, อยํ อิตฺถนฺนามา อิตฺถนฺนามาย อยฺยาย อุปสมฺปทา\-เปกฺขา, ยทิ สํฆสฺส ปตฺตกลฺลํ, อหํ อิตฺถนฺนามํ อนฺตรายิเก ธมฺเม ปุจฺเฉยฺยนฺ”ติ.}

\prose{\csromanfolio{466}สุณสิ อิตฺถนฺนาเม อยํ เต สจฺจกาโล ภูตกาโล, ยํชาตํ ตํ ปุจฺฉามิ, สนฺตํ ‘อตฺถี’ติ วตฺตพฺพํ, อสนฺตํ ‘นตฺถี’ติ วตฺตพฺพํ, นสิ อนิมิตฺตา ฯเปฯ กินฺนามาสิ, กานามา เต ปวตฺตินี”ติ. พฺยตฺตาย ภิกฺขุนิยา ปฏิพลาย สํโฆ ญาเปตพฺโพ\csromandash{}}

\proseitem{425}{“สุณาตุ เม อยฺเย สํโฆ, อยํ อิตฺถนฺนามา อิตฺถนฺนามาย อยฺยาย อุปสมฺปทา\-เปกฺขา, ปริสุทฺธา อนฺตรายิเกหิ ธมฺเมหิ, ปริปุณฺณสฺสา ปตฺตจีวรํ, อิตฺถนฺนามา สํฆํ อุปสมฺปทํ ยาจติ อิตฺถนฺนามาย อยฺยาย ปวตฺตินิยา, ยทิ สํฆสฺส ปตฺตกลฺลํ, สํโฆ อิตฺถนฺนามํ อุปสมฺปาเทยฺย อิตฺถนฺนามาย อยฺยาย ปวตฺตินิยา, เอสา ญตฺติ.}

\prose{สุณาตุ เม อยฺเย สํโฆ, อยํ อิตฺถนฺนามา อิตฺถนฺนามาย อยฺยาย อุปสมฺปาทาเปกฺขา, ปริสุทฺธา อนฺตรายิเกหิ ธมฺเมหิ, ปริปุณฺณสฺสา ปตฺตจีวรํ, อิตฺถนฺนามา สํฆํ อุปสมฺปทํ ยาจติ อิตฺถนฺนามาย อยฺยาย ปวตฺตินิยา, สํโฆ อิตฺถนฺนามํ อุปสมฺปาเทติ อิตฺถนฺนามาย อยฺยาย ปวตฺตินิยา ยสฺสา อยฺยาย ขมติ อิตฺถนฺนามาย อุปสมฺปทา อิตฺถนฺนามาย อยฺยาย ปวตฺตินิยา, สา ตุณฺหสฺส, ยสฺสา นกฺขมติ, สา ภาเสยฺย.}

\prose{ทุติยมฺปิ เอตมตฺถํ วทามิ ฯเปฯ. ตติยมฺปิ เอตมตฺถํ วทามิ. สุณาตุ เม อยฺเย สํโฆ, อยํ อิตฺถนฺนามา อิตฺถนฺนามาย อยฺยาย อุปสมฺปทา\-เปกฺขา, ปริสุทฺธา อนฺตรายิเกหิ ธมฺเมหิ, ปริปุณฺณสฺสา ปตฺตจีวรํ, อิตฺถนฺนามา สํฆํ อุปสมฺปทํ ยาจติ อิตฺถนฺนามาย อยฺยาย ปวตฺตินิยา, สํโฆ อิตฺถนฺนามํ อุปสมฺปาเทติ อิตฺถนฺนามาย อยฺยาย ปวตฺตินิยา, ยสฺสา อยฺยาย ขมติ อิตฺถนฺนามาย อุปสมฺปทา อิตฺถนฺนามาย อยฺยาย ปวตฺตินิยา, สา ตุณฺหสฺส, ยสฺสา นกฺขมติ, สา ภาเสยฺย.}

\prose{อุปสมฺปนฺนา สํเฆน อิตฺถนฺนามา อิตฺถนฺนามาย อยฺยาย ปวตฺตินิยา, ขมติ สํฆสฺส, ตสฺมา ตุณฺหี, เอวเมตํ ธาเรยามี”ติ.}

\prose{ตาวเทว ตํ อาทาย ภิกฺขุสํฆํ อุปสงฺกมิตฺวา เอกํสํ อุตฺตราสงฺคํ การาเปตฺวา ภิกฺขูนํ ปาเท วนฺทาเปตฺวา อุกฺกุฏิกํ นิสีทาเปตฺวา อญฺชลิํ ปคฺคณฺหาเปตฺวา อุปสมฺปทํ ยาจาเปตพฺพา “อหํ อยฺยา อิตฺถนฺนามา อิตฺถนฺนามาย อยฺยาย อุปสมฺปทา\-เปกฺขา เอกโต\-อุปสมฺปนฺนา ภิกฺขุนิ\-สํเฆ วิสุทฺธา, สํฆํ อยฺยา อุปสมฺปทํ ยาจามิ, อุลฺลุมฺปตุมํ อยฺยา สํโฆ อนุกมฺปํ อุปาทาย. อหํ อยฺยา อิตฺถนฺนามา อิตฺถนฺนามาย \csromanfolio{467}อยฺยาย อุปสมฺปทา\-เปกฺขา เอกโต\-อุปสมฺปนฺนา ภิกฺขุนิ\-สํเฆ วิสุทฺธา, ทุติยมฺปิ อยฺยา สํฆํ อุปสมฺปทํ ยาจามิ, อุลฺลุมฺปตุ มํ อยฺยา สํโฆ อนุกมฺปํ อุปาทาย. อหํ อยฺยา อิตฺถนฺนามา อิตฺถนฺนามาย อยฺยาย อุปสมฺปทา\-เปกฺขา เอกโต\-อุปสมฺปนฺนา ภิกฺขุนิ\-สํเฆ วิสุทฺธา, ตติยมฺปิ อยฺยา สํฆํ อุปสมฺปทํ ยาจามิ, อุลฺลุมฺปตุ มํ อยฺยา สํโฆ อนุกมฺปํ อุปาทายา”ติ. พฺยตฺเตน ภิกฺขุนา ปฏิพเลน สํโฆ ญาเปตพฺโพ\csromandash{}}

\prose{“สุณาตุ เม ภนฺเต สํโฆ, อยํ อิตฺถนฺนามา อิตฺถนฺนามาย อุปสมฺปทา\-เปกฺขา เอกโต\-อุปสมฺปนฺนา ภิกฺขุนิ\-สํเฆ วิสุทฺธา, อิตฺถนฺนามา สํฆํ อุปสมฺปทํ ยาจติ อิตฺถนฺนามาย ปวตฺตินิยา, ยทิ สํฆสฺส ปตฺตกลฺลํ, สํโฆ อิตฺถนฺนามํ อุปสมฺปาเทยฺย อิตฺถนฺนามาย ปวตฺตินิยา, เอสา ญตฺติ.}

\prose{สุณาตุ เม ภนฺเต สํโฆ, อยํ อิตฺถนฺนามา อิตฺถนฺนามาย อุปสมฺปทา\-เปกฺขา เอกโต\-อุปสมฺปนฺนา ภิกฺขุนิ\-สํเฆ วิสุทฺธา, อิตฺถนฺนามา สํฆํ อุปสมฺปทํ ยาจติ อิตฺถนฺนามาย ปวตฺตินิยา, สํโฆ อิตฺถนฺนามํ อุปสมฺปาเทติ อิตฺถนฺนามาย ปวตฺตินิยา, ยสฺสายสฺมโต ขมติ อิตฺถนฺนามาย อุปสมฺปทา อิตฺถนฺนามาย ปวตฺตินิยา, โส ตุณฺหสฺส, ยสฺส นกฺขมติ, โส ภาเสยฺย.}

\prose{ทุติยมฺปิ เอตมตฺถํ วทามิ ฯเปฯ. ตติยมฺปิ เอตมตฺถํ วทามิ. สุณาตุ เม ภนฺเต สํโฆ, อยํ อิตฺถนฺนามา อิตฺถนฺนามาย อุปสมฺปทา\-เปกฺขา เอกโต\-อุปสมฺปนฺนา ภิกฺขุนิ\-สํเฆ วิสุทฺธา, อิตฺถนฺนามา สํฆํ อุปสมฺปทํ ยาจติ อิตฺถนฺนามาย ปวตฺตินิยา, สํโฆ อิตฺถนฺนามํ อุปสมฺปาเทติ อิตฺถนฺนามาย ปวตฺตินิยา, ยสฺสายสฺมโต ขมติ อิตฺถนฺนามาย อุปสมฺปทา อิตฺถนฺนามาย ปวตฺตินิยา, โส ตุณฺหสฺส, ยสฺส นกฺขมติ, โส ภาเสยฺย.}

\prose{อุปสมฺปนฺนา สํเฆน อิตฺถนฺนามา อิตฺถนฺนามาย ปวตฺตินิยา, ขมติ สํฆสฺส, ตสฺมา ตุณฺหี, เอวเมตํ ธารยามี”ติ.}

\prose{ตาวเทว ฉายา เมตพฺพา, อุตุปฺปมาณํ อาจิกฺขิตพฺพํ, ทิวสภาโค อาจิกฺขิตพฺโพ, สงฺคีติ อาจิกฺขิตพฺพา, ภิกฺขุนิโย วตฺตพฺพา “อิมิสฺสา ตโย จ นิสฺสเย อฏฺฐ จ อกรณียานิ อาจิกฺเขยฺยาถา”ติ.}

\proseitem{426}{เตน โข ปน สมเยน ภิกฺขุนิโย ภตฺตคฺเค อาสนํ สํกสายนฺติโย กาลํ วีตินาเมสุํ. ภควโต เอตมตฺถํ \csromanfolio{468}อาโรเจสุํ. อนุชานามิ ภิกฺขเว อฏฺฐนฺนํ ภิกฺขุนีนํ ยถาวุฑฺฒํ อวเสสานํ ยถาคติกนฺติ.}

\prose{เตน โข ปน สมเยน ภิกฺขุนิโย “ภควตา อนุญฺญาตํ อฏฺฐนฺนํ ภิกฺขุนีนํ ยถาวุฑฺฒํ อวเสสานํ ยถาคติกนฺ”ติ สพฺพตฺถ อฏฺเฐว ภิกฺขุนิโย ยถาวุฑฺฒํ ปฏิพาหนฺติ อวเสสาโย ยถาคติกํ. ภควโต เอตมตฺถํ อาโรเจสุํ. อนุชานามิ ภิกฺขเว ภตฺตคฺเค อฏฺฐนฺนํ ภิกฺขุนีนํ ยถาวุฑฺฒํ อวเสสานํ ยถาคติกํ, อญฺญตฺถ สพฺพตฺถ ยถาวุฑฺฒํ\csromansharedfootnote{9f88945b9079be01}{อญฺญตฺถ ยถาวุฑฺฒํ (สฺยา)} น ปฏิพาหิตพฺพํ, ยา ปฏิพาเหยฺย, อาปตฺติ ทุกฺกฏสฺสาติ.}

\proseitem{427}{\csromansymbolfootnote{*}{วิ 2. 412-3 ปิฏฺเฐสุปิ.}เตน โข ปน สมเยน ภิกฺขุนิโย น ปวาเรนฺติ. ภควโต เอตมตฺถํ อาโรเจสุํ. น ภิกฺขเว ภิกฺขุนิยา น ปวาเรตพฺพํ, ยา น ปวาเรยฺย, ยถาธมฺโม กาเรตพฺโพติ.}

\prose{เตน โข ปน สมเยน ภิกฺขุนิโย อตฺตนา ปวาเรตฺวา ภิกฺขุสํฆํ น ปเวเรนฺติ. ภควโต เอตมตฺถํ อาโรเจสุํ. น ภิกฺขเว ภิกฺขุนิยา อตฺตนา ปวาเรตฺวา ภิกฺขุสํโฆ น ปวาเรตพฺโพ, ยา น ปวาเรยฺย, ยถาธมฺโม กาเรตพฺโพติ.}

\prose{เตน โข ปน สมเยน ภิกฺขุนิโย ภิกฺขูหิ สทฺธิํ เอกโต ปวาเรนฺติโย โกลาหลํ อกํสุ. ภควโต เอตมตฺถํ อาโรเจสุํ. น ภิกฺขเว ภิกฺขุนิยา ภิกฺขูหิ สทฺธิํ เอกโต ปวาเรตพฺพํ, ยา ปวาเรยฺย, อาปตฺติ ทุกฺกฏสฺสาติ.}

\prose{เตน โข ปน สมเยน ภิกฺขุนิโย ปุเรภตฺตํ ปวาเรนฺติโย กาลํ วีตินาเมสุํ. ภควโต เอตมตฺถํ อาโรเจสุํ. อนุชานามิ ภิกฺขเว ปจฺฉาภตฺตํ\csromansharedfootnote{ff6c864f9d25b2f0}{ภิกฺขเว ภิกฺขุนิยา ปจฺฉาภตฺตํ (สฺยา, กํ)} ปวาเรตุนฺติ. ปจฺฉาภตฺตํ ปวาเรนฺติโย วิกาเล\csromansharedfootnote{9bd1582de858d7a2}{วิกาลํ (ก)} อเหสุํ, ภควโต เอตมตฺถํ อาโรเจสุํ. อนุชานามิ ภิกฺขเว อชฺชตนา ภิกฺขุนิสํฆํ\csromansharedfootnote{dcd99e140d3dfbc1}{อชฺชตนา (สี), อชฺชตนาย (ก)} ปวาเรตฺวา อปรชฺชุ ภิกฺขุสํฆํ ปวาเรตุนฺติ.}

\prose{เตน โข ปน สมเยน สพฺโพ ภิกฺขุนิ\-สํโฆ ปวาเรนฺโต โกลาหลํ อกาสิ. ภควโต เอตมตฺถํ อาโรเจสุํ. อนุชานามิ ภิกฺขเว เอกํ ภิกฺขุนิํ พฺยตฺตํ ปฏิพลํ สมฺมนฺนิตุํ ภิกฺขุนิ\-สํฆสฺส อตฺถาย \csromanfolio{469}ภิกฺขุสํฆํ ปวาเรตุํ. เอวญฺจ ปน ภิกฺขเว สมฺมนฺนิตพฺพา, ปฐมํ ภิกฺขุนี ยาจิตพฺพา, ยาจิตฺวา พฺยตฺตาย ภิกฺขุนิยา ปฏิพลาย สํโฆ ญาเปตพฺโพ\csromandash{}}

\prose{“สุณาตุ เม อยฺเย สํโฆ, ยทิ สํฆสฺส ปตฺตกลฺลํ, สํโฆ อิตฺถนฺนามํ ภิกฺขุนิํ สมฺมนฺเนยฺย ภิกฺขุนิ\-สํฆสฺส อตฺถาย ภิกฺขุสํฆํ ปวาเรตุํ, เอสา ญตฺติ.}

\prose{สุณาตุ เม อยฺเย สํโฆ, สํโฆ อิตฺถนฺนามํ ภิกฺขุนิํ สมฺมนฺนติ ภิกฺขุนิ\-สํฆสฺส อตฺถาย ภิกฺขุสํฆํ ปวาเรตุํ, ยสฺสา อยฺยาย ขมติ อิตฺถนฺนามาย ภิกฺขุนิยา สมฺมุติ ภิกฺขุนิ\-สํฆสฺส อตฺถาย ภิกฺขุสํฆํ ปวาเรตุํ, สา ตุณฺหสฺส, ยสฺสา นกฺขมติ, สา ภาเสยฺย.}

\prose{สมฺมตา สํเฆน อิตฺถนฺนามา ภิกฺขุนี ภิกฺขุนิ\-สํฆสฺส อตฺถาย ภิกฺขุสํฆํ ปวาเรตุํ, ขมติ สํฆสฺส, ตสฺมา ตุณฺหี, เอวเมตํ ธารยามี”ติ.}

\prose{ตาย สมฺมตาย ภิกฺขุนิยา ภิกฺขุนิสํฆํ อาทาย ภิกฺขุสํฆํ อุปสงฺกมิตฺวา เอกํสํ อุตฺตราสงฺคํ กริตฺวา ภิกฺขูนํ ปาเท วนฺทิตฺวา อุกฺกุฏิกํ นิสีทิตฺวา อญฺชลิํ ปคฺคเหตฺวา เอวมสฺส วจนีโย “ภิกฺขุนิ\-สํโฆ อยฺยา ภิกฺขุสํฆํ ปวาเรติ ทิฏฺเฐน วา สุเตน วา ปริสงฺกาย วา, วทตุ\csromansharedfootnote{ab2788802a78e18c}{วทตุ มํ (ก)} อยฺยา ภิกฺขุสํโฆ ภิกฺขุนิสํฆํ อนุกมฺปํ อุปาทาย, ปสฺสนฺโต ปฏิกริสฺสติ, ทุติยมฺปิ อยฺยา ฯเปฯ. ตติยมฺปิ อยฺยา ภิกฺขุนิ\-สํโฆ ภิกฺขุสํฆํ ปวาเรติ ทิฏฺเฐน วา สุเตน วา ปริสงฺกาย วา, วทตุ อยฺยา ภิกฺขุสํโฆ ภิกฺขุนิสํฆํ อนุกมฺปํ อุปาทาย, ปสฺสนฺโต ปฏิกริสฺสตี”ติ.}

\proseitem{428}{เตน โข ปน สมเยน ภิกฺขุนิโย ภิกฺขูนํ อุโปสถํ ฐเปนฺติ. ปวารณํ ฐเปนฺติ, สวจนียํ กโรนฺติ, อนุวาทํ ปฏฺฐเปนฺติ, โอกาสํ กาเรนฺติ, โจเทนฺติ, สาเรนฺติ. ภควโต เอตมตฺถํ อาโรเจสุํ. น ภิกฺขเว ภิกฺขุนิยา ภิกฺขุสฺส อุโปสโถ ฐเปตพฺโพ, ฐปิโตปิ อฏฺฐปิโต, ฐเปนฺติยา อาปตฺติ ทุกฺกฏสฺส. น ปวารณา ฐเปตพฺพา, ฐปิตาปิ อฏฺฐปิตา, ฐเปนฺติยา อาปตฺติ ทุกฺกฏสฺส. น สวจนียํ กาตพฺพํ, กตํปิ อกตํ, กโรนฺติยา อาปตฺติ ทุกฺกฏสฺส. น อนุวาโท ปฏฺฐเปตพฺโพ, ปฏฺฐปิโตปิ อปฺปฏฺฐปิโต, ปฏฺฐเปนฺติยา อาปตฺติ ทุกฺกฏสฺส. น โอกาโส \csromanfolio{470}กาเรตพฺโพ, การิโตปิ อการิโต, กาเรนฺติยา อาปตฺติ ทุกฺกฏสฺส. น โจเทตพฺโพ, โจทิโตปิ อโจทิโต, โจเทนฺติยา อาปตฺติ ทุกฺกฏสฺส. น สาเรตพฺโพ, สาริโตปิ อสาริโต, สาเรนฺติยา อาปตฺติ ทุกฺกฏสฺสาติ.}

\prose{เตน โข ปน สมเยน ภิกฺขู ภิกฺขุนีนํ อุโปสถํ ฐเปนฺติ, ปวารณํ ฐเปนฺติ, สวจนียํ กโรนฺติ, อนุวาทํ ปฏฺฐเปนฺติ, โอกาสํ กาเรนฺติ, โจเทนฺติ, สาเรนฺติ. ภควโต เอตมตฺถํ อาโรเจสุํ. อนุชานามิ ภิกฺขเว ภิกฺขุนา ภิกฺขุนิยา อุโปสถํ ฐเปตุํ, ฐปิโตปิ สุฏฺฐปิโต, ฐเปนฺตสฺส อนาปตฺติ. ปวารณํ ฐเปตุํ, ฐปิตาปิ สุฏฺฐปิตา, ฐเปนฺตสฺส อนาปตฺติ. สวจนียํ กาตุํ, กตมฺปิ สุกตํ, กโรนฺตสฺส อนาปตฺติ. อนุวาทํ ปฏฺฐเปตุํ, ปฏฺฐปิโตปิ สุปฺปฏฺฐปิโต, ปฏฺฐเปนฺตสฺส อนาปตฺติ. โอกาสํ กาเรตุํ, การิโตปิ สุการิโต, กาเรนฺตสฺส อนาปตฺติ. โจเทตุํ, โจทิตาปิ สุโจทิตา, โจเทนฺตสฺส อนาปตฺติ. สาเรตุํ, สาริตาปิ สุสาริตา, สาเรนฺตสฺส อนาปตฺตีติ.}

\proseitem{429}{\csromansymbolfootnote{*}{วิ 2. 452 ปิฏฺเฐปิ.}เตน โข ปน สมเยน ฉพฺพคฺคิยา ภิกฺขุนิโย ยาเนน ยายนฺติ, อิตฺถิยุตฺเตนปิ ปุริสนฺตเรน ปุริสยุตฺเตนปิ อิตฺถนฺตเรน. มนุสฺสา อุชฺฌายนฺติ ขิยฺยนฺติ วิปาเจนฺติ “ฯเปฯ เสยฺยถาปิ คงฺคามหิยายา”ติ\csromansharedfootnote{809610fab96aa85d}{คงฺคามหิยาติ (สี)}. ภควโต เอตมตฺถํ อาโรเจสุํ. น ภิกฺขเว ภิกฺขุนิยา ยาเนน ยายิตพฺพํ, ยา ยาเยยฺย, ยถาธมฺโม กาเรตพฺโพติ.}

\prose{เตน โข ปน สมเยน อญฺญตรา ภิกฺขุนี คิลานา โหติ, น สกฺโกติ ปทสา คนฺตุํ. ภควโต เอตมตฺถํ อาโรเจสุํ. อนุชานามิ ภิกฺขเว คิลานาย ยานนฺติ. อถ โข ภิกฺขุนีนํ เอตทโหสิ “อิตฺถิยุตฺตํ นุ โข ปุริสยุตฺตํ นุ โข”ติ. ภควโต เอตมตฺถํ อาโรเจสุํ. อนุชานามิ ภิกฺขเว อิตฺถิยุตฺตํ ปุริสยุตฺตํ หตฺถวฏฺฏกนฺติ.}

\prose{เตน โข ปน สมเยน อญฺญตริสฺสา ภิกฺขุนิยา ยานุคฺฆาเตน พาฬฺหตรํ อผาสุ อโหสิ. ภควโต เอตมตฺถํ อาโรเจสุํ. อนุชานามิ ภิกฺขเว สิวิกํ ปาฏงฺกินฺติ.}

\proseitem{430}{\csromanfolio{471}เตน โข ปน สมเยน อฑฺฒกาสี คณิกา ภิกฺขุนีสุ ปพฺพชิตา โหติ, สา จ สาวตฺถิํ คนฺตุกามา โหติ “ภควโต สนฺติเก อุปสมฺปชฺชิสฺสามี”ติ. อสฺโสสุํ โข ธุตฺตา “อฑฺฒกาสี กิร คณิกา สาวตฺถิํ คนฺตุกามา”ติ. เต มคฺเค ปริยุฏฺฐิํสุ. อสฺโสสิ โข อฑฺฒกาสี คณิกา “ธุตฺตา กิร มคฺเค ปริยุฏฺฐิตา”ติ. ภควโต สนฺติเก ทูตํ ปาเหสิ “อหํ หิ อุปสมฺปชฺชิตุ\-กามา, กถํ นุ โข มยา ปฏิปชฺชิตพฺพนฺ”ติ. อถ โข ภควา เอตสฺมิํ นิทาเน เอตสฺมิํ ปกรเณ ธมฺมิํ กถํ กตฺวา ภิกฺขู อามนฺเตสิ “อนุชานามิ ภิกฺขเว ทูเตนปิ อุปสมฺปาเทตุนฺ”ติ.}

\prose{ภิกฺขุทูเตน อุปสมฺปาเทนฺติ. ภควโต เอตมตฺถํ อาโรเจสุํ. น ภิกฺขเว ภิกฺขุทูเตน อุปสมฺ\-ปาเท\-ตพฺพา, โย อุปสมฺปาเทยฺย, อาปตฺติ ทุกฺกฏสฺสาติ. สิกฺขามานทูเตน อุปสมฺปาเทนฺติ ฯเปฯ. สามเณรทูเตน อุปสมฺปาเทนฺติ ฯเปฯ. สามเณริทูเตน อุปสมฺปาเทนฺติ ฯเปฯ. พาลาย อพฺยตฺตาย ทูเตน อุปสมฺปาเทนฺติ. น ภิกฺขเว พาลาย อพฺยตฺตาย ทูเตน อุปสมฺ\-ปาเท\-ตพฺพา, โย อุปสมฺปาเทยฺย, อาปตฺติ ทุกฺกฏสฺส. อนุชานามิ ภิกฺขเว พฺยตฺตาย ภิกฺขุนิยา ปฏิพลาย ทูเตน อุปสมฺปาเทตุนฺติ.}

\prose{ตาย ทูตาย ภิกฺขุนิยา สํฆํ อุปสงฺกมิตฺวา เอกํสํ อุตฺตราสงฺคํ กริตฺวา ภิกฺขูนํ ปาเท วนฺทิตฺวา อุกฺกุฏิกํ นิสีทิตฺวา อญฺชลิํ ปคฺคเหตฺวา เอวมสฺส วจนีโย “อิตฺถนฺนามา อยฺยา อิตฺถนฺนามาย อยฺยาย อุปสมฺปทา\-เปกฺขา เอกโต\-อุปสมฺปนฺนา ภิกฺขุนิ\-สํเฆ วิสุทฺธา, สา เกนจิเทว อนฺตราเยน น อาคจฺฉติ, อิตฺถนฺนามา อยฺยา สํฆํ อุปสมฺปทํ ยาจติ, อุลฺลุมฺปตุ ตํ อยฺยา สํโฆ อนุกมฺปํ อุปาทาย. อิตฺถนฺนามา อยฺยา อิตฺถนฺนามาย อยฺยาย อุปสมฺปทา\-เปกฺขา เอกโต\-อุปสมฺปนฺนา ภิกฺขุนิ\-สํเฆ วิสุทฺธา, สา เกนจิเทว อนฺตราเยน น อาคจฺฉติ, ทุติยมฺปิ อยฺยา อิตฺถนฺนามา สํฆํ อุปสมฺปทํ ยาจติ, อุลฺลุมฺปตุ ตํ อยฺยา สํโฆ อนุกมฺปํ อุปาทาย. อิตฺถนฺนามา อยฺยา อิตฺถนฺนามาย อยฺยาย อุปสมฺปทา\-เปกฺขา เอกโต\-อุปสมฺปนฺนา ภิกฺขุนิ\-สํเฆ วิสุทฺธา, สา เกนจิเทว อนฺตราเยน น อาคจฺฉติ, ตติยมฺปิ อยฺยา อิตฺถนฺนามา สํฆํ อุปสมฺปทํ ยาจติ, อุลฺลุมฺปตุ ตํ อยฺยา สํโฆ อนุกมฺปํ อุปาทายา”ติ. พฺยตฺเตน ภิกฺขุนา ปฏิพเลน สํโฆ ญาเปตพฺโพ\csromandash{} \csromanfolio{472}“สุณาตุ เม ภนฺเต สํโฆ, อิตฺถนฺนามา อิตฺถนฺนามาย อุปสมฺปทา\-เปกฺขา เอกโต\-อุปสมฺปนฺนา ภิกฺขุนิ\-สํเฆ วิสุทฺธา, สา เกนจิเทว อนฺตราเยน น อาคจฺฉติ, อิตฺถนฺนามา สํฆํ อุปสมฺปทํ ยาจติ อิตฺถนฺนามาย ปวตฺตินิยา, ยทิ สํฆสฺส ปตฺตกลฺลํ, สํโฆ อิตฺถนฺนามํ อุปสมฺปาเทยฺย อิตฺถนฺนามาย ปวตฺตินิยา, เอสา ญตฺติ.}

\prose{สุณาตุ เม ภนฺเต สํโฆ, อิตฺถนฺนามาย อุปสมฺปทา\-เปกฺขา เอกโต\-อุปสมฺปนฺนา ภิกฺขุนิ\-สํเฆ วิสุทฺธา, สา เกนจิเทว อนฺตราเยน น อาคจฺฉติ, อิตฺถนฺนามา สํฆํ อุปสมฺปทํ ยาจติ อิตฺถนฺนามาย ปวตฺตินิยา, สํโฆ อิตฺถนฺนามํ อุปสมฺปาเทติ อิตฺถนฺนามาย ปวตฺตินิยา, ยสฺสายสฺมโต ขมติ อิตฺถนฺนามาย อุปสมฺปทา อิตฺถนฺนามาย ปวตฺตินิยา, โส ตุณฺหสฺส, ยสฺส นกฺขมติ, โส ภาเสยฺย.}

\prose{ทุติยมฺปิ เอตมตฺถํ วทามิ ฯเปฯ. ตติยมฺปิ เอตมตฺถํ วทามิ. สุณาตุ เม ภนฺเต สํโฆ, อิตฺถนฺนามา อิตฺถนฺนามาย อุปสมฺปทา\-เปกฺขา เอกโต\-อุปสมฺปนฺนา ภิกฺขุนิ\-สํเฆ วิสุทฺธา, สา เกนจิเทว อนฺตราเยน น อาคจฺฉติ, อิตฺถนฺนามา สํฆํ อุปสมฺปทํ ยาจติ อิตฺถนฺนามาย ปวตฺตินิยา, สํโฆ อิตฺถนฺนามํ อุปสมฺปาเทติ อิตฺถนฺนามาย ปวตฺตินิยา, ยสฺสายสฺมโต ขมติ อิตฺถนฺนามาย อุปสมฺปทา อิตฺถนฺนามาย ปวตฺตินิยา, โส ตุณฺหสฺส, ยสฺส นกฺขมติ, โส ภาเสยฺย.}

\prose{อุปสมฺปนฺนา สํเฆน อิตฺถนฺนามา อิตฺถนฺนามาย ปวตฺตินิยา, ขมติ สํฆสฺส, ตสฺมา ตุณฺหี, เอวเมตํ ธารยามี”ติ.}

\prose{ตาวเทว ฉายา เมตพฺพา, อุตุปฺปมาณํ อาจิกฺขิตพฺพํ, ทิวสภาโค อาจิกฺขิตพฺโพ, สงฺคีติ อาจิกฺขิตพฺพา, ภิกฺขุนิโย วตฺตพฺพา “ตสฺสา ตโย จ นิสฺสเย อฏฺฐ จ อกรณียานิ อาจิกฺเขยฺยาถา”ติ.}

\proseitem{431}{เตน โข ปน สมเยน ภิกฺขุนิโย อรญฺเญ วิหรนฺติ, ธุตฺตา ทูเสนฺติ. ภควโต เอตมตฺถํ อาโรเจสุํ. น ภิกฺขเว ภิกฺขุนิยา อรญฺเญ วตฺถพฺพํ, ยา วเสยฺย, อาปตฺติ ทุกฺกฏสฺสาติ.}

\prose{เตน โข ปน สมเยน อญฺญตเรน อุปาสเกน ภิกฺขุนิ\-สํฆสฺส อุโทสิโต ทินฺโน โหติ. ภควโต เอตมตฺถํ อาโรเจสุํ. อนุชานามิ ภิกฺขเว อุโทสิตนฺติ. อุโทสิโต น สมฺมติ. ภควโต \csromanfolio{473}เอตมตฺถํ อาโรเจสุํ. อนุชานามิ ภิกฺขเว อุปสฺสยนฺติ. อุปสฺสโย น สมฺมติ. ภควโต เอตมตฺถํ อาโรเจสุํ. อนุชานามิ ภิกฺขเว นวกมฺมนฺติ. นวกมฺมํ น สมฺมติ. ภควโต เอตมตฺถํ อาโรเจสุํ. อนุชานามิ ภิกฺขเว ปุคฺคลิกมฺปิ กาตุนฺติ.}

\proseitem{432}{เตน โข ปน สมเยน อญฺญตรา อิตฺถี สนฺนิสินฺน\-คพฺภา ภิกฺขุนีสุ ปพฺพชิตา โหติ, ตสฺสา ปพฺพชิตาย คพฺโภ วุฏฺฐาติ. อถ โข ตสฺสา ภิกฺขุนิยา เอตทโหสิ “กถํ นุ โข มยา อิมสฺมิํ ทารเก ปฏิปชฺชิตพฺพนฺ”ติ. ภควโต เอตมตฺถํ อาโรเจสุํ. อนุชานามิ ภิกฺขเว โปเสตุํ ยาว โส ทารโก วิญฺญุตํ ปาปุณาตีติ.}

\prose{อถ โข ตสฺสา ภิกฺขุนิยา เอตทโหสิ “มยา จ น ลพฺภา เอกิกาย วตฺถุํ, อญฺญาย จ ภิกฺขุนิยา น ลพฺภา ทารเกน สห วตฺถุํ, กถํ นุ โข มยา ปฏิปชฺชิตพฺพนฺ”ติ. ภควโต เอตมตฺถํ อาโรเจสุํ. อนุชานามิ ภิกฺขเว เอกํ ภิกฺขุนิํ สมฺมนฺนิตฺวา ตสฺสา ภิกฺขุนิยา ทุติยํ ทาตุํ. เอวญฺจ ปน ภิกฺขเว สมฺมนฺนิตพฺพา, ปฐมํ ภิกฺขุนี ยาจิตพฺพา, ยาจิตฺวา พฺยตฺตาย ภิกฺขุนิยา ปฏิพลาย สํโฆ ญาเปตพฺโพ\csromandash{}}

\prose{“สุณาตุ เม อยฺเย สํโฆ, ยทิ สํฆสฺส ปตฺตกลฺลํ, สํโฆ อิตฺถนฺนามํ ภิกฺขุนิํ สมฺมนฺเนยฺย อิตฺถนฺนามาย ภิกฺขุนิยา ทุติยํ, เอสา ญตฺติ.}

\prose{สุณาตุ เม อยฺเย สํโฆ, สํโฆ อิตฺถนฺนามํ ภิกฺขุนิํ สมฺมนฺนติ อิตฺถนฺนามาย ภิกฺขุนิยา ทุติยํ, ยสฺสา อยฺยาย ขมติ อิตฺถนฺนามาย ภิกฺขุนิยา สมฺมุติ อิตฺถนฺนามาย ภิกฺขุนิยา ทุติยาย, สา ตุณฺหสฺส, ยสฺสา นกฺขมติ, สา ภาเสยฺย.}

\prose{สมฺมตา สํเฆน อิตฺถนฺนามา ภิกฺขุนี อิตฺถนฺนามาย ภิกฺขุนิยา ทุติยา, ขมติ สํฆสฺส, ตสฺมา ตุณฺหี, เอวเมตํ ธารยามี”ติ.}

\prose{อถ โข ตสฺสา ทุติยิกาย ภิกฺขุนิยา เอตทโหสิ “กถํ นุ โข มยา อิมสฺมิํ ทารเก ปฏิปชฺชิตพฺพนฺ”ติ. ภควโต เอตมตฺถํ อาโรเจสุํ. อนุชานามิ ภิกฺขเว ฐเปตฺวา สาคารํ ยถา อญฺญสฺมิํ ปุริเส ปฏิปชฺชนฺติ\csromansharedfootnote{1f17d05cfdb8e23f}{ปฏิปชฺชติ (สฺยา)}, เอวํ ตสฺมิํ ทารเก ปฏิปชฺชิตุนฺติ.}

\proseitem{433}{\csromanfolio{474}เตน โข ปน สมเยน อญฺญตรา ภิกฺขุนี ครุธมฺมํ อชฺฌาปนฺนา โหติ มานตฺตจารินี. อถ โข ตสฺสา ภิกฺขุนิยา เอตทโหสิ “มยา จ น ลพฺภา เอกิกาย วตฺถุํ, อญฺญาย จ ภิกฺขุนิยา น ลพฺภา สห มยา วตฺถุํ, กถํ นุ โข มยา ปฏิปชฺชิตพฺพนฺ”ติ. ภควโต เอตมตฺถํ อาโรเจสุํ. อนุชานามิ ภิกฺขเว เอกํ ภิกฺขุนิํ สมฺมนฺนิตฺวา ตสฺสา ภิกฺขุนิยา ทุติยํ ทาตุํ. เอวญฺจ ปน ภิกฺขเว สมฺมนฺนิตพฺพา, ปฐมํ ภิกฺขุนี ยาจิตพฺพา, ยาจิตฺวา พฺยตฺตาย ภิกฺขุนิยา ปฏิพลาย สํโฆ ญาเปตพฺโพ\csromandash{}}

\prose{“สุณาตุ เม อยฺเย สํโฆ, ยทิ สํฆสฺส ปตฺตกลฺลํ, สํโฆ อิตฺถนฺนามํ ภิกฺขุนิํ สมฺมนฺเนยฺย อิตฺถนฺนามาย ภิกฺขุนิยา ทุติยํ, เอสา ญตฺติ.}

\prose{สุณาตุ เม อยฺเย สํโฆ, สํโฆ, สํโฆ อิตฺถนฺนามํ ภิกฺขุนิํ สมฺมนฺนติ อิตฺถนฺนามาย ภิกฺขุนิยา ทุติยํ, ยสฺสา อยฺยาย ขมติ อิตฺถนฺนามาย ภิกฺขุนิยา สมฺมุติ อิตฺถนฺนามาย ภิกฺขุนิยา ทุติยาย, สา ตุณฺหสฺส, ยสฺสา นกฺขมติ, สา ภาเสยฺย.}

\prose{สมฺมตา สํเฆน อิตฺถนฺนามา ภิกฺขุนี อิตฺถนฺนามาย ภิกฺขุนิยา ทุติยา, ขมติ สํฆสฺส, ตสฺมา ตุณฺหี, เอวเมตํ ธารยามี”ติ.}

\proseitem{434}{เตน โข ปน สมเยน อญฺญตรา ภิกฺขุนี สิกฺขํ ปจฺจกฺขาย วิพฺภมิ, สา ปุน ปจฺจาคนฺตฺวา ภิกฺขุนิโย อุปสมฺปทํ ยาจิ. ภควโต เอตมตฺถํ อาโรเจสุํ. น ภิกฺขเว ภิกฺขุนิยา สิกฺขา\-ปจฺจกฺขานํ, ยเทว สา วิพฺภนฺตา, ตเทว สา อภิกฺขุนีติ.}

\prose{เตน โข ปน สมเยน อญฺญตรา ภิกฺขุนี สกาวาสา ติตฺถายตนํ สงฺกมิ, สา ปุน ปจฺจาคนฺตฺวา ภิกฺขุนิโย อุปสมฺปทํ ยาจิ. ภควโต เอตมตฺถํ อาโรเจสุํ. ยา สา ภิกฺขเว ภิกฺขุนี สกาวาสา ติตฺถายตนํ สงฺกนฺตา, สา อาคตา น อุปสมฺ\-ปาเท\-ตพฺพาติ.}

\prose{เตน โข ปน สมเยน ภิกฺขุนิโย ปุริเสหิ อภิวาทนํ เกสจฺเฉทนํ นขจฺเฉทนํ วณปฺปฏิกมฺมํ กุกฺกุจฺจายนฺตา น สาทิยนฺติ. ภควโต เอตมตฺถํ อาโรเจสุํ. อนุชานามิ ภิกฺขเว สาทิตุนฺติ\csromansharedfootnote{be64875954517031}{สาทิยิตุํ (ก)}.}

\proseitem{435}{\csromanfolio{475}เตน โข ปน สมเยน ภิกฺขุนิโย ปลฺลงฺเกน นิสีทนฺติ ปณฺหิ\-สมฺผสฺสํ สาทิยนฺตี\csromansharedfootnote{f889413f69e0f434}{สาทิยนฺตา (ก)}. ภควโต เอตมตฺถํ อาโรเจสุํ. น ภิกฺขเว ภิกฺขุนิยา ปลฺลงฺเกน นิสีทิตพฺพํ, ยา นิสีเทยฺย, อาปตฺติ ทุกฺกฏสฺสาติ.}

\prose{เตน โข ปน สมเยน อญฺญตรา ภิกฺขุนี คิลานา โหติ, ตสฺสา วินา ปลฺลงฺเกน น ผาสุ โหติ. ภควโต เอตมตฺถํ อาโรเจสุํ. อนุชานามิ ภิกฺขเว ภิกฺขุนิยา อฑฺฒปลฺลงฺกนฺติ.}

\prose{เตน โข ปน สมเยน ภิกฺขุนิโย วจฺจกุฏิยา วจฺจํ กโรนฺติ, ฉพฺพคฺคิยา ภิกฺขุนิโย ตตฺเถว คพฺภํ ปาเตนฺติ. ภควโต เอตมตฺถํ อาโรเจสุํ. น ภิกฺขเว ภิกฺขุนิยา วจฺจกุฏิยา วจฺโจ กาตพฺโพ, ยา กเรยฺย, อาปตฺติ ทุกฺกฏสฺส. อนุชานามิ ภิกฺขเว เหฏฺฐา วิวเฏ อุปริ ปฏิจฺฉนฺเน วจฺจํ กาตุนฺติ.}

\proseitem{436}{เตน โข ปน สมเยน ภิกฺขุนิโย จุณฺเณน นหายนฺติ. มนุสฺสา อุชฺฌายนฺติ ขิยฺยนฺติ วิปาเจนฺติ. “ฯเปฯ เสยฺยถาปิ คิหินี กามโภคินิโย”ติ. ภควโต เอตมตฺถํ อาโรเจสุํ. น ภิกฺขเว ภิกฺขุนิยา จุณฺเณน นหายิตพฺพํ, ยา นหาเยยฺย, อาปตฺติ ทุกฺกฏสฺส. อนุชานามิ ภิกฺขเว กุกฺกุสํ มตฺติกนฺติ.}

\prose{เตน โข ปน สมเยน ภิกฺขุนิโย วาสิตกาย มตฺติกาย นหายนฺติ. มนุสฺสา อุชฺฌยนฺติ ขิยฺยนฺติ วิปาเจนฺติ. “ฯเปฯ เสยฺยถาปิ คิหินี กามโภคินิโย”ติ. ภควโต เอตมตฺถํ อาโรเจสุํ. น ภิกฺขเว ภิกฺขุนิยา วาสิตกาย มตฺติกาย นหายิตพฺพํ, ยา นหาเยยฺย, อาปตฺติ ทุกฺกฏสฺส. อนุชานามิ ภิกฺขเว ปกติ\-มตฺติกนฺติ.}

\prose{เตน โข ปน สมเยน ภิกฺขุนิโย ชนฺตาฆเร นหายนฺติโย โกลาหลํ อกํสุ. ภควโต เอตมตฺถํ อาโรเจสุํ. น ภิกฺขเว ภิกฺขุนิยา ชนฺตาฆเร นหายิตพฺพํ, ยา นหาเยยฺย, อาปตฺติ ทุกฺกฏสฺสาติ.}

\csromanmatikaanchor{mtk.211}
\csromantocmarkref{section}{อุทฺทานคาถา}{mtk.211}
\bookmark[dest={mtk.211},level=1]{อุทฺทานคาถา}
\prose{เตน โข ปน สมเยน ภิกฺขุนิโย ปฏิโสเต นหายนฺติ ธารา\-สมฺผสฺสํ สาทิยนฺติ. ภควโต เอตมตฺถํ อาโรเจสุํ. น ภิกฺขเว \csromanfolio{476}ภิกฺขุนิยา ปฏิโสเต นหายิตพฺพํ, ยา นหาเยยฺย, อาปตฺติ ทุกฺกฏสฺสาติ.}

\prose{เตน โข ปน สมเยน ภิกฺขุนิโย อติตฺเถ นหายนฺติ, ธุตฺตา ทูเสนฺติ. ภควโต เอตมตฺถํ อาโรเจสุํ. น ภิกฺขเว ภิกฺขุนิยา อติตฺเถ นหายิตพฺพํ, ยา นหาเยยฺย, อาปตฺติ ทุกฺกฏสฺสาติ.}

\prosewithcloser{เตน โข ปน สมเยน ภิกฺขุนิโย ปุริสติตฺเถ นหายนฺติ. มนุสฺสา อุชฺฌายนฺติ ขิยฺยนฺติ วิปาเจนฺติ. “ฯเปฯ เสยฺยถาปิ คิหินี กามโภคินิโย”ติ. ภควโต เอตมตฺถํ อาโรเจสุํ. น ภิกฺขเว ภิกฺขุนิยา ปุริสติตฺเถ นหายิตพฺพํ, ยา นหาเยยฺย, อาปตฺติ ทุกฺกฏสฺส. อนุชานามิ ภิกฺขเว มหิลาติตฺเถ นหายิตุนฺติ.}{ตติย\-ภาณวาโร นิฏฺฐิโต.}
\csromancenter{ภิกฺขุนิ\-กฺขนฺธโก ทสโม.}
\csromancenter{อิมสฺมิํ ขนฺธเก วตฺถู เอกสตํ.}
\csromancenter{\textbf{ตสฺสุทฺทานํ}}
\setlength{\csromangathapagewidth}{0pt}
\setlength{\csromangathawakwidth}{0pt}
\csromangathameasuregroup{ปพฺพชฺชํโคตมี ยาจิ, \\ กปิลวตฺถุ เวสาลิํ, \\ รโชกิณฺเณน โกฏฺฐเก, \\ ภพฺโพติ นยโต ยาจิ, \\ วสฺสสตํ ตทหุ จ, \\ ปวารณา ครุธมฺมา, \\ โอวโฏ จ อฏฺฐ ธมฺมา, \\ ครุธมฺมปฏิคฺคาโห, \\ วสฺสสหสฺสํ ปญฺเจว, \\ มญฺชิฏฺฐิกอุปมาหิ, \\ อาฬิํ พนฺเธยฺย ปาเอว, \\ อุปสมฺปาเทตุํ อยฺยา, \\ น กริสฺสนฺติ กิเมว, \\ โอวาทํ ปาติโมกฺขญฺจ, \\ น ชานนฺติ จ อาจิกฺขิ, \\ ปฏิคฺคเหตุํ ภิกฺขูหิ, \\ อาจิกฺขิ กมฺมํ ภิกฺขูหิ, \\ อาจิกฺขิตุํ ภณฺฑนญฺจ, \\ สาวตฺถิยา กทฺทโมท, \\ องฺคชาตญฺจ โอภาสํ, \\ อวนฺทิโย ทณฺฑกมฺมํ, \\ อาวรณญฺจ โอวาทํ, \\ พาลา วตฺถุวินิจฺฉยา, \\ ทุเว ติสฺโส น คณฺหนฺติ, \\ อารญฺญิโก นาโรเจนฺติ, \\ ทีฆํ วิลีวจมฺมญฺจ, \\ โจฬเวณิ จ ปฏฺฏิ จ, \\ อฏฺฐิลฺลํ โคหนุเกน, \\ อูรุํ มุขํ ทนฺตมํสํ, \\ ลญฺเฉนฺติ องฺคราคญฺจ, \\ อวงฺคํ วิเสโสโลโก, \\ เวสี ปานาคารํ สูนํ, \\ ทาสํ ทาสิํ กมฺมกรํ, \\ ติรจฺฉานหรีตกิ, \\ นีลํ ปีตํ โลหิตกํ, \\ มหารงฺคมหานาม, \\ ปุปฺผผลกญฺจุกญฺจ, \\ ภิกฺขุนี สิกฺขมานาย, \\ นิยฺยาทิเต ปริกฺขาเร, \\ ภิกฺขุสฺส สามเณรสฺส, \\ อญฺเญสญฺจ ปริกฺขาเร, \\ มลฺลี คพฺภํ ปตฺตมูลํ, \\ อุสฺสนฺนญฺจ พาฬฺหตรํ, \\ ภิกฺขูนํ ยาทิสํ โภฏฺฐํ\textsuperscript{1}, \\ เสนาสนํ อุตุนิโย, \\ ฉิชฺชนฺติ สพฺพกาลญฺจ, \\ นิมิตฺตา โลหิตา เจว, \\ ธุวโจฬปคฺฆรนฺตี, \\ เวปุริสี จ สมฺภินฺนา, \\ อนิมิตฺตาทิโต กตฺวา,}{\csromangathabat{ปพฺพชฺชํโคตมี ยาจิ,}{นานุญฺญาสิ ตถาคโต.} \\ \csromangathabat{กปิลวตฺถุ เวสาลิํ,}{อคมาสิ วินายโก.} \\ \csromangathabat{รโชกิณฺเณน โกฏฺฐเก,}{อานนฺทสฺส ปเวทยิ.} \\ \csromangathabat{ภพฺโพติ นยโต ยาจิ,}{มาตาติ โปสิกาติ จ.} \\ \csromangathabat{วสฺสสตํ ตทหุ จ,}{อภิกฺขุปจฺจาสีสนา.} \\ \csromangathabat{ปวารณา ครุธมฺมา,}{เทฺว วสฺสา อนกฺโกสนา.} \\ \csromangathabat{โอวโฏ จ อฏฺฐ ธมฺมา,}{ยาวชีวานุวตฺตนา.} \\ \csromangathabat{ครุธมฺมปฏิคฺคาโห,}{สาวสฺสา อุปสมฺปทา.} \\ \csromangathabat{วสฺสสหสฺสํ ปญฺเจว,}{กุมฺภเถนกเสตฏฺฏิ.} \\ \csromangathabat{มญฺชิฏฺฐิกอุปมาหิ,}{เอวํ สทฺธมฺมหิํสนา.} \\ \csromangathabat{อาฬิํ พนฺเธยฺย ปาเอว,}{ปุน สทฺธมฺมสณฺฐิติ.} \\ \csromangathabat{อุปสมฺปาเทตุํ อยฺยา,}{ยถาวุฑฺฒาภิวาทนา.} \\ \csromangathabat{น กริสฺสนฺติ กิเมว,}{สาธารณาสาธารณํ.} \\ \csromangathabat{โอวาทํ ปาติโมกฺขญฺจ,}{เกน นุ โข อุปสฺสยํ.} \\ \csromangathabat{น ชานนฺติ จ อาจิกฺขิ,}{น กโรนฺติ จ ภิกฺขุหิ.} \\ \csromangathabat{ปฏิคฺคเหตุํ ภิกฺขูหิ,}{ภิกฺขุนีหิ ปฏิคฺคโห.} \\ \csromangathabat{อาจิกฺขิ กมฺมํ ภิกฺขูหิ,}{อุชฺฌายนฺติ ภิกฺขุนีหิ วา.} \\ \csromangathabat{อาจิกฺขิตุํ ภณฺฑนญฺจ,}{โรเปตฺวา อุปฺปลาย จ.} \\ \csromangathabat{สาวตฺถิยา กทฺทโมท,}{อวนฺทิ กาย อูรุ จ.} \\ \csromangathabat{องฺคชาตญฺจ โอภาสํ,}{สมฺปโยเชนฺติ วคฺคิกา.} \\ \csromangathabat{อวนฺทิโย ทณฺฑกมฺมํ,}{ภิกฺขุนิโย ตถา ปุน.} \\ \csromangathabat{อาวรณญฺจ โอวาทํ,}{กปฺปติ นุ โข ปกฺกมิ.} \\ \csromangathabat{พาลา วตฺถุวินิจฺฉยา,}{โอวาทํ สํโฆ ปญฺจหิ.} \\ \csromangathabat{ทุเว ติสฺโส น คณฺหนฺติ,}{พาลา คิลานคมิกํ.} \\ \csromangathabat{อารญฺญิโก นาโรเจนฺติ,}{น ปจฺจาคจฺฉนฺติ จ.} \\ \csromangathabat{ทีฆํ วิลีวจมฺมญฺจ,}{ทุสฺสา จ เวณิวฏฺฏิ จ.} \\ \csromangathabat{โจฬเวณิ จ ปฏฺฏิ จ,}{สุตฺตเวณิ จ วฏฺฏิกา.} \\ \csromangathabat{อฏฺฐิลฺลํ โคหนุเกน,}{หตฺถโกจฺฉํ ปาทํ ตถา.} \\ \csromangathabat{อูรุํ มุขํ ทนฺตมํสํ,}{อาลิมฺโปมทฺทจุณฺณนา.} \\ \csromangathabat{ลญฺเฉนฺติ องฺคราคญฺจ,}{มุขราคํ ตถา ทุเว.} \\ \csromangathabat{อวงฺคํ วิเสโสโลโก,}{สาโลเกน นจฺเจน จ\textsuperscript{1}.} \\ \csromangathabat{เวสี ปานาคารํ สูนํ,}{อาปณํ วฑฺฒิ วณิชฺชา.} \\ \csromangathabat{ทาสํ ทาสิํ กมฺมกรํ,}{กมฺมการิํ อุปฏฺฐยฺยุํ.} \\ \csromangathabat{ติรจฺฉานหรีตกิ,}{สํธารยนฺติ นมตกํ.} \\ \csromangathabat{นีลํ ปีตํ โลหิตกํ,}{มญฺชิฏฺฐกณฺหจีวรา.} \\ \csromangathabat{มหารงฺคมหานาม,}{อจฺฉินฺนา ทีฆเมว จ.} \\ \csromangathabat{ปุปฺผผลกญฺจุกญฺจ,}{ติรีฏกญฺจ ธารยุํ.} \\ \csromangathabat{ภิกฺขุนี สิกฺขมานาย,}{สามเณราย อจฺจเย.} \\ \csromangathabat{นิยฺยาทิเต ปริกฺขาเร,}{ภิกฺขุนิโยว อิสฺสรา.} \\ \csromangathabat{ภิกฺขุสฺส สามเณรสฺส,}{อุปาสกสฺสุปาสิกา.} \\ \csromangathabat{อญฺเญสญฺจ ปริกฺขาเร,}{นิยฺยาเต ภิกฺขุอิสฺสรา.} \\ \csromangathabat{มลฺลี คพฺภํ ปตฺตมูลํ,}{พฺยญฺชนํ อามิเสน จ.} \\ \csromangathabat{อุสฺสนฺนญฺจ พาฬฺหตรํ,}{สนฺนิธิกตมามิสํ.} \\ \csromangathabat{ภิกฺขูนํ ยาทิสํ โภฏฺฐํ\textsuperscript{1},}{ภิกฺขุนีนํ ตถา กเร.} \\ \csromangathabat{เสนาสนํ อุตุนิโย,}{มกฺขียติ ปฏาณิ จ.} \\ \csromangathabat{ฉิชฺชนฺติ สพฺพกาลญฺจ,}{อนิมิตฺตาปิ ทิสฺสเร.} \\ \csromangathabat{นิมิตฺตา โลหิตา เจว,}{ตเถว ธุวโลหิตา.} \\ \csromangathabat{ธุวโจฬปคฺฆรนฺตี,}{สิขรณิตฺถิปณฺฑกา.} \\ \csromangathabat{เวปุริสี จ สมฺภินฺนา,}{อุภโตพฺยญฺชนาปิ จ.} \\ \csromangathabat{อนิมิตฺตาทิโต กตฺวา,}{ยาว อุภโตพฺยญฺชนา.} \\ เอตํ เปยฺยาลโต เหฏฺฐา, กุฏฺฐํ คณฺโฑ กิลาโส จ \\ โสสาปมาโร มานุสี, \\ อิตฺถีสิ ภุชิสฺสาสิ จ. \\ อณณา น ราชภฏี, \\ อนุญฺญาตา จ วีสติ. \\ ปริปุณฺณา จ กินฺนามา, \\ กานามา เต ปวตฺตินี. \\ จตุวีสนฺตรายานํ, \\ ปุจฺฉิตฺวา อุปสมฺปทา. \\ วิตฺถายนฺติ อนนุสิฏฺฐา, \\ สํฆมชฺเฌ ตเถว จ. \\ อุปชฺฌาคาห สํฆาฏิ, \\ อุตฺตรนฺตรวาสโก. \\ สํกจฺจุทกสาฏิ จ, \\ อาจิกฺขิตฺวาน เปสเย. \\ พาลา อสมฺมเตกโต, \\ ยาเจ ปุจฺฉนฺตรายิกา. \\ เอกโตอุปสมฺปนฺนา, \\ ภิกฺขุสํเฆ ตถา ปุน. \\ ฉายา อุตุ ทิวสา จ, \\ สงฺคีติ ตโย นิสฺสเย. \\ อฏฺฐ อกรณียานิ, \\ กาลํ สพฺพตฺถ อฏฺเฐว. \\ น ปวาเรนฺติ ภิกฺขุนี, \\ ภิกฺขุสํฆํ ตเถว จ. \\ โกลาหลํ ปุเรภตฺตํ, \\ วิกาเล จ โกลาหลํ. \\ อุโปสถํ ปวารณํ, \\ สวจนียานุวาทนํ. \\ โอกาสํ โจเท สาเรนฺติ, \\ ปฏิกฺขิตฺตํ มเหสินา. \\ ตเถว ภิกฺขุ ภิกฺขุนี, \\ อนุญฺญาตํ มเหสินา. \\ ยานํ คิลานยุตฺตญฺจ, \\ ยานุคฺฆาตฑฺฒกาสิกา. \\ ภิกฺขุ สิกฺขา สามเณร, \\ สามเณรี จ พาลาย. \\ อรญฺเญ อุปาสเกน, \\ อุโทสิโต อุปสฺสยํ. \\ น สมฺมติ นวกมฺมํ, \\ นิสินฺนคพฺภเอกิกา. \\ สาคารญฺจ ครุธมฺมํ, \\ ปจฺจกฺขาย จ สงฺกมิ. \\ อภิวาทนเกสา จ, \\ นขา จ วณกมฺมนา. \\ ปลฺลงฺเกน คิลานา จ, \\ วจฺจํ จุณฺเณน วาสิตํ. \\ ชนฺตาฆเร ปฏิโสเต, \\ อติตฺเถ ปุริเสน จ. \\ มหาโคตมี อายาจิ, \\ อานนฺโท จาปิ โยนิโส. \\ ปริสา จตสฺโส โหนฺติ, \\ ปพฺพชฺชา ชินสาสเน. \\ สํเวคชนนตฺถาย, \\ สทฺธมฺมสฺส จ วุทฺธิยา. \\ อาตุรสฺสาว เภสชฺชํ, \\ เอวํ พุทฺเธน เทสิตํ. \\ เอวํ วินีตา สทฺธมฺเม, \\ มาตุคามาปิ อิตรา. \\ ยายนฺติ\textsuperscript{1} อจฺจุตํ ฐานํ, \\ ยตฺถ คนฺตฺวา น โสจเรติ.}
\csromangathameasurewak{เอตํ เปยฺยาลโต เหฏฺฐา, กุฏฺฐํ คณฺโฑ กิลาโส จ \\ โสสาปมาโร มานุสี, \\ อิตฺถีสิ ภุชิสฺสาสิ จ. \\ อณณา น ราชภฏี, \\ อนุญฺญาตา จ วีสติ. \\ ปริปุณฺณา จ กินฺนามา, \\ กานามา เต ปวตฺตินี. \\ จตุวีสนฺตรายานํ, \\ ปุจฺฉิตฺวา อุปสมฺปทา. \\ วิตฺถายนฺติ อนนุสิฏฺฐา, \\ สํฆมชฺเฌ ตเถว จ. \\ อุปชฺฌาคาห สํฆาฏิ, \\ อุตฺตรนฺตรวาสโก. \\ สํกจฺจุทกสาฏิ จ, \\ อาจิกฺขิตฺวาน เปสเย. \\ พาลา อสมฺมเตกโต, \\ ยาเจ ปุจฺฉนฺตรายิกา. \\ เอกโตอุปสมฺปนฺนา, \\ ภิกฺขุสํเฆ ตถา ปุน. \\ ฉายา อุตุ ทิวสา จ, \\ สงฺคีติ ตโย นิสฺสเย. \\ อฏฺฐ อกรณียานิ, \\ กาลํ สพฺพตฺถ อฏฺเฐว. \\ น ปวาเรนฺติ ภิกฺขุนี, \\ ภิกฺขุสํฆํ ตเถว จ. \\ โกลาหลํ ปุเรภตฺตํ, \\ วิกาเล จ โกลาหลํ. \\ อุโปสถํ ปวารณํ, \\ สวจนียานุวาทนํ. \\ โอกาสํ โจเท สาเรนฺติ, \\ ปฏิกฺขิตฺตํ มเหสินา. \\ ตเถว ภิกฺขุ ภิกฺขุนี, \\ อนุญฺญาตํ มเหสินา. \\ ยานํ คิลานยุตฺตญฺจ, \\ ยานุคฺฆาตฑฺฒกาสิกา. \\ ภิกฺขุ สิกฺขา สามเณร, \\ สามเณรี จ พาลาย. \\ อรญฺเญ อุปาสเกน, \\ อุโทสิโต อุปสฺสยํ. \\ น สมฺมติ นวกมฺมํ, \\ นิสินฺนคพฺภเอกิกา. \\ สาคารญฺจ ครุธมฺมํ, \\ ปจฺจกฺขาย จ สงฺกมิ. \\ อภิวาทนเกสา จ, \\ นขา จ วณกมฺมนา. \\ ปลฺลงฺเกน คิลานา จ, \\ วจฺจํ จุณฺเณน วาสิตํ. \\ ชนฺตาฆเร ปฏิโสเต, \\ อติตฺเถ ปุริเสน จ. \\ มหาโคตมี อายาจิ, \\ อานนฺโท จาปิ โยนิโส. \\ ปริสา จตสฺโส โหนฺติ, \\ ปพฺพชฺชา ชินสาสเน. \\ สํเวคชนนตฺถาย, \\ สทฺธมฺมสฺส จ วุทฺธิยา. \\ อาตุรสฺสาว เภสชฺชํ, \\ เอวํ พุทฺเธน เทสิตํ. \\ เอวํ วินีตา สทฺธมฺเม, \\ มาตุคามาปิ อิตรา. \\ ยายนฺติ\textsuperscript{1} อจฺจุตํ ฐานํ,}
\csromangathasetleft{ปพฺพชฺชํโคตมี ยาจิ, \\ กปิลวตฺถุ เวสาลิํ, \\ รโชกิณฺเณน โกฏฺฐเก, \\ ภพฺโพติ นยโต ยาจิ, \\ วสฺสสตํ ตทหุ จ, \\ ปวารณา ครุธมฺมา, \\ โอวโฏ จ อฏฺฐ ธมฺมา, \\ ครุธมฺมปฏิคฺคาโห, \\ วสฺสสหสฺสํ ปญฺเจว, \\ มญฺชิฏฺฐิกอุปมาหิ, \\ อาฬิํ พนฺเธยฺย ปาเอว, \\ อุปสมฺปาเทตุํ อยฺยา, \\ น กริสฺสนฺติ กิเมว, \\ โอวาทํ ปาติโมกฺขญฺจ, \\ น ชานนฺติ จ อาจิกฺขิ, \\ ปฏิคฺคเหตุํ ภิกฺขูหิ, \\ อาจิกฺขิ กมฺมํ ภิกฺขูหิ, \\ อาจิกฺขิตุํ ภณฺฑนญฺจ, \\ สาวตฺถิยา กทฺทโมท, \\ องฺคชาตญฺจ โอภาสํ, \\ อวนฺทิโย ทณฺฑกมฺมํ, \\ อาวรณญฺจ โอวาทํ, \\ พาลา วตฺถุวินิจฺฉยา, \\ ทุเว ติสฺโส น คณฺหนฺติ, \\ อารญฺญิโก นาโรเจนฺติ, \\ ทีฆํ วิลีวจมฺมญฺจ, \\ โจฬเวณิ จ ปฏฺฏิ จ, \\ อฏฺฐิลฺลํ โคหนุเกน, \\ อูรุํ มุขํ ทนฺตมํสํ, \\ ลญฺเฉนฺติ องฺคราคญฺจ, \\ อวงฺคํ วิเสโสโลโก, \\ เวสี ปานาคารํ สูนํ, \\ ทาสํ ทาสิํ กมฺมกรํ, \\ ติรจฺฉานหรีตกิ, \\ นีลํ ปีตํ โลหิตกํ, \\ มหารงฺคมหานาม, \\ ปุปฺผผลกญฺจุกญฺจ, \\ ภิกฺขุนี สิกฺขมานาย, \\ นิยฺยาทิเต ปริกฺขาเร, \\ ภิกฺขุสฺส สามเณรสฺส, \\ อญฺเญสญฺจ ปริกฺขาเร, \\ มลฺลี คพฺภํ ปตฺตมูลํ, \\ อุสฺสนฺนญฺจ พาฬฺหตรํ, \\ ภิกฺขูนํ ยาทิสํ โภฏฺฐํ\textsuperscript{1}, \\ เสนาสนํ อุตุนิโย, \\ ฉิชฺชนฺติ สพฺพกาลญฺจ, \\ นิมิตฺตา โลหิตา เจว, \\ ธุวโจฬปคฺฆรนฺตี, \\ เวปุริสี จ สมฺภินฺนา, \\ อนิมิตฺตาทิโต กตฺวา,}
\csromangathagroupwithclosermajor{\csromangathastanza{\csromangathabat{ปพฺพชฺชํโคตมี ยาจิ,}{นานุญฺญาสิ ตถาคโต.} \\ \csromangathabat{กปิลวตฺถุ เวสาลิํ,}{อคมาสิ วินายโก.}}\csromangathastanza{\csromangathabat{รโชกิณฺเณน โกฏฺฐเก,}{อานนฺทสฺส ปเวทยิ.} \\ \csromangathabat{ภพฺโพติ นยโต ยาจิ,}{มาตาติ โปสิกาติ จ.}}\csromangathastanza{\csromangathabat{วสฺสสตํ ตทหุ จ,}{อภิกฺขุปจฺจาสีสนา.} \\ \csromangathabat{ปวารณา ครุธมฺมา,}{เทฺว วสฺสา อนกฺโกสนา.}}\csromangathastanza{\csromangathabat{โอวโฏ จ อฏฺฐ ธมฺมา,}{ยาวชีวานุวตฺตนา.} \\ \csromangathabat{ครุธมฺมปฏิคฺคาโห,}{สาวสฺสา อุปสมฺปทา.}}\csromangathastanza{\csromangathabat{วสฺสสหสฺสํ ปญฺเจว,}{กุมฺภเถนกเสตฏฺฏิ.} \\ \csromangathabat{มญฺชิฏฺฐิกอุปมาหิ,}{เอวํ สทฺธมฺมหิํสนา.}}\csromangathastanza{\csromangathabat{อาฬิํ พนฺเธยฺย ปาเอว,}{ปุน สทฺธมฺมสณฺฐิติ.} \\ \csromangathabat{อุปสมฺปาเทตุํ อยฺยา,}{ยถาวุฑฺฒาภิวาทนา.}}\csromangathastanza{\csromanfolio{477}\csromangathabat{น กริสฺสนฺติ กิเมว,}{สาธารณาสาธารณํ.} \\ \csromangathabat{โอวาทํ ปาติโมกฺขญฺจ,}{เกน นุ โข อุปสฺสยํ.}}\csromangathastanza{\csromangathabat{น ชานนฺติ จ อาจิกฺขิ,}{น กโรนฺติ จ ภิกฺขุหิ.} \\ \csromangathabat{ปฏิคฺคเหตุํ ภิกฺขูหิ,}{ภิกฺขุนีหิ ปฏิคฺคโห.}}\csromangathastanza{\csromangathabat{อาจิกฺขิ กมฺมํ ภิกฺขูหิ,}{อุชฺฌายนฺติ ภิกฺขุนีหิ วา.} \\ \csromangathabat{อาจิกฺขิตุํ ภณฺฑนญฺจ,}{โรเปตฺวา อุปฺปลาย จ.}}\csromangathastanza{\csromangathabat{สาวตฺถิยา กทฺทโมท,}{อวนฺทิ กาย อูรุ จ.} \\ \csromangathabat{องฺคชาตญฺจ โอภาสํ,}{สมฺปโยเชนฺติ วคฺคิกา.}}\csromangathastanza{\csromangathabat{อวนฺทิโย ทณฺฑกมฺมํ,}{ภิกฺขุนิโย ตถา ปุน.} \\ \csromangathabat{อาวรณญฺจ โอวาทํ,}{กปฺปติ นุ โข ปกฺกมิ.}}\csromangathastanza{\csromangathabat{พาลา วตฺถุวินิจฺฉยา,}{โอวาทํ สํโฆ ปญฺจหิ.} \\ \csromangathabat{ทุเว ติสฺโส น คณฺหนฺติ,}{พาลา คิลานคมิกํ.}}\csromangathastanza{\csromangathabat{อารญฺญิโก นาโรเจนฺติ,}{น ปจฺจาคจฺฉนฺติ จ.} \\ \csromangathabat{ทีฆํ วิลีวจมฺมญฺจ,}{ทุสฺสา จ เวณิวฏฺฏิ จ.}}\csromangathastanza{\csromangathabat{โจฬเวณิ จ ปฏฺฏิ จ,}{สุตฺตเวณิ จ วฏฺฏิกา.} \\ \csromangathabat{อฏฺฐิลฺลํ โคหนุเกน,}{หตฺถโกจฺฉํ ปาทํ ตถา.}}\csromangathastanza{\csromangathabat{อูรุํ มุขํ ทนฺตมํสํ,}{อาลิมฺโปมทฺทจุณฺณนา.} \\ \csromangathabat{ลญฺเฉนฺติ องฺคราคญฺจ,}{มุขราคํ ตถา ทุเว.}}\csromangathastanza{\csromangathabat{อวงฺคํ วิเสโสโลโก,}{สาโลเกน นจฺเจน จ\csromansharedfootnote{d190eacc79a1fe18}{สาโลเก สนจฺเจน จ (สี), สาโลเกน สนจฺจนํ (สฺยา)}.} \\ \csromangathabat{เวสี ปานาคารํ สูนํ,}{อาปณํ วฑฺฒิ วณิชฺชา.}}\csromangathastanza{\csromangathabat{ทาสํ ทาสิํ กมฺมกรํ,}{กมฺมการิํ อุปฏฺฐยฺยุํ.} \\ \csromangathabat{ติรจฺฉานหรีตกิ,}{สํธารยนฺติ นมตกํ.}}\csromangathastanza{\csromangathabat{นีลํ ปีตํ โลหิตกํ,}{มญฺชิฏฺฐกณฺหจีวรา.} \\ \csromangathabat{มหารงฺคมหานาม,}{อจฺฉินฺนา ทีฆเมว จ.}}\csromangathastanza{\csromangathabat{ปุปฺผผลกญฺจุกญฺจ,}{ติรีฏกญฺจ ธารยุํ.} \\ \csromangathabat{ภิกฺขุนี สิกฺขมานาย,}{สามเณราย อจฺจเย.}}\csromangathastanza{\csromanfolio{478}\csromangathabat{นิยฺยาทิเต ปริกฺขาเร,}{ภิกฺขุนิโยว อิสฺสรา.} \\ \csromangathabat{ภิกฺขุสฺส สามเณรสฺส,}{อุปาสกสฺสุปาสิกา.}}\csromangathastanza{\csromangathabat{อญฺเญสญฺจ ปริกฺขาเร,}{นิยฺยาเต ภิกฺขุอิสฺสรา.} \\ \csromangathabat{มลฺลี คพฺภํ ปตฺตมูลํ,}{พฺยญฺชนํ อามิเสน จ.}}\csromangathastanza{\csromangathabat{อุสฺสนฺนญฺจ พาฬฺหตรํ,}{สนฺนิธิกตมามิสํ.} \\ \csromangathabat{ภิกฺขูนํ ยาทิสํ โภฏฺฐํ\csromansharedfootnote{2910d9e2288ef36f}{เหฏฺฐา (สี), เหฏฺฐํ (สฺยา), ภุตฺติ (ก)},}{ภิกฺขุนีนํ ตถา กเร.}}\csromangathastanza{\csromangathabat{เสนาสนํ อุตุนิโย,}{มกฺขียติ ปฏาณิ จ.} \\ \csromangathabat{ฉิชฺชนฺติ สพฺพกาลญฺจ,}{อนิมิตฺตาปิ ทิสฺสเร.}}\csromangathastanza{\csromangathabat{นิมิตฺตา โลหิตา เจว,}{ตเถว ธุวโลหิตา.} \\ \csromangathabat{ธุวโจฬปคฺฆรนฺตี,}{สิขรณิตฺถิปณฺฑกา.}}\csromangathastanza{\csromangathabat{เวปุริสี จ สมฺภินฺนา,}{อุภโตพฺยญฺชนาปิ จ.} \\ \csromangathabat{อนิมิตฺตาทิโต กตฺวา,}{ยาว อุภโตพฺยญฺชนา.}}\csromangathastanza{\csromangathawak{เอตํ เปยฺยาลโต เหฏฺฐา, กุฏฺฐํ คณฺโฑ กิลาโส จ \\ โสสาปมาโร มานุสี, \\ อิตฺถีสิ ภุชิสฺสาสิ จ. \\ อณณา น ราชภฏี,}}\csromangathastanza{\csromangathawak{อนุญฺญาตา จ วีสติ. \\ ปริปุณฺณา จ กินฺนามา, \\ กานามา เต ปวตฺตินี. \\ จตุวีสนฺตรายานํ,}}\csromangathastanza{\csromangathawak{ปุจฺฉิตฺวา อุปสมฺปทา. \\ วิตฺถายนฺติ อนนุสิฏฺฐา, \\ สํฆมชฺเฌ ตเถว จ. \\ อุปชฺฌาคาห สํฆาฏิ,}}\csromangathastanza{\csromangathawak{อุตฺตรนฺตรวาสโก. \\ สํกจฺจุทกสาฏิ จ, \\ อาจิกฺขิตฺวาน เปสเย. \\ พาลา อสมฺมเตกโต,}}\csromangathastanza{\csromangathawak{ยาเจ ปุจฺฉนฺตรายิกา. \\ เอกโต\-อุปสมฺปนฺนา, \\ ภิกฺขุสํเฆ ตถา ปุน. \\ ฉายา อุตุ ทิวสา จ,}}\csromangathastanza{\csromangathawak{สงฺคีติ ตโย นิสฺสเย. \\ อฏฺฐ อกรณียานิ, \\ กาลํ สพฺพตฺถ อฏฺเฐว. \\ น ปวาเรนฺติ ภิกฺขุนี,}}\csromangathastanza{\csromangathawak{ภิกฺขุสํฆํ ตเถว จ. \\ โกลาหลํ ปุเรภตฺตํ, \\ วิกาเล จ โกลาหลํ. \\ อุโปสถํ ปวารณํ,}}\csromangathastanza{\csromangathawak{\csromanfolio{479}สวจนียานุวาทนํ. \\ โอกาสํ โจเท สาเรนฺติ, \\ ปฏิกฺขิตฺตํ มเหสินา. \\ ตเถว ภิกฺขุ ภิกฺขุนี,}}\csromangathastanza{\csromangathawak{อนุญฺญาตํ มเหสินา. \\ ยานํ คิลานยุตฺตญฺจ, \\ ยานุคฺฆาตฑฺฒกาสิกา. \\ ภิกฺขุ สิกฺขา สามเณร,}}\csromangathastanza{\csromangathawak{สามเณรี จ พาลาย. \\ อรญฺเญ อุปาสเกน, \\ อุโทสิโต อุปสฺสยํ. \\ น สมฺมติ นวกมฺมํ,}}\csromangathastanza{\csromangathawak{นิสินฺนคพฺภเอกิกา. \\ สาคารญฺจ ครุธมฺมํ, \\ ปจฺจกฺขาย จ สงฺกมิ. \\ อภิวาทนเกสา จ,}}\csromangathastanza{\csromangathawak{นขา จ วณกมฺมนา. \\ ปลฺลงฺเกน คิลานา จ, \\ วจฺจํ จุณฺเณน วาสิตํ. \\ ชนฺตาฆเร ปฏิโสเต,}}\csromangathastanza{\csromangathawak{อติตฺเถ ปุริเสน จ. \\ มหาโคตมี อายาจิ, \\ อานนฺโท จาปิ โยนิโส. \\ ปริสา จตสฺโส โหนฺติ,}}\csromangathastanza{\csromangathawak{ปพฺพชฺชา ชินสาสเน. \\ สํเวคชนนตฺถาย, \\ สทฺธมฺมสฺส จ วุทฺธิยา. \\ อาตุรสฺสาว เภสชฺชํ,}}\csromangathastanza{\csromangathawak{เอวํ พุทฺเธน เทสิตํ. \\ เอวํ วินีตา สทฺธมฺเม, \\ มาตุคามาปิ อิตรา. \\ ยายนฺติ\csromansharedfootnote{9b54d565bcb82efa}{ตายนฺติ (สี, สฺยา)} อจฺจุตํ ฐานํ,}}\csromangathastanza{ยตฺถ คนฺตฺวา น โสจเรติ.}}{ภิกฺขุนิ\-กฺขนฺธโก นิฏฺฐิโต.}
\csromanmatikaanchor{mtk.212}
\csromantocmarknopage{chapter}{11. ปญฺจสติกกฺขนฺธก}{mtk.212}
\bookmark[dest={mtk.212},level=0]{11. ปญฺจสติกกฺขนฺธก}
\chapterheadpageread{11. ปญฺจสติก\-กฺขนฺธก}
\csromanmatikaanchor{mtk.213}
\csromantocmarkref{section}{1. สงฺคีตินิทาน}{mtk.213}
\bookmark[dest={mtk.213},level=1]{1. สงฺคีตินิทาน}
\csromanheader{1}{1. \textbf{สงฺคีตินิทาน}}
\proseitem{437}{\csromanfolio{480}อถ โข อายสฺมา มหากสฺสโป ภิกฺขู อามนฺเตสิ\csromandash{}เอกมิทาหํ อาวุโส สมยํ ปาวาย กุสินารํ อทฺธาน\-มคฺค\-ปฺปฏิปนฺโน มหตา ภิกฺขุ\-สํเฆน สทฺธิํ ปญฺจมตฺเตหิ ภิกฺขุสเตหิ. อถ ขฺวาหํ อาวุโส มคฺคา โอกฺกมฺม อญฺญตรสฺมิํ รุกฺขมูเล นิสีทิํ.}

\prose{\csromansymbolfootnote{*}{ที 2. 133 ปิฏฺเฐปิ.}เตน โข ปน สมเยน อญฺญตโร อาชีวโก กุสินาราย มนฺทารว\-ปุปฺผํ คเหตฺวา ปาวํ อทฺธาน\-มคฺค\-ปฺปฏิปนฺโน โหติ. อทฺทสํ โข อหํ อาวุโส ตํ อาชีวกํ ทูรโตว อาคจฺฉนฺตํ, ทิสฺวาน ตํ อาชีวกํ เอตทโวจํ “อปาวุโส อมฺหากํ สตฺถารํ ชานาสี”ติ. อามาวุโส ชานามิ, อชฺช สตฺตาห\-ปรินิพฺพุโต สมโณ โคตโม, ตโต เม อิทํ มนฺทารว\-ปุปฺผํ คหิตนฺติ. ตตฺราวุโส เย เต ภิกฺขู อวีตราคา, อปฺเปกจฺเจ พาหา ปคฺคยฺห กนฺทนฺติ, ฉินฺนปาตํ ปปตนฺติ, อาวฏฺฏนฺติ, วิวฏฺฏนฺติ, “อติขิปฺปํ ภควา ปรินิพฺพุโต, อติขิปฺปํ สุคโต ปรินิพฺพุโต, อติขิปฺปํ จกฺขุํโลเก อนฺตรหิตนฺ”ติ. เย ปน เต ภิกฺขู วีตราคา, เต สตา สมฺปชานา อธิวาเสนฺติ “อนิจฺจา สงฺขารา, ตํ กุเตตฺถ ลพฺภา”ติ.}

\prose{อถ ขฺวาหํ อาวุโส เต ภิกฺขู เอตทโวจํ “อลํ อาวุโส มา โสจิตฺถ, มา ปริเทวิตฺถ, นเนฺวตํ อาวุโส ภควตา ปฏิกจฺเจว อกฺขาตํ สพฺเพเหว ปิเยหิ มนาเปหิ นานาภาโว วินาภาโว อญฺญถาภาโว, ตํ กุเตตฺถ อาวุโส ลพฺภา, ‘ยํ ตํ ชาตํ ภูตํ สงฺขตํ ปโลกธมฺมํ, ตํ วต มา ปลุชฺชี’ติ เนตํ ฐานํ วิชฺชตี”ติ.}

\prose{เตน โข ปนาวุโส สมเยน สุภทฺโท นาม วุฑฺฒ\-ปพฺพชิโต ตสฺสํ ปริสายํ นิสินฺโน โหติ. อถ โข อาวุโส สุภทฺโท วุฑฺฒ\-ปพฺพชิโต เต ภิกฺขู เอตทโวจ “อลํ อาวุโส มา โสจิตฺถ, มา ปริเทวิตฺถ, สุมุตฺตา มยํ เตน มหาสมเณน, อุปทฺทุตา จ มยํ โหม ‘อิทํ โว กปฺปติ, อิทํ โว น กปฺปตี’ติ, อิทานิ ปน มยํ ยํ อิจฺฉิสฺสาม ตํ กริสฺสาม, ยํ น อิจฺฉิสฺสาม น ตํ กริสฺสามา”ติ. \csromanfolio{481}หนฺท มยํ อาวุโส ธมฺมญฺจ วินยญฺจ สงฺคายาม, ปุเร อธมฺโม ทิปฺปติ\csromansharedfootnote{b85ec6811e416363}{ทิพฺพติ (ก)}, ธมฺโม ปฏิพาหิยฺยติ, ปุเร อวินโย ทิปฺปติ, วินโย ปฏิพาหิยฺยติ, ปุเร อธมฺมวาทิโน พลวนฺโต โหนฺติ, ธมฺมวาทิโน ทุพฺพลา โหนฺติ, ปุเร อวินยวาทิโน พลวนฺโต โหนฺติ, วินยวาทิโน ทุพฺพลา โหนฺตีติ.}

\prose{เตน หิ ภนฺเต เถโร ภิกฺขู อุจฺจินตูติ. อถ โข อายสฺมา มหากสฺสโป เอเกนูน\-ปญฺจ\-อรหนฺต\-สตานิ อุจฺจินิ. ภิกฺขู อายสฺมนฺตํ มหากสฺสปํ เอตทโวจุํ “อยํ ภนฺเต อายสฺมา อานนฺโท กิญฺจาปิ เสกฺโข, อภพฺโพ ฉนฺทา โทสา โมหา ภยา อคติํ คนฺตุํ, พหุ จ อเนน ภควโต สนฺติเก ธมฺโม จ วินโย จ ปริยตฺโต, เตน หิ ภนฺเต เถโร อายสฺมนฺตมฺปิ อานนฺทํ อุจฺจินตู”ติ. อถ โข อายสฺมา มหากสฺสโป อายสฺมนฺตมฺปิ อานนฺทํ อุจฺจินิ.}

\prose{อถ โข เถรานํ ภิกฺขูนํ เอตทโหสิ “กตฺถ นุ โข มยํ ธมฺมญฺจ วินยญฺจ สงฺคาเยยฺยามา”ติ. อถ โข เถรานํ ภิกฺขูนํ เอตทโหสิ “ราชคหํ โข มหาโคจรํ ปหูต\-เสนาสนํ, ยํนูน มยํ ราชคเห วสฺสํ วสนฺตา ธมฺมญฺจ วินยญฺจ สงฺคาเยยฺยาม, น อญฺเญ ภิกฺขู ราชคเห วสฺสํ อุปคจฺเฉยฺยุนฺติ.}

\prose{อถ โข อยสฺมา มหากสฺสโป สํฆํ ญาเปสิ\csromandash{}}

\proseitem{438}{“สุณาตุ เม อาวุโส สํโฆ, ยทิ สํฆสฺส ปตฺตกลฺลํ, สํโฆ อิมานิ ปญฺจ ภิกฺขุสตานิ สมฺมนฺเนยฺย ‘ราชคเห วสฺสํ วสนฺตานิ ธมฺมญฺจ วินยญฺจ สงฺคายิตุํ, น อญฺเญหิ ภิกฺขูหิ ราชคเห วสฺสํ วสิตพฺพนฺ’ติ, เอสา ญตฺติ.}

\prose{สุณาตุ เม อาวุโส สํโฆ, อิมานิ ปญฺจ ภิกฺขุสตานิ สมฺมนฺนติ ‘ราชคเห วสฺสํ วสนฺตานิ ธมฺมญฺจ วินยญฺจ สงฺคายิตุํ, น อญฺเญหิ ภิกฺขูหิ ราชคเห วสฺสํ วสิตพฺพนฺ’ติ, ยสฺสายสฺมโต ขมติ อิเมสํ ปญฺจนฺนํ ภิกฺขุสตานํ สมฺมุติ ‘ราชคเห วสฺสํ วสนฺตานํ ธมฺมญฺจ วินยญฺจ สงฺคายิตุํ น อญฺเญหิ ภิกฺขูหิ ราชคเห วสฺสํ วสิตพฺพนฺ’ติ, โส ตุณฺหสฺส, ยสฺส นกฺขมติ, โส ภาเสยฺย.}

\prose{\csromanfolio{482}สมฺมตานิ สํเฆน อิมานิ ปญฺจ ภิกฺขุสตานิ ‘ราชคเห วสฺสํ วสนฺตานิ ธมฺมญฺจ วินยญฺจ สงฺคายิตุํ น อญฺเญหิ ภิกฺขูหิ ราชคเห วสฺสํ วสิตพฺพนฺ’ติ, ขมติ สํฆสฺส, ตสฺมา ตุณฺหี, เอวเมตํ ธารยามี”ติ.}

\prose{อถ โข เถรา ภิกฺขู ราชคหํ อคมํสุ ธมฺมญฺจ วินยญฺจ สงฺคายิตุํ. อถ โข เถรานํ ภิกฺขูนํ เอตทโหสิ “ภควตา โข อาวุโส ขณฺฑผุลฺล\-ปฺปฏิสงฺขรณํ วณฺณิตํ, หนฺท มยํ อาวุโส ปฐมํ มาสํ ขณฺฑผุลฺลํ ปฏิสงฺขโรม, มชฺฌิมํ มาสํ สนฺนิปติตฺวา ธมฺมญฺจ วินยญฺจ สงฺคายิสฺสามา”ติ.}

\prose{อถ โข เถรา ภิกฺขู ปฐมํ มาสํ ขณฺฑผุลฺลํ ปฏิสงฺขาริํสุ. อถ โข อายสฺมา อานนฺโท “เสฺว สนฺนิปาโต\csromansharedfootnote{df0d60d0c129c8e6}{สนฺนิปาโตติ (ก)}, น โข เมตํ ปติรูปํ, โยหํ เสกฺโข สมาโน สนฺนิปาตํ คจฺเฉยฺยนฺ”ติ พหุเทว รตฺติํ กายคตาย สติยา วีตินาเมตฺวา รตฺติยา ปจฺจูสสมยํ “นิปชฺชิสฺสามี”ติ กายํ อาวชฺเชสิ, อปฺปตฺตญฺจ สีสํ พิพฺโพหนํ, ภูมิโต จ ปาทา มุตฺตา, เอตสฺมิํ อนฺตเร อนุปาทาย อาสเวหิ จิตฺตํ วิมุจฺจิ.}

\proseitem{439}{อถ โข อายสฺมา อานนฺโท อรหา สมาโน สนฺนิปาตํ อคมาสิ. อถ โข อายสฺมา มหากสฺสโป สํฆํ ญาเปสิ\csromandash{}}

\prose{“สุณาตุ เม อาวุโส สํโฆ, ยทิ สํฆสฺส ปตฺตกลฺลํ, อหํ อุปาลิํ วินยํ ปุจฺเฉยฺยนฺ”ติ.}

\prose{อายสฺมา อุปาลิ สํฆํ ญาเปสิ\csromandash{}}

\prose{“สุณาตุ เม ภนฺเต สํโฆ, ยทิ สํฆสฺส ปตฺตกลฺลํ, อหํ อายสฺมตา มหากสฺสเปน วินยํ ปุฏฺโฐ วิสฺสชฺเชยฺยนฺ”ติ.}

\prose{อถ โข อายสฺมา มหากสฺสโป อายสฺมนฺตํ อุปาลิํ เอตทโวจ “ปฐมํ อาวุโส อุปาลิ ปาราชิกํ กตฺถ ปญฺญตฺตนฺ”ติ. เวสลิยํ ภนฺเตติ. กํ อารพฺภาติ. สุทินฺนํ กลนฺทปุตฺตํ อารพฺภาติ. กิสฺมิํ วตฺถุสฺมินฺติ. เมถุนธมฺเมติ. อถ โข อายสฺมา มหากสฺสโป อายสฺมนฺตํ อุปาลิํ ปฐมสฺส ปาราชิกสฺส วตฺถุมฺปิ ปุจฺฉิ, นิทานมฺปิ ปุจฺฉิ, ปุคฺคลมฺปิ ปุจฺฉิ, ปญฺญตฺติมฺปิ ปุจฺฉิ, อนุปญฺญตฺติมฺปิ ปุจฺฉิ, อาปตฺติมฺปิ ปุจฺฉิ, \csromanfolio{483}อนาปตฺติมฺปิ ปุจฺฉิ. ทุติยํ ปนาวุโส อุปาลิ ปาราชิกํ กตฺถ ปญฺญตฺตนฺติ. ราชคเห ภนฺเตติ. กํ อารพฺภาติ. ธนิยํ กุมฺภการ\-ปุตฺตํ อารพฺภาติ. กิสฺมิํ วตฺถุสฺมินฺติ. อทินฺนาทาเนติ. อถ โข อายสฺมา มหากสฺสโป อายสฺมนฺตํ อุปาลิํ ทุติยสฺส ปาราชิกสฺส วตฺถุมฺปิ ปุจฺฉิ, นิทานมฺปิ ปุจฺฉิ, ปุคฺคลมฺปิ ปุจฺฉิ, ปญฺญตฺติมฺปิ ปุจฺฉิ, อนุปญฺญตฺติมฺปิ ปุจฺฉิ, อาปตฺติมฺปิ ปุจฺฉิ, อนาปตฺติมฺปิ ปุจฺฉิ. ตติยํ ปนาวุโส อุปาลิ ปาราชิกํ กตฺถ ปญฺญตฺตนฺติ. เวสาลิยํ ภนฺเตติ. กํ อารพฺภาติ. สมฺพหุเล ภิกฺขู อารพฺภาติ. กิสฺมิํ วตฺถุสฺมินฺติ. มนุสฺส\-วิคฺคเหติ. อถ โข อายสฺมา มหากสฺสโป อายสฺมนฺตํ อุปาลิํ ตติยสฺส ปาราชิกสฺส วตฺถุมฺปิ ปุจฺฉิ, นิทานมฺปิ ปุจฺฉิ, ปุคฺคลมฺปิ ปุจฺฉิ, ปญฺญตฺติมฺปิ ปุจฺฉิ, อนุปญฺญตฺติมฺปิ ปุจฺฉิ, อาปตฺติมฺปิ ปุจฺฉิ, อนาปตฺติมฺปิ ปุจฺฉิ. จตุตฺถํ ปนาวุโส อุปาลิ ปาราชิกํ กตฺถ ปญฺญตฺตนฺติ. เวสาลิยํ ภนฺเตติ. กํ อารพฺภาติ. วคฺคุมุทา\-ตีริเย ภิกฺขู อารพฺภาติ. กิสฺมิํ วตฺถุสฺมินฺติ. อุตฺตริมนุสฺสธมฺเมติ. อถ โข อายสฺมา มหากสฺสโป อายสฺมนฺตํ อุปาลิํ จตุตฺถสฺส ปาราชิกสฺส วตฺถุมฺปิ ปุจฺฉิ, นิทานมฺปิ ปุจฺฉิ, ปุคฺคลมฺปิ ปุจฺฉิ, ปญฺญตฺติมฺปิ ปุจฺฉิ, อนุปญฺญตฺติมฺปิ ปุจฺฉิ, อาปตฺติมฺปิ ปุจฺฉิ, อนาปตฺติมฺปิ ปุจฺฉิ. เอเตเนว อุปาเยน อุภโตวิภงฺเค ปุจฺฉิ, ปุฏฺโฐ ปุฏฺโฐ อายสฺมา อุปาลิ วิสฺสชฺเชสิ.}

\proseitem{440}{อถ โข อายสฺมา มหากสฺสโป สํฆํ ญาเปสิ\csromandash{}}

\prose{“สุณาตุ เม อาวุโส สํโฆ, ยทิ สํฆสฺส ปตฺตกลฺลํ, อหํ อานนฺทํ ธมฺมํ ปุจฺเฉยฺยนฺ”ติ.}

\prose{อายสฺมา อานนฺโท สํฆํ ญาเปสิ\csromandash{}}

\prose{“สุณาตุ เม ภนฺเต สํโฆ, ยทิ สํฆสฺส ปตฺตกลฺลํ, อหํ อายสฺมตา มหากสฺสเปน ธมฺมํ ปุฏฺโฐ วิสฺสชฺเชยฺยนฺ”ติ.}

\prose{อถ โข อายสฺมา มหากสฺสโป อายสฺมนฺตํ อานนฺทํ เอตทโวจ “พฺรหฺมชาลํ อาวุโส อานนฺท กตฺถ ภาสิตนฺ”ติ. อนฺตรา จ ภนฺเต ราชคหํ, อนฺตรา จ นาฬนฺทํ ราชาคารเก อมฺพลฏฺฐิกายาติ. กํ อารพฺภาติ. สุปฺปิยญฺจ ปริพฺพาชกํ พฺรหฺมทตฺตญฺจ มาณวนฺติ. อถ โข อายสฺมา มหากสฺสโป อายสฺมนฺตํ อานนฺทํ พฺรหฺมชาลสฺส นิทานมฺปิ ปุจฺฉิ, ปุคฺคลมฺปิ ปุจฺฉิ. สามญฺญผลํ ปนาวุโส อานนฺท กตฺถ ภาสิตนฺติ. ราชคเห ภนฺเต ชีวกมฺพวเนติ. เกน สทฺธินฺติ. อชาตสตฺตุนา \csromanfolio{484}เวเทหิปุตฺเตน สทฺธินฺติ. อถ โข อายสฺมา มหากสฺสโป อายสฺมนฺตํ อานนฺทํ สามญฺญ\-ผลสฺส นิทานมฺปิ ปุจฺฉิ, ปุคฺคลมฺปิ ปุจฺฉิ. เอเตเนว อุปาเยน ปญฺจปิ นิกาเย ปุจฺฉิ, ปุฏฺโฐ ปุฏฺโฐ อายสฺมา อานนฺโท วิสฺสชฺเชสิ.}

\csromanmatikaanchor{mtk.214}
\csromantocmarkref{section}{2. ขุทฺทานุขุทฺทกสิกฺขาปทกถา}{mtk.214}
\bookmark[dest={mtk.214},level=1]{2. ขุทฺทานุขุทฺทกสิกฺขาปทกถา}
\csromanheader{1}{2. ขุทฺทานุขุทฺทก\-สิกฺขาปท\-กถา}
\proseitem{441}{อถ โข อายสฺมา อานนฺโท เถเร ภิกฺขู เอตทโวจ “ภควา มํ ภนฺเต ปรินิพฺพาน\-กาเล เอวมาห ‘อากงฺขมาโน อานนฺท สํโฆ มมจฺจเยน ขุทฺทา\-นุขุทฺท\-กานิ สิกฺขาปทานิ สมูหเนยฺยา’ติ”. ปุจฺฉิ ปุน ตฺวํ อาวุโส อานนฺท ภควนฺตํ “กตมานิ ปน ภนฺเต ขุทฺทา\-นุขุทฺท\-กานิ สิกฺขาปทานี”ติ. น โขหํ ภนฺเต ภควนฺตํ ปุจฺฉิํ “กตมานิ ปน ภนฺเต ขุทฺทา\-นุขุทฺท\-กานิ สิกฺขาปทานี”ติ. เอกจฺเจ เถรา เอวมาหํสุ “จตฺตาริ ปาราชิกานิ ฐเปตฺวา อวเสสานิ ขุทฺทา\-นุขุทฺท\-กานิ สิกฺขาปทานี”ติ. เอกจฺเจ เถรา เอวมาหํสุ “จตฺตาริ ปาราชิกานิ ฐเปตฺวา เตรส สํฆาทิเสเส ฐเปตฺวา อวเสสานิ ขุทฺทา\-นุขุทฺท\-กานิ สิกฺขาปทานี”ติ. เอกจฺเจ เถรา เอวมาหํสุ “จตฺตาริ ปาราชิกานิ ฐเปตฺวา เตรส สํฆาทิเสเส ฐเปตฺวา เทฺว อนิยเต ฐเปตฺวา อวเสสานิ ขุทฺทา\-นุขุทฺท\-กานิ สิกฺขาปทานี”ติ. เอกจฺเจ เถรา เอวมาหํสุ “จตฺตาริ ปาราชิกานิ ฐเปตฺวา เตรส สํฆาทิเสเส ฐเปตฺวา เทฺว อนิยเต ฐเปตฺวา ติํส นิสฺสคฺคิเย ปาจิตฺติเย ฐเปตฺวา อวเสสานิ ขุทฺทา\-นุขุทฺท\-กานิ สิกฺขาปทานี”ติ. เอกจฺเจ เถรา เอวมาหํสุ “จตฺตาริ ปาราชิกานิ ฐเปตฺวา เตรส สํฆาทิเสเส ฐเปตฺวา เทฺว อนิยเต ฐเปตฺวา ติํส นิสฺสคฺคิเย ปาจิตฺติเย ฐเปตฺวา เทฺวนวุติ ปาจิตฺติเย ฐเปตฺวา อวเสสานิ ขุทฺทา\-นุขุทฺท\-กานิ สิกฺขาปทานี”ติ. เอกจฺเจ เถรา เอวมาหํสุ “จตฺตาริ ปาราชิกานิ ฐเปตฺวา เตรส สํฆาทิเสเส ฐเปตฺวา เทฺว อนิยเต ฐเปตฺวา ติํส นิสฺสคฺคิเย ปาจิตฺติเย ฐเปตฺวา เทฺวนวุติ ปาจิตฺติเย ฐเปตฺวา จตฺตาโร ปาฏิเทสนีเย ฐเปตฺวา อวเสสานิ ขุทฺทา\-นุขุทฺท\-กานิ สิกฺขาปทานี”ติ.}

\proseitem{442}{อถ โข อายสฺมา มหากสฺสโป สํฆํ ญาเปสิ\csromandash{}}

\prose{“สุณาตุ เม อาวุโส สํโฆ, สนฺตมฺหากํ สิกฺขาปทานิ คิหิคตานิ คิหิโนปิ ชานนฺติ ‘อิทํ โว สมณานํ สกฺย\-ปุตฺติ\-ยานํ กปฺปติ, อิทํ โว น กปฺปตี’ติ. สเจ มยํ ขุทฺทา\-นุขุทฺท\-กานิ สิกฺขาปทานิ สมูหนิสฺสาม, \csromanfolio{485}ภวิสฺสนฺติ วตฺตาโร ‘ธูมกาลิกํ สมเณน โคตเมน สาวกานํ สิกฺขาปทํ ปญฺญตฺตํ, ยาวิเมสํ สตฺถา อฏฺฐาสิ, ตาวิเม สิกฺขาปเทสุ สิกฺขิํสุ, ยโต อิเมสํ สตฺถา ปรินิพฺพุโต, น ทานิเม สิกฺขาปเทสุ สิกฺขนฺตี’ติ, ยทิ สํฆสฺส ปตฺตกลฺลํ, สํโฆ อปฺปญฺญตฺตํ นปฺปญฺญเปยฺย, ปญฺญตฺตํ น สมุจฺฉินฺเทยฺย, ยถา\-ปญฺญตฺเตสุ สิกฺขาปเทสุ สมาทาย วตฺเตยฺย, เอสา ญตฺติ.}

\prose{สุณาตุ เม อาวุโส สํโฆ, สนฺตมฺหากํ สิกฺขาปทานิ คิหิคตานิ คิหิโนปิ ชานนฺติ ‘อิทํ โว สมณานํ สกฺย\-ปุตฺติ\-ยานํ กปฺปติ, อิทํ โว น กปฺปตี”ติ, สเจ มยํ ขุทฺทา\-นุขุทฺท\-กานิ สิกฺขาปทานิ สมูหนิสฺสาม, ภวิสฺสนฺติ วตฺตาโร ‘ธูมกาลิกํ สมเณน โคตเมน สาวกานํ สิกฺขาปทํ ปญฺญตฺตํ, ยาวิเมสํ สตฺถา อฏฺฐาสิ, ตาวิเม สิกฺขาปเทสุ สิกฺขิํสุ, ยโต อิเมสํ สตฺถา ปรินิพฺพุโต, น ทานิเม สิกฺขาปเทสุ สิกฺขนฺตี’ติ, สํโฆ อปฺปญฺญตฺตํ นปฺปญฺญเปติ, ปญฺญตฺตํ น สมุจฺฉินฺทติ, ยถา\-ปญฺญตฺเตสุ สิกฺขาปเทสุ สมาทาย วตฺตติ, ยสฺสายสฺมโต ขมติ อปฺปญฺญตฺตสฺส อปฺปญฺญาปนา, ปญฺญตฺตสฺส อสมุจฺเฉโท, ยถา\-ปญฺญตฺเตสุ สิกฺขาปเทสุ สมาทาย วตฺตนา, โส ตุณฺหสฺส, ยสฺส นกฺขมติ, โส ภาเสยฺย.}

\prose{สํโฆ อปฺปญฺญตฺตํ นปฺปญฺญเปติ, ปญฺญตฺตํ น สมุจฺฉินฺทติ, ยถา\-ปญฺญตฺเตสุ สิกฺขาปเทสุ สมาทาย วตฺตติ, ขมติ สํฆสฺส, ตสฺมา ตุณฺหี, เอวเมตํ ธารยามี”ติ.}

\proseitem{443}{อถ โข เถรา ภิกฺขู อายสฺมนฺตํ อานนฺทํ เอตทโวจุํ “อิทํ เต อาวุโส อานนฺท ทุกฺกฏํ, ยํ ตฺวํ ภควนฺตํ น ปุจฺฉิ ‘กตมานิ ปน ภนฺเต ขุทฺทา\-นุขุทฺท\-กานิ สิกฺขาปทานี’ติ, เทเสหิ ตํ ทุกฺกฏนฺ”ติ. อหํ โข ภนฺเต อสฺสติยา ภควนฺตํ น ปุจฺฉิํ “กตมานิ ปน ภนฺเต ขุทฺทา\-นุขุทฺท\-กานิ สิกฺขาปทานี”ติ, นาหํ ตํ ทุกฺกฏํ ปสฺสามิ, อปิ จายสฺมนฺตานํ สทฺธาย เทเสมิ ตํ ทุกฺกฏนฺติ. อิทมฺปิ เต อาวุโส อานนฺท ทุกฺกฏํ, ยํ ตฺวํ ภควโต วสฺสิกสาฏิกํ อกฺกมิตฺวา สิพฺเพสิ, เทเสหิ ตํ ทุกฺกฏนฺติ. อหํ โข ภนฺเต น อคารเวน ภควโต วสฺสิกสาฏิกํ อกฺกมิตฺวา สิพฺเพสิํ, นาหํ ตํ ทุกฺกฏํ ปสฺสามิ, อปิ จายสฺมนฺตานํ สทฺธาย เทเสมิ ตํ ทุกฺกฏนฺติ. อิทมฺปิ เต อาวุโส อานนฺท ทุกฺกฏํ, ยํ ตฺวํ มาตุคาเมหิ \csromanfolio{486}ภควโต สรีรํ ปฐมํ วนฺทาเปสิ, ตาสํ โรทนฺตีนํ ภควโต สรีรํ อสฺสุเกน มกฺขิตํ, เทเสหิ ตํ ทุกฺกฏนฺติ. อหํ โข ภนฺเต “มายิมาสํ\csromansharedfootnote{ce07ab0cedfc4bb0}{มายิมา (สี, สฺยา)} วิกาเล อเหสุนฺ”ติ มาตุคาเมหิ ภควโต สรีรํ ปฐมํ วนฺทาเปสิํ, นาหํ ตํ ทุกฺกฏํ ปสฺสามิ, อปิ จายสฺมนฺตานํ สทฺธาย เทเสมิ ตํ ทุกฺกฏนฺติ. อิทมฺปิ เต อาวุโส อานนฺท ทุกฺกฏํ, ยํ ตฺวํ ภควโต โอฬาริเก นิมิตฺเต กยิรมาเน โอฬริเก โอภาเส กยิรมาเน น ภควนฺตํ ยาจิ “ติฏฺฐตุ ภควา กปฺปํ, ติฏฺฐตุ สุคโต กปฺปํ พหุชนหิตาย พหุชน\-สุขาย โลกานุกมฺปาย อตฺถาย หิตาย สุขาย เทวมนุสฺสานนฺ”ติ, เทเสหิ ตํ ทุกฺกฏนฺติ. อหํ โข ภนฺเต มาเรน ปริยุฏฺฐิต\-จิตฺโต น ภควนฺตํ ยาจิํ “ติฏฺฐตุ ภควา กปฺปํ, ติฏฺฐตุ สุคโต กปฺปํ พหุชนหิตาย พหุชน\-สุขาย โลกานุกมฺปาย อตฺถาย หิตาย สุขาย เทวมนุสฺสานนฺ”ติ, นาหํ ตํ ทุกฺกฏํ ปสฺสามิ, อปิ จายสฺมนฺตานํ สทฺธาย เทเสมิ ตํ ทุกฺกฏนฺติ. อิทมฺปิ เต อาวุโส อานนฺท ทุกฺกฏํ, ยํ ตฺวํ มาตุคามสฺส ตถาคต\-ปฺปเวทิเต ธมฺมวินเย ปพฺพชฺชํ อุสฺสุกฺกํ อกาสิ, เทเสหิ ตํ ทุกฺกฏนฺติ. อหํ โข ภนฺเต “อยํ มหาปชาปติ โคตมี ภควโต มาตุจฺฉา อาปาทิกา โปสิกา ขีรสฺส ทายิกา ภควนฺตํ ชเนตฺติยา กาลงฺกตาย ถญฺญํ ปาเยสี”ติ มาตุคามสฺส ตถาคต\-ปฺปเวทิเต ธมฺมวินเย ปพฺพชฺชํ อุสฺสุกฺกํ อกาสิํ, นาหํ ตํ ทุกฺกฏํ ปสฺสามิ, อปิ จายสฺมนฺตานํ สทฺธาย เทเสมิ ตํ ทุกฺกฏนฺติ.}

\proseitem{444}{เตน โข ปน สมเยน อายสฺมา ปุราโณ ทกฺขิณา\-คิริสฺมิํ จาริกํ จรติ มหตา ภิกฺขุ\-สํเฆน สทฺธิํ ปญฺจมตฺเตหิ ภิกฺขุสเตหิ. อถ โข อายสฺมา ปุราโณ เถเรหิ ภิกฺขูหิ ธมฺเม จ วินเย จ สงฺคีเต ทกฺขิณา\-คิริสฺมิํ ยถาภิรนฺตํ วิหริตฺวา เยน ราชคหํ เยน เวฬุวนํ กลนฺทก\-นิวาโป เยน เถรา ภิกฺขู เตนุปสงฺกมิ, อุปสงฺกมิตฺวา เถเรหิ ภิกฺขูหิ สทฺธิํ ปฏิสมฺโมทิตฺวา เอกมนฺตํ นิสีทิ, เอกมนฺตํ นิสินฺนํ โข อายสฺมนฺตํ ปุราณํ เถรา ภิกฺขู เอตทโวจุํ “เถเรหิ อาวุโส ปุราณ ธมฺโม จ วินโย จ สงฺกีโต, อุเปหิ ตํ สงฺคีตินฺ”ติ. สุสงฺคีตาวุโส เถเรหิ ธมฺโม จ วินโย จ, อปิ จ ยเถว มยา ภควโต สมฺมุขา สุตํ สมฺมุขา ปฏิคฺคหิตํ, ตเถวาหํ ธาเรสฺสามีติ.}

\csromanmatikaanchor{mtk.215}
\csromantocmarkref{section}{3. พฺรหฺมทณฺฑกถา}{mtk.215}
\bookmark[dest={mtk.215},level=1]{3. พฺรหฺมทณฺฑกถา}
\csromanheader{1}{3. พฺรหฺมทณฺฑ\-กถา}
\proseitem{445}{\csromanfolio{487}อถ โข อายสฺมา อานนฺโท เถเร ภิกฺขู เอตทโวจ “ภควา มํ ภนฺเต ปรินิพฺพาน\-กาเล เอวมาห ‘เตน หานนฺท สํโฆ มมจฺจเยน ฉนฺนสฺส ภิกฺขุโน พฺรหฺมทณฺฑํ อาณาเปตู’ติ”. ปุจฺฉิ ปน ตฺวํ อาวุโส อานนฺท ภควนฺตํ “กตโม ปน ภนฺเต พฺรหฺมทณฺโฑ”ติ. ปุจฺฉิํ โขหํ ภนฺเต ภควนฺตํ “กตโม ปน ภนฺเต พฺรหฺมทณฺโฑ”ติ. ฉนฺโน อานนฺท ภิกฺขุ ยํ อิจฺเฉยฺย ตํ วเทยฺย, ภิกฺขูหิ ฉนฺโน ภิกฺขุ เนว วตฺตพฺโพ, น โอวทิตพฺโพ, นานุสาสิตพฺโพติ. เตน หาวุโส อานนฺท ตฺวํเยว ฉนฺนสฺส ภิกฺขุโน พฺรหฺมทณฺฑํ อาณาเปหีติ. กถาหํ ภนฺเต ฉนฺนสฺส ภิกฺขุโน พฺรหฺมทณฺฑํ อาณาเปมิ, จณฺโฑ โส ภิกฺขุ ผรุโสติ. เตน หาวุโส อานนฺท พหุเกหิ ภิกฺขูหิ สทฺธิํ คจฺฉาหีติ. “เอวํ ภนฺเต”ติ โข อายสฺมา อานนฺโท เถรานํ ภิกฺขูนํ ปฏิสฺสุตฺวา มหตา ภิกฺขุ\-สํเฆน สทฺธิํ ปญฺจมตฺเตหิ ภิกฺขุสเตหิ นาวาย อุชฺชวนิกาย โกสมฺพิํ อุชฺชวิ, นาวาย ปจฺโจโรหิตฺวา รญฺโญ อุเทนสฺส\csromansharedfootnote{3c8ba7b277526c78}{อุเตนสฺส (ก)} อุยฺยานสฺส อวิทูเร อญฺญตรสฺมิํ รุกฺขมูเล นิสีทิ. เตน โข ปน สมเยน ราชา อุเทโน อุยฺยาเน ปริจาเรสิ สทฺธิํ โอโรเธน. อสฺโสสิ โข รญฺโญ อุเทนสฺส โอโรโธ “อมฺหากํ กิร อาจริโย อยฺโย อานนฺโท อุยฺยานสฺส อวิทูเร อญฺญตรสฺมิํ รุกฺขมูเล นิสินฺโน”ติ. อถ โข รญฺโญ อุเทนสฺส โอโรโธ ราชานํ อุเทนํ เอตทโวจ “อมฺหากํ กิร เทว อาจริโย อยฺโย อานนฺเทว อุยฺยานสฺส อวิทูเร อญฺญตรสฺมิํ รุกฺขมูเล นิสินฺโน, อิจฺฉาม มยํ เทว อยฺยํ อานนฺทํ ปสฺสิตุนฺ”ติ. เตน หิ ตุเมฺห สมณํ อานนฺทํ ปสฺสถาติ.}

\prose{อถ โข รญฺโญ อุเทนสฺส โอโรโธ เยนายสฺมา อานนฺโท เตนุปสงฺกมิ, อุปสงฺกมิตฺวา อายสฺมนฺตํ อานนฺทํ อภิวาเทตฺวา เอกมนฺตํ นิสีทิ. เอกมนฺตํ นิสินฺนํ โข รญฺโญ อุเทนสฺส โอโรธํ อายสฺมา อานนฺโท ธมฺมิยา กถาย สนฺทสฺเสสิ สมาทเปสิ สมุตฺเตเชสิ สมฺปหํเสสิ. อถ โข รญฺโญ อุเทนสฺส โอโรโธ อายสฺมตา อานนฺเทน ธมฺมิยา กถาย สนฺทสฺสิโต สมาทปิโต สมุตฺเตชิโต สมฺปหํสิโต อายสฺมโต อานนฺทสฺส ปญฺจ อุตฺตราสงฺค\-สตานิ ปาทาสิ. \csromanfolio{488}อถ โข รญฺโญ อุเทนสฺส โอโรโธ อายสฺมโต อานนฺทสฺส ภาสิตํ อภินนฺทิตฺวา อนุโมทิตฺวา อุฏฺฐายาสนา อายสฺมนฺตํ อานนฺทํ อภิวาเทตฺวา ปทกฺขิณํ กตฺวา เยน ราชา อุเทโน เตนุปสงฺกมิ. อทฺทสา โข ราชา อุเทโน โอโรธํ ทูรโตว อาคจฺฉนฺตํ, ทิสฺวาน โอโรธํ เอตทโวจ “อปิ นุ โข ตุเมฺห สมณํ อานนฺทํ ปสฺสิตฺถา”ติ. อปสฺสิมฺหา โข มยํ เทว อยฺยํ อานนฺทนฺติ. อปิ นุ ตุเมฺห สมณสฺส อานนฺทสฺส กิญฺจิ อทตฺถาติ. อทมฺหา โข มยํ เทว อยฺยสฺส อานนฺทสฺส ปญฺจ อุตฺตราสงฺคสตานีติ. ราชา อุเทโน อุชฺฌายติ ขิยฺยติ วิปาเจติ “กถํ หิ นาม สมโณ อานนฺโท ตาว พหุํ จีวรํ ปฏิคฺคเหสฺสติ, ทุสฺสวาณิชฺชํ วา สมโณ อานนฺโท กริสฺสติ, ปคฺคาหิกสาลํ วา ปสาเรสฺสตี”ติ.}

\prose{อถ โข ราชา อุเทโน เยนายสฺมา อานนฺโท เตนุปสงฺกมิ, อุปสงฺกมิตฺวา อายสฺมตา อานนฺเทน สทฺธิํ สมฺโมทิ, สมฺโมทนียํ กถํ สารณียํ วีติสาเรตฺวา เอกมนฺตํ นิสีทิ, เอกมนฺตํ นิสินฺโน โข ราชา อุเทโน อายสฺมนฺตํ อานนฺทํ เอตทโวจ “อาคมา นุ ขฺวิธ โภ อานนฺท อมฺหากํ โอโรโธ”ติ. อาคมาสิ โข เต อิธ มหาราช โอโรโธติ. อปิ ปน โภโต อานนฺทสฺส กิญฺจิ อทาสีติ. อทาสิ โข เม มหาราช ปญฺจ อุตฺตราสงฺคสตานีติ. กิํ ปน ภวํ อานนฺโท ตาว พหุํ จีวรํ กริสฺสตีติ. เย\csromansharedfootnote{57677e2a71302e46}{เย ปน (ก)} เต มหาราช ภิกฺขู ทุพฺพลจีวรา, เตหิ สทฺธิํ สํวิภชิสฺสามีติ. ยานิ โข ปน โภ อานนฺท โปราณกานิ ทุพฺพล\-จีวรานิ, ตานิ กถํ กริสฺสถาติ. ตานิ มหาราช อุตฺตร\-ตฺถรณํ กริสฺสามาติ. ยานิ ปน โภ อานนฺท โปราณกานิ อุตฺตร\-ตฺถรณานิ, ตานิ กถํ กริสฺสถาติ. ตานิ มหาราช ภิสิจฺฉวิโย กริสฺสามาติ. ยา ปน โภ อานนฺท โปราณกา ภิสิจฺฉวิโย, ตา กถํ กริสฺสถาติ. ตา มหาราช ภูมตฺถรณํ กริสฺสามาติ. ยานิ ปน โภ อานนฺท โปราณกานิ ภูมตฺถรณานิ, ตานิ กถํ กริสฺสถาติ. ตานิ มหาราช ปาทปุญฺฉนิโย กริสฺสามาติ. ยา ปน โภ อานนฺท โปราณกา ปาทปุญฺฉนิโย, ตา กถํ กริสฺสถาติ. ตา มหาราช รโชหรณํ กริสฺสมาติ. ยานิ ปน \csromanfolio{489}โภ อานนฺท โปราณกานิ รโชหรณานิ, ตานิ กถํ กริสฺสถาติ. ตานิ มหาราช โกฏฺเฏตฺวา จิกฺขลฺเลน มทฺทิตฺวา ปริภณฺฑํ ลิมฺปิสฺสามาติ.}

\prose{อถ โข ราชา อุเทโน “สพฺเพวิเม สมณา สกฺยปุตฺติยา โยนิโส อุปเนนฺติ, น กุลวํ คเมนฺตี”ติ อายสฺมโต อานนฺทสฺส อญฺญานิปิ ปญฺจ ทุสฺสสตานิ ปาทาสิ. อยญฺจรหิ อายสฺมโต อานนฺทสฺส ปฐมํ จีวรภิกฺขา อุปฺปชฺชิ จีวรสหสฺสํ. อถ โข อายสฺมา อานนฺโท เยน โฆสิตาราโม เตนุปสงฺกมิ, อุปสงฺกมิตฺวา ปญฺญตฺเต อาสเน นิสีทิ. อถ โข อายสฺมา ฉนฺโน เยนายสฺมา อานนฺโท เตนุปสงฺกมิ, อุปสงฺกมิตฺวา อายสฺมนฺตํ อานนฺทํ อภิวาเทตฺวา เอกมนฺตํ นิสีทิ, เอกมนฺตํ นิสินฺนํ โข อายสฺมนฺตํ ฉนฺนํ อายสฺมา อานนฺโท เอตทโวจ “สํเฆน เต อาวุโส ฉนฺน พฺรหฺมทณฺโฑ อาณาปิโต”ติ.}

\prose{กตโม ปน ภนฺเต อานนฺท พฺรหฺมทณฺโฑ อาณาปิโตติ. ตฺวํ อาวุโส ฉนฺน ภิกฺขู ยํ อิจฺเฉยฺยาสิ ตํ วเทยฺยาสิ, ภิกฺขูหิ ตฺวํ เนว วตฺตพฺโพ, น โอวทิตพฺโพ, นานุสาสิตพฺโพติ. “นนฺวาหํ ภนฺเต อานนฺท หโต เอตฺตาวตา, ยโตหํ ภิกฺขูหิ เนว วตฺตพฺโพ, น โอวทิตพฺโพ, นานุสาสิตพฺโพ”ติ ตตฺเถว มุจฺฉิโต ปปโต. อถ โข อายสฺมา ฉนฺโน พฺรหฺมทณฺเฑน อฏฺฏียมาโน หรายมาโน ชิคุจฺฉมาโน เอโก วูปกฏฺโฐ อปฺปมตฺโต อาตาปิ ปหิตตฺโต วิหรนฺโต นจิรสฺเสว ยสฺสตฺถาย กุลปุตฺตา สมฺมเทว อคารสฺมา อนคาริยํ ปพฺพชนฺติ, ตทนุตฺตรํ พฺรหฺมจริย\-ปริโยสานํ ทิฏฺเฐว ธมฺเม สยํ อภิณฺญา สจฺฉิกตฺวา อุปสมฺปชฺช วิหาสิ, “ขีณา ชาติ, วุสิตํ พฺรหฺมจริยํ, กตํ กรณียํ, นาปรํ อิตฺถตฺตายา”ติ อพฺภญฺญาสิ, อญฺญตโร จ ปนายสฺมา ฉนฺโน อรหตํ อโหสิ. อถ โข อายสฺมา ฉนฺโน อรหตฺตํ ปตฺโต เยนายสฺมา อานนฺโท เตนุปสงฺกมิ, อุปสงฺกมิตฺวา อายสฺมนฺตํ อานนฺทํ เอตทโวจ “ปฏิปฺปสฺสมฺเภหิ ทานิ เม ภนฺเต อานนฺท พฺรหฺมทณฺฑนฺ”ติ. ยทคฺเคน ตยา อาวุโส ฉนฺน อรหตฺตํ สจฺฉิกตํ, ตทคฺเคน เต พฺรหฺมทณฺโฑ ปฏิปฺปสฺสทฺโธติ. อิมาย โข ปน วินย\-สงฺคีติยา ปญฺจ ภิกฺขุสตานิ อนูนานิ อนธิกานิ อเหสุํ. ตสฺมา อยํ วินยสงฺคีติ “ปญฺจสติกา”ติ วุจฺจตีติ.}

\vspace{\tipitakaparskip}
\csromancenter{ปญฺจสติก\-กฺขนฺธโก เอกาทสโม.}
\csromancenter{อิมมฺหิ ขนฺธเก วตฺถู เตวีสติ.}
\csromancenter{\textbf{ตสฺสุทฺทานํ}}
\csromanmatikaanchor{mtk.216}
\csromantocmarkref{section}{อุทฺทานคาถา}{mtk.216}
\bookmark[dest={mtk.216},level=1]{อุทฺทานคาถา}
\setlength{\csromangathapagewidth}{0pt}
\setlength{\csromangathawakwidth}{0pt}
\csromangathameasuregroup{ปรินิพฺพุเต สมฺพุทฺเธ, \\ อามนฺตยิ ภิกฺขุคณํ, \\ ปาวายทฺธานมคฺคมฺหิ, \\ สงฺคายิสฺสาม สทฺธมฺมํ, \\ เอเกนูน ปญฺจสตํ, \\ ธมฺมวินยสงฺคีติํ, \\ อุปาลิํ วินยํ ปุจฺฉิ, \\ ปิฏกํ ตีณิ สงฺคีติํ, \\ ขุทฺทานุขุทฺทเก นานา, \\ น ปุจฺฉิ อกฺกมิตฺวาน, \\ ปพฺพชฺชํ มาตุคามสฺส, \\ ปุราโณ พฺรหฺมทณฺฑญฺจ, \\ ตาว พหุ ทุพฺพลญฺจ, \\ ภูมตฺถรณา ปุญฺฉนิโย, \\ สหสฺสจีวรํ อุปฺปชฺชิ, \\ ตชฺชิโต พฺรหฺมทณฺเฑน,}{\csromangathabat{ปรินิพฺพุเต สมฺพุทฺเธ,}{เถโร กสฺสปสวฺหโย.} \\ \csromangathabat{อามนฺตยิ ภิกฺขุคณํ,}{สทฺธมฺมมนุปาลโก.} \\ \csromangathabat{ปาวายทฺธานมคฺคมฺหิ,}{สุภทฺเทน ปเวทิตํ.} \\ \csromangathabat{สงฺคายิสฺสาม สทฺธมฺมํ,}{อธมฺโม ปุเร ทิปฺปติ.} \\ \csromangathabat{เอเกนูน ปญฺจสตํ,}{อานนฺทมฺปิ จ อุจฺจินิ.} \\ \csromangathabat{ธมฺมวินยสงฺคีติํ,}{วสนฺโต คุหมุตฺตเม.} \\ \csromangathabat{อุปาลิํ วินยํ ปุจฺฉิ,}{สุตฺตนฺตานนฺทปณฺฑิตํ.} \\ \csromangathabat{ปิฏกํ ตีณิ สงฺคีติํ,}{อกํสุ ชินสาวกา.} \\ \csromangathabat{ขุทฺทานุขุทฺทเก นานา,}{ยถาปญฺญตฺติวตฺตนา.} \\ \csromangathabat{น ปุจฺฉิ อกฺกมิตฺวาน,}{วนฺทาเปสิ น ยาจิ จ.} \\ \csromangathabat{ปพฺพชฺชํ มาตุคามสฺส,}{สทฺธาย ทุกฺกฏานิ เม.} \\ \csromangathabat{ปุราโณ พฺรหฺมทณฺฑญฺจ,}{โอโรโธ อุเทเนน สห.} \\ \csromangathabat{ตาว พหุ ทุพฺพลญฺจ,}{อุตฺตรตฺถรณา ภิสิ.} \\ \csromangathabat{ภูมตฺถรณา ปุญฺฉนิโย,}{รโช จิกฺขลฺลมทฺทนา.} \\ \csromangathabat{สหสฺสจีวรํ อุปฺปชฺชิ,}{ปฐมานนฺทสวฺหโย.} \\ \csromangathabat{ตชฺชิโต พฺรหฺมทณฺเฑน,}{จตุสฺสจฺจํ อปาปุณิ.}}
\csromangathasetleft{ปรินิพฺพุเต สมฺพุทฺเธ, \\ อามนฺตยิ ภิกฺขุคณํ, \\ ปาวายทฺธานมคฺคมฺหิ, \\ สงฺคายิสฺสาม สทฺธมฺมํ, \\ เอเกนูน ปญฺจสตํ, \\ ธมฺมวินยสงฺคีติํ, \\ อุปาลิํ วินยํ ปุจฺฉิ, \\ ปิฏกํ ตีณิ สงฺคีติํ, \\ ขุทฺทานุขุทฺทเก นานา, \\ น ปุจฺฉิ อกฺกมิตฺวาน, \\ ปพฺพชฺชํ มาตุคามสฺส, \\ ปุราโณ พฺรหฺมทณฺฑญฺจ, \\ ตาว พหุ ทุพฺพลญฺจ, \\ ภูมตฺถรณา ปุญฺฉนิโย, \\ สหสฺสจีวรํ อุปฺปชฺชิ, \\ ตชฺชิโต พฺรหฺมทณฺเฑน,}
\csromangathagroup{\csromangathastanza{\csromanfolio{490}\csromangathabat{ปรินิพฺพุเต สมฺพุทฺเธ,}{เถโร กสฺสปสวฺหโย.} \\ \csromangathabat{อามนฺตยิ ภิกฺขุคณํ,}{สทฺธมฺมมนุปาลโก.}}\csromangathastanza{\csromangathabat{ปาวายทฺธานมคฺคมฺหิ,}{สุภทฺเทน ปเวทิตํ.} \\ \csromangathabat{สงฺคายิสฺสาม สทฺธมฺมํ,}{อธมฺโม ปุเร ทิปฺปติ.}}\csromangathastanza{\csromangathabat{เอเกนูน ปญฺจสตํ,}{อานนฺทมฺปิ จ อุจฺจินิ.} \\ \csromangathabat{ธมฺมวินยสงฺคีติํ,}{วสนฺโต คุหมุตฺตเม.}}\csromangathastanza{\csromangathabat{อุปาลิํ วินยํ ปุจฺฉิ,}{สุตฺตนฺตานนฺทปณฺฑิตํ.} \\ \csromangathabat{ปิฏกํ ตีณิ สงฺคีติํ,}{อกํสุ ชินสาวกา.}}\csromangathastanza{\csromangathabat{ขุทฺทานุขุทฺทเก นานา,}{ยถาปญฺญตฺติวตฺตนา.} \\ \csromangathabat{น ปุจฺฉิ อกฺกมิตฺวาน,}{วนฺทาเปสิ น ยาจิ จ.}}\csromangathastanza{\csromangathabat{ปพฺพชฺชํ มาตุคามสฺส,}{สทฺธาย ทุกฺกฏานิ เม.} \\ \csromangathabat{ปุราโณ พฺรหฺมทณฺฑญฺจ,}{โอโรโธ อุเทเนน สห.}}\csromangathastanza{\csromangathabat{ตาว พหุ ทุพฺพลญฺจ,}{อุตฺตรตฺถรณา ภิสิ.} \\ \csromangathabat{ภูมตฺถรณา ปุญฺฉนิโย,}{รโช จิกฺขลฺลมทฺทนา.}}\csromangathastanza{\csromangathabat{สหสฺสจีวรํ อุปฺปชฺชิ,}{ปฐมานนฺทสวฺหโย.} \\ \csromangathabat{ตชฺชิโต พฺรหฺมทณฺเฑน,}{จตุสฺสจฺจํ อปาปุณิ.}}}

\prosewithclosermajor{วสีภูตา ปญฺจสตา, ตสฺมา “ปญฺจสตี” อิติ.}{ปญฺจสติก\-กฺขนฺธโก นิฏฺฐิโต.}
\csromanmatikaanchor{mtk.217}
\csromantocmarknopage{chapter}{12. สตฺตสติกกฺขนฺธก}{mtk.217}
\bookmark[dest={mtk.217},level=0]{12. สตฺตสติกกฺขนฺธก}
\chapterheadpageread{12. สตฺตสติก\-กฺขนฺธก}
\csromanmatikaanchor{mtk.218}
\csromantocmarkref{section}{1. ปฐมภาณวาร}{mtk.218}
\bookmark[dest={mtk.218},level=1]{1. ปฐมภาณวาร}
\csromanheader{1}{1. ปฐม\-ภาณวาร}
\proseitem{446}{\csromanfolio{491}เตน โข ปน สมเยน วสฺสสต\-ปรินิพฺพุเต ภควติ เวสาลิกา วชฺชิปุตฺตกา ภิกฺขู เวสาลิยํ ทส วตฺถูนิ ทีเปนฺติ “กปฺปติ สิงฺคิโลณกปฺโป, กปฺปติ ทฺวงฺคุลกปฺโป, กปฺปติ คามนฺตรกปฺโป, กปฺปติ อาวาสกปฺโป, กปฺปติ อนุมติกปฺโป, กปฺปติ อาจิณฺณกปฺโป, กปฺปติ อมถิตกปฺโป, กปฺปติ ชโฬคิํ ปาตุํ, กปฺปติ อทสกํ นิสีทนํ, กปฺปติ ชาตรูปรชตนฺ”ติ.}

\prose{เตน โข ปน สมเยน อายสฺมา ยโส กากณฺฑกปุตฺโตวชฺชีสุ จาริกํ จรมาโน เยน เวสาลี ตทวสริ, ตตฺร สุทํ อายสฺมา ยโส กากณฺฑกปุตฺโตเวสาลิยํ วิหรติ มหาวเน กูฏาคาร\-สาลายํ. เตน โข ปน สมเยน เวสาลิกา วชฺชิปุตฺตกา ภิกฺขู ตทหุโปสเถ กํสปาติํ\csromansharedfootnote{a4853758c123f9a4}{กํสจฏิํ (สฺยา)} อุทเกน ปูเรตฺวา มชฺเฌ ภิกฺขุ\-สํฆสฺส ฐเปตฺวา อาคตาคเต เวสาลิเก อุปาสเก เอวํ วทนฺติ “เทถาวุโส สํฆสฺส กหาปณมฺปิ อฑฺฒมฺปิ ปาทมฺปิ มาสกรูปมฺปิ, ภวิสฺสติ สํฆสฺส ปริกฺขาเรน กรณียนฺ”ติ. เอวํ วุตฺเต อายสฺมา ยโส กากณฺฑกปุตฺโต เวสาลิเก อุปาสเก เอตทโวจ “มาวุโส อทตฺถ สํฆสฺส กหาปณมฺปิ อฑฺฒมฺปิ ปาทมฺปิ มาสกรูปมฺปิ, น กปฺปติ สมณานํ สกฺย\-ปุตฺติ\-ยานํ ชาตรูป\-รชตํ, น สาทิยนฺติ สมณา สกฺยปุตฺติยา ชาตรูป\-รชตํ, น ปฏิคฺคณฺหนฺติ สมณา สกฺยปุตฺติยา ชาตรูป\-รชตํ, นิกฺขิตฺต\-มณิ\-สุวณฺณา สมณา สกฺยปุตฺติยา, อเปต\-ชาตรูปรชตา”ติ. เอวมฺปิ โข เวสาลิกา อุปาสกา อายสฺมตา ยเสน กากณฺฑก\-ปุตฺเตน วุจฺจมานา อทํสุเยว สํฆสฺส กหาปณมฺปิ อฑฺฒมฺปิ ปาทมฺปิ มาสกรูปมฺปิ.}

\prose{อถ โข เวสาลิกา วชฺชิปุตฺตกา ภิกฺขู ตสฺสา รตฺติยา อจฺจเยน ตํ หิรญฺญํ ภิกฺขคฺเคน\csromansharedfootnote{c9ea2b2ddc6914f7}{ภิกฺขุคฺเคน (สฺยา)} ปฏิวีสํ\csromansharedfootnote{fc5181e0c9b6a49b}{ปฏิวิํสํ (สี), ปฏิวิสํ (สฺยา)} ฐเปตฺวา ภาเชสุํ. อถ โข เวสาลิกา วชฺชิปุตฺตกา ภิกฺขู อายสฺมนฺตํ ยสํ กากณฺฑก\-ปุตฺตํ เอตทโวจุํ “เอโส เต อาวุโส ยส หิรญฺญสฺส ปฏิวีโส”ติ. นตฺถิ เม อาวุโส หิรญฺญสฺส ปฏิวีโส, นาหํ หิรญฺญํ สาทิยามีติ. \csromanfolio{492}อถ โข เวสาลิกา วชฺชิปุตฺตกา ภิกฺขู “อยํ อาวุโส ยโส กากณฺฑกปุตฺโต อุปาสเก สทฺเธ ปสนฺเน อกฺโกสติ ปริภาสติ อปฺปสาทํ กโรติ, หนฺทสฺส มยํ ปฏิสารณีย\-กมฺมํ กโรมา”ติ, เต ตสฺส ปฏิสารณีย\-กมฺมํ อกํสุ. อถ โข อายสฺมา ยโส กากณฺฑกปุตฺโต เวสาลิเก วชฺชิปุตฺตเก ภิกฺขู เอตทโวจ “ภควโต อาวุโส ปญฺญตฺตํ ‘ปฏิสารณีย\-กมฺมกตสฺส ภิกฺขุโน อนุทูโต ทาตพฺโพ’ติ, เทถ เม อาวุโส อนุทูตํ ภิกฺขุนฺ”ติ.}

\prose{อถ โข เวสาลิกา วชฺชิปุตฺตกา ภิกฺขู เอกํ ภิกฺขุํ สมฺมนฺนิตฺวา อายสฺมโต ยสสฺส กากณฺฑก\-ปุตฺตสฺส อนุทูตํ อทํสุ. อถ โข อายสฺมา ยโส กากณฺฑกปุตฺโต อนุทูเตน ภิกฺขุนา สทฺธิํ เวสาลิํ ปวิสิตฺวา เวสาลิเก อุปาสเก เอตทโวจ\csromandash{}อหํ กิรายสฺมนฺเต อุปาสเก สทฺเธ ปสนฺเน อกฺโกสามิ ปริภาสามิ อปฺปสาทํ กโรมิ, โยหํ อธมฺมํ “อธมฺโม”ติ วทามิ, ธมฺมํ “ธมฺโม”ติ วทามิ, อวินยํ “อวินโย”ติ วทามิ, วินยํ “วินโย”ติ วทามิ.}

\proseitem{447}{เอกมิทํ อาวุโส สมยํ ภควา สาวตฺถิยํ วิหรติ เชตวเน อนาถ\-ปิณฺฑิกสฺส อาราเม, ตตฺร โข อาวุโส ภควา ภิกฺขู อามนฺเตสิ\csromandash{} *จตฺตาโรเม ภิกฺขเว จนฺทิม\-สูริยานํ อุปกฺกิเลสา, เยหิ อุปกฺกิเลเสหิ อุปกฺกิลิฏฺฐา จนฺทิมสูริยา น ตปนฺติ, น ภาสนฺติ, น วิโรจนฺติ. กตเม จตฺตาโร, อพฺภํ ภิกฺขเว จนฺทิม\-สูริยานํ อุปกฺกิเลโส, เยน อุปกฺกิเลเสน อุปกฺกิลิฏฺฐา จนฺทิมสูริยา น ตปนฺติ, น ภาสนฺติ, น วิโรจนฺติ. มหิกา ภิกฺขเว จนฺทิม\-สูริยานํ อุปกฺกิเลโส, เยน อุปกฺกิเลเสน อุปกฺกิลิฏฺฐา จนฺทิมสูริยา น ตปนฺติ, น ภาสนฺติ, น วิโรจนฺติ. ธูมรโช ภิกฺขเว จนฺทิม\-สูริยานํ อุปกฺกิเลโส, เยน อุปกฺกิเลเสน อุปกฺกิลิฏฺฐา จนฺทิมสูริยา น ตปนฺติ, น ภาสนฺติ, น วิโรจนฺติ. ราหุ ภิกฺขเว อสุรินฺโท จนฺทิม\-สูริยานํ อุปกฺกิเลโส, เยน อุปกฺกิเลเสน อุปกฺกิลิฏฺฐา จนฺทิมสูริยา น ตปนฺติ, น ภาสนฺติ, น วิโรจนฺติ. อิเม โข ภิกฺขเว จตฺตาโร จนฺทิม\-สูริยานํ อุปกฺกิเลสา, เยหิ อุปกฺกิเลเสหิ อุปกฺกิลิฏฺฐา จนฺทิมสูริยา น ตปนฺติ, น ภาสนฺติ, น วิโรจนฺติ. เอวเมว โข ภิกฺขเว จตฺตาโรเม สมณ\-พฺราหฺมณานํ อุปกฺกิเลสา, เยหิ อุปกฺกิเลเสหิ อุปกฺกิลิฏฺฐา เอเก สมณพฺราหฺมณา \csromanfolio{493}น ตปนฺติ, น ภาสนฺติ, น วิโรจนฺติ. กตเม จตฺตาโร, สนฺติ ภิกฺขเว เอเก สมณพฺราหฺมณา สุรํ ปิวนฺติ, เมรยํ ปิวนฺติ, สุราเมรยปานา อปฺปฏิวิรตา, อยํ ภิกฺขเว ปฐโม สมณ\-พฺราหฺมณานํ อุปกฺกิเลโส, เยน อุปกฺกิเลเสน อุปกฺกิลิฏฺฐา เอเก สมณพฺราหฺมณา น ตปนฺติ, น ภาสนฺติ, น วิโรจนฺติ. ปุน จปรํ ภิกฺขเว เอเก สมณพฺรหฺมณา เมถุนํ ธมฺมํ ปฏิเสวนฺติ, เมถุนธมฺมา อปฺปฏิวิรตา, อยํ ภิกฺขเว ทุติโย สมณ\-พฺราหฺมณานํ อุปกฺกิเลโส, เยน อุปกฺกิเลเสน อุปกฺกิลิฏฺฐา เอเก สมณพฺราหฺมณา น ตปนฺติ, น ภาสนฺติ, น วิโรจนฺติ. ปุน จปรํ ภิกฺขเว เอเก สมณพฺราหฺมณา ชาตรูป\-รชตํ สาทิยนฺติ, ชาตรูปรชต\-ปฺปฏิคฺคหณา อปฺปฏิวิรตา, อยํ ภิกฺขเว ตติโย สมณ\-พฺราหฺมณานํ อุปกฺกิเลโส, เยน อุปกฺกิเลเสน อุปกฺกิลิฏฺฐา เอเก สมณพฺราหฺมณา น ตปนฺติ, น ภาสนฺติ, น วิโรจนฺติ. ปุน จปรํ ภิกฺขเว เอเก สมณ พฺราหฺมณา มิจฺฉาชีเวน ชีวิตํ กปฺเปนฺติ, มิจฺฉาชีวา อปฺปฏิวิรตา, อยํ ภิกฺขเว จตุตฺโถ สมณ\-พฺราหฺมณานํ อุปกฺกิเลโส, เยน อุปกฺกิเลเสน อุปกฺกิลิฏฺฐา เอเก สมณพฺราหฺมณา น ตปนฺติ, น ภาสนฺติ, น วิโรจนฺติ. อิเม โข ภิกฺขเว จตฺตาโร สมณ\-พฺราหฺมณานํ อุปกฺกิเลสา, เยหิ อุปกฺกิเลเสหิ อุปกฺกิลิฏฺฐา เอเก สมณพฺราหฺมณา น ตปนฺติ, น ภาสนฺติ, น วิโรจนฺติ. อิทมโวจาวุโส ภควา, อิทํ วตฺวาน สุคโต อถาปรํ เอตทโวจ สตฺถา.}

\setlength{\csromangathapagewidth}{0pt}
\setlength{\csromangathawakwidth}{0pt}
\csromangathameasuregroup{ราคโทสปริกฺลิฏฺฐา, \\ อวิชฺชานิวุฏา\textsuperscript{1} โปสา, \\ สุรํ ปิวนฺติ เมรยํ, \\ รชตํ ชาตรูปญฺจ, \\ มิจฺฉาชีเวน ชีวนฺติ, \\ เอเต อุปกฺกิเลสา วุตฺตา, \\ เยหิ อุปกฺกิเลเสหิ อุปกฺกิลิฏฺฐา, \\ น ตปนฺติ น ภาสนฺติ, \\ อนฺธกาเรน โอนทฺธา,}{\csromangathabat{ราคโทสปริกฺลิฏฺฐา,}{เอเก สมณพฺราหฺมณา.} \\ \csromangathabat{อวิชฺชานิวุฏา\textsuperscript{1} โปสา,}{ปิยรูปาภินนฺทิโน.} \\ \csromangathabat{สุรํ ปิวนฺติ เมรยํ,}{ปฏิเสวนฺติ เมถุนํ.} \\ \csromangathabat{รชตํ ชาตรูปญฺจ,}{สาทิยนฺติ อวิทฺทสู.} \\ \csromangathabat{มิจฺฉาชีเวน ชีวนฺติ,}{เอเก สมณพฺราหฺมณา.} \\ \csromangathabat{เอเต อุปกฺกิเลสา วุตฺตา,}{พุทฺเธนาทิจฺจพนฺธุนา.} \\ \csromangathabat{เยหิ อุปกฺกิเลเสหิ อุปกฺกิลิฏฺฐา,}{เอเก สมณพฺราหฺมณา.} \\ \csromangathabat{น ตปนฺติ น ภาสนฺติ,}{อสุทฺธา สรชา มคา.} \\ \csromangathabat{อนฺธกาเรน โอนทฺธา,}{ตณฺหาทาสา สเนตฺติกา.}}
\csromangathasetleft{ราคโทสปริกฺลิฏฺฐา, \\ อวิชฺชานิวุฏา\textsuperscript{1} โปสา, \\ สุรํ ปิวนฺติ เมรยํ, \\ รชตํ ชาตรูปญฺจ, \\ มิจฺฉาชีเวน ชีวนฺติ, \\ เอเต อุปกฺกิเลสา วุตฺตา, \\ เยหิ อุปกฺกิเลเสหิ อุปกฺกิลิฏฺฐา, \\ น ตปนฺติ น ภาสนฺติ, \\ อนฺธกาเรน โอนทฺธา,}
\csromangathagroup{\csromangathastanza{\csromangathabat{ราคโทสปริกฺลิฏฺฐา,}{เอเก สมณพฺราหฺมณา.} \\ \csromangathabat{อวิชฺชานิวุฏา\csromansharedfootnote{c5f02d9549809539}{อวิชฺชานีวุตา (สฺยา)} โปสา,}{ปิยรูปาภินนฺทิโน.}}\csromangathastanza{\csromangathabat{สุรํ ปิวนฺติ เมรยํ,}{ปฏิเสวนฺติ เมถุนํ.} \\ \csromangathabat{รชตํ ชาตรูปญฺจ,}{สาทิยนฺติ อวิทฺทสู.}}\csromangathastanza{\csromangathabat{มิจฺฉาชีเวน ชีวนฺติ,}{เอเก สมณพฺราหฺมณา.} \\ \csromangathabat{เอเต อุปกฺกิเลสา วุตฺตา,}{พุทฺเธนา\-ทิจฺจพนฺธุนา.}}\csromangathastanza{\csromangathabat{เยหิ อุปกฺกิเลเสหิ อุปกฺกิลิฏฺฐา,}{เอเก สมณพฺราหฺมณา.} \\ \csromangathabat{น ตปนฺติ น ภาสนฺติ,}{อสุทฺธา สรชา มคา.} \\ \csromangathabat{อนฺธกาเรน โอนทฺธา,}{ตณฺหาทาสา สเนตฺติกา.}}}

\prose{วฑฺเฒนฺติ กฏสิํ โฆรํ, อาทิยนฺติ ปุนพฺภวนฺติ}

\prose{\csromanfolio{494}เอวํวาที กิราหํ อายสฺมนฺเต อุปาสเก สทฺเธ ปสนฺเน อกฺโกสามิ ปริภาสามิ อปฺปสาทํ กโรมิ, โยหํ อธมฺมํ “อธมฺโม”ติ วทามิ, ธมฺมํ “ธมฺโม”ติ วทามิ, อวินยํ “อวินโย”ติ วทามิ, วินยํ “วินโย”ติ วทามิ.}

\proseitem{448}{\csromansymbolfootnote{*}{สํ 2. 509 ปิฏฺเฐปิ อิทํ วตฺถุ อาคตํ.}เอกมิทํ อาวุโส สมยํ ภควา ราชคเห วิหรติ เวฬุวเน กลนฺทก\-นิวาเป. เตน โข ปน สมเยน ราชนฺเตปุเร ราชปริสายํ สนฺนิสินฺนานํ สนฺนิปติตานํ อยม\-นฺตรกถา อุทปาทิ “กปฺปติ สมณานํ สกฺย\-ปุตฺติ\-ยานํ ชาตรูป\-รชตํ, สาทิยนฺติ สมณา สกฺยปุตฺติยา ชาตรูป\-รชตํ, ปฏิคฺคณฺหนฺติ สมณา สกฺยปุตฺติยา ชาตรูปรชตนฺ”ติ. เตน โข ปนาวุโส สมเยน มณิจูฬโก คามณี ตสฺสํ ปริสายํ นิสินฺโน โหติ. อถ โข อาวุโส มณิจูฬโก คามณี ตํ ปริสํ เอตทโวจ “มา อยฺยา เอวํ อวจุตฺถ, น กปฺปติ สมณานํ สกฺย\-ปุตฺติ\-ยานํ ชาตรูป\-รชตํ, น สาทิยนฺติ สมณา สกฺยปุตฺติยา ชาตรูป\-รชตํ, น ปฏิคฺคณฺหนฺติ สมณา สกฺยปุตฺติยา ชาตรูป\-รชตํ นิกฺขิตฺต\-มณิ\-สุวณฺณา สมณา สกฺยปุตฺติยา อเปต\-ชาตรูปรชตา”ติ. อสกฺขิ โข อาวุโส มณิจูฬโก คามณี ตํ ปริสํ สญฺญาเปตุํ.}

\prose{อถ โข อาวุโส มณิจูฬโก คามณี ตํ ปริสํ สญฺญาเปตฺวา เยน ภควา เตนุปสงฺกมิ, อุปสงฺกมิตฺวา ภควนฺตํ อภิวาเทตฺวา เอกมนฺตํ นิสีทิ, เอกมนฺตํ นิสินฺโน โข อาวุโส มณิจูฬโก คามณี ภควนฺตํ เอตทโวจ “อิธ ภนฺเต ราชนฺเตปุเร ราชปริสายํ สนฺนิสินฺนานํ สนฺนิปติตานํ อยม\-นฺตรกถา อุทปาทิ ‘กปฺปติ สมณานํ สกฺย\-ปุตฺติ\-ยานํ ชาตรูป\-รชตํ, สาทิยนฺติ สมณา สกฺยปุตฺติยา ชาตรูป\-รชตํ, ปฏิคฺคณฺหนฺติ สมณา สกฺยปุตฺติยา ชาตรูปรชตนฺ’ติ, เอวํ วุตฺเต อหํ ภนฺเต ตํ ปริสํ เอตทโวจํ ‘มา อยฺยา เอวํ อวจุตฺถ, น กปฺปติ สมณานํ สกฺย\-ปุตฺติ\-ยานํ ชาตรูป\-รชตํ, น สาทิยนฺติ สมณา สกฺยปุตฺติยา ชาตรูป\-รชตํ, น ปฏิคฺคณฺหนฺติ สมณา สกฺยปุตฺติยา ชาตรูป\-รชตํ, นิกฺขิตฺต\-มณิ\-สุวณฺณา สมณา สกฺยปุตฺติยา อเปตชาตรูปรชาตา’ติ. อสกฺขิํ โข อหํ ภนฺเต ตํ ปริสํ สญฺญาเปตุํ, กจฺจาหํ ภนฺเต เอวํ พฺยากรมาโน วุตฺตวาที เจว ภควโต โหมิ, น จ ภควนฺตํ อภูเตน อพฺภาจิกฺขามิ, ธมฺมสฺส จานุธมฺมํ พฺยากโรมิ, น จ โกจิ สหธมฺมิโก วาทานุวาโท คารยฺหํ ฐานํ อาคจฺฉตี”ติ. ตคฺฆ ตฺวํ คามณิ เอวํ พฺยากรมาโน วุตฺตวาที เจว เม โหสิ, น จ มํ อภูเตน \csromanfolio{495}อพฺภาจิกฺขสิ, ธมฺมสฺส จานุธมฺมํ พฺยากโรสิ, น จโกจิ สหธมฺมิโก วาทานุวาโท คารยฺหํ ฐานํ อาคจฺฉติ, น หิ คามณิ กปฺปติ สมณานํ สกฺย\-ปุตฺติ\-ยานํ ชาตรูป\-รชตํ, น สาทิยนฺติ สมณา สกฺยปุตฺติยา ชาตรูป\-รชตํ, น ปฏิคฺคณฺหนฺติ สมณา สกฺยปุตฺติยา ชาตรูป\-รชตํ, นิกฺขิตฺต\-มณิ\-สุวณฺณา สมณา สกฺยปุตฺติยา อเปต\-ชาตรูปรชตา. ยสฺส โข คามณิ ชาตรูป\-รชตํ กปฺปติ, ปญฺจปิ ตสฺส กามคุณา กปฺปนฺติ. ยสฺส ปญฺจ กามคุณา กปฺปนฺติ, เอกํเสเนตํ คามณิ ธาเรยฺยาสิ “อสฺสมณธมฺโม อสกฺยปุตฺติยธมฺโม”ติ. อปิ จาหํ คามณิ เอวํ วทามิ “ติณํ ติณตฺถิเกน ปริเยสิตพฺพํ, ทารุ ทารุตฺถิเกน ปริเยสิตพฺพํ, สกฏํ สกฏตฺถิเกน ปริเยสิตพฺพํ, ปุริโส ปุริสตฺถิเกน ปริเยสิทฺตพฺโพ, น เตฺววาหํ คามณี เกนจิ ปริยาเยน ชาตรูป\-รชตํ ‘สาทิตพฺพํ ปริเยสิตพฺพนฺ’ติ วทามี”ติ.}

\prose{เอวํวาที กิราหํ อายสฺมนฺเต อุปาสเก สทฺเธ ปสนฺเน อกฺโกสามิ ปริภาสามิ อปฺปสาทํ กโรมิ, โยหํ อธมฺมํ “อธมฺโม”ติ วทามิ, ธมฺมํ “ธมฺโม”ติ วทามิ, อวินยํ “อวินโย”ติ วทามิ, วินยํ “วินโย”ติ วทามิ.}

\proseitem{449}{\csromansymbolfootnote{*}{วิ 1. 344 ปิฏฺเฐ.}เอกมิทํ อาวุโส สมยํ ภควา ราชคเห อายสฺมนฺตํ อุปนนฺทํ สกฺยปุตฺตํ อารพฺภ ชาตรูป\-รชตํ ปฏิกฺขิปิ, สิกฺขาปทญฺจ ปญฺญเปสิ, เอวํวาที กิราหํ อายสฺมนฺเต อุปาสเก สทฺเธ ปสนฺเน อกฺโกสามิ ปริภาสามิ อปฺปสาทํ กโรมิ, โยหํ อธมฺมํ “อธมฺโม”ติ วทามิ, ธมฺมํ “ธมฺโม”ติ วทามิ, อวินยํ “อวินโย”ติ วาทามิ, วินยํ “วินโย”ติ วทามีติ.}

\prose{เอวํ วุตฺเต เวสาลิกา อุปาสกา อายสฺมนฺตํ ยสํ กากณฺฑก\-ปุตฺตํ เอตทโวจุํ “เอโกว ภนฺเต อยฺโย ยโส กากณฺฑกปุตฺโตสมโณ สกฺยปุตฺติโย, สพฺเพวิเม อสฺสมณา อสกฺยปุตฺติยา, วสตุ ภนฺเต อยฺโย ยโส กากณฺฑกปุตฺโตเวสาลิยํ, มยํ อยฺยสฺส ยสสฺส กากณฺฑก\-ปุตฺตสฺส อุสฺสุกฺกํ กริสฺสาม จีวร\-ปิณฺฑปาต\-เสนาสน\-คิลานปฺปจฺจย\-เภสชฺช\-ปริกฺขารานนฺ”ติ. อถ โข อายสฺมา ยโส กากณฺฑกปุตฺโต เวสาลิเก อุปาสเก สญฺญาเปตฺวา อนุทูเตน ภิกฺขุนา สทฺธิํ อารามํ อคมาสิ.}

\prose{อถ โข เวสาลิกา วชฺชิปุตฺตกา ภิกฺขู อนุทูตํ ภิกฺขุํ ปุจฺฉิํสุ “ขมาปิตาวุโส ยเสน กากณฺฑก\-ปุตฺเตน เวสาลิกา \csromanfolio{496}อุปาสกา”ติ. อุปาสเกหิ ปาปิกํ โน อาวุโส กตํ, เอโกว ยโส กากณฺฑกปุตฺโตสมโณ สกฺยปุตฺติโย กโต, สพฺเพว มยํ อสฺสมณา อสกฺยปุตฺติยา กตาติ. อถ โข เวสาลิกา วชฺชิปุตฺตกา ภิกฺขู “อยํ อาวุโส ยโส กากณฺฑกปุตฺโต อเมฺหหิ อสมฺมโต คิหีนํ ปกาเสสิ, หนฺทสฺส มยํ อุกฺเขปนีย\-กมฺมํ กโรมา”ติ, เต ตสฺส อุกฺเขปนีย\-กมฺมํ กตฺตุกามา สนฺนิปติํสุ. อถ โข อายสฺมา ยโส กากณฺฑกปุตฺโต เวหาสํ อพฺภุคฺคนฺตฺวา โกสมฺพิยํ ปจฺจุฏฺฐาสิ.}

\proseitem{450}{อถ โข อายสฺมา ยโส กากณฺฑกปุตฺโต ปาเวยฺยกานญฺจ\csromansharedfootnote{920e24fc16517b5b}{ปาเฐยฺยกานญฺจ (สฺยา)} อวนฺติ\-ทกฺขิณาปถ\-กานญฺจ ภิกฺขูนํ สนฺติเก ทูตํ ปาเหสิ “อาคจฺฉนฺตุ อายสฺมนฺตา, อิมํ อธิกรณํ อาทิยิสฺสาม, ปุเร อธมฺโม ทิปฺปติ, ธมฺโม ปฏิพาหิยฺยติ, อวินโย ทิปฺปติ, วินโย ปฏิพาหิยฺยติ, ปุเร อธมฺมวาทิโน พลวนฺโต โหนฺติ, ธมฺมวาทิโน ทุพฺพลา โหนฺติ, อวินยวาทิโน พลวนฺโต โหนฺติ, วินยวาทิโน ทุพฺพลา โหนฺตี”ติ.}

\prose{เตน โข ปน สมเยน อายสฺมา สมฺภูโต สาณวาสี อโหคงฺเค ปพฺพเต ปฏิวสติ. อถ โข อายสฺมา ยโส กากณฺฑกปุตฺโต เยน อโหคงฺโค ปพฺพโต เยนายสฺมา สมฺภูโต สาณวาสี เตนุปสงฺกมิ, อุปสงฺกมิตฺวา อายสฺมนฺตํ สมฺภูตํ สาณวาสิํ อภิวาเทตฺวา เอกมนฺตํ นิสีทิ, เอกมนฺตํ นิสินฺโน โข อายสฺมา ยโส กากณฺฑกปุตฺโต อายสฺมนฺตํ สมฺภูตํ สาณวาสิํ เอตทโวจ “อิเม ภนฺเต เวสาลิกา วชฺชิปุตฺตกา ภิกฺขู เวสาลิยํ ทส วตฺถูนิ ทีเปนฺติ ‘กปฺปติ สงฺคิโลณกปฺโป, กปฺปติ ทฺวงฺคุลกปฺโป, กปฺปติ คามนฺตรกปฺโป, กปฺปติ อาวาสกปฺโป, กปฺปติ อนุมติกปฺโป, กปฺปติ อาจิณฺณกปฺโป, กปฺปติ อมถิตกปฺโป, กปฺปติ ชโฬคิํ ปาตุํ, กปฺปติ อทสกํ นิสีทนํ, กปฺปติ ชาตรูปรชตนฺ’ติ, หนฺท มยํ ภนฺเต อิมํ อธิกรณํ อาทิยิสฺสาม, ปุเร อธมฺโม ทิปฺปติ, ธมฺโม ปฏิพาหิยฺยติ, อวินโย ทิปฺปติ, วินโย ปฏิคาหิยฺยติ, ปุเร อธมฺมวาทิโน พลวนฺโต โหนฺติ, ธมฺมวาทิโน ทุพฺพลา โหนฺติ, อวินยวาทิโน พลวนฺโต โหนฺติ, วินยวาทิโน ทุพฺพลา โหนฺตี”ติ. เอวมาวุโสติ โข อายสฺมา สมฺภูโต สาณวาสี อายสฺมโต ยสสฺส กากณฺฑก\-ปุตฺตสฺส ปจฺจสฺโสสิ. \csromanfolio{497}อถ โข สฏฺฐิมตฺตา ปาเวยฺยกา ภิกฺขู สพฺเพ อารญฺญิกา สพฺเพ ปิณฺณปาติกา สพฺเพ ปํสุกูลิกา สพฺเพ เตจีวริกา สพฺเพว อรหนฺโต อโหคงฺเค ปพฺพเต สนฺนิปติํสุ. อฏฺฐาสีติมตฺตา อวนฺติ\-ทกฺขิณาปถกา ภิกฺขู อปฺเปกจฺเจ อารญฺญิกา อปฺเปกจฺเจ ปิณฺฑปาติกา อปฺเปกจฺเจ ปํสุกูลิกา อปฺเปกจฺเจ เตจีวริกา สพฺเพว อรหนฺโต อโหคงฺเค ปพฺพเต สนฺนิปติํสุ. อถ โข เถรานํ ภิกฺขูนํ มนฺตยมานานํ เอตทโหสิ “อิทํ โข อธิกรณํ กกฺขฬญฺจ วาฬญฺจ, กํ นุ โข มยํ ปกฺขํ ลเภยฺยาม เยน มยํ อิมสฺมิํ อธิกรเณ พลวนฺตตรา อสฺสามา”ติ.}

\proseitem{451}{เตน โข ปน สมเยน อายสฺมา เรวโต โสเรยฺเย ปฏิวสติ พหุสฺสุโต อาคตาคโม ธมฺมธโร วินยธโร มาติกาธโร ปณฺฑิโต วิยตฺโต เมธาวี ลชฺชี กุกฺกุจฺจโก สิกฺขากาโม. อถ โข เถรานํ ภิกฺขูนํ เอตทโหสิ “อยํ โข อายสฺมา เรวโต โสเรยฺเย ปฏิวสติ พหุสฺสุโต อาคตาคโม ธมฺมธโร วินยธโร มาติกาธโร ปณฺฑิโต วิยตฺโต เมธาวี ลชฺชี กุกฺกุจฺจโก สิกฺขากาโม, สเจ มยํ อายสฺมนฺตํ เรวตํ ปกฺขํ ลภิสฺสาม, เอวํ มยํ อิมสฺมิํ อธิกรเณ พลวนฺตตรา อสฺสามา”ติ. อสฺโสสิ โข อายสฺมา เรวโต ทิพฺพาย โสตธาตุยา วิสุทฺธาย อติกฺกนฺตมานุสิกาย เถรานํ ภิกฺขูนํ มนฺตยมานานํ, สุตฺวานสฺส เอตทโหสิ “อิทํ โข อธิกรณํ กกฺขฬญฺจ วาฬญฺจ, น โข เมตํ ปติรูปํ, โยหํ เอวรูเป อธิกรเณ โอสกฺเกยฺยํ. อิทานิ จ ปน เต ภิกฺขู อาคจฺฉิสฺสนฺติ, โสหํ เตหิ อากิณฺโณ น ผาสุ คมิสฺสามิ, ยํนูนาหํ ปฏิกจฺเจว คจฺเฉยฺยนฺ”ติ. อถ โข อายสฺมา เรวโต โสเรยฺยา สงฺกสฺสํ อคมาสิ.}

\prose{อถ โข เถรา ภิกฺขู โสเรยฺยํ คนฺตฺวา ปุจฺฉิํสุ “กหํ อายสฺมา เรวโต”ติ. เต เอวมาหํสุ “เอสายสฺมา เรวโต สงฺกสฺสํ คโต”ติ. อถ โข อายสฺมา เรวโต สงฺกสฺสา กณฺณกุชฺชํ\csromansharedfootnote{5029b04292ed5043}{กนฺนกุชฺชํ (สี)} อคมาสิ. อถ โข เถรา ภิกฺขู สงฺกสฺสํ คนฺตฺวา ปุจฺฉิํสุ “กหํ อายสฺมา เรวโต”ติ. เต เอวมาหํสุ “เอสายสฺมา เรวโต กณฺณกุชฺชํ คโต”ติ. อถ โข อายสฺมา เรวโต กณฺณกุชฺชา \csromanfolio{498}อุทุมฺพรํ อคมาสิ. อถ โข เถรา ภิกฺขู กณฺณกุชฺชํ คนฺตฺวา ปุจฺฉิํสุ “กหํ อายสฺมา เรวโต”ติ. เต เอวมาหํสุ “เอสายสฺมา เรวโต อุทุมฺพรํ คโต”ติ. อถ โข อายสฺมา เรวโต อุทุมฺพรา อคฺคฬปุรํ อคมาสิ. อถ โข เถรา ภิกฺขู อุทุมฺพรํ คนฺตฺวา ปุจฺฉิํสุ “กหํ อายสฺมา เรวโต”ติ. เต เอวมาหํสุ “เอสายสฺมา เรวโต อคฺคฬปุรํ คโต”ติ. อถ โข อายสฺมา เรวโต อคฺคฬปุรา สหชาติํ\csromansharedfootnote{db04c2c5c28413c8}{สหํชาติํ (ก)} อคมาสิ. อถ โข เถรา ภิกฺขู อคฺคฬปุรํ คนฺตฺวา ปุจฺฉิํสุ “กหํ อายสฺมา เรวโต”ติ. เต เอวมาหํสุ “เอสายสฺมา เรวโต สหชาติํ คโต”ติ. อถ โข เถรา ภิกฺขู อายสฺมนฺตํ เรวตํ สหชาติยํ สมฺภาเวสุํ.}

\proseitemwithcloser{452}{อถ โข อายสฺมา สมฺภูโต สาณวาสี อายสฺมนฺตํ ยสํ กากณฺฑก\-ปุตฺตํ เอตทโวจ “อยํ อาวุโส อายสฺมา เรวโต พหุสฺสุโต อาคตาคโม ธมฺมธโร วินยธโร มาติกาธโร ปณฺฑิโต วิยตฺโต เมธาวี ลชฺชี กุกฺกุจฺจโก สิกฺขากาโม, สเจ มยํ อายสฺมนฺตํ เรวตํ ปญฺหํ ปุจฺฉิสฺสาม, ปฏิพโล อายสฺมา เรวโต เอเกเนว ปเญฺหน สกลมฺปิ รตฺติํ วีตินาเมตุํ, อิทานิ จ ปนายสฺมา เรวโต อนฺเตวาสิกํ\csromansharedfootnote{d6f4a7aecbea1090}{อนฺเตวาสิํ (สฺยา)} สรภาณกํ ภิกฺขุํ อชฺเฌสิสฺสติ, โส ตฺวํ ตสฺส ภิกฺขุโน สรภญฺญ\-ปริโยสาเน อายสฺมนฺตํ เรวตํ อุปสงฺกมิตฺวา อิมานิ ทส วตฺถูนิ ปุจฺเฉยฺยาสี”ติ. “เอวํ ภนฺเต”ติ โข อายสฺมา ยโส กากณฺฑกปุตฺโต อายสฺมโต สมฺภูตสฺส สาณวาสิสฺส ปจฺจสฺโสสิ. อถ โข อายสฺมา เรวโต อนฺเตวาสิกํ สรภาณกํ ภิกฺขุํ อชฺเฌสิ. อถ โข อายสฺมา ยโส กากณฺฑกปุตฺโตตสฺส ภิกฺขุโน สรภญฺญ\-ปริโยสาเน เยนายสฺมา เรวโต เตนุปสงฺกมิ, อุปสงฺกมิตฺวา อายสฺมนฺตํ เรวตํ อภิวาเทตฺวา เอกมนฺตํ นิสีทิ, เอกมนฺตํ นิสินฺโน โข อายสฺมา ยโส กากณฺฑกปุตฺโต อายสฺมนฺตํ เรวตํ เอตทโวจ “กปฺปติ ภนฺเต สิงฺคิโลณกปฺโป”ติ. โก โส อาวุโส สิงฺคิ\-โลณ\-กปฺโปติ. กปฺปติ ภนฺเต สิงฺคินา โลณํ ปริหริตุํ “ยตฺถ อโลณกํ ภวิสฺสติ, ตตฺถ ปริภุญฺชิสฺสามี”ติ. นาวุโส กปฺปตีติ. กปฺปติ ภนฺเต \csromanfolio{499}ทฺวงฺคุล\-กปฺโปติ. โก โส อาวุโส ทฺวงฺคุล\-กปฺโปติ. กปฺปติ ภนฺเต ทฺวงฺคุลาย ฉายาย วีติวตฺตาย วิกาเล โภชนํ ภุญฺชิตุนฺติ. นาวุโส กปฺปตีติ. กปฺปติ ภนฺเต คามนฺตร\-กปฺโปติ. โก โส อาวุโส คามนฺตร\-กปฺโปติ. กปฺปติ ภนฺเต “อิทานิ คามนฺตรํ คมิสฺสามี”ติ ภุตฺตาวินา ปวาริเตน อนติริตฺตํ โภชนํ ภุญฺชิตุนฺติ. นาวุโส กปฺปตีติ. กปฺปติ ภนฺเต อาวาสกปฺโปติ. โก โส อาวุโส อาวาสกปฺโปติ. กปฺปติ ภนฺเต สมฺพหุลา อาวาสา สมานสีมา นานุโปสถํ กาตุนฺติ. นาวุโส กปฺปตีติ. กปฺปติ ภนฺเต อนุมติกปฺโปติ. โก โส อาวุโส อนุมติกปฺโปติ. กปฺปติ ภนฺเต วคฺเคน สํเฆน กมฺมํ กาตุํ “อาคเต ภิกฺขู อนุมาเนสฺสามา”ติ\csromansharedfootnote{4d73f0c31cb63628}{อนุชานิสฺสามาติ, อนุชาเนสฺสามาติ, อนุมติํ อาเนสฺสามาติ (ก)}. นาวุโส กปฺปตีติ. กปฺปติ ภนฺเต อาจิณฺณกปฺโปติ. โก โส อาวุโส อาจิณฺณกปฺโปติ. กปฺปติ ภนฺเต อิทํ เม อุปชฺฌาเยน อชฺฌาจิณฺณํ, อิทํ เม อาจริเยน อชฺฌาจิณฺณํ, ตํ อชฺฌาจริตุนฺติ. อาจิณฺณกปฺโป โข อาวุโส เอกจฺโจ กปฺปติ, เอกจฺโจ น กปฺปตีติ. กปฺปติ ภนฺเต อมถิต\-กปฺโปติ. โก โส อาวุโส อมถิต\-กปฺโปติ. กปฺปติ ภนฺเต ยํ ตํ ขีรํ ขีรภาวํ วิชหิตํ อสมฺปตฺตํ ทธิภาวํ, ตํ ภุตฺตาวินา ปวาริเตน อนติริตฺตํ ปาตุนฺติ. นาวุโส กปฺปตีติ. กปฺปติ ภนฺเต ชโฬคิํ ปาตุนฺติ. กา สา อาวุโส ชโฬคีติ. กปฺปติ ภนฺเต ยา สา สุรา อาสุตา อสมฺปตฺตา มชฺชภาวํ สา ปาตุนฺติ. นาวุโส กปฺปตีติ. กปฺปติ ภนฺเต อทสกํ นิสีทนนฺติ. นาวุโส กปฺปตีติ. กปฺปติ ภนฺเต ชาตรูป\-รชตนฺติ. นาวุโส กปฺปตีติ. อิเม ภนฺเต เวสาลิกา วชฺชิปุตฺตกา ภิกฺขู เวสาลิยํ อิมานิ ทส วตฺถูนิ ทีเปนฺติ, หนฺท มยํ ภนฺเต อิมํ อธิกรณํ อาทิยิสฺสาม, ปุเร อธมฺโม ทิปฺปติ, ธมฺโม ปฏิพาหิยฺยติ, อวินโย ทิปฺปติ, วินโย ปฏิพาหิยฺยติ, ปุเร อธมฺมวาทิโน พลวนฺโต โหนฺติ, ธมฺมวาทิโน ทุพฺพลา โหนฺติ, อวินยวาทิโน พลวนฺโต โหนฺติ, วินยวาทิโน ทุพฺพลา โหนฺตีติ. “เอวมาวุโส”ติ โข อายสฺมา เรวโต อายสฺมโต ยสสฺส กากณฺฑก\-ปุตฺตสฺส ปจฺจสฺโสสิ.}{ปฐม\-ภาณวาโร นิฏฺฐิโต.}
\csromanmatikaanchor{mtk.219}
\csromantocmarkref{section}{2. ทุติยภาณวาร}{mtk.219}
\bookmark[dest={mtk.219},level=1]{2. ทุติยภาณวาร}
\csromanheader{1}{2. ทุติย\-ภาณวาร}
\proseitem{453}{\csromanfolio{500}อสฺโสสุํ โข เวสาลิกา วชฺชิปุตฺตกา ภิกฺขู “ยโส กิร กากณฺฑกปุตฺโต อิทํ อธิกรณํ อาทิยิตุกาโม ปกฺขํ ปริเยสติ, ลภติ จ กิร ปกฺขนฺ”ติ. อถ โข เวสาลิกานํ วชฺชิ\-ปุตฺต\-กานํ ภิกฺขูนํ เอตทโหสิ “อิทํ โข อธิกรณํ กกฺขฬญฺจ วาฬญฺจ, กํ นุ โข มยํ ปกฺขํ ลเภยฺยาม, เยน มยํ อิมสฺมิํ อธิกรเณ พลวนฺตตรา อสฺสามา”ติ.}

\prose{อถ โข เวสาลิกานํ วชฺชิ\-ปุตฺต\-กานํ ภิกฺขูนํ เอตทโหสิ “อยํ โข อายสฺมา เรวโต พหุสฺสุโต อาคตาคโม ธมฺมธโร วินยธโร มาติกาธโร ปณฺฑิโต วิยตฺโต เมธาวี ลชฺชี กุกฺกุจฺจโก สิกฺขากาโม, สเจ มยํ อายสฺมนฺตํ เรวตํ ปกฺขํ ลเภยฺยาม, เอวํ มยํ อิมสฺมิํ อธิกรเณ พลวนฺตตรา อสฺสามา”ติ.}

\prose{อถ โข เวสาลิกา วชฺชิปุตฺตก ภิกฺขู ปหูตํ สามณกํ ปริกฺขารํ ปฏิยาเทสุํ ปตฺตมฺปิ จีวรมฺปิ นิสีทนมฺปิ สูจิฆรมฺปิ กาย\-พนฺธนมฺปิ ปริสฺสาวนมฺปิ ธมฺมกรณมฺปิ, อถ โข เวสาลิกา วชฺชิปุตฺตกา ภิกฺขู ตํ สามณกํ ปริกฺขารํ อาทาย นาวาย สหชาติํ อุชฺชวิํสุ, นาวาย ปจฺโจโรหิตฺวา อญฺญตรสฺมิํ รุกฺขมูเล ภตฺตวิสฺสคฺคํ กโรนฺติ. อถ โข อายสฺมาโต สาฬฺหสฺส รโหคตสฺส ปฏิสลฺลีนสฺส เอวํ เจตโส ปริวิตกฺโก อุทปาทิ “เก นุ โข ธมฺมวาทิโน ปาจีนกา วา ภิกฺขู, ปาเวยฺยกา วา”ติ. อถ โข อายสฺมโต สาฬฺหสฺส ธมฺมญฺจ วินยญฺจ เจตสา ปจฺจเวกฺขนฺตสฺส เอตทโหสิ อธมฺมวาทิโน ปาจีนกา ภิกฺขู, ธมฺมวาทิโน ปาเวยฺยกา\csromansharedfootnote{e9a303c8f5c6b467}{ปาเฐยฺยกา (สฺยา)} ภิกฺขู”ติ.}

\prose{อถ โข อญฺญตรา สุทฺธา\-วาส\-กายิกา เทวตา อายสฺมโต สาฬฺหสฺส เจตสา เจโต\-ปริวิตกฺกม\-ญฺญาย เสยฺยถาปิ นาม พลวาปุริโส สมิญฺชิตํ วา พาหํ ปสาเรยฺย ปสาริตํ วา พาหํ สมิญฺเชยฺย เอวเมว สุทฺธาวาเสสุ เทเวสุ อนฺตรหิตา อายสฺมโต สาฬฺหสฺส สมฺมุเข ปาตุรโหสิ. อถ โข สา เทวตา อายสฺมนฺตํ สาฬฺหํ เอตทโวจ “สาธุ ภนฺเต สาฬฺห, อธมฺมวาที ปาจีนกา ภิกฺขู, ธมฺมวาที ปาเวยฺยกา ภิกฺขู, เตน หิ ภนฺเต สาฬฺห ยถาธมฺโม ตถา ติฏฺฐาหี”ติ. ปุพฺเพปิ จาหํ เทวเต เอตรหิ จ ยถาธมฺโม ตถา \csromanfolio{501}ฐิโต, อปิ จาหํ น ตาว ทิฏฺฐิํ อาวิ กโรมิ “อปฺเปว นาม มํ อิมสฺมิํ อธิกรเณ สมฺมนฺเนยฺยา”ติ.}

\proseitem{454}{อถ โข เวสาลิกา วชฺชิปุตฺตกา ภิกฺขู ตํ สามณกํ ปริกฺขารํ อาทาย เยนายสฺมา เรวโต เตนุ\-ปสงฺกมิํสุ, อุปสงฺกมิตฺวา อายสฺมนฺตํ เรวตํ เอตทโวจุํ “ปฏิคฺคณฺหาตุ ภนฺเต เถโร สามณกํ ปริกฺขารํ ปตฺตมฺปิ จีวรมฺปิ นิสีทนมฺปิ สูจิฆรมฺปิ กาย\-พนฺธนมฺปิ ปริสฺสาวนมฺปิ ธมฺมกรณมฺปี”ติ. “อลํ อาวุโส, ปริปุณฺณํ เม ปตฺตจีวรนฺ”ติ น อิจฺฉิ ปฏิคฺคเหตุํ.}

\prose{เตน โข ปน สมเยน อุตฺตโร นาม ภิกฺขุ วีสติวสฺโส อายสฺมโต เรวตสฺส อุปฏฺฐาโก โหติ. อถ โข เวสาลิกา วชฺชิปุตฺตกา ภิกฺขู เยนายสฺมา อุตฺตโร เตนุ\-ปสงฺกมิํสุ, อุปสงฺกมิตฺวา อายสฺมนฺตํ อุตฺตรํ เอตทโวจุํ “ปฏิคฺคณฺหาตุ อายสฺมา อุตฺตโร สามณกํ ปริกฺขารํ ปตฺตมฺปิ จีวรมฺปิ นิสีทนมฺปิ สูจิฆรมฺปิ กาย\-พนฺธนมฺปิ ปริสฺสาวนมฺปิ ธมฺมกรณมฺปี”ติ. “อลํ อาวุโส, ปริปุณฺณํ เม ปตฺตจีวรนฺ”ติ น อิจฺฉิ ปฏิคฺคเหตุํ. มนุสฺสา โข อาวุโส อุตฺตร ภควโต สามณกํ ปริกฺขารํ อุปนาเมนฺติ, สเจ ภควา ปฏิคฺคณฺหาติ, เตเนว อตฺตมนา โหนฺติ. โน เจ ภควา ปฏิคฺคณฺหาติ. อายสฺมโต\csromansharedfootnote{c95e26aeef92043a}{อายสฺมโต จ (สฺยา)} อานนฺทสฺส อุปนาเมนฺติ “ปฏิคฺคณฺหาตุ ภนฺเต เถโร สามณกํ ปริกฺขารํ, ยถา ภควตา ปฏิคฺคหิโต, เอวเมว โส ภวิสฺสตี”ติ. ปฏิคฺคณฺหาตุ อายสฺมา อุตฺตโร สามณกํ ปริกฺขารํ, ยถา เถเรน ปฏิคฺคหิโต, เอวเมว โส ภวิสฺสตีติ. อถ โข อายสฺมา อุตฺตโร เวสาลิเกหิ วชฺชิปุตฺเตหิ ภิกฺขูหิ นิปฺปีฬิยมาโน เอกํ จีวรํ อคฺคเหสิ. วเทยฺยาถ อาวุโส เยน อตฺโถติ. เอตฺตกํ อายสฺมา อุตฺตโร เถรํ วเทตุ “เอตฺตกญฺจ ภนฺเต เถโร สํฆมชฺเฌ วเทตุ, ปุรตฺถิเมสุ ชนปเทสุ พุทฺธา ภควนฺโต อุปฺปชฺชนฺติ. ธมฺมวาที ปาจีนกา ภิกฺขู, อธมฺมวาที ปาเวยฺยกา ภิกฺขู”ติ. “เอวมาวุโส”ติ โข อายสฺมา อุตฺตโร เวสาลิกานํ วชฺชิ\-ปุตฺต\-กานํ ภิกฺขูนํ ปฏิสฺสุตฺวา เยนายสฺมา เรวโต เตนุปสงฺกมิ, อุปสงฺกมิตฺวา อายสฺมนฺตํ เรวตํ เอตทโวจ “เอตฺตกํ ภนฺเต เถโร สํฆมชฺเฌ วเทตุ, ปุรตฺถิเมสุ ชนปเทสุ พุทฺธา ภควนฺโต \csromanfolio{502}อุปฺปชฺชนฺติ. ธมฺมวาที ปาจีนกา ภิกฺขู, อธมฺมวาที ปาเวยฺยกา ภิกฺขู”ติ. “อธมฺเม มํ ตฺวํ ภิกฺขุ นิโยเชสี”ติ เถโร อายสฺมนฺตํ อุตฺตรํ ปณาเมสิ.}

\prose{อถ โข เวสาลิกา วชฺชิปุตฺตกา ภิกฺขู อายสฺมนฺตํ อุตฺตรํ เอตทโวจุํ “กิํ อาวุโส อุตฺตร เถโร อาหา”ติ. ปาปิกํ โน อาวุโส กตํ, “อธมฺเม มํ ตฺวํ ภิกฺขุ นิโยเชสี”ติ เถโร มํ ปณาเมสีติ. นนุ ตฺวํ อาวุโส\csromansharedfootnote{e59bc760efadd6a8}{อาวุโส อุตฺตร (สฺยา, กํ)} วุฑฺโฒ วีสติวสฺโสสีติ. อามาวุโส, อปิ จ มยํ ครุนิสฺสยํ คณฺหามาติ.}

\proseitem{455}{อถ โข สํโฆ ตํ อธิกรณํ วินิจฺฉินิตุ\-กาโม สนฺนิปติ, อถ โข อายสฺมา เรวโต สํฆํ ญาเปสิ\csromandash{}}

\prose{“สุณาตุ เม อาวุโส สํโฆ, สเจ มยํ อิมํ อธิกรณํ อิธ วูปสเมสฺสาม, สิยาปิ มูลาทายกา\csromansharedfootnote{e302f91c9fd3930e}{มูลทายกา (สี)} ภิกฺขู ปุนกมฺมาย อุกฺโกเฏยฺยุํ, ยทิ สํฆสฺส ปตฺตกลฺลํ, ยตฺเถวิมํ อธิกรณํ สมุปฺปนฺนํ, สํโฆ ตตฺเถวิมํ อธิกรณํ วูปสเมยฺยา”ติ.}

\prose{อถ โข เถรา ภิกฺขู เวสาลิํ อคมํสุ ตํ อธิกรณํ วินิจฺฉินิตุ\-กามา.}

\prose{เตน โข ปน สมเยน สพฺพกามี นาม ปถพฺยา สํฆตฺเถโร วีสวสฺส\-สติโก อุปสมฺปทาย อายสฺมโต อานนฺทสฺส สทฺธิวิหาริโก เวสาลิยํ ปฏิวสติ. อถ โข อายสฺมา เรวโต อายสฺมนฺตํ สมฺภูตํ สาณวาสิํ เอตทโวจ “อหํ อาวุโส ยสฺมิํ วิหาเร สพฺพกามี เถโร วิหรติ, ตํ วิหารํ อุปคจฺฉามิ, โส ตฺวํ กาลสฺเสว อายสฺมนฺตํ สพฺพกามิํ อุปสงฺกมิตฺวา อิมานิ ทส วตฺถูนิ ปุจฺเฉยฺยาสี”ติ.}

\prose{“เอวํ ภนฺเต”ติ โข อายสฺมา สมฺภูโต สาณวาสี อายสฺมโต เรวตสฺส ปจฺจสฺโสสิ. อถ โข อายสฺมา เรวโต ยสฺมิํ วิหาเร สพฺพกามี เถโร วิหรติ ตํ วิหารํ อุปคจฺฉิ. คพฺเภ อายสฺมโต สพฺพกามิสฺส เสนาสนํ ปญฺญตฺตํ โหติ, คพฺภปฺปมุเข อายสฺมโต เรวตสฺส. อถ โข อายสฺมา เรวโต “อยํ เถโร มหลฺลโก, น นิปชฺชตี”ติ น เสยฺยํ กปฺเปสิ, อายสฺมา สพฺพกามี “อยํ ภิกฺขุ อาคนฺตุโก กิลนฺโต น นิปชฺชตี”ติ น เสยฺยํ กปฺเปสิ. อถ โข \csromanfolio{503}อายสฺมา สพฺพกามี รตฺติยา ปจฺจูสสมยํ ปจฺจุฏฺฐาย อายสฺมนฺตํ เรวตํ เอตทโวจ “กตเมน ตฺวํ ภูมิ วิหาเรน เอตรหิ พหุลํ วิหรสี”ติ. เมตฺตาวิหาเรน โข อหํ ภนฺเต เอตรหิ พหุลํ วิหรามีติ. กุลฺลกวิหาเรน กิร ตฺวํ ภูมิ เอตรหิ พหุลํ วิหรสิ, กุลฺลกวิหาโร เอโส\csromansharedfootnote{052066c943deea6c}{เหโส (สฺยา)} ภูมิ, ยทิทํ เมตฺตาติ. ปุพฺเพปิ เม ภนฺเต คิหิภูตสฺส อาจิณฺณา เมตฺตา, เตนาหํ เอตรหิปิ เมตฺตาวิหาเรน พหุลํ วิหรามิ, อปิ จ มยา จิรปฺปตฺตํ อรหตฺตนฺติ. เถโร ปน ภนฺเต กตเมน วิหาเรน เอตรหิ พหุลํ วิหรตีติ. สุญฺญตา\-วิหาเรน โข อหํ ภูมิ เอตรหิ พหุลํ วิหรามีติ. มหาปุริส\-วิหาเรน กิร ภนฺเต เถโร เอตรหิ พหุลํ วิหรติ, มหาปุริส\-วิหาโร เอโส ภนฺเต, ยทิทํ สุญฺญตาติ. ปุพฺเพปิ เม ภูมิ คิหิภูตสฺส อาจิณฺณา สุญฺญตา, เตนาหํ เอตรหิปิ สุญฺญาตาวิหาเรน พหุลํ วิหรามิ, อปิ จ มยา จิรปฺปตฺตํ อรหตฺตนฺติ. อยญฺจรหิ เถเรนํ ภิกฺขูนํ อนฺตรากถา วิปฺปกตา, อถายสฺมา สมฺภูโต สาณวาสี ตสฺมิํ อนุปฺปตฺโต โหติ. อถ โข อายสฺมา สมฺภูโต สาณวาสี เยนายสฺมา สพฺพกามี เตนุปสงฺกมิ, อุปสงฺกมิตฺวา อายสฺมนฺตํ สพฺพกามิํ อภิวาเทตฺวา เอกมนฺตํ นิสีทิ, เอกมนฺตํ นิสินฺโน โข อายสฺมา สมฺภูโต สาณวาสี อายสฺมนฺตํ สพฺพกามิํ เอตทโวจ “อิเม ภนฺเต เวสาลิกา วชฺชิปุตฺตกา ภิกฺขู เวสาลิยํ ทส วตฺถูนิ ทีเปนฺติ ‘กปฺปติ สิงฺคิโลณกปฺโป, กปฺปติ ทฺวงฺคุลกปฺโป, กปฺปติ คามนฺตรกปฺโป, กปฺปติ อาวาสกปฺโป, กปฺปติ อนุมติกปฺโป, กปฺปติ อาจิณฺณกปฺโป, กปฺปติ อมถิตกปฺโป, กปฺปติ ชโฬคิํ ปาตุํ, กปฺปติ อทสกํ นิสีทนํ, กปฺปติ ชาตรูปรชตนฺ”ติ. เถเรน ภนฺเต อุปชฺฌายสฺส มูเล พหุ ธมฺโม จ วินโย จ ปริยตฺโต, เถเรสฺส ภนฺเต ธมฺมญฺจ วินยญฺจ ปจฺจเวกฺขนฺตสฺส กถํ โหติ “เก นุ โข ธมฺมวาทิโน, ปาจีนกา วา ภิกฺขู, ปาเวยฺยกา วา”ติ. ตยาปิ โข อาวุโส อุปชฺฌายสฺส มูเล พหุ ธมฺโม จ วินโย จ ปริยตฺโต, ตุยฺหํ ปน อาวุโส ธมฺมญฺจ วินยญฺจ ปจฺจเวกฺขนฺตสฺส กถํ โหติ “เก นุ โข ธมฺมวาทิโน ปาจีนกา วา ภิกฺขู, ปาเวยฺยกา วา”ติ. มยฺหํ โข ภนฺเต ธมฺมญฺจ วินยญฺจ ปจฺจเวกฺขนฺตสฺส เอวํ โหติ “อธมฺมวาที ปาจีนกา ภิกฺขู, ธมฺมวาที ปาเวยฺยกา ภิกฺขู”ติ. อปิ จาหํ \csromanfolio{504}น ตาว ทิฏฺฐิํ อาวิ กโรมิ “อปฺเปว นาม มํ อิมสฺมิํ อธิกรเณ สมฺมนฺเนยฺยา”ติ. มยฺหมฺปิ โข อาวุโส ธมฺมญฺจ วินยญฺจ ปจฺจเวกฺขนฺตสฺส เอวํ โหติ “อธมฺมวาที ปาจีนกา ภิกฺขู, ธมฺมวาที ปาเวยฺยกา ภิกฺขู”ติ. อปิ จาหํ น ตาว ทิฏฺฐิํ อาวิ กโรมิ “อปฺเปว นาม มํ อิมสฺมิํ อธิกรเณ สมฺมนฺเนยฺยา”ติ.}

\proseitem{456}{อถ โข สํโฆ ตํ อธิกรณํ วินิจฺฉินิตุ\-กาโม สนฺนิปติ, ตสฺมิํ โข ปน อธิกรเณ วินิจฺฉิยมาเน อนคฺคานิ เจว ภสฺสานิ ชายนฺติ, น เจกสฺส ภาสิตสฺส อตฺโถ วิญฺญายติ, อถ โข อายสฺมา เรวโต สํฆํ ญาเปสิ\csromandash{}}

\prose{“สุณาตุ เม ภนฺเต สํโฆ, อมฺหากํ อิมสฺมิํ อธิกรเณ วินิจฺฉิยมาเน อนคฺคานิ เจว ภสฺสานิ ชายนฺติ, น เจกสฺส ภาสิตสฺส อตฺโถ วิญฺญายติ, ยทิ สํฆสฺส ปตฺตกลฺลํ, สํโฆ อิมํ อธิกรณํ อุพฺพาหิกาย วูปสเมยฺยา”ติ. สํโฆ จตฺตาโร ปาจีนเก ภิกฺขู จตฺตาโร ปาเวยฺยเก ภิกฺขู อุจฺจินิ, “ปาจีนกานํ ภิกฺขูนํ อายสฺมนฺตญฺจ สพฺพกามิํ อายสฺมนฺตญฺจ สาฬฺหํ อายสฺมนฺตญฺจ ขุชฺชโสภิตํ อายสฺมนฺตญฺจ วาสภคามิกํ, ปาเวยฺยกานํ ภิกฺขูนํ อายสฺมนฺตญฺจ เรวตํ อายสฺมนฺตญฺจ สมฺภูตํ สาณวาสิํ อายสฺมนฺตญฺจ ยสํ กากณฺฑก\-ปุตฺตํ อายสฺมนฺตญฺจ สุมนนฺ”ติ. อถ โข อายสฺมา เรวโต สํฆํ ญาเปสิ\csromandash{}}

\prose{“สุณาตุ เม ภนฺเต สํโฆ, อมฺหากํ อิมสฺมิํ อธิกรเณ วินิจฺฉิยมาเน อนคฺคานิ เจว ภสฺสานิ ชายนฺติ. น เจกสฺส ภาสิตสฺส อตฺโถ วิญฺญายติ, ยทิ สํฆสฺส ปตฺตกลฺลํ, สํโฆ จตฺตาโร ปาจีนเก ภิกฺขู จตฺตาโร ปาเวยฺยเก ภิกฺขู สมฺมนฺเนยฺย อุพฺพาหิกาย อิมํ อธิกรณํ วูปสเมตุํ, เอสา ญตฺติ.}

\prose{สุณาตุ เม ภนฺเต สํโฆ, อมฺหากํ อิมสฺมิํ อธิกรเณ วินิจฺฉิยมาเน อนคฺคานิ เจว ภสฺสานิ ชายนฺติ, น เจกสฺส ภาสิตสฺส อตฺโถ วิญฺญายติ, สํโฆ จตฺตาโร ปาจีนเก ภิกฺขู จตฺตาโร ปาเวยฺยเก ภิกฺขู สมฺมนฺนติ อุพฺพาหิกาย อิมํ อธิกรณํ วูปสเมตุํ, ยสฺสายสฺมโต ขมติ จตุนฺนํ ปาจีนกานํ ภิกฺขูนํ จตุนฺนํ ปาเวยฺยกานํ ภิกฺขูนํ สมฺมุติ อุพฺพาหิกาย อิมํ อธิกรณํ วูปสเมตุํ, โส ตุณฺหสฺส, ยสฺส นกฺขมติ, โส ภาเสยฺย.}

\prose{\csromanfolio{505}สมฺมตา สํเฆน จตฺตาโร ปาจีนกา ภิกฺขู จตฺตาโร ปาเวยฺยกา ภิกฺขู อุพฺพาหิกาย อิมํ อธิกรณํ วูปสเมตุํ, ขมติ สํฆสฺส, ตสฺมา ตุณฺหี, เอวเมตํ ธารยามี”ติ.}

\prose{เตน โข ปน สมเยน อชิโต นาม ภิกฺขุ ทสวสฺโส สํฆสฺส ปาติโมกฺขุ\-ทฺเทสโก โหติ. อถ โข สํโฆ อายสฺมนฺตมฺปิ อชิตํ สมฺมนฺนติ เถรานํ ภิกฺขูนํ อาสน\-ปญฺญาปกํ. อถ โข เถรานํ ภิกฺขูนํ เอตทโหสิ “กตฺถ นุ โข มยํ อิมํ อธิกรณํ วูปสเมยฺยามา”ติ. อถ โข เถรานํ ภิกฺขูนํ เอตทโหสิ “อยํ โข วาลิการาโม รมณีโย อปฺปสทฺโท อปฺปนิคฺโฆโส, ยํนูน มยํ วาลิการาเม อิมํ อธิกรณํ วูปสเมยฺยามา”ติ.}

\proseitem{457}{อถ โข เถรา ภิกฺขู วาลิการามํ อคมํสุ ตํ อธิกรณํ วินิจฺฉินิตุ\-กามา. อถ โข อายสฺมา เรวโต สํฆํ ญาเปสิ\csromandash{}}

\prose{“สุณาตุ เม ภนฺเต สํโฆ, ยทิ สํฆสฺส ปตฺตกลฺลํ, อหํ อายสฺมนฺตํ สพฺพกามิํ วินยํ ปุจฺเฉยฺยนฺ”ติ.}

\prose{อายสฺมา สพฺพกามี สํฆํ ญาเปสิ\csromandash{}}

\prose{“สุณาตุ เม อาวุโส สํโฆ, ยทิ สํฆสฺส ปตฺตกลฺลํ, อหํ เรวเตน วินยํ ปุฏฺโฐ วิสฺสชฺเชยฺยนฺ”ติ.}

\prose{อถ โข อายสฺมา เรวโต อายสฺมนฺตํ สพฺพกามิํ เอตทโวจ “กปฺปติ ภนฺเต สิงฺคิโลณกปฺโป”ติ. โก โส อาวุโส สิงฺคิ\-โลณ\-กปฺโปติ. กปฺปติ ภนฺเต สิงฺคินา โลณํ ปริหริตุํ “ยตฺถ อโลณกํ ภวิสฺสติ, ตตฺถ ปริภุญฺชิสฺสามา”ติ. นาวุโส กปฺปตีติ.\csromansymbolfootnote{*}{วิ 2. 117 ปิฏฺเฐ.}กตฺถ ปฏิกฺขิตฺตนฺติ. สาวตฺถิยํ สุตฺต\-วิภงฺเคติ. กิํ อาปชฺชตีติ. สนฺนิธิการก\-โภชเน ปาจิตฺติยนฺติ.}

\prose{สุณาตุ เม ภนฺเต สํโฆ, อิทํ ปฐมํ วตฺถุ สํเฆน วินิจฺฉิตํ, อิติปิทํ วตฺถุ อุทฺธมฺมํ อุพฺพินยํ อปคต\-สตฺถุสาสนํ, อิทํ ปฐมํ สลากํ นิกฺขิปามิ. (1)}

\prose{กปฺปติ ภนฺเต ทฺวงฺคุล\-กปฺโปติ. โก โส อาวุโส ทฺวงฺคุล\-กปฺโปติ. กปฺปติ ภนฺเต ทฺวงฺคุลาย ฉายาย วีติวตฺตาย วิกาเล โภชนํ \csromanfolio{506}ภุญฺชิตุนฺติ. นาวุโส กปฺปตีติ.\csromansharedfootnote{ea420bf0d9607ee7}{วิ 2. 114 ปิฏฺเฐ.} กตฺถ ปฏิกฺขิตฺตนฺติ. ราชคเห สุตฺต\-วิภงฺเคติ. กิํ อาปชฺชตีติ. วิกาลโภชเน ปาจิตฺติยนฺติ.}

\prose{สุณาตุ เม ภนฺเต สํโฆ, อิทํ ทุติยํ วตฺถุ สํเฆน วินิจฺฉิตํ, อิติปิทํ วตฺถุ อุทฺธมฺมํ อุพฺพินยํ อปคต\-สตฺถุสาสนํ, อิทํ ทุติยํ สลากํ นิกฺขิปามิ. (2)}

\prose{กปฺปติ ภนฺเต คามนฺตร\-กปฺโปติ. โก โส อาวุโส คามนฺตร\-กปฺโปติ. กปฺปติ ภนฺเต “อิทานิ คามนฺตรํ คมิสฺสามี”ติ ภุตฺตาวินา ปวาริเตน อนติริตฺตํ โภชนํ ภุญฺชิตุนฺติ. นาวุโส กปฺปตีติ.\csromansharedfootnote{fcf7bcbeef358a23}{วิ 2. 109 ปิฏฺเฐ.} กตฺถ ปฏิกฺขิตฺตนฺติ. สาวตฺถิยํ สุตฺต\-วิภงฺเคติ. กิํ อาปชฺชตีติ. อนติริตฺต\-โภชเน ปาจิตฺติยนฺติ.}

\prose{สุณาตุ เม ภนฺเต สํโฆ, อิทํ ตติยํ วตฺถุ สํเฆน วินิจฺฉิตํ, อิติปิทํ วตฺถุ อุทฺธมฺมํ อุพฺพินยํ อปคต\-สตฺถุสาสนํ, อิทํ ตติยํ สลากํ นิกฺขิปามิ. (3)}

\prose{กปฺปติ ภนฺเต อาวาสกปฺโปติ. โก โส อาวุโส อาวาสกปฺโปติ. กปฺปติ ภนฺเต สมฺพหุลา อาวาสา สมานสีมานานุโปสถํ กาตุนฺติ. นาวุโส กปฺปตีติ.\csromansharedfootnote{ac397c038937add8}{วิ 3. 146 ปิฏฺเฐ.} กตฺถ ปฏิกฺขิตฺตนฺติ. ราชคเห อุโปสถ\-สํยุตฺเตติ. กิํ อาปชฺชตีติ. วินยาติสาเร ทุกฺกฏนฺติ.}

\prose{สุณาตุ เม ภนฺเต สํโฆ, อิทํ จตุตฺถํ วตฺถุ สํเฆน วินิจฺฉิตํ, อิติปิทํ วตฺถุ อุทฺธมฺมํ อุพฺพินยํ อปคต\-สตฺถุสาสนํ, อิทํ จตุตฺถํ สลากํ นิกฺขิปามิ. (4)}

\prose{กปฺปติ ภนฺเต อนุมติกปฺโปติ. โก โส อาวุโส อนุมติกปฺโปติ. กปฺปติ ภนฺเต วคฺเคน สํเฆน กมฺมํ กาตุํ “อาคเต ภิกฺขู อนุมาเนสฺสามา”ติ. นาวุโส กปฺปตีติ.\csromansharedfootnote{517feb6eeb683642}{วิ 3. 435 ปิฏฺเฐ.} กตฺถ ปฏิกฺขิตฺตนฺติ. จมฺเปยฺยเก วินย\-วตฺถุสฺมินฺติ. กิํ อาปชฺชตีติ. วินยาติสาเร ทุกฺกฏนฺติ.}

\prose{สุณาตุ เม ภนฺเต สํโฆ, อิทํ ปญฺจมํ วตฺถุ สํเฆน วินิจฺฉิตํ, อิติปิทํ วตฺถุ อุทฺธมฺมํ อุพฺพินยํ อปคต\-สตฺถุสาสนํ, อิทํ ปญฺจมํ สลากํ นิกฺขิปามิ. (5)}

\prose{กปฺปติ ภนฺเต อาจิณฺณกปฺโปติ. โก โส อาวุโส อาจิณฺณกปฺโปติ. กปฺปติ ภนฺเต อิทํ เม อุปชฺฌาเยน อชฺฌาจิณฺณํ, อิทํ เม \csromanfolio{507}อาจริเยน อชฺฌาจิณฺณํ, ตํ อชฺฌาจริตุนฺติ. อาจิณฺณกปฺโป โข อาวุโส เอกจฺโจ กปฺปติ, เอกจฺโจ น กปฺปตีติ.}

\prose{สุณาตุ เม ภนฺเต สํโฆ, อิทํ ฉฏฺฐํ วตฺถุ สํเฆน วินิจฺฉิตํ, อิติปิทํ วตฺถุ อุทฺธมฺมํ อุพฺพินยํ อปคต\-สตฺถุสาสนํ, อิทํ ฉฏฺฐํ สลากํ นิกฺขิปามิ. (6)}

\prose{กปฺปติ ภนฺเต อมถิต\-กปฺโปติ. โก โส อาวุโส อมถิต\-กปฺโปติ. กปฺปติ ภนฺเต ยํ ตํ ขีรํ ขีรภาวํ วิชหิตํ อสมฺปตฺตํ ทธิภาวํ, ตํ ภุตฺตาวินา ปวาริเตน อนติริตฺตํ ปาตุนฺติ. นาวุโส กปฺปตีติ.\csromansharedfootnote{fcf7bcbeef358a23}{วิ 2. 109 ปิฏฺเฐ.} กตฺถ ปฏิกฺขิตฺตนฺติ. สาวตฺถิยํ สุตฺต\-วิภงฺเคติ. กิํ อาปชฺชตีติ. อนติริตฺต\-โภชเน ปาจิตฺติยนฺติ.}

\prose{สุณาตุ เม ภนฺเต สํโฆ, อิทํ สตฺตมํ วตฺถุ สํเฆน วินิจฺฉิตํ, อิติปิทํ วตฺถุ อุทฺธมฺมํ อุพฺพินยํ อปคต\-สตฺถุสาสนํ, อิทํ สตฺตมํ สลากํ นิกฺขิปามิ. (7)}

\prose{กปฺปติ ภนฺเต ชโฬคิํ ปาตุนฺติ. กา สา อาวุโส ชโฬคีติ. กปฺปติ ภนฺเต ยา สา สุรา อาสุตา อสมฺปตฺตา มชฺชภาวํ สา ปาตุนฺติ. นาวุโส กปฺปตีติ.\csromansharedfootnote{e2ed3392fadb8144}{วิ 2. 145 ปิฏฺเฐ. 3. วิ 2. 221 ปิฏฺเฐ.} กตฺถ ปฏิกฺขิตฺตนฺติ. โกสมฺพิยํ สุตฺต\-วิภงฺเคติ. กิํ อาปชฺชตีติ. สุราเมรยปาเน ปาจิตฺติยนฺติ.}

\prose{สุณาตุ เม ภนฺเต สํโฆ, อิทํ อฏฺฐมํ วตฺถุ สํเฆน วินิจฺฉิตํ, อิติปิทํ วตฺถุ อุทฺธมฺมํ อุพฺพินยํ อปคต\-สตฺถุสาสนํ, อิทํ อฏฺฐมํ สลากํ นิกฺขิปามิ. (8)}

\prose{กปฺปติ ภนฺเต อทสกํ นิสีทนนฺติ. นาวุโส กปฺปตีติ.\csromansharedfootnote{47654c568798548c}{วิ 2. 221 ปิฏฺเฐ.} กตฺถ ปฏิกฺขิตฺตนฺติ. สาวตฺถิยํ สุตฺต\-วิภงฺเคติ. กิํ อาปชฺชตีติ. เฉทนเก ปาจิตฺติยนฺติ.}

\prose{สุณาตุ เม ภนฺเต สํโฆ, อิทํ นวมํ วตฺถุ สํเฆน วินิจฺฉิตํ, อิติปิทํ วตฺถุ อุทฺธมฺพํ อุพฺพินยํ อปคต\-สตฺถุสาสนํ, อิทํ นวมํ สลากํ นิกฺขิปามิ. (9)}

\prose{กปฺปติ ภนฺเต ชาตรูป\-รชตนฺติ. นาวุโส กปฺปตีติ.\csromansharedfootnote{b6d972cdbee25684}{วิ 1. 344 ปิฏฺเฐ.} กตฺถ ปฏิกฺขิตฺตนฺติ. ราชคเห สุตฺต\-วิภงฺเคติ. กิํ อาปชฺชตีติ. ชาตรูป\-รชต\-ปฏิคฺคหเณ ปาจิตฺติยนฺติ.}

\csromanmatikaanchor{mtk.220}
\csromantocmarkref{section}{อุทฺทานคาถา}{mtk.220}
\bookmark[dest={mtk.220},level=1]{อุทฺทานคาถา}
\prose{\csromanfolio{508}สุณาตุ เม ภนฺเต สํโฆ, อิทํ ทสมํ วตฺถุ สํเฆน วินิจฺฉิตํ, อิติปิทํ วตฺถุ อุทฺธมฺมํ อุพฺพินยํ อปคต\-สตฺถุสาสนํ, อิทํ ทสมํ สลากํ นิกฺขิปามิ. (10)}

\prose{สุณาตุ เม ภนฺเต สํโฆ, อิมานิ ทส วตฺถูนิ สํเฆน วินิจฺฉิตานิ “อิติปิมานิ ทสวตฺถูนิ อุทฺธมฺมานิ อุพฺพินยานิ อปคตสตฺถุสาสนานี”ติ.}

\proseitem{458}{นิหตเมตํ อาวุโส อธิกรณํ สนฺตํ วูปสนฺตํ สุวูปสนฺตํ, อปิ จ มํ ตฺวํ อาวุโส สํฆมชฺเฌปิ อิมานิ ทส วตฺถูนิ ปุจฺเฉยฺยาสิ เตสํ ภิกฺขูนํ สญฺญตฺติยาติ. อถ โข อายสฺมา เรวโต อายสฺมนฺตํ สพฺพกามิํ สํฆมชฺเฌปิ อิมานิ ทส วตฺถูนิ ปุจฺฉิ, ปุฏฺโฐ ปุฏฺโฐ อายสฺมา สพฺพกามี วิสฺสชฺเชสิ. อิมายา โข ปน วินย\-สงฺคีติยา สตฺต ภิกฺขุสตานิ อนูนานิ อนธิกานิ อเหสุํ, ตสฺมายํ วินยสงฺคีติ “สตฺตสติกา”ติ วุจฺจตีติ. สตฺตสติก\-กฺขนฺธโก ทฺวาทสโม. อิมมฺหิ ขนฺธเก วตฺถู ปญฺจวีสติ.}

\vspace{\tipitakaparskip}
\csromancenter{\textbf{ตสฺสุทฺทานํ}}
\setlength{\csromangathapagewidth}{0pt}
\setlength{\csromangathawakwidth}{0pt}
\csromangathameasuregroup{ทส วตฺถูนิ ปูเรตฺวา, \\ จตฺตาโร ปุน รูปญฺจ, \\ มคฺโค โสเรยฺยํ สงฺกสฺสํ, \\ สหชาติ จ มชฺเฌสิ, \\ ปตฺตนาวาย อุชฺชวิ,}{\csromangathabat{ทส วตฺถูนิ ปูเรตฺวา,}{กมฺมํ ทูเตน ปาวิสิ.} \\ \csromangathabat{จตฺตาโร ปุน รูปญฺจ,}{โกสมฺพิ จ ปาเวยฺยโก.} \\ \csromangathabat{มคฺโค โสเรยฺยํ สงฺกสฺสํ,}{กณฺณกุชฺชํ อุทุมฺพรํ.} \\ \csromangathabat{สหชาติ จ มชฺเฌสิ,}{อสฺโสสิ กํ นุ โข มยํ.} \\ \csromangathabat{ปตฺตนาวาย อุชฺชวิ,}{รโหสิ อุปนามยํ\textsuperscript{1}.}}
\csromangathasetleft{ทส วตฺถูนิ ปูเรตฺวา, \\ จตฺตาโร ปุน รูปญฺจ, \\ มคฺโค โสเรยฺยํ สงฺกสฺสํ, \\ สหชาติ จ มชฺเฌสิ, \\ ปตฺตนาวาย อุชฺชวิ,}
\csromangathagroup{\csromangathastanza{\csromangathabat{ทส วตฺถูนิ ปูเรตฺวา,}{กมฺมํ ทูเตน ปาวิสิ.} \\ \csromangathabat{จตฺตาโร ปุน รูปญฺจ,}{โกสมฺพิ จ ปาเวยฺยโก.}}\csromangathastanza{\csromangathabat{มคฺโค โสเรยฺยํ สงฺกสฺสํ,}{กณฺณกุชฺชํ อุทุมฺพรํ.} \\ \csromangathabat{สหชาติ จ มชฺเฌสิ,}{อสฺโสสิ กํ นุ โข มยํ.} \\ \csromangathabat{ปตฺตนาวาย อุชฺชวิ,}{รโหสิ อุปนามยํ\csromansharedfootnote{3228aa94ea3dd4f1}{ทูรโตปิ อุทปาทิ (ก)}.}}}

\vspace{\tipitakaparskip}
\csromancenter{ครุ\csromansharedfootnote{acbe9cb4c910c472}{ทารุณํ (สฺยา)} สํโฆ จ เวสาลิํ, เมตฺตา สํโฆ อุพฺพาหิกาติ.}
\nitthitammajor{สตฺตสติก\-กฺขนฺธโก นิฏฺฐิโต.}
\nitthitammid{จูฬวคฺโค\csromansharedfootnote{61785d03c494c2ae}{[มิสฺสินฺคฺ โนเต 1]} นิฏฺฐิโต.}
\nitthitammajor{\textbf{จูฬวคฺคปาฬิ นิฏฺฐิตา.}}
\orphannote{มูลาย\-วิสุทฺธินวก (สี, สฺยา)}

\orphannote{อิธ ปน ภิกฺขเว ภิกฺขุ สมฺพหุลา สํฆาทิเสสา อาปตฺติโย อาปชฺชติ ปริมาณมฺปิ อปริมาณมฺปิ ฯเปฯ ววตฺถิตมฺปิ สมฺภินฺนมฺปิ. โส สํฆํ ตาสํ อาปตฺตีนํ สโมธาน\-ปริวาสํ ยาจติ, ตสฺส สํโฆ ตาสํ อาปตฺตีนํ สโมธาน\-ปริวาสํ เทติ, โส ปริวสนฺโต อนฺตรา สมฺพหุลา สํฆาทิเสสา อาปตฺติโย อาปชฺชติ ปริมาณาโย อปฺปฏิจฺฉนฺนาโย ฯเปฯ ปริมาณาโย ปฏิจฺฉนฺนาโย ฯเปฯ ปริมาณาโย}

\orphannote{มูลาย\-วิสุทฺธินวก (สี, สฺยา) ปฏิจฺฉนฺนาโยปิ อปฺปฏิจฺฉนฺนาโยปิ. โส สํฆํ อนฺตรา\-อาปตฺตีนํ มูลาย\-ปฏิกสฺสนํ ยาจติ, ตํ สํโฆ อนฺตรา\-อาปตฺตีนํ มูลาย ปฏิกสฺสติ ธมฺมิเกน กมฺเมน อกุปฺเปน ฐานารเหน, ธมฺเมน สโมธาน\-ปริวาสํ เทติ, ธมฺเมน มานตฺตํ เทติ, ธมฺเมน อพฺเภติ. โส ภิกฺขเว ภิกฺขุ วิสุทฺโธ ตาหิ อาปตฺตีหิ. \mbox{(1-3)} อิธ ปน ภิกฺขเว ภิกฺขุ สมฺพหุลา สํฆาทิเสสา อาปตฺติโย อาปชฺชติ ปริมาณมฺปิ อปริมาณมฺปิ ฯเปฯ ววตฺถิตมฺปิ สมฺภินฺนมฺปิ. โส สํฆํ ตาสํ อาปตฺตีนํ สโมธาน\-ปริวาสํ ยาจติ, ตสฺส สํโฆ ตาสํ อาปตฺตีนํ สโมธาน\-ปริวาสํ เทติ, โส ปริวสนฺโต อนฺตรา สมฺพหุลา สํฆาทิเสสา อาปตฺติโย อาปชฺชติ อปริมาณาโย}

\orphannote{มูลาย\-วิสุทฺธินวก (สี, สฺยา) อปฺปฏิจฺฉนฺนาโย ฯเปฯ อปริมาณาโย ปฏิจฺฉนฺนาโย ฯเปฯ อปริมาณาโย ปฏิจฺฉนฺนาโยปิ อปฺปฏิจฺฉนฺนาโยปิ. โส สํฆํ อนฺตรา\-อาปตฺตีนํ มูลาย\-ปฏิกสฺสนํ ยาจติ, ตํ สํโฆ อนฺตรา\-อาปตฺตีนํ มูลาย ปฏิกสฺสติ ธมฺมิเกน กมฺเมน อกุปฺเปน ฐานารเหน, ธมฺเมน สโมธาน\-ปริวาสํ เทติ, ธมฺเมน มานตฺตํ เทติ, ธมฺเมน อพฺเภติ. โส ภิกฺขเว ภิกฺขุ วิสุทฺโธ ตาหิ อาปตฺตีหิ. \mbox{(4-6)} อิธ ปน ภิกฺขเว ภิกฺขุ สมฺพหุลา สํฆาทิเสสา อาปตฺติโย อาปชฺชติ ปริมาณมฺปิ อปริมาณมฺปิ ฯเปฯ ววตฺถิตมฺปิ สมฺภินฺนมฺปิ. โส สํฆํ ตาสํ อาปตฺตีนํ สโมธาน\-ปริวาสํ ยาจติ, ตสฺส สํโฆ ตาสํ อาปตฺตีนํ สโมธาน\-ปริวาสํ เทติ, โส ปริวสนฺโต อนฺตรา สมฺพหุลา}

\orphannote{มูลาย\-วิสุทฺธินวก (สี, สฺยา) สํฆาทิเสสา อาปตฺติโย อาปชฺชติ ปริมาณาโยปิ อปริมาณาโยปิ อปฺปฏิจฺฉนฺนาโย ฯเปฯ ปริมาณาโยปิ อปริมาณาโยปิ ปฏิจฺฉนฺนาโย ฯเปฯ ปริมาณาโยปิ อปริมาณาโยปิ ปฏิจฺฉนฺนาโยปิ อปฺปฏิจฺฉนฺนาโยปิ. โส สํฆํ อนฺตรา\-อาปตฺตีนํ มูลาย\-ปฏิกสฺสนํ ยาจติ, ตํ สํโฆ อนฺตรา\-อาปตฺตีนํ มูลาย ปฏิกสฺสติ ธมฺมิเกน กมฺเมน อกุปฺเปน ฐานารเหน, ธมฺเมน สโมธาน\-ปริวาสํ เทติ, ธมฺเมน มานตฺตํ เทติ, ธมฺเมน อพฺเภติ. โส ภิกฺขเว ภิกฺขุ วิสุทฺโธ ตาหิ อาปตฺตีหิ. \mbox{(7-9)} มูลาย วิสุทฺธิ\-นวกํ นิฏฺฐิตํ.}

\orphannote{อํ 1. 362 ปิฏฺเฐปิ.}

